\documentclass[11pt,oneside]{book}
\usepackage[margin=1in]{geometry}
\usepackage{amsmath,amssymb,amsthm,mathtools}
\usepackage{booktabs,longtable,array}
\usepackage{enumitem}
\usepackage[colorlinks=true,linkcolor=blue,urlcolor=blue,citecolor=blue]{hyperref}
\usepackage[nameinlink,noabbrev]{cleveref}

\newtheorem{theorem}{Theorem}[chapter]
\newtheorem{lemma}[theorem]{Lemma}
\newtheorem{proposition}[theorem]{Proposition}
\newtheorem{corollary}[theorem]{Corollary}
\newtheorem{assessment}[theorem]{Assessment}
\newtheorem{firewall}[theorem]{Claim firewall}
\newtheorem{openproblem}[theorem]{Open problem}
\newtheorem{record}[theorem]{Record}
\newtheorem{programme}[theorem]{Programme}
\newtheorem{audit}[theorem]{Audit}
\newtheorem{correction}[theorem]{Correction}
\newtheorem{warning}[theorem]{Warning}
\theoremstyle{definition}
\newtheorem{definition}[theorem]{Definition}
\theoremstyle{remark}
\newtheorem{remark}[theorem]{Remark}

\newcommand{\MGclosed}{\textsc{closed}}
\newcommand{\MGopen}{\textsc{open}}
\newcommand{\MGpartial}{\textsc{partial}}
\newcommand{\opnorm}[1]{\lVert#1\rVert_{\mathrm{op}}}
\newcommand{\Tr}{\operatorname{Tr}}

\title{Renewal Yang--Mills Mass-Gap Research Notes\\
Sixth Series, Version V1}
\author{Compact proof-indexed continuation after archives f1--f5}
\date{September 2026}

\begin{document}
\maketitle
\tableofcontents

\chapter{Context, cumulative theorem manifest, and continuation}
\label{ch:s6v1-context}

\section{What this programme is doing}
\label{sec:s6v1-purpose}

The ultimate target is a rigorous construction of
four-dimensional quantum Yang--Mills theory with compact
nonabelian gauge group, including a continuum probability
measure, reflection positivity, Osterwalder--Schrader
reconstruction, and a strictly positive spectral gap above the
vacuum.  The present work is upstream of that target.  It studies
the concentrated one-link Wilson law and proves the
representation-uniform connected estimates needed to build a
cutoff-independent polymer expansion.

The preceding fifth series reached V100.  In accordance with the
version protocol it has been frozen, because a \(45\,000\)-line
active source is too difficult to audit reliably.  The
mathematics is not restarted: this new V1 records the principal
proved results, the corrections and no-go theorems which constrain
the proof order, the exact live implication chain, and then
continues at the vertex-accretion bottleneck.

No continuum Yang--Mills measure or Hamiltonian has yet been
constructed.  In particular, no mass gap is assumed or claimed.

\begin{record}[Immutable archive ledger]
\label{rec:s6v1-archives}
\begin{center}
\small
\begin{tabular}{p{0.62\textwidth}rr}
\toprule
archive & source lines & bytes\\
\midrule
\texttt{Renewal\_Yang\_Mills\_Mass\_Gap\_Research\_Notes\_V135\_f1.tex}
& \(88\,517\) & \(3\,083\,467\)\\
\texttt{Renewal\_Yang\_Mills\_Mass\_Gap\_Research\_Notes\_V100\_f2.tex}
& \(49\,581\) & \(1\,706\,928\)\\
\texttt{Renewal\_Yang\_Mills\_Mass\_Gap\_Research\_Notes\_V100\_f3.tex}
& \(48\,381\) & \(1\,720\,945\)\\
\texttt{Renewal\_Yang\_Mills\_Mass\_Gap\_Research\_Notes\_V100\_f4.tex}
& \(44\,712\) & \(1\,483\,037\)\\
\texttt{Renewal\_Yang\_Mills\_Mass\_Gap\_Research\_Notes\_V100\_f5.tex}
& \(45\,189\) & \(1\,467\,995\)\\
\bottomrule
\end{tabular}
\end{center}
The SHA--256 digest of the new f5 archive is
\[
 \texttt{8b9965394878a038518c180907033555ceb2494c17f2bcb0be7e46ea7e736857}.
\]
The earlier archive digests remain recorded in the opening
manifest of f5.  References to an archived version below are
descriptive proof indices; they are not local \LaTeX\
cross-references.
\end{record}

\begin{record}[Version and companion-audit protocol]
\label{rec:s6v1-protocol}
Every ordinary successor copies the complete active source,
appends one proof chapter, passes the append-prefix, terminator,
environment, label, and reference checks, and then replaces its
predecessor.  Every version divisible by five audits the newest
active Standard-Model--Einstein and Navier--Stokes notes
read-only.  At V100, the active source is frozen under the next
suffix \(\texttt{\_fX}\) and a contextual V1 starts.
\end{record}

\section{The long dependency chain}
\label{sec:s6v1-chain}

The intended implication chain is
\begin{equation}
\begin{split}
\text{small-clock nonlinear source card (SCNGC)}
&\Longrightarrow
\text{one-link Wilson connected cards}
\Longrightarrow\text{marked-tree atlas (MTA)}
\\
&\Longrightarrow\text{weighted uniform fusion control (WUFC)}
\Longrightarrow\text{connected polymer majorant (CPM)}
\\
&\Longrightarrow\text{cutoff-uniform finite-volume measures}
\Longrightarrow\text{cutoff removal and reflection positivity}
\\
&\Longrightarrow\text{Osterwalder--Schrader reconstruction}
\Longrightarrow\text{positive Hamiltonian gap}.
\end{split}                                                \label{eq:s6v1-long-chain}
\end{equation}
The active work is still in the first line.

Within the remaining intermediate thermal wedge, the current
local chain is
\begin{equation}
\begin{split}
\mathrm{VAFPC}
&\Longrightarrow\mathrm{FPC}
\Longrightarrow\mathrm{ULCC}
\Longrightarrow\mathrm{IGMC}
\Longrightarrow\mathrm{SISC}
\\
&\Longrightarrow\mathrm{TCNC}
\Longrightarrow\mathrm{TIC}
\Longrightarrow\mathrm{HLTC}
\Longrightarrow\text{full MCWC and SCNGC}.
\end{split}                                                \label{eq:s6v1-local-chain}
\end{equation}
All arrows displayed here are proved conditional implications in
archive f5.  The first hypothesis, VAFPC, is not yet proved.

\section{Cumulative manifest of the main proved mathematics}
\label{sec:s6v1-manifest}

\begin{theorem}[Representation and one-link analytic foundation]
\label{thm:s6v1-foundation}
The archives establish:
\begin{enumerate}[label=\textup{(\roman*)},leftmargin=3.5em]
\item Peter--Weyl output pinching and dimension-free
      representation normalization;
\item exact compact-group heat/L\'evy source formulae and
      dimension-free heat-increment moments;
\item the strong same-jump root estimate
      \[
       \|Q_A\|_{\rm cb}
       \le K^{|A|}
       \left(\prod_{i\in A}b_i\right)B_A^{|A|-2};
      \]
\item the clock-reserve estimate
      \[
       \|R_A\|_{\rm cb}
       \le(K|A|)^{|A|}s^{(|A|-2)/2}
       \prod_{i\in A}b_i;
      \]
\item exact endpoint-factorizing interpolation and its BKAR
      tree representation; and
\item complete percolative Duhamel contraction at every
      intermediate interpolation matrix.
\end{enumerate}
\end{theorem}

\begin{proof}
These are the result ledgers of f1 and f5-V38, V54, V60--V62,
V70, V81, and V93.  Each listed statement was proved there with
constants independent of the number of labels and the
representations.
\end{proof}

\begin{theorem}[Connected hypergraph and tree closures]
\label{thm:s6v1-tree-closures}
The following sectors satisfy the global strong target
\[
 K^m\left(\prod_{i=1}^mb_i\right)B^{m-2}:
\]
\begin{enumerate}[label=\textup{(\roman*)},leftmargin=3.5em]
\item every local connected root gas inside a light thermal
      block, including arbitrary incidence cycles;
\item every primitive incidence hypertree, with arbitrary local
      and transversal root arities;
\item every ancestry with exactly one transversal root and
      arbitrary local attachments;
\item every locally overlapped fundamental transversal cycle;
\item every fundamental transversal overlap of arity at most
      \(O(m/\log m)\); and
\item every connected simple primitive core whose incidence
      excess satisfies
      \(g\le\lfloor m/\log(m+e)\rfloor\).
\end{enumerate}
\end{theorem}

\begin{proof}
These are f5-V78, V91--V96.  The principal mechanisms are the
reserve-weighted incidence-hypertree/Cayley bijection, exact
rooted-forest formulae, union--intersection fusion, and the
excess-graded canonical skeleton.
\end{proof}

\begin{theorem}[Normalized combinatorial identities]
\label{thm:s6v1-normalizations}
The archives prove:
\begin{enumerate}[label=\textup{(\roman*)},leftmargin=3.5em]
\item the weighted Cayley--Pr\"ufer and rooted-forest identities;
\item the exact maximum-spanning-tree/Kruskal partition of unity
      for every fixed shortcut graph;
\item the exact forest-gate coefficient
      \[
       \gamma_A(E)
       ={1\over B_A^{|E|}}
       {\left(\prod_{\{i,j\}\in E}b_ib_j\right)
        \left(\prod_{C\in{\cal K}(E)}B_C\right)
        \over\prod_{i\in A}b_i};
      \]
\item exponential normalization of repeated support types before
      edge extraction;
\item exact multi-anchor tree/forest ratios and an
      \(O(w_P^2)\) anchor-partition gas; and
\item normalized minimum-edge disintegration of the BKAR cube.
\end{enumerate}
\end{theorem}

\begin{proof}
See f5-V75--V77, V83, V86--V87, V91--V92, and V100.
\end{proof}

\begin{theorem}[Optional-support and ultrametric-layer progress]
\label{thm:s6v1-layer-progress}
Let \(E\) be a nonempty owned forest.
\begin{enumerate}[label=\textup{(\roman*)},leftmargin=3.5em]
\item At a saturated gate, the complete sum over supports
      \(A\supseteq V(E)\) collapses before norms to one reserve
      activity on \(V(E)\).
\item The same collapse holds on every uniform interpolation
      line.
\item Every BKAR matrix has a nested ultrametric layer
      decomposition.
\item Expanding one layer creates a Bernoulli subprobability, not
      a Bell census.
\item After the layer expansion, every optional support label is
      removed; the survivors are forced forests with weight
      \[
       \lambda^{c(G)-1}(1-\lambda)^{|H|}
       \gamma_{V(G)}(G)\mathfrak r(V(G)).
      \]
\item All completions which introduce no new forced vertex have
      total conditional Cayley mass at most one, at every
      ultrametric depth.
\end{enumerate}
\end{theorem}

\begin{proof}
These are f5-V97--V100.  The key algebra is
\[
 I+\lambda(D^{R}(U)-I)
 =(1-\lambda)I+\lambda D^{R}(U),
\]
which is a contraction, followed by the Laplace representation
of \(B_A^{-|E|}\) and the conditional Cayley-tree probability.
\end{proof}

\section{Corrections and proof-order no-go theorems}
\label{sec:s6v1-no-go}

\begin{record}[Constraints which remain in force]
\label{rec:s6v1-no-go}
The archives rigorously rule out the following shortcuts.
\begin{enumerate}[label=\textup{(\roman*)},leftmargin=3.5em]
\item Termwise Fr\'echet differentiation creates a Bell-number
      majorant.
\item A global \(B^{2g}\) relaxation cannot sum arbitrary
      intermediate-wedge incidence excess.
\item A normalized endpoint probability cannot create a raw
      thermal port mark; the missing mark must be cut from a
      genuine global tree.
\item Block compatibility is not implied by connectedness.
\item Termwise union--intersection uncrossing has a potentially
      superexponential inverse fibre.
\item A degree-vector fibre can contain \(\exp(cm^2)\) distinct
      common-core support families, even after the known type
      factorial is granted.
\item Product-law interpolation is false on the whole entrywise
      cube; only the actual BKAR ultrametric locus is relevant.
\item Boundary rows or localized rows do not determine a missing
      interior operator estimate.
\end{enumerate}
\end{record}

\begin{assessment}[Exact live bottleneck]
\label{ass:s6v1-bottleneck}
After f5-V100, a child-layer completion forest \(H\) is already
summed whenever
\[
 V(E\cup H)=V(E).
\]
The only unsummed row has at least one edge of \(H\) incident to
a previously optional label.  Such a label becomes forced and
must be paid simultaneously by the beta factor, the conditional
Cayley probability, the clock reserve, and the rooted child-tree
sum.  This is VAFPC.
\end{assessment}

\chapter{Sixth-series V1: the first vertex-accretion estimate}
\label{ch:s6v1-accretion}

\section{Reserve and forest notation}
\label{sec:s6v1-notation}

Fix \(0<s\le1\) and marks
\[
 \sqrt s\le b_i<1,\qquad
 B_A:=\sum_{i\in A}b_i.
\]
With one fixed base \(K_0\), put
\begin{equation}
 \mathfrak r(A)
 :=
 (K_0|A|)^{|A|}
 s^{(|A|-2)/2}\prod_{i\in A}b_i,
 \qquad |A|\ge2.                                          \label{eq:s6v1-reserve}
\end{equation}
For a forest \(E\) on \(A\), \(\gamma_A(E)\) is the normalized
weighted-Cayley probability that a random tree on \(A\) contains
\(E\).

\section{Adding one new forced leaf}
\label{sec:s6v1-leaf}

Let \(E\) be a forest on \(F\), let \(r\in F\), and let \(C_r\)
be the component of \(E\) containing \(r\).  Put
\[
 f:=|F|,\qquad k:=|E|.
\]
For \(j\notin F\), set
\[
 G_j:=E\cup\{\{r,j\}\}.
\]

\begin{lemma}[Exact Cayley-probability ratio for one new leaf]
\label{lem:s6v1-gamma-leaf}
One has
\begin{equation}
 {\gamma_{F\cup\{j\}}(G_j)\over\gamma_F(E)}
 =
 b_r{B_{C_r}+b_j\over B_{C_r}}\,
 {B_F^k\over(B_F+b_j)^{k+1}}
 \le {b_r+b_j\over B_F+b_j}.                              \label{eq:s6v1-gamma-leaf}
\end{equation}
\end{lemma}

\begin{proof}
Use the exact forest formula from the manifest.  Adding the leaf
edge contributes \(b_rb_j\); the new mark \(b_j\) in the
denominator cancels it except for \(b_r\); and the component mass
\(B_{C_r}\) becomes \(B_{C_r}+b_j\).  The forest-edge count rises
from \(k\) to \(k+1\).  This proves the equality.

For the inequality, first use
\[
 {B_F^k\over(B_F+b_j)^k}\le1.
\]
Then
\[
 b_r{B_{C_r}+b_j\over B_{C_r}}
 =b_r+{b_rb_j\over B_{C_r}}
 \le b_r+b_j,
\]
because \(b_r\le B_{C_r}\).
\end{proof}

\begin{lemma}[Exact reserve ratio for one new leaf]
\label{lem:s6v1-reserve-leaf}
After enlarging a numerical base,
\begin{equation}
 {\mathfrak r(F\cup\{j\})\over\mathfrak r(F)}
 \le K(f+1)\sqrt s\,b_j
 \le K(B_F+b_j)b_j.                                      \label{eq:s6v1-reserve-leaf}
\end{equation}
\end{lemma}

\begin{proof}
Directly from \eqref{eq:s6v1-reserve},
\[
 {\mathfrak r(F\cup\{j\})\over\mathfrak r(F)}
 =
 K_0{(f+1)^{f+1}\over f^f}\sqrt s\,b_j.
\]
Since \((1+1/f)^f\le e\), the first bound follows.  The thermal
floor gives
\[
 (f+1)\sqrt s\le B_F+b_j,
\]
which proves the second.
\end{proof}

\begin{theorem}[One-leaf vertex-accretion bound]
\label{thm:s6v1-one-leaf}
For every \(j\notin F\),
\begin{equation}
 \gamma_{F\cup\{j\}}(G_j)\mathfrak r(F\cup\{j\})
 \le
 K(b_r+b_j)b_j\,
 \gamma_F(E)\mathfrak r(F).                               \label{eq:s6v1-one-leaf}
\end{equation}
Consequently, for every optional set \(J\subseteq X\setminus F\),
\begin{equation}
 \sum_{j\in J}
 \gamma_{F\cup\{j\}}(G_j)\mathfrak r(F\cup\{j\})
 \le
 K\{b_rB_J+B_J^2\}
 \gamma_F(E)\mathfrak r(F).                               \label{eq:s6v1-one-leaf-sum}
\end{equation}
\end{theorem}

\begin{proof}
Multiply
\Cref{lem:s6v1-gamma-leaf,lem:s6v1-reserve-leaf}; the factor
\(B_F+b_j\) cancels.  Sum in \(j\) and use
\[
 \sum_{j\in J}b_j^2\le\left(\sum_{j\in J}b_j\right)^2
 =B_J^2.
\]
\end{proof}

\begin{corollary}[Global single-accretion row]
\label{cor:s6v1-global-leaf}
On the intermediate wedge, the complete choice of one newly
forced leaf costs at most \(K^m\) relative to the pre-accretion
forest weight.
\end{corollary}

\begin{proof}
Take \(J=X\setminus F\) in
\eqref{eq:s6v1-one-leaf-sum}.  Since
\(b_r\le B\le m\),
\[
 b_rB_J+B_J^2\le2B^2\le2m^2\le K^m
\]
after one numerical enlargement of \(K\).
\end{proof}

\begin{corollary}[Light-block leaf gas]
\label{cor:s6v1-light-leaves}
Suppose \(r\) and every allowed new leaf lie in one light thermal
block \(P\), with \(w_P\le b_*/2\).  Then the summed one-leaf
activity is at most
\begin{equation}
 Kw_P^2\,\gamma_F(E)\mathfrak r(F).                        \label{eq:s6v1-light-leaves}
\end{equation}
After \(b_*\) is decreased, an arbitrary ordered sequence of
distinct leaf accretions in \(P\) has a bounded geometric sum.
\end{corollary}

\begin{proof}
In \eqref{eq:s6v1-one-leaf-sum},
\(b_r+B_J\le w_P\), so
\[
 b_rB_J+B_J^2=(b_r+B_J)B_J\le w_P^2.
\]
At each later accretion the same estimate applies with a smaller
unused set.  Positivity bounds the ordered sequence of length
\(\ell\) by \((Kw_P^2)^\ell\).  Choose \(b_*\) so that
\(K(b_*/2)^2\le1/2\), and sum the geometric series.
\end{proof}

\begin{assessment}[What has been closed in the restarted volume]
\label{ass:s6v1-leaf-scope}
The first VAFPC row is closed: exactly one global new leaf costs
only \(K^m\), and an arbitrary gas of new leaves is closed when
it remains inside one light thermal block.  The key cancellation
is quantitative and local to the normalized Cayley probability:
its denominator \((B_F+b_j)^{-1}\) cancels the arity factor
created by the reserve extension.

The theorem does not yet cover a rooted child tree with two or
more new vertices spanning several thermal blocks, nor an
arbitrary gas of global leaf accretions.  Replacing every global
one-leaf factor by \(KB^2\) would recreate
\(\exp(CB^2)\) and is forbidden.
\end{assessment}

\section{Next acquisition}
\label{sec:s6v1-next}

\begin{definition}[Rooted child-tree accretion card, RCTAC]
\label{def:s6v1-RCTAC}
RCTAC is the VAFPC row in which each new-vertex component of the
child-layer forest meets the old forced set in exactly one root.
It must sum the entire rooted tree by the weighted rooted-forest
identity before applying
\Cref{thm:s6v1-one-leaf}, retain the beta power
\((1-\lambda)^{|H|}\), and produce either a light-block gas or one
global \(K^m\) factor for the complete collection, not one factor
per new vertex.
\end{definition}

\begin{proposition}[RCTAC plus floating components proves VAFPC]
\label{prop:s6v1-RCTAC}
If RCTAC holds and the analogous row of child-layer components
disjoint from the old forced set is bounded by their top-level
factor \(\lambda\), then VAFPC holds.
\end{proposition}

\begin{proof}
Every component of the vertex-accreting forest either meets the
old forced set or is disjoint from it.  Contract its intersection
with the old set.  The first family is RCTAC; the second is the
stated floating-component row.  Fixed-vertex edges were already
summed in f5-V100.  These three disjoint classes exhaust the
forest.
\end{proof}

\begin{programme}[V2: rooted child trees]
\label{prog:s6v1-V2}
\begin{enumerate}[leftmargin=3em]
\item Derive the exact ratio
      \(\gamma_{F\cup J}(E\cup H)/
        \gamma_F(E)\) for a rooted tree \(H\) on
      \(\{r\}\cup J\).
\item Sum \(H\) by the weighted rooted-tree identity before
      summing \(J\).
\item Test whether the reserve arity ratio cancels against the
      full power \(B_{F\cup J}^{-|J|}\), generalizing the
      one-leaf cancellation.
\item Close the light-block row first, then test the transversal
      rooted tree.
\item Keep floating child components separate because they spend
      one top-level \(\lambda\) factor.
\end{enumerate}
\end{programme}

\begin{firewall}[Claim boundary after sixth-series V1]
\label{fire:s6v1-boundary}
This V1 gives the context, immutable archive ledger, long and
local dependency chains, a compact cumulative theorem manifest,
and the principal corrections/no-go statements.  It proves the
exact normalized Cayley-probability ratio for adding one forced
leaf, proves the corresponding reserve ratio, closes the global
single-leaf VAFPC row, and closes an arbitrary light-block leaf
gas.  It formulates RCTAC.

It does not prove RCTAC, the floating-component row, VAFPC in
full, FPC, ULCC, full IGMC, SISC, the simple high-excess row,
arbitrary repeated transversal supports, full TOCC, full TCNC,
cyclic TIC, HLTC, full MCWC, OCRC, SCNGC, WPSC, GRLTC, LZTC,
PLRC, WSDC, BRSC, HWSFC, UGCGC, thermal connected WPRAC,
all-arity RGSC, the raw Wilson aggregate strong card, all-order
strong MTA, WUFC, CPM, cutoff removal, reflection positivity,
Osterwalder--Schrader reconstruction, a continuum Yang--Mills
measure, or a positive Hamiltonian gap.  The full Yang--Mills
mass gap has not been proved.
\end{firewall}

\begin{theorem}[Sixth-series V1 result ledger]
\label{thm:s6v1-results}
This version:
\begin{enumerate}[label=\textup{(V1.\arabic*)},leftmargin=5.5em]
\item records the context, protocol, archives, and dependency
      chains;
\item records the cumulative proved-result manifest;
\item records the proof-order no-go theorems;
\item isolates VAFPC after the fixed-vertex closure;
\item proves the exact one-leaf Cayley-probability ratio;
\item proves the one-leaf reserve ratio;
\item closes the global single-leaf and local leaf-gas rows; and
\item formulates RCTAC and its relation to VAFPC.
\end{enumerate}
\end{theorem}

\begin{proof}
Use
\Cref{thm:s6v1-foundation,
thm:s6v1-tree-closures,
thm:s6v1-normalizations,
thm:s6v1-layer-progress,
lem:s6v1-gamma-leaf,
lem:s6v1-reserve-leaf,
thm:s6v1-one-leaf,
cor:s6v1-global-leaf,
cor:s6v1-light-leaves,
prop:s6v1-RCTAC}.
\end{proof}

\paragraph{Restart record.}
This sixth-series V1 follows the immutable archive
\texttt{Renewal\_Yang\_Mills\_Mass\_Gap\_Research\_Notes\_V100\_f5.tex},
with \(45\,189\) logical source lines, \(1\,467\,995\) bytes, and
SHA--256
\[
 \texttt{8b9965394878a038518c180907033555ceb2494c17f2bcb0be7e46ea7e736857}.
\]
The next compulsory companion audit is V5.

\clearpage
\chapter{Sixth-series V2: rooted child-tree accretion}
\label{ch:s6v2-rooted}

\section{Exact forest-probability ratio for a rooted child tree}
\label{sec:s6v2-ratio}

No companion audit is due at V2.

Retain the notation of V1.  Thus \(E\) is a forest on \(F\),
\(r\in F\), \(C_r\) is the component of \(E\) containing \(r\),
\[
 f=|F|,\qquad k=|E|.
\]
Let \(J\subseteq X\setminus F\), \(t:=|J|\ge1\), and let \(H\)
be a tree on \(\{r\}\cup J\).  Put
\[
 G_H:=E\cup H.
\]
The two edge sets meet only at \(r\), so \(G_H\) is a forest.

\begin{lemma}[Exact rooted-tree probability ratio]
\label{lem:s6v2-gamma-ratio}
For every such \(H\),
\begin{align}
 {\gamma_{F\cup J}(G_H)\over\gamma_F(E)}
 &=
 \left\{
  {\prod_{\{u,v\}\in H}b_ub_v\over\prod_{j\in J}b_j}
 \right\}
 {B_{C_r}+B_J\over B_{C_r}}\,
 {B_F^k\over(B_F+B_J)^{k+t}}.                             \label{eq:s6v2-gamma-ratio}
\end{align}
\end{lemma}

\begin{proof}
Use the exact forest formula for \(\gamma\).  The new edges supply
\(\prod_{\{u,v\}\in H}b_ub_v\), the new denominator marks supply
\(\prod_{j\in J}b_j\), the component mass \(B_{C_r}\) becomes
\(B_{C_r}+B_J\), and the forest-edge count rises from \(k\) to
\(k+t\).  All unchanged component and mark factors cancel.
\end{proof}

\begin{lemma}[Rooted weighted-Cayley sum]
\label{lem:s6v2-rooted-Cayley}
\begin{equation}
 \sum_{H\in\mathfrak T(\{r\}\cup J)}
 {\prod_{\{u,v\}\in H}b_ub_v\over\prod_{j\in J}b_j}
 =
 b_r(b_r+B_J)^{t-1}.                                     \label{eq:s6v2-rooted-Cayley}
\end{equation}
\end{lemma}

\begin{proof}
The weighted Cayley identity on \(\{r\}\cup J\) is
\[
 \sum_H\prod_{\{u,v\}\in H}b_ub_v
 =
 b_r\left(\prod_{j\in J}b_j\right)
 (b_r+B_J)^{t-1}.
\]
Divide by the new-vertex marks.
\end{proof}

\begin{theorem}[A complete rooted child tree does not enlarge the
Cayley probability]
\label{thm:s6v2-gamma-contraction}
For every fixed new vertex set \(J\),
\begin{equation}
 \sum_{H\in\mathfrak T(\{r\}\cup J)}
 \gamma_{F\cup J}(G_H)
 \le\gamma_F(E).                                         \label{eq:s6v2-gamma-contraction}
\end{equation}
More precisely, the ratio of the left side to the right side is
at most
\begin{equation}
 \left({b_r+B_J\over B_F+B_J}\right)^t.                   \label{eq:s6v2-gamma-sharp}
\end{equation}
\end{theorem}

\begin{proof}
Sum \eqref{eq:s6v2-gamma-ratio} and apply
\Cref{lem:s6v2-rooted-Cayley}.  The resulting ratio is
\[
 b_r(b_r+B_J)^{t-1}
 {B_{C_r}+B_J\over B_{C_r}}
 {B_F^k\over(B_F+B_J)^{k+t}}.
\]
Since \(B_F/(B_F+B_J)\le1\) and
\[
 b_r{B_{C_r}+B_J\over B_{C_r}}
 \le b_r+B_J,
\]
this is at most \eqref{eq:s6v2-gamma-sharp}.  Finally
\(b_r\le B_F\), so the latter is at most one.
\end{proof}

\section{The reserve cancels against the probability denominator}
\label{sec:s6v2-reserve}

\begin{lemma}[Multi-vertex reserve extension]
\label{lem:s6v2-reserve-ratio}
For \(J\) as above,
\begin{equation}
 {\mathfrak r(F\cup J)\over\mathfrak r(F)}
 \le
 \{K(B_F+B_J)\}^{t}\prod_{j\in J}b_j.                     \label{eq:s6v2-reserve-ratio}
\end{equation}
\end{lemma}

\begin{proof}
Directly from the reserve definition,
\[
 {\mathfrak r(F\cup J)\over\mathfrak r(F)}
 =
 K_0^t{(f+t)^{f+t}\over f^f}s^{t/2}
 \prod_{j\in J}b_j.
\]
Now
\[
 {(f+t)^{f+t}\over f^f}
 =(f+t)^t(1+t/f)^f
 \le\{e(f+t)\}^t.
\]
The thermal floor gives
\((f+t)\sqrt s\le B_F+B_J\).  Absorb \(e\) into \(K\).
\end{proof}

\begin{theorem}[Rooted child-tree accretion estimate]
\label{thm:s6v2-rooted-accretion}
For every fixed \(J\),
\begin{align}
 &\sum_{H\in\mathfrak T(\{r\}\cup J)}
 \gamma_{F\cup J}(G_H)\mathfrak r(F\cup J)
 \notag\\
 &\qquad\le
 \{K(b_r+B_J)\}^{t}
 \left(\prod_{j\in J}b_j\right)
 \gamma_F(E)\mathfrak r(F).                               \label{eq:s6v2-rooted-accretion}
\end{align}
\end{theorem}

\begin{proof}
Multiply the sharp probability ratio
\eqref{eq:s6v2-gamma-sharp} by
\eqref{eq:s6v2-reserve-ratio}.  The power
\((B_F+B_J)^t\) cancels exactly.
\end{proof}

\begin{remark}[V1 is the first coefficient]
\label{rem:s6v2-V1}
At \(t=1\), the unique rooted tree is the edge \(\{r,j\}\), and
\eqref{eq:s6v2-rooted-accretion} is precisely the one-leaf bound
of \Cref{thm:s6v1-one-leaf}, up to the harmless common base.
\end{remark}

\section{Summing the new vertex set}
\label{sec:s6v2-sets}

\begin{lemma}[Fixed-cardinality vertex-set sum]
\label{lem:s6v2-J-sum}
Let \(Y\subseteq X\setminus F\), with total mark \(B_Y\).  Then
\begin{align}
 &\sum_{\substack{J\subseteq Y\\|J|=t}}
 \sum_{H\in\mathfrak T(\{r\}\cup J)}
 \gamma_{F\cup J}(G_H)\mathfrak r(F\cup J)
 \notag\\
 &\quad\le
 { \{K(b_r+B_Y)B_Y\}^{t}\over t!}\,
 \gamma_F(E)\mathfrak r(F).                               \label{eq:s6v2-J-sum}
\end{align}
\end{lemma}

\begin{proof}
In \eqref{eq:s6v2-rooted-accretion}, use
\(B_J\le B_Y\), then apply
\[
 \sum_{\substack{J\subseteq Y\\|J|=t}}
 \prod_{j\in J}b_j
 \le {B_Y^t\over t!}.
\]
\end{proof}

\begin{theorem}[Light-block RCTAC]
\label{thm:s6v2-light-RCTAC}
Suppose \(r\) and every allowed new vertex lie in one light
thermal block \(P\), \(w_P\le b_*/2\).  Then the complete sum over
all nonempty \(J\subseteq P\setminus F\) and all rooted trees on
\(\{r\}\cup J\) is at most
\begin{equation}
 \{\exp(Kw_P^2)-1\}
 \gamma_F(E)\mathfrak r(F).                               \label{eq:s6v2-light-RCTAC}
\end{equation}
After decreasing \(b_*\), an arbitrary gas of these rooted child
trees has a bounded geometric sum.
\end{theorem}

\begin{proof}
In \Cref{lem:s6v2-J-sum},
\[
 b_r+B_Y\le w_P,\qquad B_Y\le w_P.
\]
Sum \(t\ge1\) to obtain
\(\sum_{t\ge1}(Kw_P^2)^t/t!=e^{Kw_P^2}-1\).
Choose \(b_*\) so this activity is at most \(1/2\); the gas then
sums geometrically.
\end{proof}

\begin{theorem}[Submacroscopic global RCTAC]
\label{thm:s6v2-submacro-RCTAC}
Put
\[
 T_m:=\max\left\{1,
 \left\lfloor{m\over\log(m+e)}\right\rfloor\right\}.
\]
On the intermediate wedge, the complete sum of rooted child-tree
accretions with \(1\le|J|\le T_m\) costs at most \(K^m\) relative
to \(\gamma_F(E)\mathfrak r(F)\).
\end{theorem}

\begin{proof}
Take \(Y=X\setminus F\) in
\eqref{eq:s6v2-J-sum}.  Since \(b_r+B_Y\le2B\) and \(B_Y\le B\),
the term of cardinality \(t\) is at most
\[
 {(KB^2)^t\over t!}\le(Km^2)^t.
\]
For \(t\le T_m\),
\[
 t\log(Km^2)
 \le {m\log(Km^2)\over\log(m+e)}
 \le c_0m.
\]
There are at most \(m\) terms, and this last factor is absorbed
into a fixed exponential base.
\end{proof}

\begin{assessment}[Remaining rooted-tree sector]
\label{ass:s6v2-remaining}
RCTAC is now proved for every light-block rooted child tree and
for every global rooted child tree with at most
\(m/\log(m+e)\) new vertices.  The remaining rooted row has a
transversal child tree with
\[
 |J|>{m\over\log(m+e)}.
\]
For that row the separate bound
\(\exp(KB^2)\) is unavailable.  Its large vertex set must be
summed jointly with the top-level beta integral and the thermal
block quotient, rather than one root at a time.
\end{assessment}

\section{Next acquisition}
\label{sec:s6v2-next}

\begin{definition}[Macroscopic rooted child-tree card, MRCTC]
\label{def:s6v2-MRCTC}
MRCTC is the remaining RCTAC sector
\(|J|>m/\log(m+e)\).  It must retain:
\begin{enumerate}[label=\textup{(\roman*)},leftmargin=3.5em]
\item the minimum-edge Jacobian
      \((1-\lambda)^{m-2}\);
\item the child-tree factor \((1-\lambda)^{|J|}\);
\item the exact rooted Cayley sum of
      \Cref{lem:s6v2-rooted-Cayley};
\item the thermal partition of \(J\); and
\item the block-level Kruskal probability.
\end{enumerate}
The target is \(K^m\), uniformly in the depth of the ultrametric
hierarchy.
\end{definition}

\begin{proposition}[MRCTC and the floating row prove VAFPC]
\label{prop:s6v2-MRCTC}
MRCTC, the floating-component row of
\Cref{prop:s6v1-RCTAC}, and the proved light/submacroscopic rows
complete VAFPC.
\end{proposition}

\begin{proof}
Split every rooted vertex-accreting component by whether it is
local or transversal and, in the latter case, by the threshold
\(T_m\).  The first two rows are
\Cref{thm:s6v2-light-RCTAC,thm:s6v2-submacro-RCTAC}; MRCTC is the
complement.  Add the disjoint floating components as in
\Cref{prop:s6v1-RCTAC}.
\end{proof}

\begin{programme}[V3: integrate the macroscopic rooted row]
\label{prog:s6v2-V3}
\begin{enumerate}[leftmargin=3em]
\item Combine
      \((1-\lambda)^{m-2+|J|}\) before integrating the
      minimum-edge parameter.
\item Determine the exact factorial ratio contributed by
      \(\int_0^1(1-\lambda)^{m-2+|J|}\,d\lambda\).
\item Sum the choice of the minimizing outer-tree edge only after
      that ratio is retained.
\item Test whether the resulting factor \(m/(m+|J|)\) is strong
      enough with the rooted Cayley probability; if not, use the
      full order statistics of all layer minima.
\item Keep floating components separate.
\end{enumerate}
\end{programme}

\begin{firewall}[Claim boundary after sixth-series V2]
\label{fire:s6v2-boundary}
V2 proves the exact Cayley-probability ratio for an arbitrary
rooted child tree, sums all rooted tree shapes with no probability
loss, proves exact cancellation of the reserve arity denominator,
closes the full light-block RCTAC row, and closes all global rooted
accretions with at most \(m/\log(m+e)\) new vertices.  It isolates
MRCTC.

V2 does not prove MRCTC, the floating-component row, VAFPC in
full, FPC, ULCC, full IGMC, SISC, the simple high-excess row,
arbitrary repeated transversal supports, full TOCC, full TCNC,
cyclic TIC, HLTC, full MCWC, OCRC, SCNGC, WPSC, GRLTC, LZTC,
PLRC, WSDC, BRSC, HWSFC, UGCGC, thermal connected WPRAC,
all-arity RGSC, the raw Wilson aggregate strong card, all-order
strong MTA, WUFC, CPM, cutoff removal, reflection positivity,
Osterwalder--Schrader reconstruction, a continuum Yang--Mills
measure, or a positive Hamiltonian gap.  The full Yang--Mills
mass gap has not been proved.
\end{firewall}

\begin{theorem}[Sixth-series V2 result ledger]
\label{thm:s6v2-results}
V2:
\begin{enumerate}[label=\textup{(V2.\arabic*)},leftmargin=5.5em]
\item proves the exact rooted-tree probability ratio;
\item proves the rooted weighted-Cayley sum;
\item proves probability contraction after summing all shapes;
\item proves the multi-vertex reserve ratio;
\item proves rooted child-tree accretion;
\item proves the fixed-cardinality vertex-set sum;
\item closes light-block RCTAC;
\item closes submacroscopic global RCTAC; and
\item formulates MRCTC.
\end{enumerate}
\end{theorem}

\begin{proof}
Use
\Cref{lem:s6v2-gamma-ratio,
lem:s6v2-rooted-Cayley,
thm:s6v2-gamma-contraction,
lem:s6v2-reserve-ratio,
thm:s6v2-rooted-accretion,
lem:s6v2-J-sum,
thm:s6v2-light-RCTAC,
thm:s6v2-submacro-RCTAC,
prop:s6v2-MRCTC}.
\end{proof}

\paragraph{Parent-version record.}
V2 was formed from
\texttt{Renewal\_Yang\_Mills\_Mass\_Gap\_Research\_Notes\_V1.tex}.
The V1 parent had \(612\) logical source lines, \(20\,943\) bytes,
and SHA--256
\[
 \texttt{a4f76a50a24dc218efe4254a67bad7219aeff14c621c03545077c62887c9a750}.
\]
The next compulsory companion audit is V5.

\clearpage
\chapter{Sixth-series V3: tree-valued closure of rooted
vertex accretion}
\label{ch:s6v3-tree-valued}

\section{The scalar minimum-edge gain is insufficient}
\label{sec:s6v3-scalar-no-go}

No companion audit is due at V3.

\begin{lemma}[The first minimum supplies only a rational gain]
\label{lem:s6v3-minimum-gain}
If a rooted child tree has \(t\) edges, combining its top-layer
factor \((1-\lambda)^t\) with the minimum-edge Jacobian and then
summing the \(m-1\) choices of the minimizing outer edge gives
\begin{equation}
 (m-1)\int_0^1(1-\lambda)^{m-2+t}\,d\lambda
 ={m-1\over m-1+t}.                                      \label{eq:s6v3-minimum-gain}
\end{equation}
\end{lemma}

\begin{proof}
Integrate the displayed power.
\end{proof}

\begin{assessment}[Why the V2 scalar route stops]
\label{ass:s6v3-scalar-no-go}
For \(1\le t\le m\), the factor in
\eqref{eq:s6v3-minimum-gain} lies between approximately \(1/2\)
and one.  It cannot pay the scalar series
\((KB^2)^t/t!\) when \(B\) is proportional to \(m\).
The rooted Cayley polynomial must therefore remain part of the
global output tree instead of being converted into a scalar gas.
\end{assessment}

\section{Rooted forests attached to one old component}
\label{sec:s6v3-rooted-forest}

Let \(E\) be a forest on \(F\), with \(k=|E|\).  Fix one
component \(C\) of \(E\).  Let \(J\subseteq X\setminus F\),
\(t=|J|\ge1\).

Denote by \(\mathfrak F_C(J)\) the forests \(H\) on
\(C\cup J\) such that:
\begin{enumerate}[label=\textup{(\roman*)},leftmargin=3.5em]
\item \(H\) contains no edge with both endpoints in \(C\);
\item every component of \(H\) contains exactly one vertex of
      \(C\); and
\item every vertex of \(J\) belongs to \(H\).
\end{enumerate}
Thus \(H\) has exactly \(t\) edges and \(E\cup H\) is a forest.

\begin{lemma}[Exact weighted rooted-forest polynomial]
\label{lem:s6v3-rooted-polynomial}
\begin{equation}
 \Phi_C(J)
 :=
 \sum_{H\in\mathfrak F_C(J)}
 \prod_{\{u,v\}\in H}b_ub_v
 =
 B_C\left(\prod_{j\in J}b_j\right)
 (B_C+B_J)^{t-1}.                                        \label{eq:s6v3-rooted-polynomial}
\end{equation}
\end{lemma}

\begin{proof}
This is the weighted rooted-forest identity with root set \(C\)
and nonroot set \(J\).  Equivalently, add one auxiliary vertex of
weight \(B_C\), apply the weighted Cayley identity on
\(\{\ast\}\cup J\), and expand each edge incident to \(\ast\)
by choosing its endpoint in \(C\).
\end{proof}

\begin{lemma}[Exact gamma ratio before scalar relaxation]
\label{lem:s6v3-gamma-rooted-forest}
For \(H\in\mathfrak F_C(J)\),
\begin{align}
 {\gamma_{F\cup J}(E\cup H)\over\gamma_F(E)}
 &=
 {\prod_{\{u,v\}\in H}b_ub_v\over\prod_{j\in J}b_j}
 {B_C+B_J\over B_C}
 {B_F^k\over(B_F+B_J)^{k+t}}.                             \label{eq:s6v3-gamma-rooted-forest}
\end{align}
Consequently,
\begin{equation}
 \sum_{H\in\mathfrak F_C(J)}
 {\gamma_{F\cup J}(E\cup H)\over\gamma_F(E)}
 =
 \left({B_C+B_J\over B_F+B_J}\right)^t
 \left({B_F\over B_F+B_J}\right)^k
 \le1.                                                     \label{eq:s6v3-gamma-summed}
\end{equation}
\end{lemma}

\begin{proof}
In the exact forest coefficient, the component factor \(B_C\)
is replaced by \(B_C+B_J\), the \(t\) new edges and \(t\) new
denominator marks give the first quotient, and the total-mark
denominator changes as displayed.  This proves
\eqref{eq:s6v3-gamma-rooted-forest}.

Sum \(H\) by \Cref{lem:s6v3-rooted-polynomial}.  Division by
\(\prod_{j\in J}b_j\) gives
\(B_C(B_C+B_J)^{t-1}\), which cancels the old component
denominator and produces \((B_C+B_J)^t\).  Rearrangement gives
\eqref{eq:s6v3-gamma-summed}; both ratios are at most one.
\end{proof}

\section{The tree-valued reserve estimate}
\label{sec:s6v3-reserve}

\begin{theorem}[Rooted-forest accretion with its tree weight
retained]
\label{thm:s6v3-tree-valued}
For a fixed new label set \(J\),
\begin{align}
 &\sum_{H\in\mathfrak F_C(J)}
 \gamma_{F\cup J}(E\cup H)\mathfrak r(F\cup J)
 \notag\\
 &\quad\le
 K^t\,{B_C+B_J\over B_C}\,
 \Phi_C(J)\,
 \gamma_F(E)\mathfrak r(F).                               \label{eq:s6v3-tree-valued}
\end{align}
The more precise right side also contains the factor
\((B_F/(B_F+B_J))^k\le1\).
\end{theorem}

\begin{proof}
Sum \eqref{eq:s6v3-gamma-rooted-forest} in \(H\) but retain
\(\Phi_C(J)\).  Multiply by the reserve ratio from
\Cref{lem:s6v2-reserve-ratio}:
\[
 {\mathfrak r(F\cup J)\over\mathfrak r(F)}
 \le
 \{K(B_F+B_J)\}^{t}\prod_{j\in J}b_j.
\]
The factors \(\prod_{j\in J}b_j\) and
\((B_F+B_J)^t\) cancel the corresponding denominators in the
gamma ratio, leaving
\[
 K^t\,{B_C+B_J\over B_C}\,
 \Phi_C(J)
 \left({B_F\over B_F+B_J}\right)^k.
\]
Drop only the last factor.
\end{proof}

\begin{assessment}[The factor is a component update]
\label{ass:s6v3-component-update}
The ratio \((B_C+B_J)/B_C\) in
\eqref{eq:s6v3-tree-valued} is not a new activity.  In the exact
formula for \(\gamma_F(E)\), \(B_C\) is already one numerator
factor.  Multiplication replaces that factor by the mass of the
enlarged component \(C\cup J\).  It must be updated, not bounded
by \(B/B_C\).
\end{assessment}

\begin{lemma}[Component updates telescope]
\label{lem:s6v3-telescope}
Suppose disjoint sets \(J_1,\ldots,J_q\) are attached
successively to the same evolving component,
\[
 C_\nu=C_{\nu-1}\mathbin{\dot\cup}J_\nu.
\]
Then
\begin{equation}
 \prod_{\nu=1}^q{B_{C_\nu}\over B_{C_{\nu-1}}}
 ={B_{C_q}\over B_{C_0}}.                                 \label{eq:s6v3-telescope}
\end{equation}
When multiplied by the original forest coefficient, the
denominator \(B_{C_0}\) cancels its component numerator and is
replaced by \(B_{C_q}\).
\end{lemma}

\begin{proof}
The product telescopes.  The second assertion is the component
factor in the exact gamma formula.
\end{proof}

\section{All rooted shapes and partitions are one rooted forest}
\label{sec:s6v3-gluing}

\begin{lemma}[Union of rooted accretions]
\label{lem:s6v3-union}
Let every vertex of a set \(Y\) be assigned to one child tree
which meets the old forced set \(F\) in exactly one root.  The
union of all child trees is a forest rooted at \(F\), and every
forest rooted at \(F\) arises uniquely in this way from its
connected components.
\end{lemma}

\begin{proof}
Distinct child trees have disjoint nonroot vertex sets and meet
the root set at one vertex each, so their union is acyclic and
every component contains exactly one root.  Conversely, take the
connected components of a rooted forest.
\end{proof}

\begin{theorem}[Tree-shape part of RCTAC at all arities]
\label{thm:s6v3-RCTAC-shapes}
Fix the set \(Y\) of labels first accreted at one hierarchy node
and fix the old forced forest \(E\).  The complete sum over:
\begin{enumerate}[label=\textup{(\roman*)},leftmargin=3.5em]
\item the partition of \(Y\) into rooted child components;
\item the root in the appropriate old component;
\item every rooted tree shape; and
\item every order in which components attached to the same old
      component are revealed
\end{enumerate}
has no Bell or macroscopic-arity loss.  It is bounded by \(K^{|Y|}\)
times the exact weighted forest polynomial rooted at the old
forced set, with each old component mass replaced by its final
mass.
\end{theorem}

\begin{proof}
By \Cref{lem:s6v3-union}, the apparent set partition and all
rooted tree shapes are exactly one rooted forest on \(F\cup Y\).
Sum it with the weighted rooted-forest identity, component by
component, as in \Cref{lem:s6v3-rooted-polynomial}.  Apply
\Cref{thm:s6v3-tree-valued} to each final component rather than
to each reveal order.  All numerical bases contribute once per
new vertex.  By \Cref{lem:s6v3-telescope}, the component ratios
replace existing numerator factors and introduce no inverse
mark.  Thus neither the partition nor its ordering creates an
additional fibre.
\end{proof}

\begin{corollary}[MRCTC is closed in tree-valued form]
\label{cor:s6v3-MRCTC}
The macroscopic rooted child-tree sector of V2 is bounded with
the required \(K^m\) base when its rooted forest is retained as
part of the global Cayley output.  Hence the restriction
\(|J|\le m/\log(m+e)\) is unnecessary for the tree-shape and
reserve summation.
\end{corollary}

\begin{proof}
Apply \Cref{thm:s6v3-RCTAC-shapes}.  Since the new label sets are
disjoint, the total base is
\(K^{\sum|J|}\le K^m\).
\end{proof}

\begin{correction}[The scalar MRCTC target is retired]
\label{corr:s6v3-MRCTC}
The scalar formulation of MRCTC in V2 is stronger than needed
and is retired.  Its proposed minimum-edge gain is insufficient
by \Cref{lem:s6v3-minimum-gain}.  The correct output is the
tree-valued closure \Cref{cor:s6v3-MRCTC}, which preserves the
marks required by the later global Cayley sum.
\end{correction}

\section{The remaining first-accretion allocation}
\label{sec:s6v3-allocation}

\begin{definition}[Vertex first-layer allocation card, VFLAC]
\label{def:s6v3-VFLAC}
VFLAC is the assertion that, in the complete nested ultrametric
expansion, every vertex which becomes forced is assigned to its
unique highest child edge, and that the joint sum over these
first-accretion hierarchy nodes, with all minimum-edge Jacobians
and beta factors retained, has fibre at most \(K^m\).
The rooted forest at each node is kept as tree-valued output and
is bounded by \Cref{thm:s6v3-RCTAC-shapes}.
\end{definition}

\begin{proposition}[VFLAC and the floating row prove VAFPC]
\label{prop:s6v3-VFLAC}
VFLAC together with the row of child components disjoint from the
old forced set proves VAFPC.
\end{proposition}

\begin{proof}
Every vertex-accreting component either meets the old forced
forest or is floating.  For a rooted component, choose the
highest hierarchy node at which one of its edges first makes each
new vertex forced.  VFLAC pays this unique allocation, and
\Cref{thm:s6v3-RCTAC-shapes} pays all rooted shapes at the node.
The stated floating row pays the complement.  Fixed-vertex
completions were closed in f5-V100, so the three rows exhaust
VAFPC.
\end{proof}

\begin{programme}[V4: first-accretion order statistics]
\label{prog:s6v3-V4}
\begin{enumerate}[leftmargin=3em]
\item Encode each new vertex by the highest edge level at which
      it is incident to the retained child forest.
\item Disintegrate all outer-tree parameters by their complete
      ordering, not only by the minimum.
\item Prove that the ordering simplex and the unique
      first-accretion map cancel the hierarchy-node census.
\item Insert the tree-valued rooted-forest output of V3 before
      summing thermal blocks.
\item Treat a floating child component by its first top-level
      connection and retain the corresponding power of
      \(\lambda\).
\end{enumerate}
\end{programme}

\begin{firewall}[Claim boundary after sixth-series V3]
\label{fire:s6v3-boundary}
V3 proves that the first minimum-edge integral gives no
factorial suppression, derives the exact gamma ratio for a rooted
forest attached to one old component, proves the tree-valued
reserve estimate, proves component-mass updates telescope, and
closes the tree-shape/partition part of RCTAC at every arity.  It
retires the unnecessary scalar MRCTC target and isolates VFLAC.

V3 does not prove VFLAC, the floating-component row, VAFPC in
full, FPC, ULCC, full IGMC, SISC, the simple high-excess row,
arbitrary repeated transversal supports, full TOCC, full TCNC,
cyclic TIC, HLTC, full MCWC, OCRC, SCNGC, WPSC, GRLTC, LZTC,
PLRC, WSDC, BRSC, HWSFC, UGCGC, thermal connected WPRAC,
all-arity RGSC, the raw Wilson aggregate strong card, all-order
strong MTA, WUFC, CPM, cutoff removal, reflection positivity,
Osterwalder--Schrader reconstruction, a continuum Yang--Mills
measure, or a positive Hamiltonian gap.  The full Yang--Mills
mass gap has not been proved.
\end{firewall}

\begin{theorem}[Sixth-series V3 result ledger]
\label{thm:s6v3-results}
V3:
\begin{enumerate}[label=\textup{(V3.\arabic*)},leftmargin=5.5em]
\item proves the scalar minimum-edge gain is insufficient;
\item proves the exact rooted-forest polynomial;
\item proves the exact rooted-forest gamma ratio;
\item proves the tree-valued reserve estimate;
\item proves component updates telescope;
\item proves rooted accretions glue to one rooted forest;
\item closes the tree-shape part of RCTAC at all arities;
\item closes MRCTC in tree-valued form; and
\item formulates VFLAC.
\end{enumerate}
\end{theorem}

\begin{proof}
Use
\Cref{lem:s6v3-minimum-gain,
lem:s6v3-rooted-polynomial,
lem:s6v3-gamma-rooted-forest,
thm:s6v3-tree-valued,
lem:s6v3-telescope,
lem:s6v3-union,
thm:s6v3-RCTAC-shapes,
cor:s6v3-MRCTC,
prop:s6v3-VFLAC}.
\end{proof}

\paragraph{Parent-version record.}
V3 was formed from
\texttt{Renewal\_Yang\_Mills\_Mass\_Gap\_Research\_Notes\_V2.tex}.
The V2 parent had \(977\) logical source lines, \(31\,916\) bytes,
and SHA--256
\[
 \texttt{429b940ed98b065598e89b8ff3710f34d12ed3541ae4418d766123b3d39da595}.
\]
The next compulsory companion audit is V5.

\clearpage
\chapter{Sixth-series V4: complete edge order and canonical
first accretion}
\label{ch:s6v4-order}

\section{Complete order statistics of the BKAR tree}
\label{sec:s6v4-order}

No companion audit is due at V4.

Let \(T_0\) be a tree on \(m\ge2\) labels and put
\[
 n:=|E(T_0)|=m-1.
\]
For a bijection
\(\pi:\{1,\ldots,n\}\to E(T_0)\), set
\[
 \Delta_\pi
 :=
 \left\{\boldsymbol u\in(0,1)^n:
 u_{\pi(1)}<\cdots<u_{\pi(n)}\right\}.
\]

\begin{lemma}[Normalized complete-order disintegration]
\label{lem:s6v4-order-simplex}
For every integrable \(\Psi:[0,1]^n\to\mathbb C\),
\begin{equation}
 \int_{[0,1]^n}\Psi(\boldsymbol u)\,d\boldsymbol u
 =
 \sum_{\pi\in S_n}
 \int_{\Delta_\pi}\Psi(\boldsymbol u)\,d\boldsymbol u.
                                                               \label{eq:s6v4-order-disintegration}
\end{equation}
Moreover,
\begin{equation}
 |\Delta_\pi|={1\over n!},
 \qquad
 \sum_{\pi\in S_n}|\Delta_\pi|=1.                              \label{eq:s6v4-order-mass}
\end{equation}
Equivalently, if \(\mathbb E_\pi\) denotes expectation for
Lebesgue measure conditioned on \(\Delta_\pi\), then
\begin{equation}
 \int_{[0,1]^n}\Psi(\boldsymbol u)\,d\boldsymbol u
 =
 {1\over n!}\sum_{\pi\in S_n}\mathbb E_\pi\Psi.                \label{eq:s6v4-order-average}
\end{equation}
\end{lemma}

\begin{proof}
The open simplices \(\Delta_\pi\) are disjoint and cover the
cube except for the null set on which two coordinates agree.
This proves \eqref{eq:s6v4-order-disintegration}.

On \(\Delta_\pi\), introduce spacings
\[
 a_0=u_{\pi(1)},\qquad
 a_r=u_{\pi(r+1)}-u_{\pi(r)}\ (1\le r<n),\qquad
 a_n=1-u_{\pi(n)}.
\]
They are positive, have sum one, and the triangular coordinate
map has Jacobian one.  Hence \(\Delta_\pi\) has the volume of the
standard \(n\)-simplex, namely \(1/n!\).  Summing its \(n!\)
permutations proves \eqref{eq:s6v4-order-mass}; normalization on
each simplex gives \eqref{eq:s6v4-order-average}.
\end{proof}

\begin{assessment}[What complete ordering can and cannot pay]
\label{ass:s6v4-order}
The order-simplex volume pays the \(n!\) possible strict edge
orders exactly.  It gives no additional factorial suppression
for vertex sets, tree shapes, or set partitions.  Those objects
must still be summed by the weighted rooted-forest identities of
V3.  The gain here is that an order-dependent allocation can be
made deterministically and then integrated without an extra
hierarchy census.
\end{assessment}

\section{The order determines the entire hierarchy}
\label{sec:s6v4-hierarchy}

\begin{definition}[Deletion hierarchy]
\label{def:s6v4-hierarchy}
Fix \(\boldsymbol u\in\Delta_\pi\).  Starting from \(T_0\),
delete the edges in the order
\[
 \pi(1),\pi(2),\ldots,\pi(n).
\]
The component containing a still undeleted edge is split when
the first edge of that component in the order \(\pi\) is
deleted.  The resulting binary tree of nested vertex sets is
denoted by \({\cal H}_\pi(T_0)\).
\end{definition}

\begin{lemma}[Equivalence with the recursive affine hierarchy]
\label{lem:s6v4-hierarchy}
The hierarchy \({\cal H}_\pi(T_0)\) is exactly the hierarchy
obtained by recursively applying the minimum-edge decomposition
of f5-V100.  At every child node, the affine rescaling of its
remaining edge parameters preserves their relative order.
\end{lemma}

\begin{proof}
At the root, \(\pi(1)\) is the unique minimum edge and its
deletion creates the two components in the minimum-edge
decomposition.  In either component \(Q\), the next recursive
minimum is the first member of \(E(T_0[Q])\) occurring in
\(\pi\).  The rescaling
\[
 v_e={u_e-\lambda_Q\over1-\lambda_Q}
\]
is strictly increasing in \(u_e\), so it does not change the
relative order inside \(Q\).  Induction on the number of edges
of \(Q\) proves the claim.
\end{proof}

\begin{definition}[Bottleneck owner]
\label{def:s6v4-owner}
For \(i\ne j\), define
\begin{equation}
 \beta_\pi(i,j)
 :=
 \mathop{\rm argmin}_{e\in{\rm path}_{T_0}(i,j)}u_e.
                                                               \label{eq:s6v4-owner}
\end{equation}
This is well defined on every order simplex.
\end{definition}

\begin{lemma}[Unique owner and recursive compatibility]
\label{lem:s6v4-owner}
For every pair \(i\ne j\),
\begin{equation}
 x^{T_0}_{ij}(\boldsymbol u)
 =
 u_{\beta_\pi(i,j)}.                                         \label{eq:s6v4-owned-minimum}
\end{equation}
The edge \(\beta_\pi(i,j)\) is the unique deletion edge at which
\(i\) and \(j\) first enter different hierarchy blocks.  If
\(i,j\) lie in one child node \(Q\), computing the owner in the
affinely rescaled subtree \(T_0[Q]\) gives the same original
edge.
\end{lemma}

\begin{proof}
The first equality is the definition of the BKAR path minimum.
Before \(\beta_\pi(i,j)\) is deleted, every deleted edge on the
\(i\)-to-\(j\) path would have had smaller parameter, contrary
to its minimality; after it is deleted, that path is broken.
Thus it is the first separating edge.  The last assertion follows
because the affine rescaling is strictly increasing and the
entire \(i\)-to-\(j\) path lies in \(T_0[Q]\).
\end{proof}

\begin{corollary}[Shortcut ownership has no choice]
\label{cor:s6v4-shortcut-owner}
Every completion edge \(h=\{i,j\}\), whether or not it is an
edge of \(T_0\), has one hierarchy owner
\(\beta_\pi(h)\).  Summing over possible owners is therefore
spurious: the owner is a function of \((T_0,\pi,h)\).
\end{corollary}

\begin{proof}
Apply \Cref{lem:s6v4-owner} to the endpoints of \(h\).
\end{proof}

\section{Canonical first accretion}
\label{sec:s6v4-first}

Consider one complete monomial of the recursively iterated
one-layer forest expansion.  At a hierarchy node \(Q\), let
\(F_Q\) be the forced set received from its ancestors and let
\(H_Q\) be its selected child-layer forest.  Forced sets are
passed monotonically to descendants.

\begin{lemma}[First-accretion nodes partition the new labels]
\label{lem:s6v4-first-node}
For every label \(v\) which becomes forced, there is a unique
highest hierarchy node \(q(v)\) at which
\[
 v\in V(H_{q(v)})\setminus F_{q(v)}.
\]
Consequently,
\begin{equation}
 Y_Q:=\{v:q(v)=Q\}                                         \label{eq:s6v4-YQ}
\end{equation}
are pairwise disjoint and their union is the set of all newly
forced labels.
\end{lemma}

\begin{proof}
The hierarchy blocks containing a fixed label form a chain.
Along that chain, forced sets are increasing.  Hence the nodes
at which \(v\) is incident to a selected edge have a unique
highest member, and after that node \(v\) belongs to every
descendant forced set.  This proves uniqueness.  The assertions
about the sets \(Y_Q\) are immediate.
\end{proof}

\begin{lemma}[Canonical rooted accretion card]
\label{lem:s6v4-accretion-card}
Suppose every accreting component at \(Q\) contains exactly one
component of the old forced forest.  Orient it toward that old
component, contracting the old component to one root.  Every
\(v\in Y_Q\) then has:
\begin{enumerate}[label=\textup{(\roman*)},leftmargin=3.5em]
\item one parent edge \(p(v)\) in the oriented rooted forest;
\item one first-accretion node \(q(v)\); and
\item one outer-tree bottleneck owner \(\beta_\pi(p(v))\).
\end{enumerate}
Thus
\begin{equation}
 {\cal A}(v)
 :=
 \bigl(q(v),p(v),\beta_\pi(p(v))\bigr)                    \label{eq:s6v4-card}
\end{equation}
is a function of the expanded monomial and its ordered BKAR
tree, not an additional combinatorial choice.
\end{lemma}

\begin{proof}
After the old component is contracted, each stated component is
a tree with one root.  Its orientation toward the root gives
each nonroot vertex exactly one parent edge.
\Cref{lem:s6v4-first-node,lem:s6v4-owner} give the other two
unique entries.
\end{proof}

\section{The recursive layer assignments still have unit mass}
\label{sec:s6v4-unit}

\begin{lemma}[Iterated Bernoulli contraction]
\label{lem:s6v4-iterated}
Condition at every hierarchy node on a sampled Cayley tree
containing its already owned forest.  Let \(R_Q\) be the
remaining Cayley-tree edges at node \(Q\), and let
\[
 w_Q(H)
 :=
 \lambda_Q^{|R_Q|-|H|}
 (1-\lambda_Q)^{|H|}
 \boldsymbol z_Q^H,
 \qquad H\subseteq R_Q.
\]
If every continuation below a child node has total mass at most
one, then the total mass below \(Q\), including the sum over
\(H\subseteq R_Q\), is at most one.  Consequently the complete
recursive layer-assignment mass is at most one.
\end{lemma}

\begin{proof}
Discarding continuation factors, which lie in \([0,1]\), can
only enlarge the nonnegative sum.  Then
\begin{align*}
 \sum_{H\subseteq R_Q}w_Q(H)
 &=
 \prod_{e\in R_Q}
 \{\lambda_Q+(1-\lambda_Q)z_{Q,e}\}\\
 &\le1.
\end{align*}
This is the one-layer Bernoulli identity.  Induction upward from
the terminal hierarchy nodes proves the recursive assertion.
\end{proof}

\begin{lemma}[Compatibility with Cayley disintegration]
\label{lem:s6v4-Cayley-compatibility}
The contraction in \Cref{lem:s6v4-iterated} may be applied after
the rooted-forest summation of V3.  More precisely, every factor
\(\gamma_A(G)\) can first be written as
\[
 \gamma_A(G)
 =
 \sum_{\substack{T\in\mathfrak T(A)\\G\subseteq T}}\mu_A(T),
\]
all layer choices can be summed inside the fixed \(T\), and the
remaining positive sum in \(T\) reconstructs the original
Cayley probability.
\end{lemma}

\begin{proof}
This is a finite interchange of nonnegative sums.  For fixed
\(T\), every subset of \(T\setminus G\) is a forest and is
summed by the product in \Cref{lem:s6v4-iterated}.  Summing the
surviving \(\mu_A(T)\)'s gives the displayed probability.
\end{proof}

\section{Closure of the singly rooted hierarchy allocation}
\label{sec:s6v4-VFLAC}

\begin{theorem}[Rooted VFLAC]
\label{thm:s6v4-rooted-VFLAC}
In the sector in which every vertex-accreting component meets
exactly one component of the pre-node forced forest, the complete
sum over:
\begin{enumerate}[label=\textup{(\roman*)},leftmargin=3.5em]
\item all first-accretion hierarchy nodes;
\item all strict outer-tree edge orders;
\item all rooted child-component partitions and shapes; and
\item all recursive layer assignments
\end{enumerate}
has total fibre at most \(K^m\).  All minimum-edge Jacobians,
beta factors, and Cayley probabilities may be retained until the
corresponding normalized sum is performed.
\end{theorem}

\begin{proof}
Condition first on \(T_0\), an order simplex \(\Delta_\pi\), and
the Cayley trees in the probability disintegrations.  By
\Cref{lem:s6v4-accretion-card}, every newly forced label has one
canonical allocation; there is no sum over allocation maps.

At a node \(Q\), apply the tree-valued rooted-forest closure
\Cref{thm:s6v3-RCTAC-shapes} to the set \(Y_Q\).  It costs
\(K^{|Y_Q|}\), while its component-mass quotient updates the
existing component numerator.  The updates telescope down each
hierarchy chain by \Cref{lem:s6v3-telescope}.  Since the \(Y_Q\)
partition the new labels,
\[
 \prod_QK^{|Y_Q|}
 =
 K^{\sum_Q|Y_Q|}
 \le K^m.
\]

Now sum all unused layer assignments within their sampled
Cayley trees.  Their total mass is at most one by
\Cref{lem:s6v4-iterated,lem:s6v4-Cayley-compatibility}.
Finally integrate the ordered BKAR parameters.  The sum of all
order-simplex masses is one by
\eqref{eq:s6v4-order-mass}.  No depth, node, permutation, Bell,
or rooted-tree factor remains.
\end{proof}

\begin{corollary}[The singly rooted VAFPC row is closed]
\label{cor:s6v4-rooted-row}
The VAFPC row in which every new-vertex component meets exactly
one old forced component satisfies the required \(K^m\) strong
base.
\end{corollary}

\begin{proof}
The tree-shape and component-partition sums are V3, and their
hierarchy allocation is
\Cref{thm:s6v4-rooted-VFLAC}.  These are precisely the unsummed
parts of the singly rooted row.
\end{proof}

\section{A necessary correction to the V3 exhaustion statement}
\label{sec:s6v4-correction}

\begin{correction}[There is also a multi-anchor row]
\label{corr:s6v4-multi-anchor}
The exhaustion claim in
\Cref{prop:s6v3-VFLAC} is valid when the old forced forest is
connected, or when every accreting component meets at most one
old forced component.  For a general disconnected old forest,
there is a third vertex-accreting class:
\[
 \text{a component containing new labels and at least two
 old forced components}.
\]
Thus full VAFPC requires the singly rooted row, the floating row,
and this multi-anchor coalescing row, in addition to the
fixed-vertex row.
\end{correction}

\begin{proof}
Let \(E\) be the pre-layer forest and \(G=E\cup H\).  Any
component of \(G\) containing a new label meets respectively
zero, one, or at least two components of \(E\).  These cases are
disjoint and exhaustive.  The first is floating, the second is
the singly rooted row proved above, and the third is not a
rooted forest with exactly one old root.  The fixed-vertex row
contains no new label and is separate.
\end{proof}

\begin{assessment}[Why the correction matters]
\label{ass:s6v4-correction}
Replacing several old components by one formal root before
summing would erase the component-coalescence defect and could
count the same global tree several times.  The archived
multi-anchor ratios and the Kruskal quotient normalization must
be inserted before that contraction.  No later implication is
invoked from the uncorrected exhaustion statement.
\end{assessment}

\section{Canonical bridges for floating components}
\label{sec:s6v4-floating}

Let \(G=E\cup H\) be a forest on its forced vertex set.  Call a
component of \(G\) \emph{old} if it contains a component of
\(E\), and \emph{floating} otherwise.  Write \(r\) and \(q\)
for the numbers of old and floating components.  Thus
\[
 c(G)=r+q.
\]

\begin{lemma}[Quotient-tree rooting of floating components]
\label{lem:s6v4-floating-rooting}
Fix a spanning tree \(T\supseteq G\) contributing to
\(\gamma_{V(G)}(G)\), and contract every component of \(G\).
The quotient \(\overline T\) is a tree.  Choose one old quotient
vertex \(C_*\) deterministically and orient \(\overline T\)
toward \(C_*\).  For each floating quotient vertex \(D\), let
\(a_D(T,G)\) be its parent edge.  Then:
\begin{enumerate}[label=\textup{(\roman*)},leftmargin=3.5em]
\item the \(q\) edges \(a_D(T,G)\) are distinct members of
      \(T\setminus G\);
\item adding their lifts to \(G\) gives a forest every component
      of which contains exactly one old component; and
\item the remaining \(r-1\) quotient-tree edges are precisely
      the parent edges of the old quotient vertices other than
      \(C_*\).
\end{enumerate}
\end{lemma}

\begin{proof}
Contracting connected subtrees of a tree preserves connectedness
and acyclicity, so \(\overline T\) is a tree.  Every nonroot
quotient vertex has one parent edge, and different vertices have
different parent edges.  Select the parents of the \(q\)
floating vertices.

Following selected parent edges from a floating vertex either
reaches an old vertex, where selection stops, or continues
through floating vertices and eventually reaches an old vertex.
Hence every selected-edge component contains an old vertex.
It cannot contain two: on the unique path between two old
vertices, the edge directed out of the old vertex farther from
\(C_*\) is a parent edge of an old, not a selected parent of a
floating vertex.  This proves (ii).  A tree on \(r+q\) quotient
vertices has \(r+q-1\) edges; removing the \(q\) selected parents
leaves \(r-1\), and the orientation identifies them as stated.
\end{proof}

\begin{corollary}[Exact beta split for the floating row]
\label{cor:s6v4-floating-beta}
For the fixed pair \((T,G)\), the top-layer base factor splits
canonically as
\begin{equation}
 \lambda^{c(G)-1}
 =
 \lambda^q\lambda^{r-1}.                                  \label{eq:s6v4-floating-beta}
\end{equation}
The first \(q\) copies of \(\lambda\) own the floating-parent
bridges in \Cref{lem:s6v4-floating-rooting}; the other \(r-1\)
copies own the old-component quotient edges.  This assignment
has no \(q!\) or \(r^q\) fibre after conditioning on \(T\).
\end{corollary}

\begin{proof}
The exponent identity is \(c(G)-1=q+r-1\).  The edge ownership
is the partition of the quotient-tree edges proved in
\Cref{lem:s6v4-floating-rooting}.
\end{proof}

\begin{assessment}[Exact remaining floating issue]
\label{ass:s6v4-floating-remainder}
The canonical bridges convert the floating pieces to singly
rooted pieces inside each sampled Cayley tree.  What remains is
to interchange the sum over \((G,T)\) with this
\(T\)-dependent rooting without replacing the restricted
parent event by the larger event \(T\supseteq G\cup A\).
That replacement can introduce an uncontrolled inverse fibre.
The next step must retain the quotient-tree probability, apply
the Kruskal normalization to the remaining \(r-1\) old edges,
and only then use the rooted VFLAC theorem.
\end{assessment}

\section{Next acquisition}
\label{sec:s6v4-next}

\begin{programme}[V5: audit and quotient-tree projection]
\label{prog:s6v4-V5}
\begin{enumerate}[leftmargin=3em]
\item Perform the compulsory audit of the newest active
      Standard-Model--Einstein and Navier--Stokes notes.
\item Disintegrate the conditional Cayley probability by the
      contracted quotient tree of
      \Cref{lem:s6v4-floating-rooting}.
\item Sum the \(T\)-dependent floating-parent bridges before
      forgetting \(T\).
\item Combine the remaining old-component edges with the
      archived Kruskal component-coalescence normalization.
\item Test the same quotient formula on the multi-anchor row
      identified in \Cref{corr:s6v4-multi-anchor}.
\end{enumerate}
\end{programme}

\begin{firewall}[Claim boundary after sixth-series V4]
\label{fire:s6v4-boundary}
V4 proves normalized complete-order disintegration of the BKAR
cube, identifies the exact recursive deletion hierarchy, proves
unique bottleneck ownership of every shortcut edge, constructs
the canonical first-accretion card, and proves the recursive
layer assignments have unit mass.  Together with V3, this closes
VFLAC and the full hierarchy allocation for the singly rooted
vertex-accretion row.  It corrects the V3 exhaustion statement
by isolating the multi-anchor coalescing row.  It also gives an
exact quotient-tree rooting and beta split for floating
components.

V4 does not prove the floating-component projection, the
multi-anchor coalescing row, full VAFPC, FPC, ULCC, full IGMC,
SISC, the simple high-excess row, arbitrary repeated transversal
supports, full TOCC, full TCNC, cyclic TIC, HLTC, full MCWC,
OCRC, SCNGC, WPSC, GRLTC, LZTC, PLRC, WSDC, BRSC, HWSFC,
UGCGC, thermal connected WPRAC, all-arity RGSC, the raw Wilson
aggregate strong card, all-order strong MTA, WUFC, CPM, cutoff
removal, reflection positivity, Osterwalder--Schrader
reconstruction, a continuum Yang--Mills measure, or a positive
Hamiltonian gap.  The full Yang--Mills mass gap has not been
proved.
\end{firewall}

\begin{theorem}[Sixth-series V4 result ledger]
\label{thm:s6v4-results}
V4:
\begin{enumerate}[label=\textup{(V4.\arabic*)},leftmargin=5.5em]
\item proves normalized complete-order disintegration;
\item proves equivalence of the order and affine hierarchies;
\item proves unique recursive bottleneck ownership;
\item proves the canonical first-accretion card;
\item proves iterated Bernoulli contraction;
\item proves Cayley-disintegration compatibility;
\item closes rooted VFLAC and the singly rooted VAFPC row;
\item isolates the previously omitted multi-anchor row; and
\item proves canonical quotient-tree rooting and beta splitting
      for floating components.
\end{enumerate}
\end{theorem}

\begin{proof}
Use
\Cref{lem:s6v4-order-simplex,
lem:s6v4-hierarchy,
lem:s6v4-owner,
lem:s6v4-first-node,
lem:s6v4-accretion-card,
lem:s6v4-iterated,
lem:s6v4-Cayley-compatibility,
thm:s6v4-rooted-VFLAC,
corr:s6v4-multi-anchor,
lem:s6v4-floating-rooting,
cor:s6v4-floating-beta}.
\end{proof}

\paragraph{Parent-version record.}
V4 was formed from
\texttt{Renewal\_Yang\_Mills\_Mass\_Gap\_Research\_Notes\_V3.tex}.
The V3 parent had \(1\,338\) logical source lines, \(44\,820\)
bytes, and SHA--256
\[
 \texttt{c98bf0ec6cb7010db9a6d76b3a54e530cda2952c7bd9d25bf9e9a194d7446caf}.
\]
The next compulsory companion audit is V5.

\clearpage
\chapter{Sixth-series V5: companion audit and conditional
quotient-tree projection}
\label{ch:s6v5-quotient}

\section{Compulsory companion audit}
\label{sec:s6v5-audit}

The companion directory changed while the audit began: the
Standard-Model--Einstein programme replaced V27 by V28.  The
settled files actually used below are:
\begin{center}
\small
\begin{tabular}{p{0.25\textwidth}p{0.68\textwidth}}
Standard Model--Einstein &
\texttt{Renewal\_SM\_Einstein\_Continuum\_Research\_Notes\_V28.tex},
\(468\,670\) bytes, \(11\,753\) logical source lines, SHA--256
\texttt{3AADB805D94D79F93060531FA54E9D74FF0C95EA6A7DC7F7CBF95C761E2591A0}.
\\ [0.4em]
Navier--Stokes &
\texttt{Renewal\_NS\_Stability\_Research\_Notes\_V25.tex},
\(241\,869\) bytes, \(4\,761\) logical source lines, SHA--256
\texttt{0AFECAA7C07D3D57491CE39A5D82B63FF8C37783F498E020A18C81935B09665F}.
\end{tabular}
\end{center}

\begin{audit}[Standard Model--Einstein V28]
\label{audit:s6v5-SM}
SM V28 proves the low physical momentum branch of a metric
overlap response, but leaves a separate high-branch card which
contains low--high as well as high--high divided differences.
Its proof explicitly warns that estimates of the two diagonal
branches do not control their mixed contact.

The structural lesson for the present quotient problem is exact:
the singly rooted output and the old-component coalescence output
cannot be bounded separately if child/base colours can cross the
two pieces.  One must condition on the complete Cayley tree, sum
the mixed colour assignment there, and only then split the tree
into its rooted and coalescing parts.  No Standard-Model
coefficient or operator estimate is imported.
\end{audit}

\begin{audit}[Navier--Stokes V25]
\label{audit:s6v5-NS}
NS V25 records two upstream corrections.  Concentration alone
does not give a spectral-variance floor, and the explicit floor
subsequently obtained is too small to close the desired push
budget without additional use of the equation.

The analogous warning here is that the scalar factor
\(\lambda^q\) attached to \(q\) floating components is not by
itself a useful macroscopic summability estimate.  The missing
information is the exact conditional Cayley equation.  V5
therefore does not estimate floating bridges by their number; it
derives their quotient-tree law.  No Navier--Stokes inequality is
imported.
\end{audit}

\begin{audit}[Direction selected by the V5 audit]
\label{audit:s6v5-direction}
The proof stays upstream of hierarchy iteration.  It conditions
simultaneously on the final forced set and the complete Cayley
tree, sums every child/base colour by its Bernoulli probability,
and only afterward orients the quotient tree.  This treats the
mixed bridge sector which would be lost by a separate endpoint
argument.
\end{audit}

\section{A full one-layer projection at fixed final vertex set}
\label{sec:s6v5-fixed-final}

Let \(E\) be a forest on \(F=V(E)\), let
\[
 r:=c(E)=|F|-|E|,
\]
and let \(J\) be disjoint from \(F\), with \(t=|J|\ge1\).
Put \(A:=F\mathbin{\dot\cup}J\).  For
\(0\le\boldsymbol z\le1\), define
\begin{align}
 {\cal P}_{E,J}(\lambda,\boldsymbol z)
 &:=
 \sum_{\substack{
 H\subseteq\binom A2\setminus E\\
 E\cup H\ {\rm a\ forest}\\
 V(E\cup H)=A}}
 \lambda^{c(E\cup H)-1}
 (1-\lambda)^{|H|}
 \boldsymbol z^H\,
 \gamma_A(E\cup H).                                      \label{eq:s6v5-P}
\end{align}
The vertex condition means that every member of \(J\) is incident
to \(H\); vertices of \(F\) are already forced by \(E\).

\begin{theorem}[Fixed-final-set forest projection]
\label{thm:s6v5-fixed-final}
For every \(0\le\lambda\le1\) and
\(0\le\boldsymbol z\le1\),
\begin{equation}
 {\cal P}_{E,J}(\lambda,\boldsymbol z)
 \le \gamma_A(E).                                        \label{eq:s6v5-fixed-final}
\end{equation}
This includes, in one estimate, completion components meeting
zero, one, or several old components.
\end{theorem}

\begin{proof}
Disintegrate every Cayley probability:
\[
 \gamma_A(E\cup H)
 =
 \sum_{\substack{T\in\mathfrak T(A)\\E\cup H\subseteq T}}
 \mu_A(T),
\]
and interchange the finite positive sums.  Fix
\(T\supseteq E\).  Since \(T\) is a tree,
\[
 |T\setminus E|
 =
 |A|-1-|E|
 =
 r+t-1.
\]
Every \(H\subseteq T\setminus E\) is a forest and
\[
 c(E\cup H)-1=r+t-1-|H|
 =|T\setminus E|-|H|.
\]
Therefore the inner sum, even after the requirement that \(H\)
cover \(J\) is discarded, is at most
\begin{align*}
 \sum_{H\subseteq T\setminus E}
 \lambda^{|T\setminus E|-|H|}
 (1-\lambda)^{|H|}\boldsymbol z^H
 &=
 \prod_{e\in T\setminus E}
 \{\lambda+(1-\lambda)z_e\}\\
 &\le1.
\end{align*}
The remaining sum is
\(\sum_{T\supseteq E}\mu_A(T)=\gamma_A(E)\).
\end{proof}

\begin{corollary}[Exact isolated-extension ratio]
\label{cor:s6v5-isolated-ratio}
If \(k=|E|\), then
\begin{equation}
 {\gamma_A(E)\over\gamma_F(E)}
 =
 \left({B_F\over B_F+B_J}\right)^k.                       \label{eq:s6v5-isolated-ratio}
\end{equation}
\end{corollary}

\begin{proof}
Apply the exact forest coefficient to \(E\) first on \(F\) and
then on \(A\).  On \(A\), every \(j\in J\) is an isolated
component of the declared forest \(E\); its component factor
\(b_j\) cancels the new denominator mark \(b_j\).  Only the
total-mass denominator changes, giving the displayed ratio.
\end{proof}

\begin{warning}[The scalar fixed-set bound is not the final sum]
\label{warn:s6v5-scalar-J}
Multiplying \eqref{eq:s6v5-fixed-final} by the reserve and then
summing \(J\) as an unstructured set again gives, at best, an
\(\exp(KB^2)\) majorant.  Thus the right side of
\eqref{eq:s6v5-fixed-final} must be disintegrated into a
tree-valued quotient output before the \(J\)-sum is taken.
\end{warning}

\section{The conditional Cayley law on forest components}
\label{sec:s6v5-conditional}

Let \(C_1,\ldots,C_r\) be the components of \(E\), and set
\[
 a_p:=B_{C_p}\quad(1\le p\le r),\qquad
 {\cal Q}:=\{C_1,\ldots,C_r\}\mathbin{\dot\cup}J.
\]
Give \(C_p\) weight \(a_p\), give \(j\in J\) weight \(b_j\),
and write \(w_\alpha\) for either kind of quotient weight.
Then
\[
 \sum_{\alpha\in{\cal Q}}w_\alpha=B_A.
\]

\begin{theorem}[Conditional weighted quotient-tree law]
\label{thm:s6v5-quotient-law}
Condition the weighted Cayley tree on \(A\) to contain \(E\).
After every \(C_p\) is contracted, the quotient tree
\(\tau\in\mathfrak T({\cal Q})\) has probability
\begin{equation}
 \overline\mu_{\cal Q}(\tau)
 =
 {\displaystyle\prod_{\{\alpha,\beta\}\in\tau}
        w_\alpha w_\beta
  \over
  \displaystyle
  \left(\prod_{\alpha\in{\cal Q}}w_\alpha\right)
  B_A^{r+t-2}}.                                          \label{eq:s6v5-quotient-law}
\end{equation}
Conditioned further on \(\tau\), every quotient edge
\(\{C,D\}\) is lifted independently: the lift
\(\{i,j\}\), \(i\in C,j\in D\), has probability
\begin{equation}
 {b_ib_j\over B_CB_D}.                                    \label{eq:s6v5-lift-law}
\end{equation}
Here a singleton component \(j\) has mass \(B_{\{j\}}=b_j\).
\end{theorem}

\begin{proof}
Fix a quotient tree \(\tau\).  For one quotient edge
\(\{C,D\}\), summing its lifts gives
\[
 \sum_{i\in C,\ j\in D}b_ib_j=B_CB_D.
\]
Since \(\tau\) is a tree, lifts of distinct quotient edges may be
chosen independently; together with the fixed internal forest
\(E\), every lift is a tree on \(A\), and every
\(T\supseteq E\) is obtained once.

The product of the lift sums is the numerator in
\eqref{eq:s6v5-quotient-law}.  Its sum over \(\tau\) is, by the
weighted Cayley identity,
\[
 \left(\prod_{\alpha\in{\cal Q}}w_\alpha\right)
 B_A^{r+t-2}.
\]
This proves the normalized quotient law.  Dividing an individual
lift weight \(b_ib_j\) by its total \(B_CB_D\) proves
\eqref{eq:s6v5-lift-law}, and the product factorization gives
conditional independence.
\end{proof}

\begin{corollary}[No component-shape defect in the conditional
law]
\label{cor:s6v5-no-quotient-defect}
After conditioning on \(E\), all quotient-tree shapes and all
endpoint lifts have total mass one.  In particular, a raw factor
\(r^{t}\), a Cayley count on \(r+t\) quotient vertices, or a
separate endpoint-lift census is inadmissible.
\end{corollary}

\begin{proof}
Sum \eqref{eq:s6v5-quotient-law} and then
\eqref{eq:s6v5-lift-law}.
\end{proof}

\section{A measure-preserving orientation factorization}
\label{sec:s6v5-orientation}

Choose an auxiliary old root \(\rho\in\{1,\ldots,r\}\) with
\begin{equation}
 {\mathbb P}(\rho=p)={a_p\over B_F}.                       \label{eq:s6v5-root-choice}
\end{equation}
For a quotient tree \(\tau\), orient every edge toward \(C_\rho\).
Let \(R_\rho(\tau)\) consist of the parent edges whose tails are
new labels \(j\in J\), and let \(S_\rho(\tau)\) consist of the
parent edges whose tails are old quotient vertices
\(C_p\ne C_\rho\).

\begin{lemma}[Canonical quotient split]
\label{lem:s6v5-quotient-split}
For every \((\rho,\tau)\):
\begin{enumerate}[label=\textup{(\roman*)},leftmargin=3.5em]
\item \(R_\rho(\tau)\) has exactly \(t\) edges and is a forest
      rooted at the old quotient set
      \(\{C_1,\ldots,C_r\}\);
\item every component of \(R_\rho(\tau)\) contains exactly one
      old quotient vertex;
\item \(S_\rho(\tau)\) has exactly \(r-1\) edges; and
\item \(E\), the lifted \(R_\rho(\tau)\), and the lifted
      \(S_\rho(\tau)\) reconstruct the original tree uniquely.
\end{enumerate}
\end{lemma}

\begin{proof}
Every quotient vertex except \(C_\rho\) has one parent edge.
Select the parent edges of the \(t\) new vertices.  Following
them toward the root must first reach an old vertex, at which
selection stops, so every selected component contains an old
vertex.  It cannot contain two: on the path between distinct old
vertices, the edge leaving the old vertex farther from
\(C_\rho\) is the parent of an old vertex and belongs to
\(S_\rho(\tau)\).  This proves (i) and (ii).  The remaining
nonroot tails are the \(r-1\) old vertices other than
\(C_\rho\), proving (iii).  The two edge sets partition
\(\tau\), which proves reconstruction.
\end{proof}

Fix a rooted forest \(R\) of the kind in
\Cref{lem:s6v5-quotient-split}.  Let \(D_p\) be its component
rooted at \(C_p\), and set
\begin{equation}
 X_p:=\sum_{\alpha\in D_p}w_\alpha,
 \qquad
 \sum_{p=1}^rX_p=B_A.                                    \label{eq:s6v5-branch-mass}
\end{equation}

\begin{lemma}[Exact sum of the old-tail completion]
\label{lem:s6v5-old-tail}
For fixed \(R\) and fixed auxiliary root \(C_\rho\), the sum of
the edge-weight products over all compatible
\(S_\rho\) is
\begin{equation}
 \sum_{S_\rho\ {\rm compatible}}
 \prod_{\{u,v\}\in S_\rho}w_uw_v
 =
 \left(\prod_{p\ne\rho}a_p\right)
 X_\rho B_A^{r-2}.                                      \label{eq:s6v5-old-tail}
\end{equation}
After averaging \(\rho\) with
\eqref{eq:s6v5-root-choice}, this becomes
\begin{equation}
 {\prod_{p=1}^ra_p\over B_F}\,B_A^{r-1},                 \label{eq:s6v5-old-tail-average}
\end{equation}
which is independent of the branch partition
\((D_1,\ldots,D_r)\).
\end{lemma}

\begin{proof}
Contract every \(D_p\).  A compatible completion is a tree on
the \(r\) branch indices, oriented toward \(\rho\).  For its
edge \(p\to q\), the endpoint in the tail branch must be the old
quotient vertex \(C_p\), whereas its endpoint in \(D_q\) may be
any vertex.  Its summed lift weight is therefore \(a_pX_q\).
Factor \(\prod_{p\ne\rho}a_p\).  The rooted weighted Cayley
identity for the remaining parent weights gives
\[
 \sum_{\sigma\ {\rm toward}\ \rho}
 \prod_{p\ne\rho}X_{\operatorname{parent}_\sigma(p)}
 =
 X_\rho\left(\sum_pX_p\right)^{r-2}.
\]
This proves \eqref{eq:s6v5-old-tail}.  Multiplying by
\(a_\rho/B_F\), summing \(\rho\), and using
\(\sum_\rho X_\rho=B_A\) proves
\eqref{eq:s6v5-old-tail-average}.
\end{proof}

\begin{theorem}[Measure-preserving rooted/coalescing
factorization]
\label{thm:s6v5-orientation-factorization}
The auxiliary-root orientation is an exact disintegration of the
weighted quotient-tree polynomial:
\begin{align}
 &\sum_{\rho=1}^r{a_\rho\over B_F}
 \sum_{\tau\in\mathfrak T({\cal Q})}
 \prod_{\{\alpha,\beta\}\in\tau}w_\alpha w_\beta
 \notag\\
 &\qquad=
 \left\{
 \sum_{R\in\mathfrak F_{\{C_1,\ldots,C_r\}}({\cal Q})}
 \prod_{e\in R}w_e
 \right\}
 {\prod_{p=1}^ra_p\over B_F}B_A^{r-1}
 \notag\\
 &\qquad=
 \left(\prod_{p=1}^ra_p\right)
 \left(\prod_{j\in J}b_j\right)
 B_A^{r+t-2}.                                            \label{eq:s6v5-orientation-factorization}
\end{align}
Thus the singly rooted branches and the old-component
coalescence skeleton are not two independent gases; together
they are exactly one normalized quotient Cayley tree.
\end{theorem}

\begin{proof}
For fixed \((\rho,\tau)\), use the unique split of
\Cref{lem:s6v5-quotient-split}.  Fix \(R\), sum \(S_\rho\),
and average \(\rho\) by
\Cref{lem:s6v5-old-tail}.  The factor is independent of \(R\).
The weighted rooted-forest identity gives
\[
 \sum_R\prod_{e\in R}w_e
 =
 B_F\left(\prod_{j\in J}b_j\right)B_A^{t-1}.
\]
Multiplication by
\eqref{eq:s6v5-old-tail-average} gives the last line of
\eqref{eq:s6v5-orientation-factorization}, which is also the
weighted Cayley polynomial on the \(r+t\) quotient vertices.
\end{proof}

\section{The one-node vertex-accretion kernel}
\label{sec:s6v5-one-node}

\begin{theorem}[Topology-free one-node VAFPC kernel]
\label{thm:s6v5-one-node}
Fix \(E\), \(J\), and one ultrametric hierarchy node.  Before any
scalar sum over \(J\), the complete one-layer sum over every
completion forest \(H\) with final forced set \(F\cup J\) can be
organized as follows:
\begin{enumerate}[label=\textup{(\roman*)},leftmargin=3.5em]
\item condition on one Cayley tree \(T\supseteq E\);
\item sum all child/base colours \(H\subseteq T\setminus E\),
      with the coverage restriction, with total mass at most one;
\item contract \(E\) and apply
      \Cref{thm:s6v5-orientation-factorization};
\item retain one rooted forest containing exactly one parent edge
      for every member of \(J\); and
\item retain the complementary \(r-1\) edges as the exact
      old-component coalescence skeleton.
\end{enumerate}
The rooted part costs \(K^{|J|}\) after the reserve extension,
and the two output parts together have exactly the conditional
quotient-Cayley mass.  No additional floating, multi-anchor,
partition, endpoint-lift, or child/base-colour census occurs at
this node.
\end{theorem}

\begin{proof}
For fixed \(T\), the colour sum is the product in the proof of
\Cref{thm:s6v5-fixed-final}, hence at most one.  This implements
the mixed-sector lesson of the audit before any split is made.
The conditional law and all endpoint lifts are
\Cref{thm:s6v5-quotient-law}.  Split the quotient tree by the
auxiliary orientation and use
\Cref{thm:s6v5-orientation-factorization}; this is an equality,
not a raw forest bound.

The selected new-tail edges form one rooted forest spanning
exactly \(J\).  Apply the tree-valued reserve estimate and
rooted-forest sum of
\Cref{thm:s6v3-tree-valued,thm:s6v3-RCTAC-shapes}; it contributes
one numerical factor for each new vertex.  The complementary
old-tail edges were not discarded and reconstruct the same
quotient Cayley tree, so they create no second gas.
\end{proof}

\begin{corollary}[The local topology trichotomy is eliminated]
\label{cor:s6v5-trichotomy}
At one fixed hierarchy node and final forced set, the floating,
singly rooted, and multi-anchor classes of
\Cref{corr:s6v4-multi-anchor} need not be bounded separately.
They are the child/base colourings of the single oriented
quotient tree in \Cref{thm:s6v5-one-node}.
\end{corollary}

\begin{proof}
Each original \(H\) is included in the colour sum of the theorem.
Its topology affects which quotient-tree edges receive the child
colour, but the colour product was summed before orientation.
\end{proof}

\section{The remaining cross-node compatibility}
\label{sec:s6v5-QTC}

\begin{definition}[Quotient-transition compatibility card, QTC]
\label{def:s6v5-QTC}
QTC is the assertion that the oriented quotient outputs of
\Cref{thm:s6v5-one-node}, over all nested hierarchy nodes, may be
composed with the owned outer forest and the f5-V88
component-exchange normalization so that:
\begin{enumerate}[label=\textup{(\roman*)},leftmargin=3.5em]
\item the first-accretion sets \(Y_Q\) remain disjoint;
\item every old-tail coalescence edge is owned at its unique
      highest node;
\item no final global tree is produced by more than \(K^m\)
      oriented quotient histories; and
\item all conditional quotient-tree probabilities are retained
      until that history fibre is summed.
\end{enumerate}
\end{definition}

\begin{proposition}[QTC completes VAFPC]
\label{prop:s6v5-QTC}
QTC, together with \Cref{thm:s6v4-rooted-VFLAC} and
\Cref{thm:s6v5-one-node}, proves VAFPC.
\end{proposition}

\begin{proof}
The one-node theorem removes every local topology and colour
census.  Rooted VFLAC pays the unique first-node allocation of
the disjoint new-label sets by \(K^m\).  QTC pays the only data
not contained in those two results: the compatibility of the
old-tail quotient skeletons belonging to different hierarchy
nodes.  The fixed-vertex row was proved in f5-V100.  These rows
exhaust FPC's vertex-accreting sector.
\end{proof}

\begin{assessment}[Exact live bottleneck after V5]
\label{ass:s6v5-bottleneck}
The obstruction is no longer a floating-component count or a
multi-anchor tree count.  At one node both are normalized by the
conditional quotient Cayley law.  The remaining question is
whether the old-tail edges created at several nested nodes have
a canonical highest-node ownership compatible with the existing
V88 Kruskal exchange.  Bounding each node independently would
recreate the component defect; QTC must be proved as one global
tree statement.
\end{assessment}

\begin{programme}[V6: global quotient-history gluing]
\label{prog:s6v5-V6}
\begin{enumerate}[leftmargin=3em]
\item Give every old-tail edge its bottleneck owner from
      \Cref{lem:s6v4-owner}.
\item Prove that quotient skeletons with distinct owners are
      edge-disjoint after the owned outer forest is fixed.
\item Contract all rooted new-tail branches and compare their
      old-tail union with the f5-V88 quotient graph.
\item Apply the maximum-spanning-tree/Kruskal partition of unity
      before counting surplus quotient edges.
\item Prove QTC first for a laminar quotient history; isolate any
      nonlaminar collision as a separate card.
\end{enumerate}
\end{programme}

\begin{firewall}[Claim boundary after sixth-series V5]
\label{fire:s6v5-boundary}
V5 audits SM V28 and NS V25, proves the full one-layer
fixed-final-set forest projection, derives the exact conditional
weighted quotient-tree and endpoint-lift laws, proves an exact
measure-preserving orientation factorization, and uses it to
replace the floating/singly-rooted/multi-anchor trichotomy by one
topology-free one-node VAFPC kernel.  It isolates QTC as the
remaining cross-node compatibility statement.

V5 does not prove QTC, full VAFPC, FPC, ULCC, full IGMC, SISC,
the simple high-excess row, arbitrary repeated transversal
supports, full TOCC, full TCNC, cyclic TIC, HLTC, full MCWC,
OCRC, SCNGC, WPSC, GRLTC, LZTC, PLRC, WSDC, BRSC, HWSFC,
UGCGC, thermal connected WPRAC, all-arity RGSC, the raw Wilson
aggregate strong card, all-order strong MTA, WUFC, CPM, cutoff
removal, reflection positivity, Osterwalder--Schrader
reconstruction, a continuum Yang--Mills measure, or a positive
Hamiltonian gap.  The full Yang--Mills mass gap has not been
proved.
\end{firewall}

\begin{theorem}[Sixth-series V5 result ledger]
\label{thm:s6v5-results}
V5:
\begin{enumerate}[label=\textup{(V5.\arabic*)},leftmargin=5.5em]
\item performs the compulsory SM V28 and NS V25 audit;
\item proves the fixed-final-set projection;
\item proves the isolated-extension probability ratio;
\item proves the conditional quotient-tree law;
\item proves the independent endpoint-lift law;
\item proves the canonical rooted/coalescing quotient split;
\item proves the exact old-tail completion sum;
\item proves measure-preserving orientation factorization;
\item proves the topology-free one-node VAFPC kernel; and
\item formulates QTC and proves its sufficiency for VAFPC.
\end{enumerate}
\end{theorem}

\begin{proof}
Use
\Cref{audit:s6v5-SM,
audit:s6v5-NS,
thm:s6v5-fixed-final,
cor:s6v5-isolated-ratio,
thm:s6v5-quotient-law,
lem:s6v5-quotient-split,
lem:s6v5-old-tail,
thm:s6v5-orientation-factorization,
thm:s6v5-one-node,
prop:s6v5-QTC}.
\end{proof}

\paragraph{Parent-version record.}
V5 was formed from
\texttt{Renewal\_Yang\_Mills\_Mass\_Gap\_Research\_Notes\_V4.tex}.
The V4 parent had \(1\,893\) logical source lines, \(64\,976\)
bytes, and SHA--256
\[
 \texttt{f64f3141743e3cffb20a713b714c6bcbed1ff6cf526dadfcb1a94d5d2c09c849}.
\]
The next compulsory companion audit is V10.

\clearpage
\chapter{Sixth-series V6: spacing ownership and laminar
quotient-history gluing}
\label{ch:s6v6-gluing}

\section{Correction: bottleneck owner versus affine-layer owner}
\label{sec:s6v6-owner-correction}

No companion audit is due at V6.

\begin{correction}[The bottleneck is the deepest admissible
owner]
\label{corr:s6v6-owner}
The edge \(\beta_\pi(h)\) of
\Cref{def:s6v4-owner,lem:s6v4-owner} is the unique edge whose
parameter equals the path minimum \(x_h\).  It is therefore the
deepest hierarchy level at which the shortcut \(h\) remains
inside one block.  It is not, by itself, the unique summand owner
after the affine identity
\[
 x_h=\lambda+(1-\lambda)z_h
\]
has been expanded recursively.  One monomial may select an
ancestor-layer summand.  Accordingly,
\Cref{cor:s6v4-shortcut-owner} is to be read as uniqueness of the
bottleneck edge, not uniqueness of the affine-layer summand.
\end{correction}

\begin{proof}
If the endpoints of \(h\) remain in one child block, then
\(z_h>0\) and both summands in the displayed identity may occur
in the multi-affine expansion.  Iterating gives one possible
summand at each ancestor level up to the bottleneck.  The equality
\(x_h=u_{\beta_\pi(h)}\) and uniqueness of the deepest level
remain valid.
\end{proof}

\begin{assessment}[The rooted VFLAC proof survives]
\label{ass:s6v6-VFLAC-survives}
The proof of \Cref{thm:s6v4-rooted-VFLAC} did not identify
\(q(v)\) with \(\beta_\pi(p(v))\).  It defined \(q(v)\) from the
expanded monomial and summed all affine choices by the iterated
Bernoulli contraction.  Thus its conclusion is unchanged.  The
correction only prevents the same normalization from being
misstated as pointwise uniqueness.
\end{assessment}

\section{Complete spacings give the exact layer-allocation law}
\label{sec:s6v6-spacings}

Retain the order-simplex notation of V4.  Put
\begin{align}
 a_0&:=u_{\pi(1)},\notag\\
 a_s&:=u_{\pi(s+1)}-u_{\pi(s)}
       \quad(1\le s<n),\notag\\
 a_n&:=1-u_{\pi(n)}.                                      \label{eq:s6v6-spacings}
\end{align}
Then \(a_s>0\) and \(\sum_{s=0}^na_s=1\).
For a shortcut edge \(h\), let
\[
 \kappa_\pi(h)
 \quad\hbox{be defined by}\quad
 \beta_\pi(h)=\pi(\kappa_\pi(h)).
\]

\begin{lemma}[Path minimum is a cumulative spacing]
\label{lem:s6v6-cumulative}
For every shortcut \(h\),
\begin{equation}
 x_h
 =
 \sum_{s=0}^{\kappa_\pi(h)-1}a_s.                         \label{eq:s6v6-cumulative}
\end{equation}
\end{lemma}

\begin{proof}
By \Cref{lem:s6v4-owner},
\(x_h=u_{\pi(\kappa_\pi(h))}\).  The definition
\eqref{eq:s6v6-spacings} telescopes to the right side.
\end{proof}

\begin{definition}[Admissible spacing allocation]
\label{def:s6v6-spacing-allocation}
For a finite multiset \({\cal R}\) of shortcut-edge occurrences,
an admissible allocation is a map
\[
 \ell:{\cal R}\longrightarrow\{0,\ldots,n-1\}
\]
such that
\[
 0\le\ell(h)<\kappa_\pi(h)
\]
for every occurrence \(h\).
\end{definition}

\begin{theorem}[Exact spacing-allocation normalization]
\label{thm:s6v6-spacing-normalization}
For every finite multiset \({\cal R}\),
\begin{equation}
 \prod_{h\in{\cal R}}x_h
 =
 \sum_{\ell\ {\rm admissible}}
 \prod_{h\in{\cal R}}a_{\ell(h)}
 \le1.                                                     \label{eq:s6v6-spacing-normalization}
\end{equation}
The statement includes repeated shortcut edges as distinct
occurrences.
\end{theorem}

\begin{proof}
Insert \eqref{eq:s6v6-cumulative} for each occurrence and expand
the finite product.  This gives the equality.  Since every
\(0\le x_h\le1\), their product is at most one.
\end{proof}

\begin{corollary}[No cross-node level census]
\label{cor:s6v6-no-level-census}
For a fixed ordered outer tree and a fixed collection of retained
shortcut occurrences, the complete choice of affine hierarchy
levels has total weight \(\prod_hx_h\le1\).  It is invalid to
replace this sum by a factor equal to the number of admissible
nodes per edge.
\end{corollary}

\begin{proof}
The recursive affine expansion is the repeated telescoping of
\eqref{eq:s6v6-cumulative}.  Apply
\Cref{thm:s6v6-spacing-normalization}.
\end{proof}

\section{Quotient histories on a common final vertex set}
\label{sec:s6v6-history}

Consider a complete nested history after applying the one-node
factorization of V5.  At a hierarchy node \(Q\), denote by:
\[
 Y_Q\quad\hbox{its first-accreted labels},\qquad
 R_Q\quad\hbox{its new-tail rooted forest},\qquad
 S_Q\quad\hbox{its old-tail coalescence skeleton}.
\]
Endpoint lifts are retained, so every edge occurrence is an
actual pair of final labels.  The sets \(Y_Q\) are disjoint by
\Cref{lem:s6v4-first-node}.

\begin{definition}[Laminar quotient history]
\label{def:s6v6-laminar}
A quotient history is \emph{laminar} if, after all new-tail
forests \(R_Q\) have been glued to their old components:
\begin{enumerate}[label=\textup{(\roman*)},leftmargin=3.5em]
\item the retained occurrences in the old-tail skeletons
      \(S_Q\) are pairwise distinct;
\item their union with the pre-history owned forest is a forest;
\item every occurrence is recorded at its highest selected
      affine layer; and
\item restriction to a child hierarchy block agrees with the
      skeleton recorded by that child.
\end{enumerate}
\end{definition}

\begin{lemma}[Unique reconstruction of a laminar history]
\label{lem:s6v6-laminar-reconstruction}
Fix the ordered outer tree, the final glued forest, the disjoint
sets \(Y_Q\), and the spacing allocation of every retained edge
occurrence.  There is at most one laminar quotient history
producing these data.
\end{lemma}

\begin{proof}
The spacing allocation gives the highest selected affine layer
of every retained occurrence.  Hence it assigns the occurrence
to one hierarchy node.  At that node, orientation toward the
auxiliary old root determines whether the tail is new or old.
If it is new, the edge belongs to \(R_Q\); if it is old, it
belongs to \(S_Q\).  The first-accretion partition determines the
new labels present at each node.  Child restriction then
recovers the history inductively from the root downward.
\end{proof}

\begin{theorem}[Laminar QTC]
\label{thm:s6v6-laminar-QTC}
QTC holds on the laminar quotient-history sector.  Its complete
hierarchy, level, quotient-tree, endpoint-lift, and first-accretion
fibre is at most \(K^m\).
\end{theorem}

\begin{proof}
At each node apply
\Cref{thm:s6v5-one-node}.  The reserve enlargement contributes
\(K^{|Y_Q|}\).  Since the first-accretion sets are disjoint,
\[
 \prod_QK^{|Y_Q|}\le K^m.
\]

By laminarity, the new-tail and old-tail outputs glue to one
forest, with no repeated occurrence or cycle.  Complete that
forest inside one global weighted Cayley tree and retain its
probability; the quotient-tree and endpoint-lift masses at each
step are normalized by
\Cref{thm:s6v5-quotient-law,
thm:s6v5-orientation-factorization}.

For the fixed final edge occurrences, sum all affine-level
assignments by
\Cref{thm:s6v6-spacing-normalization}; their mass is at most one.
Then sum all outer-tree orders by
\Cref{lem:s6v4-order-simplex}; their integration mass is one.
Finally,
\Cref{lem:s6v6-laminar-reconstruction} gives inverse fibre one
after these retained data are fixed.  Only the numerical
\(K^m\) reserve base remains.
\end{proof}

\begin{corollary}[Zero-collision VAFPC sector]
\label{cor:s6v6-zero-collision}
Every VAFPC history whose oriented quotient outputs have acyclic,
nonrepeated union satisfies the required strong base.
\end{corollary}

\begin{proof}
Such a history is laminar after assigning each occurrence to its
highest selected layer.  Apply
\Cref{thm:s6v6-laminar-QTC,prop:s6v5-QTC}.
\end{proof}

\section{The exact nonlaminar defect}
\label{sec:s6v6-defect}

Let \(E_0\) be the pre-history owned forest on the final quotient
vertex set \(V\); isolated vertices may be retained for this
count.  Let \({\cal M}\) be the multiset of all lifted old-tail
edge occurrences from the \(S_Q\), and form the multigraph
\[
 G:=E_0\cup{\cal M}.
\]
Parallel occurrences are retained as distinct edges.  Write
\(c_0=c(E_0)\) and \(c_G=c(G)\).

\begin{definition}[Quotient-history collision defect]
\label{def:s6v6-defect}
Set
\begin{equation}
 g_{\rm QH}
 :=
 |{\cal M}|-\{c_0-c_G\}.                                  \label{eq:s6v6-defect}
\end{equation}
\end{definition}

\begin{theorem}[Defect equals multigraph cycle excess]
\label{thm:s6v6-defect}
One has
\begin{equation}
 g_{\rm QH}
 =
 |E(G)|-|V|+c_G
 \ge0.                                                     \label{eq:s6v6-defect-excess}
\end{equation}
It is exactly the number of old-tail occurrences which must be
deleted from \(G\) to obtain a maximal forest containing \(E_0\).
In particular, \(g_{\rm QH}=0\) precisely when
\(E_0\cup{\cal M}\) is a forest.
\end{theorem}

\begin{proof}
Because \(E_0\) is a forest,
\[
 |E_0|=|V|-c_0.
\]
Since \(|E(G)|=|E_0|+|{\cal M}|\), substitution gives
\[
 |E(G)|-|V|+c_G
 =
 |{\cal M}|-(c_0-c_G)
 =
 g_{\rm QH}.
\]
A maximal forest of \(G\) has \(|V|-c_G\) edges and may be chosen
to contain the forest \(E_0\).  Therefore exactly the displayed
excess number of occurrences must be removed.  Nonnegativity and
the zero characterization follow.
\end{proof}

\begin{corollary}[Rank-defect form]
\label{cor:s6v6-rank-defect}
If \(\operatorname{rk}(A):=|V|-c(A)\) denotes graphic-matroid
rank, then
\begin{equation}
 g_{\rm QH}
 =
 \sum_Q|S_Q|
 -
 \{\operatorname{rk}(G)-\operatorname{rk}(E_0)\}.          \label{eq:s6v6-rank-defect}
\end{equation}
Thus the obstruction is failure of the nodewise old-tail forests
to add independent rank, not the depth of the hierarchy.
\end{corollary}

\begin{proof}
Use \(|{\cal M}|=\sum_Q|S_Q|\) and
\(\operatorname{rk}(G)-\operatorname{rk}(E_0)=c_0-c_G\).
\end{proof}

\begin{lemma}[Kruskal selects exactly the surplus occurrences]
\label{lem:s6v6-Kruskal-surplus}
Give the occurrences of \(G\setminus E_0\) distinct auxiliary
weights.  Decreasing-weight Kruskal, with every edge of \(E_0\)
declared earlier than the new occurrences, selects a unique
maximal forest containing \(E_0\) and rejects exactly
\(g_{\rm QH}\) occurrences.
\end{lemma}

\begin{proof}
The greedy algorithm preserves acyclicity and ends at a maximal
forest containing \(E_0\).  Its edge count is
\(|V|-c_G\), so the number rejected is
\(|E(G)|-(|V|-c_G)=g_{\rm QH}\).  Distinct auxiliary weights give
uniqueness away from a null tie set.
\end{proof}

\begin{assessment}[Kruskal identifies but does not yet pay the
defect]
\label{ass:s6v6-Kruskal}
The normalized Kruskal region prevents a spanning-forest choice
factor.  It does not authorize deletion of a rejected edge
weight.  The f5-V88--V89 analysis already proved that a weighted
exchange requires a genuine coalescing edge with fresh ownership;
a normalized endpoint probability cannot create that mark.
Therefore \Cref{lem:s6v6-Kruskal-surplus} locates the precise
surplus set but does not close it.
\end{assessment}

\section{The remaining surplus card}
\label{sec:s6v6-surplus}

\begin{definition}[Quotient-history surplus card, QHSC]
\label{def:s6v6-QHSC}
QHSC is a coefficientwise allocation for histories with
\(g_{\rm QH}>0\) such that:
\begin{enumerate}[label=\textup{(\roman*)},leftmargin=3.5em]
\item the complete child/base spacing allocation is retained;
\item the normalized Kruskal forest of
      \Cref{lem:s6v6-Kruskal-surplus} selects the
      \(g_{\rm QH}\) surplus occurrences;
\item each surplus occurrence is exchanged against a distinct
      genuinely owned coalescing edge, without reusing a root
      occurrence;
\item the exchange preserves or decreases the exact endpoint
      monomial; and
\item the inverse history fibre is at most \(K^m\).
\end{enumerate}
\end{definition}

\begin{proposition}[QHSC completes QTC]
\label{prop:s6v6-QHSC}
QHSC together with \Cref{thm:s6v6-laminar-QTC} proves QTC and
hence VAFPC.
\end{proposition}

\begin{proof}
Partition histories by \(g_{\rm QH}\).  The zero row is
\Cref{thm:s6v6-laminar-QTC}.  In every positive row, QHSC
exchanges the Kruskal surplus and produces an acyclic
nonrepeated history with an exponential inverse fibre.  Apply the
laminar theorem to its image.  Then use
\Cref{prop:s6v5-QTC}.
\end{proof}

\begin{assessment}[Exact live bottleneck after V6]
\label{ass:s6v6-bottleneck}
All hierarchy-level choices, including ancestor selections above
the bottleneck, now have an exact subprobability law.  All
zero-collision quotient histories are closed.  The remaining
problem is QHSC: source-level ownership of one coalescing payment
for each unit of the graphic-matroid rank defect
\(g_{\rm QH}\).  This is the same kind of freshness issue found
in f5-V89, but now the rejected edges and their hierarchy levels
are canonical.
\end{assessment}

\begin{programme}[V7: source ownership for the surplus set]
\label{prog:s6v6-V7}
\begin{enumerate}[leftmargin=3em]
\item Return to the pre-majorant squarefree source coefficient
      for the node histories selected by Kruskal.
\item Color root occurrences as owned-tree, rooted-accretion, or
      surplus-coalescence occurrences before taking norms.
\item Prove the three-color coefficient identity and test whether
      connectedness supplies one distinct coalescing occurrence
      per rejected edge.
\item If freshness fails, build a mixed alternating path whose
      first unused edge supplies the exchange.
\item Close first \(g_{\rm QH}=1\), then test whether the
      construction tensorizes over the Kruskal fundamental
      cycles.
\end{enumerate}
\end{programme}

\begin{firewall}[Claim boundary after sixth-series V6]
\label{fire:s6v6-boundary}
V6 corrects the distinction between bottleneck and affine-layer
ownership, proves the exact cumulative-spacing allocation law,
proves that all cross-node level choices have subprobability
mass, defines laminar quotient histories, proves laminar QTC and
the zero-collision VAFPC sector, identifies the exact nonlaminar
rank defect, and canonically selects its surplus occurrences by
Kruskal.  It formulates QHSC.

V6 does not prove QHSC, full QTC, full VAFPC, FPC, ULCC, full
IGMC, SISC, the simple high-excess row, arbitrary repeated
transversal supports, full TOCC, full TCNC, cyclic TIC, HLTC,
full MCWC, OCRC, SCNGC, WPSC, GRLTC, LZTC, PLRC, WSDC, BRSC,
HWSFC, UGCGC, thermal connected WPRAC, all-arity RGSC, the raw
Wilson aggregate strong card, all-order strong MTA, WUFC, CPM,
cutoff removal, reflection positivity, Osterwalder--Schrader
reconstruction, a continuum Yang--Mills measure, or a positive
Hamiltonian gap.  The full Yang--Mills mass gap has not been
proved.
\end{firewall}

\begin{theorem}[Sixth-series V6 result ledger]
\label{thm:s6v6-results}
V6:
\begin{enumerate}[label=\textup{(V6.\arabic*)},leftmargin=5.5em]
\item corrects bottleneck versus affine-layer ownership;
\item proves cumulative-spacing representation;
\item proves exact spacing-allocation normalization;
\item defines and reconstructs laminar quotient histories;
\item proves laminar QTC and zero-collision VAFPC;
\item identifies the quotient-history rank defect;
\item proves equality with multigraph cycle excess;
\item canonically selects all surplus occurrences by Kruskal; and
\item formulates QHSC and proves its sufficiency.
\end{enumerate}
\end{theorem}

\begin{proof}
Use
\Cref{corr:s6v6-owner,
lem:s6v6-cumulative,
thm:s6v6-spacing-normalization,
lem:s6v6-laminar-reconstruction,
thm:s6v6-laminar-QTC,
thm:s6v6-defect,
cor:s6v6-rank-defect,
lem:s6v6-Kruskal-surplus,
prop:s6v6-QHSC}.
\end{proof}

\paragraph{Parent-version record.}
V6 was formed from
\texttt{Renewal\_Yang\_Mills\_Mass\_Gap\_Research\_Notes\_V5.tex}.
The V5 parent had \(2\,464\) logical source lines, \(85\,750\)
bytes, and SHA--256
\[
 \texttt{6f6cbd7532fa75c668c3b9d1e2a96a47018542b1ea59a7dd09d11415f609e88a}.
\]
The next compulsory companion audit is V10.

\clearpage
\chapter{Sixth-series V7: marginal quotient forests and closure
of vertex accretion}
\label{ch:s6v7-marginal}

\section{The archived freshness no-go changes the route}
\label{sec:s6v7-no-freshness}

No companion audit is due at V7.

\begin{record}[Upstream result which must not be repeated]
\label{rec:s6v7-V90}
Fifth-series V90 already proves two facts relevant to the V6
programme:
\begin{enumerate}[label=\textup{(\roman*)},leftmargin=3.5em]
\item squarefreeness gives an exact disjoint colouring of root
      occurrences; but
\item connectedness does not imply that any occurrence remains
      available with a fresh background colour.
\end{enumerate}
A single differentiated transversal root can be connected and
have positive compatibility defect while the background colour
is empty.  Therefore the proposed QHSC proof cannot infer a
fresh coalescing payer merely by adding a third colour.
\end{record}

\begin{assessment}[Go upstream in the conditional probability]
\label{ass:s6v7-upstream}
The old-tail skeleton \(S_\rho\) in V5 is not a physical source
occurrence.  It is the part of a Cayley tree left after a
conditional orientation.  Requiring it to survive across
hierarchy nodes turned a normalized completion variable into a
new combinatorial history and created the V6 defect.  The proper
upstream operation is to compute its conditional marginal.
\end{assessment}

\section{The exact rooted-forest marginal}
\label{sec:s6v7-marginal}

Use the notation of V5:
\[
 A=F\mathbin{\dot\cup}J,\qquad
 t=|J|,\qquad
 C_1,\ldots,C_r\in{\cal K}(E),\qquad
 a_p=B_{C_p}.
\]
On the quotient set
\({\cal Q}=\{C_1,\ldots,C_r\}\dot\cup J\), give the old
vertices weights \(a_p\) and the new vertices weights \(b_j\).
Let
\(\mathfrak F_{\rm old}({\cal Q})\) denote the forests in which
every component contains exactly one old quotient vertex.

\begin{theorem}[Marginal of a randomly rooted quotient Cayley
tree]
\label{thm:s6v7-rooted-marginal}
Sample:
\begin{enumerate}[label=\textup{(\roman*)},leftmargin=3.5em]
\item a quotient tree \(\tau\) with law
      \(\overline\mu_{\cal Q}\) from
      \Cref{thm:s6v5-quotient-law}; and
\item independently, an auxiliary old root \(\rho\) with
      probability \(a_\rho/B_F\).
\end{enumerate}
Orient \(\tau\) toward \(C_\rho\) and retain only the parent
edges whose tails lie in \(J\).  The retained forest
\(R_\rho(\tau)\) has probability law
\begin{equation}
 \nu_{E,J}(R)
 =
 {\displaystyle\prod_{\{\alpha,\beta\}\in R}
                    w_\alpha w_\beta
  \over
  \displaystyle
  B_F\left(\prod_{j\in J}b_j\right)B_A^{t-1}},
 \qquad
 R\in\mathfrak F_{\rm old}({\cal Q}).                    \label{eq:s6v7-rooted-marginal}
\end{equation}
In particular,
\begin{equation}
 \sum_{R\in\mathfrak F_{\rm old}({\cal Q})}
 \nu_{E,J}(R)=1.                                         \label{eq:s6v7-rooted-normalized}
\end{equation}
\end{theorem}

\begin{proof}
Fix \(R\).  By
\Cref{lem:s6v5-old-tail}, the auxiliary-root average of all
compatible old-tail completions is
\[
 {\prod_{p=1}^ra_p\over B_F}B_A^{r-1}.
\]
Multiplying this by the edge-weight product of \(R\), and
dividing by the quotient Cayley normalization
\[
 \left(\prod_{p=1}^ra_p\right)
 \left(\prod_{j\in J}b_j\right)B_A^{r+t-2},
\]
gives exactly \eqref{eq:s6v7-rooted-marginal}.  The weighted
rooted-forest identity gives
\[
 \sum_{R\in\mathfrak F_{\rm old}({\cal Q})}
 \prod_{e\in R}w_e
 =
 B_F\left(\prod_{j\in J}b_j\right)B_A^{t-1},
\]
which proves \eqref{eq:s6v7-rooted-normalized}.
\end{proof}

\begin{corollary}[The old-tail completion can be marginalized]
\label{cor:s6v7-old-tail-marginal}
After the random-root average, the complete conditional sum over
\(S_\rho\) has total mass one relative to the rooted-forest
marginal.  It introduces neither a retained coalescence skeleton
nor a branch-partition factor.
\end{corollary}

\begin{proof}
Equation \eqref{eq:s6v7-rooted-marginal} was obtained by summing
all \(S_\rho\) compatible with a fixed \(R\).  Its normalization
is \eqref{eq:s6v7-rooted-normalized}.
\end{proof}

\section{Endpoint lifts preserve the marginal}
\label{sec:s6v7-lifts}

\begin{lemma}[Root-component endpoint lift]
\label{lem:s6v7-root-lift}
For every quotient edge of \(R\) joining a new vertex \(j\) to
an old component \(C_p\), choose its endpoint \(i\in C_p\) with
probability
\begin{equation}
 {b_i\over B_{C_p}}.                                      \label{eq:s6v7-root-lift}
\end{equation}
Do this independently for all such incidences.  After adjoining
the fixed internal forest \(E\), the lifted graph is a forest on
\(F\cup J\), and every one of its components contains exactly
one component of \(E\).  The lift probabilities sum to one and
restore the exact endpoint edge weights.
\end{lemma}

\begin{proof}
Splitting a quotient root vertex \(C_p\) into the connected tree
component \(E[C_p]\), and attaching its incident branches at
arbitrary vertices of \(C_p\), preserves connectedness and
acyclicity.  The quotient forest has exactly one root vertex in
each component, so the lifted forest has exactly one component of
\(E\) in each component.

For one incidence,
\[
 \sum_{i\in C_p}{b_i\over B_{C_p}}=1,
\]
and multiplying the quotient factor \(b_jB_{C_p}\) by the
probability in \eqref{eq:s6v7-root-lift} gives \(b_jb_i\).
Independence proves the joint statement.
\end{proof}

\begin{theorem}[Actual-label rooted-forest marginal]
\label{thm:s6v7-actual-marginal}
After the endpoint lifts of
\Cref{lem:s6v7-root-lift}, the marginal law is the normalized
weighted law on all forests \(R\) on \(F\cup J\) which:
\begin{enumerate}[label=\textup{(\roman*)},leftmargin=3.5em]
\item contain no new edge wholly inside a component of \(E\);
\item use exactly \(t\) new edges; and
\item make every component contain exactly one component of
      \(E\).
\end{enumerate}
Its total mass is one.
\end{theorem}

\begin{proof}
Contracting the components of \(E\) maps these actual-label
forests to
\(\mathfrak F_{\rm old}({\cal Q})\).  For each quotient forest,
\Cref{lem:s6v7-root-lift} is a normalized disintegration of all
lifts and reproduces their exact edge monomials.  Apply
\Cref{thm:s6v7-rooted-marginal}.
\end{proof}

\section{VAFPC is closed}
\label{sec:s6v7-VAFPC}

\begin{lemma}[One-node sub-Markov kernel]
\label{lem:s6v7-one-node-kernel}
Fix the old forest \(E\), the final new set \(J\), and one
hierarchy node.  The combined operation
\begin{equation}
 \text{Cayley-tree disintegration}
 \longrightarrow
 \text{child/base colour sum}
 \longrightarrow
 \text{random-root orientation}
 \longrightarrow
 \text{old-tail marginalization}                         \label{eq:s6v7-kernel}
\end{equation}
is a sub-Markov kernel whose output is only the actual-label
rooted forest of
\Cref{thm:s6v7-actual-marginal}.
\end{lemma}

\begin{proof}
For a fixed Cayley tree, the child/base colour sum, including the
coverage restriction on \(J\), has mass at most one by the proof
of \Cref{thm:s6v5-fixed-final}.  The conditional quotient-tree
law and endpoint lifts have mass one by
\Cref{thm:s6v5-quotient-law}.  The auxiliary root probabilities
sum to one.  Finally, the old-tail completion is marginalized by
\Cref{cor:s6v7-old-tail-marginal}.  Composition of these finite
positive kernels has mass at most one.
\end{proof}

\begin{theorem}[Full vertex-accreting forest-projection card]
\label{thm:s6v7-VAFPC}
VAFPC holds with total strong base \(K^m\).
\end{theorem}

\begin{proof}
Iterate \Cref{lem:s6v7-one-node-kernel} down the finite
ultrametric hierarchy.  Conditional on the past, every node
kernel has mass at most one, so the tower property gives total
probability at most one for all child/base, quotient-tree,
old-tail, endpoint-lift, and auxiliary-root data.

The only retained output at a node is a rooted forest spanning
its first-accreted set \(Y_Q\).  By
\Cref{thm:s6v3-tree-valued,thm:s6v3-RCTAC-shapes}, its reserve
extension costs \(K^{|Y_Q|}\), with component-mass updates
telescoping.  By \Cref{lem:s6v4-first-node}, the sets \(Y_Q\)
are disjoint.  Hence
\[
 \prod_QK^{|Y_Q|}\le K^m.
\]

Adding a rooted output which contains exactly one old component
per component preserves the forest property and never coalesces
two previously distinct old components.  Thus the outputs
compose without a cross-node cycle.  All hierarchy-layer
allocations are already normalized by
\Cref{lem:s6v4-iterated}; equivalently, they may be summed by
\Cref{thm:s6v6-spacing-normalization}.  This proves every item in
the VAFPC definition.
\end{proof}

\begin{corollary}[FPC, ULCC, and IGMC]
\label{cor:s6v7-FPC-chain}
The full forest-projection card, the ultrametric layer-collapse
card, and the interpolation-gate mixture card hold.  Consequently
the optional-support/dense-degree-fibre part of SISC is closed.
\end{corollary}

\begin{proof}
The fixed-vertex FPC row is f5-V100.  Its complement is VAFPC,
proved in \Cref{thm:s6v7-VAFPC}; hence FPC holds.  The proved
implications f5-V99--V100 give ULCC, and f5-V97--V98 then give
IGMC.  The f5-V97 IGMC proposition removes every optional support
membership before norms, which is exactly the support-fibre part
of SISC.
\end{proof}

\begin{correction}[QTC and QHSC are unnecessary for VAFPC]
\label{corr:s6v7-QTC}
The QTC of V5 and QHSC of V6 are retired as required
intermediate cards for VAFPC.  Their retained-history statements
remain valid for the artificial expansion in which every
old-tail conditional completion is kept.  The canonical
calculation instead marginalizes that completion by
\Cref{thm:s6v7-rooted-marginal}, so its cross-node rank defect
never enters the live VAFPC proof.
\end{correction}

\begin{proof}
The conditional marginal contains only rooted forests, whose
successive union is acyclic by the proof of
\Cref{thm:s6v7-VAFPC}.  There are no retained old-tail
occurrences from which the multigraph of
\Cref{def:s6v6-defect} could be formed.
\end{proof}

\section{The component defect which genuinely remains}
\label{sec:s6v7-SOCC}

\begin{assessment}[Do not confuse two component defects]
\label{ass:s6v7-two-defects}
The eliminated V6 defect concerned auxiliary completion edges
inside conditional Cayley probabilities.  SISC still contains a
different, source-level defect.  After IGMC removes optional
labels, distinguishable primitive root occurrences own disjoint
nonempty subforests \(E_\nu\) of the sampled outer tree.  If
\(c_\nu\) is the number of components of \(E_\nu\), their forced
incidence excess is
\[
 g_{\rm forced}=\sum_\nu(c_\nu-1).
\]
These owned edges are physical derivative/source data and cannot
be marginalized as unused Cayley completions.
\end{assessment}

\begin{definition}[Simultaneous owned-component card, SOCC]
\label{def:s6v7-SOCC}
SOCC is an exact, pre-norm reorganization of the simultaneous
source coefficient after IGMC such that:
\begin{enumerate}[label=\textup{(\roman*)},leftmargin=3.5em]
\item the sampled outer tree and its partition into nonempty
      owned forests \(E_\nu\) are retained;
\item occurrence types and Duhamel coefficients are summed
      together with that edge partition;
\item the component defect
      \(\sum_\nu(c_\nu-1)\) is paid without assuming a fresh
      background occurrence;
\item repeated and distinct support types retain their existing
      factorial/Bernoulli normalizations; and
\item the complete output is a global block-compatible Cayley or
      incidence-hypertree object with fibre at most \(K^m\).
\end{enumerate}
\end{definition}

\begin{proposition}[SOCC supplies the remaining part of SISC]
\label{prop:s6v7-SOCC}
SOCC together with \Cref{cor:s6v7-FPC-chain} proves SISC.
\end{proposition}

\begin{proof}
IGMC removes all optional support memberships and the dense
degree fibre.  The f5-V97 reduction proves that the only
remaining SISC data are the disjoint owned-edge partition and
its component-coalescence defect.  These are exactly SOCC(i)--(iv);
SOCC(v) gives the required strong output.
\end{proof}

\begin{assessment}[Exact live bottleneck after V7]
\label{ass:s6v7-bottleneck}
Vertex accretion is no longer open.  The next obstruction is
SOCC, not QHSC: a simultaneous source-level sum over partitions
of the \(m-1\) outer-tree edges into owned forests, retaining the
component defect before norms.  The archived V90 no-go forbids
deducing its payment from connectedness or from an invented fresh
colour.
\end{assessment}

\begin{programme}[V8: exact owned-edge colouring law]
\label{prog:s6v7-V8}
\begin{enumerate}[leftmargin=3em]
\item Write the simultaneous analytic coefficient with one
      colour for each distinguishable primitive root occurrence.
\item Condition on the sampled outer tree before summing its edge
      colours.
\item Derive the exact exponential/Duhamel normalization for
      nonempty colour classes.
\item Express \(\sum_\nu(c_\nu-1)\) directly from colour changes
      along the tree, rather than from a post-norm support family.
\item Close the row in which every colour class is connected and
      isolate the first disconnected-colour defect.
\end{enumerate}
\end{programme}

\begin{firewall}[Claim boundary after sixth-series V7]
\label{fire:s6v7-boundary}
V7 uses the archived no-freshness theorem to abandon the proposed
QHSC route, computes the exact rooted-forest marginal of a
randomly rooted conditional quotient Cayley tree, proves the
actual-label endpoint lift, proves the one-node sub-Markov
kernel, and closes VAFPC.  It consequently closes FPC, ULCC,
IGMC, and the optional-support/dense-degree-fibre part of SISC.
It retires QTC/QHSC as unnecessary VAFPC cards and isolates SOCC
as the genuine simultaneous source-level component defect.

V7 does not prove SOCC, SISC, the simple high-excess row,
arbitrary repeated transversal supports, full TOCC, full TCNC,
cyclic TIC, HLTC, full MCWC, OCRC, SCNGC, WPSC, GRLTC, LZTC,
PLRC, WSDC, BRSC, HWSFC, UGCGC, thermal connected WPRAC,
all-arity RGSC, the raw Wilson aggregate strong card, all-order
strong MTA, WUFC, CPM, cutoff removal, reflection positivity,
Osterwalder--Schrader reconstruction, a continuum Yang--Mills
measure, or a positive Hamiltonian gap.  The full Yang--Mills
mass gap has not been proved.
\end{firewall}

\begin{theorem}[Sixth-series V7 result ledger]
\label{thm:s6v7-results}
V7:
\begin{enumerate}[label=\textup{(V7.\arabic*)},leftmargin=5.5em]
\item imports the f5-V90 no-freshness boundary;
\item proves the normalized quotient rooted-forest marginal;
\item proves normalized actual-label endpoint lifts;
\item proves the one-node sub-Markov kernel;
\item proves VAFPC;
\item closes FPC, ULCC, and IGMC;
\item closes the optional-support part of SISC;
\item retires QTC/QHSC from the live VAFPC chain; and
\item formulates SOCC and proves its sufficiency for SISC.
\end{enumerate}
\end{theorem}

\begin{proof}
Use
\Cref{rec:s6v7-V90,
thm:s6v7-rooted-marginal,
lem:s6v7-root-lift,
thm:s6v7-actual-marginal,
lem:s6v7-one-node-kernel,
thm:s6v7-VAFPC,
cor:s6v7-FPC-chain,
corr:s6v7-QTC,
prop:s6v7-SOCC}.
\end{proof}

\paragraph{Parent-version record.}
V7 was formed from
\texttt{Renewal\_Yang\_Mills\_Mass\_Gap\_Research\_Notes\_V6.tex}.
The V6 parent had \(2\,921\) logical source lines, \(101\,564\)
bytes, and SHA--256
\[
 \texttt{4215eef6fe7f70f706da51ed42a61cae3a9d5535997ccbcfc615fcdd14e99d6b}.
\]
The next compulsory companion audit is V10.

\clearpage
\chapter{Sixth-series V8: resolved owned components and the
post-IGMC grouping obstruction}
\label{ch:s6v8-components}

\section{Owned-edge colourings resolve to an incidence tree}
\label{sec:s6v8-incidence}

No companion audit is due at V8.

Let \(T_0\) be a tree on \(X\), \(|X|=m\ge2\).  Partition its
edges into nonempty forests
\[
 E(T_0)=E_1\mathbin{\dot\cup}\cdots
        \mathbin{\dot\cup}E_q.
\]
The index \(\nu\) is the primitive root occurrence which owns the
edge.  Let \({\cal C}\) be the family of all nontrivial connected
components of all \(E_\nu\).  Every \(C\in{\cal C}\) is an
edge-subtree of \(T_0\).

Define the resolved incidence graph \({\cal I}\) to have:
\begin{enumerate}[label=\textup{(\roman*)},leftmargin=3.5em]
\item one left vertex for every \(C\in{\cal C}\);
\item one right vertex for every label \(i\in X\); and
\item an incidence edge \(C-i\) exactly when \(i\in V(C)\).
\end{enumerate}

\begin{theorem}[Resolved-component incidence tree]
\label{thm:s6v8-incidence-tree}
The bipartite graph \({\cal I}\) is a tree.
\end{theorem}

\begin{proof}
It is connected.  Indeed, every edge \(\{i,j\}\) of \(T_0\)
belongs to one component \(C\), and the length-two path
\(i-C-j\) lies in \({\cal I}\).  Replacing every edge along a
\(T_0\)-path by this length-two path connects the corresponding
label vertices, and every component vertex is incident to one of
its labels.

Let \(N:=|{\cal C}|\).  Because every \(C\) is a tree,
\[
 |V(C)|=|E(C)|+1.
\]
The component edge sets partition \(E(T_0)\), so
\[
 |E({\cal I})|
 =
 \sum_{C\in{\cal C}}|V(C)|
 =
 (m-1)+N.
\]
The graph \({\cal I}\) has \(m+N\) vertices.  Thus it is connected
and has one fewer edge than vertices, hence is a tree.
\end{proof}

\begin{corollary}[The forced defect is only forgotten grouping]
\label{cor:s6v8-forgotten-grouping}
If \(c_\nu\) is the number of components of \(E_\nu\), then
\[
 g_{\rm forced}
 =
 \sum_{\nu=1}^q(c_\nu-1)
 =
 |{\cal C}|-q.                                            \label{eq:s6v8-forgotten-grouping}
\]
The resolved component family itself has incidence excess zero.
All positive forced excess is created by remembering that
\(c_\nu\) distinct tree vertices of \({\cal I}\) belong to one
source occurrence.
\end{corollary}

\begin{proof}
The first identity is the f5-V97 forced-defect formula and
\(\sum_\nu c_\nu=|{\cal C}|\).  The second statement is
\Cref{thm:s6v8-incidence-tree}.
\end{proof}

\begin{lemma}[Two-colour calibration]
\label{lem:s6v8-two-colour}
When \(q=2\), form a bipartite graph whose vertices are the
components of \(E_1\) and \(E_2\), with an edge for every label
incident to both colours.  This graph is a tree.  Hence the
number of mixed-colour labels is
\begin{equation}
 c_1+c_2-1=g_{\rm forced}+1.                              \label{eq:s6v8-two-colour}
\end{equation}
\end{lemma}

\begin{proof}
Delete from \({\cal I}\) every label leaf, namely every label
incident to only one resolved component, and suppress nothing
else.  With two colours, a nonleaf label is incident to exactly
one component of each colour: two components of the same colour
meeting at that label would be one component.  The remaining
graph is the stated bipartite graph.  It is connected and
acyclic because it is the nontrivial core of a tree.  Its vertex
and edge counts give \eqref{eq:s6v8-two-colour}.
\end{proof}

\section{The gate coefficient contains an exact component-tree
factor}
\label{sec:s6v8-gate-split}

For \(|A|=a\ge2\), write the strong proxy as
\begin{equation}
 \mathfrak q(A)
 :=
 K_1^a\left(\prod_{i\in A}b_i\right)B_A^{a-2}.            \label{eq:s6v8-strong-proxy}
\end{equation}
Let \(E\) be a forest on \(A=V(E)\), with components
\({\cal K}(E)=\{C_1,\ldots,C_c\}\), and put
\(n_\alpha=|C_\alpha|\).

\begin{lemma}[A fixed component tree is below its Cayley
polynomial]
\label{lem:s6v8-component-tree}
For every component \(C_\alpha\),
\begin{equation}
 \prod_{\{i,j\}\in E[C_\alpha]}b_ib_j
 \le
 \left(\prod_{i\in C_\alpha}b_i\right)
 B_{C_\alpha}^{n_\alpha-2}.                              \label{eq:s6v8-component-tree}
\end{equation}
\end{lemma}

\begin{proof}
The left side is one term in the weighted Cayley sum on
\(C_\alpha\), while the right side is that complete sum.  For
\(n_\alpha=2\), both sides are \(b_ib_j\).
\end{proof}

\begin{theorem}[Strong gate splitting into resolved components]
\label{thm:s6v8-gate-splitting}
One has
\begin{equation}
 \gamma_A(E)\mathfrak q(A)
 \le
 \eta_A({\cal K}(E))
 \prod_{\alpha=1}^c\mathfrak q(C_\alpha),                 \label{eq:s6v8-gate-splitting}
\end{equation}
after one harmless adjustment of the common numerical base, where
\begin{equation}
 \eta_A({\cal K}(E))
 :=
 \left(\prod_{\alpha=1}^cB_{C_\alpha}\right)B_A^{c-2}.
                                                               \label{eq:s6v8-eta}
\end{equation}
For \(c=1\), the factor is one.  For \(c\ge2\),
\(\eta_A\) is exactly the weighted Cayley-tree polynomial on the
component set with vertex weights \(B_{C_\alpha}\).
\end{theorem}

\begin{proof}
The exact forest coefficient gives
\begin{align*}
 \gamma_A(E)\mathfrak q(A)
 &=
 K_1^a
 \left\{\prod_{\{i,j\}\in E}b_ib_j\right\}
 \left\{\prod_{\alpha=1}^cB_{C_\alpha}\right\}
 B_A^{c-2}.
\end{align*}
Apply \Cref{lem:s6v8-component-tree} in every component.  The
product of the resulting component strong polynomials is
\(\prod_\alpha\mathfrak q(C_\alpha)\), and the remaining factor
is \eqref{eq:s6v8-eta}.  The bases multiply with total exponent
\(\sum_\alpha n_\alpha=a\).

For \(c=1\),
\(B_{C_1}B_A^{-1}=1\).  For \(c\ge2\), the weighted Cayley
identity on component weights is precisely
\[
 \sum_{\sigma\in\mathfrak T(\{1,\ldots,c\})}
 \prod_{\{\alpha,\beta\}\in\sigma}
 B_{C_\alpha}B_{C_\beta}
 =
 \left(\prod_\alpha B_{C_\alpha}\right)B_A^{c-2}.
\]
\end{proof}

\begin{corollary}[Scale degree of the forgotten grouping]
\label{cor:s6v8-scale-degree}
Under \(b_i\mapsto\lambda b_i\),
\[
 \eta_A({\cal K}(E))
 \mapsto
 \lambda^{2c-2}\eta_A({\cal K}(E)).
                                                               \label{eq:s6v8-scale-degree}
\]
For a family \(E_\nu\), the product of these factors has extra
thermal degree
\[
 2\sum_\nu(c_\nu-1)=2g_{\rm forced}.                      \label{eq:s6v8-extra-degree}
\]
\end{corollary}

\begin{proof}
There are \(c\) component-mass factors and the final total-mass
factor has degree \(c-2\), giving \(2c-2\).  Multiply over
\(\nu\) and use \eqref{eq:s6v8-forgotten-grouping}.
\end{proof}

\begin{assessment}[The missing factors were not absent]
\label{ass:s6v8-missing-factors}
The crude fixed-excess estimate replaced each defect unit by a
global \(mB^2\).  The exact gate instead supplies a weighted
Cayley tree joining the disconnected owned components of the
same occurrence.  Its \(2g_{\rm forced}\) mark degree is exactly
the degree predicted by incidence excess.  The question is
whether all such component-joining trees can be summed
simultaneously; they must not be replaced by one scalar
\(B^{2g_{\rm forced}}\).
\end{assessment}

\section{The post-norm grouping still has a superexponential
fibre}
\label{sec:s6v8-matching}

The exact factor \(\eta_A\) is necessary, but a positive
post-norm sum of those factors is still insufficient.  The next
construction isolates the obstruction after optional supports
have already disappeared.

Let \(n=2r\) and let \(T_0\) be the path with edge set
\(\{e_1,\ldots,e_n\}\).  Give every label the same mark
\(b\), where \(0<b<1/8\) is fixed.  Regard each one-edge subtree
\(\{e_k\}\) as one resolved component \(C_k\); its mass is
between \(2b\) and \(2b\), hence exactly \(2b\).

Call a perfect matching \(M\) of \(\{1,\ldots,n\}\)
\emph{separated} when no matched pair consists of consecutive
path edges.  For \(\{k,\ell\}\in M\), group the two components
\(C_k,C_\ell\) into one root occurrence.  Its owned forest has
two components and its forced support has total mass at most
\(4b<1\).

\begin{lemma}[Many separated perfect matchings]
\label{lem:s6v8-separated-matchings}
For every even \(n=2r\), \(r\ge2\), the number \(N_n\) of
separated perfect matchings satisfies
\begin{equation}
 N_n\ge r!-(r-1)!=(r-1)(r-1)!.                            \label{eq:s6v8-separated-matchings}
\end{equation}
\end{lemma}

\begin{proof}
Match every member of \(\{1,\ldots,r\}\) bijectively to one
member of \(\{r+1,\ldots,2r\}\).  There are \(r!\) such
matchings.  The only consecutive pair which crosses these two
halves is \(\{r,r+1\}\).  Exactly \((r-1)!\) of the bijections
contain that pair.  Every other one is separated, proving
\eqref{eq:s6v8-separated-matchings}.
\end{proof}

\begin{remark}[A completely elementary weaker bound suffices]
\label{rem:s6v8-matching-bound}
The explicit family in
\Cref{lem:s6v8-separated-matchings} is already
superexponential in \(n\).  The weaker lower bound
\((r-2)!\) is enough for the next theorem.
\end{remark}

\begin{theorem}[Dense component-grouping obstruction]
\label{thm:s6v8-grouping-no-go}
For every separated matching above:
\begin{enumerate}[label=\textup{(\roman*)},leftmargin=3.5em]
\item every root occurrence has two disconnected owned
      components;
\item the \(r\) forced supports are pairwise distinct, so every
      repeated-support type factorial equals one; and
\item the exact product of the component-tree factors is
      \((4b^2)^r=(2b)^n\).
\end{enumerate}
Therefore
\begin{equation}
 \sum_{M\ {\rm separated}}
 \prod_{\{k,\ell\}\in M}
 \eta_{V(C_k)\cup V(C_\ell)}(\{C_k,C_\ell\})
 \ge
 (r-2)!(2b)^n,                                            \label{eq:s6v8-grouping-no-go}
\end{equation}
which is not bounded by \(K^n\) for any fixed \(K\).
Hence the exact factors \(\eta_A\), the support-type factorial,
and the resolved incidence-tree theorem do not by themselves
prove SOCC after norms.
\end{theorem}

\begin{proof}
Separated path edges have disjoint vertex sets, so their union
has two components.  Distinct matched pairs contain distinct
edge identities and hence give distinct forced supports.  With
two component masses \(2b,2b\), formula
\eqref{eq:s6v8-eta} gives
\[
 \eta=(2b)(2b)(4b)^0=4b^2.
\]
Use the \((r-2)!\) subfamily from
\Cref{rem:s6v8-matching-bound}.  Stirling's lower bound gives
\[
 \{(r-2)!(2b)^{2r}\}^{1/(2r)}
 \ge c\,b\sqrt r\longrightarrow\infty,
\]
which rules out a uniform exponential base.
\end{proof}

\begin{warning}[This is a proof-method no-go, not an operator
lower bound]
\label{warn:s6v8-not-lower}
Theorem~\ref{thm:s6v8-grouping-no-go} concerns the complete
positive proxy obtained after separate norms.  It does not assert
that the exact one-link operator coefficient has the same sign
or size.  It proves that SOCC must combine the edge grouping with
the simultaneous analytic source coefficient before separate
root norms are taken.
\end{warning}

\section{The corrected form of SOCC}
\label{sec:s6v8-SOCC}

\begin{correction}[Component resolution does not finish SOCC]
\label{corr:s6v8-SOCC}
\Cref{thm:s6v8-incidence-tree} removes the geometric incidence
cycles, and \Cref{thm:s6v8-gate-splitting} identifies their exact
component-joining weights.  It is nevertheless invalid to sum
the joining trees after taking a norm of every occurrence.  The
perfect-matching fibre in
\Cref{thm:s6v8-grouping-no-go} survives even when all occurrence
support types are distinct.
\end{correction}

\begin{definition}[Pre-norm simultaneous component-source card,
PSCSC]
\label{def:s6v8-PSCSC}
PSCSC is the remaining form of SOCC.  It must:
\begin{enumerate}[label=\textup{(\roman*)},leftmargin=3.5em]
\item retain the doubly polarized Duhamel coefficient of f5-V87;
\item retain the common sampled outer tree and all edge ownership
      variables;
\item apply the IGMC optional-label kernels without separately
      norming the resulting occurrences;
\item sum the component-grouping Cayley factors
      \(\eta_A\) jointly with the Duhamel word and one-link
      operator products; and
\item produce a \(K^m\) row normalization before the resolved
      incidence hypertree is bounded.
\end{enumerate}
\end{definition}

\begin{proposition}[PSCSC is now equivalent to the open part of
SOCC]
\label{prop:s6v8-PSCSC}
Given the proved IGMC and the exact identities
\Cref{thm:s6v8-incidence-tree,thm:s6v8-gate-splitting},
PSCSC proves SOCC.  Conversely, any SOCC proof which first takes
separate positive root norms must supply an additional
normalization not present in the known gate or type factors.
\end{proposition}

\begin{proof}
PSCSC removes the only data not already in the resolved
incidence-hypertree theorem: the grouping of component vertices
into primitive source occurrences.  After its joint \(K^m\)
normalization, \Cref{thm:s6v8-incidence-tree} supplies the
incidence-hypertree output and hence SOCC(v).

The converse is
\Cref{thm:s6v8-grouping-no-go}: all known post-norm scalar
payments are present in that example, yet their sum is
superexponential.
\end{proof}

\begin{assessment}[Exact live bottleneck after V8]
\label{ass:s6v8-bottleneck}
The geometric component defect has been resolved completely: its
components form an incidence tree and its exact missing weight is
a Cayley tree on components.  The remaining obstruction is
analytic and simultaneous.  One must obtain cancellation or a
joint probability law for the component groupings inside the
specific Duhamel one-link source; generic positive
post-norm bounds cannot do it.
\end{assessment}

\begin{programme}[V9: test the exact Duhamel component source]
\label{prog:s6v8-V9}
\begin{enumerate}[leftmargin=3em]
\item Insert independent component variables into the f5-V87
      doubly polarized exponential before support types are
      identified.
\item Compute the coefficient of the separated matching family
      in the commuting scalar calibration.
\item Determine whether the exact Duhamel simplex supplies a
      missing global occurrence factorial; if it does not, record
      a generic-source no-go.
\item Return to the noncommutative one-link increments and test a
      forest/matrix-tree determinant before taking cb norms.
\item Replace PSCSC by the smallest source-specific determinant
      or cancellation card which survives the calibration.
\end{enumerate}
\end{programme}

\begin{firewall}[Claim boundary after sixth-series V8]
\label{fire:s6v8-boundary}
V8 proves that the connected components of every owned-edge
colouring form a bipartite incidence tree, identifies the forced
defect exactly with forgotten occurrence grouping, proves the
two-colour calibration, derives the exact strong-gate splitting
factor, and identifies that factor as a component Cayley
polynomial of degree \(2g_{\rm forced}\).  It then constructs a
separated-matching family proving that the component factors and
support-type factorials still have a superexponential positive
post-norm grouping fibre.  It refines SOCC to PSCSC.

V8 does not prove PSCSC, SOCC, SISC, the simple high-excess row,
arbitrary repeated transversal supports, full TOCC, full TCNC,
cyclic TIC, HLTC, full MCWC, OCRC, SCNGC, WPSC, GRLTC, LZTC,
PLRC, WSDC, BRSC, HWSFC, UGCGC, thermal connected WPRAC,
all-arity RGSC, the raw Wilson aggregate strong card, all-order
strong MTA, WUFC, CPM, cutoff removal, reflection positivity,
Osterwalder--Schrader reconstruction, a continuum Yang--Mills
measure, or a positive Hamiltonian gap.  The full Yang--Mills
mass gap has not been proved.
\end{firewall}

\begin{theorem}[Sixth-series V8 result ledger]
\label{thm:s6v8-results}
V8:
\begin{enumerate}[label=\textup{(V8.\arabic*)},leftmargin=5.5em]
\item proves the resolved-component incidence-tree theorem;
\item identifies forced excess with forgotten grouping;
\item proves the two-colour component calibration;
\item proves exact strong-gate component splitting;
\item identifies the component Cayley factor and its scale degree;
\item constructs many separated component matchings;
\item proves the dense post-norm grouping no-go;
\item corrects the scope of component resolution; and
\item formulates PSCSC and proves its equivalence to the open
      part of SOCC.
\end{enumerate}
\end{theorem}

\begin{proof}
Use
\Cref{thm:s6v8-incidence-tree,
cor:s6v8-forgotten-grouping,
lem:s6v8-two-colour,
thm:s6v8-gate-splitting,
cor:s6v8-scale-degree,
lem:s6v8-separated-matchings,
thm:s6v8-grouping-no-go,
corr:s6v8-SOCC,
prop:s6v8-PSCSC}.
\end{proof}

\paragraph{Parent-version record.}
V8 was formed from
\texttt{Renewal\_Yang\_Mills\_Mass\_Gap\_Research\_Notes\_V7.tex}.
The V7 parent had \(3\,339\) logical source lines, \(116\,872\)
bytes, and SHA--256
\[
 \texttt{26a0cbbbcc83b8261f65a7411f518a24442a015cb4462d6738bda7ed43ece6bb}.
\]
The next compulsory companion audit is V10.

\clearpage
\chapter{Sixth-series V9: Duhamel calibration and the
tree--forest overlay}
\label{ch:s6v9-overlay}

\section{The exact distinct-type Duhamel coefficient}
\label{sec:s6v9-Duhamel}

No companion audit is due at V9.

Let \(V,H_1,\ldots,H_q\) be bounded operators.  For
\[
 \Delta_q
 :=
 \{(t_0,\ldots,t_q):t_j\ge0,\ \sum_{j=0}^qt_j=1\},
\]
write \(dt\) for its \(q\)-dimensional simplex measure.

\begin{lemma}[Polarized Duhamel coefficient]
\label{lem:s6v9-Duhamel}
One has
\begin{align}
 &[z_1\cdots z_q]\exp\left\{V+\sum_{\nu=1}^qz_\nu H_\nu\right\}
 \notag\\
 &\quad=
 \sum_{\sigma\in S_q}
 \int_{\Delta_q}
 e^{t_0V}H_{\sigma(1)}e^{t_1V}
 \cdots
 H_{\sigma(q)}e^{t_qV}\,dt.                              \label{eq:s6v9-Duhamel}
\end{align}
The simplex has volume \(1/q!\).
\end{lemma}

\begin{proof}
Iterate the first Duhamel formula
\[
 {d\over d\epsilon}e^{V+\epsilon H}\big|_{\epsilon=0}
 =
 \int_0^1e^{(1-s)V}He^{sV}\,ds
\]
in the \(q\) independent directions.  The insertion times are
distinct off a null set.  Partition their cube by the \(q!\)
orders; after passing to the \(q+1\) gaps between successive
times, each order gives the displayed simplex integral.  The
standard simplex volume is \(1/q!\).
\end{proof}

\begin{proposition}[No global factorial for distinct commuting
types]
\label{prop:s6v9-no-factorial}
If \(V=0\) and \(H_1,\ldots,H_q\) commute, then
\begin{equation}
 [z_1\cdots z_q]\exp\left\{\sum_{\nu=1}^qz_\nu H_\nu\right\}
 =
 \prod_{\nu=1}^qH_\nu.                                    \label{eq:s6v9-no-factorial}
\end{equation}
Thus the \(1/q!\) simplex volume is cancelled exactly by the
\(q!\) orders.  Only multiplicities of identical type variables
produce the factors \(1/\nu_A!\) recorded in f5-V87.
\end{proposition}

\begin{proof}
Every integrand in \eqref{eq:s6v9-Duhamel} is the same product
\(\prod_\nu H_\nu\).  Multiply it by \(q!\) simplex orders and
the volume \(1/q!\).  The last assertion follows by identifying
\(\nu_A\) variables: permutations among identical insertions no
longer produce distinct polarized monomials.
\end{proof}

\begin{corollary}[Generic Duhamel algebra does not prove PSCSC]
\label{cor:s6v9-generic-no-go}
The Duhamel simplex, cb submultiplicativity, and the known
support-type factorials cannot by themselves remove the
separated-matching fibre of
\Cref{thm:s6v8-grouping-no-go}.
\end{corollary}

\begin{proof}
In that family every support type is distinct, so all type
factorials equal one.  Take the scalar commuting calibration
\(V=0\), \(H_\nu\ge0\).  By
\Cref{prop:s6v9-no-factorial}, the exact Duhamel coefficient is
the product of the insertions and supplies no further factor.
The lower bound \eqref{eq:s6v8-grouping-no-go} therefore remains
in the corresponding positive proxy.
\end{proof}

\begin{warning}[Scope of the calibration]
\label{warn:s6v9-calibration}
\Cref{cor:s6v9-generic-no-go} is not a no-go theorem for the
specific compact-group one-link source.  It says that a proof
must use more than the universal Duhamel identity and separate
operator norms.  In particular, one may retain the component
Cayley polynomials as a global tree-valued object instead of
seeking an operator cancellation which need not exist.
\end{warning}

\section{All forgotten groupings are one forest}
\label{sec:s6v9-grouping-forest}

Let \({\cal C}\) be the resolved component set of V8, with
\[
 N:=|{\cal C}|,\qquad x_C:=B_C.
\]
A grouping is a partition
\(\Pi=\{G_1,\ldots,G_q\}\) of \({\cal C}\).  The block \(G_\nu\)
collects the connected components owned by occurrence \(\nu\).
For a nonempty \(G\subseteq{\cal C}\), define
\begin{equation}
 \eta(G)
 :=
 \left(\prod_{C\in G}x_C\right)
 \left(\sum_{C\in G}x_C\right)^{|G|-2},                  \label{eq:s6v9-eta-G}
\end{equation}
with \(\eta(\{C\})=1\).  This is the factor of
\Cref{thm:s6v8-gate-splitting}.

\begin{theorem}[Grouping-forest bijection]
\label{thm:s6v9-grouping-forest}
Data consisting of:
\begin{enumerate}[label=\textup{(\roman*)},leftmargin=3.5em]
\item a partition \(\Pi\) of \({\cal C}\); and
\item for every nonsingleton \(G\in\Pi\), a tree
      \(L_G\in\mathfrak T(G)\)
\end{enumerate}
are in bijection with forests \(L\) on vertex set \({\cal C}\).
Under this bijection,
\begin{equation}
 \prod_{G\in\Pi}\eta(G)
 =
 \sum_{\substack{L:\ {\cal K}(L)=\Pi}}
 \prod_{\{C,D\}\in L}x_Cx_D,                             \label{eq:s6v9-grouping-weight}
\end{equation}
and every such \(L\) has
\begin{equation}
 |E(L)|=N-|\Pi|
 =\sum_{\nu=1}^q(c_\nu-1)
 =g_{\rm forced}.                                        \label{eq:s6v9-grouping-size}
\end{equation}
\end{theorem}

\begin{proof}
The disjoint union \(L=\bigcup_{G\in\Pi}L_G\) is a forest whose
connected components are exactly the blocks of \(\Pi\).
Conversely, the connected components of a forest give \(\Pi\)
and their induced trees give the \(L_G\).  These constructions
are inverse.

Expand every \(\eta(G)\) by the weighted Cayley identity and
multiply; this gives
\eqref{eq:s6v9-grouping-weight}.  Finally, a forest on \(N\)
vertices with \(q\) components has \(N-q\) edges.  Use
\eqref{eq:s6v8-forgotten-grouping}.
\end{proof}

\begin{corollary}[The matching obstruction is a forest sector]
\label{cor:s6v9-matching-forest}
The separated matching family of V8 is the row in which \(L\) is
a perfect matching on the atomic resolved components.  Its
superexponential scalar fibre is therefore not an unexplained set
partition: it is the positive weight of a large family of
grouping forests which was summed without a global tree output.
\end{corollary}

\begin{proof}
Every group has two members and its Cayley tree is its unique
edge.  Thus the union of the group trees is exactly the matching.
\end{proof}

\section{Lifting the grouping forest to label edges}
\label{sec:s6v9-lift}

\begin{lemma}[Normalized endpoint lift of a component edge]
\label{lem:s6v9-component-lift}
For an edge \(\{C,D\}\in E(L)\), choose
\((i,j)\in C\times D\) with probability
\begin{equation}
 p_{CD}(i,j):={b_ib_j\over B_CB_D}.                       \label{eq:s6v9-component-lift}
\end{equation}
Then \(p_{CD}\) is a probability and
\[
 B_CB_D\,p_{CD}(i,j)=b_ib_j.
\]
Independent lifts of all edges of \(L\) have total mass one and
convert the grouping-forest monomial exactly into its label-edge
monomial.
\end{lemma}

\begin{proof}
Sum the numerator:
\[
 \sum_{i\in C,\ j\in D}b_ib_j=B_CB_D.
\]
The remaining assertions follow by multiplication over the
forest edges.
\end{proof}

\begin{remark}[The lift may create cycles with the outer tree]
\label{rem:s6v9-lift-cycles}
The forest \(L\) is acyclic on resolved component vertices, but
each pair of its vertices is already joined through the resolved
incidence tree of \Cref{thm:s6v8-incidence-tree}.  Thus every
lifted edge creates one cycle in their overlay.  This is where
the forced incidence excess resides; it must be treated as an
overlay, not by calling \(L\) itself a surplus graph.
\end{remark}

\section{The exact tree--forest overlay}
\label{sec:s6v9-overlay}

Let \({\cal I}\) be the resolved incidence tree, and regard a
grouping edge \(\{C,D\}\) as an additional edge between its two
left vertices.  Define
\[
 {\cal O}:={\cal I}\cup L.
\]

\begin{theorem}[Overlay excess identity]
\label{thm:s6v9-overlay-excess}
The overlay \({\cal O}\) is connected and has cyclomatic excess
\begin{equation}
 g({\cal O})
 =
 |E({\cal O})|-|V({\cal O})|+1
 =
 |E(L)|
 =
 g_{\rm forced}.                                         \label{eq:s6v9-overlay-excess}
\end{equation}
Deleting the grouping edges recovers the canonical spanning tree
\({\cal I}\).
\end{theorem}

\begin{proof}
\({\cal I}\) is a spanning tree of the overlay by
\Cref{thm:s6v8-incidence-tree}.  Adding \(|E(L)|\) further edges
preserves connectedness and increases the cyclomatic excess by
one per edge.  Apply \eqref{eq:s6v9-grouping-size}.
\end{proof}

\begin{lemma}[Each fundamental cycle is canonical]
\label{lem:s6v9-fundamental-cycle}
For a grouping edge \(\ell=\{C,D\}\), its fundamental cycle in
\({\cal O}\) is
\begin{equation}
 \ell\cup{\rm path}_{\cal I}(C,D).                        \label{eq:s6v9-fundamental-cycle}
\end{equation}
The collection of these cycles is a basis of the cycle space of
\({\cal O}\), indexed by \(E(L)\).
\end{lemma}

\begin{proof}
Adding one non-tree edge to the spanning tree \({\cal I}\)
creates the unique displayed cycle.  The fundamental cycles of
distinct non-tree edges are linearly independent because each
contains its own non-tree edge and no other fundamental cycle
contains that edge.  Their number equals the cycle rank from
\Cref{thm:s6v9-overlay-excess}, so they form a basis.
\end{proof}

\begin{assessment}[The right object after the no-go]
\label{ass:s6v9-right-object}
The grouping data should not be scalarized and need not be
cancelled by Duhamel order.  They are exactly one weighted forest
\(L\) laid over one canonical incidence tree \({\cal I}\), with
one canonical fundamental cycle per defect unit.  This converts
PSCSC from an unspecified operator cancellation into a
tree-valued overlay allocation problem.
\end{assessment}

\section{Tree--forest lift allocation}
\label{sec:s6v9-TFLAC}

\begin{definition}[Tree--forest lift allocation card, TFLAC]
\label{def:s6v9-TFLAC}
TFLAC is a coefficientwise allocation of the weighted overlays
\({\cal I}\cup L\), with their normalized endpoint lifts, such
that:
\begin{enumerate}[label=\textup{(\roman*)},leftmargin=3.5em]
\item the grouping forest \(L\) is retained before norms;
\item a Kruskal order treats its canonical fundamental cycles
      simultaneously;
\item the allocation outputs an ordinary global Cayley tree or
      the already proved primitive incidence-hypertree object;
\item every deleted overlay edge is paid by a retained endpoint
      monomial of the output, with no inverse small mark; and
\item the inverse fibre over one output is at most \(K^m\).
\end{enumerate}
\end{definition}

\begin{proposition}[TFLAC proves PSCSC]
\label{prop:s6v9-TFLAC}
TFLAC, together with
\Cref{thm:s6v8-gate-splitting,
thm:s6v9-grouping-forest,
lem:s6v9-component-lift}, proves PSCSC and hence SOCC and SISC.
\end{proposition}

\begin{proof}
Expand every disconnected gate into its resolved component
strong factors and \(\eta\) by
\Cref{thm:s6v8-gate-splitting}.  Expand all \(\eta\)'s jointly
as the single grouping forest of
\Cref{thm:s6v9-grouping-forest}, then lift its edges by the
normalized law of \Cref{lem:s6v9-component-lift}.  The remaining
resolved supports form the incidence tree \({\cal I}\).  TFLAC
performs the only unsummed operation, the overlay allocation,
with a \(K^m\) fibre and the correct strong output.  This is
PSCSC.  Apply \Cref{prop:s6v8-PSCSC,prop:s6v7-SOCC}.
\end{proof}

\begin{corollary}[Zero-defect row]
\label{cor:s6v9-zero-defect}
If \(g_{\rm forced}=0\), then \(L\) is empty and TFLAC holds
with fibre one.  This recovers, without a new proof, the already
closed primitive incidence-hypertree row.
\end{corollary}

\begin{proof}
By \eqref{eq:s6v9-grouping-size}, \(L\) has no edge.  The overlay
is the incidence tree \({\cal I}\), which is precisely the
hypertree object in the cumulative manifest.
\end{proof}

\begin{assessment}[Exact live bottleneck after V9]
\label{ass:s6v9-bottleneck}
The generic Duhamel factorial route is closed negatively.  The
complete forgotten occurrence grouping is now an explicit
weighted forest, and its collision with the resolved incidence
tree has a canonical fundamental-cycle basis.  TFLAC is the live
problem: a simultaneous weighted basis exchange which retains
the label endpoint lifts and has exponential inverse fibre.
\end{assessment}

\begin{programme}[V10: audit and the one-cycle overlay]
\label{prog:s6v9-V10}
\begin{enumerate}[leftmargin=3em]
\item Perform the compulsory newest-SM/newest-NS audit.
\item Prove TFLAC for \(|E(L)|=1\), using the unique
      \({\cal I}\)-path and a weighted endpoint exchange.
\item Test whether rooting that path at the lifted endpoint gives
      a probability distribution over the edge to be deleted.
\item Compare the one-cycle formula with the archived V94
      locally overlapped cycle theorem without duplicating it.
\item State the exact obstruction to iterating the exchange when
      fundamental cycles overlap.
\end{enumerate}
\end{programme}

\begin{firewall}[Claim boundary after sixth-series V9]
\label{fire:s6v9-boundary}
V9 derives the exact polarized Duhamel coefficient, proves that
distinct commuting types have no global occurrence factorial,
and closes that generic factorial route negatively.  It proves
that all occurrence groupings and their component Cayley factors
are exactly one weighted grouping forest, gives its normalized
endpoint lift, proves the tree--forest overlay excess identity
and canonical fundamental-cycle basis, and formulates TFLAC.

V9 does not prove positive-defect TFLAC, PSCSC, SOCC, SISC, the
simple high-excess row, arbitrary repeated transversal supports,
full TOCC, full TCNC, cyclic TIC, HLTC, full MCWC, OCRC, SCNGC,
WPSC, GRLTC, LZTC, PLRC, WSDC, BRSC, HWSFC, UGCGC, thermal
connected WPRAC, all-arity RGSC, the raw Wilson aggregate strong
card, all-order strong MTA, WUFC, CPM, cutoff removal,
reflection positivity, Osterwalder--Schrader reconstruction, a
continuum Yang--Mills measure, or a positive Hamiltonian gap.
The full Yang--Mills mass gap has not been proved.
\end{firewall}

\begin{theorem}[Sixth-series V9 result ledger]
\label{thm:s6v9-results}
V9:
\begin{enumerate}[label=\textup{(V9.\arabic*)},leftmargin=5.5em]
\item proves the exact polarized Duhamel coefficient;
\item proves the distinct-type factorial no-go;
\item proves the grouping-forest bijection and weight identity;
\item identifies the matching family as a forest sector;
\item proves normalized component-edge endpoint lifts;
\item proves the overlay excess identity;
\item proves the canonical fundamental-cycle basis;
\item formulates TFLAC and proves its sufficiency; and
\item closes the zero-defect TFLAC row.
\end{enumerate}
\end{theorem}

\begin{proof}
Use
\Cref{lem:s6v9-Duhamel,
prop:s6v9-no-factorial,
thm:s6v9-grouping-forest,
cor:s6v9-matching-forest,
lem:s6v9-component-lift,
thm:s6v9-overlay-excess,
lem:s6v9-fundamental-cycle,
prop:s6v9-TFLAC,
cor:s6v9-zero-defect}.
\end{proof}

\paragraph{Parent-version record.}
V9 was formed from
\texttt{Renewal\_Yang\_Mills\_Mass\_Gap\_Research\_Notes\_V8.tex}.
The V8 parent had \(3\,804\) logical source lines, \(133\,547\)
bytes, and SHA--256
\[
 \texttt{89622710a32aab0096d23b43c789e0fadd203d791688ef3d886f6d8e79db0be0}.
\]
The next compulsory companion audit is V10.

\clearpage
\chapter{Sixth-series V10: companion audit and weighted overlay
basis exchange}
\label{ch:s6v10-exchange}

\section{Compulsory companion audit}
\label{sec:s6v10-audit}

The current companion snapshots used in this audit are:
\begin{center}
\small
\begin{tabular}{p{0.25\textwidth}p{0.68\textwidth}}
Standard Model--Einstein &
\texttt{Renewal\_SM\_Einstein\_Continuum\_Research\_Notes\_V32.tex},
\(526\,117\) bytes, \(13\,100\) logical source lines, SHA--256
\texttt{5509DD0AE6F3FB30997C68FABBF0AF429F2E775A356420BD67D73F75932C2B5F}.
\\ [0.4em]
Navier--Stokes &
\texttt{Renewal\_NS\_Stability\_Research\_Notes\_V38.tex},
\(381\,843\) bytes, \(7\,551\) logical source lines, SHA--256
\texttt{59932A358FC764EB14F71CD6BF6D01A25268B114977677E22C170BFB066F6178}.
\end{tabular}
\end{center}

\begin{audit}[Standard Model--Einstein V32]
\label{audit:s6v10-SM}
SM V32 proves that convex interpolation of two gapped,
unitarily equivalent lattice stencils can close the gap.  It
therefore abandons local chart-stencil averaging and returns to
one global subdivision with a finite catalogue of frozen star
types.  Conditional on the finite star gap card, its preceding
IMS estimate gives a global overlap gap.

The applicable structural point is that separately controlled
local objects should not be interpolated across an interface
when their leading-scale difference is not small.  For TFLAC,
the analogue is to retain one global augmented overlay graph and
choose one spanning tree of it; independently exchanging each
fundamental cycle can destroy the compatibility of later cycles.
No spin-mesh estimate is imported.
\end{audit}

\begin{audit}[Navier--Stokes V38]
\label{audit:s6v10-NS}
NS V38 obtains an exact sign at an extremizer of the defect, but
its explicit estimate stops at a nonlocal pressure pairing which
cannot be controlled from the available Gaussian local norms.
The exact identity is retained and the obstruction is named
rather than bounded by an unrelated local estimate.

The corresponding rule here is to keep the entire
\(T_0\)-path in the fundamental cycle.  An endpoint estimate
alone does not determine which old edge can be exchanged without
affecting the rest of the overlay.  No fluid estimate is imported.
\end{audit}

\begin{audit}[Direction selected by the V10 audit]
\label{audit:s6v10-direction}
V10 proves the one-cycle exchange from the whole fundamental
cycle and then states its simultaneous Kruskal form on the whole
augmented graph.  It does not interpolate separate cycle
exchanges.
\end{audit}

\section{One grouping edge}
\label{sec:s6v10-one-edge}

Let \(T_0\) be the sampled outer tree on \(X\).  Suppose the
grouping forest of V9 has one edge \(\ell=\{C,D\}\).  Its exact
weight is
\[
 B_CB_D=\sum_{i\in C,\ j\in D}b_ib_j.
\]
Use the normalized endpoint lift of
\Cref{lem:s6v9-component-lift}, and call the lifted edge
\[
 h=\{i,j\}.
\]

\begin{lemma}[Coincident lift]
\label{lem:s6v10-coincident}
If \(h\in E(T_0)\), then
\begin{equation}
 \left(\prod_{e\in T_0}b_{e_-}b_{e_+}\right)b_ib_j
 \le
 \prod_{e\in T_0}b_{e_-}b_{e_+}.                         \label{eq:s6v10-coincident}
\end{equation}
Thus a grouping-edge occurrence coincident with an old tree edge
is absorbed without changing the output tree.
\end{lemma}

\begin{proof}
Every \(0<b_i,b_j<1\), so \(b_ib_j\le1\).
\end{proof}

Assume now \(h\notin E(T_0)\).  The graph \(T_0+h\) has the
unique cycle
\[
 h\cup{\rm path}_{T_0}(i,j).
\]
Fix once and for all a total order on old tree-edge occurrences
and let \(f(T_0,h)\) be the first edge of the displayed
\(T_0\)-path.

\begin{lemma}[One-cycle tree exchange]
\label{lem:s6v10-one-exchange}
Put
\begin{equation}
 U:=T_0-f(T_0,h)+h.                                      \label{eq:s6v10-U}
\end{equation}
Then \(U\) is a tree on \(X\), and
\begin{align}
 &\left(\prod_{e\in T_0}b_{e_-}b_{e_+}\right)b_ib_j
 \notag\\
 &\qquad=
 \left(\prod_{e\in U}b_{e_-}b_{e_+}\right)
 b_{f_-}b_{f_+}
 \le
 \prod_{e\in U}b_{e_-}b_{e_+}.                           \label{eq:s6v10-weight-exchange}
\end{align}
\end{lemma}

\begin{proof}
Adding \(h\) creates the unique fundamental cycle.  Deleting an
edge \(f\) on that cycle leaves a connected graph with \(m-1\)
edges, hence a tree.  Cancel the common edge factors on the two
sides of the equality.  The last inequality uses
\(b_{f_-}b_{f_+}\le1\).
\end{proof}

\begin{lemma}[One-cycle inverse fibre]
\label{lem:s6v10-one-fibre}
For fixed resolved component/source data, the map from
\((T_0,h)\) to \(U\), with the coincident case included, has
fibre at most \(m^4+m^2\), hence at most \(K^m\).
\end{lemma}

\begin{proof}
In the noncoincident row, a preimage of \(U\) is determined by
the added edge \(h\in E(U)\) and the deleted edge
\(f\in\binom X2\setminus E(U)\); then
\(T_0=U-h+f\).  There are fewer than \(m^2\) choices for each,
and the requirements that \(T_0\) be a tree, that \(h\) lift
\(\{C,D\}\), and that the deterministic rule select \(f\) can
only reduce the count.  In the coincident row, choose the
coincident edge in fewer than \(m^2\) ways.  Finally every fixed
polynomial in \(m\) is bounded by \(K^m\) for \(m\ge2\).
\end{proof}

\begin{theorem}[One-cycle TFLAC]
\label{thm:s6v10-one-cycle}
TFLAC holds for \(|E(L)|=1\), conditional on the resolved
component/source data, with one exponential fibre.  When combined
with the already proved fixed-excess source allocation, the full
simple-support forced-defect-one row is closed.
\end{theorem}

\begin{proof}
Expand \(B_CB_D\) by the normalized endpoint lift.  The coincident
row is \Cref{lem:s6v10-coincident}; the noncoincident row is
\Cref{lem:s6v10-one-exchange}.  In both cases the input monomial
is bounded by one global output-tree monomial.  The inverse
exchange fibre is \Cref{lem:s6v10-one-fibre}.  All endpoint lifts
have total mass one, proving conditional TFLAC.

The remaining ownership/source data have incidence excess one.
They are contained in the fixed-excess simple-core theorem
recorded in \Cref{thm:s6v1-tree-closures}; that theorem supplies
their already established exponential source fibre.
\end{proof}

\begin{assessment}[Relation to the archived one-defect theorem]
\label{ass:s6v10-V89}
The edge-weight identity
\eqref{eq:s6v10-weight-exchange} is the same basis-exchange
algebra used in f5-V89.  Its payer is different.  V89 required a
fresh background tree to manufacture a within-block crossing
edge.  Here \(h\) is already present in the exact component
Cayley factor \(\eta\) of the source gate, so no freshness claim
is made or needed.
\end{assessment}

\section{Simultaneous exchange on an arbitrary overlay}
\label{sec:s6v10-simultaneous}

Lift every edge of a grouping forest \(L\) independently, retaining
parallel occurrences.  Let
\[
 {\cal G}:=T_0\cup L^{\rm lift}
\]
be the resulting connected multigraph.  Put \(g=|E(L)|\).
Then \({\cal G}\) has \(m-1+g\) edge occurrences.

\begin{lemma}[Global Kruskal basis exchange]
\label{lem:s6v10-global-Kruskal}
Give all edge occurrences distinct auxiliary priorities and run
decreasing-priority Kruskal on \({\cal G}\).  It produces a
unique spanning tree \(U\) and a rejected multiset \(R\) with
\(|R|=g\).  Moreover,
\begin{equation}
 \left(\prod_{e\in T_0}b_{e_-}b_{e_+}\right)
 \left(\prod_{h\in L^{\rm lift}}b_{h_-}b_{h_+}\right)
 =
 \left(\prod_{e\in U}b_{e_-}b_{e_+}\right)
 \left(\prod_{r\in R}b_{r_-}b_{r_+}\right)
 \le
 \prod_{e\in U}b_{e_-}b_{e_+}.                           \label{eq:s6v10-global-Kruskal}
\end{equation}
\end{lemma}

\begin{proof}
Kruskal returns a maximal acyclic subset.  Since
\({\cal G}\) is connected, that subset is a spanning tree with
\(m-1\) edge occurrences, leaving \(g\) rejected occurrences.
The equality merely partitions the complete edge multiset into
\(U\) and \(R\).  Every rejected edge weight is at most one.
\end{proof}

\begin{assessment}[Why the weight inequality is not yet the full
card]
\label{ass:s6v10-fibre}
For a fixed augmented graph, the Kruskal region is normalized and
the input-tree choice is paid by the exact partition of unity
recorded in \Cref{thm:s6v1-normalizations}.  However, after the
rejected set \(R\) is dropped, many different grouping forests
can have the same output tree.  The remaining issue is not a
termwise weight comparison; it is the inverse forest fibre over
\(U\).
\end{assessment}

\begin{lemma}[Crude fixed-excess inverse bound]
\label{lem:s6v10-fixed-g-fibre}
If the rejected occurrences are remembered only by their
unordered endpoint pairs, then for fixed \(U\) there are at most
\begin{equation}
 m^{2g}                                                    \label{eq:s6v10-fixed-g-fibre}
\end{equation}
possible rejected multisets of size \(g\), before the forest,
source-type, and validity restrictions are imposed.
\end{lemma}

\begin{proof}
Order the \(g\) occurrences temporarily and choose two endpoints
for each.  This gives at most \(m^{2g}\) lists.  Forgetting the
order and imposing all restrictions can only decrease the count.
\end{proof}

\begin{theorem}[Submacroscopic overlay row]
\label{thm:s6v10-submacro}
For
\[
 0\le g\le G_m:=\left\lfloor{m\over\log(m+e)}\right\rfloor,
\]
the inverse rejected-edge factor in
\eqref{eq:s6v10-fixed-g-fibre} is at most \(K^m\).
Consequently the TFLAC overlay exchange introduces no new
obstruction in the already closed simple-support
submacroscopic-excess row.
\end{theorem}

\begin{proof}
For \(g\le G_m\),
\[
 m^{2g}
 =
 \exp\{2g\log m\}
 \le
 \exp(2m).
\]
The normalized endpoint lifts, Kruskal regions, and output
Cayley-tree sum have unit mass or their established exponential
base.  This is precisely the size range already closed by the
fixed-excess theorem in the cumulative manifest.
\end{proof}

\section{The high-excess rejected-forest card}
\label{sec:s6v10-RFAC}

\begin{definition}[Rejected-forest allocation card, RFAC]
\label{def:s6v10-RFAC}
RFAC is the high-excess continuation of TFLAC.  For
\(g>G_m\), it must retain the rejected edge multiset \(R\) and
its exact monomial in
\eqref{eq:s6v10-global-Kruskal}, sum all grouping forests and
Kruskal regions jointly, and prove that their marginal over one
output tree has mass at most \(K^m\).  It may not replace \(R\)
by its \(m^{2g}\) endpoint census.
\end{definition}

\begin{proposition}[RFAC completes TFLAC]
\label{prop:s6v10-RFAC}
RFAC, together with
\Cref{thm:s6v10-one-cycle,thm:s6v10-submacro}, proves TFLAC and
hence PSCSC, SOCC, and SISC.
\end{proposition}

\begin{proof}
Split by \(g\).  The zero row is
\Cref{cor:s6v9-zero-defect}; the positive rows through \(G_m\)
are \Cref{thm:s6v10-submacro}, including the explicit
one-cycle calibration.  RFAC treats the complement.  Apply
\Cref{prop:s6v9-TFLAC}.
\end{proof}

\begin{assessment}[Exact live bottleneck after V10]
\label{ass:s6v10-bottleneck}
The weighted basis exchange itself is closed at every excess:
the full augmented monomial is always bounded by its global
Kruskal-tree monomial.  Low excess is also closed after its crude
inverse fibre.  Only the high-excess marginal of the rejected
grouping forest remains.  It must be summed with its weights,
because forgetting \(g>m/\log m\) endpoint pairs recreates a
superexponential fibre.
\end{assessment}

\begin{programme}[V11: rejected-forest marginal]
\label{prog:s6v10-V11}
\begin{enumerate}[leftmargin=3em]
\item Condition on the output Kruskal tree \(U\) and express all
      rejected grouping edges by their fundamental \(U\)-paths.
\item Use that the pre-Kruskal grouping edges form a forest on
      resolved component vertices.
\item Test a rooted-forest identity on the rejected edges with
      path bottlenecks as roots.
\item Retain the complete product
      \(\prod_{r\in R}b_{r_-}b_{r_+}\) before summing endpoints.
\item If a dense common-path fibre survives, return to the
      maximum-spanning-tree partition before fixing \(U\).
\end{enumerate}
\end{programme}

\begin{firewall}[Claim boundary after sixth-series V10]
\label{fire:s6v10-boundary}
V10 audits SM V32 and NS V38, proves the exact one-cycle grouping
edge exchange and its exponential inverse fibre, identifies its
relation to but distinct payer from f5-V89, proves the simultaneous
global Kruskal weight exchange for every overlay, and closes the
submacroscopic-excess overlay row.  It isolates RFAC as the exact
high-excess rejected-forest marginal.

V10 does not prove RFAC, high-excess TFLAC, PSCSC, SOCC, SISC,
the simple high-excess row, arbitrary repeated transversal
supports, full TOCC, full TCNC, cyclic TIC, HLTC, full MCWC,
OCRC, SCNGC, WPSC, GRLTC, LZTC, PLRC, WSDC, BRSC, HWSFC,
UGCGC, thermal connected WPRAC, all-arity RGSC, the raw Wilson
aggregate strong card, all-order strong MTA, WUFC, CPM, cutoff
removal, reflection positivity, Osterwalder--Schrader
reconstruction, a continuum Yang--Mills measure, or a positive
Hamiltonian gap.  The full Yang--Mills mass gap has not been
proved.
\end{firewall}

\begin{theorem}[Sixth-series V10 result ledger]
\label{thm:s6v10-results}
V10:
\begin{enumerate}[label=\textup{(V10.\arabic*)},leftmargin=6em]
\item performs the compulsory SM V32 and NS V38 audit;
\item proves coincident grouping-edge absorption;
\item proves the one-cycle weighted exchange;
\item proves its inverse fibre and one-cycle TFLAC;
\item proves simultaneous global Kruskal exchange;
\item proves the crude fixed-excess inverse bound;
\item closes the submacroscopic overlay row; and
\item formulates RFAC and proves its sufficiency.
\end{enumerate}
\end{theorem}

\begin{proof}
Use
\Cref{audit:s6v10-SM,
audit:s6v10-NS,
lem:s6v10-coincident,
lem:s6v10-one-exchange,
lem:s6v10-one-fibre,
thm:s6v10-one-cycle,
lem:s6v10-global-Kruskal,
lem:s6v10-fixed-g-fibre,
thm:s6v10-submacro,
prop:s6v10-RFAC}.
\end{proof}

\paragraph{Parent-version record.}
V10 was formed from
\texttt{Renewal\_Yang\_Mills\_Mass\_Gap\_Research\_Notes\_V9.tex}.
The V9 parent had \(4\,214\) logical source lines, \(148\,168\)
bytes, and SHA--256
\[
 \texttt{df86458225b0a5c6b4138066784dc82eaf837b0d03e192d7ddcab342565f3d5a}.
\]
The next compulsory companion audit is V15.

\clearpage
\chapter{Sixth-series V11: exact grouping-forest determinants
and a diffuse-mass closure}
\label{ch:s6v11-determinant}

\section{The weighted Laplacian of resolved components}
\label{sec:s6v11-Laplacian}

No companion audit is due at V11.

Let the resolved component set be
\({\cal C}=\{1,\ldots,N\}\), with weights
\[
 x_\alpha:=B_{C_\alpha}>0,\qquad
 X:=\sum_{\alpha=1}^Nx_\alpha.
\]
Let \({\cal L}_x\) be the complete-graph weighted Laplacian
\begin{equation}
 ({\cal L}_x)_{\alpha\beta}
 =
 \begin{cases}
 x_\alpha(X-x_\alpha),&\alpha=\beta,\\
 -x_\alpha x_\beta,&\alpha\ne\beta.
 \end{cases}                                              \label{eq:s6v11-Laplacian}
\end{equation}
Thus every component edge \(\{\alpha,\beta\}\) has the exact
grouping weight \(x_\alpha x_\beta\).

\begin{lemma}[Positivity and trace]
\label{lem:s6v11-Laplacian}
\({\cal L}_x\) is positive semidefinite, has kernel spanned by
the constant vector, and
\begin{equation}
 \operatorname{Tr}{\cal L}_x
 =
 X^2-\sum_{\alpha=1}^Nx_\alpha^2
 =
 2\sum_{\alpha<\beta}x_\alpha x_\beta.                   \label{eq:s6v11-trace}
\end{equation}
\end{lemma}

\begin{proof}
For \(v\in\mathbb R^N\),
\[
 v^T{\cal L}_xv
 =
 \sum_{\alpha<\beta}x_\alpha x_\beta
 (v_\alpha-v_\beta)^2\ge0.
\]
Since all edge weights are positive, equality forces all
coordinates equal.  Summing the diagonal entries gives
\eqref{eq:s6v11-trace}.
\end{proof}

\section{Two exact forest identities}
\label{sec:s6v11-identities}

For a forest \(L\) on \({\cal C}\), write
\[
 w_x(L):=\prod_{\{\alpha,\beta\}\in L}x_\alpha x_\beta,
\qquad
 {\cal K}(L):=\text{its component partition}.
\]
For \(D\in{\cal K}(L)\), put \(X_D:=\sum_{\alpha\in D}x_\alpha\).

\begin{theorem}[Ghost-root mass identity]
\label{thm:s6v11-ghost}
For every \(t\ge0\),
\begin{equation}
 \sum_{L\ {\rm forest}}
 w_x(L)t^{|{\cal K}(L)|}
 \prod_{D\in{\cal K}(L)}X_D
 =
 t\left(\prod_{\alpha=1}^Nx_\alpha\right)(X+t)^{N-1}.
                                                               \label{eq:s6v11-ghost}
\end{equation}
Consequently, for \(0\le g\le N-1\),
\begin{equation}
 \sum_{\substack{L\ {\rm forest}\\|E(L)|=g}}
 w_x(L)\prod_{D\in{\cal K}(L)}X_D
 =
 {N-1\choose g}
 \left(\prod_{\alpha=1}^Nx_\alpha\right)X^g.             \label{eq:s6v11-ghost-grade}
\end{equation}
\end{theorem}

\begin{proof}
Add one ghost vertex \(\ast\) of weight \(t\) to the component
vertices and sum all weighted trees on
\(\{\ast,1,\ldots,N\}\).  The weighted Cayley identity gives the
right side of \eqref{eq:s6v11-ghost}.

Delete \(\ast\) from such a tree.  The remaining graph is a
forest \(L\).  Every component \(D\) of \(L\) has exactly one
edge to \(\ast\): it has at least one by connectedness, and two
would create a cycle through the path in \(D\).  Summing the
endpoint of that ghost edge gives \(tX_D\).  Multiplication over
components gives the left side.

A forest with \(g\) edges has \(N-g\) components.  Extract the
coefficient of \(t^{N-g}\) from
\(t(\prod x_\alpha)(X+t)^{N-1}\) to obtain
\eqref{eq:s6v11-ghost-grade}.
\end{proof}

\begin{assessment}[The exact root-mass ledger]
\label{ass:s6v11-root-mass}
The grouping weight \(w_x(L)\) becomes a single weighted Cayley
polynomial after multiplication by one raw component mass
\(X_D\) for every source-occurrence group \(D\).  The actual
post-IGMC grouping forest does not contain these factors.  Thus
\eqref{eq:s6v11-ghost} identifies the exact missing root-mass
ledger; it does not authorize inserting it.  When some \(X_D\)
is small, dividing by the product of group masses is forbidden.
\end{assessment}

\begin{theorem}[Matrix--forest determinant identity]
\label{thm:s6v11-matrix-forest}
\begin{equation}
 \det(tI+{\cal L}_x)
 =
 \sum_{L\ {\rm forest}}
 t^{|{\cal K}(L)|}w_x(L)
 \prod_{D\in{\cal K}(L)}|D|.                             \label{eq:s6v11-matrix-forest}
\end{equation}
\end{theorem}

\begin{proof}
Adjoin a ghost vertex and give every ghost edge
\(\{\ast,\alpha\}\) weight \(t\), independent of \(\alpha\).
Deleting the ghost from a spanning tree again produces a forest
with exactly one ghost edge into each component.  There are
\(|D|\) choices for that edge in component \(D\), producing
\(t^{|{\cal K}(L)|}\prod_D|D|\).

By the matrix-tree theorem, the sum of these augmented spanning
tree weights is the principal Laplacian minor obtained by deleting
the ghost row and column.  That minor is
\(tI+{\cal L}_x\), proving the identity.
\end{proof}

\begin{corollary}[Exact graded rooted-forest coefficient]
\label{cor:s6v11-graded}
Let \(0=\lambda_1\le\lambda_2\le\cdots\le\lambda_N\) be the
eigenvalues of \({\cal L}_x\).  Then
\begin{align}
 &\sum_{\substack{L\ {\rm forest}\\|E(L)|=g}}
 w_x(L)\prod_{D\in{\cal K}(L)}|D|
 =
 e_g(\lambda_2,\ldots,\lambda_N),                         \label{eq:s6v11-graded}\\
 &\sum_{\substack{L\ {\rm forest}\\|E(L)|=g}}
 w_x(L)
 \le
 {(\operatorname{Tr}{\cal L}_x)^g\over g!}
 \le {X^{2g}\over g!}.                                    \label{eq:s6v11-forest-bound}
\end{align}
\end{corollary}

\begin{proof}
Since \(\lambda_1=0\),
\[
 \det(tI+{\cal L}_x)
 =
 t\prod_{j=2}^N(t+\lambda_j).
\]
The coefficient of \(t^{N-g}\) is
\(e_g(\lambda_2,\ldots,\lambda_N)\), and comparison with
\Cref{thm:s6v11-matrix-forest} gives
\eqref{eq:s6v11-graded}.  Since every \(|D|\ge1\), dropping the
root-count product decreases the sum.  Finally, for nonnegative
numbers,
\[
 e_g(\lambda_2,\ldots,\lambda_N)
 \le {(\sum_{j=2}^N\lambda_j)^g\over g!},
\]
and \(\sum_{j=2}^N\lambda_j=\operatorname{Tr}{\cal L}_x\le X^2\).
\end{proof}

\begin{assessment}[The determinant supplies one factorial]
\label{ass:s6v11-one-factorial}
The crude endpoint census in V10 was \(m^{2g}\).  The exact
forest condition improves the corresponding weighted sum to
\((\operatorname{Tr}{\cal L}_x)^g/g!\).  This is a genuine
factorial, obtained from graphic acyclicity rather than Duhamel
order.  It is still insufficient when
\(\operatorname{Tr}{\cal L}_x\) is of order \(m^2\) and
\(g\) is of order \(m\).
\end{assessment}

\section{A diffuse-component closure}
\label{sec:s6v11-diffuse}

\begin{definition}[Effective component-pair mass]
\label{def:s6v11-pair-mass}
Set
\begin{equation}
 \Sigma_x
 :=
 \operatorname{Tr}{\cal L}_x
 =
 X^2-\sum_\alpha x_\alpha^2.                             \label{eq:s6v11-pair-mass}
\end{equation}
\end{definition}

\begin{theorem}[Diffuse high-excess RFAC row]
\label{thm:s6v11-diffuse}
Fix \(C_0<\infty\).  At every excess \(1\le g\le N-1\), the
complete positive grouping-forest sum obeys a \(K^m\) bound
whenever
\begin{equation}
 \Sigma_x\le C_0g.                                       \label{eq:s6v11-diffuse-condition}
\end{equation}
Consequently RFAC holds on this diffuse-pair-mass row without
discarding any endpoint weight.
\end{theorem}

\begin{proof}
By \eqref{eq:s6v11-forest-bound} and
\(g!\ge(g/e)^g\),
\[
 \sum_{\substack{L\ {\rm forest}\\|E(L)|=g}}w_x(L)
 \le
 { (C_0g)^g\over g!}
 \le(eC_0)^g
 \le K^m,
\]
because \(g\le N-1\le m-2\) for nontrivial resolved
edge-components.  Keep the sampled outer tree \(T_0\) as the
global output tree.  All endpoint lifts, IGMC kernels, and
outer-tree probabilities remain normalized, so the displayed
forest sum is the only new factor.
\end{proof}

\begin{corollary}[A simpler sufficient condition]
\label{cor:s6v11-X-condition}
The conclusion holds if
\begin{equation}
 X\le C_1\sqrt g                                          \label{eq:s6v11-X-condition}
\end{equation}
for a fixed \(C_1\).
\end{corollary}

\begin{proof}
\(\Sigma_x\le X^2\le C_1^2g\); apply
\Cref{thm:s6v11-diffuse}.
\end{proof}

\begin{corollary}[Mass concentration in the unresolved row]
\label{cor:s6v11-concentrated}
Every RFAC history not closed by
\Cref{thm:s6v11-diffuse} satisfies
\begin{equation}
 \Sigma_x>C_0g.                                          \label{eq:s6v11-concentrated}
\end{equation}
Equivalently, its total pair mass between distinct resolved
components is large compared with its forest rank.
\end{corollary}

\begin{proof}
This is the complementary case.
\end{proof}

\section{Why the component mass is not the original total mark}
\label{sec:s6v11-component-mass}

\begin{lemma}[Resolved mass as a coloured incidence degree]
\label{lem:s6v11-resolved-mass}
Let \(d_{\rm col}(i)\) be the number of resolved owned components
whose vertex sets contain the label \(i\).  Then
\begin{equation}
 X=\sum_{\alpha=1}^NB_{C_\alpha}
 =\sum_{i\in X}d_{\rm col}(i)b_i,                         \label{eq:s6v11-resolved-mass}
\end{equation}
and
\begin{equation}
 1\le d_{\rm col}(i)\le d_{T_0}(i)
 \quad\hbox{for every nonisolated label }i.               \label{eq:s6v11-col-degree}
\end{equation}
\end{lemma}

\begin{proof}
Interchange the component and label sums to obtain the first
identity.  At \(i\), all incident outer-tree edges of one colour
belong to the same resolved component, so
\(d_{\rm col}(i)\) is the number of distinct colours incident at
\(i\).  This is at most the number \(d_{T_0}(i)\) of incident
edges and at least one.
\end{proof}

\begin{warning}[Large resolved mass is not the high-\(B\) region]
\label{warn:s6v11-X-not-B}
Although \(B=\sum_i b_i\), the mass
\(X=\sum_i d_{\rm col}(i)b_i\) contains coloured incidence
degrees.  It can be much larger than \(B\) near a high-degree
tree vertex.  Therefore
\eqref{eq:s6v11-concentrated} cannot be sent to the previously
closed high-total-mark region without an additional degree
argument.
\end{warning}

\section{The concentrated component-Laplacian card}
\label{sec:s6v11-CCLC}

\begin{definition}[Concentrated component-Laplacian card, CCLC]
\label{def:s6v11-CCLC}
CCLC is the remaining RFAC row
\(\Sigma_x>C_0g\).  It must combine:
\begin{enumerate}[label=\textup{(\roman*)},leftmargin=3.5em]
\item the exact determinant coefficient
      \(e_g(\lambda_2,\ldots,\lambda_N)\);
\item the coloured incidence degrees
      \eqref{eq:s6v11-resolved-mass};
\item the outer weighted Cayley probability of \(T_0\); and
\item the thermal partition,
\end{enumerate}
before replacing \(X\) or \(\Sigma_x\) by a scalar power of \(m\).
The target is a \(K^m\) marginal on the concentrated row.
\end{definition}

\begin{proposition}[CCLC completes RFAC]
\label{prop:s6v11-CCLC}
CCLC together with
\Cref{thm:s6v10-submacro,thm:s6v11-diffuse} proves RFAC and
hence TFLAC, PSCSC, SOCC, and SISC.
\end{proposition}

\begin{proof}
Submacroscopic excess is V10.  At larger excess, split by
\(\Sigma_x\le C_0g\) and its complement.  The first row is
\Cref{thm:s6v11-diffuse}; CCLC is the second.  Apply
\Cref{prop:s6v10-RFAC,prop:s6v9-TFLAC}.
\end{proof}

\begin{assessment}[Exact live bottleneck after V11]
\label{ass:s6v11-bottleneck}
The grouping forest now has an exact determinant law.  This
closes every row in which its component-pair trace is at most
linear in the excess and supplies a \(1/g!\) factor everywhere.
The surviving CCLC row has large coloured incidence mass.  Its
resolution must use the correlation between those degrees and
the sampled outer Cayley tree; a bound on \(X\) alone discards
the only remaining normalization.
\end{assessment}

\begin{programme}[V12: average coloured incidence under the
outer Cayley law]
\label{prog:s6v11-V12}
\begin{enumerate}[leftmargin=3em]
\item Compute the moment generating function of
      \(\sum_i d_{T_0}(i)b_i\) under the weighted Cayley--Pr\"ufer
      law.
\item Dominate \(X\) by that degree statistic using
      \Cref{lem:s6v11-resolved-mass}.
\item Derive an exponential tail for \(\Sigma_x>C_0g\) before
      conditioning on edge ownership colours.
\item Determine how much of that tail survives the squarefree
      edge-colour coefficient.
\item Close CCLC if the determinant factorial and the Cayley
      degree tail combine; otherwise isolate the colour-biased
      degree card.
\end{enumerate}
\end{programme}

\begin{firewall}[Claim boundary after sixth-series V11]
\label{fire:s6v11-boundary}
V11 proves the exact ghost-root mass identity, the complete
matrix--forest determinant identity, and its graded
elementary-symmetric coefficient.  It improves the grouping
forest bound to \(\Sigma_x^g/g!\), closes RFAC whenever
\(\Sigma_x=O(g)\), identifies the resolved component mass as a
coloured outer-tree degree statistic, and isolates CCLC.

V11 does not prove CCLC, RFAC in full, high-excess TFLAC, PSCSC,
SOCC, SISC, the simple high-excess row, arbitrary repeated
transversal supports, full TOCC, full TCNC, cyclic TIC, HLTC,
full MCWC, OCRC, SCNGC, WPSC, GRLTC, LZTC, PLRC, WSDC, BRSC,
HWSFC, UGCGC, thermal connected WPRAC, all-arity RGSC, the raw
Wilson aggregate strong card, all-order strong MTA, WUFC, CPM,
cutoff removal, reflection positivity, Osterwalder--Schrader
reconstruction, a continuum Yang--Mills measure, or a positive
Hamiltonian gap.  The full Yang--Mills mass gap has not been
proved.
\end{firewall}

\begin{theorem}[Sixth-series V11 result ledger]
\label{thm:s6v11-results}
V11:
\begin{enumerate}[label=\textup{(V11.\arabic*)},leftmargin=6em]
\item proves positivity and the trace of the component Laplacian;
\item proves the ghost-root forest identity;
\item proves the matrix--forest determinant identity;
\item derives its exact fixed-rank coefficient;
\item proves the \(\Sigma_x^g/g!\) forest bound;
\item closes the diffuse-pair-mass RFAC row;
\item identifies resolved mass with coloured incidence degree;
\item proves why it is not simply the high-\(B\) row; and
\item formulates CCLC and proves its sufficiency.
\end{enumerate}
\end{theorem}

\begin{proof}
Use
\Cref{lem:s6v11-Laplacian,
thm:s6v11-ghost,
thm:s6v11-matrix-forest,
cor:s6v11-graded,
thm:s6v11-diffuse,
lem:s6v11-resolved-mass,
warn:s6v11-X-not-B,
prop:s6v11-CCLC}.
\end{proof}

\paragraph{Parent-version record.}
V11 was formed from
\texttt{Renewal\_Yang\_Mills\_Mass\_Gap\_Research\_Notes\_V10.tex}.
The V10 parent had \(4\,605\) logical source lines, \(162\,208\)
bytes, and SHA--256
\[
 \texttt{0a7418e3fbc7a5ad9d9737fcbc84cb7a65d8644655c60d68a62543ffcc47bb96}.
\]
The next compulsory companion audit is V15.

\clearpage
\chapter{Sixth-series V12: the exact Cayley--Pr\"ufer degree law
and an upstream source-aggregation route}
\label{ch:s6v12-prufer}

\section{The outer weighted Cayley law}
\label{sec:s6v12-Cayley-law}

No companion audit is due at V12.

Let \(X=\{1,\ldots,m\}\), \(m\ge2\), let \(0<b_i<1\), and put
\[
 B:=\sum_{i=1}^m b_i,\qquad b_*:=\max_i b_i.
\]
On the labelled tree set \(\mathfrak T_X\), define
\begin{equation}
 \mu_b(T)
 :=
 {\prod_{\{i,j\}\in E(T)}b_ib_j
  \over
  (\prod_{i=1}^m b_i)B^{m-2}}.                           \label{eq:s6v12-Cayley-law}
\end{equation}
The weighted Cayley identity says that this is a probability
measure.

\begin{theorem}[Exact weighted Pr\"ufer law]
\label{thm:s6v12-Prufer}
Under \(\mu_b\), the Pr\"ufer word
\((Z_1,\ldots,Z_{m-2})\) consists of independent letters with
\begin{equation}
 {\mathbb P}(Z_r=i)={b_i\over B}.                         \label{eq:s6v12-Prufer-letters}
\end{equation}
In particular, if
\begin{equation}
 {\mathsf D}_b(T):=\sum_{i=1}^m b_i d_T(i),               \label{eq:s6v12-degree-stat}
\end{equation}
then
\begin{equation}
 {\mathsf D}_b(T)
 =
 B+\sum_{r=1}^{m-2}b_{Z_r}.                              \label{eq:s6v12-degree-sum}
\end{equation}
\end{theorem}

\begin{proof}
If \(N_i\) is the number of occurrences of \(i\) in the Pr\"ufer
word, then \(d_T(i)=1+N_i\).  Therefore
\[
 \prod_{\{i,j\}\in E(T)}b_ib_j
 =
 \prod_i b_i^{d_T(i)}
 =
 \left(\prod_i b_i\right)
 \prod_{r=1}^{m-2}b_{Z_r}.
\]
The Pr\"ufer correspondence is a bijection between labelled trees
and words in \(X^{m-2}\).  Dividing the last display by the
normalizing denominator in \eqref{eq:s6v12-Cayley-law} gives the
product law \eqref{eq:s6v12-Prufer-letters}.  Finally,
\[
 \sum_i b_i d_T(i)
 =
 \sum_i b_i+\sum_i b_iN_i,
\]
which is \eqref{eq:s6v12-degree-sum}.  When \(m=2\), the word is
empty and the same formula holds.
\end{proof}

\begin{corollary}[Exact moment generating function and a
sub-Gaussian tail]
\label{cor:s6v12-MGF}
For every real \(\theta\),
\begin{equation}
 {\mathbb E}_{\mu_b}e^{\theta{\mathsf D}_b}
 =
 e^{\theta B}
 \left({1\over B}\sum_{i=1}^m b_i e^{\theta b_i}\right)^{m-2}.
                                                               \label{eq:s6v12-MGF}
\end{equation}
Moreover,
\begin{align}
 {\mathbb E}_{\mu_b}{\mathsf D}_b
 &=
 B+(m-2){\sum_i b_i^2\over B},                            \label{eq:s6v12-mean}\\
 {\mathbb P}_{\mu_b}\!\left\{
 {\mathsf D}_b-{\mathbb E}{\mathsf D}_b\ge u\right\}
 &\le
 \exp\!\left\{-{2u^2\over(m-2)b_*^2}\right\},
 \qquad u\ge0,                                            \label{eq:s6v12-tail}
\end{align}
with the evident deterministic interpretation when \(m=2\).
\end{corollary}

\begin{proof}
Equation \eqref{eq:s6v12-MGF} follows by independence in
\Cref{thm:s6v12-Prufer}.  Differentiation at zero gives
\eqref{eq:s6v12-mean}.  Each random variable \(b_{Z_r}\) lies in
an interval of length at most \(b_*\).  Hoeffding's lemma and
Chernoff optimization give \eqref{eq:s6v12-tail}.
\end{proof}

\begin{corollary}[Domination of resolved coloured mass]
\label{cor:s6v12-X-domination}
For every owned-edge colouring of \(T\), the resolved mass of
\Cref{lem:s6v11-resolved-mass} satisfies
\begin{equation}
 X_{\rm col}:=\sum_\alpha B_{C_\alpha}
 \le {\mathsf D}_b(T)\le2(m-1).                           \label{eq:s6v12-X-domination}
\end{equation}
Consequently \(\Sigma_x\le4(m-1)^2\).
\end{corollary}

\begin{proof}
The first inequality is
\(d_{\rm col}(i)\le d_T(i)\), summed against \(b_i\), from
\Cref{lem:s6v11-resolved-mass}.  Since \(b_i<1\) and a tree has
total degree \(2(m-1)\), the second inequality follows.  Finally,
\(\Sigma_x\le X_{\rm col}^2\).
\end{proof}

\section{Why the proposed tail route cannot be uniform}
\label{sec:s6v12-tail-no-go}

\begin{theorem}[Uniform marks erase all useful outer-tree
fluctuation]
\label{thm:s6v12-uniform-no-go}
Fix \(0<\beta<1\) and set \(b_i=\beta\) for all \(i\).  Then for
every labelled tree \(T\),
\begin{equation}
 {\mathsf D}_b(T)=2\beta(m-1).                            \label{eq:s6v12-uniform-D}
\end{equation}
For the path outer tree with alternating two-occurrence ownership,
every owned component is one edge.  Thus
\[
 N=m-1,\qquad x_\alpha=2\beta,\qquad q=2,\qquad g=N-q=m-3,
\]
and
\begin{equation}
 \Sigma_x=4\beta^2N(N-1).                                \label{eq:s6v12-uniform-Sigma}
\end{equation}
For every fixed \(C_0\), this history belongs to
\(\Sigma_x>C_0g\) for all sufficiently large \(m\), with
\(\mu_b\)-probability one as far as the statistic
\({\mathsf D}_b\) is concerned.
\end{theorem}

\begin{proof}
Since \(\sum_i d_T(i)=2(m-1)\),
\eqref{eq:s6v12-uniform-D} is deterministic.  On a path, colour
successive edges alternately by two occurrences.  No two edges
of one colour are adjacent, so every monochromatic component is
one edge and has mass \(2\beta\).  Hence
\[
 X_{\rm col}=2\beta N,\qquad
 \sum_\alpha x_\alpha^2=4\beta^2N,
\]
which gives \eqref{eq:s6v12-uniform-Sigma}.  Finally,
\(\Sigma_x/g\) tends to infinity linearly in \(m\).
\end{proof}

\begin{correction}[The CCLC bottleneck lies upstream of degree
concentration]
\label{corr:s6v12-CCLC}
The V11 programme proposed combining its determinant with an
outer-Cayley degree tail.  The exact tail
\eqref{eq:s6v12-tail} is valid, but
\Cref{thm:s6v12-uniform-no-go} proves that it cannot close CCLC
uniformly.  Thus the last sentence of
\Cref{ass:s6v11-bottleneck} was too narrow: the concentrated row
must retain either the source-occurrence coefficient or a
geometric locality restriction, not merely the correlation with
the sampled tree.
\end{correction}

\begin{proof}
In the uniform-mark example the proposed dominating statistic is
constant and the concentrated event occurs for every tree.
Therefore no averaging over \(\mu_b\) can reduce that row.
\end{proof}

\begin{assessment}[What the exact Pr\"ufer computation did and
did not accomplish]
\label{ass:s6v12-Prufer-scope}
The calculation closes the probabilistic question posed in V11:
the degree law and its optimal elementary concentration mechanism
are explicit.  It also proves that this mechanism is aimed at the
wrong level of the expansion.  This is a proof-method no-go, not
an operator lower bound; cancellations in the simultaneous
analytic source coefficient have still not been taken in norm.
\end{assessment}

\section{The factorial that survives finite-family aggregation}
\label{sec:s6v12-Duhamel}

The appropriate upstream object is not one independent variable
for every spatial support.  First aggregate all supports belonging
to one analytic source family and only then expand the semigroup.
The following standard coefficient statement is recorded with its
normalization because precisely this factorial was absent from the
post-norm matching proxy.

\begin{lemma}[Multi-type Duhamel coefficient]
\label{lem:s6v12-multitype}
Let \(A\) generate a semigroup satisfying
\(\|e^{tA}\|\le e^{\omega t}\), and let
\(V_1,\ldots,V_S\) be bounded operators.  For
\({\bf n}=(n_1,\ldots,n_S)\), put \(q=\sum_s n_s\).  Then, for
\(t\ge0\),
\begin{equation}
 \left\|
 [z_1^{n_1}\cdots z_S^{n_S}]
 e^{t(A+\sum_s z_sV_s)}
 \right\|
 \le
 e^{\omega t}{t^q\over\prod_{s=1}^S n_s!}
 \prod_{s=1}^S\|V_s\|^{n_s}.                             \label{eq:s6v12-multitype}
\end{equation}
\end{lemma}

\begin{proof}
The order-\(q\) Dyson term is a sum over words
\((s_1,\ldots,s_q)\).  The coefficient of \(z^{\bf n}\) contains
exactly \(q!/\prod_s n_s!\) words.  Each ordered simplex has
volume \(t^q/q!\).  Submultiplicativity and the semigroup bound
give \(e^{\omega t}\prod_s\|V_s\|^{n_s}\) for its integrand.
Multiplication proves \eqref{eq:s6v12-multitype}.
\end{proof}

\begin{corollary}[Summing a finite type catalogue]
\label{cor:s6v12-catalogue}
If \(S\le S_0\) and \(\|V_s\|\le C_V\), then the sum of the
right side of \eqref{eq:s6v12-multitype} over all
\({\bf n}\) with \(|{\bf n}|=q\) is at most
\begin{equation}
 e^{\omega t}{(tC_VS_0)^q\over q!}.                      \label{eq:s6v12-catalogue}
\end{equation}
\end{corollary}

\begin{proof}
Use
\[
 \sum_{n_1+\cdots+n_S=q}{1\over\prod_s n_s!}
 ={S^q\over q!}
 \le {S_0^q\over q!},
\]
which is the multinomial theorem.
\end{proof}

\begin{definition}[Finite-star analytic aggregation card, FSAAC]
\label{def:s6v12-FSAAC}
FSAAC is an exact reorganization of the primitive Yang--Mills
source insertions into at most \(S_0\) analytic families
\[
 V_s=\sum_{\sigma\in{\cal A}_s}V_\sigma
\]
before any separate support norm is taken, such that:
\begin{enumerate}[label=\textup{(\roman*)},leftmargin=3.5em]
\item \(S_0\) is independent of the polymer size, cutoff, shell,
      and subdivision depth;
\item the aggregate operator bounds entering
      \Cref{lem:s6v12-multitype} have one uniform strong base;
\item coefficient extraction preserves the owned outer-tree
      partition and the grouping rank \(g=N-q\); and
\item spatial support sums are paid inside \(V_s\), without
      recreating the separated-support factorial of
      \Cref{thm:s6v8-grouping-no-go}.
\end{enumerate}
\end{definition}

\begin{warning}[Finite shapes are not yet aggregate operator
bounds]
\label{warn:s6v12-shapes}
A finite list of local combinatorial shapes does not by itself
prove FSAAC.  Summing all translates of one shape may have a norm
proportional to volume, and taking their norms separately returns
the V8 obstruction.  FSAAC requires the actual simultaneous
operator estimate before the spatial sum is split.
\end{warning}

\section{A conditional closure of the rank-balanced row}
\label{sec:s6v12-balanced}

\begin{theorem}[FSAAC closes \(g\le q\)]
\label{thm:s6v12-balanced}
Assume FSAAC.  Then the complete grouping-forest marginal obeys a
\(K^m\) bound on every row
\begin{equation}
 0\le g\le q.                                             \label{eq:s6v12-balanced}
\end{equation}
This conclusion is uniform in the marks \(0<b_i<1\).
\end{theorem}

\begin{proof}
The rank-\(g\) forest sum is at most
\(\Sigma_x^g/g!\) by
\Cref{cor:s6v11-graded}.  By
\Cref{cor:s6v12-X-domination}, \(\Sigma_x\le4m^2\).
After the finite source families are summed, FSAAC and
\Cref{cor:s6v12-catalogue} supply, up to a numerical strong base,
the factor \(S_0^q/q!\).  Since \(g\le q\le m\),
\[
 { (4m^2)^g\over g!}{S_0^q\over q!}
 \le
 S_0^m{(4m^2)^g\over(g!)^2}
 \le
 S_0^m\left({2em\over g}\right)^{2g}.
\]
For \(y=g/m\in(0,1]\),
\[
 {1\over m}\log\left({2em\over g}\right)^{2g}
 =
 2y\log{2e\over y}.
\]
The right side is bounded on \(0<y\le1\) and tends to zero as
\(y\downarrow0\).  Thus the last display is at most \(K^m\).
The case \(g=0\) is immediate.
\end{proof}

\begin{corollary}[Location of the genuinely hard rank sector]
\label{cor:s6v12-defect-dominant}
Conditional on FSAAC, the only CCLC histories not covered by
\Cref{thm:s6v12-balanced} satisfy
\begin{equation}
 g>q.                                                     \label{eq:s6v12-defect-dominant}
\end{equation}
Equivalently, the number of disconnectedness defects exceeds the
number of primitive source occurrences.
\end{corollary}

\begin{proof}
The complementary ranks have \(g\le q\).  The diffuse row is
already unconditional by \Cref{thm:s6v11-diffuse}.
\end{proof}

\begin{assessment}[Exact live bottleneck after V12]
\label{ass:s6v12-bottleneck}
There are now two upstream obligations.  First, the proposed
finite family must be an aggregate operator estimate, FSAAC, not
merely a finite shape list.  Second, even its \(1/q!\) gain does
not by itself pay the defect-dominant sector \(g>q\).  The
alternating two-colour path in
\Cref{thm:s6v12-uniform-no-go} shows why: one occurrence may meet
the outer tree in linearly many separated components.  A
successful continuation must prove that physical source geometry
makes the admissible component-grouping graph sparse, or must
retain a stronger noncommutative cancellation.
\end{assessment}

\begin{programme}[V13: sparse physical grouping graphs]
\label{prog:s6v12-V13}
\begin{enumerate}[leftmargin=3em]
\item Replace the complete component graph by the graph of pairs
      which can occur together in one primitive physical source.
\item Prove an all-rank forest bound from a uniform weighted
      row-sum or bounded-overlap estimate on that graph.
\item Derive such a row-sum bound from a global subdivision with
      a finite star catalogue.
\item Separate local plaquette/vertex insertions from extended
      Wilson-boundary insertions.
\item If the extended row is not sparse, retain its cyclic order
      and formulate the precise boundary analogue.
\end{enumerate}
\end{programme}

\begin{firewall}[Claim boundary after sixth-series V12]
\label{fire:s6v12-boundary}
V12 proves the exact weighted Cayley--Pr\"ufer law, the moment
generating function and tail of the coloured-degree dominator,
and the deterministic uniform-mark obstruction to a universal
degree-tail closure.  It corrects the CCLC route accordingly.  It
also proves the exact multi-type Duhamel factorial, formulates
FSAAC with the necessary operator-level qualification, and proves
that FSAAC closes every rank-balanced grouping row \(g\le q\).

V12 does not prove FSAAC, the defect-dominant CCLC row, RFAC in
full, high-excess TFLAC, PSCSC, SOCC, SISC, the simple
high-excess row, arbitrary repeated transversal supports, full
TOCC, full TCNC, cyclic TIC, HLTC, full MCWC, OCRC, SCNGC, WPSC,
GRLTC, LZTC, PLRC, WSDC, BRSC, HWSFC, UGCGC, thermal connected
WPRAC, all-arity RGSC, the raw Wilson aggregate strong card,
all-order strong MTA, WUFC, CPM, cutoff removal, reflection
positivity, Osterwalder--Schrader reconstruction, a continuum
Yang--Mills measure, or a positive Hamiltonian gap.  The full
Yang--Mills mass gap has not been proved.
\end{firewall}

\begin{theorem}[Sixth-series V12 result ledger]
\label{thm:s6v12-results}
V12:
\begin{enumerate}[label=\textup{(V12.\arabic*)},leftmargin=6em]
\item proves the exact weighted Pr\"ufer product law;
\item computes the coloured-degree moment generating function;
\item proves its sub-Gaussian upper tail;
\item bounds resolved component mass by the degree statistic;
\item proves the uniform-mark no-go for the tail route;
\item proves the multi-type Duhamel coefficient bound;
\item identifies the finite-catalogue \(1/q!\) gain;
\item formulates FSAAC without conflating shapes and operator
      bounds; and
\item conditionally closes every grouping row \(g\le q\).
\end{enumerate}
\end{theorem}

\begin{proof}
Use
\Cref{thm:s6v12-Prufer,
cor:s6v12-MGF,
cor:s6v12-X-domination,
thm:s6v12-uniform-no-go,
corr:s6v12-CCLC,
lem:s6v12-multitype,
cor:s6v12-catalogue,
thm:s6v12-balanced}.
\end{proof}

\paragraph{Parent-version record.}
V12 was formed from
\texttt{Renewal\_Yang\_Mills\_Mass\_Gap\_Research\_Notes\_V11.tex}.
The V11 parent had \(5\,027\) logical source lines, \(176\,232\)
bytes, and SHA--256
\[
 \texttt{223bed98d69854fa8ae2aa8351edee32cb7b6cca14ac8d51a82c8ef9e8f52fcb}.
\]
The next compulsory companion audit is V15.

\clearpage
\chapter{Sixth-series V13: sparse physical grouping graphs
and closure of the strictly local source row}
\label{ch:s6v13-sparse}

\section{The physical compatibility graph}
\label{sec:s6v13-compatibility}

No companion audit is due at V13.

Fix a resolved owned-component family
\({\cal C}=\{C_1,\ldots,C_N\}\) and write
\[
 x_\alpha:=B_{C_\alpha}.
\]
The complete component graph used in V11 permits every pair
\(\{C_\alpha,C_\beta\}\).  The exact source expansion permits such
an edge only when the two components can be joined inside one
primitive physical source occurrence.  Let
\({\cal H}_{\rm phys}\) be the graph of these admissible pairs and
give its edge \(\{\alpha,\beta\}\) weight
\begin{equation}
 w_{\alpha\beta}:=x_\alpha x_\beta.                       \label{eq:s6v13-edge-weight}
\end{equation}
Every exact grouping forest is a subgraph of
\({\cal H}_{\rm phys}\), although the converse need not hold.

\begin{definition}[Weighted compatibility load]
\label{def:s6v13-load}
For a graph \({\cal H}\) on the resolved components, define
\begin{equation}
 \Omega({\cal H},x)
 :=
 \sum_{\{\alpha,\beta\}\in E({\cal H})}x_\alpha x_\beta.
                                                               \label{eq:s6v13-load}
\end{equation}
\end{definition}

\begin{lemma}[All-rank sparse-forest bound]
\label{lem:s6v13-all-rank}
For every finite graph \({\cal H}\) and nonnegative component
weights,
\begin{align}
 \sum_{\substack{L\subset{\cal H}\\L\ {\rm forest}}}
 \prod_{e\in E(L)}w_e
 &\le
 \prod_{e\in E({\cal H})}(1+w_e)                           \label{eq:s6v13-product}\\
 &\le
 \exp\{\Omega({\cal H},x)\}.                              \label{eq:s6v13-exponential}
\end{align}
In particular, \(\Omega({\cal H}_{\rm phys},x)\le C_0m\)
implies the complete grouping-forest bound \(K^m\), simultaneously
at every rank.
\end{lemma}

\begin{proof}
Dropping acyclicity enlarges the first sum to the sum over all
edge subsets, which factorizes as the product in
\eqref{eq:s6v13-product}.  The inequality \(1+u\le e^u\) gives
\eqref{eq:s6v13-exponential}.  Set \(K=e^{C_0}\).
\end{proof}

\begin{lemma}[Restricted matrix--forest refinement]
\label{lem:s6v13-restricted-determinant}
Let \({\cal L}_{\cal H}\) be the weighted Laplacian of
\({\cal H}\), with the edge weights
\eqref{eq:s6v13-edge-weight}.  Then
\begin{equation}
 \det(tI+{\cal L}_{\cal H})
 =
 \sum_{\substack{L\subset{\cal H}\\L\ {\rm forest}}}
 t^{|{\cal K}(L)|}
 \left(\prod_{e\in E(L)}w_e\right)
 \prod_{D\in{\cal K}(L)}|D|.                             \label{eq:s6v13-restricted-determinant}
\end{equation}
Moreover,
\[
 \operatorname{Tr}{\cal L}_{\cal H}=2\Omega({\cal H},x).
\]
\end{lemma}

\begin{proof}
The ghost-vertex proof of
\Cref{thm:s6v11-matrix-forest} uses only the available
non-ghost edges, so it applies verbatim to \({\cal H}\).
Every edge weight appears in two diagonal entries of its
Laplacian, giving the trace identity.
\end{proof}

\begin{correction}[The complete Laplacian is a deliberate
relaxation]
\label{corr:s6v13-complete}
The complete Laplacian \({\cal L}_x\) of V11 is an upper-bound
device, not the exact physical compatibility operator.  CCLC is
unnecessarily strong whenever
\({\cal H}_{\rm phys}\) is sparse.  The exact replacement is the
restricted determinant
\eqref{eq:s6v13-restricted-determinant}, or more directly the
all-rank product bound \eqref{eq:s6v13-exponential}.
\end{correction}

\begin{proof}
Every physical grouping forest uses only admissible component
pairs.  Positivity permits enlarging its sum to all forest
subgraphs of \({\cal H}_{\rm phys}\), but does not require adding
the missing complete-graph edges.
\end{proof}

\section{A star-cover criterion}
\label{sec:s6v13-star-cover}

The number of translated stars may grow with the cutoff volume.
What must remain finite is their local load and overlap, not their
global count.

\begin{definition}[Admissible weighted star cover]
\label{def:s6v13-star-cover}
An admissible weighted star cover of
\({\cal H}_{\rm phys}\) is a family
\(\{{\cal A}_s:s\in{\cal S}\}\) of subsets of
\(\{1,\ldots,N\}\) such that:
\begin{enumerate}[label=\textup{(\roman*)},leftmargin=3.5em]
\item every edge of \({\cal H}_{\rm phys}\) has both endpoints in
      at least one \({\cal A}_s\);
\item the star load obeys
      \(\sum_{\alpha\in{\cal A}_s}x_\alpha\le M_0\);
\item every component index belongs to at most \(R_0\) sets
      \({\cal A}_s\).
\end{enumerate}
The constants \(M_0,R_0\) must be independent of \(m\), the
cutoff, shell, and subdivision depth.
\end{definition}

\begin{theorem}[Weighted star-cover theorem]
\label{thm:s6v13-star-cover}
If \({\cal H}_{\rm phys}\) has an admissible weighted star cover,
then
\begin{equation}
 \Omega({\cal H}_{\rm phys},x)
 \le M_0R_0m,                                             \label{eq:s6v13-star-load}
\end{equation}
and its complete grouping-forest sum is at most
\(\exp(M_0R_0m)\).
\end{theorem}

\begin{proof}
Cover every physical edge by the clique on at least one
\({\cal A}_s\).  Overcounting edges is harmless, and
\begin{align*}
 \Omega({\cal H}_{\rm phys},x)
 &\le
 \sum_{s\in{\cal S}}
 \sum_{\substack{\alpha<\beta\\
                  \alpha,\beta\in{\cal A}_s}}
 x_\alpha x_\beta\\
 &\le
 {1\over2}\sum_s
 \left(\sum_{\alpha\in{\cal A}_s}x_\alpha\right)^2\\
 &\le
 {M_0\over2}\sum_s\sum_{\alpha\in{\cal A}_s}x_\alpha\\
 &\le
 {M_0R_0\over2}\sum_{\alpha=1}^Nx_\alpha.
\end{align*}
By \Cref{cor:s6v12-X-domination},
\(\sum_\alpha x_\alpha\le2(m-1)\).  This proves
\eqref{eq:s6v13-star-load}; now apply
\Cref{lem:s6v13-all-rank}.
\end{proof}

\begin{assessment}[Why this criterion hits the V8 obstruction]
\label{ass:s6v13-V8}
In the separated-matching construction of
\Cref{thm:s6v8-grouping-no-go}, arbitrarily distant path edges
may be paired.  Its compatibility graph is therefore dense and
has load of order \(m^2\) at equal marks.  A bounded-load
physical star cover forbids precisely that nonlocal positive
fibre.  It does not seek a second factorial for the same dense
graph.
\end{assessment}

\section{The strictly local primitive row}
\label{sec:s6v13-local}

\begin{definition}[Uniformly local primitive catalogue]
\label{def:s6v13-local-catalogue}
A primitive source row has a uniformly local catalogue if there
is a family of label stars \(S_s\subset X\) such that:
\begin{enumerate}[label=\textup{(\roman*)},leftmargin=3.5em]
\item \(|S_s|\le a_0\);
\item every primitive occurrence in the row has its entire forced
      support in one \(S_s\);
\item any resolved component is contained in at most \(r_0\)
      catalogue stars capable of supporting its occurrence; and
\item two components are physically compatible only if they are
      both contained in one such star.
\end{enumerate}
Here \(a_0,r_0\) are uniform cutoff-independent constants.
\end{definition}

\begin{lemma}[A local star has bounded resolved mass load]
\label{lem:s6v13-local-load}
Under \Cref{def:s6v13-local-catalogue}, for a fixed owned-edge
history the sum of \(x_\alpha\) over resolved components contained
in one star is at most
\begin{equation}
 M_0:=a_0(a_0-1).                                        \label{eq:s6v13-local-M}
\end{equation}
\end{lemma}

\begin{proof}
Every nontrivial resolved component contains at least one
outer-tree edge whose endpoints lie in the star.  Distinct
resolved components have disjoint owned edge sets.  The
restriction of an outer tree to \(a_0\) vertices is a forest and
has at most \(a_0-1\) edges.  Hence at most \(a_0-1\) resolved
components can be contained in the star.  Each such component has
at most \(a_0\) vertices and \(b_i<1\), so its mass is at most
\(a_0\).  Multiplication gives \eqref{eq:s6v13-local-M}.
\end{proof}

\begin{theorem}[Strictly local grouping closure]
\label{thm:s6v13-local-closure}
Every primitive source row with a uniformly local catalogue has
a complete all-rank grouping-forest bound
\begin{equation}
 \sum_{L\ {\rm physically\ admissible}}w_x(L)
 \le
 \exp\{a_0(a_0-1)r_0m\}.                                \label{eq:s6v13-local-closure}
\end{equation}
Thus RFAC, TFLAC, PSCSC, SOCC, and the remaining SISC estimate
hold on the strictly local primitive row.
\end{theorem}

\begin{proof}
For every label star \(S_s\), let \({\cal A}_s\) be the set of
resolved components contained in it and capable of belonging to
an occurrence supported there.  Compatibility is covered by
these sets.  Their load is bounded by
\Cref{lem:s6v13-local-load}, and their overlap is at most \(r_0\).
Apply \Cref{thm:s6v13-star-cover} and then
\Cref{prop:s6v10-RFAC,prop:s6v9-TFLAC,
prop:s6v8-PSCSC,prop:s6v7-SOCC}.
\end{proof}

\begin{corollary}[The local row does not require FSAAC]
\label{cor:s6v13-no-FSAAC}
The conclusion of \Cref{thm:s6v13-local-closure} is obtained
directly from geometric compatibility.  It does not require the
conditional aggregate-operator card FSAAC or the restriction
\(g\le q\).
\end{corollary}

\begin{proof}
The all-rank product in \Cref{lem:s6v13-all-rank} already sums
every \(g\) and every compatible grouping forest.
\end{proof}

\section{What remains outside the local catalogue}
\label{sec:s6v13-boundary}

\begin{definition}[Yang--Mills global finite-star grouping card,
YM--FSGC]
\label{def:s6v13-YMFSGC}
YM--FSGC is the statement that the actual cutoff Yang--Mills
primitive source expansion admits the star cover of
\Cref{def:s6v13-star-cover} after the global subdivision is fixed,
with uniform load and overlap constants, except for explicitly
identified extended Wilson-boundary sectors which retain their
boundary order.
\end{definition}

\begin{proposition}[YM--FSGC reduces the open SISC row to the
extended boundary]
\label{prop:s6v13-YMFSGC}
If YM--FSGC holds, then SISC is proved for every strictly local
primitive source sector.  Its only remaining SOCC contribution is
the declared extended Wilson-boundary sector.
\end{proposition}

\begin{proof}
Apply \Cref{thm:s6v13-star-cover} to each covered sector and use
\Cref{thm:s6v13-local-closure}.  By definition, the only
uncovered primitive sources are the extended boundary sources.
\end{proof}

\begin{warning}[Extended Wilson data cannot be called local]
\label{warn:s6v13-Wilson}
A Wilson loop or shell may contain a number of cutoff labels
growing with its length.  Its support therefore fails
\(|S_s|\le a_0\) unless one proves an exact subdivision into local
analytic primitives.  Merely cutting the support after taking
norms would reproduce the gapped-stencil problem seen in the SM
companion audit.
\end{warning}

\begin{definition}[Boundary-ordered grouping card, BOGC]
\label{def:s6v13-BOGC}
BOGC is the remaining assertion that the extended Wilson-boundary
source coefficient, with its cyclic order and reduced-fibre
structure retained before norms, has an admissible grouping graph
of weighted load \(O(m)\), or an equivalent \(K^m\) all-rank
forest bound.
\end{definition}

\begin{assessment}[Exact live bottleneck after V13]
\label{ass:s6v13-bottleneck}
The abstract sparse-graph estimate is complete, and the strictly
local primitive row is closed under an explicit bounded-catalogue
hypothesis.  The proof no longer needs an artificial complete
component graph.  What remains is model identification: prove
YM--FSGC for the actual local cutoff vertices and prove BOGC for
the extended Wilson sector.  This is upstream of CCLC and directly
tests the physical source geometry.
\end{assessment}

\begin{programme}[V14: use cyclic Wilson order]
\label{prog:s6v13-V14}
\begin{enumerate}[leftmargin=3em]
\item Read the reduced-fibre Wilson-shell definitions in the
      attached Yang--Mills programme and identify the exact
      primitive boundary supports.
\item Determine whether admissible component joins are interval,
      noncrossing, bounded-range, or unrestricted in boundary
      cyclic order.
\item Prove an exponential forest enumeration for the strongest
      property actually present.
\item Treat local cutoff vertices by
      \Cref{thm:s6v13-local-closure}.
\item Formulate a narrower obstruction if the boundary support
      still permits the V8 separated matching family.
\end{enumerate}
\end{programme}

\begin{firewall}[Claim boundary after sixth-series V13]
\label{fire:s6v13-boundary}
V13 defines the exact physical compatibility graph, proves its
all-rank weighted forest bound, extends the matrix--forest identity
to that graph, and corrects the complete-Laplacian relaxation.  It
proves a weighted star-cover theorem and, from it, closes RFAC
through SISC on every uniformly local primitive catalogue row.  It
formulates YM--FSGC and isolates the extended Wilson-boundary card
BOGC.

V13 does not prove YM--FSGC for all actual Yang--Mills primitives,
BOGC, FSAAC, the extended defect-dominant CCLC row, RFAC in full,
SISC in full, the simple high-excess row, arbitrary repeated
transversal supports, full TOCC, full TCNC, cyclic TIC, HLTC, full
MCWC, OCRC, SCNGC, WPSC, GRLTC, LZTC, PLRC, WSDC, BRSC, HWSFC,
UGCGC, thermal connected WPRAC, all-arity RGSC, the raw Wilson
aggregate strong card, all-order strong MTA, WUFC, CPM, cutoff
removal, reflection positivity, Osterwalder--Schrader
reconstruction, a continuum Yang--Mills measure, or a positive
Hamiltonian gap.  The full Yang--Mills mass gap has not been
proved.
\end{firewall}

\begin{theorem}[Sixth-series V13 result ledger]
\label{thm:s6v13-results}
V13:
\begin{enumerate}[label=\textup{(V13.\arabic*)},leftmargin=6em]
\item replaces the complete component graph by physical
      compatibility;
\item proves the all-rank sparse-forest product bound;
\item proves the restricted matrix--forest identity;
\item proves the weighted star-cover theorem;
\item explains exactly how it excludes the V8 obstruction;
\item derives uniform star load from bounded local supports;
\item closes the strictly local source row through SISC;
\item formulates YM--FSGC; and
\item isolates the boundary-ordered grouping card.
\end{enumerate}
\end{theorem}

\begin{proof}
Use
\Cref{lem:s6v13-all-rank,
lem:s6v13-restricted-determinant,
corr:s6v13-complete,
thm:s6v13-star-cover,
lem:s6v13-local-load,
thm:s6v13-local-closure,
prop:s6v13-YMFSGC}.
\end{proof}

\paragraph{Parent-version record.}
V13 was formed from
\texttt{Renewal\_Yang\_Mills\_Mass\_Gap\_Research\_Notes\_V12.tex}.
The V12 parent had \(5\,452\) logical source lines, \(191\,198\)
bytes, and SHA--256
\[
 \texttt{d5d429e3e963e6ea92321b9fb53e5ad339247ca04091292d8ac02350c5e07d06}.
\]
The next compulsory companion audit is V15.

\clearpage
\chapter{Sixth-series V14: the attached-programme incidence
audit and a two-anchor connected-polymer bound}
\label{ch:s6v14-polymer}

\section{Audit of the two attached Yang--Mills programmes}
\label{sec:s6v14-audit}

No SM/NS companion audit is due at V14.  This is instead the
source audit required by the V13 programme.

\begin{audit}[Reduced-fibre Wilson-shell source]
\label{audit:s6v14-YMA}
The file
\[
\texttt{Renewal\_Yang\_Mills\_Reduced\_Fibre\_Wilson\_Shell\_Symbol\_and\_Local\_Vacuum\_Programme\_V35.tex}
\]
has \(1\,560\,166\) bytes, \(37\,288\) logical source lines, and
SHA--256
\[
\texttt{7c6056d8149d3015d469b57474a75b23c88a6fe9f64fc81c7dca85c361c647a3}.
\]
Its inherited V25 theorem uses the fine-plaquette incidence graph
\(\mathscr G_{\rm plaq}\), in which two plaquettes are adjacent
when they share an integrated fine link.  In four dimensions it
proves
\begin{equation}
 \Delta_{\rm plaq}\le4(6-1)=20.                           \label{eq:s6v14-degree20}
\end{equation}
It also proves a volume-uniform connected plaquette-polymer
expansion when
\begin{equation}
 u(\beta):=e^{|\beta|}-1
 <
 {1\over2e^3\Delta_{\rm plaq}}.                          \label{eq:s6v14-KP-input}
\end{equation}
The current V35 continuation concerns a trace-shell head--tail
moment card.  It does not identify SOCC resolved components with
plaquette-polymer anchors.
\end{audit}

\begin{audit}[Boundary-sector source]
\label{audit:s6v14-YMB}
The file
\[
\texttt{Renewal\_Yang\_Mills\_Boundary\_Sector\_Cofinal\_Contraction\_Programme\_V33.tex}
\]
has \(1\,277\,180\) bytes, \(31\,148\) logical source lines, and
SHA--256
\[
\texttt{b762244c0ea25c54072c3ce3a29cf3386230f1507205c20a89ff833ea195928d}.
\]
Its current boundary object is a radial information-lifetime law,
its possible atom at infinite radius, and a diagonal remote-shell
subtraction.  It proves no cyclic order, interval order, or
noncrossing restriction on the component groupings of SOCC.
\end{audit}

\begin{correction}[BOGC was based on a semantic
misidentification]
\label{corr:s6v14-BOGC}
The boundary-ordered grouping card BOGC of V13 is retired as an
inferred property of the attached manuscripts.  In those sources,
\emph{Wilson shell} denotes the shell symbol derived from the
local Wilson plaquette action, while \emph{boundary sector}
denotes the cofinal radial-information problem.  Neither phrase
supplies an ordered extended Wilson-loop support.
\end{correction}

\begin{proof}
This is the content of
\Cref{audit:s6v14-YMA,audit:s6v14-YMB}.  Importing a noncrossing
boundary enumeration would add a hypothesis absent from both
attached programmes.
\end{proof}

\begin{assessment}[The usable imported structure]
\label{ass:s6v14-usable}
The genuinely relevant inherited facts are the bounded plaquette
degree \eqref{eq:s6v14-degree20}, the connected-support
factorization, and the small-activity polymer expansion.  These
facts can produce a summable two-anchor compatibility kernel.  A
separate incidence bridge is still required before that kernel
can pay the renewal grouping forest.
\end{assessment}

\section{A graph-animal two-anchor kernel}
\label{sec:s6v14-kernel}

Let \(\mathscr G=(P,E)\) be a locally finite graph of maximum
degree \(\Delta\ge2\).  For \(p,q\in P\) and \(a\ge0\), put
\begin{equation}
 {\cal K}_a(p,q)
 :=
 \sum_{\substack{\gamma\subset P\ {\rm finite,\ connected}\\
                  p,q\in\gamma}}
 a^{|\gamma|}.                                           \label{eq:s6v14-kernel}
\end{equation}

\begin{lemma}[Elementary rooted-animal count]
\label{lem:s6v14-animal}
For every \(p\in P\), the number of connected vertex sets
\(\gamma\) of cardinality \(n\) containing \(p\) is at most
\begin{equation}
 A_\Delta^{\,n-1},
 \qquad A_\Delta:=\Delta^2.                              \label{eq:s6v14-animal}
\end{equation}
\end{lemma}

\begin{proof}
Fix an order on the neighbors of every vertex.  To each connected
set \(\gamma\ni p\), assign the lexicographically first rooted
spanning tree of its induced graph and perform the deterministic
depth-first contour traversal of that tree.  The traversal starts
at \(p\), has \(2(n-1)\) steps, and each step chooses one of at
most \(\Delta\) neighbors.  Its visited vertex set recovers
\(\gamma\), so different sets have different traversal words.
There are at most
\(\Delta^{2(n-1)}=A_\Delta^{n-1}\) such words.
\end{proof}

\begin{theorem}[Summable two-anchor connected kernel]
\label{thm:s6v14-kernel-sum}
If \(A_\Delta a<1\), then
\begin{equation}
 \sup_{p\in P}\sum_{q\in P}{\cal K}_a(p,q)
 \le
 {a\over(1-A_\Delta a)^2}.                               \label{eq:s6v14-kernel-sum}
\end{equation}
\end{theorem}

\begin{proof}
All terms are nonnegative, so Tonelli's theorem and
\Cref{lem:s6v14-animal} give
\begin{align*}
 \sum_q{\cal K}_a(p,q)
 &=
 \sum_{\substack{\gamma\ {\rm finite,\ connected}\\p\in\gamma}}
 |\gamma|a^{|\gamma|}\\
 &\le
 \sum_{n\ge1}nA_\Delta^{n-1}a^n
 =
 {a\over(1-A_\Delta a)^2}.
\end{align*}
\end{proof}

\begin{remark}[Relation to the sharper attached count]
\label{rem:s6v14-sharper}
The attached V25 source proves the sharper rooted-animal bound
\((e\Delta_{\rm plaq})^{n-1}\).  The elementary
\(\Delta^{2(n-1)}\) bound is retained here because its proof is
self-contained.  For \(\Delta_{\rm plaq}=20\), the attached
small-activity condition
\eqref{eq:s6v14-KP-input} implies
\[
 A_{\Delta_{\rm plaq}}u(\beta)
 <
 {400\over40e^3}
 ={10\over e^3}<1,
\]
so \Cref{thm:s6v14-kernel-sum} applies throughout that disk.
\end{remark}

\section{From an anchor kernel to a grouping-forest bound}
\label{sec:s6v14-anchor}

Assign every resolved component \(C_\alpha\) an anchor
\(p(\alpha)\in P\).  Define the mass load at an anchor by
\begin{equation}
 \ell_p:=\sum_{\alpha:p(\alpha)=p}x_\alpha.               \label{eq:s6v14-anchor-load}
\end{equation}

\begin{theorem}[Kernel-load forest theorem]
\label{thm:s6v14-kernel-load}
Suppose
\begin{equation}
 \sup_p\ell_p\le L_0,\qquad
 \sup_p\sum_q{\cal K}(p,q)\le R_0,                       \label{eq:s6v14-kernel-hyp}
\end{equation}
for a nonnegative symmetric kernel \({\cal K}\).  Suppose also
that every admissible grouping edge
\(\{\alpha,\beta\}\) has coefficient bounded by
\begin{equation}
 x_\alpha x_\beta
 {\cal K}(p(\alpha),p(\beta)).                            \label{eq:s6v14-edge-domination}
\end{equation}
Then the complete positive sum over all such grouping forests is
at most
\begin{equation}
 \exp\{L_0R_0m\}.                                        \label{eq:s6v14-kernel-forest}
\end{equation}
\end{theorem}

\begin{proof}
The sum of all admissible edge weights is bounded by
\begin{align*}
 \Omega_{\cal K}
 &\le
 {1\over2}\sum_{p,q}\ell_p\ell_q{\cal K}(p,q)\\
 &\le
 {L_0R_0\over2}\sum_p\ell_p
 =
 {L_0R_0\over2}\sum_\alpha x_\alpha\\
 &\le L_0R_0m,
\end{align*}
where the final inequality is
\Cref{cor:s6v12-X-domination}.  Dropping acyclicity and summing
edge subsets as in \Cref{lem:s6v13-all-rank} gives
\(\exp\{\Omega_{\cal K}\}\), proving
\eqref{eq:s6v14-kernel-forest}.
\end{proof}

\begin{corollary}[Connected-polymer specialization]
\label{cor:s6v14-polymer-specialization}
Under the hypotheses of
\Cref{thm:s6v14-kernel-load}, take
\({\cal K}={\cal K}_a\).  If \(A_\Delta a<1\), the grouping sum
is at most
\begin{equation}
 \exp\left\{
 {L_0a\over(1-A_\Delta a)^2}\,m
 \right\}.                                               \label{eq:s6v14-polymer-specialization}
\end{equation}
\end{corollary}

\begin{proof}
Combine
\Cref{thm:s6v14-kernel-sum,thm:s6v14-kernel-load}.
\end{proof}

\begin{corollary}[Numerical bare strong-coupling kernel]
\label{cor:s6v14-strong-coupling}
For the four-dimensional plaquette graph of
\Cref{audit:s6v14-YMA}, set
\[
 a=u(\beta),\qquad A_{\Delta_{\rm plaq}}=400.
\]
On the attached V25 disk
\eqref{eq:s6v14-KP-input},
\begin{equation}
 \sup_p\sum_q{\cal K}_{u(\beta)}(p,q)
 \le
 {u(\beta)\over(1-400u(\beta))^2}<\infty.                \label{eq:s6v14-strong-kernel}
\end{equation}
The bound is uniform in the periodic volume and in the retained
coarse boundary links.
\end{corollary}

\begin{proof}
The numerical implication is
\Cref{rem:s6v14-sharper}; apply
\Cref{thm:s6v14-kernel-sum}.  Uniformity follows because the
maximum-degree and activity bounds are volume independent.
\end{proof}

\section{The missing incidence bridge}
\label{sec:s6v14-bridge}

\begin{definition}[Plaquette-incidence bridge card, PIBC]
\label{def:s6v14-PIBC}
PIBC is the assertion that, before separate source norms:
\begin{enumerate}[label=\textup{(\roman*)},leftmargin=3.5em]
\item every SOCC resolved component has a canonical plaquette or
      coarse-cell anchor;
\item the total mark mass at each anchor is bounded by a uniform
      \(L_0\);
\item every simultaneous grouping-edge coefficient is dominated
      by the two-anchor connected kernel of the same physical
      action; and
\item the domination is preserved by the IGMC endpoint lift,
      response matching, and shell transport.
\end{enumerate}
\end{definition}

\begin{proposition}[PIBC closes the polymer-controlled SOCC row]
\label{prop:s6v14-PIBC}
If PIBC holds at a shell with a connected activity \(a\) satisfying
\(A_\Delta a<1\), then RFAC, TFLAC, PSCSC, SOCC, and SISC hold at
that shell with a uniform strong base.
\end{proposition}

\begin{proof}
PIBC supplies the hypotheses of
\Cref{cor:s6v14-polymer-specialization}.  That corollary gives
the complete all-rank grouping-forest bound.  Apply the implication
chain
\Cref{prop:s6v10-RFAC,prop:s6v9-TFLAC,
prop:s6v8-PSCSC,prop:s6v7-SOCC}.
\end{proof}

\begin{warning}[The attached convergence theorem is not PIBC]
\label{warn:s6v14-not-PIBC}
The inherited V25 expansion proves that differentiated Wilson
action kernels are sums of connected plaquette clusters.  It does
not prove that the resolved components created by the renewal
outer-tree/source decomposition are the two marked ends of that
same kernel.  Treating those two incidence structures as identical
would be a new unsupported premise.
\end{warning}

\begin{assessment}[Strong coupling versus the continuum route]
\label{ass:s6v14-coupling}
The explicit disk \eqref{eq:s6v14-KP-input} is a bare
strong-coupling expansion.  The continuum Yang--Mills limit is
weakly coupled at the ultraviolet cutoff.  Therefore the numerical
kernel \eqref{eq:s6v14-strong-kernel} is a rigorous model row and
a calibration of the required mechanism, not a continuum closure.
For the continuum programme one needs a renormalized shell
activity, after response matching, which is summable in the same
two-anchor sense.
\end{assessment}

\begin{definition}[Renormalized plaquette-incidence bridge card,
RPIBC]
\label{def:s6v14-RPIBC}
RPIBC is PIBC with the bare activity \(u(\beta)\) replaced by the
actual response-matched connected shell kernel
\({\cal K}_r^{\rm conn}\), required to satisfy
\begin{equation}
 \sup_r\sup_p\sum_q{\cal K}_r^{\rm conn}(p,q)<\infty
                                                               \label{eq:s6v14-RPIBC}
\end{equation}
on the selected ultraviolet-to-infrared cutoff route.
\end{definition}

\begin{corollary}[RPIBC is the exact replacement for BOGC]
\label{cor:s6v14-replacement}
RPIBC, rather than BOGC, is the source-geometric card indicated by
the attached programmes.  If it holds, the open SISC grouping row
closes.
\end{corollary}

\begin{proof}
\Cref{corr:s6v14-BOGC} removes the unsupported cyclic premise.
The proof of \Cref{thm:s6v14-kernel-load} only needs the uniform
row sum in \eqref{eq:s6v14-RPIBC}; it does not require a bare
small-activity representation.
\end{proof}

\begin{assessment}[Exact live bottleneck after V14]
\label{ass:s6v14-bottleneck}
The finite-degree plaquette geometry and the two-anchor
connected-polymer summation are now rigorous.  The active
bottleneck is the incidence equality between that analytic kernel
and the renewal source-owned components.  At bare strong coupling
the attached notes provide the kernel side; along the continuum
route even the required uniform response-matched row sum remains
to be proved.  RPIBC states both requirements without importing a
nonexistent boundary order.
\end{assessment}

\begin{programme}[V15: compulsory companion audit and incidence
comparison]
\label{prog:s6v14-V15}
\begin{enumerate}[leftmargin=3em]
\item Audit the current SM and NS active notes, recording their
      exact versions and hashes.
\item Compare the SM global-subdivision/finite-star card with
      RPIBC's anchor and overlap requirements.
\item Compare the NS retained nonlocal pressure object with the
      need to keep the connected shell kernel before local norms.
\item Inspect the renewal/Grand Tensor source ancestry for the
      exact map from primitive source occurrences to action
      clusters.
\item Prove the map on the largest sector justified by the
      manuscripts and isolate the first failing incidence.
\end{enumerate}
\end{programme}

\begin{firewall}[Claim boundary after sixth-series V14]
\label{fire:s6v14-boundary}
V14 audits the two attached Yang--Mills programmes and corrects
the unsupported cyclic-boundary interpretation.  It proves an
elementary rooted-animal bound, a summable two-anchor connected
kernel, and a kernel-load theorem which closes every forest rank.
It specializes the result to the degree-twenty plaquette graph on
the attached strong-coupling disk and formulates PIBC and its
renormalized form RPIBC.

V14 does not prove PIBC or RPIBC, the continuum response-matched
kernel row sum, RFAC in full, SISC in full, the simple high-excess
row, arbitrary repeated transversal supports, full TOCC, full
TCNC, cyclic TIC, HLTC, full MCWC, OCRC, SCNGC, WPSC, GRLTC,
LZTC, PLRC, WSDC, BRSC, HWSFC, UGCGC, thermal connected WPRAC,
all-arity RGSC, the raw Wilson aggregate strong card, all-order
strong MTA, WUFC, CPM, cutoff removal, reflection positivity,
Osterwalder--Schrader reconstruction, a continuum Yang--Mills
measure, or a positive Hamiltonian gap.  The full Yang--Mills
mass gap has not been proved.
\end{firewall}

\begin{theorem}[Sixth-series V14 result ledger]
\label{thm:s6v14-results}
V14:
\begin{enumerate}[label=\textup{(V14.\arabic*)},leftmargin=6em]
\item audits the current attached reduced-fibre and boundary
      programmes;
\item retires the unsupported BOGC interpretation;
\item imports the exact degree-twenty plaquette incidence;
\item proves a self-contained rooted-animal count;
\item proves summability of the two-anchor connected kernel;
\item proves the kernel-load forest theorem;
\item gives an explicit bare strong-coupling specialization;
\item formulates PIBC and proves its sufficiency; and
\item replaces it by RPIBC on the continuum route.
\end{enumerate}
\end{theorem}

\begin{proof}
Use
\Cref{audit:s6v14-YMA,
audit:s6v14-YMB,
corr:s6v14-BOGC,
lem:s6v14-animal,
thm:s6v14-kernel-sum,
thm:s6v14-kernel-load,
cor:s6v14-strong-coupling,
prop:s6v14-PIBC,
cor:s6v14-replacement}.
\end{proof}

\paragraph{Parent-version record.}
V14 was formed from
\texttt{Renewal\_Yang\_Mills\_Mass\_Gap\_Research\_Notes\_V13.tex}.
The V13 parent had \(5\,846\) logical source lines, \(205\,669\)
bytes, and SHA--256
\[
 \texttt{659aac50cf47c3df5aae9ca655669834797f059ffb0ee7ff1bbdd12adf9b43ba}.
\]
The next compulsory companion audit is V15.

\clearpage
\chapter{Sixth-series V15: companion audit and the connected
operational hypergraph correction}
\label{ch:s6v15-hypergraph}

\section{Compulsory five-version companion audit}
\label{sec:s6v15-companions}

This is the audit required at V15.  The companion tasks were
advancing concurrently.  The following records identify the exact
files inspected at the time of the audit.

\begin{audit}[Standard Model--Einstein companion]
\label{audit:s6v15-SM}
The inspected file was
\[
\texttt{Renewal\_SM\_Einstein\_Continuum\_Research\_Notes\_V36.tex},
\]
with \(593\,205\) bytes, \(13\,190\) logical source lines, and
SHA--256
\[
\texttt{ee42aa4eaa7a21b4d5860aaef80a25e2e257829e28557a0d7ad4d000241978f9}.
\]
Its V32--V36 chain proves two facts relevant here.
\begin{enumerate}[label=\textup{(\roman*)},leftmargin=3.5em]
\item An explicit edgewise subdivision of one fixed finite
      complex has uniformly bounded degree and a finite normalized
      one-star catalogue.
\item One-star gaps do not imply a uniform gap on collars whose
      rescaled length diverges.  The remaining proof requires
      stratified tangent operators, absence of normal bound states,
      finite-section stability, and compatible two-tail/interface
      cards.
\end{enumerate}
Thus finite star combinatorics controls local incidence but does
not by itself control a cofinal nonlocal kernel.
\end{audit}

\begin{audit}[Navier--Stokes companion]
\label{audit:s6v15-NS}
The Navier--Stokes task advanced from V56 to V57 during this audit.
The current file actually inspected was
\[
\texttt{Renewal\_NS\_Stability\_Research\_Notes\_V57.tex},
\]
with \(571\,388\) bytes, \(10\,131\) logical source lines, and
SHA--256
\[
\texttt{6e3ebea8c9508fbf1e6dab5b2abc74f7d0d4fe2d288297de02e571b6d8ad4b1b}.
\]
It retains the nonlocal pressure as the Hessian of a Newtonian
potential.  On one cell it splits that Hessian into an exact trace
term and a traceless Riesz term.  On two separated cells it derives
the leading inverse-cube multipole pairing, whose sign depends on
the difference of velocity anisotropies.  The cell masses nearly
cancel and their first moment is fixed by the momentum identity,
but the available class does not force the total pressure-work
sign.
\end{audit}

\begin{theorem}[Companion-audit decision]
\label{thm:s6v15-audit-decision}
The YM continuation must make both of the following corrections.
\begin{enumerate}[label=\textup{(\roman*)},leftmargin=3.5em]
\item Replace a purely local finite-star conclusion by a
      stratified, cofinal kernel statement with explicit interface
      stability.
\item Retain Ward, source-short, and moment cancellations in the
      connected response before taking absolute local norms.
\end{enumerate}
Neither companion proves an estimate for the Yang--Mills renewal
grouping coefficient, so no numerical constant is imported.
\end{theorem}

\begin{proof}
Item (i) is the one-star-versus-collar distinction in
\Cref{audit:s6v15-SM}.  Item (ii) is the global
potential/Hessian and mass--dipole organization in
\Cref{audit:s6v15-NS}.  The two PDEs and their source variables
are different, so only the proof architecture, not an inequality,
is transferred.
\end{proof}

\section{Grand Tensor source audit: the primitive object is a
hypergraph}
\label{sec:s6v15-GT}

\begin{audit}[Grand Tensor V067]
\label{audit:s6v15-GT}
The attached file
\[
\texttt{Renewal\_Grand\_Tensor\_Monograph\_V067.tex}
\]
has \(2\,804\,854\) bytes, \(64\,475\) logical source lines, and
SHA--256
\[
\texttt{facd058ec8bb210ad8b296c9d7f4d78a229d08d7a2ae378808700655cb8bc0fe}.
\]
Its connected operational support construction first shorts a
mixed response by every proper-partition product, declared
mediator, lower response, and nuisance coordinate.  The remaining
positive Gram decides whether a support is a genuine connected
hyperedge.  It proves that the maximal nonzero supports form a
canonical interaction hypergraph and gives a three-sign example
whose pairwise connected panels vanish while its three-body
response is nonzero.
\end{audit}

\begin{lemma}[Self-contained pairwise no-go]
\label{lem:s6v15-pairwise-no-go}
For
\[
 p_\theta(x_1,x_2,x_3)
 ={1\over8}(1+\theta x_1x_2x_3),
 \qquad x_i\in\{-1,1\},\quad |\theta|\le1,
\]
every one- and two-coordinate marginal is independent of
\(\theta\), but
\begin{equation}
 {\mathbb E}_\theta[x_1x_2x_3]=\theta.                   \label{eq:s6v15-three-body}
\end{equation}
Hence pairwise connected data do not determine the connected
support of a simultaneous source coefficient.
\end{lemma}

\begin{proof}
Summing over any omitted sign cancels the cubic monomial, leaving
the uniform proper marginal.  Multiplication by
\(x_1x_2x_3\) and summation over all eight signs cancels the
constant term and leaves \(\theta\).
\end{proof}

\begin{correction}[RPIBC is only a pair-sector card]
\label{corr:s6v15-RPIBC}
The conclusion of \Cref{cor:s6v14-replacement} is valid only after
one proves that all higher connected operational support Grams
vanish, or controls them separately.  A summable pair kernel cannot
by itself prove full SISC, because
\Cref{lem:s6v15-pairwise-no-go} permits an invisible higher-body
source.  RPIBC remains a correct sufficient estimate for the
pairwise operational sector.
\end{correction}

\begin{proof}
The kernel-load theorem itself is unchanged.  What fails is the
unstated exhaustiveness of pair supports.  The counterexample
shows that exhaustiveness is not formal.
\end{proof}

\section{Exact hypergraph activities for owned components}
\label{sec:s6v15-activities}

Let the resolved owned components be indexed by
\([N]=\{1,\ldots,N\}\), with masses \(x_\alpha>0\).  For
\(\varnothing\ne A\subset[N]\), define the complete internal
component-joining weight
\begin{equation}
 \eta_x(A)
 :=
 \begin{cases}
  1,&|A|=1,\\
  \displaystyle
  \left(\prod_{\alpha\in A}x_\alpha\right)
  \left(\sum_{\alpha\in A}x_\alpha\right)^{|A|-2},
  &|A|\ge2.
 \end{cases}                                             \label{eq:s6v15-eta}
\end{equation}
For \(|A|\ge2\), this is the exact weighted Cayley sum over all
trees joining the component vertices in \(A\), as proved in V8.

Let \(\vartheta_A\ge0\) denote a pre-norm majorant of the actual
connected analytic source coefficient on \(A\), after all
proper-partition products and declared nuisance/mediator ranges
have been shorted in one common carrier.  Put
\begin{equation}
 z_A:=\vartheta_A\eta_x(A).                               \label{eq:s6v15-activity}
\end{equation}
Singletons carry no grouping defect and will be left implicit.

\begin{lemma}[Occurrence partitions are dominated by hyperedge
subsets]
\label{lem:s6v15-partition-subsets}
Let \(\mathfrak P_N\) be any family of partitions of \([N]\)
allowed by the physical source expansion.  Then
\begin{equation}
 \sum_{\Pi\in\mathfrak P_N}
 \prod_{\substack{A\in\Pi\\|A|\ge2}}z_A
 \le
 \prod_{\substack{\varnothing\ne A\subset[N]\\|A|\ge2}}
 (1+z_A).                                                \label{eq:s6v15-partition-product}
\end{equation}
\end{lemma}

\begin{proof}
The nonsingleton blocks of a partition form a pairwise disjoint
set of hyperedges.  Forgetting disjointness and allowing an
arbitrary subset of all nonsingleton hyperedges only enlarges the
positive sum.  The enlarged subset sum factorizes into the product
on the right.
\end{proof}

\begin{theorem}[Hypergraph activity theorem]
\label{thm:s6v15-hypergraph}
Suppose there is a cutoff- and shell-independent constant
\(\kappa_{\rm hyp}\) such that
\begin{equation}
 \sup_{\alpha\in[N]}
 \sum_{\substack{A\ni\alpha\\|A|\ge2}}
 |A|z_A
 \le\kappa_{\rm hyp}.                                    \label{eq:s6v15-KP}
\end{equation}
Then
\begin{equation}
 \sum_{\Pi\in\mathfrak P_N}
 \prod_{\substack{A\in\Pi\\|A|\ge2}}z_A
 \le e^{\kappa_{\rm hyp}N}
 \le e^{\kappa_{\rm hyp}m}.                              \label{eq:s6v15-hypergraph-bound}
\end{equation}
Thus all component-grouping ranks are bounded simultaneously by
one strong base.
\end{theorem}

\begin{proof}
By \(\log(1+u)\le u\) and
\Cref{lem:s6v15-partition-subsets}, the logarithm of the right side
of \eqref{eq:s6v15-partition-product} is at most
\(\sum_Az_A\).  On the other hand,
\[
 \sum_{\alpha=1}^N
 \sum_{\substack{A\ni\alpha\\|A|\ge2}}|A|z_A
 =
 \sum_{\substack{A\subset[N]\\|A|\ge2}}|A|^2z_A
 \ge
 \sum_Az_A.
\]
Hypothesis \eqref{eq:s6v15-KP} therefore gives
\(\sum_Az_A\le\kappa_{\rm hyp}N\).
Finally \(N\le m-1\) by the resolved incidence-tree edge count.
\end{proof}

\begin{definition}[Connected operational hypergraph activity
card, COHAC]
\label{def:s6v15-COHAC}
COHAC consists of:
\begin{enumerate}[label=\textup{(\roman*)},leftmargin=3.5em]
\item constructing every \(\vartheta_A\) from the actual
      simultaneous Yang--Mills source response after one common
      proper-partition/nuisance short;
\item proving coefficientwise domination by
      \(z_A=\vartheta_A\eta_x(A)\), with the exact occurrence type
      and Duhamel normalization retained;
\item proving the uniform activity row sum
      \eqref{eq:s6v15-KP}; and
\item transporting these activities through IGMC, response
      matching, and the cofinal shell maps without changing their
      primitive source ancestry.
\end{enumerate}
\end{definition}

\begin{proposition}[COHAC closes the live SISC row]
\label{prop:s6v15-COHAC}
COHAC proves PSCSC, SOCC, and SISC.  It therefore closes the
high-excess grouping cards RFAC and TFLAC without a separate
complete-graph determinant estimate.
\end{proposition}

\begin{proof}
The exact Cayley factor inside every primitive occurrence is
\(\eta_x(A)\), and COHAC(ii) majorizes its simultaneous analytic
coefficient by \(z_A\).  The complete sum over physical occurrence
partitions is bounded by \Cref{thm:s6v15-hypergraph}.  This is the
joint pre-norm normalization required by PSCSC.  Apply
\Cref{prop:s6v8-PSCSC,prop:s6v7-SOCC}.  Equivalently, expanding
each \(\eta_x(A)\) into its tree gives the grouping forest of V9,
so the same bound proves TFLAC and RFAC.
\end{proof}

\begin{corollary}[Earlier cards as special mechanisms]
\label{cor:s6v15-special}
\begin{enumerate}[label=\textup{(\roman*)},leftmargin=3.5em]
\item The uniformly local star-cover theorem of V13 verifies
      COHAC when all nonzero supports have uniformly bounded local
      carriers and their total activities have bounded overlap.
\item RPIBC verifies the pair part of COHAC when
      \(z_{\{\alpha,\beta\}}\) is dominated by its summable
      two-anchor kernel.
\item FSAAC supplies the analytic \(1/q!\) normalization on the
      rank-balanced part but does not replace the hypergraph row
      sum in defect-dominant sectors.
\end{enumerate}
\end{corollary}

\begin{proof}
For (i), sum the hyperedge activities inside each bounded-load star
and use the overlap estimate of
\Cref{thm:s6v13-star-cover}.  For (ii), the calculation in
\Cref{thm:s6v14-kernel-load} is the \(|A|=2\) row sum, up to a
fixed factor two.  Item (iii) is
\Cref{thm:s6v12-balanced,cor:s6v12-defect-dominant}.
\end{proof}

\section{The stratified and moment-retaining form of COHAC}
\label{sec:s6v15-stratified}

\begin{definition}[Stratified COHAC, S--COHAC]
\label{def:s6v15-SCOHAC}
S--COHAC is COHAC together with the following cofinal interface
data:
\begin{enumerate}[label=\textup{(\roman*)},leftmargin=3.5em]
\item a finite family of bulk, face, edge, and vertex tangent
      incidence hypergraphs produced by one global subdivision;
\item a uniform activity row sum on every infinite or periodic
      tangential model;
\item absence of a normal-interface activity resonance or bound
      state which destroys that row sum;
\item finite-collar stability under the declared boundary
      closure; and
\item a gluing estimate whose commutator/interface debit is
      summable in the same strong base.
\end{enumerate}
\end{definition}

\begin{assessment}[SM lesson applied without overclaim]
\label{ass:s6v15-SM-lesson}
The finite star list proves that only finitely many local incidence
types must be checked.  It does not prove their infinite tangential
row sums or their finite-collar stability.  S--COHAC is the exact
analogue of the SM companion's correction from one-star gaps to
stratified tangent and interface cards.
\end{assessment}

\begin{definition}[Moment-retaining COHAC, M--COHAC]
\label{def:s6v15-MCOHAC}
M--COHAC is the requirement that the source-short connected
response used to form \(\vartheta_A\) retain its exact Ward-null
moments, trace/traceless split, and first nonzero multipole before
an absolute activity is assigned.  Only the residual after these
orthogonal or algebraic cancellations is required to satisfy
\eqref{eq:s6v15-KP}.
\end{definition}

\begin{assessment}[NS lesson applied without importing pressure]
\label{ass:s6v15-NS-lesson}
The NS companion shows concretely that a nonlocal kernel can have
small total mass, a fixed dipole, and sign-sensitive anisotropic
pairings even when a Gaussian absolute bound is inadequate.  YM
does not inherit the Newtonian kernel or its sign.  The transferable
point is the order of operations: exact source moments and
operator decomposition precede the absolute row sum.  M--COHAC
records that order.
\end{assessment}

\begin{definition}[Cofinal moment-stratified hypergraph card,
CMSHC]
\label{def:s6v15-CMSHC}
CMSHC is S--COHAC and M--COHAC on one common response-matched
cutoff route.  It is the current full replacement for the pair-only
RPIBC proposal.
\end{definition}

\begin{proposition}[CMSHC is sufficient for the open simultaneous
source card]
\label{prop:s6v15-CMSHC}
CMSHC proves COHAC and hence the remaining SISC estimate.
\end{proposition}

\begin{proof}
M--COHAC constructs the genuine residual connected activities
without losing the exact moment cancellations.  S--COHAC proves
their uniform row sum and cofinal transport.  These are precisely
COHAC(i)--(iv); apply \Cref{prop:s6v15-COHAC}.
\end{proof}

\begin{assessment}[Exact live bottleneck after V15]
\label{ass:s6v15-bottleneck}
The combinatorial forest problem is now closed once the actual
connected operational hyperedge activities satisfy one KP-type row
sum.  This formulation survives higher-body responses and cofinal
interfaces.  The first unproved step is no longer a tree count: it
is the source-ancestry construction of
\(\vartheta_A\) from the response-matched Yang--Mills coefficient,
followed by a moment-retaining, stratified row-sum estimate.  That
is CMSHC.
\end{assessment}

\begin{programme}[V16: construct the connected activities]
\label{prog:s6v15-V16}
\begin{enumerate}[leftmargin=3em]
\item Define the common response carrier, the proper-partition
      nuisance span, and the shorted connected row for one finite
      source support \(A\).
\item Prove that its norm is invariant under source-coordinate and
      response-frame isometries.
\item Relate its squared norm to the positive connected support
      Gram in the Grand Tensor construction.
\item Attach the exact component Cayley factor
      \(\eta_x(A)\) only after the proper-partition short.
\item Derive a connected-carrier sufficient condition for the
      activity row sum and identify the first missing uniform
      shell estimate.
\end{enumerate}
\end{programme}

\begin{firewall}[Claim boundary after sixth-series V15]
\label{fire:s6v15-boundary}
V15 performs the compulsory SM V36 and NS V57 audit and records
their exact consequences for the proof order.  It audits Grand
Tensor V067, proves the pairwise support no-go, and corrects RPIBC
to a pair-sector card.  It defines exact component hyperedge
activities, proves the hypergraph activity theorem, formulates
COHAC and proves its sufficiency, and refines it to the cofinal
moment-stratified card CMSHC.

V15 does not prove COHAC or CMSHC for the actual response-matched
Yang--Mills source, RFAC in full, SISC in full, the simple
high-excess row, arbitrary repeated transversal supports, full
TOCC, full TCNC, cyclic TIC, HLTC, full MCWC, OCRC, SCNGC, WPSC,
GRLTC, LZTC, PLRC, WSDC, BRSC, HWSFC, UGCGC, thermal connected
WPRAC, all-arity RGSC, the raw Wilson aggregate strong card,
all-order strong MTA, WUFC, CPM, cutoff removal, reflection
positivity, Osterwalder--Schrader reconstruction, a continuum
Yang--Mills measure, or a positive Hamiltonian gap.  The full
Yang--Mills mass gap has not been proved.
\end{firewall}

\begin{theorem}[Sixth-series V15 result ledger]
\label{thm:s6v15-results}
V15:
\begin{enumerate}[label=\textup{(V15.\arabic*)},leftmargin=6em]
\item audits SM V36 and NS V57;
\item records the stratified-interface and retained-nonlocal
      lessons;
\item audits the Grand Tensor operational hypergraph;
\item proves a self-contained pairwise support no-go;
\item corrects RPIBC to the pair sector;
\item defines the exact Cayley-weighted hyperedge activities;
\item proves the all-rank hypergraph activity theorem;
\item formulates COHAC and proves its sufficiency; and
\item formulates CMSHC as the cofinal moment-stratified target.
\end{enumerate}
\end{theorem}

\begin{proof}
Use
\Cref{audit:s6v15-SM,
audit:s6v15-NS,
thm:s6v15-audit-decision,
audit:s6v15-GT,
lem:s6v15-pairwise-no-go,
corr:s6v15-RPIBC,
lem:s6v15-partition-subsets,
thm:s6v15-hypergraph,
prop:s6v15-COHAC,
prop:s6v15-CMSHC}.
\end{proof}

\paragraph{Parent-version record.}
V15 was formed from
\texttt{Renewal\_Yang\_Mills\_Mass\_Gap\_Research\_Notes\_V14.tex}.
The V14 parent had \(6\,277\) logical source lines, \(220\,924\)
bytes, and SHA--256
\[
 \texttt{b08312b2b8e58f5b91ba81b504899df6478cb0377873c596423ca0feed7f35bc}.
\]
The next compulsory companion audit is V20.

\clearpage
\chapter{Sixth-series V16: canonical connected source rows and
a carrier criterion for COHAC}
\label{ch:s6v16-connected-row}

\section{The proper-partition short}
\label{sec:s6v16-short}

No companion audit is due at V16.

Fix a nonempty component support \(A\).  Let \(E_A\) be its finite
polarization/coefficient Hilbert space, let \({\cal R}_A\) be one
common response Hilbert carrier, and let
\[
 F_A:E_A\longrightarrow{\cal R}_A
\]
be the full simultaneous response row before separate source
norms.  Let \({\cal N}_A\subset{\cal R}_A\) be the closed span of:
\begin{enumerate}[label=\textup{(\roman*)},leftmargin=3.5em]
\item products synthesized from proper partitions of \(A\);
\item declared lower-order response rows;
\item mediator ranges; and
\item nuisance or gauge-vertical response coordinates.
\end{enumerate}
Write \(P_A\) for the orthogonal projection onto \({\cal N}_A\).

\begin{definition}[Canonical connected row and Gram]
\label{def:s6v16-connected}
Define
\begin{equation}
 \widehat F_A:=(I-P_A)F_A,\qquad
 {\mathbb C}_A:=\widehat F_A^*\widehat F_A
 =F_A^*(I-P_A)F_A.                                      \label{eq:s6v16-connected}
\end{equation}
\end{definition}

\begin{theorem}[Positivity, vanishing, and least residual
majorant]
\label{thm:s6v16-short}
The operator \({\mathbb C}_A\) is positive semidefinite and
\begin{equation}
 {\mathbb C}_A=0
 \quad\Longleftrightarrow\quad
 \operatorname{Ran}F_A\subset{\cal N}_A.                 \label{eq:s6v16-vanishing}
\end{equation}
Moreover,
\begin{equation}
 \|\widehat F_A\|
 =
 \inf\left\{
 \|F_A-G\|:
 G:E_A\to{\cal R}_A,\ 
 \operatorname{Ran}G\subset{\cal N}_A
 \right\}.                                               \label{eq:s6v16-distance}
\end{equation}
Thus \(\|\widehat F_A\|=\|{\mathbb C}_A\|^{1/2}\) is the least
uniform operator majorant of the unexplained connected response.
\end{theorem}

\begin{proof}
Positivity and
\(\|\widehat F_A\|^2=\|\widehat F_A^*\widehat F_A\|\) are
immediate.  The Gram vanishes exactly when
\(\widehat F_A=0\), which is
\(\operatorname{Ran}F_A\subset\operatorname{Ran}P_A\).

For every admissible \(G\),
\[
 (I-P_A)(F_A-G)=(I-P_A)F_A,
\]
so \(\|F_A-G\|\ge\|\widehat F_A\|\).  Equality is attained by
\(G=P_AF_A\), proving \eqref{eq:s6v16-distance}.
\end{proof}

\begin{theorem}[Unitary covariance of the connected activity]
\label{thm:s6v16-covariance}
Let \(U:{\cal R}_A\to{\cal R}'_A\) and
\(V:E'_A\to E_A\) be unitaries.  If
\[
 F'_A=UF_AV,\qquad {\cal N}'_A=U{\cal N}_A,
\]
then
\begin{equation}
 \widehat F'_A=U\widehat F_AV,\qquad
 {\mathbb C}'_A=V^*{\mathbb C}_AV.                       \label{eq:s6v16-covariance}
\end{equation}
In particular, vanishing, rank, spectrum, and operator norm of the
connected Gram are invariant under source-coordinate and
response-frame isometries.
\end{theorem}

\begin{proof}
The projection onto \(U{\cal N}_A\) is \(UP_AU^*\).
Substitution gives the first identity; taking the Gram gives the
second.
\end{proof}

\begin{warning}[All proper partitions must be shorted
simultaneously]
\label{warn:s6v16-common-carrier}
Projecting separately against one proper partition at a time need
not equal projection against their joint span.  Likewise, Grams
formed in independently normalized carriers do not determine their
mixed relative orientation.  The single projection \(P_A\) in
\eqref{eq:s6v16-connected} is therefore part of the definition,
not a replaceable bookkeeping choice.
\end{warning}

\section{The canonical analytic activity}
\label{sec:s6v16-activity}

Let \(d_A\ge0\) contain the exact time-simplex, occurrence-type,
coupling, and shell coefficient attached to support \(A\), before
the response norm.  Let \(s_A\in E_A\) be its normalized physical
source polarization, \(\|s_A\|\le1\).

\begin{definition}[Connected response activity]
\label{def:s6v16-theta}
Put
\begin{equation}
 \vartheta_A:=d_A\|\widehat F_A\|,
 \qquad
 z_A:=\vartheta_A\eta_x(A),                              \label{eq:s6v16-activity}
\end{equation}
where \(\eta_x(A)\) is the exact component Cayley polynomial
\eqref{eq:s6v15-eta}.
\end{definition}

\begin{lemma}[Coefficientwise residual domination]
\label{lem:s6v16-domination}
The unexplained connected part of the physical coefficient obeys
\begin{equation}
 d_A\|(I-P_A)F_As_A\|
 \le\vartheta_A.                                         \label{eq:s6v16-domination}
\end{equation}
If \({\mathbb C}_A=0\), both sides vanish and no hyperedge on
\(A\) is required.
\end{lemma}

\begin{proof}
Use \(\|s_A\|\le1\) and the operator norm.  The vanishing statement
is \Cref{thm:s6v16-short}.
\end{proof}

\begin{assessment}[Why the Cayley factor is attached after the
short]
\label{ass:s6v16-order}
The proper-partition span decides whether a simultaneous support is
genuinely new.  Multiplying separate proper-subset norms by their
component Cayley factors first would recreate the V8
superexponential fibre and could turn an explained product response
into a false primitive occurrence.  Definition
\eqref{eq:s6v16-activity} first forms the canonical connected
residual and only then attaches the exact joining polynomial.
\end{assessment}

\begin{proposition}[Canonical-activity COHAC reduction]
\label{prop:s6v16-COHAC-reduction}
Suppose the actual pre-norm simultaneous coefficient is represented
in the common carriers above, and suppose
\begin{equation}
 \sup_{\alpha}
 \sum_{\substack{A\ni\alpha\\|A|\ge2}}
 |A|\,d_A\|\widehat F_A\|\eta_x(A)
 \le\kappa_{\rm hyp}.                                    \label{eq:s6v16-row-sum}
\end{equation}
Then COHAC holds and the open SISC grouping row closes.
\end{proposition}

\begin{proof}
\Cref{lem:s6v16-domination} proves coefficientwise domination by
the activities \(z_A\).  Equation
\eqref{eq:s6v16-row-sum} is precisely the hypergraph activity
condition \eqref{eq:s6v15-KP}.  Unitary covariance supplies the
coordinate requirement.  Apply
\Cref{thm:s6v15-hypergraph,prop:s6v15-COHAC}.
\end{proof}

\section{A connected-carrier sufficient condition}
\label{sec:s6v16-carrier}

Let \(\mathscr G=(P,E)\) have maximum degree \(\Delta\ge2\), and
anchor each component \(C_\alpha\) at \(p(\alpha)\in P\).  For every
nonzero activity \(z_A\), choose a finite connected carrier
\(\Gamma(A)\subset P\) which contains every \(p(\alpha)\),
\(\alpha\in A\).

\begin{definition}[Carrier multiplicity profile]
\label{def:s6v16-carrier-profile}
Fix \(C_0<\infty\), an integer \(r\ge0\), and \(a\ge0\).  The
activities have carrier profile \((C_0,r,a)\) when, for every
component index \(\alpha\) and every connected
\(\Gamma\ni p(\alpha)\),
\begin{equation}
 \sum_{\substack{A\ni\alpha\\\Gamma(A)=\Gamma}}
 |A|z_A
 \le C_0|\Gamma|^r a^{|\Gamma|}.                         \label{eq:s6v16-carrier-profile}
\end{equation}
\end{definition}

\begin{theorem}[Connected-carrier activity theorem]
\label{thm:s6v16-carrier}
Let \(A_\Delta=\Delta^2\) as in
\Cref{lem:s6v14-animal}.  If the activities have carrier profile
\((C_0,r,a)\) and
\begin{equation}
 A_\Delta a<1,                                           \label{eq:s6v16-small-carrier}
\end{equation}
then the COHAC row sum holds with
\begin{equation}
 \kappa_{\rm hyp}
 =
 C_0a\sum_{n\ge1}n^r(A_\Delta a)^{n-1}<\infty.           \label{eq:s6v16-kappa}
\end{equation}
\end{theorem}

\begin{proof}
Fix \(\alpha\).  Group the hyperedges in the left side of
\eqref{eq:s6v15-KP} by their carrier.  The carrier of any
hyperedge containing \(\alpha\) is a connected vertex set
containing \(p(\alpha)\).  By
\Cref{def:s6v16-carrier-profile,lem:s6v14-animal},
\begin{align*}
 \sum_{A\ni\alpha}|A|z_A
 &\le
 \sum_{\Gamma\ni p(\alpha)}
 C_0|\Gamma|^ra^{|\Gamma|}\\
 &\le
 C_0\sum_{n\ge1}
 A_\Delta^{n-1}n^ra^n,
\end{align*}
which is \eqref{eq:s6v16-kappa}.  The power series is finite under
\eqref{eq:s6v16-small-carrier}.
\end{proof}

\begin{corollary}[All-rank source closure from connected-carrier
decay]
\label{cor:s6v16-carrier-closure}
Under the hypotheses of
\Cref{prop:s6v16-COHAC-reduction}, it is enough to establish the
carrier profile
\eqref{eq:s6v16-carrier-profile} with
\(A_\Delta a<1\).  Then PSCSC, SOCC, SISC, RFAC, and TFLAC hold.
\end{corollary}

\begin{proof}
\Cref{thm:s6v16-carrier} proves the row sum required by
\Cref{prop:s6v16-COHAC-reduction}.
\end{proof}

\begin{remark}[Higher-body responses are included]
\label{rem:s6v16-higher}
No bound on \(|A|\) is imposed.  The sum over all hyperedges with
the same connected carrier occurs inside
\eqref{eq:s6v16-carrier-profile}.  Thus a genuine triple or
higher-body connected support is paid by the carrier activity
rather than replaced by its pairwise two-section.
\end{remark}

\section{Cofinal stratified carriers}
\label{sec:s6v16-strata}

\begin{proposition}[Finite tangent-type reduction]
\label{prop:s6v16-strata}
Suppose one global subdivision produces finitely many tangent
incidence types \(\tau\), each with degree bound
\(\Delta_\tau\le\Delta_*\).  Suppose every bulk or interface
carrier has a type \(\tau\) and satisfies
\eqref{eq:s6v16-carrier-profile} with common \(C_0,r,a\), and
\begin{equation}
 \Delta_*^2a<1.                                          \label{eq:s6v16-strata-small}
\end{equation}
If the finite-collar closure changes every activity by at most a
common multiplicative factor \(C_{\rm int}\), then S--COHAC holds
with
\begin{equation}
 \kappa_{\rm hyp}
 =
 C_{\rm int}C_0a
 \sum_{n\ge1}n^r(\Delta_*^2a)^{n-1}.                    \label{eq:s6v16-strata-kappa}
\end{equation}
\end{proposition}

\begin{proof}
Apply \Cref{thm:s6v16-carrier} on each tangent type with the common
maximum degree.  There are finitely many types, so their finite
type multiplicity is absorbed into \(C_0\).  The interface
comparison multiplies the activity row by at most
\(C_{\rm int}\).  These are exactly the tangent row-sum and
finite-collar items of S--COHAC.
\end{proof}

\begin{warning}[A local catalogue does not give
\(C_{\rm int}\)]
\label{warn:s6v16-interface}
The constant \(C_{\rm int}\) is not a consequence of finite local
degree.  A long collar can carry an interface resonance or a normal
bound state while every one-star model is regular.  This is the
precise point imported from the SM V36 audit.
\end{warning}

\section{The first missing Yang--Mills estimate}
\label{sec:s6v16-missing}

\begin{definition}[Response-matched carrier activity card, RCAC]
\label{def:s6v16-RCAC}
RCAC is the assertion that the canonical activities
\[
 z_{A,r}
 =
 d_{A,r}\| (I-P_{A,r})F_{A,r}\|\eta_x(A)
\]
on the actual response-matched Yang--Mills shell \(r\) admit:
\begin{enumerate}[label=\textup{(\roman*)},leftmargin=3.5em]
\item connected physical carriers \(\Gamma_r(A)\);
\item a shell-uniform carrier profile
      \eqref{eq:s6v16-carrier-profile};
\item a common smallness margin
      \(\Delta_*^2a<1\); and
\item the interface comparison \(C_{\rm int}<\infty\).
\end{enumerate}
\end{definition}

\begin{proposition}[RCAC completes CMSHC]
\label{prop:s6v16-RCAC}
If the source moments used in forming \(P_{A,r}\) include the
Ward-null and response-matching nuisance directions of M--COHAC,
then RCAC proves CMSHC and hence the remaining SISC row.
\end{proposition}

\begin{proof}
The moment-retaining choice of \(P_{A,r}\) gives M--COHAC.
\Cref{prop:s6v16-strata} gives S--COHAC.  Together they are CMSHC;
apply \Cref{prop:s6v15-CMSHC}.
\end{proof}

\begin{assessment}[Exact live bottleneck after V16]
\label{ass:s6v16-bottleneck}
The source activity is now canonical at every finite cutoff: it is
the distance of the simultaneous response from the common
proper-partition/nuisance span, multiplied only afterwards by the
exact component Cayley polynomial.  The remaining analytic
question is RCAC.  In particular, one must prove that the
renormalized connected residual per carrier decays fast enough that
its activity parameter beats the graph-animal entropy, and that
this margin survives every stratified interface.
\end{assessment}

\begin{programme}[V17: locate the carrier decay in the inherited
Yang--Mills bounds]
\label{prog:s6v16-V17}
\begin{enumerate}[leftmargin=3em]
\item Compare the canonical row
      \((I-P_{A,r})F_{A,r}\) with the connected score-covariance
      and local curvature-tail decompositions in the attached
      reduced-fibre programme.
\item Separate deterministic local head, connected covariance, and
      response-matching correction before taking norms.
\item Extract the best proved spatial decay for each piece.
\item Test whether their combined carrier parameter satisfies
      \(\Delta_*^2a<1\) on any nontrivial shell range.
\item If only a finite weighted row sum is available, replace the
      small-activity parameter by a direct Schur/activity estimate.
\end{enumerate}
\end{programme}

\begin{firewall}[Claim boundary after sixth-series V16]
\label{fire:s6v16-boundary}
V16 constructs the canonical connected source row as one
proper-partition/nuisance short, proves its least-residual property
and unitary covariance, and defines the corresponding analytic
hyperedge activity.  It proves a connected-carrier criterion for
the full COHAC row sum, includes arbitrary higher-body supports,
gives a finite tangent-type/interface reduction, and isolates RCAC.

V16 does not prove RCAC, CMSHC, COHAC for the actual
response-matched Yang--Mills source, RFAC in full, SISC in full,
the simple high-excess row, arbitrary repeated transversal
supports, full TOCC, full TCNC, cyclic TIC, HLTC, full MCWC, OCRC,
SCNGC, WPSC, GRLTC, LZTC, PLRC, WSDC, BRSC, HWSFC, UGCGC,
thermal connected WPRAC, all-arity RGSC, the raw Wilson aggregate
strong card, all-order strong MTA, WUFC, CPM, cutoff removal,
reflection positivity, Osterwalder--Schrader reconstruction, a
continuum Yang--Mills measure, or a positive Hamiltonian gap.  The
full Yang--Mills mass gap has not been proved.
\end{firewall}

\begin{theorem}[Sixth-series V16 result ledger]
\label{thm:s6v16-results}
V16:
\begin{enumerate}[label=\textup{(V16.\arabic*)},leftmargin=6em]
\item defines the common proper-partition short;
\item proves positivity, vanishing, and its least-residual formula;
\item proves unitary covariance;
\item defines the canonical connected response activity;
\item proves coefficientwise residual domination;
\item reduces COHAC to its canonical activity row sum;
\item proves the connected-carrier activity theorem;
\item proves the finite tangent-type/interface reduction; and
\item formulates RCAC and proves its sufficiency.
\end{enumerate}
\end{theorem}

\begin{proof}
Use
\Cref{thm:s6v16-short,
thm:s6v16-covariance,
lem:s6v16-domination,
prop:s6v16-COHAC-reduction,
thm:s6v16-carrier,
prop:s6v16-strata,
prop:s6v16-RCAC}.
\end{proof}

\paragraph{Parent-version record.}
V16 was formed from
\texttt{Renewal\_Yang\_Mills\_Mass\_Gap\_Research\_Notes\_V15.tex}.
The V15 parent had \(6\,743\) logical source lines, \(238\,512\)
bytes, and SHA--256
\[
 \texttt{bd4243291f737faf4e835607f11742554c44b4af1a781716073d021fa926531b}.
\]
The next compulsory companion audit is V20.

\clearpage
\chapter{Sixth-series V17: the exact root-cofactor connected
kernel and its conditional activity sum}
\label{ch:s6v17-root-kernel}

\section{The inherited connected-kernel prototype}
\label{sec:s6v17-prototype}

No companion audit is due at V17.

The reduced-fibre programme audited in V14 contains one source
coefficient whose ancestry can be compared directly with the
canonical short of V16.  Let \(Z_*\) be the selected partition
function, let \({\mathscr D}^{\rm norm}\) be the normalized root
equation, and let \(X\) be an action variation.  Its unnormalized
cofactor derivative has the form
\begin{equation}
 {\mathfrak c}[X]
 =
 {d\over dt}\bigl[
 Z_t^2{\mathscr D}^{\rm norm}_t
 \bigr]_{t=0}.                                           \label{eq:s6v17-cofactor}
\end{equation}

\begin{lemma}[Root cancellation removes the partition derivative]
\label{lem:s6v17-root-cancellation}
If \({\mathscr D}^{\rm norm}_0=0\), then
\begin{equation}
 {\mathfrak c}[X]
 =
 Z_*^2
 {d\over dt}{\mathscr D}^{\rm norm}_t\big|_{t=0}.         \label{eq:s6v17-root-cancellation}
\end{equation}
Consequently no derivative of the unmarked partition factor is a
new connected source occurrence at the selected root.
\end{lemma}

\begin{proof}
The product rule gives
\[
 {\mathfrak c}[X]
 =
 2Z_*Z'_*
 {\mathscr D}^{\rm norm}_0
 +
 Z_*^2({\mathscr D}^{\rm norm})'_0.
\]
The first term vanishes by the root equation.
\end{proof}

\begin{assessment}[Identification with the proper-partition
short]
\label{ass:s6v17-short}
The vanished term in
\Cref{lem:s6v17-root-cancellation} is a source-disjoint vacuum or
partition response.  It belongs to the nuisance/proper-partition
span of V16.  The surviving derivative is a linear combination of
connected covariances and is therefore an explicit finite-card
prototype of \(\widehat F_A=(I-P_A)F_A\).  This identifies one
row; it does not prove the corresponding statement for every
higher support \(A\).
\end{assessment}

\section{The signed loop--plaquette kernel}
\label{sec:s6v17-loop-plaquette}

Write the frozen marginal action writer as
\[
 P=\sum_{p\in{\mathscr P}_c}p_p
\]
and let \({\mathfrak s}_{\rm Cr}\) be the centered physical Creutz
source at the same selected root and in the same response-matched
law.  Define
\begin{equation}
 \kappa_P(p)
 :=
 -\operatorname{Cov}({\mathfrak s}_{\rm Cr},p_p).         \label{eq:s6v17-kappa}
\end{equation}

\begin{theorem}[Exact connected marginal sum]
\label{thm:s6v17-exact-sum}
On every finite selected card, the marginal response coefficient
is
\begin{equation}
 d_P=\sum_{p\in{\mathscr P}_c}\kappa_P(p).                \label{eq:s6v17-exact-sum}
\end{equation}
If
\[
 {\mathfrak s}_{\rm Cr}
 =
 \sum_{\rho=1}^{R}
 a_\rho{\cal O}_\rho,
\]
then
\begin{equation}
 \kappa_P(p)
 =
 -\sum_{\rho=1}^{R}
 a_\rho\operatorname{Cov}({\cal O}_\rho,p_p).             \label{eq:s6v17-loop-expansion}
\end{equation}
\end{theorem}

\begin{proof}
Covariance is linear in each entry, and the plaquette sum is
finite.  Hence
\[
 -\operatorname{Cov}({\mathfrak s}_{\rm Cr},P)
 =
 -\sum_p\operatorname{Cov}({\mathfrak s}_{\rm Cr},p_p),
\]
which is \eqref{eq:s6v17-exact-sum}.  Expanding the source gives
\eqref{eq:s6v17-loop-expansion}.
\end{proof}

\begin{lemma}[Exact Markov-blanket factorization]
\label{lem:s6v17-Markov}
Suppose a physical boundary record \({\cal B}\) makes
\({\cal O}_\rho\) and \(p_p\) conditionally independent.  Then
\begin{equation}
 \operatorname{Cov}({\cal O}_\rho,p_p)
 =
 \operatorname{Cov}\left(
 {\mathbb E}[{\cal O}_\rho\mid{\cal B}],
 {\mathbb E}[p_p\mid{\cal B}]
 \right).                                                \label{eq:s6v17-Markov}
\end{equation}
\end{lemma}

\begin{proof}
Conditional independence gives
\[
 {\mathbb E}[{\cal O}_\rho p_p\mid{\cal B}]
 =
 {\mathbb E}[{\cal O}_\rho\mid{\cal B}]
 {\mathbb E}[p_p\mid{\cal B}].
\]
Take expectation and subtract the product of the two means.
\end{proof}

\begin{warning}[Factorization is not decay]
\label{warn:s6v17-factorization}
Equation \eqref{eq:s6v17-Markov} identifies the only legitimate
boundary messages.  It supplies no contraction coefficient between
them.  In particular, the cofinal lifetime and diagonal-tail
identities of the attached boundary programme do not imply an
exponential covariance estimate.
\end{warning}

\section{The exact conditional row-sum constant}
\label{sec:s6v17-row-sum}

For each loop writer \({\cal O}_\rho\), let
\({\mathscr S}_\rho\) be its coarse-cell support and
\(s_\rho=|{\mathscr S}_\rho|\).  Suppose
\begin{equation}
 \left|
 \operatorname{Cov}({\cal O}_\rho,p_p)
 \right|
 \le
 C_\rho e^{-\mu\operatorname{dist}(p,{\mathscr S}_\rho)}
                                                               \label{eq:s6v17-decay}
\end{equation}
for some \(\mu>0\).

\begin{lemma}[Four-dimensional plaquette-shell sum]
\label{lem:s6v17-shell-sum}
Put \(q=e^{-\mu}\).  If the number of oriented coarse plaquettes at
distance \(n\) from one support cell is at most
\(96(n+1)^3\), then
\begin{equation}
 \sum_{n\ge0}96(n+1)^3e^{-\mu n}
 =
 96\,{1+4q+q^2\over(1-q)^4}.                             \label{eq:s6v17-shell-sum}
\end{equation}
\end{lemma}

\begin{proof}
The standard differentiated geometric series gives
\[
 \sum_{n\ge0}(n+1)^3q^n
 =
 {1+4q+q^2\over(1-q)^4},
 \qquad |q|<1.
\]
Multiply by \(96\).
\end{proof}

\begin{theorem}[Uniform absolute connected-kernel mass]
\label{thm:s6v17-l1}
Under \eqref{eq:s6v17-decay},
\begin{equation}
 \sum_{p\in{\mathscr P}_c}|\kappa_P(p)|
 \le
 96\,{1+4q+q^2\over(1-q)^4}
 \sum_{\rho=1}^{R}|a_\rho|C_\rho s_\rho.                 \label{eq:s6v17-l1}
\end{equation}
The bound is independent of the torus volume.  It is uniform along
a cutoff route if \(\mu\) has a positive lower bound and the final
sum on the right has a uniform upper bound.
\end{theorem}

\begin{proof}
Use \eqref{eq:s6v17-loop-expansion}, the triangle inequality, and
\eqref{eq:s6v17-decay}.  For a fixed support cell, sum over
plaquette shells and apply \Cref{lem:s6v17-shell-sum}.  A union
bound over the \(s_\rho\) support cells proves
\eqref{eq:s6v17-l1}.
\end{proof}

\begin{corollary}[Finite head with an explicit tail]
\label{cor:s6v17-tail}
Let \(d_P^{\le L}\) be the sum in
\eqref{eq:s6v17-exact-sum} over plaquettes within distance \(L\)
of at least one \({\mathscr S}_\rho\).  Then
\begin{equation}
 |d_P-d_P^{\le L}|
 \le
 96\sum_\rho |a_\rho|C_\rho s_\rho
 \sum_{n>L}(n+1)^3e^{-\mu n}.                            \label{eq:s6v17-tail}
\end{equation}
The right side tends to zero uniformly whenever the parameters in
\Cref{thm:s6v17-l1} are uniform.
\end{corollary}

\begin{proof}
Repeat the proof of \Cref{thm:s6v17-l1}, retaining only shells
\(n>L\).  Uniform convergence is the tail convergence of a fixed
summable positive series.
\end{proof}

\section{What this proves for RCAC}
\label{sec:s6v17-RCAC}

\begin{definition}[Root-cofactor decay card, RCDC]
\label{def:s6v17-RCDC}
RCDC is the uniform form of \eqref{eq:s6v17-decay} for the actual
four loop writers in the selected root cofactor, together with:
\begin{enumerate}[label=\textup{(\roman*)},leftmargin=3.5em]
\item a uniform positive decay rate \(\mu\);
\item uniform loop-visibility denominators, equivalently uniform
      bounds on the coefficients \(a_\rho\);
\item uniform \(C_\rho\) and support sizes \(s_\rho\); and
\item strict same-history transport of the loop and plaquette
      boundary messages.
\end{enumerate}
\end{definition}

\begin{proposition}[RCDC proves the root-cofactor pair activity
row]
\label{prop:s6v17-RCDC}
Assume RCDC and suppose the resolved pair activities belonging to
the marginal response are coefficientwise dominated by
\[
 x_\alpha x_\beta
 |\kappa_P(p(\alpha),p(\beta))|
\]
with uniform anchor mass load.  Then the pair-support part of RCAC
and RPIBC holds for that marginal row, and its complete grouping
forest has a \(K^m\) bound.
\end{proposition}

\begin{proof}
\Cref{thm:s6v17-l1} gives the uniform two-anchor kernel row sum.
The uniform anchor load and
\Cref{thm:s6v14-kernel-load} give the all-rank forest bound.
The response is already connected by
\Cref{lem:s6v17-root-cancellation,thm:s6v17-exact-sum}, so no
partition derivative is charged again.
\end{proof}

\begin{assessment}[Largest sector justified by the attached
source]
\label{ass:s6v17-largest}
The exact source ancestry, connected covariance representation, and
Markov-blanket factorization are proved for the selected root
cofactor's loop--plaquette marginal.  Its exponential decay is not
proved; it is precisely RCDC.  No inherited statement supplies the
multi-support carrier profile required for arbitrary COHAC
hyperedges.  Therefore
\Cref{prop:s6v17-RCDC} is a conditional pair-row closure, not full
RCAC.
\end{assessment}

\begin{assessment}[Exact live bottleneck after V17]
\label{ass:s6v17-bottleneck}
The first missing analytic number is now explicit: a uniform
positive decay rate for the same-root connected
loop--plaquette covariance, with uniform visibility floors.  The
Markov blanket reduces this to contraction of two physical boundary
messages, but the cofinal boundary programme currently stops before
that contraction is populated.  Even after RCDC, higher connected
source supports require their own carrier-profile theorem.
\end{assessment}

\begin{programme}[V18: go upstream to boundary-message
contraction]
\label{prog:s6v17-V18}
\begin{enumerate}[leftmargin=3em]
\item Express the covariance in
      \eqref{eq:s6v17-Markov} as an inner product of centered
      boundary messages in one conditional \(L^2\) space.
\item Derive the sharp contraction estimate supplied by a
      maximal-correlation or conditional-expectation operator.
\item Iterate nested physical shells without multiplying unrelated
      scalar norms.
\item Identify the exact product or summability condition which
      yields \(e^{-\mu r}\).
\item Compare that condition with the currently open
      fixed-radius message transport in the boundary programme.
\end{enumerate}
\end{programme}

\begin{firewall}[Claim boundary after sixth-series V17]
\label{fire:s6v17-boundary}
V17 identifies the exact selected-root cancellation, proves the
connected loop--plaquette kernel sum and Markov-blanket
factorization, evaluates its four-dimensional exponential
\(\ell^1\) constant, and derives the uniform finite-head tail.  It
formulates RCDC and proves its conditional sufficiency for the
root-cofactor pair activity row.

V17 does not prove RCDC, a positive uniform boundary-message
contraction, higher-support RCAC, CMSHC, COHAC, RFAC in full, SISC
in full, the simple high-excess row, arbitrary repeated transversal
supports, full TOCC, full TCNC, cyclic TIC, HLTC, full MCWC, OCRC,
SCNGC, WPSC, GRLTC, LZTC, PLRC, WSDC, BRSC, HWSFC, UGCGC,
thermal connected WPRAC, all-arity RGSC, the raw Wilson aggregate
strong card, all-order strong MTA, WUFC, CPM, cutoff removal,
reflection positivity, Osterwalder--Schrader reconstruction, a
continuum Yang--Mills measure, or a positive Hamiltonian gap.  The
full Yang--Mills mass gap has not been proved.
\end{firewall}

\begin{theorem}[Sixth-series V17 result ledger]
\label{thm:s6v17-results}
V17:
\begin{enumerate}[label=\textup{(V17.\arabic*)},leftmargin=6em]
\item proves the selected-root partition cancellation;
\item identifies it with one proper-partition short;
\item proves the exact loop--plaquette connected-kernel sum;
\item proves its Markov-blanket factorization;
\item evaluates the four-dimensional exponential shell series;
\item proves the conditional uniform \(\ell^1\) kernel bound;
\item proves the finite-head tail estimate;
\item formulates RCDC; and
\item proves its pair-row sufficiency.
\end{enumerate}
\end{theorem}

\begin{proof}
Use
\Cref{lem:s6v17-root-cancellation,
thm:s6v17-exact-sum,
lem:s6v17-Markov,
lem:s6v17-shell-sum,
thm:s6v17-l1,
cor:s6v17-tail,
prop:s6v17-RCDC}.
\end{proof}

\paragraph{Parent-version record.}
V17 was formed from
\texttt{Renewal\_Yang\_Mills\_Mass\_Gap\_Research\_Notes\_V16.tex}.
The V16 parent had \(7\,167\) logical source lines, \(253\,484\)
bytes, and SHA--256
\[
 \texttt{610909f687d0453d2a731a3bb1c255467c80f6e871e95672d6e8b3d522b50c8d}.
\]
The next compulsory companion audit is V20.

\clearpage
\chapter{Sixth-series V18: maximal-correlation shell transport
and the noncircular boundary-decay card}
\label{ch:s6v18-correlation}

\section{One boundary channel}
\label{sec:s6v18-one-channel}

No companion audit is due at V18.

Let \((\Omega,\mu)\) be a probability space and let
\({\cal A},{\cal B}\) be two boundary-record sigma fields.  Write
\[
 H_{\cal A}:=L^2_0({\cal A},\mu),\qquad
 H_{\cal B}:=L^2_0({\cal B},\mu)
\]
for the centered real \(L^2\) spaces.  Define the conditional
expectation channel
\begin{equation}
 T_{{\cal A}\leftarrow{\cal B}}g
 :=
 {\mathbb E}[g\mid{\cal A}],
 \qquad g\in H_{\cal B}.                                 \label{eq:s6v18-channel}
\end{equation}

\begin{lemma}[Maximal correlation is a channel norm]
\label{lem:s6v18-maxcorr}
The maximal centered correlation between the two records is
\begin{equation}
 \rho({\cal A},{\cal B})
 :=
 \sup_{\substack{f\in H_{\cal A},\,g\in H_{\cal B}\\
                  \|f\|_2=\|g\|_2=1}}
 |\langle f,g\rangle|
 =
 \|T_{{\cal A}\leftarrow{\cal B}}\|
 \le1.                                                   \label{eq:s6v18-maxcorr}
\end{equation}
\end{lemma}

\begin{proof}
For \(f\) measurable with respect to \({\cal A}\),
\[
 \langle f,g\rangle
 =
 \langle f,{\mathbb E}[g\mid{\cal A}]\rangle.
\]
Taking the supremum first over unit \(f\) gives the norm of the
conditional expectation of \(g\), and then taking the supremum
over unit \(g\) gives the operator norm.  Conditional expectation
is an orthogonal projection in \(L^2\), hence a contraction.
\end{proof}

\begin{theorem}[One-shell covariance contraction]
\label{thm:s6v18-one-shell}
Suppose observables \(F,G\) are conditionally independent given a
boundary record \({\cal B}\).  Put
\[
 f={\mathbb E}[F\mid{\cal B}]-{\mathbb E}F,\qquad
 g={\mathbb E}[G\mid{\cal B}]-{\mathbb E}G.
\]
Then
\begin{equation}
 \operatorname{Cov}(F,G)=\langle f,g\rangle_{L^2({\cal B})},
 \qquad
 |\operatorname{Cov}(F,G)|\le\|f\|_2\|g\|_2.             \label{eq:s6v18-one-shell}
\end{equation}
\end{theorem}

\begin{proof}
The first identity is
\Cref{lem:s6v17-Markov}, with the means subtracted.  The second is
Cauchy--Schwarz.
\end{proof}

\section{Nested shell transport}
\label{sec:s6v18-nested}

Consider centered Hilbert spaces \(H_0,\ldots,H_n\) of successive
physical shell records and contraction channels
\[
 T_j:H_{j+1}\longrightarrow H_j,\qquad
 \rho_j:=\|T_j\|\le1.
\]

\begin{theorem}[Exact product contraction]
\label{thm:s6v18-product}
For \(g_n\in H_n\),
\begin{equation}
 \|T_0T_1\cdots T_{n-1}g_n\|_2
 \le
 \left(\prod_{j=0}^{n-1}\rho_j\right)\|g_n\|_2.           \label{eq:s6v18-product}
\end{equation}
Consequently any same-history inner/outer covariance which factors
through these nested conditional channels obeys
\begin{equation}
 |\operatorname{Cov}(F,G)|
 \le
 \|f_0\|_2\|g_n\|_2
 \prod_{j=0}^{n-1}\rho_j.                                \label{eq:s6v18-cov-product}
\end{equation}
\end{theorem}

\begin{proof}
Submultiplicativity proves \eqref{eq:s6v18-product}.  Pair the
transported outer message with the centered inner message and use
Cauchy--Schwarz to obtain \eqref{eq:s6v18-cov-product}.
\end{proof}

\begin{corollary}[Uniform and nonuniform exponential criteria]
\label{cor:s6v18-exponential}
\begin{enumerate}[label=\textup{(\roman*)},leftmargin=3.5em]
\item If \(\rho_j\le\rho_*<1\), then the product is at most
      \(e^{-\mu n}\), where \(\mu=-\log\rho_*>0\).
\item More generally, if
      \begin{equation}
       \sum_{j=0}^{n-1}(1-\rho_j)\ge\mu n-C_0,            \label{eq:s6v18-average-gap}
      \end{equation}
      then
      \begin{equation}
       \prod_{j=0}^{n-1}\rho_j
       \le e^{C_0}e^{-\mu n}.                            \label{eq:s6v18-nonuniform}
      \end{equation}
\end{enumerate}
\end{corollary}

\begin{proof}
The first statement is immediate.  For the second, use
\(\log x\le-(1-x)\) on \(0<x\le1\).  A zero \(\rho_j\) only
strengthens the conclusion.
\end{proof}

\begin{corollary}[Square-summable defects give a positive product,
not decay]
\label{cor:s6v18-product-no-go}
If \(\sum_j(1-\rho_j)<\infty\) and every \(\rho_j>0\), the product
need not tend to zero.  For example,
\(\rho_j=1-(j+2)^{-2}\) has a strictly positive infinite product.
\end{corollary}

\begin{proof}
For sufficiently large \(j\),
\(-\log(1-(j+2)^{-2})\le2(j+2)^{-2}\).  The logarithms are
summable, so the product converges to a positive number.
\end{proof}

\begin{assessment}[Cofinal convergence is not quantitative
contraction]
\label{ass:s6v18-cofinal}
Strong convergence of each fixed-radius message, tightness of a
lifetime law, or removal of its atom at infinity gives no lower
bound like \eqref{eq:s6v18-average-gap}.  This is why the attached
boundary programme's exact cofinal identities do not yet prove
RCDC.
\end{assessment}

\section{A spectral formulation and its circularity boundary}
\label{sec:s6v18-spectral}

\begin{lemma}[Local channel gap implies maximal-correlation
contraction]
\label{lem:s6v18-local-gap}
Let \(K:H\to H\) be a self-adjoint Markov contraction with constants
as its fixed space.  If
\begin{equation}
 \langle h,(I-K)h\rangle
 \ge\lambda\|h\|_2^2
 \quad(h\perp1)                                          \label{eq:s6v18-local-gap}
\end{equation}
for some \(0<\lambda\le1\), and \(K\) is nonnegative on the
centered space, then
\begin{equation}
 \|K|_{L^2_0}\|\le1-\lambda.                             \label{eq:s6v18-local-contraction}
\end{equation}
\end{lemma}

\begin{proof}
The spectral theorem places the centered spectrum of \(K\) in
\([0,1-\lambda]\), proving the norm bound.
\end{proof}

\begin{warning}[A global transfer gap would be circular]
\label{warn:s6v18-circularity}
If \(K\) is the physical infinite-volume transfer operator whose
positive spectral gap is the desired Yang--Mills mass gap, then
assuming \eqref{eq:s6v18-local-gap} assumes the target.  It may be
used here only when \(K\) is a finite-thickness conditional block
and its \(\lambda\) is proved uniformly from local conditional
convexity, a Dobrushin estimate, a finite verified card plus stable
interfaces, or another input logically independent of the
continuum Hamiltonian gap.
\end{warning}

\begin{definition}[Local conditional boundary contraction card,
LCBCC]
\label{def:s6v18-LCBCC}
LCBCC is the construction, on one common selected
response-matched law, of nested finite-thickness shell channels
\(T_j\) such that:
\begin{enumerate}[label=\textup{(\roman*)},leftmargin=3.5em]
\item the exact Markov-blanket messages of
      \Cref{lem:s6v17-Markov} propagate by the \(T_j\);
\item their \(L^2\) endpoint norms and loop-visibility
      denominators are uniformly bounded;
\item the nonuniform gap sum
      \eqref{eq:s6v18-average-gap} holds with \(\mu>0\); and
\item the estimate is proved at finite conditional thickness,
      independently of a global Yang--Mills transfer gap.
\end{enumerate}
\end{definition}

\begin{theorem}[LCBCC proves RCDC]
\label{thm:s6v18-LCBCC}
LCBCC implies the uniform covariance decay
\eqref{eq:s6v17-decay}, with the physical shell count converted to
distance by the uniform shell thickness.  Hence it proves RCDC and
the root-cofactor pair activity row of
\Cref{prop:s6v17-RCDC}.
\end{theorem}

\begin{proof}
Apply \Cref{thm:s6v18-product,cor:s6v18-exponential} to the exact
inner and outer boundary messages.  Uniform endpoint norms supply
the constants \(C_\rho\), and uniform visibility controls the
source coefficients.  The number of separating shells is bounded
below by a fixed multiple of the plaquette distance, up to an
additive constant, so \eqref{eq:s6v18-cov-product} becomes
\eqref{eq:s6v17-decay}.  Apply
\Cref{thm:s6v17-l1,prop:s6v17-RCDC}.
\end{proof}

\begin{assessment}[Exact live bottleneck after V18]
\label{ass:s6v18-bottleneck}
RCDC is now reduced to a concrete local channel question: prove a
positive average loss of maximal correlation across the physical
conditional shells.  Fixed-radius convergence and tail subtraction
are insufficient, while a global transfer gap is circular.  The
required input is a uniform finite-thickness conditional
contraction with stable interfaces.  Higher-support COHAC remains
downstream of its multi-message analogue.
\end{assessment}

\begin{programme}[V19: finite-thickness conditional contraction]
\label{prog:s6v18-V19}
\begin{enumerate}[leftmargin=3em]
\item Write the shell channel for the compact reduced-fibre Gibbs
      law as a conditional expectation between its two boundary
      sigma fields.
\item Test a Dobrushin interdependence bound using the exact
      bounded-degree plaquette graph.
\item Separate the bare strong-coupling row, where the polymer
      expansion is available, from the continuum weak-coupling
      route.
\item On the weak-coupling route, determine whether tree gauge and
      response matching yield a finite conditional convexity
      estimate after gauge-vertical directions are removed.
\item If no uniform contraction is available, isolate an attained
      near-unit boundary-message sequence instead of assuming
      decay.
\end{enumerate}
\end{programme}

\begin{firewall}[Claim boundary after sixth-series V18]
\label{fire:s6v18-boundary}
V18 proves that maximal correlation is the norm of the conditional
boundary channel, proves exact nested-shell product contraction,
gives uniform and average-gap exponential criteria, and shows why
summable deficits do not suffice.  It proves the local spectral
criterion, records the global-gap circularity boundary, formulates
LCBCC, and proves that LCBCC implies RCDC.

V18 does not prove LCBCC, RCDC, higher-support RCAC, CMSHC, COHAC,
RFAC in full, SISC in full, the simple high-excess row, arbitrary
repeated transversal supports, full TOCC, full TCNC, cyclic TIC,
HLTC, full MCWC, OCRC, SCNGC, WPSC, GRLTC, LZTC, PLRC, WSDC,
BRSC, HWSFC, UGCGC, thermal connected WPRAC, all-arity RGSC, the
raw Wilson aggregate strong card, all-order strong MTA, WUFC, CPM,
cutoff removal, reflection positivity, Osterwalder--Schrader
reconstruction, a continuum Yang--Mills measure, or a positive
Hamiltonian gap.  The full Yang--Mills mass gap has not been
proved.
\end{firewall}

\begin{theorem}[Sixth-series V18 result ledger]
\label{thm:s6v18-results}
V18:
\begin{enumerate}[label=\textup{(V18.\arabic*)},leftmargin=6em]
\item identifies maximal correlation with a channel norm;
\item proves one-shell covariance factorization;
\item proves nested-shell product contraction;
\item proves uniform and nonuniform exponential criteria;
\item proves a positive-product no-go;
\item proves the local spectral contraction criterion;
\item establishes the global-gap circularity firewall;
\item formulates LCBCC; and
\item proves LCBCC sufficient for RCDC.
\end{enumerate}
\end{theorem}

\begin{proof}
Use
\Cref{lem:s6v18-maxcorr,
thm:s6v18-one-shell,
thm:s6v18-product,
cor:s6v18-exponential,
cor:s6v18-product-no-go,
lem:s6v18-local-gap,
thm:s6v18-LCBCC}.
\end{proof}

\paragraph{Parent-version record.}
V18 was formed from
\texttt{Renewal\_Yang\_Mills\_Mass\_Gap\_Research\_Notes\_V17.tex}.
The V17 parent had \(7\,537\) logical source lines, \(265\,685\)
bytes, and SHA--256
\[
 \texttt{fe2e855a9232d860406e4c36383a0bcfe2ea058ff9139396a8ee72af9676c406}.
\]
The next compulsory companion audit is V20.

\clearpage
\chapter{Sixth-series V19: a noncircular Dobrushin row for the
compact Wilson plaquette law}
\label{ch:s6v19-Dobrushin}

\section{One-link conditional influence}
\label{sec:s6v19-influence}

No companion audit is due at V19.

Let \(\Lambda\) be a finite four-dimensional hypercubic torus or
box, let \(E_\Lambda\) and \(P_\Lambda\) be its oriented links and
plaquettes, and let \(dH\) be normalized Haar measure on
\(\mathrm{SU}(3)\).  Consider the Wilson law
\begin{equation}
 d\mu_{\Lambda,\beta}(U)
 =
 Z_{\Lambda,\beta}^{-1}
 \exp\left\{
 {\beta\over3}\sum_{p\in P_\Lambda}\Re\operatorname{Tr}U_p
 \right\}
 \prod_{e\in E_\Lambda}dH(U_e),
 \qquad \beta\ge0.                                       \label{eq:s6v19-Wilson-law}
\end{equation}
Boundary links may be frozen; none of the estimates below depends
on their values.

\begin{lemma}[Oscillation-to-total-variation bound]
\label{lem:s6v19-osc-TV}
Let
\[
 d\nu_h={e^h\,d\nu\over\int e^h\,d\nu},
 \qquad
 d\nu_k={e^k\,d\nu\over\int e^k\,d\nu}.
\]
If
\(\operatorname{osc}(h-k)\le R\), then
\begin{equation}
 \|\nu_h-\nu_k\|_{\rm TV}
 \le\tanh(R/4).                                          \label{eq:s6v19-osc-TV}
\end{equation}
\end{lemma}

\begin{proof}
Let \(r=d\nu_h/d\nu_k\).  Then \(\int r\,d\nu_k=1\) and
\((\sup r)/(\inf r)\le e^R=:M\).  For fixed endpoints
\(a\le r\le b\), convexity of \(|r-1|\) shows that its expectation
with fixed mean one is maximized by a two-point variable at
\(a,b\).  Maximizing the resulting expression over
\(b/a\le M\) gives \(a=M^{-1/2}\), \(b=M^{1/2}\), and
\[
 {1\over2}\int|r-1|\,d\nu_k
 \le{\sqrt M-1\over\sqrt M+1}
 =\tanh(R/4).
\]
\end{proof}

For distinct links \(e,f\), let \(h_{ef}\) be the number of
plaquettes containing both.  Let \(c_{ef}\) be the supremum of the
total-variation distance between the one-link conditional laws at
\(e\) over two exterior configurations which differ only at \(f\).

\begin{lemma}[Explicit Wilson interdependence matrix]
\label{lem:s6v19-interdependence}
For the law \eqref{eq:s6v19-Wilson-law},
\begin{equation}
 c_{ef}\le\tanh(\beta h_{ef}),\qquad
 \sup_e\sum_{f\ne e}c_{ef}\le18\tanh\beta.               \label{eq:s6v19-interdependence}
\end{equation}
\end{lemma}

\begin{proof}
Only plaquettes containing both \(e\) and \(f\) change the
conditional density at \(e\).  One normalized plaquette term
\(\frac13\Re\operatorname{Tr}U_p\) lies in \([-1,1]\).
For two exterior configurations its difference, as a function of
\(U_e\), lies in \([-2,2]\), so its oscillation is at most \(4\).
Thus the conditional log-density difference has oscillation at
most \(4\beta h_{ef}\).  Apply
\Cref{lem:s6v19-osc-TV}.

Every link belongs to six plaquettes in four dimensions, and each
such plaquette has three other link slots.  Counting with
multiplicity,
\[
 \sum_{f\ne e}h_{ef}=18.
\]
Since \(\tanh(kx)\le k\tanh x\) for integers \(k\ge1\),
summing the first bound gives the second.
\end{proof}

\begin{definition}[Bare Dobrushin threshold]
\label{def:s6v19-threshold}
Put
\begin{equation}
 \beta_{\rm Dob}:=\operatorname{arctanh}(1/18).           \label{eq:s6v19-threshold}
\end{equation}
For \(0\le\beta<\beta_{\rm Dob}\), set
\[
 c_\beta:=18\tanh\beta<1.
\]
\end{definition}

\section{Volume-uniform covariance decay}
\label{sec:s6v19-decay}

For a bounded function \(F\) on link configurations, define its
single-link oscillation
\[
 \delta_e(F)
 :=
 \sup\{|F(U)-F(U')|:U_g=U'_g\text{ for }g\ne e\}.
\]
Let \(\operatorname{dist}_E\) be graph distance in the link
interaction graph: two links are adjacent when they share a
plaquette.

\begin{lemma}[Dobrushin comparison bound]
\label{lem:s6v19-comparison}
Let \(C=(c_{ef})\) be the interdependence matrix of a finite Gibbs
specification and suppose
\(\|C\|_{\ell^\infty\to\ell^\infty}\le c<1\).  Then, with
\[
 D:=(I-C)^{-1}=\sum_{n\ge0}C^n,
\]
bounded functions \(F,G\) satisfy
\begin{equation}
 |\operatorname{Cov}_\mu(F,G)|
 \le
 \sum_{e,f}\delta_e(F)D_{ef}\delta_f(G).                 \label{eq:s6v19-comparison}
\end{equation}
\end{lemma}

\begin{proof}
Expose the sites in a fixed order and write each centered function
as the sum of its martingale differences.  Couple two conditional
single-site updates maximally.  A discrepancy at \(f\) creates a
discrepancy at \(e\) with probability at most \(c_{ef}\).
After \(n\) propagated updates, the oscillation vector is therefore
bounded componentwise by \(C^n\delta(G)\).  Summing the geometric
disagreement series bounds the response of the conditional
expectation of \(G\) by \(D\delta(G)\).  Pairing with the
martingale differences of \(F\) and using their oscillation bounds
gives \eqref{eq:s6v19-comparison}.  The series converges because
\(\|C\|_\infty<1\).
\end{proof}

\begin{theorem}[Explicit strong-coupling covariance decay]
\label{thm:s6v19-decay}
Let \(0\le\beta<\beta_{\rm Dob}\).  If \(F,G\) have link supports
\(S_F,S_G\), and
\[
 r:=\operatorname{dist}_E(S_F,S_G),
\]
then
\begin{equation}
 |\operatorname{Cov}_{\mu_{\Lambda,\beta}}(F,G)|
 \le
 {c_\beta^r\over1-c_\beta}
 \left(\sum_{e\in S_F}\delta_e(F)\right)
 \max_{f\in S_G}\delta_f(G).                             \label{eq:s6v19-decay}
\end{equation}
The estimate is uniform in the volume and frozen boundary links.
\end{theorem}

\begin{proof}
The matrix \(C\) is supported on distance-one link pairs.  Hence
\((C^n)_{ef}=0\) for \(n<\operatorname{dist}_E(e,f)\), while
\(\sup_e\sum_f(C^n)_{ef}\le c_\beta^n\).  Therefore, for
\(e\in S_F\),
\[
 \sum_{f\in S_G}D_{ef}
 \le\sum_{n\ge r}c_\beta^n
 ={c_\beta^r\over1-c_\beta}.
\]
Insert this in \Cref{lem:s6v19-comparison}.
\end{proof}

\begin{corollary}[Bare Wilson loop--plaquette decay]
\label{cor:s6v19-loop-plaquette}
Let \({\cal O}\) be a bounded Wilson loop writer supported on
\(\ell({\cal O})\) links and let \(p_p\) be a normalized plaquette
writer.  After harmless normalization so that both have absolute
value at most one,
\begin{equation}
 |\operatorname{Cov}({\cal O},p_p)|
 \le
 {4\ell({\cal O})\over1-c_\beta}
 \exp\{-\mu_\beta
 \operatorname{dist}_E(\operatorname{supp}{\cal O},
                         \operatorname{supp}p_p)\},
 \quad
 \mu_\beta:=-\log c_\beta>0.                             \label{eq:s6v19-loop-plaquette}
\end{equation}
\end{corollary}

\begin{proof}
A function bounded by one has single-site oscillation at most two.
Thus the first oscillation sum is at most
\(2\ell({\cal O})\) and the second maximum at most two.  Apply
\Cref{thm:s6v19-decay}.
\end{proof}

\section{Consequences and exact scope}
\label{sec:s6v19-scope}

\begin{theorem}[RCDC decay row at bare strong coupling]
\label{thm:s6v19-RCDC-row}
On \(0\le\beta<\beta_{\rm Dob}\), the covariance-decay inequality
required in \eqref{eq:s6v17-decay} holds for every fixed finite
family of bounded Wilson loop writers, uniformly in volume and
boundary links.  If their lattice support lengths and the
root-cofactor visibility coefficients are uniformly bounded, then
RCDC and the root-cofactor pair activity row follow.
\end{theorem}

\begin{proof}
Use \Cref{cor:s6v19-loop-plaquette} for each writer.  Uniform
support length bounds the constants \(C_\rho\); the remaining RCDC
hypotheses are the stated visibility bounds and strict
same-history construction.  Apply
\Cref{thm:s6v17-l1,prop:s6v17-RCDC}.
\end{proof}

\begin{corollary}[A noncircular instance of LCBCC]
\label{cor:s6v19-noncircular}
For fixed-support observables in the bare strong-coupling region,
the exponential loss is derived entirely from the finite
one-link conditional laws and the local plaquette incidence count.
It does not assume a global transfer Hamiltonian gap.
\end{corollary}

\begin{proof}
The contraction constant is the explicit
\(c_\beta=18\tanh\beta\) from
\Cref{lem:s6v19-interdependence}; no infinite-volume spectral
quantity enters its proof.
\end{proof}

\begin{warning}[Growing physical loops are not uniform lattice
loops]
\label{warn:s6v19-growing-loop}
A loop of fixed physical size generally uses
\(\ell({\cal O})\asymp a^{-1}\) lattice links as the cutoff spacing
\(a\downarrow0\).  The prefactor in
\eqref{eq:s6v19-loop-plaquette} then grows, and the normalized
Creutz visibility may shrink.  Consequently
\Cref{thm:s6v19-RCDC-row} does not provide the cofinal RCDC
constants required for continuum Yang--Mills.
\end{warning}

\begin{warning}[The bare coupling is on the wrong continuum end]
\label{warn:s6v19-coupling}
The continuum asymptotically free route has large bare Wilson
\(\beta\), whereas \(\beta<\beta_{\rm Dob}\) is a strong-coupling
region.  The theorem is a rigorous mass-generating lattice row and
a proof that LCBCC can be noncircular; it is not ultraviolet cutoff
removal.
\end{warning}

\begin{assessment}[Exact live bottleneck after V19]
\label{ass:s6v19-bottleneck}
The local conditional-contraction mechanism is now proved with an
explicit constant for the compact Wilson law at bare strong
coupling.  Two obstacles remain on the physical continuum route:
the coupling leaves the Dobrushin region, and the selected physical
loop writers grow in lattice units with possibly collapsing
visibility.  The next upstream question is whether response
matching and block integration produce a renormalized conditional
influence matrix with norm below one on finite physical-thickness
shells, while keeping the source normalization.
\end{assessment}

\begin{programme}[V20: companion audit and renormalized
influence]
\label{prog:s6v19-V20}
\begin{enumerate}[leftmargin=3em]
\item Perform the compulsory audit of the then-current SM and NS
      active notes.
\item Define the response-matched block conditional law and its
      interdependence matrix on gauge-reduced variables.
\item Derive an oscillation bound separating local deterministic
      action, connected score covariance, and generated tails.
\item Determine which term prevents a uniform row sum below one.
\item Use the companion audits to decide between stratified
      interface control, retained nonlocal cancellation, or an
      attained near-unit influence witness.
\end{enumerate}
\end{programme}

\begin{firewall}[Claim boundary after sixth-series V19]
\label{fire:s6v19-boundary}
V19 proves the oscillation-to-total-variation lemma, derives the
explicit Wilson interdependence bound \(18\tanh\beta\), proves a
volume-uniform Dobrushin covariance estimate, and obtains
loop--plaquette exponential decay for
\(\beta<\operatorname{arctanh}(1/18)\).  It conditionally closes
the fixed-support root-cofactor pair row and proves that this
instance is noncircular.

V19 does not control growing physical-loop visibility, the
weak-coupling continuum route, a response-matched block influence
matrix, full RCDC, higher-support RCAC, CMSHC, COHAC, RFAC in full,
SISC in full, the simple high-excess row, arbitrary repeated
transversal supports, full TOCC, full TCNC, cyclic TIC, HLTC, full
MCWC, OCRC, SCNGC, WPSC, GRLTC, LZTC, PLRC, WSDC, BRSC, HWSFC,
UGCGC, thermal connected WPRAC, all-arity RGSC, the raw Wilson
aggregate strong card, all-order strong MTA, WUFC, CPM, cutoff
removal, reflection positivity, Osterwalder--Schrader
reconstruction, a continuum Yang--Mills measure, or a positive
Hamiltonian gap.  The full Yang--Mills mass gap has not been
proved.
\end{firewall}

\begin{theorem}[Sixth-series V19 result ledger]
\label{thm:s6v19-results}
V19:
\begin{enumerate}[label=\textup{(V19.\arabic*)},leftmargin=6em]
\item proves oscillation control of total variation;
\item computes the one-link Wilson influence matrix;
\item proves the row bound \(18\tanh\beta\);
\item identifies an explicit Dobrushin threshold;
\item proves the finite-volume comparison estimate;
\item proves uniform exponential covariance decay;
\item applies it to loop--plaquette covariances;
\item closes the fixed-support bare strong-coupling RCDC row; and
\item isolates growing visibility and weak coupling as distinct
      cofinal obstructions.
\end{enumerate}
\end{theorem}

\begin{proof}
Use
\Cref{lem:s6v19-osc-TV,
lem:s6v19-interdependence,
lem:s6v19-comparison,
thm:s6v19-decay,
cor:s6v19-loop-plaquette,
thm:s6v19-RCDC-row}.
\end{proof}

\paragraph{Parent-version record.}
V19 was formed from
\texttt{Renewal\_Yang\_Mills\_Mass\_Gap\_Research\_Notes\_V18.tex}.
The V18 parent had \(7\,857\) logical source lines, \(276\,967\)
bytes, and SHA--256
\[
 \texttt{4b480023dbf06452901e498154c2a11f9468412c2bcb71888ccc37c2e1bbe9aa}.
\]
The next compulsory companion audit is V20.

\clearpage
\chapter{Sixth-series V20: companion audit and whole-response
conditional influence}
\label{ch:s6v20-whole-influence}

\section{Compulsory V20 companion audit}
\label{sec:s6v20-audit}

\begin{audit}[Standard Model--Einstein V41]
\label{audit:s6v20-SM}
The inspected active file was
\[
\texttt{Renewal\_SM\_Einstein\_Continuum\_Research\_Notes\_V41.tex},
\]
with \(668\,767\) bytes, \(14\,883\) logical source lines, and
SHA--256
\[
\texttt{c54883e08fdc32e403ebc2dd79785694c4844ba55ee129d7473f50f852313e25}.
\]
Its V37--V41 chain proves an adiabatic two-tail gluing theorem,
constructs a universal gapped bulk path, reduces the interface
inverse to finite polynomial certificates, and validates the
orientation-independent incidence assembly.  It explicitly refuses
to assert the remaining face determinant before the concrete
incidence data are populated.
\end{audit}

\begin{audit}[Navier--Stokes V68]
\label{audit:s6v20-NS}
The inspected active file was
\[
\texttt{Renewal\_NS\_Stability\_Research\_Notes\_V68.tex},
\]
with \(680\,016\) bytes, \(11\,994\) logical source lines, and
SHA--256
\[
\texttt{ba094ef2ec4914c43f73beafa9ed94736ae6d9f8e2ed2321fbed8a66113c4c36}.
\]
Its current spread-anchor gate contains a Lamb pairing and a
nonlocal pressure work, both of order one.  Available bounds control
their absolute sizes but neither sign; the exact test flows with
vanishing nonlinearity are gate-neutral.  The programme therefore
does not infer the required anti-alignment from separate estimates.
\end{audit}

\begin{theorem}[V20 audit decision]
\label{thm:s6v20-audit-decision}
For the response-matched Yang--Mills block law:
\begin{enumerate}[label=\textup{(\roman*)},leftmargin=3.5em]
\item the block/link incidence and gauge reduction must be
      validated before an influence matrix is asserted; and
\item when the bare and response-generated terms are individually
      large, the conditional likelihood ratio must be assembled
      before oscillation or total variation is taken.
\end{enumerate}
Separate absolute row sums may be used only where they leave a
strict margin below one.
\end{theorem}

\begin{proof}
Item (i) is the assembly-before-determinant rule in
\Cref{audit:s6v20-SM}.  Item (ii) is the order-one
orientation/cancellation lesson in \Cref{audit:s6v20-NS}.  No SM
spectral constant or NS pressure sign is imported into Yang--Mills.
\end{proof}

\section{Exact influence perturbation calculus}
\label{sec:s6v20-perturbation}

Let a compact product specification have conditional log densities
\(h_e^0(u;\eta)\).  Add a response-matching interaction
\(\Psi\), so that
\[
 h_e(u;\eta)=h_e^0(u;\eta)+\Psi_e(u;\eta).
\]
For \(e\ne f\), define the mixed conditional oscillations
\begin{align}
 R^0_{ef}
 &:=
 \sup_{\eta\sim_f\eta'}
 \operatorname{osc}_u
 \{h_e^0(u;\eta)-h_e^0(u;\eta')\},                        \label{eq:s6v20-R0}\\
 R^\Psi_{ef}
 &:=
 \sup_{\eta\sim_f\eta'}
 \operatorname{osc}_u
 \{\Psi_e(u;\eta)-\Psi_e(u;\eta')\},                      \label{eq:s6v20-Rpsi}
\end{align}
where \(\eta\sim_f\eta'\) means equality away from \(f\).

\begin{theorem}[Response perturbation of conditional influence]
\label{thm:s6v20-perturbation}
The complete one-site influence satisfies
\begin{equation}
 c_{ef}
 \le
 \tanh\left({R^0_{ef}+R^\Psi_{ef}\over4}\right)
 \le
 \tanh(R^0_{ef}/4)+{R^\Psi_{ef}\over4}.                  \label{eq:s6v20-perturbation}
\end{equation}
Consequently, if
\begin{equation}
 c_0:=
 \sup_e\sum_{f\ne e}\tanh(R^0_{ef}/4)<1,\qquad
 \varepsilon_\Psi:=
 {1\over4}\sup_e\sum_{f\ne e}R^\Psi_{ef}<1-c_0,           \label{eq:s6v20-margin}
\end{equation}
then the full specification has Dobrushin row norm at most
\(c_0+\varepsilon_\Psi<1\).
\end{theorem}

\begin{proof}
The complete conditional log-likelihood difference has oscillation
at most \(R^0_{ef}+R^\Psi_{ef}\).  Apply
\Cref{lem:s6v19-osc-TV}.  The second inequality uses
\(\tanh(a+b)\le\tanh a+\tanh b\le\tanh a+b\) for
\(a,b\ge0\).  Sum over \(f\).
\end{proof}

\begin{corollary}[Local head and generated-tail ledger]
\label{cor:s6v20-head-tail}
If
\[
 R^\Psi_{ef}
 \le R^{\rm head}_{ef}+R^{\rm tail}_{ef}
\]
and the two row sums contribute
\(\varepsilon_{\rm head}\) and
\(\varepsilon_{\rm tail}\), then Dobrushin contraction follows
whenever
\begin{equation}
 c_0+\varepsilon_{\rm head}+\varepsilon_{\rm tail}<1.    \label{eq:s6v20-head-tail}
\end{equation}
This ledger is valid only after the exact response-matched
conditional density and its incidence have been assembled.
\end{corollary}

\begin{proof}
Insert the decomposition into
\eqref{eq:s6v20-Rpsi} and apply
\Cref{thm:s6v20-perturbation}.
\end{proof}

\section{Why perturbation around the bare Wilson law stops}
\label{sec:s6v20-no-go}

\begin{proposition}[The V19 reference loses its strict margin]
\label{prop:s6v20-bare-no-go}
For the bare Wilson reference used in V19,
\[
 c_0\le18\tanh\beta
\]
is useful only for
\(\beta<\beta_{\rm Dob}\).  When
\(\beta\ge\beta_{\rm Dob}\), the perturbative condition
\eqref{eq:s6v20-margin} cannot be verified from this upper bound
for any nonnegative response debit
\(\varepsilon_\Psi\).
\end{proposition}

\begin{proof}
At \(\beta\ge\beta_{\rm Dob}\),
\(18\tanh\beta\ge1\).  Thus the available reference estimate has
no positive residual margin \(1-c_0\).  Adding nonnegative
oscillation debits cannot restore one.
\end{proof}

\begin{warning}[This is a bound failure, not a mixing
counterexample]
\label{warn:s6v20-not-counterexample}
\Cref{prop:s6v20-bare-no-go} does not prove that the actual
weak-coupling block influence is at least one.  The estimate
\(18\tanh\beta\) separately maximizes incompatible plaquette
variations and discards gauge, curvature, and response-matching
cancellations.  It proves only that perturbing this scalar upper
bound cannot reach the continuum route.
\end{warning}

\begin{lemma}[Whole-response oscillation can be smaller than the
sum of its parts]
\label{lem:s6v20-cancellation}
For functions \(a,b\) on a common conditional fibre,
\begin{equation}
 \operatorname{osc}(a+b)
 \le\operatorname{osc}(a)+\operatorname{osc}(b),          \label{eq:s6v20-triangle}
\end{equation}
and the ratio between the right and left sides is unbounded.
\end{lemma}

\begin{proof}
The inequality is elementary.  Taking \(b=-a\) gives zero on the
left and \(2\operatorname{osc}(a)\) on the right.
\end{proof}

\begin{assessment}[The required weak-coupling object]
\label{ass:s6v20-whole}
At large bare \(\beta\), a useful influence estimate must be made
on the complete gauge-reduced response-matched conditional
log-density difference.  Bounding its deterministic curvature,
score response, and generated tail separately can erase the only
possible margin, exactly as separate order-one bounds erase the NS
spread-anchor orientation.
\end{assessment}

\section{Whole conditional influence card}
\label{sec:s6v20-WCIC}

\begin{definition}[Validated block-incidence card, VBIC]
\label{def:s6v20-VBIC}
VBIC consists of an orientation-independent global record of:
\begin{enumerate}[label=\textup{(\roman*)},leftmargin=3.5em]
\item the gauge-reduced block variables and their Haar/Jacobian
      reference measure;
\item every local action and response-generated interaction
      incident to a block;
\item exact multiplicities across bulk, face, edge, and vertex
      strata;
\item adjoint, gauge-vertical, constant-mode, and row-sum
      validation identities; and
\item strict compatibility with the source and shell transports.
\end{enumerate}
\end{definition}

\begin{definition}[Whole conditional influence card, WCIC]
\label{def:s6v20-WCIC}
Given VBIC, let \(h_{r,e}^{\rm full}\) be the complete
response-matched conditional log density on shell \(r\).  Define
\begin{equation}
 R^{\rm full}_{r,ef}
 :=
 \sup_{\eta\sim_f\eta'}
 \operatorname{osc}_u
 \{h_{r,e}^{\rm full}(u;\eta)
   -h_{r,e}^{\rm full}(u;\eta')\}.                        \label{eq:s6v20-Rfull}
\end{equation}
WCIC is the uniform strict inequality
\begin{equation}
 \sup_r\sup_e
 \sum_{f\ne e}
 \tanh(R^{\rm full}_{r,ef}/4)
 \le1-\epsilon_{\rm CI}                                  \label{eq:s6v20-WCIC}
\end{equation}
for some \(\epsilon_{\rm CI}>0\), including finite-collar and
stratified-interface blocks.
\end{definition}

\begin{theorem}[WCIC gives noncircular shell contraction]
\label{thm:s6v20-WCIC}
WCIC implies volume-uniform exponential covariance decay on every
response-matched shell, with constants depending only on
\(\epsilon_{\rm CI}\) and the uniform interaction range.  If the
source visibility and physical support normalization are uniform,
then WCIC proves LCBCC, RCDC, and the root-cofactor pair activity
row.
\end{theorem}

\begin{proof}
\Cref{lem:s6v19-osc-TV} turns
\eqref{eq:s6v20-Rfull} into the actual Dobrushin interdependence
matrix.  Equation \eqref{eq:s6v20-WCIC} gives row norm at most
\(1-\epsilon_{\rm CI}\).  The path expansion in
\Cref{thm:s6v19-decay} gives exponential spatial covariance decay,
uniformly in volume and boundary conditions.  Because the estimate
comes from finite conditional laws rather than a global transfer
Hamiltonian, it satisfies the circularity firewall
\Cref{warn:s6v18-circularity}.  Uniform visibility and support
then give RCDC by \Cref{thm:s6v18-LCBCC} and
\Cref{prop:s6v17-RCDC}.
\end{proof}

\begin{proposition}[Finite validated failure is an attained
influence witness]
\label{prop:s6v20-witness}
After VBIC, if a finite shell violates
\eqref{eq:s6v20-WCIC}, compactness of the reduced conditional
fibres implies that, for every \(\varepsilon>0\), there are a block
\(e\), finitely many neighboring blocks \(f\), exterior
configurations, and fibre points which realize the offending row
sum within \(\varepsilon\).
\end{proposition}

\begin{proof}
Every conditional log density is continuous on a finite product of
compact group fibres.  Each mixed oscillation supremum is therefore
attained.  Only finitely many \(f\) occur at finite interaction
range.  Choose the maximizers and sum them; if the supremum over
shell blocks is not attained because the shell grows, choose a
block within \(\varepsilon\) of it.
\end{proof}

\begin{assessment}[Exact live bottleneck after V20]
\label{ass:s6v20-bottleneck}
The bare strong-coupling row is rigorous, and the exact perturbation
ledger is rigorous wherever it retains a margin.  It cannot reach
large \(\beta\) from the scalar Wilson bound.  The live upstream
card is now VBIC plus WCIC for the complete gauge-reduced,
response-matched block conditional law.  This is a finite
conditional calculation at each stratum, but no attached
manuscript currently supplies its assembled mixed oscillations or
a uniform margin.
\end{assessment}

\begin{programme}[V21: construct the response-matched block
conditional]
\label{prog:s6v20-V21}
\begin{enumerate}[leftmargin=3em]
\item Extract the exact response-matched action decomposition and
      relative tree-gauge measure from the attached YM sources.
\item Write one block's full conditional density, including the
      generated plaquette marginal correction.
\item Prove the incidence multiplicities and gauge-vertical
      cancellations required by VBIC.
\item Compute the complete mixed oscillation before using triangle
      inequalities.
\item Return a strict finite conditional margin on a justified
      parameter row or an attained near-unit influence witness.
\end{enumerate}
\end{programme}

\begin{firewall}[Claim boundary after sixth-series V20]
\label{fire:s6v20-boundary}
V20 performs the compulsory SM V41 and NS V68 audit, proves the
exact conditional-influence perturbation theorem, and isolates why
the bare Wilson scalar estimate cannot be perturbed to weak
coupling.  It formulates the validated block-incidence card VBIC
and the whole conditional influence card WCIC, proves WCIC
sufficient for the noncircular RCDC route, and proves that finite
failure has an attained witness.

V20 does not prove VBIC or WCIC for the actual response-matched
block law, uniform growing-source visibility, full RCDC,
higher-support RCAC, CMSHC, COHAC, RFAC in full, SISC in full, the
simple high-excess row, arbitrary repeated transversal supports,
full TOCC, full TCNC, cyclic TIC, HLTC, full MCWC, OCRC, SCNGC,
WPSC, GRLTC, LZTC, PLRC, WSDC, BRSC, HWSFC, UGCGC, thermal
connected WPRAC, all-arity RGSC, the raw Wilson aggregate strong
card, all-order strong MTA, WUFC, CPM, cutoff removal, reflection
positivity, Osterwalder--Schrader reconstruction, a continuum
Yang--Mills measure, or a positive Hamiltonian gap.  The full
Yang--Mills mass gap has not been proved.
\end{firewall}

\begin{theorem}[Sixth-series V20 result ledger]
\label{thm:s6v20-results}
V20:
\begin{enumerate}[label=\textup{(V20.\arabic*)},leftmargin=6em]
\item audits SM V41 and NS V68;
\item derives the exact response perturbation bound;
\item gives the local-head/generated-tail margin ledger;
\item proves the bare-reference perturbative no-go;
\item distinguishes that bound failure from a mixing
      counterexample;
\item proves that separate oscillations can lose arbitrary
      cancellation;
\item formulates VBIC and WCIC;
\item proves WCIC sufficient for noncircular RCDC; and
\item proves finite failure has an attained witness.
\end{enumerate}
\end{theorem}

\begin{proof}
Use
\Cref{audit:s6v20-SM,
audit:s6v20-NS,
thm:s6v20-audit-decision,
thm:s6v20-perturbation,
prop:s6v20-bare-no-go,
lem:s6v20-cancellation,
thm:s6v20-WCIC,
prop:s6v20-witness}.
\end{proof}

\paragraph{Parent-version record.}
V20 was formed from
\texttt{Renewal\_Yang\_Mills\_Mass\_Gap\_Research\_Notes\_V19.tex}.
The V19 parent had \(8\,208\) logical source lines, \(289\,381\)
bytes, and SHA--256
\[
 \texttt{b74dab144e3d275c230c4de1eeb2cd96d903c948895f08882e97b921d286ac37}.
\]
The next compulsory companion audit is V25.

\clearpage
\chapter{Sixth-series V21: exact tree-gauge disintegration and
the upstream static-card gate}
\label{ch:s6v21-disintegration}

\section{Audit of the exact action which is actually available}
\label{sec:s6v21-action-audit}

\begin{audit}[Attached reduced-fibre programme through V35]
\label{audit:s6v21-YM35}
The inspected source was
\[
\path{Renewal_Yang_Mills_Reduced_Fibre_Wilson_Shell_Symbol_}
\]
\[
\path{and_Local_Vacuum_Programme_V35.tex}.
\]
Four facts in that source control the present continuation.
\begin{enumerate}[label=\textup{(\roman*)},leftmargin=3.5em]
\item Its Version~V2 fixes the relative parity-tree section
\[
 U_e(z,W)=
 \begin{cases}
 I,&e\in{\mathscr T}_N,\\
 W_{y,\mu},
   &e=(2y+\widehat e_\mu,2y+2\widehat e_\mu),\\
 z_e,&e\in{\mathscr E}^{\rm red}_N,
 \end{cases}
\]
and integrates the \(45M^4\) reduced links against normalized
product Haar measure.
\item The corresponding displayed action is
\[
 {\cal S}(z,W)=\beta V_M(z,W)+\Psi(z,W),
 \qquad
 V_M=\sum_p\left(1-\frac13\Re\Tr U_p\right).
\]
The symbol \(\Psi\) is declared to contain every generated or
marginally retuned interaction which is actually present.  The
source does not give an interaction list or a pointwise writer for
that symbol.
\item Its Version~V32 corrects the target order: the static shell
output is the exact trace effective action obtained by integrating
the complete fine action at fixed retained field.  The natural
coarse comparator and its relative action cocycle are not
prerequisites for differentiating this trace action.
\item Its Version~V35 corrects a possible double debit.  An
unsplit complete V32 pair card already contains every finite-shell
action and differentiated-score term.  A separate tail is added
only after one deliberately replaces that complete card by a
declared head approximation.
\end{enumerate}
\end{audit}

\begin{theorem}[Correct target of the conditional calculation]
\label{thm:s6v21-target-correction}
Let \(\widehat\beta_{2M}\) be the selected canonical fine root and
let \({\cal S}_{M,\widehat\beta_{2M}}(z;W)\) be the complete
fine action on the reduced detail fibre.  The law relevant to the
static shell symbol is
\begin{equation}
 \dd\nu_{M,W}(z)
 =
 {e^{-{\cal S}_{M,\widehat\beta_{2M}}(z;W)}
  \dd m_M(z)
  \over
  \displaystyle
  \int e^{-{\cal S}_{M,\widehat\beta_{2M}}(z;W)}
  \dd m_M(z)}.                                           \label{eq:s6v21-trace-law}
\end{equation}
Consequently:
\begin{enumerate}[label=\textup{(\alph*)},leftmargin=3.5em]
\item a WCIC intended to support the trace-shell mass route must
      be proved for \eqref{eq:s6v21-trace-law};
\item no natural-coarse likelihood cocycle is required for this
      purpose; and
\item every term already present in the complete action must occur
      once in the unsplit conditional density.
\end{enumerate}
\end{theorem}

\begin{proof}
The exact trace partition is
\[
 Z_M^{\rm tr}(W)=
 \int e^{-{\cal S}_{M,\widehat\beta_{2M}}(z;W)}
 \dd m_M(z).
\]
The probability law used when its logarithm is differentiated is
therefore exactly \eqref{eq:s6v21-trace-law}.  A natural-coarse
comparison changes neither this identity nor the conditional laws
of its integrand.  Finally, a conditional density is obtained by
disintegrating the same integrand.  Inserting a second copy of a
term already in \({\cal S}\) changes the law, so it cannot be a
tail correction to an unsplit calculation.
\end{proof}

\section{Exact conditional law in the relative tree gauge}
\label{sec:s6v21-exact-conditional}

Let \(G\) be a compact group with normalized Haar probability
\(\lambda_G\), let \(E\) be a finite set of reduced links, and set
\[
 m_E:=\lambda_G^{\otimes E}.
\]
For a nonempty \(A\subseteq E\), let
\(\Phi_A:G^A\to\mathbb R\) be continuous.  Put
\begin{equation}
 {\cal S}(z)=\sum_{A\in{\cal I}}\Phi_A(z_A),
 \qquad
 \dd\nu(z)=Z^{-1}e^{-{\cal S}(z)}\dd m_E(z),              \label{eq:s6v21-hyperaction}
\end{equation}
where \({\cal I}\) is a finite interaction hypergraph.

\begin{theorem}[Exact incident-star disintegration]
\label{thm:s6v21-star-disintegration}
For \(e\in E\) and exterior configuration
\(\eta\in G^{E\setminus\{e\}}\), define
\begin{equation}
 {\cal S}_e(u\mid\eta)
 :=
 \sum_{\substack{A\in{\cal I}\\e\in A}}
 \Phi_A\bigl((u,\eta)_A\bigr).                            \label{eq:s6v21-local-energy}
\end{equation}
A regular conditional distribution at \(e\) is
\begin{equation}
 \boxed{
 \gamma_e(\dd u\mid\eta)
 =
 {e^{-{\cal S}_e(u\mid\eta)}\lambda_G(\dd u)
  \over
  \displaystyle\int_Ge^{-{\cal S}_e(v\mid\eta)}
  \lambda_G(\dd v)}.}                                    \label{eq:s6v21-one-link-law}
\end{equation}
If \(e\ne f\) and \(\eta,\eta'\) agree away from \(f\), then
\begin{equation}
 \boxed{
 {\cal S}_e(u\mid\eta)-{\cal S}_e(u\mid\eta')
 =
 \sum_{\substack{A\in{\cal I}\\e,f\in A}}
 \left[
 \Phi_A((u,\eta)_A)-\Phi_A((u,\eta')_A)
 \right].}                                               \label{eq:s6v21-common-incidence}
\end{equation}
In particular the exact mixed conditional oscillation is
\begin{equation}
 \boxed{
 R_{ef}
 =
 \sup_{\eta\sim_f\eta'}
 \operatorname{osc}_{u\in G}
 \sum_{\substack{A\in{\cal I}\\e,f\in A}}
 \left[
 \Phi_A((u,\eta)_A)-\Phi_A((u,\eta')_A)
 \right].}                                               \label{eq:s6v21-exact-R}
\end{equation}
\end{theorem}

\begin{proof}
Split the action into the sum of interactions containing \(e\)
and those not containing \(e\).  The latter sum depends on
\(\eta\) but not on \(u\), so it cancels between numerator and
normalizing denominator in the conditional law.  This proves
\eqref{eq:s6v21-one-link-law}.  When only the \(f\)-coordinate of
the exterior is changed, every term containing \(e\) but not
\(f\) cancels in the difference.  This gives
\eqref{eq:s6v21-common-incidence}.  The difference of the two
conditional log densities has in addition a normalizing constant
independent of \(u\); oscillation removes that constant and proves
\eqref{eq:s6v21-exact-R}.
\end{proof}

\begin{corollary}[There is no hidden tree-gauge Jacobian]
\label{cor:s6v21-no-Jacobian}
For the fixed relative parity-tree section in
\Cref{audit:s6v21-YM35}, the reference conditional measure of
each reduced link is normalized Haar measure.  Tree links and
designated retained links are fixed parameters in the
disintegration.  No Faddeev--Popov determinant or additional
one-link Jacobian occurs in \eqref{eq:s6v21-one-link-law}.
\end{corollary}

\begin{proof}
The exact reduced-fibre integral in the inspected source is taken
against normalized product Haar measure on the coordinates
\((z_e)_{e\in{\mathscr E}^{\rm red}_N}\).  Disintegration of a
finite product probability measure gives Haar measure in the
selected coordinate.  Fixed coordinates create no integration
factor.  Residual simultaneous conjugation is a symmetry of this
law, not a density multiplier, because Haar measure is conjugation
invariant.
\end{proof}

\begin{corollary}[Exact pure-Wilson incidence in four dimensions]
\label{cor:s6v21-Wilson-incidence}
For the pure Wilson part of the reduced-fibre law, a reduced link
\(e\) is incident to at most six plaquette interactions.  Its
conditional law can depend on at most eighteen other reduced
links.  More precisely, only reduced links sharing a plaquette
with \(e\) occur in \eqref{eq:s6v21-exact-R}; fixed tree or
retained links can only lower these counts.
\end{corollary}

\begin{proof}
An unoriented nearest-neighbour link in four dimensions lies in
two plaquettes for each of the three transverse directions, hence
in six plaquettes.  Each such plaquette contains three other
links.  Counting with multiplicity gives eighteen candidates.
Removing fixed links and identifying candidates repeated in two
plaquettes cannot increase the number.
\end{proof}

\begin{assessment}[What part of VBIC is now exact]
\label{ass:s6v21-VBIC-progress}
The variable carrier, reference measure, Wilson plaquette
incidence, and bulk one-link multiplicities required by VBIC are
now exact in the relative tree gauge.  The remaining incidence
data are precisely the supports and pointwise writers of the
generated or marginal interaction and their boundary-stratum
multiplicities.  A symbol \(\Psi\) without those data cannot be
inserted into \eqref{eq:s6v21-exact-R}.
\end{assessment}

\section{The response-matched term must be assembled before
oscillation}
\label{sec:s6v21-response-assembly}

Suppose, in a convention which separates the terms, that
\begin{equation}
 {\cal S}
 =
 \beta V+\Psi+b^{\rm marg}P,
 \qquad
 b^{\rm marg}
 =
 -{\mathscr R_\Phi(\Psi)\over\mathscr R_\Phi(P)},         \label{eq:s6v21-marginal-action}
\end{equation}
where the last quotient is used only when
\(\mathscr R_\Phi(P)\ne0\).  Write
\(\Psi=\sum_B\Psi_B\) and \(P=\sum_QP_Q\) using their actual
interaction supports.  For an interaction writer \(X_A\), set
\[
 \Delta_fX_A(u;\eta,\eta')
 :=
 X_A((u,\eta)_A)-X_A((u,\eta')_A).
\]

\begin{theorem}[Whole response-matched mixed oscillation]
\label{thm:s6v21-whole-oscillation}
Under the explicit decomposition
\eqref{eq:s6v21-marginal-action},
\begin{equation}
 \boxed{
 \begin{split}
 R^{\rm full}_{ef}
 =\sup_{\eta\sim_f\eta'}\operatorname{osc}_u
 \bigg\{&
 \beta\sum_{p\ni e,f}\Delta_fV_p(u;\eta,\eta')\\
 &+\sum_{B\ni e,f}\Delta_f\Psi_B(u;\eta,\eta')\\
 &+b^{\rm marg}\sum_{Q\ni e,f}
       \Delta_fP_Q(u;\eta,\eta')
 \bigg\}.
 \end{split}}                                            \label{eq:s6v21-whole-R}
\end{equation}
Thus the response cofactor fixes one scalar coefficient but does
not determine the conditional influence.  The three lines in
\eqref{eq:s6v21-whole-R} must be summed pointwise before taking
oscillation whenever their cancellation is used.
\end{theorem}

\begin{proof}
Apply \Cref{thm:s6v21-star-disintegration} to the union of the
plaquette, generated, and marginal interaction hypergraphs.
Linearity gives the expression inside braces.  The cofactor
formula in \eqref{eq:s6v21-marginal-action} contains no information
about the pointwise functions
\(\Delta_f\Psi_B\) or \(\Delta_fP_Q\); it therefore cannot
determine their oscillation or their cancellation with the Wilson
line.
\end{proof}

\begin{theorem}[Exact conditional transfer by a relative action]
\label{thm:s6v21-conditional-transfer}
Let \(F,G\) be continuous actions on \(G^E\), let
\(\Delta=F-G\), and let \(\gamma_e^F,\gamma_e^G\) be their
one-link conditional laws relative to the same product Haar
measure.  Then
\begin{equation}
 \boxed{
 {\dd\gamma_e^F(\cdot\mid\eta)
  \over\dd\gamma_e^G(\cdot\mid\eta)}(u)
 =
 {e^{-\Delta(u,\eta)}
  \over
  \displaystyle
  \int e^{-\Delta(v,\eta)}
  \gamma_e^G(\dd v\mid\eta)}.}                           \label{eq:s6v21-conditional-RN}
\end{equation}
For exteriors differing only at \(f\), the additional
\(u\)-dependent log-likelihood difference is exactly
\[
 -\Delta(u,\eta)+\Delta(u,\eta').
\]
Hence an influence estimate can be transferred from \(G\) to
\(F\) only after controlling the mixed oscillation
\begin{equation}
 \sup_{\eta\sim_f\eta'}
 \operatorname{osc}_u
 \{\Delta(u,\eta)-\Delta(u,\eta')\}.                     \label{eq:s6v21-relative-mixed}
\end{equation}
\end{theorem}

\begin{proof}
Divide the two conditional densities and normalize the resulting
positive likelihood under \(\gamma_e^G\).  Comparing two
exteriors adds a ratio of normalizing constants, whose logarithm
is independent of \(u\).  It disappears under oscillation, leaving
\eqref{eq:s6v21-relative-mixed}.
\end{proof}

\begin{theorem}[Response matching does not imply WCIC]
\label{cth:s6v21-response-not-WCIC}
Equality of a macroscopic response, even together with equality
of every one-site marginal, does not imply any uniform strict
conditional-influence margin.
\end{theorem}

\begin{proof}
On \(\{-1,1\}^2\), let
\begin{equation}
 \nu_K(\sigma,\tau)
 :=
 {e^{K\sigma\tau}\over4\cosh K},
 \qquad K\ge0.                                           \label{eq:s6v21-Ising-law}
\end{equation}
Both one-site marginals are uniform for every \(K\), so the
expectation of every one-site response writer is independent of
\(K\).  On the other hand,
\[
 \nu_K(\sigma=1\mid\tau=1)
 ={1+\tanh K\over2},
 \qquad
 \nu_K(\sigma=1\mid\tau=-1)
 ={1-\tanh K\over2}.
\]
The total-variation influence of \(\tau\) on \(\sigma\) is
therefore \(\tanh K\), which tends to one.  Choosing
\(K=\operatorname{arctanh}c\) realizes every \(c\in[0,1)\) while
all one-site response data remain fixed.
\end{proof}

\begin{warning}[Logical scope of the response counterexample]
\label{warn:s6v21-response-scope}
The law \eqref{eq:s6v21-Ising-law} is not a Yang--Mills
counterexample.  It proves that response equality by itself has no
theorem-level implication for WCIC.  A Yang--Mills contraction
must use the actual action, its gauge-reduced incidence, or an
additional inequality special to that action.
\end{warning}

\section{Why the static card remains upstream of WCIC}
\label{sec:s6v21-static-first}

\begin{theorem}[Perfect conditional mixing does not fix
the effective-action sign]
\label{cth:s6v21-mixing-not-static}
There is a finite compact conditional system with Dobrushin row
norm zero whose effective action has a strictly negative Hessian
in its retained external parameter.
\end{theorem}

\begin{proof}
Take one spin \(\sigma\in\{-1,1\}\), no other conditional sites,
and the parameter-dependent action
\[
 F(\sigma,t)=-t\sigma.
\]
There are no off-diagonal influences, so the Dobrushin row norm is
zero.  With normalized counting reference,
\[
 Z(t)=\cosh t,
 \qquad
 {\cal A}(t)=-\log\cosh t,
 \qquad
 {\cal A}''(0)=-1.
\]
Thus even perfect inter-site mixing does not supply the positive
static trace-shell margin.
\end{proof}

\begin{corollary}[Upstream priority of the canonical-root pair
card]
\label{cor:s6v21-static-priority}
WCIC is a valid noncircular route to covariance decay once it is
proved for the occurring law, but it cannot replace the
canonical-root static trace-shell sign calculation.  The next
target-native Yang--Mills acquisition is therefore the complete
unsplit finite pair card of the attached V32/V35 programme.  Only
after that card passes is a weak-coupling WCIC calculation a
downstream mass-gap step on the same law.
\end{corollary}

\begin{proof}
\Cref{cth:s6v21-mixing-not-static} rules out any implication from
conditional contraction alone to the required static sign.
\Cref{thm:s6v21-target-correction} identifies the exact law and
shows that its complete unsplit card has no additional tail debit.
\end{proof}

\begin{theorem}[Strict finite VBIC/WCIC decidability after action
loading]
\label{thm:s6v21-finite-decision}
Fix a finite shell and suppose every interaction support and
writer in \eqref{eq:s6v21-hyperaction} is explicitly loaded.
Then:
\begin{enumerate}[label=\textup{(\roman*)},leftmargin=3.5em]
\item the complete conditional incidence is the finite common-star
      list in \eqref{eq:s6v21-common-incidence};
\item every \(R_{ef}\) in \eqref{eq:s6v21-exact-R} is attained;
      and
\item convergent outward enclosures of the continuous interaction
      writers terminate the strict WCIC verdict whenever the
      exact row norm is separated from one.
\end{enumerate}
\end{theorem}

\begin{proof}
The first assertion is
\Cref{thm:s6v21-star-disintegration}.  The optimization domain in
\eqref{eq:s6v21-exact-R} is a finite product of compact groups,
and the optimized function is continuous, so its maximum and
minimum are attained.  Uniform continuity implies that convergent
outward subdivisions have vanishing enclosure diameter.  Since
there are finitely many incident \(f\), a strict positive distance
between the exact row sum and one is eventually resolved.
\end{proof}

\begin{assessment}[Exact live bottleneck after V21]
\label{ass:s6v21-bottleneck}
The conditional-law architecture is no longer missing.  In the
fixed relative tree gauge its reference is product Haar, its
incidence is the exact common interaction star, and the Wilson
bulk has at most six incident plaquettes and eighteen candidate
neighbours.  The response coefficient alone cannot populate the
generated mixed oscillations, and response matching cannot imply
WCIC.

The more upstream attached V35 frontier is decisive: before
conditional contraction is used for a mass argument, one complete
canonical-root static trace-shell card must pass.  No numerical
root interval, complete unsplit pair-card value, or equivalent
four-moment enclosure is present in the supplied files.
\end{assessment}

\begin{programme}[V22: execute or bound the upstream unsplit
static card]
\label{prog:s6v21-V22}
\begin{enumerate}[leftmargin=3em]
\item Freeze the inherited physical rectangle, canonical fine
      root, and first retained nonzero momentum--colour block.
\item Use the exact trace action
      \({\cal S}_{M,\widehat\beta_{2M}}(z;W)\); do not introduce
      the natural-coarse cocycle.
\item Form the complete first and second retained-field writers
      \(J=D_W{\cal S}\) and \(K=D_W^2{\cal S}\), including every
      occurring generated and support-birth term.
\item Evaluate or rigorously enclose the unsplit two-copy writer
      whose normalized value is
      \(\mathbb EK-\operatorname{Cov}(J)\) minus the declared
      topological/reference form.
\item Add no differentiated tail.  If a head approximation is
      chosen instead, declare it first and restore exactly one
      deterministic, mixed, and tail--tail residual.
\item Return a positive static margin, an attained negative
      momentum--colour witness, or the first explicit missing
      action/root/moment interval.
\item On a positive branch, return to
      \eqref{eq:s6v21-whole-R} for WCIC on that same selected law.
\end{enumerate}
\end{programme}

\begin{firewall}[Claim boundary after sixth-series V21]
\label{fire:s6v21-boundary}
V21 audits the exact action semantics of the attached V35
programme, identifies the trace-shell law as the correct
conditional target, proves the exact tree-gauge one-link
disintegration and common-incidence cancellation, proves the
absence of a hidden tree-gauge Jacobian, certifies the pure Wilson
incidence count, derives the exact whole response-matched mixed
oscillation and relative-action conditional likelihood, and proves
two nonidentifiability countertheorems.  It also proves that WCIC,
even if available, cannot replace the unfinished static sign card.

V21 does not populate the generated interaction, a canonical-root
static pair card, a strict weak-coupling WCIC margin, uniform
growing-source visibility, full RCDC, higher-support RCAC, CMSHC,
COHAC, RFAC in full, SISC in full, the simple high-excess row,
arbitrary repeated transversal supports, full TOCC, full TCNC,
cyclic TIC, HLTC, full MCWC, OCRC, SCNGC, WPSC, GRLTC, LZTC,
PLRC, WSDC, BRSC, HWSFC, UGCGC, thermal connected WPRAC,
all-arity RGSC, the raw Wilson aggregate strong card, all-order
strong MTA, WUFC, CPM, cutoff removal, reflection positivity,
Osterwalder--Schrader reconstruction, a continuum Yang--Mills
measure, or a positive Hamiltonian gap.  The full Yang--Mills mass
gap has not been proved.
\end{firewall}

\begin{theorem}[Sixth-series V21 result ledger]
\label{thm:s6v21-results}
V21:
\begin{enumerate}[label=\textup{(V21.\arabic*)},leftmargin=6em]
\item audits the attached YM V35 action and tail semantics;
\item corrects the WCIC target to the exact trace-shell law;
\item proves exact incident-star disintegration;
\item proves common-incidence cancellation in mixed conditionals;
\item removes a spurious tree-gauge Jacobian;
\item certifies the reduced pure-Wilson incidence count;
\item derives the whole response-matched oscillation;
\item proves exact conditional likelihood transfer;
\item proves response matching does not imply WCIC;
\item proves perfect mixing does not imply a positive static
      Hessian; and
\item restores the canonical-root unsplit static card as the next
      upstream acquisition.
\end{enumerate}
\end{theorem}

\begin{proof}
Use
\Cref{audit:s6v21-YM35,
thm:s6v21-target-correction,
thm:s6v21-star-disintegration,
cor:s6v21-no-Jacobian,
cor:s6v21-Wilson-incidence,
thm:s6v21-whole-oscillation,
thm:s6v21-conditional-transfer,
cth:s6v21-response-not-WCIC,
cth:s6v21-mixing-not-static,
cor:s6v21-static-priority,
thm:s6v21-finite-decision}.
\end{proof}

\paragraph{Parent-version record.}
V21 was formed from
\texttt{Renewal\_Yang\_Mills\_Mass\_Gap\_Research\_Notes\_V20.tex}.
The V20 parent had \(8\,590\) logical source lines,
\(303\,153\) bytes, and SHA--256
\[
 \texttt{947db14da243eca91c2e3226e99b31dc6ed191b8c5303af5f1aaa344573708c6}.
\]
The next compulsory companion audit remains V25.

\clearpage
\chapter{Sixth-series V22: from whole conditional influence to
an Osterwalder--Schrader spectral gap}
\label{ch:s6v22-OS-bridge}

\section{Grand Tensor Yang--Mills trajectory audit}
\label{sec:s6v22-GT-audit}

\begin{audit}[Grand Tensor V067 Yang--Mills interface]
\label{audit:s6v22-GT}
The inspected file was
\[
 \texttt{Renewal\_Grand\_Tensor\_Monograph\_V067.tex},
\]
with \(2\,804\,854\) bytes, \(64\,475\) logical source lines, and
SHA--256
\[
 \texttt{facd058ec8bb210ad8b296c9d7f4d78a229d08d7a2ae378808700655cb8bc0fe}.
\]
The Yang--Mills interface proves four facts which change the
present mass-gap route.
\begin{enumerate}[label=\textup{(\roman*)},leftmargin=3.5em]
\item A finite Wilson transfer with strictly positive continuous
      kernel has a strict finite-card gap, but this alone gives no
      cofinal gap.
\item On a particular over-weak fixed-volume bare trajectory there
      is a centre-neutral, gauge-invariant wrapping-toron sequence
      with physical Rayleigh quotient of order
      \(K_j^{-1/3}\) and fixed-delay transfer norm tending to one.
\item The toron result obstructs a uniform gap for the full
      finite-torus sector on that bare trajectory.  It does not
      decide the infinite-volume algebra of uniformly local
      observables on a response-fixed or local-vacuum branch.
\item On such a corrected branch, a full-core gap may be promoted
      from noncollapsing finite source banks plus their strict
      delayed contraction and a collectively compact complement.
      A uniform graph--Rellich estimate is one sufficient
      complement certificate.
\end{enumerate}
\end{audit}

\begin{theorem}[Trajectory correction forced by the toron
witness]
\label{thm:s6v22-trajectory-correction}
No argument for a uniform full finite-torus Yang--Mills gap may
use the over-weak bare trajectory audited in
\Cref{audit:s6v22-GT}.  A mass-gap continuation must instead
declare one of the following before a transfer estimate is read:
\begin{enumerate}[label=\textup{(\alph*)},leftmargin=3.5em]
\item a fixed-response, adjacent-step-scaled trajectory together
      with a shell-renormalized toron sector; or
\item a projective infinite-volume local-vacuum branch on the
      uniformly local, centre-neutral gauge-invariant algebra,
      with wrapping toron writers excluded from that local target.
\end{enumerate}
\end{theorem}

\begin{proof}
On the audited bare trajectory there are normalized centred vectors
\(v_j\) in the full finite-torus gauge-invariant sector such that
\[
 \langle v_j,H_j^{\rm phys}v_j\rangle\longrightarrow0.
\]
The variational principle therefore forces the infimum of the
fixed-free spectrum to tend to zero.  The witnesses wrap a spatial
cycle, so they do not belong to any fixed uniformly local
observable carrier as the volume grows.  They consequently exclude
the first target but do not exclude the second.  The two listed
repairs are exactly the ways to keep a fixed physical scale or a
fixed local carrier before the limit is taken.
\end{proof}

\begin{warning}[The static action does not select physical time]
\label{warn:s6v22-action-not-transfer}
The exact trace-shell Hessian of V21 is an upstream branch and
local-vacuum certificate.  It does not by itself define a
reflection-positive slab kernel or its physical clock.  The
transfer law, reflection, and time normalization must occur on the
same response-fixed regulator family.
\end{warning}

\section{A spectral-support lemma for reflected correlations}
\label{sec:s6v22-spectral-support}

Let \({\cal H}\) be a Hilbert space, let \(H\ge0\) be
self-adjoint with normalized vacuum \(\Omega\), and fix a physical
delay \(\tau>0\).  Put
\[
 T=e^{-\tau H}.
\]
Let \(E\) be finite-dimensional and let
\[
 {\cal V}:E\longrightarrow\Omega^\perp
\]
be a source synthesis.  Define its delayed Gram matrices by
\begin{equation}
 G_n:={\cal V}^*T^n{\cal V},
 \qquad n\ge0.                                           \label{eq:s6v22-delayed-Gram}
\end{equation}

\begin{theorem}[Exponential delayed Grams locate the spectrum]
\label{thm:s6v22-Gram-gap}
Suppose that \(G_0\) is positive definite on \(E\) and that for
some \(C<\infty\) and \(q\in(0,1)\),
\begin{equation}
 0\preceq G_n\preceq Cq^nG_0
 \qquad(n\ge0).                                          \label{eq:s6v22-Gram-decay}
\end{equation}
Let
\begin{equation}
 {\cal C}_{\cal V}
 :=
 \overline{\operatorname{span}}
 \{p(T){\cal V}x:p\text{ a polynomial},\ x\in E\}.       \label{eq:s6v22-cyclic-bank}
\end{equation}
Then
\begin{equation}
 \boxed{
 \left\|T|_{{\cal C}_{\cal V}}\right\|\le q,
 \qquad
 H|_{{\cal C}_{\cal V}}
 \succeq -{\log q\over\tau}\,I.}                         \label{eq:s6v22-bank-gap}
\end{equation}
The constant \(C\) affects no spectral threshold.
\end{theorem}

\begin{proof}
For \(x\in E\), let \(\rho_x\) be the spectral measure of \(T\)
at \({\cal V}x\).  Then
\[
 x^*G_nx=\int_{[0,1]}\lambda^n\,\rho_x(\dd\lambda).
\]
If \(\rho_x([q+\varepsilon,1])>0\) for some
\(\varepsilon>0\), then
\[
 x^*G_nx
 \ge(q+\varepsilon)^n
 \rho_x([q+\varepsilon,1]),
\]
contradicting \eqref{eq:s6v22-Gram-decay} as
\(n\to\infty\).  Thus every \(\rho_x\) is supported in
\([0,q]\).  Polynomial functional calculus shows that the same
support bound holds on the span in
\eqref{eq:s6v22-cyclic-bank}, and closure preserves it.  Hence
\(\|T\|\le q\) on the cyclic bank.  Since
\(\lambda=e^{-\tau E}\), spectral calculus gives the stated lower
bound for \(H\).
\end{proof}

\begin{corollary}[A dense local family gives a full gap]
\label{cor:s6v22-dense-gap}
Let \({\cal D}\subset\Omega^\perp\) be dense.  Suppose that one
number \(\mu>0\) has the following property: for every
\(\psi\in{\cal D}\) there is \(C_\psi<\infty\) such that
\begin{equation}
 0\le\langle\psi,e^{-n\tau H}\psi\rangle
 \le C_\psi e^{-\mu n\tau}\|\psi\|^2
 \qquad(n\ge0).                                          \label{eq:s6v22-dense-decay}
\end{equation}
Then
\begin{equation}
 \boxed{H|_{\Omega^\perp}\succeq\mu I.}                  \label{eq:s6v22-full-gap}
\end{equation}
\end{corollary}

\begin{proof}
Apply the scalar case of
\Cref{thm:s6v22-Gram-gap} to each \(\psi\), with
\(q=e^{-\mu\tau}\).  Hence every vector in \({\cal D}\)
belongs to the closed spectral subspace
\(\mathbf1_{[\mu,\infty)}(H){\cal H}\).  Density forces this
subspace to contain all of \(\Omega^\perp\).
\end{proof}

\begin{assessment}[Why the decay prefactor is harmless]
\label{ass:s6v22-prefactor}
A mass gap is the lower endpoint of spectral support, not a
one-delay norm estimate.  Therefore an observable-dependent finite
constant in front of \(e^{-\mu t}\) disappears after large
integer delays.  What must be uniform is the exponent \(\mu\), and
the source must survive in the limiting OS Hilbert space.
\end{assessment}

\section{Dobrushin paths give the required delayed Grams}
\label{sec:s6v22-WCIC-Grams}

Let \(C=(c_{ef})\) be the actual interdependence matrix of one
response-fixed Euclidean shell specification.  Use its interaction
graph: two block variables are adjacent when some complete
interaction contains both.  Suppose
\begin{equation}
 \|C\|_{\infty}\le c:=1-\epsilon_{\rm CI}<1.             \label{eq:s6v22-uniform-c}
\end{equation}
Let \(F_1,\ldots,F_d\) be centered, reflection-admissible local
writers with common finite support \(S\), and put
\[
 F_x:=\sum_{\alpha=1}^dx_\alpha F_\alpha.
\]
Assume there are constants \(a_0,A_0,B_0>0\) such that
\begin{align}
 \|[F_x]\|_{\rm OS}^2&\ge a_0\|x\|^2,                   \label{eq:s6v22-lower-frame}\\
 \sum_{e\in S}\delta_e(F_x)&\le A_0\|x\|,               \label{eq:s6v22-osc-A}\\
 \max_{e\in S}\delta_e(F_x)&\le B_0\|x\|.               \label{eq:s6v22-osc-B}
\end{align}
Here \([F_x]\) denotes the OS equivalence class.

\begin{lemma}[Interaction-distance Dobrushin tail]
\label{lem:s6v22-Dobrushin-tail}
Let \(D=(I-C)^{-1}\).  If sets \(A,B\) have interaction-graph
distance at least \(r\), then
\begin{equation}
 \sup_{e\in A}\sum_{f\in B}D_{ef}
 \le {c^r\over1-c}.                                      \label{eq:s6v22-D-tail}
\end{equation}
Consequently bounded observables supported in \(A,B\) obey
\begin{equation}
 |\operatorname{Cov}(F,G)|
 \le {c^r\over1-c}
 \left(\sum_{e\in A}\delta_e(F)\right)
 \max_{f\in B}\delta_f(G).                              \label{eq:s6v22-cov-tail}
\end{equation}
\end{lemma}

\begin{proof}
An entry \((C^k)_{ef}\) vanishes when the graph distance from
\(e\) to \(f\) exceeds \(k\).  Thus only \(k\ge r\) contributes
from \(e\in A\) to \(f\in B\).  The row sum of \(C^k\) is at most
\(c^k\).  Summing the geometric tail proves
\eqref{eq:s6v22-D-tail}.  Insert it into the Dobrushin comparison
bound \eqref{eq:s6v19-comparison}.
\end{proof}

\begin{theorem}[WCIC-to-OS finite-bank gap]
\label{thm:s6v22-WCIC-OS}
Assume the response-fixed lattice law is reflection positive and
its physical time translation by \(\tau\) induces an OS transfer
\(T=e^{-\tau H}\).  Suppose that for some integer
\(s_\tau\ge1\) and \(r_0\ge0\),
\begin{equation}
 \operatorname{dist}_E
 (S,\vartheta_{n\tau}S)
 \ge ns_\tau-r_0
 \qquad(n\ge1),                                          \label{eq:s6v22-time-incidence}
\end{equation}
where \(\vartheta_{n\tau}\) is Euclidean-time translation.  Under
\eqref{eq:s6v22-uniform-c}--\eqref{eq:s6v22-osc-B}, the cyclic
OS bank generated by the \(F_\alpha\) has the gap
\begin{equation}
 \boxed{
 H|_{{\cal C}_F}\succeq
 \mu_{\rm CI}I,
 \qquad
 \mu_{\rm CI}:=
 {s_\tau\over\tau}
 \left|\log(1-\epsilon_{\rm CI})\right|>0.}              \label{eq:s6v22-CI-mass}
\end{equation}
\end{theorem}

\begin{proof}
Reflection positivity and the OS construction give
\begin{equation}
 \langle[F_x],T^n[F_x]\rangle_{\rm OS}
 =
 \operatorname{Cov}
 (F_x,\vartheta_{n\tau}F_x)\ge0.                         \label{eq:s6v22-OS-cov}
\end{equation}
For \(n\) with \(ns_\tau\ge r_0\),
\Cref{lem:s6v22-Dobrushin-tail} and
\eqref{eq:s6v22-osc-A}--\eqref{eq:s6v22-osc-B} give
\[
 \langle[F_x],T^n[F_x]\rangle_{\rm OS}
 \le
 {A_0B_0c^{-r_0}\over1-c}
 (c^{s_\tau})^n\|x\|^2.
\]
The lower frame \eqref{eq:s6v22-lower-frame} replaces
\(\|x\|^2\) by
\(a_0^{-1}\|[F_x]\|_{\rm OS}^2\).  Enlarging the finite
prefactor covers the remaining finitely many \(n\).  Apply
\Cref{thm:s6v22-Gram-gap} with \(q=c^{s_\tau}\).  The identity
\(-\tau^{-1}\log q=\mu_{\rm CI}\) proves the result.
\end{proof}

\begin{corollary}[Finite prefactors and support sizes may grow
with the bank]
\label{cor:s6v22-bank-growth}
For each fixed finite source bank, the constants
\(a_0^{-1},A_0,B_0,r_0\) in
\Cref{thm:s6v22-WCIC-OS} may depend on that bank.  If the same
\(\epsilon_{\rm CI}\), \(s_\tau\), and physical \(\tau\) apply
to a dense increasing union of banks, the full OS Hamiltonian
still has the common gap \(\mu_{\rm CI}\).
\end{corollary}

\begin{proof}
The bank-dependent quantities occur only in the prefactor of
\eqref{eq:s6v22-Gram-decay}; by
\Cref{ass:s6v22-prefactor} they do not change the spectral
threshold.  Apply \Cref{cor:s6v22-dense-gap} to the dense union.
\end{proof}

\section{Cofinal passage and exact remaining hypotheses}
\label{sec:s6v22-cofinal}

\begin{theorem}[Conditional response-fixed Yang--Mills mass
theorem]
\label{thm:s6v22-conditional-mass}
Consider a regulator sequence on the trajectory selected in
\Cref{thm:s6v22-trajectory-correction}.  Suppose:
\begin{enumerate}[label=\textup{(M\arabic*)},leftmargin=4.3em]
\item the exact trace-shell static card passes uniformly on the
      selected local-vacuum branch;
\item the complete gauge-reduced shell specifications satisfy
      WCIC with one \(\epsilon_{\rm CI}>0\), and the physical
      temporal incidence \eqref{eq:s6v22-time-incidence} has
      uniform \(s_\tau\ge1\) at one fixed \(\tau>0\);
\item the gauge-invariant local cylinder laws converge to a
      nontrivial reflection-positive, time-translation-invariant
      continuum measure with unique vacuum;
\item every fixed local source bank has a regulator-uniform lower
      frame and its delayed Grams converge to the corresponding
      continuum OS Grams; and
\item the centred gauge-invariant local Wilson algebra is dense
      in the continuum vacuum complement.
\end{enumerate}
Then the reconstructed continuum Yang--Mills Hamiltonian obeys
\begin{equation}
 \boxed{
 H_{\rm YM}|_{\Omega^\perp}
 \succeq
 {s_\tau\over\tau}
 \left|\log(1-\epsilon_{\rm CI})\right|I.}               \label{eq:s6v22-conditional-YM-gap}
\end{equation}
\end{theorem}

\begin{proof}
For each fixed bank, the uniform WCIC, incidence, oscillation,
and lower-frame bounds give the lattice delayed-Gram estimate in
\Cref{thm:s6v22-WCIC-OS}.  Convergence of the delayed Grams
preserves that estimate in the continuum, including its common
exponent.  The source lower frame prevents a nonzero physical
writer from disappearing into the limiting OS null space.  Apply
\Cref{cor:s6v22-bank-growth} to the dense centred local algebra.
Vacuum uniqueness identifies its orthogonal complement with the
entire positive-energy space, proving
\eqref{eq:s6v22-conditional-YM-gap}.
\end{proof}

\begin{assessment}[Relation to the Grand Tensor promotion theorem]
\label{ass:s6v22-GT-relation}
The audited Grand Tensor theorem promotes regulator finite-bank
contraction by a collectively compact complement or a
graph--Rellich bound.  \Cref{thm:s6v22-conditional-mass} supplies a
different but compatible endpoint: once a reflection-positive
continuum core and a dense noncollapsing local source algebra have
already been constructed, the WCIC exponent puts every dense
source vector in the same positive spectral subspace.  If the
continuum core is built only through finite-bank regulator limits,
the Grand Tensor collective-compactness row remains necessary to
justify that those banks exhaust the core.
\end{assessment}

\begin{firewall}[Conditional theorem versus a completed Clay
theorem]
\label{fire:s6v22-not-complete}
\Cref{thm:s6v22-conditional-mass} is a rigorous implication, not
a proof that its hypotheses hold for four-dimensional
\(\SU(3)\) Yang--Mills.  In particular, the attached sources do
not supply the weak-coupling response-fixed WCIC constant, the
nontrivial reflection-positive continuum measure, the uniform
local source lower frame, or vacuum uniqueness.
\end{firewall}

\begin{assessment}[Exact live bottleneck after V22]
\label{ass:s6v22-bottleneck}
The logical bridge from whole conditional influence to a
Hamiltonian mass is now closed: exponential Dobrushin paths give
delayed OS Grams, their asymptotic exponent gives a finite-bank
spectral gap, and density promotes the common exponent to the full
vacuum complement.

The live mathematical bottleneck is earlier.  Single-link bare
strong-coupling WCIC is known, but the selected continuum
trajectory is response-fixed and weak-facing.  The complete
block conditional on that trajectory has not been loaded with a
uniform strict row, and the canonical-root static card remains
unpopulated.  The Grand Tensor toron theorem additionally forbids
retreat to the over-weak full finite-torus bare path.
\end{assessment}

\begin{programme}[V23: test whether single-link WCIC can survive
the weak-facing limit]
\label{prog:s6v22-V23}
\begin{enumerate}[leftmargin=3em]
\item Keep the response-fixed/local-vacuum trajectory of
      \Cref{thm:s6v22-trajectory-correction}; do not return to the
      obstructed over-weak bare path.
\item On the exact tree-gauge conditional
      \eqref{eq:s6v21-one-link-law}, derive a lower as well as an
      upper bound for a one-link influence as the Wilson
      coefficient grows.
\item Determine whether adversarial exterior configurations force
      the single-link Dobrushin coefficient to one.  If so, replace
      WCIC by a physically sized block or a typical-set/local-vacuum
      contraction, rather than adding separate oscillation debits.
\item Keep the canonical-root unsplit static card as an independent
      upstream gate.
\item Preserve the physical time incidence and source lower frame
      needed by \Cref{thm:s6v22-conditional-mass}.
\end{enumerate}
\end{programme}

\begin{firewall}[Claim boundary after sixth-series V22]
\label{fire:s6v22-boundary}
V22 absorbs the Grand Tensor wrapping-toron trajectory obstruction,
proves the delayed-Gram spectral-support theorem, proves the dense
source gap corollary, derives the interaction-distance Dobrushin
tail, proves the WCIC-to-OS finite-bank gap with an explicit
physical rate, and gives a rigorous conditional continuum
Yang--Mills mass theorem.

V22 does not prove the canonical-root static card, weak-coupling
WCIC, a block or typical-set substitute, uniform source-frame
noncollapse, dense continuum occurrence, vacuum uniqueness,
collective compactness, a nontrivial reflection-positive
continuum measure, or the hypotheses of its conditional mass
theorem.  The full Yang--Mills mass gap has not been proved.
\end{firewall}

\begin{theorem}[Sixth-series V22 result ledger]
\label{thm:s6v22-results}
V22:
\begin{enumerate}[label=\textup{(V22.\arabic*)},leftmargin=6em]
\item audits the Grand Tensor Yang--Mills transfer interface;
\item imports the wrapping-toron obstruction only at its proved
      bare-trajectory scope;
\item forces a response-fixed or local-vacuum trajectory;
\item proves exponential Gram decay implies a cyclic spectral gap;
\item proves a common exponent on a dense family gives a full gap;
\item derives the Dobrushin interaction-distance tail;
\item proves WCIC gives a finite-bank OS gap;
\item identifies the explicit mass rate
      \(s_\tau|\log(1-\epsilon_{\rm CI})|/\tau\); and
\item states and proves the exact conditional continuum mass
      handoff.
\end{enumerate}
\end{theorem}

\begin{proof}
Use
\Cref{audit:s6v22-GT,
thm:s6v22-trajectory-correction,
thm:s6v22-Gram-gap,
cor:s6v22-dense-gap,
lem:s6v22-Dobrushin-tail,
thm:s6v22-WCIC-OS,
cor:s6v22-bank-growth,
thm:s6v22-conditional-mass}.
\end{proof}

\paragraph{Parent-version record.}
V22 was formed from
\texttt{Renewal\_Yang\_Mills\_Mass\_Gap\_Research\_Notes\_V21.tex}.
The V21 parent had \(9\,142\) logical source lines,
\(323\,789\) bytes, and SHA--256
\[
 \texttt{4393ea3951991d7e1bcaa67210893033d54b9b272b3260a38d1f90a3989bf150}.
\]
The next compulsory companion audit remains V25.

\clearpage
\chapter{Sixth-series V23: zero-temperature conditional
bifurcation and the single-link WCIC test}
\label{ch:s6v23-bifurcation}

\section{A quantitative compact Gibbs concentration lemma}
\label{sec:s6v23-concentration}

Let \(X\) be a compact metric space and let \(m\) be a Borel
probability measure of full support.  For a continuous function
\(\phi:X\to\mathbb R\), write
\[
 a_\phi:=\min_X\phi,
 \qquad
 {\cal M}_\phi:=\{x:\phi(x)=a_\phi\}.
\]

\begin{lemma}[Concentration near the complete minimizer set]
\label{lem:s6v23-concentration}
Let \(U\) be an open neighborhood of
\({\cal M}_\phi\).  Define
\begin{equation}
 g_U
 :=
 \inf_{x\in X\setminus U}\{\phi(x)-a_\phi\}>0,            \label{eq:s6v23-gap}
\end{equation}
and
\begin{equation}
 m_U
 :=
 m\{x:\phi(x)<a_\phi+g_U/4\}>0.                          \label{eq:s6v23-mass}
\end{equation}
For a bounded perturbation \(r_\lambda\), let
\begin{equation}
 \dd\mu_\lambda(x)
 =
 {e^{-\lambda\phi(x)-r_\lambda(x)}m(\dd x)
  \over
  \displaystyle\int_Xe^{-\lambda\phi-r_\lambda}\dd m}.
                                                                    \label{eq:s6v23-Gibbs}
\end{equation}
Then
\begin{equation}
 \boxed{
 \mu_\lambda(X\setminus U)
 \le
 m_U^{-1}
 \exp\{-3\lambda g_U/4+2\|r_\lambda\|_\infty\}.}          \label{eq:s6v23-concentration}
\end{equation}
In particular, if
\(\|r_\lambda\|_\infty\le\lambda g_U/8\), then
\begin{equation}
 \mu_\lambda(X\setminus U)
 \le m_U^{-1}e^{-\lambda g_U/2}.                         \label{eq:s6v23-simple-concentration}
\end{equation}
\end{lemma}

\begin{proof}
Compactness and the inclusion
\({\cal M}_\phi\subset U\) imply \(g_U>0\).
The strict sublevel in \eqref{eq:s6v23-mass} is a nonempty open
set, so full support gives \(m_U>0\).  On \(X\setminus U\), the
unnormalized density is at most
\[
 e^{-\lambda(a_\phi+g_U)+\|r_\lambda\|_\infty}.
\]
On the strict sublevel used in \eqref{eq:s6v23-mass}, the
partition function is at least
\[
 m_U
 e^{-\lambda(a_\phi+g_U/4)-\|r_\lambda\|_\infty}.
\]
Their ratio gives \eqref{eq:s6v23-concentration}.  The assumed
remainder bound makes
\(-3\lambda g_U/4+2\|r_\lambda\|_\infty
\le-\lambda g_U/2\).
\end{proof}

\begin{theorem}[Separated minimizers force total variation to
one]
\label{thm:s6v23-separated-TV}
Let \(\phi_0,\phi_1\) be continuous functions on \(X\) with
disjoint minimizer sets.  Then there are disjoint open
neighborhoods \(U_0,U_1\), numbers \(g>0,m_*>0\), and the
following property.  If
\[
 \dd\mu_{i,\lambda}
 \propto
 e^{-\lambda\phi_i-r_{i,\lambda}}\dd m,
 \qquad
 \|r_{i,\lambda}\|_\infty\le\lambda g/8,
\]
then
\begin{equation}
 \boxed{
 \|\mu_{0,\lambda}-\mu_{1,\lambda}\|_{\rm TV}
 \ge
 1-{2\over m_*}e^{-\lambda g/2}.}                        \label{eq:s6v23-TV-lower}
\end{equation}
Thus the total-variation distance tends to one whenever
\(\|r_{i,\lambda}\|_\infty=o(\lambda)\).
\end{theorem}

\begin{proof}
Disjoint compact subsets of a metric space have disjoint open
neighborhoods \(U_0,U_1\).  Apply
\Cref{lem:s6v23-concentration} to each pair
\((\phi_i,U_i)\), and take \(g\) to be the smaller gap and
\(m_*\) the smaller sublevel mass, decreasing \(g\) if needed.
For the event \(U_0\),
\[
 \mu_{0,\lambda}(U_0)\ge1-m_*^{-1}e^{-\lambda g/2},
 \qquad
 \mu_{1,\lambda}(U_0)
 \le\mu_{1,\lambda}(X\setminus U_1)
 \le m_*^{-1}e^{-\lambda g/2}.
\]
The variational definition of total variation gives
\eqref{eq:s6v23-TV-lower}.  An \(o(\lambda)\) remainder satisfies
the displayed bound for all sufficiently large \(\lambda\).
\end{proof}

\section{The conditional ground-state bifurcation card}
\label{sec:s6v23-card}

Use the exact tree-gauge conditional law
\eqref{eq:s6v21-one-link-law}.  Consider a weak-facing sequence
with scale \(\lambda_r\to\infty\).  For a selected reduced link
\(e_r\) and exterior \(\eta\), write its complete conditional
energy as
\begin{equation}
 {\cal S}_{r,e_r}(u\mid\eta)
 =
 \lambda_r\phi_{r,\eta}(u)+r_{r,\eta}(u).                \label{eq:s6v23-leading-energy}
\end{equation}
All Wilson, response-generated, marginal, and support-birth terms
of order \(\lambda_r\) belong to \(\phi_{r,\eta}\), not to the
remainder.

\begin{definition}[Staple-bifurcation incidence card, SBIC]
\label{def:s6v23-SBIC}
SBIC is the following finite conditional record:
\begin{enumerate}[label=\textup{(\roman*)},leftmargin=3.5em]
\item links \(e_r\ne f_r\) and exteriors
      \(\eta_r\sim_{f_r}\eta'_r\);
\item disjoint open sets \(U_r,U'_r\subset G\) containing,
      respectively, the complete minimizer sets of
      \(\phi_{r,\eta_r}\) and \(\phi_{r,\eta'_r}\);
\item uniform constants \(g>0,m_*>0\) for the two energy gaps and
      Haar sublevel masses in
      \eqref{eq:s6v23-gap}--\eqref{eq:s6v23-mass}; and
\item the lower-order condition
      \[
       \max\{\|r_{r,\eta_r}\|_\infty,
              \|r_{r,\eta'_r}\|_\infty\}
       \le\lambda_rg/8
      \]
      for all sufficiently large \(r\).
\end{enumerate}
\end{definition}

\begin{theorem}[SBIC excludes uniform single-link WCIC]
\label{thm:s6v23-SBIC-no-go}
If SBIC holds along a weak-facing sequence, then the actual
single-link interdependence coefficient satisfies
\begin{equation}
 \boxed{
 c_{e_rf_r}
 \ge
 1-{2\over m_*}e^{-\lambda_rg/2}
 \longrightarrow1.}                                     \label{eq:s6v23-influence-one}
\end{equation}
Consequently there is no cutoff-independent
\(\epsilon_{\rm CI}>0\) for the single-link WCIC row
\eqref{eq:s6v20-WCIC}.
\end{theorem}

\begin{proof}
The two conditional laws in the definition of
\(c_{e_rf_r}\) include the exterior pair in SBIC.  Apply
\Cref{thm:s6v23-separated-TV} to their complete leading energies
and remainders.  Since the Dobrushin row containing \(e_r\) is at
least \(c_{e_rf_r}\), its supremum tends to at least one.
\end{proof}

\begin{corollary}[Order-one response terms must enter the leading
card]
\label{cor:s6v23-order-one-response}
Suppose a response-generated term has conditional oscillation of
order \(\lambda_r\).  Treating it as \(r_{r,\eta}\) is invalid in
\Cref{thm:s6v23-SBIC-no-go}; it must be included in the full
leading energy \(\phi_{r,\eta}\).  A lower-order term with
uniform \(o(\lambda_r)\) norm cannot repair a uniform separated
minimizer card.
\end{corollary}

\begin{proof}
The first assertion is the definition of the decomposition
\eqref{eq:s6v23-leading-energy}.  For the second, an
\(o(\lambda_r)\) term satisfies the remainder hypothesis of
\Cref{thm:s6v23-separated-TV}, so the same total-variation lower
bound tends to one.
\end{proof}

\begin{assessment}[SBIC is a finite geometric test]
\label{ass:s6v23-geometric}
SBIC does not require a partition function, a covariance sum, or
an infinite-volume Gibbs state.  It asks for two compact
one-link minimizer problems and uniform local gaps after the
complete leading response has been assembled.  A successful SBIC
retires worst-case single-link Dobrushin contraction on that
trajectory.  A failed attempt to find separated minimizers does
not prove WCIC; it merely keeps the contraction branch open.
\end{assessment}

\section{An exact \(\SU(3)\) one-plaquette bifurcation}
\label{sec:s6v23-SU3-example}

Let
\[
 \omega=e^{2\pi i/3}.
\]
For a staple \(S\in\SU(3)\), define the one-plaquette Wilson
conditional law
\begin{equation}
 \dd\gamma_{\beta,S}(U)
 =
 { \exp\{(\beta/3)\Re\Tr(US)\}\dd\lambda_{\SU(3)}(U)
  \over
  \displaystyle
  \int_{\SU(3)}
  \exp\{(\beta/3)\Re\Tr(VS)\}\dd\lambda_{\SU(3)}(V)}.
                                                                    \label{eq:s6v23-one-plaquette-law}
\end{equation}

\begin{lemma}[The two centre staples have disjoint unique
maximizers]
\label{lem:s6v23-centre-maximizers}
The function \(\Re\Tr U\) on \(\SU(3)\) has the unique maximum
at \(U=I\).  The function \(\Re\Tr(\omega U)\) has the unique
maximum at \(U=\omega^2I\).
\end{lemma}

\begin{proof}
If \(e^{i\theta_1},e^{i\theta_2},e^{i\theta_3}\) are the
eigenvalues of \(U\), then
\[
 \Re\Tr U=\sum_{j=1}^3\cos\theta_j\le3.
\]
Equality forces every eigenvalue to be one, and a unitary matrix
with that spectrum is \(I\).  Multiplication by \(\omega I\) is
a bijection of \(\SU(3)\).  Therefore
\(\Re\Tr(\omega U)\) is maximized uniquely when
\(\omega U=I\), namely at \(U=\omega^2I\).
\end{proof}

\begin{theorem}[One-plaquette Wilson influence tends to one]
\label{thm:s6v23-one-plaquette-no-go}
The two conditional laws
\(\gamma_{\beta,I}\) and \(\gamma_{\beta,\omega I}\) satisfy
\begin{equation}
 \boxed{
 \|\gamma_{\beta,I}-\gamma_{\beta,\omega I}\|_{\rm TV}
 \longrightarrow1
 \qquad(\beta\to\infty).}                                \label{eq:s6v23-centre-TV}
\end{equation}
Hence the isolated one-plaquette two-link Wilson specification
has a single-link Dobrushin coefficient tending to one.
\end{theorem}

\begin{proof}
Apply \Cref{thm:s6v23-separated-TV} on
\(X=\SU(3)\), with Haar measure,
\[
 \phi_0(U)=-\frac13\Re\Tr U,
 \qquad
 \phi_1(U)=-\frac13\Re\Tr(\omega U),
\]
and zero remainders.  Their minimizers are the two distinct centre
elements in \Cref{lem:s6v23-centre-maximizers}.  In an isolated
plaquette, changing one neighboring link from \(I\) to
\(\omega I\) changes the staple in the displayed conditional
exactly between these two values.
\end{proof}

\begin{firewall}[The one-plaquette witness is not yet the
four-dimensional SBIC]
\label{fire:s6v23-not-embedded}
In the four-dimensional link conditional, the selected link is
incident to as many as six plaquettes.  The other five staples and
every order-one response interaction also enter the leading energy.
\Cref{thm:s6v23-one-plaquette-no-go} proves an exact local
mechanism and rules out any universal low-temperature
single-plaquette Dobrushin theorem.  It does not prove that the
full parity-tree conditional admits exteriors realizing the same
disjoint minimizers.  That is the finite SBIC still to be
executed.
\end{firewall}

\section{Consequences for the block and local-vacuum routes}
\label{sec:s6v23-consequences}

\begin{theorem}[Bifurcation persists across any insufficient
blocking]
\label{thm:s6v23-block-persistence}
Let \(B_r\) and \(B'_r\) be two neighboring conditional blocks.
Suppose their complete block conditional energies have an SBIC
with uniformly separated minimizer sets in the block fibre and
lower-order remainder \(o(\lambda_r)\).  Then the block influence
coefficient tends to one.  In particular, grouping variables
inside \(B_r\) repairs the single-link no-go only if the
ground-state distinction caused by \(B'_r\) becomes internal to
the block, is removed by the physical quotient, or is destroyed
by an order-one response term.
\end{theorem}

\begin{proof}
Apply \Cref{thm:s6v23-separated-TV} to the compact block fibre
with its product Haar reference.  The proof does not depend on the
dimension of \(X\).  The three listed repairs are the ways in which
the separated-minimizer hypothesis can fail after blocking.
\end{proof}

\begin{warning}[Worst-case exteriors may be physically rare]
\label{warn:s6v23-rare-exteriors}
Dobrushin coefficients take a supremum over all exterior
configurations.  SBIC may therefore destroy WCIC even if its two
exteriors have exponentially small probability in the selected
local vacuum.  In that case the mathematically appropriate
replacement is a typical-set, disagreement-percolation, contour,
or phase-restricted contraction with a separately paid bad-set
probability.  One may not simply discard the exteriors while
continuing to quote the worst-case WCIC theorem.
\end{warning}

\begin{theorem}[Exact V23 decision alternative]
\label{thm:s6v23-decision}
For any explicitly loaded weak-facing conditional action, exactly
one of the following outcomes is licensed by the present results.
\begin{enumerate}[label=\textup{(\roman*)},leftmargin=3.5em]
\item A uniform SBIC is proved.  Then single-link WCIC is
      impossible and the programme must pass to a sufficiently
      large block or a probability-weighted local-vacuum method.
\item Every proposed SBIC fails because the complete leading
      minimizer sets overlap or their uniform gap collapses.  No
      contraction follows; the next task is a quantitative upper
      bound for the actual conditional total variation.
\item The complete order-one response action is unavailable.
      Neither SBIC nor WCIC is defined, and action loading is the
      first missing row.
\end{enumerate}
\end{theorem}

\begin{proof}
If SBIC holds, apply
\Cref{thm:s6v23-SBIC-no-go}.  If its hypotheses fail for an
available action, the lower-bound theorem is silent and an upper
bound is needed.  If the action is unavailable,
\Cref{thm:s6v21-whole-oscillation} shows that even the relevant
mixed conditional function is undetermined.
\end{proof}

\begin{assessment}[Exact live bottleneck after V23]
\label{ass:s6v23-bottleneck}
Worst-case single-link WCIC is now subject to a sharp
zero-temperature test.  The isolated Wilson plaquette fails that
test maximally: a centre change of one neighbor moves the unique
conditional minimizer between disjoint centre elements and sends
total variation to one.  To decide the actual response-fixed
four-dimensional law, the remaining finite calculation is to
assemble all six staples and every order-one response term, then
solve their two exterior minimizer problems.

This conditional test does not remove the independent
canonical-root static-card gate, nor does it construct the
reflection-positive continuum required by
\Cref{thm:s6v22-conditional-mass}.
\end{assessment}

\begin{programme}[V24: execute the full parity-tree SBIC]
\label{prog:s6v23-V24}
\begin{enumerate}[leftmargin=3em]
\item Write the complete leading one-link energy for each
      parity-tree link type, including all six Wilson staples and
      every order-one response-generated term.
\item For each neighboring link type, impose the exact
      shared-plaquette incidence from
      \eqref{eq:s6v21-common-incidence}.
\item Search analytically for two exteriors differing only at that
      neighbor whose complete minimizer sets are disjoint.
\item If found, prove uniform energy-gap and Haar-sublevel bounds
      and retire single-link WCIC by
      \Cref{thm:s6v23-SBIC-no-go}.
\item If no bifurcation exists, compute a quantitative overlap or
      total-variation upper bound for the full conditional.
\item On the no-go branch, define the smallest physical block
      which internalizes the bifurcating centre/staple mode; do
      not assume blocking repairs it.
\end{enumerate}
\end{programme}

\begin{firewall}[Claim boundary after sixth-series V23]
\label{fire:s6v23-boundary}
V23 proves quantitative compact Gibbs concentration, proves that
separated conditional ground states force total variation to one,
defines the exact SBIC, proves SBIC incompatible with uniform
single-link WCIC, constructs an explicit centre-separated
\(\SU(3)\) one-plaquette witness, and proves the same obstruction
persists for blocks whose ground-state separation crosses the
block boundary.

V23 does not embed that witness into the full four-dimensional
response-matched parity-tree conditional, populate the
order-one generated action, prove a block or typical-set
contraction, close the canonical-root static card, construct the
continuum measure, or prove the full Yang--Mills mass gap.
\end{firewall}

\begin{theorem}[Sixth-series V23 result ledger]
\label{thm:s6v23-results}
V23:
\begin{enumerate}[label=\textup{(V23.\arabic*)},leftmargin=6em]
\item proves an explicit compact Gibbs concentration inequality;
\item proves separated minimizers force total variation to one;
\item defines the staple-bifurcation incidence card;
\item proves SBIC rules out uniform single-link WCIC;
\item shows lower-order response terms cannot repair SBIC;
\item proves the unique \(\SU(3)\) centre maximizers;
\item gives an exact one-plaquette Wilson no-go;
\item separates that local witness from the full four-dimensional
      incidence problem;
\item proves bifurcation persistence under insufficient blocking;
      and
\item gives the exact action/SBIC/upper-bound decision alternative.
\end{enumerate}
\end{theorem}

\begin{proof}
Use
\Cref{lem:s6v23-concentration,
thm:s6v23-separated-TV,
thm:s6v23-SBIC-no-go,
cor:s6v23-order-one-response,
lem:s6v23-centre-maximizers,
thm:s6v23-one-plaquette-no-go,
thm:s6v23-block-persistence,
thm:s6v23-decision}.
\end{proof}

\paragraph{Parent-version record.}
V23 was formed from
\texttt{Renewal\_Yang\_Mills\_Mass\_Gap\_Research\_Notes\_V22.tex}.
The V22 parent had \(9\,612\) logical source lines,
\(342\,023\) bytes, and SHA--256
\[
 \texttt{4fcc777a1932deaa2efc708be1d3ab6a7c13cd290b0c3ba17112a8dd3325c393}.
\]
The next compulsory companion audit remains V25.

\clearpage
\chapter{Sixth-series V24: an embedded parity-tree
low-temperature influence obstruction}
\label{ch:s6v24-embedded}

\section{The six controllable staples at an all-odd reduced link}
\label{sec:s6v24-parity-star}

Retain the exact fine-side convention \(N=2M\) and the parity
forest from \Cref{audit:s6v21-YM35}.  A fine vertex has a unique
form
\[
 x=2y+s,\qquad s\in\{0,1\}^4.
\]
Inside a parity cube, a tree edge removes the least odd coordinate
of its child.  The designated retained second constituent in
direction \(\mu\) starts at parity \(s=\widehat e_\mu\).  All
other positive links are reduced variables.

\begin{lemma}[An explicit free reduced staple star]
\label{lem:s6v24-free-star}
Let
\[
 s=(1,1,1,1),\qquad x=2y+s,
\qquad e=(x,x+\widehat e_1).
\]
Then \(e\in{\mathscr E}^{\rm red}_N\).  For
\(\nu\in\{2,3,4\}\), define
\begin{equation}
 a_{+\nu}:=(x,x+\widehat e_\nu),
 \qquad
 a_{-\nu}:=(x-\widehat e_\nu,x).                         \label{eq:s6v24-controls}
\end{equation}
The six links \(a_{\pm\nu}\) are distinct reduced links, are
different from \(e\), and each occurs in exactly one of the six
plaquette staples incident to \(e\).  By setting the other two
links in each staple to \(I\), the oriented value of that staple
can be assigned arbitrarily in the centre
\(\{I,\omega I,\omega^2I\}\) through the corresponding control
link.
\end{lemma}

\begin{proof}
The link \(e\) crosses the first parity face, but its initial
parity has four odd coordinates rather than only the first one.
It is therefore not a designated retained second constituent and
is reduced.

Each \(a_{+\nu}\) also starts at all-odd parity and crosses the
\(\nu\)-face.  Since its initial parity is not
\(\widehat e_\nu\), it is reduced.  The link \(a_{-\nu}\) is
internal to the parity cube: its child is \(x\), whose least odd
coordinate is \(1\).  Since \(\nu>1\), it is not the tree edge
which removes the least odd coordinate, and hence is reduced.

The \(+\nu\) and \(-\nu\) plaquettes through \(e\) have disjoint
transverse side links adjacent to \(x\); these are precisely the
six links in \eqref{eq:s6v24-controls}.  Direct inspection of the
plaquette words shows that each control appears in its named
staple and in none of the other five.  If its orientation in the
word is reversed, centre inversion interchanges \(\omega\) and
\(\omega^2\).  Thus an arbitrary prescribed centre value is
obtained after the other two staple links are set to the identity.
\end{proof}

\begin{corollary}[Two exteriors differing in one reduced link]
\label{cor:s6v24-two-exteriors}
There are two parity-tree exterior configurations
\(\eta\sim_f\eta'\), for one of the six control links \(f\), such
that the six oriented Wilson staples at \(e\) have total
\begin{equation}
 {\cal H}(\eta)=0,
 \qquad
 {\cal H}(\eta')=(\omega-1)I.                            \label{eq:s6v24-two-staple-sums}
\end{equation}
All tree links and all retained coarse links may simultaneously be
kept at \(I\).
\end{corollary}

\begin{proof}
Use one control link as \(f\).  Prescribe the other five oriented
staples to be
\[
 I,\quad\omega I,\quad\omega I,\quad
 \omega^2I,\quad\omega^2I.
\]
Their sum is
\[
 (1+2\omega+2\omega^2)I=-I.
\]
Set the variable staple to \(I\) under \(\eta\) and to
\(\omega I\) under \(\eta'\).  The resulting sums are those in
\eqref{eq:s6v24-two-staple-sums}.  The construction in
\Cref{lem:s6v24-free-star} uses only reduced control links, so the
tree and retained links remain fixed at the identity.  The two
exteriors differ only at \(f\).  Plaquettes not containing \(e\)
may change when \(f\) changes, but they cancel from the conditional
law at \(e\) by \Cref{thm:s6v21-star-disintegration}.
\end{proof}

\section{Haar versus a nonconstant Wilson conditional}
\label{sec:s6v24-Haar-concentration}

\begin{lemma}[A proper analytic maximum set has Haar measure zero]
\label{lem:s6v24-analytic-zero}
Let \(G\) be a connected compact real-analytic Lie group with Haar
probability \(\lambda_G\).  If \(F:G\to\mathbb R\) is
real analytic and nonconstant, then
\begin{equation}
 \lambda_G\{u:F(u)=\max_GF\}=0.                           \label{eq:s6v24-max-zero}
\end{equation}
\end{lemma}

\begin{proof}
The analytic function
\[
 u\longmapsto \max_GF-F(u)
\]
is not identically zero.  In every analytic coordinate chart, the
zero set of a nonzero real-analytic function has Lebesgue measure
zero: choose the first nonvanishing derivative in one coordinate
and apply the one-variable isolated-zero theorem on almost every
coordinate line, followed by induction and Fubini.  A finite
analytic atlas covers \(G\), and Haar measure has a smooth positive
density in each chart.  The maximum set therefore has Haar measure
zero.
\end{proof}

\begin{theorem}[A nonconstant compact Gibbs tilt separates from
Haar]
\label{thm:s6v24-Haar-TV}
Under the hypotheses of \Cref{lem:s6v24-analytic-zero}, define
\begin{equation}
 \dd\mu_\lambda(u)
 =
 {\exp\{\lambda F(u)\}\lambda_G(\dd u)
  \over
  \displaystyle\int_G\exp\{\lambda F(v)\}
  \lambda_G(\dd v)}.                                    \label{eq:s6v24-analytic-Gibbs}
\end{equation}
Then
\begin{equation}
 \boxed{
 \|\mu_\lambda-\lambda_G\|_{\rm TV}\longrightarrow1
 \qquad(\lambda\to\infty).}                              \label{eq:s6v24-Haar-TV}
\end{equation}
\end{theorem}

\begin{proof}
Let \({\cal M}\) be the maximum set.  Given
\(\varepsilon>0\), outer regularity and
\Cref{lem:s6v24-analytic-zero} give an open neighborhood
\(U\supset{\cal M}\) with
\(\lambda_G(U)<\varepsilon\).  Apply
\Cref{lem:s6v23-concentration} to \(\phi=-F\), with zero
remainder.  It gives \(\mu_\lambda(U)\to1\).  Hence
\[
 \|\mu_\lambda-\lambda_G\|_{\rm TV}
 \ge\mu_\lambda(U)-\lambda_G(U)
 \longrightarrow1-\lambda_G(U)>1-\varepsilon.
\]
Let \(\varepsilon\downarrow0\).
\end{proof}

\begin{lemma}[The nonzero centre staple gives a nonconstant tilt]
\label{lem:s6v24-nonconstant}
On \(\SU(3)\), the function
\begin{equation}
 F(U):={1\over3}\Re\Tr\bigl(U(\omega-1)I\bigr)            \label{eq:s6v24-F}
\end{equation}
is real analytic and nonconstant.
\end{lemma}

\begin{proof}
Matrix multiplication, trace, and real part are real analytic on
the compact matrix group.  At \(U=I\),
\[
 F(I)=\Re(\omega-1)=-\frac32,
\]
whereas at \(U=\omega^2I\),
\[
 F(\omega^2I)
 =\Re(\omega^2(\omega-1))
 =\Re(1-\omega^2)=\frac32.
\]
Thus \(F\) is nonconstant.
\end{proof}

\section{Actual parity-tree failure of bare single-link WCIC}
\label{sec:s6v24-WCIC-failure}

\begin{theorem}[Embedded pure-Wilson influence tends to one]
\label{thm:s6v24-embedded-no-go}
For the exact parity-tree reduced-fibre pure Wilson law at
retained background \(W=I\), there are reduced links
\(e\ne f\) such that
\begin{equation}
 \boxed{
 c_{ef}(\beta)\longrightarrow1
 \qquad(\beta\to\infty).}                                \label{eq:s6v24-cef-one}
\end{equation}
Consequently the bare single-reduced-link WCIC has no uniform
strict margin on the weak-facing branch.
\end{theorem}

\begin{proof}
Use the links and exterior pair from
\Cref{cor:s6v24-two-exteriors}.  Up to a conditional constant, the
six Wilson plaquettes incident to \(e\) give the one-link density
\begin{equation}
 \gamma_{e,\beta}(\dd U\mid\eta)
 \propto
 \exp\left\{{\beta\over3}
 \Re\Tr\bigl(U{\cal H}(\eta)\bigr)\right\}
 \lambda_{\SU(3)}(\dd U).                                \label{eq:s6v24-staple-law}
\end{equation}
When \({\cal H}(\eta)=0\), this is Haar probability.  When
\({\cal H}(\eta')=(\omega-1)I\), it is the tilt in
\eqref{eq:s6v24-analytic-Gibbs} with
\(\lambda=\beta\) and the nonconstant analytic function
\eqref{eq:s6v24-F}.  By
\Cref{thm:s6v24-Haar-TV}, the total variation between these two
conditionals tends to one.  Their exteriors differ only at \(f\),
so the defining supremum of \(c_{ef}\) is at least that distance.
\end{proof}

\begin{corollary}[The V20 bare-bound failure is an actual
worst-case failure]
\label{cor:s6v24-actual-failure}
For bare weak-coupling Wilson theory in the exact relative
parity-tree coordinates, the loss of the scalar upper margin in
\Cref{prop:s6v20-bare-no-go} is not merely a defect of the estimate:
the actual worst-case single-link influence has limsup one.
\end{corollary}

\begin{proof}
This is \eqref{eq:s6v24-cef-one}.  It concerns the supremal
Dobrushin coefficient; it does not assert slow mixing under
typical local-vacuum exteriors.
\end{proof}

\begin{theorem}[Stability under uniformly bounded generated
terms]
\label{thm:s6v24-bounded-response}
Retain the two exteriors of
\Cref{cor:s6v24-two-exteriors}.  Add conditional action terms
\(R_{\beta,\eta},R_{\beta,\eta'}\) satisfying
\begin{equation}
 \sup_{\beta\ge1}
 \max\{
 \|R_{\beta,\eta}\|_\infty,
 \|R_{\beta,\eta'}\|_\infty
 \}\le R<\infty.                                         \label{eq:s6v24-bounded-R}
\end{equation}
Then the total variation between the two complete conditionals
still tends to one.
\end{theorem}

\begin{proof}
Under the zero-staple exterior, the density with respect to Haar
is
\[
 {e^{-R_{\beta,\eta}(U)}
  \over\int e^{-R_{\beta,\eta}}\dd\lambda_G}
 \le e^{2R}.
\]
Let \({\cal M}\) be the maximum set of the nonzero-staple
function \(F\).  Given \(\varepsilon>0\), choose an open
\(U\supset{\cal M}\) with
\(\lambda_G(U)<\varepsilon e^{-2R}\).  The zero-staple
conditional assigns \(U\) mass at most \(\varepsilon\).
For the nonzero-staple conditional,
\Cref{lem:s6v23-concentration} with bounded remainder gives mass
tending to one on \(U\).  The total-variation distance is
therefore at least \(1-2\varepsilon\) for all sufficiently large
\(\beta\).  Let \(\varepsilon\downarrow0\).
\end{proof}

\begin{warning}[Sublinear but unbounded terms require the full
leading minimizer card]
\label{warn:s6v24-sublinear}
Uniform boundedness in
\Cref{thm:s6v24-bounded-response} cannot be replaced blindly by
\(\|R_\beta\|_\infty=o(\beta)\) on the zero-staple exterior.
Although sublinear, such a term may itself diverge and concentrate
that conditional on the same set selected by the nonzero Wilson
staple.  Order-\(\beta\) terms and every diverging lower-order term
must therefore be treated at their true scale in the complete
two-exterior concentration problem.
\end{warning}

\begin{assessment}[What is now ruled out]
\label{ass:s6v24-ruled-out}
The pure-Wilson single-link route is closed at weak coupling in the
actual parity-tree gauge, not merely in an isolated plaquette
model.  The obstruction uses admissible reduced coordinates, keeps
the tree and coarse background fixed, and changes exactly one
reduced exterior link.  Any response-generated term which is
uniformly bounded on this card leaves the failure intact.

This does not prove that a response term at a growing scale fails
to repair the conditional, nor that a physical block or
local-vacuum typical-set contraction fails.  Those are now the
only honest conditional-contraction branches.
\end{assessment}

\section{The corrected contraction target}
\label{sec:s6v24-corrected-target}

\begin{theorem}[Necessary escape from the single-link supremum]
\label{thm:s6v24-escape}
On a weak-facing response-fixed trajectory, a WCIC-based proof of
\Cref{thm:s6v22-conditional-mass} must establish at least one of:
\begin{enumerate}[label=\textup{(\roman*)},leftmargin=3.5em]
\item an order-one complete response interaction changes the
      two-exterior leading minimizer geometry of
      \Cref{thm:s6v24-embedded-no-go};
\item the conditional variables are grouped into blocks which
      internalize the six-staple centre rearrangement and have a
      strict inter-block row;
\item the selected local-vacuum proof replaces worst-case
      Dobrushin influence by a rigorously weighted good-exterior
      contraction plus a separately controlled bad-set or contour
      term; or
\item a different reflection-positive correlation estimate
      supplies the delayed-Gram exponent without WCIC.
\end{enumerate}
Bare single-link WCIC is not a fifth branch.
\end{theorem}

\begin{proof}
\Cref{thm:s6v24-embedded-no-go} contradicts a uniform strict
single-link row for the bare leading action, while
\Cref{thm:s6v24-bounded-response} excludes bounded perturbative
repair.  The listed alternatives exhaust changes to the leading
conditional geometry, changes to the conditional carrier,
replacement of the worst-case norm, and replacement of the
conditional-influence method itself.
\end{proof}

\begin{assessment}[Exact live bottleneck after V24]
\label{ass:s6v24-bottleneck}
The weak-coupling single-link Dobrushin route is now rigorously
eliminated for the bare parity-tree Wilson law.  The next
mathematical target is not a sharper single-link oscillation upper
bound.  It is the smallest physical block which internalizes the
explicit centre-staple rearrangement, or a local-vacuum
good/bad-exterior theorem which prices it probabilistically.

The canonical-root unsplit static card remains an independent
upstream gate.  At V25 the compulsory SM and NS audit must be
performed before choosing between block contraction and a
typical-set/contour route.
\end{assessment}

\begin{programme}[V25: companion audit and minimum repairing
block]
\label{prog:s6v24-V25}
\begin{enumerate}[leftmargin=3em]
\item Audit the then-current active SM and NS notes.
\item Determine the minimum parity-tree block containing the
      link \(e\) and all six reduced staple controls in
      \eqref{eq:s6v24-controls}.
\item Write its exact conditional reference measure and complete
      boundary interaction.
\item Test whether changing one neighboring block still produces
      separated conditional minimizer sets.
\item If the bifurcation persists, enlarge once or move to a
      rigorously weighted local-vacuum good/bad decomposition.
\item Preserve the OS time-incidence and source lower-frame rows
      of V22 and the static-card gate of V21.
\end{enumerate}
\end{programme}

\begin{firewall}[Claim boundary after sixth-series V24]
\label{fire:s6v24-boundary}
V24 proves an explicit six-control reduced staple star in the
relative parity-tree gauge, constructs two exteriors differing in
one reduced link with staple sums \(0\) and
\((\omega-1)I\), proves analytic Gibbs tilts separate from Haar,
and concludes that the actual bare reduced-link influence tends
to one as \(\beta\to\infty\).  It also proves stability under
uniformly bounded generated conditional terms and identifies the
four remaining contraction branches.

V24 does not decide an order-growing response interaction, a
minimum repairing block, a typical-set or contour estimate, the
canonical-root static card, continuum occurrence, or the full
Yang--Mills mass gap.
\end{firewall}

\begin{theorem}[Sixth-series V24 result ledger]
\label{thm:s6v24-results}
V24:
\begin{enumerate}[label=\textup{(V24.\arabic*)},leftmargin=6em]
\item constructs an explicit all-odd reduced link and six reduced
      staple controls;
\item constructs a one-coordinate exterior change with zero and
      nonzero total staples;
\item proves proper analytic maximum sets have Haar measure zero;
\item proves a nonconstant compact Gibbs tilt separates from Haar;
\item embeds the Wilson low-temperature influence obstruction in
      the exact parity-tree coordinates;
\item upgrades the V20 bare-bound failure to an actual worst-case
      failure;
\item proves bounded generated terms cannot repair it; and
\item restricts the continuation to response-leading, block,
      good/bad-exterior, or direct reflected-correlation routes.
\end{enumerate}
\end{theorem}

\begin{proof}
Use
\Cref{lem:s6v24-free-star,
cor:s6v24-two-exteriors,
lem:s6v24-analytic-zero,
thm:s6v24-Haar-TV,
lem:s6v24-nonconstant,
thm:s6v24-embedded-no-go,
thm:s6v24-bounded-response,
thm:s6v24-escape}.
\end{proof}

\paragraph{Parent-version record.}
V24 was formed from
\texttt{Renewal\_Yang\_Mills\_Mass\_Gap\_Research\_Notes\_V23.tex}.
The V23 parent had \(10\,071\) logical source lines,
\(358\,694\) bytes, and SHA--256
\[
 \texttt{608309a0886c644417fc03c567a2a6b1bfd3fb34c30db39a49a3951c9274fc15}.
\]
The next compulsory companion audit is V25.

\clearpage
\chapter{Sixth-series V25: companion audit, centre-sector
shorting, and disagreement-path absorption}
\label{ch:s6v25-disagreement}

\section{Compulsory V25 companion audit}
\label{sec:s6v25-audit}

\begin{audit}[Standard Model--Einstein V47]
\label{audit:s6v25-SM}
The inspected active file was
\[
\texttt{Renewal\_SM\_Einstein\_Continuum\_Research\_Notes\_V47.tex},
\]
with \(753\,815\) bytes, \(16\,778\) logical source lines, and
SHA--256
\[
\texttt{5852ccfeb55a3675ebb86d87bc66243c3b03d19eb81bc346232572a494694b80}.
\]
Its V44--V47 chain isolates the normalized constant transverse
subspace at each stratum, adds a self-adjoint rank-\(2q\)
completion of norm \(O(L^{-1})\), cancels the frozen zeroth moment,
and preserves a gap
\(\delta_{\rm path}-O(L^{-1/2})\).  The construction has bounded
physical first moment and shrinking physical range.  It is carried
jointly through faces, triangles, and one-simplex edges by mapping
the complete sector link into one contractible universal path
tree.  Original vertices and Gaussian locality of the growing
lattice-range completion remain open.
\end{audit}

\begin{audit}[Navier--Stokes V81]
\label{audit:s6v25-NS}
The inspected active file was
\[
\texttt{Renewal\_NS\_Stability\_Research\_Notes\_V81.tex},
\]
with \(840\,105\) bytes, \(14\,724\) logical source lines, and
SHA--256
\[
\texttt{f3a83402a09a77f11fd259d1552e6790636238dcdd7931f17fd7c6e6e89447a1}.
\]
It proves the finite-window and ancient Duhamel formulas for its
type-I ancient class, computes the heat-gradient constant
\(c_h=2/\sqrt\pi\), defines the finite Oseen constant \(c_O\),
and derives the explicit nonzero-profile floor
\(\kappa\ge1/(\pi c_O)\).  Its gradient estimate closes only after
the propagated nonlinear term is absorbed under
\[
 4c_O\kappa
 \left({\theta\over1+\theta}\right)^{1/2}<1.
\]
These are epsilon-regularity estimates, not a proof of global
regularity.
\end{audit}

\begin{theorem}[V25 audit decision]
\label{thm:s6v25-audit-decision}
The appropriate continuation after the V24 single-link no-go is:
\begin{enumerate}[label=\textup{(\roman*)},leftmargin=3.5em]
\item isolate the finite-dimensional centre/staple sector before
      estimating within-sector fluctuations;
\item assemble all incident block strata in one conditional law;
      and
\item replace a supremal one-step row by an explicit propagation
      kernel whose bad-sector contribution can be absorbed below a
      strict path threshold.
\end{enumerate}
No SM gap constant or NS regularity constant is transferred to
Yang--Mills.
\end{theorem}

\begin{proof}
Item (i) is the structural role of the normalized constant-sector
completion in \Cref{audit:s6v25-SM}.  Item (ii) is its
assembly-before-gap rule across intersecting strata.  Item (iii)
is the exact propagation-and-absorption lesson of
\Cref{audit:s6v25-NS}.  The three statements select an order of
proof, not a numerical comparison between the programmes.
\end{proof}

\section{Exact separation of sector weights and within-sector
motion}
\label{sec:s6v25-sector-TV}

\begin{lemma}[Total variation under a finite measurable sector
partition]
\label{lem:s6v25-sector-TV}
Let \(X=\bigsqcup_{z\in{\cal Z}}X_z\) be a finite measurable
partition.  Write
\[
 \mu=\sum_zp_z\mu_z,
 \qquad
 \nu=\sum_zq_z\nu_z,
\]
where \(\mu_z,\nu_z\) are probability measures supported on
\(X_z\), with arbitrary choices when the corresponding weight
vanishes.  Then
\begin{equation}
 \boxed{
 \|\mu-\nu\|_{\rm TV}
 \le
 \|p-q\|_{\rm TV}
 \sum_z\min(p_z,q_z)
       \|\mu_z-\nu_z\|_{\rm TV}.}                       \label{eq:s6v25-sector-TV}
\end{equation}
If \(\mu_z=\nu_z\) for every occupied sector, equality reduces to
\[
 \|\mu-\nu\|_{\rm TV}=\|p-q\|_{\rm TV}.
\]
\end{lemma}

\begin{proof}
On \(X_z\), set \(r_z=\min(p_z,q_z)\).  The signed measure
decomposes as
\[
 p_z\mu_z-q_z\nu_z
 =
 r_z(\mu_z-\nu_z)
 +(p_z-q_z)_+\mu_z
 -(q_z-p_z)_+\nu_z.
\]
Take total-variation norms, sum over the disjoint sectors, and
divide by two.  The two unmatched-weight sums combine to
\(\|p-q\|_{\rm TV}\).  If the conditional measures agree, the
matched term vanishes, while disjoint support makes the unmatched
weight bound exact.
\end{proof}

\begin{corollary}[Within-sector contraction cannot repair a
sector-weight jump]
\label{cor:s6v25-weight-jump}
If a sequence of exterior pairs has
\(\|p^{(r)}-q^{(r)}\|_{\rm TV}\to1\), then the corresponding
block conditional total variation tends to one, regardless of any
within-sector coupling.
\end{corollary}

\begin{proof}
Pushforward under the sector map cannot increase total variation,
so
\[
 \|\mu^{(r)}-\nu^{(r)}\|_{\rm TV}
 \ge\|p^{(r)}-q^{(r)}\|_{\rm TV}.
\]
\end{proof}

\begin{assessment}[Reading the V24 witness as a sector failure]
\label{ass:s6v25-V24-sector}
Choose a small neighborhood of the maximum set of
\eqref{eq:s6v24-F} as one sector and its complement as the other.
Under the nonzero-staple conditional its weight tends to one,
while under the zero-staple Haar conditional it can be arbitrarily
small.  Thus V24 is already a sector-weight obstruction.  Merely
proving fast motion inside each centre sector cannot repair it;
the repairing block or local-vacuum restriction must control the
sector weights themselves.
\end{assessment}

\section{A good/bad disagreement exploration}
\label{sec:s6v25-exploration}

Let \({\mathbb G}=(V,E)\) be a locally finite block interaction
graph of maximum degree \(\Delta\ge2\).  Couple two finite-volume
conditional fields whose exterior conditions differ in a source
set \(A\).  Explore the connected disagreement cluster starting
from \(A\).  At each newly exposed oriented edge, declare the
transmission gate either \emph{good} or \emph{bad}.

\begin{definition}[Conditional disagreement absorption card,
CDAC]
\label{def:s6v25-CDAC}
CDAC with constants \(a,p\in[0,1]\) means that one exploration
coupling has the following properties for every revealed history:
\begin{enumerate}[label=\textup{(\roman*)},leftmargin=3.5em]
\item the conditional probability that the next gate is bad is at
      most \(p\);
\item conditional on a good gate, the probability that a
      disagreement crosses it is at most \(a\); and
\item every disagreement at a target set \(B\) is connected to
      \(A\) by an explored self-avoiding path all of whose gates
      transmitted.
\end{enumerate}
Put
\begin{equation}
 \kappa:=a+p,
 \qquad
 \alpha:=(\Delta-1)\kappa.                              \label{eq:s6v25-alpha}
\end{equation}
\end{definition}

\begin{lemma}[Fixed disagreement-path probability]
\label{lem:s6v25-fixed-path}
Under CDAC, every predeclared self-avoiding oriented path
\(\gamma\) of \(n\) gates transmits disagreement with probability
at most
\begin{equation}
 \boxed{\mathbb P(\gamma\text{ transmits})\le\kappa^n.}  \label{eq:s6v25-path-prob}
\end{equation}
No independence of the gates is required.
\end{lemma}

\begin{proof}
At one explored gate, transmission is contained in the union of
the bad event and transmission through a good gate.  Its
conditional probability given the complete revealed history is at
most \(p+a=\kappa\).  Multiply these conditional bounds along the
ordered path by the tower property.
\end{proof}

\begin{theorem}[Subcritical disagreement-path decay]
\label{thm:s6v25-path-decay}
Assume CDAC and
\begin{equation}
 \boxed{\alpha=(\Delta-1)(a+p)<1.}                      \label{eq:s6v25-absorption}
\end{equation}
If \(B\) is at graph distance at least \(r\ge1\) from \(A\), then
\begin{equation}
 \boxed{
 \mathbb P(A\longleftrightarrow B
           \text{ by transmitted disagreement})
 \le
 {|A|\Delta\over\Delta-1}
 {\alpha^r\over1-\alpha}.}                             \label{eq:s6v25-connectivity}
\end{equation}
\end{theorem}

\begin{proof}
From one starting vertex, the number of self-avoiding paths of
length \(n\ge1\) is at most
\(\Delta(\Delta-1)^{n-1}\).  Multiply by \(|A|\), use
\Cref{lem:s6v25-fixed-path}, and sum over \(n\ge r\):
\[
 |A|\sum_{n\ge r}
 \Delta(\Delta-1)^{n-1}\kappa^n
 =
 {|A|\Delta\over\Delta-1}
 \sum_{n\ge r}\alpha^n.
\]
This is \eqref{eq:s6v25-connectivity}.
\end{proof}

\begin{assessment}[Exact analogue of absorption]
\label{ass:s6v25-absorption}
The good-gate transmission \(a\) is the propagated regular term;
the bad-sector probability \(p\) is its exceptional debit.  The
geometric factor \(\Delta-1\) counts the possible continuations.
Condition \eqref{eq:s6v25-absorption} is the exact point at which
the complete propagation contribution can be absorbed.  Estimating
\(a\), \(p\), and the path count in unrelated probability spaces
would not prove CDAC.
\end{assessment}

\section{Covariance and reflected spectral consequences}
\label{sec:s6v25-covariance}

\begin{theorem}[CDAC covariance decay]
\label{thm:s6v25-CDAC-covariance}
Suppose the disagreement exploration is a valid boundary coupling
for the Gibbs specification.  Let bounded \(F,G\) have supports
\(A,B\) at distance \(r\ge1\).  Under
\eqref{eq:s6v25-absorption},
\begin{equation}
 \boxed{
 |\operatorname{Cov}(F,G)|
 \le
 {\Delta\over\Delta-1}
 {\alpha^r\over1-\alpha}
 \left(\sum_{e\in A}\delta_e(F)\right)
 \max_{f\in B}\delta_f(G).}                            \label{eq:s6v25-CDAC-cov}
\end{equation}
\end{theorem}

\begin{proof}
Expose the field in the support of \(F\) and compare the two
conditional expectations of \(G\) produced by two values there.
Under the coupling, the value of \(G\) can differ only if a
disagreement path reaches \(B\); its change is at most the maximum
single-block oscillation in \(B\) times the number of initial
oscillation increments.  Apply
\Cref{thm:s6v25-path-decay} to each initial block and sum the
martingale increments of \(F\).  This is the usual disagreement
coupling proof of the Dobrushin comparison bound, with
\eqref{eq:s6v25-connectivity} in place of the resolvent row.
\end{proof}

\begin{theorem}[CDAC-to-OS mass rate]
\label{thm:s6v25-CDAC-OS}
Assume reflection positivity, a physical OS delay \(\tau>0\), and
the time-incidence condition
\[
 \operatorname{dist}_E(S,\vartheta_{n\tau}S)
 \ge ns_\tau-r_0
\]
for a fixed finite source bank.  Assume also the lower-frame and
oscillation bounds
\eqref{eq:s6v22-lower-frame}--\eqref{eq:s6v22-osc-B}.
If CDAC holds uniformly with
\(\alpha<1\), then the cyclic bank has Hamiltonian gap
\begin{equation}
 \boxed{
 H|_{{\cal C}_F}
 \succeq
 \mu_{\rm DA}I,
 \qquad
 \mu_{\rm DA}:=
 {s_\tau\over\tau}|\log\alpha|.}                       \label{eq:s6v25-DA-mass}
\end{equation}
\end{theorem}

\begin{proof}
Insert \(r=ns_\tau-r_0\) in
\eqref{eq:s6v25-CDAC-cov}.  The finite collar \(r_0\), source
oscillations, and lower frame contribute only an \(n\)-independent
prefactor.  Reflection positivity identifies the resulting
covariance with the positive delayed Gram as in
\eqref{eq:s6v22-OS-cov}.  Apply
\Cref{thm:s6v22-Gram-gap} with
\(q=\alpha^{s_\tau}\).
\end{proof}

\begin{corollary}[Dense-bank conditional mass handoff]
\label{cor:s6v25-dense-handoff}
On the response-fixed/local-vacuum continuum branch of
\Cref{thm:s6v22-conditional-mass}, replace WCIC by a
regulator-uniform CDAC with \(\alpha<1\).  If the same continuum,
source-frame, density, and vacuum hypotheses hold, then
\begin{equation}
 H_{\rm YM}|_{\Omega^\perp}
 \succeq\mu_{\rm DA}I.                                 \label{eq:s6v25-dense-mass}
\end{equation}
\end{corollary}

\begin{proof}
\Cref{thm:s6v25-CDAC-OS} gives the same exponent on every fixed
finite bank.  Delayed-Gram convergence transfers it to the
continuum.  Apply \Cref{cor:s6v22-bank-growth} to the dense
noncollapsing local algebra.
\end{proof}

\section{The centre-short block card}
\label{sec:s6v25-centre-short}

\begin{definition}[Centre-short block card, CSBC]
\label{def:s6v25-CSBC}
For a physical parity-tree block \(B\), CSBC consists of:
\begin{enumerate}[label=\textup{(\roman*)},leftmargin=3.5em]
\item one finite measurable sector map
      \(\pi_B:X_B\to{\cal Z}_B\) containing every centre/staple
      ground-state label exposed by boundary interactions;
\item the exact exterior-dependent sector weights and
      within-sector conditional laws;
\item one exploration coupling assembled from their common
      interaction strata;
\item uniform good-gate transmission \(a\) after the sector label
      is matched;
\item uniform conditional bad-sector probability \(p\); and
\item the strict absorption
      \((\Delta-1)(a+p)<1\).
\end{enumerate}
\end{definition}

\begin{theorem}[What a repairing block must prove]
\label{thm:s6v25-repairing-block}
A block which merely contains the six controls of
\Cref{lem:s6v24-free-star} does not yet close the conditional
route.  To yield exponential covariance decay it must additionally
prove either:
\begin{enumerate}[label=\textup{(\alph*)},leftmargin=3.5em]
\item a direct strict inter-block Dobrushin row; or
\item CSBC and its CDAC exploration.
\end{enumerate}
In the second branch the sector-weight term of
\eqref{eq:s6v25-sector-TV} and the within-sector term are both
load-bearing.
\end{theorem}

\begin{proof}
Internalizing the V24 controls removes that particular
one-coordinate exterior pair, but it gives no estimate for new
boundary interactions of the enlarged block.  The exact
sector decomposition \eqref{eq:s6v25-sector-TV} shows that
within-sector contraction alone leaves exterior-dependent sector
weights uncontrolled.  Direct inter-block Dobrushin contraction
closes by V22.  Alternatively, CSBC closes by
\Cref{thm:s6v25-path-decay,thm:s6v25-CDAC-covariance}.
\end{proof}

\begin{warning}[CSBC is not an action modification]
\label{warn:s6v25-not-modification}
The low-rank completion in the SM audit motivates isolating a
finite-dimensional dangerous sector, but CSBC does not add a term
to the Yang--Mills action.  It is an exact decomposition and
coupling of the occurring conditional law.  Changing sector
weights to obtain absorption would define a different theory.
\end{warning}

\begin{assessment}[Exact live bottleneck after V25]
\label{ass:s6v25-bottleneck}
The impossible single-link supremum has been replaced by a precise
probability-weighted route.  A finite centre-sector partition
separates exterior phase selection from within-sector motion, and
CDAC converts uniform good-gate transmission plus a conditional
bad-sector probability into exponential paths whenever
\((\Delta-1)(a+p)<1\).  The resulting exponent feeds the OS
spectral theorem without returning to single-link WCIC.

No supplied note constructs the actual response-fixed CSBC,
proves a bad-sector contour probability, or supplies its
within-sector coupling.  The canonical-root static card and
continuum landing also remain open.
\end{assessment}

\begin{programme}[V26: construct the first physical CSBC]
\label{prog:s6v25-V26}
\begin{enumerate}[leftmargin=3em]
\item Let \(B\) be the smallest parity block containing the V24
      link and all six control links; enumerate every plaquette
      crossing its boundary.
\item Define the centre/staple sector map from the complete block
      minimizer components, not from a fitted centre label.
\item Prove the exact sector-weight/within-sector decomposition of
      every neighboring-block conditional.
\item On the selected local-vacuum branch, derive a conditional
      contour or energy bound for the bad-sector probability
      \(p\).
\item Couple the matched good sectors and compute the transmission
      constant \(a\).
\item Compare the assembled value
      \((\Delta-1)(a+p)\) with one.  Enlarge the block only if the
      failed term has a demonstrated boundary origin.
\end{enumerate}
\end{programme}

\begin{firewall}[Claim boundary after sixth-series V25]
\label{fire:s6v25-boundary}
V25 performs the compulsory SM V47 and NS V81 audit, proves the
finite-sector total-variation decomposition, proves the
fixed-path and subcritical disagreement bounds without an
independence assumption, derives CDAC covariance decay, transfers
its explicit exponent to an OS finite-bank and dense-bank mass
rate, and defines the target-native centre-short block card.

V25 does not construct the physical CSBC, populate \(a\) or
\(p\), prove a response-fixed contour estimate, close the static
card, construct the continuum measure, or prove the full
Yang--Mills mass gap.
\end{firewall}

\begin{theorem}[Sixth-series V25 result ledger]
\label{thm:s6v25-results}
V25:
\begin{enumerate}[label=\textup{(V25.\arabic*)},leftmargin=6em]
\item audits SM V47 and NS V81;
\item proves an exact sector-weight/within-sector TV bound;
\item proves sector jumps survive perfect within-sector motion;
\item defines CDAC;
\item proves a conditional fixed-path product bound;
\item proves the subcritical path estimate;
\item derives covariance decay;
\item derives the OS mass rate
      \(s_\tau|\log((\Delta-1)(a+p))|/\tau\);
\item proves the dense-bank conditional handoff; and
\item identifies the complete CSBC needed from a repairing block.
\end{enumerate}
\end{theorem}

\begin{proof}
Use
\Cref{audit:s6v25-SM,
audit:s6v25-NS,
thm:s6v25-audit-decision,
lem:s6v25-sector-TV,
cor:s6v25-weight-jump,
lem:s6v25-fixed-path,
thm:s6v25-path-decay,
thm:s6v25-CDAC-covariance,
thm:s6v25-CDAC-OS,
cor:s6v25-dense-handoff,
thm:s6v25-repairing-block}.
\end{proof}

\paragraph{Parent-version record.}
V25 was formed from
\texttt{Renewal\_Yang\_Mills\_Mass\_Gap\_Research\_Notes\_V24.tex}.
The V24 parent had \(10\,504\) logical source lines,
\(374\,753\) bytes, and SHA--256
\[
 \texttt{f08520581209f4b16857620beee561c42958a166288c0c343c44729c0980c08e}.
\]
The next compulsory companion audit is V30.

\clearpage
\chapter{Sixth-series V26: plaquette-star shielding and the
multi-defect suppression card}
\label{ch:s6v26-multidefect}

\section{Two propagation graphs and their distinct thresholds}
\label{sec:s6v26-two-graphs}

Let \({\mathbb G}_{\rm link}\) be the reduced-link interaction
graph: two reduced links are adjacent if they occur in a common
fine plaquette.  Let \({\mathbb G}_{\rm plaq}\) be the plaquette
graph: two plaquettes are adjacent if they share an integrated
fine link.

\begin{lemma}[Link and plaquette degree ledger]
\label{lem:s6v26-degrees}
In four dimensions,
\begin{equation}
 \boxed{
 \Delta_{\rm link}\le18,
 \qquad
 \Delta_{\rm plaq}\le20.}                               \label{eq:s6v26-degrees}
\end{equation}
Consequently the CDAC self-avoiding continuation factors are,
respectively, at most \(17\) and \(19\).
\end{lemma}

\begin{proof}
By \Cref{cor:s6v21-Wilson-incidence}, a link is incident to at
most six plaquettes and each contributes at most three other
links, giving \(\Delta_{\rm link}\le18\).  The previously audited
plaquette geometry gives
\(\Delta_{\rm plaq}\le4(6-1)=20\), recorded in
\eqref{eq:s6v14-degree20}.  After the first step, a self-avoiding
path has at most \(\Delta-1\) forward choices.
\end{proof}

\begin{corollary}[Numerical CDAC absorption rows]
\label{cor:s6v26-absorption-rows}
For reduced-link disagreement propagation, the sufficient CDAC
condition is
\begin{equation}
 \boxed{17(a+p)<1.}                                     \label{eq:s6v26-link-absorption}
\end{equation}
For a plaquette-contour exploration with the same gate constants,
the corresponding sufficient condition is
\begin{equation}
 \boxed{19(a+p)<1.}                                     \label{eq:s6v26-plaquette-absorption}
\end{equation}
\end{corollary}

\begin{proof}
Insert the two degree bounds from
\Cref{lem:s6v26-degrees} into
\eqref{eq:s6v25-absorption}.
\end{proof}

\begin{warning}[Do not mix the two degrees]
\label{warn:s6v26-carriers}
The degree twenty bound counts neighboring plaquette defects; the
degree eighteen bound counts variables in a one-link conditional.
Using twenty in a link row is valid but wasteful.  Using eighteen
in a plaquette-animal count is unsupported.
\end{warning}

\section{Concentric plaquette-star shielding}
\label{sec:s6v26-stars}

For a reduced link \(e\), let
\[
 B_R(e):=\{f:
 \operatorname{dist}_{{\mathbb G}_{\rm link}}(e,f)\le R\},
 \qquad R\ge0.
\]
The physical fine-link star contains fixed tree and retained links
as parameters, but \(B_R(e)\) lists only integrated reduced
variables.

\begin{lemma}[The first star internalizes the V24 rearrangement]
\label{lem:s6v26-first-star}
Every reduced variable occurring in a Wilson plaquette with \(e\)
belongs to \(B_1(e)\).  Thus conditioning jointly on the variables
of \(B_1(e)\) makes the six-staple rearrangement of
\Cref{cor:s6v24-two-exteriors} internal to the block: no exterior
coordinate can directly change a plaquette containing \(e\).
Moreover,
\begin{equation}
 |B_1(e)|\le19.                                         \label{eq:s6v26-star-size}
\end{equation}
\end{lemma}

\begin{proof}
This is the definition of the link interaction graph.  The central
link contributes one vertex and has at most eighteen distinct
neighbors by \Cref{lem:s6v26-degrees}.
\end{proof}

\begin{lemma}[Exact Markov collar for a deep core]
\label{lem:s6v26-Markov-collar}
Let \(R\ge1\).  Every Wilson plaquette containing a variable in
\(B_{R-1}(e)\) has all of its reduced variables in \(B_R(e)\).
Consequently, for two exterior conditions on
\({\mathscr E}^{\rm red}_N\setminus B_R(e)\), the ratio of the two
conditional densities on \(B_R(e)\) depends only on the boundary
layer
\[
 \partial B_R(e):=B_R(e)\setminus B_{R-1}(e).
\]
An exterior disagreement can affect the marginal at \(e\) only
through a link-graph path of at least \(R\) transmissions.
\end{lemma}

\begin{proof}
If a plaquette contains \(f\in B_{R-1}(e)\), every other reduced
link in that plaquette is adjacent to \(f\), hence lies within
distance \(R\) of \(e\).  Interactions wholly inside the block are
identical under the two exterior conditions.  Interactions meeting
the exterior cannot contain a variable of \(B_{R-1}(e)\), so their
block dependence is supported in the boundary layer.  The
common-incidence cancellation
\eqref{eq:s6v21-common-incidence} then forces any effect at \(e\)
to traverse at least \(R\) adjacency steps.
\end{proof}

\begin{theorem}[Shielded-core influence under CDAC]
\label{thm:s6v26-shield}
Suppose reduced-link CDAC holds with
\(\alpha_{\rm link}:=17(a+p)<1\).  Couple two conditional laws on
\(B_R(e)\) with arbitrary exterior conditions.  Then
\begin{equation}
 \boxed{
 \mathbb P(\text{the two core values at }e\text{ disagree})
 \le
 {18\over17}
 {\alpha_{\rm link}^R\over1-\alpha_{\rm link}}.}         \label{eq:s6v26-shield}
\end{equation}
\end{theorem}

\begin{proof}
By \Cref{lem:s6v26-Markov-collar}, a transmitted disagreement
must follow a path of length at least \(R\) from the boundary
toward \(e\).  Reverse the path and apply
\Cref{thm:s6v25-path-decay} with one starting vertex and
\(\Delta=18\).  Reversal preserves self-avoidance and gate
transmission.
\end{proof}

\begin{assessment}[What blocking really buys]
\label{ass:s6v26-blocking}
The first plaquette star removes the direct V24 exterior switch,
and each additional star adds one compulsory transmission before
the exterior reaches the core.  This is a geometric shielding
statement, not a claim that the full block conditional has small
total variation.  A low-temperature boundary may still select a
whole-block sector.  CDAC or a comparable contour theorem is what
converts geometric depth into a small core influence.
\end{assessment}

\section{The available one-plaquette defect estimate}
\label{sec:s6v26-one-point}

For a plaquette \(q\), write its Wilson defect
\begin{equation}
 d_q(U):=1-\frac13\Re\Tr U_q\ge0.                       \label{eq:s6v26-defect}
\end{equation}
The attached reduced-fibre programme proves, for every corner
plaquette, every \(M\ge4\), and pure Wilson
\(\beta\ge240\),
\begin{equation}
 \mathbb E_{\beta,0,M}d_q
 \le {C_{\rm W}\over\beta},                             \label{eq:s6v26-imported-mean}
\end{equation}
where
\begin{equation}
 C_{\rm W}
 :=
 {15\over4}\log\left(
 {50\over3}e^{43/32}\pi^{15}960^4
 \right)<201.                                           \label{eq:s6v26-CW}
\end{equation}
For the symmetric generated law its stated extension is
\begin{equation}
 \mathbb E_{\beta,\Psi,M}d_q
 \le
 {C_{\rm W}\over\beta}
 +{\operatorname{osc}\Psi_M\over12\beta M^4}.           \label{eq:s6v26-generated-mean}
\end{equation}

\begin{corollary}[One-point bad-plaquette probability]
\label{cor:s6v26-one-point}
For \(\varepsilon>0\), let
\[
 {\cal B}_q(\varepsilon):=\{d_q\ge\varepsilon\}.
\]
Then
\begin{align}
 \mathbb P_{\beta,0,M}({\cal B}_q(\varepsilon))
 &\le {C_{\rm W}\over\beta\varepsilon},                 \label{eq:s6v26-pure-bad}\\
 \mathbb P_{\beta,\Psi,M}({\cal B}_q(\varepsilon))
 &\le {1\over\varepsilon}
 \left(
 {C_{\rm W}\over\beta}
 +{\operatorname{osc}\Psi_M\over12\beta M^4}
 \right).                                               \label{eq:s6v26-generated-bad}
\end{align}
\end{corollary}

\begin{proof}
Apply Markov's inequality to the nonnegative random variable
\eqref{eq:s6v26-defect} and use
\eqref{eq:s6v26-imported-mean} or
\eqref{eq:s6v26-generated-mean}.
\end{proof}

\begin{theorem}[One-point rarity does not imply path suppression]
\label{thm:s6v26-one-point-no-path}
No estimate of the form
\[
 \sup_q\mathbb P({\cal B}_q)\le p
\]
with \(0<p<1\) implies
\[
 \mathbb P\left(\bigcap_{q\in\Gamma}{\cal B}_q\right)
 \le p^{|\Gamma|}
\]
for connected plaquette sets \(\Gamma\).
\end{theorem}

\begin{proof}
Let \(Y\) be one Bernoulli variable of parameter \(p\), and set
\(\mathbf1_{{\cal B}_q}=Y\) for every plaquette \(q\).  Each
one-point probability equals \(p\), but every nonempty
intersection has probability \(p\), rather than \(p^{|\Gamma|}\).
\end{proof}

\begin{warning}[The imported mean cannot be multiplied]
\label{warn:s6v26-no-multiplication}
The explicit value in \eqref{eq:s6v26-pure-bad} is a useful
candidate scale for the CDAC bad gate.  It is not the CDAC
conditional probability and cannot be inserted into
\eqref{eq:s6v26-link-absorption} without a multi-defect or
exploration-conditional theorem.  Doing so would assume the
missing decorrelation.
\end{warning}

\section{A multi-defect suppression card}
\label{sec:s6v26-MDSC}

\begin{definition}[Multi-defect suppression card, MDSC]
\label{def:s6v26-MDSC}
Fix an exploration and bad-gate indicators
\((B_1,\ldots,B_n)\).  MDSC with parameter \(p\in[0,1]\) means
that for every path length \(n\) and every
\(J\subseteq\{1,\ldots,n\}\),
\begin{equation}
 \boxed{
 \mathbb P(B_j=1\text{ for all }j\in J)
 \le p^{|J|}.}                                          \label{eq:s6v26-MDSC}
\end{equation}
In addition, conditional on the complete bad-gate vector, the
coupling probability that every good gate on a predeclared path
transmits is at most \(a^{N_{\rm good}}\).
\end{definition}

\begin{theorem}[MDSC implies the CDAC path product]
\label{thm:s6v26-MDSC-path}
Under MDSC,
\begin{equation}
 \boxed{
 \mathbb P(\gamma\text{ transmits})
 \le(a+p)^{|\gamma|}}                                   \label{eq:s6v26-MDSC-path}
\end{equation}
for every predeclared self-avoiding path \(\gamma\).  Therefore
\eqref{eq:s6v26-link-absorption} or
\eqref{eq:s6v26-plaquette-absorption}, on the appropriate graph,
gives the covariance and OS conclusions of V25.
\end{theorem}

\begin{proof}
If \(J\) is the set of bad positions on a path of length \(n\),
bad gates are charged by one and good transmissions by \(a\).
Thus the transmission indicator has conditional expectation at
most
\[
 \prod_{j=1}^n\{a+\mathbf1_{\{B_j=1\}}\}.
\]
Expand this product and use \eqref{eq:s6v26-MDSC}:
\[
 \mathbb E\prod_{j=1}^n\{a+\mathbf1_{\{B_j=1\}}\}
 =
 \sum_{J\subseteq\{1,\ldots,n\}}
 a^{n-|J|}
 \mathbb P(B_j=1\ \forall j\in J)
 \le
 \sum_Ja^{n-|J|}p^{|J|}
 =(a+p)^n.
\]
The remaining claims are
\Cref{thm:s6v25-path-decay,
thm:s6v25-CDAC-covariance,
thm:s6v25-CDAC-OS}.
\end{proof}

\begin{corollary}[Explicit conditional threshold if the imported
scale tensorizes]
\label{cor:s6v26-formal-threshold}
If a pure-Wilson MDSC is proved with
\[
 p={C_{\rm W}\over\beta\varepsilon}
\]
and a matched-good-sector coupling has transmission \(a<1/17\)
on the reduced-link graph, then CDAC holds whenever
\begin{equation}
 \boxed{
 \beta>
 {17C_{\rm W}\over
  \varepsilon(1-17a)}.}                                \label{eq:s6v26-beta-threshold}
\end{equation}
\end{corollary}

\begin{proof}
The displayed condition is exactly
\[
 17\left(a+{C_{\rm W}\over\beta\varepsilon}\right)<1.
\]
Apply \Cref{thm:s6v26-MDSC-path}.
\end{proof}

\begin{firewall}[The threshold is conditional]
\label{fire:s6v26-threshold}
Equation \eqref{eq:s6v26-beta-threshold} is not a proved
weak-coupling Yang--Mills region.  It records the exact numerical
payoff of an MDSC which has not been established.  The one-point
bound \eqref{eq:s6v26-pure-bad} alone does not supply it.
\end{firewall}

\section{A measure-comparison route to MDSC}
\label{sec:s6v26-erasure}

\begin{definition}[Plaquette-defect erasure card, PDEC]
\label{def:s6v26-PDEC}
For each finite declared bad plaquette set
\(\Gamma\), let
\[
 {\cal E}_\Gamma
 :=
 \bigcap_{q\in\Gamma}{\cal B}_q(\varepsilon).
\]
PDEC with constants \(\sigma>0\) and \(K\ge1\) consists of a
measurable map
\[
 T_\Gamma:{\cal E}_\Gamma\to X
\]
such that:
\begin{enumerate}[label=\textup{(\roman*)},leftmargin=3.5em]
\item the complete action satisfies
      \[
       {\cal S}(U)\ge{\cal S}(T_\Gamma U)+
       \sigma|\Gamma|;
      \]
\item the reference-measure pushforward obeys
      \[
       (T_\Gamma)_*
       (\mathbf1_{{\cal E}_\Gamma}m)
       \le K^{|\Gamma|}m;
      \]
\item the map preserves the block constraint, tree gauge,
      declared boundary condition, and response/local-vacuum
      sector used by the conditional exploration.
\end{enumerate}
\end{definition}

\begin{theorem}[Erasure gives multi-defect suppression]
\label{thm:s6v26-erasure}
PDEC implies
\begin{equation}
 \boxed{
 \nu({\cal E}_\Gamma)
 \le(Ke^{-\sigma})^{|\Gamma|}.}                         \label{eq:s6v26-erasure-bound}
\end{equation}
Thus, when the bad gates are the declared plaquette defects and
PDEC is valid conditionally along the exploration, MDSC holds with
\begin{equation}
 p=Ke^{-\sigma}.                                        \label{eq:s6v26-erasure-p}
\end{equation}
\end{theorem}

\begin{proof}
On \({\cal E}_\Gamma\),
\[
 e^{-{\cal S}(U)}
 \le e^{-\sigma|\Gamma|}
 e^{-{\cal S}(T_\Gamma U)}.
\]
Integrate and use the pushforward domination:
\[
 \int_{{\cal E}_\Gamma}e^{-{\cal S}}\dd m
 \le
 e^{-\sigma|\Gamma|}
 K^{|\Gamma|}
 \int_Xe^{-{\cal S}}\dd m.
\]
Divide by the partition function.  The conditional version has
the same proof on the conditional fibre because item (iii) fixes
the exterior and the physical sector.
\end{proof}

\begin{corollary}[Plaquette-animal summability]
\label{cor:s6v26-animal-sum}
If PDEC holds on the plaquette graph and
\begin{equation}
 400Ke^{-\sigma}<1,                                     \label{eq:s6v26-animal-condition}
\end{equation}
then the total probability of an erased connected bad-plaquette
animal of size at least \(n\) containing a fixed plaquette is
bounded by a geometric tail
\begin{equation}
 {Ke^{-\sigma}
  (400Ke^{-\sigma})^{n-1}
  \over1-400Ke^{-\sigma}}.                              \label{eq:s6v26-animal-tail}
\end{equation}
\end{corollary}

\begin{proof}
By \Cref{lem:s6v14-animal}, the number of connected plaquette sets
of cardinality \(j\) containing a fixed root is at most
\(400^{j-1}\).  The union bound and
\eqref{eq:s6v26-erasure-bound} therefore give
\[
 \sum_{j\ge n}400^{j-1}(Ke^{-\sigma})^j
 =
 {Ke^{-\sigma}
  (400Ke^{-\sigma})^{n-1}
  \over1-400Ke^{-\sigma}}.
\]
\end{proof}

\begin{assessment}[First missing physical inequality]
\label{ass:s6v26-first-missing}
The one-plaquette expectation, link and plaquette degrees, star
shielding, CDAC combinatorics, and erasure-to-MDSC implication are
now exact.  The first missing physical estimate is a
gauge-compatible conditional erasure map, or an equivalent
multi-plaquette contour inequality, for the response-fixed
local-vacuum action.  It must lower the complete action by a
linear amount without leaving the block constraint or silently
changing the response sector.
\end{assessment}

\begin{programme}[V27: construct or obstruct conditional defect
erasure]
\label{prog:s6v26-V27}
\begin{enumerate}[leftmargin=3em]
\item Freeze a connected bad plaquette set in the first
      parity-star collar and its actual exterior condition.
\item Attempt an axial/tree-ordered link replacement which reduces
      each declared defect while preserving the block constraint.
\item Compute the exact Haar Jacobian or pushforward multiplicity
      \(K^{|\Gamma|}\).
\item Prove a complete-action drop
      \(\sigma|\Gamma|\), including response-generated terms, or
      return the first overlap/Bianchi obstruction.
\item If PDEC closes, combine
      \eqref{eq:s6v26-erasure-p} with a within-good-sector coupling
      and test \eqref{eq:s6v26-link-absorption}.
\item Do not substitute the one-point mean for the missing
      multi-defect estimate.
\end{enumerate}
\end{programme}

\begin{firewall}[Claim boundary after sixth-series V26]
\label{fire:s6v26-boundary}
V26 distinguishes the link and plaquette propagation degrees,
derives their exact CDAC thresholds, proves concentric
plaquette-star shielding, imports the rigorous all-volume
one-plaquette defect mean, proves that one-point rarity cannot be
multiplied, defines MDSC and PDEC, and proves the exact
erasure-to-path and animal bounds.

V26 does not prove PDEC or MDSC for Yang--Mills, populate the
within-good-sector transmission, close the canonical-root static
card, construct the continuum measure, or prove the full mass gap.
\end{firewall}

\begin{theorem}[Sixth-series V26 result ledger]
\label{thm:s6v26-results}
V26:
\begin{enumerate}[label=\textup{(V26.\arabic*)},leftmargin=6em]
\item separates the degree-eighteen link graph from the
      degree-twenty plaquette graph;
\item computes the thresholds \(17(a+p)<1\) and
      \(19(a+p)<1\);
\item proves exact plaquette-star shielding;
\item derives the one-point bad-plaquette probability;
\item proves its insufficiency for paths;
\item defines the multi-defect suppression card;
\item proves MDSC gives the CDAC path product;
\item gives the conditional explicit beta threshold;
\item defines the gauge-compatible defect-erasure card; and
\item proves its multi-defect and connected-animal consequences.
\end{enumerate}
\end{theorem}

\begin{proof}
Use
\Cref{lem:s6v26-degrees,
cor:s6v26-absorption-rows,
lem:s6v26-first-star,
lem:s6v26-Markov-collar,
thm:s6v26-shield,
cor:s6v26-one-point,
thm:s6v26-one-point-no-path,
thm:s6v26-MDSC-path,
thm:s6v26-erasure,
cor:s6v26-animal-sum}.
\end{proof}

\paragraph{Parent-version record.}
V26 was formed from
\texttt{Renewal\_Yang\_Mills\_Mass\_Gap\_Research\_Notes\_V25.tex}.
The V25 parent had \(10\,996\) logical source lines,
\(392\,293\) bytes, and SHA--256
\[
 \texttt{8d2598532555a9239fea347533e57ae802f43d4dec214f065bb8f91539d8fb08}.
\]
The next compulsory companion audit is V30.

\clearpage
\chapter{Sixth-series V27: relative centre-flux erasure and a
true Peierls sector}
\label{ch:s6v27-centre-erasure}

\section{Small centre tubes force a closed \(\mathbb Z_3\)
plaquette field}
\label{sec:s6v27-centre-cocycle}

Equip \(\SU(3)\) with Frobenius distance and write
\[
 {\cal Z}_3=\{1,\omega,\omega^2\},
 \qquad \omega=e^{2\pi i/3}.
\]
For \(0<\delta<1/2\), say that a plaquette holonomy \(U_p\)
has centre label \(z_p\in{\cal Z}_3\) when
\begin{equation}
 \|U_p-z_pI\|_F\le\delta.                               \label{eq:s6v27-centre-tube}
\end{equation}
The tubes are disjoint because distinct centre matrices have
Frobenius distance \(3\).

\begin{lemma}[Quantitative centre-label Bianchi closure]
\label{lem:s6v27-Bianchi}
Suppose every oriented face \(p\) of an elementary three-cube
\(c\) obeys \eqref{eq:s6v27-centre-tube}, after parallel transport
to a common base point.  With the face orientations induced by
\(\partial c\),
\begin{equation}
 \boxed{\prod_{p\subset\partial c}z_p=1.}               \label{eq:s6v27-centre-Bianchi}
\end{equation}
Thus the small-tube centre labels form a closed
\({\mathbb Z}_3\)-valued plaquette two-cochain.
\end{lemma}

\begin{proof}
The lattice non-Abelian Bianchi identity orders the six based,
oriented face holonomies \(V_1,\ldots,V_6\) so that
\[
 V_1\cdots V_6=I;
\]
this follows directly by expanding the plaquette words and
cancelling every oriented cube edge with its inverse after the
base-point transports.  Conjugation preserves Frobenius distance,
so \(\|V_j-z_jI\|_F\le\delta\).

For unitary matrices, telescoping the two products gives
\[
 \left\|
 V_1\cdots V_6-(z_1\cdots z_6)I
 \right\|_F
 \le\sum_{j=1}^6\|V_j-z_jI\|_F
 \le6\delta.
\]
If \(z_1\cdots z_6\ne1\), its distance from \(I\) is \(3\).
Since \(6\delta<3\), this contradicts the preceding inequality and
the exact Bianchi product.
\end{proof}

\begin{corollary}[Centre flux is a dual closed surface]
\label{cor:s6v27-dual-surface}
The support
\[
 \Sigma_z:=\{p:z_p\ne1\}
\]
has zero \({\mathbb Z}_3\) boundary at every elementary cube.
Equivalently, with the dual-cell convention, its labelled
components are closed codimension-two sheets, possibly with
nontrivial winding on a torus.
\end{corollary}

\begin{proof}
Equation \eqref{eq:s6v27-centre-Bianchi} is exactly the cubical
cocycle equation.  Cubical duality turns a closed plaquette
two-cochain into a closed dual two-chain.  On a torus a closed
cochain need not be exact, which is the stated winding
qualification.
\end{proof}

\section{Relative exactness and the Haar-preserving erasure map}
\label{sec:s6v27-relative-erasure}

Let \(K\) be a finite cubical region with fixed tree and retained
links.  A centre plaquette cochain \(z\) is called
\emph{detail-exact} when there is a centre link cochain \(k\)
such that
\begin{equation}
 k_e=1\quad\hbox{on every fixed tree, retained, and exterior link},
 \qquad
 (\dd k)_p=z_p^{-1}.                                    \label{eq:s6v27-detail-exact}
\end{equation}
Here \(\dd\) is the multiplicative cubical coboundary.

\begin{lemma}[Cohomological meaning of the detail condition]
\label{lem:s6v27-relative-cohomology}
A closed centre flux is detail-exact precisely when its class
vanishes in the relative cohomology group determined by the fixed
link subcomplex.  In particular, a contractible flux supported
inside a cubical ball and trivial on its boundary is detail-exact.
\end{lemma}

\begin{proof}
Equation \eqref{eq:s6v27-detail-exact} asks whether the closed
two-cochain \(z^{-1}\) lies in the coboundary range of one-cochains
which vanish on the fixed subcomplex.  This is exactly vanishing
of its relative cohomology class.  For a cubical ball with trivial
boundary data the relevant degree-two class vanishes, so the
relative cochain complex supplies \(k\).
\end{proof}

\begin{definition}[Centre-flux erasure]
\label{def:s6v27-erasure}
For a detail-exact cochain choose one \(k\) as in
\eqref{eq:s6v27-detail-exact} and define
\begin{equation}
 (T_kU)_e:=k_eU_e.                                      \label{eq:s6v27-Tk}
\end{equation}
\end{definition}

\begin{theorem}[Exact fibre, Haar, and plaquette action of the
erasure]
\label{thm:s6v27-erasure-map}
The map \(T_k\) is a bijection of the same reduced detail fibre,
preserves normalized product Haar measure, and leaves the tree,
retained background, and exterior links unchanged.  Its plaquette
holonomies obey
\begin{equation}
 (T_kU)_p=(\dd k)_pU_p.                                 \label{eq:s6v27-plaquette-transform}
\end{equation}
Thus plaquettes outside \(\Sigma_z\) are unchanged and, on
\(\Sigma_z\),
\[
 (T_kU)_p=z_p^{-1}U_p.
\]
\end{theorem}

\begin{proof}
Multiplication by a fixed centre element is a bijection of
\(\SU(3)\) preserving Haar measure.  The product over reduced
links therefore preserves product Haar.  The first condition in
\eqref{eq:s6v27-detail-exact} preserves every fixed coordinate.
In a plaquette word the four centre multipliers commute with the
link matrices and combine with orientation signs into the
coboundary \((\dd k)_p\), proving
\eqref{eq:s6v27-plaquette-transform}.
\end{proof}

\section{A uniform Wilson action drop}
\label{sec:s6v27-action-drop}

Recall
\[
 d(U):=1-\frac13\Re\Tr U.
\]

\begin{lemma}[Nontrivial-centre correction gains a fixed amount]
\label{lem:s6v27-centre-drop}
If \(z\in\{\omega,\omega^2\}\) and
\(\|U-zI\|_F\le\delta\), then
\begin{align}
 d(U)&\ge\frac32-{\delta\over\sqrt3},                    \label{eq:s6v27-before}\\
 d(z^{-1}U)&\le{\delta\over\sqrt3}.                      \label{eq:s6v27-after}
\end{align}
Consequently
\begin{equation}
 \boxed{
 d(U)-d(z^{-1}U)
 \ge
 \sigma_0(\delta):=
 \frac32-{2\delta\over\sqrt3}>0.}                       \label{eq:s6v27-sigma0}
\end{equation}
\end{lemma}

\begin{proof}
Cauchy--Schwarz for the Frobenius inner product gives
\[
 |\Tr A|\le\sqrt3\|A\|_F.
\]
Since \(\Re z=-1/2\),
\[
 {1\over3}\Re\Tr U
 \le-\frac12+{\delta\over\sqrt3},
\]
which proves \eqref{eq:s6v27-before}.  Also
\(\|z^{-1}U-I\|_F\le\delta\), so
\[
 {1\over3}\Re\Tr(z^{-1}U)
 \ge1-{\delta\over\sqrt3},
\]
giving \eqref{eq:s6v27-after}.  Subtract.
\end{proof}

\begin{theorem}[Relative centre-flux PDEC]
\label{thm:s6v27-centre-PDEC}
Let \(\Sigma\) be a connected detail-exact nontrivial centre-flux
sheet and let \({\cal E}_{\Sigma,\delta}\) be the event that every
\(p\in\Sigma\) lies in the \(\delta\)-tube of its assigned
nontrivial centre label.  Consider the complete action
\begin{equation}
 {\cal S}(U)=\beta\sum_pd(U_p)+\Psi(U).                  \label{eq:s6v27-complete-action}
\end{equation}
Assume the generated action has the erasure debit
\begin{equation}
 |\Psi(U)-\Psi(T_kU)|
 \le\kappa_\Psi|\Sigma|
 \qquad(U\in{\cal E}_{\Sigma,\delta}).                  \label{eq:s6v27-Psi-debit}
\end{equation}
Then
\begin{equation}
 {\cal S}(U)-{\cal S}(T_kU)
 \ge
 \sigma_\beta|\Sigma|,
 \qquad
 \sigma_\beta:=
 \beta\sigma_0(\delta)-\kappa_\Psi.                     \label{eq:s6v27-total-drop}
\end{equation}
If \(\sigma_\beta>0\), PDEC holds with
\[
 K=1,\qquad \sigma=\sigma_\beta,
\]
and
\begin{equation}
 \boxed{
 \nu({\cal E}_{\Sigma,\delta})
 \le e^{-\sigma_\beta|\Sigma|}.}                        \label{eq:s6v27-sheet-probability}
\end{equation}
\end{theorem}

\begin{proof}
By \Cref{thm:s6v27-erasure-map}, plaquettes outside \(\Sigma\)
are unchanged.  Apply \Cref{lem:s6v27-centre-drop} on each
plaquette of \(\Sigma\) and multiply by \(\beta\).  The possible
loss from the generated action is at most
\(\kappa_\Psi|\Sigma|\), proving
\eqref{eq:s6v27-total-drop}.  The erasure is Haar preserving and
injective, so the pushforward multiplicity in PDEC is one.  Apply
\Cref{thm:s6v26-erasure}.
\end{proof}

\begin{corollary}[Explicit bare centre-sheet threshold]
\label{cor:s6v27-bare-threshold}
For the pure Wilson law, \(\kappa_\Psi=0\).  Connected
detail-exact centre sheets containing a fixed plaquette have a
summable size distribution whenever
\begin{equation}
 \boxed{
 400e^{-\beta\sigma_0(\delta)}<1,
 \quad\hbox{equivalently}\quad
 \beta>{\log400\over\sigma_0(\delta)}.}                 \label{eq:s6v27-centre-threshold}
\end{equation}
For example, \(\delta=1/4\) gives
\[
 \sigma_0(1/4)=\frac32-\frac1{2\sqrt3}.
\]
\end{corollary}

\begin{proof}
Use \Cref{cor:s6v26-animal-sum} with \(K=1\) and
\(\sigma=\beta\sigma_0(\delta)\).
\end{proof}

\begin{corollary}[Generated-action threshold]
\label{cor:s6v27-generated-threshold}
Under \eqref{eq:s6v27-Psi-debit}, the same connected-sheet
summability holds if
\begin{equation}
 \boxed{
 \beta>
 {\log400+\kappa_\Psi\over\sigma_0(\delta)}.}            \label{eq:s6v27-generated-threshold}
\end{equation}
\end{corollary}

\begin{proof}
This is
\(400e^{-\sigma_\beta}<1\) with
\(\sigma_\beta\) from \eqref{eq:s6v27-total-drop}.
\end{proof}

\section{From sheet suppression to a hybrid propagation rate}
\label{sec:s6v27-hybrid}

\begin{lemma}[Centre-sheet connectivity tail]
\label{lem:s6v27-sheet-connectivity}
Assume \(q_Z:=400e^{-\sigma_\beta}<1\).  The probability that a
connected detail-exact centre-flux sheet containing a fixed
plaquette reaches plaquette-graph distance at least \(r\) is at
most
\begin{equation}
 \boxed{
 {e^{-\sigma_\beta}q_Z^r\over1-q_Z}.}                   \label{eq:s6v27-sheet-tail}
\end{equation}
\end{lemma}

\begin{proof}
A connected set reaching distance \(r\) contains at least
\(r+1\) plaquettes.  Apply the rooted-animal count
\Cref{lem:s6v14-animal} and
\eqref{eq:s6v27-sheet-probability}, then sum sizes
\(n\ge r+1\):
\[
 \sum_{n\ge r+1}400^{n-1}e^{-\sigma_\beta n}
 =
 {e^{-\sigma_\beta}q_Z^r\over1-q_Z}.
\]
\end{proof}

\begin{theorem}[Hybrid good-sector/centre-sheet decay]
\label{thm:s6v27-hybrid}
Suppose a coupling of two local-vacuum boundary laws at temporal
separation \(n\tau\) has the following exhaustive alternative:
\begin{enumerate}[label=\textup{(\roman*)},leftmargin=3.5em]
\item it follows a matched-sector reduced-link path whose
      per-gate transmission is at most \(a\) and which has at least
      \(s_gn-b_g\) gates; or
\item it is carried by a connected detail-exact centre-flux sheet
      whose plaquette-graph radius is at least \(s_Zn-b_Z\).
\end{enumerate}
Here \(s_g,s_Z>0\) and \(b_g,b_Z<\infty\) are uniform incidence
constants; integer parts are understood when necessary.  If
\begin{equation}
 q_g:=17a<1,
 \qquad
 q_Z:=400e^{-\sigma_\beta}<1,                           \label{eq:s6v27-two-rates}
\end{equation}
then the connection probability is bounded by
\begin{equation}
 \boxed{
 C'_gq_g^{s_gn}+C'_Zq_Z^{s_Zn}.}                        \label{eq:s6v27-hybrid-tail}
\end{equation}
The primed constants are explicit finite geometric prefactors
which absorb the collars \(b_g,b_Z\) and integer parts.  Under the
OS source-frame hypotheses of V25, the corresponding cyclic
Hamiltonian gap is at least
\begin{equation}
 \boxed{
 {1\over\tau}
 \min\{s_g|\log q_g|,\ s_Z|\log q_Z|\}.}                \label{eq:s6v27-hybrid-mass}
\end{equation}
\end{theorem}

\begin{proof}
The good-sector path bound is
\Cref{thm:s6v25-path-decay} with \(p=0\) and
\(\Delta_{\rm link}=18\).  The centre-sheet term is
\Cref{lem:s6v27-sheet-connectivity}.  The assumed exhaustive
alternative and the union bound give
\eqref{eq:s6v27-hybrid-tail}; the inequalities
\(\lfloor s_\bullet n-b_\bullet\rfloor
 \ge s_\bullet n-b_\bullet-1\)
are absorbed into \(C'_\bullet\).  Finite collars and prefactors do
not alter either spectral exponent.  Apply
\Cref{thm:s6v22-Gram-gap}.
\end{proof}

\begin{firewall}[What the centre erasure does not suppress]
\label{fire:s6v27-noncentral}
The centre tubes do not cover \(\SU(3)\).  A plaquette outside
all three \(\delta\)-tubes is a noncentral large-field excursion,
not a labelled centre sheet.  Nor is a winding centre sheet
detail-exact on the torus, and a response term need not satisfy
\eqref{eq:s6v27-Psi-debit}.  These three sectors require,
respectively, a noncentral excursion estimate, a declared
local-vacuum/topological convention, and a generated-action
erasure bound.  They are not included in
\eqref{eq:s6v27-sheet-probability}.
\end{firewall}

\begin{assessment}[Exact advance and live bottleneck after V27]
\label{ass:s6v27-bottleneck}
One genuinely nonperturbative bad sector now has an exact Peierls
map.  Small centre tubes force a closed \({\mathbb Z}_3\) flux
surface; on a relative-exact component, centre multiplication of
detail links is a Haar-preserving bijection which removes the
surface and lowers the pure Wilson action by
\(\beta\sigma_0(\delta)\) per plaquette.  This supplies exponential
area suppression rather than a multiplied one-point mean.

The first remaining bad sector is the noncentral complement of
the three centre tubes.  In addition, the selected
response-generated action must be shown to have a linear erasure
debit, and winding sheets must be excluded by the local target or
renormalized separately.
\end{assessment}

\begin{programme}[V28: noncentral excursion erasure]
\label{prog:s6v27-V28}
\begin{enumerate}[leftmargin=3em]
\item Partition the plaquette carrier into the three
      \(\delta\)-centre tubes and their noncentral complement.
\item On a connected noncentral component, construct a
      tree-ordered link interpolation toward the nearest centre
      tubes while preserving the detail block constraint.
\item Bound its Haar distortion and the Wilson action drop per
      newly repaired plaquette, retaining all overlap plaquettes.
\item Either obtain a PDEC for the noncentral component or exhibit
      a Bianchi/overlap configuration with no linear action drop.
\item Audit the response-generated action under both centre and
      noncentral erasures.
\item Only after the two bad sectors are exhaustive, combine them
      with the good-sector transmission in
      \Cref{thm:s6v27-hybrid}.
\end{enumerate}
\end{programme}

\begin{firewall}[Claim boundary after sixth-series V27]
\label{fire:s6v27-boundary}
V27 proves quantitative centre-label Bianchi closure, identifies
the relative cohomology condition for detail erasure, constructs
the exact Haar-preserving centre-cochain map, proves its uniform
Wilson action drop, obtains exponential centre-sheet suppression
with explicit bare and generated thresholds, and derives the
hybrid good-sector/centre-sheet OS mass rate.

V27 does not suppress noncentral excursions, erase winding flux,
bound the actual response-generated action debit, populate the
good-sector transmission, close the static card, construct the
continuum measure, or prove the full Yang--Mills mass gap.
\end{firewall}

\begin{theorem}[Sixth-series V27 result ledger]
\label{thm:s6v27-results}
V27:
\begin{enumerate}[label=\textup{(V27.\arabic*)},leftmargin=6em]
\item proves small-tube centre labels obey the Bianchi cocycle;
\item identifies their dual closed surfaces;
\item formulates the relative exactness required by the fixed
      tree and block map;
\item constructs the Haar-preserving erasure;
\item proves the per-plaquette action gain
      \(3/2-2\delta/\sqrt3\);
\item proves the centre-sheet Peierls bound;
\item gives explicit pure and generated thresholds;
\item derives the centre-sheet connectivity tail; and
\item combines it with good-sector paths into an OS spectral
      rate.
\end{enumerate}
\end{theorem}

\begin{proof}
Use
\Cref{lem:s6v27-Bianchi,
cor:s6v27-dual-surface,
lem:s6v27-relative-cohomology,
thm:s6v27-erasure-map,
lem:s6v27-centre-drop,
thm:s6v27-centre-PDEC,
cor:s6v27-bare-threshold,
cor:s6v27-generated-threshold,
lem:s6v27-sheet-connectivity,
thm:s6v27-hybrid}.
\end{proof}

\paragraph{Parent-version record.}
V27 was formed from
\texttt{Renewal\_Yang\_Mills\_Mass\_Gap\_Research\_Notes\_V26.tex}.
The V26 parent had \(11\,523\) logical source lines,
\(409\,728\) bytes, and SHA--256
\[
 \texttt{d51eb25593783101c72d9b897a112c3b9cd0f3c48a51ef378b4c3b2aec635a4e}.
\]
The next compulsory companion audit is V30.

\clearpage
\chapter{Sixth-series V28: reflection-positive suppression of
noncentral excursions}
\label{chap:s6v28}

\section{Purpose and exact claim boundary}
\label{sec:s6v28-purpose}

V27 removed detail-exact centre flux by a Haar-preserving
cochain transformation.  Its remaining large-field sector consists
of plaquettes outside the three centre tubes.  The first objective
of V28 is to decide whether the same pointwise-erasure mechanism can
handle that sector.  It cannot: a six-staple star can be frustrated
so that repairing one plaquette produces no net Wilson-action drop.

We therefore go upstream to reflection positivity.  For the
periodic pure Wilson measure, the chessboard inequality turns the
energy floor of a noncentral plaquette into a uniform
\emph{multi-plaquette} estimate.  A finite colouring then gives the
defect half of MDSC for bad reduced-link gates and an exponentially
small bad-cluster tail at sufficiently large \(\beta\).  This is a
genuine strengthening of the one-point Markov estimate in V26.
It is not yet a boundary-comparison theorem and does not supply the
matched-good-sector contraction.

\section{Noncentral tube complement and its energy floor}
\label{sec:s6v28-floor}

Let
\[
 Z_3=\{1,\omega,\omega^2\},\qquad
 \omega=e^{2\pi i/3},
\]
and, for \(0<\delta<1/2\), put
\begin{equation}
 {\cal N}_\delta
 :=
 \bigcup_{z\in Z_3}
 \{U\in\SU(3):\|U-zI\|_{\rm F}<\delta\},
 \qquad
 {\cal B}_\delta:=\SU(3)\setminus{\cal N}_\delta .
                                                               \label{eq:s6v28-tubes}
\end{equation}
For a plaquette \(p\), write
\[
 B_p(\delta):=\{U:U_p\in{\cal B}_\delta\}.
\]

\begin{lemma}[Exact Wilson-distance identity]
\label{lem:s6v28-distance}
For every \(U\in\SU(3)\),
\begin{equation}
 d(U):=1-\frac13\Re\Tr U
 =\frac16\|U-I\|_{\rm F}^2.                              \label{eq:s6v28-distance}
\end{equation}
Consequently,
\begin{equation}
 U\in{\cal B}_\delta
 \quad\Longrightarrow\quad
 d(U)\ge\varepsilon_\delta,
 \qquad
 \varepsilon_\delta:=\frac{\delta^2}{6}.                 \label{eq:s6v28-floor}
\end{equation}
\end{lemma}

\begin{proof}
Unitarity gives
\[
 \|U-I\|_{\rm F}^2
 =\Tr\{(U-I)^*(U-I)\}
 =6-2\Re\Tr U,
\]
which is \eqref{eq:s6v28-distance}.  Membership in
\({\cal B}_\delta\) implies in particular
\(\|U-I\|_{\rm F}\ge\delta\), proving
\eqref{eq:s6v28-floor}.
\end{proof}

\begin{proposition}[One-link uniform descent is impossible]
\label{prop:s6v28-local-no-go}
There is an admissible parity-tree reduced six-staple star for which
the complete incident Wilson action is independent of the retained
link \(u\), while the designated plaquette can be moved from
\({\cal B}_\delta\) into the identity tube by changing \(u\).
Hence no transformation which changes only that retained link can
have a strictly positive Wilson-action drop on every noncentral
defect configuration.
\end{proposition}

\begin{proof}
In the exact single-link conditional formula of
\Cref{thm:s6v21-star-disintegration}, the \(u\)-dependent incident
action is, up to an additive constant,
\[
 -{\beta\over3}\Re\Tr(uH),\qquad
 H=\sum_{j=1}^6s_j,
\]
where the \(s_j\in\SU(3)\) are the six oriented staples.  Take the
designated staple to be \(s_1=I\) and the remaining five to be
\[
 I,\ \omega I,\ \omega I,\ \omega^2I,\ \omega^2I.
\]
Since \(1+\omega+\omega^2=0\), their sum is \(-I\), and therefore
\(H=0\).  The parity-star realization used in
\Cref{lem:s6v24-free-star} realizes these six controls
without altering the exterior retained coordinates.

Choose \(u\in{\cal B}_\delta\) and replace it by \(I\).  The
designated plaquette changes from \(u\) to \(I\), but the sum of all
six incident Wilson terms is unchanged because \(H=0\).  A
uniformly positive total-action drop is therefore impossible for
any map supported only on this link.
\end{proof}

\section{A uniform partition-function lower bound}
\label{sec:s6v28-partition}

Let \(\Lambda\) be a four-dimensional periodic rectangular lattice
whose side lengths have the dyadic form
\(L_i=2^{m_i}\), \(m_i\ge1\).  Write \(V=|\Lambda|\),
\(E=4V\), and \(P=6V\) for its numbers of vertices, positively
oriented links, and unoriented plaquettes.  With normalized product
Haar measure \(m\), define
\begin{equation}
 Z_{\Lambda,\beta}
 :=
 \int\exp\left\{-\beta\sum_p d(U_p)\right\}\dd m(U),
 \qquad
 \mu_{\Lambda,\beta}:=Z_{\Lambda,\beta}^{-1}e^{-\beta\sum_pd(U_p)}m.
                                                               \label{eq:s6v28-Wilson}
\end{equation}
Put
\[
 h(r):=\operatorname{Haar}\{U\in\SU(3):\|U-I\|_{\rm F}<r\}.
\]

\begin{lemma}[Eight-dimensional Haar-ball lower bound]
\label{lem:s6v28-Haar-ball}
There are constants \(c_{\rm H}>0\) and \(r_{\rm H}>0\), depending
only on \(\SU(3)\), such that
\begin{equation}
 h(r)\ge c_{\rm H}r^8
 \qquad(0<r\le r_{\rm H}).                              \label{eq:s6v28-Haar-ball}
\end{equation}
\end{lemma}

\begin{proof}
Equip the eight-dimensional Lie algebra
\(\mathfrak{su}(3)\) with the Frobenius norm.  On a sufficiently
small Euclidean ball the exponential map is a diffeomorphism, its
Haar-density Jacobian is continuous and strictly positive at the
origin, and
\(\|\exp X-I\|_{\rm F}\le2\|X\|_{\rm F}\).
The exponential image of the algebra ball of radius \(r/2\) is
therefore contained in the group ball of radius \(r\), and its Haar
measure is bounded below by a fixed positive Jacobian minimum times
the Euclidean volume of that ball.  This gives
\eqref{eq:s6v28-Haar-ball}.
\end{proof}

\begin{lemma}[Small-link lower bound for the Wilson partition
function]
\label{lem:s6v28-Z-lower}
For \(0<r\le r_{\rm H}\),
\begin{equation}
 \boxed{
 Z_{\Lambda,\beta}
 \ge h(r)^E
 \exp\left\{-{8\over3}\beta Pr^2\right\}.}               \label{eq:s6v28-Z-lower}
\end{equation}
\end{lemma}

\begin{proof}
Restrict the integral to configurations satisfying
\(\|U_e-I\|_{\rm F}<r\) on every positive link.  For unitary
matrices,
\[
 \|AB-I\|_{\rm F}
 \le\|A-I\|_{\rm F}+\|B-I\|_{\rm F},
\]
and inversion preserves the distance from \(I\).  Hence every
plaquette holonomy on this set obeys
\(\|U_p-I\|_{\rm F}<4r\).  By
\eqref{eq:s6v28-distance},
\(d(U_p)<8r^2/3\).  The restricted product-Haar volume is
\(h(r)^E\), proving the claim.
\end{proof}

\section{Reflection positivity and the chessboard estimate}
\label{sec:s6v28-chessboard}

\begin{lemma}[Reflection positivity of the Wilson weight]
\label{lem:s6v28-RP}
The periodic \(\SU(3)\) Wilson measure
\eqref{eq:s6v28-Wilson} is reflection positive across every
coordinate site or bond hyperplane.  The event
\(B_p(\delta)\) is carried to the corresponding reflected event.
\end{lemma}

\begin{proof}
After discarding the constant \(e^{-\beta P}\), each plaquette
weight is
\[
 \exp\left\{{\beta\over6}
       \bigl(\chi_{\boldsymbol3}(U_p)
             +\chi_{\overline{\boldsymbol3}}(U_p)\bigr)\right\}.
\]
Its power series is a sum of products of the two fundamental
characters with nonnegative coefficients.  Tensor-product
decomposition into irreducible characters has nonnegative integer
multiplicities.  Thus every character coefficient of the
plaquette weight is nonnegative.  Peter--Weyl matrix-element
orthogonality then writes each plaquette crossing a reflection
plane as a sum
\(\sum_\alpha c_\alpha F_\alpha\vartheta F_\alpha\) with
\(c_\alpha\ge0\).  Plaquettes contained in either half are paired
by reflection.  Expanding the crossing factors and integrating
gives
\[
 \int F\,\vartheta F\,\dd\mu_{\Lambda,\beta}\ge0
\]
for every function \(F\) supported in one half.

Under reflection a plaquette holonomy is changed only by inversion
and conjugation.  Since the set of centre matrices is closed under
inversion and Frobenius distance is conjugation invariant,
\({\cal B}_\delta\) is preserved.
\end{proof}

\begin{lemma}[Block chessboard inequality]
\label{lem:s6v28-chessboard}
Tile \(\Lambda\) by translates and reflections of a
\(2\times2\times2\times2\) block, and let \(N=V/16\) be the number
of blocks.  If \(A\) is a reflection-compatible event assigned to
one block and \(A_x\) denotes its image in block \(x\), then, for
every set \(J\) of distinct blocks,
\begin{equation}
 \mu_{\Lambda,\beta}\left(\bigcap_{x\in J}A_x\right)
 \le
 \left\{
 {Z_{\Lambda,\beta}(A^{\rm disp})\over Z_{\Lambda,\beta}}
 \right\}^{|J|/N},                                    \label{eq:s6v28-chessboard}
\end{equation}
where \(A^{\rm disp}:=\bigcap_xA_x\).
\end{lemma}

\begin{proof}
Reflect successively through all bisecting coordinate
hyperplanes.  At each reflection, positivity from
\Cref{lem:s6v28-RP} and Cauchy--Schwarz replace the two half-block
event patterns by their two symmetric doubles and take a square
root.  After all dyadic reflections, each surviving factor is the
partition function with \(A\) disseminated to every block.
Counting how often each original marked block occurs gives exponent
\(|J|/N\).  This finite iteration proves
\eqref{eq:s6v28-chessboard}; no infinite-volume limit is used.
\end{proof}

\section{The multi-plaquette noncentral bound}
\label{sec:s6v28-multi}

Colour each plaquette by its unoriented coordinate pair
\(\{\mu,\nu\}\) and by the parity vector of its lower corner.
There are
\begin{equation}
 C_{\rm col}=6\cdot2^4=96                              \label{eq:s6v28-colours}
\end{equation}
colours.  For a fixed colour, a colour-dependent shifted
\(2^4\)-block tiling places exactly one such plaquette in each
block.  Reflections may reverse its orientation, which does not
change \(B_p(\delta)\).

\begin{theorem}[Uniform same-colour chessboard suppression]
\label{thm:s6v28-same-colour}
Assume
\[
 \beta\ge r_{\rm H}^{-2}
\]
and define
\begin{equation}
 Q_\delta(\beta)
 :=
 c_{\rm H}^{-64}\beta^{256}
 \exp\left\{256-{\beta\delta^2\over6}\right\}.           \label{eq:s6v28-Q}
\end{equation}
For every finite set \(J\) of plaquettes of one colour,
\begin{equation}
 \boxed{
 \mu_{\Lambda,\beta}
 \left(\bigcap_{p\in J}B_p(\delta)\right)
 \le Q_\delta(\beta)^{|J|}.}                            \label{eq:s6v28-same-colour}
\end{equation}
The constants are independent of the dyadic periodic volume.
\end{theorem}

\begin{proof}
Disseminating \(B_p(\delta)\) through its shifted block tiling
forces all \(N=V/16\) plaquettes of that colour to be noncentral.
By \eqref{eq:s6v28-floor}, their Wilson action is at least
\(\beta\varepsilon_\delta N\).  All other Wilson terms are
nonnegative and normalized product Haar has total mass one, so
\[
 Z_{\Lambda,\beta}(B^{\rm disp})
 \le e^{-\beta\varepsilon_\delta N}.
\]
Apply \Cref{lem:s6v28-chessboard,lem:s6v28-Z-lower} and take
\(r=\beta^{-1/2}\).  Since
\[
 {E\over N}=64,\qquad {P\over N}=96,
\]
the one-block chessboard factor is at most
\[
 e^{-\beta\varepsilon_\delta}
 h(\beta^{-1/2})^{-64}e^{256}
 \le
 c_{\rm H}^{-64}\beta^{256}
 e^{256-\beta\delta^2/6}
 =Q_\delta(\beta).
\]
Raising this factor to \(|J|\) proves
\eqref{eq:s6v28-same-colour}.
\end{proof}

\begin{corollary}[Arbitrary multi-plaquette suppression]
\label{cor:s6v28-all-colours}
Set
\begin{equation}
 p_{\rm plaq}(\beta,\delta)
 :=
 \min\{1,Q_\delta(\beta)^{1/96}\}.                       \label{eq:s6v28-pplaq}
\end{equation}
For every finite set \(J\) of distinct plaquettes,
\begin{equation}
 \boxed{
 \mu_{\Lambda,\beta}
 \left(\bigcap_{p\in J}B_p(\delta)\right)
 \le p_{\rm plaq}(\beta,\delta)^{|J|}.}                 \label{eq:s6v28-all-colours}
\end{equation}
\end{corollary}

\begin{proof}
When \(Q_\delta(\beta)\ge1\), the right-hand side is one.  Suppose
\(Q_\delta(\beta)<1\).  Among the \(96\) colour classes, one class
\(J_c\) contains at least \(|J|/96\) elements.  The event on the
left of \eqref{eq:s6v28-all-colours} is contained in
\(\bigcap_{p\in J_c}B_p(\delta)\).  Therefore
\[
 \mu_{\Lambda,\beta}\left(\bigcap_{p\in J}B_p(\delta)\right)
 \le Q_\delta(\beta)^{|J_c|}
 \le Q_\delta(\beta)^{|J|/96}.
\]
\end{proof}

\section{Bad reduced-link gates and a subcritical cluster}
\label{sec:s6v28-link-gates}

Call a retained or unreduced link \(e\) noncentrally bad if at least
one of its incident plaquettes is in \(B_p(\delta)\), and denote
this event by \(G_e(\delta)\).  A four-dimensional link belongs to
six plaquettes, while a plaquette contains four links.

\begin{theorem}[Defect half of link-gate MDSC]
\label{thm:s6v28-link-MDSC}
Define
\begin{equation}
 p_{\rm link}(\beta,\delta)
 :=
 \min\left\{1,\,
 6p_{\rm plaq}(\beta,\delta)^{1/4}\right\}.              \label{eq:s6v28-plink}
\end{equation}
Then every finite set \(K\) of distinct links satisfies
\begin{equation}
 \boxed{
 \mu_{\Lambda,\beta}
 \left(\bigcap_{e\in K}G_e(\delta)\right)
 \le p_{\rm link}(\beta,\delta)^{|K|}.}                 \label{eq:s6v28-link-MDSC}
\end{equation}
In particular, this proves the probability inequality in the MDSC
card for every predeclared path in a single periodic Wilson field.
\end{theorem}

\begin{proof}
Expand each event \(G_e\) as a union over the at most six
plaquettes incident to \(e\).  There are at most \(6^{|K|}\)
resulting assignments \(e\mapsto p(e)\).  Because one plaquette
contains at most four links, every assignment uses at least
\(|K|/4\) distinct plaquettes.  By
\Cref{cor:s6v28-all-colours}, each assigned intersection has
probability at most
\(p_{\rm plaq}^{|K|/4}\).  The union bound gives
\[
 \mu_{\Lambda,\beta}\left(\bigcap_{e\in K}G_e\right)
 \le
 \{6p_{\rm plaq}^{1/4}\}^{|K|}.
\]
If the expression in braces exceeds one, the trivial bound gives
the stated minimum.
\end{proof}

\begin{corollary}[Explicit large-\(\beta\) rate]
\label{cor:s6v28-explicit-rate}
Whenever \(Q_\delta(\beta)<1\),
\begin{equation}
 p_{\rm link}(\beta,\delta)
 \le
 \min\left\{1,\,
 6c_{\rm H}^{-1/6}\beta^{2/3}
 \exp\left\{{2\over3}
             -{\beta\delta^2\over2304}\right\}\right\}. \label{eq:s6v28-link-rate}
\end{equation}
Thus \(p_{\rm link}(\beta,\delta)\to0\) exponentially as
\(\beta\to\infty\), uniformly in the even periodic volume.
\end{corollary}

\begin{proof}
Under \(Q_\delta<1\),
\(p_{\rm plaq}=Q_\delta^{1/96}\).  Insert
\eqref{eq:s6v28-Q} into
\(6p_{\rm plaq}^{1/4}=6Q_\delta^{1/384}\) and simplify:
\[
 {64\over384}={1\over6},\qquad
 {256\over384}={2\over3},\qquad
 {1\over6\cdot384}={1\over2304}.
\]
\end{proof}

\begin{corollary}[Exponential noncentral bad-cluster tail]
\label{cor:s6v28-cluster}
On the degree-eighteen reduced-link dependency graph, assume
\begin{equation}
 17p_{\rm link}(\beta,\delta)<1.                         \label{eq:s6v28-subcritical}
\end{equation}
Then the probability that a fixed root link is joined by a
self-avoiding path of noncentrally bad link gates having at least
\(n\) gates is at most
\begin{equation}
 \boxed{
 {18\over17}\,
 {(17p_{\rm link})^n\over1-17p_{\rm link}}.}            \label{eq:s6v28-cluster-tail}
\end{equation}
The same bound passes to every subsequential infinite-volume limit
of the dyadic periodic Wilson measures.
\end{corollary}

\begin{proof}
There are at most \(18\cdot17^{k-1}\) self-avoiding
degree-eighteen paths of \(k\) gates from the root.  By
\Cref{thm:s6v28-link-MDSC}, the probability that all gates on a
fixed such path are bad is at most \(p_{\rm link}^k\).  Sum over
\(k\ge n\).  The bound concerns cylinder events, so it is preserved
under weak convergence on the compact product configuration space.
\end{proof}

\section{What the chessboard step does and does not close}
\label{sec:s6v28-payoff}

\begin{assessment}[A genuine noncentral-sector estimate]
\label{ass:s6v28-payoff}
For the bare periodic Wilson measure, V28 proves a volume-uniform
multi-defect estimate
\eqref{eq:s6v28-all-colours}, not an illegitimate product of
one-point probabilities.  It also proves that noncentral
bad-link clusters are subcritical for all sufficiently large
\(\beta\).  The exponent is deliberately crude: the \(96\)-colour
loss and the small-ball partition lower bound were retained
explicitly so no entropy factor is hidden.
\end{assessment}

\begin{firewall}[Why this is not yet CDAC]
\label{fire:s6v28-not-CDAC}
The chessboard estimate is proved for periodic reflection-positive
Wilson measures.  Conditioning on an arbitrary explored exterior
can destroy reflection positivity.  A response-generated action
can also destroy the nonnegative character expansion.  Moreover,
the second MDSC clause---the \(a^{N_{\rm good}}\) coupling bound
conditional on the complete bad-gate vector---is untouched.
Therefore neither
\eqref{eq:s6v28-link-MDSC} nor
\eqref{eq:s6v28-cluster-tail} may be inserted by itself into the
CDAC-to-OS theorem.
\end{firewall}

\begin{definition}[Reflection-positive comparison card, RPCC]
\label{def:s6v28-RPCC}
For a response-fixed local-vacuum measure, RPCC with debit
\(\kappa_{\rm RP}\ge0\) means that every colour-disseminated
noncentral event satisfies
\begin{equation}
 {Z_{\beta,\Psi}(B^{\rm disp})\over Z_{\beta,\Psi}}
 \le
 e^{\kappa_{\rm RP}N}
 {Z_{\Lambda,\beta}(B^{\rm disp})\over Z_{\Lambda,\beta}},
                                                               \label{eq:s6v28-RPCC}
\end{equation}
uniformly in the torus and admissible response sector.
\end{definition}

\begin{corollary}[Payoff of RPCC]
\label{cor:s6v28-RPCC-payoff}
Under RPCC, every conclusion from
\Cref{thm:s6v28-same-colour} through
\Cref{cor:s6v28-cluster} remains true after replacing
\begin{equation}
 Q_\delta(\beta)
 \quad\hbox{by}\quad
 e^{\kappa_{\rm RP}}Q_\delta(\beta).                    \label{eq:s6v28-RPCC-Q}
\end{equation}
\end{corollary}

\begin{proof}
Take the \(N\)-th root of \eqref{eq:s6v28-RPCC} in the chessboard
bound.  All subsequent colouring, incidence, and path arguments
are unchanged.
\end{proof}

\begin{programme}[V29: convert bulk defect suppression into
boundary-stable propagation control]
\label{prog:s6v28-V29}
\begin{enumerate}[leftmargin=3em]
\item Identify the smallest reflected collar whose conditional
      kernel retains a doubled reflection-positive representation.
\item Test whether the response-fixed action satisfies RPCC, or
      isolate the first generated term with a negative
      reflection-channel coefficient.
\item Compare arbitrary local-vacuum boundary laws to a doubled
      periodic collar by an explicit Radon--Nikodym debit.
\item Keep centre sheets from V27 and noncentral clusters from V28
      as separate bad carriers; do not merge their rates before
      proving an exhaustive carrier decomposition.
\item On the complement, attack the remaining matched-sector
      transmission \(a<1/17\) through the canonical-root static
      card rather than through single-link Dobrushin influence.
\end{enumerate}
\end{programme}

\begin{firewall}[Claim boundary after sixth-series V28]
\label{fire:s6v28-boundary}
V28 proves the exact noncentral Wilson energy floor, a local
one-link-erasure obstruction, a volume-uniform partition lower
bound, reflection positivity of the Wilson weight, the relevant
block chessboard inequality, uniform same-colour and arbitrary
multi-plaquette suppression, the defect half of link-gate MDSC,
and an explicit exponential noncentral bad-cluster tail.

V28 does not prove boundary-stable or response-fixed
reflection positivity, the conditional good-gate coupling, an
exhaustive disagreement-carrier decomposition, the canonical-root
static card, continuum existence, or the full Yang--Mills mass gap.
\end{firewall}

\begin{theorem}[Sixth-series V28 result ledger]
\label{thm:s6v28-results}
The claims listed in
\Cref{fire:s6v28-boundary} follow from
\Cref{lem:s6v28-distance,
prop:s6v28-local-no-go,
lem:s6v28-Haar-ball,
lem:s6v28-Z-lower,
lem:s6v28-RP,
lem:s6v28-chessboard,
thm:s6v28-same-colour,
cor:s6v28-all-colours,
thm:s6v28-link-MDSC,
cor:s6v28-explicit-rate,
cor:s6v28-cluster}.
\end{theorem}

\begin{proof}
Each item is precisely the conclusion of the cited result.
\end{proof}

\paragraph{Parent-version record.}
V28 was formed from
\texttt{Renewal\_Yang\_Mills\_Mass\_Gap\_Research\_Notes\_V27.tex}.
The V27 parent had \(11\,985\) logical source lines,
\(425\,640\) bytes, and SHA--256
\[
 \texttt{c6b950842a93b5614903bd0d4c4bff20415e3733eddb81a8fb492bf65eb46f49}.
\]
The next compulsory companion audit is V30.

\clearpage
\chapter{Sixth-series V29: vacuum-fillable collars and
boundary-stable defect suppression}
\label{chap:s6v29}

\section{Why the periodic estimate needs a new local input}
\label{sec:s6v29-purpose}

The V28 chessboard estimate is uniform in the periodic volume, but
the CDAC exploration conditions on an exterior configuration.
Reflection positivity need not survive that conditioning.  This
note replaces conditional reflection positivity by a finite local
quantity: the smallest curvature radius attainable by filling the
links of the defect collar while its exterior is held fixed.

When that radius is smaller than
\(\delta/\sqrt{24}\), a translated product-Haar ball gives an
exponentially small conditional probability for every declared set
of noncentral plaquettes.  The proof is direct and uses neither
independence nor a periodic comparison.  It also shows precisely
where the response-generated action enters.

\section{Conditional Wilson kernels and the fillability radius}
\label{sec:s6v29-fillability}

Let \({\mathscr E}\) be a finite link set in a fixed finite lattice,
let \(\Omega\subset{\mathscr E}\) be a nonempty set of variable
links, and let \(\xi\in\SU(3)^{{\mathscr E}\setminus\Omega}\) be
the frozen exterior.  Denote by
\[
 {\mathscr P}(\Omega)
 :=
 \{p:p\hbox{ contains at least one link of }\Omega\}
\]
the active plaquettes and put
\[
 L_\Omega:=|\Omega|,
 \qquad
 P_\Omega:=|{\mathscr P}(\Omega)|.
\]
The pure-Wilson conditional kernel is
\begin{equation}
 \gamma_{\Omega,\beta}^{\xi}(\dd z)
 :=
 { \exp\{-\beta\sum_{p\in{\mathscr P}(\Omega)}
                    d(U_p(z,\xi))\}
    \prod_{e\in\Omega}\dd\lambda(z_e)
  \over
   Z_{\Omega,\beta}^{\xi}}.                             \label{eq:s6v29-kernel}
\end{equation}
Terms supported outside \({\mathscr P}(\Omega)\) have cancelled
from this conditional law.

\begin{definition}[Vacuum-fillability radius]
\label{def:s6v29-radius}
Define
\begin{equation}
 \rho_\Omega(\xi)
 :=
 \inf_{z^0\in\SU(3)^\Omega}
 \max_{p\in{\mathscr P}(\Omega)}
 \|U_p(z^0,\xi)-I\|_{\rm F}.                            \label{eq:s6v29-radius}
\end{equation}
The collar \((\Omega,\xi)\) is \(\rho\)-vacuum-fillable if
\(\rho_\Omega(\xi)\le\rho\).
\end{definition}

\begin{lemma}[The fillability infimum is attained and
gauge invariant]
\label{lem:s6v29-radius-attained}
The infimum in \eqref{eq:s6v29-radius} is a minimum.  It is
unchanged by any lattice gauge transformation applied
simultaneously to the variable and exterior links.
\end{lemma}

\begin{proof}
The domain \(\SU(3)^\Omega\) is compact and the maximum of finitely
many continuous functions is continuous, so a minimizer exists.
A gauge transformation conjugates each based plaquette holonomy.
Frobenius distance from \(I\) is conjugation invariant, and the
transformation is a bijection of the filling variables.  Taking
the infimum proves invariance.
\end{proof}

\begin{lemma}[Translated Haar box and plaquette perturbation]
\label{lem:s6v29-translated-box}
Let \(z^0\) attain \(\rho_\Omega(\xi)\), and suppose
\[
 \|z_e-z_e^0\|_{\rm F}<r
 \quad(e\in\Omega).
\]
Then
\begin{equation}
 \|U_p(z,\xi)-I\|_{\rm F}
 <
 \rho_\Omega(\xi)+4r
 \quad(p\in{\mathscr P}(\Omega)).                       \label{eq:s6v29-plaquette-perturb}
\end{equation}
The product-Haar measure of this translated box is \(h(r)^{L_\Omega}\).
\end{lemma}

\begin{proof}
Write the four factors in the plaquette word for \(U_p(z,\xi)\)
and for \(U_p(z^0,\xi)\) in the same order.  Frozen factors agree,
and every changed factor, including an inversely oriented one, is
within \(r\) of its reference factor.  Telescoping the two products
and using unitary invariance of the Frobenius norm gives
\[
 \|U_p(z,\xi)-U_p(z^0,\xi)\|_{\rm F}<k_pr\le4r,
\]
where \(k_p\) is the number of variable links in \(p\).  The
triangle inequality proves
\eqref{eq:s6v29-plaquette-perturb}.  Each coordinate ball is a
left translate of the radius-\(r\) ball at \(I\), so Haar
invariance gives its measure \(h(r)\).
\end{proof}

\begin{lemma}[Conditional partition lower bound]
\label{lem:s6v29-conditional-Z}
For \(0<r\le r_{\rm H}\),
\begin{equation}
 \boxed{
 Z_{\Omega,\beta}^{\xi}
 \ge
 h(r)^{L_\Omega}
 \exp\left\{
 -{\beta P_\Omega\over6}
  \bigl(\rho_\Omega(\xi)+4r\bigr)^2
 \right\}.}                                             \label{eq:s6v29-conditional-Z}
\end{equation}
\end{lemma}

\begin{proof}
Restrict the conditional integral to the translated Haar box from
\Cref{lem:s6v29-translated-box}.  On that box,
\Cref{lem:s6v28-distance} gives
\[
 d(U_p)\le{1\over6}
 \bigl(\rho_\Omega(\xi)+4r\bigr)^2.
\]
Multiply the lower bound on the integrand by the box volume.
\end{proof}

\section{A conditional multi-defect theorem}
\label{sec:s6v29-multi}

Let \(\Gamma\subset{\mathscr P}(\Omega)\) be a set of \(k\)
distinct active plaquettes and recall the noncentral events
\(B_p(\delta)\) from \eqref{eq:s6v28-tubes}.

\begin{theorem}[Vacuum-fillable conditional suppression]
\label{thm:s6v29-conditional-suppression}
Suppose
\begin{equation}
 L_\Omega\le Ak,
 \qquad
 P_\Omega\le Bk,
 \qquad
 \rho_\Omega(\xi)\le\rho,                               \label{eq:s6v29-incidence-card}
\end{equation}
where \(A,B<\infty\).  For
\(\beta\ge r_{\rm H}^{-2}\), define
\begin{equation}
 Q_{\rm cond}(\beta,\delta,\rho;A,B)
 :=
 c_{\rm H}^{-A}\beta^{4A}
 \exp\left\{
 -{\beta\delta^2\over6}
 +{\beta B\over6}
   \bigl(\rho+4\beta^{-1/2}\bigr)^2
 \right\}.                                              \label{eq:s6v29-Qcond}
\end{equation}
Then
\begin{equation}
 \boxed{
 \gamma_{\Omega,\beta}^{\xi}
 \left(\bigcap_{p\in\Gamma}B_p(\delta)\right)
 \le
 Q_{\rm cond}(\beta,\delta,\rho;A,B)^k.}                \label{eq:s6v29-conditional-suppression}
\end{equation}
\end{theorem}

\begin{proof}
On the event in \eqref{eq:s6v29-conditional-suppression},
\Cref{lem:s6v28-distance} gives
\[
 \beta\sum_{p\in{\mathscr P}(\Omega)}d(U_p)
 \ge{\beta\delta^2\over6}k.
\]
All Wilson terms are nonnegative, and normalized product Haar has
total mass one, so the unnormalized event integral is at most
\(\exp\{-\beta\delta^2k/6\}\).  Divide by
\eqref{eq:s6v29-conditional-Z} with
\(r=\beta^{-1/2}\).  By
\Cref{lem:s6v28-Haar-ball},
\[
 h(\beta^{-1/2})^{-L_\Omega}
 \le
 \left(c_{\rm H}^{-A}\beta^{4A}\right)^k.
\]
The remaining exponential is bounded with
\eqref{eq:s6v29-incidence-card}, proving the claim.
\end{proof}

\begin{corollary}[The sharp asymptotic fillability inequality]
\label{cor:s6v29-fillability-threshold}
For fixed \(A,B,\delta,\rho\), the bound
\eqref{eq:s6v29-Qcond} tends to zero exponentially as
\(\beta\to\infty\) whenever
\begin{equation}
 \boxed{
 B\rho^2<\delta^2.}                                    \label{eq:s6v29-fillability-threshold}
\end{equation}
For the plaquette collar chosen below, \(B=24\), so the condition
is
\begin{equation}
 \boxed{\rho<{\delta\over\sqrt{24}}.}                   \label{eq:s6v29-rho-threshold}
\end{equation}
\end{corollary}

\begin{proof}
Divide \(\log Q_{\rm cond}\) by \(\beta\).  The logarithmic
Haar-entropy term vanishes, while
\[
 -{\delta^2\over6}
 +{B\over6}\bigl(\rho+4\beta^{-1/2}\bigr)^2
 \longrightarrow
 -{\delta^2-B\rho^2\over6}.
\]
This limit is strictly negative precisely under
\eqref{eq:s6v29-fillability-threshold}.
\end{proof}

\section{The canonical plaquette collar}
\label{sec:s6v29-collar}

\begin{lemma}[Uniform collar incidence]
\label{lem:s6v29-collar-incidence}
Let \(\Gamma\) be any \(k\)-plaquette set and let \(\Omega_\Gamma\)
be the set of all variable reduced links occurring in plaquettes
of \(\Gamma\).  Then
\begin{equation}
 L_{\Omega_\Gamma}\le4k,
 \qquad
 P_{\Omega_\Gamma}\le24k.                              \label{eq:s6v29-collar-incidence}
\end{equation}
\end{lemma}

\begin{proof}
Each plaquette has four links, so its union contains at most
\(4k\) variable links.  In four dimensions each link belongs to
six plaquettes.  Counting active plaquettes with multiplicity
therefore gives \(P_{\Omega_\Gamma}\le6L_{\Omega_\Gamma}\le24k\).
Fixed tree or retained links can only lower the counts.
\end{proof}

\begin{definition}[Path-local vacuum fillability card, PLVFC]
\label{def:s6v29-PLVFC}
Fix a predeclared self-avoiding reduced-link path and an exterior
sector.  PLVFC with radius \(\rho\) means that for every selection
of one incident plaquette at each of any set of distinct path
links, after duplicates are removed, the resulting plaquette set
\(\Gamma\) satisfies
\begin{equation}
 \rho_{\Omega_\Gamma}(\xi)\le\rho                       \label{eq:s6v29-PLVFC}
\end{equation}
for the actual conditional exterior \(\xi\).  If a generated
action \(\Psi\) is present, PLVFC also records the conditional
oscillation debit
\begin{equation}
 \operatorname{osc}_{z\in\SU(3)^{\Omega_\Gamma}}
 \Psi(z,\xi)\le\kappa_\Psi|\Gamma|.                     \label{eq:s6v29-Psi-osc}
\end{equation}
\end{definition}

\begin{lemma}[Exact response-action transfer]
\label{lem:s6v29-response-transfer}
Let \(\gamma_{\Omega,\beta,\Psi}^{\xi}\) be obtained from
\eqref{eq:s6v29-kernel} by adding \(\Psi(z,\xi)\) to the
conditional action.  For every event \(A\),
\begin{equation}
 \gamma_{\Omega,\beta,\Psi}^{\xi}(A)
 \le
 e^{\operatorname{osc}_\Omega\Psi(\cdot,\xi)}
 \gamma_{\Omega,\beta}^{\xi}(A).                        \label{eq:s6v29-response-transfer}
\end{equation}
\end{lemma}

\begin{proof}
The Radon--Nikodym derivative of the generated law relative to the
Wilson law is
\[
 {e^{-\Psi}\over\gamma_{\Omega,\beta}^{\xi}(e^{-\Psi})}.
\]
Its essential supremum is at most
\(e^{\sup\Psi-\inf\Psi}\).  Integrating over \(A\) proves the
claim.
\end{proof}

\begin{theorem}[Boundary-stable defect half of MDSC]
\label{thm:s6v29-boundary-MDSC}
Assume PLVFC with
\(\rho<\delta/\sqrt{24}\) and response debit
\(\kappa_\Psi\).  Put
\begin{align}
 \widehat Q(\beta,\delta,\rho,\kappa_\Psi)
 &:=
 e^{\kappa_\Psi}
 Q_{\rm cond}(\beta,\delta,\rho;4,24),                  \label{eq:s6v29-Qhat}\\
 \widehat p_{\rm link}
 &:=
 \min\{1,\,6\widehat Q^{1/4}\}.                         \label{eq:s6v29-phat}
\end{align}
For every set \(K\) of distinct path gates, conditionally on every
exterior for which PLVFC holds,
\begin{equation}
 \boxed{
 \mathbb P(B_e=1\text{ for every }e\in K\mid\xi)
 \le\widehat p_{\rm link}^{|K|}.}                       \label{eq:s6v29-boundary-MDSC}
\end{equation}
For fixed \(\delta,\rho,\kappa_\Psi\),
\(\widehat p_{\rm link}\to0\) exponentially as
\(\beta\to\infty\).
\end{theorem}

\begin{proof}
For each link gate, take the union over its at most six incident
noncentral plaquettes.  This gives at most \(6^{|K|}\)
assignments.  An assignment uses at least \(|K|/4\) distinct
plaquettes because each plaquette contains at most four links.
For its distinct set \(\Gamma\), use
\Cref{lem:s6v29-collar-incidence} and PLVFC in
\Cref{thm:s6v29-conditional-suppression}.  The response transfer
\eqref{eq:s6v29-response-transfer} and
\eqref{eq:s6v29-Psi-osc} multiply the \(k\)-plaquette bound by
\(e^{\kappa_\Psi k}\).  Hence the assignment costs at most
\(\widehat Q^{|K|/4}\).  The union bound proves
\eqref{eq:s6v29-boundary-MDSC}.  Exponential decay follows from
\Cref{cor:s6v29-fillability-threshold}; a fixed
\(\kappa_\Psi\) does not change the coefficient of \(\beta\).
\end{proof}

\begin{corollary}[Conditional CDAC threshold under the remaining
good-gate card]
\label{cor:s6v29-CDAC-threshold}
If, in addition, conditional on the bad-gate vector the matched
good gates transmit with product bound \(a^{N_{\rm good}}\), then
CDAC holds whenever
\begin{equation}
 \boxed{17(a+\widehat p_{\rm link})<1.}                 \label{eq:s6v29-CDAC-threshold}
\end{equation}
For every fixed \(a<1/17\), PLVFC with
\(\rho<\delta/\sqrt{24}\) and finite \(\kappa_\Psi\) makes
\eqref{eq:s6v29-CDAC-threshold} true for all sufficiently large
\(\beta\).
\end{corollary}

\begin{proof}
\Cref{thm:s6v29-boundary-MDSC} supplies the probability clause of
MDSC, while the additional hypothesis is its good-gate clause.
Apply \Cref{thm:s6v26-MDSC-path} and the degree-eighteen threshold
of \Cref{cor:s6v26-absorption-rows}.  The final assertion uses
\(\widehat p_{\rm link}\to0\).
\end{proof}

\section{What has been converted and what remains geometric}
\label{sec:s6v29-assessment}

\begin{theorem}[PLVFC is a finite compact optimization card]
\label{thm:s6v29-decidability}
On a fixed path collar and fixed exterior \(\xi\), every quantity
\(\rho_{\Omega_\Gamma}(\xi)\) required by PLVFC is attained.
There are finitely many plaquette assignments.  Consequently a
strict PLVFC verdict
\[
 \max_\Gamma\rho_{\Omega_\Gamma}(\xi)
 <{\delta\over\sqrt{24}}
\]
is decidable by convergent outward enclosures of the corresponding
finite-dimensional continuous functions, provided the exterior
sector itself has an explicit compact parametrization.
\end{theorem}

\begin{proof}
Attainment is \Cref{lem:s6v29-radius-attained}.  A finite path has
finitely many subsets and at most six choices per selected link.
Uniform continuity on each compact product group makes convergent
subdivision enclosures decisive whenever the strict margin is
positive.  An explicit compact exterior parametrization is needed
to take the final sector supremum.
\end{proof}

\begin{assessment}[Upstream meaning of the new radius]
\label{ass:s6v29-radius-meaning}
V29 converts boundary-stable noncentral suppression into two
auditable local statements:
\[
 \rho_*<{\delta\over\sqrt{24}},
 \qquad
 \kappa_\Psi<\infty
\]
on the occurring local-vacuum sector.  The first is geometric and
gauge invariant; the second is an exact generated-action
oscillation.  Neither is implied by the response coefficient
alone.  Unlike an arbitrary boundary Radon--Nikodym comparison,
the fillability construction pays entropy only for the
\(O(|\Gamma|)\) links in the actual defect collar.
\end{assessment}

\begin{firewall}[The local-vacuum occurrence problem remains]
\label{fire:s6v29-occurrence}
PLVFC is not proved merely because all observed plaquettes lie in
some centre tube.  A centre cocycle may be nontrivial, and even a
detail-exact centre label must be erased before the collar can be
filled near the identity.  V27 supplies that erasure only for its
declared relative-exact sector.  V28 controls bulk noncentral
clusters in periodic Wilson limits.  What is still missing is an
exhaustive theorem saying that every conditional exploration
either encounters one of those already suppressed carriers or
lies in a uniformly fillable local-vacuum collar.
\end{firewall}

\begin{programme}[V30: compulsory companion audit and the
fillable/exhaustive carrier split]
\label{prog:s6v29-V30}
\begin{enumerate}[leftmargin=3em]
\item Perform the required audit of the current SM and NS notes
      before choosing the next argument.
\item In the exact parity-tree geometry, classify the obstruction
      to \(\rho_\Omega(\xi)<\delta/\sqrt{24}\) into centre
      cohomology, a noncentral plaquette cluster, or a residual
      small-curvature holonomy.
\item Determine whether the residual small-curvature holonomy can
      be removed by a relative axial fill with an explicit radius
      amplification.
\item Load the actual generated interaction if the companion audit
      identifies it; otherwise keep \(\kappa_\Psi\) as an honest
      unpopulated row.
\item Only after the carrier split is exhaustive, return to the
      matched-good-sector constant \(a\) and the canonical-root
      static pair card.
\end{enumerate}
\end{programme}

\begin{firewall}[Claim boundary after sixth-series V29]
\label{fire:s6v29-boundary}
V29 defines a gauge-invariant, attained vacuum-fillability radius;
proves a translated-Haar conditional partition lower bound;
derives a boundary-stable multi-plaquette estimate with exact
incidence constants; identifies the sharp asymptotic condition
\(\rho<\delta/\sqrt{24}\); transfers the estimate through a
generated action by its exact local oscillation; and obtains the
defect half of conditional link-gate MDSC under PLVFC.

V29 does not prove that PLVFC covers every occurring local-vacuum
exterior, bound the actual \(\kappa_\Psi\), prove the good-gate
constant \(a<1/17\), close the canonical-root static card,
construct the continuum measure, or prove the full Yang--Mills
mass gap.
\end{firewall}

\begin{theorem}[Sixth-series V29 result ledger]
\label{thm:s6v29-results}
The claims in \Cref{fire:s6v29-boundary} follow from
\Cref{lem:s6v29-radius-attained,
lem:s6v29-translated-box,
lem:s6v29-conditional-Z,
thm:s6v29-conditional-suppression,
cor:s6v29-fillability-threshold,
lem:s6v29-collar-incidence,
lem:s6v29-response-transfer,
thm:s6v29-boundary-MDSC,
cor:s6v29-CDAC-threshold,
thm:s6v29-decidability}.
\end{theorem}

\begin{proof}
Each ledger item is the conclusion of the cited result.
\end{proof}

\paragraph{Parent-version record.}
V29 was formed from
\texttt{Renewal\_Yang\_Mills\_Mass\_Gap\_Research\_Notes\_V28.tex}.
The V28 parent had \(12\,566\) logical source lines,
\(445\,724\) bytes, and SHA--256
\[
 \texttt{32c43d66e865ba19535cd8f12cc2ca4cec4efaad6ae4987fc6c0c87f8caa0487}.
\]
The next version must perform the compulsory V30 companion audit.

\clearpage
\chapter{Sixth-series V30: companion audit, mesoscopic axial
charts, and the good-sector covariance card}
\label{chap:s6v30}

\section{Compulsory companion audit}
\label{sec:s6v30-audit}

\begin{audit}[Active Standard Model--Einstein note]
\label{audit:s6v30-SM}
The active companion was
\[
 \texttt{Renewal\_SM\_Einstein\_Continuum\_Research\_Notes\_V53.tex}.
\]
It had \(21\,361\) logical source lines, \(860\,179\) bytes,
\(1\,998\) labels, and SHA--256
\[
 \texttt{f99b1260d184d4d4f7aa257ca398f8816ef4f2c6ddfcdcaf4393c2a73f051c4b}.
\]
Its new result is a fixed-range stratified overlap kernel with a
uniform unique low cone, high-branch suppression, a covariant heat
parametrix through weight four, and logarithmic remainder
integrability at the mesoscopic endpoint
\(\varepsilon(a)=a^{1/2}\).  The scale choice
\[
 L=a^{-7/8},\qquad w=a^{1/8},\qquad R=a^{-15/16}
\]
simultaneously closes consistency, localization, and differentiated
profile-jet inequalities.  It explicitly leaves the pure
gravitational coefficient owners and the complete continuum
construction open.
\end{audit}

\begin{audit}[Active Navier--Stokes note]
\label{audit:s6v30-NS}
The active companion was
\[
 \texttt{Renewal\_NS\_Stability\_Research\_Notes\_V3.tex}.
\]
It had \(1\,173\) logical source lines, \(62\,297\) bytes,
\(71\) labels, and SHA--256
\[
 \texttt{50353bc91ed3ae22716ac135c5983bf4106eef41c64e71072d4b65a5d7ec0963}.
\]
It proves \(c_{O,2}\le1.788\), moves the initial term in the
gradient estimate to the ancient Duhamel form, and raises the
necessary homogeneous type-I floor from \(0.1815\) to \(0.3438\).
It states correctly that this is a necessary quantitative
constraint, not a singularity exclusion or a global-regularity
proof.
\end{audit}

\begin{assessment}[Effect of the audit on the Yang--Mills route]
\label{ass:s6v30-audit-effect}
The SM note contributes a usable method: choose one mesoscopic
scale only after listing every competing smallness and largeness
condition, and keep the local kernel fixed-range before transferring
to a global statement.  The NS note supplies a caution of equal
importance: improvement of a certified constant does not close a
qualitative theorem.  Accordingly V30 chooses a Yang--Mills block
scale which simultaneously makes noncentral defects rare, makes
the centre-normalized connection small in an axial chart, and
keeps the gauge-fixed Wilson Hessian coercive.  The resulting
covariance card is not called a mass gap.
\end{assessment}

\section{A quantitative non-Abelian axial chart}
\label{sec:s6v30-axial}

Let \(K_L=\{0,\ldots,L\}^4\) be a nonwrapping lattice box.  Use the
complete axial tree which reaches a vertex by moving successively
in directions \(4,3,2,1\).  Thus a direction-\(\mu\) link is in
the tree when all coordinates of smaller index vanish.  This set
has \(|K_L|-1\) links and is a spanning tree.

\begin{lemma}[Product perturbation]
\label{lem:s6v30-product}
If \(V_1,\ldots,V_m\in\SU(3)\), then
\begin{equation}
 \left\|V_1\cdots V_m-I\right\|_{\rm F}
 \le\sum_{j=1}^m\|V_j-I\|_{\rm F}.                     \label{eq:s6v30-product}
\end{equation}
The same bound holds when each \(V_j\) is first conjugated.
\end{lemma}

\begin{proof}
Telescope the product:
\[
 V_1\cdots V_m-I
 =
 \sum_{j=1}^m
 V_1\cdots V_{j-1}(V_j-I).
\]
Left multiplication and conjugation by a unitary preserve the
Frobenius norm.  Apply the triangle inequality.
\end{proof}

\begin{lemma}[Axial fundamental-loop area]
\label{lem:s6v30-axial-area}
For every non-tree link \(e\) of \(K_L\), the loop formed by \(e\)
and the two axial-tree paths from the root to its endpoints can be
tiled by at most \(3L\) elementary plaquettes.
\end{lemma}

\begin{proof}
Append \(e\) to the axial path to its initial endpoint.  To put the
result into the canonical order \(4,3,2,1\), slide this single
directed step past at most three coordinate segments.  Sliding it
past a segment of length \(x_\nu\le L\) sweeps a strip of
\(x_\nu\) elementary plaquettes in the \((\mu,\nu)\)-plane.
After the slides, the path is the axial path to the terminal
endpoint.  The total tiled area is at most
\(\sum_{\nu\ne\mu}x_\nu\le3L\).
\end{proof}

\begin{theorem}[Centre-normalized axial smallness]
\label{thm:s6v30-axial-smallness}
Suppose every plaquette of \(K_L\) lies in a centre tube of radius
\(\delta<1/2\).  Then there are a centre-valued link cochain and a
lattice gauge transformation after which:
\begin{enumerate}[label=\textup{(\roman*)},leftmargin=3.5em]
\item every plaquette holonomy \(V_p\) satisfies
      \(\|V_p-I\|_{\rm F}<\delta\);
\item every axial-tree link is \(I\); and
\item every remaining link satisfies
      \begin{equation}
       \boxed{\|V_e-I\|_{\rm F}<3L\delta.}              \label{eq:s6v30-link-smallness}
      \end{equation}
\end{enumerate}
\end{theorem}

\begin{proof}
By \Cref{lem:s6v27-Bianchi}, the nearest-centre plaquette labels
form a \(Z_3\)-valued two-cocycle.  The cubical box is contractible,
so its second cohomology vanishes and the cocycle is the coboundary
of a centre-valued link cochain.  Multiplying by the inverse
cochain changes every plaquette's nearest-centre label to one and
preserves its Frobenius distance from that centre.

Gauge-transform along the axial spanning tree to set all tree
links to \(I\).  A non-tree link is then the holonomy of its
fundamental loop.  The lattice non-Abelian Stokes cancellation
writes that holonomy as an ordered product of conjugates of the
plaquette holonomies in any elementary tiling of the loop.
\Cref{lem:s6v30-axial-area} uses at most \(3L\) plaquettes, and
\Cref{lem:s6v30-product} proves
\eqref{eq:s6v30-link-smallness}.
\end{proof}

\begin{firewall}[Relative boundary versus an absolute box]
\label{fire:s6v30-relative}
\Cref{thm:s6v30-axial-smallness} acts on every link of a
contractible box.  A conditional collar has fixed boundary links,
and its centre cochain must vanish on the declared fixed subcomplex.
That is the relative exactness condition of V27.  The theorem
therefore constructs the correct local analytic chart after the
centre carrier has been removed; it does not by itself prove the
PLVFC occurrence statement of V29.
\end{firewall}

\section{A three-condition weak-coupling scale}
\label{sec:s6v30-scales}

\begin{definition}[Mesoscopic Yang--Mills scales]
\label{def:s6v30-scales}
For \(\beta\ge2\), set
\begin{equation}
 \delta_\beta:=\beta^{-2/5},
 \qquad
 L_\beta:=2^{\lfloor(\log_2\beta)/10\rfloor}.           \label{eq:s6v30-scales}
\end{equation}
Thus \(L_\beta\) is dyadic and
\(\frac12\beta^{1/10}<L_\beta\le\beta^{1/10}\).
\end{definition}

\begin{theorem}[Mesoscopic scale ledger]
\label{thm:s6v30-scale-ledger}
The scales \eqref{eq:s6v30-scales} satisfy
\begin{align}
 \beta\delta_\beta^2&=\beta^{1/5}\longrightarrow\infty,
                                                               \label{eq:s6v30-rare-scale}\\
 L_\beta\delta_\beta&\le\beta^{-3/10}\longrightarrow0,
                                                               \label{eq:s6v30-chart-scale}\\
 L_\beta^3\delta_\beta&\le\beta^{-1/10}\longrightarrow0,
                                                               \label{eq:s6v30-Hessian-scale}\\
 {\beta\over L_\beta^2}&\ge\beta^{4/5}\longrightarrow\infty.
                                                               \label{eq:s6v30-coercive-scale}
\end{align}
\end{theorem}

\begin{proof}
Insert \(\delta_\beta=\beta^{-2/5}\) and
\(L_\beta\le\beta^{1/10}\).  The four exponents are respectively
\(1/5,-3/10,-1/10,4/5\).
\end{proof}

\begin{corollary}[Bulk noncentral defects are negligible on one
mesoscopic block]
\label{cor:s6v30-block-defects}
For the periodic Wilson measure and a fixed \(L_\beta\)-block,
\begin{equation}
 \mathbb P\{\hbox{some block plaquette lies outside all
 centre tubes of radius }\delta_\beta\}
 \le
 6(L_\beta+1)^4p_{\rm plaq}(\beta,\delta_\beta),        \label{eq:s6v30-block-defects}
\end{equation}
and the right-hand side tends to zero faster than every negative
power of \(\beta\).
\end{corollary}

\begin{proof}
Use the union bound and \Cref{cor:s6v28-all-colours}.  From
\eqref{eq:s6v28-Q},
\[
 \log p_{\rm plaq}(\beta,\delta_\beta)
 \le
 -{\beta^{1/5}\over576}+O(\log\beta).
\]
The block has at most \(6(L_\beta+1)^4=O(\beta^{2/5})\)
plaquettes, which is absorbed by the stretched exponential.
\end{proof}

\begin{corollary}[Good blocks enter a shrinking link chart]
\label{cor:s6v30-shrinking-chart}
On the complement of the event in
\eqref{eq:s6v30-block-defects}, every contractible block centre
sector admits axial coordinates satisfying
\begin{equation}
 \max_e\|V_e-I\|_{\rm F}
 \le3L_\beta\delta_\beta
 \le3\beta^{-3/10}.                                    \label{eq:s6v30-shrinking-chart}
\end{equation}
\end{corollary}

\begin{proof}
For large \(\beta\), \(\delta_\beta<1/2\).  Apply
\Cref{thm:s6v30-axial-smallness} and
\eqref{eq:s6v30-chart-scale}.
\end{proof}

\section{Coercivity in the rooted analytic chart}
\label{sec:s6v30-coercivity}

Let \(d_1\) be the linear plaquette curl on Lie-algebra-valued link
cochains, restricted to the rooted axial gauge.  All norms below
are the unweighted Euclidean sums of Frobenius norms.

\begin{lemma}[Rooted discrete one-form Poincaré inequality]
\label{lem:s6v30-Poincare}
There is a dimension-only constant \(C_{\rm P}\) such that, on
every \(K_L\),
\begin{equation}
 \|A\|_2^2\le C_{\rm P}L^2\|d_1A\|_2^2                \label{eq:s6v30-Poincare}
\end{equation}
for every Lie-algebra link cochain which vanishes on the axial
tree.
\end{lemma}

\begin{proof}
For each non-tree link, the linearized fundamental-loop identity
expresses \(A_e\) as a signed sum of plaquette curls over the strips
of \Cref{lem:s6v30-axial-area}.  Apply Cauchy--Schwarz to each
strip and then sum over links.  A plaquette in a fixed coordinate
plane occurs in at most \(C L\) such strips with a fixed transverse
endpoint, while each strip has length at most \(3L\).  The resulting
incidence count is bounded by \(C_{\rm P}L^2\).  Tensoring the scalar
identity with the eight-dimensional Lie algebra does not change
the constant.
\end{proof}

\begin{theorem}[Wilson strong convexity on the mesoscopic rooted
chart]
\label{thm:s6v30-strong-convexity}
There are constants \(c_*,C_* >0\), independent of
\(\beta\) and \(L\), with the following property.  In exponential
coordinates on the rooted axial chart
\[
 \max_e\|A_e\|_{\rm F}\le C_*L\delta,
\]
if \(L^3\delta\le c_*\), then the Hessian of the Wilson action
together with the negative logarithm of product Haar density obeys
\begin{equation}
 \boxed{
 \operatorname{Hess}{\cal V}_{\beta,L}
 \ge {c_*\beta\over L^2}I}                              \label{eq:s6v30-strong-convexity}
\end{equation}
for all sufficiently large \(\beta/L^2\).
\end{theorem}

\begin{proof}
At the identity connection, the second variation of one plaquette
term is
\[
 {1\over3}\|(d_1A)_p\|_{\rm F}^2.
\]
Thus the quadratic Wilson Hessian is
\((\beta/3)d_1^*d_1\), and
\Cref{lem:s6v30-Poincare} bounds it below by
\(\beta(3C_{\rm P}L^2)^{-1}I\).

On a fixed finite-incidence lattice, third derivatives of the
plaquette writer in exponential coordinates are bounded uniformly
on a fixed identity neighbourhood.  Hence moving every link by at
most \(C_*L\delta\) changes the Hessian, in quadratic-form norm, by
at most \(C\beta L\delta\,I\).  Relative to the preceding lower
bound this costs at most \(CL^3\delta\), which is less than one
half after decreasing \(c_*\).

The exponential-coordinate Haar Jacobian is smooth and positive
on the same neighbourhood.  The Hessian of its negative logarithm
is bounded below by \(-C I\).  The condition
\(\beta/L^2\) sufficiently large absorbs this fixed debit and
proves \eqref{eq:s6v30-strong-convexity}.
\end{proof}

\begin{corollary}[A good-sector Brascamp--Lieb card]
\label{cor:s6v30-BL}
Let \(\nu_{\beta,L}^{\rm chart}\) be the Wilson law restricted to
any convex exponential-coordinate box contained in the rooted
chart of \Cref{thm:s6v30-strong-convexity}.  For every locally
Lipschitz \(F\),
\begin{equation}
 \boxed{
 \operatorname{Var}_{\nu_{\beta,L}^{\rm chart}}(F)
 \le {C L^2\over\beta}
 \int\|\nabla F\|_2^2\dd\nu_{\beta,L}^{\rm chart}.}     \label{eq:s6v30-BL}
\end{equation}
In particular, \eqref{eq:s6v30-BL} holds at
\((L,\delta)=(L_\beta,\delta_\beta)\) for all sufficiently large
\(\beta\).
\end{corollary}

\begin{proof}
The Brascamp--Lieb variance inequality on a convex domain bounds
the variance by the inverse Hessian quadratic form.  Insert
\eqref{eq:s6v30-strong-convexity}.  The scale conditions follow
from
\eqref{eq:s6v30-Hessian-scale} and
\eqref{eq:s6v30-coercive-scale}.
\end{proof}

\begin{corollary}[Canonical pair-card covariance row]
\label{cor:s6v30-pair-card}
Let \(J_v\) be the fine-fibre score in a retained-field direction
\(v\), evaluated in one convex good-sector chart.  Then
\begin{equation}
 \operatorname{Var}(J_v)
 \le {C L_\beta^2\over\beta}
 \mathbb E\|\nabla_zJ_v\|_2^2.                          \label{eq:s6v30-score-covariance}
\end{equation}
Consequently a sufficient good-chart static inequality is
\begin{equation}
 \boxed{
 \mathbb E K_{v,v}
 -
 {C L_\beta^2\over\beta}
 \mathbb E\|\nabla_zJ_v\|_2^2
 \ge m_*\|v\|^2}                                       \label{eq:s6v30-static-sufficient}
\end{equation}
for some \(m_*>0\), where \(K_{v,v}\) is the exact second
retained-field writer of the complete action.
\end{corollary}

\begin{proof}
Apply \eqref{eq:s6v30-BL} to \(F=J_v\).  The Hessian of the trace
effective action in direction \(v\) is exactly
\(\mathbb EK_{v,v}-\operatorname{Var}(J_v)\).
Substitution gives \eqref{eq:s6v30-static-sufficient}.
\end{proof}

\begin{warning}[Classical coercivity is not a mass term]
\label{warn:s6v30-no-mass}
The lower bound \(\beta/L^2\) is the Dirichlet or rooted
Poincaré eigenvalue of the curvature operator on a finite block.
The quadratic Wilson symbol is proportional to
\(d_1^*d_1\) and has no volume-independent zeroth-order term.
On increasing periodic tori its transverse eigenvalues approach
zero with momentum.  Therefore neither
\eqref{eq:s6v30-strong-convexity} nor
\eqref{eq:s6v30-BL} proves a Yang--Mills mass gap.  Their valid role
is to control the covariance debit in the finite canonical
trace-shell pair card.
\end{warning}

\section{The next exact acquisition}
\label{sec:s6v30-next}

\begin{assessment}[State after the carrier and chart reductions]
\label{ass:s6v30-state}
At the mesoscopic scale, the bare periodic Wilson law now has:
\begin{enumerate}[label=\textup{(\roman*)},leftmargin=3.5em]
\item stretched-exponentially rare noncentral block defects;
\item exact local removal of the centre cocycle on contractible
      good blocks;
\item a shrinking rooted link chart; and
\item a strongly log-concave finite-block Wilson kernel with the
      covariance estimate \eqref{eq:s6v30-BL}.
\end{enumerate}
The response-fixed programme still requires three pieces which no
companion file supplies: relative-boundary occurrence of the chart,
the actual generated writer \(\Psi\) and its score derivatives, and
a positive evaluation of
\eqref{eq:s6v30-static-sufficient}.  These are algebraic finite-card
questions, not consequences of the bulk defect probability.
\end{assessment}

\begin{programme}[V31: populate the complete good-chart pair card]
\label{prog:s6v30-V31}
\begin{enumerate}[leftmargin=3em]
\item Write the retained-field first and second writers \(J_v\)
      and \(K_{v,v}\) for the Wilson part in the exact parity-tree
      section.
\item Separate their rooted chart bounds from the still-unloaded
      generated-action writers; assign no numerical value to
      \(\Psi\) without an interaction list.
\item Use \eqref{eq:s6v30-score-covariance} to replace the unsplit
      covariance by an explicit gradient-score row.
\item Track the bad-block complement with
      \eqref{eq:s6v30-block-defects}; require only polynomial writer
      growth so the stretched exponential absorbs it.
\item Determine whether the Wilson contribution to
      \eqref{eq:s6v30-static-sufficient} has a strict transverse
      margin after the declared reference/topological form is
      subtracted.  If it cancels, return the exact cancellation
      before touching CDAC.
\end{enumerate}
\end{programme}

\begin{firewall}[Claim boundary after sixth-series V30]
\label{fire:s6v30-boundary}
V30 performs the compulsory SM V53 and NS V3 audit; proves a
quantitative centre-normalized non-Abelian axial chart; selects
dyadic weak-coupling scales satisfying defect, chart, nonlinear
Hessian, and coercivity requirements; proves mesoscopic
noncentral-block rarity; proves rooted Wilson strong convexity and
a good-chart Brascamp--Lieb covariance bound; and turns that bound
into an explicit sufficient row for the canonical trace-shell pair
card.

V30 does not prove relative-boundary chart occurrence, load or
bound the actual generated action, establish the positive static
row, prove a good-gate transmission constant, construct the
continuum measure, or prove the full Yang--Mills mass gap.
\end{firewall}

\begin{theorem}[Sixth-series V30 result ledger]
\label{thm:s6v30-results}
The claims in \Cref{fire:s6v30-boundary} follow from
\Cref{audit:s6v30-SM,
audit:s6v30-NS,
lem:s6v30-product,
lem:s6v30-axial-area,
thm:s6v30-axial-smallness,
thm:s6v30-scale-ledger,
cor:s6v30-block-defects,
cor:s6v30-shrinking-chart,
lem:s6v30-Poincare,
thm:s6v30-strong-convexity,
cor:s6v30-BL,
cor:s6v30-pair-card}.
\end{theorem}

\begin{proof}
Each ledger item is the conclusion of the cited result.
\end{proof}

\paragraph{Parent-version record.}
V30 was formed from
\texttt{Renewal\_Yang\_Mills\_Mass\_Gap\_Research\_Notes\_V29.tex}.
The V29 parent had \(13\,056\) logical source lines,
\(462\,845\) bytes, and SHA--256
\[
 \texttt{32fef34da4b92b56cf90bfc0d14c28164fe92f6d3f7253d8dd2d5c9294f52e3d}.
\]
The next compulsory companion audit is V35.

\clearpage
\chapter{Sixth-series V31: a cofinal pure-Wilson static card}
\label{chap:s6v31}

\section{Nonduplication audit of the attached V35 programme}
\label{sec:s6v31-audit}

Before using the V30 covariance card, the complete attached file
\[
 \texttt{Renewal\_Yang\_Mills\_Reduced\_Fibre\_Wilson\_Shell\_Symbol\_and\_Local\_Vacuum\_Programme\_V35.tex}
\]
was re-read at its Wilson V2 and static V35 layers.  Four results
are already proved there and are therefore imported rather than
rederived:
\begin{enumerate}[label=\textup{(\roman*)},leftmargin=3.5em]
\item on the \(45M^4\)-link pinned reduced fibre at \(W=I\),
      \begin{equation}
       {1\over120}D(z)\le V_M(z,I)\le4D(z),
       \qquad
       D(z)=\sum_e\|z_e-I\|_{\rm F}^2;                 \label{eq:s6v31-import-coercivity}
      \end{equation}
\item the vertical unit-Wilson Hessian at the identity satisfies
      \begin{equation}
       {1\over60}I\preceq A_M\preceq8I;               \label{eq:s6v31-import-vertical}
      \end{equation}
\item on a nonzero axial coarse angle \(q\), the unit flat Schur
      reference is
      \begin{equation}
       s_0(q)
       ={t(t+8)(t+24)\over9(t^2+18t+16)},
       \qquad t=\sin^2(q/2),                           \label{eq:s6v31-import-reference}
      \end{equation}
      and
      \begin{equation}
       {5\over28}\le{s_0(q)\over4\sin^2(q/2)}
       \le{1\over3};                                   \label{eq:s6v31-import-Ward}
      \end{equation}
\item for every fixed \(M\), the pure-Wilson trace Hessian is
      \(\beta\mathsf S_M+O_M(1)\), but the source explicitly does
      not bound the growth of the remainder with \(M\).
\end{enumerate}
The attached V35 layer also proves that the complete unsplit pair
card already contains its covariance tail.  V31 therefore adds no
second tail.  Its new task is exactly the missing one in item
\textup{(iv)}: prove a cofinal margin on one explicit growing
\((\beta,M)\) schedule.

\section{Uniform concentration of the pinned reduced fibre}
\label{sec:s6v31-concentration}

Retain the explicit Haar integral constants from the attached V2
calculation:
\begin{equation}
 C_-={3e^{-43/32}\over1024\pi^8},
 \qquad
 C_+={25\pi^7\over512},
 \qquad
 {\cal C}_0:={C_+\over C_-}\,960^4.                    \label{eq:s6v31-Haar-constants}
\end{equation}
They obey
\begin{equation}
 C_-t^{-4}\le
 I(t):=\int_{\SU(3)}e^{-t\|U-I\|_{\rm F}^2}\dd U
 \le C_+t^{-4}
 \qquad(t\ge1).                                        \label{eq:s6v31-I-bounds}
\end{equation}

\begin{lemma}[A common principal logarithm chart]
\label{lem:s6v31-log-chart}
There are \(r_0,c_{\exp}>0\), depending only on \(\SU(3)\), such
that for \(0<r\le r_0\) the set
\[
 {\cal G}_{M,r}
 :=
 \{z:z_e=\exp x_e,\ \|x_e\|_{\rm F}<r
       \text{ for every reduced link }e\}
\]
is a product of principal exponential charts and
\begin{equation}
 z\notin{\cal G}_{M,r}
 \quad\Longrightarrow\quad
 D(z)\ge c_{\exp}^2r^2.                                \label{eq:s6v31-log-complement}
\end{equation}
\end{lemma}

\begin{proof}
The exponential map is a diffeomorphism on a fixed algebra ball
about zero and its differential at zero is the identity.  Hence
\(\|\exp X-I\|_{\rm F}\ge c_{\exp}\|X\|_{\rm F}\) on a smaller
ball.  By compactness, group elements outside that exponential
neighbourhood have a fixed positive distance from \(I\).  Decrease
\(c_{\exp}\) so the same implication holds for every
\(0<r\le r_0\), and sum over links.
\end{proof}

\begin{theorem}[Growing-volume chart-complement bound]
\label{thm:s6v31-chart-tail}
Let \(\nu_{\beta,M}\) be the pure-Wilson law on the pinned reduced
fibre at \(W=I\), and let \(n_M=45M^4\) be its number of group
coordinates.  If
\(\beta\ge240\) and \(0<r\le r_0\), then
\begin{equation}
 \boxed{
 \nu_{\beta,M}({\cal G}_{M,r}^{\,c})
 \le
 \exp\left\{
 -{c_{\exp}^2\beta r^2\over240}
 +n_M\log{\cal C}_0
 \right\}.}                                             \label{eq:s6v31-chart-tail}
\end{equation}
\end{theorem}

\begin{proof}
By the upper half of
\eqref{eq:s6v31-import-coercivity},
\[
 Z_{\beta,M}
 \ge\int e^{-4\beta D(z)}\dd m(z)
 =I(4\beta)^{n_M}
 \ge\{C_-(4\beta)^{-4}\}^{n_M}.
\]
On \({\cal G}_{M,r}^{\,c}\), the lower half of
\eqref{eq:s6v31-import-coercivity} and
\eqref{eq:s6v31-log-complement} give
\[
 e^{-\beta V_M}
 \le
 e^{-c_{\exp}^2\beta r^2/240}
 e^{-\beta D/240}.
\]
Its integral is at most
\[
 e^{-c_{\exp}^2\beta r^2/240}
 I(\beta/240)^{n_M}
 \le
 e^{-c_{\exp}^2\beta r^2/240}
 \{C_+(\beta/240)^{-4}\}^{n_M}.
\]
Division by the preceding lower bound produces precisely
\({\cal C}_0^{n_M}\).
\end{proof}

\section{Uniform local Schur stability}
\label{sec:s6v31-Schur}

Use product logarithm coordinates \(x\) for the reduced links and
coarse logarithm coordinates \(w\) at \(W=I\).  For the unit Wilson
action write
\begin{equation}
 A(x):=D_x^2V_M(x,0),\quad
 B(x):=D_xD_wV_M(x,0),\quad
 C(x):=D_w^2V_M(x,0),                                  \label{eq:s6v31-ABC}
\end{equation}
and
\begin{equation}
 {\cal Q}(x):=C(x)-B(x)^*A(x)^{-1}B(x)                 \label{eq:s6v31-point-Schur}
\end{equation}
whenever \(A(x)>0\).

\begin{lemma}[Volume-independent derivative bounds]
\label{lem:s6v31-derivative-bounds}
There are \(r_1,C_1,C_2>0\), independent of \(M\), such that on
\({\cal G}_{M,r_1}\),
\begin{align}
 \|A(x)-A(0)\|_{\rm op}
 +\|B(x)-B(0)\|_{\rm op}
 +\|C(x)-C(0)\|_{\rm op}
 &\le C_1\|x\|_\infty,                                 \label{eq:s6v31-block-Lipschitz}\\
 \|B(x)\|_{\rm op}+\|C(x)\|_{\rm op}&\le C_2.           \label{eq:s6v31-BC-bound}
\end{align}
Here the coarse Fourier transform and all retained polarizations
are normalized unitarily.
\end{lemma}

\begin{proof}
Each plaquette writer contains four link factors and has uniformly
bounded derivatives of every fixed order on the compact group.
Each fine link belongs to six plaquettes, and each designated
retained link has the same bounded local incidence.  Consequently
the matrices of third derivatives have row and column sums bounded
by a dimension-only constant.  The Schur test gives a
volume-independent operator norm.  Integrating the third derivative
along the segment from \(0\) to \(x\) proves
\eqref{eq:s6v31-block-Lipschitz}; the zeroth-order values and the
same incidence bound give \eqref{eq:s6v31-BC-bound}.
\end{proof}

\begin{lemma}[Pointwise Schur stability]
\label{lem:s6v31-Schur-stability}
After decreasing \(r_1\), \(A(x)\succeq I/120\) on
\({\cal G}_{M,r_1}\), and
\begin{equation}
 \boxed{
 \|{\cal Q}(x)-\mathsf S_M\|_{\rm op}
 \le C_{\rm S}\|x\|_\infty}                            \label{eq:s6v31-Schur-stability}
\end{equation}
with \(C_{\rm S}\) independent of \(M\).
\end{lemma}

\begin{proof}
\eqref{eq:s6v31-import-vertical} and
\eqref{eq:s6v31-block-Lipschitz} give the first assertion.
The resolvent identity
\[
 A(x)^{-1}-A(0)^{-1}
 =-A(x)^{-1}\{A(x)-A(0)\}A(0)^{-1}
\]
and the uniform inverse bounds make the Schur-complement map
\((A,B,C)\mapsto C-B^*A^{-1}B\) uniformly Lipschitz.
At \(x=0\) this Schur complement is exactly the imported flat
reference \(\mathsf S_M\).  This proves
\eqref{eq:s6v31-Schur-stability}.
\end{proof}

\section{The truncated trace Hessian}
\label{sec:s6v31-truncated}

Let \(\nu_{\beta,M}^{(r)}\) be the pure-Wilson law restricted to
the convex product algebra ball
\(\|x_e\|_{\rm F}<r\), with the exponential-coordinate Haar
Jacobian retained.  Its integration domain does not depend on the
coarse field \(w\).  Let \(H_{\beta,M}^{(r)}\) be the Hessian at
\(w=0\) of the corresponding truncated trace action.

\begin{lemma}[Haar-Hessian debit]
\label{lem:s6v31-Haar-debit}
On a fixed principal chart, the Hessian
\[
 R_{\rm H}(x):=D_x^2\{-\log J(x)\}
\]
of the negative logarithmic product-Haar density has
\(\|R_{\rm H}(x)\|_{\rm op}\le C_{\rm H}\), independently of
\(M\).  For \(r\le r_1\) and all sufficiently large \(\beta\),
\begin{equation}
 \beta A(x)+R_{\rm H}(x)\succeq{\beta\over240}I.        \label{eq:s6v31-total-vertical}
\end{equation}
\end{lemma}

\begin{proof}
The logarithmic Haar Jacobian is a sum of identical one-link
analytic functions, so its Hessian is block diagonal with a
uniform block norm.  Combine this with
\(A(x)\succeq I/120\) from
\Cref{lem:s6v31-Schur-stability}.
\end{proof}

\begin{theorem}[Brascamp--Lieb lower bound for the truncated
physical card]
\label{thm:s6v31-truncated-card}
There is a volume-independent constant \(C_{\rm T}\) such that,
for every normalized retained direction \(v\),
\begin{equation}
 \boxed{
 H_{\beta,M}^{(r)}[v,v]
 \ge
 \beta\,\mathsf S_M[v,v]
 -C_{\rm T}\beta r-C_{\rm T}.}                         \label{eq:s6v31-truncated-card}
\end{equation}
\end{theorem}

\begin{proof}
The retained score and deterministic second writer are
\[
 j_v=\beta D_wV_M[v],
 \qquad
 K_{v,v}=\beta D_w^2V_M[v,v].
\]
The exact trace identity gives
\[
 H_{\beta,M}^{(r)}[v,v]
 =\mathbb E K_{v,v}-\operatorname{Var}(j_v).
\]
By \eqref{eq:s6v31-total-vertical}, the truncated logarithmic
potential is strongly convex on a convex domain.  The
Brascamp--Lieb inequality therefore gives
\[
 \operatorname{Var}(j_v)
 \le
 \mathbb E\!\left[
  \beta^2
  \langle B(x)v,
   \{\beta A(x)+R_{\rm H}(x)\}^{-1}B(x)v\rangle
 \right].
\]
The inverse identity and
\Cref{lem:s6v31-derivative-bounds,lem:s6v31-Haar-debit}
imply
\[
 \beta^2B^*
 \{\beta A+R_{\rm H}\}^{-1}B
 \preceq
 \beta B^*A^{-1}B+C I.
\]
Consequently
\[
 H_{\beta,M}^{(r)}[v,v]
 \ge\beta\,\mathbb E{\cal Q}(x)[v,v]-C.
\]
Use \eqref{eq:s6v31-Schur-stability} and
\(\|x\|_\infty<r\).
\end{proof}

\section{Restoring the chart complement}
\label{sec:s6v31-restore}

\begin{lemma}[Uniform writer envelopes]
\label{lem:s6v31-writer-envelopes}
For every normalized coarse Fourier polarization \(v\),
\begin{equation}
 \sup_z|K_{v,v}(z)|\le C\beta,
 \qquad
 \sup_z|j_v(z)|\le C\beta M^2,                         \label{eq:s6v31-writer-envelopes}
\end{equation}
with \(C\) independent of \(M,\beta\).
\end{lemma}

\begin{proof}
A normalized plane wave on \(M^4\) coarse links has pointwise
amplitude \(M^{-2}\), squared \(\ell^2\)-norm one, and
\(\ell^1\)-norm \(M^2\).  Every retained link meets six
plaquettes.  The second writer is a bounded finite-incidence
quadratic form, so its operator norm is \(C\beta\).  The first
writer is a sum of uniformly bounded local terms times the
polarization; its absolute value is at most
\(C\beta\|v\|_1=C\beta M^2\).
\end{proof}

\begin{lemma}[Mixture restoration bound]
\label{lem:s6v31-mixture}
Let
\(\epsilon_{\beta,M,r}:=\nu_{\beta,M}({\cal G}_{M,r}^{\,c})\).
Then
\begin{equation}
 \boxed{
 H_{\beta,M}[v,v]
 \ge
 H_{\beta,M}^{(r)}[v,v]
 -C\epsilon_{\beta,M,r}\beta^2M^4.}                   \label{eq:s6v31-mixture}
\end{equation}
\end{lemma}

\begin{proof}
Disintegrate the full law between \({\cal G}_{M,r}\) and its
complement.  The deterministic expectation changes by at most
\(2\epsilon\sup|K_{v,v}|\).  The variance decomposition
\[
 \operatorname{Var}(j_v)
 =p\operatorname{Var}_{\cal G}(j_v)
 +(1-p)\operatorname{Var}_{{\cal G}^c}(j_v)
 +p(1-p)(m_{\cal G}-m_{{\cal G}^c})^2
\]
shows that the full variance exceeds the good conditional variance
by at most \(C\epsilon\sup|j_v|^2\).  Insert
\eqref{eq:s6v31-writer-envelopes} and enlarge \(C\).
\end{proof}

\section{An explicit cofinal schedule}
\label{sec:s6v31-schedule}

\begin{definition}[Cofinal pure-Wilson card schedule]
\label{def:s6v31-schedule}
For sufficiently large \(\beta\), set
\begin{equation}
 M_\beta
 :=
 2^{\max\{2,\lfloor(\log_2\beta)/20\rfloor\}},
 \qquad
 r_\beta:=\beta^{-1/5}.                                \label{eq:s6v31-schedule}
\end{equation}
\end{definition}

\begin{lemma}[Schedule ledger]
\label{lem:s6v31-schedule-ledger}
Along \eqref{eq:s6v31-schedule},
\begin{align}
 M_\beta&\longrightarrow\infty,
 &M_\beta^2r_\beta&\longrightarrow0,                   \label{eq:s6v31-margin-scale}\\
 n_{M_\beta}&\le45\beta^{1/5},
 &\beta r_\beta^2&=\beta^{3/5}.                        \label{eq:s6v31-tail-scale}
\end{align}
Furthermore,
\begin{equation}
 \epsilon_{\beta,M_\beta,r_\beta}
 \le
 \exp\{-c\beta^{3/5}+C\beta^{1/5}\}                   \label{eq:s6v31-stretched-tail}
\end{equation}
for constants \(c,C>0\).
\end{lemma}

\begin{proof}
The dyadic definition gives
\(\frac12\beta^{1/20}<M_\beta\le\beta^{1/20}\) eventually.
Thus \(M_\beta^2r_\beta\le\beta^{-1/10}\).
The remaining algebra is immediate, and
\eqref{eq:s6v31-stretched-tail} is
\Cref{thm:s6v31-chart-tail}.
\end{proof}

\begin{theorem}[Cofinal pure-Wilson axial static margin]
\label{thm:s6v31-cofinal-margin}
There is \(\beta_0<\infty\) such that, for
\(\beta\ge\beta_0\), \(M=M_\beta\), and every nonzero allowed axial
angle \(q\),
\begin{equation}
 \boxed{
 H_{\beta,M}(q)
 \succeq {1\over2}\,\beta s_0(q)I}                    \label{eq:s6v31-cofinal-margin}
\end{equation}
on the physical transverse--colour block.  Equivalently, the
complete pure-Wilson full trace card has a uniform normalized
margin at least \(1/2\) on this growing dyadic schedule.
\end{theorem}

\begin{proof}
By \eqref{eq:s6v31-import-Ward} and
\(\sin(\pi/M)\ge2/M\) for \(M\ge2\),
\begin{equation}
 \min_{q\ne0}s_0(q)
 \ge{5\over7}\sin^2(\pi/M)
 \ge{20\over7M^2}.                                    \label{eq:s6v31-smin}
\end{equation}
Combine
\Cref{thm:s6v31-truncated-card,lem:s6v31-mixture}:
\[
 H_{\beta,M}[v,v]
 \ge
 \beta s_0(q)
 -C\beta r_\beta-C
 -C\epsilon_{\beta,M,r_\beta}\beta^2M^4.
\]
Relative to \(\beta s_0(q)\), the three debits are bounded by
\[
 CM_\beta^2r_\beta,\qquad
 C{M_\beta^2\over\beta},\qquad
 C\epsilon_{\beta,M_\beta,r_\beta}\beta M_\beta^6.
\]
The first two vanish by
\Cref{lem:s6v31-schedule-ledger}; the third vanishes by the
stretched exponential \eqref{eq:s6v31-stretched-tail}.
Their sum is below \(1/2\) for all sufficiently large \(\beta\).
The pure Wilson trace action has the central axial symmetry, so
the physical block scalarizes and the polarization estimate proves
the matrix inequality.
\end{proof}

\begin{assessment}[What V31 closes]
\label{ass:s6v31-closes}
The attached V35 source stopped because its fixed-\(M\)
\(O_M(1)\) remainder had no cofinal growth bound.  V31 avoids a
dimension-blind Laplace remainder: global fibre coercivity gives a
stretched-exponential chart tail, while the exact vertical Hessian
and finite incidence give a direct Brascamp--Lieb Schur estimate
inside the chart.  Choosing
\(M_\beta\asymp\beta^{1/20}\) makes every error small relative to
the smallest axial reference \(\beta/M_\beta^2\).  Thus the
\emph{pure-Wilson full trace} static card is now cofinally positive
on an explicit growing schedule.
\end{assessment}

\begin{firewall}[Why the response-matched static card is still open]
\label{fire:s6v31-response}
\Cref{thm:s6v31-cofinal-margin} is not yet the selected
response-matched local card.  The actual trajectory may contain a
generated action \(\Psi_{\beta,M}\); its vertical Hessian, mixed
score, deterministic retained jet, and head--tail covariance have
not been bounded.  The full finite-torus trace may also contain a
topological packet which must be subtracted on the same carrier
before declaring a local contractible margin.  Finally, the
canonical response root must be shown to enter a schedule on which
the present inequalities hold.  No one of these rows follows from
the pure-Wilson theorem.
\end{firewall}

\begin{corollary}[Exact remaining perturbation budget]
\label{cor:s6v31-budget}
On the schedule \eqref{eq:s6v31-schedule}, let
\({\cal R}_{\Psi,\rm top}(q)\) denote the single exact signed
difference between the complete response-matched local trace
Hessian and the pure-Wilson full trace Hessian, including the
generated action, all mixed score covariances, and the topological
subtraction once.  If
\begin{equation}
 \boxed{
 \|{\cal R}_{\Psi,\rm top}(q)\|_{\rm op}
 \le{1\over4}\beta s_0(q)
 \quad\hbox{uniformly for }q\ne0,}                     \label{eq:s6v31-budget}
\end{equation}
then the response-matched local static card has normalized margin
at least \(1/4\).
\end{corollary}

\begin{proof}
Subtract \eqref{eq:s6v31-budget} from
\eqref{eq:s6v31-cofinal-margin}.  The residual is defined as one
complete signed difference, so the no-double-charging rule of the
attached V35 programme is respected.
\end{proof}

\begin{programme}[V32: attack the generated/topological residual]
\label{prog:s6v31-V32}
\begin{enumerate}[leftmargin=3em]
\item Express the response-matched action as the pure-Wilson head
      plus one declared generated/topological residual on the same
      canonical root and carrier.
\item Translate the V29 local oscillation debit into vertical,
      mixed, and retained second-jet norms on the
      \(M_\beta\)-schedule.
\item Use the complete head--tail score Gram from the attached V35
      source; do not set the mixed covariance to zero.
\item Prove \eqref{eq:s6v31-budget}, or identify the first
      generated writer whose Ward-relative norm is not
      \(o(\beta s_0)\).
\item Only after the local static card passes, return to the
      good-gate transmission and OS continuum landing.
\end{enumerate}
\end{programme}

\begin{firewall}[Claim boundary after sixth-series V31]
\label{fire:s6v31-boundary}
V31 audits and imports the stronger nonduplicated V35 pinned-fibre
results; proves a growing-volume chart-complement estimate; proves
volume-independent joint derivative and Schur stability; derives a
truncated trace-Hessian lower bound by the matrix
Brascamp--Lieb inequality; restores the compact-fibre complement;
and proves a cofinal normalized pure-Wilson axial trace margin of
\(1/2\) on the explicit dyadic schedule
\(M_\beta\asymp\beta^{1/20}\).

V31 does not bound the actual generated/topological residual, place
the canonical response root on this schedule, prove the local
response-matched static card, establish good-gate transmission,
construct a reflection-positive continuum measure, or prove the
full Yang--Mills mass gap.
\end{firewall}

\begin{theorem}[Sixth-series V31 result ledger]
\label{thm:s6v31-results}
The claims in \Cref{fire:s6v31-boundary} follow from
\Cref{thm:s6v31-chart-tail,
lem:s6v31-derivative-bounds,
lem:s6v31-Schur-stability,
lem:s6v31-Haar-debit,
thm:s6v31-truncated-card,
lem:s6v31-writer-envelopes,
lem:s6v31-mixture,
lem:s6v31-schedule-ledger,
thm:s6v31-cofinal-margin,
cor:s6v31-budget}.
\end{theorem}

\begin{proof}
Each ledger item is the conclusion of the cited result.
\end{proof}

\paragraph{Parent-version record.}
V31 was formed from
\texttt{Renewal\_Yang\_Mills\_Mass\_Gap\_Research\_Notes\_V30.tex}.
The V30 parent had \(13\,532\) logical source lines,
\(480\,739\) bytes, and SHA--256
\[
 \texttt{f0e862e04e70b666ecee578638c4cc7e937d12d1b88fd9687741bc8c5d55e066}.
\]
The next compulsory companion audit is V35.

\clearpage
\chapter{Sixth-series V32: unsplit stability under the actual
generated action}
\label{chap:s6v32}

\section{Source audit and choice of route}
\label{sec:s6v32-audit}

The V31 follow-up audit searched both the attached reduced-fibre
V35 programme and the Grand Tensor V067 manuscript for a literal
interaction dictionary
\[
 \Psi_{\beta,M}(z,W)=\sum_A\Psi_{\beta,M,A}(z_A,W)
\]
with regulator-uniform first and second writers.  Neither supplied
such a dictionary.  Both instead require the complete finite action
to be retained symbolically.  Consequently no numerical
\(\Psi\)-bound is inferred from the words \emph{generated} or
\emph{response matched}.

V32 uses the admissible unsplit route of the attached V35
programme.  It gives sufficient norms directly on the complete
\(\Psi\), inserts them into the full joint Schur complement, and
uses one Brascamp--Lieb estimate for the complete score.  There is
no separately added covariance tail.

\section{The generated joint-jet debit}
\label{sec:s6v32-jets}

On the principal chart of V31, write
\begin{equation}
 P_\Psi(x):=D_x^2\Psi(x,0),\qquad
 Q_\Psi(x):=D_xD_w\Psi(x,0),\qquad
 R_\Psi(x):=D_w^2\Psi(x,0).                            \label{eq:s6v32-PQR}
\end{equation}
For \(0<r\le r_1\), define
\begin{equation}
 \begin{split}
 {\mathfrak d}_{\Psi}(M,\beta,r)
 :=\sup_{\|x\|_\infty<r}\bigg\{&
 \|P_\Psi(x)\|_{\rm op}
 +\|Q_\Psi(x)\|_{\rm op}
 +\|R_\Psi(x)\|_{\rm op}\\
 &+{\|Q_\Psi(x)\|_{\rm op}^2\over\beta}
 \bigg\}.
 \end{split}                                           \label{eq:s6v32-jet-debit}
\end{equation}
All operators use the same unitary coarse Fourier normalization as
V31.

\begin{lemma}[Stability of a joint Schur complement]
\label{lem:s6v32-Schur-perturbation}
Let \(A,B,C\) be the unit-Wilson blocks in
\eqref{eq:s6v31-ABC}.  Suppose
\begin{equation}
 \|P_\Psi(x)\|_{\rm op}\le{\beta\over240}               \label{eq:s6v32-vertical-small}
\end{equation}
throughout the chart.  Put
\[
 A_F=\beta A+P_\Psi,\qquad
 B_F=\beta B+Q_\Psi,\qquad
 C_F=\beta C+R_\Psi.
\]
Then \(A_F\succeq\beta I/240\), and
\begin{equation}
 \boxed{
 \left\|
 C_F-B_F^*A_F^{-1}B_F
 -\beta\{C-B^*A^{-1}B\}
 \right\|_{\rm op}
 \le C{\mathfrak d}_{\Psi}(M,\beta,r)}                 \label{eq:s6v32-Schur-perturbation}
\end{equation}
with \(C\) independent of \(M,\beta\).
\end{lemma}

\begin{proof}
\Cref{lem:s6v31-Schur-stability} gives
\(A\succeq I/120\); hence
\eqref{eq:s6v32-vertical-small} gives
\(A_F\succeq\beta I/240\).  The inverse identity yields
\[
 A_F^{-1}-(\beta A)^{-1}
 =-A_F^{-1}P_\Psi(\beta A)^{-1},
\]
so its norm is at most
\(C\|P_\Psi\|/\beta^2\).  Expand
\[
 (\beta B+Q_\Psi)^*A_F^{-1}(\beta B+Q_\Psi)
\]
into the Wilson--Wilson, two mixed, and generated--generated
terms.  Use the inverse estimate,
\(\|B\|\le C_2\), and
\(\|A_F^{-1}\|\le C/\beta\).  The four respective debits are
\(C\|P_\Psi\|\), \(C\|Q_\Psi\|\),
\(C\|Q_\Psi\|\), and
\(C\|Q_\Psi\|^2/\beta\).  Add
\(\|R_\Psi\|\).
\end{proof}

\begin{lemma}[The Haar inverse costs only order one]
\label{lem:s6v32-Haar-inverse}
Under \eqref{eq:s6v32-vertical-small}, for all sufficiently large
\(\beta\),
\[
 A_F+R_{\rm H}\succeq{\beta\over480}I,
\]
and
\begin{equation}
 \left\|
 B_F^*(A_F+R_{\rm H})^{-1}B_F
 -B_F^*A_F^{-1}B_F
 \right\|_{\rm op}
 \le
 C\left(1+{\|Q_\Psi\|_{\rm op}\over\beta}\right)^2.    \label{eq:s6v32-Haar-inverse}
\end{equation}
\end{lemma}

\begin{proof}
The first statement follows from the uniform Haar-Hessian bound in
\Cref{lem:s6v31-Haar-debit}.  The resolvent identity makes the
difference of inverses \(O(\beta^{-2})\).  Since
\(\|B_F\|\le C\beta+\|Q_\Psi\|\), multiplication on both sides
gives \eqref{eq:s6v32-Haar-inverse}.
\end{proof}

\section{Unsplit complete-score control in the chart}
\label{sec:s6v32-unsplit}

Let
\[
 F_{\beta,M}(x,w)=\beta V_M(x,w)+\Psi_{\beta,M}(x,w)
\]
and let \(H_{\beta,\Psi,M}^{(r)}\) be the trace Hessian obtained by
integrating \(e^{-F_{\beta,M}}\) over the fixed convex product
algebra ball \(\|x_e\|<r\), with the Haar Jacobian included.

\begin{theorem}[Complete generated-action chart bound]
\label{thm:s6v32-chart-card}
Assume \eqref{eq:s6v32-vertical-small}.  For every normalized
retained direction \(v\),
\begin{equation}
 \boxed{
 \begin{split}
 H_{\beta,\Psi,M}^{(r)}[v,v]
 \ge{}&
 \beta\mathsf S_M[v,v]
 -C\beta r
 -C{\mathfrak d}_{\Psi}(M,\beta,r)\\
 &-C\left(1+
 {\sup_{\|x\|_\infty<r}\|Q_\Psi(x)\|_{\rm op}\over\beta}
 \right)^2.
 \end{split}}                                           \label{eq:s6v32-chart-card}
\end{equation}
This inequality contains the covariance of the complete score
\(\beta D_wV_M+D_w\Psi\) exactly once.
\end{theorem}

\begin{proof}
The vertical Hessian of the complete logarithmic density is
\(A_F+R_{\rm H}\).  Its strong convexity and the matrix
Brascamp--Lieb inequality give
\[
 \operatorname{Var}(D_wF[v])
 \le
 \mathbb E
 \langle B_Fv,(A_F+R_{\rm H})^{-1}B_Fv\rangle.
\]
Subtract this single covariance bound from
\(\mathbb E C_F[v,v]\).  By
\Cref{lem:s6v32-Haar-inverse}, the result is bounded below by the
pointwise complete-action Schur complement, up to the last line of
\eqref{eq:s6v32-chart-card}.  Apply
\Cref{lem:s6v32-Schur-perturbation} and then
\eqref{eq:s6v31-Schur-stability}.  No head score and tail score
have been separated, so no mixed covariance is omitted or charged
twice.
\end{proof}

\section{Concentration transfer and complement restoration}
\label{sec:s6v32-complement}

\begin{lemma}[Generated-action chart probability]
\label{lem:s6v32-chart-probability}
Let
\[
 \operatorname{osc}\Psi_{\beta,M}
 :=\sup_z\Psi_{\beta,M}(z,0)-\inf_z\Psi_{\beta,M}(z,0).
\]
For the full generated law at \(W=I\),
\begin{equation}
 \boxed{
 \nu_{\beta,\Psi,M}({\cal G}_{M,r}^{\,c})
 \le
 e^{\operatorname{osc}\Psi_{\beta,M}}
 \nu_{\beta,M}({\cal G}_{M,r}^{\,c}).}                 \label{eq:s6v32-chart-probability}
\end{equation}
\end{lemma}

\begin{proof}
The likelihood relative to the pure-Wilson law is
\[
 {e^{-\Psi}\over\mathbb E_{\beta,M}e^{-\Psi}}.
\]
Its supremum is at most
\(e^{\sup\Psi-\inf\Psi}\).  Integrate over the chart complement.
\end{proof}

\begin{definition}[Polynomial global writer card, PGWC]
\label{def:s6v32-PGWC}
PGWC means that there are fixed exponents \(a,b<\infty\) and a
constant \(C\) such that, for every normalized selected axial
polarization \(v\),
\begin{equation}
 \sup_z\left\{
 |D_wF_{\beta,M}(z,0)[v]|+
 |D_w^2F_{\beta,M}(z,0)[v,v]|
 \right\}
 \le C\beta^aM^b.                                      \label{eq:s6v32-PGWC}
\end{equation}
\end{definition}

\begin{lemma}[Complete-action mixture debit]
\label{lem:s6v32-mixture}
Under PGWC, if
\[
 \epsilon_{\beta,\Psi,M,r}
 :=\nu_{\beta,\Psi,M}({\cal G}_{M,r}^{\,c}),
\]
then
\begin{equation}
 H_{\beta,\Psi,M}[v,v]
 \ge
 H_{\beta,\Psi,M}^{(r)}[v,v]
 -C\epsilon_{\beta,\Psi,M,r}\beta^{2a}M^{2b}.          \label{eq:s6v32-mixture}
\end{equation}
\end{lemma}

\begin{proof}
Repeat the exact two-component variance decomposition in
\Cref{lem:s6v31-mixture}, now for the complete first and second
writers.  The deterministic difference is linear in the bad mass;
the within-bad and between-component covariance debits are bounded
by the square of the global first-writer envelope.  PGWC absorbs
both after increasing \(C,a,b\) if necessary.
\end{proof}

\section{A cofinal generated-action stability theorem}
\label{sec:s6v32-cofinal}

Use the schedule
\((M_\beta,r_\beta)\) from
\eqref{eq:s6v31-schedule}.

\begin{definition}[Generated cofinal short card, GCSC]
\label{def:s6v32-GCSC}
The actual generated action satisfies GCSC on the selected root
sequence if:
\begin{align}
 \operatorname{osc}\Psi_{\beta,M_\beta}
 &=o(\beta r_\beta^2),                                 \label{eq:s6v32-GCSC-osc}\\
 {M_\beta^2\over\beta}\,
 {\mathfrak d}_{\Psi}
 (M_\beta,\beta,r_\beta)
 &\longrightarrow0,                                   \label{eq:s6v32-GCSC-jets}\\
 {1\over\beta}
 \sup_{\|x\|_\infty<r_\beta}
 \|Q_\Psi(x)\|_{\rm op}
 &\longrightarrow0,                                   \label{eq:s6v32-GCSC-mixed}\\
 \Psi_{\beta,M_\beta}
 &\text{ satisfies PGWC}.                              \label{eq:s6v32-GCSC-global}
\end{align}
\end{definition}

\begin{theorem}[Cofinal lower stability of the full trace card]
\label{thm:s6v32-generated-margin}
Assume GCSC.  Then there is
\(\varepsilon_\beta=o(\beta/M_\beta^2)\),
\(\varepsilon_\beta\ge0\), such that, uniformly on every nonzero
selected axial block,
\begin{equation}
 \boxed{
 H_{\beta,\Psi,M_\beta}(q)
 \succeq
 \beta s_0(q)I-\varepsilon_\beta I,}                    \label{eq:s6v32-generated-asymptotic}
\end{equation}
and therefore
\begin{equation}
 \boxed{
 H_{\beta,\Psi,M_\beta}(q)
 \succeq{1\over2}\beta s_0(q)I}                        \label{eq:s6v32-generated-margin}
\end{equation}
for all sufficiently large \(\beta\).
\end{theorem}

\begin{proof}
\eqref{eq:s6v32-GCSC-jets} implies
\(\|P_\Psi\|/\beta\to0\), so
\eqref{eq:s6v32-vertical-small} holds eventually.
By
\Cref{lem:s6v32-chart-probability,thm:s6v31-chart-tail},
\[
 \epsilon_{\beta,\Psi,M_\beta,r_\beta}
 \le
 \exp\{-c\beta^{3/5}+C\beta^{1/5}
       +o(\beta^{3/5})\}.
\]
PGWC makes the mixture debit
\eqref{eq:s6v32-mixture} smaller than every negative power of
\(\beta\).

In \eqref{eq:s6v32-chart-card}, the Wilson chart debit satisfies
\[
 \beta r_\beta
 =o(\beta/M_\beta^2)
\]
by \eqref{eq:s6v31-margin-scale}.  The generated jet debit is
little-oh of \(\beta/M_\beta^2\) by
\eqref{eq:s6v32-GCSC-jets}.  The Haar-inverse term is
\(o(\beta/M_\beta^2)\) by
\eqref{eq:s6v32-GCSC-mixed} and
\(\beta/M_\beta^2\to\infty\).
This proves
\eqref{eq:s6v32-generated-asymptotic}.  No matching covariance
lower bound, and hence no upper asymptotic for the trace Hessian,
is asserted.  Finally use
\eqref{eq:s6v31-smin}; the little-oh error is uniformly smaller
than \(\frac12\beta s_0(q)\).
\end{proof}

\begin{definition}[Topological short card, TSC]
\label{def:s6v32-TSC}
If the full finite-torus trace Hessian decomposes on the same
carrier as
\[
 H_{\beta,\Psi,M}^{\rm full}
 =
 H_{\beta,\Psi,M}^{\rm loc}
 +H_{\beta,\Psi,M}^{\rm top},
\]
TSC means
\begin{equation}
 {M_\beta^2\over\beta}
 \sup_{q\ne0}
 \|H_{\beta,\Psi,M_\beta}^{\rm top}(q)\|_{\rm op}
 \longrightarrow0.                                    \label{eq:s6v32-TSC}
\end{equation}
\end{definition}

\begin{corollary}[Local response-matched static margin]
\label{cor:s6v32-local-margin}
Under GCSC and TSC,
\begin{equation}
 \boxed{
 H_{\beta,\Psi,M_\beta}^{\rm loc}(q)
 \succeq{1\over4}\beta s_0(q)I}                        \label{eq:s6v32-local-margin}
\end{equation}
uniformly on every nonzero selected axial block for all
sufficiently large \(\beta\).
\end{corollary}

\begin{proof}
By TSC and \eqref{eq:s6v31-smin}, the topological subtraction is
eventually at most \(\frac14\beta s_0(q)\).
Subtract it from \eqref{eq:s6v32-generated-margin}.
\end{proof}

\section{Canonical-root landing}
\label{sec:s6v32-root}

\begin{definition}[Canonical-root landing card, CRLC]
\label{def:s6v32-CRLC}
For a cofinal sequence of dyadic shells \(M_j\to\infty\), CRLC
means that the selected canonical roots
\(\widehat\beta_{2M_j}\) satisfy
\begin{equation}
 \widehat\beta_{2M_j}\longrightarrow\infty,
 \qquad
 M_j\le\widehat\beta_{2M_j}^{1/20}                     \label{eq:s6v32-CRLC-scale}
\end{equation}
eventually, and that GCSC and TSC hold with
\((\beta,M_\beta)\) replaced by
\((\widehat\beta_{2M_j},M_j)\).
\end{definition}

\begin{corollary}[Static-card closure under the three remaining
source rows]
\label{cor:s6v32-static-closure}
CRLC implies a cofinally uniform normalized local axial static
margin of at least \(1/4\) on the canonical response-root
sequence.
\end{corollary}

\begin{proof}
The estimates of
\Cref{thm:s6v31-chart-tail,thm:s6v32-chart-card} require only
\[
 M^4=o(\beta r^2),\qquad M^2r=o(1),\qquad
 \beta/M^2\to\infty
\]
with \(r=\beta^{-1/5}\).  The stronger corridor
\eqref{eq:s6v32-CRLC-scale} supplies all three.  Repeat
\Cref{thm:s6v32-generated-margin,cor:s6v32-local-margin}.
\end{proof}

\begin{assessment}[Exact residual after V32]
\label{ass:s6v32-residual}
The response-matched static problem is now reduced to three
source-specific acquisitions, with no hidden covariance:
\begin{enumerate}[label=\textup{(\roman*)},leftmargin=3.5em]
\item GCSC: oscillation, joint-jet, mixed-jet, and polynomial
      global bounds for the actual generated action;
\item TSC: a Ward-relative topological Hessian short; and
\item CRLC: placement of the canonical fine root in the explicit
      weak-coupling corridor.
\end{enumerate}
If these pass, the canonical local static card passes by
\Cref{cor:s6v32-static-closure}.  No attached source currently
populates these three rows.  Their absence is data absence, not a
negative Yang--Mills result.
\end{assessment}

\begin{programme}[V33: derive GCSC from a support-local generated
interaction dictionary]
\label{prog:s6v32-V33}
\begin{enumerate}[leftmargin=3em]
\item State a support-local norm on interaction writers which
      implies each row of GCSC by a finite-incidence Schur test.
\item Separate an extensive zero-jet oscillation from the first and
      second retained-field writers; do not use one bound for all
      three.
\item Determine the maximal support radius and multiplicity growth
      allowed on \(M\le\beta^{1/20}\).
\item Apply the criterion to every generated interaction which is
      actually specified in the supplied manuscripts; leave
      unnamed terms unpopulated.
\item If no dictionary exists, return the minimal required data
      schema and move upstream to the response-construction step
      which must produce it.
\end{enumerate}
\end{programme}

\begin{firewall}[Claim boundary after sixth-series V32]
\label{fire:s6v32-boundary}
V32 proves a uniform perturbation theorem for the complete joint
Schur complement; controls the covariance of the complete
generated score once by an unsplit Brascamp--Lieb argument;
transfers V31 concentration through the generated-action
oscillation; restores the compact-fibre complement under a
polynomial writer card; and proves that GCSC, TSC, and CRLC imply a
canonical local axial static margin of \(1/4\).

V32 does not prove GCSC, TSC, or CRLC for the actual response
construction, establish good-gate transmission, build the
reflection-positive continuum measure, or prove the full
Yang--Mills mass gap.
\end{firewall}

\begin{theorem}[Sixth-series V32 result ledger]
\label{thm:s6v32-results}
The claims in \Cref{fire:s6v32-boundary} follow from
\Cref{lem:s6v32-Schur-perturbation,
lem:s6v32-Haar-inverse,
thm:s6v32-chart-card,
lem:s6v32-chart-probability,
lem:s6v32-mixture,
thm:s6v32-generated-margin,
cor:s6v32-local-margin,
cor:s6v32-static-closure}.
\end{theorem}

\begin{proof}
Each ledger item is the conclusion of the cited result.
\end{proof}

\paragraph{Parent-version record.}
V32 was formed from
\texttt{Renewal\_Yang\_Mills\_Mass\_Gap\_Research\_Notes\_V31.tex}.
The V31 parent had \(14\,096\) logical source lines,
\(499\,462\) bytes, and SHA--256
\[
 \texttt{5e719036dce3cb7a8e9691c4aa54eba99177ceb34e3ca8822b03d1964643b2b6}.
\]
The next compulsory companion audit is V35.

\clearpage
\chapter{Sixth-series V33: the optimal polynomial scale corridor
and a generated-interaction compiler}
\label{chap:s6v33}

\section{Optimizing the cofinal chart scale}
\label{sec:s6v33-scales}

V31 used the deliberately conservative choice
\(M\asymp\beta^{1/20}\), \(r=\beta^{-1/5}\).  The proof itself
contains two competing strict exponent inequalities.  Making them
explicit both enlarges the canonical-root corridor and identifies
the point where this particular concentration--Schur method ceases
to decide.

\begin{definition}[Polynomial shell and chart exponents]
\label{def:s6v33-exponents}
Let \(\gamma,\alpha>0\), and define
\begin{equation}
 M_{\beta,\gamma}
 :=
 2^{\max\{2,\lfloor\gamma\log_2\beta\rfloor\}},
 \qquad
 r_{\beta,\alpha}:=\beta^{-\alpha}.                    \label{eq:s6v33-exponents}
\end{equation}
\end{definition}

\begin{lemma}[The two strict scale inequalities]
\label{lem:s6v33-two-inequalities}
The chart-tail argument of V31 is stretched-exponentially effective
when
\begin{equation}
 1-2\alpha>4\gamma,                                    \label{eq:s6v33-tail-inequality}
\end{equation}
and its Schur error is negligible relative to the smallest axial
reference when
\begin{equation}
 \alpha>2\gamma.                                       \label{eq:s6v33-Schur-inequality}
\end{equation}
\end{lemma}

\begin{proof}
\Cref{thm:s6v31-chart-tail} has negative exponent of order
\(\beta r^2=\beta^{1-2\alpha}\) and entropy exponent of order
\(M^4=\beta^{4\gamma}\), proving
\eqref{eq:s6v33-tail-inequality}.  The chart debit in
\eqref{eq:s6v31-truncated-card} is \(C\beta r\), while the smallest
reference is \(c\beta/M^2\).  Their ratio is
\(CM^2r=O(\beta^{2\gamma-\alpha})\), proving
\eqref{eq:s6v33-Schur-inequality}.
\end{proof}

\begin{theorem}[Exact exponent feasibility interval]
\label{thm:s6v33-feasibility}
There exists an \(\alpha\) satisfying both
\eqref{eq:s6v33-tail-inequality} and
\eqref{eq:s6v33-Schur-inequality} if and only if
\begin{equation}
 \boxed{0<\gamma<{1\over8}.}                           \label{eq:s6v33-gamma-range}
\end{equation}
For such a \(\gamma\), every
\begin{equation}
 \boxed{
 2\gamma<\alpha<{1-4\gamma\over2}}                     \label{eq:s6v33-alpha-range}
\end{equation}
is admissible.
\end{theorem}

\begin{proof}
The two strict inequalities require
\[
 2\gamma<\alpha<{1-4\gamma\over2}.
\]
This open interval is nonempty precisely when
\[
 2\gamma<{1\over2}-2\gamma,
 \quad\hbox{that is}\quad
 \gamma<{1\over8}.
\]
Positivity of \(\gamma\) is the cofinal condition \(M\to\infty\).
\end{proof}

\begin{theorem}[General cofinal pure-Wilson margin]
\label{thm:s6v33-general-Wilson}
Fix \(0<\gamma<1/8\) and choose \(\alpha\) in
\eqref{eq:s6v33-alpha-range}.  Then for all sufficiently large
\(\beta\), every nonzero allowed axial block at
\(M=M_{\beta,\gamma}\) obeys
\begin{equation}
 \boxed{
 H_{\beta,M}(q)\succeq{1\over2}\beta s_0(q)I.}         \label{eq:s6v33-general-Wilson}
\end{equation}
\end{theorem}

\begin{proof}
Use \(r=r_{\beta,\alpha}\) in
\Cref{thm:s6v31-chart-tail}.  The chart-complement probability is
\[
 \epsilon_{\beta,M,r}
 \le
 \exp\{-c\beta^{1-2\alpha}+C\beta^{4\gamma}\},
\]
which is stretched-exponentially small by
\eqref{eq:s6v33-tail-inequality}.  The normalized chart debit is
\(O(M^2r)\), which vanishes by
\eqref{eq:s6v33-Schur-inequality}.  The Haar debit is
\(O(M^2/\beta)\), and the polynomial writer envelope multiplied by
the chart tail also vanishes.  The proof of
\Cref{thm:s6v31-cofinal-margin} now applies verbatim.
\end{proof}

\begin{corollary}[A convenient enlarged schedule]
\label{cor:s6v33-convenient}
The choice
\begin{equation}
 \gamma={1\over10},
 \qquad
 \alpha={1\over4}                                      \label{eq:s6v33-convenient}
\end{equation}
is admissible.  It gives
\[
 M\asymp\beta^{1/10},\qquad
 \epsilon_{\beta,M,r}
 \le e^{-c\beta^{1/2}+C\beta^{2/5}},
 \qquad
 M^2r=O(\beta^{-1/20}).
\]
\end{corollary}

\begin{proof}
\(1/5<1/4<3/10\), which is
\eqref{eq:s6v33-alpha-range} at \(\gamma=1/10\).
\end{proof}

\begin{warning}[Meaning of the \(1/8\) endpoint]
\label{warn:s6v33-endpoint}
The endpoint \(\gamma=1/8\) is a limitation of the present
product-Haar entropy bound and uniform Schur perturbation, not a
physical obstruction.  At the endpoint there is no power
\(\alpha\) which makes both strict estimates close.  A sharper
concentration theorem or a cancellation-sensitive Schur estimate
could enlarge the corridor; V33 does not infer failure of the
static card there.
\end{warning}

\section{A support-local interaction dictionary}
\label{sec:s6v33-dictionary}

Let \(X_M\) be the reduced-link coordinate set and \(W_M\) the
retained coarse-link coordinate set.  Suppose an actual generated
action is presented as
\begin{equation}
 \Psi_{\beta,M}(x,w)
 =
 \sum_{A\in{\cal I}_{\beta,M}}\psi_A(x_{A_X},w_{A_W}), \label{eq:s6v33-dictionary}
\end{equation}
where \(A_X\subset X_M\), \(A_W\subset W_M\), and every writer is
twice continuously differentiable on the declared chart.

For a block matrix \(T=(T_{ij})\), write
\[
 \|T\|_{\rm row}:=\sup_i\sum_j\|T_{ij}\|_{\rm op}.
\]
For rectangular matrices both the row norm and the corresponding
column norm will be recorded.

\begin{definition}[Generated interaction incidence stocks]
\label{def:s6v33-stocks}
On a radius-\(r\) chart define
\begin{align}
 {\cal O}_0
 &:=
 \sum_{A\in{\cal I}_{\beta,M}}
 \operatorname{osc}_{x}\psi_A(x,0),                    \label{eq:s6v33-O0}\\
 {\cal P}_2
 &:=
 \sup_{i\in X_M}
 \sum_A\sum_{j\in A_X}
 \sup_{\|x\|_\infty<r}
 \|\partial_{x_i}\partial_{x_j}\psi_A\|_{\rm op},
                                                               \label{eq:s6v33-P2}\\
 {\cal Q}_{XW}
 &:=
 \sup_{i\in X_M}
 \sum_A\sum_{\lambda\in A_W}
 \sup_{\|x\|_\infty<r}
 \|\partial_{x_i}\partial_{w_\lambda}\psi_A\|_{\rm op},
                                                               \label{eq:s6v33-QXW}\\
 {\cal Q}_{WX}
 &:=
 \sup_{\lambda\in W_M}
 \sum_A\sum_{i\in A_X}
 \sup_{\|x\|_\infty<r}
 \|\partial_{w_\lambda}\partial_{x_i}\psi_A\|_{\rm op},
                                                               \label{eq:s6v33-QWX}\\
 {\cal R}_2
 &:=
 \sup_{\lambda\in W_M}
 \sum_A\sum_{\mu\in A_W}
 \sup_{\|x\|_\infty<r}
 \|\partial_{w_\lambda}\partial_{w_\mu}\psi_A\|_{\rm op}.
                                                               \label{eq:s6v33-R2}
\end{align}
Terms whose displayed differentiated coordinate is not in the
support of \(A\) are zero.
\end{definition}

\begin{lemma}[Incidence stocks control the complete joint jets]
\label{lem:s6v33-stock-control}
For the operators in \eqref{eq:s6v32-PQR},
\begin{align}
 \|P_\Psi\|_{\rm op}
 &\le{\cal P}_2,                                       \label{eq:s6v33-P-control}\\
 \|Q_\Psi\|_{\rm op}
 &\le({\cal Q}_{XW}{\cal Q}_{WX})^{1/2},               \label{eq:s6v33-Q-control}\\
 \|R_\Psi\|_{\rm op}
 &\le{\cal R}_2,                                       \label{eq:s6v33-R-control}\\
 \operatorname{osc}\Psi_{\beta,M}(\cdot,0)
 &\le{\cal O}_0.                                       \label{eq:s6v33-osc-control}
\end{align}
\end{lemma}

\begin{proof}
The oscillation bound is subadditivity of oscillation.  The two
square Hessians are self-adjoint, so the Schur row-sum bound gives
\eqref{eq:s6v33-P-control} and
\eqref{eq:s6v33-R-control}.  For the rectangular mixed Hessian,
\[
 \|Q_\Psi\|_2
 \le\sqrt{\|Q_\Psi\|_1\|Q_\Psi\|_\infty}
 \le({\cal Q}_{XW}{\cal Q}_{WX})^{1/2}.
\]
\end{proof}

\begin{definition}[First-writer incidence stocks]
\label{def:s6v33-first-stocks}
Let
\begin{align}
 {\cal J}_1
 &:=
 \sup_{\lambda\in W_M}
 \sum_A
 \sup_{\|x\|_\infty<r}
 \|\partial_{w_\lambda}\psi_A\|,                       \label{eq:s6v33-J1}\\
 {\cal K}_2
 &:=
 \sup_{\lambda\in W_M}
 \sum_A\sum_{\mu\in A_W}
 \sup_{x\in\SU(3)^{X_M}}
 \|\partial_{w_\lambda}\partial_{w_\mu}\psi_A\|.
                                                               \label{eq:s6v33-K2}
\end{align}
For PGWC, the first supremum is also taken on the full compact
fibre rather than only the chart.
\end{definition}

\begin{theorem}[Dictionary-to-GCSC compiler]
\label{thm:s6v33-dictionary-compiler}
Fix an admissible exponent pair
\((\gamma,\alpha)\) from
\Cref{thm:s6v33-feasibility}, put
\[
 M=M_{\beta,\gamma},\qquad r=r_{\beta,\alpha},
\]
and define
\begin{equation}
 {\cal D}_2
 :=
 {\cal P}_2+{\cal R}_2+
 ({\cal Q}_{XW}{\cal Q}_{WX})^{1/2}
 +{{\cal Q}_{XW}{\cal Q}_{WX}\over\beta}.              \label{eq:s6v33-D2}
\end{equation}
The generated action satisfies the scale-adjusted GCSC whenever
\begin{align}
 {\cal O}_0&=o(\beta r^2),                              \label{eq:s6v33-compiler-osc}\\
 {M^2\over\beta}{\cal D}_2&\longrightarrow0,            \label{eq:s6v33-compiler-jets}\\
 {\cal J}_1+{\cal K}_2&\le C\beta^aM^b                 \label{eq:s6v33-compiler-global}
\end{align}
for some fixed \(a,b,C\), with the last bound valid on the full
compact fibre.  Under the analogous topological short and
canonical-root landing rows, the local static margin of
\Cref{cor:s6v32-local-margin} follows.
\end{theorem}

\begin{proof}
\Cref{lem:s6v33-stock-control} inserts
\({\cal D}_2\) into the debit
\eqref{eq:s6v32-jet-debit}.  Since
\(({\cal Q}_{XW}{\cal Q}_{WX})^{1/2}\le{\cal D}_2\),
\eqref{eq:s6v33-compiler-jets} also makes the mixed jet
little-oh of \(\beta\).  Equation
\eqref{eq:s6v33-compiler-osc} is the chart-probability row.
Finite incidence and the full-fibre first/second writer bounds in
\eqref{eq:s6v33-compiler-global} give PGWC after contraction with a
unitary Fourier polarization.  The proof of
\Cref{thm:s6v32-generated-margin} then applies with the general
tail and margin inequalities from
\Cref{lem:s6v33-two-inequalities}.
\end{proof}

\section{A finite-range coefficient corollary}
\label{sec:s6v33-finite-range}

\begin{corollary}[Uniformly local generated terms]
\label{cor:s6v33-local-terms}
Suppose:
\begin{enumerate}[label=\textup{(\roman*)},leftmargin=3.5em]
\item at most \(N_\beta\) generated interactions contain any one
      coordinate;
\item each support contains at most \(S_\beta\) reduced and
      retained coordinates;
\item every second partial derivative is bounded by \(b_{2,\beta}\);
\item every interaction oscillation is bounded by \(b_{0,\beta}\);
      and
\item the total number of interactions is at most
      \(C M^4N_\beta\).
\end{enumerate}
Then the compiler rows follow from
\begin{align}
 M^4N_\beta b_{0,\beta}
 &=o(\beta r^2),                                      \label{eq:s6v33-local-osc}\\
 N_\beta S_\beta b_{2,\beta}
 &=o(\beta/M^2),                                      \label{eq:s6v33-local-jets}
\end{align}
together with polynomial full-fibre first-writer bounds.
\end{corollary}

\begin{proof}
The hypotheses give
\[
 {\cal O}_0\le CM^4N_\beta b_{0,\beta}
\]
and bound each of
\({\cal P}_2,{\cal Q}_{XW},{\cal Q}_{WX},{\cal R}_2\)
by \(CN_\beta S_\beta b_{2,\beta}\).
The quadratic mixed term divided by \(\beta\) is smaller under
\eqref{eq:s6v33-local-jets}.  Apply
\Cref{thm:s6v33-dictionary-compiler}.
\end{proof}

\begin{corollary}[A bounded plaquette-local retuning passes]
\label{cor:s6v33-plaquette-retuning}
If the generated action is a sum of fixed-radius plaquette-local
writers with uniformly bounded coefficients and uniformly bounded
first two derivatives, then it satisfies the dictionary compiler
on every admissible exponent pair
\((\gamma,\alpha)\).  More generally, a common coefficient
\(b_{\beta,M}\) passes the second-jet row whenever
\begin{equation}
 |b_{\beta,M}|=o(\beta/M^2),                            \label{eq:s6v33-b-condition}
\end{equation}
and passes the oscillation row whenever
\begin{equation}
 M^4|b_{\beta,M}|=o(\beta r^2).                        \label{eq:s6v33-b-osc}
\end{equation}
\end{corollary}

\begin{proof}
Fixed-radius support makes \(N_\beta,S_\beta=O(1)\).
The stated coefficient bounds insert directly into
\eqref{eq:s6v33-local-osc} and
\eqref{eq:s6v33-local-jets}.  For bounded coefficients,
\eqref{eq:s6v33-tail-inequality} gives
\(M^4=o(\beta r^2)\), and
\(\beta/M^2\to\infty\).
\end{proof}

\begin{assessment}[Minimal missing data schema]
\label{ass:s6v33-schema}
The generated-action bottleneck no longer requires an unspecified
global phrase.  A submitted interaction dictionary need only
populate:
\[
 ({\cal O}_0,\ {\cal P}_2,\ {\cal Q}_{XW},\
  {\cal Q}_{WX},\ {\cal R}_2,\ {\cal J}_1,\ {\cal K}_2)
\]
with support multiplicities and the canonical-root growth law.
Bounded finite-range terms pass automatically.  A term can fail
only by extensive zero-jet oscillation faster than the chart energy,
joint-jet incidence of order \(\beta/M^2\) or larger, nonpolynomial
global writers, a nonshort topological Hessian, or a root outside
the exponent corridor.  The supplied sources name
\(\Psi_{\beta,M}\) but do not provide these seven stocks.
\end{assessment}

\begin{programme}[V34: move upstream to the response construction]
\label{prog:s6v33-V34}
\begin{enumerate}[leftmargin=3em]
\item Locate the exact step which constructs or selects the
      response-matched generated action before it is denoted by
      \(\Psi\).
\item Require that step to emit the seven-stock dictionary of
      \Cref{ass:s6v33-schema}, rather than only a scalar response
      coefficient.
\item Test the marginal plaquette retuning separately using
      \eqref{eq:s6v33-b-condition}--\eqref{eq:s6v33-b-osc}.
\item Audit whether the topological packet has exponentially small
      axial Hessian on \(M\asymp\beta^\gamma\).
\item If the response construction has no such writer output,
      record the exact missing interface and return to the
      pure-Wilson local-vacuum/good-gate route without treating
      GCSC as proved.
\end{enumerate}
\end{programme}

\begin{firewall}[Claim boundary after sixth-series V33]
\label{fire:s6v33-boundary}
V33 proves that the concentration--Schur method supports every
polynomial cofinal shell exponent \(\gamma<1/8\), with the exact
chart interval
\(2\gamma<\alpha<(1-4\gamma)/2\); proves the corresponding
pure-Wilson cofinal margin; defines seven support-local generated
interaction stocks; proves their operator-norm control; and gives
a finite-incidence compiler which implies the generated-action
short card.

V33 does not populate those stocks for the symbolic generated
action, prove the topological short or root-landing cards,
establish good-gate transmission, construct the continuum measure,
or prove the full Yang--Mills mass gap.
\end{firewall}

\begin{theorem}[Sixth-series V33 result ledger]
\label{thm:s6v33-results}
The claims in \Cref{fire:s6v33-boundary} follow from
\Cref{lem:s6v33-two-inequalities,
thm:s6v33-feasibility,
thm:s6v33-general-Wilson,
cor:s6v33-convenient,
lem:s6v33-stock-control,
thm:s6v33-dictionary-compiler,
cor:s6v33-local-terms,
cor:s6v33-plaquette-retuning}.
\end{theorem}

\begin{proof}
Each ledger item is the conclusion of the cited result.
\end{proof}

\paragraph{Parent-version record.}
V33 was formed from
\texttt{Renewal\_Yang\_Mills\_Mass\_Gap\_Research\_Notes\_V32.tex}.
The V32 parent had \(14\,580\) logical source lines,
\(514\,627\) bytes, and SHA--256
\[
 \texttt{4db29068cab48c617ffce4a5d5a339c80ac487669e750ff5169923c9c94b47a4}.
\]
The next compulsory companion audit is V35.

\clearpage
\chapter{Sixth-series V34: response regularity implies a bounded
marginal compiler}
\label{chap:s6v34}

\section{The actual upstream response interface}
\label{sec:s6v34-interface}

The response-chart layer of the attached V35 programme defines four
positive Wilson-loop numerators \(N_\alpha\), the Creutz response
\(\chi\), and, at a selected root \(\beta_*\), the two derivatives
\begin{equation}
 d:=\partial_\beta\chi(\beta_*,0),
 \qquad
 n:=\partial_s\chi(\beta_*,0).                          \label{eq:s6v34-dn}
\end{equation}
On a regular root, the exact marginal coefficient is
\begin{equation}
 \boxed{b^{\rm marg}=-{n\over d}.}                      \label{eq:s6v34-b}
\end{equation}
The first-order response-matched generated action is therefore
\begin{equation}
 \boxed{
 \Theta_{\beta,M}
 :=
 \Psi_{\beta,M}^{\rm gen}+b^{\rm marg}_{\beta,M}P_M,}  \label{eq:s6v34-Theta}
\end{equation}
where \(P_M\) is the declared plaquette marginal writer.  The
coefficient is selected before coarse-field differentiation and is
held fixed inside that directional Hessian.

\begin{definition}[Uniform response-regularity card, URRC]
\label{def:s6v34-URRC}
A selected cofinal root family satisfies URRC if every root is
simple and there are constants \(d_*,n^*>0\) such that
\begin{equation}
 \boxed{
 |d_{\beta,M}|\ge d_*,
 \qquad
 |n_{\beta,M}|\le n^*.}                                \label{eq:s6v34-URRC}
\end{equation}
\end{definition}

\begin{lemma}[URRC bounds the selected marginal coefficient]
\label{lem:s6v34-b-bound}
Under URRC,
\begin{equation}
 \boxed{|b^{\rm marg}_{\beta,M}|\le B_*:=n^*/d_*.}     \label{eq:s6v34-b-bound}
\end{equation}
\end{lemma}

\begin{proof}
Insert \eqref{eq:s6v34-URRC} into
\eqref{eq:s6v34-b}.
\end{proof}

\section{Stock addition for the matched action}
\label{sec:s6v34-stocks}

Assume the plaquette writer has a fixed-radius decomposition
\[
 P_M=\sum_{p}P_p
\]
whose one-plaquette oscillations and first two link/coarse
derivatives are bounded by a regulator-independent constant.  This
is true for the ordinary Wilson plaquette writer.

\begin{lemma}[Plaquette marginal stock bounds]
\label{lem:s6v34-P-stocks}
There is a constant \(C_P\), independent of \(M,\beta,b\), such
that the seven stocks of
\Cref{def:s6v33-stocks,def:s6v33-first-stocks} for \(bP_M\)
satisfy
\begin{align}
 {\cal O}_0[bP_M]&\le C_P|b|M^4,                       \label{eq:s6v34-P-O}\\
 {\cal P}_2[bP_M]+{\cal Q}_{XW}[bP_M]
 +{\cal Q}_{WX}[bP_M]+{\cal R}_2[bP_M]
 &\le C_P|b|,                                          \label{eq:s6v34-P-D}\\
 {\cal J}_1[bP_M]+{\cal K}_2[bP_M]
 &\le C_P|b|                                           \label{eq:s6v34-P-J}
\end{align}
after unitary coarse normalization; the global first writer has
the polynomial envelope \(C_P|b|M^2\).
\end{lemma}

\begin{proof}
Each plaquette contains four links and has a fixed derivative
bound on the compact group.  A link belongs to six plaquettes, so
every differentiated coordinate row meets only a fixed number of
terms.  This proves \eqref{eq:s6v34-P-D} and
\eqref{eq:s6v34-P-J}.  There are \(O(M^4)\) plaquettes, which gives
\eqref{eq:s6v34-P-O}.  Contracting the first coarse writer with a
unit normalized plane wave uses its \(\ell^1\)-norm \(M^2\).
\end{proof}

\begin{lemma}[Matched-action stock addition]
\label{lem:s6v34-stock-addition}
Let superscript \({\rm gen}\) denote the stocks of
\(\Psi^{\rm gen}\), and superscript \({\rm mat}\) those of
\(\Theta\).  Then
\begin{align}
 {\cal O}_0^{\rm mat}
 &\le{\cal O}_0^{\rm gen}+C_P|b^{\rm marg}|M^4,
                                                               \label{eq:s6v34-O-add}\\
 {\cal P}_2^{\rm mat}+{\cal R}_2^{\rm mat}
 &\le{\cal P}_2^{\rm gen}+{\cal R}_2^{\rm gen}
      +C_P|b^{\rm marg}|,                              \label{eq:s6v34-PR-add}\\
 {\cal Q}_{XW}^{\rm mat}
 &\le{\cal Q}_{XW}^{\rm gen}+C_P|b^{\rm marg}|,        \label{eq:s6v34-Q1-add}\\
 {\cal Q}_{WX}^{\rm mat}
 &\le{\cal Q}_{WX}^{\rm gen}+C_P|b^{\rm marg}|.        \label{eq:s6v34-Q2-add}
\end{align}
The first-writer stocks obey the analogous triangle inequalities.
\end{lemma}

\begin{proof}
Differentiate \eqref{eq:s6v34-Theta} with \(b^{\rm marg}\) held
fixed on the field tangent.  Every stock is a sum of absolute row
contributions, so the triangle inequality and
\Cref{lem:s6v34-P-stocks} apply.
\end{proof}

\section{Response numerators to the cofinal short card}
\label{sec:s6v34-compiler}

Fix an admissible scale pair
\[
 0<\gamma<1/8,
 \qquad
 2\gamma<\alpha<(1-4\gamma)/2,
\]
and put \(M=M_{\beta,\gamma}\), \(r=\beta^{-\alpha}\).

\begin{theorem}[General ratio conditions for marginal matching]
\label{thm:s6v34-ratio-compiler}
Suppose the base generated action
\(\Psi^{\rm gen}\) passes the dictionary compiler of
\Cref{thm:s6v33-dictionary-compiler}.  Then the matched action
\(\Theta\) also passes whenever
\begin{align}
 M^4\left|{n\over d}\right|
 &=o(\beta r^2),                                       \label{eq:s6v34-ratio-osc}\\
 \left|{n\over d}\right|
 &=o(\beta/M^2),                                       \label{eq:s6v34-ratio-jet}
\end{align}
and \(n/d\) has at most polynomial growth for the global writer
card.
\end{theorem}

\begin{proof}
By \eqref{eq:s6v34-b}, the added coefficient has magnitude
\(|n/d|\).  Equations
\eqref{eq:s6v34-ratio-osc} and
\eqref{eq:s6v34-ratio-jet} are exactly the two sufficient
plaquette-retuning conditions
\eqref{eq:s6v33-b-osc} and
\eqref{eq:s6v33-b-condition}.  Stock addition in
\Cref{lem:s6v34-stock-addition} preserves each little-oh row, and
polynomial growth preserves PGWC.
\end{proof}

\begin{corollary}[Uniform response regularity closes the marginal
rows]
\label{cor:s6v34-URRC-compiler}
If the base generated action passes the V33 dictionary compiler
and URRC holds, then the complete matched action
\(\Theta=\Psi^{\rm gen}+b^{\rm marg}P\) satisfies GCSC on every
admissible exponent pair.
\end{corollary}

\begin{proof}
By \Cref{lem:s6v34-b-bound}, \(b^{\rm marg}=O(1)\).
The scale inequalities give
\[
 M^4=o(\beta r^2),
 \qquad
 1=o(\beta/M^2).
\]
Thus
\eqref{eq:s6v34-ratio-osc} and
\eqref{eq:s6v34-ratio-jet} hold.
\end{proof}

\begin{corollary}[Pure generated head with bounded marginal
retuning]
\label{cor:s6v34-pure-marginal}
If \(\Psi^{\rm gen}=0\) and URRC holds, the response-matched action
\(b^{\rm marg}P\) satisfies GCSC.  Under the topological short and
root-landing cards, the local static margin is at least \(1/4\) for
all sufficiently large selected roots.
\end{corollary}

\begin{proof}
The zero generated head has zero stocks.
\Cref{cor:s6v34-URRC-compiler} gives GCSC, and
\Cref{cor:s6v32-static-closure} gives the static conclusion.
\end{proof}

\section{Root multiplicity is a domain condition}
\label{sec:s6v34-ramification}

\begin{theorem}[Regular/ramified response fork]
\label{thm:s6v34-response-fork}
For every selected response root, exactly one of the following
applies.
\begin{enumerate}[label=\textup{(\roman*)},leftmargin=3.5em]
\item \(d\ne0\): the ordinary matched coefficient is uniquely
      \(b=-n/d\), and
      \Cref{thm:s6v34-ratio-compiler} applies.
\item \(d=0\), \(n\ne0\): no finite first-order plaquette
      coefficient preserves the response level.
\item \(d=n=0\): first order does not select \(b\), and the least
      nonzero mixed response jet is required.
\end{enumerate}
The latter two branches do not define the action
\eqref{eq:s6v34-Theta} and therefore cannot be inserted into the
static theorem as if \(b\) were finite.
\end{theorem}

\begin{proof}
The first-order response constraint is
\[
 n+bd=0.
\]
Elementary linear algebra gives the three disjoint alternatives.
The existence of an analytic ordinary root graph in the first case
is the implicit-function theorem.  In the other cases that theorem
does not provide a unique first derivative.
\end{proof}

\begin{corollary}[Near-ramification quantitative requirement]
\label{cor:s6v34-near-ramification}
On a regular sequence with \(d\to0\), the marginal compiler remains
valid only if the actual ratio \(n/d\) satisfies
\eqref{eq:s6v34-ratio-osc}--\eqref{eq:s6v34-ratio-jet}.
A bound on \(n\) alone is insufficient.
\end{corollary}

\begin{proof}
This is \Cref{thm:s6v34-ratio-compiler}.  If \(n\) stays separated
from zero while \(d\to0\), the ratio diverges and may violate either
row.  If \(n\to0\) as well, its relative order against \(d\) is
still needed.
\end{proof}

\section{The complete response-to-static handoff}
\label{sec:s6v34-handoff}

\begin{definition}[Response-root landing and regularity card,
RRLRC]
\label{def:s6v34-RRLRC}
RRLRC consists of:
\begin{enumerate}[label=\textup{(\roman*)},leftmargin=3.5em]
\item selected roots \((\widehat\beta_{2M},M)\) satisfying an
      exponent corridor with some fixed \(\gamma<1/8\);
\item a simple-root certificate and URRC;
\item the base generated-action dictionary rows of V33;
\item the same-carrier topological short TSC.
\end{enumerate}
\end{definition}

\begin{theorem}[RRLRC closes the canonical local static card]
\label{thm:s6v34-RRLRC}
RRLRC implies, for every sufficiently large selected root and every
nonzero retained axial block,
\begin{equation}
 \boxed{
 H_{\widehat\beta,\Theta,M}^{\rm loc}(q)
 \succeq{1\over4}\widehat\beta\,s_0(q)I.}              \label{eq:s6v34-RRLRC}
\end{equation}
\end{theorem}

\begin{proof}
URRC and the base dictionary give matched GCSC by
\Cref{cor:s6v34-URRC-compiler}.  The exponent corridor is the
generalized canonical-root landing condition of V33.  TSC removes
the topological packet.  Apply
\Cref{thm:s6v32-generated-margin,cor:s6v32-local-margin}.
\end{proof}

\begin{assessment}[Exact upstream frontier]
\label{ass:s6v34-frontier}
The scalar response coefficient is no longer an unstructured
addition to the generated action.  On a regular root it is
\(-n/d\), and a uniform slope floor plus bounded numerator makes
the plaquette marginal automatically harmless for the cofinal
static card.  The remaining upstream data are now:
\begin{enumerate}[label=\textup{(\alph*)},leftmargin=3.5em]
\item a cofinal lower bound for the response slope \(d\);
\item a cofinal upper bound for the generated response numerator
      \(n\);
\item the seven-stock dictionary for the base
      \(\Psi^{\rm gen}\);
\item the root growth exponent; and
\item the topological Hessian short.
\end{enumerate}
The attached response programme gives finite evaluators and
interval alternatives for these quantities, but no cofinal
analytic estimates which populate them.
\end{assessment}

\begin{programme}[V35: compulsory audit and response-root
asymptotics]
\label{prog:s6v34-V35}
\begin{enumerate}[leftmargin=3em]
\item Perform the required audit of the then-current SM and NS
      notes.
\item Use the exact four-loop response numerator formulas to test
      whether the pure-Wilson weak-coupling root can satisfy a
      power corridor \(M=O(\beta^\gamma)\), \(\gamma<1/8\).
\item Seek a reflection-positive or character-expansion estimate
      for the response slope \(d\) and numerator \(n\), keeping
      their ratio signed.
\item Do not infer URRC from root existence or monotonicity unless
      a uniform derivative floor is proved.
\item Keep TSC independent from the response-slope calculation.
\end{enumerate}
\end{programme}

\begin{firewall}[Claim boundary after sixth-series V34]
\label{fire:s6v34-boundary}
V34 identifies the exact upstream matched action
\(\Psi^{\rm gen}-(n/d)P\); proves the seven-stock contribution of a
plaquette marginal; gives explicit sufficient scale
conditions on \(n/d\) for the absolute-stock marginal compiler; proves that
uniform response regularity makes bounded marginal retuning
automatic; records the regular/ramified domain fork; and proves
that RRLRC closes the canonical local static card.

V34 does not prove URRC, the base generated-action dictionary
bounds, root landing, TSC, good-gate transmission, continuum
existence, or the full Yang--Mills mass gap.
\end{firewall}

\begin{theorem}[Sixth-series V34 result ledger]
\label{thm:s6v34-results}
The claims in \Cref{fire:s6v34-boundary} follow from
\Cref{lem:s6v34-b-bound,
lem:s6v34-P-stocks,
lem:s6v34-stock-addition,
thm:s6v34-ratio-compiler,
cor:s6v34-URRC-compiler,
thm:s6v34-response-fork,
thm:s6v34-RRLRC}.
\end{theorem}

\begin{proof}
Each ledger item is the conclusion of the cited result.
\end{proof}

\paragraph{Parent-version record.}
V34 was formed from
\texttt{Renewal\_Yang\_Mills\_Mass\_Gap\_Research\_Notes\_V33.tex}.
The V33 parent had \(15\,029\) logical source lines,
\(529\,887\) bytes, and SHA--256
\[
 \texttt{fdd4ffcaf4cdb9fe321e0f854b7872a6e6449bdc8c9afe55a35d5002a49b2d94}.
\]
The next version must perform the compulsory V35 companion audit.

\clearpage
\chapter{Sixth-series V35: companion audit, response-root scale,
and physical-clock no-escape}
\label{chap:s6v35}

\section{Compulsory audit of the current companion programmes}
\label{sec:s6v35-audit}

The synchronized read-only audit was made against the newest complete
companion files present when this continuation was written:
\begin{align}
 &\texttt{Renewal\_SM\_Einstein\_Continuum\_Research\_Notes\_V59.tex},
 \notag\\
 &\qquad 954\,504\ \hbox{bytes},\quad 23\,743\ \hbox{logical lines},
 \notag\\
 &\qquad
 \texttt{c0892b99f05245f80ae108987c1bf69676278c8e4dce67d9519986a8bf535fa4},
 \label{eq:s6v35-SM-checkpoint}\\
 &\texttt{Renewal\_NS\_Stability\_Research\_Notes\_V21.tex},
 \notag\\
 &\qquad 222\,679\ \hbox{bytes},\quad 4\,316\ \hbox{logical lines},
 \notag\\
 &\qquad
 \texttt{9159a5cf188d664d233382a3b1ba0b64b74bf6a93eb91c756943d1cefd87fe08}.
 \label{eq:s6v35-NS-checkpoint}
\end{align}
The files are parallel working notes, not authorities for the present
proof.

SM V59 proves that the real Euclidean Einstein--Weyl functional has a
uniformly elliptic high-frequency conformal escape, gives a deterministic
curvature--matter lower bound after positive \(R^2+C^2\) completion, and
then separates that lower bound from the entropy, exponential-moment,
compactness, and reflection-positivity rows needed for its measure limit.
The directly reusable methodological point is exact: deterministic
coercivity does not itself pay configuration-space entropy or prove
no-escape for normalized measures.

NS V21 proves that its loop constraints become identities on part of its
parameter region and are unavailable on the remainder, and tabulates
derivative bounds which become too loose away from the endpoint where they
were designed.  Its sign and conditional-window statements are specific to
Navier--Stokes and are not imported.  The relevant audit warning is only
logical: a magnitude or an inequality which is automatic on the selected
region cannot decide a signed physical gate.

\begin{theorem}[V35 companion-audit decision]
\label{thm:s6v35-audit-decision}
No SM or NS theorem is imported as a Yang--Mills estimate.  The audit changes
the order of attack in two ways:
\begin{enumerate}[label=\textup{(\roman*)},leftmargin=3.5em]
\item the \(M^4\) entropy term in the V31 global-chart estimate is treated as
      a genuine scale debit, rather than being hidden behind finite-card
      coercivity; and
\item the signed response numerator is retained when the response-slope
      denominator is bounded, rather than replacing the ratio by unsigned
      size data.
\end{enumerate}
\end{theorem}

\begin{proof}
The first point is the distinction between an action lower bound and a
normalized probability estimate.  The second is forced algebraically by
\(b^{\rm marg}=-n/d\): the sign of the matched marginal is the sign of a
ratio, not a consequence of separate magnitude bounds.  Neither statement
uses an equation from the companion theories.
\end{proof}

\section{What the V31--V33 scale corridor demands of the response root}
\label{sec:s6v35-root-scale}

Write \(\widehat\beta_{2M}\) for the selected weak-facing root of the exact
four-loop function
\[
 D_{2M,u_*}(\beta)
 =N_{11}(\beta)N_{00}(\beta)
  -e^{-u_*}N_{10}(\beta)N_{01}(\beta).
\]
The attached response programme proves finite isolation of each fixed-\(M\)
root but supplies no cofinal growth estimate.  The V33 global-chart argument
requires \(M=O(\widehat\beta_{2M}^{\,\gamma})\) for some
\(\gamma<1/8\).  The next theorem records the exact content of that
requirement.

\begin{theorem}[Super-octic root-growth equivalence]
\label{thm:s6v35-super-octic}
Let \(M\) tend to infinity through the selected dyadic sequence and suppose
\(\widehat\beta_{2M}>0\).  The following are equivalent:
\begin{enumerate}[label=\textup{(\roman*)},leftmargin=3.5em]
\item there are \(C<\infty\) and \(0<\gamma<1/8\) such that
      \[
       M\le C\widehat\beta_{2M}^{\,\gamma}
      \]
      eventually;
\item there are \(c>0\) and \(p>8\) such that
      \[
       \widehat\beta_{2M}\ge cM^p
      \]
      eventually.
\end{enumerate}
Consequently any cofinal root law satisfying
\(\widehat\beta_{2M}=O(M^8)\), and in particular any polylogarithmic root
law, fails the V33 global-chart corridor.
\end{theorem}

\begin{proof}
If \textup{(i)} holds, raise its inequality to the power
\(p=1/\gamma>8\):
\[
 \widehat\beta_{2M}\ge C^{-p}M^p.
\]
This is \textup{(ii)}.  Conversely, if \textup{(ii)} holds, take
\(\gamma=1/p<1/8\).  Then
\[
 M\le c^{-1/p}\widehat\beta_{2M}^{\,1/p},
\]
which is \textup{(i)}.  If
\(\widehat\beta_{2M}\le C_8M^8\) along an unbounded subsequence, combining it
with \textup{(ii)} would give
\(cM^{p-8}\le C_8\), impossible for \(p>8\).  Every polylogarithm is
\(o(M^8)\).
\end{proof}

\begin{corollary}[A fixed-card root certificate does not land the cofinal
static theorem]
\label{cor:s6v35-root-certificate}
Argument-principle isolation, simplicity, and arbitrary interval refinement
of every fixed-\(M\) zero of \(D_{2M,u_*}\) do not imply the root landing
required by V31--V34.  To use that global-chart theorem one must additionally
prove a super-octic lower growth law for the selected roots, or replace the
global all-link chart estimate.
\end{corollary}

\begin{proof}
The fixed-card assertions concern one analytic function at one \(M\) and
contain no inequality relating different \(M\).  The missing relation is
exactly either condition of
\Cref{thm:s6v35-super-octic}.  The stated alternative changes the proof
which generated that condition and therefore does not infer it from root
isolation.
\end{proof}

\begin{assessment}[Origin of the exponent eight]
\label{ass:s6v35-eight}
The threshold is not a four-loop response theorem.  It comes from using the
single event that all \(45M^4\) reduced links lie in one principal chart.
The V31 tail pays an entropy \(CM^4\), while the local Schur error must obey
\(M^2r\to0\).  With \(M\asymp\beta^\gamma\) and
\(r=\beta^{-\alpha}\), these rows are
\[
 1-2\alpha>4\gamma,
 \qquad
 \alpha>2\gamma.
\]
They have a common solution precisely when \(\gamma<1/8\).  Thus a
local-density or polymer restoration which does not demand a simultaneous
chart for all links is the upstream way to remove the super-octic debit.
\end{assessment}

\section{A signed cofinal compiler for the four-loop response ratio}
\label{sec:s6v35-response}

At a selected regular root write
\[
 W_\alpha={N_\alpha\over Z}>0,
 \qquad
 d=\partial_\beta\chi,
 \qquad
 n=\partial_s\chi,
 \qquad
 b^{\rm marg}=-{n\over d},
\]
with the four Creutz signs
\(\varepsilon_{11}=\varepsilon_{00}=-1\) and
\(\varepsilon_{10}=\varepsilon_{01}=1\).  For a bounded real insertion
\(\Psi\), put
\[
 \operatorname{osc}(\Psi)=\sup\Psi-\inf\Psi.
\]

\begin{lemma}[Partition-free oscillation bound for the generated numerator]
\label{lem:s6v35-n-bound}
Assume the normalized real loop writers obey
\(\|\mathcal O_\alpha\|_\infty\le1\), and that at the selected root
\[
 \min_\alpha W_\alpha\ge\kappa>0.
\]
Then the exact four-loop response numerator satisfies
\begin{equation}
 \boxed{
 |n|\le {2\over\kappa}\operatorname{osc}(\Psi).}
 \label{eq:s6v35-n-bound}
\end{equation}
This estimate is invariant under adding a constant to \(\Psi\).
\end{lemma}

\begin{proof}
The response differential is
\[
 n=-\sum_\alpha\varepsilon_\alpha
 {N_{\alpha;\Psi}\over N_\alpha}.
\]
Because \(\sum_\alpha\varepsilon_\alpha=0\), replace \(\Psi\) by
\(\Psi-c\), where
\(c=(\sup\Psi+\inf\Psi)/2\).  Then
\(\|\Psi-c\|_\infty\le\operatorname{osc}(\Psi)/2\), and
\[
 |N_{\alpha;\Psi-c}|
 \le Z\|\Psi-c\|_\infty,
 \qquad
 N_\alpha=ZW_\alpha\ge\kappa Z.
\]
Each of the four summands has absolute value at most
\(\operatorname{osc}(\Psi)/(2\kappa)\).  Summation proves
\eqref{eq:s6v35-n-bound}.  The cancellation of \(c\) proves the last
claim.
\end{proof}

\begin{theorem}[Two-sided response bracket gives a signed slope floor]
\label{thm:s6v35-bracket-slope}
Let \(\chi\) be \(C^2\) on
\([\beta_*-h,\beta_*+h]\), where
\(\chi(\beta_*)=u_*\) and \(h>0\).  Suppose
\begin{align}
 \chi(\beta_*-h)-u_*&\ge\eta_-,&
 u_*-\chi(\beta_*+h)&\ge\eta_+,                     \label{eq:s6v35-bracket}\\
 \sup_{|\beta-\beta_*|\le h}|\chi''(\beta)|&\le K,    \label{eq:s6v35-second}
\end{align}
where \(\eta_\pm>0\).  Put
\begin{equation}
 L:={\max\{\eta_-,\eta_+\}\over h}-{Kh\over2}.
 \label{eq:s6v35-L}
\end{equation}
If \(L>0\), then the selected root is transverse in the down-crossing
orientation and
\begin{equation}
 \boxed{d=\chi'(\beta_*)\le-L<0.}
 \label{eq:s6v35-d-floor}
\end{equation}
Under the hypotheses of \Cref{lem:s6v35-n-bound},
\begin{equation}
 \boxed{
 |b^{\rm marg}|
 \le {2\,\operatorname{osc}(\Psi)\over\kappa L},
 \qquad
 \operatorname{sgn}b^{\rm marg}=\operatorname{sgn}n
 \quad(n\ne0).}
 \label{eq:s6v35-b-bound}
\end{equation}
\end{theorem}

\begin{proof}
Taylor's theorem at \(\beta_*\) gives
\[
 \chi(\beta_*-h)-u_*=-hd+R_-,
 \qquad |R_-|\le Kh^2/2.
\]
Therefore
\(-d\ge\eta_-/h-Kh/2\).  Likewise
\[
 \chi(\beta_*+h)-u_*=hd+R_+,
 \qquad |R_+|\le Kh^2/2,
\]
and the second bracket gives
\(-d\ge\eta_+/h-Kh/2\).  Taking the larger lower bound proves
\eqref{eq:s6v35-d-floor}.  Combine it with
\eqref{eq:s6v35-n-bound} and \(b^{\rm marg}=-n/d\).
Since \(d<0\), the signs of \(b^{\rm marg}\) and \(n\) agree.
\end{proof}

\begin{corollary}[Scale-dependent response transversality is enough]
\label{cor:s6v35-scaled-response}
Along selected roots let \(\kappa_M,h_M,\eta_{\pm,M},K_M\) satisfy
\Cref{thm:s6v35-bracket-slope}, and let
\[
 L_M={\max\{\eta_{-,M},\eta_{+,M}\}\over h_M}
       -{K_Mh_M\over2}>0,
 \qquad
 B_M={2\,\operatorname{osc}(\Psi_M)\over\kappa_ML_M}.
\]
For any admissible V33 scale pair \((M,\beta,r)\), the plaquette marginal
passes the two V34 ratio rows whenever
\begin{equation}
 {M^4B_M\over\beta r^2}\longrightarrow0,
 \qquad
 {M^2B_M\over\beta}\longrightarrow0.
 \label{eq:s6v35-scaled-rows}
\end{equation}
Thus the uniform response-regularity card
\(|d|\ge d_*>0\), \(|n|\le n^*\) is sufficient but not necessary.
\end{corollary}

\begin{proof}
\Cref{thm:s6v35-bracket-slope} gives
\(|n/d|=|b^{\rm marg}|\le B_M\).  The two limits in
\eqref{eq:s6v35-scaled-rows} imply respectively
\[
 M^4|n/d|=o(\beta r^2),
 \qquad
 |n/d|=o(\beta/M^2),
\]
which are exactly the marginal hypotheses of the V34 ratio compiler.
\end{proof}

\begin{firewall}[What the four-loop formula does not decide]
\label{fire:s6v35-response-limit}
The attached exact Haar evaluator can certify all entries of
\Cref{thm:s6v35-bracket-slope} at a fixed cutoff.  It does not supply
uniform cofinal bounds on \(\kappa_M\), \(K_M\), the bracket opening, or
\(\operatorname{osc}(\Psi_M)\).  In particular, finite simplicity of every
root does not imply that \(L_M\) stays bounded away from zero.  The signed
quantity \(n_M\), not merely its upper bound
\eqref{eq:s6v35-n-bound}, remains necessary if the sign of the matched
action is used.
\end{firewall}

\section{The physical-clock correction to CDAC}
\label{sec:s6v35-clock}

The CDAC theorem in V25 used one fixed physical delay \(\tau>0\).
When a block scale varies with the regulator, the contraction per block and
the physical duration of that block must be renormalized together.  The
following statement makes this dependence explicit.

\begin{theorem}[Variable-clock OS spectral transfer]
\label{thm:s6v35-variable-clock}
Let \(\tau_j\downarrow0\).  For one fixed finite centred source bank, let
\(G_j(k)\) be its positive reflected Gram after \(k\) regulator transfer
steps.  Suppose that for one \(C_B<\infty\) and \(q_j\in(0,1)\),
\begin{equation}
 0\le x^*G_j(k)x
 \le C_Bq_j^k\,x^*G_j(0)x
 \qquad(k\ge1)                                        \label{eq:s6v35-lattice-Gram}
\end{equation}
for every coefficient vector \(x\).  Suppose also that for
\[
 k_j(t)=\lceil t/\tau_j\rceil
\]
the zero-delay and delayed Grams converge,
\begin{align}
 x^*G_j(0)x&\longrightarrow\|\psi_x\|^2,              \label{eq:s6v35-Gram-zero}\\
 x^*G_j(k_j(t))x
 &\longrightarrow
 \langle\psi_x,e^{-tH}\psi_x\rangle
 \quad(t>0),                                           \label{eq:s6v35-Gram-time}
\end{align}
where \(H\ge0\) is strongly continuous and the limiting bank is
nonzero.  Define
\[
 \mu_j={|\log q_j|\over\tau_j}.
\]
If
\begin{equation}
 0<\mu_*:=\liminf_{j\to\infty}\mu_j<\infty,
 \label{eq:s6v35-rate-liminf}
\end{equation}
then the cyclic subspace generated by the limiting bank obeys
\begin{equation}
 \boxed{H|_{{\cal C}_{\rm bank}}\succeq\mu_*I.}
 \label{eq:s6v35-variable-gap}
\end{equation}
If instead \(\mu_j\to\infty\), the hypotheses are incompatible with a
nonzero limiting bank.
\end{theorem}

\begin{proof}
Fix \(0<\mu<\mu_*\).  Eventually \(\mu_j\ge\mu\), and
\[
 q_j^{k_j(t)}
 =\exp\{-\mu_j\tau_j k_j(t)\}
 \le e^{-\mu t}.
\]
Pass to the limit in \eqref{eq:s6v35-lattice-Gram} using
\eqref{eq:s6v35-Gram-zero}--\eqref{eq:s6v35-Gram-time}:
\[
 0\le
 \langle\psi_x,e^{-tH}\psi_x\rangle
 \le C_Be^{-\mu t}\|\psi_x\|^2.
\]
The positive spectral measure of \(\psi_x\) can have no mass in
\([0,\mu)\): otherwise its Laplace transform times \(e^{\mu t}\) would be
unbounded as \(t\to\infty\).  Letting \(\mu\uparrow\mu_*\) proves
\eqref{eq:s6v35-variable-gap} for every bank vector and hence for its
cyclic closure.

If \(\mu_j\to\infty\), the same bound tends to zero for every fixed \(t>0\).
Thus
\(\langle\psi_x,e^{-tH}\psi_x\rangle=0\) for all \(t>0\).
Strong continuity at zero gives
\[
 \|\psi_x\|^2
 =\lim_{t\downarrow0}
   \langle\psi_x,e^{-tH}\psi_x\rangle=0.
\]
Every limiting bank vector is zero, contrary to noncollapse.
\end{proof}

\begin{corollary}[Fixed block contraction on a shrinking clock is
ultralocal]
\label{cor:s6v35-ultralocal}
Under the convergence and noncollapse hypotheses of
\Cref{thm:s6v35-variable-clock}, if
\(\tau_j\downarrow0\) and \(q_j\le q_*<1\), no nonzero strongly continuous
continuum bank exists.  A finite nonzero physical rate
\(\mu_j\to\mu\) instead requires
\begin{equation}
 q_j=e^{-\mu_j\tau_j}
     =1-\mu\tau_j+o(\tau_j).
 \label{eq:s6v35-critical-contraction}
\end{equation}
\end{corollary}

\begin{proof}
The fixed upper bound gives
\[
 \mu_j\ge{|\log q_*|\over\tau_j}\longrightarrow\infty,
\]
so the last branch of
\Cref{thm:s6v35-variable-clock} applies.  If
\(\mu_j\to\mu<\infty\), Taylor expansion of
\(e^{-\mu_j\tau_j}\) proves
\eqref{eq:s6v35-critical-contraction}.
\end{proof}

\begin{definition}[Yang--Mills Gram no-escape card, YMGNC]
\label{def:s6v35-YMGNC}
For a regulator sequence with block clock \(\tau_j\), YMGNC consists of:
\begin{enumerate}[label=\textup{(\roman*)},leftmargin=3.5em]
\item reflection-positive normalized cutoff laws and compatible physical
      time translations;
\item for every fixed gauge-invariant centred source bank, convergence of
      the zero-delay and moving-delay Grams in
      \eqref{eq:s6v35-Gram-zero}--\eqref{eq:s6v35-Gram-time};
\item a regulator-uniform lower frame, or an equivalent noncollapse
      certificate, and uniform integrability for every rescaled unbounded
      source used in the limiting bank;
\item a complete-law CDAC or other correlation estimate of the form
      \eqref{eq:s6v35-lattice-Gram}, with uniform bank prefactor and
      \[
       \liminf_j{|\log q_j|\over\tau_j}\ge\mu_0>0;
      \]
\item a dense increasing union of the limiting centred banks and uniqueness
      of the vacuum.
\end{enumerate}
\end{definition}

\begin{theorem}[YMGNC closes the physical mass handoff]
\label{thm:s6v35-YMGNC-gap}
If YMGNC holds and its limiting OS measure is nontrivial, then
\begin{equation}
 \boxed{
 H_{\rm YM}|_{\Omega^\perp}\succeq\mu_0I.}
 \label{eq:s6v35-YMGNC-gap}
\end{equation}
If the block clock is \(\tau_j=a_jL_j\downarrow0\), the contraction row
cannot be populated by one fixed \(q<1\); it must be calibrated at the
critical scale \eqref{eq:s6v35-critical-contraction}, or the source
noncollapse/convergence rows fail.
\end{theorem}

\begin{proof}
Apply \Cref{thm:s6v35-variable-clock} to each fixed bank.  Uniform
integrability is the exact condition needed to pass the rescaled
source-products to the limiting Grams when they are not bounded.  The common
lower rate \(\mu_0\) therefore places every bank's cyclic subspace above
\(\mu_0\).  Density and vacuum uniqueness identify the closure of their
union with \(\Omega^\perp\), giving
\eqref{eq:s6v35-YMGNC-gap}.  The final statement is
\Cref{cor:s6v35-ultralocal}.
\end{proof}

\begin{assessment}[The scale fork exposed by V35]
\label{ass:s6v35-fork}
Two independent scale mismatches are now explicit.
\begin{enumerate}[label=\textup{(\alph*)},leftmargin=3.5em]
\item The present full-chart static proof accepts selected roots only if
      they grow at least like \(M^p\) for some \(p>8\).
\item A CDAC factor bounded strictly below one per block cannot be carried
      through blocks whose physical clock tends to zero; a nontrivial
      continuum requires a near-critical factor
      \(q_j=1-\mu\tau_j+o(\tau_j)\).
\end{enumerate}
These are not contradictions in Yang--Mills theory.  They show that the
global all-link chart and fixed-factor disagreement estimates are too
coarse to be simultaneously read as a physical continuum proof on a
shrinking block scale.  The upstream repair is to prove an intensive
local-density/cluster no-escape estimate and a scale-calibrated
correlation loss.
\end{assessment}

\begin{programme}[V36: remove the volume-entropy bottleneck]
\label{prog:s6v35-V36}
\begin{enumerate}[leftmargin=3em]
\item Replace the event that all \(45M^4\) reduced links share one chart by
      a local bad-set density or polymer expansion compatible with the
      normalized axial Fourier quadratic form.
\item Prove a local Schur/restoration estimate in which the debit is the
      expected bad density and cluster diameter, not the probability of one
      global good event.
\item Recompute the admissible \((M,\beta,r)\) corridor.  Only then compare
      it with the actual four-loop root growth.
\item For the response marginal, populate
      \(\kappa_M,L_M,\operatorname{osc}(\Psi_M)\) in
      \eqref{eq:s6v35-scaled-rows}, retaining the sign of \(n_M\).
\item Express CDAC in physical-rate form
      \(|\log q_j|/\tau_j\), and require the Gram no-escape rows before
      invoking the continuum mass theorem.
\end{enumerate}
\end{programme}

\begin{firewall}[Claim boundary after sixth-series V35]
\label{fire:s6v35-boundary}
V35 performs the compulsory SM/NS audit; proves that the existing
global-chart root landing is exactly a super-octic root-growth requirement;
derives a partition-free oscillation bound for the generated four-loop
response numerator; proves a two-sided response-bracket slope floor and a
scale-dependent replacement for uniform response regularity; proves the
variable-clock OS spectral theorem and the ultralocality obstruction for a
fixed contraction on shrinking physical blocks; and formulates YMGNC, which
conditionally closes the continuum mass handoff.

V35 does not prove the actual cofinal root asymptotics, the local-density
replacement of the global chart, the base generated-action stocks, TSC,
scale-calibrated CDAC, YMGNC, existence of the nontrivial reflection-positive
continuum measure, vacuum uniqueness, or the full Yang--Mills mass gap.
\end{firewall}

\begin{theorem}[Sixth-series V35 result ledger]
\label{thm:s6v35-results}
The claims in \Cref{fire:s6v35-boundary} follow from
\Cref{thm:s6v35-audit-decision,
thm:s6v35-super-octic,
cor:s6v35-root-certificate,
lem:s6v35-n-bound,
thm:s6v35-bracket-slope,
cor:s6v35-scaled-response,
thm:s6v35-variable-clock,
cor:s6v35-ultralocal,
thm:s6v35-YMGNC-gap}.
\end{theorem}

\begin{proof}
Each ledger item is the conclusion of the cited result.
\end{proof}

\paragraph{Parent-version record.}
V35 was formed from
\texttt{Renewal\_Yang\_Mills\_Mass\_Gap\_Research\_Notes\_V34.tex}.
The V34 parent had \(15\,389\) logical source lines,
\(542\,585\) bytes, and SHA--256
\[
 \texttt{9e3cd3090f7b839a7f62f545e3ebf78bd5bf9d1a5e1638b5f056375453510098}.
\]
The next compulsory companion audit is V40.

\clearpage
\chapter{Sixth-series V36: intensive Wilson concentration and a
local covariance target}
\label{chap:s6v36}

\section{Why the partition-function entropy can be divided by volume}
\label{sec:s6v36-energy}

Retain the pinned pure-Wilson reduced fibre at \(W=I\), with
\[
 n_M=45M^4,\qquad
 D(z)=\sum_e d_e(z),\qquad
 d_e(z)=\|z_e-I\|_{\rm F}^2,
\]
and the exact comparison imported in V31,
\[
 {1\over120}D(z)\le V_M(z,I)\le4D(z).
\]
Also retain
\[
 I(t)=\int_{\SU(3)}e^{-t\|U-I\|_{\rm F}^2}\,dU,
 \qquad
 C_-t^{-4}\le I(t)\le C_+t^{-4}\quad(t\ge1),
\]
and
\[
 {\cal C}_0={C_+\over C_-}\,960^4.
\]

\begin{theorem}[Volume-uniform mean Wilson energy]
\label{thm:s6v36-mean-energy}
For \(\beta\ge240\), the pinned pure-Wilson law
\[
 d\nu_{\beta,M}=Z_{\beta,M}^{-1}e^{-\beta V_M}\,dm
\]
satisfies
\begin{equation}
 \boxed{
 \mathbb E_{\beta,M}V_M
 \le {2n_M\log{\cal C}_0\over\beta},
 \qquad
 \mathbb E_{\beta,M}D
 \le {240n_M\log{\cal C}_0\over\beta}.}
 \label{eq:s6v36-mean-energy}
\end{equation}
The constants are independent of \(M\).
\end{theorem}

\begin{proof}
Put \(f(\beta)=\log Z_{\beta,M}\).  Differentiation under the compact
integral gives
\[
 f'(\beta)=-\mathbb E_{\beta,M}V_M,
 \qquad
 f''(\beta)=\operatorname{Var}_{\beta,M}(V_M)\ge0.
\]
Convexity therefore implies
\[
 \mathbb E_{\beta,M}V_M
 =-f'(\beta)
 \le {2\over\beta}\{f(\beta/2)-f(\beta)\}.
\]
The action comparison gives
\[
 Z_{\beta/2,M}
 \le I(\beta/240)^{n_M},
 \qquad
 Z_{\beta,M}
 \ge I(4\beta)^{n_M}.
\]
The one-link integral bounds are applicable because \(\beta\ge240\), and
hence
\[
 {Z_{\beta/2,M}\over Z_{\beta,M}}
 \le
 \left\{
 {C_+(\beta/240)^{-4}\over C_-(4\beta)^{-4}}
 \right\}^{n_M}
 ={\cal C}_0^{n_M}.
\]
This proves the first inequality in
\eqref{eq:s6v36-mean-energy}.  Since \(D\le120V_M\) pointwise, the second
follows.
\end{proof}

\begin{corollary}[Intensive bad-link density]
\label{cor:s6v36-bad-density}
For \(\rho>0\), define
\[
 {\mathfrak b}_{M,\rho}(z)
 ={1\over n_M}\sum_e
 {\bf1}_{\{\|z_e-I\|_{\rm F}\ge\rho\}}.
\]
Then
\begin{align}
 \boxed{
 \mathbb E_{\beta,M}{\mathfrak b}_{M,\rho}
 \le {240\log{\cal C}_0\over\beta\rho^2},}
 \label{eq:s6v36-bad-density}\\
 \boxed{
 \mathbb E_{\beta,M}
 \left[{1\over n_M}\sum_e\|z_e-I\|_{\rm F}\right]
 \le
 \left({240\log{\cal C}_0\over\beta}\right)^{1/2}.}
 \label{eq:s6v36-mean-distance}
\end{align}
In particular, if \(\rho_\beta=\beta^{-\alpha}\) with
\(\alpha<1/2\), the expected bad density tends to zero for arbitrary
growth of \(M\).
\end{corollary}

\begin{proof}
The pointwise inequality
\[
 {\bf1}_{\{\|z_e-I\|_{\rm F}\ge\rho\}}
 \le {d_e(z)\over\rho^2}
\]
and \Cref{thm:s6v36-mean-energy} give
\eqref{eq:s6v36-bad-density}.  Cauchy--Schwarz gives
\[
 {1\over n_M}\sum_e\sqrt{d_e}
 \le\sqrt{D/n_M}.
\]
Apply Jensen's inequality followed by
\eqref{eq:s6v36-mean-energy} to obtain
\eqref{eq:s6v36-mean-distance}.  The last assertion follows because
\(\beta\rho_\beta^2=\beta^{1-2\alpha}\to\infty\).
\end{proof}

\begin{assessment}[What happened to the \(M^4\) entropy]
\label{ass:s6v36-entropy}
The global event \(\max_e\|z_e-I\|<\rho\) pays for all
\(n_M\) coordinates at once and led to the V31 term
\(n_M\log{\cal C}_0\).  The derivative of the logarithmic partition
function measures total energy.  Dividing it by \(n_M\) before applying
Markov or Cauchy--Schwarz gives
\eqref{eq:s6v36-bad-density}--\eqref{eq:s6v36-mean-distance}, where the
volume cancels.  This is the first rigorous local-density replacement of
the global chart.  It is not yet a covariance estimate.
\end{assessment}

\section{Finite-incidence writers see the mean distance}
\label{sec:s6v36-writers}

\begin{lemma}[Local-writer averaging]
\label{lem:s6v36-local-writer}
Let \(E_M\) be a graph with \(n_M\) sites and uniformly bounded local
degree.  Suppose a Hermitian quadratic writer has a decomposition
\[
 R_M(z)[v,v]=\sum_X r_X(z_X;v_X),
\]
where every \(X\) has at most \(s\) sites, every site belongs to at most
\(\nu\) sets, and
\begin{equation}
 |r_X(z_X;v_X)-r_X(I;v_X)|
 \le L\left(\sum_{e\in X}\sqrt{d_e(z)}\right)
        \left(\sum_{f\in X}|v_f|^2\right).
 \label{eq:s6v36-local-Lipschitz}
\end{equation}
If \(\|v\|_2=1\) and
\(\max_f|v_f|^2\le c_v/n_M\), then
\begin{equation}
 \boxed{
 |R_M(z)[v,v]-R_M(I)[v,v]|
 \le {Lsc_v\nu\over n_M}\sum_e\sqrt{d_e(z)}.}
 \label{eq:s6v36-writer-point}
\end{equation}
Under the Wilson law,
\begin{equation}
 \boxed{
 \left|
 \mathbb E_{\beta,M}R_M(z)[v,v]-R_M(I)[v,v]
 \right|
 \le C_R\beta^{-1/2},}
 \label{eq:s6v36-writer-mean}
\end{equation}
where \(C_R=Lsc_v\nu(240\log{\cal C}_0)^{1/2}\) is volume
independent.
\end{lemma}

\begin{proof}
In \eqref{eq:s6v36-local-Lipschitz},
\(\sum_{f\in X}|v_f|^2\le sc_v/n_M\).  After summing over \(X\), every
\(\sqrt{d_e}\) occurs at most \(\nu\) times.  This proves
\eqref{eq:s6v36-writer-point}.  Take expectations and use
\eqref{eq:s6v36-mean-distance}.
\end{proof}

\begin{corollary}[The deterministic retained Wilson writer has a
super-quartic corridor]
\label{cor:s6v36-deterministic}
Let \(C_M(z)=D_w^2V_M(z,0)\) be the unit-Wilson retained second writer
on a normalized axial Fourier polarization \(v_q\).  There is
\(C_{\rm det}<\infty\), independent of \(M,\beta,q\), such that
\begin{equation}
 \mathbb E_{\beta,M}\{\beta C_M(z)[v_q,v_q]\}
 \ge
 \beta\,\mathsf S_M[v_q,v_q]-C_{\rm det}\sqrt\beta.
 \label{eq:s6v36-deterministic}
\end{equation}
Consequently the deterministic debit is
\(o(\beta s_0(q))\), uniformly over \(q\ne0\), whenever
\begin{equation}
 {M^2\over\sqrt\beta}\longrightarrow0.
 \label{eq:s6v36-super-quartic}
\end{equation}
\end{corollary}

\begin{proof}
Every plaquette contains four links and has globally bounded fixed-order
derivatives on the compact group.  Thus \(C_M\) has a decomposition
satisfying \Cref{lem:s6v36-local-writer}; a normalized plane wave has
pointwise squared amplitude \(O(M^{-4})=O(n_M^{-1})\).  Hence
\[
 \mathbb E C_M(z)[v_q,v_q]
 \ge C_M(I)[v_q,v_q]-C\beta^{-1/2}.
\]
At the identity,
\[
 \mathsf S_M
 =C_M(I)-B_M(I)^*A_M(I)^{-1}B_M(I)
 \preceq C_M(I),
\]
which proves \eqref{eq:s6v36-deterministic}.  The V31 Ward bound gives
\[
 s_0(q)\ge{20\over7M^2}
\qquad(q\ne0).
\]
The error relative to \(\beta s_0(q)\) is therefore at most
\(CM^2/\sqrt\beta\), which vanishes under
\eqref{eq:s6v36-super-quartic}.
\end{proof}

\begin{firewall}[The covariance subtraction has not disappeared]
\label{fire:s6v36-covariance}
The full trace Hessian is
\[
 H_{\beta,M}[v,v]
 =\mathbb E\{\beta C_M(z)[v,v]\}
  -\operatorname{Var}\{\beta D_wV_M(z,0)[v]\}.
\]
\Cref{cor:s6v36-deterministic} controls only the first term.  A
bad-density bound does not justify applying Brascamp--Lieb on a nonconvex
configuration set and does not control long coherent correlations among
the exceptional links.
\end{firewall}

\begin{proposition}[Counterexample: vanishing bad density need not control a normalized
score variance]
\label{cth:s6v36-density-not-covariance}
There are bounded random arrays
\((X_1^{(n)},\ldots,X_n^{(n)})\), \(|X_i^{(n)}|\le1\), such that
\[
 {1\over n}\sum_i{\bf1}_{\{X_i^{(n)}\ne0\}}
 \longrightarrow0
\quad\hbox{almost surely},
\]
while the flat normalized score
\[
 J_n={1\over\sqrt n}\sum_iX_i^{(n)}
\]
has \(\operatorname{Var}(J_n)\to\infty\).
\end{proposition}

\begin{proof}
Let \(m_n=\lfloor n^{3/4}\rfloor\), let \(S\) take the values
\(\pm1\) with equal probability, and set
\[
 X_i^{(n)}=
 \begin{cases}
 S,&i\le m_n,\\
 0,&i>m_n.
 \end{cases}
\]
The bad density is \(m_n/n\to0\).  On the other hand
\[
 J_n={m_n\over\sqrt n}S,
 \qquad
 \operatorname{Var}(J_n)={m_n^2\over n}\sim n^{1/2}.
\]
The example uses coherent exceptional sites.  It proves that density must
be supplemented by a cluster-diameter or connected-covariance estimate.
\end{proof}

\section{A local connected-kernel replacement for the global chart}
\label{sec:s6v36-LCKC}

Index retained coarse-link components by a bounded-degree four-dimensional
torus graph.  Write the unit-action retained score as
\[
 D_wV_M(z,0)[v]=\sum_a v_aY_a(z),
\]
and let
\[
 R_M^0=B_M(I)^*A_M(I)^{-1}B_M(I).
\]
All colour and direction blocks below are included in the indices.

\begin{definition}[Local covariance-kernel card, LCKC]
\label{def:s6v36-LCKC}
LCKC holds if there are \(c,C>0\), independent of \(M,\beta\), such that
\begin{equation}
 \left\|
 \operatorname{Cov}_{\beta,M}(Y_a,Y_b)
 -{1\over\beta}(R_M^0)_{ab}
 \right\|_{\rm op}
 \le C\beta^{-3/2}e^{-c\,\operatorname{dist}(a,b)}
 \label{eq:s6v36-LCKC}
\end{equation}
for every pair of retained local components.
\end{definition}

\begin{lemma}[LCKC gives the missing intensive covariance error]
\label{lem:s6v36-LCKC-variance}
Under LCKC, every normalized retained vector \(v\) satisfies
\begin{equation}
 \boxed{
 \left|
 \operatorname{Var}_{\beta,M}
 \{\beta D_wV_M[v]\}
 -\beta R_M^0[v,v]
 \right|
 \le C_{\rm cov}\sqrt\beta,}
 \label{eq:s6v36-LCKC-variance}
\end{equation}
with \(C_{\rm cov}\) independent of \(M\).
\end{lemma}

\begin{proof}
On a bounded-degree four-dimensional graph,
\[
 \sup_a\sum_b e^{-c\,\operatorname{dist}(a,b)}<\infty
\]
uniformly in the torus size.  The Schur test applied to
\eqref{eq:s6v36-LCKC} therefore gives operator norm at most
\(C'\beta^{-3/2}\) for the covariance-kernel difference.  The score has
the prefactor \(\beta\), so its variance difference has norm at most
\(\beta^2C'\beta^{-3/2}=C'\sqrt\beta\).
\end{proof}

\begin{theorem}[Intensive pure-Wilson static theorem under LCKC]
\label{thm:s6v36-intensive-static}
Assume LCKC along a cofinal sequence with
\[
 M\longrightarrow\infty,\qquad
 \beta\longrightarrow\infty,\qquad
 M^2/\sqrt\beta\longrightarrow0.
\]
Then for all sufficiently large members of the sequence and every nonzero
axial momentum,
\begin{equation}
 \boxed{
 H_{\beta,M}(q)\succeq{1\over2}\beta s_0(q)I.}
 \label{eq:s6v36-intensive-static}
\end{equation}
\end{theorem}

\begin{proof}
Combine
\Cref{cor:s6v36-deterministic,lem:s6v36-LCKC-variance}:
\begin{align*}
 H_{\beta,M}[v_q,v_q]
 &\ge
 \beta C_M(I)[v_q,v_q]
 -\beta R_M^0[v_q,v_q]
 -C\sqrt\beta\\
 &=\beta\mathsf S_M[v_q,v_q]-C\sqrt\beta.
\end{align*}
Since \(s_0(q)\ge20/(7M^2)\), the relative error is at most
\(CM^2/\sqrt\beta\), which is below \(1/2\) eventually.  Central colour
symmetry scalarizes the axial block exactly as in V31.
\end{proof}

\begin{corollary}[Root-growth debit after an LCKC replacement]
\label{cor:s6v36-super-quartic-root}
If LCKC is proved and the selected roots are used in
\Cref{thm:s6v36-intensive-static}, root landing requires only
\[
 M=o(\widehat\beta_{2M}^{\,1/4}).
\]
In particular, a sufficient power law is
\(\widehat\beta_{2M}\ge cM^p\) for any \(p>4\), in place of the
super-octic global-chart condition of V35.
\end{corollary}

\begin{proof}
The first condition is
\eqref{eq:s6v36-super-quartic} with
\(\beta=\widehat\beta_{2M}\).  If
\(\widehat\beta_{2M}\ge cM^p\), then
\[
 {M^2\over\sqrt{\widehat\beta_{2M}}}
 \le c^{-1/2}M^{2-p/2}\longrightarrow0
\]
exactly when \(p>4\).
\end{proof}

\begin{assessment}[The new first missing inequality]
\label{ass:s6v36-frontier}
\Cref{thm:s6v36-mean-energy,cor:s6v36-bad-density} remove volume entropy
from one-point and local-writer control.  They also show that the
deterministic retained curvature has the larger corridor
\(M=o(\beta^{1/4})\).  The first unsupplied pure-Wilson statement is now
the connected covariance asymptotic
\eqref{eq:s6v36-LCKC}, including its \(\beta^{-3/2}\) remainder and
uniform exponential spatial summability.  It cannot be inferred from
vanishing bad density by
\Cref{cth:s6v36-density-not-covariance}.  A Wilson-specific pinned-fibre
cluster expansion, a Helffer--Sj\"ostrand representation with controlled
bad polymers, or an equivalent conditional mixing theorem is required.
\end{assessment}

\begin{programme}[V37: attack LCKC at its first non-Gaussian order]
\label{prog:s6v36-V37}
\begin{enumerate}[leftmargin=3em]
\item Expand the pinned detail field at scale
      \(x=\beta^{-1/2}y\) and identify the exact Gaussian covariance
      \(A_M(I)^{-1}\).
\item Use the uniform vertical gap and finite range to prove exponential
      decay of \(A_M(I)^{-1}\) with constants independent of \(M\).
\item Bound the cubic and higher local vertices in a weighted connected
      norm, separating the small-field region from bad polymers using
      \eqref{eq:s6v36-bad-density}.
\item Derive the score covariance remainder at order
      \(O(\beta^{-3/2})\), or record the first polymer activity whose
      uniform summability is missing.
\item Only after LCKC is proved, combine it with the response and
      topological stocks; do not label the deterministic result a full
      trace Hessian.
\end{enumerate}
\end{programme}

\begin{firewall}[Claim boundary after sixth-series V36]
\label{fire:s6v36-boundary}
V36 proves a volume-uniform \(O(\beta^{-1})\) mean energy per reduced
link, an intensive bad-link-density bound, and an
\(O(\beta^{-1/2})\) local-writer expectation bound.  It applies these
estimates to the deterministic retained Wilson Hessian, proves by
counterexample that bad density alone does not control the score
covariance, formulates the precise local covariance-kernel card LCKC, and
proves that LCKC would enlarge the pure-Wilson static corridor from
\(M=o(\beta^{1/8})\) to \(M=o(\beta^{1/4})\).

V36 does not prove LCKC, the response-root growth needed even by the
enlarged corridor, the generated-action version of the covariance kernel,
TSC, YMGNC, continuum existence, or the full Yang--Mills mass gap.
\end{firewall}

\begin{theorem}[Sixth-series V36 result ledger]
\label{thm:s6v36-results}
The claims in \Cref{fire:s6v36-boundary} follow from
\Cref{thm:s6v36-mean-energy,
cor:s6v36-bad-density,
lem:s6v36-local-writer,
cor:s6v36-deterministic,
cth:s6v36-density-not-covariance,
lem:s6v36-LCKC-variance,
thm:s6v36-intensive-static,
cor:s6v36-super-quartic-root}.
\end{theorem}

\begin{proof}
Each ledger item is the conclusion of the cited result.
\end{proof}

\paragraph{Parent-version record.}
V36 was formed from
\texttt{Renewal\_Yang\_Mills\_Mass\_Gap\_Research\_Notes\_V35.tex}.
The V35 parent had \(15\,944\) logical source lines,
\(562\,659\) bytes, and SHA--256
\[
 \texttt{12e332699e7819595ab4550e9457796d0a8700d2fb9e4d2df63c9ef6c7551909}.
\]
The next compulsory companion audit is V40.

\clearpage
\chapter{Sixth-series V37: the pinned Gaussian kernel and the
connected nonlinear remainder}
\label{chap:s6v37}

\section{Uniform locality of the vertical Gaussian inverse}
\label{sec:s6v37-inverse}

Let \({\cal E}_M\) be the reduced-detail link graph.  Matrix blocks carry
the eight colour components and any fixed direction multiplicity.  The
unit-Wilson vertical Hessian
\[
 A_M=A_M(I)
\]
has finite graph range \(R_A\), independent of \(M\), and the attached
reduced-fibre calculation imported in V31 gives
\begin{equation}
 \lambda I\preceq A_M\preceq\Lambda I,
 \qquad
 \lambda={1\over60},
 \qquad
 \Lambda=8.
 \label{eq:s6v37-spectrum}
\end{equation}

\begin{theorem}[Explicit exponential inverse bound]
\label{thm:s6v37-inverse}
Put
\[
 q_A=1-{\lambda\over\Lambda}={479\over480}.
\]
For reduced-link blocks \(e,f\),
\begin{equation}
 \boxed{
 \|(A_M^{-1})_{ef}\|_{\rm op}
 \le {1\over\lambda}\,
 q_A^{\,\lceil\operatorname{dist}(e,f)/R_A\rceil}.}
 \label{eq:s6v37-inverse}
\end{equation}
The estimate is uniform in \(M\).  In particular, for every
\(0<c_A<R_A^{-1}|\log q_A|\),
\begin{equation}
 \sup_e\sum_f e^{c_A\operatorname{dist}(e,f)}
 \|(A_M^{-1})_{ef}\|_{\rm op}<\infty
 \label{eq:s6v37-weighted-inverse}
\end{equation}
uniformly in the torus size.
\end{theorem}

\begin{proof}
Set \(K_M=I-\Lambda^{-1}A_M\).  By
\eqref{eq:s6v37-spectrum},
\[
 0\preceq K_M\preceq q_AI.
\]
The Neumann series converges in operator norm:
\[
 A_M^{-1}
 ={1\over\Lambda}\sum_{k=0}^\infty K_M^k.
\]
Since \(K_M\) has graph range \(R_A\), its \(ef\) block vanishes when
\(kR_A<\operatorname{dist}(e,f)\).  If
\(k_0=\lceil\operatorname{dist}(e,f)/R_A\rceil\), then
\[
 \|(A_M^{-1})_{ef}\|
 \le {1\over\Lambda}\sum_{k\ge k_0}\|K_M\|^k
 \le {q_A^{k_0}\over\Lambda(1-q_A)}
 ={q_A^{k_0}\over\lambda}.
\]
This is \eqref{eq:s6v37-inverse}.  A bounded-degree
four-dimensional graph has at most \(C(1+r)^3\) vertices at distance
\(r\).  The polynomial growth is summable against
\(\exp\{-(R_A^{-1}|\log q_A|-c_A)r\}\), proving
\eqref{eq:s6v37-weighted-inverse}.
\end{proof}

Let \(B_M=B_M(I)=D_xD_wV_M(0,0)\).  It also has fixed graph range and
uniform row and column sums.  Recall
\[
 R_M^0=B_M^*A_M^{-1}B_M.
\]

\begin{corollary}[Exponential locality of the flat Schur covariance]
\label{cor:s6v37-flat-kernel}
There are \(c_0,C_0>0\), independent of \(M\), such that
\begin{equation}
 \boxed{
 \|(R_M^0)_{ab}\|_{\rm op}
 \le C_0e^{-c_0\operatorname{dist}(a,b)}.}
 \label{eq:s6v37-flat-kernel}
\end{equation}
\end{corollary}

\begin{proof}
The \(ab\) block is a sum of
\[
 (B_M)_{ea}^*(A_M^{-1})_{ef}(B_M)_{fb}
\]
over \(e\) in a fixed-radius neighbourhood of \(a\) and \(f\) in a
fixed-radius neighbourhood of \(b\).  The two neighbourhood cardinalities
and the \(B_M\) blocks are uniformly bounded.  Apply
\Cref{thm:s6v37-inverse} and absorb the fixed distance shift into the
constant.
\end{proof}

\section{The exact Gaussian card}
\label{sec:s6v37-Gaussian}

Let \(\gamma_M\) be the centred Gaussian law on the vertical Lie-algebra
fibre with density proportional to
\[
 \exp\{-\tfrac12\langle y,A_My\rangle\}\,dy.
\]
For a retained coordinate \(a\), let
\[
 L_a(y)=\langle B_Me_a,y\rangle.
\]

\begin{proposition}[The quadratic model satisfies LCKC exactly]
\label{prop:s6v37-Gaussian}
For the quadratic action at inverse coupling \(\beta\), with
\(x=\beta^{-1/2}y\),
\begin{equation}
 \boxed{
 \operatorname{Cov}_{\gamma_M}
 \bigl(\beta^{-1/2}L_a,\beta^{-1/2}L_b\bigr)
 ={1\over\beta}(R_M^0)_{ab}.}
 \label{eq:s6v37-Gaussian-cov}
\end{equation}
The covariance kernel is exponentially summable uniformly in \(M\).
\end{proposition}

\begin{proof}
The Gaussian covariance of \(y\) is \(A_M^{-1}\).  Therefore
\[
 \operatorname{Cov}_{\gamma_M}(L_a,L_b)
 =\langle B_Me_a,A_M^{-1}B_Me_b\rangle
 =(R_M^0)_{ab}.
\]
Scaling proves \eqref{eq:s6v37-Gaussian-cov}, and
\Cref{cor:s6v37-flat-kernel} gives the last claim.
\end{proof}

\begin{assessment}[A massless coarse symbol does not make the detail
inverse long range]
\label{ass:s6v37-detail-gap}
The smallest retained axial Schur eigenvalue is \(O(M^{-2})\), but the
vertical detail operator has the fixed gap \(1/60\).  Hence
\(A_M^{-1}\) and \(R_M^0\) are exponentially local even though the
coarse physical symbol has a Ward zero at \(q=0\).  The LCKC target is
therefore compatible with the exact pinned geometry; its difficulty is
the nonlinear full law, not the Gaussian inverse.
\end{assessment}

\section{An exact connected interpolation identity}
\label{sec:s6v37-interpolation}

The next lemma is finite dimensional and will be applied only after a
small-field chart or a polymer decomposition has supplied a common
integration space.  Let
\[
 d\mu_t=Z_t^{-1}e^{-Q-t{\cal R}}\,d\rho,
 \qquad 0\le t\le1,
\]
where all displayed expectations are finite.  For random variables
\(F,G,H\), write
\[
 \operatorname{Cum}_t(F,G,H)
 =\mathbb E_t
 [(F-\mathbb E_tF)(G-\mathbb E_tG)(H-\mathbb E_tH)].
\]

\begin{lemma}[Covariance Duhamel identity]
\label{lem:s6v37-cov-Duhamel}
If differentiation under the integral is justified, then
\begin{equation}
 \boxed{
 \operatorname{Cov}_{\mu_1}(F,G)
 -\operatorname{Cov}_{\mu_0}(F,G)
 =-\int_0^1
 \operatorname{Cum}_t(F,G,{\cal R})\,dt.}
 \label{eq:s6v37-cov-Duhamel}
\end{equation}
\end{lemma}

\begin{proof}
For every \(K\),
\[
 {d\over dt}\mathbb E_tK
 =-\operatorname{Cov}_t(K,{\cal R}).
\]
Apply this once to \(K=FG\), once to \(K=F\), and once to \(K=G\);
differentiate
\(\operatorname{Cov}_t(F,G)=\mathbb E_t(FG)-\mathbb E_tF\,\mathbb E_tG\).
The resulting five terms are exactly
\(-\operatorname{Cum}_t(F,G,{\cal R})\).  Integrate from zero to one.
\end{proof}

In the rescaled Wilson chart, Taylor expansion has the local form
\begin{align}
 \beta V_M(\beta^{-1/2}y,0)
 &=\tfrac12\langle y,A_My\rangle+{\cal R}_{\beta,M}(y),
 \label{eq:s6v37-action-expand}\\
 Y_a(\beta^{-1/2}y)
 &=\beta^{-1/2}L_a(y)+\beta^{-1}Q_{a,\beta,M}(y),
 \label{eq:s6v37-score-expand}
\end{align}
where the coefficients of each fixed local Taylor jet are bounded
uniformly in \(M\).  The factors of \(\beta\) are exact consequences of
the rescaling; they are not yet expectation bounds.

\begin{definition}[Connected perturbation and tail card, CPTC]
\label{def:s6v37-CPTC}
CPTC with exponent \(\sigma>0\) consists of a chart/polymer
representation of the full pinned Wilson law for which, uniformly in
\(M,\beta\),
\begin{enumerate}[label=\textup{(\roman*)},leftmargin=3.5em]
\item the covariance contributions containing at least one score remainder
      \(\beta^{-1}Q_{a,\beta,M}\) differ from their declared Gaussian
      contractions by at most
      \(C\beta^{-1-\sigma}e^{-c\operatorname{dist}(a,b)}\);
\item the connected interpolation term satisfies
      \[
       \int_0^1
       |\operatorname{Cum}_t(L_a,L_b,{\cal R}_{\beta,M})|\,dt
       \le
       C\beta^{-\sigma}e^{-c\operatorname{dist}(a,b)};
      \]
      after the two exterior factors \(\beta^{-1/2}\), this contributes
      \(C\beta^{-1-\sigma}e^{-c\operatorname{dist}(a,b)}\);
\item polymers meeting the complement of the chosen local charts change
      the two-point score kernel by at most
      \(C\beta^{-1-\sigma}e^{-c\operatorname{dist}(a,b)}\).
\end{enumerate}
\end{definition}

\begin{theorem}[CPTC compiles to a generalized LCKC]
\label{thm:s6v37-CPTC}
CPTC implies
\begin{equation}
 \boxed{
 \left\|
 \operatorname{Cov}_{\beta,M}(Y_a,Y_b)
 -{1\over\beta}(R_M^0)_{ab}
 \right\|_{\rm op}
 \le
 C'\beta^{-1-\sigma}
 e^{-c'\operatorname{dist}(a,b)}.}
 \label{eq:s6v37-general-LCKC}
\end{equation}
\end{theorem}

\begin{proof}
Expand both score writers by
\eqref{eq:s6v37-score-expand}.  The linear--linear Gaussian term is
\(\beta^{-1}R_M^0\) by
\Cref{prop:s6v37-Gaussian}.  Apply
\Cref{lem:s6v37-cov-Duhamel} to the linear--linear change of law; CPTC
item \textup{(ii)} supplies its bound after multiplication by
\(\beta^{-1}\).  Item \textup{(i)} bounds every term containing a score
remainder.  Item \textup{(iii)} restores the chart complement.  Adding
the three exponentially decaying kernels and decreasing the exponent if
necessary proves \eqref{eq:s6v37-general-LCKC}.
\end{proof}

\begin{theorem}[Static corridor for a general connected exponent]
\label{thm:s6v37-sigma-corridor}
Assume \eqref{eq:s6v37-general-LCKC} with \(\sigma>0\), together with the
deterministic estimate of V36.  Put
\[
 \theta=\min\{\sigma,1/2\}.
\]
Then the pure-Wilson full trace axial margin obeys
\begin{equation}
 H_{\beta,M}(q)
 \succeq\beta s_0(q)I-C\beta^{1-\theta}I.
 \label{eq:s6v37-sigma-margin}
\end{equation}
It is at least \(\frac12\beta s_0(q)I\), uniformly for \(q\ne0\), on
every sequence satisfying
\begin{equation}
 M^2\beta^{-\theta}\longrightarrow0.
 \label{eq:s6v37-sigma-corridor}
\end{equation}
If \(\beta\ge cM^p\), a sufficient root-growth exponent is
\begin{equation}
 p>{2\over\theta}.
 \label{eq:s6v37-root-exponent}
\end{equation}
\end{theorem}

\begin{proof}
The Schur test applied to
\eqref{eq:s6v37-general-LCKC} gives covariance-kernel operator error
\(O(\beta^{-1-\sigma})\).  Multiplication by the two score prefactors
\(\beta\) gives the full score-variance error
\(O(\beta^{1-\sigma})\).  The deterministic error from V36 is
\(O(\beta^{1/2})\).  Their sum is
\(O(\beta^{1-\theta})\), proving
\eqref{eq:s6v37-sigma-margin}.  Divide by
\(\beta s_0(q)\ge20\beta/(7M^2)\) to obtain
\eqref{eq:s6v37-sigma-corridor}.  If
\(\beta\ge cM^p\), then
\[
 M^2\beta^{-\theta}
 \le c^{-\theta}M^{2-p\theta}\to0
\]
whenever \eqref{eq:s6v37-root-exponent} holds.
\end{proof}

\begin{firewall}[An extensive Taylor remainder is not CPTC]
\label{fire:s6v37-extensive}
The rescaled remainder \({\cal R}_{\beta,M}\) is a sum of
\(O(M^4)\) local vertices.  Bounding its absolute size before taking the
connected cumulant in \eqref{eq:s6v37-cov-Duhamel} produces a volume
factor and does not prove CPTC.  The two marked scores must remain in one
connected cluster, while vacuum polymers cancel against the normalizer.
Likewise, the bad-density estimate of V36 controls the mean number of bad
links but not the diameter or coherent connectivity of a bad polymer.
\end{firewall}

\begin{assessment}[Exact first nonlinear target]
\label{ass:s6v37-frontier}
The flat part of LCKC is now closed with explicit exponential constants.
The first nonlinear target is the connected three-point quantity
\[
 \operatorname{Cum}_t(L_a,L_b,{\cal R}_{\beta,M})
\]
in a norm which sums its spatial tail uniformly in \(M\), together with
the same connected control for the score Taylor remainder and bad
polymers.  A bound with \(\sigma=1/2\) reproduces the V36
\(M=o(\beta^{1/4})\) corridor.  Any \(\sigma>0\) already improves on a
dimension-dependent absolute expansion, though the deterministic
\(O(\sqrt\beta)\) debit caps this particular compiler at
\(\theta=1/2\).
\end{assessment}

\begin{programme}[V38: a pinned-fibre polymer norm]
\label{prog:s6v37-V38}
\begin{enumerate}[leftmargin=3em]
\item Freeze local balls of radius
      \(\rho_\beta=\beta^{-\alpha}\), \(0<\alpha<1/2\), and represent
      connected chart failures by polymers on the reduced-link incidence
      graph.
\item Derive a conditional energy cost for a connected bad polymer which
      is stable under its exterior, or prove that the frustrated-star
      obstruction of V28 prevents such a cost.
\item On the good complement, use the vertical gap and
      \Cref{thm:s6v37-inverse} to build a weighted connected expansion for
      the cubic and higher local vertices.
\item Test the Koteck\'y--Preiss sum for the marked two-score clusters.
      Keep vacuum clusters cancelled before estimating volume.
\item If worst-case exterior stability fails, return to a typical-exterior
      coupling and state the additional no-escape input explicitly.
\end{enumerate}
\end{programme}

\begin{firewall}[Claim boundary after sixth-series V37]
\label{fire:s6v37-boundary}
V37 proves a uniform exponential kernel bound for the exact pinned
vertical inverse and the flat Schur covariance, proves that the quadratic
detail model satisfies LCKC exactly, derives an exact covariance Duhamel
identity, formulates the connected perturbation and tail card CPTC, proves
that CPTC implies a generalized LCKC, and computes the resulting static and
root-growth corridors.

V37 does not prove CPTC for the nonlinear Wilson law, a bad-polymer
activity bound, the selected root asymptotics, the generated-action
extension, TSC, YMGNC, continuum existence, or the full Yang--Mills mass
gap.
\end{firewall}

\begin{theorem}[Sixth-series V37 result ledger]
\label{thm:s6v37-results}
The claims in \Cref{fire:s6v37-boundary} follow from
\Cref{thm:s6v37-inverse,
cor:s6v37-flat-kernel,
prop:s6v37-Gaussian,
lem:s6v37-cov-Duhamel,
thm:s6v37-CPTC,
thm:s6v37-sigma-corridor}.
\end{theorem}

\begin{proof}
Each ledger item is the conclusion of the cited result.
\end{proof}

\paragraph{Parent-version record.}
V37 was formed from
\texttt{Renewal\_Yang\_Mills\_Mass\_Gap\_Research\_Notes\_V36.tex}.
The V36 parent had \(16\,419\) logical source lines,
\(576\,945\) bytes, and SHA--256
\[
 \texttt{8481dc206975d38a45459bc9bb44a7322cc36f6b9f8149b97437be037b0884d0}.
\]
The next compulsory companion audit is V40.

\clearpage
\chapter{Sixth-series V38: bad-density large deviations and
staple-aligned conditional convexity}
\label{chap:s6v38}

\section{An exponential large-deviation bound for the bad density}
\label{sec:s6v38-density}

Continue with the pinned pure-Wilson law and notation of V36.  The mean
estimate there can be upgraded to an exponential bound for a macroscopic
bad density without reintroducing a volume-dependent threshold.

\begin{lemma}[Exponential moment of the total pinned distance]
\label{lem:s6v38-D-mgf}
For \(\beta\ge240\),
\begin{equation}
 \boxed{
 \mathbb E_{\beta,M}
 \exp\left\{{\beta\over240}D\right\}
 \le {\cal C}_0^{\,n_M}.}
 \label{eq:s6v38-D-mgf}
\end{equation}
\end{lemma}

\begin{proof}
The coercive half \(V_M\ge D/120\) gives
\[
 -\beta V_M+{\beta\over240}D
 \le-{\beta\over240}D.
\]
The upper half \(V_M\le4D\) gives
\[
 Z_{\beta,M}\ge I(4\beta)^{n_M}.
\]
Consequently
\[
 \mathbb E_{\beta,M}e^{\beta D/240}
 \le {I(\beta/240)^{n_M}\over I(4\beta)^{n_M}}
 \le
 \left\{
 {C_+(\beta/240)^{-4}\over C_-(4\beta)^{-4}}
 \right\}^{n_M}
 ={\cal C}_0^{n_M}.
\]
\end{proof}

\begin{theorem}[Bad-density large deviations]
\label{thm:s6v38-density-LD}
For the bad-link density
\({\mathfrak b}_{M,\rho}\) of V36 and every \(\eta,\rho>0\),
\begin{equation}
 \boxed{
 \nu_{\beta,M}\{
 {\mathfrak b}_{M,\rho}\ge\eta\}
 \le
 \exp\left\{
 -n_M\left(
 {\beta\eta\rho^2\over240}-\log{\cal C}_0
 \right)\right\}.}
 \label{eq:s6v38-density-LD}
\end{equation}
In particular, for any \(K>240\log{\cal C}_0\),
\begin{equation}
 \nu_{\beta,M}\left\{
 {\mathfrak b}_{M,\rho}\ge{K\over\beta\rho^2}
 \right\}
 \le
 \exp\{-n_M(K/240-\log{\cal C}_0)\}.
 \label{eq:s6v38-density-scale}
\end{equation}
\end{theorem}

\begin{proof}
On
\(\{{\mathfrak b}_{M,\rho}\ge\eta\}\),
\[
 D=\sum_ed_e\ge n_M\eta\rho^2.
\]
Exponential Markov inequality and
\Cref{lem:s6v38-D-mgf} give
\[
 \nu\{D\ge n_M\eta\rho^2\}
 \le
 e^{-\beta n_M\eta\rho^2/240}
 \mathbb Ee^{\beta D/240},
\]
which is \eqref{eq:s6v38-density-LD}.  Substitute
\(\eta=K/(\beta\rho^2)\) for
\eqref{eq:s6v38-density-scale}.
\end{proof}

\begin{corollary}[High-probability sparse chart failures]
\label{cor:s6v38-sparse}
If \(\rho_\beta=\beta^{-\alpha}\) with \(0<\alpha<1/2\), then for
every \(K>240\log{\cal C}_0\),
\[
 {\mathfrak b}_{M,\rho_\beta}
 \le K\beta^{2\alpha-1}
\]
outside an event of probability at most \(e^{-c_Kn_M}\), uniformly in
the growth of \(M\).
\end{corollary}

\begin{proof}
This is \eqref{eq:s6v38-density-scale}, since
\((\beta\rho_\beta^2)^{-1}=\beta^{2\alpha-1}\).
\end{proof}

\section{Why global coercivity does not give a conditional polymer cost}
\label{sec:s6v38-conditional-no}

\begin{proposition}[Counterexample to conditional coercivity from a global
quadratic comparison]
\label{prop:s6v38-global-not-conditional}
For every \(0<\epsilon<1\), the function
\[
 V_\epsilon(x,y)=(x-y)^2+\epsilon(x^2+y^2),
 \qquad (x,y)\in[-1,1]^2,
\]
obeys
\begin{equation}
 \epsilon(x^2+y^2)
 \le V_\epsilon(x,y)
 \le(2+\epsilon)(x^2+y^2).
 \label{eq:s6v38-global-comparison}
\end{equation}
Nevertheless, with exterior value \(y=1\), the conditional minimizer is
\[
 x_\epsilon={1\over1+\epsilon},
\]
which stays far from the global vacuum \(x=0\).  In particular there is no
\(c>0\) for which
\[
 V_\epsilon(x,1)-V_\epsilon(0,1)\ge cx^2
\]
holds for all \(x\in[-1,1]\).
\end{proposition}

\begin{proof}
The lower inequality follows by discarding \((x-y)^2\), and
\((x-y)^2\le2(x^2+y^2)\) gives the upper inequality.  Differentiating the
conditional function gives
\[
 \partial_xV_\epsilon(x,1)=2(x-1)+2\epsilon x,
\]
whose unique zero is \(x_\epsilon=(1+\epsilon)^{-1}\).  Moreover
\[
 V_\epsilon(x_\epsilon,1)<V_\epsilon(0,1),
\]
so the proposed conditional lower bound fails at \(x_\epsilon\).
\end{proof}

\begin{assessment}[The V28 frustrated star is the gauge version]
\label{ass:s6v38-frustration}
The proposition explains why
\(D/120\le V_M\le4D\) cannot by itself produce a Koteck\'y--Preiss
activity for a specified link polymer after its exterior is frozen.  In the
Wilson fibre, the six exterior staples can prefer incompatible link values;
V28 exhibited exactly this obstruction to universal one-link action
descent.  The correct local good event must therefore concern the
\emph{relative alignment of the staples}, not the absolute distance of
the integrated link from the identity.
\end{assessment}

\section{A gauge-covariant staple-alignment lemma}
\label{sec:s6v38-staples}

Every Wilson plaquette incident to one link \(U\) can, after cyclic
reordering and using the real part for the opposite orientation, be written
as
\[
 1-\frac13\operatorname{Re}\Tr(UA_i),
\qquad A_i\in\SU(3).
\]
There are \(k=6\) such terms for an interior four-dimensional link.
For arbitrary \(k\), put
\[
 \Phi_A(U)
 =\sum_{i=1}^k
 \left(1-\frac13\operatorname{Re}\Tr(UA_i)\right).
\]

\begin{definition}[Staple-alignment event]
\label{def:s6v38-SAE}
The exterior staples are \((\epsilon,U_0)\)-aligned if
\begin{equation}
 \max_{1\le i\le k}\|U_0A_i-I\|_{\rm F}\le\epsilon
 \label{eq:s6v38-alignment}
\end{equation}
for some \(U_0\in\SU(3)\).
\end{definition}

\begin{theorem}[Aligned staples give a uniform conditional Hessian]
\label{thm:s6v38-staple-Hessian}
Assume \eqref{eq:s6v38-alignment}.  Let
\[
 U=e^XU_0,\qquad \|X\|_{\rm F}\le r,
\]
inside a geodesically convex exponential ball, and set
\[
 \delta_{\epsilon,r}=\epsilon+e^r-1.
\]
For every anti-Hermitian traceless tangent \(H\),
\begin{equation}
 \boxed{
 {d^2\over dt^2}\Phi_A(e^{tH}U)\bigg|_{t=0}
 \ge {k\over3}(1-\delta_{\epsilon,r})
       \|H\|_{\rm F}^2.}
 \label{eq:s6v38-staple-Hessian}
\end{equation}
Also
\begin{equation}
 \|\nabla\Phi_A(U_0)\|_{\rm F}\le{k\epsilon\over3}.
 \label{eq:s6v38-staple-gradient}
\end{equation}
Thus if \(\delta_{\epsilon,r}<1\), the conditional Wilson action
\(\beta\Phi_A\) is strongly convex on this ball with constant
\[
 {\beta k\over3}(1-\delta_{\epsilon,r}).
\]
\end{theorem}

\begin{proof}
Let \(B_i=UA_i=e^XU_0A_i\).  The elementary exponential bound gives
\[
 \|B_i-I\|_{\rm F}
 \le\|e^X-I\|_{\rm F}+\|U_0A_i-I\|_{\rm F}
 \le e^r-1+\epsilon.
\]
Along the left geodesic \(e^{tH}U\),
\[
 {d^2\over dt^2}
 \left(1-\frac13\operatorname{Re}\Tr(e^{tH}B_i)\right)_{t=0}
 =-\frac13\operatorname{Re}\Tr(H^2B_i).
\]
Because \(H^*=-H\),
\[
 -\operatorname{Re}\Tr H^2=\|H\|_{\rm F}^2.
\]
Moreover,
\[
 |\Tr\{H^2(B_i-I)\}|
 \le\|H^2\|_{\rm F}\|B_i-I\|_{\rm F}
 \le\|H\|_{\rm F}^2\delta_{\epsilon,r}.
\]
Sum over \(i\) to obtain
\eqref{eq:s6v38-staple-Hessian}.  At \(U_0\), the first derivative in a
unit direction \(H\) is
\[
 -{1\over3}\sum_i
 \operatorname{Re}\Tr\{H(U_0A_i-I)\},
\]
because \(\operatorname{Re}\Tr H=0\).  Cauchy--Schwarz bounds its
absolute value by \(k\epsilon/3\), proving
\eqref{eq:s6v38-staple-gradient}.
\end{proof}

\begin{corollary}[Existence and uniqueness of the local conditional vacuum]
\label{cor:s6v38-local-vacuum}
Under the hypotheses of
\Cref{thm:s6v38-staple-Hessian}, suppose
\(\delta_{\epsilon,r}<1\) and
\begin{equation}
 r(1-\delta_{\epsilon,r})>\epsilon.
 \label{eq:s6v38-interior-condition}
\end{equation}
Then \(\Phi_A\) has a unique critical point \(U_*\) in the radius-\(r\)
ball about \(U_0\), it is the unique minimizer there, and
\begin{equation}
 \operatorname{dist}(U_*,U_0)
 \le{\epsilon\over1-\delta_{\epsilon,r}}.
 \label{eq:s6v38-minimizer-shift}
\end{equation}
\end{corollary}

\begin{proof}
Let
\(\lambda_A=(k/3)(1-\delta_{\epsilon,r})\) and
\(g_0=k\epsilon/3\).  Strong geodesic convexity gives, along every radial
unit geodesic of length \(s\),
\[
 {d\over ds}\Phi_A
 \ge-g_0+\lambda_As.
\]
At the boundary \(s=r\), this is positive by
\eqref{eq:s6v38-interior-condition}.  Hence a minimizer on the closed
ball cannot lie on its boundary and is an interior critical point.
Strong convexity makes it unique.  If \(s_*=\operatorname{dist}(U_*,U_0)\),
strong monotonicity of the gradient gives
\(\lambda_As_*\le g_0\), which is
\eqref{eq:s6v38-minimizer-shift}.
\end{proof}

\begin{corollary}[An exact local good/bad fork for one Wilson link]
\label{cor:s6v38-link-fork}
Fix \(\epsilon,r\) satisfying
\eqref{eq:s6v38-interior-condition}.  For every exterior configuration,
exactly one of the following occurs:
\begin{enumerate}[label=\textup{(\roman*)},leftmargin=3.5em]
\item the six staples are \((\epsilon,U_0)\)-aligned for some \(U_0\);
      the one-link conditional has the strongly convex local vacuum of
      \Cref{cor:s6v38-local-vacuum};
\item no such \(U_0\) exists; the exterior belongs to the
      staple-misalignment event.
\end{enumerate}
The second event, not the absolute event
\(\|U-I\|\ge\rho\), is the local polymer type which must be controlled to
prove CPTC.
\end{corollary}

\begin{proof}
The two cases are a statement and its negation.  The conclusion in the
first case is
\Cref{thm:s6v38-staple-Hessian,cor:s6v38-local-vacuum}.
\end{proof}

\section{From plaquette sectors to staple misalignment}
\label{sec:s6v38-sector}

\begin{lemma}[Pairwise criterion for staple alignment]
\label{lem:s6v38-pairwise}
If for some index \(i_0\)
\begin{equation}
 \max_i\|A_{i_0}^{-1}A_i-I\|_{\rm F}\le\epsilon,
 \label{eq:s6v38-pairwise}
\end{equation}
then the staples are \((\epsilon,U_0)\)-aligned with
\(U_0=A_{i_0}^{-1}\).
Conversely, \((\epsilon,U_0)\)-alignment implies
\begin{equation}
 \|A_i^{-1}A_j-I\|_{\rm F}\le2\epsilon
 \label{eq:s6v38-pairwise-converse}
\end{equation}
for all \(i,j\).
\end{lemma}

\begin{proof}
The first assertion is the definition after taking
\(U_0=A_{i_0}^{-1}\).  For the converse, unitary invariance and the triangle
inequality give
\[
 \|A_i^{-1}A_j-I\|_{\rm F}
 =\|(U_0A_i)^{-1}(U_0A_j)-I\|_{\rm F}
 \le
 \|U_0A_j-I\|_{\rm F}
 +\|U_0A_i-I\|_{\rm F}.
\]
\end{proof}

\begin{assessment}[What a bad-polymer estimate must now measure]
\label{ass:s6v38-polymer-target}
Staple misalignment is a finite, gauge-covariant relative holonomy event.
It can be generated by a noncentral plaquette defect, by incompatible
near-centre labels on incident plaquettes, or by a connected centre sheet
passing through the star.  V28 controls unconditional noncentral
multi-plaquette events, while V27 records the centre-sheet branch.  What is
still missing is a same-law connected estimate showing that clusters of
misaligned stars have a summable activity, with the two marked score
supports kept in the connected component.  The absolute bad-link density
of \Cref{thm:s6v38-density-LD} is useful for global capture but does not
replace this relative-staple estimate.
\end{assessment}

\begin{programme}[V39: aligned-star conditional mixing]
\label{prog:s6v38-V39}
\begin{enumerate}[leftmargin=3em]
\item Use the strong convexity in
      \eqref{eq:s6v38-staple-Hessian} to bound the one-link conditional
      covariance and its Lipschitz response to aligned staple data.
\item Express failure of pairwise staple alignment through a fixed set of
      short relative Wilson loops around the star.
\item Combine the V28 noncentral chessboard estimate and the V27
      centre-sheet split to seek a connected misaligned-star activity.
\item Test whether the resulting activity controls the marked cumulant in
      CPTC rather than merely the unmarked bad probability.
\item If the centre-sheet branch remains unsummable, isolate it as a
      separate \(\mathbb Z_3\) contour gas and do not charge it to the
      Gaussian detail field.
\end{enumerate}
\end{programme}

\begin{firewall}[Claim boundary after sixth-series V38]
\label{fire:s6v38-boundary}
V38 proves an exponential large-deviation bound for the density of
absolute bad links, proves that global quadratic action coercivity does not
imply conditional coercivity, proves a gauge-covariant staple-alignment
criterion with an explicit uniform conditional Hessian and local-vacuum
shift, and identifies staple misalignment as the correct local bad-polymer
event for the nonlinear covariance problem.

V38 does not prove a connected misaligned-star activity, CPTC, LCKC, the
selected root asymptotics, the generated-action extension, TSC, YMGNC,
continuum existence, or the full Yang--Mills mass gap.
\end{firewall}

\begin{theorem}[Sixth-series V38 result ledger]
\label{thm:s6v38-results}
The claims in \Cref{fire:s6v38-boundary} follow from
\Cref{lem:s6v38-D-mgf,
thm:s6v38-density-LD,
cor:s6v38-sparse,
prop:s6v38-global-not-conditional,
thm:s6v38-staple-Hessian,
cor:s6v38-local-vacuum,
cor:s6v38-link-fork,
lem:s6v38-pairwise}.
\end{theorem}

\begin{proof}
Each ledger item is the conclusion of the cited result.
\end{proof}

\paragraph{Parent-version record.}
V38 was formed from
\texttt{Renewal\_Yang\_Mills\_Mass\_Gap\_Research\_Notes\_V37.tex}.
The V37 parent had \(16\,816\) logical source lines,
\(590\,146\) bytes, and SHA--256
\[
 \texttt{6b2a98bef53bcad8b6eca6d27e2131fa11a3056af7b80261512d298aad9e93f6}.
\]
The next compulsory companion audit is V40.

\clearpage
\chapter{Sixth-series V39: aligned-star conditional concentration
and the coherent-transmission obstruction}
\label{chap:s6v39}

\section{A global conditional tail from local staple alignment}
\label{sec:s6v39-tail}

Use the bi-invariant Riemannian distance on \(\SU(3)\) and define
\[
 \varphi(V)=1-\frac13\operatorname{Re}\Tr V.
\]
Its unique zero is \(I\).  For radii below the injectivity radius put
\begin{align}
 h(s)&=\sup_{\operatorname{dist}(V,I)\le s}\varphi(V),
 \label{eq:s6v39-h}\\
 g(r)&=\inf_{\operatorname{dist}(V,I)\ge r}\varphi(V).
 \label{eq:s6v39-g}
\end{align}
Then \(h(s)\downarrow0\) as \(s\downarrow0\), while \(g(r)>0\) for
every fixed \(r>0\).

\begin{lemma}[Aligned-star energy gap outside a ball]
\label{lem:s6v39-energy-gap}
Suppose the \(k\) staples are \((\epsilon,U_0)\)-aligned in the sense of
V38.  For \(0<s<r\), set
\begin{equation}
 \Delta_{s,r,\epsilon}
 =k\left\{
 g(r)-h(s)-{2\epsilon\over\sqrt3}
 \right\}.
 \label{eq:s6v39-Delta}
\end{equation}
If \(\Delta_{s,r,\epsilon}>0\), then
\begin{align}
 \sup_{\operatorname{dist}(U,U_0)\le s}\Phi_A(U)
 &\le k\{h(s)+\epsilon/\sqrt3\},
 \label{eq:s6v39-inner-energy}\\
 \inf_{\operatorname{dist}(U,U_0)\ge r}\Phi_A(U)
 &\ge k\{g(r)-\epsilon/\sqrt3\}.
 \label{eq:s6v39-outer-energy}
\end{align}
\end{lemma}

\begin{proof}
Write \(U=VU_0\) and \(B_i=U_0A_i\), so
\(\|B_i-I\|_{\rm F}\le\epsilon\).  Unitary invariance and
Cauchy--Schwarz give
\[
 \left|
 {1\over3}\operatorname{Re}\Tr\{V(B_i-I)\}
 \right|
 \le {1\over3}\|V\|_{\rm F}\|B_i-I\|_{\rm F}
 \le{\epsilon\over\sqrt3}.
\]
Thus
\[
 \varphi(V)-{\epsilon\over\sqrt3}
 \le
 1-\frac13\operatorname{Re}\Tr(VB_i)
 \le
 \varphi(V)+{\epsilon\over\sqrt3}.
\]
Sum over \(i\), then use the definitions of \(h(s)\) and \(g(r)\).
\end{proof}

Let \(m_s\) denote normalized Haar measure of the radius-\(s\) ball about
the identity.  By invariance it is also the measure of the ball about
\(U_0\).  The full one-link conditional is
\[
 d\mu_{\beta,A}(U)
 ={e^{-\beta\Phi_A(U)}\,dU
   \over\int_{\SU(3)}e^{-\beta\Phi_A(V)}\,dV}.
\]

\begin{theorem}[Uniform conditional chart-complement bound]
\label{thm:s6v39-conditional-tail}
Under the hypotheses of
\Cref{lem:s6v39-energy-gap},
\begin{equation}
 \boxed{
 \mu_{\beta,A}\{
 \operatorname{dist}(U,U_0)\ge r\}
 \le m_s^{-1}e^{-\beta\Delta_{s,r,\epsilon}}.}
 \label{eq:s6v39-conditional-tail}
\end{equation}
The constants are uniform over all exterior staple configurations with the
same alignment parameters.
\end{theorem}

\begin{proof}
The numerator over the outer event is at most
\[
 \exp\{-\beta k(g(r)-\epsilon/\sqrt3)\}
\]
because Haar probability is one.  Restricting the denominator to the
radius-\(s\) ball and using
\eqref{eq:s6v39-inner-energy} gives the lower bound
\[
 m_s\exp\{-\beta k(h(s)+\epsilon/\sqrt3)\}.
\]
Their ratio is \eqref{eq:s6v39-conditional-tail}.
\end{proof}

\section{A one-link Poincar\'e and response card}
\label{sec:s6v39-Poincare}

Fix \(r\) satisfying the V38 strong-convexity conditions and smaller than
the convexity radius, so the radius-\(r\) Riemannian ball is strongly
geodesically convex and has convex boundary.
Put
\begin{equation}
 \lambda_{\epsilon,r}
 ={k\over3}(1-\delta_{\epsilon,r}),
 \qquad
 m_{\beta,\epsilon,r}
 =\beta\lambda_{\epsilon,r}.
 \label{eq:s6v39-m}
\end{equation}

\begin{theorem}[Aligned-star conditional Poincar\'e inequality]
\label{thm:s6v39-Poincare}
Assume \(m_{\beta,\epsilon,r}>0\), and let
\(\mu_{\beta,A}^{(r)}\) be the one-link conditional restricted and
renormalized to the radius-\(r\) ball about \(U_0\).  Every smooth real
\(f\) satisfies
\begin{equation}
 \boxed{
 \operatorname{Var}_{\mu_{\beta,A}^{(r)}}(f)
 \le {1\over m_{\beta,\epsilon,r}}
 \mathbb E_{\mu_{\beta,A}^{(r)}}\|\nabla f\|^2.}
 \label{eq:s6v39-Poincare}
\end{equation}
For the full conditional,
\begin{equation}
 \boxed{
 \operatorname{Var}_{\mu_{\beta,A}}(f)
 \le {\|\nabla f\|_\infty^2\over m_{\beta,\epsilon,r}}
 +2m_s^{-1}e^{-\beta\Delta_{s,r,\epsilon}}
   \operatorname{osc}(f)^2.}
 \label{eq:s6v39-full-variance}
\end{equation}
\end{theorem}

\begin{proof}
Equip \(\SU(3)\) with its bi-invariant metric, whose Ricci tensor is
nonnegative.  \Cref{thm:s6v38-staple-Hessian} gives
\(\operatorname{Hess}(\beta\Phi_A)\succeq
m_{\beta,\epsilon,r}I\) on the ball.  The Bakry--\'Emery
Poincar\'e inequality with convex Neumann boundary therefore proves
\eqref{eq:s6v39-Poincare}.

Let \(p\) be the full conditional mass outside the ball.  Mixture variance
decomposition between the ball and its complement gives
\[
 \operatorname{Var}_{\mu_{\beta,A}}(f)
 \le
 \operatorname{Var}_{\mu_{\beta,A}^{(r)}}(f)
 +2p\,\operatorname{osc}(f)^2.
\]
Use \eqref{eq:s6v39-Poincare} and
\Cref{thm:s6v39-conditional-tail}.
\end{proof}

\begin{theorem}[Lipschitz response to an aligned staple path]
\label{thm:s6v39-response}
Let \(A_i(t)\), \(0\le t\le1\), be \(C^1\) staple paths which remain
\((\epsilon,U_0)\)-aligned with the same centre and parameters.  Let
\(\mu_t^{(r)}\) be their restricted one-link conditionals.  For every
smooth \(f\) independent of \(t\),
\begin{equation}
 \boxed{
 \left|
 {d\over dt}\mathbb E_{\mu_t^{(r)}}f
 \right|
 \le
 {\beta\|\nabla f\|_\infty
  \over3m_{\beta,\epsilon,r}}
 \sum_{i=1}^k\|\dot A_i(t)\|_{\rm F}.}
 \label{eq:s6v39-response}
\end{equation}
Consequently the full conditional expectations at the endpoints obey
\begin{align}
 |\mathbb E_{\mu_{\beta,A(1)}}f
  -\mathbb E_{\mu_{\beta,A(0)}}f|
 \le{}&
 {\beta\|\nabla f\|_\infty
  \over3m_{\beta,\epsilon,r}}
 \sum_i\int_0^1\|\dot A_i(t)\|_{\rm F}\,dt
 \notag\\
 &+2\operatorname{osc}(f)
 \sum_{t\in\{0,1\}}
 m_s^{-1}e^{-\beta\Delta_{s,r,\epsilon}}.
 \label{eq:s6v39-full-response}
\end{align}
\end{theorem}

\begin{proof}
Differentiate the normalized restricted expectation:
\[
 {d\over dt}\mathbb E_{\mu_t^{(r)}}f
 =-\operatorname{Cov}_{\mu_t^{(r)}}
  (f,\beta\dot\Phi_t).
\]
The covariance form of Brascamp--Lieb and
\eqref{eq:s6v39-m} gives
\[
 |\operatorname{Cov}(f,\beta\dot\Phi_t)|
 \le {1\over m_{\beta,\epsilon,r}}
 \|\nabla f\|_\infty
 \|\nabla(\beta\dot\Phi_t)\|_\infty.
\]
For a unit tangent \(H\),
\[
 |D_U\dot\Phi_t[H]|
 =\left|
 -{1\over3}\sum_i
 \operatorname{Re}\Tr(HU\dot A_i)
 \right|
 \le{1\over3}\sum_i\|\dot A_i\|_{\rm F}.
\]
This proves \eqref{eq:s6v39-response}.  Integrate it.  At either endpoint,
replacing the full expectation by the restricted one changes it by at most
\(2\operatorname{osc}(f)\) times the outer probability.  Apply
\Cref{thm:s6v39-conditional-tail} at both endpoints.
\end{proof}

\begin{assessment}[Sharp scale of the aligned conditional]
\label{ass:s6v39-scale}
For fixed alignment parameters,
\[
 m_{\beta,\epsilon,r}
 =\beta\lambda_{\epsilon,r}.
\]
Thus local fluctuations have variance \(O(\beta^{-1})\), while the
response coefficient in \eqref{eq:s6v39-response} tends to
\[
 {1\over3\lambda_{\epsilon,r}}
 ={1\over k(1-\delta_{\epsilon,r})}.
\]
Concentration becomes stronger with \(\beta\), but deterministic movement
of the conditional vacuum under coherent staple motion does not acquire an
extra factor \(\beta^{-1}\).
\end{assessment}

\section{Coherent rotations forbid weak-coupling one-link contraction}
\label{sec:s6v39-coherent}

\begin{theorem}[Exact translation covariance of a coherent staple star]
\label{thm:s6v39-translation}
For \(G\in\SU(3)\), take
\[
 A_i^{(G)}=G^{-1}\qquad(1\le i\le k).
\]
Then
\[
 \Phi_{A^{(G)}}(U)=k\varphi(UG^{-1}),
\]
and the corresponding conditional law is the left translate by \(G\) of
the law with all staples equal to \(I\):
\begin{equation}
 \boxed{
 \mu_{\beta,A^{(G)}}=(L_G)_*\mu_{\beta,A^{(I)}}.}
 \label{eq:s6v39-translation}
\end{equation}
Moreover
\begin{equation}
 \mu_{\beta,A^{(G)}}\Longrightarrow\delta_G
 \qquad(\beta\to\infty).
 \label{eq:s6v39-zero-temperature}
\end{equation}
\end{theorem}

\begin{proof}
The action identity is immediate.  In the integral defining the
conditional, substitute \(U=GV\) and use left Haar invariance; this proves
\eqref{eq:s6v39-translation}.  The continuous function
\(\varphi\) is nonnegative and has the unique zero \(I\).  For every
neighbourhood of \(I\), compactness gives a strictly positive action gap
on its complement, while the neighbourhood has positive Haar measure.
The same numerator/denominator comparison as in
\Cref{thm:s6v39-conditional-tail} shows concentration at \(I\).
Translation gives \eqref{eq:s6v39-zero-temperature}.
\end{proof}

\begin{corollary}[No regulator-uniform strict contraction on the unreduced
coherent mode]
\label{cor:s6v39-no-one-link-contraction}
Let \(G\ne I\).  For every continuous \(1\)-Lipschitz function \(f\),
\[
 \mathbb E_{\mu_{\beta,A^{(G)}}}f
 -\mathbb E_{\mu_{\beta,A^{(I)}}}f
 \longrightarrow f(G)-f(I).
\]
Taking
\[
 f(U)=\operatorname{dist}(U,I)
\]
gives the limiting difference
\(\operatorname{dist}(G,I)\).  Hence no one-link influence estimate which
measures coherent staple displacement by
\(\operatorname{dist}(G,I)\) can have a regulator-uniform contraction
coefficient strictly below one as \(\beta\to\infty\).
\end{corollary}

\begin{proof}
Weak convergence in
\eqref{eq:s6v39-zero-temperature} gives the first limit.  The distance
function is continuous and \(1\)-Lipschitz, with values \(0\) and
\(\operatorname{dist}(G,I)\) at \(I\) and \(G\).  A strict influence
coefficient \(c<1\) would bound the left side by
\(c\operatorname{dist}(G,I)\) for every \(\beta\); passing to the limit
contradicts the displayed equality.
\end{proof}

\begin{assessment}[The contraction must live on the physical quotient]
\label{ass:s6v39-quotient}
The coherent mode in
\Cref{thm:s6v39-translation} is a transported local gauge orientation.
The conditional is sharply concentrated, but its centre follows the
exterior exactly.  Therefore a weak-coupling Dobrushin row on unreduced
links is the wrong target.  A strict loss can arise only after a parity
tree, Hodge, or relative-boundary quotient removes coherent transport, or
from a block observable insensitive to it.  This is compatible with the
physical CDAC formulation, but it prevents the aligned one-link Poincar\'e
constant from being advertised as CDAC.
\end{assessment}

\section{The exact V40 entrance}
\label{sec:s6v39-entrance}

\begin{definition}[Aligned quotient influence card, AQIC]
\label{def:s6v39-AQIC}
For one physical block, AQIC consists of:
\begin{enumerate}[label=\textup{(\roman*)},leftmargin=3.5em]
\item a declared quotient map which annihilates the common coherent staple
      rotation of \Cref{thm:s6v39-translation};
\item the aligned-star tail and Poincar\'e bounds
      \eqref{eq:s6v39-conditional-tail} and
      \eqref{eq:s6v39-full-variance} on every good star;
\item a Lipschitz estimate for the quotient observable under noncoherent
      staple perturbations, with a strict summed row;
\item a connected contour activity for staple-misaligned stars, including
      the \(\mathbb Z_3\) sheet branch; and
\item a physical-clock rate satisfying YMGNC rather than a fixed
      contraction on a shrinking clock.
\end{enumerate}
\end{definition}

\begin{programme}[V40: compulsory audit and quotient incidence]
\label{prog:s6v39-V40}
\begin{enumerate}[leftmargin=3em]
\item Perform the compulsory audit of the then-current SM and NS notes.
\item Write the exact differential of the parity-tree or Hodge quotient
      with respect to the six staple inputs.
\item Prove that the common coherent direction lies in its kernel.
\item Compute the remaining incidence row and test whether
      \eqref{eq:s6v39-response} becomes strictly contractive on the
      quotient.
\item Keep the misaligned-star contour activity as a separate unresolved
      row if the audit supplies no estimate for it.
\end{enumerate}
\end{programme}

\begin{firewall}[Claim boundary after sixth-series V39]
\label{fire:s6v39-boundary}
V39 proves a uniform exponential complement bound for every aligned
one-link conditional, a restricted and full conditional Poincar\'e
inequality, and a Lipschitz response estimate along aligned staple paths.
It also proves an exact coherent-rotation covariance and shows that
weak-coupling one-link influence cannot contract uniformly on the
unreduced coherent mode.  It therefore locates the next contraction target
on the physical block quotient.

V39 does not prove AQIC, a strict quotient influence row, a connected
misaligned-star activity, CPTC, LCKC, response-root asymptotics, TSC,
YMGNC, continuum existence, or the full Yang--Mills mass gap.
\end{firewall}

\begin{theorem}[Sixth-series V39 result ledger]
\label{thm:s6v39-results}
The claims in \Cref{fire:s6v39-boundary} follow from
\Cref{lem:s6v39-energy-gap,
thm:s6v39-conditional-tail,
thm:s6v39-Poincare,
thm:s6v39-response,
thm:s6v39-translation,
cor:s6v39-no-one-link-contraction}.
\end{theorem}

\begin{proof}
Each ledger item is the conclusion of the cited result.
\end{proof}

\paragraph{Parent-version record.}
V39 was formed from
\texttt{Renewal\_Yang\_Mills\_Mass\_Gap\_Research\_Notes\_V38.tex}.
The V38 parent had \(17\,238\) logical source lines,
\(603\,035\) bytes, and SHA--256
\[
 \texttt{bc8b7bdb204d4aeb43c03b9fce2531d2c5b6a59bb844cdba8b1cd6c21f980d71}.
\]
The next compulsory companion audit is V40.

\clearpage
\chapter{Sixth-series V40: compulsory audit and the exact
aligned-star quotient incidence}
\label{chap:s6v40}

\section{Compulsory audit of the current companion programmes}
\label{sec:s6v40-audit}

The read-only audit was frozen at the following complete files:
\begin{align}
 &\texttt{Renewal\_SM\_Einstein\_Continuum\_Research\_Notes\_V66.tex},
 \notag\\
 &\qquad 1\,055\,461\ \hbox{bytes},\quad
 26\,332\ \hbox{logical lines},
 \notag\\
 &\qquad
 \texttt{fe6523301385afce2e05d7bf8c30d2f918b74d816a9ac4966d649d3550c61957},
 \label{eq:s6v40-SM-checkpoint}\\
 &\texttt{Renewal\_NS\_Stability\_Research\_Notes\_V26.tex},
 \notag\\
 &\qquad 278\,283\ \hbox{bytes},\quad
 5\,320\ \hbox{logical lines},
 \notag\\
 &\qquad
 \texttt{82c887f8c2c69e4f28fad7c3dc8de5b9099cb5a4bf431eba1fff3e91935978ad}.
 \label{eq:s6v40-NS-checkpoint}
\end{align}

SM V65--V66 separates pointwise positivity and reversibility from positive
transfer spectrum, proves preservation of a positive transfer seed under
an exact conditional extension, constructs the finite free transverse
traceless Mehler seed, and finds that its coupled sourced constraint has a
volume-growing inverse Laplacian and needs connected neutral-cluster
control.  Wilson lattice gauge theory already has the required
reflection-positive seed, so no gravitational theorem is imported.  The
audit reinforces two proof obligations already present here: the quotient
must be an actual orthogonal or reduced-coordinate construction, and an
extensive finite-volume bound must be replaced by a connected local norm.

NS V26 computes a rank-one restriction of two kernel derivative norms.
The first norm is unchanged because its extremizer is already rank one;
the second improves by at most eleven percent pointwise.  No
Navier--Stokes estimate is imported.  The relevant warning is that
projecting to a physically smaller input class does not by itself make an
influence norm strictly smaller; the quotient incidence must be evaluated.

\begin{theorem}[V40 companion-audit decision]
\label{thm:s6v40-audit}
The companion audit supplies no missing Yang--Mills probability or
correlation estimate.  It confirms that the next noncircular calculation
is the differential of the aligned conditional after the coherent staple
direction has been removed.  That calculation is performed below rather
than inferred from the existence of a quotient.
\end{theorem}

\begin{proof}
The SM statements concern a gravitational TT transfer system and the NS
statements concern Oseen-kernel tensor norms.  Neither bounds a Wilson
staple conditional.  Their methodological conclusions identify an
interface but provide no numerical row for it.
\end{proof}

\section{Linear response of the aligned local vacuum}
\label{sec:s6v40-minimizer}

Work first at the perfectly aligned star, translated so that
\(U_0=I\) and \(A_i=I\).  Let
\[
 A_i(s)=e^{sX_i},
 \qquad X_i\in\mathfrak{su}(3),
 \qquad
 \overline X={1\over k}\sum_{i=1}^kX_i.
\]
For small \(s\), let \(U_*(s)=e^{Y(s)}\) be the unique local conditional
vacuum supplied by V38.

\begin{theorem}[Exact differential of the conditional vacuum]
\label{thm:s6v40-vacuum-differential}
At the perfectly aligned star,
\begin{equation}
 \boxed{
 Y'(0)=-\overline X.}
 \label{eq:s6v40-vacuum-differential}
\end{equation}
Thus the coherent input subspace
\[
 {\cal C}=\{(X,\ldots,X):X\in\mathfrak{su}(3)\}
\]
translates the link vacuum isometrically, while the transverse subspace
\[
 {\cal C}^{\perp}
 =\{(X_1,\ldots,X_k):\sum_iX_i=0\}
\]
does not move it to first order.
\end{theorem}

\begin{proof}
For \(Y,X_i\) small,
\[
 1-\frac13\operatorname{Re}\Tr(e^Ye^{sX_i})
 ={1\over6}\|Y+sX_i\|_{\rm F}^2
 +O\bigl((\|Y\|+|s|\|X_i\|)^3\bigr).
\]
Consequently the derivative at the origin of the gradient equation
\(D_Y\Phi_A(e^Y)=0\) is
\[
 {k\over3}Y'(0)+{1\over3}\sum_iX_i=0.
\]
The vertical derivative \(kI/3\) is invertible, so the implicit-function
theorem applies and gives
\eqref{eq:s6v40-vacuum-differential}.  The two subspace statements are
immediate.
\end{proof}

\section{Recentering kills the complete first variation}
\label{sec:s6v40-recenter}

Define the recentered link variable by
\[
 U=e^{-s\overline X}V.
\]
At perfect alignment the recentered action is
\[
 \widehat\Phi_s(V)
 =\sum_{i=1}^k
 \left[
 1-\frac13\operatorname{Re}\Tr
 \{e^{-s\overline X}Ve^{sX_i}\}
 \right].
\]

\begin{theorem}[Perfect-star quotient cancellation]
\label{thm:s6v40-perfect-cancellation}
For every \(V\in\SU(3)\),
\begin{equation}
 \boxed{
 {d\over ds}\widehat\Phi_s(V)\bigg|_{s=0}=0.}
 \label{eq:s6v40-perfect-cancellation}
\end{equation}
Hence the derivative at \(s=0\) of the entire normalized recentered
one-link conditional law is zero, not merely the derivative of its
minimizer.
\end{theorem}

\begin{proof}
Differentiate and use cyclicity of the trace:
\begin{align*}
 {d\over ds}\widehat\Phi_s(V)\bigg|_{s=0}
 &=-{1\over3}\operatorname{Re}
   \sum_i\Tr(-\overline XV+VX_i)\\
 &=-{1\over3}\operatorname{Re}
   \Tr\left\{V\left(\sum_iX_i-k\overline X\right)\right\}=0.
\end{align*}
The Haar measure is left invariant, so the recentering has unit Jacobian.
The unnormalized density and its normalizer both have zero derivative.
\end{proof}

\begin{definition}[Coherent quotient coordinates]
\label{def:s6v40-quotient}
The linear input quotient is
\[
 {\cal Q}_{\rm st}(X_1,\ldots,X_k)
 =(X_1-\overline X,\ldots,X_k-\overline X).
\]
The output quotient is the recentered variable
\[
 V=e^{s\overline X}U.
\]
Together they annihilate the coherent gauge-transport direction at the
level of the complete conditional law by
\Cref{thm:s6v40-perfect-cancellation}.
\end{definition}

\section{Near alignment: the quotient influence is proportional to
misalignment}
\label{sec:s6v40-near}

Translate a general aligned star so that \(U_0=I\), and write
\[
 B_i=A_i,\qquad \|B_i-I\|_{\rm F}\le\epsilon.
\]
Perturb it by
\[
 A_i(s)=B_i e^{sX_i}
\]
and retain the recentering \(U=e^{-s\overline X}V\).  Denote the
derivative at zero of the recentered action by
\(\dot{\widehat\Phi}\).

\begin{lemma}[Pointwise quotient insertion bounds]
\label{lem:s6v40-insertion}
With
\[
 S_X=\sum_i\|X_i\|_{\rm F},
\]
one has
\begin{equation}
 \boxed{
 \|\nabla_V\dot{\widehat\Phi}\|_\infty
 \le {2\epsilon\over3}S_X,
 \qquad
 \operatorname{osc}(\dot{\widehat\Phi})
 \le {4\epsilon\over3}S_X.}
 \label{eq:s6v40-insertion}
\end{equation}
\end{lemma}

\begin{proof}
At \(s=0\),
\[
 \dot{\widehat\Phi}(V)
 =-{1\over3}\operatorname{Re}\sum_i
 \Tr(-\overline XVB_i+VB_iX_i).
\]
The same expression with every \(B_i=I\) vanishes by
\(\sum_iX_i=k\overline X\).  Subtract it.  In a unit tangent \(H\), the
derivative is a sum of
\[
 \Tr\{-\overline XHV(B_i-I)+HV(B_i-I)X_i\}.
\]
Cauchy--Schwarz, unitarity, and
\[
 k\|\overline X\|_{\rm F}\le S_X
\]
give the first bound in
\eqref{eq:s6v40-insertion}.  The same estimate without \(H\) bounds
\(\|\dot{\widehat\Phi}\|_\infty\) by
\(2\epsilon S_X/3\); twice this bound controls the oscillation.
\end{proof}

\begin{theorem}[Full conditional quotient response]
\label{thm:s6v40-full-quotient-response}
Choose \(s<r\) and alignment parameters so that the V39 gap
\(\Delta=\Delta_{s,r,\epsilon}\) and convexity
\(\lambda=\lambda_{\epsilon,r}\) are positive.  Let \(f\) be a smooth
function of the recentered link.  At the aligned star,
\begin{align}
 \left|
 {d\over ds}\mathbb E_{\mu_{\beta,A(s)}}f(e^{s\overline X}U)
 \right|_{s=0}
 \le{}&
 {2\epsilon\over3\lambda}
 \|\nabla f\|_\infty S_X
 \notag\\
 &+
 C_{\epsilon,r,s}\,
 \epsilon\beta e^{-\beta\Delta/2}
 \{\|\nabla f\|_\infty+\operatorname{osc}(f)\}S_X.
 \label{eq:s6v40-full-quotient-response}
\end{align}
At perfect alignment \(\epsilon=0\), the left side is exactly zero.
\end{theorem}

\begin{proof}
Differentiation of the normalized recentered expectation gives
\[
 -\beta\operatorname{Cov}
 (f,\dot{\widehat\Phi}).
\]
Let
\[
 p\le m_s^{-1}e^{-\beta\Delta}
\]
be the conditional mass outside the radius-\(r\) ball.  The full variance
bound of V39, applied once to \(f\) and once to
\(\dot{\widehat\Phi}\), gives
\begin{align*}
 \operatorname{Var}(f)
 &\le{\|\nabla f\|_\infty^2\over\beta\lambda}
     +2p\operatorname{osc}(f)^2,\\
 \operatorname{Var}(\dot{\widehat\Phi})
 &\le{(2\epsilon S_X/3)^2\over\beta\lambda}
     +2p(4\epsilon S_X/3)^2.
\end{align*}
Use Cauchy--Schwarz for the covariance and
\(\sqrt{a+b}\le\sqrt a+\sqrt b\).  The product of the two main terms is
\[
 {2\epsilon\over3\beta\lambda}
 \|\nabla f\|_\infty S_X,
\]
which becomes the first term of
\eqref{eq:s6v40-full-quotient-response} after multiplication by
\(\beta\).  Every remaining product contains \(p^{1/2}\) or \(p\);
for \(\beta\ge1\) it is bounded by the displayed exponential remainder
after enlarging \(C_{\epsilon,r,s}\).  The perfect case also follows
directly from
\Cref{thm:s6v40-perfect-cancellation}.
\end{proof}

\section{The exact six-staple incidence row}
\label{sec:s6v40-incidence}

An interior four-dimensional link has \(k=6\) staples, each an ordered
product of three exterior links.  A unit perturbation of one exterior-link
slot induces a staple tangent of the same Frobenius norm.  Counting with
multiplicity, the full star therefore has
\begin{equation}
 N_{\rm slot}=3k=18
 \label{eq:s6v40-slots}
\end{equation}
incoming slots.

\begin{corollary}[Strict aligned quotient row]
\label{cor:s6v40-strict-row}
For the six-staple star, the summed leading quotient influence from all
exterior-link slots is at most
\begin{equation}
 \boxed{
 \rho_{\rm AQ}^{(0)}
 ={6\epsilon\over1-\delta_{\epsilon,r}},
 \qquad
 \delta_{\epsilon,r}=\epsilon+e^r-1.}
 \label{eq:s6v40-strict-row}
\end{equation}
Including the chart-complement restoration,
\begin{equation}
 \rho_{\rm AQ}(\beta)
 \le
 {6\epsilon\over1-\delta_{\epsilon,r}}
 +C_{\epsilon,r,s}\epsilon\beta e^{-\beta\Delta/2}.
 \label{eq:s6v40-full-row}
\end{equation}
Hence the good aligned quotient row is strictly contractive for all
sufficiently large \(\beta\) whenever
\begin{equation}
 \boxed{
 {6\epsilon\over1-\delta_{\epsilon,r}}<1.}
 \label{eq:s6v40-row-condition}
\end{equation}
\end{corollary}

\begin{proof}
In
\eqref{eq:s6v40-full-quotient-response}, the leading coefficient per
unit of \(S_X\) is
\[
 {2\epsilon\over3\lambda}
 ={2\epsilon\over k(1-\delta_{\epsilon,r})}.
\]
Summing over the \(3k=18\) exterior slots and setting \(k=6\) gives
\eqref{eq:s6v40-strict-row}.  The same incidence factor is absorbed into
the remainder constant in
\eqref{eq:s6v40-full-row}.  Its second term tends to zero exponentially,
so \eqref{eq:s6v40-row-condition} implies strictness eventually.
\end{proof}

\begin{assessment}[What AQIC now contains]
\label{ass:s6v40-AQIC}
V40 populates the first three good-star rows of AQIC:
\begin{enumerate}[label=\textup{(\roman*)},leftmargin=3.5em]
\item the common coherent input and output directions are explicit;
\item the complete conditional law, not only its minimizer, has zero
      perfect-star quotient derivative; and
\item a near-aligned star has the strict row
      \eqref{eq:s6v40-full-row} under
      \eqref{eq:s6v40-row-condition}.
\end{enumerate}
The audit warning from NS V26 is resolved by calculation: here the leading
unquotiented coherent norm is one, while the recentered perfect-star norm
is zero.  The remaining AQIC row is probabilistic and connected---a
summable activity for staple-misaligned stars and centre sheets.
\end{assessment}

\begin{programme}[V41: the relative-loop description of misalignment]
\label{prog:s6v40-V41}
\begin{enumerate}[leftmargin=3em]
\item Express \(A_i^{-1}A_j\) as the holonomy of an explicit short loop in
      the union of the two incident plaquettes.
\item Show that staple misalignment forces at least one such loop outside a
      fixed identity tube.
\item Split that event into a noncentral defect and incompatible
      \(\mathbb Z_3\) centre labels.
\item Apply the V28 same-colour chessboard bound to separated relative-loop
      defects and keep the centre-sheet branch separate.
\item Determine whether the resulting unconditional contour activity can
      be upgraded to the marked connected activity required by CPTC.
\end{enumerate}
\end{programme}

\begin{firewall}[Claim boundary after sixth-series V40]
\label{fire:s6v40-boundary}
V40 performs the compulsory SM/NS audit; computes the exact differential
of the perfectly aligned conditional vacuum; defines coherent quotient
coordinates; proves cancellation of the first variation of the complete
recentered conditional law; proves an \(O(\epsilon)\) full conditional
response bound near alignment; and computes the 18-slot incidence,
obtaining the explicit strict good-star row
\eqref{eq:s6v40-row-condition}.

V40 does not prove a connected activity for misaligned stars or centre
sheets, complete AQIC, CPTC, LCKC, selected root asymptotics, the
generated-action extension, TSC, YMGNC, continuum existence, or the full
Yang--Mills mass gap.
\end{firewall}

\begin{theorem}[Sixth-series V40 result ledger]
\label{thm:s6v40-results}
The claims in \Cref{fire:s6v40-boundary} follow from
\Cref{thm:s6v40-audit,
thm:s6v40-vacuum-differential,
thm:s6v40-perfect-cancellation,
lem:s6v40-insertion,
thm:s6v40-full-quotient-response,
cor:s6v40-strict-row}.
\end{theorem}

\begin{proof}
Each ledger item is the conclusion of the cited result.
\end{proof}

\paragraph{Parent-version record.}
V40 was formed from
\texttt{Renewal\_Yang\_Mills\_Mass\_Gap\_Research\_Notes\_V39.tex}.
The V39 parent had \(17\,642\) logical source lines,
\(616\,228\) bytes, and SHA--256
\[
 \texttt{198f5917253ea95b4ee32bbbd7b3dbc2829bedaff87bab26fcf9b6854b097f56}.
\]
The next compulsory companion audit is V45.

\clearpage
\chapter{Sixth-series V41: relative staple loops and an explicit
misalignment-cluster rate}
\label{chap:s6v41}

\section{The relative holonomy of two staples}
\label{sec:s6v41-relative}

Let \(e\) be an oriented link from \(x\) to \(y\).  Each oriented staple
\(A_i\) is a three-link path from \(y\) back to \(x\), chosen so the
corresponding incident plaquette holonomy based at \(x\) is
\[
 P_i=U_eA_i.
\]
Oppositely oriented plaquettes are put in this form using inversion and
the real trace.

\begin{lemma}[Two-staple relative loop identity]
\label{lem:s6v41-relative-loop}
For every pair of staples,
\begin{equation}
 \boxed{
 A_i^{-1}A_j=P_i^{-1}P_j.}
 \label{eq:s6v41-relative-loop}
\end{equation}
It is the holonomy of a closed path of length at most six contained in the
union of the two incident plaquettes.  Under a gauge transformation it is
conjugated at its base point, so its Frobenius distance from \(I\) is
gauge invariant.
\end{lemma}

\begin{proof}
Since \(A_i=U_e^{-1}P_i\),
\[
 A_i^{-1}A_j
 =P_i^{-1}U_eU_e^{-1}P_j=P_i^{-1}P_j.
\]
The path \(A_i^{-1}A_j\) traverses one three-link staple in reverse and
the other forward.  Gauge factors cancel at intermediate vertices and
leave conjugation at the common endpoint.
\end{proof}

\begin{lemma}[Misalignment forces an identity-bad incident plaquette]
\label{lem:s6v41-misalignment-plaquette}
Let \({\cal M}_e(\epsilon)\) be the event that the \(k=6\) staples at
\(e\) are not \((\epsilon,U_0)\)-aligned for any \(U_0\).  Then
\begin{equation}
 \boxed{
 {\cal M}_e(\epsilon)
 \subseteq
 \bigcup_{p\ni e}
 \{\|U_p-I\|_{\rm F}>\epsilon/2\}.}
 \label{eq:s6v41-misalignment-inclusion}
\end{equation}
\end{lemma}

\begin{proof}
Argue by contraposition.  If all incident plaquette holonomies satisfy
\(\|P_i-I\|_{\rm F}\le\epsilon/2\), then
\[
 \|A_{i_0}^{-1}A_i-I\|_{\rm F}
 =\|P_{i_0}^{-1}P_i-I\|_{\rm F}
 \le\|P_{i_0}-I\|_{\rm F}+\|P_i-I\|_{\rm F}
 \le\epsilon.
\]
The pairwise criterion of V38, with
\(U_0=A_{i_0}^{-1}\), makes all staples
\((\epsilon,U_0)\)-aligned.  This is the complement of
\({\cal M}_e(\epsilon)\).
\end{proof}

\begin{assessment}[No separate centre exception in the pure Wilson star]
\label{ass:s6v41-centre}
A plaquette near a nontrivial element of \(Z_3\) is far from \(I\) and is
therefore included on the right of
\eqref{eq:s6v41-misalignment-inclusion}.  For the bare Wilson measure it
has a large positive action cost.  The centre-sheet split remains useful
for observables defined modulo centre and for response-generated actions,
but it is not needed to make the present pure-Wilson misalignment carrier
exhaustive.
\end{assessment}

\section{Identity-defect chessboard suppression}
\label{sec:s6v41-chessboard}

Define the identity-defect event
\[
 C_p(\delta)=\{\|U_p-I\|_{\rm F}\ge\delta\}.
\]
The exact Wilson defect identity is
\begin{equation}
 d(U)=1-\frac13\operatorname{Re}\Tr U
 ={1\over6}\|U-I\|_{\rm F}^2.
 \label{eq:s6v41-Wilson-distance}
\end{equation}
Unlike the V28 event \(B_p(\delta)\), \(C_p(\delta)\) does not quotient
by the centre.

\begin{theorem}[Same-colour identity-defect suppression]
\label{thm:s6v41-same-colour}
Assume \(\beta\ge r_{\rm H}^{-2}\), and define
\begin{equation}
 Q_\delta^{I}(\beta)
 =c_{\rm H}^{-64}\beta^{256}
  \exp\left\{256-{\beta\delta^2\over6}\right\}.
 \label{eq:s6v41-QI}
\end{equation}
For every finite set \(J\) of plaquettes of one V28 colour,
\begin{equation}
 \boxed{
 \mu_{\Lambda,\beta}
 \left(\bigcap_{p\in J}C_p(\delta)\right)
 \le\{Q_\delta^{I}(\beta)\}^{|J|}.}
 \label{eq:s6v41-same-colour}
\end{equation}
The constants are independent of the even dyadic periodic volume.
\end{theorem}

\begin{proof}
The event \(C_p(\delta)\) is preserved when a reflection changes the
plaquette holonomy by inversion and conjugation.  Thus the V28 block
chessboard inequality applies.  Dissemination through the shifted
\(2^4\)-block tiling forces \(N=V/16\) plaquettes of that colour to
satisfy \(d(U_p)\ge\delta^2/6\) by
\eqref{eq:s6v41-Wilson-distance}.  Hence the disseminated numerator is at
most \(e^{-\beta\delta^2N/6}\).  The V28 small-link partition lower bound,
with \(r=\beta^{-1/2}\), contributes
\[
 h(\beta^{-1/2})^{-E/N}e^{(8/3)\beta(P/N)\beta^{-1}}
 \le c_{\rm H}^{-64}\beta^{256}e^{256}.
\]
Taking the \(N\)-th root gives
\eqref{eq:s6v41-QI}, and the chessboard inequality raises it to
\(|J|\).
\end{proof}

\begin{corollary}[Arbitrary identity-defect family]
\label{cor:s6v41-all-colours}
Put
\begin{equation}
 p_{\rm plaq}^{I}(\beta,\delta)
 =\min\{1,(Q_\delta^{I}(\beta))^{1/96}\}.
 \label{eq:s6v41-pplaq}
\end{equation}
Then every finite set \(J\) of distinct plaquettes obeys
\begin{equation}
 \boxed{
 \mu_{\Lambda,\beta}
 \left(\bigcap_{p\in J}C_p(\delta)\right)
 \le\{p_{\rm plaq}^{I}(\beta,\delta)\}^{|J|}.}
 \label{eq:s6v41-all-colours}
\end{equation}
\end{corollary}

\begin{proof}
If \(Q_\delta^I\ge1\), the bound is trivial.  Otherwise one of the
ninety-six colour classes contains at least \(|J|/96\) members.  The
intersection over all of \(J\) is contained in the intersection over that
class.  Apply \Cref{thm:s6v41-same-colour}.
\end{proof}

\section{Multi-star suppression and the contour tail}
\label{sec:s6v41-stars}

\begin{theorem}[Misaligned-star multi-defect bound]
\label{thm:s6v41-star-MDSC}
Define
\begin{equation}
 p_{\rm star}(\beta,\epsilon)
 =\min\left\{1,\,
 6\{p_{\rm plaq}^{I}(\beta,\epsilon/2)\}^{1/4}\right\}.
 \label{eq:s6v41-pstar}
\end{equation}
For every finite set \(K\) of distinct links,
\begin{equation}
 \boxed{
 \mu_{\Lambda,\beta}
 \left(\bigcap_{e\in K}{\cal M}_e(\epsilon)\right)
 \le p_{\rm star}(\beta,\epsilon)^{|K|}.}
 \label{eq:s6v41-star-MDSC}
\end{equation}
\end{theorem}

\begin{proof}
By \Cref{lem:s6v41-misalignment-plaquette}, expand every
\({\cal M}_e\) into a union over its at most six incident plaquettes.
There are at most \(6^{|K|}\) assignments.  Since one plaquette contains
four links, every assignment uses at least \(|K|/4\) distinct
plaquettes.  Apply
\Cref{cor:s6v41-all-colours} to each assigned intersection and sum the
union bound:
\[
 \mu\left(\bigcap_{e\in K}{\cal M}_e\right)
 \le
 \left[
 6\{p_{\rm plaq}^{I}(\beta,\epsilon/2)\}^{1/4}
 \right]^{|K|}.
\]
The minimum with one gives the stated formula.
\end{proof}

\begin{corollary}[Explicit weak-coupling star rate]
\label{cor:s6v41-star-rate}
Whenever \(Q_{\epsilon/2}^{I}(\beta)<1\),
\begin{equation}
 \boxed{
 p_{\rm star}(\beta,\epsilon)
 \le
 \min\left\{1,\,
 6c_{\rm H}^{-1/6}\beta^{2/3}
 \exp\left({2\over3}
 -{\beta\epsilon^2\over9216}\right)\right\}.}
 \label{eq:s6v41-star-rate}
\end{equation}
\end{corollary}

\begin{proof}
Insert \(\delta=\epsilon/2\) into
\eqref{eq:s6v41-QI} and then into
\eqref{eq:s6v41-pstar}.  The exponent is divided by
\(96\cdot4=384\):
\[
 {1\over6\cdot384}\left({\epsilon\over2}\right)^2
 ={\epsilon^2\over9216}.
\]
The remaining powers are the same arithmetic as in V28.
\end{proof}

\begin{corollary}[Exponential misalignment-cluster tail]
\label{cor:s6v41-cluster}
On the degree-eighteen link dependency graph, if
\begin{equation}
 17p_{\rm star}(\beta,\epsilon)<1,
 \label{eq:s6v41-subcritical}
\end{equation}
then the probability that a fixed root belongs to a self-avoiding path of
at least \(n\) misaligned stars is at most
\begin{equation}
 \boxed{
 {18\over17}
 {(17p_{\rm star})^n\over1-17p_{\rm star}}.}
 \label{eq:s6v41-cluster}
\end{equation}
The bound passes to every subsequential infinite-volume limit of the
periodic pure-Wilson measures.
\end{corollary}

\begin{proof}
There are at most \(18\cdot17^{j-1}\) self-avoiding paths of \(j\)
links from a root.  By
\Cref{thm:s6v41-star-MDSC}, the probability that all stars of one path
are misaligned is at most \(p_{\rm star}^j\).  Sum over \(j\ge n\).
The events are cylinders on the compact product space, so weak limits
preserve the inequality.
\end{proof}

\section{What this adds to AQIC and CPTC}
\label{sec:s6v41-handoff}

\begin{theorem}[Bare periodic AQIC probability half]
\label{thm:s6v41-AQIC-half}
Fix \(\epsilon,r,s\) satisfying the strict aligned quotient condition
\eqref{eq:s6v40-row-condition} and the V39 local energy-gap conditions.
For the bare periodic pure-Wilson law, all of the following hold for
sufficiently large \(\beta\):
\begin{enumerate}[label=\textup{(\roman*)},leftmargin=3.5em]
\item every aligned star has the strict quotient influence row
      \eqref{eq:s6v40-full-row};
\item every finite family of misaligned stars satisfies
      \eqref{eq:s6v41-star-MDSC}; and
\item misaligned-star paths have the exponential tail
      \eqref{eq:s6v41-cluster}.
\end{enumerate}
\end{theorem}

\begin{proof}
The first item is V40.  Since
\eqref{eq:s6v41-star-rate} tends to zero exponentially at fixed
\(\epsilon>0\), the subcritical condition
\eqref{eq:s6v41-subcritical} holds eventually.  The last two items are
\Cref{thm:s6v41-star-MDSC,cor:s6v41-cluster}.
\end{proof}

\begin{firewall}[Unconditional contours are not a marked connected
expansion]
\label{fire:s6v41-not-CPTC}
\Cref{thm:s6v41-AQIC-half} is proved for one periodic pure-Wilson law.
It does not show that the same multi-star bound survives conditioning on
an explored exterior, nor that a response-generated action preserves the
nonnegative character expansion.  It also does not prove cancellation of
two disjoint unmarked bad clusters in the normalized two-score cumulant.
Consequently it supplies the probability half of the bare AQIC card but
not CPTC, CDAC, or the response-matched covariance kernel.
\end{firewall}

\begin{assessment}[Centre sheets have been removed from this particular
bottleneck]
\label{ass:s6v41-centre-removed}
The V40 programme anticipated a split between noncentral defects and
incompatible centre labels.  The relative-loop identity gives a stronger
route: failure of staple alignment forces an incident plaquette away from
the identity itself.  The Wilson defect
\eqref{eq:s6v41-Wilson-distance} then charges both noncentral and
nontrivial-centre plaquettes.  Thus the bare pure-Wilson
misaligned-star contour is exhaustive without a separate \(Z_3\) gas.
Centre sheets may re-enter only after a quotient or generated interaction
makes identity distance no longer the correct energy carrier.
\end{assessment}

\begin{programme}[V42: boundary-stable star contours]
\label{prog:s6v41-V42}
\begin{enumerate}[leftmargin=3em]
\item Condition on an explored exterior and isolate the minimal doubled
      collar containing every plaquette used by the relative-star event.
\item Compare the collar conditional to a periodic reflected collar by an
      explicit boundary oscillation or vacuum-fillability debit.
\item Determine whether the V29 threshold
      \(\rho_\Omega<\delta/\sqrt{24}\) can be expressed directly in the
      staple-alignment parameters \((\epsilon,r)\).
\item If so, derive a conditional version of
      \eqref{eq:s6v41-star-MDSC} and combine it with the V40 good quotient
      row.
\item If not, retain the present unconditional contour for CPTC and seek a
      separate disagreement coupling for CDAC.
\end{enumerate}
\end{programme}

\begin{firewall}[Claim boundary after sixth-series V41]
\label{fire:s6v41-boundary}
V41 proves the exact two-staple relative-loop identity, proves that staple
misalignment forces an identity-bad incident plaquette, establishes a
reflection-positive chessboard bound for identity defects, and derives
uniform multi-star suppression and an explicit exponential
misalignment-cluster tail for the periodic pure-Wilson law.  Together
with V40 it closes the bare periodic probability half of AQIC.

V41 does not prove boundary-stable or response-fixed star suppression,
the marked connected cancellations of CPTC, LCKC for the nonlinear law,
selected root asymptotics, the generated-action extension, TSC, YMGNC,
continuum existence, or the full Yang--Mills mass gap.
\end{firewall}

\begin{theorem}[Sixth-series V41 result ledger]
\label{thm:s6v41-results}
The claims in \Cref{fire:s6v41-boundary} follow from
\Cref{lem:s6v41-relative-loop,
lem:s6v41-misalignment-plaquette,
thm:s6v41-same-colour,
cor:s6v41-all-colours,
thm:s6v41-star-MDSC,
cor:s6v41-star-rate,
cor:s6v41-cluster,
thm:s6v41-AQIC-half}.
\end{theorem}

\begin{proof}
Each ledger item is the conclusion of the cited result.
\end{proof}

\paragraph{Parent-version record.}
V41 was formed from
\texttt{Renewal\_Yang\_Mills\_Mass\_Gap\_Research\_Notes\_V40.tex}.
The V40 parent had \(18\,062\) logical source lines,
\(629\,878\) bytes, and SHA--256
\[
 \texttt{d4c85164f1d7f46d2725a526e44606e6b5d28d26a1e3396d8fc96b857030cfdf}.
\]
The next compulsory companion audit is V45.

\clearpage
\chapter{Sixth-series V42: boundary-stable identity defects and
the conditional star card}
\label{chap:s6v42}

\section{The V29 conditional theorem for identity defects}
\label{sec:s6v42-conditional}

Let \(\Omega\) be a finite variable-link collar, let \(\xi\) be its frozen
exterior, and retain the V29 vacuum-fillability radius
\[
 \rho_\Omega(\xi)
 =\min_{z^0}
 \max_{p\in{\mathscr P}(\Omega)}
 \|U_p(z^0,\xi)-I\|_{\rm F}.
\]
For \(\Gamma\subset{\mathscr P}(\Omega)\), use the identity-defect events
\[
 C_p(\delta)=\{\|U_p-I\|_{\rm F}\ge\delta\}.
\]

\begin{theorem}[Vacuum-fillable conditional identity-defect suppression]
\label{thm:s6v42-conditional-identity}
Suppose
\[
 |\Omega|\le A|\Gamma|,
 \qquad
 |{\mathscr P}(\Omega)|\le B|\Gamma|,
 \qquad
 \rho_\Omega(\xi)\le\rho.
\]
For \(\beta\ge r_{\rm H}^{-2}\), put
\begin{equation}
 Q_{\rm cond}^{I}(\beta,\delta,\rho;A,B)
 =c_{\rm H}^{-A}\beta^{4A}
 \exp\left\{
 -{\beta\delta^2\over6}
 +{\beta B\over6}
  (\rho+4\beta^{-1/2})^2
 \right\}.
 \label{eq:s6v42-Qcond}
\end{equation}
Then the pure-Wilson conditional satisfies
\begin{equation}
 \boxed{
 \gamma_{\Omega,\beta}^{\xi}
 \left(\bigcap_{p\in\Gamma}C_p(\delta)\right)
 \le
 \{Q_{\rm cond}^{I}(\beta,\delta,\rho;A,B)\}^{|\Gamma|}.}
 \label{eq:s6v42-conditional-identity}
\end{equation}
If a generated action \(\Psi\) has conditional oscillation at most
\(\kappa_\Psi|\Gamma|\), the right side is multiplied by
\(e^{\kappa_\Psi|\Gamma|}\).
\end{theorem}

\begin{proof}
On the event,
\[
 \sum_{p\in{\mathscr P}(\Omega)}d(U_p)
 \ge {|\Gamma|\delta^2\over6}
\]
by the exact identity
\(d(U)=\|U-I\|_{\rm F}^2/6\).  Hence its unnormalized conditional
integral is at most
\(e^{-\beta|\Gamma|\delta^2/6}\).
The V29 translated-Haar box at a minimizing filling gives
\[
 Z_{\Omega,\beta}^{\xi}
 \ge
 h(\beta^{-1/2})^{|\Omega|}
 \exp\left\{
 -{\beta|{\mathscr P}(\Omega)|\over6}
 (\rho+4\beta^{-1/2})^2
 \right\}.
\]
Divide, use
\(h(\beta^{-1/2})\ge c_{\rm H}\beta^{-4}\), and apply the incidence
bounds.  This gives
\eqref{eq:s6v42-conditional-identity}.  The response-action statement is
the V29 Radon--Nikodym oscillation inequality.
\end{proof}

\begin{corollary}[Sharp identity-defect fillability threshold]
\label{cor:s6v42-fillability}
For fixed \(A,B,\delta,\rho,\kappa_\Psi\), the conditional bound decays
exponentially in \(\beta|\Gamma|\) whenever
\begin{equation}
 B\rho^2<\delta^2.
 \label{eq:s6v42-fillability}
\end{equation}
For the canonical plaquette collar \(A=4,B=24\), and for the
misalignment reduction \(\delta=\epsilon/2\), this is
\begin{equation}
 \boxed{
 \rho<{\epsilon\over4\sqrt6}.}
 \label{eq:s6v42-star-fillability}
\end{equation}
\end{corollary}

\begin{proof}
Divide the logarithm of
\eqref{eq:s6v42-Qcond} by \(\beta\).  Its limit is
\[
 -{\delta^2\over6}+{B\rho^2\over6},
\]
which is negative precisely under
\eqref{eq:s6v42-fillability}.  Substituting
\(B=24\) and \(\delta=\epsilon/2\) gives
\(24\rho^2<\epsilon^2/4\), equivalent to
\eqref{eq:s6v42-star-fillability}.
\end{proof}

\section{Conditional multi-star suppression}
\label{sec:s6v42-stars}

\begin{definition}[Star path fillability card, SPFC]
\label{def:s6v42-SPFC}
Fix a predeclared self-avoiding link path.  SPFC with parameters
\((\epsilon,\rho,\kappa_\Psi)\) means that for every set of path links
and every selection of one incident identity-bad plaquette at each
misaligned star, the canonical variable-link collar of the distinct
selected plaquettes has
\[
 \rho_\Omega(\xi)\le\rho
\]
for the actual explored exterior, and the generated-action oscillation on
that collar is at most
\(\kappa_\Psi|\Gamma|\).
\end{definition}

\begin{theorem}[Boundary-stable misaligned-star MDSC]
\label{thm:s6v42-star-MDSC}
Assume SPFC and
\(\rho<\epsilon/(4\sqrt6)\).  Define
\begin{align}
 \widehat Q_{\rm star}
 &=
 e^{\kappa_\Psi}
 Q_{\rm cond}^{I}(\beta,\epsilon/2,\rho;4,24),
 \label{eq:s6v42-Qstar}\\
 \widehat p_{\rm star}
 &=
 \min\{1,6\widehat Q_{\rm star}^{1/4}\}.
 \label{eq:s6v42-phat}
\end{align}
Then for every set \(K\) of distinct path links, conditionally on every
exterior satisfying SPFC,
\begin{equation}
 \boxed{
 \mathbb P\{
 {\cal M}_e(\epsilon)\text{ for every }e\in K
 \mid\xi\}
 \le
 \widehat p_{\rm star}^{|K|}.}
 \label{eq:s6v42-star-MDSC}
\end{equation}
For fixed \(\epsilon,\rho,\kappa_\Psi\) in this range,
\(\widehat p_{\rm star}\to0\) exponentially as
\(\beta\to\infty\).
\end{theorem}

\begin{proof}
By V41, every misaligned star forces at least one incident
\(C_p(\epsilon/2)\) event.  Expanding the union gives at most
\(6^{|K|}\) assignments.  Each plaquette contains at most four links, so
an assignment uses at least \(|K|/4\) distinct plaquettes.  Their canonical
collar has \(A=4,B=24\).  Apply SPFC and
\Cref{thm:s6v42-conditional-identity}, including the response
oscillation.  Each assignment costs at most
\(\widehat Q_{\rm star}^{|K|/4}\); summing proves
\eqref{eq:s6v42-star-MDSC}.  Exponential decay follows from
\Cref{cor:s6v42-fillability}.
\end{proof}

\begin{corollary}[Conditional star-path tail]
\label{cor:s6v42-path-tail}
If
\[
 17\widehat p_{\rm star}<1,
\]
the conditional probability that a fixed root is connected by a
self-avoiding path of at least \(n\) misaligned stars is at most
\begin{equation}
 \boxed{
 {18\over17}
 {(17\widehat p_{\rm star})^n
  \over1-17\widehat p_{\rm star}}.}
 \label{eq:s6v42-path-tail}
\end{equation}
\end{corollary}

\begin{proof}
Condition on the exterior.  There are at most
\(18\cdot17^{j-1}\) self-avoiding paths of length \(j\), and
\Cref{thm:s6v42-star-MDSC} bounds the event that every star of one path is
misaligned by \(\widehat p_{\rm star}^j\).  Sum the geometric series over
\(j\ge n\).
\end{proof}

\section{Compatibility with the strict aligned quotient row}
\label{sec:s6v42-compatibility}

\begin{lemma}[The local parameter window is nonempty]
\label{lem:s6v42-window}
There exist \(0<s<r\) and \(\epsilon,\rho>0\) such that simultaneously:
\begin{enumerate}[label=\textup{(\roman*)},leftmargin=3.5em]
\item \(\delta_{\epsilon,r}=\epsilon+e^r-1<1\);
\item \(r(1-\delta_{\epsilon,r})>\epsilon\);
\item \(\Delta_{s,r,\epsilon}
      =6\{g(r)-h(s)-2\epsilon/\sqrt3\}>0\);
\item \(6\epsilon/(1-\delta_{\epsilon,r})<1/17\); and
\item \(\rho<\epsilon/(4\sqrt6)\).
\end{enumerate}
\end{lemma}

\begin{proof}
Choose any sufficiently small fixed \(r>0\) with \(e^r-1<1/2\).
Since \(g(r)>0\) and \(h(s)\to0\), choose \(s<r\) with
\(h(s)<g(r)/2\).  Now choose \(\epsilon>0\) small enough that
\[
 \epsilon<\min\left\{
 {r(2-e^r)\over1+r},
 {\sqrt3\,g(r)\over4},
 {1-(e^r-1)\over109}
 \right\}.
\]
The first bound is equivalent to
\(r(1-\epsilon-e^r+1)>\epsilon\); the second gives
\(h(s)+2\epsilon/\sqrt3<g(r)\); the third gives
\(6\epsilon/(1-\delta_{\epsilon,r})<1/18<1/17\).
It also makes \(\delta_{\epsilon,r}<1\).  Finally choose any
\(0<\rho<\epsilon/(4\sqrt6)\).
\end{proof}

\begin{theorem}[Conditional pure-Wilson star-card handoff]
\label{thm:s6v42-handoff}
Choose parameters from
\Cref{lem:s6v42-window}, assume SPFC, and let
\[
 a_{\rm W}(\beta)
 ={6\epsilon\over1-\delta_{\epsilon,r}}
 +C_{\epsilon,r,s}\epsilon\beta e^{-\beta\Delta/2}
\]
be the V40 aligned quotient row.  Then, for all sufficiently large
\(\beta\),
\begin{equation}
 \boxed{
 17\{a_{\rm W}(\beta)+\widehat p_{\rm star}(\beta)\}<1.}
 \label{eq:s6v42-numerical-absorption}
\end{equation}
Thus the Wilson good-star quotient coefficient and the conditional
misaligned-star probability are numerically compatible with the
degree-eighteen path absorption threshold.
\end{theorem}

\begin{proof}
By construction,
\[
 17{6\epsilon\over1-\delta_{\epsilon,r}}<1.
\]
The remaining term in \(a_{\rm W}\) tends to zero exponentially by
\(\Delta>0\).  The conditional star probability tends to zero
exponentially by
\Cref{thm:s6v42-star-MDSC}.  Their sum therefore satisfies
\eqref{eq:s6v42-numerical-absorption} eventually.
\end{proof}

\begin{firewall}[Numerical compatibility is not yet CDAC]
\label{fire:s6v42-not-CDAC}
The coefficient \(a_{\rm W}\) is a Wasserstein/Lipschitz influence on the
recentered one-link quotient.  The earlier CDAC theorem was stated using a
specific disagreement coupling and its product event.  A coupling theorem
which carries the quotient metric through a sequence of good and bad
stars has not yet been proved.  Moreover, the generated action was used
above only through its collar oscillation in the bad-star probability.
Its quotient directional derivative may add a good-star influence debit.
Therefore
\eqref{eq:s6v42-numerical-absorption} is the exact feasible input to, not
a proof of, metric CDAC.
\end{firewall}

\begin{definition}[Generated quotient influence debit]
\label{def:s6v42-generated-debit}
For a generated conditional action \(\Psi\), define
\[
 a_\Psi
 =\sup_{\rm good\ exteriors}
 \sum_{f\leadsto e}
 \sup_{\operatorname{Lip}(f_e)\le1}
 \left|
 D_f\,
 \mathbb E_{\gamma_e^{\rm quotient}}f_e
 \right|_{\Psi\text{-part}},
\]
where the coherent staple direction is removed before the derivative.
The response-matched absorption budget is
\begin{equation}
 \boxed{
 17\{a_{\rm W}+a_\Psi+\widehat p_{\rm star}\}<1.}
 \label{eq:s6v42-generated-budget}
\end{equation}
The seven static stocks of V33 do not bound \(a_\Psi\); this is a
first-derivative conditional specification norm, not a retained Hessian.
\end{definition}

\begin{programme}[V43: metric disagreement coupling]
\label{prog:s6v42-V43}
\begin{enumerate}[leftmargin=3em]
\item Formulate the quotient single-site kernels on one complete common
      metric space and use Kantorovich duality to convert the V40
      Lipschitz row into an optimal Wasserstein coupling.
\item Explore the lattice sequentially and prove the path product with
      factor \(a_{\rm W}+a_\Psi+\widehat p_{\rm star}\).
\item Derive covariance decay for quotient-Lipschitz Wilson sources.
\item Keep the physical clock explicit as in V35; do not use a fixed
      lattice contraction as a continuum mass if the block duration
      shrinks.
\item Identify the generated writer which first contributes to
      \(a_\Psi\).
\end{enumerate}
\end{programme}

\begin{firewall}[Claim boundary after sixth-series V42]
\label{fire:s6v42-boundary}
V42 proves boundary-stable conditional suppression for identity-defect
plaquettes, derives the exact star fillability threshold
\(\rho<\epsilon/(4\sqrt6)\), proves conditional multi-star MDSC and a
star-path tail under SPFC, proves that the fillability, local-vacuum,
strict quotient, and path-absorption numerical parameters have a common
nonempty range, and isolates the generated quotient influence debit.

V42 does not prove SPFC occurrence for every exploration, the metric
disagreement product, the generated quotient debit, complete CDAC, CPTC,
LCKC, selected root asymptotics, TSC, YMGNC, continuum existence, or the
full Yang--Mills mass gap.
\end{firewall}

\begin{theorem}[Sixth-series V42 result ledger]
\label{thm:s6v42-results}
The claims in \Cref{fire:s6v42-boundary} follow from
\Cref{thm:s6v42-conditional-identity,
cor:s6v42-fillability,
thm:s6v42-star-MDSC,
cor:s6v42-path-tail,
lem:s6v42-window,
thm:s6v42-handoff}.
\end{theorem}

\begin{proof}
Each ledger item is the conclusion of the cited result.
\end{proof}

\paragraph{Parent-version record.}
V42 was formed from
\texttt{Renewal\_Yang\_Mills\_Mass\_Gap\_Research\_Notes\_V41.tex}.
The V41 parent had \(18\,431\) logical source lines,
\(642\,499\) bytes, and SHA--256
\[
 \texttt{845dcfe971a22db495f669224e046ded32ecf5d9d41ffb796282494d9443094d}.
\]
The next compulsory companion audit is V45.

\clearpage
\chapter{Sixth-series V43: metric disagreement coupling on the
aligned quotient}
\label{chap:s6v43}

\section{The compact quotient-kernel setup}
\label{sec:s6v43-setup}

Let \(({\cal X},d)\) be the compact one-block quotient state space obtained
by the coherent recentering of V40, and normalize
\(\operatorname{diam}{\cal X}\le1\).  For a probability kernel \(K\) on
\({\cal X}\), write
\[
 W_1(K,K')
 =\sup_{\operatorname{Lip}(f)\le1}|Kf-K'f|
\]
for the Kantorovich distance.

\begin{lemma}[Lipschitz response gives an optimal metric coupling]
\label{lem:s6v43-Kantorovich}
Suppose two good-star one-site conditionals \(K_e^\xi,K_e^{\xi'}\) obey
\begin{equation}
 W_1(K_e^\xi,K_e^{\xi'})
 \le\sum_{f\sim e}c_{ef}d(\xi_f,\xi_f').
 \label{eq:s6v43-W1-row}
\end{equation}
Then there is a coupling \((X_e,X'_e)\) of these conditionals such that
\begin{equation}
 \mathbb E d(X_e,X'_e)
 \le\sum_{f\sim e}c_{ef}d(\xi_f,\xi_f').
 \label{eq:s6v43-optimal-coupling}
\end{equation}
For a measurable family of finite-dimensional compact kernels, the
couplings may be selected measurably in the exterior data.
\end{lemma}

\begin{proof}
Kantorovich--Rubinstein duality identifies the displayed supremum with the
minimum of \(\int d(x,y)\,d\pi(x,y)\) over couplings.  Compactness makes
the coupling set compact and the cost continuous, so the minimum is
attained.  The graph of the set of minimizers is measurable for a
measurable finite-dimensional kernel family; the measurable-selection
theorem supplies a measurable minimizer.
\end{proof}

\begin{definition}[Metric star exploration card, MSEC]
\label{def:s6v43-MSEC}
MSEC with parameters \(a,p\in[0,1]\) means that one common sequential
exploration of two boundary conditions has:
\begin{enumerate}[label=\textup{(\roman*)},leftmargin=3.5em]
\item bad-star indicators \(B_e\) satisfying, for every set \(S\) on every
      predeclared self-avoiding exploration path,
      \[
       \mathbb E\prod_{e\in S}B_e\le p^{|S|};
      \]
\item at a good star \(B_e=0\), an optimal quotient coupling whose expected
      outgoing distance is at most \(a\) times the incoming path distance;
      at a bad star the normalized outgoing distance is bounded by one;
      and
\item the same quotient coordinates and metric are used at every step.
\end{enumerate}
\end{definition}

\section{The exact mixed good/bad path factor}
\label{sec:s6v43-path}

\begin{lemma}[MDSC moment expansion]
\label{lem:s6v43-moment}
Let \(B_1,\ldots,B_\ell\in\{0,1\}\) satisfy
\[
 \mathbb E\prod_{i\in S}B_i\le p^{|S|}
\]
for every \(S\subseteq\{1,\ldots,\ell\}\).  Then for \(0\le a\le1\),
\begin{equation}
 \boxed{
 \mathbb E\prod_{i=1}^{\ell}
 \{a+(1-a)B_i\}
 \le\{a+(1-a)p\}^{\ell}
 \le(a+p)^\ell.}
 \label{eq:s6v43-moment}
\end{equation}
\end{lemma}

\begin{proof}
Expand the product:
\[
 \sum_{S\subseteq\{1,\ldots,\ell\}}
 a^{\ell-|S|}(1-a)^{|S|}
 \mathbb E\prod_{i\in S}B_i.
\]
Insert the moment bound and sum the binomial series.  The second inequality
uses \((1-a)p\le p\).
\end{proof}

\begin{theorem}[Metric path transmission]
\label{thm:s6v43-path}
Under MSEC, the expected quotient distance carried along any fixed
self-avoiding exploration path of \(\ell\) gates is at most
\begin{equation}
 \boxed{
 \theta^\ell,
 \qquad
 \theta:=a+(1-a)p.}
 \label{eq:s6v43-theta}
\end{equation}
\end{theorem}

\begin{proof}
Condition on the bad-star vector.  At a good gate the sequential optimal
coupling multiplies the incoming expected distance by at most \(a\); at a
bad gate it multiplies by at most one because the metric diameter is one.
Thus the conditional path factor is bounded by
\[
 \prod_{i=1}^{\ell}\{a+(1-a)B_i\}.
\]
Take expectation and apply
\Cref{lem:s6v43-moment}.
\end{proof}

\begin{corollary}[Connectivity tail on a degree-\(\Delta\) graph]
\label{cor:s6v43-connectivity}
Let
\begin{equation}
 \alpha=(\Delta-1)\theta.
 \label{eq:s6v43-alpha}
\end{equation}
If \(\alpha<1\), the total expected metric influence carried from a root
to graph distance at least \(r\ge1\) is at most
\begin{equation}
 \boxed{
 {\Delta\over\Delta-1}
 {\alpha^r\over1-\alpha}.}
 \label{eq:s6v43-connectivity}
\end{equation}
\end{corollary}

\begin{proof}
There are at most
\(\Delta(\Delta-1)^{\ell-1}\) self-avoiding paths of length \(\ell\)
from the root.  Apply
\Cref{thm:s6v43-path} to each and sum over \(\ell\ge r\).
\end{proof}

\section{Covariance decay for quotient-Lipschitz observables}
\label{sec:s6v43-covariance}

For a function \(F\) of finitely many quotient blocks, define its
single-block Lipschitz oscillation
\[
 \delta_e^{\rm Lip}(F)
 =\sup_{\xi,\xi':\,\xi_{e^c}=\xi'_{e^c}}
 { |F(\xi)-F(\xi')|\over d(\xi_e,\xi'_e)}.
\]

\begin{theorem}[MSEC covariance decay]
\label{thm:s6v43-covariance}
Assume the MSEC exploration is a valid coupling of the two Gibbs
specifications.  Let bounded quotient-Lipschitz \(F,G\) have supports
\(A,B\) at graph distance \(r\ge1\).  If \(\alpha<1\), then
\begin{equation}
 \boxed{
 |\operatorname{Cov}(F,G)|
 \le
 {\Delta\over\Delta-1}
 {\alpha^r\over1-\alpha}
 \left(\sum_{e\in A}\delta_e^{\rm Lip}(F)\right)
 \max_{f\in B}\delta_f^{\rm Lip}(G).}
 \label{eq:s6v43-covariance}
\end{equation}
\end{theorem}

\begin{proof}
Reveal the blocks in the support of \(F\) one at a time and write the
resulting martingale difference decomposition.  Couple the two conditional
laws of the unrevealed exterior after changing one revealed block.
The change in the conditional expectation of \(G\) is bounded by the
Lipschitz oscillation of \(G\) times the expected quotient distance which
reaches its support.  By
\Cref{cor:s6v43-connectivity}, this is at most the geometric factor in
\eqref{eq:s6v43-covariance}.  Sum the martingale increments of \(F\).
\end{proof}

\section{Population by the V40--V42 pure-Wilson rows}
\label{sec:s6v43-population}

\begin{theorem}[SPFC gives pure-Wilson quotient MSEC]
\label{thm:s6v43-SPFC-MSEC}
Fix parameters in the nonempty V42 window and assume SPFC along the
exploration.  For the pure-Wilson quotient specification, MSEC holds with
\begin{align}
 a&=a_{\rm W}(\beta)
 ={6\epsilon\over1-\delta_{\epsilon,r}}
 +C_{\epsilon,r,s}\epsilon\beta e^{-\beta\Delta/2},
 \label{eq:s6v43-aW}\\
 p&=\widehat p_{\rm star}(\beta).
 \label{eq:s6v43-pstar}
\end{align}
Consequently, on the degree-eighteen graph,
\begin{equation}
 \alpha_{\rm W}(\beta)
 =17\{a_{\rm W}+(1-a_{\rm W})\widehat p_{\rm star}\}<1
 \label{eq:s6v43-alphaW}
\end{equation}
for all sufficiently large \(\beta\), and the covariance bound
\eqref{eq:s6v43-covariance} holds with
\(\Delta=18\) and \(\alpha=\alpha_{\rm W}\).
\end{theorem}

\begin{proof}
V42 gives the bad-star moment bound conditionally on every SPFC exterior.
V40 bounds the \(W_1\) response of every aligned good-star kernel in the
common recentered quotient.  By
\Cref{lem:s6v43-Kantorovich}, those response bounds are realized by
sequential optimal couplings.  The metric was normalized to diameter one,
so a bad gate has the trivial factor one.  These are precisely the three
MSEC rows.  The V42 numerical absorption theorem gives
\(17(a_{\rm W}+\widehat p_{\rm star})<1\) eventually; since
\(a+(1-a)p\le a+p\), it implies
\eqref{eq:s6v43-alphaW}.  Apply
\Cref{thm:s6v43-covariance}.
\end{proof}

\begin{corollary}[Generated-action extension with a quotient debit]
\label{cor:s6v43-generated}
Suppose the response-generated conditional kernels preserve the same
quotient, SPFC bad-star bound, and common metric, and their good-star
Wasserstein row is at most \(a_{\rm W}+a_\Psi\le1\).  Then MSEC and
\Cref{thm:s6v43-covariance} hold with
\begin{equation}
 \alpha_\Psi
 =17\left[
 a_{\rm W}+a_\Psi+
 \{1-a_{\rm W}-a_\Psi\}\widehat p_{\rm star}
 \right],
 \label{eq:s6v43-alphaPsi}
\end{equation}
provided \(\alpha_\Psi<1\).
\end{corollary}

\begin{proof}
Replace the good-gate factor \(a\) in
\Cref{thm:s6v43-path} by \(a_{\rm W}+a_\Psi\).  The bad-star moment row is
unchanged by hypothesis.  The path counting and covariance proof are
identical.
\end{proof}

\section{Physical-clock spectral consequence}
\label{sec:s6v43-OS}

\begin{theorem}[Metric-CDAC finite-regulator OS rate]
\label{thm:s6v43-OS}
Assume reflection positivity, a quotient-Lipschitz finite source bank with
a regulator-uniform lower frame, and that one exploration step advances
physical time by \(\tau_j>0\).  If MSEC holds at regulator \(j\) with
\(\alpha_j<1\), then the bank's delayed reflected Grams have exponential
rate
\begin{equation}
 \boxed{
 \mu_j={|\log\alpha_j|\over\tau_j}.}
 \label{eq:s6v43-muj}
\end{equation}
If the moving-delay Grams converge and
\[
 0<\liminf_j\mu_j<\infty,
\]
the limiting cyclic bank has Hamiltonian gap at least
\(\liminf_j\mu_j\).
\end{theorem}

\begin{proof}
\Cref{thm:s6v43-covariance} applied after \(k\) time steps gives a
bank-dependent prefactor times \(\alpha_j^k\).  Reflection positivity
identifies this covariance with the positive delayed OS Gram, and the
lower frame converts coefficient norms to OS norms.  The exponent per
physical time is \eqref{eq:s6v43-muj}.  Apply the variable-clock spectral
transfer theorem of V35.
\end{proof}

\begin{firewall}[The fixed-factor continuum trap remains]
\label{fire:s6v43-clock}
At fixed \(\epsilon,r\), the pure-Wilson factor
\(\alpha_{\rm W}(\beta)\) is eventually bounded strictly below one.  If
the physical duration \(\tau_j\) of one such star step tends to zero, then
\(\mu_j\to\infty\).  By the V35 ultralocality theorem this is incompatible
with a nonzero strongly continuous limiting source bank under uniform
Gram convergence.  Thus the present result is a valid lattice
correlation theorem, but its continuum use requires either blocks of
nonvanishing physical duration or a scale-dependent near-critical
\(\alpha_j=1-\mu\tau_j+o(\tau_j)\).
\end{firewall}

\begin{assessment}[What is now closed and what remains]
\label{ass:s6v43-frontier}
For the pure-Wilson quotient, SPFC now gives an actual metric
disagreement coupling and exponential covariance decay, rather than only
numerical compatibility.  Two independent inputs still prevent a
response-matched continuum mass claim:
\begin{enumerate}[label=\textup{(\roman*)},leftmargin=3.5em]
\item SPFC occurrence has not been proved along every exploration, and the
      generated quotient debit \(a_\Psi\) is unbounded; and
\item the lattice contraction has not been calibrated to a finite
      physical rate on the continuum clock.
\end{enumerate}
The static LCKC and the dynamic MSEC are complementary.  MSEC gives decay
but not the Gaussian covariance asymptotic required by the trace-shell
static compiler.
\end{assessment}

\begin{programme}[V44: populate the generated quotient debit]
\label{prog:s6v43-V44}
\begin{enumerate}[leftmargin=3em]
\item Differentiate a general finite-range generated interaction under the
      V40 coherent recentering.
\item Express its good-star \(W_1\) influence through support-local first
      and mixed conditional derivatives.
\item Define a finite interaction stock whose row sum is exactly
      \(a_\Psi\), keeping it separate from the seven second-jet stocks of
      V33.
\item Test plaquette marginal retuning
      \(b^{\rm marg}P\) first; its coherent cancellation should be
      computable explicitly from the six incident plaquettes.
\item State the remaining base-generated writer which first lacks a
      cofinal bound.
\end{enumerate}
\end{programme}

\begin{firewall}[Claim boundary after sixth-series V43]
\label{fire:s6v43-boundary}
V43 proves the optimal-coupling realization of a quotient Lipschitz row,
the exact mixed path factor \(a+(1-a)p\), metric connectivity and
covariance decay, and population of MSEC by the V40--V42 pure-Wilson rows
under SPFC.  It extends the theorem conditionally through a generated
quotient debit and gives the regulator and variable-clock OS rates.

V43 does not prove unconditional SPFC occurrence, the generated quotient
debit, CPTC, LCKC, selected root asymptotics, TSC, YMGNC, a finite
continuum physical rate, continuum existence, or the full Yang--Mills
mass gap.
\end{firewall}

\begin{theorem}[Sixth-series V43 result ledger]
\label{thm:s6v43-results}
The claims in \Cref{fire:s6v43-boundary} follow from
\Cref{lem:s6v43-Kantorovich,
lem:s6v43-moment,
thm:s6v43-path,
cor:s6v43-connectivity,
thm:s6v43-covariance,
thm:s6v43-SPFC-MSEC,
cor:s6v43-generated,
thm:s6v43-OS}.
\end{theorem}

\begin{proof}
Each ledger item is the conclusion of the cited result.
\end{proof}

\paragraph{Parent-version record.}
V43 was formed from
\texttt{Renewal\_Yang\_Mills\_Mass\_Gap\_Research\_Notes\_V42.tex}.
The V42 parent had \(18\,794\) logical source lines,
\(654\,064\) bytes, and SHA--256
\[
 \texttt{34b3edb49881f68cc40d1eec8ebeab7352410d960ffa3cf8f4138ccc1425f4c4}.
\]
The next compulsory companion audit is V45.

\clearpage
\chapter{Sixth-series V44: generated quotient-influence stocks and
the plaquette-marginal cancellation}
\label{chap:s6v44}

\section{The matched one-link conditional}
\label{sec:s6v44-conditional}

On one link star, write the response-matched conditional potential as
\begin{equation}
 {\cal U}(U;\xi)
 =(\beta+b^{\rm marg})\Phi_A(U;\xi)
  +\Psi^{\rm gen}(U;\xi).
 \label{eq:s6v44-potential}
\end{equation}
Set
\[
 \beta_{\rm eff}=\beta+b^{\rm marg}.
\]
The decomposition is exact because the marginal writer \(P\) is the same
plaquette-defect sum as the Wilson head.

\begin{theorem}[Plaquette marginal is an effective-coupling shift]
\label{thm:s6v44-effective-beta}
If \(\Psi^{\rm gen}=0\) and \(\beta_{\rm eff}>0\), every result of
V38--V43 for the pure-Wilson one-link conditional remains valid after the
single substitution
\[
 \beta\longmapsto\beta_{\rm eff}.
\]
In particular:
\begin{enumerate}[label=\textup{(\roman*)},leftmargin=3.5em]
\item the coherent recentered first variation still vanishes exactly;
\item the leading good-star quotient row remains
      \(6\epsilon/(1-\delta_{\epsilon,r})\), independent of
      \(b^{\rm marg}\); and
\item the aligned-star complement and conditional identity-defect
      exponents use \(\beta_{\rm eff}\).
\end{enumerate}
\end{theorem}

\begin{proof}
With \(\Psi^{\rm gen}=0\),
\eqref{eq:s6v44-potential} is precisely the Wilson star potential at
inverse coupling \(\beta_{\rm eff}\).  The perfect-star cancellation of
V40 is an algebraic identity in the plaquette sum and is unchanged by a
scalar multiple.  In the near-aligned response estimate, both the
insertion gradient and the strong-convexity Hessian are multiplied by
\(\beta_{\rm eff}\); their ratio is therefore unchanged.  Every energy
tail is obtained by exponentiating the action and hence uses
\(\beta_{\rm eff}\).
\end{proof}

\begin{corollary}[The V34 marginal scale rows imply positivity of the
effective coupling]
\label{cor:s6v44-effective-positive}
If the V34 ratio compiler holds along a cofinal sequence, then
\[
 |b^{\rm marg}|=o(\beta/M^2)=o(\beta).
\]
Consequently
\[
 \beta_{\rm eff}/\beta\longrightarrow1
\]
and \(\beta_{\rm eff}>0\) eventually.  The plaquette marginal contributes
no separate generated quotient debit \(a_\Psi\).
\end{corollary}

\begin{proof}
Since \(M\to\infty\), the V34 row
\(|b^{\rm marg}|=o(\beta/M^2)\) implies
\(|b^{\rm marg}|/\beta\to0\).  Apply
\Cref{thm:s6v44-effective-beta}.
\end{proof}

\begin{firewall}[A near-cancelling marginal is a genuine obstruction]
\label{fire:s6v44-negative-beta}
If \(b^{\rm marg}\le-\beta\), the effective Wilson coefficient is
nonpositive.  The aligned-star convexity and energy-tail arguments cannot
be used.  Root existence or a finite value of \(b^{\rm marg}\) at one
cutoff does not exclude this branch; it is excluded cofinally only by a
ratio estimate such as V34 or by a direct positive lower bound on
\(\beta_{\rm eff}\).
\end{firewall}

\section{Three dynamic stocks for the genuinely generated interaction}
\label{sec:s6v44-stocks}

Fix the aligned quotient coordinates of V40.  For one good star, let
\({\cal B}_r(U_0)\) be the convex radius-\(r\) ball, and let
\(\zeta_f\) denote one exterior-link coordinate feeding a staple or a
generated term.  The derivative \(D_{\zeta_f}^{Q}\) is taken after the
coherent recentering, so the common staple rotation is removed.

\begin{definition}[Dynamic quotient-interaction stocks, DQIS]
\label{def:s6v44-DQIS}
The three one-star stocks of \(\Psi^{\rm gen}\) are
\begin{align}
 {\mathfrak h}_\Psi
 &:=
 \sup_{\rm good\ star}
 \sup_{U\in{\cal B}_r}
 \bigl\|(D_U^2\Psi^{\rm gen})_-\bigr\|_{\rm op},
 \label{eq:s6v44-h}\\
 {\mathfrak j}_\Psi
 &:=
 \sup_{\rm good\ star}
 \sum_{f\leadsto e}
 \sup_{U\in SU(3)}
 \|D_UD_{\zeta_f}^{Q}\Psi^{\rm gen}\|_{\rm op},
 \label{eq:s6v44-j}\\
 {\mathfrak o}_\Psi
 &:=
 \sup_{\rm declared\ star\ collar}
 \operatorname{osc}_{U,\mathrm{collar}}
 \Psi^{\rm gen}.
 \label{eq:s6v44-o}
\end{align}
Here \(H_-\) is the negative part of a Hermitian form.  The first stock
debits conditional convexity, the second is the summed quotient response
insertion, and the third transfers Wilson complement and defect
probabilities through the generated law.  The supremum in
\eqref{eq:s6v44-j} is deliberately global on the compact one-link fibre:
the convex-ball argument uses only its restriction, but the
full-versus-restricted covariance comparison also needs the oscillation
of the exterior insertion on the complement.
\end{definition}

\begin{assessment}[DQIS is not the V33 static stock list]
\label{ass:s6v44-stock-distinction}
The V33 stocks control vertical and retained second jets of a trace action
for a static Schur complement.  The new stock
\({\mathfrak j}_\Psi\) differentiates the \emph{one-site conditional
specification} with respect to an exterior quotient coordinate.  It is a
mixed spatial influence norm and is not bounded merely by knowing the
retained field Hessian.  The oscillation stock has appeared in V29, while
the negative vertical stock is the sign-sensitive part needed to preserve
strong convexity.
\end{assessment}

\section{Generated good-star response}
\label{sec:s6v44-response}

Let
\[
 \lambda=\lambda_{\epsilon,r}
 ={k\over3}(1-\delta_{\epsilon,r})
\]
be the Wilson aligned-star Hessian constant and suppose
\begin{equation}
 m_\Psi
 :=\beta_{\rm eff}\lambda-{\mathfrak h}_\Psi>0.
 \label{eq:s6v44-mPsi}
\end{equation}

\begin{theorem}[Generated quotient response bound]
\label{thm:s6v44-generated-response}
For a quotient-Lipschitz one-link observable \(f\), the summed derivative
of its conditional expectation over all exterior-link slots satisfies
\begin{align}
 \sum_{f_0\leadsto e}
 \left|
 D_{\zeta_{f_0}}
 \mathbb E_{{\cal U}}f
 \right|
 \le{}&
 \left\{
 {12\epsilon\beta_{\rm eff}\over m_\Psi}
 +{{\mathfrak j}_\Psi\over m_\Psi}
 \right\}\|\nabla f\|_\infty
 \notag\\
 &+
 C\bigl(\beta_{\rm eff}+{\mathfrak j}_\Psi+
          {\mathfrak o}_\Psi\bigr)
 \exp\left\{
 -{1\over2}
 (\beta_{\rm eff}\Delta-{\mathfrak o}_\Psi)
 \right\}
 \{\|\nabla f\|_\infty+\operatorname{osc}(f)\}.
 \label{eq:s6v44-generated-response}
\end{align}
The constant \(C\) depends only on the fixed local geometry and alignment
parameters.  In particular, the good-star Wasserstein row may be taken as
the expression in braces plus the displayed exponential remainder.
\end{theorem}

\begin{proof}
On the convex ball, the Hessian of the complete potential
\eqref{eq:s6v44-potential} is bounded below by
\[
 \beta_{\rm eff}\lambda I-{\mathfrak h}_\Psi I
 =m_\Psi I.
\]
The covariance form of the Bakry--\'Emery inequality therefore bounds the
derivative of a normalized conditional expectation by
\[
 {1\over m_\Psi}\|\nabla f\|_\infty
 \|\nabla(\hbox{exterior insertion})\|_\infty.
\]
For the Wilson part, the V40 recentered insertion bound is
\((2\epsilon/3)\) per unit of staple tangent.  Summing the
three slots in each of six staples gives
\[
 18\,{2\epsilon\beta_{\rm eff}\over3}
 =12\epsilon\beta_{\rm eff}.
\]
The generated insertion gradients sum to
\({\mathfrak j}_\Psi\) by definition.

The generated conditional is absolutely continuous with respect to the
\(\beta_{\rm eff}\)-Wilson conditional.  On the declared collar, its
Radon--Nikodym likelihood ratio has essential supremum at most
\(e^{{\mathfrak o}_\Psi}\).  Thus the V39 complement probability is at
most
\[
 C\exp\{-\beta_{\rm eff}\Delta+{\mathfrak o}_\Psi\}.
\]
Repeat the full-versus-restricted covariance decomposition of V40 for
both \(f\) and the complete exterior insertion.  Every complement term
contains the square root of this probability or the probability itself.
For a generated exterior insertion
\(g_f(U)=D_{\zeta_f}^{Q}\Psi^{\rm gen}(U)\), choose any fixed
\(U_\circ\in SU(3)\).  Compactness and the geodesic fundamental theorem
of calculus give
\[
 \operatorname{osc}_{U\in SU(3)}g_f
 \le
 \operatorname{diam}(SU(3))
 \sup_{U\in SU(3)}\|D_Ug_f(U)\|.
\]
Consequently the sum of the relevant insertion oscillations is bounded
by a fixed geometric multiple of \({\mathfrak j}_\Psi\); the constant
values \(g_f(U_\circ)\) vanish from normalized covariances.  The Wilson
insertion oscillations are bounded by a fixed incidence multiple of
\(\beta_{\rm eff}\).  Enlarging \(C\) gives the second line of
\eqref{eq:s6v44-generated-response}.
\end{proof}

\begin{corollary}[A cofinal sufficient DQIS row]
\label{cor:s6v44-DQIS-row}
Suppose
\begin{equation}
 {b^{\rm marg}\over\beta}\to0,
 \qquad
 {{\mathfrak h}_\Psi\over\beta}\to0,
 \qquad
 {{\mathfrak j}_\Psi\over\beta}\to0,
 \qquad
 {{\mathfrak o}_\Psi\over\beta}\to0.
 \label{eq:s6v44-DQIS-small}
\end{equation}
Then the generated good-star row satisfies
\begin{equation}
 a_{\rm gen}(\beta)
 \le
 {6\epsilon\over1-\delta_{\epsilon,r}}+o(1).
 \label{eq:s6v44-row-limit}
\end{equation}
For parameters in the V42 window, it is strictly below \(1/17\)
eventually.  If SPFC also holds with a response oscillation per selected
plaquette bounded uniformly, then
\[
 17\left[
 a_{\rm gen}+(1-a_{\rm gen})\widehat p_{\rm star}
 \right]<1
\]
eventually, and the MSEC covariance theorem of V43 applies to the
response-generated quotient law.
\end{corollary}

\begin{proof}
The first two limits give
\[
 {12\epsilon\beta_{\rm eff}\over m_\Psi}
 \longrightarrow
 {12\epsilon\over\lambda}
 ={6\epsilon\over1-\delta_{\epsilon,r}}.
\]
The third limit makes
\({\mathfrak j}_\Psi/m_\Psi\to0\), while the fourth makes
\(\beta_{\rm eff}\Delta-{\mathfrak o}_\Psi\) a positive linear function
of \(\beta\) eventually.  Hence the exponential remainder vanishes.
The V42 parameter choice makes the limit strictly below \(1/17\).
SPFC makes \(\widehat p_{\rm star}\to0\), so the last inequality follows.
Apply V43.
\end{proof}

\begin{corollary}[Bounded finite-range generated terms pass the local
dynamic row]
\label{cor:s6v44-bounded}
If the number of generated terms meeting one star is uniformly bounded,
their coefficients and first two local derivatives are uniformly bounded,
and their collar oscillation is uniformly bounded per selected plaquette,
then
\[
 {\mathfrak h}_\Psi+{\mathfrak j}_\Psi+
 {\mathfrak o}_\Psi=O(1).
\]
Together with the V34 marginal ratio row, these interactions satisfy
\eqref{eq:s6v44-DQIS-small}.
\end{corollary}

\begin{proof}
Finite range and bounded incidence reduce every stock to a uniformly
bounded finite sum of the assumed local derivative or oscillation bounds.
The V34 ratio row gives
\(b^{\rm marg}/\beta\to0\) by
\Cref{cor:s6v44-effective-positive}.
\end{proof}

\begin{assessment}[The first missing generated dynamic writer]
\label{ass:s6v44-frontier}
The plaquette marginal is now removed from the dynamic bottleneck: it is an
effective-coupling shift and cancels from the leading quotient row.
Likewise, any genuinely bounded finite-range generated family passes DQIS.
The unresolved input is the actual cofinal generated action produced by
the exact block pushforward.  The attached manuscripts do not give a
support-local decomposition with uniform bounds on
\({\mathfrak j}_\Psi\), nor do they prove that its collar oscillation is
sublinear in \(\beta\).  These are the first dynamic writers which must be
loaded; the static seven-stock dictionary alone cannot substitute for
them.
\end{assessment}

\begin{programme}[V45: compulsory audit and generated interaction
dictionary]
\label{prog:s6v44-V45}
\begin{enumerate}[leftmargin=3em]
\item Perform the compulsory audit of the then-current SM and NS notes.
\item Search the attached Yang--Mills and Grand Tensor sources for an
      actual finite-range/quasilocal expansion of
      \(\Psi^{\rm gen}\) with first mixed conditional derivatives.
\item Translate every supplied writer into
      \(({\mathfrak h}_\Psi,{\mathfrak j}_\Psi,
      {\mathfrak o}_\Psi)\) without double-counting the plaquette marginal.
\item If no dictionary exists, prove the strongest abstract weighted
      interaction-norm compiler and leave only its numerical population
      open.
\item Keep the physical-clock rate and SPFC occurrence separate from the
      generated dynamic stocks.
\end{enumerate}
\end{programme}

\begin{firewall}[Claim boundary after sixth-series V44]
\label{fire:s6v44-boundary}
V44 proves that response matching by the plaquette marginal is exactly an
effective Wilson-coupling shift and adds no leading quotient influence
when \(\beta+b^{\rm marg}>0\).  It defines the three dynamic
quotient-interaction stocks, proves a generated good-star response bound,
and gives cofinal sufficient conditions under which the response-generated
quotient law satisfies the V43 MSEC covariance theorem, conditional on
SPFC.

V44 does not populate DQIS for the actual generated block action, prove
SPFC occurrence, CPTC, LCKC, selected root asymptotics, TSC, a finite
physical-clock rate, YMGNC, continuum existence, or the full Yang--Mills
mass gap.
\end{firewall}

\begin{theorem}[Sixth-series V44 result ledger]
\label{thm:s6v44-results}
The claims in \Cref{fire:s6v44-boundary} follow from
\Cref{thm:s6v44-effective-beta,
cor:s6v44-effective-positive,
thm:s6v44-generated-response,
cor:s6v44-DQIS-row,
cor:s6v44-bounded}.
\end{theorem}

\begin{proof}
Each ledger item is the conclusion of the cited result.
\end{proof}

\paragraph{Parent-version record.}
V44 was formed from
\texttt{Renewal\_Yang\_Mills\_Mass\_Gap\_Research\_Notes\_V43.tex}.
The V43 parent had \(19\,161\) logical source lines,
\(666\,832\) bytes, and SHA--256
\[
 \texttt{8dc98c24ba7f6ae6d1c10db375af58da0aeba09d4662d790bac5784353239cf3}.
\]
The next compulsory companion audit is V45.

% =============================================================================
% APPENDED SIXTH-SERIES V45 MATERIAL
% =============================================================================

\section{Compulsory companion audit, a weighted generated-interaction
compiler, and the upstream marginal-covariance identity}
\label{sec:s6v45}

Version V44 separated the Wilson plaquette marginal from the genuinely
generated interaction and reduced the good-star dynamic problem to three
stocks: a negative vertical Hessian, a global mixed quotient derivative,
and a collar oscillation.  This note performs the compulsory fifth-version
audit, searches the supplied sources for those exact inputs, and proves two
routes by which they could be loaded.  The first route starts from an
absolutely summable interaction dictionary.  The second differentiates the
fine-shell marginal exactly and identifies the conditional connected
covariances which must be controlled upstream.  Neither route is silently
declared for the physical block action.

\subsection{The V45 compulsory SM/NS audit}
\label{sec:s6v45-audit}

\begin{definition}[Frozen V45 audit packet]
\label{def:s6v45-audit-packet}
The companion files inspected read-only for this audit are
\begin{align}
 &\texttt{Renewal\_SM\_Einstein\_Continuum\_Research\_Notes\_V72.tex},
 \notag\\
 &\qquad 28\,921\ {\rm lines},\quad 1\,146\,496\ {\rm bytes},\quad
 \texttt{43e761121a335ce077dcd86e7054baaf9bfcbe0c7e8d3857f7ccc47945496fd5},
 \label{eq:s6v45-SM-snapshot}\\
 &\texttt{Renewal\_NS\_Stability\_Research\_Notes\_V26.tex},
 \notag\\
 &\qquad 5\,319\ {\rm lines},\quad 278\,283\ {\rm bytes},\quad
 \texttt{82c887f8c2c69e4f28fad7c3dc8de5b9099cb5a4bf431eba1fff3e91935978ad}.
 \label{eq:s6v45-NS-snapshot}
\end{align}
The hashes are SHA--256 digests and freeze the exact audit state even if the
parallel programmes subsequently advance.
\end{definition}

\begin{assessment}[SM V72 lesson]
\label{ass:s6v45-SM}
SM V72 proves positivity of a finite filled-sea Fock transfer after retaining
the fermionic carrier, while SM V71 proves that integrating that carrier out
produces an indefinite determinant Hessian.  The transferable lesson is
architectural: positivity of a joint local transfer must not be confused with
positivity or locality of its logarithmic marginal.  For Yang--Mills this
directs attention to the covariance term generated when fine shell variables
are eliminated.  The SM construction supplies no numerical Yang--Mills
interaction or covariance bound.
\end{assessment}

\begin{assessment}[NS V26 lesson]
\label{ass:s6v45-NS}
NS V26 obtains a smaller tensor norm only after proving that the actual input
is rank one; its first-derivative norm receives no improvement because the
extremal channel is already rank one.  The corresponding warning here is that
one may restrict the mixed quotient norm only after proving that all exterior
gauge tangents occupy the proposed smaller cone.  No such restriction is
available for the eighteen one-link exterior slots, so V44's full operator
norm is retained.
\end{assessment}

\begin{theorem}[V45 audit decision]
\label{thm:s6v45-audit-decision}
The audit changes the order of attack but imports no unproved companion
conclusion.  The generated logarithmic marginal is first differentiated
exactly; its connected covariance is then either bounded by a same-law
conditional mixing estimate or avoided by retaining the detail carrier.
No rank-restricted quotient norm is used without an invariant tangent-cone
proof.
\end{theorem}

\begin{proof}
The first two sentences are the precise lessons of
\Cref{ass:s6v45-SM,ass:s6v45-NS}.  They concern proof architecture rather
than a shared field equation, so they do not populate a Yang--Mills stock.
\end{proof}

\subsection{Source search and its exact negative result}
\label{sec:s6v45-source-search}

\begin{definition}[Frozen source-search packet]
\label{def:s6v45-source-packet}
The two searched source files are
\begin{align}
 &\texttt{Renewal\_Yang\_Mills\_Reduced\_Fibre\_Wilson\_Shell\_Symbol\_and\_}
 \notag\\
 &\qquad
 \texttt{Local\_Vacuum\_Programme\_V35.tex},
 \quad 39\,841\ {\rm lines},\quad 1\,560\,166\ {\rm bytes},
 \notag\\
 &\qquad
 \texttt{7c6056d8149d3015d469b57474a75b23c88a6fe9f64fc81c7dca85c361c647a3},
 \label{eq:s6v45-YM-source}\\
 &\texttt{Renewal\_Grand\_Tensor\_Monograph\_V067.tex},
 \quad 68\,938\ {\rm lines},\quad 2\,804\,854\ {\rm bytes},
 \notag\\
 &\qquad
 \texttt{facd058ec8bb210ad8b296c9d7f4d78a229d08d7a2ae378808700655cb8bc0fe}.
 \label{eq:s6v45-GT-source}
\end{align}
\end{definition}

\begin{proposition}[What the supplied sources do and do not load]
\label{prop:s6v45-source-result}
The supplied Yang--Mills V35 source proves a rooted plaquette-polymer
expansion only on its explicit strong-coupling disk
\[
 |\beta|<\beta_{\rm KP}=0.0012439027\ldots .
\]
It also states explicitly that the complete submitted generated tail has not
been populated in the required interaction norm and that the missing
weak-facing quantity is the law-aligned response of the complete generated
action.  Its V35 head--tail covariance identities bound a declared static
tail after the corresponding moments have been supplied, but do not supply
the moments.

The Grand Tensor V067 source proves abstract weighted-influence and
Dobrushin compilers.  Its directly checkable row assumes bounds on every
mixed local energy and a strict row sum; it does not evaluate those
quantities for the present Yang--Mills block pushforward.  Hence neither
source populates V44's DQIS for the cofinal weak-coupling branch
\(\beta\to\infty\).
\end{proposition}

\begin{proof}
This is a finite textual and formula audit of the exact files frozen in
\Cref{def:s6v45-source-packet}.  The polymer disk and the rooted marked
derivative expansion occur in Yang--Mills V25, whereas its V21 ledger marks
actual generated interaction summability as open and V30 identifies the
law-aligned complete response as missing.  V35 itself proves algebraic
head--tail enclosures conditional on the head, mixed, tail, and deterministic
stocks.  The Grand Tensor Dobrushin theorem starts from an already bounded
influence matrix.  Finally, a disk bounded by \(0.0012439027\ldots\) cannot
contain a sequence tending to \(+\infty\).  Thus none of these statements is
the required weak-coupling loading theorem.
\end{proof}

\subsection{A support-local interaction norm}
\label{sec:s6v45-interaction}

Let \({\mathcal E}\) be the link graph of a finite periodic regulator or its
infinite-volume cover, equipped with graph distance.  Let
\begin{equation}
 \Psi^{\rm gen}(U)=\sum_{X\Subset{\mathcal E}}\psi_X(U_X)
 \label{eq:s6v45-interaction}
\end{equation}
be a gauge-invariant \(C^2\) interaction, where additive constants may be
discarded.  All derivatives below use the fixed bi-invariant metrics and the
coherent quotient derivative of V40.  The positive plaquette marginal
\(b^{\rm marg}\Phi_A\) is excluded from this sum and remains in
\(\beta_{\rm eff}\).

\begin{definition}[Weighted dynamic interaction stocks]
\label{def:s6v45-weighted-stocks}
For \(\mu>0\), define
\begin{align}
 {\cal N}_{0,\mu}
 &:=
 \sup_{e\in{\mathcal E}}
 \sum_{X\ni e}
 e^{\mu\operatorname{diam}X}
 \operatorname{osc}_X\psi_X,
 \label{eq:s6v45-N0}\\
 {\cal N}_{-,\mu}
 &:=
 \sup_{e\in{\mathcal E}}
 \sum_{X\ni e}
 e^{\mu\operatorname{diam}X}
 \sup_{U_X}
 \bigl\|(\nabla_e^2\psi_X)_-\bigr\|_{\rm op},
 \label{eq:s6v45-Nminus}\\
 {\cal N}_{\times,\mu}
 &:=
 \sup_{e\in{\mathcal E}}
 \sum_{f\in{\mathcal E}}
 \sum_{X\ni e,f}
 e^{\mu\operatorname{diam}X}
 \sup_{U_X}
 \|\nabla_e\nabla_f^Q\psi_X\|_{\rm op},
 \label{eq:s6v45-Ncross}\\
 {\cal N}_{{\rm DQ},\mu}
 &:={\cal N}_{0,\mu}+{\cal N}_{-,\mu}
    +{\cal N}_{\times,\mu}.
 \label{eq:s6v45-NDQ}
\end{align}
The definition includes absolute \(C^2\) convergence of the displayed
series.  A finite-range interaction of bounded incidence and bounded local
\(C^2\) norm has finite \({\cal N}_{{\rm DQ},\mu}\).
\end{definition}

\begin{theorem}[Interaction-to-DQIS compiler]
\label{thm:s6v45-interaction-DQIS}
Let \({\cal C}_e\) be the fixed link collar used by one declared star and set
\[
 q_{\rm col}:=\sup_e|{\cal C}_e|<\infty.
\]
For the interaction \eqref{eq:s6v45-interaction}, the V44 stocks obey
\begin{equation}
 \boxed{
 {\mathfrak h}_\Psi\le{\cal N}_{-,\mu},
 \qquad
 {\mathfrak j}_\Psi\le{\cal N}_{\times,\mu},
 \qquad
 {\mathfrak o}_\Psi\le q_{\rm col}{\cal N}_{0,\mu}.}
 \label{eq:s6v45-DQIS-compiler}
\end{equation}
Moreover, if \(\Psi_{\le R}\) is the sum of terms with
\(\operatorname{diam}X\le R\), the same three stocks of the omitted tail are
bounded respectively by
\begin{equation}
 e^{-\mu R}{\cal N}_{-,\mu},\qquad
 e^{-\mu R}{\cal N}_{\times,\mu},\qquad
 q_{\rm col}e^{-\mu R}{\cal N}_{0,\mu}.
 \label{eq:s6v45-tail-stocks}
\end{equation}
\end{theorem}

\begin{proof}
Only terms containing \(e\) contribute to \(\nabla_e^2\Psi^{\rm gen}\).
For Hermitian forms \(H_i\),
\[
 \|( \sum_i H_i)_-\|_{\rm op}
 \le\sum_i\|(H_i)_-\|_{\rm op},
\]
because the minimum eigenvalue is superadditive.  Dropping the exponential
weights in \eqref{eq:s6v45-Nminus} proves the first bound.

Only terms containing both \(e\) and \(f\) contribute to
\(\nabla_e\nabla_f^Q\Psi^{\rm gen}\).  Absolute \(C^2\) convergence permits
termwise differentiation; summing first over \(f\) gives the second bound.

When the variables in a finite collar \({\cal C}\) vary, every interaction
term disjoint from \({\cal C}\) is constant.  Hence
\[
 \operatorname{osc}_{\cal C}\Psi^{\rm gen}
 \le
 \sum_{X:X\cap{\cal C}\ne\varnothing}\operatorname{osc}_X\psi_X
 \le
 \sum_{e\in{\cal C}}\sum_{X\ni e}\operatorname{osc}_X\psi_X
 \le |{\cal C}|{\cal N}_{0,\mu}.
\]
This proves the third bound.  If \(\operatorname{diam}X>R\), then
\(1\le e^{-\mu R}e^{\mu\operatorname{diam}X}\); inserting this inequality
in the same three sums proves \eqref{eq:s6v45-tail-stocks}.
\end{proof}

\begin{lemma}[Exact local perturbation comparison]
\label{lem:s6v45-perturbation}
Let
\[
 \dd\nu_0=Z_0^{-1}e^{-H_0}\dd m,\qquad
 \dd\nu=Z^{-1}e^{-H_0-G}\dd m
\]
be probability measures on one conditional cylinder.  For every event
\(A\),
\begin{equation}
 \boxed{\nu(A)\le
 e^{\operatorname{osc}G}\nu_0(A).}
 \label{eq:s6v45-event-comparison}
\end{equation}
\end{lemma}

\begin{proof}
The numerator of \(\nu(A)\) is at most
\(e^{-\inf G}Z_0\nu_0(A)\), whereas
\(Z\ge e^{-\sup G}Z_0\).  Their ratio is
\eqref{eq:s6v45-event-comparison}.
\end{proof}

\begin{corollary}[Generated conditional defect-path bound]
\label{cor:s6v45-generated-path}
Assume the boundary-fillability hypotheses of V42 and let
\(\Gamma\) be a selected star path.  If the union of its comparison collars
has at most \(q_{\rm path}|\Gamma|\) links, then
\begin{equation}
 \mathbb P_{\rm gen}(\Gamma\ {\rm is\ misaligned}\mid{\cal F}_{\partial})
 \le
 \left[
  6c_H^{-1/6}\beta_{\rm eff}^{2/3}
  \exp\left\{
    {2\over3}-{\beta_{\rm eff}\epsilon^2\over9216}
    +q_{\rm path}{\cal N}_{0,\mu}
  \right\}
 \right]^{|\Gamma|}.
 \label{eq:s6v45-generated-path}
\end{equation}
In particular, if
\begin{equation}
 {\cal N}_{0,\mu}=o(\beta),
 \qquad
 {\beta_{\rm eff}\over\beta}\longrightarrow1,
 \label{eq:s6v45-N0-small}
\end{equation}
the bracket tends to zero exponentially in \(\beta\).
\end{corollary}

\begin{proof}
V41--V42 give the displayed Wilson factor, with \(\beta\) replaced by
\(\beta_{\rm eff}\), for every conditionally fillable selected path.
In the likelihood ratio between the generated and Wilson conditionals,
only interaction terms meeting the path collar survive.  The proof of
\Cref{thm:s6v45-interaction-DQIS} bounds their oscillation by
\(q_{\rm path}|\Gamma|{\cal N}_{0,\mu}\).  Apply
\Cref{lem:s6v45-perturbation} and absorb the result into the
\(|\Gamma|\)-th power.  Under \eqref{eq:s6v45-N0-small}, the negative
linear term in \(\beta\) dominates the sublinear interaction term and the
logarithmic polynomial prefactor.
\end{proof}

\begin{theorem}[Quasilocal generated-action closure of the local dynamic
row]
\label{thm:s6v45-quasilocal-closure}
Fix V42 parameters \((\epsilon,r)\) for which
\begin{equation}
 17\,{6\epsilon\over1-\delta_{\epsilon,r}}<1.
 \label{eq:s6v45-strict-window}
\end{equation}
Suppose along the selected weak-coupling sequence
\begin{equation}
 {b^{\rm marg}\over\beta}\to0,
 \qquad
 {{\cal N}_{{\rm DQ},\mu}\over\beta}\to0
 \label{eq:s6v45-interaction-small}
\end{equation}
for one fixed \(\mu>0\), uniformly in volume and admissible boundary
conditions.  Assume in addition the geometric half of SPFC: every canonical
comparison collar used by the exploration has
\(\rho_\Omega(\xi)\le\rho<\epsilon/(4\sqrt6)\).  Then eventually
\[
 m_\Psi>0,\qquad
 a_{\rm gen}
 \le {6\epsilon\over1-\delta_{\epsilon,r}}+o(1),
\qquad
 \widehat p_{\rm star}^{\rm gen}=o(1),
\]
and the V43 mixed good/bad exploration coefficient satisfies
\begin{equation}
 \boxed{
 17\left[
 a_{\rm gen}+
 (1-a_{\rm gen})\widehat p_{\rm star}^{\rm gen}
 \right]<1.}
 \label{eq:s6v45-MSEC}
\end{equation}
Consequently the finite-regulator quotient Gibbs law has the V43
volume-uniform exponential covariance bound.
\end{theorem}

\begin{proof}
The V34 marginal ratio gives
\(\beta_{\rm eff}/\beta\to1\).  The compiler
\eqref{eq:s6v45-DQIS-compiler} turns
\eqref{eq:s6v45-interaction-small} into every DQIS smallness row of V44.
Thus V44 gives the asserted \(m_\Psi\) and \(a_{\rm gen}\).  The stipulated
fillability and the oscillation estimate from
\Cref{thm:s6v45-interaction-DQIS} supply the two SPFC halves, so the path
estimate is \Cref{cor:s6v45-generated-path}.  Insert both limits in the
left side of \eqref{eq:s6v45-MSEC}; its limit is strictly below one by
\eqref{eq:s6v45-strict-window}.  The V43 connectivity-to-covariance theorem
then applies.
\end{proof}

\begin{assessment}[What this theorem closes]
\label{ass:s6v45-closure-scope}
\Cref{thm:s6v45-quasilocal-closure} removes the \emph{generated oscillation}
half of SPFC as a separate hypothesis whenever the same support-local
interaction norm loads DQIS: its oscillation component transfers the Wilson
path estimate and its two derivative components control the good-star
response.  The geometric vacuum-fillability half remains an independent
assumption.  This is a theorem about a submitted interaction dictionary.  By
\Cref{prop:s6v45-source-result}, the physical block pushforward has not yet
been proved to possess the required uniform norm.
\end{assessment}

\subsection{Going upstream: exact differentiation of the eliminated shell}
\label{sec:s6v45-upstream}

Let \(U\) be a retained one-link variable, \(\zeta\) an exterior quotient
coordinate, \(Y\) the eliminated detail variables, and
\begin{equation}
 Z(U,\zeta)=\int e^{-S(U,\zeta,Y)}\,\dd\lambda(Y),
 \qquad
 \Psi(U,\zeta)=-\log Z(U,\zeta).
 \label{eq:s6v45-marginal}
\end{equation}
Assume the displayed derivatives are dominated so differentiation under the
integral is valid.  Denote expectation under the normalized conditional
detail law by \(\mathbb E_{U,\zeta}\).

\begin{theorem}[Exact marginal first and second derivatives]
\label{thm:s6v45-marginal-derivatives}
For tangent vectors \(u,v\) and an exterior tangent \(\xi\),
\begin{align}
 D_U\Psi[u]
 &=\mathbb E[D_US[u]],
 \label{eq:s6v45-first}\\
 D_U^2\Psi[u,v]
 &=\mathbb E[D_U^2S[u,v]]
   -\operatorname{Cov}(D_US[u],D_US[v]),
 \label{eq:s6v45-Hessian}\\
 D_UD_\zeta\Psi[u,\xi]
 &=\mathbb E[D_UD_\zeta S[u,\xi]]
   -\operatorname{Cov}(D_US[u],D_\zeta S[\xi]).
 \label{eq:s6v45-mixed}
\end{align}
Thus the negative vertical stock and the mixed quotient stock of the
generated logarithmic marginal contain connected conditional covariances of
the eliminated shell.
\end{theorem}

\begin{proof}
From \(D_UZ[u]=-\int D_US[u]e^{-S}\dd\lambda\) one obtains
\eqref{eq:s6v45-first}.  For every observable \(F\),
\[
 D\mathbb E[F]
 =\mathbb E[DF]-\operatorname{Cov}(F,DS).
\]
Apply this identity first with \(F=D_US[u]\) in direction \(v\), and then in
the exterior direction \(\xi\).  This gives
\eqref{eq:s6v45-Hessian} and \eqref{eq:s6v45-mixed}.
\end{proof}

\begin{proposition}[A positive joint action need not give a positive
marginal Hessian]
\label{cth:s6v45-marginal-negative}
There is a compact eliminated fibre for which \(D_U^2S=0\) but
\(D_U^2\Psi<0\).
\end{proposition}

\begin{proof}
Let \(Y\) be any compact probability space carrying a bounded nonconstant
real function \(F\), restrict \(t\) to a compact interval about zero, and set
\(S(t,Y)=S_0(Y)+tF(Y)\).  At \(t=0\),
\[
 \Psi''(0)=-\operatorname{Var}_0(F)<0
\]
by \eqref{eq:s6v45-Hessian}, although \(S_{tt}=0\).  On the compact
\((t,Y)\)-cylinder, adding a sufficiently large constant to \(S\) makes the
joint action nonnegative without changing
either derivative.  Hence joint positivity is irrelevant to the sign of the
logarithmic marginal Hessian.
\end{proof}

\begin{definition}[Conditional-detail mixing card, CDMC]
\label{def:s6v45-CDMC}
Index the eliminated detail variables by \(x\in{\mathcal D}\).  CDMC consists
of a nonnegative kernel \(K(x,y)\), uniform in volume, boundary condition,
and the declared \(U,\zeta\) range, such that
\begin{equation}
 \|\operatorname{Cov}(F,G)\|_{\rm op}
 \le
 \sum_{x,y\in{\mathcal D}}
 K(x,y)L_x(F)L_y(G)
 \label{eq:s6v45-CDMC}
\end{equation}
for the vector-valued score observables below.  Here \(L_x\) is the
coordinate Lipschitz seminorm.  Define
\begin{align}
 H_{\rm dir}
 &:=
 \sup\bigl\|(\mathbb E D_U^2S)_-\bigr\|_{\rm op},
 \qquad
 J_{\rm dir}
 :=
 \sup_e\sum_{f\leadsto e}
 \|\mathbb E D_UD_{\zeta_f}^{Q}S\|_{\rm op},
 \label{eq:s6v45-direct}\\
 A_e(x)&:=L_x(D_US),\qquad
 C_f(y):=L_y(D_{\zeta_f}^{Q}S),
 \label{eq:s6v45-score-Lip}\\
 V_U&:=\sup_e\sum_{x,y}K(x,y)A_e(x)A_e(y),
 \notag\\
 V_\times&:=
 \sup_e\sum_{f\leadsto e}\sum_{x,y}
 K(x,y)A_e(x)C_f(y).
 \label{eq:s6v45-cov-stocks}
\end{align}
All suprema use the same domains as the corresponding V44 stocks; in
particular the mixed row is global in \(U\).
\end{definition}

\begin{theorem}[Conditional-detail loader for the two derivative stocks]
\label{thm:s6v45-CDMC-loader}
Under CDMC,
\begin{equation}
 \boxed{
 {\mathfrak h}_\Psi\le H_{\rm dir}+V_U,
 \qquad
 {\mathfrak j}_\Psi\le J_{\rm dir}+V_\times.}
 \label{eq:s6v45-CDMC-loader}
\end{equation}
If any two collar boundary configurations can be joined by coordinate
geodesics of total length at most \(\ell_{\rm col}\) and
\begin{equation}
 B_{\rm col}
 :=
 \sup_{\rm path}\sum_{f\in{\cal C}}
 \|\mathbb E D_{\zeta_f}S\|<\infty,
 \label{eq:s6v45-Bcol}
\end{equation}
then also
\begin{equation}
 \boxed{{\mathfrak o}_\Psi\le
 \ell_{\rm col}B_{\rm col}.}
 \label{eq:s6v45-CDMC-osc}
\end{equation}
\end{theorem}

\begin{proof}
In \eqref{eq:s6v45-Hessian}, the covariance of a real score with itself is
positive semidefinite.  The norm of that covariance is bounded by CDMC with
\(F=G=D_US\), giving the first inequality.  Apply CDMC to the two scores in
\eqref{eq:s6v45-mixed}, then sum over exterior slots, to obtain the second.

For the last claim, \eqref{eq:s6v45-first} in a collar coordinate gives
\(D_{\zeta_f}\Psi=\mathbb E D_{\zeta_f}S\).  Integrate this derivative along
the stipulated coordinate-geodesic path and use its total length together
with \eqref{eq:s6v45-Bcol}.
\end{proof}

\begin{corollary}[Upstream sufficient row for the V45 closure theorem]
\label{cor:s6v45-upstream-row}
If, uniformly on the selected branch,
\[
 H_{\rm dir}+V_U=o(\beta),\qquad
 J_{\rm dir}+V_\times=o(\beta),\qquad
 \ell_{\rm col}B_{\rm col}=o(\beta),
\]
and the same estimates apply to every selected defect-path collar with a
linear-in-length constant, and if those collars satisfy the geometric
fillability inequality
\(\rho_\Omega(\xi)\le\rho<\epsilon/(4\sqrt6)\), then the conclusions of
\Cref{thm:s6v45-quasilocal-closure} hold without first choosing an
interaction decomposition of \(\Psi\).
\end{corollary}

\begin{proof}
\Cref{thm:s6v45-CDMC-loader} supplies DQIS directly.  The path-collar version
and \Cref{lem:s6v45-perturbation} supply the generated SPFC bound.  The proof
of \Cref{thm:s6v45-quasilocal-closure} then uses only those stocks, not the
particular decomposition \eqref{eq:s6v45-interaction}.
\end{proof}

\begin{assessment}[The new precise upstream frontier]
\label{ass:s6v45-frontier}
The generated-action problem is no longer an undifferentiated request for
\emph{locality}.  There are now two exact admissible routes:
\begin{enumerate}[label=\textup{(\alph*)},leftmargin=3.5em]
\item prove
      \({\cal N}_{{\rm DQ},\mu}=o(\beta)\) for one support-local interaction
      dictionary; or
\item prove CDMC and the three sublinear rows in
      \Cref{cor:s6v45-upstream-row} directly for the conditional eliminated
      shell.
\end{enumerate}
The second route exposes the first real missing estimate:
volume- and boundary-uniform decay of connected fine-shell score
covariances.  If that estimate becomes circular because it presupposes the
coarse mass gap, the SM audit dictates the upstream alternative: retain the
detail variables in a joint positive transfer dilation and prove scale
projectivity there before taking the logarithmic marginal.
\end{assessment}

\begin{programme}[V46: conditional detail mixing before another static
compiler]
\label{prog:s6v45-V46}
\begin{enumerate}[leftmargin=3em]
\item Form the actual fine-detail dependency graph of one block, keeping
      gauge fixing, torons, centre sectors, and boundary links explicit.
\item Combine the V38 aligned-star conditional convexity with the V41--V42
      misalignment-cluster estimate to seek a boundary-uniform CDMC kernel.
\item Prove the covariance inequality
      \eqref{eq:s6v45-CDMC} by a stopped good-cluster exploration; do not use
      the desired coarse covariance theorem inside this proof.
\item Insert the actual local score incidences in
      \eqref{eq:s6v45-cov-stocks}.  If a score is not local or the row diverges,
      return that support as the exact witness and move to the retained-detail
      transfer dilation.
\item Defer further root or static-symbol compilers until one of the two
      routes in \Cref{ass:s6v45-frontier} is populated.
\end{enumerate}
\end{programme}

\begin{firewall}[Claim boundary after sixth-series V45]
\label{fire:s6v45-boundary}
V45 performs the compulsory SM/NS audit and the exact source search; proves
the weighted interaction-to-DQIS compiler and exponential truncation
estimate; transfers the Wilson conditional defect-path bound through a
generated interaction on geometrically fillable collars; proves a
conditional finite-regulator MSEC covariance closure under one explicit
sublinear weighted norm and that fillability condition; derives the exact
vertical and mixed derivatives of the logarithmic block marginal; proves
that their covariance term can destroy positivity; and gives a
conditional-detail mixing loader for all three DQIS stocks.

V45 does not prove the required interaction norm or CDMC for the physical
block pushforward, execute the V32 physical static card, prove CPTC or LCKC,
establish selected-root asymptotics or TSC, obtain a nonzero physical-clock
continuum transfer rate, prove YMGNC, construct the continuum measure, or
prove the full Yang--Mills mass gap.
\end{firewall}

\begin{theorem}[Sixth-series V45 result ledger]
\label{thm:s6v45-results}
Every positive conclusion in \Cref{fire:s6v45-boundary} follows from
\Cref{thm:s6v45-audit-decision,
prop:s6v45-source-result,
thm:s6v45-interaction-DQIS,
lem:s6v45-perturbation,
cor:s6v45-generated-path,
thm:s6v45-quasilocal-closure,
thm:s6v45-marginal-derivatives,
cth:s6v45-marginal-negative,
thm:s6v45-CDMC-loader,
cor:s6v45-upstream-row}.
\end{theorem}

\begin{proof}
Each ledger statement is exactly the conclusion of the cited result, with
the limitations retained from their hypotheses.
\end{proof}

\paragraph{Parent-version record.}
V45 was formed from
\texttt{Renewal\_Yang\_Mills\_Mass\_Gap\_Research\_Notes\_V44.tex}.
The V44 parent had \(19\,544\) logical source lines,
\(680\,399\) bytes, and SHA--256
\[
 \texttt{a58ef4b86d8e2ba521fe2fc5ca357e39d9302c5e2a26a7aa6f278f0d003ef9b7}.
\]
The next compulsory companion audit is V50.

% =============================================================================
% END APPENDED SIXTH-SERIES V45 MATERIAL
% =============================================================================

% =============================================================================
% APPENDED SIXTH-SERIES V46 MATERIAL
% =============================================================================

\section{Conditional-shell mixing, the fillability obstruction, and the
Gaussian Schur owner}
\label{sec:s6v46}

V45 reduced the generated logarithmic marginal to connected covariances of
the eliminated shell.  The present note applies the V43 exploration more
precisely.  It produces an explicit row-summable covariance kernel on every
boundary class for which MSEC is valid, but it also proves that small
vacuum-fillability cannot hold uniformly over arbitrary boundary holonomy.
Finally, an exact Gaussian calculation shows why qualitative mixing alone is
still one power of \(\beta\) too weak: the leading conditional covariance is
the Schur complement owner and must be subtracted or absorbed before the
generated remainder can be \(o(\beta)\).

\subsection{The path-resolved MSEC covariance kernel}
\label{sec:s6v46-kernel}

Let \({\mathcal G}=({\mathcal D},{\mathcal A})\) be a finite or locally
finite graph of maximum degree \(\Delta\).  Assume MSEC with parameters
\(a,p\), and write
\begin{equation}
 \theta=a+(1-a)p,\qquad
 \alpha=(\Delta-1)\theta<1.
 \label{eq:s6v46-alpha}
\end{equation}
For \(x,y\in{\mathcal D}\), define the self-avoiding path kernel
\begin{equation}
 {\cal K}_\theta(x,y)
 :=
 \mathbf1_{\{x=y\}}+
 \sum_{\substack{\gamma:x\to y\ {\rm self\!-\!avoiding}\\|\gamma|\ge1}}
 \theta^{|\gamma|}.
 \label{eq:s6v46-path-kernel}
\end{equation}
The sum is over unoriented graph paths with the displayed endpoints, read in
the direction \(x\to y\).

\begin{lemma}[Row sum and exponential row tail]
\label{lem:s6v46-row-tail}
The kernel \({\cal K}_\theta\) is symmetric and
\begin{align}
 \sup_x\sum_y{\cal K}_\theta(x,y)
 &\le
 1+{\Delta\theta\over1-\alpha},
 \label{eq:s6v46-row-sum}\\
 \sup_x\sum_{y:\,d(x,y)\ge R}{\cal K}_\theta(x,y)
 &\le
 {\Delta\over\Delta-1}{\alpha^R\over1-\alpha},
 \qquad R\ge1.
 \label{eq:s6v46-row-tail}
\end{align}
\end{lemma}

\begin{proof}
Path reversal proves symmetry.  From a fixed root there are at most
\(\Delta(\Delta-1)^{\ell-1}\) self-avoiding paths of length \(\ell\ge1\).
Summing their weights over all endpoints gives
\[
 1+\sum_{\ell\ge1}\Delta(\Delta-1)^{\ell-1}\theta^\ell
 =1+{\Delta\theta\over1-\alpha}.
\]
For endpoints at distance at least \(R\), every contributing path has length
at least \(R\).  Starting the same geometric sum at \(R\) gives
\eqref{eq:s6v46-row-tail}.
\end{proof}

\begin{theorem}[Path-resolved MSEC covariance inequality]
\label{thm:s6v46-path-CDMC}
Suppose the MSEC exploration is valid for every sequential conditional used
to compare two boundary conditions.  Then bounded quotient-Lipschitz
observables \(F,G\) satisfy
\begin{equation}
 \boxed{
 |\operatorname{Cov}(F,G)|
 \le
 \sum_{x,y\in{\mathcal D}}
 {\cal K}_\theta(x,y)
 \delta_x^{\rm Lip}(F)\delta_y^{\rm Lip}(G).}
 \label{eq:s6v46-CDMC}
\end{equation}
Thus \({\cal K}_\theta\) is an explicit CDMC kernel in the sense of
\Cref{def:s6v45-CDMC}.  Unlike the support-distance form of V43, it is
row-summable and may be inserted directly into score-incidence sums.
\end{theorem}

\begin{proof}
Reveal the coordinates on which \(F\) depends and use the same martingale
difference decomposition as in the proof of
\Cref{thm:s6v43-covariance}.  Change one revealed coordinate \(x\).  The
effect reaching a target coordinate \(y\ne x\) is carried by a
self-avoiding disagreement path from \(x\) to \(y\).  By
\Cref{thm:s6v43-path}, the expected metric transmitted along a fixed path
\(\gamma\) is at most \(\theta^{|\gamma|}\).  The union bound over all such
paths therefore gives at most
\({\cal K}_\theta(x,y)\) times the initial normalized displacement.  At
\(y=x\), the trivial bound is one, which is the diagonal term in
\eqref{eq:s6v46-path-kernel}.

The change in the conditional expectation of \(G\) is consequently at most
\[
 \sum_y{\cal K}_\theta(x,y)\delta_y^{\rm Lip}(G)
\]
times the change at \(x\).  The martingale increment of \(F\) is bounded by
\(\delta_x^{\rm Lip}(F)\) times that change.  Sum over revealed \(x\).
This proves \eqref{eq:s6v46-CDMC}.  The proof uses only sequential
conditionals and therefore also applies to vector-valued scores after
polarization.
\end{proof}

\begin{corollary}[Inheritance by a fixed-boundary detail domain]
\label{cor:s6v46-domain}
Let \(\Lambda\subset{\mathcal D}\) be an eliminated-detail domain and freeze
all variables in \(\Lambda^c\).  If the one-site kernels encountered inside
\(\Lambda\) satisfy the same MSEC rows for that frozen exterior and every
later revealed exterior, then its conditional Gibbs law satisfies
\eqref{eq:s6v46-CDMC} with \({\cal K}_\theta\) restricted to \(\Lambda\).
The constants are independent of \(|\Lambda|\).
\end{corollary}

\begin{proof}
The induced graph has degree at most \(\Delta\), and fixing exterior
coordinates deletes paths rather than creating them.  Apply
\Cref{thm:s6v46-path-CDMC}; the row estimates remain those of
\Cref{lem:s6v46-row-tail}.
\end{proof}

\begin{corollary}[Conditional pure-Wilson CDMC on the admissible class]
\label{cor:s6v46-Wilson-CDMC}
Fix parameters in the V42 window.  For all sufficiently large \(\beta\),
every pure-Wilson detail domain whose successive explored exteriors satisfy
the geometric SPFC fillability row has CDMC kernel
\[
 {\cal K}_{\theta_{\rm W}},\qquad
 \theta_{\rm W}
 =a_{\rm W}+(1-a_{\rm W})\widehat p_{\rm star},
\]
and its row sum is bounded uniformly in the domain and volume.
\end{corollary}

\begin{proof}
\Cref{thm:s6v43-SPFC-MSEC} supplies MSEC with
\(\alpha_{\rm W}=17\theta_{\rm W}<1\) on the degree-eighteen graph.
Apply \Cref{cor:s6v46-domain} and then
\Cref{lem:s6v46-row-tail}.
\end{proof}

\subsection{Why arbitrary-boundary fillability is false}
\label{sec:s6v46-fillability}

\begin{lemma}[Boundary holonomy forces a plaquette defect]
\label{lem:s6v46-holonomy}
Let a frozen boundary loop \(C\) bound a simply connected tiled surface
\(\Sigma\) of \(m\) plaquettes in a variable-link collar.  For every filling,
\begin{equation}
 \max_{p\in\Sigma}\|U_p-I\|_{\rm F}
 \ge {\,\|W_C-I\|_{\rm F}\over m},
 \label{eq:s6v46-holonomy}
\end{equation}
where \(W_C\) is the boundary holonomy.
\end{lemma}

\begin{proof}
Choose a base vertex and a spanning tree in the tiled disk.  Cancellation of
interior edges expresses \(W_C\) as an ordered product of the \(m\)
plaquette holonomies, each conjugated by a tree transporter.  For unitaries,
\[
 \|AB-I\|_{\rm F}
 \le\|A-I\|_{\rm F}+\|B-I\|_{\rm F},
\]
and the Frobenius norm is invariant under unitary conjugation.  Iteration
gives
\(\|W_C-I\|_{\rm F}\le\sum_{p\in\Sigma}\|U_p-I\|_{\rm F}\), which implies
\eqref{eq:s6v46-holonomy}.
\end{proof}

\begin{proposition}[No uniform small-fillability theorem over all exteriors]
\label{prop:s6v46-no-uniform-fill}
If the admissible frozen exteriors include a loop with
\(W_C=zI\), \(z=e^{2\pi i/3}\), then
\begin{equation}
 \rho_\Omega(\xi)\ge {3\over m}.
 \label{eq:s6v46-centre-lower}
\end{equation}
Consequently a fixed collar with
\(3/m\ge\epsilon/(4\sqrt6)\) cannot satisfy the V42 strict fillability row
for every boundary configuration.
\end{proposition}

\begin{proof}
For \(SU(3)\),
\[
 \|zI-I\|_{\rm F}
 =\sqrt3\,|z-1|=3.
\]
Apply \Cref{lem:s6v46-holonomy} and then minimize over fillings.  The final
claim is immediate from the V42 requirement
\(\rho_\Omega(\xi)<\epsilon/(4\sqrt6)\).
\end{proof}

\begin{assessment}[The correct role of SPFC]
\label{ass:s6v46-SPFC}
SPFC is not a deterministic all-boundary property.  It can hold on a
protected exterior chart, or its failures can be charged as additional bad
gates under the actual Gibbs law.  Requiring it uniformly over all boundary
holonomies would contradict \Cref{prop:s6v46-no-uniform-fill}.  Hence the
conditional CDMC of \Cref{cor:s6v46-Wilson-CDMC} is rigorous but restricted
to the admissible class; an unconditional marginal theorem needs a
probabilistic fillability estimate.
\end{assessment}

\begin{lemma}[Absorbing probabilistic nonfillability into MSEC]
\label{lem:s6v46-prob-fill}
Let \(M_e\) be the ordinary misalignment indicator and \(N_e\) the indicator
that the required fillability row fails at gate \(e\).  Suppose that for
every finite gate set \(S\) and every \(A\subseteq S\),
\begin{equation}
 \mathbb E\left[
  \prod_{e\in A}M_e\prod_{e\in S\setminus A}N_e
 \right]
 \le p_M^{|A|}p_N^{|S\setminus A|}.
 \label{eq:s6v46-mixed-MDSC}
\end{equation}
Then the enlarged bad indicator
\(B_e=\mathbf1_{\{M_e=1\ {\rm or}\ N_e=1\}}\) satisfies
\begin{equation}
 \mathbb E\prod_{e\in S}B_e
 \le(p_M+p_N)^{|S|}.
 \label{eq:s6v46-combined-MDSC}
\end{equation}
If good gates retain contraction \(a\), all V43 path and covariance results
hold with \(p=\min\{1,p_M+p_N\}\).
\end{lemma}

\begin{proof}
Use \(B_e\le M_e+N_e\), expand the product, apply
\eqref{eq:s6v46-mixed-MDSC} term by term, and sum the binomial series.  The
last assertion is exactly the MSEC definition with the enlarged bad event.
\end{proof}

\begin{definition}[Probabilistic fillability stock, PFS]
\label{def:s6v46-PFS}
PFS is the smallest \(p_N\) for which
\eqref{eq:s6v46-mixed-MDSC} holds jointly with the V42 misalignment stock on
every predeclared exploration.  It is a same-law, boundary-holonomy stock;
a one-collar marginal probability without the mixed moment bound does not
populate PFS.
\end{definition}

\begin{corollary}[A nonuniform route to conditional-shell CDMC]
\label{cor:s6v46-PFS-CDMC}
If PFS holds and
\begin{equation}
 (\Delta-1)\left[
 a+(1-a)\min\{1,p_M+p_N\}
 \right]<1,
 \label{eq:s6v46-PFS-margin}
\end{equation}
then the full conditional detail law has the path kernel
\eqref{eq:s6v46-path-kernel} and CDMC, without deterministic fillability for
every exterior.
\end{corollary}

\begin{proof}
Apply \Cref{lem:s6v46-prob-fill} and then
\Cref{thm:s6v46-path-CDMC}.
\end{proof}

\subsection{The scale audit: mixing is not yet a sublinear generated stock}
\label{sec:s6v46-scale}

\begin{proposition}[What the MSEC kernel gives for Wilson-sized scores]
\label{prop:s6v46-beta-count}
Let \(F_\beta,G_\beta\) be local score observables with
\[
 \sum_x\delta_x^{\rm Lip}(F_\beta)\le C_F\beta,
 \qquad
 \sum_y\delta_y^{\rm Lip}(G_\beta)\le C_G\beta.
\]
Then \eqref{eq:s6v46-CDMC} gives only
\begin{equation}
 |\operatorname{Cov}(F_\beta,G_\beta)|
 \le
 \left(1+{\Delta\theta\over1-\alpha}\right)
 C_FC_G\beta^2.
 \label{eq:s6v46-beta2}
\end{equation}
Even a strong-convexity covariance inequality with conditional Hessian
floor \(\kappa\beta\) and score gradients of order \(\beta\) gives in
general only \(O(\beta)\).  Neither conclusion is \(o(\beta)\).
\end{proposition}

\begin{proof}
The double sum in \eqref{eq:s6v46-CDMC} is bounded by the kernel row norm
times the two \(\ell^1\) Lipschitz totals, giving
\eqref{eq:s6v46-beta2}.  On a convex chart, the Brascamp--Lieb inequality
has the form
\[
 |\operatorname{Cov}(F_\beta,G_\beta)|
 \le {1\over\kappa\beta}
 \mathbb E\|\nabla F_\beta\|\,\|\nabla G_\beta\|.
\]
Gradients of order \(\beta\) make the right side order \(\beta\).  The final
statement follows.
\end{proof}

\begin{theorem}[Exact Gaussian Schur owner]
\label{thm:s6v46-Gaussian-Schur}
Let \(u\in\mathbb R^r\), \(y\in\mathbb R^n\), \(C\succ0\), and
\begin{equation}
 S_\beta(u,y)
 ={\,\beta\over2}
 \left(u^TAu+2u^TBy+y^TCy\right).
 \label{eq:s6v46-quadratic-action}
\end{equation}
Then, up to a constant independent of \(u\),
\begin{equation}
 -\log\int_{\mathbb R^n}e^{-S_\beta(u,y)}\,\dd y
 ={\,\beta\over2}
 u^T(A-BC^{-1}B^T)u.
 \label{eq:s6v46-Schur}
\end{equation}
In the exact marginal identity of V45,
\begin{equation}
 \mathbb E[D_u^2S_\beta]=\beta A,
 \qquad
 \operatorname{Cov}(D_uS_\beta,D_uS_\beta)
 =\beta BC^{-1}B^T.
 \label{eq:s6v46-two-beta-terms}
\end{equation}
Thus the conditional law is perfectly mixing and strongly log-concave, but
both constituents of the marginal Hessian are order \(\beta\).  If the
Schur form in \eqref{eq:s6v46-Schur} is assigned to the declared Gaussian
or plaquette marginal owner, the genuinely generated quadratic remainder is
exactly zero.
\end{theorem}

\begin{proof}
Complete the square:
\[
 y^TCy+2u^TBy
 =(y+C^{-1}B^Tu)^TC(y+C^{-1}B^Tu)
  -u^TBC^{-1}B^Tu.
\]
Translation invariance of Lebesgue measure gives
\eqref{eq:s6v46-Schur}.  Under the conditional Gaussian,
\[
 \mathbb E[y]=-C^{-1}B^Tu,\qquad
 \operatorname{Cov}(y,y)=(\beta C)^{-1}.
\]
Since \(D_uS_\beta=\beta(Au+By)\), its covariance is
\(\beta BC^{-1}B^T\), proving
\eqref{eq:s6v46-two-beta-terms}.  Subtracting the complete Schur owner leaves
zero.
\end{proof}

\begin{proposition}[Qualitative CDMC cannot identify the generated owner]
\label{prop:s6v46-CDMC-no-scale}
There is no implication
\[
 \text{uniform CDMC and a }\kappa\beta\text{ detail Hessian floor}
 \quad\Longrightarrow\quad
 {\mathfrak h}_\Psi+{\mathfrak j}_\Psi=o(\beta).
\]
\end{proposition}

\begin{proof}
In \Cref{thm:s6v46-Gaussian-Schur}, choose fixed matrices with
\(A-BC^{-1}B^T\ne0\).  The conditional Gaussian has a volume-independent
Poincar\'e constant \(O(\beta^{-1})\) and zero long-range influence, yet its
marginal Hessian is the nonzero order-\(\beta\) matrix in
\eqref{eq:s6v46-Schur}.  Therefore qualitative mixing and strong convexity
do not imply the displayed sublinear conclusion.
\end{proof}

\begin{definition}[Gaussian Schur owner matching card, GSOMC]
\label{def:s6v46-GSOMC}
At the selected vacuum and scale, let \(A_M,B_M,C_M\) be the retained,
mixed, and detail blocks of the gauge-quotiented quadratic Wilson Hessian.
GSOMC requires:
\begin{enumerate}[label=\textup{(\roman*)},leftmargin=3.5em]
\item \(C_M\) is positive on the pinned detail quotient and
      \(B_MC_M^{-1}B_M^*\) is formed on that same carrier;
\item the full quadratic Schur form
      \(\beta(A_M-B_MC_M^{-1}B_M^*)\), including every non-plaquette
      exponentially local component, is assigned once to the baseline
      marginal action;
\item the response coefficient \(b^{\rm marg}\) is used for this owner only
      to the extent proved by the response-matching identity; and
\item DQIS and CDMC are applied to the residual logarithmic marginal after
      this quadratic owner is removed.
\end{enumerate}
\end{definition}

\begin{corollary}[Quadratic exactness of GSOMC]
\label{cor:s6v46-GSOMC}
For a purely quadratic block action, GSOMC makes every residual DQIS stock
zero.  For a nonlinear block action, any \(o(\beta)\) claim must control the
difference between the exact connected score covariance and its Gaussian
Schur value, not either order-\(\beta\) term separately.
\end{corollary}

\begin{proof}
The first statement is \Cref{thm:s6v46-Gaussian-Schur}.  The second follows
by subtracting \eqref{eq:s6v46-two-beta-terms} from the exact V45 marginal
identity before estimating the remainder.
\end{proof}

\begin{assessment}[V46 bottleneck after going upstream]
\label{ass:s6v46-frontier}
The conditional mixing question has split into two independent rows.
First, arbitrary exterior fillability is false; one needs either a protected
boundary chart or PFS.  Second, even successful MSEC/CDMC controls the size
of conditional covariances but not their order-\(\beta\) Gaussian
cancellation.  The next useful object is therefore a \emph{renormalized}
conditional perturbation theorem comparing the exact marginal Hessian and
mixed response with the GSOMC Gaussian Schur owner.  This is the
noncircular content required from CPTC/LCKC.
\end{assessment}

\begin{programme}[V47: renormalized conditional perturbation]
\label{prog:s6v46-V47}
\begin{enumerate}[leftmargin=3em]
\item Write the gauge-pinned fine block as its vacuum quadratic form plus a
      nonlinear remainder on the aligned chart.
\item Use the V37 exponential inverse bound for \(C_M^{-1}\) to record the
      exact GSOMC kernel and its support tail.
\item Prove a finite-dimensional Laplace--Schur remainder bound, with all
      dimension dependence explicit, for the exact marginal second and mixed
      derivatives after subtracting the Gaussian owner.
\item Combine the good-chart estimate with the V41 defect-cluster tail and
      PFS; do not invoke the desired coarse covariance decay in this step.
\item Test whether the resulting debit is \(O(M^4)\).  On the existing
      corridor \(M=o(\beta^{1/4})\), such a debit would be \(o(\beta)\) and
      would populate the V45 local dynamic row.
\end{enumerate}
\end{programme}

\begin{firewall}[Claim boundary after sixth-series V46]
\label{fire:s6v46-boundary}
V46 proves a path-resolved, row-summable MSEC covariance kernel and its
inheritance by fixed-boundary detail domains; proves a boundary-holonomy
lower bound which rules out uniform small fillability over arbitrary
exteriors; introduces the same-law probabilistic fillability stock and its
exact absorption into MSEC; proves that unrenormalized MSEC and
Brascamp--Lieb estimates are not sublinear in \(\beta\); computes the exact
Gaussian marginal Schur owner; and formulates GSOMC.

V46 does not populate PFS, prove GSOMC for the physical response-matched
action, bound the nonlinear renormalized covariance remainder, prove the
physical generated DQIS, CPTC, LCKC, selected-root asymptotics, TSC, YMGNC,
continuum existence, a finite physical-clock mass, or the full Yang--Mills
mass gap.
\end{firewall}

\begin{theorem}[Sixth-series V46 result ledger]
\label{thm:s6v46-results}
The positive statements in \Cref{fire:s6v46-boundary} follow from
\Cref{lem:s6v46-row-tail,
thm:s6v46-path-CDMC,
cor:s6v46-domain,
cor:s6v46-Wilson-CDMC,
lem:s6v46-holonomy,
prop:s6v46-no-uniform-fill,
lem:s6v46-prob-fill,
cor:s6v46-PFS-CDMC,
prop:s6v46-beta-count,
thm:s6v46-Gaussian-Schur,
prop:s6v46-CDMC-no-scale,
cor:s6v46-GSOMC}.
\end{theorem}

\begin{proof}
Each item is exactly the conclusion of the cited result.
\end{proof}

\paragraph{Parent-version record.}
V46 was formed from
\texttt{Renewal\_Yang\_Mills\_Mass\_Gap\_Research\_Notes\_V45.tex}.
The V45 parent had \(20\,183\) logical source lines,
\(704\,720\) bytes, and SHA--256
\[
 \texttt{f0a666df5fd9d6b5515bce23fac291b5fb25c95d3048557426b07ef401ab1b45}.
\]
The next compulsory companion audit is V50.

% =============================================================================
% END APPENDED SIXTH-SERIES V46 MATERIAL
% =============================================================================

% =============================================================================
% APPENDED SIXTH-SERIES V47 MATERIAL
% =============================================================================

\section{Renormalized Laplace--Schur control and the block-volume corridor}
\label{sec:s6v47}

V46 showed that the order-\(\beta\) Gaussian Schur form must be assigned to
the baseline marginal owner before the generated remainder is estimated.
This note carries out the corresponding nonlinear calculation.  It first
separates the mixed derivative needed inside the convex ball from the much
weaker amplitude needed on its exponentially unlikely complement.  It then
proves an exact perturbation theorem around the Gaussian Schur owner and an
\(O(M^4)\) compiler under explicit local moment stocks.

\subsection{Interior and complement mixed-derivative stocks}
\label{sec:s6v47-split-stock}

\begin{definition}[Split mixed quotient stocks]
\label{def:s6v47-split-j}
For the V44 generated remainder define
\begin{align}
 {\mathfrak j}_{\Psi}^{\rm in}
 &:=
 \sup_{\rm good\ star}
 \sum_{f\leadsto e}
 \sup_{U\in{\cal B}_r}
 \|D_UD_{\zeta_f}^{Q}\Psi^{\rm gen}\|_{\rm op},
 \label{eq:s6v47-j-in}\\
 {\mathfrak j}_{\Psi}^{\rm all}
 &:=
 \sup_{\rm good\ star}
 \sum_{f\leadsto e}
 \sup_{U\in SU(3)}
 \|D_UD_{\zeta_f}^{Q}\Psi^{\rm gen}\|_{\rm op}.
 \label{eq:s6v47-j-all}
\end{align}
The V44 stock is
\({\mathfrak j}_{\Psi}^{\rm all}\), and
\({\mathfrak j}_{\Psi}^{\rm in}\le
{\mathfrak j}_{\Psi}^{\rm all}\).
\end{definition}

\begin{theorem}[Split-stock generated response bound]
\label{thm:s6v47-split-response}
With
\[
 m_\Psi=\beta_{\rm eff}\lambda-{\mathfrak h}_\Psi>0,
\]
the V44 response theorem remains valid in the sharper form
\begin{align}
 \sum_{f_0\leadsto e}
 \left|
 D_{\zeta_{f_0}}\mathbb E_{\cal U}f
 \right|
 \le{}&
 \left\{
 {12\epsilon\beta_{\rm eff}\over m_\Psi}
 +{{\mathfrak j}_{\Psi}^{\rm in}\over m_\Psi}
 \right\}\|\nabla f\|_\infty
 \notag\\
 &+
 C\bigl(\beta_{\rm eff}
        +{\mathfrak j}_{\Psi}^{\rm all}
        +{\mathfrak o}_\Psi\bigr)
 \exp\left\{
 -{1\over2}
 (\beta_{\rm eff}\Delta-{\mathfrak o}_\Psi)
 \right\}
 \{\|\nabla f\|_\infty+\operatorname{osc}f\}.
 \label{eq:s6v47-split-response}
\end{align}
Consequently the leading good-star row needs only
\[
 {\mathfrak h}_\Psi=o(\beta),\qquad
 {\mathfrak j}_{\Psi}^{\rm in}=o(\beta),
 \qquad
 {\mathfrak o}_\Psi=o(\beta),
 \label{eq:s6v47-leading-small}
\]
while the complement term vanishes provided additionally
\begin{equation}
 \log\bigl(
  2+\beta_{\rm eff}
  +{\mathfrak j}_{\Psi}^{\rm all}
  +{\mathfrak o}_\Psi
 \bigr)=o(\beta).
 \label{eq:s6v47-tail-subexp}
\end{equation}
\end{theorem}

\begin{proof}
Repeat the proof of \Cref{thm:s6v44-generated-response}.  The
Bakry--\'Emery covariance is taken entirely inside \({\cal B}_r\), so its
generated insertion gradient is bounded by
\({\mathfrak j}_{\Psi}^{\rm in}\).  In the full-versus-restricted
comparison, the insertion oscillation on \(SU(3)\) is bounded by the
geodesic fundamental theorem of calculus using
\({\mathfrak j}_{\Psi}^{\rm all}\).  This gives
\eqref{eq:s6v47-split-response}.

Under \eqref{eq:s6v47-leading-small},
\(\beta_{\rm eff}\Delta-{\mathfrak o}_\Psi\) is a positive linear function
of \(\beta\) eventually.  Taking logarithms of the prefactor and using
\eqref{eq:s6v47-tail-subexp} shows that the second line tends to zero.
\end{proof}

\begin{assessment}[Why the split matters]
\label{ass:s6v47-split}
The global stock introduced in the V44 correction was sufficient but
stronger than the leading response calculation needs.  V47 retains it only
as a complement-amplitude stock, for which subexponential growth is enough.
All \(o(\beta)\) cancellation work may now be confined to the aligned convex
chart where the Gaussian recentering and conditional concentration are
available.
\end{assessment}

\subsection{A fixed-fibre perturbation of the Gaussian Schur owner}
\label{sec:s6v47-RLS}

Let \(q\in{\mathcal Q}\subset\mathbb R^d\) collect the retained one-link and
exterior quotient coordinates.  Let \(y\in\mathbb R^n\) be pinned detail
coordinates, and take
\begin{equation}
 S_\beta(q,y)
 ={\,\beta\over2}
 \left(q^TAq+2q^TBy+y^TCy\right)+R(q,y),
 \qquad C\succ0.
 \label{eq:s6v47-action}
\end{equation}
Put
\[
 m(q)=-C^{-1}B^Tq,\qquad
 \xi=y-m(q),\qquad
 \widetilde R(q,\xi)=R(q,m(q)+\xi).
\]
The Schur-transported derivative in a retained direction \(v\) is
\begin{equation}
 {\cal D}_vR
 :=
 D_vR-D_yR[C^{-1}B^Tv]
 =D_v\widetilde R.
 \label{eq:s6v47-transported}
\end{equation}
Assume the recentered fibre \(\Xi\) and reference measure are independent of
\(q\).  This holds on \(\mathbb R^n\); on a compact group chart it is a
separate recentering requirement.

Let \(\gamma_{\beta,C}\) be the normalized Gaussian measure with density
proportional to
\(\exp\{-\beta\xi^TC\xi/2\}\), restricted to \(\Xi\) when appropriate, and
define
\begin{equation}
 G_\beta(q)
 :=
 -\log\int_\Xi e^{-\widetilde R(q,\xi)}
 \,\gamma_{\beta,C}(\dd\xi).
 \label{eq:s6v47-G}
\end{equation}

\begin{lemma}[Exact Schur-owner factorization]
\label{lem:s6v47-factorization}
Up to a constant independent of \(q\), the marginal action is
\begin{equation}
 \Psi_\beta(q)
 ={\,\beta\over2}q^T(A-BC^{-1}B^T)q+G_\beta(q).
 \label{eq:s6v47-factorization}
\end{equation}
\end{lemma}

\begin{proof}
Complete the square exactly as in
\Cref{thm:s6v46-Gaussian-Schur}, translate \(y=m(q)+\xi\), and use the
fixed-fibre hypothesis.  The Gaussian normalizer is independent of \(q\);
the remaining logarithm is \eqref{eq:s6v47-G}.
\end{proof}

Let \(\nu_q\) be the probability measure on \(\Xi\) proportional to
\[
 \exp\left\{
 -{\beta\over2}\xi^TC\xi-\widetilde R(q,\xi)
 \right\}\dd\xi .
\]
For a retained unit direction \(v\), set
\begin{align}
 r_{vw}
 &:=
 \sup_{q\in{\mathcal Q}}
 \left\|
  \mathbb E_{\nu_q}
  D_vD_w\widetilde R
 \right\|_{\rm op},
 \label{eq:s6v47-rvw}\\
 g_v^2
 &:=
 \sup_{q\in{\mathcal Q}}
 \mathbb E_{\nu_q}
 \left\|
  \nabla_\xi D_v\widetilde R
 \right\|^2.
 \label{eq:s6v47-gv}
\end{align}

\begin{theorem}[Renormalized Laplace--Schur remainder bound]
\label{thm:s6v47-RLS}
Assume \(\Xi\) is convex and the recentered conditional potential satisfies
\begin{equation}
 \nabla_\xi^2
 \left(
  {\beta\over2}\xi^TC\xi+\widetilde R(q,\xi)
 \right)
 \succeq\kappa\beta I
 \label{eq:s6v47-convexity}
\end{equation}
uniformly in \(q\), for some \(\kappa>0\).  Then
\begin{align}
 D_vG_\beta
 &=\mathbb E_{\nu_q}D_v\widetilde R,
 \label{eq:s6v47-Gfirst}\\
 D_vD_wG_\beta
 &=
 \mathbb E_{\nu_q}D_vD_w\widetilde R
 -\operatorname{Cov}_{\nu_q}
  (D_v\widetilde R,D_w\widetilde R),
 \label{eq:s6v47-Gsecond}\\
 \|D_vD_wG_\beta\|_{\rm op}
 &\le
 r_{vw}+{g_vg_w\over\kappa\beta}.
 \label{eq:s6v47-RLS-bound}
\end{align}
The same result holds on a geodesically convex chart for which the
Brascamp--Lieb inequality associated with
\eqref{eq:s6v47-convexity} is valid.
\end{theorem}

\begin{proof}
The first two identities follow by differentiating the normalized integral
\eqref{eq:s6v47-G}; they are the V45 marginal identities applied after the
Gaussian owner has been removed.  Strong convexity and the
Brascamp--Lieb inequality give, for scalar polarizations \(X,Y\),
\[
 |\operatorname{Cov}_{\nu_q}(X,Y)|
 \le {1\over\kappa\beta}
 \left(\mathbb E_{\nu_q}\|\nabla_\xi X\|^2\right)^{1/2}
 \left(\mathbb E_{\nu_q}\|\nabla_\xi Y\|^2\right)^{1/2}.
\]
Take \(X=D_v\widetilde R\), \(Y=D_w\widetilde R\), polarize the retained
and target components, and use
\eqref{eq:s6v47-rvw}--\eqref{eq:s6v47-gv}.  This proves
\eqref{eq:s6v47-RLS-bound}.
\end{proof}

\begin{lemma}[Oscillation of a logarithmic perturbation]
\label{lem:s6v47-log-osc}
For two retained configurations \(q,q'\),
\begin{equation}
 |G_\beta(q)-G_\beta(q')|
 \le
 \sup_{\xi\in\Xi}
 |\widetilde R(q,\xi)-\widetilde R(q',\xi)|.
 \label{eq:s6v47-log-osc}
\end{equation}
\end{lemma}

\begin{proof}
If the supremum on the right is \(\omega\), then pointwise
\[
 e^{-\omega}e^{-\widetilde R(q',\xi)}
 \le e^{-\widetilde R(q,\xi)}
 \le e^\omega e^{-\widetilde R(q',\xi)}.
\]
Integrate against \(\gamma_{\beta,C}\), take negative logarithms, and obtain
\eqref{eq:s6v47-log-osc}.
\end{proof}

\subsection{The explicit block-volume compiler}
\label{sec:s6v47-volume}

Split retained directions into a one-link direction \(u\) and the exterior
quotient slots \(\zeta_f\).  Define
\begin{align}
 R_{2,U}
 &:=
 \sup_{|u|=1}r_{uu},
 &
 G_U^2
 &:=
 \sup_{|u|=1}g_u^2,
 \label{eq:s6v47-U-stocks}\\
 R_{2,\times}
 &:=
 \sup_{|u|=1}
 \sum_{f\leadsto e}\sup_{|\zeta_f|=1}r_{u\zeta_f},
 &
 G_\times
 &:=
 \sup_{|u|=1}
 \sum_{f\leadsto e}\sup_{|\zeta_f|=1}g_ug_{\zeta_f}.
 \label{eq:s6v47-cross-stocks}
\end{align}
Let \(O_{\rm col}\) be the right side of
\eqref{eq:s6v47-log-osc}, maximized over each declared collar pair.

\begin{theorem}[RLS-to-DQIS compiler]
\label{thm:s6v47-RLS-DQIS}
Assume GSOMC, so the quadratic term in
\eqref{eq:s6v47-factorization} belongs to the baseline owner and
\(G_\beta\) is the generated remainder.  On the aligned convex chart,
\begin{equation}
 \boxed{
 {\mathfrak h}_{G}
 \le R_{2,U}+{G_U^2\over\kappa\beta},
 \qquad
 {\mathfrak j}_{G}^{\rm in}
 \le R_{2,\times}+{G_\times\over\kappa\beta},
 \qquad
 {\mathfrak o}_{G}\le O_{\rm col}.}
 \label{eq:s6v47-RLS-DQIS}
\end{equation}
\end{theorem}

\begin{proof}
The first inequality follows from
\(\|(D_u^2G)_-\|\le\|D_u^2G\|\) and
\eqref{eq:s6v47-RLS-bound} with \(v=w=u\).  Sum the same estimate with
\(v=u,w=\zeta_f\) to obtain the second.  The third is
\Cref{lem:s6v47-log-osc}.
\end{proof}

\begin{corollary}[Block-volume sublinearity]
\label{cor:s6v47-volume}
Let \(N_M\le c_{\rm vol}M^4\) be the number of local fine cells in a
four-dimensional block.  Suppose, uniformly in the aligned chart, volume,
and admissible exterior,
\begin{align}
 R_{2,U}+R_{2,\times}+O_{\rm col}
 &\le c_2N_M,
 \label{eq:s6v47-local-second}\\
 G_U^2+G_\times
 &\le c_1\beta N_M,
 \label{eq:s6v47-local-gradient}\\
 \log\bigl(2+{\mathfrak j}_{G}^{\rm all}\bigr)
 &\le c_3N_M+c_4\log(2+\beta).
 \label{eq:s6v47-global-growth}
\end{align}
If
\begin{equation}
 M=o(\beta^{1/4}),
 \label{eq:s6v47-corridor}
\end{equation}
then
\begin{equation}
 {\mathfrak h}_{G}
 +{\mathfrak j}_{G}^{\rm in}
 +{\mathfrak o}_{G}
 =o(\beta),
 \qquad
 \log(2+{\mathfrak j}_{G}^{\rm all})=o(\beta).
 \label{eq:s6v47-sublinear}
\end{equation}
Thus the generated remainder passes the split-stock response row of
\Cref{thm:s6v47-split-response}.
\end{corollary}

\begin{proof}
\Cref{thm:s6v47-RLS-DQIS} and
\eqref{eq:s6v47-local-second}--\eqref{eq:s6v47-local-gradient} bound the
three leading stocks by \(CN_M\).  The growth row has the same order, up to
\(\log(2+\beta)=o(\beta)\).  Since
\(N_M=O(M^4)=o(\beta)\) by
\eqref{eq:s6v47-corridor}, every conclusion in
\eqref{eq:s6v47-sublinear} follows.  Apply
\Cref{thm:s6v47-split-response}.
\end{proof}

\begin{assessment}[Relation to the earlier corridor]
\label{ass:s6v47-corridor}
The exponent \(1/4\) is not inserted retroactively.  V36 obtained the same
corridor from the deterministic retained-writer debit.  V47 shows that a
local nonlinear Laplace--Schur remainder with one \(O(1)\) debit per
four-dimensional fine cell is compatible with that corridor as well.  A
larger \(O(M^p)\) packet would instead require \(M=o(\beta^{1/p})\).
\end{assessment}

\begin{assessment}[What remains to populate]
\label{ass:s6v47-frontier}
The RLS theorem is exact, but the physical Wilson block still needs four
model-specific verifications:
\begin{enumerate}[label=\textup{(\roman*)},leftmargin=3.5em]
\item a fixed-fibre coherent chart after the Gaussian conditional mean is
      removed;
\item the uniform strong-convexity row
      \eqref{eq:s6v47-convexity} on its good component;
\item the local moment bounds
      \eqref{eq:s6v47-local-second}--\eqref{eq:s6v47-local-gradient} for
      the Taylor remainder; and
\item the global subexponential complement amplitude
      \eqref{eq:s6v47-global-growth}, together with PFS or a protected
      exterior chart.
\end{enumerate}
These are strictly upstream of another static trace-shell evaluation and do
not use the desired continuum mass gap.
\end{assessment}

\begin{programme}[V48: the compact Wilson Taylor packet]
\label{prog:s6v47-V48}
\begin{enumerate}[leftmargin=3em]
\item Use exponential coordinates on each pinned fine link and write the
      Wilson plaquette action through its exact quadratic part plus an
      integral Taylor remainder.
\item Bound the transported first, second, and fine-gradient derivatives of
      that remainder by local incidence, keeping the powers of the chart
      radius, \(\beta\), and \(N_M\) explicit.
\item Use the aligned-star Hessian of V38 to choose a shrinking chart radius
      for which \eqref{eq:s6v47-convexity} holds, and compare its complement
      with the V41--V42 defect exponent.
\item Determine whether the shrinking-radius moments give the
      \(O(N_M)\) rows of
      \eqref{eq:s6v47-local-second}--\eqref{eq:s6v47-local-gradient}.
      If not, record the exact excess power and revise the \(M\)-corridor.
\item Keep the nonfillable boundary event separate as PFS; do not replace it
      by an all-boundary assertion excluded in V46.
\end{enumerate}
\end{programme}

\begin{firewall}[Claim boundary after sixth-series V47]
\label{fire:s6v47-boundary}
V47 proves the interior/complement split of the generated quotient response,
reducing the global mixed-derivative requirement to subexponential growth;
proves the exact fixed-fibre Gaussian-owner factorization; proves the
renormalized Laplace--Schur second-derivative bound by Brascamp--Lieb;
proves the logarithmic perturbation oscillation bound; and proves that the
explicit local moment packet
\eqref{eq:s6v47-local-second}--\eqref{eq:s6v47-global-growth} loads all
generated dynamic stocks on \(M=o(\beta^{1/4})\).

V47 does not verify that packet for the compact physical Wilson block,
populate PFS or GSOMC, prove the physical generated DQIS, CPTC, LCKC,
selected-root asymptotics, TSC, YMGNC, continuum existence, a finite
physical-clock mass, or the full Yang--Mills mass gap.
\end{firewall}

\begin{theorem}[Sixth-series V47 result ledger]
\label{thm:s6v47-results}
The positive conclusions in \Cref{fire:s6v47-boundary} follow from
\Cref{thm:s6v47-split-response,
lem:s6v47-factorization,
thm:s6v47-RLS,
lem:s6v47-log-osc,
thm:s6v47-RLS-DQIS,
cor:s6v47-volume}.
\end{theorem}

\begin{proof}
Each ledger statement is the corresponding cited conclusion.
\end{proof}

\paragraph{Parent-version record.}
V47 was formed from
\texttt{Renewal\_Yang\_Mills\_Mass\_Gap\_Research\_Notes\_V46.tex}.
The V46 parent had \(20\,676\) logical source lines,
\(723\,148\) bytes, and SHA--256
\[
 \texttt{1f8916a969f8118299c04b9eb461cf7525f5ce322c4426c372f714eaed740ce9}.
\]
The next compulsory companion audit is V50.

% =============================================================================
% END APPENDED SIXTH-SERIES V47 MATERIAL
% =============================================================================

% =============================================================================
% APPENDED SIXTH-SERIES V48 MATERIAL
% =============================================================================

\section{Parametric Morse reduction, classical minimization, and the
entropy owner}
\label{sec:s6v48}

The direct Taylor route proposed at the end of V47 has a structural
weakness: in link exponential coordinates, the plaquette holonomy contains
quadratic Baker--Campbell--Hausdorff terms, so a supremum estimate on the
cubic remainder can retain a spurious factor of \(\beta\).  This note replaces
that route by an exact local quadraticization.  The block marginal then
splits into an order-\(\beta\) classical minimized action, an entropy term
coming from the Morse Jacobian, and an exponentially small chart-complement
term.  Only the latter two should be estimated as fluctuation corrections.

\subsection{Exact parametric quadraticization of the detail fibre}
\label{sec:s6v48-Morse}

Let \(q\) range in a finite-dimensional retained chart \({\mathcal Q}\), let
\(y\) range in a pinned detail chart, and let
\(\phi(q,y)\) be \(C^4\).  Suppose that for every \(q\) in the declared
retained chart there is a unique critical point \(y_*(q)\) in the detail
chart and
\begin{equation}
 D_y\phi(q,y_*(q))=0,\qquad
 D_y^2\phi(q,y)\succeq c_*I
 \label{eq:s6v48-detail-floor}
\end{equation}
throughout a common neighbourhood, with \(c_*>0\).

\begin{lemma}[Parametric stationary section]
\label{lem:s6v48-stationary}
The map \(q\mapsto y_*(q)\) is \(C^3\), and for a retained tangent \(v\),
\begin{equation}
 D_vy_*
 =-\bigl(D_y^2\phi\bigr)^{-1}D_yD_v\phi
 \quad\hbox{at }(q,y_*(q)).
 \label{eq:s6v48-yprime}
\end{equation}
\end{lemma}

\begin{proof}
The derivative of \(D_y\phi\) in the \(y\)-variable is invertible by
\eqref{eq:s6v48-detail-floor}.  The parameter-dependent implicit function
theorem gives \(C^3\) regularity.  Differentiate
\(D_y\phi(q,y_*(q))=0\) in direction \(v\) and solve for \(D_vy_*\).
\end{proof}

Put \(z=y-y_*(q)\) and define the averaged detail Hessian
\begin{equation}
 H(q,z)
 :=
 2\int_0^1(1-t)
 D_y^2\phi(q,y_*(q)+tz)\,\dd t .
 \label{eq:s6v48-averaged-H}
\end{equation}

\begin{theorem}[Explicit parametric Morse map]
\label{thm:s6v48-Morse}
On a sufficiently small common neighbourhood of \(z=0\),
\begin{equation}
 \eta=\Theta_q(z):=H(q,z)^{1/2}z
 \label{eq:s6v48-Morse-map}
\end{equation}
is a \(C^2\) diffeomorphism onto its image and
\begin{equation}
 \boxed{
 \phi(q,y_*(q)+z)
 =\phi_*(q)+{1\over2}\|\eta\|^2,
 \qquad
 \phi_*(q):=\phi(q,y_*(q)).}
 \label{eq:s6v48-exact-quadratic}
\end{equation}
The neighbourhoods may be chosen uniformly on a compact retained chart on
which \eqref{eq:s6v48-detail-floor} and the \(C^4\) bounds are uniform.
\end{theorem}

\begin{proof}
Taylor's formula with integral remainder and
\(D_y\phi(q,y_*)=0\) gives
\[
 \phi(q,y_*+z)-\phi(q,y_*)
 =
 \int_0^1(1-t)
 D_y^2\phi(q,y_*+tz)[z,z]\,\dd t
 ={1\over2}z^TH(q,z)z.
\]
The floor \eqref{eq:s6v48-detail-floor} implies
\(H(q,z)\succeq c_*I\), so its positive square root is \(C^2\).
At \(z=0\),
\[
 D_z\Theta_q(0)=
 \bigl(D_y^2\phi(q,y_*(q))\bigr)^{1/2},
\]
which is invertible.  The inverse function theorem proves the local
diffeomorphism and the preceding Taylor identity proves
\eqref{eq:s6v48-exact-quadratic}.  Uniformity follows by compactness and the
uniform lower and derivative bounds.
\end{proof}

\begin{remark}[Pinned quotient and Haar density]
\label{rem:s6v48-Haar}
On the gauge-pinned \(SU(3)\) detail fibre, apply
\Cref{thm:s6v48-Morse} in exponential coordinates on the physical quotient.
The exponential-coordinate Haar density and the determinant of
\(\Theta_q^{-1}\) multiply into one positive Jacobian \(J(q,\eta)\).
Gauge zero modes must be removed before
\eqref{eq:s6v48-detail-floor}; otherwise the square root is not invertible.
\end{remark}

\subsection{Exact classical, entropy, and complement decomposition}
\label{sec:s6v48-decomposition}

Choose a fixed ball \(B_R\) contained in every Morse image and let
\({\mathcal U}_q=\Theta_q^{-1}(B_R)+y_*(q)\).  For the block integral
\[
 Z_\beta(q)=\int e^{-\beta\phi(q,y)}\,\dd m(y),
\]
write \(Z_\beta^{\rm loc}\) for the integral over \({\mathcal U}_q\) and
\(Z_\beta^{\rm out}=Z_\beta-Z_\beta^{\rm loc}\).  Let
\(\gamma_{\beta,R}\) be normalized Gaussian measure on \(B_R\), and set
\begin{align}
 {\mathcal E}_\beta(q)
 &:=
 -\log\int_{B_R}J(q,\eta)\,
 \gamma_{\beta,R}(\dd\eta),
 \label{eq:s6v48-entropy}\\
 \rho_\beta(q)
 &:={Z_\beta^{\rm out}(q)\over Z_\beta^{\rm loc}(q)},
 \qquad
 {\mathcal T}_\beta(q):=-\log(1+\rho_\beta(q)).
 \label{eq:s6v48-tail}
\end{align}

\begin{theorem}[Exact three-owner marginal decomposition]
\label{thm:s6v48-three-owner}
Up to a \(q\)-independent Gaussian normalization,
\begin{equation}
 \boxed{
 -\log Z_\beta(q)
 =\beta\phi_*(q)
  +{\mathcal E}_\beta(q)
  +{\mathcal T}_\beta(q).}
 \label{eq:s6v48-three-owner}
\end{equation}
\end{theorem}

\begin{proof}
The change of variables in \Cref{thm:s6v48-Morse} gives
\[
 Z_\beta^{\rm loc}(q)
 =
 e^{-\beta\phi_*(q)}
 \int_{B_R}e^{-\beta\|\eta\|^2/2}J(q,\eta)\,\dd\eta.
\]
Factor out the \(q\)-independent Gaussian normalizer and take negative
logarithms.  Since
\(Z_\beta=Z_\beta^{\rm loc}(1+\rho_\beta)\), the last logarithm is
\({\mathcal T}_\beta\).
\end{proof}

\begin{theorem}[Exact nonlinear classical Schur Hessian]
\label{thm:s6v48-classical-Schur}
For retained tangents \(v,w\),
\begin{align}
 D_v\phi_*
 &=D_v\phi(q,y_*(q)),
 \label{eq:s6v48-envelope-first}\\
 D_vD_w\phi_*
 &=
 D_vD_w\phi
 -D_vD_y\phi\,
  (D_y^2\phi)^{-1}
  D_yD_w\phi
 \quad\hbox{at }(q,y_*(q)).
 \label{eq:s6v48-envelope-second}
\end{align}
\end{theorem}

\begin{proof}
Differentiate \(\phi(q,y_*(q))\).  The term containing \(D_y\phi\) vanishes
at the stationary section, proving
\eqref{eq:s6v48-envelope-first}.  Differentiate once more and substitute
\eqref{eq:s6v48-yprime}; this gives
\eqref{eq:s6v48-envelope-second}.
\end{proof}

\begin{assessment}[The correct order-\(\beta\) owner]
\label{ass:s6v48-classical-owner}
The full nonlinear classical action \(\beta\phi_*(q)\), not merely its
vacuum quadratic Taylor polynomial, is the order-\(\beta\) owner.  GSOMC is
its quadratic germ.  A claim that the generated remainder is \(o(\beta)\)
therefore requires a classical minimizer approximation theorem comparing
\(\phi_*\) with the declared coarse Wilson or response-matched baseline.
Mixing of fluctuations cannot prove that deterministic comparison.
\end{assessment}

\subsection{The entropy Jacobian as an \(O(N_M)\) candidate}
\label{sec:s6v48-entropy-bound}

Put
\[
 L(q,\eta):=-\log J(q,\eta).
\]
Then \({\mathcal E}_\beta\) is, up to a constant,
\[
 -\log\int_{B_R}
 e^{-\beta\|\eta\|^2/2-L(q,\eta)}\,\dd\eta.
\]
Let \(\nu_q^J\) be the corresponding normalized law.

\begin{theorem}[Morse-Jacobian Laplace bound]
\label{thm:s6v48-Jacobian}
Assume
\begin{equation}
 \beta I+D_\eta^2L(q,\eta)\succeq\kappa\beta I
 \label{eq:s6v48-J-convexity}
\end{equation}
on \(B_R\), uniformly in \(q\).  For retained tangents \(v,w\), define
\begin{align}
 j_{vw}
 &:=
 \sup_q
 \left\|
  \mathbb E_{\nu_q^J}D_vD_wL
 \right\|_{\rm op},
 \label{eq:s6v48-j2}\\
 \ell_v^2
 &:=
 \sup_q
 \mathbb E_{\nu_q^J}
 \|\nabla_\eta D_vL\|^2.
 \label{eq:s6v48-jgrad}
\end{align}
Then
\begin{equation}
 \boxed{
 \|D_vD_w{\mathcal E}_\beta\|_{\rm op}
 \le
 j_{vw}+{\ell_v\ell_w\over\kappa\beta}.}
 \label{eq:s6v48-J-bound}
\end{equation}
Moreover,
\begin{equation}
 |{\mathcal E}_\beta(q)-{\mathcal E}_\beta(q')|
 \le
 \sup_{\eta\in B_R}|L(q,\eta)-L(q',\eta)|.
 \label{eq:s6v48-J-osc}
\end{equation}
\end{theorem}

\begin{proof}
This is \Cref{thm:s6v47-RLS,lem:s6v47-log-osc} with
\(\widetilde R=L\), \(C=I\), and the fixed fibre \(B_R\).  The density
condition \eqref{eq:s6v48-J-convexity} is exactly the strong-convexity
hypothesis required there.
\end{proof}

\begin{corollary}[Local Jacobian incidence gives a volume debit]
\label{cor:s6v48-J-volume}
Suppose the pinned block has \(N_M\le c_{\rm vol}M^4\) local coordinates and
the Jacobian stocks satisfy
\begin{align}
 \sup_{|u|=1}j_{uu}
 +\sup_{|u|=1}\sum_{f\leadsto e}
  \sup_{|\zeta_f|=1}j_{u\zeta_f}
 +\operatorname{osc}_{\rm col}L
 &\le c_JN_M,
 \label{eq:s6v48-J-local}\\
 \sup_{|u|=1}\ell_u^2
 +\sup_{|u|=1}\sum_{f\leadsto e}
  \sup_{|\zeta_f|=1}\ell_u\ell_{\zeta_f}
 &\le c'_J\beta N_M.
 \label{eq:s6v48-J-gradient}
\end{align}
Then the interior Hessian, mixed, and collar-oscillation stocks of
\({\mathcal E}_\beta\) are bounded by \(C_JN_M\).
\end{corollary}

\begin{proof}
Apply \eqref{eq:s6v48-J-bound} in the one-link and mixed directions, sum the
eighteen exterior slots, and use \eqref{eq:s6v48-J-osc}.
\end{proof}

\begin{lemma}[Chart-complement logarithm]
\label{lem:s6v48-tail}
Suppose \(0\le\rho_\beta\le1/2\) and, for all one-link and exterior
directions \(v,w\) used below,
\begin{equation}
 |\rho_\beta|+\|D_v\rho_\beta\|+
 \|D_vD_w\rho_\beta\|
 \le P(\beta,N_M)e^{-\beta\Delta_*},
 \label{eq:s6v48-rho-bound}
\end{equation}
where \(\log(2+P(\beta,N_M))=o(\beta)\) and \(\Delta_*>0\).  Then the
corresponding first, second, mixed, and oscillation stocks of
\({\mathcal T}_\beta=-\log(1+\rho_\beta)\) are \(o(1)\).
\end{lemma}

\begin{proof}
The first two derivatives are
\[
 D_v{\mathcal T}_\beta
 =-{D_v\rho_\beta\over1+\rho_\beta},
 \qquad
 D_vD_w{\mathcal T}_\beta
 =-{D_vD_w\rho_\beta\over1+\rho_\beta}
  +{D_v\rho_\beta D_w\rho_\beta\over(1+\rho_\beta)^2}.
\]
The denominators are bounded below by one, and
\(|\log(1+\rho_\beta)|\le\rho_\beta\).  Insert
\eqref{eq:s6v48-rho-bound}; the exponential dominates the subexponential
prefactor.
\end{proof}

\subsection{A two-scale closing card}
\label{sec:s6v48-closing}

Let \(\Phi_M^{\rm base}(q)\) be a declared full classical/perfect-action
baseline whose vacuum quadratic germ is the complete GSOMC Schur kernel,
including its exponentially local non-plaquette part.  A coarse Wilson
action plus one matched plaquette coefficient is admissible here only if a
separate uniform comparison theorem proves that it has this full germ.
Define the classical discrepancy
\begin{equation}
 {\mathcal C}_{\beta,M}(q)
 :=\beta\{\phi_*(q)-\Phi_M^{\rm base}(q)\}.
 \label{eq:s6v48-classical-discrepancy}
\end{equation}

\begin{definition}[Classical minimizer approximation card, CMAC]
\label{def:s6v48-CMAC}
CMAC with exponent two means that, on the aligned retained chart,
\begin{equation}
 {\mathfrak h}_{\cal C}
 +{\mathfrak j}_{\cal C}^{\rm in}
 +{\mathfrak o}_{\cal C}
 \le C_{\rm cl}{\beta\over M^2},
 \qquad
 \log(2+{\mathfrak j}_{\cal C}^{\rm all})
 \le C_{\rm cl}M^4+C'_{\rm cl}\log(2+\beta).
 \label{eq:s6v48-CMAC}
\end{equation}
The first inequality is a deterministic approximation of the nonlinear
minimized action, not a fluctuation estimate.
\end{definition}

\begin{theorem}[Classical-plus-entropy DQIS closure]
\label{thm:s6v48-closing}
Assume CMAC, the Jacobian hypotheses
\eqref{eq:s6v48-J-convexity}--\eqref{eq:s6v48-J-gradient}, and the complement
bound \eqref{eq:s6v48-rho-bound}.  If
\begin{equation}
 M\longrightarrow\infty,
 \qquad
 M=o(\beta^{1/4}),
 \label{eq:s6v48-two-scale}
\end{equation}
then, after the baseline owner is removed from
\eqref{eq:s6v48-three-owner}, the generated remainder
\[
 {\mathcal G}_{\beta,M}
 ={\mathcal C}_{\beta,M}
  +{\mathcal E}_\beta
  +{\mathcal T}_\beta
\]
satisfies
\begin{equation}
 {\mathfrak h}_{\cal G}
 +{\mathfrak j}_{\cal G}^{\rm in}
 +{\mathfrak o}_{\cal G}
 =o(\beta),
 \qquad
 \log(2+{\mathfrak j}_{\cal G}^{\rm all})=o(\beta).
 \label{eq:s6v48-generated-small}
\end{equation}
Hence it passes the split-stock response theorem of V47.
\end{theorem}

\begin{proof}
CMAC contributes \(O(\beta/M^2)=o(\beta)\).  By
\Cref{cor:s6v48-J-volume}, the entropy contributes
\(O(N_M)=O(M^4)=o(\beta)\).  The complement contribution is \(o(1)\) by
\Cref{lem:s6v48-tail}.  The global logarithmic stock is
\(O(M^4+\log(2+\beta))=o(\beta)\), after enlarging the subexponential
prefactor in \eqref{eq:s6v48-rho-bound}.  Summing the three owners gives
\eqref{eq:s6v48-generated-small}; apply
\Cref{thm:s6v47-split-response}.
\end{proof}

\begin{assessment}[What V48 changes]
\label{ass:s6v48-change}
The nonlinear fluctuation problem is no longer assigned to a raw Taylor
remainder multiplied by \(\beta\).  Exact Morse reduction moves it into the
Jacobian entropy, for which one local debit per fine coordinate gives
\(O(M^4)\).  The remaining order-\(\beta\) task is CMAC: compare the exact
classical minimized fine action with the declared coarse baseline.  This
separation prevents a fluctuation estimate from being asked to prove a
deterministic renormalization statement.
\end{assessment}

\begin{programme}[V49: classical minimizer approximation before the V50
audit]
\label{prog:s6v48-V49}
\begin{enumerate}[leftmargin=3em]
\item Express the pinned vacuum minimizer by discrete harmonic extension
      using the V37 inverse detail Hessian.
\item Compare its quadratic energy exactly with the full GSOMC owner and
      identify every exponentially local non-plaquette component; test
      separately whether a scalar coarse-Wilson owner can represent it.
\item Use the plaquette Taylor formula along the minimizing extension to
      seek the CMAC rate \(O(\beta/M^2)\), separating boundary, toron, and
      topological sectors.
\item If the exponent two fails, record the actual deterministic rate and
      recompute the nonempty two-scale corridor.
\item Keep the Jacobian jet packet and PFS as separate next inputs; neither
      may be inferred from CMAC.
\end{enumerate}
\end{programme}

\begin{firewall}[Claim boundary after sixth-series V48]
\label{fire:s6v48-boundary}
V48 proves an explicit parameter-dependent Morse quadraticization on a
pinned strongly convex detail chart; proves the exact decomposition of the
block marginal into classical minimized, Jacobian entropy, and chart-tail
owners; proves the nonlinear classical Schur formula; proves
Brascamp--Lieb bounds for the entropy Hessian and mixed response; proves the
chart-tail derivative lemma; and shows that CMAC plus \(O(M^4)\) Jacobian
stocks closes the generated split-DQIS row on
\(1\ll M\ll\beta^{1/4}\).

V48 does not prove CMAC or identify a response-matched full perfect-action
baseline, the physical Morse-Jacobian jet stocks, PFS, the
physical chart-tail bound, CPTC, LCKC, selected-root asymptotics, TSC,
YMGNC, continuum existence, a finite physical-clock mass, or the full
Yang--Mills mass gap.
\end{firewall}

\begin{theorem}[Sixth-series V48 result ledger]
\label{thm:s6v48-results}
The positive conclusions in \Cref{fire:s6v48-boundary} follow from
\Cref{lem:s6v48-stationary,
thm:s6v48-Morse,
thm:s6v48-three-owner,
thm:s6v48-classical-Schur,
thm:s6v48-Jacobian,
cor:s6v48-J-volume,
lem:s6v48-tail,
thm:s6v48-closing}.
\end{theorem}

\begin{proof}
Each item is the exact conclusion of the cited result.
\end{proof}

\paragraph{Parent-version record.}
V48 was formed from
\texttt{Renewal\_Yang\_Mills\_Mass\_Gap\_Research\_Notes\_V47.tex}.
The V47 parent had \(21\,147\) logical source lines,
\(738\,211\) bytes, and SHA--256
\[
 \texttt{528aac9da91e6a4353a7144b7d17a128de04dae1b3cb93f68115b0b51b552f37}.
\]
The next compulsory companion audit is V50.

% =============================================================================
% END APPENDED SIXTH-SERIES V48 MATERIAL
% =============================================================================

% =============================================================================
% APPENDED SIXTH-SERIES V49 MATERIAL
% =============================================================================

\section{The flat-Schur symbol obstruction and the Ward-critical influence
row}
\label{sec:s6v49}

V48 made the full GSOMC kernel the quadratic germ of any admissible
classical baseline.  The exact axial symbol imported in V31 now decides
whether one scalar coarse-Wilson coefficient can be that germ.  It cannot:
the Ward-normalized Schur ratio varies by a fixed amount across the Brillouin
zone.  The full exponentially local Schur kernel is therefore unavoidable.
That repair reveals a second, deeper fact: because the kernel has a Ward
zero, its Gaussian one-site absolute-influence row cannot be strictly below
one.  Thus the fixed Dobrushin contraction used for a finite-regulator
estimate cannot be the cofinal physical continuum mechanism.

\subsection{One scalar plaquette owner does not equal the flat Schur owner}
\label{sec:s6v49-scalar}

On the central axial transverse block, V31 imported the exact unit flat
Schur symbol
\begin{equation}
 s_0(q)
 ={t(t+8)(t+24)\over9(t^2+18t+16)},
 \qquad
 t=\sin^2(q/2),
 \label{eq:s6v49-s0}
\end{equation}
and the Wilson Ward symbol is
\begin{equation}
 \omega(q)=4\sin^2(q/2)=4t.
 \label{eq:s6v49-omega}
\end{equation}
For \(q\ne0\), put
\begin{equation}
 r(t):={s_0(q)\over\omega(q)}
 ={(t+8)(t+24)\over36(t^2+18t+16)}.
 \label{eq:s6v49-ratio}
\end{equation}

\begin{lemma}[Exact endpoint separation]
\label{lem:s6v49-endpoints}
The continuous extension of \(r\) to \(t=0\) satisfies
\begin{equation}
 r(0)={1\over3},
 \qquad
 r(1)={5\over28},
 \qquad
 r(0)-r(1)={13\over84}.
 \label{eq:s6v49-endpoints}
\end{equation}
\end{lemma}

\begin{proof}
Substitute \(t=0\) and \(t=1\) in
\eqref{eq:s6v49-ratio}:
\[
 r(0)={8\cdot24\over36\cdot16}={1\over3},
 \qquad
 r(1)={9\cdot25\over36\cdot35}={5\over28}.
\]
Their difference is \(28/84-15/84=13/84\).
\end{proof}

\begin{theorem}[Scalar-baseline obstruction]
\label{thm:s6v49-scalar-obstruction}
For every real scalar \(c\),
\begin{align}
 \boxed{
 \sup_{q\ne0}
 {\,|s_0(q)-c\omega(q)|\over\omega(q)}
 \ge {13\over168}.}
 \label{eq:s6v49-scalar-obstruction}
 \\
 \boxed{
 \sup_{q}
 |s_0(q)-c\omega(q)|
 \ge {1286\over19089}.}
 \label{eq:s6v49-absolute-obstruction}
\end{align}
Consequently no scalar adjustment of a coarse Wilson plaquette action,
including one response-matched coefficient \(b^{\rm marg}\), approximates
the complete flat Schur Hessian with either a vanishing Ward-relative error
or a vanishing uniform absolute symbol error.
\end{theorem}

\begin{proof}
For any \(a,b,c\in\mathbb R\),
\(\max\{|a-c|,|b-c|\}\ge|a-b|/2\).  Apply this with
\(a=r(0)\), \(b=r(1)\), and use
\Cref{lem:s6v49-endpoints}.  The value at \(t=0\) is reached as
\(q\to0\) through nonzero momenta, while \(t=1\) is an allowed axial
endpoint.  For the absolute bound, direct substitution gives
\[
 r(1/2)={833\over3636},
 \qquad
 r(1/2)-r(1)={643\over12726}.
\]
If \(M_c\) is the maximum of the absolute symbol errors at
\(t=1/2\) and \(t=1\), then
\[
 |r(1/2)-r(1)|
 \le {M_c\over2}+{M_c\over4},
\]
because \(\omega=2\) and \(4\) at those two points.  Hence
\[
 M_c\ge {4\over3}{643\over12726}
 ={1286\over19089}.
\]
The last assertion follows because a scalar plaquette adjustment changes
\(c\) but not the nonconstant function \(r(t)\).
\end{proof}

\begin{corollary}[Quadratic failure of scalar CMAC]
\label{cor:s6v49-scalar-CMAC}
If \(\Phi_M^{\rm base}\) in V48 is restricted to a scalar Wilson plaquette
owner, its quadratic classical discrepancy has Ward-relative norm at least
\(13\beta/168\) and uniform absolute symbol norm at least
\(1286\beta/19089\).  It therefore cannot satisfy the \(o(\beta)\)
conclusion of CMAC uniformly over the axial momentum block.
\end{corollary}

\begin{proof}
The quadratic Hessian of the exact minimized action is
\(\beta s_0(q)\), whereas that of a scalar Wilson owner is
\(\beta c\omega(q)\).  Multiply
\eqref{eq:s6v49-scalar-obstruction} by \(\beta\).
\end{proof}

\begin{assessment}[What response matching actually fixes]
\label{ass:s6v49-response}
The response cofactor \(b^{\rm marg}=-n/d\) fixes one scalar projection of
the action.  \Cref{thm:s6v49-scalar-obstruction} proves that it cannot also
identify the complete momentum-dependent flat Schur kernel.  The
non-plaquette exponentially local component is not a small fluctuation
tail; it is part of the leading quadratic owner and must be retained once.
\end{assessment}

\subsection{The full flat owner is quasilocal}
\label{sec:s6v49-full-owner}

Let
\begin{equation}
 {\mathsf S}_M^0
 =K_M^0-(B_M^0)^*(A_M^0)^{-1}B_M^0
 \label{eq:s6v49-full-Schur}
\end{equation}
denote the pinned unit-Wilson retained Schur Hessian at the vacuum, in the
notation of V37.

\begin{theorem}[Perfect quadratic baseline]
\label{thm:s6v49-perfect-quadratic}
The quadratic action
\begin{equation}
 \Phi_M^{\rm PQ}(x)
 ={1\over2}\langle x,{\mathsf S}_M^0x\rangle
 \label{eq:s6v49-PQ}
\end{equation}
has exactly the GSOMC quadratic germ.  Its kernel is exponentially
summable uniformly in \(M\).  Hence it is a legitimate quasilocal
quadratic baseline, although it is not a finite-range scalar Wilson
plaquette action.
\end{theorem}

\begin{proof}
The first claim is the definition of the Schur complement
\eqref{eq:s6v49-full-Schur}.  The coarse block \(K_M^0\) has fixed range.
The nonlocal term
\((B_M^0)^*(A_M^0)^{-1}B_M^0\) has the exponential kernel bound of
\Cref{cor:s6v37-flat-kernel}.  Adding a finite-range kernel preserves
exponential summability uniformly in \(M\).  The scalar obstruction is
\Cref{thm:s6v49-scalar-obstruction}.
\end{proof}

\begin{corollary}[Quadratic CMAC is exact for the perfect owner]
\label{cor:s6v49-PQ-CMAC}
If the V48 classical baseline uses
\(\Phi_M^{\rm PQ}\), its vacuum quadratic discrepancy vanishes exactly.
The remaining classical discrepancy begins at third retained derivative.
\end{corollary}

\begin{proof}
\Cref{thm:s6v48-classical-Schur} evaluated at the vacuum gives
\({\mathsf S}_M^0\), which is the Hessian of
\eqref{eq:s6v49-PQ}.  Subtraction cancels the second jet.
\end{proof}

\begin{lemma}[Local cubic remainder on a shrinking retained chart]
\label{lem:s6v49-cubic}
Let
\[
 {\cal R}_M(x)
 :=\phi_*(x)-\phi_*(0)-D\phi_*(0)[x]
  -{1\over2}\langle x,{\mathsf S}_M^0x\rangle .
\]
If
\begin{equation}
 \sup_{\|x\|\le r_M}\|D^3\phi_*(x)\|_{\rm op}\le L_M,
 \label{eq:s6v49-third}
\end{equation}
then on that ball
\begin{align}
 \|D^2{\cal R}_M(x)\|_{\rm op}
 &\le L_Mr_M,
 \label{eq:s6v49-H-rem}\\
 \|D{\cal R}_M(x)\|
 &\le {L_M\over2}r_M^2,
 \label{eq:s6v49-D-rem}\\
 |{\cal R}_M(x)|
 &\le {L_M\over6}r_M^3.
 \label{eq:s6v49-R-rem}
\end{align}
\end{lemma}

\begin{proof}
The value and first two derivatives of \({\cal R}_M\) vanish at zero by
\Cref{cor:s6v49-PQ-CMAC}.  Integrate \(D^3\phi_*\) once, twice, and three
times along the segment \(t\mapsto tx\).
\end{proof}

\begin{proposition}[Weighted locality propagates to the classical third
jet]
\label{prop:s6v49-third-locality}
Fix a weighted row norm with exponent below the V37 inverse exponent.
Suppose on a pinned vacuum chart:
\begin{enumerate}[label=\textup{(\roman*)},leftmargin=3.5em]
\item the detail Hessian inverse has uniformly bounded weighted row norm;
\item the retained--detail Hessian and the third derivatives of the
      unit-coupling Wilson action have uniformly bounded finite-range
      incidence norms; and
\item the nonlinear detail Hessian differs from its vacuum value by less
      than the inverse weighted-norm radius.
\end{enumerate}
Then the stationary-section derivatives through order two and
\(D^3\phi_*\) have uniformly bounded weighted row norms.  In particular,
\(L_M\le L_*\) in the corresponding local retained norm, independently of
\(M\).
\end{proposition}

\begin{proof}
V37 supplies item~\textup{(i)} at the vacuum.  Under item~\textup{(iii)},
the resolvent identity and Neumann series preserve the weighted inverse
bound.  Formula \eqref{eq:s6v48-yprime} writes \(Dy_*\) as one product of
that inverse with the mixed Hessian.  Differentiating the stationary
equation once more writes \(D^2y_*\) as the same inverse applied to a finite
sum of contractions of \(D^3\phi\) with \(Dy_*\).  Differentiating the
classical Schur formula \eqref{eq:s6v48-envelope-second} expresses
\(D^3\phi_*\) as a finite sum of the same factors.

Weighted row norms are submultiplicative by the graph triangle inequality.
The number of contractions and local incidences is fixed, so the resulting
bounds are uniform in \(M\).
\end{proof}

\begin{corollary}[A shrinking-chart classical row]
\label{cor:s6v49-shrinking-CMAC}
Under \Cref{prop:s6v49-third-locality}, if
\[
 r_M\longrightarrow0,
\]
then on the retained ball of radius \(r_M\), the Hessian and mixed
influence stocks of
\(\beta{\cal R}_M\) are \(O(\beta r_M)=o(\beta)\).  Its local oscillation is
\(O(\beta r_M^3)\).
\end{corollary}

\begin{proof}
Use \Cref{lem:s6v49-cubic} with \(L_M\le L_*\).  Mixed second derivatives
are components of the same Hessian bound.
\end{proof}

\begin{assessment}[The price of the perfect quadratic baseline]
\label{ass:s6v49-PQ-price}
The shrinking-chart result repairs the deterministic nonlinear remainder,
but the perfect quadratic baseline itself is massless at the Ward mode.
Its conditional influence cannot have a cofinal row bounded strictly below
one.  This prevents the finite-regulator fixed-contraction argument from
being promoted directly to a continuum mass theorem.
\end{assessment}

\subsection{The Ward zero forbids strict absolute one-site contraction}
\label{sec:s6v49-Ward-row}

\begin{theorem}[Zero-row-sum influence obstruction]
\label{thm:s6v49-zero-row}
Let \(Q=(Q_{xy})\) be a translation-invariant scalar precision kernel with
\[
 Q_{00}>0,\qquad
 \sum_yQ_{0y}=0.
\]
For the Gaussian one-site conditional, the absolute influence row is
\begin{equation}
 \sum_{y\ne0}{|Q_{0y}|\over Q_{00}}\ge1.
 \label{eq:s6v49-Gaussian-row}
\end{equation}
Thus no Dobrushin estimate based on the absolute one-site Gaussian response
has a strict row sum below one.
\end{theorem}

\begin{proof}
The conditional mean at the origin is
\[
 \mathbb E[X_0\mid X_{0^c}]
 =-{1\over Q_{00}}\sum_{y\ne0}Q_{0y}X_y,
\]
so its coordinate Lipschitz coefficients are
\(|Q_{0y}|/Q_{00}\).  The zero-row-sum identity gives
\(\sum_{y\ne0}Q_{0y}=-Q_{00}\).  Hence
\[
 \sum_{y\ne0}|Q_{0y}|
 \ge\left|\sum_{y\ne0}Q_{0y}\right|=Q_{00},
\]
which proves \eqref{eq:s6v49-Gaussian-row}.
\end{proof}

\begin{corollary}[Application to the flat Schur owner]
\label{cor:s6v49-Schur-row}
On every scalarized physical polarization for which
\(\widehat{\mathsf S}_M^0(0)=0\) and the bulk diagonal is positive, the
perfect quadratic baseline obeys the obstruction
\eqref{eq:s6v49-Gaussian-row}.  On growing pinned tori the same lower bound
holds on bulk rows up to a boundary error tending to zero.
\end{corollary}

\begin{proof}
The Ward zero is equivalent to the zero sum of the convolution kernel.
Apply \Cref{thm:s6v49-zero-row}.  Pinning changes rows only through the
boundary correction; exponential locality from
\Cref{thm:s6v49-perfect-quadratic} makes that correction tend to zero for
rows whose distance from the pinning boundary tends to infinity.
\end{proof}

\begin{theorem}[Fixed-contraction versus Ward-critical continuum
alternative]
\label{thm:s6v49-critical-alternative}
For a cofinal family whose quadratic germ converges to the Ward-critical
Schur owner, the following two claims cannot both hold:
\begin{enumerate}[label=\textup{(\roman*)},leftmargin=3.5em]
\item the linearized absolute one-site quotient response row of the actual
      specification converges to the Gaussian response row of the Schur
      owner; and
\item its Dobrushin/MSEC row has a uniform strict margin below one.
\end{enumerate}
Therefore a nontrivial continuum route must use a scale-dependent
near-critical row, a non-one-site block of nonvanishing physical duration,
or a cancellation norm not dominated by absolute coordinate influences.
\end{theorem}

\begin{proof}
Under item~\textup{(i)}, the limiting actual linearized row is bounded below
by one by \Cref{cor:s6v49-Schur-row}.  Item~\textup{(ii)} would bound the
same row strictly below one, a contradiction.  The three listed alternatives
are precisely ways of abandoning the conjunction: allow the row to approach
one at the physical clock rate, keep a macroscopic block duration, or avoid
the absolute one-site majorant.
\end{proof}

\begin{assessment}[Revised location of the mass mechanism]
\label{ass:s6v49-frontier}
The full flat owner can be made quasilocal and its nonlinear classical
remainder can be small on a shrinking chart, but its Ward zero rules out a
cofinal fixed one-site contraction.  This agrees with the V35 and V43
ultralocality firewall.  The physical mass gap must arise from the
scale-dependent reflected transfer or an infrared block contraction after
the massless ultraviolet Ward structure has been retained, not from a
uniform ultraviolet one-link Dobrushin margin.
\end{assessment}

\begin{programme}[V50: compulsory audit and physical-clock critical row]
\label{prog:s6v49-V50}
\begin{enumerate}[leftmargin=3em]
\item Perform the compulsory audit of the then-current SM and NS notes.
\item Replace the fixed one-link margin by a scale-dependent block
      coefficient
      \(\alpha_j=1-\mu\tau_j+o(\tau_j)\), keeping the Ward-critical
      quadratic owner explicit.
\item Derive the exact finite-block condition under which MSEC covariance
      decay over \(k\) microscopic gates produces that physical-time row.
\item Separate ultraviolet spatial mixing from the reflected temporal
      transfer; do not identify graph distance with physical Euclidean time
      without the declared block map.
\item Retain CMAC, the Jacobian packet, PFS, root asymptotics, and transfer
      projectivity as independent inputs.
\end{enumerate}
\end{programme}

\begin{firewall}[Claim boundary after sixth-series V49]
\label{fire:s6v49-boundary}
V49 proves a quantitative Ward-normalized obstruction to representing the
flat Schur symbol by one scalar Wilson plaquette coefficient; constructs the
full exponentially local perfect quadratic owner; proves local third-jet
control of its nonlinear classical remainder under the V37 weighted inverse
packet; proves the zero-row-sum obstruction to strict absolute Gaussian
one-site influence; and derives the fixed-contraction versus Ward-critical
continuum alternative.

V49 does not construct the response-matched full perfect-action baseline,
prove a shrinking-chart probability estimate or PFS, populate the physical
Morse-Jacobian packet, obtain the scale-dependent physical-clock row,
prove transfer projectivity, CPTC, LCKC, root asymptotics, TSC, YMGNC,
continuum existence, or the full Yang--Mills mass gap.
\end{firewall}

\begin{theorem}[Sixth-series V49 result ledger]
\label{thm:s6v49-results}
The positive statements in \Cref{fire:s6v49-boundary} follow from
\Cref{lem:s6v49-endpoints,
thm:s6v49-scalar-obstruction,
cor:s6v49-scalar-CMAC,
thm:s6v49-perfect-quadratic,
cor:s6v49-PQ-CMAC,
lem:s6v49-cubic,
prop:s6v49-third-locality,
cor:s6v49-shrinking-CMAC,
thm:s6v49-zero-row,
cor:s6v49-Schur-row,
thm:s6v49-critical-alternative}.
\end{theorem}

\begin{proof}
Each ledger item is the exact conclusion of the cited result.
\end{proof}

\paragraph{Parent-version record.}
V49 was formed from
\texttt{Renewal\_Yang\_Mills\_Mass\_Gap\_Research\_Notes\_V48.tex}.
The V48 parent had \(21\,618\) logical source lines,
\(753\,674\) bytes, and SHA--256
\[
 \texttt{b5cca4da1d9829e8e425d6d6d55769a129f4e07cc8e89a75f74b8ffa8d5699d6}.
\]
The next compulsory companion audit is V50.

% =============================================================================
% END APPENDED SIXTH-SERIES V49 MATERIAL
% =============================================================================

% =============================================================================
% APPENDED SIXTH-SERIES V50 MATERIAL
% =============================================================================

\section{Compulsory companion audit and the physical-clock critical row}
\label{sec:s6v50}

V49 proved that the Ward-critical flat Schur owner cannot support a
cofinal fixed absolute one-site contraction.  Version V50 performs the
required companion audit and replaces that route by a physical-time slab
coefficient.  The correct finite positive mass regime is
\(q_j=1-\mu\delta_j+o(\delta_j)\), not a fixed \(q<1\) as the slab duration
\(\delta_j\) vanishes.  The same calculation exposes the strengthened
projectivity requirement: adjacent kernel defects are amplified by
\((1-q_j)^{-1}\asymp\delta_j^{-1}\).

\subsection{The V50 compulsory audit}
\label{sec:s6v50-audit}

\begin{definition}[Frozen V50 companion packet]
\label{def:s6v50-audit}
The exact files inspected read-only are
\begin{align}
 &\texttt{Renewal\_SM\_Einstein\_Continuum\_Research\_Notes\_V99.tex},
 \notag\\
 &\qquad
 39\,816\ {\rm lines},\quad1\,553\,794\ {\rm bytes},\quad
 \texttt{2ec68514815f62ecea1ea478676e89870599c6d334df1d628c31ab2395dae014},
 \label{eq:s6v50-SM}\\
 &\texttt{Renewal\_NS\_Stability\_Research\_Notes\_V26.tex},
 \notag\\
 &\qquad
 5\,319\ {\rm lines},\quad278\,283\ {\rm bytes},\quad
 \texttt{82c887f8c2c69e4f28fad7c3dc8de5b9099cb5a4bf431eba1fff3e91935978ad}.
 \label{eq:s6v50-NS}
\end{align}
\end{definition}

\begin{assessment}[SM V99 lesson]
\label{ass:s6v50-SM}
SM V97--V99 prove, at an abstract channel level, that a mixed response is
amplified by a resolvent factor \((1-c)^{-1}\), that one summable scalar
ledger can close countably many fixed local rows, and that convergence along
one regulator chain is not regulator independence without a quotient
cocycle or common-refinement bridge.  These are directly relevant to a
near-critical Yang--Mills transfer, where \(1-q_j\) is of order
\(\delta_j\).  The SM notes do not populate any Yang--Mills kernel,
reflection-positive slab, or bridge estimate.
\end{assessment}

\begin{assessment}[NS V26 unchanged]
\label{ass:s6v50-NS}
The NS companion is byte-for-byte unchanged from the V45 audit.  Its
rank-one norm lesson remains applicable only after the physical input cone
has been proved rank one.  It supplies no physical-clock or Yang--Mills
projectivity row.
\end{assessment}

\begin{theorem}[V50 audit decision]
\label{thm:s6v50-audit-decision}
The Yang--Mills programme adopts two pieces of companion proof logic:
near-critical response estimates must retain their inverse-margin
amplification, and continuum uniqueness must be checked on common
refinements rather than one preferred chain.  No companion field estimate
is imported.
\end{theorem}

\begin{proof}
This is exactly the distinction stated in
\Cref{ass:s6v50-SM,ass:s6v50-NS}: the imported statements are functional
analytic implications, while their physical hypotheses remain to be proved
for Yang--Mills.
\end{proof}

\subsection{Time-slab contraction and its rate trichotomy}
\label{sec:s6v50-clock}

Let \(P_j\) be a positive self-adjoint Markov contraction on the
finite-regulator physical Hilbert space, normalized so constants have
eigenvalue one.  Let \({\cal H}_{j,0}\) be the orthogonal complement of the
vacuum and suppose
\begin{equation}
 \|P_j|_{{\cal H}_{j,0}}\|\le q_j<1.
 \label{eq:s6v50-q}
\end{equation}
One application of \(P_j\) advances physical Euclidean time by
\(\delta_j>0\).

\begin{theorem}[Exact finite-slab spectral rate]
\label{thm:s6v50-finite-rate}
The normalized finite-regulator Hamiltonian
\[
 H_j=-\delta_j^{-1}\log P_j
\]
has vacuum-complement gap at least
\begin{equation}
 \boxed{\mu_j={-\log q_j\over\delta_j}.}
 \label{eq:s6v50-muj}
\end{equation}
\end{theorem}

\begin{proof}
By the spectral theorem, the spectrum of \(P_j\) on
\({\cal H}_{j,0}\) lies in \([0,q_j]\).  The decreasing spectral map
\(-\delta_j^{-1}\log\lambda\) is at least
\(-\delta_j^{-1}\log q_j\) there.
\end{proof}

\begin{theorem}[Physical-clock trichotomy]
\label{thm:s6v50-trichotomy}
Assume \(\delta_j\to0\) and \(q_j\to1\).  Then:
\begin{enumerate}[label=\textup{(\roman*)},leftmargin=3.5em]
\item if \((1-q_j)/\delta_j\to\mu\in(0,\infty)\), then
      \(\mu_j\to\mu\);
\item if \((1-q_j)/\delta_j\to0\), then \(\mu_j\to0\);
\item if \((1-q_j)/\delta_j\to\infty\), then
      \(\mu_j\to\infty\).
\end{enumerate}
Equivalently, a finite nonzero physical rate is
\begin{equation}
 \boxed{q_j=1-\mu\delta_j+o(\delta_j).}
 \label{eq:s6v50-critical-row}
\end{equation}
A fixed \(q_j\le q_*<1\) with \(\delta_j\to0\) lies in the third,
ultralocal regime.
\end{theorem}

\begin{proof}
For \(q\to1\),
\(-\log q=(1-q)+O((1-q)^2)\).  Divide by \(\delta_j\).  In the first two
cases \(1-q_j\to0\) makes the quadratic remainder negligible relative to
the asserted limit.  In the third case, the elementary inequality
\(-\log q\ge1-q\) gives divergence.  The fixed-\(q_*\) statement follows
directly from \eqref{eq:s6v50-muj}.
\end{proof}

\begin{corollary}[Microscopic gate iteration does not cure the clock]
\label{cor:s6v50-micro}
Suppose a time slab contains \(m_j\) microscopic gates of duration
\(\tau_j\), so \(\delta_j=m_j\tau_j\), and the only available estimate is
\[
 q_j\le\alpha_j^{m_j}.
\]
Then
\begin{equation}
 \mu_j\ge{-\log\alpha_j\over\tau_j}.
 \label{eq:s6v50-micro-rate}
\end{equation}
If \(\alpha_j\le\alpha_*<1\) and \(\tau_j\to0\), this lower bound diverges,
regardless of \(m_j\).  A finite physical rate therefore requires a
near-critical microscopic coefficient or a genuinely different
cancellation-compatible block estimate.
\end{corollary}

\begin{proof}
Insert \(q_j\le\alpha_j^{m_j}\) and
\(\delta_j=m_j\tau_j\) in \eqref{eq:s6v50-muj}.
\end{proof}

\begin{assessment}[Spatial MSEC is not automatically a time slab]
\label{ass:s6v50-incidence}
The V43 MSEC coefficient is obtained by exploring a degree-eighteen link
graph.  It becomes a \(P_j\) coefficient in
\eqref{eq:s6v50-q} only after an explicit reflected time-cylinder map proves
that one exploration block is the positive physical transfer step and that
its quotient metric controls the OS Hilbert norm on a dense source core.
Graph distance, link-update count, and Euclidean time are not interchangeable
by notation.
\end{assessment}

\subsection{Continuum spectral handoff}
\label{sec:s6v50-handoff}

\begin{theorem}[Near-critical OS gap handoff]
\label{thm:s6v50-gap-handoff}
Assume:
\begin{enumerate}[label=\textup{(\roman*)},leftmargin=3.5em]
\item the reflected finite-regulator theories have positive normalized
      transfer steps \(P_j\) of physical duration \(\delta_j\to0\);
\item a common local OS source core is transported projectively and its
      delayed diagonal Grams converge for every fixed physical time \(t\);
\item for every vacuum-orthogonal core vector \(F\),
      \[
       \langle F,P_j^nF\rangle
       \le C_Fq_j^n
      \]
      whenever \(n\delta_j\) remains in a compact time interval;
\item \(q_j=1-\mu\delta_j+o(\delta_j)\) for one \(\mu>0\); and
\item the transported source core is dense in the limiting
      vacuum-orthogonal OS Hilbert space.
\end{enumerate}
Then the limiting positive semigroup \(e^{-tH}\) has
\begin{equation}
 \operatorname{spec}(H|_{{\cal H}_0})\subset[\mu,\infty).
 \label{eq:s6v50-limit-gap}
\end{equation}
\end{theorem}

\begin{proof}
For \(n_j=\lfloor t/\delta_j\rfloor\),
\[
 q_j^{n_j}
 =\exp\left\{
  {t+o(1)\over\delta_j}\log q_j
 \right\}
 \longrightarrow e^{-\mu t}.
\]
Gram convergence gives
\[
 \langle F,e^{-tH}F\rangle\le C_Fe^{-\mu t}
\]
for every core vector.  Let \(E_H\) be the spectral measure of \(H\).
If a nonzero core vector had spectral mass in
\([0,\mu-\varepsilon]\), its diagonal semigroup matrix element would be
bounded below by a positive multiple of
\(e^{-(\mu-\varepsilon)t}\), contradicting the preceding estimate as
\(t\to\infty\).  Density of the core and closedness of spectral subspaces
then exclude spectrum below \(\mu\) on all of \({\cal H}_0\).
\end{proof}

\begin{remark}[Uniformity in physical time]
\label{rem:s6v50-time-uniform}
Item~\textup{(iii)} must be applicable for arbitrarily large fixed \(t\)
before the spectral contradiction is taken.  Convergence on one bounded
time interval alone proves no spectral gap.
\end{remark}

\subsection{Inverse-margin amplification and adjacent projectivity}
\label{sec:s6v50-projectivity}

Let \(d\) be a complete Wasserstein-type metric on probability laws over a
common transported slab state space.

\begin{lemma}[Invariant-law perturbation with the exact margin]
\label{lem:s6v50-invariant}
Let \(P,Q\) be Markov kernels with invariant laws \(\pi,\sigma\).  Suppose
\begin{align}
 d(\nu P,\nu'P)&\le q\,d(\nu,\nu'),
 \notag\\
 d(\nu Q,\nu'Q)&\le q\,d(\nu,\nu'),
 \qquad 0\le q<1,
 \label{eq:s6v50-kernel-contract}\\
 \sup_x d(P(x,\cdot),Q(x,\cdot))&\le\varepsilon.
 \label{eq:s6v50-kernel-defect}
\end{align}
Then
\begin{equation}
 \boxed{d(\pi,\sigma)\le{\varepsilon\over1-q}.}
 \label{eq:s6v50-invariant-bound}
\end{equation}
\end{lemma}

\begin{proof}
Invariance, the triangle inequality, and the two hypotheses give
\[
 d(\pi,\sigma)
 =d(\pi P,\sigma Q)
 \le d(\pi P,\sigma P)+d(\sigma P,\sigma Q)
 \le qd(\pi,\sigma)+\varepsilon.
\]
Rearrange.
\end{proof}

\begin{corollary}[Near-critical projectivity debit]
\label{cor:s6v50-projectivity}
For adjacent transported Yang--Mills slab kernels with contraction
\(q_j\) and kernel defect \(\varepsilon_j\), define
\begin{equation}
 \Lambda_j^{\rm YM}
 :={\varepsilon_j\over1-q_j}.
 \label{eq:s6v50-ledger}
\end{equation}
If
\begin{equation}
 \sum_j\Lambda_j^{\rm YM}<\infty,
 \label{eq:s6v50-ledger-sum}
\end{equation}
then the transported invariant laws form a Cauchy sequence.  In the
physical regime \eqref{eq:s6v50-critical-row}, the required summability is
equivalent up to bounded factors to
\begin{equation}
 \sum_j{\varepsilon_j\over\delta_j}<\infty.
 \label{eq:s6v50-scaled-ledger}
\end{equation}
\end{corollary}

\begin{proof}
\Cref{lem:s6v50-invariant} bounds the adjacent law distance by
\(\Lambda_j^{\rm YM}\).  Summing adjacent distances proves the Cauchy
claim.  Under \eqref{eq:s6v50-critical-row},
\((1-q_j)/\delta_j\to\mu>0\), proving the equivalence of the two positive
series up to fixed factors.
\end{proof}

\begin{proposition}[Unweighted adjacent convergence is insufficient]
\label{prop:s6v50-unweighted}
There exist sequences \(\delta_j\downarrow0\),
\(q_j=1-\delta_j\), and \(\varepsilon_j>0\) for which
\(\sum_j\varepsilon_j<\infty\) but
\(\sum_j\varepsilon_j/(1-q_j)=\infty\).
\end{proposition}

\begin{proof}
Take \(\delta_j=2^{-j}\) and
\(\varepsilon_j=2^{-j}/j\).  Then
\(\sum_j\varepsilon_j<\infty\), while
\[
 {\varepsilon_j\over1-q_j}={1\over j},
\]
whose series diverges.
\end{proof}

\subsection{Regulator independence requires bridges}
\label{sec:s6v50-bridges}

\begin{definition}[Yang--Mills projective bridge card, YMPBC]
\label{def:s6v50-YMPBC}
For two cofinal regulator chains \(r_j,s_j\), YMPBC requires common
refinements \(t_j\succeq r_j,s_j\) and transported law distances
\begin{equation}
 b_j^{(r)}=d(\pi_{r_j}\!\to t_j,\pi_{t_j}),
 \qquad
 b_j^{(s)}=d(\pi_{s_j}\!\to t_j,\pi_{t_j})
 \label{eq:s6v50-bridge}
\end{equation}
such that \(b_j^{(r)}+b_j^{(s)}\to0\).  For differentiated source or stress
rows, the same statement is required in the corresponding complete
weighted dual norm.
\end{definition}

\begin{theorem}[Common-refinement uniqueness]
\label{thm:s6v50-bridge}
If the transported laws along \(r_j\) and \(s_j\) converge and YMPBC holds,
then their limits agree on every common local observable.
\end{theorem}

\begin{proof}
For every unit-Lipschitz common observable, the difference of its two
expectations is bounded by
\[
 b_j^{(r)}+b_j^{(s)}
\]
plus the two distances from the chain laws to their limits.  Let
\(j\to\infty\).
\end{proof}

\begin{assessment}[One chain is not a continuum theory]
\label{ass:s6v50-one-chain}
The ledger \eqref{eq:s6v50-ledger-sum} closes one chosen adjacent sequence.
It does not compare a second block map, gauge pinning, anisotropy, or
subsequence.  YMPBC, or an equivalent directed tail-diameter theorem, is a
separate part of continuum existence.
\end{assessment}

\subsection{The physical critical contraction card}
\label{sec:s6v50-card}

\begin{definition}[Physical critical contraction card, PCCC]
\label{def:s6v50-PCCC}
PCCC consists of:
\begin{enumerate}[label=\textup{(\roman*)},leftmargin=3.5em]
\item an actual reflection-positive time-slab kernel \(P_j\), its physical
      duration \(\delta_j\), and a common vacuum-orthogonal source core;
\item the near-critical norm row
      \(q_j=1-\mu\delta_j+o(\delta_j)\), \(\mu>0\);
\item delayed-Gram convergence for every fixed physical time and density of
      the limiting source core;
\item the inverse-margin projectivity ledger
      \(\sum_j\varepsilon_j/(1-q_j)<\infty\);
\item YMPBC for every admitted regulator route;
\item noncollapse of at least one local source norm and the required
      Osterwalder--Schrader axioms; and
\item compatibility of the full perfect-action owner, CMAC/Jacobian
      remainder, topological sectors, PFS, and response matching with the
      same slab kernel.
\end{enumerate}
\end{definition}

\begin{theorem}[Conditional PCCC mass-gap handoff]
\label{thm:s6v50-PCCC}
If PCCC is populated and its projective limit is a nontrivial
Osterwalder--Schrader Yang--Mills theory whose local source core is cyclic,
then its reconstructed Hamiltonian has vacuum-complement gap at least
\(\mu\).
\end{theorem}

\begin{proof}
Rows~\textup{(i)--(v)} construct a regulator-independent limiting delayed
Gram family.  Row~\textup{(vi)} excludes the zero Hilbert space and permits
OS reconstruction.  Row~\textup{(ii)} and the dense cyclic core give
\Cref{thm:s6v50-gap-handoff}.  Row~\textup{(vii)} ensures that the slab law
is the same physical response-matched Yang--Mills law rather than an
auxiliary comparison kernel.
\end{proof}

\begin{assessment}[Exact frontier after V50]
\label{ass:s6v50-frontier}
The continuum mass mechanism is now correctly scaled, but no physical
\(q_j\) has been constructed.  The existing fixed MSEC factor is a valid
finite-regulator spatial covariance bound and an invalid candidate for a
shrinking-time physical row.  The next work must construct the actual
reflected slab, calculate its near-critical deficit \(1-q_j\), and estimate
the adjacent kernel defect at the stronger \(o(\delta_j)\)-summable scale.
\end{assessment}

\begin{programme}[V51: time-cylinder ownership]
\label{prog:s6v50-V51}
\begin{enumerate}[leftmargin=3em]
\item Identify the reflected half-slab kernel in the attached Yang--Mills
      and Grand Tensor source layers and freeze its physical time
      normalization.
\item Prove which exact action owner generates that kernel: bare Wilson,
      trace-shell perfect action, or response-matched action.
\item Express its vacuum-complement coefficient \(q_j\) as a Rayleigh or
      conductance quantity on the physical OS carrier.
\item Test the near-critical expansion
      \(1-q_j\sim\mu\delta_j\); if the source supplies only a fixed graph
      contraction, return the missing time-incidence map rather than reuse
      it.
\item Write the adjacent kernel defect \(\varepsilon_j\) in the same norm
      and test the amplified ledger
      \eqref{eq:s6v50-scaled-ledger}.
\end{enumerate}
\end{programme}

\begin{firewall}[Claim boundary after sixth-series V50]
\label{fire:s6v50-boundary}
V50 performs the compulsory SM/NS audit; proves the exact finite-slab rate
and physical-clock trichotomy; proves the near-critical OS spectral handoff;
proves the invariant-law inverse-margin perturbation estimate and its
strengthened adjacent projectivity ledger; proves common-refinement
uniqueness; and formulates PCCC with a conditional mass-gap handoff.

V50 does not construct the physical reflected slab or its time-incidence
map, prove the near-critical coefficient, populate the amplified ledger or
YMPBC, close the perfect-action/CMAC/Jacobian/PFS rows, construct the
continuum OS theory, or prove the full Yang--Mills mass gap.
\end{firewall}

\begin{theorem}[Sixth-series V50 result ledger]
\label{thm:s6v50-results}
The positive conclusions in \Cref{fire:s6v50-boundary} follow from
\Cref{thm:s6v50-audit-decision,
thm:s6v50-finite-rate,
thm:s6v50-trichotomy,
cor:s6v50-micro,
thm:s6v50-gap-handoff,
lem:s6v50-invariant,
cor:s6v50-projectivity,
prop:s6v50-unweighted,
thm:s6v50-bridge,
thm:s6v50-PCCC}.
\end{theorem}

\begin{proof}
Each ledger item is the conclusion of the cited result with its hypotheses
unchanged.
\end{proof}

\paragraph{Parent-version record.}
V50 was formed from
\texttt{Renewal\_Yang\_Mills\_Mass\_Gap\_Research\_Notes\_V49.tex}.
The V49 parent had \(22\,057\) logical source lines,
\(769\,656\) bytes, and SHA--256
\[
 \texttt{fbb6ef594d5ac09f3d02906c53288e09589001489650cd6834e6435c9b0c8c6a}.
\]
The next compulsory companion audit is V55.

% =============================================================================
% END APPENDED SIXTH-SERIES V50 MATERIAL
% =============================================================================

% =============================================================================
% BEGIN APPENDED SIXTH-SERIES V51 MATERIAL
% =============================================================================

\section{V51: the exact half-slab edge, clock normalization, and the
Gibbs/transfer separation}
\label{sec:s6v51}

\subsection{Why the V50 frontier needs one correction}
\label{sec:s6v51-correction}

\begin{assessment}[Correction to the V50 source assessment]
\label{ass:s6v51-correction}
The Grand Tensor source does contain a finite-regulator half-slab packet.
In its definition of the physical pure-colour half-slab amplitude, a
Hilbert--Schmidt integral operator \(\mathcal A\) is formed from the finite
Wilson half-slab kernel, the normalized full-slab transfer is obtained from
\(\mathcal A^*\mathcal A\), and a physical full-slab duration \(t>0\) is
declared.  Its half-slab Doob--Schmidt theorem then proves
\[
 \rho_{1/2}={s_2(\mathcal A)\over s_1(\mathcal A)},\qquad
 \|T_{\rm norm}|_{\Omega^\perp}\|=\rho_{1/2}^2,\qquad
 m_{\rm FV}=-{2\over t}\log\rho_{1/2}.
 \tag{51.1}
\]
Thus the first sentence of
\Cref{ass:s6v50-frontier} is too broad if read at one fixed finite Wilson
regulator.

The unresolved ownership statement is narrower and more consequential:
the sources do not prove that this bare finite Wilson amplitude is the same
kernel as the response-matched perfect-action owner after the cofinal
shell integrations used in V37--V49; nor do they prove a regulator-uniform
near-critical estimate for its first \emph{relative} singular edge, the
amplified adjacent-defect ledger, or common-refinement independence.
\end{assessment}

\begin{proof}
The finite statement is the direct content of the source theorem labelled
\texttt{thm:YM-half-slab-Doob}.  The same source explicitly says that every
finite compact Wilson half-slab is Hilbert--Schmidt and that its full-slab
transfer is positive trace class.  By contrast, the reduced-fibre Yang--Mills
source states that its exact Witten lift does not identify an arbitrary
half-slab amplitude, boundary instrument, or spatial route.  Its later
same-history incidence theorem begins with the conditional clause that one
reflected slab carries both commands; it proves the resulting Gram identity
but not the existence of that common slab or its cofinal compatibility.
These statements establish exactly the distinction above.
\end{proof}

\subsection{Operator ownership of the finite half-slab mass}
\label{sec:s6v51-operator}

Let \({\cal H}_0,{\cal H}_{1/2}\) be Hilbert spaces and let
\(\mathcal A:{\cal H}_0\to{\cal H}_{1/2}\) be a nonzero compact operator.
Write its singular values in decreasing order as
\[
 s_1\ge s_2\ge\cdots\ge0.
\]
Assume \(s_1\) is simple and let \(u_0\) be its normalized right singular
vector.  Define
\begin{equation}
 T_{\mathcal A}:=s_1^{-2}\mathcal A^*\mathcal A,
 \qquad {\cal H}_{0,\circ}:=u_0^\perp.
 \label{eq:s6v51-normal-transfer}
\end{equation}

\begin{theorem}[Exact half-slab to full-slab spectral dictionary]
\label{thm:s6v51-dictionary}
The operator \(T_{\mathcal A}\) is a positive compact contraction,
\(T_{\mathcal A}u_0=u_0\), and
\begin{equation}
 \boxed{
 q_{\mathcal A}:=
 \|T_{\mathcal A}|_{{\cal H}_{0,\circ}}\|
 =\left({s_2\over s_1}\right)^2.}
 \label{eq:s6v51-q-rho}
\end{equation}
If one full slab has physical duration \(t>0\), then
\begin{equation}
 \boxed{
 -{1\over t}\log q_{\mathcal A}
 =-{2\over t}\log{s_2\over s_1}.}
 \label{eq:s6v51-mass}
\end{equation}
When \(\mathcal A\) is Hilbert--Schmidt,
\(T_{\mathcal A}\) is trace class.
\end{theorem}

\begin{proof}
The singular-value decomposition diagonalizes
\(\mathcal A^*\mathcal A\) with eigenvalues \(s_k^2\).  Division by
\(s_1^2\) gives eigenvalue one on \(u_0\), and simplicity of \(s_1\)
removes every further unit eigenvector.  The largest eigenvalue on
\(u_0^\perp\) is \(s_2^2/s_1^2\), proving
\eqref{eq:s6v51-q-rho}.  The spectral definition
\(H=-t^{-1}\log T_{\mathcal A}\) gives
\eqref{eq:s6v51-mass}.  Finally, a product of two Hilbert--Schmidt
operators is trace class.
\end{proof}

\begin{remark}[Sector-relative version]
\label{rem:s6v51-relative}
If deterministic centre, topology, harmonic, or boundary data are readable
from both half-boundaries, the unshorted carrier can contain additional
unit correlations.  Then \({\cal H}_{0,\circ}\) in
\Cref{thm:s6v51-dictionary} must be replaced by the relative carrier after
the maximal common factor has been shorted.  Small probability of a common
record does not change its operator norm from one.
\end{remark}

\subsection{The Gibbs gap is a diagnostic, not the transfer Hamiltonian}
\label{sec:s6v51-Gibbs}

Let \(\Lambda\) be a two-boundary probability law on
\((X_0,X_{1/2})\), let \(\mathsf P_i\) be conditional expectation onto the
\(i\)-th boundary sigma-field, and remove their common range.  On the
relative carrier put
\[
 \rho:=\|\mathsf P_1\mathsf P_0\|,
 \qquad
 G:={1\over2}(\mathsf P_0+\mathsf P_1).
\]

\begin{theorem}[Exact transfer/Gibbs conversion]
\label{thm:s6v51-Gibbs-conversion}
Suppose \(\Lambda\) is the Doob law of a positive half-slab amplitude
\(\mathcal A\), and \(\rho=s_2(\mathcal A)/s_1(\mathcal A)\) on the
shorted carrier.  Then
\begin{align}
 \|T_{\mathcal A}|_{\rm rel}\|&=\rho^2,
 \label{eq:s6v51-transfer-edge}\\
 \inf\operatorname{spec}(I-G)|_{\rm rel}
 &=:\lambda_{01}={1-\rho\over2},
 \label{eq:s6v51-Gibbs-gap}\\
 m_{\rm FV}
 &=-{2\over t}\log(1-2\lambda_{01}).
 \label{eq:s6v51-Gibbs-mass}
\end{align}
Nevertheless \(G\) is an alternating-boundary Gibbs diagnostic on
\(L^2(\Lambda)\), whereas \(T_{\mathcal A}\) is the physical full-slab
transfer on the initial-boundary Hilbert space.  Equality of the numerical
edges through \eqref{eq:s6v51-Gibbs-mass} does not identify the two
operators or their clocks.
\end{theorem}

\begin{proof}
The first identity is \Cref{thm:s6v51-dictionary}.  On each generic
two-projection block the canonical representation is
\[
 \mathsf P_0=
 \begin{pmatrix}1&0\\0&0\end{pmatrix},
 \qquad
 \mathsf P_1=
 \begin{pmatrix}c^2&c(1-c^2)^{1/2}\\
 c(1-c^2)^{1/2}&1-c^2\end{pmatrix},
 \qquad 0\le c\le\rho.
\]
The two eigenvalues of
\((\mathsf P_0+\mathsf P_1)/2\) are
\((1\pm c)/2\).  Taking the supremum over the relative blocks gives
\(\|G\|_{\rm rel}=(1+\rho)/2\), which is
\eqref{eq:s6v51-Gibbs-gap}.  Since
\(\rho=1-2\lambda_{01}\), substitution in
\eqref{eq:s6v51-mass} proves \eqref{eq:s6v51-Gibbs-mass}.
The operators act on the displayed different Hilbert spaces and have
different spectra even on a single generic block; only their extremal
numbers are related as stated.
\end{proof}

\begin{corollary}[Near-critical conversion factors]
\label{cor:s6v51-critical-conversion}
For a sequence of full-slab durations \(\delta_j\downarrow0\), the following
are equivalent for every fixed \(\mu\in(0,\infty)\):
\begin{align}
 q_j&=1-\mu\delta_j+o(\delta_j),
 \label{eq:s6v51-q-critical}\\
 \rho_j&=1-{\mu\over2}\delta_j+o(\delta_j),
 \label{eq:s6v51-rho-critical}\\
 \lambda_{01,j}&={\mu\over4}\delta_j+o(\delta_j).
 \label{eq:s6v51-lambda-critical}
\end{align}
In particular, a regulator-independent positive lower bound on the
dimensionless Gibbs gap \(\lambda_{01,j}\) is not the desired shrinking-time
regime; it would yield an infinite physical mass scale.
\end{corollary}

\begin{proof}
Use \(q_j=\rho_j^2\) and
\(\lambda_{01,j}=(1-\rho_j)/2\).  Expanding the square and square root at
one proves all equivalences.  The last statement follows from
\eqref{eq:s6v51-Gibbs-mass} divided by \(\delta_j\).
\end{proof}

\subsection{Stable singular edges and the correct perturbative scale}
\label{sec:s6v51-stability}

\begin{lemma}[Relative singular-ratio stability]
\label{lem:s6v51-ratio-stability}
Let \(A,B\) be compact operators between the same Hilbert spaces.  If
\(s_1(A)\ge\sigma>0\) and
\(\varepsilon:=\|A-B\|<\sigma\), then
\begin{equation}
 \left|
 {s_2(A)\over s_1(A)}-{s_2(B)\over s_1(B)}
 \right|
 \le {2\varepsilon\over\sigma}.
 \label{eq:s6v51-ratio-stability}
\end{equation}
Consequently their squared ratios differ by at most
\begin{equation}
 \left|
 \left({s_2(A)\over s_1(A)}\right)^2
 -\left({s_2(B)\over s_1(B)}\right)^2
 \right|
 \le {4\varepsilon\over\sigma}.
 \label{eq:s6v51-q-stability}
\end{equation}
\end{lemma}

\begin{proof}
The min--max principle for singular values gives
\(\lvert s_k(A)-s_k(B)\rvert\le\varepsilon\) for \(k=1,2\).
Since \(s_1(B)>0\), insert the intermediate ratio
\(s_2(B)/s_1(A)\):
\[
 \left|{s_2(A)\over s_1(A)}
       -{s_2(B)\over s_1(B)}\right|
 \le {\varepsilon\over s_1(A)}
 +{s_2(B)\lvert s_1(B)-s_1(A)\rvert
       \over s_1(A)s_1(B)}
 \le {2\varepsilon\over\sigma}.
\]
Both ratios lie in \([0,1]\), so
\(\lvert x^2-y^2\rvert\le2\lvert x-y\rvert\), proving
\eqref{eq:s6v51-q-stability}.
\end{proof}

\begin{assessment}[The normalization is part of the estimate]
\label{ass:s6v51-normalization}
The small parameter in \Cref{lem:s6v51-ratio-stability} is the
dimensionless relative defect
\(\|A-B\|/s_1(A)\), not the unnormalized amplitude defect.  A scalar
partition-function factor may make \(\|A-B\|\) small or large without
changing either normalized channel.  Any amplitude comparison used in the
mass programme must therefore control the defect after leading-singular
normalization, or state an equivalent relative estimate.
\end{assessment}

\begin{lemma}[Telescoping at fixed physical time]
\label{lem:s6v51-telescoping}
Let \(P_j,Q_j\) be contractions and let one step have duration
\(\delta_j\downarrow0\).  If
\[
 \|P_j-Q_j\|=\eta_j,
 \qquad n_j=\lfloor t/\delta_j\rfloor,
\]
then
\begin{equation}
 \boxed{\|P_j^{n_j}-Q_j^{n_j}\|\le n_j\eta_j.}
 \label{eq:s6v51-telescope}
\end{equation}
Hence \(\eta_j=o(\delta_j)\) makes the two fixed-\(t\) channels
operator-norm indistinguishable, while an \(O(\delta_j)\) one-step
deformation can accumulate to an order-one fixed-time change.
\end{lemma}

\begin{proof}
The identity
\[
 P^n-Q^n=\sum_{\ell=0}^{n-1}
 P^{n-1-\ell}(P-Q)Q^\ell
\]
and contractivity give \(\|P^n-Q^n\|\le n\|P-Q\|\).
Now \(n_j\delta_j\to t\), proving both conclusions.
\end{proof}

\begin{theorem}[No mass generation from a sub-generator perturbation]
\label{thm:s6v51-no-subgenerator}
Let \(Q_j\) be reference transfer contractions with physical steps
\(\delta_j\downarrow0\), and suppose a transported normalized mode
\(f_j\) satisfies
\[
 \|f_j\|=1,\qquad
 \|Q_j^{\lfloor t/\delta_j\rfloor}f_j\|\longrightarrow1
\]
for every fixed \(t>0\).  If \(P_j\) are physical transfer contractions
with
\begin{equation}
 \|P_j-Q_j\|=o(\delta_j),
 \label{eq:s6v51-subgenerator}
\end{equation}
then
\[
 \|P_j^{\lfloor t/\delta_j\rfloor}f_j\|\longrightarrow1
\]
for every fixed \(t>0\).  Such a comparison cannot prove a positive
limiting mass on this infrared mode.
\end{theorem}

\begin{proof}
By the reverse triangle inequality and
\Cref{lem:s6v51-telescoping},
\[
 \left|
 \|P_j^{n_j}f_j\|-\|Q_j^{n_j}f_j\|
 \right|
 \le \|P_j^{n_j}-Q_j^{n_j}\|
 \le n_j o(\delta_j)=o(1).
\]
The asserted limit follows.  A semigroup with gap \(\mu>0\) on the
vacuum-orthogonal carrier would instead contract every such normalized
mode by at most \(e^{-\mu t}<1\), giving the contradiction.
\end{proof}

\begin{assessment}[Reconciling fixed-slab infrared saturation with V50]
\label{ass:s6v51-infrared}
The Grand Tensor Gaussian saturation theorem is correct at fixed positive
half-slab duration: along \(k\to0\), the massless Gaussian channel norm
tends to one, so a uniformly massive fixed-time channel must differ from it
by an order-one norm.  For a shrinking elementary duration
\(\delta_j\), however, a finite mass changes a one-step edge only by
order \(\delta_j\).  The two statements are reconciled by
\Cref{lem:s6v51-telescoping}: an order-\(\delta_j\) generator-level change
accumulates to the required order-one fixed-time reorganization, whereas an
\(o(\delta_j)\) change cannot.
\end{assessment}

\subsection{The exact owner and the remaining cofinal rows}
\label{sec:s6v51-owner}

\begin{definition}[Half-slab ownership and clock card, HSOCC]
\label{def:s6v51-HSOCC}
At regulator \(j\), HSOCC contains:
\begin{enumerate}[label=\textup{(\roman*)},leftmargin=3.5em]
\item the initial and half-slab physical boundary spaces, measures, and
      gauge/common-sector quotient;
\item the positive reflected amplitude \(A_j\), its leading singular
      normalization, and proof that \(A_j^*A_j\) is the physical transfer;
\item the exact action producing \(A_j\), including every shell return,
      Jacobian, counterterm, and response-matching term;
\item the full-slab physical duration \(\delta_j\), with its anisotropic
      lattice units and refinement law;
\item the relative half-slab edge
      \(\rho_j=s_2(A_j)/s_1(A_j)\) after common-factor extraction;
\item the estimate
      \(1-\rho_j=(\mu/2)\delta_j+o(\delta_j)\), uniformly in volume and
      every admitted cofinal route;
\item an adjacent normalized-amplitude or Markov-kernel defect strong
      enough for the inverse-margin projectivity ledger; and
\item fixed-time delayed-Gram convergence on a common dense physical
      source core.
\end{enumerate}
\end{definition}

\begin{proposition}[What the attached half-slab theorem populates]
\label{prop:s6v51-source-population}
For each fixed finite compact Wilson regulator with a declared duration
\(t\), the attached Grand Tensor half-slab theorem populates the
Hilbert--Schmidt/trace-class part of HSOCC rows~\textup{(i)--(ii)}, the bare
Wilson instance of row~\textup{(iii)}, the declared finite clock in
row~\textup{(iv)}, and the exact spectral dictionary in row~\textup{(v)}.
It does not by itself populate rows~\textup{(vi)--(viii)} or identify the
bare amplitude with the response-matched perfect-action amplitude.
\end{proposition}

\begin{proof}
Compactness and bounded continuity of the finite Wilson action make its
kernel square-integrable; hence its integral operator is Hilbert--Schmidt
and its reflected square is positive trace class.  Positivity improvement
at fixed regulator gives the simple leading singular pair and the
Doob--Schmidt dictionary, proving the positive assertions.  The theorem's
hypotheses and conclusion contain no uniform volume/refinement asymptotic,
no estimate of \(1-\rho_j\), no adjacent-regulator defect, and no
fixed-time cofinal Gram convergence.  The reduced-fibre source's explicit
incidence caveat rules out inferring the perfect-action identification from
its Witten lift alone.
\end{proof}

\begin{theorem}[Conditional half-slab completion criterion]
\label{thm:s6v51-HSOCC}
Assume HSOCC is populated along every admitted cofinal regulator route,
the common-refinement bridges of V50 hold, and the limiting OS source core
is nontrivial and cyclic.  Then the reconstructed Yang--Mills Hamiltonian
has vacuum-complement gap at least \(\mu\).
\end{theorem}

\begin{proof}
By rows~\textup{(ii)} and \textup{(v)},
\(q_j=\rho_j^2\) is the physical full-slab coefficient.  Row~\textup{(vi)}
and \Cref{cor:s6v51-critical-conversion} give
\[
 q_j=1-\mu\delta_j+o(\delta_j).
\]
Rows~\textup{(vii)--(viii)} and the common-refinement bridges produce a
regulator-independent limiting delayed-Gram semigroup on the common core.
The V50 near-critical OS gap handoff applies, and cyclic density extends
the spectral lower bound from the core to the full vacuum-orthogonal
physical Hilbert space.
\end{proof}

\begin{assessment}[Upstream bottleneck after V51]
\label{ass:s6v51-frontier}
The half-slab operator and its finite-regulator mass formula are no longer
missing.  The first genuinely unpopulated quantitative row is
\[
 {1-\rho_j\over\delta_j}\longrightarrow{\mu\over2}>0
\]
uniformly after common-factor extraction and volume removal, for the
\emph{same} response-matched amplitude that obeys the projective defect
ledger.  Neither a fixed spatial Dobrushin coefficient nor finite-volume
positivity improvement estimates this ratio.  The next iteration should
move upstream to the variational derivative of the singular edge under the
exact signed Wilsonian head, because that is where an order-\(\delta_j\)
generator correction can be distinguished from an \(o(\delta_j)\)
perturbation.
\end{assessment}

\begin{programme}[V52: singular-edge response and owner matching]
\label{prog:s6v51-V52}
\begin{enumerate}[leftmargin=3em]
\item Prove a Feynman--Hellmann formula for the leading and first relative
      eigenvalues of \(A_\theta^*A_\theta\), with an explicit degeneracy
      firewall for \(s_2\).
\item Express the derivative of
      \(-2\log(s_2/s_1)\) as a difference of physical half-slab score
      expectations.
\item Insert the exact signed Wilsonian head and identify which shell
      returns cancel from the ratio and which survive at generator order.
\item Determine whether the available response-matching bounds are
      \(O(\delta_j)\) or \(o(\delta_j)\) in normalized operator norm.
\item If differentiability or edge isolation fails, replace the scalar
      derivative by min--max one-sided bounds on the complete relative
      spectral cluster.
\end{enumerate}
\end{programme}

\begin{firewall}[Claim boundary after sixth-series V51]
\label{fire:s6v51-boundary}
V51 corrects the finite-regulator source assessment; proves the exact
half-slab/full-slab spectral dictionary; separates the two-boundary Gibbs
diagnostic from the physical transfer; computes all near-critical
conversion factors; proves normalized singular-ratio stability; proves the
fixed-time telescoping threshold and the no-sub-generator mass theorem; and
defines the precise HSOCC completion criterion.

V51 does not estimate the relative singular edge of four-dimensional
Yang--Mills uniformly in volume or regulator, identify the response-matched
perfect amplitude with the bare Wilson amplitude, populate the amplified
projectivity ledger, prove the OS continuum limit, or prove the full
Yang--Mills mass gap.
\end{firewall}

\begin{theorem}[Sixth-series V51 result ledger]
\label{thm:s6v51-results}
The positive conclusions in \Cref{fire:s6v51-boundary} are exactly those of
\Cref{ass:s6v51-correction,
thm:s6v51-dictionary,
thm:s6v51-Gibbs-conversion,
cor:s6v51-critical-conversion,
lem:s6v51-ratio-stability,
lem:s6v51-telescoping,
thm:s6v51-no-subgenerator,
prop:s6v51-source-population,
thm:s6v51-HSOCC}.
\end{theorem}

\begin{proof}
Each item is the conclusion of the cited result, with its stated
hypotheses and sector qualifications unchanged.
\end{proof}

\paragraph{Parent-version record.}
V51 was formed from
\texttt{Renewal\_Yang\_Mills\_Mass\_Gap\_Research\_Notes\_V50.tex}.
The V50 parent had \(22\,530\) logical source lines,
\(786\,893\) bytes, and SHA--256
\[
 \texttt{05097b755029d69f6d7c08aa87e6b98b28e91cd68473864dc730e9adb32018d6}.
\]
The next compulsory companion audit remains V55.

% =============================================================================
% END APPENDED SIXTH-SERIES V51 MATERIAL
% =============================================================================

% =============================================================================
% BEGIN APPENDED SIXTH-SERIES V52 MATERIAL
% =============================================================================

\section{V52: singular-edge response, signed mode scores, and the
degenerate-clock firewall}
\label{sec:s6v52}

\subsection{Differentiable half-slab amplitudes}
\label{sec:s6v52-amplitudes}

Let \(I\subset\mathbb R\) be open and let
\[
 A_\theta:{\cal H}_0\longrightarrow{\cal H}_{1/2}
\]
be a norm-\(C^1\) family of compact operators.  Put
\[
 B_\theta=A_\theta^*A_\theta.
\]
Whenever \(\lambda_k(\theta)>0\) is a simple eigenvalue of \(B_\theta\),
choose normalized singular vectors \(u_k(\theta),v_k(\theta)\) so that
\[
 A_\theta u_k=s_kv_k,\qquad
 A_\theta^*v_k=s_ku_k,\qquad
 \lambda_k=s_k^2.
\]

\begin{theorem}[Feynman--Hellmann formula for a simple singular edge]
\label{thm:s6v52-FH}
At every parameter where \(\lambda_k\) is simple,
\begin{align}
 \lambda_k'
 &=\langle u_k,B_\theta'u_k\rangle
 =2\operatorname{Re}\langle A_\theta u_k,A_\theta'u_k\rangle,
 \label{eq:s6v52-lambda-prime}\\
 {d\over d\theta}\log s_k
 &={\operatorname{Re}\langle v_k,A_\theta'u_k\rangle\over s_k}.
 \label{eq:s6v52-log-s}
\end{align}
If \(s_1\) and the first relative singular value \(s_2\) are both simple,
then the dimensionless one-slab mass
\[
 M_\theta:=-2\log{s_2(\theta)\over s_1(\theta)}
\]
satisfies
\begin{equation}
 M_\theta'
 =2\operatorname{Re}\left\{
 {\langle v_1,A_\theta'u_1\rangle\over s_1}
 -{\langle v_2,A_\theta'u_2\rangle\over s_2}
 \right\}.
 \label{eq:s6v52-M-prime}
\end{equation}
\end{theorem}

\begin{proof}
Norm-\(C^1\) regularity gives
\[
 B_\theta'=A_\theta'^*A_\theta+A_\theta^*A_\theta'.
\]
The ordinary Feynman--Hellmann theorem for the simple eigenvalue of the
self-adjoint compact family \(B_\theta\) gives the first equality in
\eqref{eq:s6v52-lambda-prime}; substitution of \(B_\theta'\) gives the
second.  Since \(\lambda_k=s_k^2>0\),
\[
 {s_k'\over s_k}={\lambda_k'\over2\lambda_k}
 ={\operatorname{Re}\langle s_kv_k,A_\theta'u_k\rangle\over s_k^2},
\]
which is \eqref{eq:s6v52-log-s}.  Subtract the formulas for
\(k=2\) and \(k=1\).
\end{proof}

\subsection{The physical score is signed on the excited edge}
\label{sec:s6v52-score}

Suppose the amplitude is an integral operator with a real positive kernel
\(F_\theta(y,x)\), and assume differentiation under the integral sign is
valid with
\begin{equation}
 \partial_\theta F_\theta(y,x)
 =-{1\over2}J_\theta(y,x)F_\theta(y,x).
 \label{eq:s6v52-score-kernel}
\end{equation}
For a positive singular value \(s_k\), define the mode-score functional
\begin{equation}
 {\mathfrak E}_{k,\theta}(G)
 :={1\over s_k}
 \iint \overline{v_k(y)}G(y,x)F_\theta(y,x)u_k(x)
 \,\mu_{1/2}(\dd y)\mu_0(\dd x).
 \label{eq:s6v52-mode-functional}
\end{equation}

\begin{lemma}[Normalization and sign of the mode functional]
\label{lem:s6v52-mode-functional}
For every \(k\) in \eqref{eq:s6v52-mode-functional},
\[
 {\mathfrak E}_{k,\theta}(1)=1.
\]
For the positivity-improving leading pair, the phases may be chosen so
that \({\mathfrak E}_{1,\theta}\) is expectation under the positive
Doob half-slab law.  For \(k\ge2\), the functional has total mass one but
is generally signed or complex and is not a probability expectation.
\end{lemma}

\begin{proof}
The first identity is
\[
 s_k^{-1}\langle v_k,A_\theta u_k\rangle=1.
\]
Positivity improvement gives strictly positive leading left and right
singular vectors, so the \(k=1\) integrand is a positive probability
density after the displayed normalization.  Every \(u_k\) with \(k\ge2\)
is orthogonal to the strictly positive \(u_1\), and hence cannot itself be
strictly positive or strictly negative.  Thus positivity of the associated
functional is not available in general; complex phases give the analogous
statement in the complex case.
\end{proof}

\begin{theorem}[Exact score-difference response]
\label{thm:s6v52-score-response}
Under \eqref{eq:s6v52-score-kernel}, if \(s_k\) is simple then
\begin{equation}
 {d\over d\theta}\log s_k
 =-{1\over2}\operatorname{Re}{\mathfrak E}_{k,\theta}(J_\theta).
 \label{eq:s6v52-score-log}
\end{equation}
If \(s_1,s_2\) are simple, the one-slab mass response is
\begin{equation}
 \boxed{
 M_\theta'
 =\operatorname{Re}\left[
 {\mathfrak E}_{2,\theta}(J_\theta)
 -{\mathfrak E}_{1,\theta}(J_\theta)\right].}
 \label{eq:s6v52-score-mass}
\end{equation}
In particular, multiplication of \(F_\theta\) by any positive scalar
\(c_\theta\) leaves \eqref{eq:s6v52-score-mass} unchanged.
\end{theorem}

\begin{proof}
Equation \eqref{eq:s6v52-score-kernel} says that
\[
 \langle v_k,A_\theta'u_k\rangle
 =-{s_k\over2}{\mathfrak E}_{k,\theta}(J_\theta).
\]
Insert this identity in \eqref{eq:s6v52-log-s} and then in
\eqref{eq:s6v52-M-prime}.  Replacing \(F_\theta\) by \(c_\theta F_\theta\)
adds the constant
\(-2c_\theta'/c_\theta\) to \(J_\theta\).  It cancels between the two mode
functionals because both send the constant one to one.
\end{proof}

\begin{corollary}[Exact integrated-owner score]
\label{cor:s6v52-integrated-score}
Let
\[
 F_\theta(y,x)=\int
 \exp\{-S_\theta(y,x;\eta)/2\}\,\nu(\dd\eta)
\]
and suppose differentiation is dominated.  Under the conditional interior
law
\[
 \pi_{\theta,y,x}(\dd\eta)
 ={e^{-S_\theta(y,x;\eta)/2}\nu(\dd\eta)
   \over F_\theta(y,x)},
\]
the amplitude score in \eqref{eq:s6v52-score-kernel} is exactly
\begin{equation}
 J_\theta(y,x)
 =\mathbb E_{\pi_{\theta,y,x}}
   [\partial_\theta S_\theta(y,x;\eta)].
 \label{eq:s6v52-integrated-J}
\end{equation}
Therefore every integrated shell return, Jacobian, and counterterm enters
the singular-edge response through the same conditional action owner; it
cannot be fitted afterward from a separate covariance.
\end{corollary}

\begin{proof}
Differentiate the defining integral:
\[
 \partial_\theta F_\theta
 =-{1\over2}\int(\partial_\theta S_\theta)
 e^{-S_\theta/2}\,\nu(\dd\eta)
 =-{1\over2}F_\theta
 \mathbb E_{\pi_{\theta,y,x}}[\partial_\theta S_\theta].
\]
Comparison with \eqref{eq:s6v52-score-kernel} proves the formula.
\end{proof}

\begin{assessment}[Why concentration of the vacuum law is insufficient]
\label{ass:s6v52-signed}
The leading term in \eqref{eq:s6v52-score-mass} is an ordinary expectation
under the positive Doob law, but the excited term is a signed mode matrix
element.  Bounds on probability tails, Dobrushin oscillations, or positive
covariances of the vacuum law do not by themselves give a positive lower
bound on their difference.  A proof must control the first relative
spectral subspace or establish an operator inequality on the entire
vacuum-orthogonal carrier.
\end{assessment}

\subsection{Degenerate first-relative clusters}
\label{sec:s6v52-cluster}

Fix \(\theta_0\).  Suppose the leading eigenvalue
\(\lambda_1(\theta_0)\) of \(B_{\theta_0}\) is simple and the largest
relative eigenvalue \(\lambda_\circ\) is isolated from both
\(\lambda_1\) and the remainder, but may have finite multiplicity \(r\).
Let \(P_\circ\) be its spectral projection and set
\[
 C_\circ:=P_\circ B_{\theta_0}'P_\circ
 \quad\hbox{on }\operatorname{Ran}P_\circ.
\]

\begin{theorem}[One-sided response of a degenerate relative edge]
\label{thm:s6v52-degenerate}
Let
\[
 M(\theta)=\log\lambda_1(\theta)
 -\log\lambda_\circ^{\max}(\theta),
\]
where \(\lambda_\circ^{\max}\) denotes the largest eigenvalue in the
continuing isolated relative cluster.  Then
\begin{align}
 D_+M(\theta_0)
 &={\lambda_1'(\theta_0)\over\lambda_1(\theta_0)}
   -{\lambda_{\max}(C_\circ)\over\lambda_\circ},
 \label{eq:s6v52-Dplus}\\
 D_-M(\theta_0)
 &={\lambda_1'(\theta_0)\over\lambda_1(\theta_0)}
   -{\lambda_{\min}(C_\circ)\over\lambda_\circ}.
 \label{eq:s6v52-Dminus}
\end{align}
Thus differentiability of the ordered edge is equivalent to
\(C_\circ\) being scalar on the cluster, unless an accidental equality
among its extreme eigenvalues occurs.
\end{theorem}

\begin{proof}
The Riesz projection of the isolated cluster varies continuously in
operator norm and reduces the local perturbation problem to an
\(r\)-dimensional Hermitian matrix family.  Its eigenvalue branches have
first derivatives equal to the eigenvalues of \(C_\circ\).  For positive
increments the largest branch therefore has derivative
\(\lambda_{\max}(C_\circ)\); for negative increments it has derivative
\(\lambda_{\min}(C_\circ)\).  The leading eigenvalue is simple, so ordinary
Feynman--Hellmann applies to it.  Differentiate the two logarithms
one-sidedly to obtain \eqref{eq:s6v52-Dplus} and
\eqref{eq:s6v52-Dminus}.
\end{proof}

When the amplitude score representation
\eqref{eq:s6v52-score-kernel} holds, choose orthonormal right singular
vectors \(u_a\) for the relative cluster and put
\(v_a=A_{\theta_0}u_a/s_\circ\).  Define
\[
 [R_\circ]_{ab}
 ={1\over s_\circ}
 \iint\overline{v_a(y)}J_{\theta_0}(y,x)
 F_{\theta_0}(y,x)u_b(x)\,
 \mu_{1/2}(\dd y)\mu_0(\dd x),
 \qquad
 S_\circ={R_\circ+R_\circ^*\over2}.
\]

\begin{corollary}[Complete cluster-score bounds]
\label{cor:s6v52-cluster-score}
With \(J_1=\operatorname{Re}{\mathfrak E}_{1,\theta_0}
(J_{\theta_0})\),
\begin{align}
 D_+M(\theta_0)&=\lambda_{\min}(S_\circ)-J_1,
 \label{eq:s6v52-cluster-plus}\\
 D_-M(\theta_0)&=\lambda_{\max}(S_\circ)-J_1.
 \label{eq:s6v52-cluster-minus}
\end{align}
Consequently a basis-free sufficient lower response bound
\(D_+M\ge c\) is the operator inequality
\begin{equation}
 \boxed{S_\circ\ \succeq\ (J_1+c)I_{\operatorname{Ran}P_\circ}.}
 \label{eq:s6v52-cluster-lower}
\end{equation}
\end{corollary}

\begin{proof}
Using \(A_{\theta_0}'=-J_{\theta_0}F_{\theta_0}/2\) at kernel level,
direct expansion gives
\[
 {1\over\lambda_\circ}C_\circ=-S_\circ.
\]
Also \(\lambda_1'/\lambda_1=-J_1\).
Substitute these identities in
\eqref{eq:s6v52-Dplus}--\eqref{eq:s6v52-Dminus}.  The last assertion is
the min--max characterization of
\(\lambda_{\min}(S_\circ)\).
\end{proof}

\begin{remark}[Symmetry multiplets]
\label{rem:s6v52-symmetry}
Colour, cubic rotation, and momentum symmetries can force a nontrivial
relative multiplicity.  The correct target is then the complete
cluster inequality \eqref{eq:s6v52-cluster-lower}, not a selected
eigenvector expectation.  Selecting one vector before proving that the
perturbation respects the full physical symmetry can hide a descending
branch and overstate the gap.
\end{remark}

\subsection{Why differentiation at zero slab width is singular}
\label{sec:s6v52-zero-width}

\begin{theorem}[Simple-edge perturbation cannot start at the
identity clock]
\label{cth:s6v52-identity}
Let \({\cal H}\) have dimension at least two and let
\(P_\delta\to I\) in operator norm as \(\delta\downarrow0\).
At \(\delta=0\), the leading and first-relative eigenvalues of
\(P_0=I\) coincide.  Hence no hypothesis asserting an isolated simple
\(s_2\) at \(\delta=0\) can justify differentiation of the physical
near-critical edge there.  In infinite dimension the unit eigenvalue has
infinite multiplicity.
\end{theorem}

\begin{proof}
Every vector is a unit eigenvector of the identity.  Thus the eigenvalue
one has multiplicity \(\dim{\cal H}\), so neither the leading nor the
first-relative branch is spectrally isolated at zero width.
\end{proof}

\begin{assessment}[Parameter derivatives versus clock derivatives]
\label{ass:s6v52-clock}
The exact formulas of
\Cref{thm:s6v52-score-response,thm:s6v52-degenerate} are valid for a
coupling or shell deformation at a fixed positive slab.  They do not by
themselves compute
\[
 \lim_{\delta\downarrow0}{1-\rho_\delta\over\delta},
\]
because the zero-width amplitude has a completely degenerate edge.
The clock limit must instead be controlled at the level of the rescaled
Dirichlet forms
\[
 {\cal E}_\delta(f)
 ={1\over\delta}\langle f,(I-P_\delta)f\rangle,
\]
or through strong resolvent/semigroup convergence of the generators.
\end{assessment}

\subsection{Owner-matching consequence and revised target}
\label{sec:s6v52-owner}

\begin{theorem}[Exact normalization-free response target]
\label{thm:s6v52-owner-target}
At fixed positive slab width, any proposed bare/perfect amplitude
identification sufficient for singular-edge transport must control the
matrix elements of the conditional owner
\[
 J_\theta(y,x)
 =\mathbb E[\partial_\theta S_\theta\mid y,x]
\]
on the leading singular pair and on the complete first-relative cluster.
Scalar normalization terms cancel exactly.  Terms whose induced normalized
amplitude defect is \(o(\delta)\) cannot alter the limiting physical
generator, while an \(O(\delta)\) term may do so and therefore cannot be
discarded as a small remainder.
\end{theorem}

\begin{proof}
The matrix-element requirement is
\Cref{thm:s6v52-score-response,cor:s6v52-cluster-score}; scalar cancellation
was proved in \Cref{thm:s6v52-score-response}.  The last two statements are
the fixed-time telescoping conclusions of
\Cref{lem:s6v51-telescoping,thm:s6v51-no-subgenerator}.
\end{proof}

\begin{assessment}[Actual result of the V52 owner test]
\label{ass:s6v52-owner-test}
The accumulated DQIS, CMAC, perfect quadratic symbol, and weighted
third-jet estimates are not expressed in normalized half-slab operator norm
with an explicit physical-width factor.  They therefore do not currently
decide whether the bare/perfect amplitude difference is \(O(\delta)\),
\(o(\delta)\), or larger.  The score formula identifies the missing
translation: shell-response bounds must be converted into a uniform
rescaled Dirichlet-form comparison on the common physical carrier.
\end{assessment}

\begin{programme}[V53: rescaled Dirichlet-form route]
\label{prog:s6v52-V53}
\begin{enumerate}[leftmargin=3em]
\item Prove the exact equivalence between a uniform Poincar\'e lower bound
      for \({\cal E}_{\delta,j}\) and the near-critical transfer edge.
\item Formulate a Mosco-type liminf/recovery packet on the transported OS
      source core that survives the identity-edge degeneracy.
\item Determine the precise form perturbation scale: an unscaled
      \(o(\delta)\) transfer defect becomes an \(o(1)\) generator-form
      defect, whereas an \(O(\delta)\) defect must be retained.
\item Insert the perfect quadratic owner and isolate the infrared
      transverse form from gauge, centre, harmonic, and topological common
      factors.
\item State the remaining nonlinear form inequality without replacing it
      by a one-site Dobrushin estimate.
\end{enumerate}
\end{programme}

\begin{firewall}[Claim boundary after sixth-series V52]
\label{fire:s6v52-boundary}
V52 proves simple and degenerate singular-edge response formulas; constructs
the exact signed mode-score functional; proves cancellation of scalar
normalizers; identifies the conditional integrated action as the unique
score owner; gives a basis-free cluster lower-bound target; and proves that
simple-edge differentiation cannot be initiated at the zero-width identity
transfer.

V52 does not prove the cluster score inequality for Yang--Mills, convert the
existing shell estimates to a rescaled OS Dirichlet form, obtain a uniform
near-critical edge, close projectivity, construct the continuum theory, or
prove the full Yang--Mills mass gap.
\end{firewall}

\begin{theorem}[Sixth-series V52 result ledger]
\label{thm:s6v52-results}
The positive conclusions in \Cref{fire:s6v52-boundary} are exactly those of
\Cref{thm:s6v52-FH,
lem:s6v52-mode-functional,
thm:s6v52-score-response,
cor:s6v52-integrated-score,
thm:s6v52-degenerate,
cor:s6v52-cluster-score,
cth:s6v52-identity,
thm:s6v52-owner-target}.
\end{theorem}

\begin{proof}
Each assertion is the conclusion of the cited result with all isolation,
regularity, sector, and clock qualifications retained.
\end{proof}

\paragraph{Parent-version record.}
V52 was formed from
\texttt{Renewal\_Yang\_Mills\_Mass\_Gap\_Research\_Notes\_V51.tex}.
The V51 parent had \(23\,016\) logical source lines,
\(805\,865\) bytes, and SHA--256
\[
 \texttt{c7883b4094f3c201cb7cbf19e45237ccb355d6dc94cd41564aa547cdc0221a7a}.
\]
The next compulsory companion audit remains V55.

% =============================================================================
% END APPENDED SIXTH-SERIES V52 MATERIAL
% =============================================================================

% =============================================================================
% BEGIN APPENDED SIXTH-SERIES V53 MATERIAL
% =============================================================================

\section{V53: rescaled transfer forms, Mosco gap transport, and the
infrared coercivity test}
\label{sec:s6v53}

\subsection{The exact finite-step form dictionary}
\label{sec:s6v53-dictionary}

Let \({\cal H}\) be a Hilbert space, let \(P\) be a positive
self-adjoint contraction, and let \(\Pi\) be the orthogonal projection onto
a closed protected space such that
\[
 P\Pi=\Pi P=\Pi.
\]
The protected space contains the vacuum and may also contain every common
sector deliberately shorted before a relative mass is assigned.  For a
physical time step \(\delta>0\), define the bounded nonnegative form
\begin{equation}
 {\cal E}_{P,\delta}(f,g)
 :={1\over\delta}\langle f,(I-P)g\rangle,
 \qquad
 {\cal E}_{P,\delta}(f):={\cal E}_{P,\delta}(f,f).
 \label{eq:s6v53-form}
\end{equation}

\begin{theorem}[Exact Poincar\'e/edge equivalence]
\label{thm:s6v53-edge-form}
Put
\[
 q=\|P|_{(I-\Pi){\cal H}}\|.
\]
Then
\begin{equation}
 \boxed{
 \inf_{\substack{f\in(I-\Pi){\cal H}\\f\ne0}}
 {{\cal E}_{P,\delta}(f)\over\|f\|^2}
 ={1-q\over\delta}.}
 \label{eq:s6v53-form-gap}
\end{equation}
Consequently
\[
 {\cal E}_{P,\delta}(f)\ge\mu\|(I-\Pi)f\|^2
 \quad\hbox{for all }f
\]
if and only if \(q\le1-\mu\delta\).  The Hamiltonian
\(H_\delta=-\delta^{-1}\log P\) satisfies
\begin{equation}
 \inf\operatorname{spec}
 (H_\delta|_{(I-\Pi){\cal H}})
 ={-\log q\over\delta}
 \ge {1-q\over\delta}.
 \label{eq:s6v53-H-vs-form}
\end{equation}
\end{theorem}

\begin{proof}
On \((I-\Pi){\cal H}\), positivity and self-adjointness give
\(\sup\operatorname{spec}P=q\).  The spectral theorem therefore gives
\[
 \inf\operatorname{spec}(I-P)=1-q,
\]
which is exactly \eqref{eq:s6v53-form-gap}.  The Poincar\'e equivalence
follows.  Applying the decreasing spectral map
\(-\delta^{-1}\log\lambda\) and using
\(-\log q\ge1-q\) proves \eqref{eq:s6v53-H-vs-form}.
\end{proof}

\begin{corollary}[Near-critical form normalization]
\label{cor:s6v53-critical}
For \(\delta_j\downarrow0\) and relative edges \(q_j\to1\),
\[
 q_j=1-\mu\delta_j+o(\delta_j)
\]
is equivalent to
\[
 \inf_{f\perp\operatorname{Ran}\Pi_j}
 {{\cal E}_{P_j,\delta_j}(f)\over\|f\|^2}
 =\mu+o(1).
\]
Thus the identity-edge degeneracy found in V52 disappears after the unique
physical rescaling by \(\delta_j^{-1}\).
\end{corollary}

\begin{proof}
This is the identity \eqref{eq:s6v53-form-gap} applied to each \(j\).
\end{proof}

\subsection{The exact perturbation scale}
\label{sec:s6v53-perturbation}

\begin{lemma}[Transfer norm to generator-form norm]
\label{lem:s6v53-transfer-form}
Let \(P,Q\) be positive self-adjoint contractions on the same Hilbert
space and use the same time step \(\delta\).  Then
\begin{equation}
 \left|{\cal E}_{P,\delta}(f)
       -{\cal E}_{Q,\delta}(f)\right|
 \le{\|P-Q\|\over\delta}\|f\|^2.
 \label{eq:s6v53-form-defect}
\end{equation}
If \(Q\) has a form gap \(\gamma\) on one common relative carrier and
\(\|P-Q\|\le\varepsilon\delta\), then \(P\) has form gap at least
\(\gamma-\varepsilon\) there.
\end{lemma}

\begin{proof}
The form difference is
\(\delta^{-1}\langle f,(Q-P)f\rangle\); the operator-norm bound proves
\eqref{eq:s6v53-form-defect}.  Apply it to unit vectors in the relative
carrier and take the infimum.
\end{proof}

\begin{corollary}[Sharp order bookkeeping]
\label{cor:s6v53-order}
An \(o(\delta_j)\) normalized transfer defect gives an \(o(1)\)
generator-form defect and preserves every strict limiting form gap.
An \(O(\delta_j)\) transfer defect produces an \(O(1)\) form contribution
and must be retained in the limiting generator.  A bound with no explicit
factor \(\delta_j\) is insufficient for continuum gap transport.
\end{corollary}

\begin{proof}
Divide the respective operator defects by \(\delta_j\) in
\eqref{eq:s6v53-form-defect}.
\end{proof}

\begin{proposition}[Carrier motion is a separate debit]
\label{prop:s6v53-carrier}
Let \(\Pi_P,\Pi_Q\) be the protected projections of \(P,Q\).  Even if
\(\|P-Q\|=o(\delta)\), a gap comparison on
\((I-\Pi_P){\cal H}\) and \((I-\Pi_Q){\cal H}\) also requires control of
\(\|\Pi_P-\Pi_Q\|\), or a common-factor quotient constructed before the
comparison.  Operator closeness alone does not identify which unit
eigenvectors are physical vacua and which are protected common sectors.
\end{proposition}

\begin{proof}
Take \(P=Q=I\).  Then \(\|P-Q\|=0\), while arbitrary orthogonal
projections commute with both operators and can declare different
protected carriers.  Hence no conclusion comparing the relative spaces
can follow from the operator defect alone.
\end{proof}

\subsection{Gap transport by variational form convergence}
\label{sec:s6v53-Mosco}

We use the following fixed-carrier formulation after all regulator spaces
have been embedded isometrically into one transported Hilbert space
\({\cal H}\).  A sequence of closed nonnegative forms
\({\cal E}_j\) Mosco-converges to \({\cal E}\) if:
\begin{enumerate}[label=\textup{(M\arabic*)},leftmargin=3.5em]
\item \(f_j\rightharpoonup f\) implies
      \({\cal E}(f)\le\liminf_j{\cal E}_j(f_j)\);
\item for every \(f\) there exist \(f_j\to f\) with
      \(\limsup_j{\cal E}_j(f_j)\le{\cal E}(f)\).
\end{enumerate}
Values outside a form domain are understood as \(+\infty\).

\begin{theorem}[Mosco transport of a protected-sector gap]
\label{thm:s6v53-Mosco-gap}
Let \({\cal E}_j\) Mosco-converge to \({\cal E}\).  Let
\(\Pi_j,\Pi\) be orthogonal projections with
\(\Pi_j\to\Pi\) strongly.  Assume:
\begin{enumerate}[label=\textup{(\roman*)},leftmargin=3.5em]
\item \(\Pi_j\operatorname{Dom}{\cal E}_j\subseteq
      \operatorname{Dom}{\cal E}_j\) and
      \({\cal E}_j(\Pi_j f,g)=0\) for all form-domain vectors;
\item for one \(\mu>0\),
      \[
       {\cal E}_j(f)\ge
       \mu\|(I-\Pi_j)f\|^2
       \quad\hbox{for all }f\in\operatorname{Dom}{\cal E}_j.
      \]
\end{enumerate}
Then
\begin{equation}
 \boxed{
 {\cal E}(f)\ge\mu\|f\|^2
 \quad\hbox{for all }f\in
 \operatorname{Dom}{\cal E}\cap(I-\Pi){\cal H}.}
 \label{eq:s6v53-limit-gap}
\end{equation}
\end{theorem}

\begin{proof}
Take \(f\perp\operatorname{Ran}\Pi\).
Choose a Mosco recovery sequence \(h_j\to f\) satisfying
\(\limsup_j{\cal E}_j(h_j)\le{\cal E}(f)\), and put
\[
 g_j=(I-\Pi_j)h_j.
\]
Strong convergence of the projections, their uniform norm bound, and
\(h_j\to f\) give \(g_j\to f\).  Assumption~\textup{(i)} gives
\({\cal E}_j(g_j)={\cal E}_j(h_j)\), while
assumption~\textup{(ii)} gives
\[
 \mu\|g_j\|^2\le{\cal E}_j(g_j).
\]
Taking the upper limit yields
\(\mu\|f\|^2\le{\cal E}(f)\), as required.
\end{proof}

\begin{remark}[Extension to arbitrary vectors]
\label{rem:s6v53-protected-limit}
If one wants the stronger formula
\({\cal E}(f)\ge\mu\|(I-\Pi)f\|^2\) for arbitrary form-domain vectors,
one must additionally include
\({\cal E}(\Pi f,g)=0\) as an explicit limiting protected-sector
hypothesis.  The spectral gap conclusion in
\Cref{thm:s6v53-Mosco-gap} concerns only
\((I-\Pi){\cal H}\) and needs no such extension.
\end{remark}

\begin{corollary}[Core version]
\label{cor:s6v53-core}
Let \({\cal D}\) be a form core for a closed nonnegative form
\({\cal E}\), invariant under \(I-\Pi\).  If
\[
 {\cal E}(f)\ge\mu\|f\|^2
 \qquad(f\in{\cal D}\cap(I-\Pi){\cal H}),
\]
then the same inequality holds on the full relative form domain.
\end{corollary}

\begin{proof}
Approximate a relative form-domain vector in the form norm by vectors in
\({\cal D}\), apply \(I-\Pi\), and pass to the limit using closedness.
\end{proof}

\begin{assessment}[What Mosco convergence does and does not supply]
\label{ass:s6v53-Mosco}
Mosco convergence transports a gap already proved uniformly at the
regulators; it does not manufacture the inequality in
\Cref{thm:s6v53-Mosco-gap}.  Conversely, convergence of a few Wilson-loop
expectations is weaker than form convergence and does not identify a
limiting Hamiltonian.  The Yang--Mills programme therefore needs both a
uniform relative form inequality and a recovery/liminf packet for the same
response-matched transfer forms.
\end{assessment}

\subsection{The perfect quadratic owner fails the infrared coercivity test}
\label{sec:s6v53-infrared}

\begin{theorem}[Relative-form perturbations cannot open a massless edge]
\label{thm:s6v53-relative-no-gap}
Let \({\cal E}_0\) be a nonnegative form on a relative physical carrier
and suppose there are unit vectors \(f_n\) in its domain with
\[
 {\cal E}_0(f_n)\longrightarrow0.
\]
Let \({\cal R}\) be a symmetric form satisfying
\begin{equation}
 |{\cal R}(f)|\le a{\cal E}_0(f)
 \quad\hbox{for all }f
\label{eq:s6v53-relative-bound}
\end{equation}
for some finite \(a\), and suppose
\({\cal E}={\cal E}_0+{\cal R}\ge0\).  Then \({\cal E}\) has no positive
Poincar\'e gap on that carrier.
\end{theorem}

\begin{proof}
For the infrared sequence,
\[
 0\le{\cal E}(f_n)
 \le(1+a){\cal E}_0(f_n)\longrightarrow0.
\]
Since \(\|f_n\|=1\), the infimum of the Rayleigh quotient of
\({\cal E}\) is zero.
\end{proof}

\begin{corollary}[Vanishing absolute remainder version]
\label{cor:s6v53-vanishing}
For regulator forms
\({\cal E}_j={\cal E}_{0,j}+{\cal R}_j\), assume there are relative unit
vectors \(f_j\) such that
\({\cal E}_{0,j}(f_j)\to0\), and
\[
 |{\cal R}_j(f)|
 \le a{\cal E}_{0,j}(f)+b_j\|f\|^2,
 \qquad b_j\to0.
\]
Then the regulator gaps of \({\cal E}_j\) tend to zero along a subsequence.
\end{corollary}

\begin{proof}
Evaluation on \(f_j\) gives
\[
 0\le{\cal E}_j(f_j)
 \le(1+a){\cal E}_{0,j}(f_j)+b_j\longrightarrow0.
\]
\end{proof}

\begin{assessment}[Application to the V49 perfect quadratic symbol]
\label{ass:s6v53-perfect}
The V49 exact Schur/perfect quadratic owner remains a massless transverse
Gaussian form: its low-momentum Rayleigh quotient tends to zero after the
gauge and common harmonic directions are treated correctly.  Therefore a
nonlinear estimate that is only relatively bounded by that quadratic form,
as in \eqref{eq:s6v53-relative-bound}, cannot open the Yang--Mills mass
gap.  The missing nonperturbative statement must contribute an
\(L^2\)-coercive relative term on the infrared carrier, or reorganize the
physical carrier so that the massless trial sequence is no longer a valid
low-energy sequence.  This is the form version of the fixed-time
order-one channel reorganization found in V51.
\end{assessment}

\begin{theorem}[Necessary infrared coercivity alternative]
\label{thm:s6v53-coercivity-alternative}
Let \({\cal E}_{0,j}\) be the response-matched quadratic baseline and
\({\cal E}_j\) the exact relative transfer form.  If
\(\inf\operatorname{gap}{\cal E}_j\ge\mu>0\) while there are normalized
infrared vectors \(f_j\) with
\({\cal E}_{0,j}(f_j)\to0\), then necessarily
\begin{equation}
 \liminf_j
 \bigl[{\cal E}_j(f_j)-{\cal E}_{0,j}(f_j)\bigr]\ge\mu.
 \label{eq:s6v53-coercive-deformation}
\end{equation}
In particular the nonlinear exact-owner correction is not
form-small on every Gaussian infrared sequence.
\end{theorem}

\begin{proof}
The exact gap gives \({\cal E}_j(f_j)\ge\mu\).  Subtract
\({\cal E}_{0,j}(f_j)\to0\) and take the lower limit.
\end{proof}

\subsection{The rescaled physical form card}
\label{sec:s6v53-card}

\begin{definition}[Rescaled OS Dirichlet-form card, RODFC]
\label{def:s6v53-RODFC}
RODFC consists of:
\begin{enumerate}[label=\textup{(\roman*)},leftmargin=3.5em]
\item the exact response-matched positive transfers \(P_j\), physical
      durations \(\delta_j\), and protected projections \(\Pi_j\);
\item a common transported Hilbert carrier and strong convergence
      \(\Pi_j\to\Pi\);
\item the rescaled forms
      \({\cal E}_j=\delta_j^{-1}\langle\cdot,(I-P_j)\cdot\rangle\);
\item a regulator-, volume-, and route-uniform inequality
      \[
       {\cal E}_j(f)\ge\mu\|(I-\Pi_j)f\|^2;
      \]
\item Mosco liminf and recovery rows for \({\cal E}_j\);
\item an explicit nonlinear infrared estimate compatible with
      \eqref{eq:s6v53-coercive-deformation}, rather than a purely relative
      bound by the massless quadratic owner;
\item the amplified adjacent defect ledger and common-refinement bridges;
      and
\item noncollapse, OS reconstruction, and cyclicity of the limiting local
      source core.
\end{enumerate}
\end{definition}

\begin{theorem}[Conditional RODFC mass-gap handoff]
\label{thm:s6v53-RODFC}
If RODFC is populated, the limiting reconstructed Hamiltonian has
vacuum-complement gap at least \(\mu\).
\end{theorem}

\begin{proof}
Rows~\textup{(i)--(v)} and
\Cref{thm:s6v53-Mosco-gap} produce a limiting closed form with the
relative lower bound \(\mu\).  Rows~\textup{(vii)--(viii)} make the limit
regulator independent and identify its form with the Hamiltonian form of
the nontrivial reconstructed OS theory.  The representation theorem for
closed nonnegative forms then places the relative Hamiltonian spectrum in
\([\mu,\infty)\).
\end{proof}

\begin{assessment}[Exact frontier after V53]
\label{ass:s6v53-frontier}
The unstable scalar singular-edge derivative has been replaced by the
correct rescaled variational target.  The first missing field-theoretic
theorem is now sharp: prove a uniform \(L^2\)-coercive inequality for the
exact response-matched transfer form after gauge and common-sector
shorting.  Existing CMAC and third-jet smallness estimates are useful for
constructing the form and comparing ultraviolet owners, but if they remain
only relative to the massless perfect quadratic form they cannot prove the
infrared gap.
\end{assessment}

\begin{programme}[V54: locate a nonperturbative coercive mechanism]
\label{prog:s6v53-V54}
\begin{enumerate}[leftmargin=3em]
\item Search the attached Yang--Mills and Grand Tensor sources for a
      reflection-compatible inequality that is \(L^2\)-coercive on the
      transverse infrared carrier: character expansion, centre-sector
      tunnelling, magnetic flux, or a Wilson-loop surface inequality.
\item Separate finite-volume positivity from a volume-uniform form bound.
\item Test whether any claimed polymer or shell estimate is absolute on
      infrared trial vectors or merely relative to the massless quadratic
      energy.
\item If no mechanism is present, prove the precise missing inequality as
      an explicit conjectural card rather than recycling ultraviolet
      contraction.
\item At V55, perform the compulsory companion audit before changing the
      route.
\end{enumerate}
\end{programme}

\begin{firewall}[Claim boundary after sixth-series V53]
\label{fire:s6v53-boundary}
V53 proves the exact transfer-edge/form-gap equivalence; proves the sharp
transfer-defect scaling; isolates carrier motion as a separate debit;
proves Mosco transport of a protected-sector gap; proves that a perturbation
only relatively bounded by the massless quadratic form cannot open a gap;
and states the necessary nonlinear infrared coercivity alternative and
RODFC handoff.

V53 does not prove the Yang--Mills uniform form inequality, the required
nonperturbative infrared deformation, Mosco convergence of the exact
Yang--Mills forms, projective regulator independence, the continuum OS
theory, or the full Yang--Mills mass gap.
\end{firewall}

\begin{theorem}[Sixth-series V53 result ledger]
\label{thm:s6v53-results}
The positive conclusions in \Cref{fire:s6v53-boundary} are exactly those of
\Cref{thm:s6v53-edge-form,
cor:s6v53-critical,
lem:s6v53-transfer-form,
cor:s6v53-order,
prop:s6v53-carrier,
thm:s6v53-Mosco-gap,
cor:s6v53-core,
thm:s6v53-relative-no-gap,
cor:s6v53-vanishing,
thm:s6v53-coercivity-alternative,
thm:s6v53-RODFC}.
\end{theorem}

\begin{proof}
Each item follows from the cited variational result with its carrier,
convergence, and uniformity assumptions unchanged.
\end{proof}

\paragraph{Parent-version record.}
V53 was formed from
\texttt{Renewal\_Yang\_Mills\_Mass\_Gap\_Research\_Notes\_V52.tex}.
The V52 parent had \(23\,465\) logical source lines,
\(822\,064\) bytes, and SHA--256
\[
 \texttt{02503f43327d3e1563a91a7644fdd2607c793d434effd0e67f01c3eca0c2ca45}.
\]
The next compulsory companion audit remains V55.

% =============================================================================
% END APPENDED SIXTH-SERIES V53 MATERIAL
% =============================================================================

% =============================================================================
% BEGIN APPENDED SIXTH-SERIES V54 MATERIAL
% =============================================================================

\section{V54: temporal conductance, complete two-boundary coercivity, and
the static-data obstruction}
\label{sec:s6v54}

\subsection{Targeted source audit}
\label{sec:s6v54-audit}

\begin{assessment}[What the attached coercivity statements actually prove]
\label{ass:s6v54-source-audit}
The targeted search of the attached Yang--Mills and Grand Tensor layers
finds four relevant but logically distinct statements.
\begin{enumerate}[label=\textup{(\roman*)},leftmargin=3.5em]
\item The reduced-fibre theorem labelled
      \texttt{thm:YMA2-global-fibre-coercivity} proves global coercivity of
      one pinned finite Wilson fibre and a positive vertical Hessian.  Its
      own scope statement says this is not coarse coercivity and does not
      prove a uniform Gibbs Poincar\'e inequality.
\item The boundary-sector theorem labelled
      \texttt{thm:V27-strict-one-link} proves strict contraction of every
      nontrivial irreducible under one finite-coupling one-link conditional
      law.  Its descendants concatenate those insertions along a
      \emph{spatial} rectangle side.  Their continuum contraction requires
      the unpopulated hypothesis \(r_N^{\rm hb}\to0\), and the source
      explicitly refuses to identify this spatial estimate with OS time.
\item The reduced-fibre result labelled
      \texttt{cor:YMA30-area-law-macrocell} is an exact calibration:
      \emph{if} a rectangular free energy already has an area term
      \(\sigma xy\), its macrocell mixed difference extracts
      \(\sigma\Delta_1\Delta_2\).  It neither proves \(\sigma>0\) nor
      controls a temporal transfer.
\item The Grand Tensor half-slab theorem proves finite-regulator
      Hilbert--Schmidt nuclearity, positivity, and an exact finite-volume
      spectral formula.  Its finite-volume firewall leaves uniform
      relative boundary correlation and continuum transport open.
\end{enumerate}
No item supplies an \(L^2\)-coercive, volume-uniform inequality for the
actual reflected time-boundary pair law.
\end{assessment}

\begin{proof}
Each assertion is the hypothesis/conclusion boundary stated in the named
source result.  In particular, the words \emph{if on the selected
rectangle} precede the area-law calibration; the boundary-sector firewalls
state that their spatial Peter--Weyl estimate is not a temporal mass gap;
and the reduced-fibre firewall following its pinning comparison says that
the surviving covariance and uniform Poincar\'e inequality are not
controlled.  Therefore none can populate the V53 RODFC coercivity row.
\end{proof}

\subsection{The exact two-boundary Dirichlet identity}
\label{sec:s6v54-pair}

Let \(\pi\) be a probability law, let \(P\) be a positive self-adjoint
Markov contraction on \(L^2(\pi)\), and let
\[
 {\mathsf Q}(\dd x,\dd y)=\pi(\dd x)P(x,\dd y)
\]
be its stationary reversible pair law.  Let \(\Pi\) project onto the
protected invariant factor, with \(P\Pi=\Pi\).

\begin{theorem}[Reflected pair-law identity]
\label{thm:s6v54-pair-identity}
For every \(f\in L^2(\pi)\),
\begin{equation}
 \boxed{
 \langle f,(I-P)f\rangle
 ={1\over2}\iint |f(x)-f(y)|^2\,
 {\mathsf Q}(\dd x,\dd y).}
 \label{eq:s6v54-pair-identity}
\end{equation}
Consequently, for a physical step \(\delta\), the relative form inequality
with constant \(\mu\),
\[
 {1\over\delta}\langle f,(I-P)f\rangle
 \ge\mu\|(I-\Pi)f\|_{L^2(\pi)}^2,
\]
is equivalent to the complete two-boundary coercivity inequality
\begin{equation}
 \boxed{
 \iint |f(x)-f(y)|^2\,\dd{\mathsf Q}
 \ge2\mu\delta\,
 \|(I-\Pi)f\|_{L^2(\pi)}^2
 \quad(f\in L^2(\pi)).}
 \label{eq:s6v54-complete-coercivity}
\end{equation}
\end{theorem}

\begin{proof}
Expand the square.  Both marginals of \({\mathsf Q}\) are \(\pi\), while
reversibility gives
\[
 \iint\overline{f(x)}f(y)\,\dd{\mathsf Q}
 =\langle f,Pf\rangle\in\mathbb R.
\]
Thus the right side of \eqref{eq:s6v54-pair-identity} is
\(\|f\|^2-\langle f,Pf\rangle\).  Multiply the rescaled form inequality by
\(2\delta\) to obtain \eqref{eq:s6v54-complete-coercivity}; the reverse
implication is the same calculation.
\end{proof}

\begin{corollary}[Near-critical maximal-correlation form]
\label{cor:s6v54-correlation}
On the relative carrier,
\eqref{eq:s6v54-complete-coercivity} holds if and only if
\[
 \|P|_{\rm rel}\|\le1-\mu\delta.
\]
For a Doob half-slab law with relative maximal correlation \(\rho\),
this is
\begin{equation}
 \rho^2\le1-\mu\delta.
 \label{eq:s6v54-rho-coercivity}
\end{equation}
Hence the desired Yang--Mills mass is exactly a complete-source
two-boundary variance loss of order \(\delta\), not an order-one
one-step contraction.
\end{corollary}

\begin{proof}
Apply \Cref{thm:s6v53-edge-form} and the half-slab identity
\(q=\rho^2\) from V51.
\end{proof}

\subsection{Temporal conductance at the correct scale}
\label{sec:s6v54-conductance}

For the moment suppose the protected factor consists only of constants.
Define the one-step conductance
\begin{equation}
 \Phi(P):=
 \inf_{\substack{A\ {\rm measurable}\\0<\pi(A)\le1/2}}
 {{\mathsf Q}(A,A^c)\over\pi(A)}.
 \label{eq:s6v54-conductance}
\end{equation}
For a protected factor, the same definition is applied fibrewise after
common-factor disintegration, with an essential uniform lower bound over
the admitted fibres.

\begin{theorem}[Cheeger lower bound for a positive reflected transfer]
\label{thm:s6v54-Cheeger}
Let \(P\) be a positive self-adjoint Markov contraction.  Then
\begin{equation}
 \boxed{
 1-\|P|_{\mathbf1^\perp}\|
 \ge{\Phi(P)^2\over2}.}
 \label{eq:s6v54-Cheeger}
\end{equation}
Consequently, if a family of physical steps obeys
\begin{equation}
 \Phi(P_j)\ge c\sqrt{\delta_j}
 \label{eq:s6v54-conductance-scale}
\end{equation}
uniformly in volume and regulator, then its rescaled transfer forms have
gap at least \(c^2/2\).
\end{theorem}

\begin{proof}
Let \(f\) be real, subtract a median, and split it into positive and
negative parts, whose supports each have measure at most \(1/2\).
For a nonnegative part \(g\), the layer-cake and coarea identities give
\[
 \iint|g(x)^2-g(y)^2|\,\dd{\mathsf Q}
 \ge\Phi(P)\int g^2\,\dd\pi.
\]
On the other hand, Cauchy--Schwarz, stationarity, and
\((a+b)^2\le2(a^2+b^2)\) give
\[
 \iint|g(x)^2-g(y)^2|\,\dd{\mathsf Q}
 \le
 \left(\iint|g(x)-g(y)|^2\,\dd{\mathsf Q}\right)^{1/2}
 \left(4\int g^2\,\dd\pi\right)^{1/2}.
\]
After squaring and summing the two median parts, the cross-sign pair
contribution is nonnegative, yielding
\[
 {1\over2}\iint|f(x)-f(y)|^2\,\dd{\mathsf Q}
 \ge{\Phi(P)^2\over2}\operatorname{Var}_\pi(f).
\]
By \Cref{thm:s6v54-pair-identity}, the best constant on the left is
\(1-\|P|_{\mathbf1^\perp}\|\), proving
\eqref{eq:s6v54-Cheeger}.  Insert
\eqref{eq:s6v54-conductance-scale} and divide by \(\delta_j\).
\end{proof}

\begin{remark}[Why the square-root scale is natural]
\label{rem:s6v54-sqrt}
A local diffusion crosses a codimension-one boundary over a distance of
order \(\sqrt{\delta}\) in time \(\delta\), so its one-step conductance is
typically of order \(\sqrt{\delta}\), while its spectral deficit is of
order \(\delta\).  Demanding a fixed positive conductance from a shrinking
time step would again impose an ultralocal infinite generator scale.
\end{remark}

\begin{assessment}[Conductance is sufficient, not yet populated]
\label{ass:s6v54-conductance}
The source's one-link heat-bath action and spatial side telescope resemble
a conductance calculation, but they are taken along a spatial Wilson
rectangle after conditioning on its complement.  To use
\Cref{thm:s6v54-Cheeger}, one must instead estimate cuts of the complete
initial-boundary law under the \emph{physical reflected time transfer},
after gauge and common-sector extraction.  No available incidence theorem
identifies these two cut problems.
\end{assessment}

\subsection{Static Yang--Mills data do not determine a temporal gap}
\label{sec:s6v54-static}

\begin{theorem}[Stationary-law non-identifiability of the mass]
\label{thm:s6v54-static-obstruction}
Let \(\pi\) be any probability law on a configuration space and let
\(\Pi_0 f=\int f\,\dd\pi\).  For every \(0<\varepsilon\le1\), define
\begin{equation}
 P_\varepsilon=(1-\varepsilon)I+\varepsilon\Pi_0.
 \label{eq:s6v54-refresh}
\end{equation}
Then \(P_\varepsilon\) is a positive self-adjoint Markov contraction with
stationary law \(\pi\), and
\begin{equation}
 \|P_\varepsilon|_{\mathbf1^\perp}\|=1-\varepsilon,
 \qquad
 \inf_{\mathbf1^\perp}
 {{\cal E}_{P_\varepsilon,\delta}(f)\over\|f\|^2}
 ={\varepsilon\over\delta}.
 \label{eq:s6v54-refresh-gap}
\end{equation}
Thus all equal-time expectations, all spatial Wilson-loop area data, and
all spatial conditional expectations determined solely by \(\pi\) can be
held fixed while the physical mass rate is made to tend to zero, to a
finite constant, or to infinity by choosing respectively
\[
 \varepsilon_\delta=o(\delta),\qquad
 \varepsilon_\delta\sim\mu\delta,\qquad
 \varepsilon_\delta/\delta\to\infty.
\]
\end{theorem}

\begin{proof}
Both \(I\) and \(\Pi_0\) are positive self-adjoint Markov contractions and
fix \(\pi\); their convex combination has the same properties.  On
constants \(P_\varepsilon=I\), while on the mean-zero subspace
\(P_\varepsilon=(1-\varepsilon)I\).  This proves
\eqref{eq:s6v54-refresh-gap}.  The law \(\pi\) is identical for every
\(\varepsilon\), so every statistic and spatial conditional expectation
computed from that one law is identical as well.  The three rate regimes
follow by division by \(\delta\).
\end{proof}

\begin{corollary}[An area law alone is not a mass-gap proof]
\label{cor:s6v54-area-not-gap}
Even an exact positive string tension for every loop in a fixed-time
spatial law does not, without a physical reflected pair law or an
equivalent time-incidence theorem, determine a Hamiltonian mass gap.
\end{corollary}

\begin{proof}
Choose \(\pi\) to be any law having the asserted spatial loop statistics
and apply \Cref{thm:s6v54-static-obstruction}.  The static data remain fixed
while the temporal rate varies arbitrarily.
\end{proof}

\begin{remark}[Scope of the logical counterexample]
\label{rem:s6v54-static-scope}
\Cref{cor:s6v54-area-not-gap} does not say that confinement and a mass gap
are unrelated in pure Yang--Mills theory.  It says that the static
observable data alone omit the time-transfer incidence needed to prove the
spectral statement.  A Yang--Mills-specific theorem relating the actual
area-law measure to its reflected slab kernel could supply that missing
information, but it must be proved.
\end{remark}

\subsection{A complete temporal cut card}
\label{sec:s6v54-card}

\begin{definition}[Temporal half-slab conductance card, THSCC]
\label{def:s6v54-THSCC}
For every regulator and volume, THSCC records:
\begin{enumerate}[label=\textup{(\roman*)},leftmargin=3.5em]
\item the exact response-matched half-slab amplitude and its Doob pair law
      \({\mathsf Q}_j\);
\item the physical full-slab duration \(\delta_j\);
\item the maximal common-factor disintegration and protected projection
      \(\Pi_j\);
\item either the complete coercivity inequality
      \eqref{eq:s6v54-complete-coercivity} with one \(\mu>0\), or a
      fibre-uniform conductance lower bound
      \(\Phi_j\ge\sqrt{2\mu\delta_j}\);
\item proof that the bound holds for the complete centre-neutral
      gauge-invariant relative carrier, not only a finite Wilson-loop bank;
\item stability of the cut/coercivity constant under exact shell
      integration and adjacent regulator transport;
\item the inverse-margin projectivity ledger and common-refinement bridge;
      and
\item Mosco recovery and liminf on the limiting cyclic OS source core.
\end{enumerate}
\end{definition}

\begin{theorem}[Conditional THSCC completion]
\label{thm:s6v54-THSCC}
If THSCC is populated and the limiting OS theory is nontrivial, then its
Hamiltonian has relative vacuum gap at least \(\mu\).
\end{theorem}

\begin{proof}
Row~\textup{(iv)} gives the uniform regulator form gap by
\Cref{thm:s6v54-pair-identity}, directly in the coercivity branch, or by
\Cref{thm:s6v54-Cheeger}, in the conductance branch.
Rows~\textup{(v)--(viii)} give the complete carrier, regulator-independent
Mosco limit, and cyclic OS identification.  Apply
\Cref{thm:s6v53-Mosco-gap,thm:s6v53-RODFC}.
\end{proof}

\begin{assessment}[Exact frontier after V54]
\label{ass:s6v54-frontier}
The bottleneck is now an explicit cut or variance inequality on an existing
object: the response-matched reflected half-slab pair law.  This target is
strictly stronger than the attached finite-fibre coercivity, spatial
Peter--Weyl attenuation, and conditional area-law calibration, and it has
the correct \(\sqrt{\delta}\) conductance or \(\delta\) variance-loss
scaling.  The next compulsory SM/NS audit should be used only to search for
a transferable method proving this temporal inequality or its projective
stability.
\end{assessment}

\begin{programme}[V55: compulsory companion audit at the temporal-cut
frontier]
\label{prog:s6v54-V55}
\begin{enumerate}[leftmargin=3em]
\item Hash and inspect the current SM and NS companion files.
\item Search them for complete-source Poincar\'e, Cheeger, reflected
      two-boundary variance, Mosco, or generator-form stability arguments.
\item Import only abstract proof mechanisms whose Yang--Mills hypotheses
      are separately restated and verified.
\item Re-test the current exact Yang--Mills sources for any newly exposed
      time-incidence map.
\item If no temporal inequality is populated, proceed upstream to a
      canonical-path or block-dynamics comparison for the exact half-slab
      pair law, with gauge/common-sector fibres explicit.
\end{enumerate}
\end{programme}

\begin{firewall}[Claim boundary after sixth-series V54]
\label{fire:s6v54-boundary}
V54 audits the attached coercivity claims; proves the exact reflected
pair-law Dirichlet identity and complete two-boundary coercivity
equivalence; proves the correctly scaled Cheeger sufficient condition;
proves that static area-law and spatial conditional data cannot determine a
temporal gap; and formulates THSCC with a conditional mass-gap handoff.

V54 does not prove temporal conductance or complete two-boundary coercivity
for Yang--Mills, relate the spatial character telescope to OS time,
establish a positive string tension, close projectivity or Mosco
convergence, construct the continuum OS theory, or prove the full
Yang--Mills mass gap.
\end{firewall}

\begin{theorem}[Sixth-series V54 result ledger]
\label{thm:s6v54-results}
The positive conclusions in \Cref{fire:s6v54-boundary} are exactly those of
\Cref{ass:s6v54-source-audit,
thm:s6v54-pair-identity,
cor:s6v54-correlation,
thm:s6v54-Cheeger,
thm:s6v54-static-obstruction,
cor:s6v54-area-not-gap,
thm:s6v54-THSCC}.
\end{theorem}

\begin{proof}
Each statement is the conclusion of the cited theorem or audited source
boundary, with all time-incidence, carrier, and uniformity qualifications
retained.
\end{proof}

\paragraph{Parent-version record.}
V54 was formed from
\texttt{Renewal\_Yang\_Mills\_Mass\_Gap\_Research\_Notes\_V53.tex}.
The V53 parent had \(23\,921\) logical source lines,
\(838\,387\) bytes, and SHA--256
\[
 \texttt{266aeb29d72031b3b96ca8d4e9e22127140c2f3e3b4f5fe62bc8438a152249d4}.
\]
The next version performs the compulsory companion audit.

% =============================================================================
% END APPENDED SIXTH-SERIES V54 MATERIAL
% =============================================================================

% =============================================================================
% BEGIN APPENDED SIXTH-SERIES V55 MATERIAL
% =============================================================================

\section{V55 compulsory companion audit: clock factors, slow-sector
projectors, and the Yang--Mills import boundary}
\label{sec:s6v55}

\subsection{Audited files and exact provenance}
\label{sec:s6v55-files}

\begin{definition}[V55 companion snapshot]
\label{def:s6v55-snapshot}
The compulsory audit used the following current active files:
\begin{center}
\begin{tabular}{p{0.43\linewidth}rr}
\toprule
file & logical lines & bytes\\
\midrule
\texttt{Renewal\_SM\_Einstein\_Continuum\_Research\_Notes\_V3.tex}
& \(1\,191\) & \(46\,246\)\\
\texttt{Renewal\_NS\_Stability\_Research\_Notes\_V26.tex}
& \(5\,319\) & \(278\,283\)\\
\bottomrule
\end{tabular}
\end{center}
Their SHA--256 digests are respectively
\[
 \texttt{aa524965426043b76e7fba2e170928e1e74642493cd1a8386898abce2fa6fd61}
\]
and
\[
 \texttt{82c887f8c2c69e4f28fad7c3dc8de5b9099cb5a4bf431eba1fff3e91935978ad}.
\]
\end{definition}

\begin{assessment}[SM V3 audit]
\label{ass:s6v55-SM}
The restarted SM series now contains three abstract mechanisms relevant to
the V54 frontier:
\begin{enumerate}[label=\textup{(\roman*)},leftmargin=3.5em]
\item exact Duhamel formulas force every short-time channel derivative to
      carry at least one factor of the physical step;
\item a strict generator logarithmic norm yields
      \(e^{-m\tau}=1-m\tau+O(\tau^2)\), while any stationary Ward mode
      detected by the norm forbids such strict dissipativity;
\item quotienting a slow sector is valid only for observables in its
      annihilator and with a proved local dual norm; approximate reduction
      creates an explicit Duhamel leakage, and a complementary gap closing
      in volume forces the block time inferred from that gap to diverge.
\end{enumerate}
These are reusable functional-analytic warnings and estimates.  SM V3 does
not construct a Yang--Mills half-slab, common-sector projector, temporal
conductance, or volume-uniform complementary generator gap.
\end{assessment}

\begin{assessment}[NS V26 audit]
\label{ass:s6v55-NS}
NS V26 is byte-for-byte unchanged from the V45 and V50 audits.  Its
Dirichlet forms concern the Navier--Stokes blow-up programme and are not
forms of a positive reflected Markov transfer.  It supplies no
half-slab pair law, clock-scaled Yang--Mills conductance, protected-sector
projection, or Mosco comparison for the present problem.
\end{assessment}

\begin{theorem}[V55 audit decision]
\label{thm:s6v55-audit-decision}
The Yang--Mills programme imports from SM V3 only:
\begin{enumerate}[label=\textup{(\roman*)},leftmargin=3.5em]
\item the Duhamel origin of a physical-step factor;
\item the rule that stationary/common modes must be removed by a proved
      invariant quotient with observable duality; and
\item the accounting of approximate slow/fast leakage as a separate
      error.
\end{enumerate}
No companion contraction constant, generator, field estimate, or
continuum conclusion is imported.
\end{theorem}

\begin{proof}
The imported items are general operator identities whose Yang--Mills
objects are restated below.  By
\Cref{ass:s6v55-SM,ass:s6v55-NS}, neither companion supplies those physical
objects or their required uniform estimates.
\end{proof}

\subsection{The clock factor from the actual generator}
\label{sec:s6v55-Duhamel}

\begin{lemma}[Contractive Duhamel comparison]
\label{lem:s6v55-Duhamel}
Let \(H,K\) be nonnegative self-adjoint operators on one Hilbert space and
assume \(V:=H-K\) extends to a bounded operator.  Then, for every
\(\delta\ge0\),
\begin{equation}
 e^{-\delta H}-e^{-\delta K}
 =-\int_0^\delta e^{-(\delta-s)H}V e^{-sK}\,\dd s
 \label{eq:s6v55-Duhamel}
\end{equation}
in operator norm, and
\begin{equation}
 \boxed{
 \|e^{-\delta H}-e^{-\delta K}\|
 \le\delta\|V\|.}
 \label{eq:s6v55-Duhamel-bound}
\end{equation}
\end{lemma}

\begin{proof}
Differentiate
\(s\mapsto e^{-(\delta-s)H}e^{-sK}\) on a common operator core and
integrate.  Boundedness of \(V\) extends the identity to the whole Hilbert
space.  Both semigroups are contractions, so the integrand norm is at most
\(\|V\|\); integration gives \eqref{eq:s6v55-Duhamel-bound}.
\end{proof}

\begin{corollary}[Generator-sized corrections survive]
\label{cor:s6v55-generator-sized}
If the response-matched Yang--Mills Hamiltonian can be written
\(H_j=H_{0,j}+V_j\) with
\(\sup_j\|V_j\|<\infty\), then its one-step transfer differs from the
reference transfer by \(O(\delta_j)\).  This is exactly the scale that
survives in the V53 rescaled form.  It cannot be dropped under an
ultraviolet small-remainder label.
\end{corollary}

\begin{proof}
Apply \Cref{lem:s6v55-Duhamel} and then
\Cref{lem:s6v53-transfer-form}.
\end{proof}

\begin{remark}[Unbounded interactions require form Duhamel theory]
\label{rem:s6v55-unbounded}
The bounded-perturbation hypothesis is not automatic for continuum
Yang--Mills.  An unbounded electric or magnetic contribution must instead
be handled by relative form bounds and strong-resolvent convergence.
The explicit factor \(\delta_j\) may not be asserted merely by writing a
formal exponential.
\end{remark}

\subsection{Near-critical stability of the protected projection}
\label{sec:s6v55-projector}

\begin{theorem}[Spectral projector perturbation at a closing margin]
\label{thm:s6v55-projector}
Let \(P,Q\) be positive self-adjoint contractions on a Hilbert space.
Suppose \(1\) is an isolated eigenvalue of \(P\), with spectral projection
\(\Pi_P\), and
\begin{equation}
 \operatorname{spec}(P)\setminus\{1\}\subset[0,1-g],
 \qquad g>0.
 \label{eq:s6v55-P-gap}
\end{equation}
If
\(\varepsilon:=\|P-Q\|<g/4\), let \(\Pi_Q^{\rm top}\) be the spectral
projection of \(Q\) for the interval \((1-g/2,1]\).  Then
\begin{equation}
 \boxed{
 \|\Pi_Q^{\rm top}-\Pi_P\|
 \le {4\varepsilon\over g}.}
 \label{eq:s6v55-projector-bound}
\end{equation}
In particular, if \(g_j=\mu\delta_j+o(\delta_j)\), convergence of the
protected projectors by this route requires
\begin{equation}
 {\|P_j-Q_j\|\over\delta_j}\longrightarrow0.
 \label{eq:s6v55-projector-scale}
\end{equation}
\end{theorem}

\begin{proof}
Let \(\Gamma\) be the circle of radius \(g/2\) around \(1\), oriented
positively.  It separates \(\{1\}\) from the rest of
\(\operatorname{spec}P\).  Since \(\varepsilon<g/4\), spectral inclusion
for self-adjoint perturbations shows that the same contour separates the
top cluster of \(Q\).  The Riesz formulas and the resolvent identity give
\[
 \Pi_Q^{\rm top}-\Pi_P
 ={1\over2\pi i}\int_\Gamma
 (z-Q)^{-1}(Q-P)(z-P)^{-1}\,\dd z.
\]
On \(\Gamma\),
\(\|(z-P)^{-1}\|\le2/g\), while
\(\|(z-Q)^{-1}\|\le1/(g/2-\varepsilon)\le4/g\).
The contour length is \(\pi g\), hence the displayed integral is at most
\(4\varepsilon/g\).  Substituting
\(g_j=\mu\delta_j+o(\delta_j)\) proves
\eqref{eq:s6v55-projector-scale}.
\end{proof}

\begin{assessment}[Top cluster versus declared common factor]
\label{ass:s6v55-top-cluster}
The projection \(\Pi_Q^{\rm top}\) in
\Cref{thm:s6v55-projector} is determined spectrally.  It equals the desired
vacuum/common-sector projection only after a physical sector theorem says
which top states are protected.  Projector perturbation controls the angle
between already identified clusters; it does not decide whether a centre
or topological label should be shorted.
\end{assessment}

\begin{proposition}[Cost of using an approximate protected carrier]
\label{prop:s6v55-approx-carrier}
Let \(P\Pi=\Pi\), let
\(\|P|_{(I-\Pi){\cal H}}\|\le q<1\), and let
\(\widetilde\Pi\) be another orthogonal projection with
\(\|\Pi-\widetilde\Pi\|\le\kappa<1\).  Then
\begin{equation}
 \sup_{\substack{f\perp\operatorname{Ran}\widetilde\Pi\\\|f\|=1}}
 \langle f,Pf\rangle
 \le q+(1-q)\kappa^2,
 \label{eq:s6v55-approx-edge}
\end{equation}
so the form deficit on the approximate relative carrier is at least
\begin{equation}
 (1-q)(1-\kappa^2).
 \label{eq:s6v55-approx-deficit}
\end{equation}
\end{proposition}

\begin{proof}
For \(f\perp\operatorname{Ran}\widetilde\Pi\),
\[
 \|\Pi f\|
 =\|(\Pi-\widetilde\Pi)f\|\le\kappa\|f\|.
\]
The reducing decomposition and positivity give, for unit \(f\),
\[
 \langle f,Pf\rangle
 \le\|\Pi f\|^2+q\|(I-\Pi)f\|^2
 =q+(1-q)\|\Pi f\|^2.
\]
This proves both displayed bounds.
\end{proof}

\begin{corollary}[Angle control and incidence control are different]
\label{cor:s6v55-angle-incidence}
A bounded projector angle \(\kappa<1\) preserves the order-\(\delta\)
spectral deficit once the true gap is known.  To prove that transported
local observables belong to the approximate quotient, one still needs the
annihilator and local-dual estimates emphasized by SM V3.
\end{corollary}

\begin{proof}
The first statement follows from
\eqref{eq:s6v55-approx-deficit}; the second is independent because the
projector norm contains no support or observable-functional information.
\end{proof}

\subsection{Leakage in one shrinking physical step}
\label{sec:s6v55-leakage}

\begin{lemma}[Off-diagonal generator leakage]
\label{lem:s6v55-leakage}
Let \(H=H_0+V\ge0\), where \(H_0\ge0\) commutes with an orthogonal
projection \(\Pi\), and \(V\) is bounded.  Put \(R=I-\Pi\).  Then
\begin{equation}
 \boxed{
 \|\Pi e^{-\delta H}R\|
 \le\delta\|\Pi VR\|+\delta^2\|V\|^2 e^{\delta\|V\|}/2.}
 \label{eq:s6v55-leakage}
\end{equation}
Thus an off-diagonal generator block of order one produces an
order-\(\delta\) one-step leakage source.
\end{lemma}

\begin{proof}
Use the Dyson expansion around \(H_0\).  Since
\(\Pi e^{-\delta H_0}R=0\), the one-insertion term is
\[
 -\int_0^\delta
 \Pi e^{-(\delta-s)H_0}V e^{-sH_0}R\,\dd s,
\]
whose norm is at most \(\delta\|\Pi VR\|\).
The sum of all terms with at least two insertions is bounded by
\[
 \sum_{n\ge2}{\delta^n\|V\|^n\over n!}
 \le{\delta^2\|V\|^2\over2}e^{\delta\|V\|}.
\]
\end{proof}

\begin{assessment}[Leakage cannot be hidden in conductance]
\label{ass:s6v55-leakage}
If the protected sector is only approximately invariant, the off-diagonal
term in \eqref{eq:s6v55-leakage} is an inhomogeneous source coupling
protected and relative carriers.  A conductance or Poincar\'e estimate for
the compressed relative block does not control this source.  It must vanish
after division by the physical step, be absorbed into an enlarged exact
protected cluster, or be retained in the limiting generator.
\end{assessment}

\subsection{The audited Yang--Mills packet}
\label{sec:s6v55-packet}

\begin{definition}[Audited temporal projector packet, ATPP]
\label{def:s6v55-ATPP}
ATPP strengthens THSCC by requiring:
\begin{enumerate}[label=\textup{(\roman*)},leftmargin=3.5em]
\item an exactly identified physical protected projection \(\Pi_j\) for
      the response-matched transfer \(P_j\);
\item if a comparison projection \(\widetilde\Pi_j\) is used, a bound
      \(\|\Pi_j-\widetilde\Pi_j\|\le\kappa<1\) and a proved
      observable-annihilator/local-dual row;
\item if projectors are transported from adjacent kernels, a normalized
      transfer defect \(o(\delta_j)\) strong enough for
      \eqref{eq:s6v55-projector-scale};
\item an off-diagonal leakage estimate whose rescaled limit is either zero
      or an explicit term of the limiting generator;
\item the complete temporal coercivity or conductance row on the resulting
      exact relative carrier; and
\item the THSCC projectivity, Mosco, noncollapse, and cyclicity rows.
\end{enumerate}
\end{definition}

\begin{theorem}[Conditional ATPP handoff]
\label{thm:s6v55-ATPP}
If ATPP is populated with temporal form constant \(\mu>0\), the
regulator-independent reconstructed Yang--Mills Hamiltonian has relative
vacuum gap at least \(\mu\).
\end{theorem}

\begin{proof}
Rows~\textup{(i)--(iv)} identify one exact transported protected/relative
decomposition and prevent uncharged leakage.  If an approximate carrier is
used as an intermediate comparison,
\Cref{prop:s6v55-approx-carrier} quantifies its diagnostic loss, but
row~\textup{(v)} still requires the coercivity conclusion on the resulting
exact physical relative carrier.  That row supplies the regulator form
gap, and row~\textup{(vi)} permits the V53 Mosco and OS handoff.
\end{proof}

\begin{assessment}[Direction after the audit]
\label{ass:s6v55-frontier}
The audit does not populate temporal conductance, but it closes a hidden
logical gap in the proposed cut route: the protected/common-sector
projection must be transported at the same inverse-margin scale as the
kernel, and approximate slow/fast leakage must be charged at generator
order.  The next upstream attack is a decomposition theorem for the actual
half-slab pair law: local block coercivity plus a coarse mode chain,
with every constant scaled in physical time and every common factor
removed before tensorization.
\end{assessment}

\begin{programme}[V56: block decomposition of temporal coercivity]
\label{prog:s6v55-V56}
\begin{enumerate}[leftmargin=3em]
\item Prove an exact variance decomposition through a time-slab block
      sigma-field.
\item Combine conditional block Poincar\'e constants with a coarse
      boundary-chain constant, keeping the physical \(\delta\) scale.
\item Identify the overlap or comparison debit when conditional blocks are
      not independent.
\item Show why one-link finite-card gaps can deteriorate with block count
      unless a parallel/tensorized physical update is proved.
\item Formulate the concrete Yang--Mills block inputs without importing an
      SM generator or an NS form.
\end{enumerate}
\end{programme}

\begin{firewall}[Claim boundary after sixth-series V55]
\label{fire:s6v55-boundary}
V55 performs the compulsory SM/NS audit; imports only abstract Duhamel,
quotient, and leakage logic; proves the bounded-generator clock factor;
proves near-critical spectral-projector stability and its inverse-margin
cost; quantifies the loss from an approximate protected carrier; and proves
an off-diagonal generator leakage bound.

V55 does not identify or control the exact Yang--Mills protected projector,
prove temporal conductance or block coercivity, populate the required
projector/kernel defect scales, close Mosco/projective transport, construct
the continuum OS theory, or prove the full Yang--Mills mass gap.
\end{firewall}

\begin{theorem}[Sixth-series V55 result ledger]
\label{thm:s6v55-results}
The positive conclusions in \Cref{fire:s6v55-boundary} are exactly those of
\Cref{def:s6v55-snapshot,
thm:s6v55-audit-decision,
lem:s6v55-Duhamel,
cor:s6v55-generator-sized,
thm:s6v55-projector,
prop:s6v55-approx-carrier,
cor:s6v55-angle-incidence,
lem:s6v55-leakage,
thm:s6v55-ATPP}.
\end{theorem}

\begin{proof}
Each ledger item is the cited result with its boundedness, spectral
separation, carrier, and physical-incidence hypotheses unchanged.
\end{proof}

\paragraph{Parent-version record.}
V55 was formed from
\texttt{Renewal\_Yang\_Mills\_Mass\_Gap\_Research\_Notes\_V54.tex}.
The V54 parent had \(24\,312\) logical source lines,
\(854\,038\) bytes, and SHA--256
\[
 \texttt{42d821d2211f1e17857c7ab3e544f3f62b9e11be2ea42713ff696a1c915c3ff3}.
\]
The next compulsory companion audit is V60.

% =============================================================================
% END APPENDED SIXTH-SERIES V55 MATERIAL
% =============================================================================

% =============================================================================
% BEGIN APPENDED SIXTH-SERIES V56 MATERIAL
% =============================================================================

\section{V56: two-level temporal block coercivity and the update-schedule
debit}
\label{sec:s6v56}

\subsection{Exact coarse/fine Hilbert decomposition}
\label{sec:s6v56-decomposition}

Let \(({\mathsf X},\pi)\) be a probability space and let
\(r:{\mathsf X}\to{\mathsf Z}\) be a measurable coarse boundary map.
Write \(\bar\pi=r_*\pi\).  Define
\[
 J:L^2(\bar\pi)\to L^2(\pi),\qquad
 (Jm)(x)=m(r(x)),
\]
and let \(K=J^*\), so \(Kf=\mathbb E_\pi[f\mid r]\).  Then \(J\) is an
isometry, \(KJ=I\), and \(JK\) is the orthogonal projection onto the
coarse-measurable subspace.

\begin{theorem}[Exact variance decomposition]
\label{thm:s6v56-variance}
For every \(f\in L^2(\pi)\), put \(m=Kf\) and \(g=f-Jm\).  Then
\begin{align}
 \|f-\pi f\|_{L^2(\pi)}^2
 &=\|g\|_{L^2(\pi)}^2
   +\|m-\bar\pi m\|_{L^2(\bar\pi)}^2,
 \label{eq:s6v56-variance}\\
 \|g\|_{L^2(\pi)}^2
 &=\mathbb E_\pi\operatorname{Var}(f\mid r).
 \label{eq:s6v56-fine-variance}
\end{align}
\end{theorem}

\begin{proof}
The vector \(g=(I-JK)f\) is orthogonal to the range of \(J\), while
\[
 f-\pi f=g+J(m-\bar\pi m).
\]
Pythagoras and isometry of \(J\) prove
\eqref{eq:s6v56-variance}.  The second identity is the standard
conditional-variance formula, obtained by integrating
\(|f-\mathbb E[f\mid r]|^2\).
\end{proof}

\subsection{The canonical two-level comparison transfer}
\label{sec:s6v56-comparison}

Let \(Q\) be a positive self-adjoint Markov contraction on
\(L^2(\bar\pi)\), and define
\begin{equation}
 R:=JQK
 \quad\hbox{on }L^2(\pi).
 \label{eq:s6v56-R}
\end{equation}
Probabilistically, \(R\) reads the coarse variable, advances it with \(Q\),
and samples the fine variable from the target conditional law at the new
coarse value.

\begin{theorem}[Exact two-level form identity]
\label{thm:s6v56-two-level}
The operator \(R\) is a positive self-adjoint Markov contraction.  For
\(f=Jm+g\) as in \Cref{thm:s6v56-variance},
\begin{equation}
 \boxed{
 \langle f,(I-R)f\rangle
 =\|g\|^2+\langle m,(I-Q)m\rangle_{L^2(\bar\pi)}.}
 \label{eq:s6v56-two-level-form}
\end{equation}
If constants are the only protected modes and
\[
 q_Q=\|Q|_{\mathbf1^\perp}\|,
\]
then
\begin{equation}
 \boxed{
 1-\|R|_{\mathbf1^\perp}\|=1-q_Q.}
 \label{eq:s6v56-R-gap}
\end{equation}
\end{theorem}

\begin{proof}
The adjoint identity \(K=J^*\) and self-adjointness of \(Q\) make \(R\)
self-adjoint; positivity follows from
\(\langle f,Rf\rangle=\langle Kf,QKf\rangle\ge0\).
Moreover \(R1=JQ1=1\), and positivity preservation follows from that of
conditional expectation, \(Q\), and pullback, so \(R\) is Markov.
Since \(Kg=0\),
\[
 R(Jm+g)=JQm.
\]
Orthogonality of \(g\) to the range of \(J\) gives
\[
 \langle Jm+g,Jm+g-JQm\rangle
 =\|g\|^2+\langle m,(I-Q)m\rangle,
\]
proving \eqref{eq:s6v56-two-level-form}.  On \(\ker K\), \(R=0\);
on \(\operatorname{Ran}J\), it is unitarily equivalent to \(Q\).
Thus the largest nonconstant eigenvalue is \(q_Q\).
\end{proof}

\begin{remark}[Protected coarse factors]
\label{rem:s6v56-protected}
If \(Q\) has a protected projection \(\bar\Pi\), then
\(\Pi=J\bar\Pi K\) is protected for \(R\), while the fine kernel
\(\ker K\) remains an eigenvalue-zero subspace.  All statements below hold
on the corresponding relative carriers after this replacement.
\end{remark}

\subsection{A block criterion for the physical transfer}
\label{sec:s6v56-criterion}

Let \(P\) be the actual positive self-adjoint physical transfer on
\(L^2(\pi)\).  Its dimensionless Dirichlet form is
\[
 {\cal D}_P(f)=\langle f,(I-P)f\rangle.
\]

\begin{theorem}[Two-level temporal coercivity compiler]
\label{thm:s6v56-compiler}
Suppose there are constants \(a,b\ge0\) such that, for every relative
\(f=Jm+g\),
\begin{equation}
 {\cal D}_P(f)
 \ge a\|g\|^2+b\langle m,(I-Q)m\rangle.
 \label{eq:s6v56-block-domination}
\end{equation}
If the relative coarse gap is
\[
 \langle m,(I-Q)m\rangle
 \ge\gamma_Q\|m\|^2,
\]
then
\begin{equation}
 \boxed{
 {\cal D}_P(f)\ge
 \min\{a,b\gamma_Q\}\|f\|^2.}
 \label{eq:s6v56-compiled-gap}
\end{equation}
For a physical step \(\delta\), a mass lower bound \(\mu\) follows whenever
\begin{equation}
 a\ge\mu\delta,\qquad
 b\gamma_Q\ge\mu\delta.
 \label{eq:s6v56-physical-rows}
\end{equation}
\end{theorem}

\begin{proof}
Apply the coarse gap in
\eqref{eq:s6v56-block-domination}, then use orthogonality:
\[
 {\cal D}_P(f)\ge a\|g\|^2+b\gamma_Q\|Jm\|^2
 \ge\min\{a,b\gamma_Q\}(\|g\|^2+\|Jm\|^2).
\]
The last bracket is \(\|f\|^2\) on the relative carrier.
Divide \eqref{eq:s6v56-compiled-gap} by \(\delta\) and apply V53.
\end{proof}

\begin{corollary}[Direct comparison with the canonical block transfer]
\label{cor:s6v56-direct-comparison}
If
\begin{equation}
 I-P\succeq c(I-R)
 \label{eq:s6v56-operator-domination}
\end{equation}
as quadratic forms, then
\[
 1-\|P|_{\rm rel}\|\ge c(1-q_Q).
\]
Equivalently, \Cref{thm:s6v56-compiler} holds with \(a=b=c\).
\end{corollary}

\begin{proof}
Insert the exact decomposition
\eqref{eq:s6v56-two-level-form} into
\eqref{eq:s6v56-operator-domination} and use
\eqref{eq:s6v56-R-gap}.
\end{proof}

\begin{lemma}[Additive comparison debit]
\label{lem:s6v56-additive-debit}
If \(P,R\) act on the same exact relative carrier, then
\begin{equation}
 1-\|P|_{\rm rel}\|
 \ge 1-\|R|_{\rm rel}\|-\|P-R\|.
 \label{eq:s6v56-additive-debit}
\end{equation}
Thus a comparison transfer with deficit \(\mu\delta+o(\delta)\) transports
that coefficient only when
\(\|P-R\|=o(\delta)\), or when a sharper signed form comparison retains a
strict positive part of the coefficient.
\end{lemma}

\begin{proof}
The triangle inequality gives
\(\|P|_{\rm rel}\|\le\|R|_{\rm rel}\|+\|P-R\|\).
Rearrange.
\end{proof}

\subsection{Parallel versus random-scan updates}
\label{sec:s6v56-schedule}

\begin{theorem}[Tensorized parallel update]
\label{thm:s6v56-parallel}
For \(N\) probability spaces, let \(P_i\) be positive self-adjoint Markov
contractions with mean-zero norms \(q_i<1\).  The parallel product transfer
\[
 P_{\rm par}=\bigotimes_{i=1}^N P_i
\]
has
\begin{equation}
 \boxed{
 \|P_{\rm par}|_{\mathbf1^\perp}\|=\max_iq_i,
 \qquad
 \operatorname{gap}(I-P_{\rm par})=\min_i(1-q_i).}
 \label{eq:s6v56-parallel-gap}
\end{equation}
No factor \(N^{-1}\) is lost.
\end{theorem}

\begin{proof}
Decompose each \(L^2\) space into constants plus its mean-zero part.
The product Hilbert space is the orthogonal sum indexed by nonempty sets of
mean-zero factors.  On the component indexed by \(S\), the operator norm is
\(\prod_{i\in S}q_i\), whose maximum over nonempty \(S\) is
\(\max_iq_i\).
\end{proof}

\begin{theorem}[Random-scan volume debit]
\label{thm:s6v56-random-scan}
Let \(E_i\) be the operator that applies \(P_i\) to coordinate \(i\) and
the identity to all other coordinates, and put
\[
 P_{\rm rs}={1\over N}\sum_{i=1}^NE_i.
\]
Then
\begin{equation}
 \boxed{
 \operatorname{gap}(I-P_{\rm rs})
 ={1\over N}\min_i(1-q_i).}
 \label{eq:s6v56-random-gap}
\end{equation}
Thus a proof based on one uniformly chosen one-link refresh per physical
step loses a factor \(N^{-1}\).
\end{theorem}

\begin{proof}
The same tensor decomposition diagonalizes the commuting \(E_i\).
On a component with mean-zero index set \(S\), its largest eigenvalue is
at most
\[
 {1\over N}\left(\sum_{i\notin S}1+\sum_{i\in S}q_i\right)
 =1-{1\over N}\sum_{i\in S}(1-q_i).
\]
The largest nonconstant value is obtained by a singleton \(S=\{i\}\)
whose deficit is minimal, proving
\eqref{eq:s6v56-random-gap}.
\end{proof}

\begin{corollary}[Sweep clocks must be declared]
\label{cor:s6v56-sweep}
If \(N\) single-coordinate updates are called one sweep, then converting
\eqref{eq:s6v56-random-gap} into a physical rate requires declaring whether
one update has duration \(\delta/N\) or one sweep has duration \(\delta\).
Changing this convention after estimating the gap changes the generator
by a factor \(N\).
\end{corollary}

\begin{proof}
The rate is the dimensionless deficit divided by the physical duration of
the operator to which the deficit belongs.  Substitute the two possible
durations.
\end{proof}

\begin{assessment}[Why the attached one-link theorem does not yet load
the fine row]
\label{ass:s6v56-one-link}
The boundary-sector source proves strict finite-card contraction for each
one-link conditional representation and concatenates these matrices along
a spatial loop side.  It does not prove that the physical half-slab
transfer is the parallel product in \Cref{thm:s6v56-parallel}, a
random-scan chain with a calibrated sweep clock, or the canonical
coarse/fine transfer \(R=JQK\).  Therefore its one-link constants cannot be
inserted as \(a\) in \eqref{eq:s6v56-physical-rows} without a new
time-incidence and comparison theorem.
\end{assessment}

\subsection{Overlap and coarse-mode firewall}
\label{sec:s6v56-overlap}

\begin{proposition}[Local block gaps leave the coarse mode untouched]
\label{prop:s6v56-coarse-firewall}
For the canonical transfer \(R=JQK\), arbitrarily strong fine resampling
does not improve the coarse deficit \(1-q_Q\).  In particular, if
\(q_Q\to1\) in volume, then the full gap of \(R\) closes even though every
conditional fine variable is refreshed exactly.
\end{proposition}

\begin{proof}
By \eqref{eq:s6v56-R-gap}, the nonconstant edge of \(R\) is exactly
\(q_Q\).  The eigenfunctions \(Jm\) depend only on the coarse variable and
are unaffected by the exact conditional fine resampling.
\end{proof}

\begin{assessment}[The remaining Yang--Mills content]
\label{ass:s6v56-YM-content}
The two-level theorem separates the missing field estimate into two honest
parts:
\begin{enumerate}[label=\textup{(\roman*)},leftmargin=3.5em]
\item a physical half-slab form domination that pays at least
      \(a\asymp\delta\) for every conditional fine fluctuation; and
\item a volume-uniform coarse boundary transfer with
      \(b\gamma_Q\asymp\delta\) after centre, gauge, harmonic, and
      topological common factors are removed.
\end{enumerate}
Finite-fibre Wilson coercivity is relevant only to the first candidate
row.  The massless transverse infrared tower lives in the second and
cannot be removed by exact local refresh alone.
\end{assessment}

\subsection{Temporal block coercivity card}
\label{sec:s6v56-card}

\begin{definition}[Two-level temporal block card, TTBC]
\label{def:s6v56-TTBC}
TTBC consists of:
\begin{enumerate}[label=\textup{(\roman*)},leftmargin=3.5em]
\item an exact coarse boundary map \(r_j\) compatible with reflection,
      gauge quotient, and the protected projector;
\item the conditional expectation pair \(J_j,K_j\) and a positive
      reversible coarse transfer \(Q_j\);
\item proof that the physical transfer form satisfies
      \eqref{eq:s6v56-block-domination};
\item uniform physical rows
      \(a_j/\delta_j\ge\mu_f>0\) and
      \(b_j\gamma_{Q,j}/\delta_j\ge\mu_c>0\);
\item an update schedule whose operator and physical duration are declared
      before any tensorization;
\item overlap/comparison errors charged in signed form order
      \(o(\delta_j)\), or absorbed while retaining a strict coefficient;
\item common-refinement projectivity and Mosco convergence for the same
      decomposed forms; and
\item density and noncollapse of the complete relative OS source core.
\end{enumerate}
\end{definition}

\begin{theorem}[Conditional TTBC mass-gap handoff]
\label{thm:s6v56-TTBC}
If TTBC is populated, the reconstructed Yang--Mills Hamiltonian has
relative vacuum gap at least
\[
 \min\{\mu_f,\mu_c\}.
\]
\end{theorem}

\begin{proof}
Rows~\textup{(i)--(vi)} and
\Cref{thm:s6v56-compiler} give the uniform regulator form inequality with
constant \(\min\{\mu_f,\mu_c\}\).  Rows~\textup{(vii)--(viii)} permit the
V53 Mosco/OS handoff.
\end{proof}

\begin{assessment}[Exact frontier after V56]
\label{ass:s6v56-frontier}
V56 converts the temporal conductance bottleneck into a local/coarse
decomposition.  The fine row can plausibly use the existing pinned-fibre
geometry only after a reflected-transfer comparison is proved.  The coarse
row is the irreducible infrared problem: it must control the transverse
centre-neutral boundary mode uniformly in volume.  The next iteration
should derive a Schur-complement formula for this coarse form and determine
whether the perfect quadratic owner leaves a positive nonlinear remainder
or merely the massless symbol already isolated in V49.
\end{assessment}

\begin{programme}[V57: coarse-form Schur complement]
\label{prog:s6v56-V57}
\begin{enumerate}[leftmargin=3em]
\item Write the exact infimal Schur complement of a two-level quadratic
      form and prove its gap comparison.
\item Separate gauge/common null directions before inversion.
\item Insert the V49 perfect quadratic symbol and recover its massless
      transverse low-momentum edge.
\item State the exact nonlinear remainder inequality that would be needed
      to lift that edge by an \(L^2\) amount.
\item Test whether CMAC, PFS, or the signed Wilsonian head controls that
      remainder with the required sign.
\end{enumerate}
\end{programme}

\begin{firewall}[Claim boundary after sixth-series V56]
\label{fire:s6v56-boundary}
V56 proves the exact variance and two-level transfer-form decompositions;
proves a temporal fine/coarse coercivity compiler; quantifies additive
comparison error; proves parallel tensorization without volume loss and
the exact random-scan \(N^{-1}\) debit; and isolates the coarse infrared
mode left untouched by exact local refresh.

V56 does not prove the physical Yang--Mills block-form domination, identify
the update schedule of the reflected half-slab, establish either uniform
fine or coarse physical row, control overlaps, close Mosco/projectivity,
construct the continuum OS theory, or prove the full Yang--Mills mass gap.
\end{firewall}

\begin{theorem}[Sixth-series V56 result ledger]
\label{thm:s6v56-results}
The positive conclusions in \Cref{fire:s6v56-boundary} are exactly those of
\Cref{thm:s6v56-variance,
thm:s6v56-two-level,
thm:s6v56-compiler,
cor:s6v56-direct-comparison,
lem:s6v56-additive-debit,
thm:s6v56-parallel,
thm:s6v56-random-scan,
cor:s6v56-sweep,
prop:s6v56-coarse-firewall,
thm:s6v56-TTBC}.
\end{theorem}

\begin{proof}
Each ledger statement is the cited result with its exact carrier,
schedule, comparison, and physical-clock hypotheses unchanged.
\end{proof}

\paragraph{Parent-version record.}
V56 was formed from
\texttt{Renewal\_Yang\_Mills\_Mass\_Gap\_Research\_Notes\_V55.tex}.
The V55 parent had \(24\,717\) logical source lines,
\(869\,402\) bytes, and SHA--256
\[
 \texttt{c5556d6f8a55f21ce1951d3b742a90453d7a164f2a5e0c3574237ac1621742b2}.
\]
The next compulsory companion audit is V60.

% =============================================================================
% END APPENDED SIXTH-SERIES V56 MATERIAL
% =============================================================================

% =============================================================================
% BEGIN APPENDED SIXTH-SERIES V57 MATERIAL
% =============================================================================

\section{V57: the transfer-form Schur complement and the coarse
infrared remainder}
\label{sec:s6v57}

\subsection{Exact Schur reduction of a nonnegative form}
\label{sec:s6v57-Schur}

Let \({\cal X},{\cal Y}\) be Hilbert spaces and consider the bounded
self-adjoint block operator
\begin{equation}
 L=
 \begin{pmatrix}
 A&B^*\\
 B&C
 \end{pmatrix}\succeq0
 \quad\hbox{on }{\cal X}\oplus{\cal Y}.
 \label{eq:s6v57-block}
\end{equation}
Here \({\cal X}\) is the retained coarse relative carrier and
\({\cal Y}\) is the eliminated conditional fine carrier.

\begin{theorem}[Strict-fibre Schur identity]
\label{thm:s6v57-Schur}
Assume \(C\succeq cI_{\cal Y}\) for some \(c>0\).  Define
\begin{equation}
 S:=A-B^*C^{-1}B.
 \label{eq:s6v57-S}
\end{equation}
Then \(S\succeq0\), and for every \(x\in{\cal X}\),
\begin{equation}
 \boxed{
 \inf_{y\in{\cal Y}}
 \left\langle\binom{x}{y},
 L\binom{x}{y}\right\rangle
 =\langle x,Sx\rangle.}
 \label{eq:s6v57-infimum}
\end{equation}
The unique minimizer is \(y_*=-C^{-1}Bx\), and the exact completed-square
identity is
\begin{equation}
 \left\langle\binom{x}{y},
 L\binom{x}{y}\right\rangle
 =\langle x,Sx\rangle
 +\langle y+C^{-1}Bx,C(y+C^{-1}Bx)\rangle.
 \label{eq:s6v57-square}
\end{equation}
\end{theorem}

\begin{proof}
Expand the right side of \eqref{eq:s6v57-square}.  Its cross term is
\(2\operatorname{Re}\langle Bx,y\rangle\), its \(x\)-quadratic return is
\(\langle x,B^*C^{-1}Bx\rangle\), and this cancels the subtraction in
\(S\), recovering \eqref{eq:s6v57-block}.  The second term is nonnegative
and vanishes exactly at \(y_*\), proving the infimum.  Since \(L\succeq0\),
the infimum is nonnegative for every \(x\), hence \(S\succeq0\).
\end{proof}

\begin{theorem}[Null-fibre compatibility firewall]
\label{thm:s6v57-null}
Suppose only \(C\succeq0\) and \(L\succeq0\).  Then necessarily
\begin{equation}
 B^*n=0\qquad(n\in\ker C).
 \label{eq:s6v57-null-compatibility}
\end{equation}
If \(C\) has closed range and is bounded below on
\((\ker C)^\perp\), the infimal formula remains valid with the
Moore--Penrose inverse:
\begin{equation}
 S=A-B^*C^\dagger B,
 \qquad
 \inf_y\left\langle\binom{x}{y},L\binom{x}{y}\right\rangle
 =\langle x,Sx\rangle.
 \label{eq:s6v57-pseudo}
\end{equation}
The minimizers are
\(-C^\dagger Bx+\ker C\).
\end{theorem}

\begin{proof}
For \(n\in\ker C\), positivity of the quadratic polynomial
\[
 \left\langle\binom{x}{tn},
 L\binom{x}{tn}\right\rangle
 =\langle x,Ax\rangle
 +2t\operatorname{Re}\langle Bx,n\rangle
\]
for every real \(t\) forces the real part of
\(\langle Bx,n\rangle\) to vanish.  Replacing \(x\) by \(ix\) forces the
imaginary part to vanish, proving
\eqref{eq:s6v57-null-compatibility}.  Thus \(Bx\in(\ker C)^\perp\), where
\(C\) is invertible by the closed-range lower bound.  Complete the square
there as in \Cref{thm:s6v57-Schur}; the null component of \(y\) contributes
nothing.
\end{proof}

\begin{assessment}[Gauge directions must be removed before inversion]
\label{ass:s6v57-gauge}
Equation \eqref{eq:s6v57-null-compatibility} is automatic for an exact
positive transfer form, but only after the block decomposition represents
the true physical Hilbert carrier.  In a coordinate Hessian with unfixed
gauge directions, an apparent failure of compatibility can signal that the
block is not a physical nonnegative operator.  A raw inverse of \(C\) may
neither exist nor respect the gauge/common-sector quotient.
\end{assessment}

\subsection{From coarse and fine gaps to a full form gap}
\label{sec:s6v57-gap}

\begin{theorem}[Quantitative reconstruction from the Schur gap]
\label{thm:s6v57-reconstruction}
Under \Cref{thm:s6v57-Schur}, suppose
\[
 S\succeq sI_{\cal X},\qquad
 C\succeq cI_{\cal Y},\qquad
 \|C^{-1}B\|\le r.
\]
Then
\begin{equation}
 \boxed{
 L\succeq
 \min\left\{{s\over1+2r^2},{c\over2}\right\}
 I_{{\cal X}\oplus{\cal Y}}.}
 \label{eq:s6v57-full-gap}
\end{equation}
\end{theorem}

\begin{proof}
Put \(z=y+C^{-1}Bx\).  Then
\[
 \|x\|^2+\|y\|^2
 \le(1+2r^2)\|x\|^2+2\|z\|^2.
\]
By \eqref{eq:s6v57-square},
\[
 \left\langle\binom{x}{y},L\binom{x}{y}\right\rangle
 \ge s\|x\|^2+c\|z\|^2.
\]
Comparison of the two displayed inequalities gives
\eqref{eq:s6v57-full-gap}.
\end{proof}

\begin{theorem}[Massless Schur edge forces a massless full form]
\label{thm:s6v57-massless-Schur}
Under \Cref{thm:s6v57-Schur}, if there are unit vectors \(x_n\in{\cal X}\)
such that
\[
 \langle x_n,Sx_n\rangle\longrightarrow0,
\]
then \(L\) has no positive lower spectral bound on
\({\cal X}\oplus{\cal Y}\).
\end{theorem}

\begin{proof}
Choose \(y_n=-C^{-1}Bx_n\).  The vector
\(\xi_n=(x_n,y_n)\) has norm at least one, and
\[
 \langle\xi_n,L\xi_n\rangle
 =\langle x_n,Sx_n\rangle\to0.
\]
After normalizing \(\xi_n\), its Rayleigh quotient still tends to zero.
\end{proof}

\begin{corollary}[Fine coercivity alone cannot create the gap]
\label{cor:s6v57-fine-not-enough}
Even an arbitrarily large uniform lower bound on the eliminated block
\(C\) cannot produce a full gap when the effective coarse Schur form
\(S\) has an infrared sequence of vanishing Rayleigh quotient.
\end{corollary}

\begin{proof}
This is exactly \Cref{thm:s6v57-massless-Schur}; the minimizing fine vector
tracks the coarse mode and cancels the return \(B^*C^{-1}B\).
\end{proof}

\subsection{Stability of the exact coarse operator}
\label{sec:s6v57-stability}

\begin{lemma}[Schur-complement perturbation bound]
\label{lem:s6v57-Schur-stability}
Let
\[
 S=A-B^*C^{-1}B,\qquad
 S_0=A_0-B_0^*C_0^{-1}B_0,
\]
where \(C,C_0\succeq cI\) and
\(\|B\|,\|B_0\|\le M\).  Then
\begin{equation}
 \boxed{
 \|S-S_0\|
 \le\|A-A_0\|
 +{2M\over c}\|B-B_0\|
 +{M^2\over c^2}\|C-C_0\|.}
 \label{eq:s6v57-Schur-stability}
\end{equation}
\end{lemma}

\begin{proof}
Use
\[
 C^{-1}-C_0^{-1}
 =C^{-1}(C_0-C)C_0^{-1},
\]
so its norm is at most \(c^{-2}\|C-C_0\|\).  Add and subtract
\(B_0^*C^{-1}B\) and \(B_0^*C^{-1}B_0\) in the quadratic return.  The two
outer-factor differences cost at most
\(M c^{-1}\|B-B_0\|\) each, and the inverse difference costs at most
\(M^2c^{-2}\|C-C_0\|\).  Add \(\|A-A_0\|\).
\end{proof}

\begin{proposition}[A norm-small Schur correction cannot hide a fixed
mass]
\label{prop:s6v57-small-correction}
Suppose \(S_0\succeq0\) has unit infrared vectors \(x_n\) with
\(\langle x_n,S_0x_n\rangle\to0\), and
\(\|S-S_0\|\le\varepsilon\).  Then
\begin{equation}
 \inf\operatorname{spec}S\le\varepsilon.
 \label{eq:s6v57-small-gap}
\end{equation}
Consequently, proving \(S\succeq\mu I\) requires a coarse correction of
operator or form size at least \(\mu\) on every such infrared sequence.
\end{proposition}

\begin{proof}
For each \(x_n\),
\[
 \langle x_n,Sx_n\rangle
 \le\langle x_n,S_0x_n\rangle+\varepsilon.
\]
Take the lower spectral infimum and then \(n\to\infty\).
\end{proof}

\subsection{Two different Schur complements}
\label{sec:s6v57-two-Schur}

\begin{assessment}[Classical action Schur versus transfer-form Schur]
\label{ass:s6v57-two-Schur}
The V49 symbol
\[
 s_0(t)
 ={t(t+8)(t+24)\over9(t^2+18t+16)}
 ={4\over3}t+O(t^2)
 \qquad(t\downarrow0)
 \tag{57.1}
\]
is the Schur complement of the \emph{quadratic classical action Hessian}
after eliminating fine lattice variables.  The operator \(S\) in
\eqref{eq:s6v57-S} is instead the Schur complement of the
\emph{rescaled OS transfer Dirichlet operator}
\(\delta^{-1}(I-P)\) on a decomposition of the physical Hilbert space.
They are not identical by algebra.  An incidence theorem deriving the
physical half-slab transfer form from the response-matched action is
required before the V49 symbol can be inserted into
\eqref{eq:s6v57-S}.
\end{assessment}

\begin{proof}
The small-\(t\) expansion in \((57.1)\) follows by differentiating the
displayed rational function at zero:
\[
 s_0'(0)={8\cdot24\over9\cdot16}={4\over3}.
\]
The two operators act on different spaces and arise from different
quadratic variations: field-coordinate variation of the action versus
Hilbert-space variation of a Markov transfer.  Equality would therefore be
an additional theorem, not a formal consequence of sharing the word
Schur.
\end{proof}

\begin{corollary}[Even a successful incidence leaves a massless baseline]
\label{cor:s6v57-baseline}
If a future incidence theorem identifies the low-momentum transfer coarse
baseline with a positive multiple of \(s_0(t(k))\), and
\(t(k)\to0\) along physical transverse momenta, then that baseline has no
uniform positive gap.  The exact nonlinear Schur remainder must satisfy,
on every normalized infrared sequence \(x_j\),
\begin{equation}
 \liminf_j
 \langle x_j,(S_j-S_{0,j})x_j\rangle\ge\mu
 \label{eq:s6v57-required-remainder}
\end{equation}
to prove a gap \(\mu\).
\end{corollary}

\begin{proof}
Equation \((57.1)\) makes the baseline Rayleigh quotient tend to zero.
The required lower bound is then
\Cref{thm:s6v53-coercivity-alternative} applied to the exact coarse Schur
form.
\end{proof}

\subsection{Testing the accumulated nonlinear estimates}
\label{sec:s6v57-test}

\begin{assessment}[CMAC]
\label{ass:s6v57-CMAC}
CMAC controls the Jacobian/entropy derivatives of a nonlinear Morse
coordinate map around the full perfect quadratic baseline.  It is a local
smoothness and owner-comparison packet.  No current CMAC row gives the
positive infrared matrix element
\eqref{eq:s6v57-required-remainder}; if its remainder remains norm-small,
\Cref{prop:s6v57-small-correction} shows that it cannot open a fixed gap.
\end{assessment}

\begin{assessment}[PFS]
\label{ass:s6v57-PFS}
PFS replaces impossible uniform boundary fillability by a probabilistic
good-fill event plus a mixed misalignment debit.  It can control exceptional
geometry in an action comparison, but it neither identifies the
transfer-form Schur operator nor supplies a signed positive lower bound on
its transverse infrared remainder.
\end{assessment}

\begin{assessment}[Signed Wilsonian head]
\label{ass:s6v57-head}
The signed Wilsonian head is the only accumulated object retaining the
signs of direct second derivatives and integrated covariance returns.
It is therefore the correct candidate from which
\eqref{eq:s6v57-required-remainder} might be derived.  Present estimates
bound its tails and local jets but do not prove that the completed
coarse transfer form is \(L^2\)-positive by a uniform amount on every
centre-neutral infrared mode.
\end{assessment}

\begin{theorem}[Coarse Schur completion criterion]
\label{thm:s6v57-completion}
Suppose the exact response-matched transfer generator admits the physical
block form \eqref{eq:s6v57-block} after common-sector removal, uniformly in
regulator and volume, and assume:
\begin{enumerate}[label=\textup{(\roman*)},leftmargin=3.5em]
\item \(C_j\succeq cI\) and
      \(\|C_j^{-1}B_j\|\le r\);
\item \(S_j=A_j-B_j^*C_j^{-1}B_j\succeq sI\);
\item the block forms Mosco-converge projectively to the OS Hamiltonian
      form.
\end{enumerate}
Then the limiting Hamiltonian gap is at least
\[
 \min\left\{{s\over1+2r^2},{c\over2}\right\}.
\]
\end{theorem}

\begin{proof}
\Cref{thm:s6v57-reconstruction} gives the displayed uniform regulator
bound.  V53 transports it through the projective Mosco limit.
\end{proof}

\begin{assessment}[Exact frontier after V57]
\label{ass:s6v57-frontier}
The upstream target is no longer an unspecified nonlinear correction.  It
is the operator inequality
\[
 A_j-B_j^*C_j^{-1}B_j\succeq sI
\]
for the rescaled physical transfer form on the complete transverse
centre-neutral coarse carrier.  V49 proves that the natural quadratic
baseline is massless.  CMAC and PFS do not provide the missing sign.  The
signed Wilsonian head preserves the information needed to attempt it, but
an action-to-transfer incidence theorem must come first.
\end{assessment}

\begin{programme}[V58: action-to-transfer form incidence]
\label{prog:s6v57-V58}
\begin{enumerate}[leftmargin=3em]
\item Starting from a positive half-slab kernel \(F=e^{-W/2}\), derive the
      exact Markov pair-law Dirichlet form in terms of \(W\) without
      replacing the boundary difference by a local Hessian.
\item Prove the small-time diffusion expansion under explicit regularity
      and identify its carré-du-champ owner.
\item Determine which action Hessian information enters the leading
      transfer form and which enters only higher-order drift or measure
      terms.
\item Apply the result to compact group boundary variables after gauge
      shorting.
\item If the required small-time kernel asymptotic is absent from the
      sources, isolate it as the next physical acquisition card.
\end{enumerate}
\end{programme}

\begin{firewall}[Claim boundary after sixth-series V57]
\label{fire:s6v57-boundary}
V57 proves strict and null-compatible Schur identities; reconstructs a
full form gap from fine and coarse Schur gaps; proves that a massless coarse
Schur sequence defeats arbitrary fine coercivity; proves a quantitative
Schur perturbation bound; distinguishes the classical action and OS
transfer Schur complements; and identifies the precise signed nonlinear
coarse remainder required for a gap.

V57 does not prove action-to-transfer incidence, the positive nonlinear
Schur remainder, a uniform Yang--Mills coarse or fine form gap,
Mosco/projective convergence, the continuum OS theory, or the full
Yang--Mills mass gap.
\end{firewall}

\begin{theorem}[Sixth-series V57 result ledger]
\label{thm:s6v57-results}
The positive conclusions in \Cref{fire:s6v57-boundary} are exactly those of
\Cref{thm:s6v57-Schur,
thm:s6v57-null,
thm:s6v57-reconstruction,
thm:s6v57-massless-Schur,
cor:s6v57-fine-not-enough,
lem:s6v57-Schur-stability,
prop:s6v57-small-correction,
ass:s6v57-two-Schur,
cor:s6v57-baseline,
thm:s6v57-completion}.
\end{theorem}

\begin{proof}
Each conclusion is the cited result with its positivity, gauge-short,
closed-range, incidence, and convergence hypotheses retained.
\end{proof}

\paragraph{Parent-version record.}
V57 was formed from
\texttt{Renewal\_Yang\_Mills\_Mass\_Gap\_Research\_Notes\_V56.tex}.
The V56 parent had \(25\,156\) logical source lines,
\(884\,377\) bytes, and SHA--256
\[
 \texttt{259211bbe50bf7cf53c1b828b7a7e10555a896270f02643cf275f0241c0ff5b6}.
\]
The next compulsory companion audit is V60.

% =============================================================================
% END APPENDED SIXTH-SERIES V57 MATERIAL
% =============================================================================

% =============================================================================
% BEGIN APPENDED SIXTH-SERIES V58 MATERIAL
% =============================================================================

\section{V58: exact action-to-transfer incidence through the Doob
conditional variance}
\label{sec:s6v58}

\subsection{The half-slab law determines the full transfer form}
\label{sec:s6v58-Doob}

Let \(({\mathsf X},\mu_0)\) and
\(({\mathsf Y},\mu_{1/2})\) be finite-regulator boundary spaces, and let
\[
 F(y,x)>0,\qquad
 (Af)(y)=\int F(y,x)f(x)\,\mu_0(\dd x)
\]
be a Hilbert--Schmidt half-slab amplitude.  Let
\[
 Au_0=s_1v_0,\qquad A^*v_0=s_1u_0
\]
be its normalized positive leading singular pair.  Put
\[
 \pi_0(\dd x)=u_0(x)^2\mu_0(\dd x),\qquad
 \pi_{1/2}(\dd y)=v_0(y)^2\mu_{1/2}(\dd y)
\]
and
\begin{equation}
 \Lambda(\dd x,\dd y)
 ={1\over s_1}u_0(x)F(y,x)v_0(y)
 \,\mu_0(\dd x)\mu_{1/2}(\dd y).
 \label{eq:s6v58-Doob-law}
\end{equation}
Let
\[
 (Kf)(y)=\mathbb E_\Lambda[f(X)\mid Y=y].
\]

\begin{theorem}[Doob conditional-variance identity]
\label{thm:s6v58-conditional-variance}
The operator
\[
 P:=K^*K
\]
is the normalized physical full-slab transfer, unitarily equivalent to
\(s_1^{-2}A^*A\).  For every \(f\in L^2(\pi_0)\),
\begin{equation}
 \boxed{
 \langle f,(I-P)f\rangle_{L^2(\pi_0)}
 =\mathbb E_\Lambda
   \operatorname{Var}_\Lambda(f(X)\mid Y).}
 \label{eq:s6v58-conditional-form}
\end{equation}
Equivalently,
\begin{equation}
 \langle f,(I-P)f\rangle
 =\inf_{h\in L^2(\pi_{1/2})}
 \mathbb E_\Lambda|f(X)-h(Y)|^2,
 \label{eq:s6v58-infimum}
\end{equation}
with unique minimizer \(h=Kf\) modulo null sets.
\end{theorem}

\begin{proof}
The source half-slab unitary equivalence follows directly from
\eqref{eq:s6v58-Doob-law}:
\[
 v_0(y)(Kf)(y)
 ={1\over s_1}\int F(y,x)u_0(x)f(x)\,\mu_0(\dd x).
\]
Thus multiplication by \(u_0,v_0\) conjugates \(K\) to \(A/s_1\), and
conjugates \(K^*K\) to \(A^*A/s_1^2\).
Conditional expectation is the orthogonal projection of \(f(X)\) onto
the \(Y\)-measurable subspace of \(L^2(\Lambda)\), so
\[
 \mathbb E|f(X)-Kf(Y)|^2
 =\|f\|_{L^2(\pi_0)}^2-\|Kf\|_{L^2(\pi_{1/2})}^2
 =\langle f,(I-K^*K)f\rangle.
\]
The projection theorem also gives the infimum statement and its
uniqueness.
\end{proof}

\begin{corollary}[Complete-source half-slab criterion]
\label{cor:s6v58-complete}
After removing the maximal common factor with protected projection
\(\Pi\), the physical rescaled form has gap at least \(\mu\) if and only if
\begin{equation}
 \boxed{
 \mathbb E_\Lambda
 \operatorname{Var}(f(X)\mid Y)
 \ge\mu\delta\,
 \|(I-\Pi)f\|_{L^2(\pi_0)}^2
 \quad(f\in L^2(\pi_0)).}
 \label{eq:s6v58-half-coercivity}
\end{equation}
This is the half-slab version of the V54 two-boundary inequality.
\end{corollary}

\begin{proof}
Divide \eqref{eq:s6v58-conditional-form} by the full-slab duration
\(\delta\) and use \Cref{thm:s6v53-edge-form}.  Composing the two
conditional half-channels produces the V54 full stationary pair law, so
the criteria have the same full-transfer edge.
\end{proof}

\subsection{What the action contributes exactly}
\label{sec:s6v58-action}

Suppose
\begin{equation}
 F(y,x)=e^{-W(y,x)/2}
 \label{eq:s6v58-action-kernel}
\end{equation}
relative to the fixed reference boundary measures.  Then
\eqref{eq:s6v58-Doob-law} becomes
\begin{equation}
 \Lambda(\dd x,\dd y)
 ={1\over s_1}
 \exp\left\{
 -{1\over2}W(y,x)+\log u_0(x)+\log v_0(y)
 \right\}
 \mu_0(\dd x)\mu_{1/2}(\dd y).
 \label{eq:s6v58-effective-pair}
\end{equation}

\begin{theorem}[Exact conditional kernel from the action owner]
\label{thm:s6v58-action-owner}
For \(\pi_{1/2}\)-almost every \(y\), the conditional law entering
\eqref{eq:s6v58-half-coercivity} is
\begin{equation}
 \boxed{
 \Lambda(\dd x\mid y)
 ={u_0(x)e^{-W(y,x)/2}\mu_0(\dd x)
   \over
   \int u_0(x')e^{-W(y,x')/2}\mu_0(\dd x')}.}
 \label{eq:s6v58-conditional-kernel}
\end{equation}
Thus \(W\) alone, and in particular its local Hessian alone, is not the
conditional-variance owner: the leading ground singular function \(u_0\)
and the global normalization are part of the exact conditional law.
\end{theorem}

\begin{proof}
Disintegrate \eqref{eq:s6v58-effective-pair} with respect to \(y\).
The factors depending only on \(y\), including \(v_0(y)\) and the global
\(s_1^{-1}\), cancel from the conditional ratio.  The remaining numerator
and denominator are exactly
\eqref{eq:s6v58-conditional-kernel}.
\end{proof}

\begin{corollary}[Scalar and one-boundary gauge of the amplitude]
\label{cor:s6v58-boundary-gauge}
Multiplying \(F(y,x)\) by a positive scalar changes neither
\(\Lambda(\dd x\mid y)\) nor the normalized transfer.  By contrast, adding a
nonconstant one-boundary term to \(W\) changes the singular ground function
and cannot be removed from \eqref{eq:s6v58-conditional-kernel} without
solving the new singular problem.
\end{corollary}

\begin{proof}
A scalar multiplies \(A\) and all its singular values by the same factor,
and cancels from the conditional ratio.  A nonconstant boundary factor is
a multiplication operator, not a scalar; it changes \(A^*A\) and its
leading eigenfunction in general.
\end{proof}

\begin{assessment}[Precise failure of the local-Hessian shortcut]
\label{ass:s6v58-Hessian}
The V49 classical Hessian controls the quadratic variation of \(W\) near
its local minimum.  The physical form
\eqref{eq:s6v58-conditional-form} instead measures the best
\(Y\)-prediction error for \emph{all} relative \(L^2(\pi_0)\) functions
under the global law \eqref{eq:s6v58-effective-pair}.  Passing from the
former to the latter requires global control of \(u_0\), the tails and
topological components of \(W\), and the common-factor quotient.  No local
Morse chart or Hessian identity performs this passage by itself.
\end{assessment}

\subsection{The small-time generator limit}
\label{sec:s6v58-small-time}

\begin{theorem}[Spectral small-time form limit]
\label{thm:s6v58-spectral-limit}
Let \(H\ge0\) be self-adjoint and \(P_\delta=e^{-\delta H}\).  For every
\(f\in\operatorname{Dom}H^{1/2}\),
\begin{equation}
 \boxed{
 \lim_{\delta\downarrow0}
 {1\over\delta}\langle f,(I-P_\delta)f\rangle
 =\|H^{1/2}f\|^2.}
 \label{eq:s6v58-spectral-limit}
\end{equation}
For arbitrary \(f\in{\cal H}\), the left side increases to \(+\infty\) if
\(f\notin\operatorname{Dom}H^{1/2}\).
\end{theorem}

\begin{proof}
Let \(\nu_f\) be the spectral measure of \(H\).  Then
\[
 {1\over\delta}\langle f,(I-e^{-\delta H})f\rangle
 =\int_{[0,\infty)}
 {1-e^{-\delta\lambda}\over\delta}\,\nu_f(\dd\lambda).
\]
As \(\delta\downarrow0\), the integrand increases pointwise to
\(\lambda\).  Monotone convergence gives
\(\int\lambda\,\dd\nu_f=\|H^{1/2}f\|^2\), finite exactly on the form
domain.
\end{proof}

\begin{corollary}[Conditional variance converges to the Hamiltonian form]
\label{cor:s6v58-conditional-limit}
If the response-matched half-slab laws \(\Lambda_\delta\) generate one
semigroup \(P_\delta=e^{-\delta H}\), then
\begin{equation}
 \lim_{\delta\downarrow0}{1\over\delta}
 \mathbb E_{\Lambda_\delta}
 \operatorname{Var}(f(X)\mid Y)
 =\|H^{1/2}f\|^2.
 \label{eq:s6v58-conditional-limit}
\end{equation}
\end{corollary}

\begin{proof}
Combine \Cref{thm:s6v58-conditional-variance} with
\Cref{thm:s6v58-spectral-limit}.
\end{proof}

\begin{assessment}[Semigroup consistency is an acquisition row]
\label{ass:s6v58-semigroup}
At independent regulators, a declared half-slab duration does not prove
that the kernels form one semigroup in \(\delta\), nor that their rescaled
forms converge to a common \(H\).  Equation
\eqref{eq:s6v58-conditional-limit} is therefore a theorem after
semigroup or Mosco consistency is established, not a definition that can
replace that work.
\end{assessment}

\subsection{Diffusion incidence and its limitation}
\label{sec:s6v58-diffusion}

\begin{theorem}[Carr\'e-du-champ form of a reversible boundary diffusion]
\label{thm:s6v58-diffusion}
Let \(M\) be a compact Riemannian manifold without boundary, let
\(\pi=\rho\,\mathrm{vol}\) with smooth \(\rho>0\), and let \(a(x)\) be a
smooth symmetric positive bundle endomorphism.  Define on smooth functions
\[
 {\cal L}f={1\over\rho}\operatorname{div}(\rho\,a\nabla f).
\]
Then \({\cal L}\) is symmetric and nonpositive in \(L^2(\pi)\), and the
closed form of \(H=-{\cal L}\) is
\begin{equation}
 \boxed{
 {\cal E}(f)
 =\int_M\langle\nabla f,a\nabla f\rangle\,\dd\pi.}
 \label{eq:s6v58-carre}
\end{equation}
If \(a\succeq\alpha I\) and
\[
 \operatorname{Var}_\pi(f)
 \le C_{\rm P}\int_M|\nabla f|^2\,\dd\pi,
\]
then \(H\) has gap at least \(\alpha/C_{\rm P}\).
\end{theorem}

\begin{proof}
Integration by parts on the compact boundaryless manifold gives
\[
 -\langle f,{\cal L}g\rangle_{L^2(\pi)}
 =\int_M\langle\nabla f,a\nabla g\rangle\,\dd\pi.
\]
Closure gives \eqref{eq:s6v58-carre}.  Ellipticity and the Poincar\'e
inequality imply
\[
 {\cal E}(f)\ge{\alpha\over C_{\rm P}}
 \operatorname{Var}_\pi(f),
\]
which is the spectral-gap variational principle.
\end{proof}

\begin{assessment}[The action Hessian is not the carr\'e du champ]
\label{ass:s6v58-carre}
In a reversible diffusion, the principal tensor \(a\) owns the local
energy in \eqref{eq:s6v58-carre}; the potential hidden in \(\rho\) changes
the global Poincar\'e constant.  A Hessian of the Wilson action can help
prove a Poincar\'e inequality only through an additional global theorem
such as a valid convexity/curvature criterion on the physical quotient.
The compact non-Abelian configuration space, gauge null directions,
multiple sectors, and increasing volume prevent importing the pinned local
Hessian as a uniform \(C_{\rm P}\) bound.
\end{assessment}

\begin{proposition}[Uniform ellipticity without uniform Poincar\'e is
insufficient]
\label{prop:s6v58-ellipticity}
There are reversible diffusions with \(a=I\) on circles of length \(L\)
whose gaps tend to zero as \(L\to\infty\).
\end{proposition}

\begin{proof}
Take normalized Lebesgue measure on the circle and
\({\cal L}=\partial_x^2\).  The first nonzero eigenvalue of
\(-{\cal L}\) is \((2\pi/L)^2\), which tends to zero.
\end{proof}

\subsection{A global conditional-variance comparison}
\label{sec:s6v58-comparison}

\begin{lemma}[Density-ratio comparison of half-slab coercivity]
\label{lem:s6v58-density-comparison}
Let \(\Lambda\) and \(\Lambda_0\) be two probability laws on
\({\mathsf X}\times{\mathsf Y}\) with the same measurable spaces, and
assume
\begin{equation}
 0<m\le {d\Lambda\over d\Lambda_0}\le M<\infty.
 \label{eq:s6v58-density-ratio}
\end{equation}
Let \(\pi,\pi_0\) be their \(X\)-marginals.  Then, for every \(f\),
\begin{align}
 \inf_h\mathbb E_\Lambda|f(X)-h(Y)|^2
 &\ge m\inf_h\mathbb E_{\Lambda_0}|f(X)-h(Y)|^2,
 \label{eq:s6v58-numerator-comparison}\\
 \operatorname{Var}_\pi(f)
 &\le M\inf_c\mathbb E_{\pi_0}|f-c|^2
 =M\operatorname{Var}_{\pi_0}(f).
 \label{eq:s6v58-denominator-comparison}
\end{align}
Consequently a reference half-slab coefficient
\[
 \inf_h\mathbb E_{\Lambda_0}|f-h|^2
 \ge\gamma_0\operatorname{Var}_{\pi_0}(f)
\]
implies the target coefficient at least
\begin{equation}
 \gamma\ge{m\over M}\gamma_0.
 \label{eq:s6v58-comparison-gap}
\end{equation}
\end{lemma}

\begin{proof}
The lower density bound holds for every \(h\), so it survives taking the
infimum, proving \eqref{eq:s6v58-numerator-comparison}.  For every constant
\(c\), the upper marginal density bound gives
\[
 \mathbb E_\pi|f-c|^2\le
 M\mathbb E_{\pi_0}|f-c|^2.
\]
Take the infimum over \(c\).  Combine the two estimates.
\end{proof}

\begin{assessment}[Volume oscillation kills a naive comparison]
\label{ass:s6v58-volume-comparison}
For Gibbs weights, a pointwise ratio bound of the form
\eqref{eq:s6v58-density-ratio} typically costs
\(M/m\sim e^{\operatorname{osc}(W-W_0)}\).  If the action difference is a
sum over the volume, this comparison can decay exponentially with volume
even when every local summand is small.  A useful Yang--Mills argument must
therefore use local conditional comparison, a tensorized curvature
estimate, or a signed form cancellation; a global Holley--Stroock ratio is
not automatically uniform.
\end{assessment}

\subsection{Action-to-transfer incidence card}
\label{sec:s6v58-card}

\begin{definition}[Doob conditional-form card, DCFC]
\label{def:s6v58-DCFC}
DCFC consists of:
\begin{enumerate}[label=\textup{(\roman*)},leftmargin=3.5em]
\item the exact response-matched action \(W_j\), reference boundary
      measures, and positive half-slab amplitude
      \(F_j=e^{-W_j/2}\);
\item the leading singular functions \(u_{0,j},v_{0,j}\), maximal common
      factor, and exact physical relative carrier;
\item the conditional law
      \eqref{eq:s6v58-conditional-kernel} and the complete inequality
      \eqref{eq:s6v58-half-coercivity};
\item a proof that the half-slab family composes at its declared physical
      durations, or Mosco convergence of the rescaled conditional-variance
      forms;
\item a volume-uniform global Poincar\'e/conductance or signed
      conditional-form comparison, not merely a local action Hessian;
\item adjacent singular-ground-function and projector transport at the
      inverse-margin scale;
\item common-refinement projectivity; and
\item noncollapse, OS reconstruction, and cyclicity.
\end{enumerate}
\end{definition}

\begin{theorem}[Conditional DCFC completion]
\label{thm:s6v58-DCFC}
If DCFC is populated with constant \(\mu>0\), the reconstructed
Yang--Mills Hamiltonian has relative vacuum gap at least \(\mu\).
\end{theorem}

\begin{proof}
Rows~\textup{(i)--(iii)} and
\Cref{thm:s6v58-conditional-variance} give the uniform regulator
Dirichlet-form inequality for the actual response-matched transfer.
Rows~\textup{(iv)--(vii)} give a unique projective Mosco limit of those
forms.  Row~\textup{(viii)} identifies the nontrivial limiting form with
the reconstructed OS Hamiltonian, to which V53 applies.
\end{proof}

\begin{assessment}[Exact frontier after V58]
\label{ass:s6v58-frontier}
The action-to-transfer incidence is now exact and exposes the missing
global variables: the ground singular function and the conditional
half-slab Poincar\'e coefficient.  The remaining mass problem cannot be
settled by differentiating the Wilson action at its minimum.  It requires
a volume-uniform lower bound in
\eqref{eq:s6v58-half-coercivity} for the complete centre-neutral carrier,
plus cofinal control of the singular ground functions.  The next attack
should test conditional tensorization: whether local half-slab conditional
variances and a coarse boundary law can yield the global coefficient
without an exponential density-ratio loss.
\end{assessment}

\begin{programme}[V59: conditional tensorization without volume
oscillation]
\label{prog:s6v58-V59}
\begin{enumerate}[leftmargin=3em]
\item Prove the Efron--Stein variance decomposition for a block cover and
      state the exact dependency-matrix correction.
\item Derive a global Poincar\'e bound from uniform conditional block gaps
      and a strict inter-block influence norm.
\item Track the physical \(\delta\) factor and the Ward/common-sector
      quotient.
\item Test the existing MSEC coefficient against the dependency matrix,
      respecting V49's proof that a fixed one-site margin is incompatible
      with the critical Gaussian owner.
\item Decide whether a multiscale rather than one-site tensorization is
      mathematically necessary.
\end{enumerate}
\end{programme}

\begin{firewall}[Claim boundary after sixth-series V58]
\label{fire:s6v58-boundary}
V58 proves that the full-slab transfer form is exactly a conditional
variance under the Doob half-slab law; derives the exact conditional kernel
from the response-matched action and its ground singular function; proves
the spectral small-time form limit; states the reversible diffusion
carr\'e-du-champ incidence and its global Poincar\'e requirement; and proves
a density-ratio comparison with its volume-loss firewall.

V58 does not prove the Yang--Mills complete conditional-variance
inequality, control the singular ground functions cofinally, establish a
uniform global Poincar\'e or conductance bound, prove semigroup/Mosco
consistency or projectivity, construct the continuum OS theory, or prove
the full Yang--Mills mass gap.
\end{firewall}

\begin{theorem}[Sixth-series V58 result ledger]
\label{thm:s6v58-results}
The positive conclusions in \Cref{fire:s6v58-boundary} are exactly those of
\Cref{thm:s6v58-conditional-variance,
cor:s6v58-complete,
thm:s6v58-action-owner,
cor:s6v58-boundary-gauge,
thm:s6v58-spectral-limit,
cor:s6v58-conditional-limit,
thm:s6v58-diffusion,
prop:s6v58-ellipticity,
lem:s6v58-density-comparison,
thm:s6v58-DCFC}.
\end{theorem}

\begin{proof}
Each ledger item is the conclusion of the cited result with its
positivity, singular-ground-state, clock, carrier, and uniformity
hypotheses retained.
\end{proof}

\paragraph{Parent-version record.}
V58 was formed from
\texttt{Renewal\_Yang\_Mills\_Mass\_Gap\_Research\_Notes\_V57.tex}.
The V57 parent had \(25\,588\) logical source lines,
\(898\,940\) bytes, and SHA--256
\[
 \texttt{5afa9c0321832ad4fe04962ec465d94a13937cd4678911f08a360fdc1c7cafe1}.
\]
The next compulsory companion audit is V60.

% =============================================================================
% END APPENDED SIXTH-SERIES V58 MATERIAL
% =============================================================================

% =============================================================================
% BEGIN APPENDED SIXTH-SERIES V59 MATERIAL
% =============================================================================

\section{V59: conditional tensorization and the critical dependency
matrix}
\label{sec:s6v59}

\subsection{Product variance and its exact martingale proof}
\label{sec:s6v59-product}

Let
\[
 ({\mathsf X},\pi)=\prod_{i=1}^N({\mathsf X}_i,\pi_i)
\]
and let \(\mathcal F_i=\sigma(X_1,\ldots,X_i)\).  Put
\[
 M_i=\mathbb E[f\mid\mathcal F_i],
 \qquad d_i=M_i-M_{i-1},
 \qquad M_0=\mathbb Ef.
\]

\begin{theorem}[Efron--Stein tensorization]
\label{thm:s6v59-Efron-Stein}
For every \(f\in L^2(\pi)\),
\begin{equation}
 \boxed{
 \operatorname{Var}_\pi(f)
 =\sum_{i=1}^N\|d_i\|_2^2
 \le\sum_{i=1}^N
 \mathbb E_\pi\operatorname{Var}
 (f\mid X_1,\ldots,X_{i-1},X_{i+1},\ldots,X_N).}
 \label{eq:s6v59-Efron-Stein}
\end{equation}
\end{theorem}

\begin{proof}
The martingale differences \(d_i\) are mutually orthogonal and sum to
\(f-\mathbb Ef\), proving the equality.  Let \(X_i'\) be an independent
copy of \(X_i\), with all other coordinates fixed, and let \(f^{(i)}\) be
the resulting value.  Conditional Jensen gives
\[
 \|d_i\|_2^2
 \le {1\over2}\mathbb E|f-f^{(i)}|^2
 =\mathbb E\operatorname{Var}
 (f\mid X_{-i}).
\]
Summing proves the inequality.
\end{proof}

\begin{assessment}[Static tensorization is not yet channel tensorization]
\label{ass:s6v59-static}
Equation \eqref{eq:s6v59-Efron-Stein} decomposes variance under one product
boundary law.  The Yang--Mills numerator is the cross-boundary prediction
loss
\(\mathbb E\operatorname{Var}(f(X)\mid Y)\).  To obtain a mass, the
martingale pieces must be seen with a quantitative loss by the same
half-slab channel.  Efron--Stein alone does not supply that incidence.
\end{assessment}

\subsection{An abstract dependency-matrix compiler}
\label{sec:s6v59-compiler}

Let \({\cal H}_{\rm rel}\) be a relative Hilbert carrier and let
\[
 P=K^*K,\qquad
 T=(I-P)^{1/2}.
\]
Thus the dimensionless half-slab prediction loss is
\[
 {\cal D}(f)=\|Tf\|^2.
\]
Assume the carrier has an orthogonal decomposition
\begin{equation}
 I_{\rm rel}=\sum_{i=1}^ND_i,
 \qquad D_iD_j=0\quad(i\ne j),
 \label{eq:s6v59-input-arms}
\end{equation}
and the loss space has analysis maps \(R_i\) satisfying the Bessel bound
\begin{equation}
 \sum_{i=1}^N\|R_i z\|^2\le\|z\|^2.
 \label{eq:s6v59-output-arms}
\end{equation}

\begin{theorem}[Orthogonal-arm dependency compiler]
\label{thm:s6v59-dependency}
Suppose there are \(\gamma>0\) and a nonnegative matrix
\(C=(c_{ij})\) such that, for every \(f\in{\cal H}_{\rm rel}\) and every
\(i\),
\begin{equation}
 \sqrt{\gamma}\,\|D_i f\|
 \le\|R_iTf\|+\sum_{j=1}^Nc_{ij}\|D_jf\|.
 \label{eq:s6v59-local-loss}
\end{equation}
If \(c:=\|C\|_{\ell^2\to\ell^2}<\sqrt{\gamma}\), then
\begin{equation}
 \boxed{
 {\cal D}(f)\ge
 (\sqrt{\gamma}-c)^2\|f\|^2
 \qquad(f\in{\cal H}_{\rm rel}).}
 \label{eq:s6v59-global-loss}
\end{equation}
\end{theorem}

\begin{proof}
Set
\[
 a_i=\|D_if\|,\qquad b_i=\|R_iTf\|.
\]
The componentwise hypothesis implies
\[
 \sqrt{\gamma}\|a\|_{\ell^2}
 \le\|b\|_{\ell^2}+\|Ca\|_{\ell^2}
 \le\|Tf\|+c\|a\|_{\ell^2},
\]
where \eqref{eq:s6v59-output-arms} gives the second inequality.  By
orthogonality in \eqref{eq:s6v59-input-arms},
\(\|a\|_{\ell^2}=\|f\|\).  Rearranging and squaring proves
\eqref{eq:s6v59-global-loss}.
\end{proof}

\begin{corollary}[The critical dependency scale]
\label{cor:s6v59-critical-scale}
Suppose the local channel loss and dependency norm obey
\begin{align}
 \gamma_j&=\mu_0\delta_j+o(\delta_j),
 \label{eq:s6v59-local-critical}\\
 c_j&=\kappa\sqrt{\delta_j}+o(\sqrt{\delta_j})
 \label{eq:s6v59-dependency-critical}
\end{align}
with \(0\le\kappa<\sqrt{\mu_0}\).  Then
\begin{equation}
 { {\cal D}_j(f)\over\delta_j}
 \ge\left[
  (\sqrt{\mu_0}-\kappa)^2+o(1)
 \right]\|f\|^2.
 \label{eq:s6v59-physical-gap}
\end{equation}
A fixed dependency norm \(c_j\to c_*>0\) cannot satisfy the compiler when
\(\gamma_j\asymp\delta_j\to0\).
\end{corollary}

\begin{proof}
Insert \eqref{eq:s6v59-local-critical} and
\eqref{eq:s6v59-dependency-critical} into
\eqref{eq:s6v59-global-loss}, use
\(\sqrt{\gamma_j}=\sqrt{\mu_0}\sqrt{\delta_j}
+o(\sqrt{\delta_j})\), and divide by \(\delta_j\).
The last assertion follows because eventually
\(c_j>\sqrt{\gamma_j}\).
\end{proof}

\begin{remark}[Why the square root appears]
\label{rem:s6v59-square-root}
The matrix \(C\) in \eqref{eq:s6v59-local-loss} compares
\emph{amplitudes}, while \(\gamma\) and \({\cal D}\) are squared losses.
Thus the natural critical debit is \(O(\sqrt{\delta})\).  Squaring a
dimensionless influence coefficient only after subtracting it from a local
loss amplitude is essential; replacing this by a fixed Dobrushin margin
changes the physical scaling.
\end{remark}

\subsection{Nonorthogonal output arms and frame loss}
\label{sec:s6v59-frame}

\begin{proposition}[Output-frame correction]
\label{prop:s6v59-frame}
If \eqref{eq:s6v59-output-arms} is replaced by
\[
 \sum_i\|R_i z\|^2\le B\|z\|^2
\]
for \(B\ge1\), then the conclusion becomes
\begin{equation}
 {\cal D}(f)\ge
 {(\sqrt{\gamma}-c)^2\over B}\|f\|^2.
 \label{eq:s6v59-frame-loss}
\end{equation}
\end{proposition}

\begin{proof}
The proof of \Cref{thm:s6v59-dependency} now gives
\[
 (\sqrt{\gamma}-c)\|f\|
 \le\|b\|_{\ell^2}\le\sqrt B\|Tf\|.
\]
Square and rearrange.
\end{proof}

\begin{proposition}[Overlap multiplicity is not automatically harmless]
\label{prop:s6v59-overlap}
Suppose each physical loss coordinate is counted in at most \(m\) output
arms, and each \(R_i\) is a coordinate restriction.  Then one may take
\(B\le m\) in \Cref{prop:s6v59-frame}.  If \(m\) diverges with regulator
or volume, the compiler loses uniformity unless a sharper signed
orthogonality identity is proved.
\end{proposition}

\begin{proof}
On summing \(\|R_i z\|^2\), every squared coordinate of \(z\) occurs at
most \(m\) times.  Hence the sum is at most \(m\|z\|^2\).
\end{proof}

\subsection{A recursion form useful for Gibbs dependencies}
\label{sec:s6v59-recursion}

\begin{lemma}[Nonnegative dependency recursion]
\label{lem:s6v59-recursion}
Let \(a,b\in[0,\infty)^N\) and let \(C\) have nonnegative entries.  If
\[
 a\le b+Ca
\]
componentwise and \(\|C\|_{\ell^2\to\ell^2}=c<1\), then
\begin{equation}
 \boxed{\|a\|_{\ell^2}\le{1\over1-c}\|b\|_{\ell^2}.}
 \label{eq:s6v59-recursion}
\end{equation}
\end{lemma}

\begin{proof}
Monotonicity and the componentwise inequality give
\[
 \|a\|_2\le\|b\|_2+\|Ca\|_2
 \le\|b\|_2+c\|a\|_2.
\]
Rearrange.
\end{proof}

\begin{assessment}[What must be proved to call \(C\) a Yang--Mills
dependency matrix]
\label{ass:s6v59-C-meaning}
The entries \(c_{ij}\) must be derived from the actual Doob conditional law
\eqref{eq:s6v58-conditional-kernel}, after its ground singular factor and
common-sector quotient are included.  A matrix of action-Hessian
couplings, one-link heat-bath sensitivities, or spatial shell influences is
not automatically the matrix in
\eqref{eq:s6v59-local-loss}.  The proof must construct the input
martingale arms \(D_i\), the observed loss arms \(R_i\), and the displayed
inequality on the same physical half-slab carrier.
\end{assessment}

\subsection{Testing the existing MSEC route}
\label{sec:s6v59-MSEC}

\begin{theorem}[Fixed one-site influence cannot load the critical
compiler]
\label{thm:s6v59-fixed-influence}
Suppose a proposed one-site dependency matrix satisfies
\[
 \liminf_j\|C_j\|_{\ell^2\to\ell^2}\ge c_*>0
\]
while the local physical loss is
\(\gamma_j=\mu_0\delta_j+o(\delta_j)\) with
\(\delta_j\to0\).  Then the strict hypothesis of
\Cref{thm:s6v59-dependency} fails for all sufficiently large \(j\).
\end{theorem}

\begin{proof}
\(\sqrt{\gamma_j}\to0\), whereas
\(\|C_j\|\ge c_*/2\) eventually.  Thus
\(\|C_j\|<\sqrt{\gamma_j}\) is impossible.
\end{proof}

\begin{assessment}[MSEC verdict]
\label{ass:s6v59-MSEC}
V49 proves that the linearized one-site influence row of the Ward-critical
Gaussian owner approaches a row with norm at least one.  V45--V47 provide
useful spatial and interaction-tail bounds but no
\(O(\sqrt{\delta_j})\) temporal dependency matrix on the Doob loss arms.
Therefore the fixed one-site MSEC route cannot populate
\eqref{eq:s6v59-local-loss} in the shrinking-time regime.  Reapplying it
with a fixed margin would reproduce the ultralocal clock error already
excluded in V50.
\end{assessment}

\begin{corollary}[A multiscale split is necessary for this route]
\label{cor:s6v59-multiscale}
Within the dependency-compiler strategy, the Ward-critical long modes must
be moved out of the one-site fast block and into an explicit coarse
carrier.  Only the remaining fast arms may be assigned a strict local
loss; the coarse carrier then requires its own temporal form estimate.
\end{corollary}

\begin{proof}
The one-site matrix fails by
\Cref{thm:s6v59-fixed-influence,ass:s6v59-MSEC}.  The V56 exact
coarse/fine decomposition is the available way to isolate the modes
responsible for that failure without discarding them.
\end{proof}

\subsection{Critical conditional tensorization card}
\label{sec:s6v59-card}

\begin{definition}[Critical half-slab dependency card, CHDC]
\label{def:s6v59-CHDC}
At each regulator, CHDC consists of:
\begin{enumerate}[label=\textup{(\roman*)},leftmargin=3.5em]
\item the exact Doob half-slab loss operator
      \(T_j=(I-K_j^*K_j)^{1/2}\) on the protected relative carrier;
\item orthogonal fast input arms \(D_{i,j}\) and loss analysis arms
      \(R_{i,j}\) with a volume-uniform frame bound \(B\);
\item the local loss/dependency inequality
      \eqref{eq:s6v59-local-loss};
\item \(\gamma_j/\delta_j\to\mu_0>0\) and
      \(\|C_j\|/\sqrt{\delta_j}\to\kappa<\sqrt{\mu_0}\);
\item an explicit coarse carrier containing all Ward-critical infrared
      modes excluded from the fast arms;
\item a coarse temporal form gap and fine/coarse Schur comparison;
\item adjacent projectivity for the arms, ground singular functions, and
      protected projectors; and
\item Mosco/OS convergence and source-core cyclicity.
\end{enumerate}
\end{definition}

\begin{theorem}[Conditional CHDC handoff]
\label{thm:s6v59-CHDC}
If CHDC is populated and the coarse form gap is \(\mu_c>0\), then the
fast half-slab form has asymptotic gap at least
\[
 {(\sqrt{\mu_0}-\kappa)^2\over B}.
\]
Together with a V56/V57 fine--coarse form domination, the limiting
Hamiltonian gap is bounded below by the resulting positive minimum of this
constant and \(\mu_c\).
\end{theorem}

\begin{proof}
\Cref{thm:s6v59-dependency,prop:s6v59-frame,
cor:s6v59-critical-scale} give the fast constant.  Rows~\textup{(v)--(vi)}
and the V56/V57 compilers combine it with the coarse gap.  Rows
\textup{(vii)--(viii)} permit the V53 limiting handoff.
\end{proof}

\begin{assessment}[Exact frontier after V59]
\label{ass:s6v59-frontier}
Conditional tensorization does not close the mass gap with the existing
one-site MSEC stock: its dependency margin has the wrong critical scaling,
and the Ward-critical infrared tower must remain in a coarse carrier.
The proof has nevertheless isolated a quantitative viable route.  A fast
block needs an \(O(\delta)\) local prediction loss with only
\(O(\sqrt\delta)\) inter-arm amplitude debit, while the coarse
centre-neutral carrier needs an independent positive rescaled transfer
form.  V60 must audit the parallel programmes before deciding how to attack
that coarse form.
\end{assessment}

\begin{programme}[V60: compulsory audit and coarse-carrier decision]
\label{prog:s6v59-V60}
\begin{enumerate}[leftmargin=3em]
\item Audit the then-current SM and NS notes for a multiscale
      slow/fast split, infrared block estimate, or form-Schur comparison.
\item Import only exact operator mechanisms, not companion physical
      constants.
\item Reassess whether the Yang--Mills sources contain a coarse
      centre-neutral temporal observable family complete enough for a
      Rayleigh bound.
\item If not, formulate the minimal infrared trial-space exclusion theorem
      needed to prove a positive coarse form.
\item Continue with a Fourier-block or wavelet decomposition rather than
      retrying a forbidden one-site margin.
\end{enumerate}
\end{programme}

\begin{firewall}[Claim boundary after sixth-series V59]
\label{fire:s6v59-boundary}
V59 proves product Efron--Stein tensorization; proves an orthogonal-arm
half-slab dependency compiler with frame correction; derives the exact
\(O(\sqrt\delta)\) critical dependency scale; proves a nonnegative
dependency recursion; and shows that the existing fixed one-site influence
route cannot load the shrinking-time compiler.

V59 does not construct the Yang--Mills loss arms or dependency matrix,
prove their critical scaling, establish the fast or coarse temporal form
gap, close projectivity/Mosco convergence, construct the continuum OS
theory, or prove the full Yang--Mills mass gap.
\end{firewall}

\begin{theorem}[Sixth-series V59 result ledger]
\label{thm:s6v59-results}
The positive conclusions in \Cref{fire:s6v59-boundary} are exactly those of
\Cref{thm:s6v59-Efron-Stein,
thm:s6v59-dependency,
cor:s6v59-critical-scale,
prop:s6v59-frame,
prop:s6v59-overlap,
lem:s6v59-recursion,
thm:s6v59-fixed-influence,
cor:s6v59-multiscale,
thm:s6v59-CHDC}.
\end{theorem}

\begin{proof}
Each ledger item is the cited theorem with its exact arm-incidence, frame,
critical-scaling, coarse-carrier, and convergence hypotheses retained.
\end{proof}

\paragraph{Parent-version record.}
V59 was formed from
\texttt{Renewal\_Yang\_Mills\_Mass\_Gap\_Research\_Notes\_V58.tex}.
The V58 parent had \(26\,075\) logical source lines,
\(916\,290\) bytes, and SHA--256
\[
 \texttt{150575545276900e419df8b87fffc4af09df40f0a5c0e5659b5b372c6f4d4602}.
\]
The next version performs the compulsory companion audit.

% =============================================================================
% END APPENDED SIXTH-SERIES V59 MATERIAL
% =============================================================================

% =============================================================================
% BEGIN APPENDED SIXTH-SERIES V60 MATERIAL
% =============================================================================

\section{V60 compulsory companion audit: finite-horizon locality and the
physical singlet infrared carrier}
\label{sec:s6v60}

\subsection{Audited companion snapshot}
\label{sec:s6v60-audit}

\begin{definition}[V60 companion files]
\label{def:s6v60-files}
The current active companion notes at this audit are:
\begin{center}
\begin{tabular}{p{0.43\linewidth}rr}
\toprule
file & logical lines & bytes\\
\midrule
\texttt{Renewal\_SM\_Einstein\_Continuum\_Research\_Notes\_V7.tex}
& \(2\,537\) & \(96\,410\)\\
\texttt{Renewal\_NS\_Stability\_Research\_Notes\_V26.tex}
& \(5\,319\) & \(278\,283\)\\
\bottomrule
\end{tabular}
\end{center}
Their SHA--256 digests are respectively
\[
 \texttt{ec93b11558e7300f728b76b976abbacbbd2a0386b081d87607676ef4e0652362}
\]
and
\[
 \texttt{82c887f8c2c69e4f28fad7c3dc8de5b9099cb5a4bf431eba1fff3e91935978ad}.
\]
\end{definition}

\begin{assessment}[SM V7 audit]
\label{ass:s6v60-SM}
Since V55, the SM programme has:
\begin{enumerate}[label=\textup{(\roman*)},leftmargin=3.5em]
\item computed the exact mean-zero periodic lattice-Laplacian gap,
      which tends to \(4\pi^2/L^2\) at fixed physical side \(L\) and
      therefore closes when \(L\to\infty\);
\item shown that a fixed contraction inferred from this gap requires block
      time of order \(L^2\);
\item proved finite-time heat-kernel tails that can make remote refinement
      sources summable without a global gap;
\item separated local algebraic state/projectivity landing from uniqueness,
      exponential clustering, and a transfer gap; and
\item proved bounded-generator near/far Duhamel comparisons and their
      differentiated two-source form, while retaining an explicit
      unbounded-form-domain firewall.
\end{enumerate}
These results validate a local finite-horizon route for projectivity but
supply no positive coarse Yang--Mills mass form.
\end{assessment}

\begin{assessment}[NS V26 audit]
\label{ass:s6v60-NS}
NS V26 remains byte-for-byte unchanged.  As in V55, its Dirichlet
functionals and finite-dimensional norm lessons do not provide a positive
reflected Yang--Mills transfer, a physical time incidence, or a
centre-neutral coarse form estimate.
\end{assessment}

\begin{theorem}[V60 audit decision]
\label{thm:s6v60-audit-decision}
The Yang--Mills programme imports from SM V7 only:
\begin{enumerate}[label=\textup{(\roman*)},leftmargin=3.5em]
\item the exact \(L^{-2}\) infrared block-time warning;
\item finite-horizon heat locality as a possible projectivity tool that is
      logically independent of a mass gap; and
\item the requirement of form rather than bounded-operator Duhamel theory
      for unbounded limiting generators.
\end{enumerate}
No companion spectral constant or field estimate is imported.
\end{theorem}

\begin{proof}
The three items are abstract consequences of the companion theorems.  By
\Cref{ass:s6v60-SM,ass:s6v60-NS}, neither companion constructs or estimates
the response-matched Yang--Mills half-slab coarse form.
\end{proof}

\subsection{Local projectivity may proceed without the mass}
\label{sec:s6v60-locality}

\begin{theorem}[Finite-horizon projectivity is gap independent]
\label{thm:s6v60-local-projectivity}
Let \(A\) be a local observable.  Suppose adjacent regulator defects
decompose as
\[
 |\omega_{j+1}(A)-\omega_j(A)|
 \le n_j(A)+r_j(A),
\]
where \(\sum_jn_j(A)<\infty\), and the remote term obeys
\begin{equation}
 r_j(A)\le
 C_A\varepsilon_j
 \exp\left\{-{R_j(A)^2\over C T_j}\right\}
 \label{eq:s6v60-heat-tail}
\end{equation}
for finite physical horizons \(T_j\).  If the right side is summable for
every member of a dense local algebra, the expectations define a limiting
algebraic state on that local algebra.  No volume-uniform Hamiltonian gap
is required for this conclusion.
\end{theorem}

\begin{proof}
The assumed series makes \(\omega_j(A)\) Cauchy for every observable in
the dense local algebra.  Positivity and normalization pass to pointwise
limits, giving a state; boundedness extends it to the norm closure.  The
argument uses neither long-time decay nor a lower spectral bound.
\end{proof}

\begin{firewall}[State existence is not the V58 form inequality]
\label{fire:s6v60-state-not-gap}
\Cref{thm:s6v60-local-projectivity} can help populate DCFC projectivity and
common-refinement rows.  It does not imply uniqueness, clustering, or the
conditional-variance lower bound
\eqref{eq:s6v58-half-coercivity}.  The latter remains a separate
infinite-volume temporal theorem.
\end{firewall}

\subsection{The coloured one-particle caveat}
\label{sec:s6v60-colour}

Let \(G=\SU(3)\) and let \(V_{\rm ad}\) be its adjoint representation.
For any nonzero spatial momentum \(k\), a linearized transverse coloured
mode has one-particle carrier
\[
 {\cal K}_k\otimes V_{\rm ad}.
\]

\begin{proposition}[A single adjoint mode is not globally gauge invariant]
\label{prop:s6v60-one-colour}
The global-colour invariant subspace of \(V_{\rm ad}\) is zero.  Hence the
Haar gauge projection annihilates every one-particle adjoint state:
\begin{equation}
 \int_{\SU(3)}\operatorname{Ad}(g)v\,\dd g=0
 \qquad(v\in V_{\rm ad}).
 \label{eq:s6v60-adjoint-average}
\end{equation}
\end{proposition}

\begin{proof}
The Haar average is the orthogonal projection onto invariant vectors.  An
invariant vector in the adjoint representation commutes with every element
of \(\mathfrak{su}(3)\), hence lies in the centre of the Lie algebra.
That centre is zero.
\end{proof}

\begin{assessment}[Correction to an overly quick infrared trial]
\label{ass:s6v60-one-particle}
A single transverse colour Fourier mode cannot by itself be used as a
Rayleigh trial vector in the gauge-invariant OS Hilbert space.  Every
argument in V53 and V57 was conditional on the infrared vectors belonging
to the physical relative carrier; \Cref{prop:s6v60-one-colour} shows that
this condition fails for the naive one-particle adjoint vector.  The
massless-baseline test must be performed on gauge-singlet composites.
\end{assessment}

\subsection{The Gaussian singlet sector is still gapless}
\label{sec:s6v60-singlet}

Let \({\cal K}\) be a one-particle momentum-polarization space with
nonnegative dispersion \(\omega\), and consider the bosonic Fock space
over \({\cal K}\otimes V_{\rm ad}\).  Let
\[
 H_0=d\Gamma(\omega)
\]
be the free Hamiltonian.  The Killing form gives a nonzero invariant tensor
\[
 \Xi={1\over\sqrt8}\sum_{a=1}^8 e_a\otimes e_a
 \in(V_{\rm ad}\otimes V_{\rm ad})^{\SU(3)}.
\]

\begin{theorem}[Gauge-singlet two-particle infrared sequence]
\label{thm:s6v60-singlet-sequence}
Assume there are normalized one-particle wave packets
\(\varphi_n,\psi_n\in{\cal K}\) with disjoint supports and
\[
 \langle\varphi_n,\omega\varphi_n\rangle
 +\langle\psi_n,\omega\psi_n\rangle\longrightarrow0.
\]
Then the global-colour-invariant vacuum-orthogonal Fock sector contains
normalized two-particle vectors \(\Psi_n\) such that
\begin{equation}
 \langle\Psi_n,H_0\Psi_n\rangle\longrightarrow0.
 \label{eq:s6v60-singlet-energy}
\end{equation}
Therefore the massless Gaussian Hamiltonian has no positive gap even after
global \(\SU(3)\) singlet projection.
\end{theorem}

\begin{proof}
Let \(a_a^*(\varphi)\) create the one-particle vector
\(\varphi\otimes e_a\), and set
\[
 \Psi_n={1\over\sqrt8}\sum_{a=1}^8
 a_a^*(\varphi_n)a_a^*(\psi_n)\Omega.
\]
Disjoint supports make the summands orthonormal and ensure
\(\|\Psi_n\|=1\).  The contraction of the two adjoint indices by the
Killing tensor is invariant under global colour, so \(\Psi_n\) is a
singlet.  It is a two-particle vector and hence orthogonal to the vacuum.
Second quantization gives
\[
 \langle\Psi_n,H_0\Psi_n\rangle
 =\langle\varphi_n,\omega\varphi_n\rangle
 +\langle\psi_n,\omega\psi_n\rangle\to0.
\]
The variational principle proves the last statement.
\end{proof}

\begin{corollary}[Finite-torus scale]
\label{cor:s6v60-torus}
For a massless dispersion
\(\omega(k)=|k|+o(|k|)\) on a three-torus of side \(L\), choose opposite
lowest nonzero momenta.  The singlet two-particle energy is
\begin{equation}
 {4\pi\over L}+o(L^{-1}),
 \label{eq:s6v60-torus-energy}
\end{equation}
and therefore closes as \(L\to\infty\).
\end{corollary}

\begin{proof}
Each lowest nonzero momentum has magnitude \(2\pi/L\).  The two-particle
energy is their sum.
\end{proof}

\begin{remark}[Local gauge invariance remains a separate incidence]
\label{rem:s6v60-local-gauge}
Global-colour singlet projection is necessary but weaker than the full
Gauss-law/local-gauge physical projection.  In the linearized continuum
theory, local gauge-invariant quadratic fields such as colour-contracted
field-strength bilinears have the same two-massless-quantum threshold.
To use a particular lattice sequence in the interacting OS carrier, one
must construct it from gauge-invariant Wilson or field-strength writers
and prove noncollapse of its norm.  V60 does not assume that incidence.
\end{remark}

\begin{theorem}[Singlet-relative perturbations still cannot open a gap]
\label{thm:s6v60-singlet-relative}
Let \({\cal E}_0\) be the Gaussian form on the singlet carrier and let
\(\Psi_n\) be the sequence of
\Cref{thm:s6v60-singlet-sequence}.  If a nonnegative exact form
\({\cal E}={\cal E}_0+{\cal R}\) satisfies
\[
 |{\cal R}(\Psi)|\le a{\cal E}_0(\Psi)
\]
on the singlet form domain, then \({\cal E}\) has no positive singlet
gap.
\end{theorem}

\begin{proof}
This is \Cref{thm:s6v53-relative-no-gap} applied within the invariant
closed subspace, using
\({\cal E}_0(\Psi_n)\to0\).
\end{proof}

\begin{corollary}[Physical nonlinear remainder target]
\label{cor:s6v60-physical-remainder}
If the exact gauge-invariant Yang--Mills coarse form has gap
\(\mu>0\) and a noncollapsing lattice/OS realization of the singlet
sequence satisfies
\({\cal E}_{0,j}(\Psi_j)\to0\), then necessarily
\begin{equation}
 \liminf_j
 \bigl[{\cal E}_j(\Psi_j)-{\cal E}_{0,j}(\Psi_j)\bigr]\ge\mu.
 \label{eq:s6v60-physical-remainder}
\end{equation}
\end{corollary}

\begin{proof}
Apply \Cref{thm:s6v53-coercivity-alternative} to the physical singlet
sequence.
\end{proof}

\subsection{The correct coarse observable family}
\label{sec:s6v60-writers}

\begin{definition}[Gauge-singlet infrared writer card, GSIWC]
\label{def:s6v60-GSIWC}
At regulator \(j\), GSIWC requires:
\begin{enumerate}[label=\textup{(\roman*)},leftmargin=3.5em]
\item gauge-invariant, centre-neutral positive-time writers
      \(W_{j,\varphi}\) built from small-momentum field-strength bilinears
      or controlled small Wilson-loop combinations;
\item subtraction of the vacuum component and proof that their OS norms
      do not collapse;
\item a Fourier/wave-packet map showing that the perfect quadratic owner
      gives their baseline energy tending to zero in the infinite-volume
      limit;
\item inclusion of these writers in the exact Doob conditional-form and
      projective source core;
\item a lower bound of the form
      \eqref{eq:s6v60-physical-remainder} for the exact signed nonlinear
      owner; and
\item uniform control of all other gauge-invariant relative vectors, not
      only this diagnostic family, to conclude a full gap.
\end{enumerate}
\end{definition}

\begin{assessment}[Why a diagnostic sequence does not prove the gap]
\label{ass:s6v60-diagnostic}
GSIWC rows~\textup{(i)--(iv)} would prove that the Gaussian baseline has a
physical singlet infrared sequence.  Row~\textup{(v)} is necessary for a
mass but only on that sequence.  A mass-gap lower bound is universal over
the vacuum-orthogonal gauge-invariant carrier, so row~\textup{(vi)} remains
essential.  Demonstrating that one candidate glueball writer is massive
does not exclude a different lower-energy sequence.
\end{assessment}

\subsection{Audit-adjusted direction}
\label{sec:s6v60-direction}

\begin{theorem}[Separated existence/gap workstreams]
\label{thm:s6v60-separated}
The following two conclusions are logically compatible and require
different estimates:
\begin{enumerate}[label=\textup{(\roman*)},leftmargin=3.5em]
\item local regulator projectivity may be proved by finite-horizon heat or
      Lieb--Robinson tails as in
      \eqref{eq:s6v60-heat-tail}, even while the Gaussian infrared carrier
      is massless;
\item a Yang--Mills mass gap requires the global conditional-variance or
      coarse Schur lower bound on the complete gauge-invariant carrier,
      including the singlet sequence of
      \Cref{thm:s6v60-singlet-sequence}.
\end{enumerate}
Neither conclusion implies the other.
\end{theorem}

\begin{proof}
Item~\textup{(i)} follows by the Cauchy argument of
\Cref{thm:s6v60-local-projectivity} and contains no large-time estimate.
Item~\textup{(ii)} is the variational characterization of the Hamiltonian
gap and is challenged by
\Cref{thm:s6v60-singlet-sequence}.  A state may exist without being unique
or gapped, while a spectral inequality on an unidentified nonprojective
family does not construct a continuum state.
\end{proof}

\begin{assessment}[Exact frontier after V60]
\label{ass:s6v60-frontier}
The compulsory audit changes the coarse-carrier attack in two ways.
Projectivity should use finite-horizon locality and should no longer wait
for a global gap.  The gap calculation must abandon a single coloured
Fourier trial but does not escape the massless baseline: a gauge-singlet
two-particle infrared sequence remains.  The next rigorous task is to
construct its lattice gauge-invariant writer realization and compute its
OS norm/energy under the perfect quadratic half-slab, before asking the
signed nonlinear owner for the order-one lift
\eqref{eq:s6v60-physical-remainder}.
\end{assessment}

\begin{programme}[V61: lattice gauge-singlet infrared writers]
\label{prog:s6v60-V61}
\begin{enumerate}[leftmargin=3em]
\item Linearize a small plaquette or clover field-strength writer and
      form a colour-contracted two-writer wave packet at opposite momenta.
\item Verify exact lattice gauge invariance by parallel transport to a
      common basepoint or by a closed Wilson network.
\item Compute the quadratic OS two-point Gram and show noncollapse after
      normalization.
\item Calculate its Gaussian transfer Rayleigh quotient and recover the
      two-particle energy \(2\omega(k)\).
\item Identify the signed nonlinear matrix element whose uniform positive
      lower bound would lift the sequence.
\end{enumerate}
\end{programme}

\begin{firewall}[Claim boundary after sixth-series V60]
\label{fire:s6v60-boundary}
V60 performs the compulsory SM/NS audit; separates finite-horizon
projectivity from a global transfer gap; proves that a single adjoint mode
is removed by global gauge projection; constructs a gauge-singlet
two-particle infrared sequence showing that the Gaussian singlet sector is
still gapless; and formulates the physical nonlinear remainder and GSIWC
targets.

V60 does not construct the lattice gauge-invariant realization of the
singlet sequence, prove its noncollapsing OS norm, estimate the exact
nonlinear remainder, establish a uniform global form gap, close
projectivity/Mosco convergence, construct the continuum OS theory, or
prove the full Yang--Mills mass gap.
\end{firewall}

\begin{theorem}[Sixth-series V60 result ledger]
\label{thm:s6v60-results}
The positive conclusions in \Cref{fire:s6v60-boundary} are exactly those of
\Cref{def:s6v60-files,
thm:s6v60-audit-decision,
thm:s6v60-local-projectivity,
prop:s6v60-one-colour,
thm:s6v60-singlet-sequence,
cor:s6v60-torus,
thm:s6v60-singlet-relative,
cor:s6v60-physical-remainder,
thm:s6v60-separated}.
\end{theorem}

\begin{proof}
Each ledger item is the cited result with its finite-horizon, global-colour,
physical-incidence, noncollapse, and carrier-completeness qualifications
unchanged.
\end{proof}

\paragraph{Parent-version record.}
V60 was formed from
\texttt{Renewal\_Yang\_Mills\_Mass\_Gap\_Research\_Notes\_V59.tex}.
The V59 parent had \(26\,484\) logical source lines,
\(930\,310\) bytes, and SHA--256
\[
 \texttt{1966a6ba2e46161d7d78ec5a0fe76035d2d9734830533be6ee6de3f0857d37cf}.
\]
The next compulsory companion audit is V65.

% =============================================================================
% END APPENDED SIXTH-SERIES V60 MATERIAL
% =============================================================================

% =============================================================================
% BEGIN APPENDED SIXTH-SERIES V61 MATERIAL
% =============================================================================

\section{V61: exact lattice gauge-singlet infrared writers and their
quadratic transfer energy}
\label{sec:s6v61}

\subsection{A bounded adjoint plaquette field}
\label{sec:s6v61-plaquette}

For an oriented plaquette \(p\) based at a vertex \(x\), let \(U_p\) be
its holonomy and define
\begin{equation}
 {\mathfrak F}_p
 :={U_p-U_p^*\over2i}
 -{1\over3}\operatorname{Tr}
   \left({U_p-U_p^*\over2i}\right)I_3.
 \label{eq:s6v61-Fp}
\end{equation}
This is a bounded traceless Hermitian matrix.

\begin{lemma}[Exact gauge covariance]
\label{lem:s6v61-covariance}
Under a lattice gauge transformation \(g\),
\begin{equation}
 {\mathfrak F}_p\longmapsto
 g_x{\mathfrak F}_pg_x^{-1}.
 \label{eq:s6v61-F-transform}
\end{equation}
\end{lemma}

\begin{proof}
The based plaquette holonomy transforms as
\(U_p\mapsto g_xU_pg_x^{-1}\).  Taking the adjoint, anti-Hermitian part,
trace, and traceless projection commutes with unitary conjugation.
\end{proof}

Let \(p\) and \(q\) be spatial plaquettes on one positive-time slice,
based at \(x\) and \(y\), and let \(\gamma:x\to y\) be a spatial lattice
path.  With \(U_\gamma\) its path holonomy, define the bilocal Wilson-network
writer
\begin{equation}
 {\cal O}_{p,q;\gamma}
 :=\operatorname{Tr}
 \left({\mathfrak F}_pU_\gamma
       {\mathfrak F}_qU_\gamma^*\right).
 \label{eq:s6v61-network}
\end{equation}

\begin{theorem}[Exact gauge invariance and centre neutrality]
\label{thm:s6v61-network-invariance}
The writer \({\cal O}_{p,q;\gamma}\) is real, bounded, and invariant under
every lattice gauge transformation.  If \(p,q,\gamma\) lie in a
contractible spatial region, it is neutral under the one-form centre
action.
\end{theorem}

\begin{proof}
The path transforms as
\(U_\gamma\mapsto g_xU_\gamma g_y^{-1}\).  Using
\eqref{eq:s6v61-F-transform}, all \(g_y\) factors cancel internally and
the matrix inside the trace is conjugated by \(g_x\); its trace is
unchanged.  Both
\({\mathfrak F}_p\) and
\(U_\gamma{\mathfrak F}_qU_\gamma^*\) are Hermitian, so the trace of their
product is real.  Boundedness follows from compactness of \(\SU(3)\).
Under a one-form centre action, contractible plaquette holonomies are
neutral and the centre phases of \(U_\gamma\) and \(U_\gamma^*\) cancel.
\end{proof}

\begin{remark}[Path dependence is physical data]
\label{rem:s6v61-path}
Different connecting paths give different bounded gauge-invariant writers
at finite lattice spacing.  A continuum comparison must either fix a
translation-covariant path rule, average with declared weights, or prove
path-independence up to a controlled curvature insertion.  Gauge
invariance does not remove this choice.
\end{remark}

\subsection{Quadratic colour-singlet linearization}
\label{sec:s6v61-linearization}

Let a smooth weak field be inserted through
\[
 U_e(\varepsilon)=\exp\{i\varepsilon aA_e\}.
\]
For a plaquette in the \(ij\)-plane, write
\({\cal F}_{ij}^{\rm lin}(x)\) for the linearized continuum curvature.

\begin{theorem}[Second-order network jet]
\label{thm:s6v61-network-jet}
At the trivial connection,
\begin{align}
 {\mathfrak F}_{p_{ij}(x)}(\varepsilon)
 &=\varepsilon a^2{\cal F}_{ij}^{\rm lin}(x)+O(\varepsilon^2),
 \label{eq:s6v61-F-jet}\\
 {\cal O}_{p_{ij}(x),p_{ij}(y);\gamma}(\varepsilon)
 &=\varepsilon^2a^4
 \operatorname{Tr}\bigl(
 {\cal F}_{ij}^{\rm lin}(x)
 {\cal F}_{ij}^{\rm lin}(y)\bigr)
 +O(\varepsilon^3).
 \label{eq:s6v61-network-jet}
\end{align}
For a fixed finite path, the remainder is uniform on bounded weak-field
sets.  The quadratic coefficient is independent of the connecting path.
\end{theorem}

\begin{proof}
The ordered plaquette product has the standard Taylor expansion
\[
 U_{p_{ij}(x)}(\varepsilon)
 =I+i\varepsilon a^2{\cal F}_{ij}^{\rm lin}(x)
 +O(\varepsilon^2).
\]
Taking \eqref{eq:s6v61-Fp} yields
\eqref{eq:s6v61-F-jet}.  Also
\(U_\gamma(\varepsilon)=I+O(\varepsilon)\).
Substitution in \eqref{eq:s6v61-network} shows that the two plaquette
fields supply the entire order-\(\varepsilon^2\) term; every path
correction has at least one additional power of \(\varepsilon\).
\end{proof}

\begin{corollary}[Exact singlet colour contraction]
\label{cor:s6v61-colour}
Writing
\({\cal F}^{\rm lin}=F^aT^a\) in an orthonormal Lie-algebra basis, the
quadratic jet in \eqref{eq:s6v61-network-jet} is a positive constant
multiple of
\[
 \sum_{a=1}^8F^a_{ij}(x)F^a_{ij}(y).
\]
It is precisely the invariant adjoint two-particle contraction used in
V60.
\end{corollary}

\begin{proof}
Use
\(\operatorname{Tr}(T^aT^b)=c_{\mathfrak g}\delta^{ab}\) with
\(c_{\mathfrak g}>0\).
\end{proof}

\subsection{A momentum-resolved closed network}
\label{sec:s6v61-momentum}

On a finite spatial torus \(\Lambda\), fix a translation-covariant
nonwinding path rule \(\gamma_{x,r}:x\to x+r\).  For a lattice momentum
\(k\), define
\begin{equation}
 {\cal W}_{k}^{(ij)}
 :={1\over|\Lambda|}
 \sum_{x,r\in\Lambda}e^{ik\cdot r}
 {\cal O}_{p_{ij}(x),p_{ij}(x+r);\gamma_{x,r}}.
 \label{eq:s6v61-Wk}
\end{equation}
Replace it by its real part if \(k\) is not identified with \(-k\).

\begin{theorem}[Quadratic Fourier selection]
\label{thm:s6v61-Fourier}
At quadratic order around the trivial connection,
\begin{equation}
 {\cal W}_{k}^{(ij)}
 =\varepsilon^2a^4c_{\mathfrak g}
 \sum_{a=1}^8
 \widehat F_{ij}^a(k)\widehat F_{ij}^a(-k)
 +O(\varepsilon^3),
 \label{eq:s6v61-Fourier}
\end{equation}
up to the fixed Fourier normalization.  Thus the exact gauge-invariant
network has the desired opposite-momentum colour-singlet two-particle
component.
\end{theorem}

\begin{proof}
Insert \eqref{eq:s6v61-network-jet} into
\eqref{eq:s6v61-Wk}.  With
\[
 F(x)=|\Lambda|^{-1/2}\sum_p e^{ip\cdot x}\widehat F(p),
\]
the sum over \(x\) imposes \(p+q=0\), and the sum over \(r\) with
\(e^{ik\cdot r}\) imposes \(q=-k\), hence \(p=k\).
\end{proof}

\begin{assessment}[Nonlocality of exact momentum]
\label{ass:s6v61-nonlocal}
The exact plane-wave writer \eqref{eq:s6v61-Wk} spans the whole spatial
volume and uses long connecting paths.  It is suitable for a finite-volume
Rayleigh test but is not one fixed local observable for projective
continuum landing.  The correct order is to use finite-radius wave packets,
take the regulator limit at fixed radius, and only then let the radius
tend to infinity.
\end{assessment}

\subsection{Local wave-packet exhaustion}
\label{sec:s6v61-wavepacket}

\begin{theorem}[Local singlet packets with vanishing Gaussian energy]
\label{thm:s6v61-local-packets}
Let the quadratic physical theory have transverse one-particle dispersion
\(\omega(p)=|p|+o(|p|)\) near zero.  There exist normalized,
vacuum-orthogonal, global-colour-singlet two-particle vectors
\(\Psi_R\), generated by gauge-singlet bilocal writers supported in a
spatial ball of radius \(O(R)\), such that
\begin{equation}
 \langle\Psi_R,H_0\Psi_R\rangle=O(R^{-1}).
 \label{eq:s6v61-local-energy}
\end{equation}
\end{theorem}

\begin{proof}
Choose two orthogonal smooth compactly supported spatial profiles
\(\phi,\psi\), and scale
\[
 \phi_R(x)=R^{-3/2}\phi(x/R),\qquad
 \psi_R(x)=R^{-3/2}\psi(x/R).
\]
Their Fourier transforms are
\(R^{3/2}\widehat\phi(Rp)\) and
\(R^{3/2}\widehat\psi(Rp)\).  Hence
\[
 \langle\phi_R,\omega\phi_R\rangle
 =R^{-1}\int|\widehat\phi(q)|^2|q|\,\dd q+o(R^{-1}),
\]
and similarly for \(\psi_R\).  Contract their adjoint colour indices by
the invariant Killing tensor as in
\Cref{thm:s6v60-singlet-sequence}.  The resulting normalized two-particle
state is a singlet, is orthogonal to the vacuum, and has energy equal to
the sum of the two one-particle expectations, proving
\eqref{eq:s6v61-local-energy}.  In position space the bilinear is
supported where the two compact profiles are supported; connecting their
field-strength insertions by paths inside the enclosing ball gives the
exact lattice writer \eqref{eq:s6v61-network}.
\end{proof}

\begin{remark}[Field-strength coefficient and normalization]
\label{rem:s6v61-normalization}
The field-strength insertion contributes momentum factors to the raw
quadratic writer, so its unnormalized Gaussian norm may shrink as
\(R\to\infty\).  The Rayleigh vector is normalized after the writer acts
on the vacuum.  To transport this sequence into the interacting continuum,
one must prove that this normalization is compatible with the projective
OS norms and does not collapse to the zero vector.
\end{remark}

\subsection{Exact quadratic transfer quotient}
\label{sec:s6v61-transfer}

\begin{theorem}[Gaussian two-particle transfer edge]
\label{thm:s6v61-Gaussian-transfer}
Assume the perfect quadratic half-slab reconstructs the free transverse
Fock Hamiltonian \(H_0=d\Gamma(\omega_a)\).  For a nonzero momentum
\(k\) and the normalized singlet state selected by
\eqref{eq:s6v61-Fourier},
\begin{align}
 \langle\Psi_k,e^{-\delta H_0}\Psi_k\rangle
 &=e^{-2\delta\omega_a(k)},
 \label{eq:s6v61-transfer-edge}\\
 {1\over\delta}
 \langle\Psi_k,(I-e^{-\delta H_0})\Psi_k\rangle
 &={1-e^{-2\delta\omega_a(k)}\over\delta}
 \longrightarrow2\omega_a(k)
 \quad(\delta\downarrow0).
 \label{eq:s6v61-form-edge}
\end{align}
On a torus, this quotient tends to zero as the side length tends to
infinity.
\end{theorem}

\begin{proof}
The state contains two quanta of momenta \(k,-k\), so it is an eigenvector
of \(d\Gamma(\omega_a)\) with eigenvalue
\(\omega_a(k)+\omega_a(-k)=2\omega_a(k)\).
Functional calculus proves both identities.  The first nonzero momentum
has magnitude \(2\pi/L\), and the massless dispersion tends to zero with
it.
\end{proof}

\begin{corollary}[Wave-packet Rayleigh closure]
\label{cor:s6v61-wavepacket-transfer}
For the local packets of \Cref{thm:s6v61-local-packets},
\[
 \limsup_{\delta\downarrow0}
 {1\over\delta}
 \langle\Psi_R,(I-e^{-\delta H_0})\Psi_R\rangle
 =O(R^{-1}).
\]
Thus the gauge-invariant local exhaustion has no Gaussian gap.
\end{corollary}

\begin{proof}
Apply the spectral small-time identity
\Cref{thm:s6v58-spectral-limit} and
\eqref{eq:s6v61-local-energy}.
\end{proof}

\subsection{The exact interacting lift that remains}
\label{sec:s6v61-lift}

Let \({\cal E}_j\) be the exact response-matched rescaled transfer form and
\({\cal E}_{0,j}\) its perfect quadratic baseline on a common
gauge-invariant writer realization.  For a normalized centred network
writer state \(\widehat{\cal W}_{j,R}\), define
\begin{equation}
 \Delta_{j,R}
 :={\cal E}_j(\widehat{\cal W}_{j,R})
   -{\cal E}_{0,j}(\widehat{\cal W}_{j,R}).
 \label{eq:s6v61-lift}
\end{equation}

\begin{theorem}[Necessary singlet-network lift]
\label{thm:s6v61-necessary-lift}
If the exact continuum Yang--Mills form has gap at least \(\mu>0\), the
network states land without norm collapse, and the regulator limit is
taken first at fixed \(R\), then
\begin{equation}
 \boxed{
 \liminf_{R\to\infty}\liminf_{j\to\infty}
 \Delta_{j,R}\ge\mu.}
 \label{eq:s6v61-lift-lower}
\end{equation}
\end{theorem}

\begin{proof}
The exact gap gives
\({\cal E}_j(\widehat{\cal W}_{j,R})\ge\mu-o_j(1)\) after form landing.
By \Cref{cor:s6v61-wavepacket-transfer}, the quadratic energy tends to
zero as \(R\to\infty\).  Subtract it in
\eqref{eq:s6v61-lift} and take the limits in the stated order.
\end{proof}

\begin{assessment}[No available positive lift estimate]
\label{ass:s6v61-lift-status}
The exact Wilson-network writer and its Gaussian energy are now explicit.
The accumulated signed Wilsonian head gives formal derivatives of the
action owner, while CMAC and PFS control local charts and exceptional
geometry.  None currently proves the positive order-one lower bound
\eqref{eq:s6v61-lift-lower}.  Bounds making
\(\Delta_{j,R}\) vanish uniformly with the coupling would instead preserve
the massless singlet sequence and refute that proof route.
\end{assessment}

\begin{definition}[Singlet-network lift card, SNLC]
\label{def:s6v61-SNLC}
SNLC consists of:
\begin{enumerate}[label=\textup{(\roman*)},leftmargin=3.5em]
\item the exact network family \eqref{eq:s6v61-network} with a coherent
      path and smearing rule;
\item positive-time support, centering, and noncollapsing projective OS
      normalization at every fixed packet radius;
\item identification of the quadratic limit and
      \eqref{eq:s6v61-local-energy};
\item a signed exact-owner representation of \(\Delta_{j,R}\);
\item the uniform lower lift \eqref{eq:s6v61-lift-lower};
\item a universal coercivity extension from the diagnostic network family
      to the complete gauge-invariant vacuum-orthogonal carrier; and
\item DCFC/RODFC projectivity, Mosco, and OS reconstruction.
\end{enumerate}
\end{definition}

\begin{assessment}[Exact frontier after V61]
\label{ass:s6v61-frontier}
Gauge incidence no longer blocks the diagnostic infrared sequence:
\eqref{eq:s6v61-network} is an exact bounded gauge-invariant,
centre-neutral lattice writer whose quadratic jet is the colour-singlet
two-particle mode.  The new sharp bottleneck is the signed order-one
nonlinear lift \eqref{eq:s6v61-lift-lower}, followed by its extension from
this family to every physical state.  This is the point at which dimensional
transmutation must enter; ultraviolet norm-smallness alone cannot supply
it.
\end{assessment}

\begin{programme}[V62: signed lift representation]
\label{prog:s6v61-V62}
\begin{enumerate}[leftmargin=3em]
\item Express \(\Delta_{j,R}\) through an interpolation between the perfect
      quadratic and exact response-matched half-slab amplitudes.
\item Use the V52 singular-cluster score formula at positive slab width,
      or a form Duhamel formula, to retain the sign of every integrated
      Wilsonian return.
\item Separate local ultraviolet terms from the radius-dependent
      connected surface/network contribution.
\item Prove which terms vanish on the centred singlet writer and isolate
      the remaining candidate for an order-one lift.
\item If only absolute smallness is available, record the resulting
      no-lift theorem and move upstream to the connected network
      coefficient rather than claim a mass.
\end{enumerate}
\end{programme}

\begin{firewall}[Claim boundary after sixth-series V61]
\label{fire:s6v61-boundary}
V61 constructs an exact bounded gauge-invariant centre-neutral bilocal
plaquette network; computes its quadratic colour-singlet jet and
momentum selection; constructs a local wave-packet exhaustion with
Gaussian energy \(O(R^{-1})\); computes its exact quadratic transfer
Rayleigh quotient; and states the necessary interacting lift.

V61 does not prove projective OS noncollapse of the network writers,
estimate the signed nonlinear lift, extend coercivity to the complete
physical carrier, close Mosco/projectivity, construct the continuum OS
theory, or prove the full Yang--Mills mass gap.
\end{firewall}

\begin{theorem}[Sixth-series V61 result ledger]
\label{thm:s6v61-results}
The positive conclusions in \Cref{fire:s6v61-boundary} are exactly those of
\Cref{lem:s6v61-covariance,
thm:s6v61-network-invariance,
thm:s6v61-network-jet,
cor:s6v61-colour,
thm:s6v61-Fourier,
thm:s6v61-local-packets,
thm:s6v61-Gaussian-transfer,
cor:s6v61-wavepacket-transfer,
thm:s6v61-necessary-lift}.
\end{theorem}

\begin{proof}
Each result is cited with its weak-field, path, quadratic-owner,
normalization, order-of-limits, and projective-incidence hypotheses
unchanged.
\end{proof}

\paragraph{Parent-version record.}
V61 was formed from
\texttt{Renewal\_Yang\_Mills\_Mass\_Gap\_Research\_Notes\_V60.tex}.
The V60 parent had \(26\,908\) logical source lines,
\(946\,730\) bytes, and SHA--256
\[
 \texttt{6900a79a7744bd1580f67af6d8c3eefad46e98acc90147f8e61d3dd8441c4ed6}.
\]
The next compulsory companion audit is V65.

% =============================================================================
% END APPENDED SIXTH-SERIES V61 MATERIAL
% =============================================================================

% =============================================================================
% BEGIN APPENDED SIXTH-SERIES V62 MATERIAL
% =============================================================================

\section{V62: signed interpolation of the Doob prediction loss and the
ground-state response debit}
\label{sec:s6v62}

\subsection{Differentiating the conditional-variance numerator}
\label{sec:s6v62-envelope}

Let \(\Lambda_\theta\), \(0\le\theta\le1\), be a \(C^1\) family of
probability laws on \({\mathsf X}\times{\mathsf Y}\), all equivalent to a
common reference law.  Write
\begin{equation}
 S_\theta(x,y)
 :=\partial_\theta\log{d\Lambda_\theta\over d\Lambda_*}(x,y),
 \qquad
 \mathbb E_{\Lambda_\theta}S_\theta=0.
 \label{eq:s6v62-pair-score}
\end{equation}
For a fixed bounded real writer \(f=f(x)\), define
\begin{align}
 h_\theta(y)&:=\mathbb E_{\Lambda_\theta}[f(X)\mid Y=y],
 \label{eq:s6v62-predictor}\\
 r_\theta(x,y)&:=f(x)-h_\theta(y),
 \label{eq:s6v62-residual}\\
 N_\theta(f)&:=\mathbb E_{\Lambda_\theta}r_\theta^2
 =\inf_h\mathbb E_{\Lambda_\theta}|f(X)-h(Y)|^2.
 \label{eq:s6v62-N}
\end{align}
Assume the scores and derivatives are dominated sufficiently to
differentiate these quantities.

\begin{theorem}[Conditional-variance envelope formula]
\label{thm:s6v62-envelope}
For every such fixed writer,
\begin{equation}
 \boxed{
 {d\over d\theta}N_\theta(f)
 =\mathbb E_{\Lambda_\theta}
  [S_\theta r_\theta^2]
 =\operatorname{Cov}_{\Lambda_\theta}
  (S_\theta,r_\theta^2).}
 \label{eq:s6v62-N-prime}
\end{equation}
No derivative of the optimal predictor \(h_\theta\) remains.
\end{theorem}

\begin{proof}
Differentiate
\(N_\theta=\mathbb E_\theta(f-h_\theta)^2\):
\[
 N_\theta'
 =\mathbb E_\theta[S_\theta r_\theta^2]
 -2\mathbb E_\theta[r_\theta h_\theta'(Y)].
\]
By conditional expectation,
\(\mathbb E_\theta[r_\theta\mid Y]=0\), so the second term vanishes.
The score is centered by normalization of \(\Lambda_\theta\), which gives
the covariance form.
\end{proof}

Let \(\pi_\theta\) be the \(X\)-marginal, with score
\[
 s_\theta^X(x)
 =\partial_\theta\log{d\pi_\theta\over d\pi_*}(x),
 \qquad\mathbb E_{\pi_\theta}s_\theta^X=0.
\]
Put
\begin{equation}
 \bar f_\theta=\mathbb E_{\pi_\theta}f,\qquad
 V_\theta(f)=\operatorname{Var}_{\pi_\theta}(f).
 \label{eq:s6v62-V}
\end{equation}

\begin{lemma}[Marginal variance envelope formula]
\label{lem:s6v62-V-prime}
\begin{equation}
 \boxed{
 V_\theta'(f)
 =\mathbb E_{\pi_\theta}
 [s_\theta^X(f-\bar f_\theta)^2].}
 \label{eq:s6v62-V-prime}
\end{equation}
\end{lemma}

\begin{proof}
Differentiate
\(\mathbb E_{\pi_\theta}(f-c_\theta)^2\) at the minimizing constant
\(c_\theta=\bar f_\theta\).  The derivative of \(c_\theta\) is multiplied
by \(\mathbb E(f-\bar f_\theta)=0\), leaving the displayed score term.
\end{proof}

\begin{theorem}[Exact Rayleigh-loss interpolation]
\label{thm:s6v62-Rayleigh}
If \(V_\theta(f)>0\), define the dimensionless transfer loss
\[
 \Gamma_\theta(f):={N_\theta(f)\over V_\theta(f)}.
\]
Then
\begin{equation}
 \boxed{
 \Gamma_\theta'(f)
 ={V_\theta\,
   \mathbb E_{\Lambda_\theta}[S_\theta r_\theta^2]
  -N_\theta\,
   \mathbb E_{\pi_\theta}
       [s_\theta^X(f-\bar f_\theta)^2]
  \over V_\theta^2}.}
 \label{eq:s6v62-Gamma-prime}
\end{equation}
Consequently
\begin{equation}
 {\Gamma_1(f)-\Gamma_0(f)\over\delta}
 ={1\over\delta}\int_0^1\Gamma_\theta'(f)\,\dd\theta
 \label{eq:s6v62-lift-integral}
\end{equation}
is the exact rescaled nonlinear lift of the writer's transfer Rayleigh
quotient.
\end{theorem}

\begin{proof}
Apply the quotient rule to \(N_\theta/V_\theta\) and use
\Cref{thm:s6v62-envelope,lem:s6v62-V-prime}.  Integrate the derivative
from zero to one and divide by the physical full-slab duration.
\end{proof}

\subsection{The score of an interpolated half-slab action}
\label{sec:s6v62-action-score}

Let
\begin{equation}
 F_\theta(y,x)
 =\exp\left\{-{1\over2}
 \bigl(W_0(y,x)+\theta R(y,x)\bigr)\right\}
 \label{eq:s6v62-amplitude}
\end{equation}
relative to fixed reference boundary measures.  Let
\(s_{1,\theta},u_\theta,v_\theta\) be its positive normalized leading
singular data, and let \(\Lambda_\theta\) be its Doob law.

\begin{theorem}[Complete Doob pair score]
\label{thm:s6v62-Doob-score}
Choose the real normalization gauge
\(\langle u_\theta,u_\theta'\rangle
=\langle v_\theta,v_\theta'\rangle=0\).  Then
\begin{equation}
 \boxed{
 S_\theta(x,y)
 =-{1\over2}R(y,x)
  +{u_\theta'(x)\over u_\theta(x)}
  +{v_\theta'(y)\over v_\theta(y)}
  +{1\over2}\mathbb E_{\Lambda_\theta}R.}
 \label{eq:s6v62-complete-score}
\end{equation}
The \(X\)-marginal score is
\begin{equation}
 \boxed{s_\theta^X(x)=2{u_\theta'(x)\over u_\theta(x)}.}
 \label{eq:s6v62-X-score}
\end{equation}
\end{theorem}

\begin{proof}
Differentiate the logarithm of
\[
 d\Lambda_\theta
 =s_{1,\theta}^{-1}u_\theta(x)F_\theta(y,x)v_\theta(y)
 \,\mu_0(\dd x)\mu_{1/2}(\dd y).
\]
The V52 leading singular-value score formula gives
\[
 {s_{1,\theta}'\over s_{1,\theta}}
 =-{1\over2}\mathbb E_{\Lambda_\theta}R.
\]
Substitution proves \eqref{eq:s6v62-complete-score}.  Since
\(\pi_\theta(\dd x)=u_\theta(x)^2\mu_0(\dd x)\), differentiation gives
\eqref{eq:s6v62-X-score}; its mean is zero because
\(\|u_\theta\|_2=1\).
\end{proof}

\begin{corollary}[The action score is not the Doob score]
\label{cor:s6v62-not-action-only}
The signed Wilsonian insertion \(-R/2\) is only one term in the exact
Rayleigh response \eqref{eq:s6v62-Gamma-prime}.  The left and right ground
singular responses, together with the marginal normalization term, are
equally part of the physical transfer incidence.
\end{corollary}

\begin{proof}
Insert \eqref{eq:s6v62-complete-score} and
\eqref{eq:s6v62-X-score} into
\eqref{eq:s6v62-Gamma-prime}.
\end{proof}

\subsection{The inverse-margin ground response}
\label{sec:s6v62-ground}

Let
\[
 B_\theta=A_\theta^*A_\theta,\qquad
 \lambda_{1,\theta}=s_{1,\theta}^2,
\]
and let \(\lambda_{2,\theta}\) be the next eigenvalue on the relative
carrier.  Put
\[
 G_\theta=\lambda_{1,\theta}-\lambda_{2,\theta}>0,\qquad
 Q_\theta=I-|u_\theta\rangle\langle u_\theta|.
\]

\begin{theorem}[Leading singular-vector response]
\label{thm:s6v62-ground-response}
In the normalization gauge of
\Cref{thm:s6v62-Doob-score},
\begin{equation}
 \boxed{
 u_\theta'
 =(\lambda_{1,\theta}-B_\theta)^{-1}_{Q_\theta}
   Q_\theta B_\theta'u_\theta,}
 \label{eq:s6v62-u-prime}
\end{equation}
and therefore
\begin{equation}
 \boxed{
 \|u_\theta'\|_2
 \le{\|B_\theta'\|\over G_\theta}.}
 \label{eq:s6v62-u-bound}
\end{equation}
Moreover the marginal score has the exact Fisher norm
\begin{equation}
 \boxed{
 \|s_\theta^X\|_{L^2(\pi_\theta)}
 =2\|u_\theta'\|_{L^2(\mu_0)}.}
 \label{eq:s6v62-score-Fisher}
\end{equation}
The analogous statements hold for the left singular vector using
\(A_\theta A_\theta^*\).
\end{theorem}

\begin{proof}
Differentiate
\(B_\theta u_\theta=\lambda_{1,\theta}u_\theta\):
\[
 (B_\theta-\lambda_{1,\theta})u_\theta'
 =-(B_\theta'-\lambda_{1,\theta}')u_\theta.
\]
Apply \(Q_\theta\); the scalar eigenvalue-derivative term vanishes, and
\(u_\theta'\in\operatorname{Ran}Q_\theta\) by the normalization gauge.
Inversion on that range proves \eqref{eq:s6v62-u-prime}.  The inverse norm
is \(G_\theta^{-1}\), proving \eqref{eq:s6v62-u-bound}.  Finally,
\[
 \int\left|2{u_\theta'\over u_\theta}\right|^2
 u_\theta^2\,\dd\mu_0=4\|u_\theta'\|_2^2.
\]
\end{proof}

\begin{corollary}[Near-critical ground-state debit]
\label{cor:s6v62-ground-debit}
For a normalized transfer with
\[
 1-{\lambda_{2,\theta}\over\lambda_{1,\theta}}
 \asymp\mu\delta,
\]
the absolute spectral separation is
\(G_\theta\asymp s_{1,\theta}^2\mu\delta\).
Thus a direct perturbative bound on the ground score generally contains an
inverse factor \(\delta^{-1}\).  To make that score small by
\eqref{eq:s6v62-u-bound}, one would need
\[
 \|B_\theta'\|=o(s_{1,\theta}^2\delta),
\]
unless a sharper cancellation with the action and left-ground terms is
proved.
\end{corollary}

\begin{proof}
Factor
\[
 G_\theta=\lambda_{1,\theta}
 \left(1-{\lambda_{2,\theta}\over\lambda_{1,\theta}}\right)
\]
and apply \eqref{eq:s6v62-u-bound}.
\end{proof}

\begin{assessment}[A new circularity firewall]
\label{ass:s6v62-circularity}
Bounding the ground singular response by its reduced resolvent uses the
very near-critical relative gap that the programme is trying to prove.
This formula is valid as a conditional stability estimate but cannot serve
as an independent proof of the gap.  A viable interpolation must either
obtain the complete score covariance without separately estimating
\(u_\theta'\) by \(G_\theta^{-1}\), or establish a different uniform
ground-state regularity theorem from positivity and locality.
\end{assessment}

\subsection{A signed no-lift criterion}
\label{sec:s6v62-no-lift}

\begin{lemma}[Rayleigh response bound]
\label{lem:s6v62-response-bound}
Assume \(V_\theta(f)\ge v_*>0\).  Since
\(0\le N_\theta(f)\le V_\theta(f)\),
\begin{equation}
 |\Gamma_\theta'(f)|
 \le{1\over v_*}
 \left(
 \left|\mathbb E_{\Lambda_\theta}
 [S_\theta r_\theta^2]\right|
 +\left|\mathbb E_{\pi_\theta}
 [s_\theta^X(f-\bar f_\theta)^2]\right|
 \right).
 \label{eq:s6v62-response-bound}
\end{equation}
\end{lemma}

\begin{proof}
Take absolute values in \eqref{eq:s6v62-Gamma-prime} and use
\(N_\theta/V_\theta\le1\) and \(V_\theta\ge v_*\).
\end{proof}

\begin{theorem}[No lift from sub-clock score covariances]
\label{thm:s6v62-no-lift}
Let \(f_j\) be a diagnostic writer family with
\(\inf_{\theta,j}V_{\theta,j}(f_j)\ge v_*>0\).  Suppose
\[
 \int_0^1\left(
 \left|\mathbb E_{\Lambda_{\theta,j}}
 [S_{\theta,j}r_{\theta,j}^2]\right|
 +\left|\mathbb E_{\pi_{\theta,j}}
 [s_{\theta,j}^X(f_j-\bar f_{\theta,j})^2]\right|
 \right)\dd\theta
 =o(\delta_j).
 \label{eq:s6v62-subclock-score}
\]
Then
\[
 {\Gamma_{1,j}(f_j)-\Gamma_{0,j}(f_j)\over\delta_j}
 \longrightarrow0.
\]
If the baseline rescaled Rayleigh quotient tends to zero, the exact forms
cannot have a uniform positive gap on this writer sequence.
\end{theorem}

\begin{proof}
Integrate \eqref{eq:s6v62-response-bound}, divide by \(\delta_j\), and use
\eqref{eq:s6v62-subclock-score}.  Add the vanishing baseline quotient and
apply the variational principle.
\end{proof}

\begin{corollary}[Necessary signed score stock]
\label{cor:s6v62-necessary-score}
For the V61 singlet-network packets to acquire a limiting lift at least
\(\mu\), the signed integral in
\eqref{eq:s6v62-lift-integral} must have lower limit at least \(\mu\).
It is insufficient to show that every absolute score covariance is small;
such a result at the \(o(\delta)\) scale would instead prove
\Cref{thm:s6v62-no-lift}.
\end{corollary}

\begin{proof}
This is the contrapositive of
\Cref{thm:s6v62-no-lift} together with the necessary lift
\eqref{eq:s6v61-lift-lower}.
\end{proof}

\subsection{Connected localization of the score response}
\label{sec:s6v62-connected}

Suppose the complete pair score has an \(L^2\)-convergent decomposition
\[
 S_\theta=\sum_ZS_{\theta,Z},\qquad
 \mathbb E_\theta S_{\theta,Z}=0.
\]

\begin{lemma}[Exact connected score decomposition]
\label{lem:s6v62-connected}
\begin{equation}
 N_\theta'(f)
 =\sum_Z\operatorname{Cov}_{\Lambda_\theta}
 (S_{\theta,Z},r_\theta^2).
 \label{eq:s6v62-connected}
\end{equation}
If \(S_{\theta,Z}\) is independent of \(r_\theta^2\), its summand vanishes
exactly.
\end{lemma}

\begin{proof}
Linearity and \(L^2\) convergence permit the sum to pass through the
covariance in \eqref{eq:s6v62-N-prime}.  Independence and centering make
the corresponding covariance zero.
\end{proof}

\begin{assessment}[What locality can and cannot do]
\label{ass:s6v62-locality}
Finite-horizon locality may show that score components far from a
fixed-radius network have summable covariance with its squared prediction
residual.  This can close projectivity and reduce the lift to a connected
neighbourhood.  It does not determine the sign of the remaining connected
sum or make it order one after division by \(\delta\).  That connected
signed term is the precise location at which a nonperturbative mass
mechanism must be proved.
\end{assessment}

\subsection{Complete signed lift card}
\label{sec:s6v62-card}

\begin{definition}[Doob score lift card, DSLC]
\label{def:s6v62-DSLC}
For the V61 network packets, DSLC consists of:
\begin{enumerate}[label=\textup{(\roman*)},leftmargin=3.5em]
\item a \(C^1\) path of positive response-matched half-slab amplitudes from
      the perfect quadratic owner to the exact owner;
\item the complete Doob pair and marginal scores
      \eqref{eq:s6v62-complete-score}--\eqref{eq:s6v62-X-score};
\item uniform noncollapse \(V_{\theta,j}(f_{j,R})\ge v_R>0\) at every
      fixed radius;
\item a connected localization of every action and ground-score term;
\item a noncircular estimate of the combined score covariance, not a
      separate reduced-resolvent bound assuming the desired gap;
\item the positive signed lift
      \[
       \liminf_{R\to\infty}\liminf_{j\to\infty}
       {1\over\delta_j}\int_0^1
       \Gamma_{\theta,j}'(f_{j,R})\,\dd\theta\ge\mu;
      \]
\item a universal extension from the network diagnostics to the complete
      relative carrier; and
\item projective Mosco/OS landing of the same form path.
\end{enumerate}
\end{definition}

\begin{assessment}[Exact frontier after V62]
\label{ass:s6v62-frontier}
The nonlinear lift is now represented by an exact signed covariance
integral.  This exposes an omission in any action-only calculation: the
Doob ground singular functions and the changing marginal normalization
contribute at the same order, and naive control of them is amplified by
the inverse near-critical margin.  The next upstream task is to avoid that
circularity by exploiting ground-state transformed operator identities or
by proving a direct joint-law comparison for the complete score covariance.
\end{assessment}

\begin{programme}[V63: ground-transform cancellation]
\label{prog:s6v62-V63}
\begin{enumerate}[leftmargin=3em]
\item Rewrite the Rayleigh quotient on the fixed reference boundary space
      using \(A_\theta^*A_\theta/s_{1,\theta}^2\), so no pointwise
      derivative of \(u_\theta\) appears in the operator derivative.
\item Track the moving vacuum projection by a constrained min--max
      principle rather than a reduced-resolvent estimate.
\item Derive one-sided cluster formulas if the singlet edge is degenerate.
\item Compare the resulting operator insertion with the signed Wilsonian
      action score.
\item Decide whether the ground-response terms cancel on centred
      gauge-invariant network packets or remain an independent acquisition.
\end{enumerate}
\end{programme}

\begin{firewall}[Claim boundary after sixth-series V62]
\label{fire:s6v62-boundary}
V62 proves exact envelope derivatives for the Doob conditional-variance
numerator, marginal variance, and their Rayleigh quotient; derives the
complete interpolated Doob score; proves the leading singular-vector
response and its inverse-margin Fisher cost; proves a signed sub-clock
no-lift theorem; and localizes the remaining lift into connected score
covariances.

V62 does not control the complete Yang--Mills Doob score, remove the
inverse-margin circularity, prove a positive network lift or its universal
coercive extension, close projectivity/Mosco convergence, construct the
continuum OS theory, or prove the full Yang--Mills mass gap.
\end{firewall}

\begin{theorem}[Sixth-series V62 result ledger]
\label{thm:s6v62-results}
The positive conclusions in \Cref{fire:s6v62-boundary} are exactly those of
\Cref{thm:s6v62-envelope,
lem:s6v62-V-prime,
thm:s6v62-Rayleigh,
thm:s6v62-Doob-score,
cor:s6v62-not-action-only,
thm:s6v62-ground-response,
cor:s6v62-ground-debit,
lem:s6v62-response-bound,
thm:s6v62-no-lift,
cor:s6v62-necessary-score,
lem:s6v62-connected}.
\end{theorem}

\begin{proof}
Each ledger assertion is the cited result with its differentiability,
noncollapse, complete-score, spectral-separation, locality, and
order-of-limits hypotheses retained.
\end{proof}

\paragraph{Parent-version record.}
V62 was formed from
\texttt{Renewal\_Yang\_Mills\_Mass\_Gap\_Research\_Notes\_V61.tex}.
The V61 parent had \(27\,346\) logical source lines,
\(962\,682\) bytes, and SHA--256
\[
 \texttt{40ca8e0f9aad41eb60e65c07df231f9450211c4a524fe958542f005767bada58}.
\]
The next compulsory companion audit is V65.

% =============================================================================
% END APPENDED SIXTH-SERIES V62 MATERIAL
% =============================================================================

% =============================================================================
% BEGIN APPENDED SIXTH-SERIES V63 MATERIAL
% =============================================================================

\section{V63: fixed-reference ground transform and signed spectral-edge
separation}
\label{sec:s6v63}

\subsection{The normalized transfer on a fixed Hilbert space}
\label{sec:s6v63-fixed}

Retain the norm-\(C^1\) compact half-slab family
\[
 A_\theta:{\cal H}_0\to{\cal H}_{1/2},\qquad
 B_\theta=A_\theta^*A_\theta.
\]
Let \(\lambda_{1,\theta}>0\) be the simple leading eigenvalue, with
normalized positive eigenvector \(u_\theta\), and define
\begin{equation}
 T_\theta={B_\theta\over\lambda_{1,\theta}},
 \qquad
 L_\theta=I-T_\theta,
 \qquad
 Q_\theta=I-|u_\theta\rangle\langle u_\theta|.
 \label{eq:s6v63-normalized}
\end{equation}

\begin{theorem}[Normalized operator derivative]
\label{thm:s6v63-operator-derivative}
The normalized transfer and dimensionless form derivatives are
\begin{align}
 T_\theta'
 &={B_\theta'\over\lambda_{1,\theta}}
 -{\lambda_{1,\theta}'\over\lambda_{1,\theta}}T_\theta,
 \label{eq:s6v63-T-prime}\\
 L_\theta'
 &=-{B_\theta'\over\lambda_{1,\theta}}
 +{\lambda_{1,\theta}'\over\lambda_{1,\theta}}T_\theta.
 \label{eq:s6v63-L-prime}
\end{align}
Neither formula contains \(u_\theta'\) or \(v_\theta'\).
\end{theorem}

\begin{proof}
Differentiate \(T_\theta=\lambda_{1,\theta}^{-1}B_\theta\), and use
\(L_\theta=I-T_\theta\).
\end{proof}

\begin{theorem}[Constrained min--max formula]
\label{thm:s6v63-minmax}
Let
\[
 q_\theta=\|T_\theta|_{Q_\theta{\cal H}_0}\|,
 \qquad g_\theta=1-q_\theta.
\]
Then
\begin{equation}
 \boxed{
 g_\theta
 =\inf_{\substack{f\perp u_\theta\\f\ne0}}
 {\langle f,(\lambda_{1,\theta}I-B_\theta)f\rangle
  \over\lambda_{1,\theta}\|f\|^2}.}
 \label{eq:s6v63-gap-minmax}
\end{equation}
Consequently a physical step \(\delta\) has form gap at least \(\mu\) if
and only if
\begin{equation}
 \boxed{
 Q_\theta(\lambda_{1,\theta}I-B_\theta)Q_\theta
 \succeq
 \mu\delta\,\lambda_{1,\theta}Q_\theta.}
 \label{eq:s6v63-edge-operator}
\end{equation}
\end{theorem}

\begin{proof}
The largest spectral value of \(B_\theta\) on \(u_\theta^\perp\) is
\(\lambda_{2,\theta}=q_\theta\lambda_{1,\theta}\).  The Rayleigh--Ritz
principle gives the infimum in
\eqref{eq:s6v63-gap-minmax}.  The operator inequality is exactly the
statement that this infimum is at least \(\mu\delta\).
\end{proof}

\begin{assessment}[Where the ground derivative went]
\label{ass:s6v63-cancellation}
The moving physical marginal in V62 forces \(u_\theta'\) into the
derivative of a fixed configuration writer.  The spectral gap is instead
a constrained extremum over the entire moving subspace
\(u_\theta^\perp\).  On the fixed reference Hilbert space its value is
given by \eqref{eq:s6v63-gap-minmax}; differentiating the eigenvalue
extremum uses \(B_\theta'\) and \(\lambda_{1,\theta}'\), but not the
derivative of the optimizing eigenvector.  This removes the
inverse-margin circularity for a direct spectral-edge proof.  It does not
remove the ground/projector debit from projective transport of prescribed
physical writers.
\end{assessment}

\subsection{The exact amplitude-score insertion}
\label{sec:s6v63-score}

Assume the action interpolation
\[
 \partial_\theta F_\theta(y,x)
 =-{1\over2}R_\theta(y,x)F_\theta(y,x)
\]
and let \(A_{\theta,R}\) be the inserted integral operator
\begin{equation}
 (A_{\theta,R}f)(y)
 =\int R_\theta(y,x)F_\theta(y,x)f(x)\,\mu_0(\dd x).
 \label{eq:s6v63-inserted-A}
\end{equation}

\begin{lemma}[Inserted derivative of the reflected square]
\label{lem:s6v63-B-prime}
\begin{equation}
 A_\theta'=-{1\over2}A_{\theta,R},
 \qquad
 B_\theta'
 =-{1\over2}
 \left(A_{\theta,R}^*A_\theta
       +A_\theta^*A_{\theta,R}\right).
 \label{eq:s6v63-B-prime}
\end{equation}
For every \(f\) with \(A_\theta f\ne0\), define
\begin{equation}
 {\cal J}_\theta(f)
 :={\operatorname{Re}
 \langle A_\theta f,A_{\theta,R}f\rangle
 \over\|A_\theta f\|^2}.
 \label{eq:s6v63-Jf}
\end{equation}
Then
\begin{equation}
 \langle f,B_\theta'f\rangle
 =-{\cal J}_\theta(f)\langle f,B_\theta f\rangle.
 \label{eq:s6v63-B-quadratic}
\end{equation}
\end{lemma}

\begin{proof}
Differentiate the kernel and then \(B_\theta=A_\theta^*A_\theta\).
Taking the quadratic form of the resulting symmetrized operator gives
\[
 \langle f,B_\theta'f\rangle
 =-\operatorname{Re}
 \langle A_\theta f,A_{\theta,R}f\rangle,
\]
which is \eqref{eq:s6v63-B-quadratic}.
\end{proof}

\begin{lemma}[Leading action score]
\label{lem:s6v63-leading-score}
For the leading singular vector,
\begin{equation}
 {\cal J}_\theta(u_\theta)
 =\mathbb E_{\Lambda_\theta}R_\theta
 =:{\cal J}_{1,\theta},
 \qquad
 {\lambda_{1,\theta}'\over\lambda_{1,\theta}}
 =-{\cal J}_{1,\theta}.
 \label{eq:s6v63-leading-score}
\end{equation}
\end{lemma}

\begin{proof}
Use \(A_\theta u_\theta=s_{1,\theta}v_\theta\) in
\eqref{eq:s6v63-Jf}.  The resulting normalized integral is precisely the
Doob expectation.  The second identity follows from
\eqref{eq:s6v63-B-quadratic} and Feynman--Hellmann.
\end{proof}

\begin{theorem}[Ground-transform score cancellation]
\label{thm:s6v63-score-cancellation}
For every fixed \(f\) with \(A_\theta f\ne0\),
\begin{align}
 \langle f,T_\theta'f\rangle
 &=\bigl({\cal J}_{1,\theta}-{\cal J}_\theta(f)\bigr)
   \langle f,T_\theta f\rangle,
 \label{eq:s6v63-T-score}\\
 \langle f,L_\theta'f\rangle
 &=\bigl({\cal J}_\theta(f)-{\cal J}_{1,\theta}\bigr)
   \langle f,T_\theta f\rangle.
 \label{eq:s6v63-L-score}
\end{align}
Thus the normalized spectral response is a difference between an inserted
relative-mode score and the leading Doob score; all explicit derivatives
of the singular ground functions cancel.
\end{theorem}

\begin{proof}
Insert \eqref{eq:s6v63-B-quadratic} and
\eqref{eq:s6v63-leading-score} into
\eqref{eq:s6v63-T-prime} and negate for \(L_\theta'\).
\end{proof}

\begin{corollary}[Scalar owner terms cancel]
\label{cor:s6v63-scalar}
Adding a scalar \(c_\theta\) to \(R_\theta(y,x)\) leaves both score
differences in \eqref{eq:s6v63-T-score}--\eqref{eq:s6v63-L-score}
unchanged.
\end{corollary}

\begin{proof}
The inserted operator changes by \(c_\theta A_\theta\), so both
\({\cal J}_\theta(f)\) and \({\cal J}_{1,\theta}\) increase by
\(c_\theta\).
\end{proof}

\subsection{Degenerate relative edges and one-sided separation}
\label{sec:s6v63-cluster}

At a fixed \(\theta\), suppose the largest relative eigenvalue
\(\lambda_{2,\theta}\) is isolated as a finite-multiplicity cluster
\(E_\theta\), and choose an orthonormal basis \(u_a\) of that cluster.
Put \(s_2=\sqrt{\lambda_2}\), \(v_a=A_\theta u_a/s_2\), and define the
Hermitian inserted score matrix
\begin{equation}
 [\mathbb J_{2,\theta}]_{ab}
 :={1\over2s_2}
 \left(
 \langle v_a,A_{\theta,R}u_b\rangle
 +\overline{\langle v_b,A_{\theta,R}u_a\rangle}
 \right).
 \label{eq:s6v63-cluster-score}
\end{equation}

\begin{theorem}[One-sided dimensionless-gap derivative]
\label{thm:s6v63-cluster-derivative}
Let
\[
 q_\theta={\lambda_{2,\theta}\over\lambda_{1,\theta}},
 \qquad g_\theta=1-q_\theta.
\]
Then the right and left derivatives at the cluster point are
\begin{align}
 D_+g_\theta
 &=q_\theta\left(
 \lambda_{\min}(\mathbb J_{2,\theta})
 -{\cal J}_{1,\theta}\right),
 \label{eq:s6v63-g-plus}\\
 D_-g_\theta
 &=q_\theta\left(
 \lambda_{\max}(\mathbb J_{2,\theta})
 -{\cal J}_{1,\theta}\right).
 \label{eq:s6v63-g-minus}
\end{align}
\end{theorem}

\begin{proof}
On \(E_\theta\), direct expansion of
\eqref{eq:s6v63-B-prime} gives
\[
 {1\over\lambda_{2,\theta}}
 P_{E_\theta}B_\theta'P_{E_\theta}
 =-\mathbb J_{2,\theta}.
\]
For positive parameter increments, the derivative of the largest
eigenvalue in the cluster is the largest eigenvalue of the compressed
derivative, namely
\(-\lambda_2\lambda_{\min}(\mathbb J_2)\).  For negative increments it is
\(-\lambda_2\lambda_{\max}(\mathbb J_2)\).
Differentiate \(g=1-\lambda_2/\lambda_1\), and use
\(\lambda_1'/\lambda_1=-{\cal J}_1\).
\end{proof}

\begin{remark}[Simple-edge recovery]
\label{rem:s6v63-simple}
For a simple relative edge,
\(\mathbb J_{2,\theta}\) is the scalar
\({\cal J}_\theta(u_2)\), and
\[
 g_\theta'
 =q_\theta\bigl(
 {\cal J}_\theta(u_2)-{\cal J}_{1,\theta}\bigr).
\]
\end{remark}

\subsection{Integrated edge opening without a ground resolvent}
\label{sec:s6v63-integrated}

\begin{theorem}[Dini edge-opening compiler]
\label{thm:s6v63-Dini}
Fix the physical step \(\delta>0\).  Suppose \(g_\theta\) is continuous on
\([0,1]\), its top relative cluster is isolated at every \(\theta\), and
for one \(m>0\),
\begin{equation}
 q_\theta\left(
 \lambda_{\min}(\mathbb J_{2,\theta})
 -{\cal J}_{1,\theta}\right)
 \ge m\delta
 \qquad(0\le\theta\le1).
 \label{eq:s6v63-score-separation}
\end{equation}
Then
\begin{equation}
 \boxed{g_1\ge g_0+m\delta.}
 \label{eq:s6v63-opened-gap}
\end{equation}
In particular the exact endpoint has rescaled form gap at least \(m\) if
the baseline gap is nonnegative.
\end{theorem}

\begin{proof}
By \Cref{thm:s6v63-cluster-derivative},
\(D_+g_\theta\ge m\delta\) at every point.  Set
\(h(\theta)=g_\theta-m\delta\,\theta\).  Its lower right Dini derivative is
nonnegative everywhere.  The elementary Dini monotonicity theorem implies
that \(h\) is nondecreasing.  Hence
\(g_1-m\delta\ge g_0\).
\end{proof}

\begin{corollary}[Volume-uniform score target]
\label{cor:s6v63-uniform-target}
To open a volume-uniform mass from a massless Gaussian baseline by this
interpolation, it is sufficient to prove
\eqref{eq:s6v63-score-separation} uniformly in volume, regulator,
top-relative symmetry cluster, and interpolation parameter, after all
common factors have been removed.
\end{corollary}

\begin{proof}
Apply \Cref{thm:s6v63-Dini} at every finite regulator and then the V53
projective Mosco handoff.
\end{proof}

\begin{assessment}[What was gained and what remains]
\label{ass:s6v63-gain}
Unlike the fixed-writer formula in V62, the spectral-edge compiler requires
no bound on \(u_\theta'\) and hence no division by the unknown gap.
However, the leading Doob expectation and the complete worst relative
score matrix still depend on the global singular data of the half-slab
amplitude.  The cancellation removes a circular \emph{derivative}
estimate; it does not turn the signed score separation into a local action
Hessian bound.
\end{assessment}

\subsection{The physical score-separation card}
\label{sec:s6v63-card}

\begin{definition}[Spectral action-score separation card, SASSC]
\label{def:s6v63-SASSC}
SASSC consists of:
\begin{enumerate}[label=\textup{(\roman*)},leftmargin=3.5em]
\item a positive \(C^1\) half-slab interpolation on fixed reference
      boundary spaces, with exact physical clock \(\delta_j\);
\item maximal-common-factor removal and isolation of the complete top
      relative cluster at every interpolation point;
\item the inserted action operator \(A_{\theta,R}\) containing the full
      signed Wilsonian head, Jacobian, shell returns, and counterterms;
\item uniform leading and cluster score construction
      \({\cal J}_{1,\theta},\mathbb J_{2,\theta}\);
\item the volume- and regulator-uniform signed separation
      \eqref{eq:s6v63-score-separation};
\item adjacent transport of the endpoint forms and protected projections
      at the inverse-margin scale;
\item common-refinement Mosco convergence; and
\item noncollapse and cyclicity of the limiting gauge-invariant OS core.
\end{enumerate}
\end{definition}

\begin{theorem}[Conditional SASSC mass-gap handoff]
\label{thm:s6v63-SASSC}
If SASSC is populated with constant \(m>0\), the reconstructed
Yang--Mills Hamiltonian has relative vacuum gap at least \(m\).
\end{theorem}

\begin{proof}
Rows~\textup{(i)--(v)} and
\Cref{thm:s6v63-Dini} give a dimensionless endpoint deficit at least
\(m\delta_j\).  Rows~\textup{(vi)--(viii)} supply the V53 projective
Mosco/OS handoff.
\end{proof}

\begin{assessment}[Exact frontier after V63]
\label{ass:s6v63-frontier}
The ground-response circularity has been removed from the direct gap
calculation.  The exact remaining inequality is now
\[
 \lambda_{\min}(\mathbb J_{2,\theta})
 -{\cal J}_{1,\theta}
 \ge {m\delta\over q_\theta}.
\]
It compares the full signed action insertion in the worst
vacuum-orthogonal half-slab mode with the same insertion in the vacuum
Doob law.  No accumulated CMAC, PFS, MSEC, character, or local Hessian
estimate proves this global spectral score separation.  V64 should derive
variational lower bounds for the cluster matrix from a declared network
source bank and expose whether finite-bank control can ever be
source-complete.
\end{assessment}

\begin{programme}[V64: from network scores to the full relative cluster]
\label{prog:s6v63-V64}
\begin{enumerate}[leftmargin=3em]
\item Express \(\lambda_{\min}(\mathbb J_{2,\theta})\) by a min--max over
      the complete top relative singular cluster.
\item Compare that min--max with the V61 gauge-invariant network writer
      bank.
\item Prove the direction of every finite-bank inequality and prevent a
      lower bound on a compression from being promoted to the full
      carrier.
\item Formulate a quantitative density/frame condition under which an
      expanding writer bank controls the cluster.
\item Determine the projective and clock-scaled error required for that
      frame exhaustion.
\end{enumerate}
\end{programme}

\begin{firewall}[Claim boundary after sixth-series V63]
\label{fire:s6v63-boundary}
V63 rewrites the normalized reflected transfer and its form on fixed
reference spaces; proves the constrained min--max gap formula; derives the
exact inserted action-score operator; proves cancellation of all explicit
ground-vector derivatives; treats degenerate relative clusters with
one-sided derivatives; and gives a Dini score-separation compiler for
opening an order-\(\delta\) edge.

V63 does not prove the Yang--Mills signed score separation, isolate its
relative cluster uniformly, construct a source-complete writer frame,
close projectivity/Mosco convergence, construct the continuum OS theory,
or prove the full Yang--Mills mass gap.
\end{firewall}

\begin{theorem}[Sixth-series V63 result ledger]
\label{thm:s6v63-results}
The positive conclusions in \Cref{fire:s6v63-boundary} are exactly those of
\Cref{thm:s6v63-operator-derivative,
thm:s6v63-minmax,
lem:s6v63-B-prime,
lem:s6v63-leading-score,
thm:s6v63-score-cancellation,
cor:s6v63-scalar,
thm:s6v63-cluster-derivative,
thm:s6v63-Dini,
cor:s6v63-uniform-target,
thm:s6v63-SASSC}.
\end{theorem}

\begin{proof}
Each assertion is the cited result with its positivity, fixed-reference,
cluster-isolation, signed-insertion, common-factor, clock, and convergence
hypotheses retained.
\end{proof}

\paragraph{Parent-version record.}
V63 was formed from
\texttt{Renewal\_Yang\_Mills\_Mass\_Gap\_Research\_Notes\_V62.tex}.
The V62 parent had \(27\,855\) logical source lines,
\(979\,227\) bytes, and SHA--256
\[
 \texttt{d8b6d5269c0591f88a2f7ca89733115752b972e081ace766eafa164053b56f20}.
\]
The next compulsory companion audit is V65.

% =============================================================================
% END APPENDED SIXTH-SERIES V63 MATERIAL
% =============================================================================

% =============================================================================
% BEGIN APPENDED SIXTH-SERIES V64 MATERIAL
% =============================================================================

\section{V64: source-complete spin-network frames for the relative score
cluster}
\label{sec:s6v64}

\subsection{The compression direction}
\label{sec:s6v64-compression}

Let \(E\) be a Hilbert space and let \(J=J^*\) be a bounded score operator
on \(E\).  For a closed subspace \(W\subset E\), let \(P_W\) be its
orthogonal projection.

\begin{proposition}[A positive compression is only a partial result]
\label{prop:s6v64-compression}
If
\[
 P_WJP_W\succeq cP_W,
\]
then
\[
 \inf\operatorname{spec}J\le
 \inf\operatorname{spec}(P_WJP_W|_W).
\]
Thus a lower bound on the compression gives no lower bound on \(J\) unless
\(W\) is dense in \(E\), or an independent estimate controls
\(W^\perp\).
\end{proposition}

\begin{proof}
The Rayleigh infimum for \(J\) is taken over all unit vectors of \(E\);
the compression infimum restricts that set to unit vectors of \(W\).
Restricting an infimum can only increase it.
\end{proof}

\begin{proposition}[Two-dimensional compression counterexample]
\label{prop:s6v64-counterexample}
For every \(c>0\) and \(M>0\), take
\[
 J=\begin{pmatrix}c&0\\0&-M\end{pmatrix},
 \qquad W=\operatorname{span}\{(1,0)\}.
\]
Then \(P_WJP_W=cP_W\), while
\(\lambda_{\min}(J)=-M\).
\end{proposition}

\begin{proof}
Both claims follow by direct diagonalization.
\end{proof}

\begin{assessment}[Finite Wilson banks cannot be promoted silently]
\label{ass:s6v64-finite-bank}
The V61 singlet network family probes a physically relevant infrared
direction, but even a positive signed lift on every member of one finite
network bank would be a compression statement.  The SASSC cluster minimum
in V63 ranges over the complete top relative singular eigenspace.
Source completeness or an explicit orthogonal-complement estimate is
therefore mandatory.
\end{assessment}

\subsection{The generalized frame pencil}
\label{sec:s6v64-frame}

Let \(w_1,\ldots,w_m\in E\), and define the synthesis map
\[
 W:\mathbb C^m\to E,\qquad
 Wa=\sum_{\alpha=1}^ma_\alpha w_\alpha.
\]
Put
\begin{equation}
 G=W^*W,\qquad H=W^*JW.
 \label{eq:s6v64-GH}
\end{equation}

\begin{theorem}[Exact source-frame criterion]
\label{thm:s6v64-frame}
If \(\operatorname{Ran}W=E\), then the following are equivalent:
\begin{align}
 J&\succeq cI_E,
 \label{eq:s6v64-J-lower}\\
 H&\succeq cG
 \quad\hbox{on }\mathbb C^m.
 \label{eq:s6v64-pencil-lower}
\end{align}
Equivalently,
\begin{equation}
 \lambda_{\min}(J)
 =\inf_{\substack{a\notin\ker W}}
 {a^*Ha\over a^*Ga}.
 \label{eq:s6v64-generalized-eigenvalue}
\end{equation}
\end{theorem}

\begin{proof}
For every coefficient vector \(a\),
\[
 a^*Ha=\langle Wa,JWa\rangle,\qquad
 a^*Ga=\|Wa\|^2.
\]
Thus \eqref{eq:s6v64-J-lower} implies
\eqref{eq:s6v64-pencil-lower}.  Conversely, surjectivity writes every
\(x\in E\) as \(x=Wa\); the pencil inequality then gives
\(\langle x,Jx\rangle\ge c\|x\|^2\).  Taking Rayleigh infima proves
\eqref{eq:s6v64-generalized-eigenvalue}.
\end{proof}

\begin{corollary}[Application to the V63 score cluster]
\label{cor:s6v64-cluster}
Let \(E_\theta\) be the complete top relative cluster and set
\[
 J_\theta^{\rm sep}
 =\mathbb J_{2,\theta}-{\cal J}_{1,\theta}I_{E_\theta}.
\]
If projected physical writers \(w_\alpha=P_{E_\theta}\varphi_\alpha\)
span \(E_\theta\), then the V63 score-separation condition is equivalent to
\begin{equation}
 H_\theta^{\rm sep}
 \succeq {m\delta\over q_\theta}G_\theta,
 \label{eq:s6v64-cluster-pencil}
\end{equation}
where
\[
 G_\theta=W_\theta^*W_\theta,\qquad
 H_\theta^{\rm sep}
 =W_\theta^*J_\theta^{\rm sep}W_\theta.
\]
\end{corollary}

\begin{proof}
Apply \Cref{thm:s6v64-frame} with
\(c=m\delta/q_\theta\).
\end{proof}

\subsection{Finite-lattice spin-network density}
\label{sec:s6v64-density}

Let \({\cal E}\) and \({\cal V}\) be the finite edge and vertex sets of a
spatial lattice.  The boundary configuration space is
\[
 {\mathsf X}=G^{\cal E},\qquad G=\SU(3),
\]
and the gauge group \(G^{\cal V}\) acts by endpoint multiplication.
Let \(\mu_H\) be product Haar measure.  A spin-network function means a
finite product of irreducible matrix coefficients on edges with all vertex
indices contracted by invariant tensors.

\begin{theorem}[Density of gauge-invariant spin networks]
\label{thm:s6v64-spin-density}
The linear span \({\cal S}_{\rm spin}\) of spin-network functions is dense
in
\[
 L^2(G^{\cal E},\mu_H)^{G^{\cal V}}.
\]
If \(\pi\) is a probability law with strictly positive continuous density
with respect to \(\mu_H\), the same span is dense in
\[
 L^2(G^{\cal E},\pi)^{G^{\cal V}}.
\]
\end{theorem}

\begin{proof}
Peter--Weyl implies that finite products of irreducible matrix
coefficients are uniformly dense in \(C(G^{\cal E})\).  Haar averaging
over the compact gauge group is a norm-one projection
\[
 {\cal P}_{\rm gauge}f(U)
 =\int_{G^{\cal V}}f(g\cdot U)\,\dd g
\]
from continuous functions onto continuous gauge-invariant functions.
Applying this projection to products of matrix coefficients contracts
their vertex indices by invariant tensors and produces precisely finite
linear combinations of spin networks.  Hence spin networks are uniformly
dense in \(C(G^{\cal E})^{G^{\cal V}}\), and continuous invariant
functions are dense in invariant \(L^2(\mu_H)\).

For the second statement, strict positive continuity on compact
\(G^{\cal E}\) gives constants \(0<c\le d\pi/d\mu_H\le C<\infty\).
The two \(L^2\) norms are equivalent, so the same dense span works.
\end{proof}

\begin{corollary}[Doob-boundary density]
\label{cor:s6v64-Doob-density}
At every finite compact positivity-improving Wilson regulator, the leading
singular function \(u_0\) is strictly positive and continuous.  Therefore
the Doob boundary law
\(\pi_0=u_0^2\mu_H\) has a strictly positive continuous density, and
gauge-invariant spin networks are dense in its physical boundary
\(L^2\) space.  Orthogonal projection away from the protected common
factor preserves density in the relative carrier.
\end{corollary}

\begin{proof}
The source Jentzsch--Perron theorem gives the stated regularity of \(u_0\).
Apply \Cref{thm:s6v64-spin-density}.  The image of a dense set under a
bounded orthogonal projection is dense in the range of that projection.
\end{proof}

\begin{remark}[Topological sectors]
\label{rem:s6v64-topology}
If the finite configuration space is decomposed into protected
superselection components rather than connected by a positivity-improving
kernel, the density theorem is applied in each declared component.
It does not decide whether components are protected zero-energy sectors or
dynamically connected states.
\end{remark}

\subsection{A finite frame exists at each isolated finite cluster}
\label{sec:s6v64-finite-frame}

\begin{theorem}[Finite spin-network frame extraction]
\label{thm:s6v64-frame-extraction}
Let \(E_\theta\) be a finite-dimensional isolated relative spectral
cluster at one finite regulator.  Then there exist finitely many
gauge-invariant spin networks
\(\varphi_1,\ldots,\varphi_m\) whose projections
\[
 w_\alpha=P_{E_\theta}\varphi_\alpha
\]
span \(E_\theta\).  Hence the exact cluster lower bound can be tested by
the finite generalized pencil of
\Cref{cor:s6v64-cluster}.
\end{theorem}

\begin{proof}
By \Cref{cor:s6v64-Doob-density}, projected spin networks are dense in
\(E_\theta\).  A dense linear subspace of a finite-dimensional space is
the whole space.  Choose a basis from finitely many projected elements and
apply \Cref{thm:s6v64-frame}.
\end{proof}

\begin{assessment}[Existence is not a cofinal construction]
\label{ass:s6v64-nonuniform}
\Cref{thm:s6v64-frame-extraction} is a finite-regulator existence theorem.
It gives neither a canonical list of networks nor a frame lower bound
uniform in volume, regulator, interpolation parameter, or symmetry-cluster
crossing.  The necessary network size and path range can diverge.  Those
are precisely the uniform data required before the finite generalized
pencil becomes a continuum proof.
\end{assessment}

\subsection{Frame conditioning and the physical clock}
\label{sec:s6v64-conditioning}

Let \(W:\mathbb C^m\to E\) be surjective and restrict coefficient space to
\((\ker W)^\perp\).  Assume
\begin{equation}
 G=W^*W\succeq aI,\qquad a>0.
 \label{eq:s6v64-frame-lower}
\end{equation}

\begin{lemma}[Absolute matrix error becomes a generalized error]
\label{lem:s6v64-frame-error}
If an approximation \(\widehat H\) satisfies
\(\|H-\widehat H\|\le\varepsilon\), then
\begin{equation}
 \left|
 {z^*Hz\over z^*Gz}
 -{z^*\widehat Hz\over z^*Gz}
 \right|
 \le{\varepsilon\over a}
 \qquad(z\in(\ker W)^\perp\setminus\{0\}).
 \label{eq:s6v64-generalized-error}
\end{equation}
If
\[
 \widehat H\succeq(c+\varepsilon/a)G,
\]
then \(J\succeq cI_E\).
\end{lemma}

\begin{proof}
The numerator error is at most
\(\varepsilon\|z\|^2\), while
\(z^*Gz\ge a\|z\|^2\), proving the ratio estimate.  The last statement
follows from \Cref{thm:s6v64-frame}.
\end{proof}

\begin{corollary}[Clock-scaled frame accuracy]
\label{cor:s6v64-clock-error}
To certify the score separation
\(c_j=m\delta_j/q_j\) from an absolute matrix approximation, it is
sufficient that
\begin{equation}
 {\varepsilon_j\over a_j}=o(\delta_j)
 \label{eq:s6v64-clock-frame}
\end{equation}
and that the approximate generalized eigenvalue exceeds
\((m+o(1))\delta_j/q_j\).  If \(a_j\to0\), an \(o(\delta_j)\) absolute
matrix error is not enough by itself.
\end{corollary}

\begin{proof}
Apply \Cref{lem:s6v64-frame-error} with
\(c_j=m\delta_j/q_j\).
\end{proof}

\begin{proposition}[Nested dense banks can extend a common lower bound]
\label{prop:s6v64-dense-extension}
Let \(W_1\subset W_2\subset\cdots\) be subspaces with dense union in a
Hilbert space \(E\), and let \(J=J^*\) be bounded.  If
\[
 \langle w,Jw\rangle\ge c\|w\|^2
 \qquad(w\in W_n,\ n\ge1)
\]
with one common \(c\), then \(J\succeq cI_E\).
\end{proposition}

\begin{proof}
Approximate any \(x\in E\) in norm by \(w_n\) from the union.  Boundedness
of \(J\) makes its quadratic form norm-continuous, so the inequality passes
to \(x\).
\end{proof}

\begin{assessment}[Pointwise bank convergence is weaker]
\label{ass:s6v64-pointwise}
It is not enough that every fixed network writer eventually has a positive
score.  The common lower constant must hold on every linear combination in
an expanding dense bank, including cancellations whose coefficients can
grow as the frame lower bound collapses.  The generalized pencil and
\eqref{eq:s6v64-clock-frame} record this conditioning exactly.
\end{assessment}

\subsection{Recovering the cluster projector from the transfer}
\label{sec:s6v64-projector}

\begin{theorem}[Riesz extraction of the relative cluster]
\label{thm:s6v64-Riesz}
Let \(T\) be compact self-adjoint, with simple top eigenvalue one and an
isolated relative eigenvalue \(q<1\) separated from the remaining spectrum
by a contour \(\Gamma\).  Then the complete relative cluster projection is
\begin{equation}
 \boxed{
 P_E={1\over2\pi i}\int_\Gamma(z-T)^{-1}\,\dd z.}
 \label{eq:s6v64-Riesz}
\end{equation}
For every physical writer \(\varphi\), the projected frame vector
\(P_E\varphi\) is therefore determined by the full transfer resolvent.
\end{theorem}

\begin{proof}
This is the spectral Riesz projection theorem for an isolated compact
self-adjoint spectral cluster.
\end{proof}

\begin{assessment}[Projection is not a local observable]
\label{ass:s6v64-Riesz-locality}
Formula \eqref{eq:s6v64-Riesz} makes finite-frame extraction exact, but its
resolvent is global and its norm is controlled by the isolation margin.
Replacing it with a finite correlation or transfer-power filter creates a
new inverse-isolation error.  Thus the spin-network density theorem solves
algebraic source completeness at a fixed regulator, not uniform cluster
identification.
\end{assessment}

\subsection{The source-complete spectral score card}
\label{sec:s6v64-card}

\begin{definition}[Spin-network cluster frame card, SNCFC]
\label{def:s6v64-SNCFC}
SNCFC consists of:
\begin{enumerate}[label=\textup{(\roman*)},leftmargin=3.5em]
\item the complete gauge-invariant spin-network algebra on every finite
      boundary regulator and its protected relative projection;
\item a specified finite or nested bank whose projection spans or is dense
      in every top relative score cluster;
\item uniform control of the synthesis Gram \(G_{\theta,j}\), including a
      lower frame stock \(a_{\theta,j}\);
\item an exact or approximated signed score pencil
      \(H_{\theta,j}^{\rm sep}\);
\item the generalized lower bound
      \[
       H_{\theta,j}^{\rm sep}
       \succeq {m\delta_j\over q_{\theta,j}}G_{\theta,j}
      \]
      with approximation debit
      \(\varepsilon_{\theta,j}/a_{\theta,j}=o(\delta_j)\);
\item uniform cluster isolation or a cluster-crossing formulation that
      retains the complete maximal relative space;
\item SASSC endpoint projectivity and Mosco convergence; and
\item noncollapse and cyclicity of the limiting spin-network OS algebra.
\end{enumerate}
\end{definition}

\begin{theorem}[Conditional SNCFC completion]
\label{thm:s6v64-SNCFC}
If SNCFC is populated with \(m>0\), the reconstructed Yang--Mills
Hamiltonian has relative vacuum gap at least \(m\).
\end{theorem}

\begin{proof}
Rows~\textup{(i)--(vi)} and
\Cref{thm:s6v64-frame,cor:s6v64-cluster,
lem:s6v64-frame-error,prop:s6v64-dense-extension}
establish the full-cluster SASSC separation, not merely a compression
bound.  Rows~\textup{(vii)--(viii)} permit the V63 and V53
spectral-edge/Mosco handoffs.
\end{proof}

\begin{assessment}[Exact frontier after V64]
\label{ass:s6v64-frontier}
At each finite Wilson regulator, gauge-invariant spin networks are
source-complete and a finite projected frame for the top relative cluster
exists.  The unresolved theorem is now uniform and quantitative: construct
those frames cofinally, keep their conditioning under control, and prove
the signed generalized score pencil lower bound at order \(\delta\).
Neither density nor positivity improvement gives that lower bound.  The
V65 companion audit should look specifically for a uniform frame,
Riesz-cluster, or generalized-pencil mechanism before the programme
chooses the next attack.
\end{assessment}

\begin{programme}[V65: compulsory audit at the source-complete frontier]
\label{prog:s6v64-V65}
\begin{enumerate}[leftmargin=3em]
\item Audit the current SM and NS files for uniform frame conditioning,
      cluster projection, generalized eigenvalue, or dense-core coercivity
      mechanisms.
\item Preserve the distinction between finite-regulator density and
      cofinal uniformity.
\item Reinspect the Yang--Mills source layers for a complete Wilson-network
      bank rather than a selected four-loop source.
\item If no signed pencil estimate exists, derive a dual obstruction:
      construct the worst score vector from the negative part of the
      cluster operator and identify which observable would detect it.
\item Continue without promoting finite-bank numerics to a full mass gap.
\end{enumerate}
\end{programme}

\begin{firewall}[Claim boundary after sixth-series V64]
\label{fire:s6v64-boundary}
V64 proves the direction of compression inequalities and an explicit
counterexample; proves the exact generalized frame-pencil criterion;
proves density of gauge-invariant spin networks for finite Wilson/Doob
boundary laws; extracts finite frames for isolated finite clusters;
quantifies frame-conditioning and clock-scaled matrix errors; and gives
Riesz extraction of the full cluster.

V64 does not construct a cofinally uniform spin-network frame, control its
conditioning or cluster isolation, prove the signed generalized score
pencil, close projectivity/Mosco convergence, construct the continuum OS
theory, or prove the full Yang--Mills mass gap.
\end{firewall}

\begin{theorem}[Sixth-series V64 result ledger]
\label{thm:s6v64-results}
The positive conclusions in \Cref{fire:s6v64-boundary} are exactly those of
\Cref{prop:s6v64-compression,
prop:s6v64-counterexample,
thm:s6v64-frame,
cor:s6v64-cluster,
thm:s6v64-spin-density,
cor:s6v64-Doob-density,
thm:s6v64-frame-extraction,
lem:s6v64-frame-error,
cor:s6v64-clock-error,
prop:s6v64-dense-extension,
thm:s6v64-Riesz,
thm:s6v64-SNCFC}.
\end{theorem}

\begin{proof}
Each conclusion is the cited result with its gauge-invariance, finite
regulator, frame-surjectivity, conditioning, cluster-isolation, physical
clock, and projective-limit qualifications retained.
\end{proof}

\paragraph{Parent-version record.}
V64 was formed from
\texttt{Renewal\_Yang\_Mills\_Mass\_Gap\_Research\_Notes\_V63.tex}.
The V63 parent had \(28\,313\) logical source lines,
\(994\,460\) bytes, and SHA--256
\[
 \texttt{739654baa2e95fc286deeb727754b9b1531258c0b14412880ca0a8643d8ec191}.
\]
The next version performs the compulsory companion audit.

% =============================================================================
% END APPENDED SIXTH-SERIES V64 MATERIAL
% =============================================================================

% =============================================================================
% BEGIN APPENDED SIXTH-SERIES V65 MATERIAL
% =============================================================================

\section{V65 compulsory companion audit: bad-mode witnesses and the
cluster-filtering time}
\label{sec:s6v65}

\subsection{Companion snapshot and decision}
\label{sec:s6v65-audit}

\begin{definition}[V65 companion files]
\label{def:s6v65-files}
The active companion files at this audit are:
\begin{center}
\begin{tabular}{p{0.43\linewidth}rr}
\toprule
file & logical lines & bytes\\
\midrule
\texttt{Renewal\_SM\_Einstein\_Continuum\_Research\_Notes\_V11.tex}
& \(3\,876\) & \(146\,117\)\\
\texttt{Renewal\_NS\_Stability\_Research\_Notes\_V26.tex}
& \(5\,319\) & \(278\,283\)\\
\bottomrule
\end{tabular}
\end{center}
Their SHA--256 digests are
\[
 \texttt{bfd7ba8ff3df159f1840ee414e19b6110944da228dc97e4a77cfedd0feeae4eb}
\]
and
\[
 \texttt{82c887f8c2c69e4f28fad7c3dc8de5b9099cb5a4bf431eba1fff3e91935978ad}.
\]
\end{definition}

\begin{assessment}[SM V11 audit]
\label{ass:s6v65-SM}
Since the V60 audit, SM V8--V11 has developed:
\begin{enumerate}[label=\textup{(\roman*)},leftmargin=3.5em]
\item one- and two-marked Dyson-path locality from a weighted finite-range
      influence matrix;
\item a sector small-gain compiler that yields finite-horizon propagation
      under finite absolute influence, but cannot yield strict contraction
      when a Ward sector has zero diagonal margin and nonnegative cross
      influences enter it;
\item explicit local readout bounds for Hamiltonian and Lindblad
      interactions; and
\item a correction that support-reset two-source locality does not
      automatically apply to generic Hamiltonian commutators, which require
      connected nested-commutator paths.
\end{enumerate}
These mechanisms can support finite-horizon projectivity of spin-network
writers.  They do not provide a source-complete top-relative cluster
frame, a uniform frame lower bound, a Riesz isolation margin, or a signed
Yang--Mills score pencil.
\end{assessment}

\begin{assessment}[NS V26 audit]
\label{ass:s6v65-NS}
NS V26 is again byte-for-byte unchanged.  Its finite-dimensional source
and Dirichlet calculations contain no compact half-slab spectral cluster,
spin-network frame, or clock-scaled generalized pencil that can populate
SNCFC.
\end{assessment}

\begin{theorem}[V65 audit decision]
\label{thm:s6v65-audit-decision}
The Yang--Mills programme imports marked finite-horizon path locality only
as a possible projectivity tool for fixed-support writers.  It does not
import the companion absolute influence matrix as a gap mechanism and does
not change the V64 signed cluster-pencil target.
\end{theorem}

\begin{proof}
The locality conclusions require only finite weighted influence, while the
companion itself proves that a zero-margin Ward sector obstructs strict
absolute small gain.  By
\Cref{ass:s6v65-SM,ass:s6v65-NS}, neither source supplies the missing
cluster frame or signed lower bound.
\end{proof}

\subsection{Every failed cluster inequality has a source witness}
\label{sec:s6v65-witness}

\begin{theorem}[Dense-source detection of a bad score direction]
\label{thm:s6v65-bad-witness}
Let \(E\) be an isolated finite-dimensional relative cluster, let
\(J=J^*\in{\cal B}(E)\), and let \({\cal D}\subset E\) be the set of
projected gauge-invariant spin-network linear combinations.  By V64,
\({\cal D}\) is dense in \(E\).  If \(J\not\succeq cI_E\), then for every
\(\varepsilon>0\) there exists \(0\ne w\in{\cal D}\) such that
\begin{equation}
 {\langle w,Jw\rangle\over\|w\|^2}<c+\varepsilon.
 \label{eq:s6v65-bad-witness}
\end{equation}
\end{theorem}

\begin{proof}
The variational principle supplies a unit vector \(e\in E\) with
\(\langle e,Je\rangle<c+\varepsilon/2\).  Choose
\(w_n\in{\cal D}\) with \(w_n\to e\).  Since \(J\) is bounded,
\[
 {\langle w_n,Jw_n\rangle\over\|w_n\|^2}
 \longrightarrow\langle e,Je\rangle.
\]
For sufficiently large \(n\), the quotient is below
\(c+\varepsilon\).
\end{proof}

\begin{lemma}[Quantitative witness approximation]
\label{lem:s6v65-quantitative-witness}
Let \(\|e\|=1\), \(\|J\|\le M\), and
\(\|w-e\|\le\eta<1/2\).  Then
\begin{equation}
 \left|
 {\langle w,Jw\rangle\over\|w\|^2}
 -\langle e,Je\rangle
 \right|
 \le {6M\eta\over(1-\eta)^2}.
 \label{eq:s6v65-Rayleigh-error}
\end{equation}
\end{lemma}

\begin{proof}
First,
\[
 |\langle w,Jw\rangle-\langle e,Je\rangle|
 \le M\|w-e\|(\|w\|+\|e\|)
 \le M\eta(2+\eta).
\]
Also
\[
 |\|w\|^2-1|\le\eta(2+\eta),\qquad
 \|w\|\ge1-\eta.
\]
Insert and subtract
\(\langle e,Je\rangle/\|w\|^2\), use
\(|\langle e,Je\rangle|\le M\), and bound
\(2+\eta\le5/2\); the displayed constant \(6\) dominates the resulting
sum.
\end{proof}

\begin{corollary}[No algebraic hiding place at fixed regulator]
\label{cor:s6v65-no-hiding}
At a fixed compact Wilson regulator, failure of the full SASSC cluster
inequality can be detected arbitrarily accurately by a finite linear
combination of gauge-invariant spin networks after exact cluster
projection.
\end{corollary}

\begin{proof}
Apply \Cref{thm:s6v65-bad-witness} and the V64 spin-network density
theorem.
\end{proof}

\begin{assessment}[The witness may escape cofinally]
\label{ass:s6v65-escape}
The number, representation labels, spatial range, and coefficient norm of
the detecting spin-network combination may grow with regulator or volume.
Thus \Cref{cor:s6v65-no-hiding} is not a uniform finite-bank theorem.
A bad mode can escape every fixed writer list while remaining detectable
by the complete algebra.
\end{assessment}

\subsection{Transfer-power extraction of the relative cluster}
\label{sec:s6v65-filter}

Let \(T\) be a positive compact self-adjoint contraction with spectral
decomposition
\begin{equation}
 T=P_1+qP_2+T_<,
 \label{eq:s6v65-spectral-decomposition}
\end{equation}
where \(P_1\) is the vacuum/common-factor projection,
\(P_2\) is the complete top relative cluster,
\[
 0\le T_<\le r(I-P_1-P_2),\qquad 0\le r<q<1.
\]

\begin{theorem}[Exact transfer-power cluster filter]
\label{thm:s6v65-power-filter}
Define
\begin{equation}
 {\cal F}_n=q^{-n}(T^n-P_1).
 \label{eq:s6v65-filter}
\end{equation}
Then
\begin{equation}
 \boxed{
 \|{\cal F}_n-P_2\|
 \le\left({r\over q}\right)^n.}
 \label{eq:s6v65-filter-error}
\end{equation}
Consequently, for every source \(\varphi\),
\[
 {\cal F}_n\varphi\longrightarrow P_2\varphi
\]
in norm.
\end{theorem}

\begin{proof}
Functional calculus gives
\[
 {\cal F}_n
 =P_2+q^{-n}T_<^n.
\]
The remainder norm is at most \((r/q)^n\).
\end{proof}

\begin{corollary}[Filtered witness accuracy]
\label{cor:s6v65-filtered-witness}
Let \(\|J\|\le M\).  If a source \(\varphi\) has
\(\|P_2\varphi\|\ge b>0\), then the normalized filtered vector
\[
 w_n={{\cal F}_n\varphi\over\|{\cal F}_n\varphi\|}
\]
approximates the normalized cluster vector
\[
 e={P_2\varphi\over\|P_2\varphi\|}
\]
with error
\begin{equation}
 \|w_n-e\|
 \le {2\|\varphi\|\over b}
 \left({r\over q}\right)^n
 \label{eq:s6v65-vector-filter-error}
\end{equation}
whenever the right side is below one.  Its score Rayleigh error is then
controlled by \Cref{lem:s6v65-quantitative-witness}.
\end{corollary}

\begin{proof}
The unnormalized difference is bounded by
\[
 \|({\cal F}_n-P_2)\varphi\|
 \le(r/q)^n\|\varphi\|.
\]
Normalization of two nonzero vectors changes their distance by at most
twice the unnormalized error divided by the norm \(b\) of the target,
which gives \eqref{eq:s6v65-vector-filter-error}.  Apply
\Cref{lem:s6v65-quantitative-witness}.
\end{proof}

\subsection{The physical time cost of resolving the cluster}
\label{sec:s6v65-time}

\begin{theorem}[Logarithmic physical filtering time]
\label{thm:s6v65-filter-time}
Suppose one transfer step has duration \(\delta_j\), and
\begin{equation}
 q_j=e^{-\delta_jE_{1,j}},\qquad
 r_j\le e^{-\delta_jE_{2,j}},
 \qquad
 E_{2,j}-E_{1,j}\ge\Delta>0.
 \label{eq:s6v65-energy-separation}
\end{equation}
To make the operator filter error
\((r_j/q_j)^{n_j}=o(\delta_j)\), it is sufficient and, at the exponential
rate in \eqref{eq:s6v65-energy-separation}, necessary that
\begin{equation}
 n_j\delta_j\Delta-\log(1/\delta_j)\longrightarrow+\infty.
 \label{eq:s6v65-filter-time}
\end{equation}
Thus the physical filter duration \(t_j=n_j\delta_j\) grows at least
logarithmically in \(1/\delta_j\).
\end{theorem}

\begin{proof}
Equation \eqref{eq:s6v65-energy-separation} gives
\[
 \left({r_j\over q_j}\right)^{n_j}
 \le e^{-n_j\delta_j\Delta}.
\]
The ratio of this bound to \(\delta_j\) tends to zero exactly when
\[
 -n_j\delta_j\Delta+\log(1/\delta_j)\to-\infty,
\]
which is \eqref{eq:s6v65-filter-time}.  If equality holds in the
exponential spectral rates, the same calculation is necessary.
\end{proof}

\begin{corollary}[Frame-conditioned filter scale]
\label{cor:s6v65-frame-filter}
If the generalized pencil requires cluster-vector accuracy
\(o(a_j\delta_j)\), where \(a_j\) is the frame lower stock, then the
transfer-power filter must satisfy
\begin{equation}
 \left({r_j\over q_j}\right)^{n_j}
 {\|\varphi_j\|\over b_j}
 =o(a_j\delta_j).
 \label{eq:s6v65-conditioned-filter}
\end{equation}
Collapsing source overlap \(b_j\), collapsing frame stock \(a_j\), or a
closing energy separation \(E_{2,j}-E_{1,j}\) lengthens the required
physical filtering time.
\end{corollary}

\begin{proof}
Combine the vector error
\eqref{eq:s6v65-vector-filter-error} with the V64 clock-scaled generalized
pencil accuracy.
\end{proof}

\begin{assessment}[Finite-horizon locality does not control the filter]
\label{ass:s6v65-horizon}
The marked-path estimates imported from SM V11 apply on fixed physical
horizons with constants growing in that horizon.  The filter required by
\eqref{eq:s6v65-filter-time} has \(t_j\to\infty\).  Therefore fixed-horizon
projectivity does not automatically make the Riesz or power-filtered
cluster writers projective.  A long-time uniform locality/clustering
theorem would already contain mass-like information and cannot be assumed
for free.
\end{assessment}

\subsection{A dual failure certificate}
\label{sec:s6v65-certificate}

\begin{definition}[Cluster failure certificate, CFC]
\label{def:s6v65-CFC}
At a finite regulator and interpolation parameter, CFC consists of:
\begin{enumerate}[label=\textup{(\roman*)},leftmargin=3.5em]
\item a gauge-invariant spin-network source \(\varphi\);
\item a certified nonzero cluster overlap
      \(b\le\|P_2\varphi\|\);
\item a Riesz or transfer-power filter with a proved approximation error;
\item a score Rayleigh value below the required threshold
      \(m\delta/q\), with all frame/filter errors included.
\end{enumerate}
Such a certificate rigorously disproves the proposed SASSC constant at
that regulator; it does not disprove the existence of a smaller positive
constant unless the violation is cofinal.
\end{definition}

\begin{theorem}[Completeness of finite-regulator failure certificates]
\label{thm:s6v65-CFC-complete}
If the SASSC cluster inequality fails at a finite compact regulator, then
for every positive tolerance there exists a CFC whose certified Rayleigh
value violates the threshold up to that tolerance.
\end{theorem}

\begin{proof}
\Cref{thm:s6v65-bad-witness} supplies a projected spin-network witness.
\Cref{thm:s6v65-power-filter} approximates its cluster projection
arbitrarily accurately.  The quantitative Rayleigh estimate
\Cref{lem:s6v65-quantitative-witness} transfers the strict violation after
the tolerance is chosen smaller than its margin.
\end{proof}

\begin{assessment}[Exact frontier after V65]
\label{ass:s6v65-frontier}
The finite-regulator cluster problem is now decidable in principle by a
source-complete spin-network failure certificate or generalized-pencil
proof.  The continuum obstruction has moved to uniformity: source overlap,
frame conditioning, cluster energy separation, and signed score values
must remain controlled while the physical filter horizon diverges.  No
companion programme supplies these rows.  The next attack should avoid
long-time cluster extraction by seeking a direct quadratic-form lower
bound on the whole relative carrier, for example through a local-to-global
uncertainty principle for the signed insertion.
\end{assessment}

\begin{programme}[V66: direct form domination without cluster filtering]
\label{prog:s6v65-V66}
\begin{enumerate}[leftmargin=3em]
\item Rewrite SASSC as an operator inequality on
      \(Q_\theta{\cal H}\) without projecting onto the top cluster.
\item Seek a decomposition of the signed inserted form into local positive
      pieces plus a controlled overlap matrix.
\item Prove a local-to-global frame inequality on the complete
      gauge-invariant spin-network core.
\item Track the order-\(\delta\) constant and prevent absolute interaction
      majorants from erasing signed cancellations.
\item If the local pieces are not positive, isolate the exact negative
      return responsible and move upstream to its Wilsonian covariance
      representation.
\end{enumerate}
\end{programme}

\begin{firewall}[Claim boundary after sixth-series V65]
\label{fire:s6v65-boundary}
V65 performs the compulsory SM/NS audit; proves that every failed
finite-regulator cluster inequality has a projected spin-network witness;
quantifies witness approximation; proves transfer-power extraction of the
complete relative cluster; computes the logarithmically diverging physical
filter time needed for \(o(\delta)\) accuracy; and gives complete
finite-regulator failure certificates.

V65 does not produce a cofinally uniform cluster frame or filter, prove
the signed score lower bound, establish long-time locality, close
projectivity/Mosco convergence, construct the continuum OS theory, or
prove the full Yang--Mills mass gap.
\end{firewall}

\begin{theorem}[Sixth-series V65 result ledger]
\label{thm:s6v65-results}
The positive conclusions in \Cref{fire:s6v65-boundary} are exactly those of
\Cref{def:s6v65-files,
thm:s6v65-audit-decision,
thm:s6v65-bad-witness,
lem:s6v65-quantitative-witness,
cor:s6v65-no-hiding,
thm:s6v65-power-filter,
cor:s6v65-filtered-witness,
thm:s6v65-filter-time,
cor:s6v65-frame-filter,
thm:s6v65-CFC-complete}.
\end{theorem}

\begin{proof}
Each item is the cited result with its finite-regulator density, cluster
isolation, overlap, frame conditioning, physical-time, and cofinal
uniformity qualifications unchanged.
\end{proof}

\paragraph{Parent-version record.}
V65 was formed from
\texttt{Renewal\_Yang\_Mills\_Mass\_Gap\_Research\_Notes\_V64.tex}.
The V64 parent had \(28\,795\) logical source lines,
\(1\,011\,601\) bytes, and SHA--256
\[
 \texttt{f28a12bb8d70028f01e500053e0cf8cfffd90d1690d1e3f463ac9acec4f00478}.
\]
The next compulsory companion audit is V70.

% =============================================================================
% END APPENDED SIXTH-SERIES V65 MATERIAL
% =============================================================================

% =============================================================================
% BEGIN APPENDED SIXTH-SERIES V66 MATERIAL
% =============================================================================

\section{V66: direct signed form domination without cluster filtering}
\label{sec:s6v66}

\subsection{Purpose and fixed-reference notation}
\label{sec:s6v66-purpose}

V65 showed that exact extraction of the top relative transfer cluster
costs a physical time of order \(\log(1/\delta)\).  This is incompatible
with using only the fixed-horizon locality estimates presently available
from the companion programmes.  The purpose of V66 is therefore to remove
the cluster projector from the sufficient hypothesis.  The replacement
is a signed G\aa rding inequality for the derivative of the normalized
transfer form on the \emph{entire} vacuum-orthogonal carrier.  The
argument below is finite-regulator and exact; its cofinal input is stated
as a separate card.

Retain the norm-\(C^1\) family of V63,
\[
 A_\theta:{\cal H}_0\longrightarrow{\cal H}_{1/2},\qquad
 B_\theta=A_\theta^*A_\theta,
\]
with simple leading eigenpair
\((\lambda_{1,\theta},u_\theta)\), and put
\begin{equation}
 T_\theta={B_\theta\over\lambda_{1,\theta}},\qquad
 L_\theta=I-T_\theta,\qquad
 Q_\theta=I-|u_\theta\rangle\langle u_\theta|.
 \label{eq:s6v66-normalized}
\end{equation}
Thus \(L_\theta u_\theta=0\), \(0\le T_\theta\le I\), and
\(L_\theta=Q_\theta L_\theta Q_\theta\).  The one-step relative form
gap is
\begin{equation}
 g_\theta
 =\inf_{\substack{f\in Q_\theta{\cal H}_0\\\|f\|=1}}
 \langle f,L_\theta f\rangle
 =1-\|T_\theta|_{Q_\theta{\cal H}_0}\|.
 \label{eq:s6v66-gap}
\end{equation}

\subsection{The whole-carrier signed insertion}
\label{sec:s6v66-insertion}

Assume, as in V63, that the positive half-slab amplitude satisfies
\[
 \partial_\theta F_\theta(y,x)
 =-{1\over2}R_\theta(y,x)F_\theta(y,x),
\]
and write \(A_{\theta,R}\) for insertion of \(R_\theta\).  Let
\({\cal J}_{1,\theta}\) be the leading Doob score of
\Cref{lem:s6v63-leading-score}.

\begin{proposition}[Exact whole-carrier derivative form]
\label{prop:s6v66-derivative-form}
For every \(f\in{\cal H}_0\),
\begin{equation}
 \boxed{
 \langle f,L_\theta'f\rangle
 ={1\over\lambda_{1,\theta}}
 \operatorname{Re}
 \langle A_\theta f,
 (A_{\theta,R}-{\cal J}_{1,\theta}A_\theta)f\rangle.}
 \label{eq:s6v66-derivative-form}
\end{equation}
Equivalently,
\begin{equation}
 L_\theta'
 ={1\over2\lambda_{1,\theta}}
 \left(A_{\theta,R}^*A_\theta+
       A_\theta^*A_{\theta,R}\right)
 -{\cal J}_{1,\theta}T_\theta .
 \label{eq:s6v66-derivative-operator}
\end{equation}
Neither expression differentiates the singular ground vectors.
\end{proposition}

\begin{proof}
By \Cref{lem:s6v63-B-prime},
\[
 B_\theta'
 =-{1\over2}
 (A_{\theta,R}^*A_\theta+A_\theta^*A_{\theta,R}),
\]
whereas
\(\lambda_{1,\theta}'/\lambda_{1,\theta}
=-{\cal J}_{1,\theta}\).  Substitution into
\[
 L_\theta'
 =-{B_\theta'\over\lambda_{1,\theta}}
 +{\lambda_{1,\theta}'\over\lambda_{1,\theta}}T_\theta
\]
gives \eqref{eq:s6v66-derivative-operator}.  Taking its quadratic form
gives \eqref{eq:s6v66-derivative-form}.
\end{proof}

\begin{proposition}[Exact centered two-arm kernel]
\label{prop:s6v66-two-arm}
Suppose \(F_\theta\) and \(R_\theta\) are real, and that Fubini's theorem
applies to the displayed integrals.  Then
\begin{align}
 \langle f,L_\theta'f\rangle
 &={1\over2\lambda_{1,\theta}}
 \int\mu_{1/2}(\dd y)
 \int\mu_0(\dd x)\int\mu_0(\dd x')\,
 \overline{f(x')}f(x)                                      \notag\\
 &\quad\cdot F_\theta(y,x')F_\theta(y,x)
 \bigl[
 R_\theta(y,x)+R_\theta(y,x')
 -2{\cal J}_{1,\theta}
 \bigr].
 \label{eq:s6v66-two-arm}
\end{align}
If \(R_\theta=\sum_ZR_{\theta,Z}\) with convergence in a class allowing
termwise integration, then the form in
\eqref{eq:s6v66-two-arm} is the sum of the corresponding centered local
two-arm forms.
\end{proposition}

\begin{proof}
Expand the first inner product in
\eqref{eq:s6v66-derivative-form}.  Its real part equals one half of that
integral plus its complex conjugate.  In the conjugate exchange \(x\)
and \(x'\).  This replaces the single factor \(R_\theta(y,x)\) by
\(\frac12(R_\theta(y,x)+R_\theta(y,x'))\).  The term
\(-{\cal J}_{1,\theta}\|A_\theta f\|^2\) supplies the remaining
\(-2{\cal J}_{1,\theta}\) inside the same symmetrized kernel.
Linearity and the stated convergence justify the last assertion.
\end{proof}

\begin{remark}[The sign has not been discarded]
\label{rem:s6v66-sign}
The bracket in \eqref{eq:s6v66-two-arm} is centered and generally changes
sign.  Replacing it by its absolute value gives only an upper bound on the
magnitude of \(L_\theta'\); it cannot prove a positive spectral response.
Any successful local decomposition must retain a positive diagonal
contribution and charge only the overlap terms to an absolute majorant.
\end{remark}

\subsection{A direct G\aa rding-to-gap compiler}
\label{sec:s6v66-garding}

\begin{definition}[Direct signed form inequality]
\label{def:s6v66-DSFI}
For nonnegative functions \(a_\theta,b_\theta\), the direct signed form
inequality, abbreviated DSFI, is
\begin{equation}
 \boxed{
 Q_\theta L_\theta'Q_\theta
 \succeq
 b_\theta\delta\,Q_\theta
 -a_\theta Q_\theta L_\theta Q_\theta}
 \label{eq:s6v66-DSFI}
\end{equation}
as a quadratic-form inequality on the whole relative carrier
\(Q_\theta{\cal H}_0\).
\end{definition}

\begin{theorem}[Whole-carrier G\aa rding edge-opening theorem]
\label{thm:s6v66-garding-gap}
Fix a finite regulator and \(\delta>0\).  Suppose:
\begin{enumerate}[label=\textup{(\roman*)},leftmargin=3.5em]
\item \(T_\theta\) is a norm-\(C^1\), compact, nonnegative self-adjoint
      family on a fixed Hilbert space;
\item its eigenvalue \(1\) is simple for every \(\theta\in[0,1]\);
\item the top relative eigenvalue is an isolated finite-multiplicity
      cluster, allowing crossings inside that cluster;
\item \(a,b:[0,1]\to[0,\infty)\) are continuous and
      \eqref{eq:s6v66-DSFI} holds for every \(\theta\).
\end{enumerate}
Set
\[
 {\cal A}(t)=\int_0^t a_s\,\dd s.
\]
Then
\begin{equation}
 \boxed{
 g_1\ge
 e^{-{\cal A}(1)}g_0
 +\delta\int_0^1
 e^{-({\cal A}(1)-{\cal A}(s))}b_s\,\dd s.}
 \label{eq:s6v66-integrated-gap}
\end{equation}
In particular, a regulator-uniform positive lower bound on the integral
in \eqref{eq:s6v66-integrated-gap} opens an order-\(\delta\) transfer gap
without constructing the top relative cluster projector.
\end{theorem}

\begin{proof}
Let \(E_\theta\subset Q_\theta{\cal H}_0\) be the eigenspace of
\(L_\theta\) at its smallest positive eigenvalue \(g_\theta\).
The right derivative of the bottom eigenvalue of this isolated cluster
is the smallest eigenvalue of
\(P_{E_\theta}L_\theta'P_{E_\theta}\).  Therefore DSFI gives
\begin{equation}
 D_+g_\theta\ge b_\theta\delta-a_\theta g_\theta.
 \label{eq:s6v66-dini}
\end{equation}
This conclusion is unchanged at an internal multiplicity crossing,
because the compressed derivative formula already takes the smallest
branch derivative.

Define
\[
 h(\theta)
 =e^{{\cal A}(\theta)}g_\theta
 -\delta\int_0^\theta e^{{\cal A}(s)}b_s\,\dd s.
\]
The product rule for lower right Dini derivatives and
\eqref{eq:s6v66-dini} give \(D_+h(\theta)\ge0\).  Continuity of
\(g,a,b\) and the Dini monotonicity theorem imply
\(h(1)\ge h(0)=g_0\).  Divide by \(e^{{\cal A}(1)}\) to obtain
\eqref{eq:s6v66-integrated-gap}.
\end{proof}

\begin{corollary}[Constant-stock form]
\label{cor:s6v66-constant-stock}
If \(a_\theta\le a_*\) and \(b_\theta\ge b_*>0\), then
\begin{equation}
 g_1\ge e^{-a_*}g_0+
 \begin{cases}
 \displaystyle
 \delta b_*{1-e^{-a_*}\over a_*},&a_*>0,\\
 \delta b_*,&a_*=0.
 \end{cases}
 \label{eq:s6v66-constant-gap}
\end{equation}
\end{corollary}

\begin{proof}
Use \({\cal A}(1)-{\cal A}(s)\le a_*(1-s)\) in
\eqref{eq:s6v66-integrated-gap} and integrate.
\end{proof}

\begin{assessment}[Why this removes the V65 long-time debit]
\label{ass:s6v66-no-filter}
The hypothesis \eqref{eq:s6v66-DSFI} is imposed directly on
\(Q_\theta{\cal H}_0\), not merely on its top transfer cluster.  Its
verification therefore needs neither the Riesz contour from V64 nor the
transfer-power filter from V65.  The theorem preserves the one-step
\(O(\delta)\) clock.  The price is stronger: DSFI must control every
relative direction, including high-energy directions, with their
possible loss charged to \(a_\theta L_\theta\).
\end{assessment}

\subsection{Local positive pieces and overlap control}
\label{sec:s6v66-local-global}

\begin{theorem}[Local-frame DSFI compiler]
\label{thm:s6v66-local-frame}
For each \(\theta\), let
\[
 D_{\theta,z}:Q_\theta{\cal H}_0\longrightarrow{\cal K}_{\theta,z},
 \qquad z\in{\cal Z}_\theta,
\]
be bounded analysis maps satisfying the lower frame estimate
\begin{equation}
 \sum_z\|D_{\theta,z}f\|^2
 \ge A_\theta\|f\|^2,
 \qquad A_\theta>0.
 \label{eq:s6v66-lower-frame}
\end{equation}
Suppose there are \(\mu_\theta>0\), \(a_\theta\ge0\), and a bounded
operator \(C_\theta\) on \(\ell^2({\cal Z}_\theta)\) such that
\(\|C_\theta\|\le\kappa_\theta<\mu_\theta\) and
\begin{align}
 \langle f,L_\theta'f\rangle
 &\ge
 \delta\mu_\theta\sum_z\|D_{\theta,z}f\|^2
 -\delta\sum_{z,w}
 (C_\theta)_{zw}
 \|D_{\theta,z}f\|\|D_{\theta,w}f\| \notag\\
 &\quad-a_\theta\langle f,L_\theta f\rangle
 \label{eq:s6v66-local-decomposition}
\end{align}
for every \(f\in Q_\theta{\cal H}_0\), where the middle quadratic
expression is real.  Then DSFI holds with
\begin{equation}
 b_\theta=A_\theta(\mu_\theta-\kappa_\theta).
 \label{eq:s6v66-b-stock}
\end{equation}
\end{theorem}

\begin{proof}
Put \(d_z=\|D_{\theta,z}f\|\).  The operator-norm bound gives
\[
 \left|\sum_{z,w}(C_\theta)_{zw}d_zd_w\right|
 \le\kappa_\theta\sum_zd_z^2.
\]
Consequently \eqref{eq:s6v66-local-decomposition} is bounded below by
\[
 \delta(\mu_\theta-\kappa_\theta)\sum_zd_z^2
 -a_\theta\langle f,L_\theta f\rangle.
\]
Apply \eqref{eq:s6v66-lower-frame}.  This is
\eqref{eq:s6v66-DSFI} with \eqref{eq:s6v66-b-stock}.
\end{proof}

\begin{corollary}[Cofinal stock sufficient for an endpoint gap]
\label{cor:s6v66-cofinal-stock}
For a cofinal regulator sequence \(j\), suppose the hypotheses of
\Cref{thm:s6v66-local-frame} hold and
\begin{equation}
 \inf_{j,\theta}A_{j,\theta}
       (\mu_{j,\theta}-\kappa_{j,\theta})>0,
 \qquad
 \sup_j\int_0^1a_{j,\theta}\,\dd\theta<\infty.
 \label{eq:s6v66-cofinal-stocks}
\end{equation}
Then the endpoint dimensionless gaps satisfy
\[
 \inf_j{g_{j,1}\over\delta_j}>0
\]
provided \(g_{j,0}\ge0\).  This is the finite-regulator spectral row
needed by the V53 Mosco handoff.
\end{corollary}

\begin{proof}
Apply \Cref{thm:s6v66-garding-gap} and use
\eqref{eq:s6v66-cofinal-stocks} to bound its integral uniformly below.
\end{proof}

\begin{remark}[The three stocks are logically distinct]
\label{rem:s6v66-three-stocks}
The frame stock \(A\) asserts that the local probes see every relative
direction.  The diagonal stock \(\mu\) is the positive signed response of
each local piece.  The overlap stock \(\kappa\) measures destructive
interference between pieces.  Locality estimates can help bound
\(\kappa\), but they do not by themselves produce either \(A>0\) or
\(\mu>0\).
\end{remark}

\subsection{A complete spin-network pencil with no spectral filter}
\label{sec:s6v66-pencil}

Let \(\{\phi_n\}_{n\ge1}\) be a countable gauge-invariant spin-network
family whose linear span is dense in the finite-regulator invariant
Hilbert space.  Put
\[
 \psi_{\theta,n}=Q_\theta\phi_n.
\]
The projected family is dense in \(Q_\theta{\cal H}_0\), although it need
not be linearly independent.

\begin{theorem}[Infinite-pencil characterization of DSFI]
\label{thm:s6v66-pencil}
Fix \(\theta\), \(a\ge0\), and \(b\in\mathbb R\).  For \(N\ge1\), define
\begin{align}
 [G_\theta^{(N)}]_{mn}
 &=\langle\psi_{\theta,m},\psi_{\theta,n}\rangle,
 \label{eq:s6v66-pencil-G}\\
 [H_{\theta,a}^{(N)}]_{mn}
 &=\langle\psi_{\theta,m},
 (L_\theta'+aL_\theta)\psi_{\theta,n}\rangle.
 \label{eq:s6v66-pencil-H}
\end{align}
If \(L_\theta'\) is bounded, then
\begin{equation}
 Q_\theta(L_\theta'+aL_\theta)Q_\theta
 \succeq b\delta Q_\theta
 \label{eq:s6v66-whole-inequality}
\end{equation}
if and only if, for every \(N\),
\begin{equation}
 \boxed{H_{\theta,a}^{(N)}
 \succeq b\delta\,G_\theta^{(N)}.}
 \label{eq:s6v66-all-pencils}
\end{equation}
No inverse Gram matrix and no top-cluster approximation enters this
criterion.
\end{theorem}

\begin{proof}
If \eqref{eq:s6v66-whole-inequality} holds, apply it to
\(f=\sum_{n=1}^Nc_n\psi_{\theta,n}\); its quadratic form is exactly
\[
 c^*H_{\theta,a}^{(N)}c
 -b\delta\,c^*G_\theta^{(N)}c\ge0.
\]
Conversely, \eqref{eq:s6v66-all-pencils} proves the same inequality on
every finite projected spin-network combination.  Their union is dense
in \(Q_\theta{\cal H}_0\).  Boundedness of
\(L_\theta'+aL_\theta-b\delta I\) lets the inequality pass to the norm
closure.
\end{proof}

\begin{corollary}[Failure has a finite unfiltered witness]
\label{cor:s6v66-finite-witness}
If \eqref{eq:s6v66-whole-inequality} fails, then it fails on some finite
projected spin-network combination.  Equivalently, one of the finite
matrix inequalities \eqref{eq:s6v66-all-pencils} has a negative
generalized direction.
\end{corollary}

\begin{proof}
If the bounded quadratic form is negative on some unit vector, density
and continuity make it negative on a sufficiently close finite
combination.  Its coefficient vector violates the corresponding pencil.
\end{proof}

\begin{lemma}[Weighted-Haar Riesz stock]
\label{lem:s6v66-Haar-stock}
Let \(\{\phi_n\}\) be an orthonormal invariant spin-network basis in
\(L^2(\mu_{\rm H})\), and let
\(\pi=w\mu_{\rm H}\) with
\[
 0<m\le w\le M<\infty.
\]
Then every finite coefficient family satisfies
\begin{equation}
 m\sum_n|c_n|^2
 \le\left\|\sum_nc_n\phi_n\right\|_{L^2(\pi)}^2
 \le M\sum_n|c_n|^2.
 \label{eq:s6v66-Haar-Riesz}
\end{equation}
\end{lemma}

\begin{proof}
Bound \(w\) above and below in
\[
 \int\left|\sum_nc_n\phi_n\right|^2w\,\dd\mu_{\rm H}
\]
and use Haar orthonormality.
\end{proof}

\begin{assessment}[Conditioning warning]
\label{ass:s6v66-conditioning}
At each compact finite regulator a strictly positive continuous Doob
density has finite \(m,M\), so
\eqref{eq:s6v66-Haar-Riesz} supplies a finite-regulator Riesz stock.
It does not give a cofinal stock: \(m\) may tend to zero and \(M/m\) may
grow exponentially with volume or refinement.  The infinite-pencil
criterion itself avoids inverting this stock, but any truncation with a
quantified tail must still control conditioning.
\end{assessment}

\subsection{Testing the presently available Wilsonian input}
\label{sec:s6v66-test}

\begin{proposition}[What the local action decomposition actually gives]
\label{prop:s6v66-local-action}
If the Wilsonian interpolation head has an absolutely convergent local
expansion \(R_\theta=\sum_ZR_{\theta,Z}\), then
\[
 L_\theta'=\sum_Z{\cal C}_{\theta,Z}
\]
as a bounded quadratic-form sum, where each \({\cal C}_{\theta,Z}\) has
the centered two-arm kernel obtained from
\eqref{eq:s6v66-two-arm} by replacing \(R_\theta\) with
\(R_{\theta,Z}\) and centering its leading score.  Absolute polymer or
marked-path majorants control
\(\sum_Z\|{\cal C}_{\theta,Z}\|\), but do not imply
\({\cal C}_{\theta,Z}\succeq0\).
\end{proposition}

\begin{proof}
The decomposition is the termwise linearity in
\Cref{prop:s6v66-two-arm}.  An absolute majorant bounds the norm of the
sum by the sum of the norms.  Positivity is not preserved by a statement
about absolute values, and the centered bracket in
\eqref{eq:s6v66-two-arm} changes sign unless an additional correlation
inequality is proved.
\end{proof}

\begin{assessment}[Upstream bottleneck after V66]
\label{ass:s6v66-bottleneck}
The cluster-filter bottleneck has been bypassed, but not the signed
inequality.  The current sources provide local summability, boundary
stability, and finite-horizon marked paths.  They do not provide:
\begin{enumerate}[label=\textup{(\roman*)},leftmargin=3.5em]
\item a positive local diagonal stock \(\mu_\theta\) for the centered
      two-arm insertion;
\item a frame lower stock \(A_\theta\) uniform in volume and refinement;
\item an overlap norm \(\kappa_\theta<\mu_\theta\).
\end{enumerate}
Thus \Cref{thm:s6v66-local-frame} is a new exact compiler, not a proof
that the Wilsonian action satisfies its hypotheses.  The next upstream
question is whether the centered two-arm kernel has a covariance or
conditional-variance representation with a positive principal part, or
whether explicit null/negative directions rule out that route.
\end{assessment}

\begin{definition}[Direct signed form card, DSFC]
\label{def:s6v66-DSFC}
For a cofinal regulator family, DSFC consists of:
\begin{enumerate}[label=\textup{(\roman*)},leftmargin=3.5em]
\item a norm-\(C^1\) normalized transfer family with simple vacuum;
\item the exact inserted kernel \eqref{eq:s6v66-two-arm};
\item local analysis maps satisfying
      \eqref{eq:s6v66-lower-frame};
\item a signed local decomposition
      \eqref{eq:s6v66-local-decomposition};
\item cofinal stocks \eqref{eq:s6v66-cofinal-stocks};
\item the already separated V53 Mosco/projectivity hypotheses.
\end{enumerate}
\end{definition}

\begin{theorem}[Conditional mass-gap handoff from DSFC]
\label{thm:s6v66-DSFC-handoff}
If DSFC holds and the V53 form-limit hypotheses identify the limiting OS
Hamiltonian, then that Hamiltonian has a strictly positive
vacuum-orthogonal spectral gap.
\end{theorem}

\begin{proof}
\Cref{thm:s6v66-local-frame} gives DSFI with a uniform positive \(b\).
\Cref{thm:s6v66-garding-gap} gives
\(g_{j,1}\ge m\delta_j\) for one \(m>0\) independent of \(j\).
The V53 Mosco gap-transport theorem passes the rescaled form gap to the
limiting vacuum-orthogonal Hamiltonian form.  Its spectral infimum is
therefore at least the transported positive constant.
\end{proof}

\begin{programme}[V67: covariance test and obstruction search]
\label{prog:s6v66-V67}
\begin{enumerate}[leftmargin=3em]
\item Polarize the centered two-arm form into conditional covariance
      and commutator pieces relative to the half-slab return law.
\item Test whether action terms depending on only finitely many local
      coordinates leave a large kernel in the complete spin-network
      pencil.
\item Construct explicit sign-changing finite pencils when positivity
      fails, instead of weakening the desired inequality by absolute
      values.
\item If a positive principal covariance exists, identify its analysis
      maps and prove the sharp overlap matrix bound needed in
      \Cref{thm:s6v66-local-frame}.
\end{enumerate}
\end{programme}

\begin{firewall}[Claim boundary after sixth-series V66]
\label{fire:s6v66-boundary}
V66 derives the exact whole-carrier derivative and centered two-arm
kernel; proves a direct signed G\aa rding inequality opens an
order-\(\delta\) transfer gap without cluster filtering; proves a
local-frame/overlap sufficient criterion; gives an exact complete
spin-network pencil characterization with finite failure witnesses; and
states the conditional continuum handoff.

V66 does not prove DSFI for the Yang--Mills Wilsonian head, establish
cofinal frame or overlap stocks, close the V53 projective/Mosco
hypotheses, construct the continuum OS theory, or prove the full
Yang--Mills mass gap.
\end{firewall}

\begin{theorem}[Sixth-series V66 result ledger]
\label{thm:s6v66-results}
The positive conclusions in \Cref{fire:s6v66-boundary} are exactly those
of
\Cref{prop:s6v66-derivative-form,
prop:s6v66-two-arm,
thm:s6v66-garding-gap,
cor:s6v66-constant-stock,
thm:s6v66-local-frame,
cor:s6v66-cofinal-stock,
thm:s6v66-pencil,
cor:s6v66-finite-witness,
lem:s6v66-Haar-stock,
prop:s6v66-local-action,
thm:s6v66-DSFC-handoff}.
\end{theorem}

\begin{proof}
Each item is the cited theorem with its finite-regulator, compactness,
isolation, boundedness, clock-scaling, and cofinal-limit qualifications
unchanged.
\end{proof}

\paragraph{Parent-version record.}
V66 was formed from
\texttt{Renewal\_Yang\_Mills\_Mass\_Gap\_Research\_Notes\_V65.tex}.
The V65 parent had \(29\,224\) logical source lines,
\(1\,026\,536\) bytes, and SHA--256
\[
 \texttt{7adea9601872c80a0591b08234d5f1ccf9765ac4410baf4b33eaebc50142fe8d}.
\]
The next compulsory companion audit remains V70.

% =============================================================================
% END APPENDED SIXTH-SERIES V66 MATERIAL
% =============================================================================

% =============================================================================
% BEGIN APPENDED SIXTH-SERIES V67 MATERIAL
% =============================================================================

\section{V67: Doob conditional covariance and the sign obstruction}
\label{sec:s6v67}

\subsection{Ground-transformed probability space}
\label{sec:s6v67-ground}

The signed two-arm form of V66 is exact but its positivity is opaque in
reference-amplitude coordinates.  This note conjugates it by the leading
singular functions.  The resulting formula separates a conditional-mean
multiplier from a within-fibre covariance.  That separation both yields a
new sufficient criterion and proves that the conditional-mean term alone
cannot be the sought positive mechanism under a natural dense-range
hypothesis.

Fix \(\theta\) and suppress it temporarily.  Write
\[
 Au=s_1v,\qquad A^*v=s_1u,\qquad s_1^2=\lambda_1,
\]
with \(u,v>0\) normalized in their reference Hilbert spaces.  Define
\begin{align}
 \pi_0(\dd x)&=u(x)^2\mu_0(\dd x),&
 \pi_{1/2}(\dd y)&=v(y)^2\mu_{1/2}(\dd y),                 \label{eq:s6v67-marginals}\\
 \Lambda(\dd x,\dd y)
 &={u(x)F(y,x)v(y)\over s_1}\,
   \mu_0(\dd x)\mu_{1/2}(\dd y).
 \label{eq:s6v67-Doob-law}
\end{align}

\begin{lemma}[Doob law and conditional expectation]
\label{lem:s6v67-Doob}
\(\Lambda\) is a probability measure with marginals
\(\pi_0,\pi_{1/2}\).  The operator
\begin{equation}
 ({\mathsf P}\varphi)(y)
 :=\mathbb E_\Lambda[\varphi(X)\mid Y=y]
 ={A(u\varphi)(y)\over s_1v(y)}
 \label{eq:s6v67-P}
\end{equation}
is a contraction
\[
 {\mathsf P}:L^2(\pi_0)\longrightarrow L^2(\pi_{1/2})
\]
and preserves constants and means.
\end{lemma}

\begin{proof}
Integrating \eqref{eq:s6v67-Doob-law} in \(y\) and using
\(A^*v=s_1u\) gives \(u^2\mu_0=\pi_0\).  Integrating in \(x\) and using
\(Au=s_1v\) gives \(v^2\mu_{1/2}=\pi_{1/2}\).  Formula
\eqref{eq:s6v67-P} follows by division by the \(Y\)-marginal density.
Conditional Jensen gives the contraction property.  Conditional
expectation fixes \(1\), and the tower property preserves the integral.
\end{proof}

\begin{lemma}[Exact transfer and loss forms after the ground transform]
\label{lem:s6v67-transfer-loss}
For \(f=u\varphi\),
\begin{align}
 \|f\|_{{\cal H}_0}^2
 &=\|\varphi\|_{L^2(\pi_0)}^2,                              \label{eq:s6v67-norm}\\
 \langle f,Tf\rangle
 &=\|{\mathsf P}\varphi\|_{L^2(\pi_{1/2})}^2,               \label{eq:s6v67-T}\\
 \langle f,Lf\rangle
 &=\mathbb E_\Lambda
   \left|\varphi(X)-\mathbb E[\varphi(X)\mid Y]\right|^2.  \label{eq:s6v67-L}
\end{align}
Moreover,
\[
 f\perp u
 \quad\Longleftrightarrow\quad
 \mathbb E_{\pi_0}\varphi=0.
\]
\end{lemma}

\begin{proof}
The norm identity is the definition of \(\pi_0\).  Since
\(A(u\varphi)=s_1v\,{\mathsf P}\varphi\),
\[
 \langle f,Tf\rangle
 ={1\over s_1^2}\|Af\|^2
 =\int|{\mathsf P}\varphi|^2\,\dd\pi_{1/2}.
\]
Subtract this from \(\|f\|^2\) and use the conditional variance identity
to obtain \eqref{eq:s6v67-L}.  Finally
\(\langle u,u\varphi\rangle=\int\varphi\,\dd\pi_0\).
\end{proof}

\subsection{Exact covariance polarization}
\label{sec:s6v67-covariance}

Assume \(R\) is real and square integrable under \(\Lambda\).  Define
\begin{align}
 m(y)&=\mathbb E_\Lambda[R(X,Y)\mid Y=y],                  \label{eq:s6v67-m}\\
 J&=\mathbb E_\Lambda R,                                  \label{eq:s6v67-J}\\
 \operatorname{Cov}_Y(R,\varphi)
 &:=\mathbb E_\Lambda[(R-m(Y))\varphi(X)\mid Y].
 \label{eq:s6v67-cov}
\end{align}

\begin{theorem}[Exact conditional-covariance formula]
\label{thm:s6v67-covariance}
For every \(f=u\varphi\) in the form domain of the insertion,
\begin{align}
 \langle f,L'f\rangle
 &=\mathbb E_{\pi_{1/2}}
 \left[
 (m-J)|{\mathsf P}\varphi|^2
 +\operatorname{Re}\left\{
 \overline{{\mathsf P}\varphi}\,
 \operatorname{Cov}_Y(R,\varphi)
 \right\}
 \right].
 \label{eq:s6v67-covariance-form}
\end{align}
Thus the signed response has exactly two parts:
\begin{enumerate}[label=\textup{(\roman*)},leftmargin=3.5em]
\item a centered conditional-mean multiplier acting on the returned
      amplitude \({\mathsf P}\varphi\);
\item a within-\(Y\)-fibre covariance coupling the returned amplitude to
      the unresolved conditional fluctuation of \(\varphi\).
\end{enumerate}
\end{theorem}

\begin{proof}
The inserted analogue of \eqref{eq:s6v67-P} is
\[
 {A_R(u\varphi)(y)\over s_1v(y)}
 =\mathbb E_\Lambda[R(X,Y)\varphi(X)\mid Y=y].
\]
Insert this and \(A(u\varphi)=s_1v{\mathsf P}\varphi\) into
\Cref{prop:s6v66-derivative-form}.  Since
\[
 \mathbb E[R\varphi\mid Y]-J{\mathsf P}\varphi
 =(m-J){\mathsf P}\varphi+\operatorname{Cov}_Y(R,\varphi),
\]
the displayed formula follows.
\end{proof}

\begin{corollary}[Independent-return triple form]
\label{cor:s6v67-triple}
Let \((X,X',Y)\) have law in which \(Y\sim\pi_{1/2}\) and, conditional
on \(Y\), \(X,X'\) are independent with the conditional law induced by
\(\Lambda\).  Then
\begin{align}
 \langle u\varphi,L'u\varphi\rangle
 ={1\over2}\mathbb E\big[
 \overline{\varphi(X')}\varphi(X)
 \{R(X,Y)+R(X',Y)-2J\}\big].
 \label{eq:s6v67-triple}
\end{align}
\end{corollary}

\begin{proof}
Condition first on \(Y\), use conditional independence, and compare with
\eqref{eq:s6v67-covariance-form}.  Equivalently this is
\Cref{prop:s6v66-two-arm} divided by the Doob marginal factors.
\end{proof}

\begin{remark}[A genuine covariance, but not a variance]
\label{rem:s6v67-not-variance}
The second term in \eqref{eq:s6v67-covariance-form} is a covariance
paired with \({\mathsf P}\varphi\), not a square.  It has no fixed sign.
The first term is multiplication by the centered function \(m-J\), and
also has no fixed sign.  The Doob transform therefore clarifies the
sign problem but does not solve it by a hidden positivity identity.
\end{remark}

\subsection{Sharp conditional-variance control}
\label{sec:s6v67-control}

Put
\begin{align}
 t(\varphi)&=\|{\mathsf P}\varphi\|_{L^2(\pi_{1/2})}^2,
 \label{eq:s6v67-t}\\
 \ell(\varphi)&=
 \|\varphi\|_{L^2(\pi_0)}^2-t(\varphi),
 \label{eq:s6v67-ell}\\
 \sigma^2&=\mathop{\rm ess\,sup}_{y}
 \operatorname{Var}_\Lambda(R\mid Y=y).
 \label{eq:s6v67-sigma}
\end{align}

\begin{lemma}[Covariance debit]
\label{lem:s6v67-covariance-debit}
If \(\sigma<\infty\), then
\begin{equation}
 \left|
 \mathbb E_{\pi_{1/2}}
 \overline{{\mathsf P}\varphi}\,
 \operatorname{Cov}_Y(R,\varphi)
 \right|
 \le \sigma\sqrt{t(\varphi)\ell(\varphi)}.
 \label{eq:s6v67-covariance-debit}
\end{equation}
Consequently, for every \(\eta>0\),
\begin{equation}
 \langle u\varphi,L'u\varphi\rangle
 \ge
 \mathbb E_{\pi_{1/2}}(m-J)|{\mathsf P}\varphi|^2
 -{\eta\over2}\|\varphi\|^2
 -{\sigma^2\over2\eta}\ell(\varphi).
 \label{eq:s6v67-Young}
\end{equation}
\end{lemma}

\begin{proof}
Conditional Cauchy--Schwarz gives
\[
 |\operatorname{Cov}_Y(R,\varphi)|^2
 \le\operatorname{Var}(R\mid Y)
    \operatorname{Var}(\varphi\mid Y).
\]
Multiply by \(|{\mathsf P}\varphi|\), integrate, and use
Cauchy--Schwarz in \(Y\).  The two squared factors integrate to at most
\(\sigma^2t(\varphi)\) and exactly \(\ell(\varphi)\), respectively.
This proves \eqref{eq:s6v67-covariance-debit}.  Young's inequality
\(\sigma\sqrt{t\ell}\le\eta t/2+\sigma^2\ell/(2\eta)\), followed by
\(t\le\|\varphi\|^2\), proves \eqref{eq:s6v67-Young}.
\end{proof}

\begin{definition}[Conditional-mean range inequality]
\label{def:s6v67-CMRI}
The conditional-mean range inequality, abbreviated CMRI, with stocks
\(\beta>0\) and \(\alpha\ge0\), is
\begin{equation}
 \mathbb E_{\pi_{1/2}}
 (m-J)|{\mathsf P}\varphi|^2
 \ge
 \beta\delta\|\varphi\|_{L^2(\pi_0)}^2
 -\alpha\,\ell(\varphi)
 \label{eq:s6v67-CMRI}
\end{equation}
for every mean-zero \(\varphi\).
\end{definition}

\begin{theorem}[CMRI plus sub-clock fibre variance implies DSFI]
\label{thm:s6v67-CMRI-DSFI}
Suppose CMRI holds and
\begin{equation}
 \sigma^2\le s\delta
 \label{eq:s6v67-subclock-variance}
\end{equation}
for some \(s<\infty\).  Then on the whole relative carrier,
\begin{equation}
 QL'Q
 \succeq{\beta\over2}\delta Q
 -\left(\alpha+{s\over2\beta}\right)QLQ.
 \label{eq:s6v67-DSFI}
\end{equation}
Hence the V66 edge-opening theorem applies with
\[
 b={\beta\over2},\qquad
 a=\alpha+{s\over2\beta}.
\]
\end{theorem}

\begin{proof}
Apply \eqref{eq:s6v67-Young} with \(\eta=\beta\delta\), then use CMRI
and \eqref{eq:s6v67-subclock-variance}.  By
\Cref{lem:s6v67-transfer-loss},
\(\|\varphi\|^2=\|u\varphi\|^2\) and
\(\ell(\varphi)=\langle u\varphi,Lu\varphi\rangle\).  Mean-zero
\(\varphi\) parametrizes \(Q{\cal H}_0\), so the resulting form
inequality is \eqref{eq:s6v67-DSFI}.
\end{proof}

\begin{assessment}[Clock interpretation]
\label{ass:s6v67-clock}
The within-fibre standard deviation need only be \(O(\sqrt\delta)\), not
\(O(\delta)\).  Its square is then an \(O(\delta)\) covariance debit that
can be absorbed into the G\aa rding loss with an order-one coefficient.
This reproduces, now in an exact transfer identity, the critical
\(\sqrt\delta\) dependency amplitude found abstractly in V59.
\end{assessment}

\subsection{The centered-multiplier obstruction}
\label{sec:s6v67-obstruction}

\begin{lemma}[A centered nonconstant multiplier has a negative
mean-zero direction]
\label{lem:s6v67-centered-negative}
Let \((Y,\pi)\) be a probability space, let \(c\in L^\infty(\pi)\) be
real with \(\int c\,\dd\pi=0\), and suppose there are \(\varepsilon>0\)
and a measurable set
\[
 N_\varepsilon=\{y:c(y)\le-\varepsilon\}
\]
such that \(L^2(N_\varepsilon,\pi)\) contains a nonzero function
orthogonal to constants.  Then there is \(h\in L^2(\pi)\) with
\(\int h\,\dd\pi=0\) and
\begin{equation}
 \int c|h|^2\,\dd\pi
 \le-\varepsilon\|h\|^2<0.
 \label{eq:s6v67-negative}
\end{equation}
\end{lemma}

\begin{proof}
Choose nonzero \(h\) supported in \(N_\varepsilon\) and orthogonal to
the constant function there.  Extended by zero, it has global mean zero.
On its support \(c\le-\varepsilon\), which gives
\eqref{eq:s6v67-negative}.
\end{proof}

\begin{theorem}[Dense-return no-go for an output-only insertion]
\label{thm:s6v67-output-no-go}
Suppose:
\begin{enumerate}[label=\textup{(\roman*)},leftmargin=3.5em]
\item \(R(X,Y)=r(Y)\) is output-only and \(r\in L^\infty(\pi_{1/2})\);
\item \(r-J\) satisfies the nontrivial negative-set hypothesis of
      \Cref{lem:s6v67-centered-negative};
\item \({\mathsf P}\) maps
      \(L^2_0(\pi_0)\) densely into \(L^2_0(\pi_{1/2})\).
\end{enumerate}
Then \(QL'Q\) has a negative quadratic direction.  In particular there
is no \(b\ge0\) for which
\[
 QL'Q\succeq b\delta Q.
\]
\end{theorem}

\begin{proof}
For an output-only insertion,
\(\operatorname{Cov}_Y(R,\varphi)=0\), so
\[
 \langle u\varphi,L'u\varphi\rangle
 =\int(r-J)|{\mathsf P}\varphi|^2\,\dd\pi_{1/2}.
\]
By \Cref{lem:s6v67-centered-negative} there is a mean-zero \(h\) on
which the bounded multiplier form is strictly negative.  Dense range
provides mean-zero \(\varphi_n\) with
\({\mathsf P}\varphi_n\to h\) in \(L^2(\pi_{1/2})\).  Continuity of the
bounded multiplier form makes it negative for all sufficiently large
\(n\).  Then \(u\varphi_n\in Q{\cal H}_0\) is the required direction.
\end{proof}

\begin{corollary}[Where a positive mechanism must reside]
\label{cor:s6v67-positive-location}
Under the hypotheses of \Cref{thm:s6v67-output-no-go}, a positive direct
response cannot be obtained from the centered conditional mean alone.
For a general Wilsonian insertion, any successful DSFI proof must use at
least one of:
\begin{enumerate}[label=\textup{(\roman*)},leftmargin=3.5em]
\item the within-fibre covariance in
      \eqref{eq:s6v67-covariance-form};
\item a G\aa rding loss \(-aL\) that charges negative returned modes to
      conditional variance;
\item a failure of dense return range caused by a proved structural
      constraint.
\end{enumerate}
\end{corollary}

\begin{proof}
The theorem rules out the covariance-free, loss-free, dense-range
alternative.  The listed mechanisms are precisely the hypotheses used
to escape one of those three conditions.
\end{proof}

\begin{remark}[CMRI is not pointwise positivity]
\label{rem:s6v67-CMRI-not-pointwise}
Since \(\int(m-J)\,\dd\pi_{1/2}=0\), a pointwise lower bound
\(m-J\ge\beta\delta>0\) is impossible.  CMRI can hold only because the
negative multiplier region is controlled by the loss
\(\ell(\varphi)\), or because the returned range avoids it.  This is the
precise upstream inequality to seek; a pointwise sign argument is
circular or false.
\end{remark}

\subsection{A finite certificate for the new obstruction}
\label{sec:s6v67-certificate}

\begin{definition}[Conditional sign certificate, CSC]
\label{def:s6v67-CSC}
At a finite regulator, a CSC consists of a finite projected
spin-network combination \(f=u\varphi\in Q{\cal H}_0\) and certified
bounds on
\[
 \mathbb E(m-J)|{\mathsf P}\varphi|^2,\quad
 \mathbb E\operatorname{Re}
 \{\overline{{\mathsf P}\varphi}\operatorname{Cov}_Y(R,\varphi)\},
 \quad
 \ell(\varphi),
\]
which either verify or violate a proposed DSFI constant pair \((a,b)\).
\end{definition}

\begin{proposition}[Completeness of finite CSC failures]
\label{prop:s6v67-CSC-complete}
If a bounded proposed DSFI fails at a finite regulator, then for every
tolerance there is a CSC on a finite spin-network combination whose
margin differs from the true negative margin by less than that
tolerance.
\end{proposition}

\begin{proof}
\Cref{cor:s6v66-finite-witness} supplies a finite projected
spin-network witness.  Apply the unitary multiplication map
\(\varphi\mapsto u\varphi\) and evaluate its derivative form by
\Cref{thm:s6v67-covariance}.  Bounded-form continuity gives arbitrary
quantitative tolerance.
\end{proof}

\begin{assessment}[Exact frontier after V67]
\label{ass:s6v67-frontier}
The signed insertion has now been reduced to two explicit probabilistic
stocks.  A sufficient route is:
\[
 \boxed{\text{CMRI with }\beta>0,\quad
        \operatorname*{ess\,sup}\operatorname{Var}(R\mid Y)=O(\delta),
        \quad\int a_\theta\,\dd\theta=O(1).}
\]
The second condition is compatible with the V59 critical clock, but no
current source proves it for the full Wilsonian head.  More importantly,
the centered-multiplier no-go shows that positivity cannot come from a
conditional-mean sign alone.  V68 should derive an exact commutator or
gradient representation of the within-fibre covariance and test whether
the Yang--Mills local action supplies the needed positive contribution,
with negative regions charged to the conditional-variance loss.
\end{assessment}

\begin{programme}[V68: fibre-gradient representation]
\label{prog:s6v67-V68}
\begin{enumerate}[leftmargin=3em]
\item Express
      \(\operatorname{Cov}_Y(R,\varphi)\) through the conditional
      generator or a Helffer--Sj\"ostrand inverse when such a structure
      is available.
\item Separate the symmetric positive gradient term from commutators
      created by gauge constraints and Wilson-loop interactions.
\item Determine whether the covariance term is intrinsically signed by
      constructing a two-state and a compact-group test model.
\item State the exact additional conditional Poincar\'e or curvature
      estimate needed for a cofinal CMRI proof.
\end{enumerate}
\end{programme}

\begin{firewall}[Claim boundary after sixth-series V67]
\label{fire:s6v67-boundary}
V67 constructs the exact Doob probability law; identifies the normalized
transfer loss with conditional variance; proves the exact
conditional-covariance polarization of \(L'\); proves its sharp
\(\sqrt{t\ell}\) debit; derives CMRI plus an \(O(\delta)\) conditional
variance stock as a sufficient DSFI criterion; and proves a dense-return
sign obstruction for nonconstant output-only insertions.

V67 does not prove CMRI or the conditional variance stock for the
Yang--Mills Wilsonian head, prove that its return operator has dense
range, establish cofinal DSFI, close the V53 projective/Mosco hypotheses,
construct the continuum OS theory, or prove the full Yang--Mills mass
gap.
\end{firewall}

\begin{theorem}[Sixth-series V67 result ledger]
\label{thm:s6v67-results}
The positive conclusions in \Cref{fire:s6v67-boundary} are exactly those
of
\Cref{lem:s6v67-Doob,
lem:s6v67-transfer-loss,
thm:s6v67-covariance,
cor:s6v67-triple,
lem:s6v67-covariance-debit,
thm:s6v67-CMRI-DSFI,
lem:s6v67-centered-negative,
thm:s6v67-output-no-go,
cor:s6v67-positive-location,
prop:s6v67-CSC-complete}.
\end{theorem}

\begin{proof}
Each item is the cited result with its positivity, integrability,
dense-range, finite-regulator, clock, and cofinal qualifications
unchanged.
\end{proof}

\paragraph{Parent-version record.}
V67 was formed from
\texttt{Renewal\_Yang\_Mills\_Mass\_Gap\_Research\_Notes\_V66.tex}.
The V66 parent had \(29\,809\) logical source lines,
\(1\,047\,024\) bytes, and SHA--256
\[
 \texttt{b234f7ab93a4245159933f94a1c1378648ea8839d0e42e71f3795330c8eb50e1}.
\]
The next compulsory companion audit remains V70.

% =============================================================================
% END APPENDED SIXTH-SERIES V67 MATERIAL
% =============================================================================

% =============================================================================
% BEGIN APPENDED SIXTH-SERIES V68 MATERIAL
% =============================================================================

\section{V68: fibre block geometry and an exactly soluble sign model}
\label{sec:s6v68}

\subsection{The conditional fluctuation embedding}
\label{sec:s6v68-embedding}

V67 proved that the only possible new positive mechanism, beyond a
G\aa rding loss, is the within-return-fibre covariance.  We now expose
the geometry of that term.  Continue on the Doob probability space of
\Cref{lem:s6v67-Doob}.  Disintegrate
\[
 \Lambda(\dd x,\dd y)
 =\nu_y(\dd x)\pi_{1/2}(\dd y)
\]
and define the mean-zero fibre
\[
 {\cal K}_y=L^2_0(\nu_y).
\]
For \(\varphi\in L^2(\pi_0)\), put
\begin{align}
 g_\varphi(y)&=({\mathsf P}\varphi)(y),&
 \zeta_\varphi(y,x)&=\varphi(x)-g_\varphi(y).
 \label{eq:s6v68-g-zeta}
\end{align}
Then \(\zeta_\varphi(y,\cdot)\in{\cal K}_y\) for almost every \(y\).

\begin{lemma}[Isometric return--fluctuation embedding]
\label{lem:s6v68-isometry}
The map
\begin{equation}
 {\mathfrak J}\varphi
 :=(g_\varphi,\zeta_\varphi)
 \label{eq:s6v68-J}
\end{equation}
is an isometry from \(L^2(\pi_0)\) into
\begin{equation}
 {\cal K}
 :=L^2(\pi_{1/2})
 \oplus\int_Y^\oplus{\cal K}_y\,\pi_{1/2}(\dd y).
 \label{eq:s6v68-K}
\end{equation}
Its range \({\cal R}={\mathfrak J}L^2(\pi_0)\) is closed, and
\begin{align}
 \|\varphi\|^2&=\|g_\varphi\|^2+\|\zeta_\varphi\|^2,        \label{eq:s6v68-Pythagoras}\\
 \langle u\varphi,Lu\varphi\rangle&=\|\zeta_\varphi\|^2.  \label{eq:s6v68-loss}
\end{align}
\end{lemma}

\begin{proof}
Conditional orthogonality gives, for almost every \(y\),
\[
 \mathbb E[|\varphi(X)|^2\mid Y=y]
 =|g_\varphi(y)|^2+
   \|\zeta_\varphi(y,\cdot)\|_{{\cal K}_y}^2.
\]
Integrate in \(y\) to obtain \eqref{eq:s6v68-Pythagoras}.
Thus \({\mathfrak J}\) is an isometry and has closed range.
Equation \eqref{eq:s6v68-loss} is
\Cref{lem:s6v67-transfer-loss}.
\end{proof}

Let
\begin{align}
 c(y)&=m(y)-J,&
 r_y(x)&=R(x,y)-m(y),&
 v(y)&=\|r_y\|_{{\cal K}_y}^2
      =\operatorname{Var}(R\mid Y=y).
 \label{eq:s6v68-c-r-v}
\end{align}

\begin{theorem}[Exact fibre block compression]
\label{thm:s6v68-block}
On the ambient direct integral \({\cal K}\), define the measurable
self-adjoint fibre block
\begin{equation}
 {\mathbb C}_y
 =
 \begin{pmatrix}
 c(y)&\frac12\langle r_y,\,\cdot\,\rangle_{{\cal K}_y}\\
 \frac12 r_y&0
 \end{pmatrix}.
 \label{eq:s6v68-block}
\end{equation}
Whenever this field defines a bounded decomposable operator
\({\mathbb C}\), the signed transfer derivative is its compression to the
compatibility range:
\begin{equation}
 \boxed{
 \langle u\varphi,L'u\varphi\rangle
 =\langle{\mathfrak J}\varphi,
 {\mathbb C}{\mathfrak J}\varphi\rangle_{\cal K}.}
 \label{eq:s6v68-compression}
\end{equation}
Moreover the loss form is the compression of
\[
 {\mathbb D}_y=
 \begin{pmatrix}0&0\\0&I_{{\cal K}_y}\end{pmatrix}.
\]
\end{theorem}

\begin{proof}
Since \(r_y\) has conditional mean zero,
\[
 \operatorname{Cov}_Y(R,\varphi)
 =\langle r_y,\zeta_\varphi(y,\cdot)\rangle_{{\cal K}_y},
\]
with the inner-product convention adjusted by complex conjugation in the
first entry.  The quadratic form of
\eqref{eq:s6v68-block} is therefore
\[
 c(y)|g_\varphi(y)|^2+
 \operatorname{Re}\left\{
 \overline{g_\varphi(y)}
 \operatorname{Cov}_Y(R,\varphi)(y)\right\}.
\]
Integration and \Cref{thm:s6v67-covariance} give
\eqref{eq:s6v68-compression}.  The loss statement is
\eqref{eq:s6v68-loss}.
\end{proof}

\begin{assessment}[Compatibility is the only possible rescue]
\label{ass:s6v68-compatibility}
The ambient block \({\mathbb C}_y\) is indefinite whenever \(r_y\ne0\):
its lower-right diagonal is zero and its off-diagonal entry is nonzero.
The physical pairs \((g,\zeta)\), however, do not fill the ambient direct
integral; they lie in the closed compatibility range \({\cal R}\).
Consequently any positive result must use a quantitative geometric
property of \({\cal R}\), not merely fibrewise Cauchy--Schwarz.
\end{assessment}

\subsection{The exact Schur debit}
\label{sec:s6v68-Schur}

\begin{lemma}[Fibrewise completion of the square]
\label{lem:s6v68-square}
For \(a>0\), \(g\in\mathbb C\), and \(\zeta\in{\cal K}_y\),
\begin{align}
 &c(y)|g|^2+
 \operatorname{Re}\{\overline g\langle r_y,\zeta\rangle\}
 +a\|\zeta\|^2                                      \notag\\
 &\qquad
 =a\left\|\zeta+{g\over2a}r_y\right\|^2
 +\left(c(y)-{v(y)\over4a}\right)|g|^2.
 \label{eq:s6v68-square}
\end{align}
\end{lemma}

\begin{proof}
Expand the square and use \(\|r_y\|^2=v(y)\).
\end{proof}

\begin{corollary}[Pointwise ambient G\aa rding criterion]
\label{cor:s6v68-pointwise}
Let \(a>b\delta\).  The ambient block inequality
\[
 {\mathbb C}+a{\mathbb D}\succeq b\delta I_{\cal K}
\]
holds if and only if, almost everywhere,
\begin{equation}
 c(y)-b\delta
 -{v(y)\over4(a-b\delta)}
 \ge0.
 \label{eq:s6v68-pointwise}
\end{equation}
If it holds, its compression proves DSFI.
\end{corollary}

\begin{proof}
Subtract \(b\delta(|g|^2+\|\zeta\|^2)\) and apply
\Cref{lem:s6v68-square} with \(a-b\delta\).  The resulting square is
nonnegative, and minimizing it in \(\zeta\) shows that
\eqref{eq:s6v68-pointwise} is also necessary for the ambient inequality.
Compression to \({\cal R}\) preserves a positive inequality.
\end{proof}

\begin{proposition}[Pointwise ambient positivity is impossible for a
nontrivial centered head]
\label{prop:s6v68-pointwise-no-go}
Suppose \(c=m-J\) is integrable,
\(\int c\,\dd\pi_{1/2}=0\), and is not zero almost everywhere.  Then
\eqref{eq:s6v68-pointwise} cannot hold for any \(b\ge0\) and finite
\(a>b\delta\).
\end{proposition}

\begin{proof}
The rightmost two terms subtracted in
\eqref{eq:s6v68-pointwise} are nonnegative, so that inequality would
imply \(c(y)\ge0\) almost everywhere.  A nonnegative integrable function
with zero integral vanishes almost everywhere, contrary to the
hypothesis.
\end{proof}

\begin{corollary}[No proof by unrestricted fibre relaxation]
\label{cor:s6v68-no-relaxation}
A DSFI proof cannot replace the physical compatibility range
\({\cal R}\) by the full direct integral and then apply a pointwise Schur
bound.  Such a relaxation erases exactly the correlation between
returned amplitudes and fibre fluctuations that would have to create
the mass response.
\end{corollary}

\begin{proof}
The relaxed proof would require the ambient inequality in
\Cref{cor:s6v68-pointwise}, which
\Cref{prop:s6v68-pointwise-no-go} rules out for a nontrivial centered
head.
\end{proof}

\subsection{An exactly soluble binary return channel}
\label{sec:s6v68-binary}

Let \(X,Y\in\{-1,+1\}\), both with uniform marginal law, and let
\begin{equation}
 \Lambda_\rho(x,y)={1\over4}(1+\rho xy),
 \qquad |\rho|<1.
 \label{eq:s6v68-binary-law}
\end{equation}
This is the Doob law of the positive amplitude
\[
 F_\rho(y,x)=1+\rho xy
\]
on the two uniform reference spaces, with \(u=v=1\) and \(s_1=1\).
The nonconstant input mode \(\varphi(x)=x\) spans the complete
vacuum-orthogonal carrier.

\begin{lemma}[Binary transfer data]
\label{lem:s6v68-binary-transfer}
For the law \eqref{eq:s6v68-binary-law},
\begin{align}
 {\mathsf P}\varphi(y)&=\rho y,                            \label{eq:s6v68-binary-P}\\
 \langle\varphi,T\varphi\rangle&=\rho^2,                  \label{eq:s6v68-binary-T}\\
 \langle\varphi,L\varphi\rangle&=1-\rho^2.                \label{eq:s6v68-binary-L}
\end{align}
\end{lemma}

\begin{proof}
The conditional law is
\(\mathbb P(X=x\mid Y=y)=\frac12(1+\rho xy)\), hence
\(\mathbb E[X\mid Y=y]=\rho y\).  The last two identities follow from
\Cref{lem:s6v67-transfer-loss} and \(\|\varphi\|=1\).
\end{proof}

Choose the real insertion
\begin{equation}
 R_{\alpha,\beta}(x,y)=\alpha xy+\beta y.
 \label{eq:s6v68-binary-R}
\end{equation}

\begin{theorem}[Covariance controls the sign at the clock scale]
\label{thm:s6v68-binary-sign}
For \eqref{eq:s6v68-binary-R},
\begin{align}
 J&=\alpha\rho,                                           \label{eq:s6v68-binary-J}\\
 m(y)-J&=\beta y,                                         \label{eq:s6v68-binary-c}\\
 \operatorname{Cov}_Y(R,\varphi)(y)
 &=\alpha y(1-\rho^2),                                    \label{eq:s6v68-binary-cov}\\
 \operatorname{Var}(R\mid Y=y)
 &=\alpha^2(1-\rho^2),                                    \label{eq:s6v68-binary-var}\\
 \boxed{\langle\varphi,L'\varphi\rangle
 =\alpha\rho(1-\rho^2).}                                  \label{eq:s6v68-binary-response}
\end{align}
The output-only coefficient \(\beta\) makes no contribution to the
relative response.  The input--output coupling \(\alpha\) makes the
response positive or negative according to the sign of
\(\alpha\rho\).
\end{theorem}

\begin{proof}
Given \(Y=y\),
\[
 \mathbb E[XY\mid Y=y]=\rho,\qquad
 \mathbb E[R\mid Y=y]=\alpha\rho+\beta y,
\]
which proves \eqref{eq:s6v68-binary-J} and
\eqref{eq:s6v68-binary-c}.  Since the \(\beta y\) term is
conditionally constant,
\begin{align*}
 \operatorname{Cov}_Y(R,X)
 &=\alpha y\,
   \mathbb E[(X-\rho y)X\mid Y=y]\\
 &=\alpha y(1-\rho^2),
\end{align*}
and the same Bernoulli computation gives
\eqref{eq:s6v68-binary-var}.  Insert
\eqref{eq:s6v68-binary-P}--\eqref{eq:s6v68-binary-cov} in
\Cref{thm:s6v67-covariance}.  The centered-multiplier term is
\(\mathbb E[\beta Y\,\rho^2]=0\), while the covariance term is
\(\alpha\rho(1-\rho^2)\).
\end{proof}

\begin{corollary}[Exact critical scaling in the binary model]
\label{cor:s6v68-binary-clock}
If
\[
 1-\rho_\delta^2=\mu\delta+o(\delta),\qquad
 \rho_\delta\longrightarrow1,
\]
and \(\alpha_\delta\to\alpha_*>0\), then
\begin{align}
 \langle\varphi,L_\delta'\varphi\rangle
 &=\alpha_*\mu\delta+o(\delta),                            \label{eq:s6v68-binary-clock-response}\\
 \operatorname{Var}(R\mid Y)
 &=\alpha_*^2\mu\delta+o(\delta).                          \label{eq:s6v68-binary-clock-var}
\end{align}
Thus a covariance mechanism can simultaneously have a positive
order-\(\delta\) gap response and the V67 sub-clock fibre-variance stock.
\end{corollary}

\begin{proof}
Substitute the assumed asymptotics in
\eqref{eq:s6v68-binary-response} and
\eqref{eq:s6v68-binary-var}.
\end{proof}

\begin{assessment}[What the binary model proves and does not prove]
\label{ass:s6v68-binary}
The model proves that the within-fibre covariance is not merely an error:
an input--output score aligned with the transmitted mode produces the
correct positive clock-scale response, whereas reversing its sign closes
the gap.  It does not show that the Yang--Mills Wilsonian score is aligned
with every low relative mode.  That is now the precise nonabelian
spectral question.
\end{assessment}

\subsection{The compatibility-angle target}
\label{sec:s6v68-angle}

For \(v(y)>0\), define the normalized score direction
\[
 e_y={r_y\over\sqrt{v(y)}}\in{\cal K}_y
\]
and the aligned fluctuation coefficient
\begin{equation}
 a_\varphi(y)=\langle e_y,\zeta_\varphi(y,\cdot)\rangle.
 \label{eq:s6v68-alignment}
\end{equation}
Then
\begin{equation}
 \operatorname{Cov}_Y(R,\varphi)
 =\sqrt{v}\,a_\varphi.
 \label{eq:s6v68-cov-alignment}
\end{equation}

\begin{definition}[Positive compatibility-angle inequality]
\label{def:s6v68-PCAI}
PCAI with stocks \(\gamma>0\) and \(a_0\ge0\) is
\begin{equation}
 \mathbb E_{\pi_{1/2}}
 \left[
 c|g_\varphi|^2+
 \sqrt v\,\operatorname{Re}
 (\overline{g_\varphi}a_\varphi)
 \right]
 \ge
 \gamma\delta\|\varphi\|^2-a_0\|\zeta_\varphi\|^2
 \label{eq:s6v68-PCAI}
\end{equation}
for every mean-zero \(\varphi\).
\end{definition}

\begin{proposition}[PCAI is exactly DSFI in conditional coordinates]
\label{prop:s6v68-PCAI}
PCAI holds if and only if
\[
 QL'Q\succeq\gamma\delta Q-a_0QLQ.
\]
\end{proposition}

\begin{proof}
Use \Cref{thm:s6v68-block} for the left side,
\eqref{eq:s6v68-Pythagoras} for the norm, and
\eqref{eq:s6v68-loss} for the loss.
\end{proof}

\begin{assessment}[Exact upstream target after V68]
\label{ass:s6v68-frontier}
The abstract covariance route is no longer ambiguous.  Fibrewise
positivity is impossible, while compression to the physical
compatibility range can be positive and has the correct clock scaling.
For Yang--Mills one must prove PCAI uniformly on the gauge-invariant
relative carrier.  Equivalently, one must show that every weakly lost
mode has a positively aligned score fluctuation.  Modes with large
misalignment may be charged to \(L\), but modes with
\(\|\zeta_\varphi\|^2=O(\delta)\|\varphi\|^2\) must have the positive
alignment explicitly.  Neither the area-law macrocell estimate nor an
absolute polymer majorant supplies this signed angle.
\end{assessment}

\begin{programme}[V69: low-loss compatibility and score alignment]
\label{prog:s6v68-V69}
\begin{enumerate}[leftmargin=3em]
\item Restrict PCAI to the low-loss cone
      \(\|\zeta_\varphi\|^2\le C\delta\|\varphi\|^2\), since its
      complement can be paid for by the G\aa rding loss.
\item Derive an exact low-loss/high-loss cone lemma with optimized
      constants.
\item Test Wilson plaquette insertions on the bilocal gauge-invariant
      soft modes constructed in V61.
\item Decide whether locality forces the alignment to vanish on
      delocalized infrared packets; if so, abandon interpolation opening
      and return to a genuinely nonperturbative transfer inequality.
\end{enumerate}
\end{programme}

\begin{firewall}[Claim boundary after sixth-series V68]
\label{fire:s6v68-boundary}
V68 gives the isometric return--fluctuation embedding; represents the
signed derivative as a fibre-block compression; proves the exact Schur
debit and the impossibility of unrestricted pointwise fibre positivity;
solves a binary return model in which covariance produces either sign
and the correct \(O(\delta)\) response; and identifies PCAI as precisely
the DSFI target in conditional coordinates.

V68 does not prove PCAI for Yang--Mills, establish uniform score
alignment on its low-loss cone, close the V53 projective/Mosco
hypotheses, construct the continuum OS theory, or prove the full
Yang--Mills mass gap.
\end{firewall}

\begin{theorem}[Sixth-series V68 result ledger]
\label{thm:s6v68-results}
The positive conclusions in \Cref{fire:s6v68-boundary} are exactly those
of
\Cref{lem:s6v68-isometry,
thm:s6v68-block,
lem:s6v68-square,
cor:s6v68-pointwise,
prop:s6v68-pointwise-no-go,
cor:s6v68-no-relaxation,
lem:s6v68-binary-transfer,
thm:s6v68-binary-sign,
cor:s6v68-binary-clock,
prop:s6v68-PCAI}.
\end{theorem}

\begin{proof}
Each item is the cited result with its boundedness, disintegration,
compatibility-range, finite-model, clock-scaling, and cofinal
qualifications unchanged.
\end{proof}

\paragraph{Parent-version record.}
V68 was formed from
\texttt{Renewal\_Yang\_Mills\_Mass\_Gap\_Research\_Notes\_V67.tex}.
The V67 parent had \(30\,301\) logical source lines,
\(1\,064\,104\) bytes, and SHA--256
\[
 \texttt{734fcaa1214f0575d0a66eca6aa6eb7ac7a62a99a9001ebe4b5c29b0a56005d7}.
\]
The next compulsory companion audit remains V70.

% =============================================================================
% END APPENDED SIXTH-SERIES V68 MATERIAL
% =============================================================================

% =============================================================================
% BEGIN APPENDED SIXTH-SERIES V69 MATERIAL
% =============================================================================

\section{V69: low-loss cone reduction and directional score alignment}
\label{sec:s6v69}

\subsection{Why only the low-loss cone is difficult}
\label{sec:s6v69-cone}

V68 identified the physical compatibility-angle inequality, but stated it
on the whole relative carrier.  The G\aa rding loss already controls
directions on which \(L\) is large.  The genuinely difficult set is the
near-isometric cone of the return channel.  This note makes that
reduction exact and then places the V61 gauge-singlet infrared packets
inside the cone.

Fix a regulator, interpolation parameter, and step \(\delta>0\).  Write
\begin{equation}
 {\cal C}(f)=\langle f,L'f\rangle,\qquad
 {\cal L}(f)=\langle f,Lf\rangle
 \label{eq:s6v69-C-L}
\end{equation}
for \(f\in Q{\cal H}_0\).

\begin{definition}[Low-loss cone]
\label{def:s6v69-cone}
For \(K>0\), define
\begin{equation}
 {\mathfrak C}_{K,\delta}
 =\left\{f\in Q{\cal H}_0:
 {\cal L}(f)\le K\delta\|f\|^2\right\}.
 \label{eq:s6v69-cone}
\end{equation}
\end{definition}

\begin{theorem}[Low-loss cone compiler]
\label{thm:s6v69-cone}
Suppose there are constants \(K,\gamma,M>0\) such that
\begin{align}
 {\cal C}(f)&\ge\gamma\delta\|f\|^2
 &&(f\in{\mathfrak C}_{K,\delta}),                         \label{eq:s6v69-low-positive}\\
 {\cal C}(f)&\ge-M\delta\|f\|^2
 &&(f\in Q{\cal H}_0).                                    \label{eq:s6v69-global-lower}
\end{align}
Then for every \(b\) with \(0<b\le\gamma\),
\begin{equation}
 \boxed{
 QL'Q\succeq b\delta Q-aQLQ,\qquad
 a={M+b\over K}.}
 \label{eq:s6v69-cone-DSFI}
\end{equation}
\end{theorem}

\begin{proof}
If \(f\in{\mathfrak C}_{K,\delta}\), then
\[
 {\cal C}(f)\ge\gamma\delta\|f\|^2
 \ge b\delta\|f\|^2-a{\cal L}(f).
\]
If \(f\notin{\mathfrak C}_{K,\delta}\), then
\[
 b\delta\|f\|^2-a{\cal L}(f)
 <\left(b-aK\right)\delta\|f\|^2
 =-M\delta\|f\|^2
 \le{\cal C}(f).
\]
The two cases cover the relative carrier.
\end{proof}

\begin{corollary}[Cofinal low-loss sufficient condition]
\label{cor:s6v69-cofinal-cone}
If \eqref{eq:s6v69-low-positive}--\eqref{eq:s6v69-global-lower}
hold uniformly in volume, refinement, and \(\theta\), with
\(K,\gamma,M\) independent of the regulator, then the V66 integrated
edge-opening theorem yields a regulator-uniform positive rescaled
endpoint gap.
\end{corollary}

\begin{proof}
\Cref{thm:s6v69-cone} gives DSFI with uniform
\(b=\gamma\) and \(a=(M+\gamma)/K\).  Substitute these constants in
\Cref{cor:s6v66-constant-stock}.
\end{proof}

\begin{assessment}[A necessary clock-scale lower control]
\label{ass:s6v69-global-clock}
The global lower bound \eqref{eq:s6v69-global-lower} must be of order
\(\delta\).  If only
\({\cal C}(f)\ge-M\|f\|^2\) were known, the same proof would require
\(a\asymp\delta^{-1}\); the integrating factor in
\Cref{thm:s6v66-garding-gap} would then erase the endpoint lower bound.
Thus absolute boundedness of the action score is not enough.  One needs
a normalized-transfer cancellation or relative-form estimate at the
physical clock scale.
\end{assessment}

\subsection{A volume-robust directional covariance stock}
\label{sec:s6v69-directional}

In the V67 conditional coordinates, define
\begin{equation}
 ({\mathsf K}\varphi)(y)
 :=\operatorname{Cov}_Y(R,\varphi)(y).
 \label{eq:s6v69-K}
\end{equation}
It annihilates constants.  The essential bound on
\(\operatorname{Var}(R\mid Y)\) used in V67 can grow with volume when
\(R\) is an extensive action score.  The following relative operator
bound is the correct directional replacement.

\begin{definition}[Score-fluctuation relative bound]
\label{def:s6v69-SFRB}
SFRB with stock \(s\ge0\) is
\begin{equation}
 \boxed{
 \|{\mathsf K}\varphi\|_{L^2(\pi_{1/2})}^2
 \le s\delta\,\ell(\varphi)}
 \label{eq:s6v69-SFRB}
\end{equation}
for every mean-zero \(\varphi\), where
\(\ell(\varphi)=\|\varphi\|^2-\|{\mathsf P}\varphi\|^2\).
\end{definition}

\begin{lemma}[Directional covariance debit]
\label{lem:s6v69-directional-debit}
Under SFRB,
\begin{equation}
 \left|\langle{\mathsf P}\varphi,{\mathsf K}\varphi\rangle\right|
 \le\sqrt{s\delta\,t(\varphi)\ell(\varphi)},
 \qquad
 t(\varphi)=\|{\mathsf P}\varphi\|^2.
 \label{eq:s6v69-directional-debit}
\end{equation}
\end{lemma}

\begin{proof}
Cauchy--Schwarz gives
\[
 |\langle{\mathsf P}\varphi,{\mathsf K}\varphi\rangle|
 \le\sqrt{t(\varphi)}\|{\mathsf K}\varphi\|.
\]
Apply \eqref{eq:s6v69-SFRB}.
\end{proof}

\begin{theorem}[CMRI plus SFRB implies DSFI]
\label{thm:s6v69-CMRI-SFRB}
Suppose the conditional-mean range inequality
\eqref{eq:s6v67-CMRI} holds with stocks \((\beta,\alpha)\), and SFRB
holds with stock \(s\).  Then
\begin{equation}
 QL'Q
 \succeq{\beta\over2}\delta Q
 -\left(\alpha+{s\over2\beta}\right)QLQ.
 \label{eq:s6v69-CMRI-SFRB}
\end{equation}
\end{theorem}

\begin{proof}
By \Cref{thm:s6v67-covariance} and
\Cref{lem:s6v69-directional-debit},
\[
 \langle u\varphi,L'u\varphi\rangle
 \ge
 \mathbb E(m-J)|{\mathsf P}\varphi|^2
 -\sqrt{s\delta\,t\ell}.
\]
Young's inequality at scale \(\beta\delta\) gives
\[
 \sqrt{s\delta\,t\ell}
 \le{\beta\delta\over2}t+{s\over2\beta}\ell.
\]
Use CMRI and \(t\le\|\varphi\|^2\), then conjugate back by
\(f=u\varphi\) using \Cref{lem:s6v67-transfer-loss}.
\end{proof}

\begin{corollary}[SFRB is weaker than the V67 variance stock]
\label{cor:s6v69-SFRB-weaker}
If
\[
 \mathop{\rm ess\,sup}_y\operatorname{Var}(R\mid Y=y)
 \le s\delta,
\]
then SFRB holds with the same \(s\).  The converse need not hold.
\end{corollary}

\begin{proof}
The conditional Cauchy--Schwarz step in
\Cref{lem:s6v67-covariance-debit} gives
\[
 |{\mathsf K}\varphi(y)|^2
 \le\operatorname{Var}(R\mid Y=y)
    \operatorname{Var}(\varphi\mid Y=y).
\]
Integrate and use the essential bound.  SFRB controls only the action of
the score covariance on admissible input functions; it does not bound
the full fibre norm of the extensive score.
\end{proof}

\subsection{Product-channel test of extensivity}
\label{sec:s6v69-product}

Take \(N\) independent copies of the binary return channel of V68, with
common correlation \(0<\rho<1\), and let
\begin{equation}
 \varphi_N(X)={1\over\sqrt N}\sum_{i=1}^NX_i,\qquad
 R_N(X,Y)=\alpha\sum_{i=1}^NX_iY_i.
 \label{eq:s6v69-product-data}
\end{equation}

\begin{theorem}[An extensive local score has a finite collective
response]
\label{thm:s6v69-product}
For \eqref{eq:s6v69-product-data},
\begin{align}
 \|\varphi_N\|^2&=1,                                      \label{eq:s6v69-product-norm}\\
 {\mathsf P}\varphi_N
 &={\rho\over\sqrt N}\sum_{i=1}^NY_i,                     \label{eq:s6v69-product-P}\\
 \ell(\varphi_N)&=1-\rho^2,                               \label{eq:s6v69-product-loss}\\
 {\mathsf K}\varphi_N
 &={\alpha(1-\rho^2)\over\sqrt N}\sum_{i=1}^NY_i,          \label{eq:s6v69-product-K}\\
 \langle\varphi_N,L'\varphi_N\rangle
 &=\alpha\rho(1-\rho^2),                                  \label{eq:s6v69-product-response}\\
 \operatorname{Var}(R_N\mid Y)
 &=N\alpha^2(1-\rho^2).                                   \label{eq:s6v69-product-score-var}
\end{align}
Thus the total conditional score variance diverges linearly in volume,
while the normalized delocalized mode has a volume-independent positive
response and satisfies
\begin{equation}
 {\|{\mathsf K}\varphi_N\|^2\over\ell(\varphi_N)}
 =\alpha^2(1-\rho^2).
 \label{eq:s6v69-product-relative}
\end{equation}
\end{theorem}

\begin{proof}
Independence and centering give \(\|\varphi_N\|^2=1\).
Conditionally on \(Y\),
\(\mathbb E[X_i\mid Y]=\rho Y_i\), proving
\eqref{eq:s6v69-product-P} and hence
\eqref{eq:s6v69-product-loss}.  Conditional cross-covariances vanish,
while the diagonal calculation of V68 gives
\[
 \operatorname{Cov}_Y(X_iY_i,X_j)
 =\delta_{ij}Y_i(1-\rho^2).
\]
This proves \eqref{eq:s6v69-product-K}.  Pair it with
\eqref{eq:s6v69-product-P} and use
\(\mathbb E(\sum_iY_i)^2=N\) to obtain
\eqref{eq:s6v69-product-response}.  Conditional independence also gives
\eqref{eq:s6v69-product-score-var}.  Taking the norm in
\eqref{eq:s6v69-product-K} yields
\(\alpha^2(1-\rho^2)^2\), and division by
\eqref{eq:s6v69-product-loss} gives
\eqref{eq:s6v69-product-relative}.
\end{proof}

\begin{assessment}[Locality is compatible with a collective lift]
\label{ass:s6v69-collective}
The product model rules out a tempting but false shortcut: an extensive
sum of local action insertions need not have vanishing matrix element on
a normalized delocalized mode.  Translation-coherent local
contributions add before normalization and can leave an order-one
generator response.  At the same time,
\eqref{eq:s6v69-product-score-var} shows why a volume-uniform essential
variance bound is the wrong stock.  SFRB retains the directional
cancellation needed by a collective physical mode.
\end{assessment}

\subsection{The V61 infrared packets are decisive low-loss tests}
\label{sec:s6v69-infrared}

\begin{theorem}[Soft-packet alignment consequence of DSFI]
\label{thm:s6v69-soft-packet}
Let \(f_{j,R}\in Q_{j,0}{\cal H}_{j,0}\) be normalized regulator
realizations of the V61 gauge-singlet packet at the quadratic endpoint.
Assume
\begin{equation}
 \limsup_{j\to\infty}
 {{\cal L}_{j,0}(f_{j,R})\over\delta_j}
 \le {C\over R}+o_R(1).
 \label{eq:s6v69-soft-loss}
\end{equation}
If DSFI holds at \(\theta=0\) with cofinal constants \(a<\infty\) and
\(b>0\), then
\begin{equation}
 \liminf_{R\to\infty}\liminf_{j\to\infty}
 {{\cal C}_{j,0}(f_{j,R})\over\delta_j}
 \ge b.
 \label{eq:s6v69-soft-response}
\end{equation}
\end{theorem}

\begin{proof}
DSFI gives
\[
 {{\cal C}_{j,0}(f_{j,R})\over\delta_j}
 \ge b-a\,{{\cal L}_{j,0}(f_{j,R})\over\delta_j}.
\]
Take the regulator lower limit, use
\eqref{eq:s6v69-soft-loss}, and then let \(R\to\infty\).
\end{proof}

\begin{corollary}[Soft-response no-go certificate]
\label{cor:s6v69-soft-no-go}
If a noncollapsing cofinal packet family satisfies
\eqref{eq:s6v69-soft-loss} and
\[
 \lim_{R\to\infty}\liminf_{j\to\infty}
 {{\cal C}_{j,0}(f_{j,R})\over\delta_j}=0,
\]
then no DSFI with \(b>0\) and bounded \(a\) holds at the quadratic
endpoint.
\end{corollary}

\begin{proof}
This is the contrapositive of
\Cref{thm:s6v69-soft-packet}.
\end{proof}

\begin{corollary}[Conditional-angle form of the soft test]
\label{cor:s6v69-soft-angle}
Writing \(f_{j,R}=u_{j,0}\varphi_{j,R}\), a positive DSFI requires
\begin{align}
 \liminf_{R\to\infty}\liminf_{j\to\infty}{1\over\delta_j}
 \mathbb E\big[
 c_{j,0}|g_{j,R}|^2+
 \sqrt{v_{j,0}}\,
 \operatorname{Re}(\overline{g_{j,R}}a_{j,R})
 \big]\ge b,
 \label{eq:s6v69-soft-angle}
\end{align}
where \(g,a,c,v\) are the compatibility variables of V68.
\end{corollary}

\begin{proof}
Use \Cref{prop:s6v68-PCAI} and the fact that the loss term divided by
\(\delta_j\) vanishes after \(j\to\infty\) and then \(R\to\infty\).
\end{proof}

\begin{assessment}[Connection with the V61 nonlinear lift]
\label{ass:s6v69-lift}
Equation \eqref{eq:s6v69-soft-response} is the infinitesimal,
fixed-reference version of the integrated nonlinear lift required in
\Cref{thm:s6v61-necessary-lift}.  It shows where dimensional
transmutation must appear in the present interpolation route: the
conditional score must retain a positive order-\(\delta\) alignment on
arbitrarily large, exactly gauge-invariant colour-singlet packets whose
quadratic loss tends to zero.  Absolute locality estimates neither prove
nor refute that signed alignment.
\end{assessment}

\subsection{Low-loss alignment card}
\label{sec:s6v69-card}

\begin{definition}[Low-loss alignment card, LLAC]
\label{def:s6v69-LLAC}
For a cofinal regulator family, LLAC consists of:
\begin{enumerate}[label=\textup{(\roman*)},leftmargin=3.5em]
\item a fixed \(K>0\) and the low-loss cones
      \(\mathfrak C_{j,\theta;K,\delta_j}\);
\item the uniform signed lower bound
      \eqref{eq:s6v69-low-positive} with \(\gamma>0\);
\item the uniform global clock lower bound
      \eqref{eq:s6v69-global-lower};
\item a direct spin-network verification or CMRI--SFRB proof of the
      low-loss bound;
\item the V61 packet test \eqref{eq:s6v69-soft-angle};
\item the V53 projective/Mosco rows.
\end{enumerate}
\end{definition}

\begin{theorem}[Conditional mass-gap handoff from LLAC]
\label{thm:s6v69-LLAC}
LLAC implies a regulator-uniform endpoint rescaled transfer gap and,
under the V53 limit hypotheses, a positive gap for the limiting OS
Hamiltonian.
\end{theorem}

\begin{proof}
\Cref{thm:s6v69-cone} proves DSFI.
\Cref{thm:s6v66-garding-gap} integrates it.
The V53 Mosco handoff transports the uniform rescaled form inequality to
the limiting Hamiltonian.
\end{proof}

\begin{assessment}[Exact frontier after V69]
\label{ass:s6v69-frontier}
The whole-carrier problem is reduced to two clock-scaled statements: a
positive signed response on the low-loss cone and a global negative-part
bound.  The extensive-score variance bottleneck has been sharpened to
the directional SFRB.  The first unavoidable Yang--Mills diagnostic is
now \eqref{eq:s6v69-soft-angle} on the exact V61 bilocal packets.
Evaluating it requires a signed connected Wilsonian coefficient, not an
absolute polymer majorant.  V70 must first perform the compulsory SM/NS
audit and then determine whether either companion programme has acquired
a directional covariance or low-loss-cone estimate that changes this
route.
\end{assessment}

\begin{programme}[V70: companion audit and low-loss transfer]
\label{prog:s6v69-V70}
\begin{enumerate}[leftmargin=3em]
\item Audit the latest SM and NS notes with line counts, hashes, and
      targeted searches for conditional covariance, relative form
      bounds, low-loss cones, and signed connected coefficients.
\item Import only proved estimates with matching hypotheses and clocks.
\item Test any new estimate against LLAC and the V61 soft-packet
      inequality.
\item If no companion estimate fills the signed row, retain the present
      route and isolate the smallest Wilson-network coefficient whose
      sign would decide it.
\end{enumerate}
\end{programme}

\begin{firewall}[Claim boundary after sixth-series V69]
\label{fire:s6v69-boundary}
V69 proves that a positive low-loss-cone response plus a global
clock-scale lower bound implies DSFI; replaces a volume-divergent total
score variance by the directional SFRB; proves CMRI plus SFRB implies
DSFI; solves an extensive product model showing coherent local scores
can lift a delocalized mode; and proves that the V61 infrared packets are
necessary signed alignment tests.

V69 does not prove the low-loss response, global clock lower bound,
CMRI, or SFRB for Yang--Mills; evaluate the V61 signed connected
coefficient; close the V53 projective/Mosco hypotheses; construct the
continuum OS theory; or prove the full Yang--Mills mass gap.
\end{firewall}

\begin{theorem}[Sixth-series V69 result ledger]
\label{thm:s6v69-results}
The positive conclusions in \Cref{fire:s6v69-boundary} are exactly those
of
\Cref{thm:s6v69-cone,
cor:s6v69-cofinal-cone,
lem:s6v69-directional-debit,
thm:s6v69-CMRI-SFRB,
cor:s6v69-SFRB-weaker,
thm:s6v69-product,
thm:s6v69-soft-packet,
cor:s6v69-soft-no-go,
cor:s6v69-soft-angle,
thm:s6v69-LLAC}.
\end{theorem}

\begin{proof}
Each item is the cited result with its low-loss, clock, directional,
finite-regulator, packet-landing, and cofinal qualifications unchanged.
\end{proof}

\paragraph{Parent-version record.}
V69 was formed from
\texttt{Renewal\_Yang\_Mills\_Mass\_Gap\_Research\_Notes\_V68.tex}.
The V68 parent had \(30\,760\) logical source lines,
\(1\,079\,283\) bytes, and SHA--256
\[
 \texttt{51c60ddeea9bf627397ec05fead958b0d3399eccc90f8a769dd983fd91c8e3b5}.
\]
The next compulsory companion audit is V70.

% =============================================================================
% END APPENDED SIXTH-SERIES V69 MATERIAL
% =============================================================================

% =============================================================================
% BEGIN APPENDED SIXTH-SERIES V70 MATERIAL
% =============================================================================

\section{V70: compulsory companion audit and the loss-factorization target}
\label{sec:s6v70}

\subsection{Audited files and integrity record}
\label{sec:s6v70-files}

The two current companion notes in \texttt{latex\_notes} were inspected
read-only before continuing the Yang--Mills argument:
\begin{align}
 &\texttt{Renewal\_SM\_Einstein\_Continuum\_Research\_Notes\_V14.tex},
 \notag\\
 &\quad 4\,979\ {\rm lines},\quad184\,018\ {\rm bytes},\quad
 \texttt{0eb5b4e34244f97e4a080624bbf80fc602fe0983dff5289edcc8a686f753b4a3},
 \label{eq:s6v70-SM-file}\\
 &\texttt{Renewal\_NS\_Stability\_Research\_Notes\_V26.tex},
 \notag\\
 &\quad 5\,319\ {\rm lines},\quad278\,283\ {\rm bytes},\quad
 \texttt{82c887f8c2c69e4f28fad7c3dc8de5b9099cb5a4bf431eba1fff3e91935978ad}.
 \label{eq:s6v70-NS-file}
\end{align}
The NS file and digest are unchanged from the V65 audit.  The SM file has
advanced from V11 to V14.

\begin{definition}[V70 audit question]
\label{def:s6v70-question}
The audit asks only whether a companion theorem proves one of the current
Yang--Mills rows:
\begin{enumerate}[label=\textup{(\roman*)},leftmargin=3.5em]
\item a positive signed response on the low-loss cone;
\item a global order-\(\delta\) lower bound for \(L'\);
\item the score-fluctuation relative bound SFRB;
\item projective/Mosco landing of the resulting transfer forms.
\end{enumerate}
Similarity of vocabulary is not counted as an import.
\end{definition}

\subsection{What is new in SM V12--V14}
\label{sec:s6v70-SM}

\paragraph{SM V12.}
For bounded finite-regulator Hamiltonian interactions, conditional
oscillations obey a Lieb--Robinson estimate controlled by the exponential
influence row
\[
 k_\mu=\sup_i\sum_j e^{\mu d(i,j)}
 4\sum_{Z\ni i,j}\|h_Z\|.
\]
Two Hamiltonian insertions satisfy a connected three-set tree bound.  The
proof uses two complementary commutator estimates and Jacobi's identity;
the spatial exponent is reduced from \(\mu\) to \(\mu/2\).

\paragraph{SM V13.}
For bounded local Lindblad insertions, which have no Jacobi identity, the
same kind of tree geometry is obtained by decomposing an insertion into a
persistent old oscillation and a localized new bump.  This removes the
earlier support-reset assumption but remains a finite-horizon absolute
majorant.

\paragraph{SM V14.}
For nonnegative generators \(H_0,H_1\) with common closed form domain and
relative form difference
\[
 |(q_1-q_0)[z,z]|\le a q_0[z,z]+b\|z\|^2,\qquad a<1,
\]
SM V14 proves the weak Duhamel estimate
\begin{equation}
 |\langle z,(e^{-tH_1}-e^{-tH_0})w\rangle|
 \le t\gamma_\lambda M_\lambda
 \|z\|_{{\cal Q},0}\|w\|_{{\cal Q},0},
 \label{eq:s6v70-SM-Duhamel}
\end{equation}
with explicit form-equivalence constants.  It passes this estimate
through Mosco-convergent bounded truncations when the transported form
norms stay bounded.  Two form insertions additionally require an
integrable \({\cal Q}^*\to{\cal Q}\) smoothing kernel.

\begin{proposition}[Exact use of the SM import]
\label{prop:s6v70-SM-use}
The SM V12--V14 results can supply:
\begin{enumerate}[label=\textup{(\roman*)},leftmargin=3.5em]
\item absolute connected-tree decay for two bounded local marked
      insertions at fixed physical horizon;
\item an order-\(t\) weak comparison for prescribed vectors with
      uniformly bounded form energy;
\item a conditional truncation-to-form-limit passage for those matrix
      elements.
\end{enumerate}
They do not imply any item of \Cref{def:s6v70-question} without additional
Yang--Mills-specific sign, energy-moment, frame, and relative-loss
estimates.
\end{proposition}

\begin{proof}
The tree estimates bound norms and hence lose the sign required in the
PCAI expression \eqref{eq:s6v68-PCAI}.  The weak Duhamel estimate
\eqref{eq:s6v70-SM-Duhamel} is clock-scaled only on a form-norm-bounded
set, whereas the V69 global lower row is required on the whole relative
carrier.  Its Mosco theorem assumes, rather than proves, the relative
form constants and bounded transported energy moments.  Finally none of
SM V12--V14 identifies the Yang--Mills score covariance
\({\mathsf K}\), proves its factorization through the Yang--Mills loss,
or provides a positive diagonal response.  Therefore each possible use
is conditional on an additional row listed in the statement.
\end{proof}

\begin{firewall}[Finite horizon is not the return spectral edge]
\label{fire:s6v70-finite-horizon}
The connected-tree estimates are valuable for summing off-diagonal local
matrix elements of a covariance operator.  They are not an equilibrium
mixing estimate and do not control the smallest relative eigenvalue of
the reflected transfer.  Importing their exponential spatial decay as a
positive mass would repeat the locality/gap conflation excluded in V60
and in the SM notes themselves.
\end{firewall}

\subsection{Low transfer loss does not supply the SM energy moment}
\label{sec:s6v70-energy}

\begin{proposition}[Low-loss/form-energy separation]
\label{prop:s6v70-loss-energy}
Let \(H\ge0\) have orthonormal vacuum-orthogonal eigenvectors
\(e_0,e_n\) with
\[
 He_0=E_0e_0,\qquad He_n=E_ne_n,\qquad E_n\longrightarrow\infty.
\]
Fix \(0<\delta<1\), set \(L_\delta=I-e^{-\delta H}\), and define
\[
 f_n=\sqrt{1-\delta}\,e_0+\sqrt\delta\,e_n.
\]
Then
\begin{equation}
 \langle f_n,L_\delta f_n\rangle
 \le\delta(E_0+1)
 \label{eq:s6v70-low-loss}
\end{equation}
uniformly in \(n\), while
\begin{equation}
 \langle f_n,Hf_n\rangle
 =(1-\delta)E_0+\delta E_n\longrightarrow\infty.
 \label{eq:s6v70-high-energy}
\end{equation}
Thus membership in a fixed V69 low-loss cone does not imply a uniform
SM V14 form-energy moment.
\end{proposition}

\begin{proof}
Functional calculus and \(1-e^{-x}\le x\) give
\begin{align*}
 \langle f_n,L_\delta f_n\rangle
 &=(1-\delta)(1-e^{-\delta E_0})
   +\delta(1-e^{-\delta E_n})\\
 &\le\delta E_0+\delta.
\end{align*}
The Hamiltonian expectation is
\eqref{eq:s6v70-high-energy} and diverges with \(E_n\).
\end{proof}

\begin{corollary}[Required refinement of the cone]
\label{cor:s6v70-energy-cone}
The weak SM V14 estimate can be used uniformly on the hard V69 cone only
after either:
\begin{enumerate}[label=\textup{(\roman*)},leftmargin=3.5em]
\item intersecting it with a cofinally bounded form-energy ball and
      controlling the complement separately; or
\item proving a Yang--Mills ultraviolet inequality that bounds the
      relevant form energy by loss plus a uniform local moment.
\end{enumerate}
Neither condition follows from low loss alone.
\end{corollary}

\begin{proof}
\Cref{prop:s6v70-loss-energy} rules out the missing implication.  Each
listed alternative supplies exactly the form-norm boundedness used on the
right side of \eqref{eq:s6v70-SM-Duhamel}.
\end{proof}

\subsection{NS V26 decision}
\label{sec:s6v70-NS}

NS V26 computes pointwise rank-one norms for the first two derivatives of
the Oseen kernel and compares them with symmetric-traceless tensor norms.
Its first-derivative refinement is empty; its second-derivative
rank-one norm improves the earlier bound modestly on an intermediate
radial range.  These are deterministic kernel constants in an ancient
Navier--Stokes profile argument.

\begin{theorem}[V70 companion-audit decision]
\label{thm:s6v70-audit}
No companion theorem proves a Yang--Mills low-loss positive response,
global clock lower bound, SFRB, or continuum mass-gap handoff.
The SM tree estimates may be used later as an absolute off-diagonal
summation device, and the SM weak Duhamel estimate may be used on
form-energy-bounded Yang--Mills writers.  Nothing is imported from NS
V26.
\end{theorem}

\begin{proof}
\Cref{prop:s6v70-SM-use,prop:s6v70-loss-energy} gives the SM decision.
The NS quantities concern Euclidean derivatives and rank-one contractions
of the Oseen kernel.  They contain no reflected transfer, Doob law,
Yang--Mills score, gauge-invariant carrier, or conditional covariance, so
they do not populate any row of \Cref{def:s6v70-question}.
\end{proof}

\subsection{The operator factorization hidden in SFRB}
\label{sec:s6v70-factorization}

The audit suggests using the SM tree estimates only after identifying the
correct operator to which they should apply.  We now give that
identification exactly.  Work on the mean-zero input space
\[
 {\cal H}_{0,\circ}=L^2_0(\pi_0).
\]
Let
\begin{equation}
 D=I-{\mathsf P}^*{\mathsf P}\succeq0,
 \label{eq:s6v70-D}
\end{equation}
so that
\(\ell(\varphi)=\langle\varphi,D\varphi\rangle\), and retain the
score-covariance operator
\({\mathsf K}\varphi=\operatorname{Cov}_Y(R,\varphi)\).

\begin{theorem}[Equivalent forms of the score-fluctuation relative bound]
\label{thm:s6v70-factorization}
Assume \({\mathsf K}:{\cal H}_{0,\circ}\to L^2(\pi_{1/2})\) is bounded.
For \(s\ge0\), the following are equivalent:
\begin{enumerate}[label=\textup{(\roman*)},leftmargin=3.5em]
\item
\(\|{\mathsf K}\varphi\|^2\le s\delta
\langle\varphi,D\varphi\rangle\) for all \(\varphi\);
\item
\begin{equation}
 {\mathsf K}^*{\mathsf K}\preceq s\delta D;
 \label{eq:s6v70-operator-SFRB}
\end{equation}
\item for every \(h\in L^2(\pi_{1/2})\) and
      \(\varphi\in{\cal H}_{0,\circ}\),
\begin{equation}
 |\langle h,{\mathsf K}\varphi\rangle|^2
 \le s\delta\|h\|^2\langle\varphi,D\varphi\rangle;
 \label{eq:s6v70-dual-SFRB}
\end{equation}
\item there is a contraction
\[
 {\mathsf C}:\overline{\operatorname{Ran}D^{1/2}}
 \longrightarrow L^2(\pi_{1/2})
\]
such that
\begin{equation}
 \boxed{
 {\mathsf K}=\sqrt{s\delta}\,{\mathsf C}D^{1/2}.}
 \label{eq:s6v70-factorization}
\end{equation}
\end{enumerate}
\end{theorem}

\begin{proof}
The equality
\(\|{\mathsf K}\varphi\|^2
=\langle\varphi,{\mathsf K}^*{\mathsf K}\varphi\rangle\)
shows that \textup{(i)} and \textup{(ii)} are equivalent.  Taking the
supremum of
\(|\langle h,{\mathsf K}\varphi\rangle|\) over unit \(h\) proves the
equivalence of \textup{(i)} and \textup{(iii)}.

If \textup{(iv)} holds, contraction of \({\mathsf C}\) gives
\textup{(i)}.  Conversely assume \textup{(i)}.  Define initially on
\(\operatorname{Ran}D^{1/2}\)
\[
 {\mathsf C}_0(D^{1/2}\varphi)
 ={1\over\sqrt{s\delta}}{\mathsf K}\varphi
\]
when \(s\delta>0\).  If \(D^{1/2}\varphi=0\), \textup{(i)} gives
\({\mathsf K}\varphi=0\), so the definition is independent of the
representative.  Again by \textup{(i)},
\(\|{\mathsf C}_0D^{1/2}\varphi\|\le\|D^{1/2}\varphi\|\).
It therefore extends uniquely to a contraction on the closure of the
range and proves \eqref{eq:s6v70-factorization}.  The case \(s=0\) is
\({\mathsf K}=0\).
\end{proof}

\begin{corollary}[Exact-return modes must annihilate the score covariance]
\label{cor:s6v70-null}
SFRB implies
\begin{equation}
 \ker(I-{\mathsf P}^*{\mathsf P})\subseteq\ker{\mathsf K}.
 \label{eq:s6v70-null}
\end{equation}
More generally the covariance output is Lipschitz in the square root of
the prediction loss, uniformly at scale \(\sqrt\delta\).
\end{corollary}

\begin{proof}
If \(D\varphi=0\), the factorization
\eqref{eq:s6v70-factorization} gives \({\mathsf K}\varphi=0\).
The norm statement is item~\textup{(i)} of
\Cref{thm:s6v70-factorization}.
\end{proof}

\begin{assessment}[How the SM tree bound could enter]
\label{ass:s6v70-tree-use}
Equation \eqref{eq:s6v70-dual-SFRB} is a two-marked matrix-element
estimate: one mark is the centered action score and the other is the
output test \(h\).  The SM V12--V13 tree technology can plausibly control
its off-diagonal local entries after a gauge-invariant local frame is
chosen.  It still does not provide the essential factor \(D^{1/2}\), the
\(\sqrt\delta\) clock, or the sign in CMRI.  Those must arise from the
reflected return law.  This prevents importing a spatial tree majorant as
if it were already SFRB.
\end{assessment}

\subsection{Audit-adjusted next target}
\label{sec:s6v70-next}

\begin{definition}[Loss-factorization card, LFC]
\label{def:s6v70-LFC}
At a finite regulator and interpolation parameter, LFC consists of:
\begin{enumerate}[label=\textup{(\roman*)},leftmargin=3.5em]
\item a gauge-invariant local frame for the input loss operator
      \(D=I-{\mathsf P}^*{\mathsf P}\);
\item a frame for the output of the score covariance \({\mathsf K}\);
\item a connected two-marked estimate for the matrix of
      \({\mathsf K}\), with the V70 SM tree bound used only for
      off-diagonal summation;
\item a diagonal return estimate proving
      \({\mathsf K}^*{\mathsf K}\preceq s\delta D\);
\item cofinal control of the frame constants and \(s\).
\end{enumerate}
\end{definition}

\begin{programme}[V71: exact-return nullspace and local factorization]
\label{prog:s6v70-V71}
\begin{enumerate}[leftmargin=3em]
\item Characterize \(\ker D\) through equality in conditional Jensen:
      these are input functions measurable through the return variable.
\item Prove the necessary cancellation
      \(\ker D\subseteq\ker{\mathsf K}\) directly from the Wilsonian
      score, or construct a counterexample.
\item On the orthogonal complement, seek a local square-root
      factorization of \({\mathsf K}\) through the conditional residual
      \(\varphi-\mathsf P\varphi\).
\item Use connected-tree decay only after the diagonal factorization is
      established.
\item Keep the positive CMRI/soft-packet alignment as a separate signed
      row.
\end{enumerate}
\end{programme}

\begin{firewall}[Claim boundary after sixth-series V70]
\label{fire:s6v70-boundary}
V70 performs the compulsory SM/NS audit; records the new SM connected
two-marked locality and weak form-Duhamel tools; proves that low transfer
loss does not imply the form-energy moment those tools require; imports no
unsupported gap claim; and proves SFRB is exactly a Douglas-type
factorization of the score covariance through the square root of the
return loss.

V70 does not prove that factorization, CMRI, the low-loss positive
response, or the global clock lower bound for Yang--Mills; close the V53
projective/Mosco hypotheses; construct the continuum OS theory; or prove
the full Yang--Mills mass gap.
\end{firewall}

\begin{theorem}[Sixth-series V70 result ledger]
\label{thm:s6v70-results}
The positive conclusions in \Cref{fire:s6v70-boundary} are exactly those
of
\Cref{prop:s6v70-SM-use,
prop:s6v70-loss-energy,
cor:s6v70-energy-cone,
thm:s6v70-audit,
thm:s6v70-factorization,
cor:s6v70-null}.
\end{theorem}

\begin{proof}
Each item is the cited result with its finite-horizon, form-energy,
boundedness, clock, directional, gauge, and cofinal qualifications
unchanged.
\end{proof}

\paragraph{Parent-version record.}
V70 was formed from
\texttt{Renewal\_Yang\_Mills\_Mass\_Gap\_Research\_Notes\_V69.tex}.
The V69 parent had \(31\,226\) logical source lines,
\(1\,095\,222\) bytes, and SHA--256
\[
 \texttt{14850424cb5682fc9ce2c8a3728e4991b90d62c8529db89ff1e2e15a94162b41}.
\]
The next compulsory companion audit is V75.

% =============================================================================
% END APPENDED SIXTH-SERIES V70 MATERIAL
% =============================================================================

% =============================================================================
% BEGIN APPENDED SIXTH-SERIES V71 MATERIAL
% =============================================================================

\section{V71: physical residual factorization of the score covariance}
\label{sec:s6v71}

\subsection{The residual operator}
\label{sec:s6v71-residual}

V70 identified SFRB with an abstract Douglas factorization through
\(D^{1/2}\), where \(D=I-\mathsf P^*\mathsf P\).  The Doob conditional
law contains a canonical square root of this operator.  Making it explicit
separates an automatic algebraic cancellation from the genuinely hard
uniform norm estimate.

Let
\[
 {\cal K}_{\rm fib}
 =\int_Y^\oplus L^2_0(\nu_y)\,\pi_{1/2}(\dd y)
\]
be the conditional mean-zero fibre space of V68, and define
\begin{equation}
 ({\mathsf Z}\varphi)(y,x)
 :=\varphi(x)-({\mathsf P}\varphi)(y).
 \label{eq:s6v71-Z}
\end{equation}

\begin{theorem}[The return loss has an exact residual square root]
\label{thm:s6v71-residual}
\({\mathsf Z}:L^2(\pi_0)\to{\cal K}_{\rm fib}\) is a contraction and
\begin{equation}
 \boxed{
 {\mathsf Z}^*{\mathsf Z}
 =I-{\mathsf P}^*{\mathsf P}=D.}
 \label{eq:s6v71-ZstarZ}
\end{equation}
Consequently its polar decomposition is
\begin{equation}
 {\mathsf Z}={\mathsf U}D^{1/2},
 \label{eq:s6v71-polar}
\end{equation}
where \({\mathsf U}\) is a partial isometry from
\(\overline{\operatorname{Ran}D^{1/2}}\) onto
\(\overline{\operatorname{Ran}{\mathsf Z}}\).
\end{theorem}

\begin{proof}
For \(\varphi,\psi\in L^2(\pi_0)\), conditional orthogonality gives
\begin{align*}
 \langle{\mathsf Z}\varphi,{\mathsf Z}\psi\rangle
 &=\langle\varphi,\psi\rangle_{L^2(\pi_0)}
 -\langle{\mathsf P}\varphi,
          {\mathsf P}\psi\rangle_{L^2(\pi_{1/2})}\\
 &=\langle\varphi,
 (I-{\mathsf P}^*{\mathsf P})\psi\rangle.
\end{align*}
This proves \eqref{eq:s6v71-ZstarZ}; positivity and
\(\|{\mathsf P}\|\le1\) show that \(\|{\mathsf Z}\|\le1\).
The standard polar decomposition of a bounded operator gives
\eqref{eq:s6v71-polar}.
\end{proof}

\begin{corollary}[Equality in conditional Jensen]
\label{cor:s6v71-Jensen}
For \(\varphi\in L^2(\pi_0)\), the following are equivalent:
\begin{enumerate}[label=\textup{(\roman*)},leftmargin=3.5em]
\item \(\langle\varphi,D\varphi\rangle=0\);
\item \({\mathsf Z}\varphi=0\);
\item
\begin{equation}
 \varphi(X)=({\mathsf P}\varphi)(Y)
 \quad\Lambda\text{-almost surely};
 \label{eq:s6v71-common-info}
\end{equation}
\item equality holds in
\(\|{\mathsf P}\varphi\|\le\|\varphi\|\).
\end{enumerate}
Thus \(\ker D\) is the common information measurable from both the input
and the return variable.
\end{corollary}

\begin{proof}
\Cref{thm:s6v71-residual} makes \textup{(i)}, \textup{(ii)}, and
\textup{(iv)} equivalent.  Definition \eqref{eq:s6v71-Z} makes
\textup{(ii)} and \textup{(iii)} equivalent.
\end{proof}

\begin{proposition}[Strictly positive kernels have no nonconstant exact
return mode]
\label{prop:s6v71-positive-kernel}
Suppose the Doob law \(\Lambda\) is equivalent to the product
\(\pi_0\otimes\pi_{1/2}\).  Then
\begin{equation}
 \ker D=\operatorname{span}\{1\}.
 \label{eq:s6v71-kernel-D}
\end{equation}
In particular \(D\) is injective on the mean-zero carrier.
\end{proposition}

\begin{proof}
If \(D\varphi=0\), \Cref{cor:s6v71-Jensen} gives measurable \(g\) with
\(\varphi(x)=g(y)\) for \(\Lambda\)-almost every \((x,y)\).  Equivalence
to the product law makes the same equality hold for
\(\pi_0\otimes\pi_{1/2}\)-almost every pair.  Fubini then supplies
full-measure sets of \(x\)'s and \(y\)'s on which any two admissible
values are equal through a common pair.  Hence both \(\varphi\) and \(g\)
are almost surely constant.  Constants plainly lie in \(\ker D\).
\end{proof}

\begin{remark}[Injectivity is not a gap]
\label{rem:s6v71-no-gap}
At a fixed positivity-improving compact regulator the amplitude is
strictly positive and \eqref{eq:s6v71-kernel-D} is typical.  It gives no
cofinal lower bound on \(D\).  In the near-critical regime the smallest
relative eigenvalue of \(D\) is precisely the closing transfer gap.
Using injectivity as coercivity would therefore be circular.
\end{remark}

\subsection{The physical score-readout operator}
\label{sec:s6v71-readout}

Retain
\[
 r_y(x)=R(x,y)-\mathbb E[R\mid Y=y]
\]
and define on \(\operatorname{Ran}{\mathsf Z}\)
\begin{equation}
 ({\mathsf R}_{\rm phys}{\mathsf Z}\varphi)(y)
 :=\langle r_y,
 ({\mathsf Z}\varphi)(y,\cdot)\rangle_{L^2(\nu_y)}.
 \label{eq:s6v71-Rphys}
\end{equation}

\begin{theorem}[Automatic loss-null cancellation and exact
factorization]
\label{thm:s6v71-score-factor}
The definition \eqref{eq:s6v71-Rphys} is independent of the
representative \(\varphi\), and
\begin{equation}
 \boxed{
 {\mathsf K}
 ={\mathsf R}_{\rm phys}{\mathsf Z}
 ={\mathsf R}_{\rm phys}{\mathsf U}D^{1/2}}
 \label{eq:s6v71-score-factor}
\end{equation}
on its natural domain.  In particular,
\begin{equation}
 \ker D\subseteq\ker{\mathsf K}
 \label{eq:s6v71-automatic-null}
\end{equation}
holds for every square-integrable score, without an additional
Yang--Mills hypothesis.
\end{theorem}

\begin{proof}
If \({\mathsf Z}\varphi=0\), then
\(\varphi(X)={\mathsf P}\varphi(Y)\) conditionally and
\[
 \operatorname{Cov}_Y(R,\varphi)
 =\mathbb E[(R-m(Y)){\mathsf P}\varphi(Y)\mid Y]=0.
\]
Thus two representatives with the same residual have the same covariance
output, proving well-definedness.  For arbitrary \(\varphi\),
\[
 ({\mathsf R}_{\rm phys}{\mathsf Z}\varphi)(y)
 =\mathbb E[(R-m(Y))
 (\varphi-{\mathsf P}\varphi(Y))\mid Y=y]
 =({\mathsf K}\varphi)(y).
\]
This gives the first factorization.  Insert
\eqref{eq:s6v71-polar} for the second.  If \(D\varphi=0\), then
\({\mathsf Z}\varphi=0\), proving
\eqref{eq:s6v71-automatic-null}.
\end{proof}

\begin{theorem}[SFRB is a restricted score-readout norm]
\label{thm:s6v71-restricted-norm}
SFRB with stock \(s\) holds if and only if
\({\mathsf R}_{\rm phys}\) extends continuously to
\(\overline{\operatorname{Ran}{\mathsf Z}}\) and
\begin{equation}
 \boxed{
 \|{\mathsf R}_{\rm phys}\|_
 {\overline{\operatorname{Ran}{\mathsf Z}}\to L^2(\pi_{1/2})}
 \le\sqrt{s\delta}.}
 \label{eq:s6v71-restricted-norm}
\end{equation}
Equivalently,
\({\mathsf R}_{\rm phys}{\mathsf U}\) is the contraction factor in
\Cref{thm:s6v70-factorization}, up to the scale \(\sqrt{s\delta}\).
\end{theorem}

\begin{proof}
For every \(\varphi\),
\[
 \|{\mathsf K}\varphi\|
 =\|{\mathsf R}_{\rm phys}{\mathsf Z}\varphi\|,
 \qquad
 \|{\mathsf Z}\varphi\|^2
 =\ell(\varphi).
\]
Therefore \eqref{eq:s6v69-SFRB} is exactly the assertion that the
densely defined readout in \eqref{eq:s6v71-Rphys} is bounded by
\(\sqrt{s\delta}\) on its range.  Such a bound is equivalent to its
unique continuous extension to the closure.  The polar factorization
\eqref{eq:s6v71-score-factor} identifies the Douglas factor.
\end{proof}

\begin{corollary}[Why the extensive fibre variance may be harmless]
\label{cor:s6v71-extensive}
The unrestricted fibre readout
\[
 \zeta\longmapsto\langle r_y,\zeta_y\rangle
\]
has norm
\(\mathop{\rm ess\,sup}_y\sqrt{\operatorname{Var}(R\mid Y=y)}\), which
may diverge with volume.  SFRB requires only its restriction to
\(\overline{\operatorname{Ran}{\mathsf Z}}\) to have norm
\(O(\sqrt\delta)\).  The product-channel calculation of V69 is an
explicit example in which a collective physical residual has the latter
scaling while the total score variance is extensive.
\end{corollary}

\begin{proof}
The norm of a fibre functional is \(\|r_y\|\), giving the unrestricted
norm.  The restricted statement is
\Cref{thm:s6v71-restricted-norm}.  The claimed example is
\eqref{eq:s6v69-product-score-var} compared with
\eqref{eq:s6v69-product-relative}.
\end{proof}

\subsection{An exact finite-pencil criterion for SFRB}
\label{sec:s6v71-pencil}

Let \(\{\phi_n\}_{n\ge1}\) have dense span in
\(L^2_0(\pi_0)\).  For \(N\ge1\), define the residual and score Gram
matrices
\begin{align}
 [D^{(N)}]_{mn}
 &=\langle{\mathsf Z}\phi_m,{\mathsf Z}\phi_n\rangle
 =\langle\phi_m,D\phi_n\rangle,                            \label{eq:s6v71-D-pencil}\\
 [S^{(N)}]_{mn}
 &=\langle{\mathsf K}\phi_m,{\mathsf K}\phi_n\rangle.
 \label{eq:s6v71-S-pencil}
\end{align}

\begin{theorem}[Loss-covariance pencil]
\label{thm:s6v71-pencil}
If \({\mathsf K}\) is bounded, SFRB holds with stock \(s\) if and only if
\begin{equation}
 \boxed{
 S^{(N)}\preceq s\delta D^{(N)}
 \quad\hbox{for every }N.}
 \label{eq:s6v71-pencil}
\end{equation}
The matrices may be singular; no inverse or condition number is needed.
\end{theorem}

\begin{proof}
For \(c\in\mathbb C^N\) and
\(\varphi=\sum_{n=1}^Nc_n\phi_n\),
\[
 c^*S^{(N)}c=\|{\mathsf K}\varphi\|^2,\qquad
 c^*D^{(N)}c=\|{\mathsf Z}\varphi\|^2.
\]
Thus SFRB implies every finite pencil.  Conversely the pencils give SFRB
on the dense algebraic span.  Boundedness of \({\mathsf K}\) and
\({\mathsf Z}\) passes the inequality to its closure.
\end{proof}

\begin{corollary}[Finite failure witness]
\label{cor:s6v71-pencil-failure}
If SFRB fails for a proposed \(s\), then some finite spin-network
combination violates \eqref{eq:s6v71-pencil}.  Conversely every such
violation disproves that stock at the given regulator.
\end{corollary}

\begin{proof}
The forward statement follows by approximating a strict violating vector
from the dense span; continuity preserves a sufficiently small negative
margin.  The reverse statement is immediate from the corresponding
quadratic form.
\end{proof}

\subsection{Local diagonal and overlap stocks}
\label{sec:s6v71-local}

Suppose the centered score is decomposed into local terms and hence
\[
 {\mathsf K}=\sum_z{\mathsf K}_z.
\]
Define
\[
 {\boldsymbol{\mathsf K}}\varphi
 =({\mathsf K}_z\varphi)_z
 \in\ell^2({\cal Z};L^2(\pi_{1/2}))
\]
whenever the vector is square summable, and let
\({\mathsf S}_{\rm add}(h_z)_z=\sum_zh_z\) on the relevant range.

\begin{theorem}[Local score-readout compiler]
\label{thm:s6v71-local-compiler}
Suppose, for all mean-zero \(\varphi\),
\begin{align}
 \sum_z\|{\mathsf K}_z\varphi\|^2
 &\le d\delta\,\|{\mathsf Z}\varphi\|^2,                   \label{eq:s6v71-local-diagonal}\\
 \|{\mathsf S}_{\rm add}{\boldsymbol{\mathsf K}}\varphi\|^2
 &\le\kappa
 \|{\boldsymbol{\mathsf K}}\varphi\|^2.                   \label{eq:s6v71-local-overlap}
\end{align}
Then SFRB holds with
\begin{equation}
 s=\kappa d.
 \label{eq:s6v71-local-stock}
\end{equation}
\end{theorem}

\begin{proof}
Since
\({\mathsf K}\varphi
={\mathsf S}_{\rm add}{\boldsymbol{\mathsf K}}\varphi\),
\eqref{eq:s6v71-local-overlap} followed by
\eqref{eq:s6v71-local-diagonal} gives
\[
 \|{\mathsf K}\varphi\|^2
 \le\kappa d\delta\,\|{\mathsf Z}\varphi\|^2
 =\kappa d\delta\,\ell(\varphi).
\]
\end{proof}

\begin{assessment}[Division of labour after the V70 audit]
\label{ass:s6v71-division}
The local diagonal estimate
\eqref{eq:s6v71-local-diagonal} is a return-specific statement: every
local centered score readout must be small relative to the same
conditional residual that defines the transfer loss.  The overlap
estimate \eqref{eq:s6v71-local-overlap} is where the SM V12--V13
connected-tree technology may help, after bounded or form-controlled
local insertions and a gauge-invariant frame are supplied.  Absolute
tree decay cannot prove the diagonal factor \(\delta\), while the
diagonal estimate alone cannot prevent volume growth of the coherent
sum.
\end{assessment}

\subsection{Revised bottleneck}
\label{sec:s6v71-frontier}

\begin{assessment}[Exact frontier after V71]
\label{ass:s6v71-frontier}
The nullspace obstruction anticipated in V70 is automatically resolved:
every exact return-measurable mode has zero score covariance.  The hard
part is quantitative, not algebraic.  The decisive unsigned row is now
the restricted norm
\eqref{eq:s6v71-restricted-norm}, or its local compiler
\eqref{eq:s6v71-local-diagonal}--\eqref{eq:s6v71-local-overlap}.
Even if this row is proved, CMRI or the V69 low-loss signed alignment
remains independently necessary.  Hence the programme has separated the
remaining mass mechanism into:
\[
 \boxed{
 \text{loss-relative score readout}
 \quad+\quad
 \text{positive low-loss alignment}.}
\]
\end{assessment}

\begin{programme}[V72: local plaquette score versus return residual]
\label{prog:s6v71-V72}
\begin{enumerate}[leftmargin=3em]
\item Write the elementary Wilson plaquette score readout
      \({\mathsf K}_z\) on the half-slab Doob law.
\item Test \eqref{eq:s6v71-local-diagonal} by conditioning successively
      on all variables except the plaquette neighbourhood.
\item Determine the exact \(\delta\)-factor from the temporal plaquette
      heat-kernel or exhibit a finite spin-network violation.
\item Only after the diagonal estimate, use connected-tree decay to
      control \(\kappa\).
\item Evaluate the sign of the resulting operator on the V61 bilocal
      soft packets.
\end{enumerate}
\end{programme}

\begin{firewall}[Claim boundary after sixth-series V71]
\label{fire:s6v71-boundary}
V71 constructs the canonical conditional residual operator and proves
\(D={\mathsf Z}^*{\mathsf Z}\); characterizes exact-return modes and
proves strict positive kernels leave only constants; proves the score
covariance factors automatically through the physical residual; reduces
SFRB to a restricted readout norm; gives a complete finite loss-covariance
pencil; and proves a local diagonal/overlap compiler.

V71 does not prove the restricted readout bound, its local diagonal or
overlap stocks, positive low-loss alignment, CMRI, the global clock lower
bound, the V53 projective/Mosco hypotheses, the continuum OS theory, or
the full Yang--Mills mass gap.
\end{firewall}

\begin{theorem}[Sixth-series V71 result ledger]
\label{thm:s6v71-results}
The positive conclusions in \Cref{fire:s6v71-boundary} are exactly those
of
\Cref{thm:s6v71-residual,
cor:s6v71-Jensen,
prop:s6v71-positive-kernel,
thm:s6v71-score-factor,
thm:s6v71-restricted-norm,
cor:s6v71-extensive,
thm:s6v71-pencil,
cor:s6v71-pencil-failure,
thm:s6v71-local-compiler}.
\end{theorem}

\begin{proof}
Each item is the cited result with its conditional-law, positivity,
boundedness, compatibility-range, local-decomposition, clock, and
cofinal qualifications unchanged.
\end{proof}

\paragraph{Parent-version record.}
V71 was formed from
\texttt{Renewal\_Yang\_Mills\_Mass\_Gap\_Research\_Notes\_V70.tex}.
The V70 parent had \(31\,626\) logical source lines,
\(1\,110\,390\) bytes, and SHA--256
\[
 \texttt{7bd53b480498654998c811ced6e1d0253c56bc89ad48205088dfd8d3fbf0556a}.
\]
The next compulsory companion audit is V75.

% =============================================================================
% END APPENDED SIXTH-SERIES V71 MATERIAL
% =============================================================================

% =============================================================================
% BEGIN APPENDED SIXTH-SERIES V72 MATERIAL
% =============================================================================

\section{V72: compact-group heat return and positive covariance splitting}
\label{sec:s6v72}

\subsection{Why the full covariance must not be treated as an error}
\label{sec:s6v72-purpose}

V71 reduced SFRB to the norm of the physical score readout on conditional
residuals.  That is the correct estimate for a \emph{remainder} score.
It need not hold for the positive principal score itself.  An exact
compact-group heat channel proves this point: its entire spectral response
is a positive within-fibre covariance, it obeys DSFI with explicit
constants, yet its full score covariance can violate a cofinal SFRB on
ultraviolet modes.  This forces a positive-principal/error splitting
before absolute covariance estimates are applied.

\subsection{The compact-group heat channel}
\label{sec:s6v72-heat}

Let \(G\) be a compact connected Lie group with normalized Haar measure,
and let \(H\ge0\) be a bi-invariant Laplace-type operator on
\(L^2(G)\), normalized so that
\[
 \ker H=\operatorname{span}\{1\}.
\]
Write \(P_t=e^{-tH}\) and let \(p_t(yx^{-1})>0\) be its heat kernel for
\(t>0\).  Take the half-slab amplitude
\begin{equation}
 F_t(y,x)=p_t(yx^{-1}),\qquad A_t=P_t.
 \label{eq:s6v72-heat-amplitude}
\end{equation}

\begin{lemma}[Heat-channel Doob data]
\label{lem:s6v72-heat-data}
The leading singular data are
\[
 s_1=1,\qquad u=v=1.
\]
The Doob law and normalized full transfer are
\begin{align}
 \Lambda_t(\dd x,\dd y)
 &=p_t(yx^{-1})\,\dd x\,\dd y,                            \label{eq:s6v72-heat-Doob}\\
 T_t&=P_{2t}=e^{-2tH},&
 L_t&=I-e^{-2tH}.                                        \label{eq:s6v72-heat-transfer}
\end{align}
If \(\lambda_*>0\) is the smallest nonzero eigenvalue of \(H\), the
one-step relative gap is \(1-e^{-2t\lambda_*}\).
\end{lemma}

\begin{proof}
The heat semigroup preserves constants and Haar measure, so
\(P_t1=P_t^*1=1\).  Positivity improvement makes this leading eigenvalue
simple.  The amplitude is self-adjoint and the semigroup law gives
\(A_t^*A_t=P_{2t}\).  Peter--Weyl diagonalizes \(H\); on the
vacuum-orthogonal carrier the lowest eigenvalue is \(\lambda_*\), proving
the last assertion.
\end{proof}

Let \(t=t(\theta)>0\) be \(C^1\) and nondecreasing.  Define the amplitude
score by
\begin{equation}
 \partial_\theta F_{t(\theta)}
 =-{1\over2}R_\theta F_{t(\theta)}.
 \label{eq:s6v72-score-definition}
\end{equation}

\begin{theorem}[The heat score is a purely aligned covariance]
\label{thm:s6v72-heat-score}
For the heat channel,
\begin{align}
 {\cal J}_{1,\theta}&=0,&
 m_\theta(y)&=\mathbb E[R_\theta\mid Y=y]=0,               \label{eq:s6v72-zero-means}\\
 A_{\theta,R}
 &=2\dot t(\theta)H e^{-t(\theta)H},                       \label{eq:s6v72-inserted-heat}\\
 {\mathsf K}_\theta
 &=2\dot t(\theta)H e^{-t(\theta)H},                       \label{eq:s6v72-K-heat}\\
 L_\theta'
 &=2\dot t(\theta)H e^{-2t(\theta)H}\succeq0.              \label{eq:s6v72-Lprime-heat}
\end{align}
For every mean-zero \(\varphi\),
\begin{equation}
 \boxed{
 \langle\varphi,L_\theta'\varphi\rangle
 =\langle P_t\varphi,{\mathsf K}_\theta\varphi\rangle
 =2\dot t\,\|H^{1/2}e^{-tH}\varphi\|^2\ge0.}
 \label{eq:s6v72-positive-covariance}
\end{equation}
Thus the entire response is the within-return-fibre covariance; the
centered conditional-mean multiplier vanishes.
\end{theorem}

\begin{proof}
Differentiate \(A_t=e^{-tH}\):
\[
 A_\theta'=-\dot t\,He^{-tH}.
\]
Since \(A_{\theta,R}=-2A_\theta'\), this proves
\eqref{eq:s6v72-inserted-heat}.  The leading eigenvalue is identically
one, so its logarithmic derivative and hence \({\cal J}_{1,\theta}\)
vanish.  Also
\[
 m_\theta(y)
 ={(A_{\theta,R}1)(y)\over(A_\theta1)(y)}=0
\]
because \(H1=0\).  With \(u=v=s_1=1\), the conditional covariance
operator equals \(A_{\theta,R}\), proving
\eqref{eq:s6v72-K-heat}.  Differentiating
\(L_\theta=I-e^{-2tH}\) proves
\eqref{eq:s6v72-Lprime-heat}.  Finally, functional calculus and
self-adjointness give
\[
 \langle e^{-tH}\varphi,
 2\dot t\,He^{-tH}\varphi\rangle
 =2\dot t\,\|H^{1/2}e^{-tH}\varphi\|^2.
\]
\end{proof}

\subsection{An explicit heat-channel G\aa rding inequality}
\label{sec:s6v72-heat-Garding}

\begin{theorem}[Clock-scaled heat DSFI]
\label{thm:s6v72-heat-DSFI}
Assume
\begin{equation}
 \dot t(\theta)=c_\theta\delta,\qquad c_\theta>0.
 \label{eq:s6v72-clock-path}
\end{equation}
For every \(\alpha>0\), define
\begin{equation}
 b_\theta
 :=\min\left\{
 \alpha,\,
 2c_\theta\lambda_*e^{-2t(\theta)\lambda_*}
 +\alpha(1-e^{-2t(\theta)\lambda_*})
 \right\}>0.
 \label{eq:s6v72-b}
\end{equation}
Then on \(1^\perp\),
\begin{equation}
 \boxed{
 L_\theta'
 \succeq b_\theta\delta I-\alpha\delta L_\theta.}
 \label{eq:s6v72-heat-DSFI}
\end{equation}
If \(c_\theta\ge c_*>0\) and \(t(\theta)\le t^*<\infty\), then
\(\inf_\theta b_\theta>0\).
\end{theorem}

\begin{proof}
On an \(H\)-eigenvector of eigenvalue \(\lambda\ge\lambda_*\),
the left side of
\(L_\theta'+\alpha\delta L_\theta\) divided by \(\delta\) is
\begin{equation}
 f_\theta(\lambda)
 =2c_\theta\lambda e^{-2t\lambda}
 +\alpha(1-e^{-2t\lambda}).
 \label{eq:s6v72-f}
\end{equation}
Its derivative is
\[
 f_\theta'(\lambda)
 =2e^{-2t\lambda}
 \{c_\theta+\alpha t-2c_\theta t\lambda\}.
\]
It has at most one critical point, and that point is a maximum.
Therefore the minimum on \([\lambda_*,\infty)\) is the smaller of its
value at \(\lambda_*\) and its limit \(\alpha\) at infinity.  This is
exactly \(b_\theta\).  Functional calculus proves
\eqref{eq:s6v72-heat-DSFI}.  The uniform statement follows by minimizing
the positive continuous endpoint expression over the compact range
\(c_\theta\ge c_*\), \(0<t\le t^*\); decreasing \(c\) to \(c_*\) gives
the worst first term.
\end{proof}

\begin{assessment}[What the single-link heat gap means]
\label{ass:s6v72-heat-meaning}
The theorem is a complete direct-form proof for one compact heat channel,
and tensor products of independent such channels retain a one-link
Casimir gap.  It is not the four-dimensional Yang--Mills mass gap:
spatial plaquette couplings, gauge constraints, continuum scaling, and
the V61 delocalized colour-singlet modes change the physical low-energy
carrier.  Its role is to identify the correct sign mechanism and the
correct treatment of ultraviolet directions.
\end{assessment}

\subsection{Full SFRB fails although the response is positive}
\label{sec:s6v72-SFRB-failure}

\begin{proposition}[Ultraviolet obstruction to full-score SFRB]
\label{prop:s6v72-SFRB-failure}
Let
\[
 t_\delta=\tau\delta,\qquad
 \dot t_\delta=c\delta,\qquad \tau,c>0.
\]
On an \(H\)-eigenvector of eigenvalue \(\lambda\),
\begin{equation}
 {\|{\mathsf K}_\delta\varphi\|^2
  \over\langle\varphi,D_\delta\varphi\rangle}
 ={4c^2\delta^2\lambda^2\over e^{2\tau\delta\lambda}-1},
 \qquad
 D_\delta=I-e^{-2\tau\delta H}.
 \label{eq:s6v72-SFRB-ratio}
\end{equation}
If a cofinal sequence of eigenvalues satisfies
\[
 2\tau\delta\lambda_\delta\longrightarrow x\in(0,\infty),
\]
then the ratio in \eqref{eq:s6v72-SFRB-ratio} converges to
\begin{equation}
 {c^2\over\tau^2}{x^2\over e^x-1}>0.
 \label{eq:s6v72-SFRB-positive-limit}
\end{equation}
Hence no regulator-independent \(s\) can make the \emph{full} heat score
obey
\({\mathsf K}_\delta^*{\mathsf K}_\delta\preceq
s\delta D_\delta\) along such a sequence.
\end{proposition}

\begin{proof}
Use \eqref{eq:s6v72-K-heat} and the spectral values
\[
 {\mathsf K}_\delta:
 2c\delta\lambda e^{-\tau\delta\lambda},
 \qquad
 D_\delta:1-e^{-2\tau\delta\lambda}.
\]
Their squared ratio simplifies to
\eqref{eq:s6v72-SFRB-ratio}.  Substitute
\(\lambda_\delta\sim x/(2\tau\delta)\) to obtain
\eqref{eq:s6v72-SFRB-positive-limit}.  SFRB would force the same ratio
to be at most \(s\delta\to0\), a contradiction.
\end{proof}

\begin{corollary}[SFRB is an error estimate, not a positivity estimate]
\label{cor:s6v72-SFRB-role}
The positive heat covariance in
\eqref{eq:s6v72-positive-covariance} must be retained as a principal
signed term.  Estimating all of it by absolute SFRB is both unnecessary
and, cofinally, false on the ultraviolet band
\(\lambda\asymp\delta^{-1}\).
\end{corollary}

\begin{proof}
Positivity is \Cref{thm:s6v72-heat-score}; failure of the absolute
relative estimate is \Cref{prop:s6v72-SFRB-failure}.
\end{proof}

\subsection{Positive-principal/error covariance compiler}
\label{sec:s6v72-splitting}

Return to a general Doob half-slab law.  Split the covariance operator as
\begin{equation}
 {\mathsf K}={\mathsf K}_+ +{\mathsf K}_{\rm err}.
 \label{eq:s6v72-K-split}
\end{equation}
The notation \(+\) records a term to be treated with its sign; it does
not assume the operator itself is positive between different Hilbert
spaces.

\begin{definition}[Aligned principal response inequality]
\label{def:s6v72-APRI}
APRI with stocks \(\beta>0\) and \(\alpha\ge0\) is
\begin{align}
 \mathbb E_{\pi_{1/2}}(m-J)|{\mathsf P}\varphi|^2
 +\operatorname{Re}
 \langle{\mathsf P}\varphi,{\mathsf K}_+\varphi\rangle
 \ge
 \beta\delta\|\varphi\|^2-\alpha\ell(\varphi)
 \label{eq:s6v72-APRI}
\end{align}
for every mean-zero \(\varphi\).
\end{definition}

\begin{theorem}[APRI plus error-SFRB implies DSFI]
\label{thm:s6v72-split-compiler}
Suppose APRI holds and
\begin{equation}
 \|{\mathsf K}_{\rm err}\varphi\|^2
 \le s\delta\,\ell(\varphi)
 \label{eq:s6v72-error-SFRB}
\end{equation}
for every mean-zero \(\varphi\).  Then
\begin{equation}
 \boxed{
 QL'Q\succeq{\beta\over2}\delta Q
 -\left(\alpha+{s\over2\beta}\right)QLQ.}
 \label{eq:s6v72-split-DSFI}
\end{equation}
\end{theorem}

\begin{proof}
The exact covariance formula and
\eqref{eq:s6v72-K-split} give the APRI left side plus
\(\operatorname{Re}
\langle{\mathsf P}\varphi,{\mathsf K}_{\rm err}\varphi\rangle\).
By Cauchy--Schwarz and \eqref{eq:s6v72-error-SFRB},
\[
 |\langle{\mathsf P}\varphi,{\mathsf K}_{\rm err}\varphi\rangle|
 \le\sqrt{s\delta\,t(\varphi)\ell(\varphi)}
 \le{\beta\delta\over2}t(\varphi)
    +{s\over2\beta}\ell(\varphi).
\]
Use \(t(\varphi)\le\|\varphi\|^2\), then conjugate by
\(f=u\varphi\) as in
\Cref{thm:s6v69-CMRI-SFRB}.
\end{proof}

\begin{corollary}[Heat channel as the zero-error case]
\label{cor:s6v72-heat-split}
For the heat channel take
\[
 {\mathsf K}_+={\mathsf K},\qquad
 {\mathsf K}_{\rm err}=0,\qquad m-J=0.
\]
The principal response is the positive form
\eqref{eq:s6v72-positive-covariance}; no SFRB is imposed on it.
\end{corollary}

\begin{proof}
This is \Cref{thm:s6v72-heat-score} inserted in
\eqref{eq:s6v72-K-split}.
\end{proof}

\subsection{Yang--Mills interpretation and new target}
\label{sec:s6v72-YM}

\begin{assessment}[Plaquette heat part versus interacting remainder]
\label{ass:s6v72-YM-split}
The temporal plaquette heat kernel suggests a noncircular split of the
Yang--Mills score:
\[
 \text{single-plaquette heat/Casimir response}
 \quad+\quad
 \text{spatial, gauge-incidence, and Wilsonian return error}.
\]
The first term should be kept signed and compared with
\eqref{eq:s6v72-positive-covariance}; only the second should be tested by
the V71 residual readout norm and the SM connected-tree overlap bounds.
This does not yet prove APRI.  In particular, the V61 soft singlet packets
show that a sum of local Casimir responses must survive gauge projection
and continuum scaling with a uniform positive coefficient; a one-link
Casimir gap alone cannot establish that.
\end{assessment}

\begin{definition}[Heat-principal score card, HPSC]
\label{def:s6v72-HPSC}
For the response-matched Yang--Mills half-slab, HPSC consists of:
\begin{enumerate}[label=\textup{(\roman*)},leftmargin=3.5em]
\item an exact decomposition of the amplitude tangent into temporal
      heat/Casimir and interacting-return pieces;
\item identification of the induced
      \({\mathsf K}_+\) and \({\mathsf K}_{\rm err}\) without changing
      the ground transform;
\item APRI for the heat/Casimir principal part on the full
      gauge-invariant low-loss carrier, including the V61 packets;
\item error-SFRB for the remainder, using the V71 local
      diagonal/overlap compiler;
\item the global clock lower bound and V53 limit rows.
\end{enumerate}
\end{definition}

\begin{assessment}[Exact frontier after V72]
\label{ass:s6v72-frontier}
The covariance route has been corrected: forcing the whole score through
SFRB can discard the mass-producing sign and is false even for the exact
heat channel.  The viable sufficient theorem is now
\Cref{thm:s6v72-split-compiler}.  The next task is not another absolute
bound; it is to extract an exact heat/Casimir principal tangent from the
physical response-matched Wilsonian half-slab and prove APRI after gauge
projection.  If no such decomposition exists for the attached amplitude,
the heat analogy must remain diagnostic and the programme must return to
the complete signed pencil.
\end{assessment}

\begin{programme}[V73: exact Wilsonian tangent split]
\label{prog:s6v72-V73}
\begin{enumerate}[leftmargin=3em]
\item Locate the temporal heat-kernel/Casimir term in the physical
      response-matched half-slab action, including its exact coefficient
      and clock.
\item Differentiate the same amplitude owner to define
      \({\mathsf K}_+\) and \({\mathsf K}_{\rm err}\); do not splice
      kernels from different measures.
\item Compute the principal response on gauge-invariant spin networks
      and the V61 bilocal packets.
\item Test whether gauge projection preserves a cofinal Casimir floor or
      whether spatial soft modes make it vanish.
\item Apply SFRB only to the residual interacting tangent.
\end{enumerate}
\end{programme}

\begin{firewall}[Claim boundary after sixth-series V72]
\label{fire:s6v72-boundary}
V72 solves the compact-group heat return exactly; proves its action score
is a purely positive conditional covariance; proves an explicit DSFI;
shows full-score SFRB can fail on the ultraviolet clock band despite
positivity; and proves the positive-principal/error covariance compiler.

V72 does not identify the heat-principal split in the physical
response-matched Yang--Mills amplitude, prove APRI or error-SFRB there,
establish a gauge-invariant cofinal Casimir floor, close the V53
projective/Mosco hypotheses, construct the continuum OS theory, or prove
the full Yang--Mills mass gap.
\end{firewall}

\begin{theorem}[Sixth-series V72 result ledger]
\label{thm:s6v72-results}
The positive conclusions in \Cref{fire:s6v72-boundary} are exactly those
of
\Cref{lem:s6v72-heat-data,
thm:s6v72-heat-score,
thm:s6v72-heat-DSFI,
prop:s6v72-SFRB-failure,
cor:s6v72-SFRB-role,
thm:s6v72-split-compiler,
cor:s6v72-heat-split}.
\end{theorem}

\begin{proof}
Each item is the cited result with its compact-group, heat-kernel,
spectral, clock, splitting, gauge, and cofinal qualifications unchanged.
\end{proof}

\paragraph{Parent-version record.}
V72 was formed from
\texttt{Renewal\_Yang\_Mills\_Mass\_Gap\_Research\_Notes\_V71.tex}.
The V71 parent had \(32\,059\) logical source lines,
\(1\,125\,154\) bytes, and SHA--256
\[
 \texttt{c694fc9decef34c5b8479249411dbbfe6d5c1f01461cb4073428096e2f07c7eb}.
\]
The next compulsory companion audit is V75.

% =============================================================================
% END APPENDED SIXTH-SERIES V72 MATERIAL
% =============================================================================

% =============================================================================
% BEGIN APPENDED SIXTH-SERIES V73 MATERIAL
% =============================================================================

\section{V73: same-owner Casimir tangent and the source-separation firewall}
\label{sec:s6v73}

\subsection{Targeted source audit}
\label{sec:s6v73-audit}

V72 showed that a heat/Casimir principal covariance can open the transfer
edge, but only if it is a tangent of the same physical half-slab
amplitude.  The attached sources were therefore inspected at the exact
points where such an identity could occur.

\begin{proposition}[What the attached sources actually identify]
\label{prop:s6v73-source-findings}
The attached sources establish the following:
\begin{enumerate}[label=\textup{(\roman*)},leftmargin=3.5em]
\item The Grand Tensor manuscript defines, at each finite compact Wilson
      regulator, the physical pure-colour half-slab amplitude
      \[
      F(U_+,U_0)=c\exp[-S_+(U_+,U_0)/2]
      \]
      and proves its Hilbert--Schmidt, positivity-improving, and
      Doob--Schmidt properties.
\item Its reflected-boundary theorem identifies a positive cross-source
      Gram and a strict local electric-Casimir contact in a boundary
      Hessian.  It calls that contact an audit term, not an independent
      gauge-invariant mass, and states that the coherent remainder cancels
      it on boundary gauge gradients.
\item The reduced-fibre manuscript explicitly says that its Witten lift
      does not identify an arbitrary half-slab amplitude, boundary
      instrument, or spatial route.
\item The reduced-fibre manuscript requires the actual response-matched
      coarse-link family with every generated, marginal, and support term
      retained, and distinguishes exact pushforwards from terms that are
      merely named.
\end{enumerate}
None of these statements gives an operator identity
\begin{equation}
 A_\theta'=-c_\theta\delta H_\theta A_\theta+E_\theta
 \label{eq:s6v73-missing-identity}
\end{equation}
for the response-matched half-slab amplitude.
\end{proposition}

\begin{proof}
Item~\textup{(i)} is the definition and Doob--Schmidt theorem in the
Grand Tensor half-slab hierarchy.  Item~\textup{(ii)} is its theorem on
the reflection correction and positive cross source; both qualifications
occur in that theorem's conclusion.  Item~\textup{(iii)} is the scope
assessment following the reduced-fibre exact-form Witten lift.
Item~\textup{(iv)} is the reduced-fibre colour-germ gate and its later
assessment of what is specified by exact fibre pushforward.  None of
those displayed results differentiates the physical half-slab amplitude
in the form \eqref{eq:s6v73-missing-identity}.
\end{proof}

\begin{assessment}[Source consequence]
\label{ass:s6v73-source-consequence}
The Casimir contact is evidence for a possible principal term, but it is
not yet that term.  Importing it directly into V72 would splice a
boundary-Hessian statement into a first interpolation derivative of a
possibly different amplitude owner.  The cancellation on gauge gradients
also prevents promotion of the contact before the coherent remainder and
physical Hodge short are included.
\end{assessment}

\subsection{The exact same-owner tangent theorem}
\label{sec:s6v73-tangent}

Let \(A_\theta:{\cal H}_0\to{\cal H}_{1/2}\) be the physical
norm-\(C^1\) half-slab family, with
\[
 A_\theta u_\theta=s_{1,\theta}v_\theta,\qquad
 A_\theta^*v_\theta=s_{1,\theta}u_\theta,\qquad
 \lambda_{1,\theta}=s_{1,\theta}^2.
\]
Fix \(\theta\) and suppress it below.

\begin{definition}[Same-owner Casimir tangent]
\label{def:s6v73-SCT}
A same-owner Casimir tangent with stocks \(c,\kappa>0\) is a decomposition
\begin{equation}
 A'=-c\delta HA+E
 \label{eq:s6v73-tangent-split}
\end{equation}
on the actual amplitude spaces, where:
\begin{enumerate}[label=\textup{(\roman*)},leftmargin=3.5em]
\item \(H\ge0\) is self-adjoint on the output space and \(HA\) is
      bounded;
\item \(Hv=0\);
\item
\begin{equation}
 H\succeq\kappa(I-|v\rangle\langle v|)
 \label{eq:s6v73-output-gap}
\end{equation}
      on the relevant output carrier;
\item \(E\) is the exact residual tangent of this same amplitude.
\end{enumerate}
\end{definition}

\begin{theorem}[Same-owner Casimir response]
\label{thm:s6v73-Casimir-response}
Under \Cref{def:s6v73-SCT}, split the normalized derivative linearly as
\[
 L'=L'_+ +L'_{\rm err},
\]
using \(-c\delta HA\) and \(E\), respectively, in the derivative of
\(B=A^*A\) and of \(\lambda_1\).  Then
\begin{equation}
 \boxed{
 L'_+={2c\delta\over\lambda_1}A^*HA.}
 \label{eq:s6v73-Lplus}
\end{equation}
On the vacuum-orthogonal input carrier,
\begin{align}
 \langle f,L'_+f\rangle
 &\ge2c\kappa\delta\langle f,Tf\rangle                    \label{eq:s6v73-principal-T}\\
 &=2c\kappa\delta
   \{\|f\|^2-\langle f,Lf\rangle\}.                       \label{eq:s6v73-principal-L}
\end{align}
Thus the principal tangent satisfies DSFI with
\[
 b_+=2c\kappa,\qquad a_+=2c\kappa\delta.
\]
\end{theorem}

\begin{proof}
The principal derivative is \(A'_+=-c\delta HA\).  Since \(H\) is
self-adjoint,
\[
 B'_+=-2c\delta A^*HA.
\]
Its contribution to the leading eigenvalue derivative is
\[
 \lambda'_{1,+}
 =2\operatorname{Re}\langle Au,A'_+u\rangle
 =-2c\delta\lambda_1\langle v,Hv\rangle=0.
\]
The normalized derivative formula of V63 gives
\eqref{eq:s6v73-Lplus}.

If \(f\perp u\), then
\[
 \langle v,Af\rangle
 =\langle A^*v,f\rangle=s_1\langle u,f\rangle=0.
\]
Apply \eqref{eq:s6v73-output-gap} to \(Af\):
\[
 \langle f,L'_+f\rangle
 ={2c\delta\over\lambda_1}\langle Af,HAf\rangle
 \ge{2c\kappa\delta\over\lambda_1}\|Af\|^2
 =2c\kappa\delta\langle f,Tf\rangle.
\]
Finally \(T=I-L\), proving
\eqref{eq:s6v73-principal-L}.
\end{proof}

\begin{corollary}[Residual debit allowed by the principal stock]
\label{cor:s6v73-residual-debit}
If
\begin{equation}
 QL'_{\rm err}Q
 \succeq-d\delta Q-a_{\rm err}QLQ,
 \qquad 0\le d<2c\kappa,
 \label{eq:s6v73-error-lower}
\end{equation}
then
\begin{equation}
 QL'Q
 \succeq(2c\kappa-d)\delta Q
 -(a_{\rm err}+2c\kappa\delta)QLQ.
 \label{eq:s6v73-full-DSFI}
\end{equation}
\end{corollary}

\begin{proof}
Add \eqref{eq:s6v73-principal-L} and
\eqref{eq:s6v73-error-lower}.
\end{proof}

\begin{corollary}[Cofinal Casimir handoff]
\label{cor:s6v73-cofinal}
Suppose the preceding hypotheses hold cofinally and
\[
 \inf_{j,\theta}(2c_{j,\theta}\kappa_{j,\theta}
 -d_{j,\theta})>0,\qquad
 \sup_j\int_0^1(a_{{\rm err},j,\theta}
 +2c_{j,\theta}\kappa_{j,\theta}\delta_j)\,\dd\theta<\infty.
\]
Then the V66 theorem gives a regulator-uniform positive rescaled endpoint
gap.
\end{corollary}

\begin{proof}
\Cref{cor:s6v73-residual-debit} gives DSFI with the displayed uniform
stocks.  Apply \Cref{thm:s6v66-garding-gap}.
\end{proof}

\begin{assessment}[Content of the same-owner theorem]
\label{ass:s6v73-content}
The output gap \(\kappa\) is not by itself a Yang--Mills mass.  It becomes
a transfer response only through the exact identity
\eqref{eq:s6v73-tangent-split}.  The theorem then uses \(Af\perp v\) to
convert output Casimir coercivity into a positive response on every
transmitted relative direction.  The loss term pays directions poorly
transmitted by the half-slab.
\end{assessment}

\subsection{A Hessian cannot determine the interpolation tangent}
\label{sec:s6v73-Hessian}

\begin{lemma}[Tangent independence from a fixed boundary Hessian]
\label{lem:s6v73-Hessian-independent}
Let \(q\) be a finite-dimensional boundary coordinate, let \(W_0(q)\) be
smooth with any prescribed positive Hessian at \(q=0\), and let \(r(q)\)
be an arbitrary smooth real function.  The two positive amplitude paths
\begin{equation}
 F_\theta^\pm(q)
 =\exp\left\{-{1\over2}[W_0(q)\pm\theta r(q)]\right\}
 \label{eq:s6v73-two-paths}
\end{equation}
have the same amplitude and boundary Hessian at \(\theta=0\), but their
interpolation scores and normalized transfer derivatives there are
opposites, apart from the common scalar normalization line.
\end{lemma}

\begin{proof}
At \(\theta=0\), both paths equal \(e^{-W_0/2}\), so every derivative in
the boundary coordinate \(q\), including its Hessian, agrees.  Their
interpolation derivatives are
\[
 \left.\partial_\theta F_\theta^\pm\right|_{\theta=0}
 =\mp{1\over2}rF_0,
\]
and hence their scores are \(\pm r\).  The exact derivative formula
\Cref{prop:s6v66-derivative-form} is linear in the centered score, so the
non-scalar normalized derivative forms have opposite signs.
\end{proof}

\begin{corollary}[No Hessian-to-tangent inference]
\label{cor:s6v73-no-Hessian-splice}
A positive local boundary Hessian or Casimir contact does not imply the
same-owner tangent identity \eqref{eq:s6v73-tangent-split}, nor even the
sign of \(L'\).  An explicit amplitude derivative or an exact
pushforward/Duhamel identity is indispensable.
\end{corollary}

\begin{proof}
\Cref{lem:s6v73-Hessian-independent} keeps the Hessian fixed while
reversing the tangent response.
\end{proof}

\subsection{A complete verification criterion}
\label{sec:s6v73-verification}

For a proposed output generator \(H\) and coefficient \(c\), define the
same-owner residual
\begin{equation}
 E_{c,H}:=A'+c\delta HA.
 \label{eq:s6v73-residual}
\end{equation}
Let \(\{\phi_n\}\) and \(\{\chi_m\}\) be complete input and output
spin-network/Peter--Weyl cores.

\begin{proposition}[Complete matrix test for a claimed Casimir tangent]
\label{prop:s6v73-matrix-test}
If \(A'\) and \(HA\) are bounded, then
\[
 A'=-c\delta HA
\]
if and only if
\begin{equation}
 \langle\chi_m,E_{c,H}\phi_n\rangle=0
 \quad\hbox{for every }m,n.
 \label{eq:s6v73-matrix-test}
\end{equation}
Every quantitative error bound on \(E_{c,H}\) must likewise be proved for
the matrix of this same residual.
\end{proposition}

\begin{proof}
The forward implication is immediate.  Conversely
\eqref{eq:s6v73-matrix-test} says \(E_{c,H}\phi_n\) is orthogonal to a
dense output core for every \(n\), hence vanishes.  Boundedness extends
the equality from the dense input core to all inputs.
\end{proof}

\begin{definition}[Same-owner Casimir card, SOCC]
\label{def:s6v73-SOCC}
For a cofinal response-matched half-slab family, SOCC consists of:
\begin{enumerate}[label=\textup{(\roman*)},leftmargin=3.5em]
\item the actual \(A_\theta\) and its derivative, including all generated
      Wilsonian returns and normalization terms;
\item an output \(H_\theta\) with \(H_\theta v_\theta=0\) and a cofinal
      relative gap \(\kappa>0\);
\item the exact split \eqref{eq:s6v73-tangent-split}, verified on complete
      cores by \eqref{eq:s6v73-matrix-test};
\item the residual lower bound \eqref{eq:s6v73-error-lower} with
      \(d<2c\kappa\);
\item common-sector shorting and V53 projective/Mosco landing.
\end{enumerate}
\end{definition}

\begin{assessment}[Exact frontier after V73]
\label{ass:s6v73-frontier}
SOCC would close the finite-regulator spectral row by
\Cref{cor:s6v73-cofinal}, but the attached sources do not presently
supply its tangent split.  The correct next move is upstream and finite:
differentiate the literal Wilson half-slab \(F=c e^{-S_+/2}\), choose the
electric output Casimir proposed by the reflected-head theorem, and
compute the residual \eqref{eq:s6v73-residual} on the first
gauge-invariant Peter--Weyl sectors.  If that residual is not smaller than
the principal stock, the heat-principal route fails before any cofinal
claim.
\end{assessment}

\begin{programme}[V74: first-sector same-owner residual]
\label{prog:s6v73-V74}
\begin{enumerate}[leftmargin=3em]
\item Fix one finite Wilson half-slab and write its literal action
      derivative under a declared coupling/time interpolation.
\item Evaluate \(A'\), \(HA\), and \(E_{c,H}\) on the vacuum and first
      nontrivial gauge-invariant Peter--Weyl sectors.
\item Separate normalization, gauge-gradient, and physical Hodge
      components before taking a sign.
\item Test whether one coefficient \(c>0\) cancels the leading temporal
      Casimir tangent across those sectors.
\item If no common coefficient exists, return to the complete signed
      pencil rather than forcing SOCC.
\end{enumerate}
\end{programme}

\begin{firewall}[Claim boundary after sixth-series V73]
\label{fire:s6v73-boundary}
V73 audits the exact half-slab and reflected-Casimir statements in the
attached sources; proves the same-owner Casimir tangent theorem and its
DSFI consequence; proves a boundary Hessian cannot determine an
interpolation tangent; and gives a complete core-matrix verification
criterion.

V73 does not prove the physical Yang--Mills amplitude has the same-owner
Casimir split, bound its residual below the principal stock, establish
cofinal SOCC, close the V53 projective/Mosco hypotheses, construct the
continuum OS theory, or prove the full Yang--Mills mass gap.
\end{firewall}

\begin{theorem}[Sixth-series V73 result ledger]
\label{thm:s6v73-results}
The positive conclusions in \Cref{fire:s6v73-boundary} are exactly those
of
\Cref{prop:s6v73-source-findings,
thm:s6v73-Casimir-response,
cor:s6v73-residual-debit,
cor:s6v73-cofinal,
lem:s6v73-Hessian-independent,
cor:s6v73-no-Hessian-splice,
prop:s6v73-matrix-test}.
\end{theorem}

\begin{proof}
Each item is the cited result with its same-owner, differentiability,
output-gap, residual, gauge, finite-regulator, and cofinal qualifications
unchanged.
\end{proof}

\paragraph{Parent-version record.}
V73 was formed from
\texttt{Renewal\_Yang\_Mills\_Mass\_Gap\_Research\_Notes\_V72.tex}.
The V72 parent had \(32\,500\) logical source lines,
\(1\,140\,586\) bytes, and SHA--256
\[
 \texttt{c908857b64ba0492c25218097a5023835f583213fae51de49b455acdfff3a4e3}.
\]
The next compulsory companion audit is V75.

% =============================================================================
% END APPENDED SIXTH-SERIES V73 MATERIAL
% =============================================================================

% =============================================================================
% BEGIN APPENDED SIXTH-SERIES V74 MATERIAL
% =============================================================================

\section{V74: representation-ratio test for a Wilson/Casimir tangent}
\label{sec:s6v74}

\subsection{Central one-block diagnostic}
\label{sec:s6v74-central}

V73 reduces the heat-principal route to a same-owner operator identity.
Before attempting the full coupled half-slab, that identity can be tested
on any central convolution block.  This is only a necessary diagnostic:
the physical Wilson half-slab is not asserted to factor into independent
central blocks.

Let \(G\) be a compact group with normalized Haar measure and let
\(k_\theta:G\to\mathbb R\) be a central \(L^2\) kernel.  Define
\begin{equation}
 (A_\theta f)(y)
 =\int_G k_\theta(yx^{-1})f(x)\,\dd x.
 \label{eq:s6v74-convolution}
\end{equation}
For an irreducible representation \(\pi\) of dimension \(d_\pi\), put
\begin{equation}
 a_\pi(\theta)
 ={1\over d_\pi}\int_G
 k_\theta(g)\overline{\chi_\pi(g)}\,\dd g.
 \label{eq:s6v74-Fourier-coefficient}
\end{equation}

\begin{lemma}[Peter--Weyl scalarization]
\label{lem:s6v74-scalarization}
Every matrix coefficient \(D^\pi_{mn}\) is an eigenvector of
\(A_\theta\) with eigenvalue \(a_\pi(\theta)\).  If \(k_\theta\) is real
and inversion invariant, then \(A_\theta\) is self-adjoint and all
\(a_\pi\) are real.
\end{lemma}

\begin{proof}
Centrality makes the Fourier transform
\[
 \widehat k_\theta(\pi)
 =\int k_\theta(g)\pi(g)^*\,\dd g
\]
an intertwiner of the irreducible representation \(\pi\).  Schur's lemma
makes it a scalar, and taking its trace gives
\eqref{eq:s6v74-Fourier-coefficient}.  The convolution identity then
gives the eigenvalue assertion.  A real inversion-invariant convolution
kernel equals its Hilbert-space adjoint.
\end{proof}

Let \(H_C\) be the bi-invariant Casimir operator, normalized by
\begin{equation}
 H_CD^\pi_{mn}=C_\pi D^\pi_{mn},\qquad
 C_{\bf1}=0,\quad C_\pi>0\ (\pi\ne{\bf1}).
 \label{eq:s6v74-Casimir}
\end{equation}

\subsection{Exact projective tangent criterion}
\label{sec:s6v74-criterion}

An overall scalar multiple of \(A_\theta\) cancels from the normalized
transfer.  The correct central version of a heat tangent therefore allows
one scalar normalization velocity.

\begin{theorem}[Representation-ratio criterion]
\label{thm:s6v74-ratio-criterion}
Assume \(a_{\bf1}\ne0\) and \(a_\pi\ne0\) on an interval.  The projective
same-owner identity
\begin{equation}
 A_\theta'
 =\sigma_\theta A_\theta
 -c_\theta\delta H_CA_\theta
 \label{eq:s6v74-projective-heat}
\end{equation}
holds on the Peter--Weyl core if and only if
\begin{align}
 \sigma_\theta
 &={a_{\bf1}'\over a_{\bf1}},                             \label{eq:s6v74-sigma}\\
 \partial_\theta\log{a_\pi\over a_{\bf1}}
 &=-c_\theta\delta C_\pi
 \qquad(\pi\ne{\bf1}).                                    \label{eq:s6v74-ratio}
\end{align}
Equivalently, the quantities
\begin{equation}
 -{1\over\delta C_\pi}
 \partial_\theta\log{a_\pi\over a_{\bf1}}
 \label{eq:s6v74-common-c}
\end{equation}
must have one common value for every nontrivial representation.
\end{theorem}

\begin{proof}
Apply \eqref{eq:s6v74-projective-heat} to each matrix coefficient and use
\Cref{lem:s6v74-scalarization} and \eqref{eq:s6v74-Casimir}.  The trivial
representation gives \eqref{eq:s6v74-sigma}; subtracting its logarithmic
derivative from the nontrivial equations gives
\eqref{eq:s6v74-ratio}.  Conversely those scalar equations verify
\eqref{eq:s6v74-projective-heat} on every Peter--Weyl coefficient, hence
on their algebraic span.
\end{proof}

\begin{corollary}[Two-representation obstruction determinant]
\label{cor:s6v74-two-rep}
Define
\[
 d_\pi=\partial_\theta\log(a_\pi/a_{\bf1}).
\]
For two nontrivial representations \(\pi,\sigma\), an exact projective
heat tangent requires
\begin{equation}
 \boxed{
 \Delta_{\pi,\sigma}
 :=C_\sigma d_\pi-C_\pi d_\sigma=0.}
 \label{eq:s6v74-obstruction}
\end{equation}
For any proposed scalar \(\tau\),
\begin{equation}
 \max\{|d_\pi+\tau C_\pi|,
       |d_\sigma+\tau C_\sigma|\}
 \ge{|\Delta_{\pi,\sigma}|\over C_\pi+C_\sigma}.
 \label{eq:s6v74-residual-lower}
\end{equation}
\end{corollary}

\begin{proof}
Equation \eqref{eq:s6v74-ratio} makes
\(d_\rho=-c\delta C_\rho\), proving the necessary determinant.
For arbitrary \(\tau\),
\[
 \Delta_{\pi,\sigma}
 =C_\sigma(d_\pi+\tau C_\pi)
 -C_\pi(d_\sigma+\tau C_\sigma).
\]
The triangle inequality gives
\[
 |\Delta_{\pi,\sigma}|
 \le(C_\pi+C_\sigma)
 \max\{|d_\pi+\tau C_\pi|,
       |d_\sigma+\tau C_\sigma|\}.
\]
\end{proof}

\begin{remark}[This is a same-owner test]
\label{rem:s6v74-same-owner}
All coefficients in \eqref{eq:s6v74-obstruction} come from the same
kernel \(k_\theta\) and the same derivative.  A Casimir coefficient
computed from a boundary Hessian of another law does not enter the
determinant.
\end{remark}

\subsection{Modewise normalized-transfer response}
\label{sec:s6v74-response}

Assume for simplicity that the relevant coefficients are positive and
that the trivial coefficient is the leading one.  On the \(\pi\)-sector,
the normalized full transfer eigenvalue and gap are
\begin{equation}
 q_\pi=\left({a_\pi\over a_{\bf1}}\right)^2,\qquad
 g_\pi=1-q_\pi.
 \label{eq:s6v74-q}
\end{equation}

\begin{proposition}[Ratio slope is the exact spectral response]
\label{prop:s6v74-mode-response}
Define
\begin{equation}
 r_\pi:=-{1\over\delta}
 \partial_\theta\log{a_\pi\over a_{\bf1}}.
 \label{eq:s6v74-rpi}
\end{equation}
Then
\begin{equation}
 \boxed{g_\pi'=2\delta q_\pi r_\pi.}
 \label{eq:s6v74-gap-response}
\end{equation}
The projective heat tangent is exactly
\(r_\pi=c_\theta C_\pi\) for every \(\pi\ne{\bf1}\).
\end{proposition}

\begin{proof}
Differentiate \eqref{eq:s6v74-q}:
\[
 q_\pi'
 =2q_\pi\partial_\theta\log(a_\pi/a_{\bf1})
 =-2\delta q_\pi r_\pi.
\]
Since \(g_\pi=1-q_\pi\), this proves
\eqref{eq:s6v74-gap-response}.  The last assertion is
\eqref{eq:s6v74-ratio}.
\end{proof}

\begin{corollary}[First-sector positivity is insufficient]
\label{cor:s6v74-one-sector}
Matching one nontrivial representation fixes a candidate \(c_\theta\),
but it proves no whole-carrier DSFI unless every other low-loss
representation has a controlled residual.  A nonzero
\(\Delta_{\pi,\sigma}\) quantitatively lower-bounds the failure of the
single coefficient on at least one of the two sectors.
\end{corollary}

\begin{proof}
Choose \(c_\theta\) from one ratio in
\eqref{eq:s6v74-common-c} and apply
\eqref{eq:s6v74-residual-lower} to the other.
\end{proof}

\subsection{An \(\SU(3)\) Wilson character obstruction at small coupling}
\label{sec:s6v74-SU3}

Let
\[
 X(U)=\operatorname{Re}\operatorname{Tr}U
 ={1\over2}(\chi_{\bf3}(U)+\chi_{\overline{\bf3}}(U))
\]
and consider the central one-plaquette Wilson kernel
\begin{equation}
 k_s(U)=e^{sX(U)},\qquad s>0.
 \label{eq:s6v74-Wilson-kernel}
\end{equation}
Let \(a_{\bf1},a_{\bf3},a_{\bf8}\) be
\eqref{eq:s6v74-Fourier-coefficient}.

\begin{lemma}[First Wilson character coefficients]
\label{lem:s6v74-SU3-expansion}
As \(s\downarrow0\),
\begin{align}
 a_{\bf1}(s)&=1+{s^2\over4}+O(s^3),                       \label{eq:s6v74-a1}\\
 a_{\bf3}(s)&={s\over6}+O(s^2),                           \label{eq:s6v74-a3}\\
 a_{\bf8}(s)&={s^2\over32}+O(s^3).                        \label{eq:s6v74-a8}
\end{align}
Consequently,
\begin{align}
 {d\over ds}\log{a_{\bf3}\over a_{\bf1}}
 &={1\over s}+O(1),                                      \label{eq:s6v74-d3}\\
 {d\over ds}\log{a_{\bf8}\over a_{\bf1}}
 &={2\over s}+O(1).                                      \label{eq:s6v74-d8}
\end{align}
\end{lemma}

\begin{proof}
Expand
\[
 e^{sX}=1+sX+{s^2\over2}X^2+O(s^3)
\]
uniformly on the compact group.  Character orthogonality and
\[
 {\bf3}\otimes\overline{\bf3}={\bf1}\oplus{\bf8}
\]
give
\[
 \int X^2\,\dd U={1\over2}.
\]
Thus the trivial coefficient is
\(1+(s^2/2)(1/2)+O(s^3)\).
For the fundamental coefficient,
\[
 {1\over3}\int sX\,\overline{\chi_{\bf3}}\,\dd U
 ={s\over6},
\]
because only
\(\chi_{\bf3}\chi_{\overline{\bf3}}\) contains the trivial
representation at first order.  In
\[
 X^2={1\over4}
 (\chi_{\bf3}^2+
 2\chi_{\bf3}\chi_{\overline{\bf3}}+
 \chi_{\overline{\bf3}}^2),
\]
the coefficient of \(\chi_{\bf8}\) is \(1/2\).  Multiplication by
\(s^2/2\), integration against \(\chi_{\bf8}\), and division by
\(d_{\bf8}=8\) give \(s^2/32\).
The logarithmic derivatives follow from the three expansions.
\end{proof}

\begin{theorem}[The literal Wilson convolution is not an exact Casimir
semigroup near \(s=0\)]
\label{thm:s6v74-SU3-obstruction}
Use the standard \(\SU(3)\) Casimirs
\[
 C_{\bf3}={4\over3},\qquad C_{\bf8}=3.
\]
For all sufficiently small \(s>0\), the Wilson kernel
\eqref{eq:s6v74-Wilson-kernel} fails the exact projective heat criterion
\eqref{eq:s6v74-ratio}.
\end{theorem}

\begin{proof}
If the criterion held, the ratio of the two nonzero logarithmic
derivatives would equal
\[
 {C_{\bf8}\over C_{\bf3}}={9\over4}.
\]
By \eqref{eq:s6v74-d3}--\eqref{eq:s6v74-d8}, that ratio tends to \(2\)
as \(s\downarrow0\).  Since \(2\ne9/4\), equality is impossible on a
sufficiently small punctured interval.
\end{proof}

\begin{assessment}[Scope of the \(\SU(3)\) obstruction]
\label{ass:s6v74-SU3-scope}
The theorem rules out an \emph{exact} universal heat-semigroup identity
for the literal central Wilson factor.  It does not rule out an
asymptotic heat principal part at weak coupling, nor a useful residual
bound after the full half-slab pushforward.  The physical continuum
trajectory is a weak-coupling trajectory, whereas
\Cref{thm:s6v74-SU3-obstruction} is a small-\(s\) diagnostic.  Its role
is to show that the residual in V73 is structural and must be estimated,
not declared zero.
\end{assessment}

\subsection{The finite Wilson ratio card}
\label{sec:s6v74-card}

\begin{definition}[Wilson representation-ratio card, WRRC]
\label{def:s6v74-WRRC}
At a selected finite half-slab regulator, WRRC consists of:
\begin{enumerate}[label=\textup{(\roman*)},leftmargin=3.5em]
\item one declared central conditional block of the actual amplitude;
\item its coefficients \(a_\pi\) and logarithmic slopes \(d_\pi\);
\item the Casimirs \(C_\pi\) in the same normalization;
\item the mismatch determinants
      \(\Delta_{\pi,\sigma}\) on a representation bank containing every
      low-loss sector;
\item a residual lower-form estimate that charges mismatched
      high-loss sectors to \(L\);
\item proof that the block calculation survives gauge incidence and the
      full Wilsonian pushforward.
\end{enumerate}
\end{definition}

\begin{assessment}[Exact frontier after V74]
\label{ass:s6v74-frontier}
The same-owner question now has a computable local diagnostic:
\eqref{eq:s6v74-obstruction}.  A literal Wilson factor is not exactly a
Casimir heat semigroup, already by the \({\bf3}\)/\({\bf8}\) small-coupling
test.  The viable route is therefore approximate: choose the principal
coefficient from the first physical low-loss sector and prove that all
representation-ratio mismatches belong to the residual debit allowed by
\Cref{cor:s6v73-residual-debit}.  This must be done on the actual
response-matched half-slab; the central block alone does not prove it.
\end{assessment}

\begin{programme}[V75: companion audit and weak-coupling ratio residual]
\label{prog:s6v74-V75}
\begin{enumerate}[leftmargin=3em]
\item Perform the compulsory SM/NS audit.
\item Derive the large-coupling stationary-phase expansion of
      \(d_\pi/C_\pi\) for the \(\SU(3)\) Wilson kernel on a fixed
      low-representation bank.
\item Track the remainder uniformly over the representations that can
      lie in the low-loss cone.
\item Compare the ratio mismatch with the principal stock
      \(2c\kappa\) and the allowed G\aa rding debit.
\item Keep the full Wilsonian pushforward and gauge-incidence transfer as
      separate required rows.
\end{enumerate}
\end{programme}

\begin{firewall}[Claim boundary after sixth-series V74]
\label{fire:s6v74-boundary}
V74 diagonalizes a central same-owner tangent by Peter--Weyl theory;
proves the exact representation-ratio criterion and a quantitative
two-representation obstruction; identifies the ratio slope with the
normalized transfer response; and proves from \(\SU(3)\) character
algebra that the literal one-plaquette Wilson kernel is not exactly a
Casimir semigroup near zero coupling.

V74 does not control the weak-coupling ratio residual, transfer the
central diagnostic through the response-matched half-slab, prove the
residual debit or cofinal SOCC, close the V53 projective/Mosco hypotheses,
construct the continuum OS theory, or prove the full Yang--Mills mass
gap.
\end{firewall}

\begin{theorem}[Sixth-series V74 result ledger]
\label{thm:s6v74-results}
The positive conclusions in \Cref{fire:s6v74-boundary} are exactly those
of
\Cref{lem:s6v74-scalarization,
thm:s6v74-ratio-criterion,
cor:s6v74-two-rep,
prop:s6v74-mode-response,
cor:s6v74-one-sector,
lem:s6v74-SU3-expansion,
thm:s6v74-SU3-obstruction}.
\end{theorem}

\begin{proof}
Each item is the cited result with its centrality, nonvanishing,
representation, coupling-regime, same-owner, gauge, and cofinal
qualifications unchanged.
\end{proof}

\paragraph{Parent-version record.}
V74 was formed from
\texttt{Renewal\_Yang\_Mills\_Mass\_Gap\_Research\_Notes\_V73.tex}.
The V73 parent had \(32\,882\) logical source lines,
\(1\,154\,463\) bytes, and SHA--256
\[
 \texttt{3bbadc3b4f758c4892383b4d33ccebe3699b9299caa247de1df98f03badd851d}.
\]
The next compulsory companion audit is V75.

% =============================================================================
% END APPENDED SIXTH-SERIES V74 MATERIAL
% =============================================================================

% =============================================================================
% BEGIN APPENDED SIXTH-SERIES V75 MATERIAL
% =============================================================================

\section{V75: compulsory companion audit and weak-coupling Wilson ratios}
\label{sec:s6v75}

\subsection{Compulsory companion audit}
\label{sec:s6v75-audit}

The current companion files were inspected read-only:
\begin{align}
 &\texttt{Renewal\_SM\_Einstein\_Continuum\_Research\_Notes\_V17.tex},
 \notag\\
 &\quad 6\,287\ {\rm lines},\quad233\,407\ {\rm bytes},\quad
 \texttt{3ed9f8701c5fe0304260dc386add01aa0a01e5ced70c839e23117a94ad57949f},
 \label{eq:s6v75-SM-file}\\
 &\texttt{Renewal\_NS\_Stability\_Research\_Notes\_V26.tex},
 \notag\\
 &\quad 5\,319\ {\rm lines},\quad278\,283\ {\rm bytes},\quad
 \texttt{82c887f8c2c69e4f28fad7c3dc8de5b9099cb5a4bf431eba1fff3e91935978ad}.
 \label{eq:s6v75-NS-file}
\end{align}
The NS file is unchanged from V70.  The SM programme has advanced from
V14 to V17.

\paragraph{SM V15.}
It proves a conditional \(C^1\) landing theorem for local transfer matrix
elements under a form-locality acquisition card.  The card requires a
common form domain, relative form bounds, local energy moments, one- and
two-mark estimates, spatial and contact rows, cofinal summability, and
physical incidence.  It separately proves that absolute form-locality
rows cannot imply signed spectral coercivity.

\paragraph{SM V16.}
It obtains local vacuum form-energy moments from the exact identity
\[
 \|H^{1/2}A\Omega\|^2=\langle A\Omega,[H,A]\Omega\rangle,
 \qquad H\Omega=0,
\]
and from bounded or form-local commutator estimates.  It also proves that
a raw finite-temperature Gibbs moment retains an extensive background
energy and is not the required vacuum-local row.

\paragraph{SM V17.}
It computes
\[
 \|e^{-tH}\|_{{\cal H}_{-s}\to{\cal H}_s}
 =\sup_{E\in\sigma(H)}(E+\lambda)^s e^{-tE}.
\]
For the natural form exponent \(s=1\), the two-mark smoothing integral
diverges for every unbounded \(H\).  It is integrable for fixed
\(0<s<1\).  The surviving branches are fractional subcriticality, a
fixed positive collar with an explicit logarithmic debit, or a
model-specific contact cancellation proved before absolute values.

\begin{theorem}[V75 companion-audit decision]
\label{thm:s6v75-audit}
The SM V15--V17 results may support future Yang--Mills local transfer
landing in the following conditional ways:
\begin{enumerate}[label=\textup{(\roman*)},leftmargin=3.5em]
\item the V16 commutator identity can populate a local energy row for a
      declared Yang--Mills vacuum carrier;
\item the V15 card can pass already proved local matrix-element and
      parameter-derivative bounds to a route-independent limit;
\item the V17 fractional or cancellation branches can control two
      unbounded form insertions.
\end{enumerate}
They do not prove a same-owner Wilson/Casimir tangent, a positive
representation-ratio slope, a residual debit, or a Yang--Mills transfer
gap.  Nothing is imported from NS V26.
\end{theorem}

\begin{proof}
The V15 sign counterexample and gap firewall explicitly separate its
absolute landing theorem from signed coercivity.  V16 assumes a
zero-energy vacuum and a uniform local commutator row; neither is supplied
for the present cofinal Yang--Mills half-slab.  V17 proves only the
consequences of fractional regularity, a collar, or a combined contact
cancellation, and states those as additional hypotheses.  None of these
statements computes the Wilson character slopes of V74.  The NS file
still concerns Oseen-kernel rank-one norms and contains no reflected
gauge transfer or Wilson character coefficient.
\end{proof}

\begin{assessment}[Audit adjustment]
\label{ass:s6v75-audit-adjustment}
If the Wilson residual is later represented by two coincident
order-one form insertions, the natural absolute smoothing route is now
closed by SM V17.  One must either prove a fractional-order residual,
retain a positive collar, or combine the Casimir contact with the coherent
remainder before taking norms.  This reinforces the same-owner rule of
V73 rather than filling it.
\end{assessment}

\subsection{Large-coupling Laplace expansion for a fixed representation
bank}
\label{sec:s6v75-Laplace}

Retain the central \(\SU(3)\) Wilson kernel
\[
 k_s(U)=e^{s\operatorname{Re}\operatorname{Tr}U}
\]
and the coefficients \(a_\pi(s)\) of
\eqref{eq:s6v74-Fourier-coefficient}.  Normalize Hermitian generators by
\[
 \operatorname{Tr}_{\bf3}(T_aT_b)={1\over2}\delta_{ab},
 \qquad a,b=1,\ldots,8,
\]
so that
\(\sum_aT_a^{(\pi)}T_a^{(\pi)}=C_\pi I\).

\begin{theorem}[Fixed-bank Wilson-to-heat asymptotics]
\label{thm:s6v75-fixed-bank}
Let \({\cal B}\) be a finite set of irreducible \(\SU(3)\)
representations.  As \(s\to\infty\), uniformly for \(\pi\in{\cal B}\),
\begin{align}
 {a_\pi(s)\over a_{\bf1}(s)}
 &=1-{C_\pi\over s}+O_{\cal B}(s^{-2}),                   \label{eq:s6v75-ratio-asymptotic}\\
 \log{a_\pi(s)\over a_{\bf1}(s)}
 &=-{C_\pi\over s}+O_{\cal B}(s^{-2}),                    \label{eq:s6v75-log-asymptotic}\\
 {d\over ds}\log{a_\pi(s)\over a_{\bf1}(s)}
 &={C_\pi\over s^2}+O_{\cal B}(s^{-3}).                   \label{eq:s6v75-slope-asymptotic}
\end{align}
Thus the literal Wilson convolution has the Casimir heat time
\(t_{\rm eff}=s^{-1}\) to leading order on every fixed representation
bank.
\end{theorem}

\begin{proof}
The phase \(\Phi(U)=\operatorname{Re}\operatorname{Tr}U\) has its unique
global maximum \(3\) at \(I\).  In exponential coordinates
\[
 U=\exp(iX),\qquad X=x^aT_a,
\]
one has
\begin{equation}
 \Phi(e^{iX})=3-{1\over4}|x|^2+O(|x|^4),
 \qquad
 \dd U=(1+O(|x|^2))\,\dd x
 \label{eq:s6v75-phase}
\end{equation}
near zero.  Outside a fixed exponential-coordinate neighbourhood,
\(\Phi\le3-\varepsilon\), so that part of every coefficient is
exponentially smaller than the identity saddle.

For a nontrivial irreducible \(\pi\), invariance and the Casimir identity
give
\[
 {1\over d_\pi}
 \operatorname{Tr}_\pi(T_a^{(\pi)}T_b^{(\pi)})
 ={C_\pi\over8}\delta_{ab}.
\]
Therefore
\begin{equation}
 {\overline{\chi_\pi(e^{iX})}\over d_\pi}
 =1-{C_\pi\over16}|x|^2+O_\pi(|x|^3).
 \label{eq:s6v75-character-Taylor}
\end{equation}
The odd contribution integrates to zero after pairing \(X\) and \(-X\)
against the real inversion-invariant Wilson phase and Haar Jacobian.

Scale \(x=s^{-1/2}z\).  The leading Gaussian density is proportional to
\(e^{-|z|^2/4}\), for which
\[
 \mathbb E|z|^2=16.
\]
Dividing the character-weighted saddle integral by the trivial saddle
integral cancels their common Haar and phase corrections.  The
representation-dependent quadratic term in
\eqref{eq:s6v75-character-Taylor} contributes
\[
 -{C_\pi\over16s}\mathbb E|z|^2=-{C_\pi\over s}.
\]
Taylor's theorem on the fixed finite bank, Gaussian domination on a
shrinking coordinate ball, and the exponentially small exterior give a
uniform \(O_{\cal B}(s^{-2})\) remainder.  This proves
\eqref{eq:s6v75-ratio-asymptotic}; logarithmic expansion gives
\eqref{eq:s6v75-log-asymptotic}.  The same dominated saddle expansion
may be differentiated in \(s\); differentiating its first two orders
gives \eqref{eq:s6v75-slope-asymptotic}.
\end{proof}

\begin{corollary}[Fixed-bank same-owner residual]
\label{cor:s6v75-fixed-residual}
Let \(s=s(\theta)\to\infty\) and suppose the interpolation moves toward
larger effective heat time:
\begin{equation}
 -{\dot s(\theta)\over s(\theta)^2}=c_\theta\delta,
 \qquad c_\theta>0.
 \label{eq:s6v75-path}
\end{equation}
For the V74 ratio slope
\[
 d_\pi=\partial_\theta\log(a_\pi/a_{\bf1}),
\]
one has, uniformly on a fixed finite bank,
\begin{equation}
 d_\pi=-c_\theta\delta C_\pi
 +O_{\cal B}\left({c_\theta\delta\over s}\right).
 \label{eq:s6v75-residual}
\end{equation}
Equivalently,
\[
 r_\pi=c_\theta C_\pi+O_{\cal B}(c_\theta/s)
\]
for the exact normalized-transfer response coefficient of V74.
\end{corollary}

\begin{proof}
Multiply \eqref{eq:s6v75-slope-asymptotic} by
\(\dot s=-c_\theta\delta s^2\), then use
\eqref{eq:s6v74-rpi}.
\end{proof}

\begin{corollary}[Positive fixed-bank response]
\label{cor:s6v75-fixed-positive}
Suppose \(c_\theta\ge c_*>0\) and \({\cal B}\) contains no trivial
representation.  For all sufficiently large \(s\), there is
\(m_{\cal B}>0\) such that
\begin{equation}
 2q_\pi r_\pi\ge m_{\cal B}
 \qquad(\pi\in{\cal B}).
 \label{eq:s6v75-fixed-positive}
\end{equation}
\end{corollary}

\begin{proof}
On a fixed bank, \(q_\pi=(a_\pi/a_{\bf1})^2\to1\) uniformly and
\(r_\pi=c_\theta C_\pi+O_{\cal B}(c_\theta/s)\).
The positive finite minimum of the nonzero Casimirs and \(c_\theta\ge c_*\)
give the claim.
\end{proof}

\begin{assessment}[What the saddle theorem does not cover]
\label{ass:s6v75-fixed-only}
The theorem is uniform only on a fixed finite representation bank.  The
whole relative carrier contains arbitrarily high representations, and
the bank relevant to a refining regulator may grow.  A fixed-bank
stationary-phase expansion cannot by itself prove SOCC or DSFI.
\end{assessment}

\subsection{An exact low/high representation-band compiler}
\label{sec:s6v75-band}

For a central diagonal family, decompose a relative vector into its
Peter--Weyl sectors:
\[
 f=\sum_{\pi\ne{\bf1}}f_\pi.
\]
Let \(q_\pi\) be the normalized transfer eigenvalue and \(r_\pi\) the
response coefficient \eqref{eq:s6v74-rpi}.  Then
\begin{align}
 \langle f,Lf\rangle
 &=\sum_\pi(1-q_\pi)\|f_\pi\|^2,                           \label{eq:s6v75-diagonal-L}\\
 \langle f,L'f\rangle
 &=2\delta\sum_\pi q_\pi r_\pi\|f_\pi\|^2.                \label{eq:s6v75-diagonal-C}
\end{align}

\begin{theorem}[Representation-band DSFI compiler]
\label{thm:s6v75-band-compiler}
Partition the nontrivial representations into a low band
\({\cal B}_{\rm lo}\) and high band \({\cal B}_{\rm hi}\).  Suppose
\begin{align}
 2q_\pi r_\pi&\ge m>0
 &&(\pi\in{\cal B}_{\rm lo}),                             \label{eq:s6v75-low-slope}\\
 2q_\pi r_\pi&\ge-M
 &&(\pi\in{\cal B}_{\rm hi}),                             \label{eq:s6v75-high-slope}\\
 1-q_\pi&\ge\ell_{\rm hi}>0
 &&(\pi\in{\cal B}_{\rm hi}).                             \label{eq:s6v75-high-loss}
\end{align}
Then
\begin{equation}
 \boxed{
 L'\succeq m\delta I
 -{(m+M)\delta\over\ell_{\rm hi}}L}
 \label{eq:s6v75-band-DSFI}
\end{equation}
on the relative carrier.
\end{theorem}

\begin{proof}
Let
\[
 h=\sum_{\pi\in{\cal B}_{\rm hi}}\|f_\pi\|^2.
\]
Equations
\eqref{eq:s6v75-low-slope}--\eqref{eq:s6v75-high-slope} give
\[
 \langle f,L'f\rangle
 \ge\delta\{m(\|f\|^2-h)-Mh\}
 =m\delta\|f\|^2-(m+M)\delta h.
\]
By \eqref{eq:s6v75-high-loss},
\[
 \langle f,Lf\rangle\ge\ell_{\rm hi}h.
\]
Substitution proves \eqref{eq:s6v75-band-DSFI}.
\end{proof}

\begin{corollary}[Heat-comparison choice of the high band]
\label{cor:s6v75-heat-band}
If a Wilson block obeys
\begin{equation}
 q_\pi\le e^{-2C_\pi/s+\varepsilon_s}
 \label{eq:s6v75-heat-comparison}
\end{equation}
and the high band is \(C_\pi\ge\eta s\), then
\begin{equation}
 \ell_{\rm hi}\ge1-e^{-2\eta+\varepsilon_s}.
 \label{eq:s6v75-high-loss-value}
\end{equation}
For fixed \(\eta>0\) and \(\varepsilon_s\to0\), this is uniformly
positive.
\end{corollary}

\begin{proof}
Apply \eqref{eq:s6v75-heat-comparison} on the high band and subtract from
one.
\end{proof}

\begin{corollary}[High-band mass inside a low-loss vector]
\label{cor:s6v75-high-mass}
Under \eqref{eq:s6v75-high-loss}, every normalized
\(f\in{\mathfrak C}_{K,\delta}\) satisfies
\begin{equation}
 \sum_{\pi\in{\cal B}_{\rm hi}}\|f_\pi\|^2
 \le{K\delta\over\ell_{\rm hi}}.
 \label{eq:s6v75-high-mass}
\end{equation}
\end{corollary}

\begin{proof}
The high-band contribution to \(\langle f,Lf\rangle\) is at least
\(\ell_{\rm hi}\) times its squared mass, while the cone definition bounds
the whole loss by \(K\delta\).
\end{proof}

\begin{assessment}[The new cofinal target]
\label{ass:s6v75-cofinal-target}
The fixed-bank saddle expansion supplies
\eqref{eq:s6v75-low-slope} on every preassigned collection of
representations.  The band compiler shows what remains:
\begin{enumerate}[label=\textup{(\roman*)},leftmargin=3.5em]
\item a representation-uniform weak-coupling expansion on a growing
      low band, naturally \(C_\pi/s\ll1\);
\item a heat-comparison loss floor on its complement;
\item a clock-scaled negative lower bound for the exact response there;
\item transfer of all three rows from the central block to the actual
      response-matched, gauge-constrained half-slab.
\end{enumerate}
Only after these rows may the V66 G\aa rding theorem be invoked.
\end{assessment}

\subsection{A quantitative growing-bank obligation}
\label{sec:s6v75-growing}

\begin{definition}[Growing representation-ratio card, GRRC]
\label{def:s6v75-GRRC}
For a cofinal sequence \(s_j\to\infty\), GRRC consists of a threshold
\(\eta_j\downarrow0\) with \(\eta_js_j\to\infty\) and:
\begin{align}
 \sup_{0<C_\pi\le\eta_js_j}
 \left|
 {1\over C_\pi}{d\over ds}
 \log{a_\pi(s_j)\over a_{\bf1}(s_j)}
 -{1\over s_j^2}
 \right|
 &\le{\varepsilon_j\over s_j^2},
 \qquad\varepsilon_j\to0,                                \label{eq:s6v75-GRRC-low}\\
 \sup_{C_\pi>\eta_js_j}q_{\pi,j}
 &\le1-\ell_j,                                            \label{eq:s6v75-GRRC-high}\\
 \inf_{C_\pi>\eta_js_j}2q_{\pi,j}r_{\pi,j}
 &\ge-M.                                                  \label{eq:s6v75-GRRC-lower}
\end{align}
The usable cofinal form requires
\[
 \inf_j\ell_j>0
\]
or, more generally, bounded
\(\delta_j/\ell_j\) in the G\aa rding debit after the interpolation is
integrated.
\end{definition}

\begin{proposition}[GRRC closes the central spectral row]
\label{prop:s6v75-GRRC}
Assume the path clock
\eqref{eq:s6v75-path} with \(c_\theta\ge c_*>0\), and suppose GRRC holds
uniformly in \(\theta\).  If the smallest nonzero Casimir is \(C_*>0\),
then for all sufficiently large \(j\) the central convolution transfer
satisfies DSFI with a positive \(j\)-independent \(b\), provided the
high-band debit in \eqref{eq:s6v75-band-DSFI} has bounded integrated
coefficient.
\end{proposition}

\begin{proof}
Equation \eqref{eq:s6v75-GRRC-low} and the path clock give
\[
 r_{\pi,j}
 \ge c_*(1-\varepsilon_j)C_\pi
 \]
on the low band.  Its coefficients also obey \(q_{\pi,j}\to1\) uniformly
on every subband with \(C_\pi/s_j\to0\); this uniform convergence is part
of the same GRRC weak-coupling control.  Hence
\[
 2q_{\pi,j}r_{\pi,j}\ge c_*C_*
\]
for large \(j\), after decreasing the constant once.  The high rows of
GRRC give \eqref{eq:s6v75-high-slope}--\eqref{eq:s6v75-high-loss}.
Apply \Cref{thm:s6v75-band-compiler}; its \(b\) is positive and its loss
coefficient is integrable by hypothesis.
\end{proof}

\begin{firewall}[Central GRRC is not the physical half-slab]
\label{fire:s6v75-central}
Even a proof of GRRC for the one-plaquette central Wilson kernel would not
establish the Yang--Mills mass gap.  The response-matched half-slab
contains shared links, spatial plaquettes, gauge projection, generated
Wilsonian returns, and moving ground factors.  V73's same-owner residual
must control the difference between that amplitude and the central
diagnostic without consuming the positive stock.
\end{firewall}

\begin{assessment}[Exact frontier after V75]
\label{ass:s6v75-frontier}
The weak-coupling calculation repairs the small-coupling obstruction of
V74 on every fixed representation bank: Wilson character ratios become
Casimir heat ratios to relative error \(O(1/s)\).  The infinite-carrier
problem has not disappeared; it has become the growing-bank card GRRC and
the same-owner transfer from a central block to the full half-slab.
The next upstream task is to make the stationary-phase remainder uniform
for \(C_\pi/s\le\eta\), using Casimir-controlled character derivatives,
and to determine whether a high-band loss floor survives the Wilsonian
pushforward.
\end{assessment}

\begin{programme}[V76: Casimir-uniform character remainder]
\label{prog:s6v75-V76}
\begin{enumerate}[leftmargin=3em]
\item Bound normalized character derivatives near the identity by
      explicit polynomials in \(C_\pi\).
\item Run the Wilson saddle expansion uniformly on
      \(C_\pi\le\eta s\), retaining the dependence on \(\eta\).
\item Optimize the coordinate radius and exponential exterior tail.
\item Derive a uniform ratio-slope error or construct a representation
      sequence that violates GRRC.
\item Keep the full half-slab transfer as a subsequent, separate row.
\end{enumerate}
\end{programme}

\begin{firewall}[Claim boundary after sixth-series V75]
\label{fire:s6v75-boundary}
V75 performs the compulsory SM/NS audit; proves the large-coupling
Wilson-to-Casimir asymptotics on every fixed \(\SU(3)\) representation
bank; derives the clock-scaled same-owner residual there; proves an exact
low/high representation-band DSFI compiler; and formulates the growing
representation-ratio card.

V75 does not prove the growing-bank uniform estimates, transfer the
central calculation to the response-matched half-slab, prove the
same-owner residual debit, close the V53 projective/Mosco hypotheses,
construct the continuum OS theory, or prove the full Yang--Mills mass
gap.
\end{firewall}

\begin{theorem}[Sixth-series V75 result ledger]
\label{thm:s6v75-results}
The positive conclusions in \Cref{fire:s6v75-boundary} are exactly those
of
\Cref{thm:s6v75-audit,
thm:s6v75-fixed-bank,
cor:s6v75-fixed-residual,
cor:s6v75-fixed-positive,
thm:s6v75-band-compiler,
cor:s6v75-heat-band,
cor:s6v75-high-mass,
prop:s6v75-GRRC}.
\end{theorem}

\begin{proof}
Each item is the cited result with its audit, finite-bank,
stationary-phase, centrality, clock, high-band, same-owner, and cofinal
qualifications unchanged.
\end{proof}

\paragraph{Parent-version record.}
V75 was formed from
\texttt{Renewal\_Yang\_Mills\_Mass\_Gap\_Research\_Notes\_V74.tex}.
The V74 parent had \(33\,290\) logical source lines,
\(1\,168\,260\) bytes, and SHA--256
\[
 \texttt{dfbe042a92ab18a566979199013dc609a1052f4f3ff7a6a4ef7713ba5881bf49}.
\]
The next compulsory companion audit is V80.

% =============================================================================
% END APPENDED SIXTH-SERIES V75 MATERIAL
% =============================================================================

% =============================================================================
% BEGIN APPENDED SIXTH-SERIES V76 MATERIAL
% =============================================================================

\section{V76: Casimir-uniform Wilson saddle control and the true band scale}
\label{sec:s6v76}

\subsection{Purpose and source-scope audit}
\label{sec:s6v76-purpose}

V75 reduced the central spectral row to a growing representation problem.
This iteration proves the low-band part of that problem with constants
independent of the irreducible representation.  It also corrects an
overstrong possible reading of the V75 high-band card: when the boundary
between the bands is \(C_\pi=\eta_s s\) with \(\eta_s\downarrow0\), even
the exact heat kernel has high-band loss only of order \(\eta_s\), not a
positive constant.

Two earlier Yang--Mills source results were checked before proceeding.
The boundary-sector theorem labelled
\texttt{thm:V27-strict-one-link} proves, at each finite coupling and
fixed exterior configuration,
\[
 \|A_{e,\rho}^{\pm}\|_{\rm op}<1
\]
for every nontrivial finite-dimensional irreducible \(\rho\).  Its proof
uses full support and the mean-unitary variance identity.  It does not
state a bound uniform in \(\rho\), the exterior configuration, or the
cofinal weak-coupling parameter.  The reduced-fibre source theorem
labelled \texttt{thm:YMA25-positive-character} proves an entire
nonnegative character expansion of the Wilson plaquette weight, but its
ledger explicitly leaves the weak-coupling uniform remainder open.

\begin{theorem}[Exact scope of the imported source facts]
\label{thm:s6v76-source-scope}
For the central Wilson kernel used below, every normalized Fourier
coefficient is nonnegative.  At each fixed \(s<\infty\), every
nontrivial normalized coefficient is strictly smaller than one in
absolute value.  Neither source theorem supplies a cofinal number
\(\ell>0\) such that
\[
 1-q_\pi\ge\ell
\]
on a representation band depending on \(s\).
\end{theorem}

\begin{proof}
The positive character expansion has nonnegative scalar coefficients.
Character orthogonality identifies those coefficients, up to the
positive dimension convention in
\eqref{eq:s6v74-Fourier-coefficient}, with the \(a_\pi(s)\); hence
\(a_\pi(s)\ge0\).  Normalize the Wilson density to a probability.  If a
nontrivial unitary irreducible had normalized mean of modulus one, the
mean-unitary variance identity and the strictly positive density would
make a nonzero vector invariant under every group element, a
contradiction.  This is a strict statement at the given \(s\).  Its
proof contains no quantitative constant uniform along \(s\to\infty\),
so the last assertion is a statement about logical scope, not a new
estimate.
\end{proof}

\subsection{A representation-uniform character remainder}
\label{sec:s6v76-character}

Retain the generator normalization of V75 and write
\[
 X=x^aT_a,\qquad |x|^2=\sum_{a=1}^{8}(x^a)^2.
\]
For an irreducible representation \(\pi\), let
\[
 A_\pi(x):=x^aT_a^{(\pi)},\qquad
 h_\pi(x):={1\over d_\pi}\Re\chi_\pi(e^{iX}).
\]

\begin{lemma}[Operator Casimir bound]
\label{lem:s6v76-operator-Casimir}
For every irreducible \(\pi\) and every \(x\in\mathbb R^8\),
\begin{equation}
 \boxed{\|A_\pi(x)\|_{\rm op}\le\sqrt{C_\pi}\,|x|.}
 \label{eq:s6v76-operator-Casimir}
\end{equation}
\end{lemma}

\begin{proof}
For a unit vector \(v\), Cauchy--Schwarz gives
\[
 \|A_\pi(x)v\|^2
 \le |x|^2\sum_{a=1}^{8}\|T_a^{(\pi)}v\|^2
 =|x|^2
 \left\langle v,\sum_a(T_a^{(\pi)})^2v\right\rangle
 =C_\pi|x|^2.
\]
Take the supremum over \(v\).
\end{proof}

\begin{theorem}[Uniform normalized-character Taylor inequality]
\label{thm:s6v76-character-remainder}
For every irreducible unitary representation of \(\SU(3)\) and every
\(x\in\mathbb R^8\),
\begin{equation}
 \boxed{
 \left|
 h_\pi(x)-1+{C_\pi\over16}|x|^2
 \right|
 \le {C_\pi^2\over24}|x|^4.}
 \label{eq:s6v76-character-remainder}
\end{equation}
No constant in this inequality depends on \(d_\pi\).
\end{theorem}

\begin{proof}
Since \(A_\pi(x)\) is Hermitian,
\[
 h_\pi(x)={1\over d_\pi}\Tr_\pi\cos A_\pi(x).
\]
Invariance of the trace form and the Casimir identity give
\[
 {1\over d_\pi}\Tr_\pi A_\pi(x)^2
 ={C_\pi\over8}|x|^2.
\]
For every real \(t\),
\[
 |\cos t-1+t^2/2|\le |t|^4/24.
\]
Apply this scalar inequality through the spectral calculus, average the
eigenvalues, and use
\[
 {1\over d_\pi}\Tr |A_\pi(x)|^4
 \le\|A_\pi(x)\|_{\rm op}^4
 \le C_\pi^2|x|^4
\]
from \Cref{lem:s6v76-operator-Casimir}.  The quadratic trace gives the
coefficient \(C_\pi/16\).
\end{proof}

\begin{remark}[Why real parts are essential]
\label{rem:s6v76-real-part}
The complex normalized character has a representation-dependent cubic
term.  The Wilson density is inversion invariant, and the coefficient
\(a_\pi\) is real, so its ratio is the expectation of the real part.
Passing to \(h_\pi\) removes all odd powers before an estimate is made;
no cancellation of an uncontrolled cubic remainder is being assumed.
\end{remark}

\subsection{Representation-free Wilson saddle moments}
\label{sec:s6v76-moments}

Let
\begin{equation}
 \nu_s(\dd U)
 ={e^{s\Phi(U)}\over Z_s}\,\dd U,\qquad
 \Phi(U)=\Re\Tr_{\bf3}U,\qquad
 Z_s=\int_{\SU(3)}e^{s\Phi(V)}\,\dd V.
 \label{eq:s6v76-Wilson-law}
\end{equation}
Choose a fixed exponential chart \({\cal N}\) about \(I\), write
\(U=e^{iX}\) there, and put \(Q(U)=|x|^2\) on \({\cal N}\).

\begin{lemma}[Eight-dimensional saddle card]
\label{lem:s6v76-saddle-card}
There are constants \(s_0,K,c>0\), depending only on the fixed chart,
such that for \(s\ge s_0\),
\begin{align}
 \nu_s({\cal N}^{\,c})&\le Ke^{-cs},                     \label{eq:s6v76-tail}\\
 \mathbb E_s[Q\,1_{\cal N}]
 &= {16\over s}+O(s^{-2}),                               \label{eq:s6v76-Q1}\\
 \mathbb E_s[Q^m1_{\cal N}]
 &\le K_m s^{-m}\qquad(m=2,3),                           \label{eq:s6v76-Qm}\\
 \operatorname{Cov}_s(Q1_{\cal N},\Phi)
 &= -{16\over s^2}+O(s^{-3}).                            \label{eq:s6v76-Qcov}
\end{align}
The constants in the two \(O(\cdot)\) terms are numerical chart
constants and have no representation index.
\end{lemma}

\begin{proof}
The maximum of \(\Phi\) is uniquely attained at \(I\).  Shrinking the
chart if necessary, there are \(r,\alpha,\varepsilon>0\) such that
\[
 \Phi(e^{iX})
 =3-\frac14Q+P_4(x)+O(|x|^6),\qquad
 \Phi(e^{iX})\le3-\alpha Q\quad(|x|\le r),
\]
where \(P_4\) is homogeneous of degree four, while
\(\Phi\le3-\varepsilon\) outside the chart.  The exponential-coordinate
Haar density is a smooth even function
\[
 J(x)=1+J_2(x)+O(|x|^4).
\]
The lower bound obtained by integrating over \(|x|\le s^{-1/2}\) and
the exterior phase gap prove \eqref{eq:s6v76-tail}.

On the chart set \(x=s^{-1/2}z\).  After removing the factor \(e^{3s}\),
the numerator and denominator are dominated by a fixed multiple of
\(e^{-\alpha|z|^2}\), and on every ball
\(|z|\le r\sqrt{s}\) their density ratio has the expansion
\[
 e^{-|z|^2/4}
 \left[
 1+s^{-1}\{P_4(z)+J_2(z)\}
 +O\!\left(s^{-2}(1+|z|^{12})\right)
 \right].
\]
The complement of \(|z|\le s^{1/10}\) is smaller than every displayed
power by Gaussian domination; on the inner ball the remainder is
uniform.  Termwise integration therefore gives a full differentiable
Laplace expansion through the orders used below.

For the normalized Gaussian
\(\gamma(\dd z)\propto e^{-|z|^2/4}\dd z\) in eight dimensions,
\[
 \mathbb E_\gamma|z|^2=16,\qquad
 \operatorname{Var}_\gamma(|z|^2)=64.
\]
These values give \eqref{eq:s6v76-Q1} and
\eqref{eq:s6v76-Qm}.  Finally,
\[
 Q=s^{-1}|z|^2,\qquad
 \Phi=3-\frac1{4s}|z|^2+O(s^{-2}|z|^4).
\]
Thus the leading covariance is
\[
 -{1\over4s^2}\operatorname{Var}_\gamma(|z|^2)
 =-{16\over s^2}.
\]
The next term of the dominated expansion is \(O(s^{-3})\), proving
\eqref{eq:s6v76-Qcov}.
\end{proof}

\subsection{The uniform growing-low-band theorem}
\label{sec:s6v76-low-band}

Define the normalized Wilson ratio
\begin{equation}
 R_\pi(s):={a_\pi(s)\over a_{\bf1}(s)}
 =\int_{\SU(3)}{1\over d_\pi}\Re\chi_\pi(U)\,\nu_s(\dd U).
 \label{eq:s6v76-R}
\end{equation}

\begin{theorem}[Casimir-uniform ratio and slope bounds]
\label{thm:s6v76-uniform-ratio}
There are universal constants \(K,c,s_0>0\) such that, for every
irreducible \(\pi\) and \(s\ge s_0\),
\begin{align}
 \left|R_\pi(s)-1+{C_\pi\over s}\right|
 &\le K\left\{
 {C_\pi+C_\pi^2\over s^2}+e^{-cs}\right\},               \label{eq:s6v76-R-bound}\\
 \left|R_\pi'(s)-{C_\pi\over s^2}\right|
 &\le K\left\{
 {C_\pi+C_\pi^2\over s^3}+e^{-cs}\right\}.               \label{eq:s6v76-Rprime-bound}
\end{align}
The estimates are useful, rather than merely true, when
\(C_\pi/s\) is small.
\end{theorem}

\begin{proof}
On \({\cal N}\), insert
\eqref{eq:s6v76-character-remainder} into
\eqref{eq:s6v76-R}.  Since normalized characters have modulus at most
one, \eqref{eq:s6v76-tail} controls the exterior.  Equations
\eqref{eq:s6v76-Q1} and \eqref{eq:s6v76-Qm} then give
\[
 R_\pi
 =1-{C_\pi\over16}\mathbb E_s[Q1_{\cal N}]
 +O(C_\pi^2s^{-2})+O(e^{-cs}),
\]
which is \eqref{eq:s6v76-R-bound}.

Differentiation under the compact-group integral gives the exact
identity
\begin{equation}
 R_\pi'(s)=\operatorname{Cov}_s(h_\pi,\Phi).
 \label{eq:s6v76-cov-identity}
\end{equation}
The quadratic character term and
\eqref{eq:s6v76-Qcov} contribute
\[
 -{C_\pi\over16}
 \operatorname{Cov}_s(Q1_{\cal N},\Phi)
 ={C_\pi\over s^2}+O(C_\pi s^{-3}).
\]
Write the local character remainder as \(\rho_\pi\).  On the chart,
\[
 |\rho_\pi|\le {C_\pi^2\over24}Q^2.
\]
The saddle expansion also gives
\[
 |\Phi-\mathbb E_s\Phi|
 \le K(Q+s^{-1})
\]
in expectation on the chart, up to an exponentially small exterior
term.  Consequently,
\[
 |\operatorname{Cov}_s(\rho_\pi1_{\cal N},\Phi)|
 \le K C_\pi^2
 \mathbb E_s[Q^2(Q+s^{-1})1_{\cal N}]
 \le K C_\pi^2s^{-3}
\]
by \eqref{eq:s6v76-Qm}.  The bounded exterior character and the phase
gap cost \(O(e^{-cs})\).  This proves
\eqref{eq:s6v76-Rprime-bound}.
\end{proof}

\begin{corollary}[Growing-bank logarithmic slope]
\label{cor:s6v76-growing-slope}
There are \(\eta_0>0\), \(K_0<\infty\), and \(s_1\) such that, whenever
\(\pi\ne{\bf1}\), \(s\ge s_1\), and \(C_\pi\le\eta_0s\),
\[
 R_\pi(s)\ge\frac12
\]
and
\begin{equation}
 \boxed{
 \left|
 {1\over C_\pi}{\dd\over\dd s}\log R_\pi(s)-{1\over s^2}
 \right|
 \le {K_0\over s^2}
 \left({1\over s}+{C_\pi\over s}\right).}
 \label{eq:s6v76-log-slope}
\end{equation}
Therefore, for every \(\eta_s\downarrow0\) with
\(\eta_ss\to\infty\), the low row
\eqref{eq:s6v75-GRRC-low} holds on
\(0<C_\pi\le\eta_ss\), with
\[
 \varepsilon_s\le K_0(\eta_s+s^{-1})\longrightarrow0.
\]
\end{corollary}

\begin{proof}
For \(\SU(3)\), nontrivial Casimirs have a positive minimum
\(C_*=4/3\).  Choose \(\eta_0\) small and then \(s_1\) large in
\eqref{eq:s6v76-R-bound}; it gives \(R_\pi\ge1/2\) throughout the stated
band.  Now
\[
 \left|{R_\pi'\over R_\pi}-{C_\pi\over s^2}\right|
 \le 2\left|R_\pi'-{C_\pi\over s^2}\right|
 +{2C_\pi\over s^2}|1-R_\pi|.
\]
Use \eqref{eq:s6v76-R-bound}--\eqref{eq:s6v76-Rprime-bound}, divide by
\(C_\pi\), and absorb the exponential term using \(C_\pi\ge C_*\).
After reducing \(\eta_0\), all quadratic products are bounded by the
right side of \eqref{eq:s6v76-log-slope}.  Taking the supremum over
\(C_\pi\le\eta_ss\) proves the last assertion.
\end{proof}

\begin{corollary}[Low-band same-owner response]
\label{cor:s6v76-low-response}
Under the V75 clock
\(-\dot s/s^2=c_\theta\delta\), with
\(0<c_*\le c_\theta\le c^*<\infty\), one has uniformly for
\(0<C_\pi\le\eta_ss\),
\begin{equation}
 r_\pi
 =c_\theta C_\pi
 \left[1+O(\eta_s+s^{-1})\right].
 \label{eq:s6v76-low-response}
\end{equation}
In particular the low-slope row
\eqref{eq:s6v75-low-slope} holds for all large \(s\), with a positive
constant independent of the growing bank.
\end{corollary}

\begin{proof}
Multiply \eqref{eq:s6v76-log-slope} by
\(-\dot s/\delta=c_\theta s^2\) and use
\eqref{eq:s6v74-rpi}.  Also
\eqref{eq:s6v76-R-bound} gives \(q_\pi=R_\pi^2\to1\) uniformly on the
band.  Since \(C_\pi\ge C_*\), the quantity
\(2q_\pi r_\pi\) has a positive uniform lower bound for large \(s\).
\end{proof}

\subsection{The high-band scale cannot be constant}
\label{sec:s6v76-high-scale}

\begin{proposition}[Exact heat model forces a vanishing band loss]
\label{prop:s6v76-heat-band-obstruction}
Let the normalized half-transfer be the exact heat multiplier
\[
 R_\pi^{\rm heat}(s)=e^{-C_\pi/s},
 \qquad q_\pi^{\rm heat}(s)=e^{-2C_\pi/s}.
\]
If \(\eta_s\downarrow0\) and \(\eta_ss\to\infty\), then
\begin{equation}
 \sup_{C_\pi>\eta_ss}q_\pi^{\rm heat}(s)\longrightarrow1.
 \label{eq:s6v76-heat-sup}
\end{equation}
Thus no high-band loss floor bounded away from zero can coexist with
this growing-low-band choice, even for the ideal heat channel.
\end{proposition}

\begin{proof}
The irreducible \(\SU(3)\) representation of highest weight \((n,0)\)
has
\[
 C_{(n,0)}={n^2+3n\over3}.
\]
Let \(n(s)\) be the least integer for which
\(C_{(n(s),0)}>\eta_ss\).  Since
\[
 C_{(n+1,0)}-C_{(n,0)}={2n+4\over3}
 =O(\sqrt{\eta_ss}),
\]
the hypothesis \(\eta_ss\to\infty\) implies
\[
 {C_{(n(s),0)}\over\eta_ss}\longrightarrow1.
\]
Consequently
\[
 q_{(n(s),0)}^{\rm heat}
 =\exp\!\left[
 -2\eta_s{C_{(n(s),0)}\over\eta_ss}
 \right]\longrightarrow1.
\]
This member belongs to the high band and proves
\eqref{eq:s6v76-heat-sup}.
\end{proof}

\begin{proposition}[Moderate-band Wilson loss]
\label{prop:s6v76-moderate-loss}
After decreasing the constant \(\eta_0\) in
\Cref{cor:s6v76-growing-slope}, there is \(s_2\) such that
\begin{equation}
 1-q_\pi(s)\ge {C_\pi\over2s}
 \label{eq:s6v76-moderate-loss}
\end{equation}
whenever \(s\ge s_2\) and
\[
 0<C_\pi\le\eta_0s.
\]
Hence the intermediate band
\(\eta_ss<C_\pi\le\eta_0s\) has
\[
 1-q_\pi(s)\ge{\eta_s\over2}.
\]
\end{proposition}

\begin{proof}
Positivity from \Cref{thm:s6v76-source-scope} and
\eqref{eq:s6v76-R-bound} imply, after choosing \(\eta_0\) small and
\(s_2\) large,
\[
 0\le R_\pi(s)\le1-{C_\pi\over2s}.
\]
Because \(q_\pi=R_\pi^2\) and \(0\le R_\pi\le1\),
\[
 1-q_\pi=(1-R_\pi)(1+R_\pi)\ge1-R_\pi
 \ge{C_\pi\over2s}.
\]
\end{proof}

\begin{lemma}[What fixed-card compactness actually gives]
\label{lem:s6v76-fixed-compact-gap}
For every fixed \(s<\infty\),
\begin{equation}
 \ell(s):=
 1-\sup_{\pi\ne{\bf1}}R_\pi(s)^2>0.
 \label{eq:s6v76-fixed-compact-gap}
\end{equation}
Nevertheless \(\ell(s)=O(s^{-1})\) as \(s\to\infty\).
\end{lemma}

\begin{proof}
The normalized Wilson convolution has an \(L^2\) kernel on the compact
group, hence is Hilbert--Schmidt and compact.  Its nontrivial
Peter--Weyl eigenvalues therefore tend to zero outside a finite set.
Each is strictly smaller than one in modulus by
\Cref{thm:s6v76-source-scope}; taking the maximum over the remaining
finite set proves \eqref{eq:s6v76-fixed-compact-gap}.

For the fundamental representation,
\eqref{eq:s6v76-R-bound} gives
\[
 1-R_{\bf3}(s)^2={2C_{\bf3}\over s}+O(s^{-2}).
\]
Since \(\ell(s)\le1-R_{\bf3}(s)^2\), the asserted upper order follows.
\end{proof}

\begin{assessment}[Corrected GRRC debit]
\label{ass:s6v76-corrected-GRRC}
The V75 alternative allowing a controlled
\(\delta_s/\ell_s\) is the viable one.  A constant
\(\inf_s\ell_s>0\) is not compatible with a boundary
\(\eta_s\downarrow0\), even in the exact heat model.  V76 proves
\(\ell_s\gtrsim\eta_s\) only on
\[
 \eta_ss<C_\pi\le\eta_0s.
\]
Thus this portion of the V75 G\aa rding debit is controlled if
\(\delta_s/\eta_s\) is bounded, and vanishes if
\(\delta_s/\eta_s\to0\).  The ultraviolet band
\[
 C_\pi>\eta_0s
\]
still requires a cofinal Fourier-decay or heat-comparison estimate; the
fixed-card strict theorem cannot replace it.
\end{assessment}

\begin{firewall}[Central estimates are not yet half-slab estimates]
\label{fire:s6v76-central-only}
The proof in V76 is for the literal one-plaquette central Wilson
convolution.  It neither shows that the response-matched half-slab is
central nor bounds the difference caused by shared links, spatial
plaquettes, gauge projection, generated Wilsonian interactions, and
moving ground factors.  Its role is to populate the exact low-band row
of the V75 compiler and to identify the scale of the remaining
ultraviolet row.
\end{firewall}

\begin{assessment}[Exact frontier after V76]
\label{ass:s6v76-frontier}
The representation-dependent \(O_\pi(\cdot)\) in V75 has been replaced
by explicit Casimir dependence.  This closes the growing-low-band
ratio-slope card for every choice
\(\eta_s\downarrow0\), \(\eta_ss\to\infty\).  It also proves a loss of
order \(\eta_s\) on the moderate band.  The next central obstruction is
now sharply confined to \(C_\pi>\eta_0s\): one needs both a uniform loss
floor there and a clock-scaled lower bound for the exact logarithmic
response.  After that, the still harder task is transfer to the
response-matched physical half-slab.
\end{assessment}

\begin{programme}[V77: ultraviolet Fourier loss and response]
\label{prog:s6v76-V77}
\begin{enumerate}[leftmargin=3em]
\item Use highest-weight or Sobolev Fourier estimates to bound
      \(R_\pi(s)\) uniformly when \(C_\pi>\eta_0s\), retaining the
      \(s\)-dependence of every derivative norm.
\item Differentiate the same estimate and determine whether
      \(2q_\pi r_\pi\) has a clock-independent negative lower bound.
\item If absolute Sobolev estimates lose exponentially in \(s\), replace
      them by a local oscillatory estimate at the Wilson saddle.
\item Combine the ultraviolet and moderate bands only after both loss
      and response are owned by the same Wilson kernel.
\end{enumerate}
\end{programme}

\begin{firewall}[Claim boundary after sixth-series V76]
\label{fire:s6v76-boundary}
V76 proves an operator-Casimir bound and a dimension-free normalized
character remainder; proves representation-uniform Wilson ratio and
ratio-slope estimates; closes the growing-low-band GRRC row; proves an
order-\(\eta_s\) moderate-band loss; proves that a constant high-band
floor is impossible for a shrinking band threshold even in the exact
heat model; and distinguishes fixed-card strictness from cofinal
quantitative contraction.

V76 does not prove ultraviolet Wilson loss or response, transfer the
central calculation to the response-matched half-slab, prove the
same-owner residual debit, close the V53 projective/Mosco hypotheses,
construct the continuum OS theory, or prove the full Yang--Mills mass
gap.
\end{firewall}

\begin{theorem}[Sixth-series V76 result ledger]
\label{thm:s6v76-results}
The positive conclusions in \Cref{fire:s6v76-boundary} are exactly those
of
\Cref{thm:s6v76-source-scope,
lem:s6v76-operator-Casimir,
thm:s6v76-character-remainder,
lem:s6v76-saddle-card,
thm:s6v76-uniform-ratio,
cor:s6v76-growing-slope,
cor:s6v76-low-response,
prop:s6v76-heat-band-obstruction,
prop:s6v76-moderate-loss,
lem:s6v76-fixed-compact-gap}.
\end{theorem}

\begin{proof}
Each item is the cited statement with its central-kernel, Casimir-band,
large-\(s\), clock, and source-scope qualifications unchanged.
\end{proof}

\paragraph{Parent-version record.}
V76 was formed from
\texttt{Renewal\_Yang\_Mills\_Mass\_Gap\_Research\_Notes\_V75.tex}.
The V75 parent had \(33\,789\) logical source lines,
\(1\,186\,484\) bytes, and SHA--256
\[
 \texttt{2a3d9ed012ec77d3fafe136cad6dafc6c610a9cc288e7777332029c23782a5e7}.
\]
The next compulsory companion audit remains V80.

% =============================================================================
% END APPENDED SIXTH-SERIES V76 MATERIAL
% =============================================================================

% =============================================================================
% BEGIN APPENDED SIXTH-SERIES V77 MATERIAL
% =============================================================================

\section{V77: heat minorization, ultraviolet loss, and response localization}
\label{sec:s6v77}

\subsection{A uniform heat component inside the Wilson law}
\label{sec:s6v77-minorization}

Let \(d(\cdot,\cdot)\) be the bi-invariant Riemannian distance whose
positive Laplacian is the Casimir \(H_C\) of
\eqref{eq:s6v74-Casimir}.  Denote by \(p_t(U)\) the normalized heat
kernel:
\[
 e^{-tH_C}f(U)=\int_{\SU(3)}p_t(UV^{-1})f(V)\,\dd V.
\]

\begin{lemma}[Compact-group Gaussian heat bound]
\label{lem:s6v77-heat-Gaussian}
There are constants \(A,B,t_0>0\) such that
\begin{equation}
 p_t(U)\le A t^{-4}
 \exp\!\left[-{d(I,U)^2\over Bt}\right]
 \label{eq:s6v77-heat-Gaussian}
\end{equation}
for \(0<t\le t_0\) and every \(U\in\SU(3)\).
\end{lemma}

\begin{proof}
In a normal neighbourhood of \(I\), the heat parametrix has leading
term
\[
 (4\pi t)^{-4}
 e^{-d(I,U)^2/(4t)}
\]
times a smooth bounded Jacobian factor, while every differentiated
remainder satisfies the same estimate with a slightly enlarged
Gaussian denominator.  On the complement of a smaller normal
neighbourhood, apply the semigroup identity at time \(t/2\), the local
bound in two balls, and compactness; the positive separation of the
balls supplies an exponential \(e^{-c/t}\).  Increasing \(A\) and \(B\)
combines the two regions.  This is the usual finite-parametrix proof,
and all constants are fixed because the group and metric are fixed.
\end{proof}

Put
\[
 V(U):=3-\Re\Tr_{\bf3}U.
\]
Then \(V\ge0\), with equality only at \(I\).

\begin{lemma}[Distance comparison and Wilson normalization]
\label{lem:s6v77-Wilson-geometry}
There are constants \(\rho,K_1,K_2>0\) such that
\begin{align}
 d(I,U)^2&\ge\rho V(U),                                  \label{eq:s6v77-distance-V}\\
 K_1e^{3s}s^{-4}
 \le Z_s&\le K_2e^{3s}s^{-4}                             \label{eq:s6v77-Z-bounds}
\end{align}
for every \(U\in\SU(3)\) and all sufficiently large \(s\).
\end{lemma}

\begin{proof}
Near \(I\), both \(d(I,e^{iX})^2\) and \(V(e^{iX})\) are positive
definite quadratic forms plus \(O(|X|^3)\); in the normalization of V75,
\(V(e^{iX})=|x|^2/4+O(|x|^4)\).  Their quotient therefore has a positive
lower limit.  Away from \(I\), compactness and the strict positivity of
both functions give a positive lower bound.  This proves
\eqref{eq:s6v77-distance-V}.

For the lower normalization bound, integrate \(e^{-sV}\) over a
coordinate ball of radius \(s^{-1/2}\), on which \(sV\) is bounded.
For the upper bound, use \(V(e^{iX})\ge c|x|^2\) in a fixed chart and
the positive phase gap outside it.  The eight-dimensional Gaussian
integral is of order \(s^{-4}\).  Restoring \(e^{3s}\) proves
\eqref{eq:s6v77-Z-bounds}.
\end{proof}

\begin{theorem}[Uniform heat minorization of the Wilson probability]
\label{thm:s6v77-heat-minorization}
There exist constants
\[
 \kappa>0,\qquad \epsilon\in(0,1),\qquad s_0<\infty
\]
such that, for all \(s\ge s_0\),
\begin{equation}
 \boxed{\nu_s(\dd U)\ge
 \epsilon\,p_{\kappa/s}(U)\,\dd U.}
 \label{eq:s6v77-heat-minorization}
\end{equation}
Consequently there is a central probability measure \(\rho_s\) with
\begin{equation}
 \nu_s=\epsilon\,p_{\kappa/s}\dd U+(1-\epsilon)\rho_s.
 \label{eq:s6v77-mixture}
\end{equation}
\end{theorem}

\begin{proof}
Choose
\[
 0<\kappa\le{\rho\over2B},
\]
where \(\rho\) is from \eqref{eq:s6v77-distance-V} and \(B\) from
\eqref{eq:s6v77-heat-Gaussian}.  The Gaussian heat bound gives
\[
 p_{\kappa/s}(U)
 \le A\kappa^{-4}s^4
 \exp\!\left[-{s\,d(I,U)^2\over B\kappa}\right]
 \le A\kappa^{-4}s^4e^{-2sV(U)}.
\]
On the other hand, the upper bound for \(Z_s\) gives the pointwise lower
bound
\[
 {e^{s\Phi(U)}\over Z_s}
 \ge K_2^{-1}s^4e^{-sV(U)}.
\]
Since \(e^{-2sV}\le e^{-sV}\), the heat density is bounded above by a
fixed multiple of the Wilson probability density.  Take \(\epsilon\)
to be the smaller of the reciprocal multiple and \(1/2\).
Subtracting the heat component and dividing by \(1-\epsilon\) gives a
nonnegative density of total mass one.  Centrality is inherited from
the two central densities.
\end{proof}

\subsection{A cofinal ultraviolet loss floor}
\label{sec:s6v77-UV-loss}

\begin{theorem}[Heat minorization implies Fourier loss]
\label{thm:s6v77-Fourier-loss}
For every irreducible \(\pi\) and \(s\ge s_0\),
\begin{equation}
 0\le R_\pi(s)
 \le 1-\epsilon\left(1-e^{-\kappa C_\pi/s}\right).
 \label{eq:s6v77-R-upper}
\end{equation}
In particular, for every fixed \(\eta_0>0\),
\begin{equation}
 C_\pi\ge\eta_0s
 \quad\Longrightarrow\quad
 1-q_\pi(s)\ge
 \ell_{\rm UV}(\eta_0)>0,                                \label{eq:s6v77-UV-loss}
\end{equation}
where
\[
 \ell_{\rm UV}(\eta_0)
 :=
 1-\left[
 1-\epsilon(1-e^{-\kappa\eta_0})
 \right]^2.
\]
\end{theorem}

\begin{proof}
The heat component in \eqref{eq:s6v77-mixture} has normalized
\(\pi\)-Fourier coefficient \(e^{-\kappa C_\pi/s}\).  Every Fourier
coefficient of the probability \(\rho_s\) has operator norm at most
one.  Since \(\rho_s\) is central, its normalized scalar coefficient is
at most one.  Equation \eqref{eq:s6v77-mixture} therefore gives the
upper bound in \eqref{eq:s6v77-R-upper}; the lower bound is the positive
Wilson character coefficient proved in
\Cref{thm:s6v76-source-scope}.  Square the upper bound and subtract from
one to obtain \eqref{eq:s6v77-UV-loss}.
\end{proof}

\begin{corollary}[The full high-band loss row]
\label{cor:s6v77-full-high-loss}
Let \(\eta_s\downarrow0\) and \(\eta_ss\to\infty\).  After reducing the
fixed \(\eta_0\) used in V76, there is \(c_{\rm hi}>0\) such that
\begin{equation}
 C_\pi>\eta_ss
 \quad\Longrightarrow\quad
 1-q_\pi(s)\ge c_{\rm hi}\eta_s                         \label{eq:s6v77-full-high-loss}
\end{equation}
for all sufficiently large \(s\).
\end{corollary}

\begin{proof}
If \(C_\pi\le\eta_0s\), apply
\Cref{prop:s6v76-moderate-loss}.  If \(C_\pi>\eta_0s\),
apply \eqref{eq:s6v77-UV-loss}.  Since \(\eta_s\to0\), the fixed
ultraviolet number \(\ell_{\rm UV}(\eta_0)\) eventually dominates a
fixed multiple of \(\eta_s\).
\end{proof}

\begin{assessment}[High loss is now sharp at the correct order]
\label{ass:s6v77-loss-sharp}
The heat-model obstruction
\Cref{prop:s6v76-heat-band-obstruction} shows that the order
\(\eta_s\) in \eqref{eq:s6v77-full-high-loss} cannot generally be
replaced by a positive constant.  Thus the central Wilson high-loss row
is now complete at the scale needed by the shrinking low-band cutoff.
What remains is the response lower bound, not contraction.
\end{assessment}

\subsection{Absolute response control and its limitation}
\label{sec:s6v77-response}

\begin{lemma}[Wilson phase variance]
\label{lem:s6v77-phase-variance}
There are \(K,s_1>0\) such that
\begin{equation}
 \operatorname{Var}_{\nu_s}(\Phi)\le {K\over s^2}
 \label{eq:s6v77-phase-variance}
\end{equation}
for \(s\ge s_1\).
\end{lemma}

\begin{proof}
In the exponential chart,
\[
 \Phi=3-\frac14Q+O(Q^2).
\]
The moment bounds
\eqref{eq:s6v76-Qm} therefore make its local variance \(O(s^{-2})\).
The exterior probability is exponentially small by
\eqref{eq:s6v76-tail}, while \(\Phi\) is bounded on the compact group.
\end{proof}

\begin{proposition}[Variance response inequality]
\label{prop:s6v77-response-variance}
For every irreducible \(\pi\),
\begin{equation}
 |R_\pi'(s)|
 \le {K\over s}\sqrt{1-R_\pi(s)^2}.                       \label{eq:s6v77-response-variance}
\end{equation}
Under the V75 clock and \(c_\theta\le c^*\),
\begin{equation}
 2q_\pi r_\pi
 \ge-K_1s\,R_\pi\sqrt{1-R_\pi^2}
 \ge-{K_1\over2}s.                                       \label{eq:s6v77-crude-response}
\end{equation}
\end{proposition}

\begin{proof}
Centrality gives
\[
 \int\pi(U)\,\nu_s(\dd U)=R_\pi I.
\]
For a unit vector \(v\), the mean-unitary variance identity yields
\[
 \int\|\pi(U)v-R_\pi v\|^2\,\nu_s(\dd U)
 =1-R_\pi^2.
\]
Differentiate the matrix mean:
\[
 R_\pi'I
 =\int(\Phi-\mathbb E_s\Phi)\pi(U)\,\nu_s(\dd U).
\]
The constant vector \(R_\pi v\) drops out against the centered phase.
Cauchy--Schwarz and \eqref{eq:s6v77-phase-variance} prove
\eqref{eq:s6v77-response-variance}.  Under
\(-\dot s/s^2=c_\theta\delta\),
\[
 2q_\pi r_\pi=2c_\theta s^2R_\pi R_\pi'.
\]
Use \eqref{eq:s6v77-response-variance} and
\(\max_{0\le x\le1}x\sqrt{1-x^2}=1/2\).
\end{proof}

\begin{corollary}[Positive response through the moderate band]
\label{cor:s6v77-moderate-response}
After reducing \(\eta_0>0\), for all sufficiently large \(s\),
\begin{equation}
 2q_\pi r_\pi\ge0
 \qquad(0<C_\pi\le\eta_0s).                              \label{eq:s6v77-moderate-response}
\end{equation}
On the smaller band \(C_\pi\le\eta_ss\), the uniform positive lower
bound of \Cref{cor:s6v76-low-response} still holds.
\end{corollary}

\begin{proof}
Equation \eqref{eq:s6v76-Rprime-bound} gives
\[
 R_\pi'
 \ge{C_\pi\over s^2}
 \left[
 1-K\left({1\over s}+{C_\pi\over s}\right)
 \right]-Ke^{-cs}.
\]
Use \(C_\pi\ge C_*>0\), choose \(\eta_0\) small, and then take \(s\)
large.  The bracket is positive.  Since \(R_\pi\ge0\), the clock
identity in the preceding proof gives
\eqref{eq:s6v77-moderate-response}.
\end{proof}

\subsection{Far-ultraviolet Fourier integration}
\label{sec:s6v77-Fourier-tail}

\begin{lemma}[Uniform differentiated Fourier tails]
\label{lem:s6v77-Fourier-tail}
For every integer \(N\ge1\), there are \(K_N,s_N<\infty\) such that
\begin{align}
 |R_\pi(s)|
 &\le K_N\left({s\over C_\pi}\right)^N,                  \label{eq:s6v77-R-tail}\\
 |R_\pi'(s)|
 &\le {K_N\over s}
 \left({s\over C_\pi}\right)^N                           \label{eq:s6v77-Rprime-tail}
\end{align}
for all nontrivial \(\pi\) and \(s\ge s_N\).
\end{lemma}

\begin{proof}
Let \(H_C\) act on the \(U\)-variable.  Self-adjointness and
\(H_C\chi_\pi=C_\pi\chi_\pi\) give
\[
 C_\pi^N R_\pi
 ={1\over d_\pi Z_s}
 \int H_C^N(e^{s\Phi})\Re\chi_\pi\,\dd U.
\]
The scaled saddle calculation used in
\Cref{lem:s6v76-saddle-card} gives
\begin{equation}
 {1\over Z_s}\int
 |H_C^N(e^{s\Phi})|\,\dd U\le K_Ns^N.                   \label{eq:s6v77-HN}
\end{equation}
Indeed, in \(x=s^{-1/2}z\) coordinates each conjugated second-order
operator \(e^{-s\Phi}H_Ce^{s\Phi}\) has weighted degree one in \(s\);
all coefficient polynomials have Gaussian moments.  Outside the chart,
derivatives grow only polynomially in \(s\) and the phase gap absorbs
that polynomial.  Since
\(|\chi_\pi|/d_\pi\le1\), this proves
\eqref{eq:s6v77-R-tail}.

For the derivative, use
\eqref{eq:s6v76-cov-identity} and set
\[
 F_s=(\Phi-\mathbb E_s\Phi)e^{s\Phi}.
\]
The same integration by parts gives
\[
 C_\pi^N R_\pi'
 ={1\over d_\pi Z_s}
 \int H_C^NF_s\,\Re\chi_\pi\,\dd U.
\]
At the saddle,
\(\Phi-\mathbb E_s\Phi=s^{-1}\psi_s(z)\), where
\(\psi_s\) and all derivatives required at fixed \(N\) have
Gaussian-polynomial bounds.  Hence
\[
 {1\over Z_s}\int|H_C^NF_s|\,\dd U
 \le K_Ns^{N-1}.
\]
The exterior is again exponentially small after any fixed polynomial
loss.  This proves \eqref{eq:s6v77-Rprime-tail}.
\end{proof}

\begin{corollary}[Far-ultraviolet response is harmless]
\label{cor:s6v77-far-UV-response}
Fix \(\alpha>0\).  Choose an integer \(N>1/(2\alpha)\).  Uniformly on
\[
 C_\pi\ge s^{1+\alpha},
\]
the clock-scaled response obeys
\begin{equation}
 |2q_\pi r_\pi|
 \le K_{N,\alpha}s^{1-2N\alpha}=o(1).                    \label{eq:s6v77-far-UV-response}
\end{equation}
\end{corollary}

\begin{proof}
The clock identity gives
\[
 |2q_\pi r_\pi|
 =2c_\theta s^2|R_\pi R_\pi'|.
\]
Apply both estimates in \Cref{lem:s6v77-Fourier-tail}:
\[
 |2q_\pi r_\pi|
 \le K_Ns
 \left({s\over C_\pi}\right)^{2N}
 \le K_{N,\alpha}s^{1-2N\alpha}.
\]
\end{proof}

\subsection{A three-band conditional central compiler}
\label{sec:s6v77-three-band}

\begin{theorem}[Central Wilson DSFI under the clock separation
\(s\delta\to0\)]
\label{thm:s6v77-central-conditional-DSFI}
Let \(\eta_s\downarrow0\), \(\eta_ss\to\infty\), and fix
\(\alpha>0\).  Split the nontrivial carrier into
\begin{align*}
 {\cal B}_0&=\{C_\pi\le\eta_ss\},\\
 {\cal B}_1&=\{\eta_ss<C_\pi\le\eta_0s\},\\
 {\cal B}_2&=\{\eta_0s<C_\pi<s^{1+\alpha}\},\\
 {\cal B}_3&=\{C_\pi\ge s^{1+\alpha}\}.
\end{align*}
Then, for all large \(s\), there are constants
\(m,\ell_0,M_0>0\), independent of \(s\), such that
\begin{align}
 2q_\pi r_\pi&\ge m &&(\pi\in{\cal B}_0),                 \label{eq:s6v77-B0}\\
 2q_\pi r_\pi&\ge0 &&(\pi\in{\cal B}_1),                  \label{eq:s6v77-B1}\\
 2q_\pi r_\pi&\ge-M_0s &&(\pi\in{\cal B}_2),              \label{eq:s6v77-B2}\\
 2q_\pi r_\pi&\ge-M_0 &&(\pi\in{\cal B}_3),               \label{eq:s6v77-B3}\\
 1-q_\pi&\ge\ell_0 &&(\pi\in{\cal B}_2\cup{\cal B}_3).
                                                                    \label{eq:s6v77-B23-loss}
\end{align}
Consequently
\begin{equation}
 L'\succeq m\delta P_0
 -{M_0s\delta\over\ell_0}L,                              \label{eq:s6v77-three-band-form}
\end{equation}
where \(P_0\) is the projection onto \({\cal B}_0\).
On the low-loss cone
\(\langle f,Lf\rangle\le K_{\rm cone}\delta\|f\|^2\), the
\({\cal B}_1\)-mass is at most
\(2K_{\rm cone}\delta\eta_s^{-1}\|f\|^2\).  Hence if
\begin{equation}
 {\delta\over\eta_s}\longrightarrow0,
 \qquad s\delta\longrightarrow0,                         \label{eq:s6v77-clock-separation}
\end{equation}
then
\begin{equation}
 \langle f,L'f\rangle
 \ge {m\over2}\delta\|f\|^2                              \label{eq:s6v77-central-DSFI}
\end{equation}
on that cone for all sufficiently large \(s\).
\end{theorem}

\begin{proof}
Rows \eqref{eq:s6v77-B0} and \eqref{eq:s6v77-B1} are
\Cref{cor:s6v76-low-response,cor:s6v77-moderate-response}.
Equation \eqref{eq:s6v77-crude-response} gives
\eqref{eq:s6v77-B2}, and
\Cref{cor:s6v77-far-UV-response} gives
\eqref{eq:s6v77-B3} after increasing a constant.
The ultraviolet loss
\eqref{eq:s6v77-UV-loss} gives
\eqref{eq:s6v77-B23-loss}.

Discard the nonnegative \({\cal B}_1\) contribution and compare the
negative mass on \({\cal B}_2\cup{\cal B}_3\) with \(L\) using
\eqref{eq:s6v77-B23-loss}; this proves
\eqref{eq:s6v77-three-band-form}.  By
\Cref{prop:s6v76-moderate-loss}, the cone bound gives
\[
 \|P_1f\|^2
 \le {2K_{\rm cone}\delta\over\eta_s}\|f\|^2.
\]
Thus
\[
 \|P_0f\|^2
 \ge
 \left[
 1-{2K_{\rm cone}\delta\over\eta_s}
   -{K_{\rm cone}\delta\over\ell_0}
 \right]\|f\|^2.
\]
Insert this into \eqref{eq:s6v77-three-band-form}.  The positive-stock
loss is \(o(\delta)\) by \(\delta/\eta_s\to0\), and the negative debit is
\[
 O(s\delta)\langle f,Lf\rangle
 =O(s\delta^2)\|f\|^2=o(\delta)\|f\|^2.
\]
For large \(s\), both errors consume less than half the positive stock,
giving \eqref{eq:s6v77-central-DSFI}.
\end{proof}

\begin{firewall}[The clock separation is an added hypothesis]
\label{fire:s6v77-clock-hypothesis}
V77 has not proved
\eqref{eq:s6v77-clock-separation} for the physical cofinal regulator.
The parameter \(\delta\) in the spectral interpolation and the Wilson
coupling \(s\) must be related by the response-matched construction,
not chosen independently after the estimate.  The theorem is therefore
a conditional closure of the literal central block, not a closure of
the physical half-slab.
\end{firewall}

\begin{assessment}[Exact frontier after V77]
\label{ass:s6v77-frontier}
The central high-loss problem is solved: a fixed heat component inside
the Wilson probability gives a constant loss for \(C_\pi\gtrsim s\),
and V76 supplies the sharp order-\(\eta_s\) loss below that scale.
The response is positive through \(C_\pi\le\eta_0s\), harmless beyond
\(s^{1+\alpha}\), and only bounded below by \(-O(s)\) in the transition
band.  This already closes the central DSFI if
\(\delta/\eta_s\to0\) and \(s\delta\to0\).  Without that clock
separation, the next upstream question is a one-sided response estimate
on
\[
 \eta_0s<C_\pi<s^{1+\alpha}.
\]
\end{assessment}

\begin{programme}[V78: character-graph response sign]
\label{prog:s6v77-V78}
\begin{enumerate}[leftmargin=3em]
\item Rewrite the Wilson coefficients as the positive semigroup
      generated by tensoring with
      \({\bf3}\oplus\overline{\bf3}\) on the \(\SU(3)\) representation
      graph.
\item Express \(R_\pi'/R_\pi\) as a difference of two adjacent
      representation-graph averages.
\item Prove monotonicity of \(R_\pi(s)\), or isolate the precise
      log-concavity inequality whose failure permits a negative
      transition-band response.
\item If monotonicity fails, retain the quantitative \(O(s)\) debit and
      derive the regulator relation needed to absorb it in the physical
      half-slab.
\end{enumerate}
\end{programme}

\begin{firewall}[Claim boundary after sixth-series V77]
\label{fire:s6v77-boundary}
V77 proves a uniform heat-kernel minorization of the central Wilson
probability; derives a fixed ultraviolet Fourier loss and the sharp
order-\(\eta_s\) complete high-band loss; proves an exact
variance-response inequality; proves positive response through the
moderate band and negligible response in a far-ultraviolet band; and
gives a conditional central DSFI under explicit clock separation.

V77 does not prove transition-band response monotonicity, derive the
clock separation from the physical regulator, transfer the central
calculation to the response-matched half-slab, prove the same-owner
residual debit, close the V53 projective/Mosco hypotheses, construct the
continuum OS theory, or prove the full Yang--Mills mass gap.
\end{firewall}

\begin{theorem}[Sixth-series V77 result ledger]
\label{thm:s6v77-results}
The positive conclusions in \Cref{fire:s6v77-boundary} are exactly those
of
\Cref{lem:s6v77-heat-Gaussian,
lem:s6v77-Wilson-geometry,
thm:s6v77-heat-minorization,
thm:s6v77-Fourier-loss,
cor:s6v77-full-high-loss,
lem:s6v77-phase-variance,
prop:s6v77-response-variance,
cor:s6v77-moderate-response,
lem:s6v77-Fourier-tail,
cor:s6v77-far-UV-response,
thm:s6v77-central-conditional-DSFI}.
\end{theorem}

\begin{proof}
Each item is the cited statement with its centrality, large-\(s\),
Casimir-band, clock, and cone qualifications unchanged.
\end{proof}

\paragraph{Parent-version record.}
V77 was formed from
\texttt{Renewal\_Yang\_Mills\_Mass\_Gap\_Research\_Notes\_V76.tex}.
The V76 parent had \(34\,378\) logical source lines,
\(1\,206\,045\) bytes, and SHA--256
\[
 \texttt{f4e3acd3c34645791e3cdde48024908480f24a3b969f22a82caa49f12b33885c}.
\]
The next compulsory companion audit remains V80.

% =============================================================================
% END APPENDED SIXTH-SERIES V77 MATERIAL
% =============================================================================

% =============================================================================
% BEGIN APPENDED SIXTH-SERIES V78 MATERIAL
% =============================================================================

\section{V78: the exact \(\SU(3)\) fusion-graph response equation}
\label{sec:s6v78}

\subsection{Why the transition-band sign is a representation-ring problem}
\label{sec:s6v78-purpose}

V77 leaves only one central analytic ambiguity if the clock separation
\(s\delta\to0\) is unavailable: the sign of \(R_\pi'(s)\) for
\[
 \eta_0s<C_\pi<s^{1+\alpha}.
\]
The positive character expansion does not by itself settle that sign,
because a ratio of two positive entire functions need not be increasing.
This iteration removes all group integrals from the question and rewrites
it as one explicit parabolic inequality on the \(\SU(3)\) fusion graph.

Let
\[
 W={\bf3}\oplus\overline{\bf3},\qquad
 \chi_W=\chi_{\bf3}+\chi_{\overline{\bf3}}=2\Phi.
\]
For irreducibles \(\pi,\mu\), write
\[
 N_{W\pi}^{\mu}
 :=\dim\operatorname{Hom}_{\SU(3)}(\mu,W\otimes\pi).
\]

\begin{definition}[Fusion adjacency and coefficient vector]
\label{def:s6v78-fusion}
On the algebraic span of irreducible labels, define
\begin{equation}
 M_{\pi\mu}:=N_{W\pi}^{\mu}.
 \label{eq:s6v78-M}
\end{equation}
Expand
\begin{equation}
 e^{s\Phi(U)}
 =\sum_{\pi}b_\pi(s)\chi_\pi(U).
 \label{eq:s6v78-b-expansion}
\end{equation}
\end{definition}

\begin{lemma}[Exact fusion semigroup]
\label{lem:s6v78-fusion-semigroup}
The adjacency \(M\) is symmetric, extends to a bounded self-adjoint
operator of norm \(6\) on the character \(\ell^2\)-space, and
\begin{equation}
 \boxed{
 b'(s)=\frac12Mb(s),\qquad b(0)=e_{\bf1},\qquad
 b(s)=e^{sM/2}e_{\bf1}.}
 \label{eq:s6v78-fusion-semigroup}
\end{equation}
Moreover
\begin{equation}
 R_\pi(s)={b_\pi(s)\over d_\pi b_{\bf1}(s)}.
 \label{eq:s6v78-R-b}
\end{equation}
\end{lemma}

\begin{proof}
Multiplication by the real character \(\chi_W\) is a bounded
self-adjoint operator on the Hilbert space of square-integrable class
functions.  In the orthonormal irreducible-character basis its matrix is
\(M\).  Since \(\|\chi_W\|_\infty=6\), its operator norm is \(6\).
Differentiate \eqref{eq:s6v78-b-expansion}:
\[
 {d\over ds}e^{s\Phi}
 ={1\over2}\chi_We^{s\Phi}.
\]
The tensor-product rule for characters gives the differential equation
and initial condition in \eqref{eq:s6v78-fusion-semigroup}; boundedness
gives its unique exponential solution.  Character orthogonality and
\eqref{eq:s6v74-Fourier-coefficient} give
\(a_\pi=b_\pi/d_\pi\) and \(a_{\bf1}=b_{\bf1}\), proving
\eqref{eq:s6v78-R-b}.
\end{proof}

\subsection{The reversible Markov form}
\label{sec:s6v78-Markov}

The dimension vector satisfies
\[
 \sum_\mu N_{W\pi}^{\mu}d_\mu=6d_\pi.
\]
It therefore gives a Doob transform of the fusion adjacency.

\begin{definition}[Dimension-biased fusion walk]
\label{def:s6v78-P}
Set
\begin{equation}
 P_{\pi\mu}
 :={N_{W\pi}^{\mu}d_\mu\over6d_\pi}.
 \label{eq:s6v78-P}
\end{equation}
\end{definition}

\begin{lemma}[Reversibility and heat-ratio representation]
\label{lem:s6v78-Markov-ratio}
The matrix \(P\) is Markov and reversible with respect to the
\(\sigma\)-finite measure
\[
 \mathfrak m(\pi)=d_\pi^2.
\]
If \(K_s=e^{3s(P-I)}\), then
\begin{equation}
 \boxed{
 R_\pi(s)={K_s(\pi,{\bf1})\over K_s({\bf1},{\bf1})}.}
 \label{eq:s6v78-heat-ratio}
\end{equation}
\end{lemma}

\begin{proof}
The dimension identity makes every row sum of \(P\) equal to one.
Self-duality of \(W\) gives
\[
 N_{W\pi}^{\mu}=N_{W\mu}^{\pi}.
\]
Hence
\[
 d_\pi^2P_{\pi\mu}
 ={d_\pi d_\mu\over6}N_{W\pi}^{\mu}
 =d_\mu^2P_{\mu\pi}.
\]
Put \(u_\pi=b_\pi/d_\pi\).  The fusion equation gives
\[
 u_\pi'=3\sum_\mu P_{\pi\mu}u_\mu,\qquad
 u(0)=e_{\bf1}.
\]
Thus
\[
 u(s)=e^{3s}K_se_{\bf1}.
\]
Divide its \(\pi\)-entry by its trivial entry and use
\eqref{eq:s6v78-R-b}.
\end{proof}

\begin{remark}[The ratio is not the heat kernel itself]
\label{rem:s6v78-ratio-not-kernel}
The denominator \(K_s({\bf1},{\bf1})\) is essential.  Positivity of the
Markov heat kernel proves \(R_\pi\ge0\), but it does not imply that the
ratio in \eqref{eq:s6v78-heat-ratio} is increasing.  A parabolic
comparison or a graph-specific path inequality is still needed.
\end{remark}

\subsection{Coordinates on the \(\SU(3)\) Weyl chamber}
\label{sec:s6v78-coordinates}

Label an irreducible by its highest weight \((p,q)\in\mathbb N_0^2\).
Its dimension and Casimir are
\begin{align}
 d_{p,q}
 &={(p+1)(q+1)(p+q+2)\over2},                            \label{eq:s6v78-dimension}\\
 C_{p,q}
 &={p^2+q^2+pq+3p+3q\over3}.                             \label{eq:s6v78-Casimir}
\end{align}

\begin{lemma}[Six-neighbour fusion rule]
\label{lem:s6v78-six-neighbour}
The neighbours of \((p,q)\) under tensoring with \(W\) are
\begin{equation}
 \begin{split}
 &(p+1,q),\quad(p-1,q+1),\quad(p,q-1),\\
 &(p,q+1),\quad(p+1,q-1),\quad(p-1,q),
 \end{split}
 \label{eq:s6v78-neighbours}
\end{equation}
with every term having a negative coordinate omitted.  Each retained
edge has multiplicity one.
\end{lemma}

\begin{proof}
The first row is the standard highest-weight decomposition of
\((p,q)\otimes(1,0)\), and the second is its conjugate decomposition
with \((0,1)\).  Boundary terms are absent because a negative Dynkin
coordinate is not dominant.
\end{proof}

\subsection{The exact graph Harnack card}
\label{sec:s6v78-Harnack}

\begin{theorem}[Fusion-graph response equation]
\label{thm:s6v78-response-equation}
Let \(R(s)=(R_\pi(s))_\pi\), with \(R_{\bf1}=1\).  Then
\begin{equation}
 \boxed{
 R_\pi'(s)
 =3\left[(PR(s))_\pi-R_{\bf3}(s)R_\pi(s)\right].}
 \label{eq:s6v78-response-equation}
\end{equation}
Consequently,
\begin{equation}
 R_\pi'(s)\ge0
 \quad\Longleftrightarrow\quad
 (PR(s))_\pi\ge R_{\bf3}(s)R_\pi(s).                     \label{eq:s6v78-Harnack-equivalence}
\end{equation}
\end{theorem}

\begin{proof}
Differentiate \(R_\pi=u_\pi/u_{\bf1}\) and use \(u'=3Pu\):
\[
 R_\pi'
 =3\{(PR)_\pi-R_\pi(PR)_{\bf1}\}.
\]
From the root, the fusion walk moves to \({\bf3}\) and
\(\overline{\bf3}\), each with probability \(1/2\).  Charge conjugation
of the Wilson density gives
\[
 R_{\overline{\bf3}}=R_{\bf3},
\]
so \((PR)_{\bf1}=R_{\bf3}\).  This proves
\eqref{eq:s6v78-response-equation}, and
\eqref{eq:s6v78-Harnack-equivalence} follows by division by \(3\).
\end{proof}

\begin{definition}[Fusion heat-ratio Harnack inequality, FHRH]
\label{def:s6v78-FHRH}
FHRH is the assertion that, for every \(s>0\) and every irreducible
\(\pi\),
\begin{equation}
 (PR(s))_\pi\ge R_{\bf3}(s)R_\pi(s).                     \label{eq:s6v78-FHRH}
\end{equation}
\end{definition}

\begin{corollary}[What FHRH would close]
\label{cor:s6v78-FHRH-closes}
If FHRH holds, then every central Wilson ratio is nondecreasing,
\[
 R_\pi'(s)\ge0,
\]
and every central response coefficient under the V75 clock is
nonnegative:
\[
 2q_\pi r_\pi=2c_\theta s^2R_\pi R_\pi'\ge0.
\]
Together with the low-band strict lower bound of V76, the V75 band
compiler yields central DSFI without the additional condition
\(s\delta\to0\).
\end{corollary}

\begin{proof}
Apply \Cref{thm:s6v78-response-equation}; the clock identity is the one
used in \Cref{prop:s6v77-response-variance}.  Since there is no negative
high-band response, take \(M=0\) in
\Cref{thm:s6v75-band-compiler}.  The high loss is
\eqref{eq:s6v77-full-high-loss}; on a low-loss cone its order
\(\eta_s\) is absorbable whenever \(\delta/\eta_s\to0\).
\end{proof}

\begin{proposition}[The fundamental FHRH row is strict]
\label{prop:s6v78-fundamental}
For \(s>0\),
\begin{equation}
 R_{\bf3}(s)={1\over3}\mathbb E_s\Phi,\qquad
 R_{\bf3}'(s)={1\over3}\operatorname{Var}_{\nu_s}(\Phi)>0.
 \label{eq:s6v78-fundamental}
\end{equation}
Thus \eqref{eq:s6v78-FHRH} holds strictly for
\(\pi={\bf3},\overline{\bf3}\).
\end{proposition}

\begin{proof}
Charge conjugation makes the expectation of
\(\Im\chi_{\bf3}\) vanish.  Since
\(\Phi=\Re\chi_{\bf3}\), definition \eqref{eq:s6v76-R} gives the first
identity.  Differentiate the exponential-family expectation to obtain
the covariance with \(\Phi\), which is its variance divided by three.
The Wilson law has full support and \(\Phi\) is not constant, so the
variance is strictly positive.  The equivalence
\eqref{eq:s6v78-Harnack-equivalence} proves the last assertion.
\end{proof}

\subsection{An exact path-count formulation}
\label{sec:s6v78-paths}

Let
\[
 m_\pi(n):=(M^n)_{\pi{\bf1}}.
\]
By \Cref{lem:s6v78-six-neighbour}, this is the number of length-\(n\)
fusion paths from the root to \(\pi\), with boundary omissions.

\begin{proposition}[Response determinant and its coefficients]
\label{prop:s6v78-response-determinant}
Define
\begin{equation}
 D_\pi(s)
 :=b_\pi'(s)b_{\bf1}(s)-b_\pi(s)b_{\bf1}'(s).
 \label{eq:s6v78-D}
\end{equation}
Then
\begin{equation}
 R_\pi'(s)={D_\pi(s)\over d_\pi b_{\bf1}(s)^2},           \label{eq:s6v78-D-response}
\end{equation}
and the coefficient of \(s^k\) in \(D_\pi\) is
\begin{equation}
 {1\over2^{k+1}k!}
 \sum_{i=0}^{k}\binom{k}{i}
 \left[
 m_\pi(i+1)m_{\bf1}(k-i)
 -m_\pi(i)m_{\bf1}(k-i+1)
 \right].
 \label{eq:s6v78-D-coefficient}
\end{equation}
Thus coefficientwise nonnegativity of
\eqref{eq:s6v78-D-coefficient} is a sufficient, entirely integral,
certificate for FHRH at \(\pi\).
\end{proposition}

\begin{proof}
Differentiate \eqref{eq:s6v78-R-b} to obtain
\eqref{eq:s6v78-D-response}.  The semigroup series is
\[
 b_\pi(s)=\sum_{n=0}^{\infty}
 m_\pi(n){(s/2)^n\over n!}.
\]
Multiply the derivative series for \(b_\pi'\) and the series for
\(b_{\bf1}\), subtract the reversed product, and group terms of total
degree \(k\).  Replacing \(1/(i!(k-i)!)\) by
\(\binom{k}{i}/k!\) gives \eqref{eq:s6v78-D-coefficient}.
\end{proof}

\begin{lemma}[A tempting coefficient-ratio shortcut fails]
\label{lem:s6v78-ratio-shortcut-fails}
The sequence
\[
 {m_{\bf3}(n)\over3m_{\bf1}(n)}
\]
is not nondecreasing on the indices where the denominator is nonzero.
Specifically,
\begin{equation}
 {m_{\bf3}(2)\over3m_{\bf1}(2)}={1\over6},\qquad
 {m_{\bf3}(3)\over3m_{\bf1}(3)}=1,\qquad
 {m_{\bf3}(4)\over3m_{\bf1}(4)}={5\over12}.
 \label{eq:s6v78-ratio-counter}
\end{equation}
\end{lemma}

\begin{proof}
Successive application of the six-neighbour rule gives
\[
 (m_{\bf1}(2),m_{\bf3}(2))=(2,1),\quad
 (m_{\bf1}(3),m_{\bf3}(3))=(2,6),\quad
 (m_{\bf1}(4),m_{\bf3}(4))=(12,15).
\]
Substitution proves \eqref{eq:s6v78-ratio-counter}.
\end{proof}

\begin{assessment}[Meaning of the failed shortcut]
\label{ass:s6v78-shortcut}
A standard ratio-of-power-series lemma cannot prove FHRH by comparing
the raw coefficients \(m_\pi(n)\) with the return counts
\(m_{\bf1}(n)\): the required monotonicity already fails for the
fundamental representation, even though its full analytic ratio is
strictly increasing by \Cref{prop:s6v78-fundamental}.  Any successful
proof must use the binomially convolved determinant
\eqref{eq:s6v78-D-coefficient}, a parabolic comparison, or an equivalent
path-switching argument.  Positivity of the character expansion alone
is insufficient.
\end{assessment}

\subsection{A rigorous finite-window certificate}
\label{sec:s6v78-finite-certificate}

\begin{lemma}[Uniform exponential-series remainder]
\label{lem:s6v78-series-remainder}
For \(0\le s\le S\), let
\[
 b^{[N]}(s):=\sum_{n=0}^{N}{(sM/2)^n\over n!}e_{\bf1}.
\]
Then
\begin{align}
 \|b(s)-b^{[N]}(s)\|_{\ell^2}
 &\le e^{3S}{(3S)^{N+1}\over(N+1)!},                    \label{eq:s6v78-b-tail}\\
 \|b'(s)-(b^{[N]})'(s)\|_{\ell^2}
 &\le3e^{3S}{(3S)^N\over N!}.                            \label{eq:s6v78-bprime-tail}
\end{align}
\end{lemma}

\begin{proof}
Use \(\|M/2\|\le3\) in the exponential series.  The scalar Taylor
remainder satisfies
\[
 \sum_{n=N+1}^{\infty}{x^n\over n!}
 \le e^x{x^{N+1}\over(N+1)!},
\]
which proves \eqref{eq:s6v78-b-tail}.  Differentiate the series, or
apply the same estimate after the bounded factor \(M/2\), to obtain
\eqref{eq:s6v78-bprime-tail}.
\end{proof}

\begin{proposition}[Finite-card FHRH certification protocol]
\label{prop:s6v78-finite-certificate}
Fix a finite representation set \({\cal F}\), a compact coupling
interval \([0,S]\), and \(N\).  Form the rational polynomials
\[
 D_\pi^{[N]}
 :=(b_\pi^{[N]})'b_{\bf1}^{[N]}
   -b_\pi^{[N]}(b_{\bf1}^{[N]})'.
\]
The bounds
\eqref{eq:s6v78-b-tail}--\eqref{eq:s6v78-bprime-tail} give an explicit
uniform error \(E_{N,S}\) with
\[
 |D_\pi-D_\pi^{[N]}|\le E_{N,S}
 \qquad(\pi\in{\cal F},\ 0\le s\le S).
\]
If exact rational interval arithmetic, a Bernstein-basis test, or a
Sturm test proves
\[
 D_\pi^{[N]}(s)>E_{N,S}
\]
on the interval for every \(\pi\in{\cal F}\), then FHRH is rigorously
certified on that finite card.
\end{proposition}

\begin{proof}
Each component of \(b,b'\) is bounded respectively by \(e^{3S}\) and
\(3e^{3S}\).  Expanding
\(D_\pi-D_\pi^{[N]}\) into four product errors and applying
\Cref{lem:s6v78-series-remainder} gives, for example, the explicit
choice
\[
 E_{N,S}
 =6e^{6S}{(3S)^{N+1}\over(N+1)!}
 +2e^{6S}{3(3S)^N\over N!}
 +2e^{6S}
 {3(3S)^{2N+1}\over N!(N+1)!},
\]
where harmless overestimation keeps the formula uniform in \(\pi\).
The stated exact polynomial test then implies \(D_\pi>0\), and
\eqref{eq:s6v78-D-response} gives FHRH.
\end{proof}

\begin{firewall}[Finite-window evidence is not the cofinal theorem]
\label{fire:s6v78-finite-only}
The transition band contains representations whose highest weights and
coupling interval grow with \(s\).  A list of numerical checks, or even
an exact certificate on each preassigned finite card, does not prove
FHRH there.  The purpose of
\Cref{prop:s6v78-finite-certificate} is falsification and local
certification; the needed cofinal theorem remains
\eqref{eq:s6v78-FHRH} or a weaker transition-band lower bound with
controlled clock debit.
\end{firewall}

\begin{assessment}[Exact frontier after V78]
\label{ass:s6v78-frontier}
The transition response is now an explicit graph inequality:
\[
 (PR)_\pi\ge R_{\bf3}R_\pi
\]
for the dimension-biased six-neighbour walk on the dominant
\(\SU(3)\) Weyl chamber.  The fundamental row is strictly proved, and
every finite card has a rigorous, computable certificate.  Raw
coefficient-ratio monotonicity is disproved as a proof route.  The next
upstream step is therefore a path-switching or maximum-principle proof
for the convolved determinant \(D_\pi\), not another saddle expansion.
\end{assessment}

\begin{programme}[V79: parabolic comparison or obstruction]
\label{prog:s6v78-V79}
\begin{enumerate}[leftmargin=3em]
\item Derive the evolution equation for the Harnack defect
      \(F_\pi=(PR)_\pi-R_{\bf3}R_\pi\).
\item Test whether the first hypothetical zero of \(F\) is excluded by
      the six-neighbour geometry and the boundary transition weights.
\item In parallel, seek a sign-reversing path injection proving
      nonnegativity of the coefficients
      \eqref{eq:s6v78-D-coefficient}.
\item If neither closes, extract the weakest lower bound on \(F_\pi\)
      sufficient to improve the \(O(s)\) transition debit in V77.
\end{enumerate}
\end{programme}

\begin{firewall}[Claim boundary after sixth-series V78]
\label{fire:s6v78-boundary}
V78 proves the exact fusion-semigroup and reversible Markov
representations of every central Wilson coefficient; derives the
six-neighbour fusion-walk response equation and the exact FHRH
equivalence; proves the fundamental response is strictly positive;
derives the integral path-count determinant; rules out a naive
coefficient-ratio proof; and gives a rigorous finite-window
certification protocol.

V78 does not prove FHRH on the cofinal transition band, remove the
clock-separation hypothesis of V77, transfer the central calculation to
the response-matched half-slab, prove the same-owner residual debit,
close the V53 projective/Mosco hypotheses, construct the continuum OS
theory, or prove the full Yang--Mills mass gap.
\end{firewall}

\begin{theorem}[Sixth-series V78 result ledger]
\label{thm:s6v78-results}
The positive conclusions in \Cref{fire:s6v78-boundary} are exactly those
of
\Cref{lem:s6v78-fusion-semigroup,
lem:s6v78-Markov-ratio,
lem:s6v78-six-neighbour,
thm:s6v78-response-equation,
cor:s6v78-FHRH-closes,
prop:s6v78-fundamental,
prop:s6v78-response-determinant,
lem:s6v78-ratio-shortcut-fails,
lem:s6v78-series-remainder,
prop:s6v78-finite-certificate}.
\end{theorem}

\begin{proof}
Each item is the cited statement with its central-Wilson,
fusion-graph, finite-card, and cofinal qualifications unchanged.
\end{proof}

\paragraph{Parent-version record.}
V78 was formed from
\texttt{Renewal\_Yang\_Mills\_Mass\_Gap\_Research\_Notes\_V77.tex}.
The V77 parent had \(34\,935\) logical source lines,
\(1\,224\,729\) bytes, and SHA--256
\[
 \texttt{671a2296e2d2b74afab01204f47f7a90997a80e897a1ace295895a0ed9adc38c}.
\]
The next compulsory companion audit remains V80.

% =============================================================================
% END APPENDED SIXTH-SERIES V78 MATERIAL
% =============================================================================

% =============================================================================
% BEGIN APPENDED SIXTH-SERIES V79 MATERIAL
% =============================================================================

\section{V79: direct central gap asymptotics and the finite-mass clock}
\label{sec:s6v79}

\subsection{Upstream pivot from response sign to the endpoint spectrum}
\label{sec:s6v79-pivot}

V78 isolates the still-open fusion-graph inequality FHRH.  That
inequality would make one interpolation especially clean, but it is not
needed to determine the spectrum of the endpoint central Wilson
transfer.  V76 and V77 already control every Casimir band strongly
enough to compute the first relative edge directly.  This is the
appropriate upstream move because the continuum mass is determined by
the endpoint quantity
\[
 -\delta_s^{-1}\log q_{*,s},
\]
not by the sign of an arbitrarily chosen path at every intermediate
parameter.

Let
\begin{equation}
 q_{*,s}:=\sup_{\pi\ne{\bf1}}q_\pi(s),\qquad
 g_s:=1-q_{*,s}.
 \label{eq:s6v79-global-edge}
\end{equation}
The smallest nonzero \(\SU(3)\) Casimir is
\[
 C_*={4\over3},
\]
attained exactly by \({\bf3}\) and \(\overline{\bf3}\); the next
Casimir is \(3\), attained by the adjoint \({\bf8}\).

\begin{lemma}[Uniform low-Casimir loss asymptotic]
\label{lem:s6v79-low-loss-asymptotic}
Let \(\eta_s\downarrow0\) and \(\eta_ss\to\infty\).  Uniformly for
\[
 0<C_\pi\le\eta_ss,
\]
one has
\begin{equation}
 1-q_\pi(s)
 ={2C_\pi\over s}
 \left[1+O\!\left(\eta_s+s^{-1}\right)\right].
 \label{eq:s6v79-low-loss-asymptotic}
\end{equation}
\end{lemma}

\begin{proof}
By \eqref{eq:s6v76-R-bound},
\[
 R_\pi=1-{C_\pi\over s}+E_{\pi,s},\qquad
 |E_{\pi,s}|
 \le K{C_\pi\over s}
 \left({1\over s}+{C_\pi\over s}\right)+Ke^{-cs}.
\]
Since \(C_\pi\ge C_*\), the exponential term is absorbed into the same
right side.  On the stated band this is
\[
 E_{\pi,s}
 ={C_\pi\over s}O(\eta_s+s^{-1}).
\]
Now
\[
 1-q_\pi
 =1-R_\pi^2
 =(1-R_\pi)(1+R_\pi).
\]
The first factor is
\((C_\pi/s)[1+O(\eta_s+s^{-1})]\), and the second is
\(2+O(\eta_s)\), uniformly on the band.  Multiplication proves
\eqref{eq:s6v79-low-loss-asymptotic}.
\end{proof}

\begin{theorem}[The central Wilson spectral edge is fundamental]
\label{thm:s6v79-central-edge}
For all sufficiently large \(s\),
\begin{equation}
 \boxed{
 q_{*,s}=q_{\bf3}(s)=q_{\overline{\bf3}}(s)}
 \label{eq:s6v79-edge-sector}
\end{equation}
and
\begin{equation}
 \boxed{
 g_s={2C_*\over s}+O(s^{-2})
 ={8\over3s}+O(s^{-2}).}
 \label{eq:s6v79-gap-asymptotic}
\end{equation}
\end{theorem}

\begin{proof}
Choose any \(\eta_s\downarrow0\) with
\(\eta_ss\to\infty\), for example \(\eta_s=s^{-1/3}\).
For the fundamental representation, the fixed-Casimir estimate
\eqref{eq:s6v76-R-bound} gives
\begin{equation}
 1-q_{\bf3}(s)
 ={2C_*\over s}+O(s^{-2}).                               \label{eq:s6v79-fund-gap}
\end{equation}
Charge conjugation gives the same value for
\(\overline{\bf3}\).

If \(0<C_\pi\le\eta_ss\) and
\(\pi\notin\{{\bf3},\overline{\bf3}\}\), then
\(C_\pi\ge3\).  By
\Cref{lem:s6v79-low-loss-asymptotic},
\[
 1-q_\pi(s)\ge{6\over s}[1-o(1)],
\]
which is strictly larger than
\eqref{eq:s6v79-fund-gap} for large \(s\).

If \(\eta_ss<C_\pi\le\eta_0s\), then
\Cref{prop:s6v76-moderate-loss} gives
\[
 1-q_\pi(s)\ge{\eta_s\over2}.
\]
Because \(s\eta_s\to\infty\), this is much larger than \(s^{-1}\).
Finally, if \(C_\pi>\eta_0s\), the fixed ultraviolet loss
\eqref{eq:s6v77-UV-loss} is again much larger than \(s^{-1}\).
Thus only the two fundamental sectors attain the first nontrivial edge,
and \eqref{eq:s6v79-fund-gap} proves
\eqref{eq:s6v79-gap-asymptotic}.
\end{proof}

\begin{remark}[Multiplicity of the edge]
\label{rem:s6v79-edge-multiplicity}
The central convolution acts by one scalar on every matrix coefficient.
Thus \eqref{eq:s6v79-edge-sector} identifies the representation sectors,
not a simple eigenvector.  The fundamental and antifundamental
Peter--Weyl blocks carry the same first relative eigenvalue.  This
degeneracy is harmless for the spectral edge but must be respected by
any Feynman--Hellmann argument along a noncentral perturbation.
\end{remark}

\subsection{The physical duration forced by a finite generator}
\label{sec:s6v79-clock}

Let the central full-transfer generator at physical duration
\(\delta_s>0\) be
\[
 H_s=-\delta_s^{-1}\log P_s
\]
on the vacuum complement.  Its gap is
\begin{equation}
 \mu_s:=-{1\over\delta_s}\log q_{*,s}.
 \label{eq:s6v79-mus}
\end{equation}

\begin{theorem}[Central clock trichotomy]
\label{thm:s6v79-clock-trichotomy}
As \(s\to\infty\),
\begin{equation}
 \boxed{
 \mu_s={2C_*\over s\delta_s}
 \left[1+O(s^{-1})\right].}
 \label{eq:s6v79-mass-clock}
\end{equation}
Consequently:
\begin{enumerate}[label=\textup{(\roman*)},leftmargin=3.5em]
\item if \(s\delta_s\to\lambda\in(0,\infty)\), then
      \[
      \mu_s\longrightarrow{2C_*\over\lambda}
      ={8\over3\lambda};
      \]
\item if \(s\delta_s\to0\), then \(\mu_s\to+\infty\);
\item if \(s\delta_s\to+\infty\), then \(\mu_s\to0\).
\end{enumerate}
\end{theorem}

\begin{proof}
By \Cref{thm:s6v79-central-edge},
\[
 q_{*,s}=1-{2C_*\over s}+O(s^{-2}).
\]
The Taylor formula
\[
 -\log(1-x)=x+O(x^2)
\]
with \(x=g_s=O(s^{-1})\) gives
\[
 -\log q_{*,s}={2C_*\over s}+O(s^{-2}).
\]
Divide by \(\delta_s\).  The three alternatives follow immediately.
\end{proof}

\begin{assessment}[Correction to the provisional V77 branch]
\label{ass:s6v79-clock-correction}
The condition \(s\delta_s\to0\) in
\Cref{thm:s6v77-central-conditional-DSFI} is a mathematically valid
absorption regime for that pathwise estimate, but
\Cref{thm:s6v79-clock-trichotomy} shows that it is not the
finite-central-mass regime.  It drives the central generator edge to
infinity.  The physically relevant central scaling is instead
\[
 s\delta_s\to\lambda\in(0,\infty).
\]
At that scale the \(O(s)\) transition-band debit from V77 cannot be
discarded by the factor \(\delta_s\); either FHRH must control its sign,
or the endpoint spectrum must be transferred directly without using
that interpolation.
\end{assessment}

\subsection{Endpoint transfer stability at generator scale}
\label{sec:s6v79-stability}

\begin{lemma}[Relative-edge perturbation]
\label{lem:s6v79-edge-perturbation}
Let \(P_s\) and \(\widetilde P_s\) be self-adjoint positive contractions
on one Hilbert space with the same normalized vacuum vector and vacuum
eigenvalue one.  Let \(q_{*,s}\) and \(\widetilde q_{*,s}\) be their
vacuum-complement spectral suprema.  Then
\begin{equation}
 |\widetilde q_{*,s}-q_{*,s}|
 \le\|\widetilde P_s-P_s\|.                              \label{eq:s6v79-edge-perturbation}
\end{equation}
If
\begin{equation}
 \|\widetilde P_s-P_s\|=o(s^{-1}),                       \label{eq:s6v79-os-perturbation}
\end{equation}
then
\begin{equation}
 1-\widetilde q_{*,s}
 ={2C_*\over s}+o(s^{-1}).                               \label{eq:s6v79-transferred-edge}
\end{equation}
\end{lemma}

\begin{proof}
On the common vacuum complement, the variational characterization of
the spectral supremum gives
\[
 \widetilde q_{*,s}
 \le q_{*,s}+\|\widetilde P_s-P_s\|.
\]
Interchanging the operators proves
\eqref{eq:s6v79-edge-perturbation}.  Combine it with
\eqref{eq:s6v79-gap-asymptotic} and
\eqref{eq:s6v79-os-perturbation}.
\end{proof}

\begin{theorem}[Conditional endpoint handoff to a physical half-slab]
\label{thm:s6v79-endpoint-handoff}
Suppose a response-matched, gauge-reduced physical full transfer
\(\widetilde P_j\) can be identified on a common boundary Hilbert space
with a central Wilson reference \(P_{s_j}\) so that:
\begin{enumerate}[label=\textup{(\roman*)},leftmargin=3.5em]
\item both are positive self-adjoint contractions with the same
      normalized vacuum;
\item \(s_j\to\infty\) and
      \(s_j\delta_j\to\lambda\in(0,\infty)\);
\item
      \[
      \|\widetilde P_j-P_{s_j}\|=o(\delta_j);
      \]
\item the projective/common-refinement and fixed-time convergence rows
      of HSOCC hold for \(\widetilde P_j\).
\end{enumerate}
Then the continuum Hamiltonian reconstructed from
\(\widetilde P_j\) has vacuum-complement gap at least
\[
 {2C_*\over\lambda}={8\over3\lambda}.
\]
\end{theorem}

\begin{proof}
The clock assumption makes
\[
 o(\delta_j)=o(s_j^{-1}).
\]
Apply \Cref{lem:s6v79-edge-perturbation}:
\[
 1-\widetilde q_{*,j}
 ={2C_*\over s_j}+o(s_j^{-1})
 ={2C_*\over\lambda}\delta_j+o(\delta_j).
\]
The logarithmic conversion of V50--V51 gives the same limiting
generator gap.  The HSOCC projective and fixed-time rows transfer that
edge to the regulator-independent OS semigroup and then to its
Hamiltonian.
\end{proof}

\begin{corollary}[A quantitative \(O(\delta)\) variant]
\label{cor:s6v79-Odelta-handoff}
In the setting of
\Cref{thm:s6v79-endpoint-handoff}, replace row~\textup{(iii)} by
\[
 \|\widetilde P_j-P_{s_j}\|
 \le\varepsilon\delta_j+o(\delta_j)
\]
for a constant
\[
 0\le\varepsilon<{2C_*\over\lambda}.
\]
Then the limiting physical gap is at least
\begin{equation}
 {2C_*\over\lambda}-\varepsilon>0.
 \label{eq:s6v79-Odelta-gap}
\end{equation}
\end{corollary}

\begin{proof}
The edge perturbation lemma gives
\[
 1-\widetilde q_{*,j}
 \ge
 \left({2C_*\over\lambda}-\varepsilon\right)\delta_j
o(\delta_j).
\]
Apply the same logarithmic and HSOCC handoff.
\end{proof}

\begin{firewall}[The comparison operator must have the same owner]
\label{fire:s6v79-owner}
The norm estimate in
\Cref{thm:s6v79-endpoint-handoff} is not supplied by matching one
plaquette Hessian or one fundamental character coefficient.  It must
compare normalized full transfers after shared-link integration, gauge
projection, all spatial plaquettes, generated/marginal returns,
Jacobians, counterterms, and moving ground factors have been included.
It must also identify the vacuum complements.  Without those rows,
\(\widetilde P_j-P_{s_j}\) is not even a defined same-space operator.
\end{firewall}

\subsection{What an amplitude estimate must deliver}
\label{sec:s6v79-amplitude}

Let \(A_s\) and \(\widetilde A_s\) be normalized half-transfer
amplitudes with
\[
 P_s=A_s^*A_s,\qquad
 \widetilde P_s=\widetilde A_s^*\widetilde A_s,
\qquad
 \|A_s\|,\|\widetilde A_s\|\le1.
\]

\begin{lemma}[Half-amplitude to full-transfer defect]
\label{lem:s6v79-amplitude-defect}
\begin{equation}
 \|\widetilde P_s-P_s\|
 \le2\|\widetilde A_s-A_s\|.                             \label{eq:s6v79-amplitude-defect}
\end{equation}
Consequently an \(o(s^{-1})\) normalized half-amplitude defect suffices
for the endpoint handoff.
\end{lemma}

\begin{proof}
Insert and subtract \(\widetilde A_s^*A_s\):
\[
 \widetilde A_s^*\widetilde A_s-A_s^*A_s
 =\widetilde A_s^*(\widetilde A_s-A_s)
  +(\widetilde A_s^*-A_s^*)A_s.
\]
Take norms and use the contraction bounds.
\end{proof}

\begin{proposition}[Ground-factor normalization ledger]
\label{prop:s6v79-ground-ledger}
Suppose unnormalized amplitudes \(B_s,\widetilde B_s\) have leading
singular values \(\sigma_s,\widetilde\sigma_s>0\), and set
\[
 A_s={B_s\over\sigma_s},\qquad
 \widetilde A_s={\widetilde B_s\over\widetilde\sigma_s}.
\]
Then
\begin{equation}
 \|\widetilde A_s-A_s\|
 \le{\|\widetilde B_s-B_s\|\over\widetilde\sigma_s}
 +{\|B_s\|\over\sigma_s\widetilde\sigma_s}
 |\widetilde\sigma_s-\sigma_s|.                          \label{eq:s6v79-ground-ledger}
\end{equation}
In particular, an unnormalized \(o(s^{-1})\) estimate is useful only
after lower bounds on both leading singular values and the normalized
ground-factor difference have been kept at the same scale.
\end{proposition}

\begin{proof}
Write
\[
 {\widetilde B_s\over\widetilde\sigma_s}
 -{B_s\over\sigma_s}
 ={\widetilde B_s-B_s\over\widetilde\sigma_s}
 +B_s\left({1\over\widetilde\sigma_s}
           -{1\over\sigma_s}\right)
\]
and take norms.
\end{proof}

\begin{assessment}[Exact frontier after V79]
\label{ass:s6v79-frontier}
For the literal central Wilson transfer, the first relative sector and
the full asymptotic gap are now proved:
\[
 g_s={8\over3s}+O(s^{-2}).
\]
The finite-mass clock is \(s\delta_s\to\lambda\), and a physical
endpoint transfer within \(o(\delta_s)\) of this reference inherits the
positive mass \(8/(3\lambda)\).  This route bypasses the unresolved FHRH
sign but exposes the next true bottleneck: a same-space, normalized,
response-matched half-slab comparison at generator scale.  V73's
same-owner warning and V51's \(o(\delta)\) firewall are therefore
confirmed, not avoided.
\end{assessment}

\begin{programme}[V80: companion audit and physical endpoint defect]
\label{prog:s6v79-V80}
\begin{enumerate}[leftmargin=3em]
\item Perform the compulsory SM/NS companion audit and record only
      changes since V75.
\item Inventory the exact response-matched half-slab terms absent from
      the central reference.
\item Express their normalized amplitude difference by a common
      Duhamel formula, including the leading singular normalization.
\item Test the resulting insertion against the V71 residual readout
      and the SM V17 two-form smoothing obstruction.
\item Either prove \(o(s^{-1})\) or isolate the first term that remains
      of generator order.
\end{enumerate}
\end{programme}

\begin{firewall}[Claim boundary after sixth-series V79]
\label{fire:s6v79-boundary}
V79 proves that the fundamental and antifundamental sectors are the
first nontrivial central Wilson sectors at large coupling; computes the
global central gap as \(8/(3s)+O(s^{-2})\); proves the exact
finite-mass clock trichotomy; corrects the provisional clock branch of
V77; and proves endpoint transfer and half-amplitude perturbation
criteria at generator scale.

V79 does not prove the physical endpoint defect, identify a common
gauge-reduced boundary space for the central and response-matched
transfers, prove FHRH, close the HSOCC projective/Mosco rows, construct
the continuum OS theory, or prove the full Yang--Mills mass gap.
\end{firewall}

\begin{theorem}[Sixth-series V79 result ledger]
\label{thm:s6v79-results}
The positive conclusions in \Cref{fire:s6v79-boundary} are exactly those
of
\Cref{lem:s6v79-low-loss-asymptotic,
thm:s6v79-central-edge,
thm:s6v79-clock-trichotomy,
lem:s6v79-edge-perturbation,
thm:s6v79-endpoint-handoff,
cor:s6v79-Odelta-handoff,
lem:s6v79-amplitude-defect,
prop:s6v79-ground-ledger}.
\end{theorem}

\begin{proof}
Each item is the cited statement with its central-reference,
normalization, common-space, clock, projective, and endpoint-defect
qualifications unchanged.
\end{proof}

\paragraph{Parent-version record.}
V79 was formed from
\texttt{Renewal\_Yang\_Mills\_Mass\_Gap\_Research\_Notes\_V78.tex}.
The V78 parent had \(35\,465\) logical source lines,
\(1\,241\,726\) bytes, and SHA--256
\[
 \texttt{ac435a032221b5ad20f453c1fcd01cfba9719d130be78515d62a4c546b70450d}.
\]
The next iteration, V80, requires the compulsory companion audit.

% =============================================================================
% END APPENDED SIXTH-SERIES V79 MATERIAL
% =============================================================================

% =============================================================================
% BEGIN APPENDED SIXTH-SERIES V80 MATERIAL
% =============================================================================

\section{V80: companion audit, gauge-product reference, and the
generator-order interaction}
\label{sec:s6v80}

\subsection{Compulsory V80 companion audit}
\label{sec:s6v80-audit}

At this checkpoint the current Standard Model--Einstein companion was
\[
 \texttt{Renewal\_SM\_Einstein\_Continuum\_Research\_Notes\_V21.tex}.
\]
It had \(7\,823\) logical source lines, \(291\,914\) bytes, one live
document terminator, and SHA--256
\[
 \texttt{2a7061a4eb082502476da3178b74ad90d452a6653288609aa7639dc180084d90}.
\]
The current Navier--Stokes companion remained
\[
 \texttt{Renewal\_NS\_Stability\_Research\_Notes\_V26.tex},
\]
with \(5\,319\) logical source lines, \(278\,283\) bytes, one live
document terminator, and SHA--256
\[
 \texttt{82c887f8c2c69e4f28fad7c3dc8de5b9099cb5a4bf431eba1fff3e91935978ad}.
\]

The SM file advanced from V17 to V21 after the preceding Yang--Mills
audit.  Its four new conclusions are:
\begin{enumerate}[label=\textup{(\roman*)},leftmargin=3.5em]
\item V18 proves that an energy-form congruence is not a physical
      semigroup conjugacy; a local physical unitary freezes a principal
      kinetic form only along an isometric orbit and creates a required
      connection commutator.
\item V19 proves exact simplex cancellation for full-order insertions
      that are spectral multipliers of the same Hamiltonian.  Generic
      local metric or interaction insertions are not aligned by that
      theorem.
\item V20 proves a combined, noncommuting Cauchy estimate for a
      common-domain holomorphic sectorial family.  It controls the sum of
      both time orders and the direct contact without splitting the
      nonintegrable natural-form bridge.
\item V21 constructs an explicit complex sectorial radius from uniform
      ellipticity and a relative-form margin, but keeps a uniform
      operator-and-adjoint square-root estimate as an additional
      hypothesis.
\end{enumerate}
NS V26 is unchanged and still concerns pointwise rank-one Oseen-kernel
norms rather than reflected gauge transfer.

\begin{theorem}[V80 audit decision]
\label{thm:s6v80-audit}
No model-specific estimate is imported from either companion.  The SM
results do, however, force the following division for the Yang--Mills
endpoint problem:
\begin{enumerate}[label=\textup{(\alph*)},leftmargin=3.5em]
\item a scalar change of the electric clock may use the aligned V19
      branch;
\item a noncommuting magnetic or generated-action term must remain in a
      combined signed principal form;
\item holomorphic Cauchy control may be used for its complete contact or
      adjacent residual only after a uniform complex radius,
      square-root row, and local component sum are proved;
\item a form congruence cannot be used as though it identified the
      physical transfer spectrum.
\end{enumerate}
\end{theorem}

\begin{proof}
Items~(a)--(d) are precisely the scope restrictions in SM V18--V21.
None of those theorems contains the Yang--Mills boundary Hilbert space,
gauge projection, response-matched action, or a signed spectral lower
bound.  The NS note contains none of these objects.  The audit therefore
changes the analytic routing but imports no numerical constant.
\end{proof}

\subsection{From one central link to the full boundary product}
\label{sec:s6v80-product}

Let \(E\) be a finite oriented boundary-edge set and put
\[
 {\cal H}_E=L^2(\SU(3)^E,\dd U_E).
\]
Let \(P_s^{(1)}\) denote the normalized one-edge full transfer of
V74--V79, with eigenvalues \(q_\pi(s)=R_\pi(s)^2\).  Define
\begin{equation}
 Q_{E,s}:=\bigotimes_{e\in E}P_s^{(1)}.
 \label{eq:s6v80-product-transfer}
\end{equation}

\begin{theorem}[Volume-uniform product edge]
\label{thm:s6v80-product-edge}
The constant function is the unique eigenvector of \(Q_{E,s}\) with
eigenvalue one, and
\begin{equation}
 \sup\sigma\!\left(
 Q_{E,s}\big|_{{\bf1}^{\perp}}\right)=q_{*,s}.             \label{eq:s6v80-product-edge}
\end{equation}
Consequently
\begin{equation}
 \operatorname{gap}(I-Q_{E,s})
 ={8\over3s}+O(s^{-2}),                                  \label{eq:s6v80-product-gap}
\end{equation}
with constants independent of \(|E|\).
\end{theorem}

\begin{proof}
The Peter--Weyl tensor basis labels an eigenvector by a representation
\(\pi_e\) and matrix indices on every edge.  Its eigenvalue is
\[
 \prod_{e\in E}q_{\pi_e}(s).
\]
It equals one only when every \(\pi_e\) is trivial, because all
nontrivial one-edge coefficients are strictly below one.  On the
orthogonal complement at least one label is nontrivial, so the product
is at most \(q_{*,s}\).  Equality is attained by putting a fundamental
label on one edge and trivial labels on all others.  Apply
\Cref{thm:s6v79-central-edge}.
\end{proof}

Let \({\cal G}_E\) be the boundary gauge group, acting at vertices by
left and right multiplication on incident edges, and let
\(\Pi_E\) be the orthogonal projection onto the invariant subspace
\({\cal H}_E^{\cal G}\).

\begin{lemma}[Gauge projection commutes with the product reference]
\label{lem:s6v80-gauge-commutes}
\begin{equation}
 [Q_{E,s},\Pi_E]=0.
 \label{eq:s6v80-gauge-commutes}
\end{equation}
\end{lemma}

\begin{proof}
A central convolution commutes with both the left and right regular
actions on one edge.  Each vertex gauge transformation is a product of
such left or right actions on incident edges.  Hence
\(Q_{E,s}\) commutes with every element of \({\cal G}_E\), and therefore
with its Haar average \(\Pi_E\).
\end{proof}

\begin{corollary}[Gauge restriction cannot lower the reference gap]
\label{cor:s6v80-gauge-gap}
On the gauge-invariant vacuum complement,
\begin{equation}
 \operatorname{gap}\!\left(
 I-Q_{E,s}\big|_{{\cal H}_E^{\cal G}}\right)
 \ge {8\over3s}+O(s^{-2}),                               \label{eq:s6v80-gauge-gap}
\end{equation}
uniformly in the boundary volume.
\end{corollary}

\begin{proof}
By \Cref{lem:s6v80-gauge-commutes}, the gauge-invariant space is a
reducing subspace.  Its vacuum complement is contained in the full
vacuum complement, so its spectral supremum is no larger than
\eqref{eq:s6v80-product-edge}.  The constant remains the unique
eigenvector at one because it was unique on the larger space.
\end{proof}

\begin{remark}[The lower bound need not be sharp after gauge reduction]
\label{rem:s6v80-gauge-sharp}
A single fundamental edge is not gauge invariant.  The first admitted
spin network may therefore have a larger electric Casimir sum than
\(C_*\).  V80 uses only the volume-uniform lower bound; identifying the
first physical spin network requires the boundary topology and Gauss-law
sector and is a separate finite-regulator question.
\end{remark}

\subsection{Why an \(o(\delta)\) comparison would remove the interaction}
\label{sec:s6v80-no-subgenerator}

Let
\[
 Q_j:=Q_{E_j,s_j}\big|_{{\cal H}_{E_j}^{\cal G}}
\]
be the gauge-product reference, transported when necessary to the
declared physical boundary space.  Let \(\widetilde P_j\) be the exact
response-matched physical transfer of duration \(\delta_j\).

\begin{theorem}[Sub-generator comparison gives the same continuum
semigroup]
\label{thm:s6v80-no-subgenerator}
Suppose both families are contractions on identified Hilbert spaces and
\begin{equation}
 \|\widetilde P_j-Q_j\|=o(\delta_j).                      \label{eq:s6v80-subgenerator}
\end{equation}
For \(n_j=\lfloor t/\delta_j\rfloor\),
\begin{equation}
 \|\widetilde P_j^{\,n_j}-Q_j^{\,n_j}\|\longrightarrow0
 \label{eq:s6v80-fixed-time-same}
\end{equation}
for every fixed \(t>0\).  Hence any common strong continuum limit has
the same semigroup for the two families.
\end{theorem}

\begin{proof}
The telescoping identity gives
\[
 \widetilde P_j^{\,n}-Q_j^{\,n}
 =\sum_{k=0}^{n-1}
 \widetilde P_j^{\,n-1-k}
 (\widetilde P_j-Q_j)Q_j^k.
\]
Taking norms and using contractivity,
\[
 \|\widetilde P_j^{\,n_j}-Q_j^{\,n_j}\|
 \le n_j\|\widetilde P_j-Q_j\|
 ={t+o(1)\over\delta_j}o(\delta_j)=o(1).
\]
\end{proof}

\begin{assessment}[Endpoint comparison corrected]
\label{ass:s6v80-endpoint-correction}
The sufficient \(o(\delta_j)\) handoff in
\Cref{thm:s6v79-endpoint-handoff} is appropriate only if the physical
and reference transfers are intended to have the same continuum
generator.  A magnetic Yang--Mills term, a genuine response-matched
return, or a counterterm that survives in the continuum is of order
\(\delta_j\), not \(o(\delta_j)\).  Requiring all of it to satisfy
\eqref{eq:s6v80-subgenerator} would prove convergence to the product
electric reference and erase the very interaction whose mass gap is at
issue.
\end{assessment}

\subsection{The generator-order signed principal form}
\label{sec:s6v80-principal}

For a positive contraction \(P_j\), write its rescaled transfer form
\[
 {\cal E}_{P,j}(f)
 :={1\over\delta_j}\langle f,(I-P_j)f\rangle.
\]
After the physical ground transform and common-space identification,
suppose
\begin{equation}
 {\cal E}_{\widetilde P,j}
 ={\cal E}_{Q,j}+{\cal B}_j+{\cal R}_j                 \label{eq:s6v80-form-split}
\end{equation}
on a common vacuum-orthogonal form core.  Here \({\cal B}_j\) contains
every generator-order magnetic, spatial, gauge-incidence, and
response-matched term; \({\cal R}_j\) is reserved for a genuinely small
remainder.

\begin{theorem}[Signed principal-form gap compiler]
\label{thm:s6v80-principal-compiler}
Assume \(s_j\delta_j\to\lambda\in(0,\infty)\) and, on the common
vacuum-orthogonal core,
\begin{align}
 {\cal E}_{Q,j}(f)
 &\ge\mu_{E,j}\|f\|^2,\qquad
 \mu_{E,j}\longrightarrow\mu_E\ge{8\over3\lambda},
                                                               \label{eq:s6v80-E-gap}\\
 {\cal B}_j(f)
 &\ge-\alpha\,{\cal E}_{Q,j}(f)-b\|f\|^2,
 \qquad 0\le\alpha<1,                                    \label{eq:s6v80-B-Garding}\\
 |{\cal R}_j(f)|
 &\le\varepsilon_j\{
 {\cal E}_{Q,j}(f)+\|f\|^2\},
 \qquad\varepsilon_j\to0.                                \label{eq:s6v80-R-small}
\end{align}
If
\begin{equation}
 (1-\alpha)\mu_E-b>0,                                    \label{eq:s6v80-positive-margin}
\end{equation}
then
\begin{equation}
 \liminf_j{\cal E}_{\widetilde P,j}(f)
 \ge\{(1-\alpha)\mu_E-b\}\|f\|^2                         \label{eq:s6v80-physical-form-gap}
\end{equation}
along normalized common-core sequences for which the reference energies
remain bounded.  With the HSOCC Mosco/projective rows, the continuum
Hamiltonian has at least this gap.
\end{theorem}

\begin{proof}
Insert \eqref{eq:s6v80-B-Garding} and
\eqref{eq:s6v80-R-small} in
\eqref{eq:s6v80-form-split}:
\[
 {\cal E}_{\widetilde P,j}(f)
 \ge(1-\alpha-\varepsilon_j){\cal E}_{Q,j}(f)
 -(b+\varepsilon_j)\|f\|^2.
\]
Use \eqref{eq:s6v80-E-gap} and pass to the lower limit.  The common-core
bounded-energy condition and the HSOCC Mosco/projective rows justify
passage from core forms to the reconstructed closed form.  The strict
margin \eqref{eq:s6v80-positive-margin} gives the positive spectral
lower bound.
\end{proof}

\begin{remark}[Why positivity of the bare magnetic potential is not enough]
\label{rem:s6v80-ground-shift}
Before ground subtraction, a magnetic plaquette potential may be
nonnegative.  The physical generator is normalized by its own ground
energy and ground vector.  Transporting it to the reference
ground space creates a scalar subtraction, a density connection, and
commutator/contact terms.  Therefore
\({\cal B}_j\ge0\) cannot be inferred from bare pointwise positivity.
The signed inequality \eqref{eq:s6v80-B-Garding} must be proved after the
same ground transform that defines \(\widetilde P_j\).
\end{remark}

\begin{corollary}[Small \(O(\delta)\) operator perturbations]
\label{cor:s6v80-small-generator}
If, on the common vacuum complement,
\[
 \|\widetilde P_j-Q_j\|
 \le\varepsilon\delta_j+o(\delta_j)
\]
with \(\varepsilon<8/(3\lambda)\), then
\Cref{thm:s6v80-principal-compiler} applies with
\(\alpha=0\), \(b=\varepsilon\), and vanishing remainder, reproducing
the V79 quantitative endpoint margin.
\end{corollary}

\begin{proof}
The operator bound implies
\[
 |{\cal E}_{\widetilde P,j}(f)-{\cal E}_{Q,j}(f)|
 \le(\varepsilon+o(1))\|f\|^2.
\]
Use \(\mu_E\ge8/(3\lambda)\).
\end{proof}

\subsection{Duhamel ownership of the physical difference}
\label{sec:s6v80-Duhamel}

At a finite regulator let \(A_j(z)\), \(0\le z\le1\), be a common-space
half-amplitude path from the gauge-product reference to the exact
response-matched physical amplitude.  Let its unnormalized kernel be
\[
 B_j(z)=e^{-S_{0,j}/2}\,{\cal T}_j(z),
\]
where \({\cal T}_j(z)\) contains all shared-link integrations and the
complete action difference.  Normalize by the leading singular value
\(\sigma_j(z)\):
\[
 A_j(z)=B_j(z)/\sigma_j(z).
\]

\begin{lemma}[Exact normalized amplitude derivative]
\label{lem:s6v80-normalized-derivative}
Whenever \(B_j(z)\) and its simple leading singular value are
differentiable,
\begin{equation}
 A_j'
 ={B_j'\over\sigma_j}
 -{\sigma_j'\over\sigma_j}A_j.                           \label{eq:s6v80-normalized-derivative}
\end{equation}
Consequently
\begin{equation}
 A_j(1)-A_j(0)
 =\int_0^1
 \left\{
 {B_j'(z)\over\sigma_j(z)}
 -{\sigma_j'(z)\over\sigma_j(z)}A_j(z)
 \right\}\,\dd z.                                        \label{eq:s6v80-normalized-Duhamel}
\end{equation}
\end{lemma}

\begin{proof}
Differentiate the quotient and apply the fundamental theorem of
calculus in operator norm, or first on the common finite-regulator
Hilbert--Schmidt class and then in the controlled topology.
\end{proof}

\begin{assessment}[Term inventory]
\label{ass:s6v80-term-inventory}
For the physical half-slab, \(B_j'\) must include in one owner:
\begin{enumerate}[label=\textup{(\roman*)},leftmargin=3.5em]
\item temporal plaquette differences from the product central kernel;
\item spatial plaquettes and every boundary plaquette sharing a link;
\item the Gauss-law/gauge projection and any gauge-fixing Jacobian used
      to identify the boundary space;
\item generated or marginal Wilsonian returns and response-matching
      cofactors;
\item counterterms and anisotropic clock factors;
\item derivatives of all moving boundary measures and carriers.
\end{enumerate}
The second term in
\eqref{eq:s6v80-normalized-derivative} is the compulsory ground-factor
subtraction.  Omitting it compares unnormalized amplitudes and does not
control the physical relative spectrum.
\end{assessment}

\begin{theorem}[Audit-guided analytic routing]
\label{thm:s6v80-routing}
In the decomposition of
\eqref{eq:s6v80-normalized-Duhamel}:
\begin{enumerate}[label=\textup{(\alph*)},leftmargin=3.5em]
\item a scalar electric-clock derivative aligned with the product
      Casimir reference may be combined before norms, as in SM V19;
\item the complete generator-order noncommuting insertion and its
      ground subtraction belong to \({\cal B}_j\) in
      \eqref{eq:s6v80-form-split} and require a signed inequality;
\item only a remainder already proved small in a common form norm may
      use the SM V20 holomorphic combined-contact estimate;
\item use of that estimate additionally requires the V21 uniform
      complex-radius, square-root, and spatial component rows.
\end{enumerate}
\end{theorem}

\begin{proof}
Item~(a) is the aligned spectral-multiplier hypothesis.  The terms in
item~(b) survive after division by \(\delta_j\), so
\Cref{thm:s6v80-no-subgenerator} forbids classifying all of them as an
\(o(\delta_j)\) error.  SM V20 controls a complete analytic variation
but supplies no sign, proving item~(c).  SM V21 states the additional
uniform hypotheses needed for that complex form argument, proving
item~(d).
\end{proof}

\begin{firewall}[What the companion audit does not prove]
\label{fire:s6v80-audit-scope}
The holomorphic SM branch is a convergence and contact-control device.
It does not prove
\eqref{eq:s6v80-B-Garding}.  Conversely, the gauge-product electric
gap does not make the magnetic and response-matched principal form a
small residual.  The signed principal estimate and the analytic
remainder estimate are distinct rows and must not be substituted for
one another.
\end{firewall}

\begin{assessment}[Exact frontier after V80]
\label{ass:s6v80-frontier}
The one-link central result now yields a gauge-invariant product
reference with a volume-uniform generator gap at the finite-mass clock.
The hoped-for blanket \(o(\delta)\) comparison has been rejected for a
physical reason: it would erase every surviving Yang--Mills interaction.
The first genuinely unproved spectral row is now the signed
ground-transformed G\aa rding estimate
\[
 {\cal B}_j\ge-\alpha{\cal E}_{Q,j}-b,
 \qquad
 (1-\alpha){8\over3\lambda}-b>0,
\]
for the complete generator-order magnetic, shared-link, and
response-matched principal form.  SM V20--V21 may control the remainder
after this row, but cannot create its sign.
\end{assessment}

\begin{programme}[V81: ground-transformed magnetic principal form]
\label{prog:s6v80-V81}
\begin{enumerate}[leftmargin=3em]
\item Write the exact finite-regulator product-reference ground
      transform and the response-matched physical ground transform on
      one gauge-invariant boundary space.
\item Compute the generator-order form difference, retaining the scalar
      ground-energy subtraction and density connection.
\item Separate local nonnegative plaquette oscillations from the
      sign-indefinite ground/contact terms.
\item Test the latter with a local form bound relative to the electric
      Casimir form, keeping constants uniform in volume and refinement.
\item Apply the signed compiler only if the resulting
      \((\alpha,b)\) leave a strict part of \(8/(3\lambda)\).
\end{enumerate}
\end{programme}

\begin{firewall}[Claim boundary after sixth-series V80]
\label{fire:s6v80-boundary}
V80 performs the compulsory SM/NS audit; proves a volume-uniform product
Wilson edge and its preservation under gauge restriction; proves that
an \(o(\delta)\) physical comparison would give the same continuum
generator and therefore cannot contain a genuine Yang--Mills
interaction; proves a signed generator-order gap compiler; and gives
the exact normalized-amplitude Duhamel and term-ownership ledger.

V80 does not prove the signed physical principal-form estimate, construct
the common response-matched ground transform, populate the uniform
holomorphic/square-root remainder rows, close HSOCC or the V53
projective/Mosco hypotheses, construct the continuum OS theory, or prove
the full Yang--Mills mass gap.
\end{firewall}

\begin{theorem}[Sixth-series V80 result ledger]
\label{thm:s6v80-results}
The positive conclusions in \Cref{fire:s6v80-boundary} are exactly those
of
\Cref{thm:s6v80-audit,
thm:s6v80-product-edge,
lem:s6v80-gauge-commutes,
cor:s6v80-gauge-gap,
thm:s6v80-no-subgenerator,
thm:s6v80-principal-compiler,
cor:s6v80-small-generator,
lem:s6v80-normalized-derivative,
thm:s6v80-routing}.
\end{theorem}

\begin{proof}
Each item is the cited statement with its product-reference,
gauge-restriction, finite-mass-clock, same-ground, signed-form, and
companion-scope qualifications unchanged.
\end{proof}

\paragraph{Parent-version record.}
V80 was formed from
\texttt{Renewal\_Yang\_Mills\_Mass\_Gap\_Research\_Notes\_V79.tex}.
The V79 parent had \(35\,937\) logical source lines,
\(1\,256\,727\) bytes, and SHA--256
\[
 \texttt{d5c4b3c5a059d099493bc9e63fc82589c5b27e213b650a8603c00f63a501dc35}.
\]
The next compulsory companion audit is V85.

% =============================================================================
% END APPENDED SIXTH-SERIES V80 MATERIAL
% =============================================================================

% =============================================================================
% BEGIN APPENDED SIXTH-SERIES V81 MATERIAL
% =============================================================================

\section{V81: moving vacua and the intrinsic ground-state Dirichlet form}
\label{sec:s6v81}

\subsection{A finite-volume moving-vacuum compiler}
\label{sec:s6v81-angle}

V80 formulated a signed comparison after a common ground transform.
Before applying it, one must account for the fact that the interacting
vacuum need not equal the product-reference vacuum.  The following
elementary angle estimate makes that dependence explicit.

Let \({\cal E}_0\) be a closed nonnegative quadratic form on
\({\cal H}\), with normalized null vector \(\Omega_0\) and gap
\begin{equation}
 {\cal E}_0(f)\ge\mu_0\|f\|^2
 \qquad(f\perp\Omega_0).                                  \label{eq:s6v81-ref-gap}
\end{equation}
Let \(h\) be another normalized vector and put
\begin{equation}
 \kappa^2:=1-|\langle h,\Omega_0\rangle|^2.
 \label{eq:s6v81-kappa}
\end{equation}

\begin{lemma}[Vacuum-angle loss]
\label{lem:s6v81-vacuum-angle}
For every \(f\perp h\),
\begin{equation}
 {\cal E}_0(f)\ge
 \mu_0(1-\kappa^2)\|f\|^2.                               \label{eq:s6v81-angle-gap}
\end{equation}
\end{lemma}

\begin{proof}
Choose the phase of \(h\) so that
\[
 h=\sqrt{1-\kappa^2}\,\Omega_0+h_\perp,\qquad
 h_\perp\perp\Omega_0,\quad\|h_\perp\|=\kappa.
\]
Write \(f=c\Omega_0+f_\perp\).  Orthogonality to \(h\) gives
\[
 \sqrt{1-\kappa^2}|c|
 \le\kappa\|f_\perp\|.
\]
Since \(\|f\|^2=|c|^2+\|f_\perp\|^2\), this inequality is equivalent to
\[
 |c|^2\le\kappa^2\|f\|^2,
 \qquad
 \|f_\perp\|^2\ge(1-\kappa^2)\|f\|^2.
\]
Now use \eqref{eq:s6v81-ref-gap} and
\({\cal E}_0(f)={\cal E}_0(f_\perp)\).
\end{proof}

\begin{theorem}[Common-space moving-vacuum G\aa rding compiler]
\label{thm:s6v81-moving-vacuum}
Let
\[
 {\cal E}={\cal E}_0+{\cal B}+{\cal R}
\]
be a nonnegative closed form with normalized null vector \(h\).  Suppose
\begin{align}
 {\cal B}(f)&\ge-\alpha{\cal E}_0(f)-b\|f\|^2,
 \qquad 0\le\alpha<1,                                    \label{eq:s6v81-B}\\
 |{\cal R}(f)|&\le\varepsilon\{
 {\cal E}_0(f)+\|f\|^2\}.                                \label{eq:s6v81-R}
\end{align}
Then for \(f\perp h\),
\begin{equation}
 {\cal E}(f)\ge
 \left[
 (1-\alpha-\varepsilon)\mu_0(1-\kappa^2)
 -b-\varepsilon
 \right]\|f\|^2.                                         \label{eq:s6v81-moving-gap}
\end{equation}
\end{theorem}

\begin{proof}
The form bounds give
\[
 {\cal E}(f)
 \ge(1-\alpha-\varepsilon){\cal E}_0(f)
 -(b+\varepsilon)\|f\|^2.
\]
Apply \Cref{lem:s6v81-vacuum-angle}.
\end{proof}

\subsection{Why a global vacuum angle is not cofinal}
\label{sec:s6v81-orthogonality}

\begin{theorem}[Product orthogonality catastrophe]
\label{thm:s6v81-orthogonality-catastrophe}
Let \(\mu_0\) and \(\mu_1=h^2\mu_0\) be distinct probability measures on
one compact one-site space, with \(h>0\) and
\(\int h^2\,\dd\mu_0=1\).  Put
\[
 r:=\int h\,\dd\mu_0.
\]
Then \(0<r<1\).  On \(N\) sites, the normalized ground-density vector
\[
 h_N(x_1,\ldots,x_N)=\prod_{i=1}^Nh(x_i)
\]
satisfies
\begin{equation}
 \langle h_N,1\rangle_{L^2(\mu_0^{\otimes N})}=r^N
 \longrightarrow0.                                      \label{eq:s6v81-product-overlap}
\end{equation}
This can occur even when both the reference and tilted product
Dirichlet forms have spectral gaps bounded below independently of \(N\).
\end{theorem}

\begin{proof}
Cauchy--Schwarz gives \(r\le1\).  Equality would make \(h\) proportional
to the constant function, and normalization would give
\(\mu_1=\mu_0\), contrary to hypothesis.  Positivity gives \(r>0\).
Fubini proves \eqref{eq:s6v81-product-overlap}.  For the last assertion,
take any two one-site reversible Markov generators having positive
gaps in their respective measures and tensorize them.  The spectral gap
of a product generator is the minimum of the one-site gaps, independent
of \(N\), while the overlap still equals \(r^N\).
\end{proof}

\begin{assessment}[Global-angle comparison is the wrong volume tool]
\label{ass:s6v81-angle-fails}
In \Cref{thm:s6v81-moving-vacuum}, the factor
\(1-\kappa^2=|\langle h,\Omega_0\rangle|^2\) can vanish exponentially
with volume even when both phases have perfectly uniform gaps.  Thus a
global unitary alignment or vacuum-overlap estimate cannot be the
thermodynamic bridge from the product electric reference to the
interacting Yang--Mills transfer.  The ground transform must be used
intrinsically, in its own measure, and controlled by local conditional
or multiscale estimates.
\end{assessment}

\subsection{Exact ground-state transform}
\label{sec:s6v81-Doob}

Let \(({\cal X},\mu_0)\) be a finite-dimensional compact configuration
space at one regulator.  Let the reference form be strongly local:
\begin{equation}
 {\cal E}_0(u,v)=\int_{\cal X}\Gamma(u,v)\,\dd\mu_0,
 \label{eq:s6v81-reference-form}
\end{equation}
where \(\Gamma\) is a symmetric carré du champ satisfying the product
rule
\[
 \Gamma(fg,k)
 =f\Gamma(g,k)+g\Gamma(f,k).
\]
Let \(H_0\) be the associated nonnegative self-adjoint operator, with
constant ground state.  Let \(W\) be a real multiplication potential,
and suppose
\begin{equation}
 (H_0+W)h=eh,\qquad h>0,\qquad
 \int h^2\,\dd\mu_0=1.                                   \label{eq:s6v81-ground-equation}
\end{equation}
Define the physical ground-state measure
\begin{equation}
 \dd\mu_h=h^2\,\dd\mu_0.                                 \label{eq:s6v81-muh}
\end{equation}

\begin{theorem}[Ground-state representation]
\label{thm:s6v81-ground-representation}
For every form-core function \(f\),
\begin{equation}
 \boxed{
 \left\langle hf,(H_0+W-e)hf\right\rangle_{L^2(\mu_0)}
 =\int_{\cal X}\Gamma(f,f)\,\dd\mu_h.}
 \label{eq:s6v81-ground-representation}
\end{equation}
By polarization, the unitary map
\[
 U_h:L^2(\mu_h)\to L^2(\mu_0),\qquad U_hf=hf,
\]
conjugates \(H_0+W-e\) to the nonnegative operator associated with the
intrinsic Dirichlet form
\begin{equation}
 {\cal E}_h(f,g)
 =\int\Gamma(f,g)\,\dd\mu_h.                              \label{eq:s6v81-intrinsic-form}
\end{equation}
Its vacuum is the constant function in \(L^2(\mu_h)\).
\end{theorem}

\begin{proof}
The weak ground equation states that, for every core function \(\varphi\),
\begin{equation}
 {\cal E}_0(h,\varphi)
 +\int(W-e)h\varphi\,\dd\mu_0=0.                         \label{eq:s6v81-weak-ground}
\end{equation}
Use the product rule twice:
\[
 \Gamma(hf,hf)
 =h^2\Gamma(f,f)
  +2hf\Gamma(h,f)+f^2\Gamma(h,h),
\]
and
\[
 \Gamma(h,hf^2)
 =2hf\Gamma(h,f)+f^2\Gamma(h,h).
\]
Therefore
\[
 {\cal E}_0(hf,hf)
 =\int h^2\Gamma(f,f)\,\dd\mu_0
 +{\cal E}_0(h,hf^2).
\]
Take \(\varphi=hf^2\) in
\eqref{eq:s6v81-weak-ground}.  Its last two terms cancel the potential
and energy contribution in
\(\langle hf,(H_0+W-e)hf\rangle\), leaving precisely
\eqref{eq:s6v81-ground-representation}.  Closure and polarization give
\eqref{eq:s6v81-intrinsic-form}; \(U_h\) is unitary by
\eqref{eq:s6v81-muh}.
\end{proof}

\begin{corollary}[Mass gap equals an intrinsic Poincaré constant]
\label{cor:s6v81-Poincare}
In the setting of
\Cref{thm:s6v81-ground-representation}, the spectral gap of
\(H_0+W-e\) is the largest \(\mu\ge0\) for which
\begin{equation}
 \boxed{
 \operatorname{Var}_{\mu_h}(f)
 \le{1\over\mu}
 \int\Gamma(f,f)\,\dd\mu_h}
 \label{eq:s6v81-Poincare}
\end{equation}
holds on the form domain.
\end{corollary}

\begin{proof}
Under \(U_h\), the physical vacuum becomes the constant function and
the physical quadratic form becomes
\eqref{eq:s6v81-intrinsic-form}.  The variational characterization of
the first nonzero spectral value is exactly
\eqref{eq:s6v81-Poincare}.
\end{proof}

\begin{assessment}[The magnetic potential has not disappeared]
\label{ass:s6v81-potential-measure}
The pointwise term \(W-e\), including its extensive ground-energy
subtraction, cancels exactly from the intrinsic form.  Its full effect
has moved into the correlated density \(h^2\).  Thus V80's signed
principal-form card has an equivalent and often safer formulation:
prove a uniform Poincaré inequality for the actual
ground-state measure.  Declaring \(W\ge0\) gives no such inequality,
while estimating \(W-e\ge-e\) before the ground transform manufactures
an extensive false debit.
\end{assessment}

\subsection{A scalar-shift firewall}
\label{sec:s6v81-scalar}

\begin{proposition}[Extensive scalar shifts do not change a gap]
\label{prop:s6v81-scalar-shift}
Let \(H_N\ge0\) have ground energy zero and gap \(\mu_N\).  For any real
number \(c_N\), the operator
\[
 \widetilde H_N=H_N+c_NI
\]
has ground energy \(c_N\) and, after physical ground subtraction, the
same gap \(\mu_N\).  In particular \(c_N\) may grow proportionally to
the volume without affecting the mass gap.
\end{proposition}

\begin{proof}
Every spectral value is translated by \(c_N\).  Subtracting the new
ground energy translates the spectrum back exactly.
\end{proof}

\begin{firewall}[Do not bound before centering]
\label{fire:s6v81-center-first}
A lower bound of the form
\[
 W-e\ge-e
\]
for a nonnegative extensive potential charges the entire ground energy
as a negative scalar and can destroy a volume-uniform estimate even
when \(W\) is itself a pure scalar.  The exact ground representation
\eqref{eq:s6v81-ground-representation}, or an equivalent connected
centering identity, must be applied before absolute values.
\end{firewall}

\subsection{Yang--Mills ground-measure card}
\label{sec:s6v81-YM-card}

\begin{definition}[Gauge ground-measure Poincaré card, GMPC]
\label{def:s6v81-GMPC}
At regulator \(j\), GMPC consists of:
\begin{enumerate}[label=\textup{(G\arabic*)},leftmargin=4em]
\item a common gauge-invariant boundary configuration space
      \({\cal X}_j\) and reference carré du champ \(\Gamma_j\);
\item a positive physical ground density \(h_j\), including the exact
      response-matched action, such that
      \(\mu_j=h_j^2\mu_{0,j}\) is normalized;
\item a proof of the exact ground representation
      \[
      {\cal E}_j(f)=\int\Gamma_j(f,f)\,\dd\mu_j
      \]
      after gauge quotient and all boundary contacts;
\item a volume- and regulator-uniform Poincaré inequality
      \[
      \operatorname{Var}_{\mu_j}(f)
      \le\mu^{-1}\int\Gamma_j(f,f)\,\dd\mu_j
      \]
      for one \(\mu>0\);
\item HSOCC projective, Mosco, and fixed-time convergence for the same
      ground-transformed family.
\end{enumerate}
\end{definition}

\begin{theorem}[GMPC completion criterion]
\label{thm:s6v81-GMPC}
If GMPC holds with \(\mu>0\) along every admitted cofinal route and the
limiting physical source core is cyclic, then the reconstructed
Yang--Mills Hamiltonian has vacuum-complement gap at least \(\mu\).
\end{theorem}

\begin{proof}
Rows (G2)--(G4) and
\Cref{cor:s6v81-Poincare} give the uniform finite-regulator spectral
gap for the actual physical ground-transformed generator.  Row (G5)
passes its closed forms and semigroups to the route-independent OS
limit.  Cyclicity extends the core lower bound to the full physical
vacuum complement.
\end{proof}

\begin{firewall}[Schrödinger form is itself a hypothesis]
\label{fire:s6v81-Schrodinger}
The exact identity
\eqref{eq:s6v81-ground-representation} applies when the physical
generator differs from the reference strongly local generator by a
multiplication potential on the same configuration space.  A
response-matched transfer may also change the kinetic carré du champ,
introduce nonlocal Wilsonian returns, or change the boundary carrier.
Those changes must either be included in a more general Dirichlet form
or retained as signed principal terms.  V81 has not shown that the full
half-slab satisfies row (G3).
\end{firewall}

\begin{assessment}[Exact frontier after V81]
\label{ass:s6v81-frontier}
The extensive scalar and global vacuum-overlap obstructions have been
removed from the viable route.  For a Schrödinger-type physical
generator, the complete magnetic potential and ground-energy
subtraction are exactly equivalent to changing the reversible measure;
the mass gap is its intrinsic Poincaré constant.  Global density or
vacuum-overlap comparison cannot be uniform in volume, even for benign
product tilts.  The next upstream target is therefore a local
conditional or block Poincaré inequality for the actual
gauge-invariant ground-state measure, with gauge-coherent rotations
removed before any Dobrushin estimate.
\end{assessment}

\begin{programme}[V82: intrinsic block Poincaré compiler]
\label{prog:s6v81-V82}
\begin{enumerate}[leftmargin=3em]
\item Prove a block variance decomposition that retains gauge-invariant
      conditional sigma-algebras.
\item Combine uniform conditional block gaps with a quantitative
      approximate-tensorization constant.
\item Compare the required inter-block constant with the earlier
      coherent-rotation obstruction to one-link Dobrushin contraction.
\item Identify a minimal gauge-reduced block geometry on which the
      obstruction is absent, or prove that the present block choice is
      still insufficient.
\end{enumerate}
\end{programme}

\begin{firewall}[Claim boundary after sixth-series V81]
\label{fire:s6v81-boundary}
V81 proves a moving-vacuum angle compiler; proves that global vacuum
overlap can vanish exponentially despite uniform product gaps; proves
the exact ground-state Dirichlet-form representation and its Poincaré
equivalence; removes extensive scalar shifts before estimation; and
formulates the intrinsic GMPC completion card.

V81 does not prove that the response-matched Yang--Mills generator has
the required Schrödinger/Dirichlet form, prove a uniform Poincaré
inequality for its ground measure, close HSOCC, construct the continuum
OS theory, or prove the full Yang--Mills mass gap.
\end{firewall}

\begin{theorem}[Sixth-series V81 result ledger]
\label{thm:s6v81-results}
The positive conclusions in \Cref{fire:s6v81-boundary} are exactly those
of
\Cref{lem:s6v81-vacuum-angle,
thm:s6v81-moving-vacuum,
thm:s6v81-orthogonality-catastrophe,
thm:s6v81-ground-representation,
cor:s6v81-Poincare,
prop:s6v81-scalar-shift,
thm:s6v81-GMPC}.
\end{theorem}

\begin{proof}
Each item is the cited statement with its finite-volume,
common-space, ground-density, strongly-local, gauge, and cofinal
qualifications unchanged.
\end{proof}

\paragraph{Parent-version record.}
V81 was formed from
\texttt{Renewal\_Yang\_Mills\_Mass\_Gap\_Research\_Notes\_V80.tex}.
The V80 parent had \(36\,450\) logical source lines,
\(1\,276\,461\) bytes, and SHA--256
\[
 \texttt{93f110dab8c8c0a1b5a25dab7fb4dc9499fe929df64119ad3c976d087c921dc9}.
\]
The next compulsory companion audit is V85.

% =============================================================================
% END APPENDED SIXTH-SERIES V81 MATERIAL
% =============================================================================

% =============================================================================
% BEGIN APPENDED SIXTH-SERIES V82 MATERIAL
% =============================================================================

\section{V82: an intrinsic gauge-block Poincaré compiler}
\label{sec:s6v82}

\subsection{Relation to the earlier V59 theorem}
\label{sec:s6v82-relation}

V59 already proved an orthogonal-arm dependency compiler for the
\emph{half-slab prediction loss} and identified its critical
\(O(\sqrt\delta)\) amplitude scale.  It also proved that the existing
fixed one-site influence matrix does not populate that compiler.  V82
does not repeat those statements.  It treats the different object
created by V81: the intrinsic, continuous-time ground-state Dirichlet
form
\[
 {\cal E}_\mu(f)=\int\Gamma(f,f)\,\dd\mu
\]
and asks when conditional Poincaré inequalities for gauge-reduced
blocks imply its global Poincaré inequality.

\subsection{Conditional block variances}
\label{sec:s6v82-conditionals}

Let \(({\cal X},\mu)\) be a probability space with coordinates indexed
by a finite set \(E\).  For a block \(B\subset E\), let
\({\cal F}_{B^c}\) be the exterior sigma-algebra and define
\begin{equation}
 {\cal V}_B(f)
 :=\mathbb E_\mu\!\left[
 \operatorname{Var}_\mu(f\mid{\cal F}_{B^c})
 \right].
 \label{eq:s6v82-block-variance}
\end{equation}
Suppose the carré du champ decomposes into nonnegative local pieces
\[
 \Gamma=\sum_{a\in E}\Gamma_a.
\]
For a block \(B\), put
\begin{equation}
 {\cal E}_B(f)
 :=\int\sum_{a\in B}\Gamma_a(f,f)\,\dd\mu.
 \label{eq:s6v82-block-energy}
\end{equation}

\begin{definition}[Weighted approximate tensorization]
\label{def:s6v82-AT}
For blocks \({\cal B}\), weights \(w_B>0\), and \(C_{\rm AT}<\infty\),
the measure \(\mu\) satisfies weighted approximate tensorization if
\begin{equation}
 \operatorname{Var}_\mu(f)
 \le C_{\rm AT}\sum_{B\in{\cal B}}w_B{\cal V}_B(f)
 \label{eq:s6v82-AT}
\end{equation}
for every \(f\in L^2(\mu)\).
\end{definition}

\begin{definition}[Conditional block Poincaré row]
\label{def:s6v82-conditional-gap}
The block \(B\) has conditional gap \(\lambda_B>0\) if, for
\(\mu\)-almost every exterior configuration,
\begin{equation}
 \operatorname{Var}_\mu(f\mid{\cal F}_{B^c})
 \le {1\over\lambda_B}
 \mathbb E_\mu\!\left[
 \sum_{a\in B}\Gamma_a(f,f)
 \,\middle|\,{\cal F}_{B^c}\right].                      \label{eq:s6v82-conditional-gap}
\end{equation}
\end{definition}

\begin{theorem}[Intrinsic block Poincaré compiler]
\label{thm:s6v82-block-compiler}
Assume weighted approximate tensorization
\eqref{eq:s6v82-AT}, conditional gaps
\[
 \lambda_B\ge\lambda_{\rm loc}>0,
\]
and the weighted incidence bound
\begin{equation}
 \sum_{\substack{B\in{\cal B}\\a\in B}}w_B
 \le m_{\rm inc}
 \qquad(a\in E).                                         \label{eq:s6v82-incidence}
\end{equation}
Then
\begin{equation}
 \boxed{
 \operatorname{Var}_\mu(f)
 \le {C_{\rm AT}m_{\rm inc}\over\lambda_{\rm loc}}
 {\cal E}_\mu(f).}
 \label{eq:s6v82-global-Poincare}
\end{equation}
Thus the intrinsic generator has spectral gap at least
\begin{equation}
 \mu_{\rm glob}
 ={ \lambda_{\rm loc}\over C_{\rm AT}m_{\rm inc}}.
 \label{eq:s6v82-global-gap}
\end{equation}
\end{theorem}

\begin{proof}
Average \eqref{eq:s6v82-conditional-gap} over the exterior:
\[
 {\cal V}_B(f)\le\lambda_{\rm loc}^{-1}{\cal E}_B(f).
\]
Insert this in \eqref{eq:s6v82-AT} and exchange the two finite sums:
\[
 \operatorname{Var}_\mu(f)
 \le {C_{\rm AT}\over\lambda_{\rm loc}}
 \sum_{a\in E}
 \left(\sum_{B\ni a}w_B\right)
 \int\Gamma_a(f,f)\,\dd\mu.
\]
Apply \eqref{eq:s6v82-incidence}.
\end{proof}

\begin{remark}[The constants have distinct meanings]
\label{rem:s6v82-constants}
The local spectral number \(\lambda_{\rm loc}\), the global dependency
number \(C_{\rm AT}\), and the geometric overlap
\(m_{\rm inc}\) are separate.  A large block gap does not compensate an
unbounded approximate-tensorization constant, and bounded overlap does
not prove either analytic row.
\end{remark}

\subsection{Exact quotient descent}
\label{sec:s6v82-quotient}

Let a compact gauge group \({\cal G}\) act measurably on
\({\cal X}\), preserve \(\mu\), and preserve the carré du champ:
\[
 \Gamma(f\circ g,f\circ g)
 =\Gamma(f,f)\circ g.
\]
Let \(q:{\cal X}\to\overline{\cal X}={\cal X}/{\cal G}\) be the orbit
map and \(\overline\mu=q_\#\mu\).

\begin{lemma}[Invariant variance and energy descend exactly]
\label{lem:s6v82-quotient}
For every invariant \(f=\overline f\circ q\),
\begin{align}
 \operatorname{Var}_\mu(f)
 &=\operatorname{Var}_{\overline\mu}(\overline f),
                                                               \label{eq:s6v82-quotient-var}\\
 \int_{\cal X}\Gamma(f,f)\,\dd\mu
 &=\int_{\overline{\cal X}}
   \overline\Gamma(\overline f,\overline f)\,\dd\overline\mu,
                                                               \label{eq:s6v82-quotient-energy}
\end{align}
where the second identity defines the descended carré du champ on its
invariant form core.
\end{lemma}

\begin{proof}
The pushforward identity gives equality of every moment of
\(f=\overline f\circ q\), proving
\eqref{eq:s6v82-quotient-var}.  Invariance makes
\(\Gamma(f,f)\) constant on orbits, so it equals a measurable function
\(\overline\Gamma(\overline f,\overline f)\circ q\).
Pushforward again gives \eqref{eq:s6v82-quotient-energy}.
\end{proof}

\begin{corollary}[Quotient Poincaré equivalence]
\label{cor:s6v82-quotient-Poincare}
A Poincaré inequality on the quotient is exactly the same inequality,
with the same constant, on the invariant subspace of the unreduced
form.  Pure gauge-orbit fluctuations are absent rather than estimated.
\end{corollary}

\begin{proof}
Apply both identities in \Cref{lem:s6v82-quotient}.
\end{proof}

\begin{firewall}[Gauge fixing needs its Jacobian]
\label{fire:s6v82-gauge-fixing}
An explicit tree gauge or slice may coordinatize
\(\overline{\cal X}\), but its pushforward density and any residual
stabilizer must be included.  Replacing the quotient measure by product
Haar measure on the retained links changes both conditional variances
and local gaps.  The exact descent theorem authorizes quotienting; it
does not authorize dropping the Faddeev--Popov or orbit-volume factor.
\end{firewall}

\subsection{Heat-bath form and approximate tensorization}
\label{sec:s6v82-heat-bath}

Let
\[
 E_Bf:=\mathbb E_\mu[f\mid{\cal F}_{B^c}]
\]
and define the block heat-bath form
\begin{equation}
 {\cal D}_{\cal B}(f)
 :=\sum_{B\in{\cal B}}w_B\langle f,(I-E_B)f\rangle.
 \label{eq:s6v82-heat-bath-form}
\end{equation}

\begin{lemma}[Conditional variance is the heat-bath loss]
\label{lem:s6v82-heat-bath}
Each \(E_B\) is an orthogonal projection on \(L^2(\mu)\), and
\begin{equation}
 \langle f,(I-E_B)f\rangle={\cal V}_B(f).                 \label{eq:s6v82-heat-bath-loss}
\end{equation}
Consequently \eqref{eq:s6v82-AT} is equivalent to a spectral gap at
least \(C_{\rm AT}^{-1}\) for the weighted block heat-bath generator
\(\sum_Bw_B(I-E_B)\).
\end{lemma}

\begin{proof}
Conditional expectation is the orthogonal projection onto
\({\cal F}_{B^c}\)-measurable functions.  Therefore
\[
 \langle f,(I-E_B)f\rangle
 =\|f-E_Bf\|_2^2
 =\mathbb E\operatorname{Var}(f\mid{\cal F}_{B^c}).
\]
Sum with the weights.  The variational characterization of the first
nonzero eigenvalue gives the last assertion.
\end{proof}

\begin{assessment}[Static versus half-slab dependency]
\label{ass:s6v82-static-v-channel}
The operator in \Cref{lem:s6v82-heat-bath} is an auxiliary reversible
sampler for the \emph{ground-state measure}.  It is not the physical
half-slab transfer and is not the V59 loss operator
\((I-K^*K)^{1/2}\).  Its use is legitimate only through
\Cref{thm:s6v82-block-compiler}: prove a Poincaré inequality for the
intrinsic form, then invoke V81's exact ground-state representation.
One may not identify its Monte Carlo time with the physical clock.
\end{assessment}

\subsection{Interaction with the earlier coherent-rotation obstruction}
\label{sec:s6v82-coherent}

\begin{theorem}[What the existing one-link result does and does not rule
out]
\label{thm:s6v82-one-link-verdict}
The earlier weak-coupling coherent-rotation construction proves that a
worst-exterior one-link Dobrushin influence approaches one on unreduced
link variables.  Therefore the existing Dobrushin comparison theorem
does not provide a regulator-uniform \(C_{\rm AT}\) for one-link blocks.
It does not, by itself, disprove approximate tensorization
\eqref{eq:s6v82-AT}, because approximate tensorization is an averaged
\(L^2(\mu)\) inequality rather than a worst-exterior total-variation
bound.
\end{theorem}

\begin{proof}
The first assertion is the scope of the coherent-rotation theorem
recorded in V39: it constructs exterior configurations whose one-link
conditional laws are nearly rigid translates, forcing the supremum
influence to one.  Standard Dobrushin-to-tensorization estimates require
a strict norm below one and therefore cannot be invoked.  However,
\eqref{eq:s6v82-AT} integrates conditional variances under \(\mu\) and
contains no supremum over all exterior configurations.  Failure of a
sufficient worst-case criterion is not failure of the averaged
inequality.
\end{proof}

\begin{assessment}[Why quotienting comes first]
\label{ass:s6v82-quotient-first}
Some rigid conditional translations are gauge-orbit motions.  By
\Cref{cor:s6v82-quotient-Poincare}, those directions should be removed
exactly, not paid for in a dependency matrix.  Any remaining near-rigid
motion on the quotient is physical and must be assigned either to a
larger block or to an explicit coarse carrier.  This is consistent with
the V59 conclusion that Ward-critical modes cannot be hidden inside a
one-site fast block.
\end{assessment}

\subsection{A cofinal intrinsic block card}
\label{sec:s6v82-card}

\begin{definition}[Gauge-reduced block approximate tensorization card,
GRBAT]
\label{def:s6v82-GRBAT}
For the physical ground measure \(\mu_j\), GRBAT consists of:
\begin{enumerate}[label=\textup{(B\arabic*)},leftmargin=4em]
\item an exact quotient or gauge slice, including its pushforward
      density and residual stabilizer;
\item a covering by physical blocks \({\cal B}_j\) with weights satisfying
      a uniform incidence bound \(m_{\rm inc}\);
\item conditional quotient measures for every block and almost every
      exterior;
\item a uniform conditional Poincaré number
      \(\lambda_{\rm loc}>0\);
\item weighted approximate tensorization with
      \(C_{\rm AT}\) uniform in volume and regulator;
\item identification of the block carré du champ with the corresponding
      pieces of the physical ground-state form;
\item projective compatibility of the quotient, blocks, and ground
      measures.
\end{enumerate}
\end{definition}

\begin{theorem}[GRBAT populates the GMPC gap row]
\label{thm:s6v82-GRBAT}
If GRBAT holds, then GMPC row (G4) holds with
\begin{equation}
 \mu={\lambda_{\rm loc}\over
 C_{\rm AT}m_{\rm inc}}>0.                               \label{eq:s6v82-GRBAT-gap}
\end{equation}
If the remaining GMPC rows hold, the continuum Yang--Mills Hamiltonian
has gap at least this number.
\end{theorem}

\begin{proof}
Rows (B2)--(B6) are exactly the hypotheses of
\Cref{thm:s6v82-block-compiler} on the descended gauge-invariant form.
Use \Cref{thm:s6v81-GMPC} for the continuum handoff.
\end{proof}

\begin{firewall}[Uniform conditional gap has not been proved]
\label{fire:s6v82-local-gap}
Finite-volume smooth positivity gives a positive conditional block gap
for each fixed exterior configuration.  It does not give the uniform
number in GRBAT row (B4).  At weak coupling, conditional actions may
develop several separated maximizers for adversarial exterior
configurations.  The next iteration must test this possibility before a
worst-exterior block theorem is used.
\end{firewall}

\begin{assessment}[Exact frontier after V82]
\label{ass:s6v82-frontier}
The V81 intrinsic Poincaré target now has a noncircular local-to-global
compiler that is distinct from V59's half-slab dependency theorem.
Gauge directions descend without loss, and the remaining constants are
separated into conditional block gap, approximate tensorization, and
incidence.  The old one-link Dobrushin estimate cannot prove the
tensorization row, but does not logically refute it.  The first row to
test is more basic: whether worst-exterior conditional gaps remain
uniform at weak coupling after the chosen quotient.
\end{assessment}

\begin{programme}[V83: conditional multimodality test]
\label{prog:s6v82-V83}
\begin{enumerate}[leftmargin=3em]
\item Analyze the exact one-link matrix-Fisher conditional law for
      staple matrices with degenerate polar data.
\item Construct a separated-maximizer exterior if it is admissible and
      test its Poincaré constant by a barrier function.
\item Determine which degeneracy is pure gauge and which survives the
      quotient.
\item If worst-exterior uniformity fails, replace GRBAT row (B4) by a
      good-exterior averaged inequality with an explicit bad-set debit.
\end{enumerate}
\end{programme}

\begin{firewall}[Claim boundary after sixth-series V82]
\label{fire:s6v82-boundary}
V82 proves an intrinsic conditional-block Poincaré compiler with exact
incidence constant; proves exact descent of invariant variance and
energy to the gauge quotient; identifies conditional variance with an
auxiliary heat-bath spectral loss; distinguishes static approximate
tensorization from the physical V59 channel compiler; and formulates
the GRBAT-to-GMPC handoff.

V82 does not prove a uniform conditional block gap, approximate
tensorization for the Yang--Mills ground measure, the exact physical
ground-state representation, HSOCC, the continuum OS theory, or the full
Yang--Mills mass gap.
\end{firewall}

\begin{theorem}[Sixth-series V82 result ledger]
\label{thm:s6v82-results}
The positive conclusions in \Cref{fire:s6v82-boundary} are exactly those
of
\Cref{thm:s6v82-block-compiler,
lem:s6v82-quotient,
cor:s6v82-quotient-Poincare,
lem:s6v82-heat-bath,
thm:s6v82-one-link-verdict,
thm:s6v82-GRBAT}.
\end{theorem}

\begin{proof}
Each item is the cited statement with its intrinsic-measure,
conditional-gap, quotient, incidence, static-time, and cofinal
qualifications unchanged.
\end{proof}

\paragraph{Parent-version record.}
V82 was formed from
\texttt{Renewal\_Yang\_Mills\_Mass\_Gap\_Research\_Notes\_V81.tex}.
The V81 parent had \(36\,869\) logical source lines,
\(1\,291\,327\) bytes, and SHA--256
\[
 \texttt{8240ebfb99d2a0781c0a199350c93c9b2ccbe7d6c304337bc0733fa24987e690}.
\]
The next compulsory companion audit is V85.

% =============================================================================
% END APPENDED SIXTH-SERIES V82 MATERIAL
% =============================================================================

% =============================================================================
% BEGIN APPENDED SIXTH-SERIES V83 MATERIAL
% =============================================================================

\section{V83: an admissible center-well obstruction to worst-exterior
block gaps}
\label{sec:s6v83}

\subsection{The exact one-link conditional family}
\label{sec:s6v83-family}

The boundary-sector source gives, after all exterior links around one
selected edge are fixed, the conditional law
\begin{equation}
 \nu_{\tau,H}(\dd U)
 ={1\over Z_{\tau,H}}
 \exp\{\tau\Re\Tr(UH)\}\,\dd U,
 \qquad U\in\SU(3),                                      \label{eq:s6v83-matrix-Fisher}
\end{equation}
where \(\tau>0\) is the coupling factor and \(H\) is the sum of the
oriented staple matrices.  The fixed-card strict-contraction theorem
uses only that this density is positive.  GRBAT row (B4), by contrast,
would require a Poincaré number uniform in both \(\tau\) and the
admissible exterior \(H\).

\subsection{Two isolated center wells}
\label{sec:s6v83-wells}

Let
\[
 z=e^{2\pi i/3},\qquad
 U_+=zI,\qquad U_-=\overline z I.
\]

\begin{lemma}[Minimum real trace on \(\SU(3)\)]
\label{lem:s6v83-trace-min}
For every \(U\in\SU(3)\),
\begin{equation}
 \Re\Tr U\ge-{3\over2},                                  \label{eq:s6v83-trace-min}
\end{equation}
with equality exactly at \(U_+\) and \(U_-\).
\end{lemma}

\begin{proof}
Write the eigenangles as
\[
 \theta_1=x+y,\qquad\theta_2=x-y,\qquad\theta_3=-2x.
\]
Then
\[
 \Re\Tr U
 =2\cos x\cos y+\cos2x.
\]
For fixed \(x\), minimizing over \(y\) gives
\[
 \cos2x-2|\cos x|.
\]
With \(t=|\cos x|\in[0,1]\), this is
\[
 2t^2-2t-1
 =2\left(t-\frac12\right)^2-\frac32
 \ge-\frac32.
\]
Equality requires \(t=1/2\) and
\(\cos y=-\operatorname{sgn}(\cos x)\).  Substitution modulo
\(2\pi\) makes all three eigenvalues equal to \(z\), or all three equal
to \(\overline z\).  A unitary matrix with three equal eigenvalues is
that scalar center matrix, proving the equality statement.
\end{proof}

\begin{corollary}[Double-well matrix-Fisher law]
\label{cor:s6v83-double-well}
For \(H=-hI\), \(h>0\), the phase
\[
 \Psi_H(U):=\Re\Tr(UH)=-h\Re\Tr U
\]
has exactly two global maxima \(U_\pm\), each of value \(3h/2\).
\end{corollary}

\begin{proof}
Apply \Cref{lem:s6v83-trace-min}.
\end{proof}

\begin{lemma}[Positive barrier between the center wells]
\label{lem:s6v83-barrier}
There are disjoint conjugation-invariant neighbourhoods
\({\cal N}_\pm\) of \(U_\pm\) and a number \(\Delta>0\) such that
\begin{equation}
 \sup_{U\notin{\cal N}_+\cup{\cal N}_-}\Psi_H(U)
 \le {3h\over2}-\Delta.                                  \label{eq:s6v83-barrier}
\end{equation}
\end{lemma}

\begin{proof}
The two maximizers are isolated by
\Cref{cor:s6v83-double-well}.  Choose disjoint invariant
neighbourhoods with compact complement.  The continuous phase has a
strictly smaller maximum on that complement; their difference is
\(\Delta>0\).
\end{proof}

\subsection{Exponential collapse of the conditional Poincaré gap}
\label{sec:s6v83-gap-collapse}

Let \(\lambda(\tau,-hI)\) be the optimal constant in
\begin{equation}
 \operatorname{Var}_{\nu_{\tau,-hI}}(f)
 \le {1\over\lambda(\tau,-hI)}
 \int|\nabla f|^2\,\dd\nu_{\tau,-hI}.                    \label{eq:s6v83-Poincare}
\end{equation}

\begin{theorem}[Center-well upper bound]
\label{thm:s6v83-center-well-gap}
There are constants \(K,c>0\) such that
\begin{equation}
 \boxed{
 \lambda(\tau,-hI)\le K\tau^4e^{-c\tau}}
 \label{eq:s6v83-gap-upper}
\end{equation}
for all sufficiently large \(\tau\).  In particular no positive
Poincaré lower bound can hold uniformly over all matrix-Fisher
parameters \(H\) and weak couplings.
\end{theorem}

\begin{proof}
Choose the neighbourhoods in
\Cref{lem:s6v83-barrier}.  There is a smooth function
\(f:\SU(3)\to[-1,1]\) which equals \(1\) on a smaller neighbourhood of
\(U_+\), equals \(-1\) on a smaller neighbourhood of \(U_-\), and whose
gradient is supported in the barrier region
\[
 {\cal K}:=\SU(3)\setminus({\cal N}_+\cup{\cal N}_-).
\]
Complex conjugation preserves Haar measure and the phase and exchanges
the two wells.  Choose \(f(\overline U)=-f(U)\).  Then
\[
 \mathbb E_{\nu_{\tau,-hI}}f=0.
\]

At either nondegenerate central maximum, a coordinate ball of radius
\(\tau^{-1/2}\) has phase at least
\[
 {3h\over2}-{C\over\tau}.
\]
Its eight-dimensional Haar volume is bounded below by
\(c_0\tau^{-4}\).  Therefore
\begin{equation}
 Z_{\tau,-hI}
 \ge c_1\tau^{-4}e^{3h\tau/2}.                            \label{eq:s6v83-Z-lower}
\end{equation}
On the support of \(\nabla f\),
\eqref{eq:s6v83-barrier} gives
\[
 \int|\nabla f|^2e^{\tau\Psi_H}\,\dd U
 \le\|\nabla f\|_\infty^2
 e^{\tau(3h/2-\Delta)}.
\]
After division by \eqref{eq:s6v83-Z-lower},
\begin{equation}
 \int|\nabla f|^2\,\dd\nu_{\tau,-hI}
 \le K\tau^4e^{-\Delta\tau}.                             \label{eq:s6v83-test-energy}
\end{equation}
Laplace concentration in the two wells and the values \(f=\pm1\) imply
\(\operatorname{Var}_{\nu_{\tau,-hI}}(f)\to1\); in particular it is at
least \(1/2\) for large \(\tau\).  The Rayleigh quotient in
\eqref{eq:s6v83-Poincare} proves
\eqref{eq:s6v83-gap-upper}, with \(c=\Delta\) after changing \(K\).
\end{proof}

\begin{remark}[Only an upper bound is needed for the obstruction]
\label{rem:s6v83-upper-enough}
The theorem does not claim an Eyring--Kramers asymptotic.  A single
smooth barrier test function proves that the gap tends to zero
exponentially, which is sufficient to refute a worst-exterior uniform
lower bound.
\end{remark}

\subsection{Admissibility on a four-dimensional one-link star}
\label{sec:s6v83-admissible}

\begin{lemma}[Center-staple realization]
\label{lem:s6v83-staple-realization}
On a four-dimensional hypercubic one-link star, there is an exterior
link configuration for which the six oriented staple matrices consist
of three copies of \(zI\) and three copies of \(\overline zI\).
Consequently their sum is
\begin{equation}
 H=3(z+\overline z)I=-3I.                                \label{eq:s6v83-H-minus}
\end{equation}
\end{lemma}

\begin{proof}
For each of the six plaquettes incident to the selected link, its staple
contains the parallel link on the opposite side of that plaquette.
Those six parallel links are distinct.  Set all transverse exterior
links in the star to \(I\).  Then assign three of the opposite parallel
links the center value \(zI\) and the other three
\(\overline zI\), with inverses chosen when the orientation requires.
Every resulting oriented staple has the prescribed center value.
Because center elements commute, no further compatibility condition is
created inside the isolated star.  Their sum is
\eqref{eq:s6v83-H-minus}.
\end{proof}

\begin{corollary}[Worst-exterior one-link GRBAT row fails]
\label{cor:s6v83-GRBAT-fails}
If GRBAT row (B4) is required for every exterior configuration on
unreduced one-link blocks, then it fails on the weak-coupling cofinal
sequence.
\end{corollary}

\begin{proof}
Use the admissible exterior of
\Cref{lem:s6v83-staple-realization} in the conditional law
\eqref{eq:s6v83-matrix-Fisher}.  It has \(H=-3I\), so
\Cref{thm:s6v83-center-well-gap} makes its conditional Poincaré number
tend to zero.
\end{proof}

\subsection{The two wells are visible to plaquette invariants}
\label{sec:s6v83-gauge}

Partition the six labelled staples into
\[
 {\cal S}_+=\{zI,zI,zI\},\qquad
 {\cal S}_-=\{\overline zI,\overline zI,\overline zI\},
\]
and define
\begin{equation}
 {\cal O}(U)
 :=\sum_{S\in{\cal S}_-}\Re\Tr(US)
  -\sum_{S\in{\cal S}_+}\Re\Tr(US).                      \label{eq:s6v83-separator}
\end{equation}

\begin{lemma}[Gauge-invariant well separator]
\label{lem:s6v83-well-separator}
The summands in \eqref{eq:s6v83-separator} are labelled plaquette-trace
observables.  They are gauge invariant when the exterior links are
transformed together with \(U\), and
\begin{equation}
 {\cal O}(U_+)= {27\over2},\qquad
 {\cal O}(U_-)=-{27\over2}.                              \label{eq:s6v83-separator-values}
\end{equation}
Thus the two conditional wells are not identified by gauge conjugation
of the complete labelled star.
\end{lemma}

\begin{proof}
Each product \(US\), with the orientation fixed, is the holonomy of an
incident plaquette up to cyclic order, so its trace is gauge invariant.
At \(U_+=zI\), multiplication by a staple in
\({\cal S}_-\) gives \(I\), with real trace \(3\), while multiplication
by a staple in \({\cal S}_+\) gives \(\overline zI\), with real trace
\(-3/2\).  Hence
\[
 {\cal O}(U_+)=3(3)-3(-3/2)=27/2.
\]
Complex conjugation gives the other value.  Since a gauge orbit cannot
change a gauge-invariant observable, the two labelled-star
configurations are distinct on the quotient.
\end{proof}

\begin{firewall}[Rarity has not yet been estimated]
\label{fire:s6v83-rarity}
The center-staple exterior is highly frustrated for the positive Wilson
action and may have exponentially small probability under the physical
ground measure.  The theorem proves failure of a \emph{worst-exterior}
conditional gap; it does not prove failure of an averaged block
Poincaré inequality or of GMPC.  A good-exterior theorem must quantify
the bad-set contribution as an operator or variance debit, not merely
say that the configuration is rare.
\end{firewall}

\subsection{A corrected good-exterior compiler}
\label{sec:s6v83-good}

For each block \(B\), let \(G_B\in{\cal F}_{B^c}\) be a good-exterior
event.  Decompose
\begin{align}
 {\cal V}_B^{\rm good}(f)
 &:=\mathbb E[1_{G_B}
 \operatorname{Var}(f\mid{\cal F}_{B^c})],                \label{eq:s6v83-Vgood}\\
 {\cal V}_B^{\rm bad}(f)
 &:=\mathbb E[1_{G_B^c}
 \operatorname{Var}(f\mid{\cal F}_{B^c})].                \label{eq:s6v83-Vbad}
\end{align}

\begin{theorem}[Good-exterior approximate-tensorization compiler]
\label{thm:s6v83-good-compiler}
Assume weighted approximate tensorization with constant \(C_{\rm AT}\),
the incidence bound \eqref{eq:s6v82-incidence}, and:
\begin{align}
 {\cal V}_B^{\rm good}(f)
 &\le\lambda_{\rm good}^{-1}{\cal E}_B(f),                \label{eq:s6v83-good-gap}\\
 \sum_Bw_B{\cal V}_B^{\rm bad}(f)
 &\le\varepsilon_{\rm bad}\operatorname{Var}_\mu(f)
   +\kappa_{\rm bad}{\cal E}_\mu(f).                     \label{eq:s6v83-bad-debit}
\end{align}
If
\[
 C_{\rm AT}\varepsilon_{\rm bad}<1,
\]
then
\begin{equation}
 \boxed{
 \operatorname{Var}_\mu(f)
 \le
 {C_{\rm AT}
 \left(m_{\rm inc}/\lambda_{\rm good}
       +\kappa_{\rm bad}\right)
 \over1-C_{\rm AT}\varepsilon_{\rm bad}}
 {\cal E}_\mu(f).}
 \label{eq:s6v83-good-global}
\end{equation}
\end{theorem}

\begin{proof}
Insert
\({\cal V}_B={\cal V}_B^{\rm good}+{\cal V}_B^{\rm bad}\)
in \eqref{eq:s6v82-AT}.  Sum
\eqref{eq:s6v83-good-gap} and use the incidence bound:
\[
 \sum_Bw_B{\cal V}_B^{\rm good}
 \le {m_{\rm inc}\over\lambda_{\rm good}}{\cal E}_\mu.
\]
Use \eqref{eq:s6v83-bad-debit} for the other sum.  Move
\(C_{\rm AT}\varepsilon_{\rm bad}\operatorname{Var}_\mu(f)\)
to the left and divide by its positive coefficient.
\end{proof}

\begin{remark}[Small probability alone is insufficient on \(L^2\)]
\label{rem:s6v83-probability-insufficient}
A bound \(\mu(G_B^c)\ll1\) controls
\({\cal V}_B^{\rm bad}(f)\) only for bounded \(f\), not uniformly on the
full \(L^2\) form domain.  Row
\eqref{eq:s6v83-bad-debit} needs a weighted spectral, entropy, or
capacity estimate.  This is the precise replacement for an informal
rare-bad-exterior argument.
\end{remark}

\begin{assessment}[Exact frontier after V83]
\label{ass:s6v83-frontier}
Worst-exterior conditional Poincaré control is now rigorously closed as
a route: an admissible center-staple exterior produces two separated,
gauge-visible wells and an exponentially small one-link gap.  The
failure does not settle the averaged physical problem.  The viable
replacement is
\Cref{thm:s6v83-good-compiler}, whose new substantive row is the
bad-exterior variance/capacity debit
\eqref{eq:s6v83-bad-debit}.  The next upstream step is to derive that
debit from Wilson action cost and a local repair map, or to show by a
positive-density family of frustrated exteriors that no such estimate
can be uniform.
\end{assessment}

\begin{programme}[V84: bad-exterior capacity and repair map]
\label{prog:s6v83-V84}
\begin{enumerate}[leftmargin=3em]
\item Define a gauge-invariant frustration functional measuring distance
      of the staple sum from the unique-polar-maximizer region.
\item Construct a local exterior repair which lowers the Wilson action
      by a quantified amount while changing only boundedly many links.
\item Convert the repair injection into a probability and capacity
      bound for the bad exterior.
\item Test whether the bound controls
      \eqref{eq:s6v83-bad-debit} on the full form domain; if it controls
      only bounded observables, retain that limitation explicitly.
\end{enumerate}
\end{programme}

\begin{firewall}[Claim boundary after sixth-series V83]
\label{fire:s6v83-boundary}
V83 proves the two-center-well structure of a matrix-Fisher conditional
law; proves an exponentially small conditional Poincaré upper bound;
constructs the corresponding admissible four-dimensional staple
exterior; proves that labelled plaquette traces distinguish the wells
after gauge quotient; rules out worst-exterior one-link GRBAT; and
proves a corrected good-exterior compiler with an explicit bad-set
debit.

V83 does not estimate the physical probability or capacity of the bad
exteriors, prove averaged approximate tensorization, populate GMPC or
HSOCC, construct the continuum OS theory, or prove the full
Yang--Mills mass gap.
\end{firewall}

\begin{theorem}[Sixth-series V83 result ledger]
\label{thm:s6v83-results}
The positive conclusions in \Cref{fire:s6v83-boundary} are exactly those
of
\Cref{lem:s6v83-trace-min,
cor:s6v83-double-well,
lem:s6v83-barrier,
thm:s6v83-center-well-gap,
lem:s6v83-staple-realization,
cor:s6v83-GRBAT-fails,
lem:s6v83-well-separator,
thm:s6v83-good-compiler}.
\end{theorem}

\begin{proof}
Each item is the cited statement with its matrix-Fisher,
fixed-exterior, star-admissibility, gauge-label, good/bad, and averaged
qualifications unchanged.
\end{proof}

\paragraph{Parent-version record.}
V83 was formed from
\texttt{Renewal\_Yang\_Mills\_Mass\_Gap\_Research\_Notes\_V82.tex}.
The V82 parent had \(37\,265\) logical source lines,
\(1\,306\,163\) bytes, and SHA--256
\[
 \texttt{3e7342489942619509b9e7ccc72af17f9d47e5e1c07def09109b2b1278d67016}.
\]
The next compulsory companion audit is V85.

% =============================================================================
% END APPENDED SIXTH-SERIES V83 MATERIAL
% =============================================================================

% =============================================================================
% BEGIN APPENDED SIXTH-SERIES V84 MATERIAL
% =============================================================================

\section{V84: local frustration cost and a capacity route for bad
exteriors}
\label{sec:s6v84}

\subsection{The isolated star has an explicit exponential penalty}
\label{sec:s6v84-star-cost}

For a staple matrix \(H\), retain the matrix-Fisher partition function
\[
 Z_\tau(H)
 :=\int_{\SU(3)}e^{\tau\Re\Tr(UH)}\,\dd U.
\]
The V83 bad star has \(H_-=-3I\).  The completely aligned six-staple
star has \(H_+=6I\).

\begin{lemma}[Aligned and frustrated phase maxima]
\label{lem:s6v84-phase-maxima}
\begin{align}
 \max_{U\in\SU(3)}\Re\Tr(UH_+)&=18,                      \label{eq:s6v84-good-max}\\
 \max_{U\in\SU(3)}\Re\Tr(UH_-)&={9\over2}.               \label{eq:s6v84-bad-max}
\end{align}
The first maximum is attained only at \(I\); the second is attained
only at \(zI,\overline zI\).
\end{lemma}

\begin{proof}
For \(H_+=6I\), maximize \(6\Re\Tr U\).  Since
\(\Re\Tr U\le3\), with equality only at \(I\), this gives
\eqref{eq:s6v84-good-max}.  For \(H_-=-3I\), use
\Cref{lem:s6v83-trace-min}:
\[
 -3\Re\Tr U\le-3(-3/2)=9/2,
\]
with the stated equality cases.
\end{proof}

\begin{theorem}[Isolated-star partition penalty]
\label{thm:s6v84-star-penalty}
There are constants \(K,\tau_0<\infty\) such that
\begin{equation}
 \boxed{
 {Z_\tau(H_-)\over Z_\tau(H_+)}
 \le K e^{-27\tau/2}}
 \label{eq:s6v84-star-penalty}
\end{equation}
for \(\tau\ge\tau_0\).
\end{theorem}

\begin{proof}
Every phase is bounded by its maximum, so
\[
 Z_\tau(H_-)\le e^{(9/2)\tau}.
\]
Near the nondegenerate maximum \(I\) of the aligned phase, its Taylor
expansion is
\[
 \Re\Tr(e^{iX}H_+)
 =18-c_0|x|^2+O(|x|^4)
\]
with \(c_0>0\).  Integrating over \(|x|\le\tau^{-1/2}\) gives
\[
 Z_\tau(H_+)\ge c_1\tau^{-4}e^{18\tau}.
\]
Thus the ratio is at most
\[
 c_1^{-1}\tau^4e^{-(27/2)\tau}.
\]
For \(\tau\ge\tau_0\), absorb the polynomial into
\(K e^{-(27/2-\varepsilon)\tau}\) for any fixed
\(\varepsilon>0\).  To retain the displayed exponent exactly, use the
full Laplace lower bound
\[
 Z_\tau(H_+)\ge c_2\tau^{-4}e^{18\tau}
\]
and the corresponding Laplace upper bound
\[
 Z_\tau(H_-)\le C_2\tau^{-4}e^{(9/2)\tau},
\]
obtained by the two nondegenerate saddles and the exterior phase gap.
Their polynomial factors cancel, giving
\eqref{eq:s6v84-star-penalty}.
\end{proof}

\begin{remark}[Local suppression is not a full-lattice probability]
\label{rem:s6v84-local-only}
Equation \eqref{eq:s6v84-star-penalty} compares two conditional
one-link partition functions with prescribed staples.  In a full
lattice, changing the six opposite links also changes other plaquettes.
Their action difference can offset part of the local \(27\tau/2\)
gain.  A probability estimate therefore needs a repair map on the full
configuration and an accounting of every changed plaquette.
\end{remark}

\subsection{A repair-map probability lemma}
\label{sec:s6v84-repair}

Let
\[
 \mu_s(\dd x)=Z_s^{-1}e^{-S_s(x)}\,\rho(\dd x)
\]
on a finite-regulator configuration space, and let \(B_s\) be a bad
set.

\begin{definition}[Finite-multiplicity repair]
\label{def:s6v84-repair}
A measurable map \(R_s:B_s\to{\cal X}\) is an
\((\Delta_s,J_s,M_s)\)-repair if:
\begin{enumerate}[label=\textup{(\roman*)},leftmargin=3.5em]
\item \(S_s(R_sx)\le S_s(x)-\Delta_s\);
\item under each injective branch of \(R_s\), the reference measure
      satisfies
      \[
      \rho(\dd x)\le J_s\,\rho(\dd R_sx);
      \]
\item every point has at most \(M_s\) preimages.
\end{enumerate}
\end{definition}

\begin{lemma}[Repair suppression]
\label{lem:s6v84-repair-suppression}
If \(R_s\) is an \((\Delta_s,J_s,M_s)\)-repair, then
\begin{equation}
 \boxed{
 \mu_s(B_s)\le M_sJ_se^{-\Delta_s}.}
 \label{eq:s6v84-repair-suppression}
\end{equation}
\end{lemma}

\begin{proof}
On each injective branch, change variables \(y=R_sx\).  The action
improvement and Jacobian bound give
\[
 \int_{\text{branch}}e^{-S_s(x)}\,\rho(\dd x)
 \le J_se^{-\Delta_s}
 \int_{R_s(\text{branch})}e^{-S_s(y)}\,\rho(\dd y).
\]
Sum the branches.  The multiplicity bound counts each image point at
most \(M_s\) times.  Divide by \(Z_s\).
\end{proof}

\begin{corollary}[Cofinal exponential rarity]
\label{cor:s6v84-exponential-rarity}
If
\[
 \Delta_s\ge cs,\qquad
 \log(M_sJ_s)=o(s)
\]
for some \(c>0\), then
\[
 \mu_s(B_s)\le e^{-(c-o(1))s}.
\]
\end{corollary}

\begin{proof}
Take logarithms in \eqref{eq:s6v84-repair-suppression}.
\end{proof}

\begin{assessment}[The missing geometric repair]
\label{ass:s6v84-repair-missing}
The isolated-star calculation supplies a candidate local action gain,
but no full-lattice repair satisfying
\Cref{def:s6v84-repair} has been constructed.  A valid map must retain
Gauss-law incidence, periodic/topological sectors, response-matched
returns, and all plaquettes touching the repaired links.  The number of
branches and the Haar Jacobian must be subexponential in the coupling
uniformly in volume.
\end{assessment}

\subsection{Why probability is not capacity}
\label{sec:s6v84-capacity}

\begin{lemma}[Small measure does not control an \(L^2\) bad part]
\label{lem:s6v84-small-probability}
On every nonatomic probability space and for every
\(p\in(0,1)\), there are an event \(B\) of probability \(p\) and a
unit vector \(g\in L^2\) supported on \(B\).  Hence
\begin{equation}
 \int_B|g|^2\,\dd\mu=1                                   \label{eq:s6v84-L2-bad}
\end{equation}
even as \(p\downarrow0\).
\end{lemma}

\begin{proof}
Choose \(\mu(B)=p\) and take
\[
 g=p^{-1/2}1_B.
\]
\end{proof}

\begin{assessment}[Required strengthening]
\label{ass:s6v84-capacity-needed}
A repair estimate for \(\mu(B_s)\) cannot by itself prove the bad-set
debit \eqref{eq:s6v83-bad-debit} on the full form domain.  One needs
integrability improvement, a capacity inequality, or a spectral
penalty preventing normalized low-energy functions from concentrating
on \(B_s\).
\end{assessment}

\subsection{A response-moment conversion}
\label{sec:s6v84-moment}

For a block \(B\), define the exterior conditional-variance field
\begin{equation}
 v_B(\omega)
 :=\operatorname{Var}_\mu(f\mid{\cal F}_{B^c})(\omega).
 \label{eq:s6v84-vB}
\end{equation}

\begin{theorem}[Rarity plus an \(L^r\) response bound]
\label{thm:s6v84-rarity-response}
Fix \(r>1\).  Suppose
\begin{equation}
 \sup_{B\in{\cal B}}\mu(G_B^c)\le p                       \label{eq:s6v84-p}
\end{equation}
and there are constants \(S_V,S_E\) such that
\begin{equation}
 \sum_{B\in{\cal B}}w_B
 \|v_B\|_{L^r(\mu|_{{\cal F}_{B^c}})}
 \le
 S_V\operatorname{Var}_\mu(f)+S_E{\cal E}_\mu(f).
 \label{eq:s6v84-response-moment}
\end{equation}
Then
\begin{equation}
 \sum_Bw_B{\cal V}_B^{\rm bad}(f)
 \le p^{1-1/r}
 \left\{
 S_V\operatorname{Var}_\mu(f)+S_E{\cal E}_\mu(f)
 \right\}.                                               \label{eq:s6v84-bad-from-moment}
\end{equation}
Thus the V83 bad debit holds with
\[
 \varepsilon_{\rm bad}=p^{1-1/r}S_V,\qquad
 \kappa_{\rm bad}=p^{1-1/r}S_E.
\]
\end{theorem}

\begin{proof}
Since \(G_B^c\) is exterior-measurable, Hölder gives
\[
 {\cal V}_B^{\rm bad}(f)
 =\mathbb E[1_{G_B^c}v_B]
 \le\mu(G_B^c)^{1-1/r}\|v_B\|_{L^r}
 \le p^{1-1/r}\|v_B\|_{L^r}.
\]
Multiply by \(w_B\), sum, and use
\eqref{eq:s6v84-response-moment}.
\end{proof}

\begin{corollary}[Exponential repair absorbs polynomial response growth]
\label{cor:s6v84-exp-absorbs}
Suppose a repair gives \(p_s\le e^{-cs}\), while for one \(r>1\)
\[
 S_V(s)+S_E(s)\le K(1+s)^m
\]
for fixed \(K,m\).  Then the bad debit coefficients in
\Cref{thm:s6v84-rarity-response} tend to zero exponentially.
\end{corollary}

\begin{proof}
The factor \(p_s^{1-1/r}\) is exponential in \(-s\) and dominates every
fixed polynomial.
\end{proof}

\begin{firewall}[The moment bound is substantive]
\label{fire:s6v84-moment}
Equation \eqref{eq:s6v84-response-moment} is not a consequence of
\(v_B\in L^1\).  It is an \(L^2\)-to-\(L^{2r}\) improvement for
conditional fluctuations, summed without volume loss.  A log-Sobolev,
super-Poincaré, or local hypercontractive estimate could supply it, but
none has yet been proved for the response-matched Yang--Mills ground
measure.
\end{firewall}

\subsection{A direct capacity alternative}
\label{sec:s6v84-direct-capacity}

\begin{definition}[Bad-exterior form capacity]
\label{def:s6v84-bad-capacity}
The family \((G_B^c)\) has weighted form-capacity pair
\((\varepsilon,\kappa)\) if
\begin{equation}
 \sum_Bw_B
 \mathbb E[1_{G_B^c}v_B]
 \le\varepsilon\operatorname{Var}_\mu(f)
  +\kappa{\cal E}_\mu(f)                                 \label{eq:s6v84-capacity}
\end{equation}
for every form-domain \(f\).
\end{definition}

\begin{theorem}[Capacity closes the good-exterior branch]
\label{thm:s6v84-capacity-closes}
Assume approximate tensorization with \(C_{\rm AT}\), a good-exterior
conditional gap \(\lambda_{\rm good}\), incidence
\(m_{\rm inc}\), and the capacity bound
\eqref{eq:s6v84-capacity}.  If
\[
 C_{\rm AT}\varepsilon<1,
\]
then the intrinsic physical gap is at least
\begin{equation}
 \boxed{
 {1-C_{\rm AT}\varepsilon\over
 C_{\rm AT}(m_{\rm inc}/\lambda_{\rm good}+\kappa)}.}
 \label{eq:s6v84-capacity-gap}
\end{equation}
\end{theorem}

\begin{proof}
This is \Cref{thm:s6v83-good-compiler} with the bad debit renamed as
the capacity pair, followed by inversion of its Poincaré constant.
\end{proof}

\subsection{Revised acquisition card}
\label{sec:s6v84-card}

\begin{definition}[Frustrated-exterior repair/capacity card, FERC]
\label{def:s6v84-FERC}
FERC consists of:
\begin{enumerate}[label=\textup{(F\arabic*)},leftmargin=4em]
\item a gauge-invariant definition of the bad exterior containing every
      separated-maximizer staple configuration;
\item a full-lattice repair map with action gain \(cs\), controlled
      multiplicity and Jacobian;
\item either the response-moment bound
      \eqref{eq:s6v84-response-moment} or the direct capacity bound
      \eqref{eq:s6v84-capacity};
\item a good-exterior conditional Poincaré gap;
\item approximate tensorization for the actual ground measure;
\item compatibility with the gauge quotient and cofinal bridges.
\end{enumerate}
\end{definition}

\begin{theorem}[FERC-to-GMPC handoff]
\label{thm:s6v84-FERC}
If FERC holds with uniform
\(C_{\rm AT},m_{\rm inc},\lambda_{\rm good}\), and if its bad capacity
pair satisfies
\[
 \limsup C_{\rm AT}\varepsilon<1,\qquad
 \sup\kappa<\infty,
\]
then GMPC has a positive uniform finite-regulator gap.  If
\(\varepsilon,\kappa\to0\), the limiting lower bound is
\(\lambda_{\rm good}/(C_{\rm AT}m_{\rm inc})\).
\end{theorem}

\begin{proof}
Use the repair lemma and response-moment theorem to obtain the capacity
pair, or use the direct pair in FERC row (F3).  Apply
\Cref{thm:s6v84-capacity-closes} and then
\Cref{thm:s6v81-GMPC}.
\end{proof}

\begin{assessment}[Exact frontier after V84]
\label{ass:s6v84-frontier}
The V83 bad star has a local free-energy penalty of
\(27\tau/2\) relative to the aligned star.  A full-lattice repair with
only subexponential entropy loss would make such exteriors
exponentially rare.  V84 also proves exactly why rarity alone is
insufficient and gives two rigorous conversions to the needed
bad-variance debit: an \(L^r\) response bound or a direct form-capacity
bound.  The unresolved rows are now geometric repair on the complete
response-matched action, integrability/capacity of conditional
fluctuations, and approximate tensorization.
\end{assessment}

\begin{programme}[V85: companion audit and repair feasibility]
\label{prog:s6v84-V85}
\begin{enumerate}[leftmargin=3em]
\item Perform the compulsory SM/NS companion audit and record only
      changes since V80.
\item Check whether the SM holomorphic/square-root branch supplies the
      \(L^r\) response improvement in
      \eqref{eq:s6v84-response-moment}.
\item Audit the Yang--Mills sources for a full-lattice local repair or
      Peierls map that includes generated action terms.
\item If neither exists, retain FERC as an open theorem card and move
      upstream to a multiscale rather than worst-block decomposition.
\end{enumerate}
\end{programme}

\begin{firewall}[Claim boundary after sixth-series V84]
\label{fire:s6v84-boundary}
V84 computes the isolated-star exponential partition penalty; proves a
finite-multiplicity repair suppression lemma; proves that small
probability alone does not control \(L^2\) bad mass; proves conversion
of exponential rarity plus an \(L^r\) response bound into the V83 bad
debit; and proves the direct capacity and FERC completion compilers.

V84 does not construct the full-lattice response-matched repair, prove
the \(L^r\) response or capacity bound, establish approximate
tensorization, populate GMPC or HSOCC, construct the continuum OS
theory, or prove the full Yang--Mills mass gap.
\end{firewall}

\begin{theorem}[Sixth-series V84 result ledger]
\label{thm:s6v84-results}
The positive conclusions in \Cref{fire:s6v84-boundary} are exactly those
of
\Cref{lem:s6v84-phase-maxima,
thm:s6v84-star-penalty,
lem:s6v84-repair-suppression,
cor:s6v84-exponential-rarity,
lem:s6v84-small-probability,
thm:s6v84-rarity-response,
cor:s6v84-exp-absorbs,
thm:s6v84-capacity-closes,
thm:s6v84-FERC}.
\end{theorem}

\begin{proof}
Each item is the cited statement with its isolated-star,
full-lattice-repair, probability-versus-capacity, response-moment,
ground-measure, and cofinal qualifications unchanged.
\end{proof}

\paragraph{Parent-version record.}
V84 was formed from
\texttt{Renewal\_Yang\_Mills\_Mass\_Gap\_Research\_Notes\_V83.tex}.
The V83 parent had \(37\,685\) logical source lines,
\(1\,320\,869\) bytes, and SHA--256
\[
 \texttt{a5f0e704d129ff3249f8b862e4f193aa65b281ea7421098eebe234cb7b5d0123}.
\]
The next iteration, V85, requires the compulsory companion audit.

% =============================================================================
% END APPENDED SIXTH-SERIES V84 MATERIAL
% =============================================================================

% =============================================================================
% BEGIN APPENDED SIXTH-SERIES V85 MATERIAL
% =============================================================================

\section{V85: companion audit and the sparse bad-cluster correction}
\label{sec:s6v85}

\subsection{Compulsory V85 companion audit}
\label{sec:s6v85-audit}

At the completed audit read, the current Standard Model--Einstein
companion was
\[
 \texttt{Renewal\_SM\_Einstein\_Continuum\_Research\_Notes\_V25.tex}.
\]
It had \(9\,134\) logical source lines, \(341\,078\) bytes, one live
document terminator, and SHA--256
\[
 \texttt{85687530d16cca4378a21dd8882fffde47f36bed727756bdc6667f418d65c29c}.
\]
The occurrences of the terminator string earlier in that file are
verbatim prose, not live terminators.  The Navier--Stokes companion
remained
\[
 \texttt{Renewal\_NS\_Stability\_Research\_Notes\_V26.tex},
\]
with \(5\,319\) logical source lines, \(278\,283\) bytes, one live
terminator, and SHA--256
\[
 \texttt{82c887f8c2c69e4f28fad7c3dc8de5b9099cb5a4bf431eba1fff3e91935978ad}.
\]

The new SM conclusions since the V80 audit are:
\begin{enumerate}[label=\textup{(\roman*)},leftmargin=3.5em]
\item V22 proves cutoff-independent square-root constants on the real
      self-adjoint and complex energy-commuting branches, but leaves the
      generic nonnormal complex graph row as DQSRC.
\item V23 proves a geometrically summable
      \(O(h_n^\alpha)\) adjacent complex metric-form defect for uniformly
      \(C^{0,\alpha}\) coefficients.
\item V24 proves exact nested \(P_1\) Galerkin bridge commutation and zero
      or \(O(h_n^\alpha)\) kinetic quadrature error, while retaining
      nonlinear and independently blocked gauge-link contacts.
\item V25 adopts the Yang--Mills good-chart/bad-capacity distinction and
      formulates a bad-geometry form debit; it does not prove that debit.
\end{enumerate}

The attached Yang--Mills sources were also searched for a full-lattice
repair, Peierls, or bad-exterior capacity theorem.  They contain finite
Wilson geometry, source-visible Witten capacities, and fixed-card
conditional contraction, but no theorem with the repair-map hypotheses
of V84 for the complete response-matched action.

\begin{theorem}[V85 audit decision]
\label{thm:s6v85-audit}
Neither companion supplies the response-moment estimate
\eqref{eq:s6v84-response-moment}, the direct bad-exterior capacity
\eqref{eq:s6v84-capacity}, or a full-lattice Yang--Mills repair.
The SM square-root and Hölder-bridge results may be used later to land a
remainder already controlled on a good chart, but they do not improve
the \(L^2\) bad-set debit.  No NS result is imported.
\end{theorem}

\begin{proof}
SM V22 controls form-domain square roots, not an
\(L^2\)-to-\(L^{2r}\) norm.  V23--V24 control adjacent discretization
errors on a uniformly regular chart.  V25 explicitly retains its
bad-geometry debit as a hypothesis.  None is a probability-to-capacity
conversion for the Yang--Mills ground measure.  The NS file is
unchanged and has no gauge ground measure.  The source search found no
map satisfying the complete action, multiplicity, and Jacobian rows of
\Cref{def:s6v84-repair}.
\end{proof}

\subsection{Why a globally good exterior is impossible in infinite
volume}
\label{sec:s6v85-global-good}

\begin{theorem}[Positive local defect density destroys the all-good
event]
\label{thm:s6v85-all-good}
Let \(B_1,\ldots,B_N\) be local bad-block events revealed in that order.
Suppose
\begin{equation}
 \mathbb P(B_k\mid B_1^c,\ldots,B_{k-1}^c)\ge p>0
 \label{eq:s6v85-conditional-lower}
\end{equation}
whenever the conditioning event has positive probability.  Then
\begin{equation}
 \mathbb P\!\left(\bigcap_{k=1}^NB_k^c\right)
 \le(1-p)^N.                                             \label{eq:s6v85-all-good}
\end{equation}
In particular the probability that every block is good tends to zero
as \(N\to\infty\).
\end{theorem}

\begin{proof}
Iterated conditioning gives
\[
 \mathbb P\!\left(\bigcap_{k=1}^NB_k^c\right)
 =\prod_{k=1}^N
 \mathbb P\!\left(
 B_k^c\,\middle|\,\bigcap_{i<k}B_i^c
 \right).
\]
Each factor is at most \(1-p\) by
\eqref{eq:s6v85-conditional-lower}.
\end{proof}

\begin{corollary}[Exponential rarity per block is not global exclusion]
\label{cor:s6v85-local-not-global}
Even if \(p=p_s\le e^{-cs}\), the thermodynamic limit at a fixed
finite \(s\) contains infinitely many opportunities for a bad block.
A proof based on the event \emph{no bad block anywhere} cannot be
uniform in volume.
\end{corollary}

\begin{proof}
For each fixed \(s\), \(p_s>0\) unless the bad event is impossible.
Apply \Cref{thm:s6v85-all-good} along arbitrarily many separated blocks.
\end{proof}

\begin{assessment}[Correct thermodynamic meaning of a repair]
\label{ass:s6v85-repair-meaning}
A repair estimate should make bad blocks sparse and their connected
clusters short; it should not attempt to remove all bad blocks.  The
local form/capacity argument must tolerate infinitely many separated
defects while charging only bounded overlap near each one.
\end{assessment}

\subsection{Connected-set counting}
\label{sec:s6v85-counting}

Let \({\mathbb G}=({\mathbb V},{\mathbb E})\) be the block adjacency
graph and assume its maximum degree is at most \(D\).

\begin{lemma}[Rooted connected-set bound]
\label{lem:s6v85-lattice-animal}
For every vertex \(o\) and integer \(n\ge1\), the number
\(A_n(o)\) of connected vertex sets of cardinality \(n\) containing
\(o\) satisfies
\begin{equation}
 A_n(o)\le(eD)^{n-1}.                                    \label{eq:s6v85-animal}
\end{equation}
\end{lemma}

\begin{proof}
Every connected set has a spanning tree rooted at \(o\).  The standard
exponential generating-function bound for rooted labelled trees,
together with at most \(D\) choices for each tree edge, gives
\[
 A_n(o)\le {n^{n-2}\over(n-1)!}D^{n-1}\le(eD)^{n-1}.
\]
Multiple tree encodings only overcount the connected sets, so the bound
is valid.
\end{proof}

\begin{remark}[A coarser constant would suffice]
\label{rem:s6v85-count-constant}
Nothing below depends on the sharp factor \(eD\).  A depth-first
encoding gives the elementary replacement \((4D)^{n-1}\), which leaves
the criterion unchanged in kind.  Only exponential, rather than
volume-dependent, growth in cluster size matters.
\end{remark}

\begin{definition}[Bernoulli domination of bad blocks]
\label{def:s6v85-domination}
Bad indicators \((\xi_v)_{v\in{\mathbb V}}\) are \(p\)-dominated if, for
every finite \(F\subset{\mathbb V}\),
\begin{equation}
 \mathbb P(\xi_v=1\ \text{for all }v\in F)\le p^{|F|}.
 \label{eq:s6v85-domination}
\end{equation}
\end{definition}

\begin{theorem}[Bad-cluster tail]
\label{thm:s6v85-cluster-tail}
Under \(p\)-domination,
\begin{equation}
 \mathbb P\bigl(
 |\mathcal C(o)|\ge n,\ \xi_o=1
 \bigr)
 \le
 \sum_{m\ge n}(eD)^{m-1}p^m,                             \label{eq:s6v85-cluster-tail}
\end{equation}
where \(\mathcal C(o)\) is the bad connected component of \(o\).
If \(eDp<1\), this is at most
\begin{equation}
 {p(eDp)^{n-1}\over1-eDp}.                               \label{eq:s6v85-cluster-geometric}
\end{equation}
\end{theorem}

\begin{proof}
If the bad cluster of \(o\) has at least \(n\) vertices, it contains a
connected bad set of some size \(m\ge n\) containing \(o\).
For every such set, domination bounds the probability that all its
vertices are bad by \(p^m\).  Sum over the sets and use
\Cref{lem:s6v85-lattice-animal}.  The geometric-series formula gives
\eqref{eq:s6v85-cluster-geometric}.
\end{proof}

\subsection{A weighted sparse-cluster summability test}
\label{sec:s6v85-weighted}

\begin{proposition}[Exponential defect cost versus cluster entropy]
\label{prop:s6v85-cluster-sum}
Suppose a connected bad cluster of size \(n\) carries a local repair,
response, or capacity coefficient at most
\[
 A e^{an}
\]
for constants \(A,a\ge0\).  Under \(p\)-domination, the expected
rooted coefficient is bounded by
\begin{equation}
 \sum_{n\ge1}A e^{an}A_n(o)p^n
 \le
 {A e^ap\over1-e^{a+1}Dp}                               \label{eq:s6v85-weighted-sum}
\end{equation}
provided
\begin{equation}
 e^{a+1}Dp<1.                                            \label{eq:s6v85-subcritical}
\end{equation}
\end{proposition}

\begin{proof}
Use \eqref{eq:s6v85-animal}:
\[
 \sum_{n\ge1}Ae^{an}(eD)^{n-1}p^n
 =Ae^ap\sum_{n\ge1}(e^{a+1}Dp)^{n-1}.
\]
Condition \eqref{eq:s6v85-subcritical} permits summation.
\end{proof}

\begin{corollary}[Weak-coupling cluster margin]
\label{cor:s6v85-weak-cluster}
If a complete-action repair gives
\[
 p_s\le e^{-cs+o(s)}
\]
and the cluster capacity cost has
\[
 a_s=o(s),
\]
then \eqref{eq:s6v85-subcritical} holds for all sufficiently large
\(s\), for every fixed bounded-degree block graph.
\end{corollary}

\begin{proof}
Take logarithms:
\[
 \log(e^{a_s+1}Dp_s)
 \le a_s+1+\log D-cs+o(s)\longrightarrow-\infty.
\]
\end{proof}

\begin{firewall}[Cluster counting is not a capacity theorem]
\label{fire:s6v85-counting-only}
The factor \(Ae^{an}\) in
\Cref{prop:s6v85-cluster-sum} must already bound the relevant
form-domain response or capacity of one specified cluster.  Counting
does not manufacture that bound from probability.  V84's
\(L^r\)-response or direct-capacity row is still required, now locally
on connected clusters rather than globally on the event of any bad
block.
\end{firewall}

\subsection{Sparse frustrated-cluster card}
\label{sec:s6v85-card}

\begin{definition}[Sparse frustrated-cluster card, SFCC]
\label{def:s6v85-SFCC}
For the physical ground measure and a bounded-degree gauge-reduced block
graph, SFCC consists of:
\begin{enumerate}[label=\textup{(S\arabic*)},leftmargin=4em]
\item a local bad-block event containing every conditional
      separated-well configuration;
\item \(p_s\)-domination of all finite bad-block sets, derived from the
      complete response-matched action;
\item a good-component conditional Poincaré inequality;
\item for each connected bad cluster \(C\), a repair/capacity estimate
      with cost at most \(Ae^{a_s|C|}\);
\item the strict entropy margin
      \(e^{a_s+1}Dp_s<1\);
\item bounded overlap of enlarged cluster neighbourhoods in the
      assembled form estimate;
\item approximate tensorization or a multiscale martingale
      decomposition for the good components and repaired clusters.
\end{enumerate}
\end{definition}

\begin{theorem}[What SFCC would contribute]
\label{thm:s6v85-SFCC}
Rows (S2), (S4), and (S5) make the rooted sum of all connected bad-cluster
capacity coefficients finite and volume independent, with the explicit
bound \eqref{eq:s6v85-weighted-sum}.  If rows (S3), (S6), and (S7)
insert that sum into the V83 good-exterior compiler with absorption
coefficient below one, then GRBAT and hence GMPC hold.
\end{theorem}

\begin{proof}
\Cref{prop:s6v85-cluster-sum} proves the volume-independent rooted sum.
Bounded overlap turns rooted local bounds into a global form debit.
The assumed strict absorption row permits
\Cref{thm:s6v83-good-compiler}; then apply
\Cref{thm:s6v82-GRBAT,thm:s6v81-GMPC}.
\end{proof}

\begin{assessment}[Exact frontier after V85]
\label{ass:s6v85-frontier}
The companion audit supplies no missing Yang--Mills capacity estimate.
The thermodynamic mistake in a global good-exterior event has now been
removed: rare frustrated blocks must be treated as a dilute random
cluster gas.  The combinatorial part is closed by
\Cref{thm:s6v85-cluster-tail,prop:s6v85-cluster-sum}.  Two analytic
rows remain genuinely open: complete-action stochastic domination of
bad clusters and a connected-cluster form-capacity bound whose
exponential cost loses more slowly than the Wilson action gain.
\end{assessment}

\begin{programme}[V86: complete-action bad-cluster domination]
\label{prog:s6v85-V86}
\begin{enumerate}[leftmargin=3em]
\item Define bad blocks by a gauge-invariant plaquette-frustration
      threshold rather than by the staple sum alone.
\item Attempt a disjoint local repair on a separated family of blocks,
      accounting for every plaquette altered by the repair.
\item Derive the finite-set bound
      \(\mathbb P(\cap_{v\in F}B_v)\le p_s^{|F|}\), with colouring if
      repairs overlap.
\item Compare the action gain with the entropy and Jacobian losses.
\item If generated response-matched terms have no controlled local
      oscillation, record that term as the first obstruction and move
      upstream to its locality theorem.
\end{enumerate}
\end{programme}

\begin{firewall}[Claim boundary after sixth-series V85]
\label{fire:s6v85-boundary}
V85 performs the compulsory SM/NS and source audit; proves that an
all-good-block event cannot be the thermodynamic mechanism; proves a
rooted connected-set bound, a bad-cluster tail under Bernoulli
domination, and the exact entropy-versus-capacity summability criterion;
and formulates the sparse frustrated-cluster card.

V85 does not prove complete-action bad-block domination, a connected
cluster capacity estimate, approximate tensorization, GMPC or HSOCC,
the continuum OS theory, or the full Yang--Mills mass gap.
\end{firewall}

\begin{theorem}[Sixth-series V85 result ledger]
\label{thm:s6v85-results}
The positive conclusions in \Cref{fire:s6v85-boundary} are exactly those
of
\Cref{thm:s6v85-audit,
thm:s6v85-all-good,
cor:s6v85-local-not-global,
lem:s6v85-lattice-animal,
thm:s6v85-cluster-tail,
prop:s6v85-cluster-sum,
cor:s6v85-weak-cluster,
thm:s6v85-SFCC}.
\end{theorem}

\begin{proof}
Each item is the cited statement with its companion-scope,
thermodynamic, stochastic-domination, bounded-degree, cluster-capacity,
and cofinal qualifications unchanged.
\end{proof}

\paragraph{Parent-version record.}
V85 was formed from
\texttt{Renewal\_Yang\_Mills\_Mass\_Gap\_Research\_Notes\_V84.tex}.
The V84 parent had \(38\,120\) logical source lines,
\(1\,334\,905\) bytes, and SHA--256
\[
 \texttt{0fc9cb0e118e707cc7b51299d196fa8861b4fdde8e72ba339e26b4bef0825c7c}.
\]
The next compulsory companion audit is V90.

% =============================================================================
% END APPENDED SIXTH-SERIES V85 MATERIAL
% =============================================================================

% =============================================================================
% BEGIN APPENDED SIXTH-SERIES V86 MATERIAL
% =============================================================================

\section{V86: reflection-positive domination of Wilson
star-frustration}
\label{sec:s6v86}

\subsection{Purpose and change of route}
\label{sec:s6v86-purpose}

V85 asked for a complete-action local repair producing Bernoulli
domination of bad blocks.  There is a preliminary obstruction to that
literal route.  The V83 bad event is an event of the \emph{exterior}
staples of a selected link.  Once those exterior links are conditioned
upon, its indicator is deterministic.  Consequently no nontrivial
upper bound on its conditional probability can be obtained by
resampling only the selected link.

This note replaces that invalid one-link conditioning step by a
global-to-local argument.  First, a gauge-invariant star-frustration
functional converts each V83 exterior into a definite amount of
actual Wilson plaquette energy, regardless of the value of the central
link.  Second, simultaneous occurrence on a reflected tiling forces an
extensive action cost.  Third, reflection positivity and the
chessboard inequality convert the disseminated cost into the required
finite-set domination.  This closes SFCC row (S2) for the
\emph{pure Wilson measure}.  The extension to the physical
response-matched ground measure is reduced to an exact
reflection-positivity and relative-stability card.

\subsection{Gauge-invariant star frustration}
\label{sec:s6v86-star}

Let
\[
 d(V):=3-\Re\Tr V,\qquad V\in\SU(3).
\]
On an oriented four-dimensional hypercubic lattice, fix an oriented
edge \(e\).  Write the six incident oriented plaquette holonomies in
the form
\[
 U_p=U_eS_{e,p},\qquad p\ni e,
\]
after cyclically orienting their traces.  The staple sum and the
star-frustration functional are
\begin{align}
 H_e&:=\sum_{p\ni e}S_{e,p},                              \label{eq:s6v86-H}\\
 {\mathfrak F}_e
 &:=\min_{V\in\SU(3)}\sum_{p\ni e}d(VS_{e,p})
   =18-\max_{V\in\SU(3)}\Re\Tr(VH_e).                    \label{eq:s6v86-F}
\end{align}

\begin{lemma}[Gauge invariance and energetic meaning]
\label{lem:s6v86-F-basic}
The quantity \({\mathfrak F}_e\) is gauge invariant, depends only on
the exterior staples, and satisfies
\begin{equation}
 0\le {\mathfrak F}_e
 \le\sum_{p\ni e}d(U_p)                                  \label{eq:s6v86-F-energy}
\end{equation}
for every value of the central link \(U_e\).
\end{lemma}

\begin{proof}
A gauge transformation at the initial and terminal vertices of \(e\)
sends
\[
 U_e\longmapsto g_xU_eg_y^{-1},
 \qquad
 H_e\longmapsto g_yH_eg_x^{-1}.
\]
The substitution \(V\mapsto g_x^{-1}Vg_y\) is a bijection of
\(\SU(3)\), and cyclicity of the trace proves invariance of the maximum
in \eqref{eq:s6v86-F}.  The formula uses only the staples and therefore
does not involve \(U_e\).  Nonnegativity follows from \(d\ge0\).
Finally, \(U_e\) is one admissible value of \(V\) in the minimum in
\eqref{eq:s6v86-F}; this gives \eqref{eq:s6v86-F-energy}.
\end{proof}

\begin{lemma}[Exact V83 frustration value]
\label{lem:s6v86-center-F}
For the V83 center-staple exterior, consisting of three staples \(zI\)
and three staples \(\overline zI\),
\begin{equation}
 H_e=-3I,\qquad {\mathfrak F}_e={27\over2}.               \label{eq:s6v86-center-F}
\end{equation}
Hence, for every fixed
\(\kappa\in(0,27/2)\), the gauge-invariant event
\begin{equation}
 B_e(\kappa):=\{ {\mathfrak F}_e\ge\kappa\}               \label{eq:s6v86-Be}
\end{equation}
contains an open neighbourhood of the V83 frustrated exterior.
\end{lemma}

\begin{proof}
The identity \(H_e=-3I\) is
\eqref{eq:s6v83-H-minus}.  By
\Cref{lem:s6v83-trace-min},
\[
 \max_{V\in\SU(3)}\Re\Tr(-3V)=9/2.
\]
Insert this into \eqref{eq:s6v86-F}.  The maximum over the compact
group depends continuously on \(H_e\), so
\({\mathfrak F}_e\) is continuous in the exterior staples.  The final
claim follows from the strict inequality \(\kappa<27/2\).
\end{proof}

\begin{corollary}[A frustrated star cannot be healed at its centre]
\label{cor:s6v86-central-no-repair}
On \(B_e(\kappa)\),
\begin{equation}
 \sum_{p\ni e}d(U_p)\ge\kappa                            \label{eq:s6v86-star-cost}
\end{equation}
for every \(U_e\).  In particular, changing \(U_e\) alone cannot remove
the Wilson cost carried by this exterior.
\end{corollary}

\begin{proof}
This is \eqref{eq:s6v86-F-energy}.
\end{proof}

\begin{firewall}[Why a one-link finite-energy proof is impossible]
\label{fire:s6v86-one-link}
Let \({\cal F}_{e^c}\) be the sigma-field generated by every link
except \(U_e\).  Since \(B_e(\kappa)\in{\cal F}_{e^c}\),
\[
 {\mathbb P}\bigl(B_e(\kappa)\mid{\cal F}_{e^c}\bigr)
 =1_{B_e(\kappa)}
\]
under every measure.  Thus no estimate
\[
 {\mathbb P}\bigl(B_e(\kappa)\mid{\cal F}_{e^c}\bigr)
 \le p_s<1
\]
can hold almost surely when the event has positive probability.
Any conditional repair must resample a collar of exterior links, and
an arbitrary outside boundary can frustrate that collar in turn.  A
Peierls, multiscale, or reflection-positive dissemination argument is
therefore required.
\end{firewall}

\subsection{Incidence converts disseminated frustration into action}
\label{sec:s6v86-incidence}

For a finite set \(T\) of selected edges define
\[
 m(T):=\max_p\#\{e\in T:p\ni e\}.
\]
On a hypercubic lattice \(m(T)\le4\), since a plaquette has four
boundary edges.

\begin{lemma}[Star-to-plaquette incidence inequality]
\label{lem:s6v86-incidence}
For every link configuration,
\begin{equation}
 \sum_{e\in T}{\mathfrak F}_e
 \le m(T)\sum_p d(U_p).                                  \label{eq:s6v86-incidence}
\end{equation}
Consequently,
\begin{equation}
 \bigcap_{e\in T}B_e(\kappa)
 \subset
 \left\{\sum_p d(U_p)\ge{\kappa|T|\over m(T)}\right\}.    \label{eq:s6v86-extensive}
\end{equation}
\end{lemma}

\begin{proof}
Sum \eqref{eq:s6v86-F-energy} over \(e\in T\).  A given plaquette
defect is counted once for each selected edge on its boundary, hence at
most \(m(T)\) times.  On the intersection in
\eqref{eq:s6v86-extensive}, the left side of
\eqref{eq:s6v86-incidence} is at least \(\kappa|T|\).
\end{proof}

\begin{remark}[No independence was used]
\label{rem:s6v86-no-independence}
The inequality is pointwise.  It accounts for every plaquette touched
by every star, including overlaps, and makes no probabilistic
independence or conditional-boundary assumption.
\end{remark}

\subsection{A volume-uniform lower bound on the Wilson partition
function}
\label{sec:s6v86-partition}

Let \(\Lambda\) be a finite periodic lattice with \(L_\Lambda\)
unfixed links and \(P_\Lambda\) plaquettes, and let product Haar
measure be normalized to one.  Put
\begin{equation}
 A_\Lambda(U):=\sum_{p\subset\Lambda}d(U_p),\qquad
 Z_{\Lambda,s}:=\int e^{-sA_\Lambda(U)}\,\dd U.           \label{eq:s6v86-Z}
\end{equation}

\begin{lemma}[Near-flat tube bound]
\label{lem:s6v86-tube}
There are group constants \(c_0,C_0>0\) and \(s_0\) such that, for
every finite periodic lattice and \(s\ge s_0\),
\begin{equation}
 Z_{\Lambda,s}
 \ge
 \bigl(c_0s^{-4}\bigr)^{L_\Lambda}
 e^{-C_0P_\Lambda}.                                     \label{eq:s6v86-Zlower}
\end{equation}
\end{lemma}

\begin{proof}
Fix a bi-invariant Riemannian metric on \(\SU(3)\).  Since
\(\dim\SU(3)=8\), there are \(r_0,c_0>0\) such that normalized Haar
measure satisfies
\[
 |B(I,r)|\ge c_0r^8,\qquad 0<r\le r_0.
\]
The plaquette defect is smooth, vanishes at the identity, and has
vanishing first derivative there.  On a sufficiently small product
chart there is \(C_0\) such that
\[
 d(U_1U_2U_3^{-1}U_4^{-1})
 \le C_0\sum_{j=1}^4\operatorname{dist}(U_j,I)^2.
\]
Restrict every link to \(B(I,s^{-1/2})\), increasing \(s_0\) so this
ball lies in the chart.  The product Haar volume of that tube is at
least \((c_0s^{-4})^{L_\Lambda}\), while
\[
 sA_\Lambda(U)\le4C_0P_\Lambda
\]
there.  Absorb the harmless factor \(4\) into \(C_0\).
\end{proof}

\begin{corollary}[Block-normalized free-energy cost]
\label{cor:s6v86-free-energy}
Suppose a periodic block tiling has \(N\) tiles and
\[
 L_\Lambda\le\ell N,\qquad P_\Lambda\le qN              \label{eq:s6v86-count-ratios}
\]
for fixed \(\ell,q\).  Then
\begin{equation}
 Z_{\Lambda,s}\ge
 \exp\{-N(C_1\log s+C_2)\},                              \label{eq:s6v86-Zblock}
\end{equation}
where \(C_1=4\ell\) and \(C_2\) depends only on
\(\ell,q,\SU(3)\).
\end{corollary}

\begin{proof}
Take logarithms in \eqref{eq:s6v86-Zlower} and use
\eqref{eq:s6v86-count-ratios}.
\end{proof}

\subsection{Reflection positivity and the chessboard step}
\label{sec:s6v86-chessboard}

\begin{lemma}[Positive character expansion of the Wilson weight]
\label{lem:s6v86-character-positive}
For \(s\ge0\), the central function
\[
 w_s(U):=\exp\{s\Re\Tr U\}
\]
has an irreducible-character expansion
\begin{equation}
 w_s(U)=\sum_{\pi\in\widehat{\SU(3)}}a_\pi(s)\chi_\pi(U),
 \qquad a_\pi(s)\ge0.                                   \label{eq:s6v86-positive-expansion}
\end{equation}
Consequently the periodic Wilson lattice measure is reflection
positive across every coordinate reflection plane compatible with the
lattice.
\end{lemma}

\begin{proof}
Since \(\Tr U=\chi_{\mathbf3}(U)\) and
\(\overline{\Tr U}=\chi_{\overline{\mathbf3}}(U)\),
\[
 w_s
 =\sum_{n,m\ge0}
 { (s/2)^{n+m}\over n!\,m!}
 \chi_{\mathbf3}^{\,n}\chi_{\overline{\mathbf3}}^{\,m}.
\]
Every product of characters decomposes into irreducible characters
with nonnegative integer tensor-product multiplicities.  Absolute
uniform convergence of the exponential series permits regrouping and
proves \eqref{eq:s6v86-positive-expansion}.

For a plaquette crossing a reflection plane, apply Peter--Weyl to each
character in \eqref{eq:s6v86-positive-expansion}.  The crossing weight
is a sum
\[
 \sum_{\pi,i,j}c_{\pi}\,
 F_{\pi,ij}\,\Theta F_{\pi,ij},
 \qquad c_\pi\ge0,
\]
where \(\Theta\) is reflected complex conjugation and the factors on
the two sides are matrix coefficients of the two half-plaquette
holonomies.  Products of such sums retain this form.  Integrating
links lying on the reflection plane with Haar measure therefore gives
\[
 \int F\,\Theta F\,\dd\mu_{\Lambda,s}\ge0
\]
for every observable \(F\) supported on one side.  This is reflection
positivity.
\end{proof}

\begin{theorem}[Chessboard domination from one disseminated event]
\label{thm:s6v86-chessboard}
Let a periodic lattice be tiled by \(N\) congruent reflection blocks,
and let a reflection-positive probability measure be invariant under
the generated block reflections.  If \(B\) is an event supported in
one block, denote its reflected translate in tile \(t\) by \(B_t\) and
its full dissemination by
\[
 B^\#:=\bigcap_{t=1}^N B_t.
\]
Then, for every set \(F\) of tiles,
\begin{equation}
 \mathbb P\left(\bigcap_{t\in F}B_t\right)
 \le\mathbb P(B^\#)^{|F|/N}.                             \label{eq:s6v86-chessboard}
\end{equation}
\end{theorem}

\begin{proof}
Write the left side as the expectation of the product of the
corresponding indicator functions, inserting the constant function
one in unoccupied tiles.  Reflect successively across the midplanes of
the block tiling.  At each reflection, reflection positivity and
Cauchy--Schwarz replace the two half-patterns by their two symmetric
duplicates and take a square root.  After all coordinate reflections,
each originally occupied tile has been disseminated to every tile.
Multiplying the accumulated powers gives exponent \(1/N\) for each of
the \(|F|\) original indicators, hence
\eqref{eq:s6v86-chessboard}.  Indicators follow from bounded
nonnegative block observables by monotone approximation if a boundary
of the event has positive measure.
\end{proof}

\begin{remark}[Scope of the chessboard theorem]
\label{rem:s6v86-chessboard-scope}
The result requires a periodic reflection tiling.  It is not an
assertion for arbitrary boundary conditions.  This is precisely why
it avoids the false worst-exterior conditional estimate ruled out in
\Cref{fire:s6v86-one-link}.
\end{remark}

\subsection{Pure-Wilson finite-set domination}
\label{sec:s6v86-Wilson-domination}

Choose a reflection block large enough that the six staples defining a
distinguished \(B_e(\kappa)\) lie in its interior.  Reflected copies
choose one distinguished edge in every tile.  Let \(m_\star\le4\) be
the maximum number of those selected edges incident on one plaquette.

\begin{theorem}[Wilson star-frustration activity]
\label{thm:s6v86-Wilson-activity}
For the pure Wilson measure
\[
 \mu_{\Lambda,s}(\dd U)
 =Z_{\Lambda,s}^{-1}e^{-sA_\Lambda(U)}\,\dd U,
\]
there are constants \(C_1,C_2\), independent of the periodic volume,
such that
\begin{equation}
 \mu_{\Lambda,s}\bigl(B(\kappa)^\#\bigr)
 \le
 \exp\left\{
 -N\left({\kappa s\over m_\star}
          -C_1\log s-C_2\right)
 \right\}.                                               \label{eq:s6v86-disseminated}
\end{equation}
Therefore, for every set \(F\) of tiles,
\begin{equation}
 \mu_{\Lambda,s}\left(\bigcap_{t\in F}B_t(\kappa)\right)
 \le p_{\rm W}(s,\kappa)^{|F|},                          \label{eq:s6v86-W-domination}
\end{equation}
where
\begin{equation}
 p_{\rm W}(s,\kappa)
 :=
 \min\left\{1,
 \exp\left[-{\kappa s\over m_\star}
            +C_1\log s+C_2\right]\right\}.               \label{eq:s6v86-pW}
\end{equation}
In particular
\[
 p_{\rm W}(s,\kappa)
 =\exp\{-\kappa s/m_\star+O(\log s)\}.
\]
\end{theorem}

\begin{proof}
On the disseminated event, apply
\Cref{lem:s6v86-incidence} to the \(N\) distinguished edges:
\[
 A_\Lambda(U)\ge{\kappa N\over m_\star}.
\]
Since product Haar measure has total mass one,
\[
 \int_{B(\kappa)^\#}e^{-sA_\Lambda(U)}\,\dd U
 \le e^{-\kappa sN/m_\star}.
\]
Divide by the partition-function lower bound
\eqref{eq:s6v86-Zblock}; this proves
\eqref{eq:s6v86-disseminated}.  Reflection positivity is
\Cref{lem:s6v86-character-positive}.  Insert
\eqref{eq:s6v86-disseminated} in
\eqref{eq:s6v86-chessboard} and cap the resulting activity at one.
\end{proof}

\begin{corollary}[Finite colouring gives domination for all star
blocks]
\label{cor:s6v86-colouring}
Suppose the full family of oriented star blocks can be partitioned into
\(\chi<\infty\) reflection-tiling classes of the preceding kind, with
the same upper constants and \(m_\star\).  Then for every finite family
\(F\) of star blocks,
\begin{equation}
 \mu_{\Lambda,s}\left(\bigcap_{e\in F}B_e(\kappa)\right)
 \le p_{\rm all}(s,\kappa)^{|F|},                        \label{eq:s6v86-all-domination}
\end{equation}
where
\begin{equation}
 p_{\rm all}(s,\kappa)
 :=p_{\rm W}(s,\kappa)^{1/\chi}
 =\exp\left\{-{\kappa s\over\chi m_\star}
              +O(\log s)\right\}                        \label{eq:s6v86-pall}
\end{equation}
for large \(s\).
\end{corollary}

\begin{proof}
One colour class contains a subset \(F_0\subset F\) with
\(|F_0|\ge|F|/\chi\).  Since an intersection only becomes larger when
conditions are discarded,
\[
 \mu_{\Lambda,s}\left(\bigcap_{e\in F}B_e\right)
 \le
 \mu_{\Lambda,s}\left(\bigcap_{e\in F_0}B_e\right)
 \le p_{\rm W}^{|F_0|}
 \le p_{\rm W}^{|F|/\chi}.
\]
\end{proof}

\begin{corollary}[Pure-Wilson cluster subcriticality]
\label{cor:s6v86-W-clusters}
For every fixed block adjacency degree \(D\) and fixed
\(\kappa>0\),
\[
 eD\,p_{\rm all}(s,\kappa)<1
\]
for all sufficiently large \(s\).  Thus the bad-cluster tail
\eqref{eq:s6v85-cluster-geometric} holds for the pure Wilson
star-frustration indicators, uniformly in periodic volume.
Moreover, every cluster capacity exponent \(a_s=o(s)\) satisfies the
V85 entropy margin
\[
 e^{a_s+1}D\,p_{\rm all}(s,\kappa)<1
\]
for sufficiently large \(s\).
\end{corollary}

\begin{proof}
Use the exponential asymptotic
\eqref{eq:s6v86-pall} in
\Cref{thm:s6v85-cluster-tail,cor:s6v85-weak-cluster}.
\end{proof}

\subsection{The complete-action extension and exact upstream card}
\label{sec:s6v86-complete}

The physical ground density need not equal the pure Wilson density.
The following theorem records exactly what is sufficient to retain the
argument without incorrectly treating a generator-order return as a
small perturbation.

\begin{definition}[Complete-action reflection/stability card, CARSC]
\label{def:s6v86-CARSC}
On a periodic \(N\)-tile lattice, write the exact normalized physical
ground measure, before or after an exact gauge quotient, as
\[
 \mu^{\rm ph}_{\Lambda,s}(\dd U)
 ={1\over Z^{\rm ph}_{\Lambda,s}}
 e^{-S^{\rm ph}_{\Lambda,s}(U)}\,\rho_\Lambda(\dd U).
\]
CARSC at the star threshold \(\kappa\) consists of:
\begin{enumerate}[label=\textup{(C\arabic*)},leftmargin=4em]
\item \(\rho_\Lambda\) is normalized and the complete measure is
      reflection positive for the chosen block reflections;
\item there is a scalar \(c_{\Lambda,s}\), allowed to be extensive,
      and numbers \(r_s,b_s\ge0\), independent of volume, such that on
      the disseminated bad event
      \begin{equation}
       S^{\rm ph}_{\Lambda,s}(U)
       \ge {s\kappa N\over m_\star}
           +c_{\Lambda,s}-r_sN;                          \label{eq:s6v86-bad-stability}
      \end{equation}
\item there is a measurable near-flat tube \({\cal T}_{\Lambda,s}\)
      with
      \begin{align}
       \rho_\Lambda({\cal T}_{\Lambda,s})
       &\ge e^{-N(C_1\log s+C_2)},                        \label{eq:s6v86-tube-volume}\\
       S^{\rm ph}_{\Lambda,s}(U)
       &\le c_{\Lambda,s}+b_sN
       \quad\text{on }{\cal T}_{\Lambda,s};               \label{eq:s6v86-tube-action}
      \end{align}
\item \(r_s+b_s=o(s)\).
\end{enumerate}
\end{definition}

\begin{theorem}[CARSC gives complete-action domination]
\label{thm:s6v86-CARSC}
Under CARSC,
\begin{equation}
 \mu^{\rm ph}_{\Lambda,s}\bigl(B(\kappa)^\#\bigr)
 \le
 \exp\left\{-N\left(
 {\kappa s\over m_\star}
 -r_s-b_s-C_1\log s-C_2\right)\right\}.                  \label{eq:s6v86-complete-disseminated}
\end{equation}
Hence the chessboard and finite-colouring arguments give
\begin{equation}
 \mu^{\rm ph}_{\Lambda,s}
 \left(\bigcap_{e\in F}B_e(\kappa)\right)
 \le p_{\rm ph}(s,\kappa)^{|F|},                         \label{eq:s6v86-complete-domination}
\end{equation}
with
\begin{equation}
 p_{\rm ph}(s,\kappa)
 \le
 \exp\left\{-{1\over\chi}\left(
 {\kappa s\over m_\star}
 -r_s-b_s-C_1\log s-C_2\right)\right\}                  \label{eq:s6v86-pph}
\end{equation}
for all sufficiently large \(s\).  In particular SFCC row (S2) holds
for this family of bad events with
\[
 p_{\rm ph}(s,\kappa)=e^{-c_\kappa s+o(s)},
 \qquad c_\kappa={\kappa\over\chi m_\star}>0.
\]
\end{theorem}

\begin{proof}
The bad-set numerator is at most
\[
 \exp\left\{-{s\kappa N\over m_\star}
              -c_{\Lambda,s}+r_sN\right\}.
\]
By integrating only over the tube,
\[
 Z^{\rm ph}_{\Lambda,s}
 \ge
 \exp\{-c_{\Lambda,s}
        -N(b_s+C_1\log s+C_2)\}.
\]
Division cancels the arbitrary extensive scalar
\(c_{\Lambda,s}\) exactly and proves
\eqref{eq:s6v86-complete-disseminated}.  CARSC row (C1) permits
\Cref{thm:s6v86-chessboard}; then repeat
\Cref{cor:s6v86-colouring}.  Row (C4) gives the stated cofinal
asymptotic.
\end{proof}

\begin{remark}[Extensive vacuum energies are harmless]
\label{rem:s6v86-scalar-cancels}
The explicit \(c_{\Lambda,s}\) is essential.  Bounding it separately
would manufacture a volume loss.  Its cancellation is the static
measure analogue of
\Cref{prop:s6v81-scalar-shift}.
\end{remark}

\begin{lemma}[A checkable locality bound for generated terms]
\label{lem:s6v86-locality}
Suppose, on the unreduced link space,
\[
 S^{\rm ph}_{\Lambda,s}=sA_\Lambda+c_{\Lambda,s}+\Psi_{\Lambda,s},
 \qquad
 \Psi_{\Lambda,s}=\sum_X\psi_{X,s},
 \]
where \(X\) runs over finite link sets.  If
\begin{equation}
 \sup_{e}
 \sum_{X:X\cap{\cal N}(e)\ne\varnothing}
 2\|\psi_{X,s}\|_\infty
 \le\omega_s                                             \label{eq:s6v86-local-osc}
\end{equation}
for a fixed star collar \({\cal N}(e)\), then changing all variables
in that collar changes \(\Psi_{\Lambda,s}\) by at most \(\omega_s\).
If, in addition, every term is assigned to at most \(q_{\rm inc}\)
collars, then on a union of \(n\) collars its total oscillation is at
most \(q_{\rm inc}\omega_sn\).
\end{lemma}

\begin{proof}
Terms disjoint from the changed variables cancel.  Each remaining
term changes by at most twice its uniform norm.  Summing gives the
first claim.  The assignment multiplicity gives the second.
\end{proof}

\begin{firewall}[Local oscillation does not prove reflection
positivity]
\label{fire:s6v86-locality-RP}
The bound \eqref{eq:s6v86-local-osc} can supply the
\(o(s)\) relative-stability debit in CARSC if
\(\omega_s=o(s)\), but it does not imply row (C1).  Reflection
positivity requires a positive half-space factorization of every term
crossing a reflection plane.  A nonlocal or signed
response-matching cofactor can satisfy a norm bound while destroying
that factorization.
\end{firewall}

\begin{proposition}[Positive-character sufficient condition for the
complete action]
\label{prop:s6v86-complete-RP}
Suppose every complete-action factor lying across a block reflection
plane is a central kernel
\[
 K(UV^{-1})=\sum_\pi k_\pi\chi_\pi(UV^{-1}),
 \qquad k_\pi\ge0,
\]
or a product and locally uniform limit of such kernels, while all
remaining factors occur in reflected pairs.  Then CARSC row (C1)
holds.
\end{proposition}

\begin{proof}
Peter--Weyl writes each crossing kernel as a nonnegative sum of
matrix-coefficient squares across the plane, exactly as in the second
part of \Cref{lem:s6v86-character-positive}.  Products, reflected
pairs, positive sums, integration over plane variables, and locally
uniform limits preserve the reflection-positive cone.
\end{proof}

\subsection{What has and has not been closed}
\label{sec:s6v86-frontier}

\begin{theorem}[V86 domination decision]
\label{thm:s6v86-decision}
For the periodic pure Wilson \(\SU(3)\) measure, the V83
center-frustrated star neighbourhoods satisfy volume-uniform
finite-set domination with activity
\[
 p_s=e^{-c_\kappa s+O(\log s)}.
\]
Thus the geometric full-lattice probability row left open in V84 and
V85 is closed for the Wilson part without constructing a
finite-multiplicity local repair.

For the exact physical ground measure, the same conclusion follows
from CARSC.  The first unresolved complete-action input is now
specific: prove that all response-matched and generated terms preserve
the required reflection-positive factorization and have
sublinear-in-\(s\) relative stability between the disseminated bad
sector and a near-flat tube.  If reflection positivity fails, one must
replace the chessboard step by a polymer/repair comparison for those
terms.
\end{theorem}

\begin{proof}
The Wilson assertion is
\Cref{thm:s6v86-Wilson-activity,cor:s6v86-colouring}.  The physical
assertion and its exact hypotheses are \Cref{thm:s6v86-CARSC}.
\Cref{fire:s6v86-locality-RP} separates the two required properties
and prevents either from being inferred from the other.
\end{proof}

\begin{assessment}[Upstream bottleneck after V86]
\label{ass:s6v86-frontier}
The probability half of SFCC is no longer an unspecified Peierls-map
problem for the Wilson action.  It is a proved reflection-positive
chessboard estimate.  Two genuinely analytic tasks remain before the
mass-gap compiler can use it:
\begin{enumerate}[label=\textup{(\roman*)},leftmargin=3.5em]
\item derive the exact finite-regulator physical ground density from
      the response-matched half-amplitude and test CARSC term by term;
\item prove a connected-cluster form-capacity estimate, not merely a
      probability estimate.
\end{enumerate}
The second item remains the \(L^2\) obstruction identified in V84.
The first must be settled upstream before assigning the Wilson
domination estimate to the physical measure.
\end{assessment}

\begin{programme}[V87: response-matched reflection-positivity audit]
\label{prog:s6v86-V87}
\begin{enumerate}[leftmargin=3em]
\item Return to the exact half-amplitude and list every factor in the
      physical ground density after shared-link integration.
\item For each cross-plane factor, compute its character or
      matrix-coefficient kernel and test coefficient positivity.
\item Separate extensive scalar normalization from the relative
      bad-sector and near-flat-tube bounds in CARSC.
\item Prove \(r_s+b_s=o(s)\) for local positive factors, or identify
      the first generated cofactor for which it fails.
\item If CARSC closes, combine
      \Cref{thm:s6v86-CARSC} with the V85 cluster sum and attack the
      connected-cluster capacity row.  If it fails, construct the
      corresponding signed polymer comparison rather than assuming
      reflection positivity.
\end{enumerate}
\end{programme}

\begin{firewall}[Claim boundary after sixth-series V86]
\label{fire:s6v86-boundary}
V86 defines a gauge-invariant star-frustration functional; proves its
exact value \(27/2\) on the V83 center exterior; proves the
star-to-plaquette incidence bound; proves a volume-uniform near-flat
partition lower bound; proves positive character coefficients and
reflection positivity for the Wilson factor; derives the chessboard
finite-set domination
\[
 p_s=e^{-c_\kappa s+O(\log s)}
\]
for periodic pure-Wilson star blocks; and gives the exact CARSC
extension theorem for a complete physical ground action.

V86 does not prove CARSC for the response-matched physical measure,
does not show that its chosen star event contains every possible
bad conditional geometry, does not prove the connected-cluster
capacity estimate or approximate tensorization, does not populate
GMPC or HSOCC, does not construct the continuum OS theory, and does
not prove the full Yang--Mills mass gap.
\end{firewall}

\begin{theorem}[Sixth-series V86 result ledger]
\label{thm:s6v86-results}
The positive conclusions in \Cref{fire:s6v86-boundary} are exactly
those of
\Cref{lem:s6v86-F-basic,
lem:s6v86-center-F,
cor:s6v86-central-no-repair,
lem:s6v86-incidence,
lem:s6v86-tube,
cor:s6v86-free-energy,
lem:s6v86-character-positive,
thm:s6v86-chessboard,
thm:s6v86-Wilson-activity,
cor:s6v86-colouring,
cor:s6v86-W-clusters,
thm:s6v86-CARSC,
lem:s6v86-locality,
prop:s6v86-complete-RP,
thm:s6v86-decision}.
\end{theorem}

\begin{proof}
Each assertion is the cited result with its periodic-volume,
reflection-tiling, pure-Wilson, complete-action, relative-stability,
and form-capacity qualifications unchanged.
\end{proof}

\paragraph{Parent-version record.}
V86 was formed from
\texttt{Renewal\_Yang\_Mills\_Mass\_Gap\_Research\_Notes\_V85.tex}.
The V85 parent had \(38\,503\) logical source lines,
\(1\,349\,086\) bytes, and SHA--256
\[
 \texttt{1c90611e3407ebd0ecddf58bb213497158189cccedc9665b958239a1174870f0}.
\]
The next compulsory companion audit is V90.

% =============================================================================
% END APPENDED SIXTH-SERIES V86 MATERIAL
% =============================================================================

% =============================================================================
% BEGIN APPENDED SIXTH-SERIES V87 MATERIAL
% =============================================================================

\section{V87: singular-ground ownership and the spatial
reflection-positive bridge}
\label{sec:s6v87}

\subsection{The measure mismatch that must be corrected}
\label{sec:s6v87-mismatch}

V86 proves a chessboard domination theorem for a periodic Wilson Gibbs
measure.  The gap compiler of V81--V82, however, is written for the
physical boundary ground measure.  These measures are not
automatically the same.

Retain the exact half-slab notation of V58.  A positive amplitude
\[
 (Af)(y)=\int F(y,x)f(x)\,\mu_0(\dd x)
\]
has normalized leading singular pair
\[
 Au_0=s_1v_0,\qquad A^*v_0=s_1u_0,
\]
and the physical initial-boundary law is
\begin{equation}
 \pi_0(\dd x)=u_0(x)^2\mu_0(\dd x).                      \label{eq:s6v87-pi0}
\end{equation}

\begin{proposition}[Exact one-link ground correction]
\label{prop:s6v87-ground-conditional}
Assume that, conditional on the exterior \(\xi\), the reference law
\(\mu_0\) has the matrix-Fisher conditional
\[
 \mu_0(\dd U\mid\xi)
 ={1\over Z_{\tau,H(\xi)}}
 e^{\tau\Re\Tr(UH(\xi))}\,\dd U.
\]
Then the physical conditional is
\begin{equation}
 \boxed{
 \pi_0(\dd U\mid\xi)
 =
 {u_0(U,\xi)^2e^{\tau\Re\Tr(UH(\xi))}\,\dd U
  \over
  \int_{\SU(3)}
  u_0(V,\xi)^2e^{\tau\Re\Tr(VH(\xi))}\,\dd V}.}
                                                               \label{eq:s6v87-physical-conditional}
\end{equation}
In particular, the matrix-Fisher law
\eqref{eq:s6v83-matrix-Fisher} is the physical conditional only when
the conditional ground factor is constant in \(U\), or when it has
already been included in the declared potential.
\end{proposition}

\begin{proof}
Disintegrate \eqref{eq:s6v87-pi0} with respect to the exterior
sigma-field.  The exterior-only factors cancel between numerator and
denominator.  The remaining density is
\eqref{eq:s6v87-physical-conditional}.  This is the one-boundary
version of \Cref{thm:s6v58-action-owner}.
\end{proof}

\begin{assessment}[Scope correction for V83--V86]
\label{ass:s6v87-scope-correction}
The V83 double-well theorem and V86 Wilson chessboard theorem remain
valid for their stated matrix-Fisher and periodic Wilson measures.
They do not, without an estimate on \(u_0\), prove the same conditional
gap obstruction or bad-star activity for the physical ground measure
\eqref{eq:s6v87-pi0}.  This is an ownership correction, not a
retraction of either finite-measure theorem.
\end{assessment}

\begin{lemma}[A ground factor can erase the bare barrier]
\label{lem:s6v87-barrier-erasure}
Let
\[
 \nu_\tau(\dd U)=Z_\tau^{-1}e^{\tau\Phi(U)}\,\dd U
\]
be any positive compact-group Gibbs law.  If
\[
 h_\tau(U)=e^{-\tau\Phi(U)/2},
\]
then the ground-weighted probability proportional to
\(h_\tau^2\nu_\tau\) is normalized Haar measure.  Thus a bare
double-well potential alone does not determine the Poincar\'e gap of a
ground-weighted law.
\end{lemma}

\begin{proof}
The two factors cancel pointwise:
\[
 h_\tau(U)^2e^{\tau\Phi(U)}=1.
\]
\end{proof}

\begin{firewall}[Positivity of the singular function is not a
regularity estimate]
\label{fire:s6v87-positive-not-flat}
Finite-regulator positivity improvement gives \(u_0>0\), but it gives
no volume- or coupling-uniform bound on
\[
 \operatorname*{osc}_{U}\log u_0(U,\xi).
\]
The example in \Cref{lem:s6v87-barrier-erasure} has a strictly positive
factor whose logarithmic oscillation is of order \(\tau\), exactly the
order needed to change the well geometry.
\end{firewall}

\subsection{When the two-well obstruction survives the ground factor}
\label{sec:s6v87-robust-well}

Let \(G\) be a compact eight-dimensional Riemannian manifold with
volume measure \(\dd U\).  Let \(\Phi\) have two symmetry-related,
nondegenerate maxima \(U_\pm\) of value \(M\).  Choose disjoint
neighbourhoods \({\cal N}_\pm\), exchanged by the symmetry, so that
the closed complement
\({\cal K}:=G\setminus({\cal N}_+^\circ\cup{\cal N}_-^\circ)\)
contains the support of the gradient of a smooth odd separator \(f\)
and satisfies
\begin{equation}
 \sup_{\cal K}\Phi\le M-\Delta,\qquad \Delta>0.           \label{eq:s6v87-barrier}
\end{equation}
Let \(\iota\) be the measure-preserving involution exchanging the two
wells and satisfying \(\Phi\circ\iota=\Phi\).

\begin{theorem}[Ground-profile-stable two-well upper bound]
\label{thm:s6v87-robust-well}
Consider
\begin{equation}
 \nu_\tau^g(\dd U)
 ={1\over Z_\tau^g}g_\tau(U)e^{\tau\Phi(U)}\,\dd U,
 \qquad g_\tau>0.                                       \label{eq:s6v87-weighted-well}
\end{equation}
Assume \(g_\tau\circ\iota=g_\tau\).  Suppose there are scales
\(G_\tau>0\) and numbers \(a_\tau,b_\tau\ge0\) such that:
\begin{align}
 \inf_{B(U_\pm,\tau^{-1/2})}g_\tau
 &\ge G_\tau e^{-a_\tau\tau},                           \label{eq:s6v87-well-lower}\\
 \sup_{\cal K}g_\tau
 &\le G_\tau e^{b_\tau\tau},                            \label{eq:s6v87-barrier-upper}\\
 a_\tau+b_\tau&\le\Delta-\varepsilon                    \label{eq:s6v87-profile-margin}
\end{align}
for some fixed \(\varepsilon>0\) and all large \(\tau\).  Then the
optimal Poincar\'e number satisfies
\begin{equation}
 \boxed{
 \lambda(\nu_\tau^g)
 \le C\tau^4e^{-\varepsilon\tau}}                       \label{eq:s6v87-weighted-gap}
\end{equation}
for all large \(\tau\).
\end{theorem}

\begin{proof}
Choose the separator \(f\) from the V83 proof so that
\[
 f\circ\iota=-f,\qquad
 f=1\text{ near }U_+,\qquad
 f=-1\text{ near }U_-,
\]
and \(\nabla f\) is supported in \({\cal K}\).  Symmetry gives
\(\int f\,\dd\nu_\tau^g=0\).

The nondegenerate maximum and the eight-dimensional ball volume give
\[
 Z_\tau^g
 \ge c\tau^{-4}G_\tau
     e^{\tau M-a_\tau\tau-C}.
\]
On the gradient support,
\eqref{eq:s6v87-barrier} and
\eqref{eq:s6v87-barrier-upper} give
\[
 \int|\nabla f|^2g_\tau e^{\tau\Phi}\,\dd U
 \le C G_\tau e^{\tau(M-\Delta+b_\tau)}.
\]
After normalization,
\[
 \int|\nabla f|^2\,\dd\nu_\tau^g
 \le C\tau^4
 e^{-\tau(\Delta-a_\tau-b_\tau)}
 \le C\tau^4e^{-\varepsilon\tau}.                       \label{eq:s6v87-test-energy}
\]
The same estimates show that the barrier mass is exponentially small
relative to the two well balls.  Symmetry gives equal mass to the two
balls, and \(f^2=1\) there, so
\(\operatorname{Var}_{\nu_\tau^g}(f)\ge1/2\) for large \(\tau\).
Its Rayleigh quotient proves \eqref{eq:s6v87-weighted-gap}.
\end{proof}

\begin{corollary}[Sublinear ground logarithms cannot heal the V83
wells]
\label{cor:s6v87-sublinear-ground}
In the V83 exterior \(H=-hI\), put
\[
 g_\tau(U)=u_0(U,\xi)^2.
\]
If complex conjugation preserves \(g_\tau\) and if, after multiplication
by an arbitrary scalar,
\[
 \sup_{{\cal N}_+\cup{\cal N}_-\cup{\cal K}}
 |\log g_\tau|=o(\tau),                                 \label{eq:s6v87-sublinear-log}
\]
then the physical conditional Poincar\'e number still tends to zero
exponentially.
\end{corollary}

\begin{proof}
Choose \(G_\tau=1\) after the permitted scalar normalization.
Equation \eqref{eq:s6v87-sublinear-log} gives
\(a_\tau,b_\tau=o(1)\).  Apply
\Cref{thm:s6v87-robust-well} with any fixed
\(\varepsilon<\Delta\) for all sufficiently large \(\tau\).
\end{proof}

\begin{corollary}[Ground-healing dichotomy]
\label{cor:s6v87-healing-dichotomy}
To avoid the V83 exponential conditional-gap collapse at a
symmetry-preserving center exterior, the physical singular factor must
violate the sub-barrier profile hypotheses of
\Cref{thm:s6v87-robust-well}.  In particular, at least one of the
following must occur along a subsequence:
\begin{enumerate}[label=\textup{(\roman*)},leftmargin=3.5em]
\item the ground law breaks the involutive balance between the wells;
\item the logarithmic ground contrast between a well and the barrier
      is of order \(\tau\);
\item the physical conditional is not owned by the declared
      matrix-Fisher reference.
\end{enumerate}
\end{corollary}

\begin{proof}
This is the contrapositive of
\Cref{thm:s6v87-robust-well,prop:s6v87-ground-conditional}.
\end{proof}

\subsection{Time-transfer positivity does not imply spatial
reflection positivity}
\label{sec:s6v87-two-positivities}

\begin{proposition}[Finite counterexample separating the two notions]
\label{prop:s6v87-positivity-counterexample}
There is a finite positivity-improving, positive self-adjoint transfer
operator which commutes with a spatial reflection, but whose unique
ground measure is not spatially reflection positive.
\end{proposition}

\begin{proof}
Let one spatial half have two states and identify a full boundary
configuration with \((i,j)\in\{1,2\}^2\).  Spatial reflection exchanges
\(i\) and \(j\).  Define the strictly positive symmetric ground vector
by the matrix
\[
 h=
 \begin{pmatrix}
 1&2\\
 2&1
 \end{pmatrix}.
\]
On the four-dimensional boundary Hilbert space let
\[
 T={|h\rangle\langle h|\over\langle h,h\rangle}.
\]
Its kernel is strictly positive, \(T\) is positive self-adjoint and
positivity improving, its leading eigenvalue is one and simple, and it
commutes with reflection because \(h(i,j)=h(j,i)\).

Relative to counting measure, the ground probability has reflection
matrix proportional to
\[
 \bigl(h(i,j)^2\bigr)_{i,j}
 =
 \begin{pmatrix}
 1&4\\
 4&1
 \end{pmatrix}.
\]
This matrix has eigenvalues \(5\) and \(-3\).  Taking a left-half
function in the negative eigendirection makes the spatial
reflection form negative.  Hence the ground law is not spatially
reflection positive.
\end{proof}

\begin{firewall}[HSOCC does not populate CARSC row (C1)]
\label{fire:s6v87-HSOCC-CARSC}
Positivity of \(A^*A\) is reflection positivity in the Euclidean
\emph{time} construction of the transfer.  CARSC row (C1) uses
reflection positivity across \emph{spatial block planes} of the
boundary law.  By
\Cref{prop:s6v87-positivity-counterexample}, the former, even with
spatial symmetry and a unique positive ground vector, does not imply
the latter.
\end{firewall}

\subsection{A valid route from spatially positive cylinders}
\label{sec:s6v87-cylinder}

Let \(\Theta\) be spatial reflected conjugation on boundary
observables.  A probability \(\pi\) is spatially reflection positive
if
\[
 \int F\,\Theta F\,\dd\pi\ge0                           \label{eq:s6v87-RP}
\]
for every bounded continuous \(F\) supported on the positive spatial
half.

\begin{lemma}[Reflection positivity is closed under weak limits]
\label{lem:s6v87-RP-closed}
On a compact boundary configuration space, if
\(\pi_n\Rightarrow\pi\) weakly and every \(\pi_n\) satisfies
\eqref{eq:s6v87-RP}, then \(\pi\) also satisfies
\eqref{eq:s6v87-RP}.
\end{lemma}

\begin{proof}
For each admissible bounded continuous \(F\), the function
\(F\Theta F\) is bounded and continuous.  Therefore
\[
 \int F\Theta F\,\dd\pi
 =\lim_{n\to\infty}\int F\Theta F\,\dd\pi_n\ge0.
\]
\end{proof}

\begin{theorem}[Ground law as a midpoint-cylinder limit]
\label{thm:s6v87-midpoint}
Let \(T\) be a compact positive self-adjoint operator on
\(L^2(\mu_0)\) with simple leading eigenvalue
\(\lambda_1>0\), normalized positive eigenfunction \(u_0\), and
\[
 \sup\bigl(\operatorname{spec}T\setminus\{\lambda_1\}\bigr)
 =\lambda_2<\lambda_1.
\]
For \(b\ge0\) with \(\langle b,u_0\rangle>0\), define the midpoint law
\begin{equation}
 \pi_n(\dd x)
 :={|T^nb(x)|^2\over\|T^nb\|_2^2}\,\mu_0(\dd x).         \label{eq:s6v87-midpoint-law}
\end{equation}
Then
\begin{equation}
 \|\pi_n-u_0^2\mu_0\|_{\rm TV}
 \le C_b\left({\lambda_2\over\lambda_1}\right)^n.        \label{eq:s6v87-midpoint-convergence}
\end{equation}
If every midpoint cylinder law \(\pi_n\) is spatially reflection
positive, then the physical ground law \(u_0^2\mu_0\) is spatially
reflection positive.
\end{theorem}

\begin{proof}
After division by its norm, the spectral theorem gives
\[
 {T^nb\over\|T^nb\|_2}\longrightarrow u_0
\]
in \(L^2(\mu_0)\) at the displayed geometric rate, after choosing the
positive phase.  For unit \(L^2\) functions \(a,c\),
\[
 \|a^2-c^2\|_{L^1}
 \le\|a-c\|_2\|a+c\|_2
 \le2\|a-c\|_2.
\]
This gives total-variation convergence, with an inessential convention
factor in its definition absorbed into \(C_b\).  Apply
\Cref{lem:s6v87-RP-closed}.
\end{proof}

\begin{corollary}[Spatially positive Euclidean construction populates
the ground row]
\label{cor:s6v87-Euclidean-RP}
Suppose the exact finite-cylinder action and boundary vector producing
\eqref{eq:s6v87-midpoint-law} have a positive half-space
factorization under every required spatial block reflection, uniformly
for all cylinder lengths.  Then CARSC row (C1) holds for the
finite-regulator ground measure.
\end{corollary}

\begin{proof}
The factorization proves reflection positivity of every cylinder
law by the same square-sum argument as
\Cref{lem:s6v86-character-positive}.  Apply
\Cref{thm:s6v87-midpoint}.
\end{proof}

\begin{remark}[Uniformity in cylinder length is substantive]
\label{rem:s6v87-cylinder-uniformity}
A spatially positive bare Wilson cylinder suffices for the bare Wilson
ground.  It does not prove the claim for a different
response-matched amplitude.  Every generated action, marginal return,
boundary Jacobian, and counterterm producing that amplitude must occur
inside the positive cylinder factorization or be shown not to alter
the normalized ground law.
\end{remark}

\subsection{Ground-profile transfer of the V86 activity}
\label{sec:s6v87-profile-transfer}

Let \(\mu_{\rm W}\) be a normalized periodic Wilson reference on a
boundary configuration space, and let
\[
 \pi_h(\dd x)
 ={h(x)^2\over\int h^2\,\dd\mu_{\rm W}}\,
 \mu_{\rm W}(\dd x).                                    \label{eq:s6v87-pih}
\]
The scalar normalization of \(h\) is deliberately left free.

\begin{theorem}[Ground-profile chessboard transfer]
\label{thm:s6v87-profile-transfer}
Let \(B^\#\) be the dissemination of \(N\) star-frustration events.
Assume
\begin{equation}
 \mu_{\rm W}(B^\#)\le e^{-a_sN}.                         \label{eq:s6v87-W-activity}
\end{equation}
Suppose there is a measurable comparison set \(G\), a scale \(m>0\),
and \(r_s,b_s\ge0\) such that
\begin{align}
 \mu_{\rm W}(G)&\ge e^{-b_sN},                           \label{eq:s6v87-G-mass}\\
 h^2&\ge m\quad\mu_{\rm W}\text{-a.e. on }G,             \label{eq:s6v87-h-lower}\\
 h^2&\le me^{r_sN}
       \quad\mu_{\rm W}\text{-a.e. on }B^\#.             \label{eq:s6v87-h-upper}
\end{align}
Then
\begin{equation}
 \boxed{
 \pi_h(B^\#)\le e^{-N(a_s-r_s-b_s)}.}                    \label{eq:s6v87-ground-activity}
\end{equation}
If \(\pi_h\) is spatially reflection positive, a \(\chi\)-colour
version of the V86 argument gives, for all finite \(F\),
\begin{equation}
 \pi_h\left(\bigcap_{e\in F}B_e\right)
 \le
 \exp\left\{-{|F|\over\chi}
 (a_s-r_s-b_s)\right\}.                                 \label{eq:s6v87-ground-domination}
\end{equation}
\end{theorem}

\begin{proof}
The numerator in \eqref{eq:s6v87-pih} obeys
\[
 \int_{B^\#}h^2\,\dd\mu_{\rm W}
 \le me^{r_sN}\mu_{\rm W}(B^\#)
 \le me^{-N(a_s-r_s)}.
\]
The denominator is at least
\[
 \int_Gh^2\,\dd\mu_{\rm W}
 \ge m\mu_{\rm W}(G)\ge me^{-b_sN}.
\]
Division proves \eqref{eq:s6v87-ground-activity}; the arbitrary scale
\(m\) cancels.  Apply
\Cref{thm:s6v86-chessboard,cor:s6v86-colouring} to
\(\pi_h\) for the final assertion.
\end{proof}

\begin{corollary}[Sublinear profile debit preserves exponential
domination]
\label{cor:s6v87-sublinear-profile}
For the V86 Wilson activity,
\[
 a_s={\kappa s\over m_\star}-O(\log s).
\]
If \(r_s+b_s=o(s)\) and the ground law is spatially reflection
positive, then
\[
 p_h(s,\kappa)
 \le\exp\left\{-{\kappa s\over\chi m_\star}+o(s)\right\}.
\]
Thus the V85 cluster-entropy margin survives the exact ground
reweighting.
\end{corollary}

\begin{proof}
Insert \eqref{eq:s6v86-disseminated} in
\Cref{thm:s6v87-profile-transfer}.
\end{proof}

\begin{definition}[Singular-ground reflection/profile card, SGRPC]
\label{def:s6v87-SGRPC}
For the exact response-matched amplitude, SGRPC consists of:
\begin{enumerate}[label=\textup{(R\arabic*)},leftmargin=4em]
\item exact identification of the reference measure
      \(\mu_0=\mu_{\rm W}\) or an explicit Radon--Nikodym bridge to
      the V86 Wilson reference;
\item midpoint-cylinder realization of \(u_0^2\mu_0\);
\item positive half-space factorization of every spatially crossing
      complete-action factor, so that
      \Cref{cor:s6v87-Euclidean-RP} applies;
\item a comparison set \(G\) and ground-profile bounds
      \eqref{eq:s6v87-G-mass}--\eqref{eq:s6v87-h-upper};
\item the cofinal debit \(r_s+b_s=o(s)\);
\item the same quotient and block identification used in GRBAT.
\end{enumerate}
\end{definition}

\begin{theorem}[SGRPC supplies the physical SFCC probability row]
\label{thm:s6v87-SGRPC}
If SGRPC holds, the physical ground star-frustration indicators satisfy
SFCC row (S2) with exponentially small activity.  The V85 rooted
cluster and entropy-versus-capacity conclusions then apply to the
physical ground measure.
\end{theorem}

\begin{proof}
Rows (R2)--(R3) and
\Cref{cor:s6v87-Euclidean-RP} give spatial reflection positivity.
Rows (R1), (R4), and (R5), followed by
\Cref{cor:s6v87-sublinear-profile}, give finite-set domination.
Row (R6) puts the events on the GRBAT block graph.  Apply
\Cref{thm:s6v85-cluster-tail,prop:s6v85-cluster-sum}.
\end{proof}

\subsection{A direct route to the singular profile}
\label{sec:s6v87-row-comparison}

\begin{lemma}[Kernel-row comparison controls the ground profile]
\label{lem:s6v87-row-comparison}
Let \(T\) be a positivity-preserving integral operator
\[
 (Tf)(x)=\int K(x,y)f(y)\,\mu(\dd y),
 \qquad K\ge0,
\]
with strictly positive leading eigenfunction \(h\) and eigenvalue
\(\lambda>0\).  If, for two configurations \(x,z\),
\begin{equation}
 K(x,y)\le R\,K(z,y)
 \quad\text{for }\mu\text{-almost every }y,              \label{eq:s6v87-row-order}
\end{equation}
then
\begin{equation}
 h(x)\le R\,h(z).                                       \label{eq:s6v87-ground-order}
\end{equation}
\end{lemma}

\begin{proof}
Use the eigenfunction equation and positivity:
\[
 \lambda h(x)
 =\int K(x,y)h(y)\,\mu(\dd y)
 \le R\int K(z,y)h(y)\,\mu(\dd y)
 =R\lambda h(z).
\]
Divide by \(\lambda\).
\end{proof}

\begin{corollary}[Uniform response-row suppression would populate the
profile debit]
\label{cor:s6v87-row-profile}
Suppose every \(x\in B^\#\) can be paired with \(z\in G\) so that the
complete normalized transfer kernel obeys
\[
 K(x,y)\le
 \exp\{-\alpha_sN/2\}K(z,y)
\]
for almost every \(y\), and suppose the pairing and the lower profile
on \(G\) give \eqref{eq:s6v87-h-lower}.  Then
\[
 h(x)^2\le m e^{-\alpha_sN}\qquad(x\in B^\#).
\]
Thus \eqref{eq:s6v87-h-upper} holds with \(r_s=0\); retaining the
sharper displayed estimate improves the activity exponent \(a_s\) to
\(a_s+\alpha_s\) before any probability estimate.
\end{corollary}

\begin{proof}
Square \eqref{eq:s6v87-ground-order}.
\end{proof}

\begin{firewall}[The row comparison is deliberately strong]
\label{fire:s6v87-row-strong}
An exterior half-boundary \(y\) can favor a frustrated \(x\), so the
pointwise row order in
\eqref{eq:s6v87-row-order} may fail exactly as the one-link
conditional route failed in V86.  The lemma is a valid acquisition
route, not an assertion that the response-matched kernel has the
order.  If it fails, one needs an averaged projective-kernel,
reflection, or polymer estimate for \(h\).
\end{firewall}

\subsection{Revised frontier}
\label{sec:s6v87-frontier}

\begin{theorem}[V87 ownership decision]
\label{thm:s6v87-decision}
The correct object connecting the V86 Wilson activity to the V81
physical Poincar\'e programme is the leading singular ground profile
\(u_0\), not a symbolic local response action.

More precisely:
\begin{enumerate}[label=\textup{(\roman*)},leftmargin=3.5em]
\item time-transfer positivity and uniqueness of \(u_0\) do not imply
      the spatial reflection positivity needed by the chessboard
      estimate;
\item spatially reflection-positive midpoint cylinders do imply that
      property for \(u_0^2\mu_0\);
\item the ground-profile comparison
      \eqref{eq:s6v87-G-mass}--\eqref{eq:s6v87-h-upper} transfers the
      V86 exponential activity with an exact debit \(r_s+b_s\);
\item a sublinear logarithmic ground profile cannot erase the V83
      two-well obstruction.
\end{enumerate}
\end{theorem}

\begin{proof}
The four assertions are, respectively,
\Cref{prop:s6v87-positivity-counterexample,
thm:s6v87-midpoint,
thm:s6v87-profile-transfer,
cor:s6v87-sublinear-ground}.
\end{proof}

\begin{assessment}[Exact bottleneck after V87]
\label{ass:s6v87-frontier}
The response-matched sources supply a symbolic generated action and
finite response coefficients, while the bare Grand Tensor packet
supplies a finite Wilson half-slab singular problem.  They do not yet
identify one cofinal response-matched amplitude whose leading singular
function satisfies SGRPC.

The next upstream calculation must therefore act on the
\emph{half-slab kernel}: either prove a spatially positive complete
cylinder factorization and sublinear ground-profile debit, or prove an
averaged projective/polymer substitute.  Estimating only the local
Wilson Hessian or only the scalar response coefficient cannot populate
this row.
\end{assessment}

\begin{programme}[V88: projective ground-profile control]
\label{prog:s6v87-V88}
\begin{enumerate}[leftmargin=3em]
\item Express the leading singular equation for the complete
      half-slab kernel on the gauge quotient.
\item Replace the generally false pointwise row order by an integrated
      domination or Hilbert-projective cross-ratio estimate.
\item Determine whether local kernel factors give an
      \(O(\log s)\), \(o(s)\), or order-\(s\) profile debit per
      disseminated star.
\item Keep spatial reflection positivity separate from profile
      magnitude; prove each or route around the failed one.
\item If SGRPC closes, return to the connected-cluster capacity row.
\end{enumerate}
\end{programme}

\begin{firewall}[Claim boundary after sixth-series V87]
\label{fire:s6v87-boundary}
V87 proves the exact singular-ground correction to a matrix-Fisher
conditional; gives an explicit barrier-erasure example; proves a
ground-profile-stable two-well upper bound; proves that transfer
positivity does not imply spatial reflection positivity; proves
inheritance of spatial reflection positivity from midpoint cylinders;
proves a ground-profile transfer theorem for the V86 activity; and
formulates SGRPC.

V87 does not construct the exact response-matched amplitude, prove
spatial reflection positivity or the profile debit for its singular
ground law, prove that the V83 event exhausts all bad conditional
geometries, prove connected-cluster capacity or approximate
tensorization, populate GMPC or HSOCC, construct the continuum OS
theory, or prove the full Yang--Mills mass gap.
\end{firewall}

\begin{theorem}[Sixth-series V87 result ledger]
\label{thm:s6v87-results}
The positive conclusions in \Cref{fire:s6v87-boundary} are exactly
those of
\Cref{prop:s6v87-ground-conditional,
lem:s6v87-barrier-erasure,
thm:s6v87-robust-well,
cor:s6v87-sublinear-ground,
cor:s6v87-healing-dichotomy,
prop:s6v87-positivity-counterexample,
lem:s6v87-RP-closed,
thm:s6v87-midpoint,
cor:s6v87-Euclidean-RP,
thm:s6v87-profile-transfer,
cor:s6v87-sublinear-profile,
thm:s6v87-SGRPC,
lem:s6v87-row-comparison,
cor:s6v87-row-profile,
thm:s6v87-decision}.
\end{theorem}

\begin{proof}
Each conclusion is the cited result with its reference-measure,
singular-ground, spatial-reflection, profile-debit, and
complete-amplitude qualifications unchanged.
\end{proof}

\paragraph{Parent-version record.}
V87 was formed from
\texttt{Renewal\_Yang\_Mills\_Mass\_Gap\_Research\_Notes\_V86.tex}.
The V86 parent had \(39\,263\) logical source lines,
\(1\,375\,989\) bytes, and SHA--256
\[
 \texttt{7f6bb87317a9be5befb8e1afd4416e52aefed60b59e237a76c92003fe3b6d9d2}.
\]
The next compulsory companion audit is V90.

% =============================================================================
% END APPENDED SIXTH-SERIES V87 MATERIAL
% =============================================================================

% =============================================================================
% BEGIN APPENDED SIXTH-SERIES V88 MATERIAL
% =============================================================================

\section{V88: restricted Hilbert--Schmidt control of the physical
ground activity}
\label{sec:s6v88}

\subsection{Why the singular function need not be estimated
pointwise}
\label{sec:s6v88-purpose}

V87 placed the exact leading singular function \(u_0\) between the
bare Wilson law and the physical boundary law.  A direct Harnack or
row-ratio estimate of \(u_0\) is sufficient but unnecessarily strong.
The singular equation itself gives an integrated estimate in which the
unknown function disappears.

Let \(({\mathsf X},\mu_0)\) and
\(({\mathsf Y},\mu_{1/2})\) be probability spaces, let
\[
 (Af)(y)=\int_{\mathsf X}F(y,x)f(x)\,\mu_0(\dd x)
\]
be Hilbert--Schmidt, and assume a simple leading singular pair
\[
 Au_0=s_1v_0,\qquad A^*v_0=s_1u_0,\qquad
 \|u_0\|_2=\|v_0\|_2=1.
\]
The physical initial-boundary measure is
\[
 \pi_0(\dd x)=u_0(x)^2\mu_0(\dd x).
\]

\begin{theorem}[Restricted Hilbert--Schmidt ground bound]
\label{thm:s6v88-restricted-HS}
For every measurable \({\cal B}\subset{\mathsf X}\),
\begin{equation}
 \boxed{
 \pi_0({\cal B})
 \le {1\over s_1^2}
 \int_{\cal B}\int_{\mathsf Y}|F(y,x)|^2\,
 \mu_{1/2}(\dd y)\mu_0(\dd x).}                          \label{eq:s6v88-restricted-HS}
\end{equation}
The analogous final-boundary estimate is
\begin{equation}
 \pi_{1/2}({\cal C})
 \le {1\over s_1^2}
 \int_{\cal C}\int_{\mathsf X}|F(y,x)|^2\,
 \mu_0(\dd x)\mu_{1/2}(\dd y).                           \label{eq:s6v88-left-HS}
\end{equation}
\end{theorem}

\begin{proof}
Let \(M_{\cal B}\) be multiplication by \(1_{\cal B}\).  The singular
equation gives
\[
 \pi_0({\cal B})
 =\|M_{\cal B}u_0\|_2^2
 ={1\over s_1^2}\|M_{\cal B}A^*v_0\|_2^2.
\]
Since \(\|v_0\|_2=1\),
\[
 \|M_{\cal B}A^*v_0\|_2
 \le\|M_{\cal B}A^*\|_{\rm op}
 \le\|M_{\cal B}A^*\|_{\rm HS}.
\]
The square of the last norm is precisely the restricted integral in
\eqref{eq:s6v88-restricted-HS}.  Exchange \(A\) and \(A^*\) for
\eqref{eq:s6v88-left-HS}.
\end{proof}

\begin{corollary}[Stable-rank form]
\label{cor:s6v88-stable-rank}
Define the squared-amplitude probability
\begin{equation}
 \widehat\Lambda_F(\dd x,\dd y)
 :={|F(y,x)|^2\over\|A\|_{\rm HS}^2}
 \,\mu_0(\dd x)\mu_{1/2}(\dd y)                         \label{eq:s6v88-square-law}
\end{equation}
and the effective stable rank
\begin{equation}
 r_{\rm eff}(A):={\|A\|_{\rm HS}^2\over s_1^2}.          \label{eq:s6v88-effective-rank}
\end{equation}
Then
\begin{equation}
 \boxed{
 \pi_0({\cal B})
 \le r_{\rm eff}(A)\,
 \widehat\Lambda_F({\cal B}\times{\mathsf Y}).}          \label{eq:s6v88-stable-rank-bound}
\end{equation}
\end{corollary}

\begin{proof}
Multiply and divide the right side of
\eqref{eq:s6v88-restricted-HS} by \(\|A\|_{\rm HS}^2\).
\end{proof}

\begin{remark}[Exact source ownership]
\label{rem:s6v88-source-ownership}
Equation \eqref{eq:s6v88-restricted-HS} uses exactly the
Hilbert--Schmidt integral printed as \textup{(YMH.2)} in the attached
Grand Tensor source.  It neither replaces the physical ground measure
by the square law \eqref{eq:s6v88-square-law} nor assumes
\(u_0\) is constant.  The stable-rank factor is the exact price of the
replacement.
\end{remark}

\subsection{A lower bound for the leading singular value}
\label{sec:s6v88-s1}

\begin{lemma}[Positive rectangle lower bound]
\label{lem:s6v88-rectangle}
Suppose \(F\ge0\).  Let
\({\cal G}_0\subset{\mathsf X}\) and
\({\cal G}_{1/2}\subset{\mathsf Y}\) have positive measure.  If
\begin{align}
 \mu_0({\cal G}_0)&\ge e^{-b_0N},                        \label{eq:s6v88-G0}\\
 \mu_{1/2}({\cal G}_{1/2})&\ge e^{-b_1N},                \label{eq:s6v88-G1}\\
 F(y,x)&\ge c_F e^{-dN/2}
 \quad\text{on }{\cal G}_{1/2}\times{\cal G}_0,          \label{eq:s6v88-F-lower}
\end{align}
where \(c_F>0\) is an arbitrary global scalar, then
\begin{equation}
 \boxed{
 s_1^2\ge
 c_F^2e^{-N(d+b_0+b_1)}.}                               \label{eq:s6v88-s1-lower}
\end{equation}
\end{lemma}

\begin{proof}
Use the normalized positive test functions
\[
 f_0={1_{{\cal G}_0}\over\sqrt{\mu_0({\cal G}_0)}},
 \qquad
 f_1={1_{{\cal G}_{1/2}}\over
             \sqrt{\mu_{1/2}({\cal G}_{1/2})}}.
\]
The variational definition of \(s_1\) and positivity give
\begin{align*}
 s_1
 &\ge\langle f_1,Af_0\rangle\\
 &\ge c_Fe^{-dN/2}
 \sqrt{\mu_0({\cal G}_0)\mu_{1/2}({\cal G}_{1/2})}\\
 &\ge c_F
 e^{-N(d+b_0+b_1)/2}.
\end{align*}
Square the result.
\end{proof}

\begin{corollary}[Near-flat control of effective rank]
\label{cor:s6v88-rank}
If, in addition, \(0\le F\le c_F\), then
\begin{equation}
 r_{\rm eff}(A)
 \le e^{N(d+b_0+b_1)}.                                  \label{eq:s6v88-rank-bound}
\end{equation}
\end{corollary}

\begin{proof}
Normalized reference measures give
\(\|A\|_{\rm HS}^2\le c_F^2\).  Divide by
\eqref{eq:s6v88-s1-lower}.
\end{proof}

\begin{remark}[No physical normalizer appears]
\label{rem:s6v88-normalizer}
Multiplication of \(F\) by \(c_F\) multiplies both
\(\|M_{\cal B}A^*\|_{\rm HS}^2\) and \(s_1^2\) by \(c_F^2\).
It cancels from every conclusion, as required by
\Cref{ass:s6v51-normalization}.
\end{remark}

\subsection{Action form of the ground-activity estimate}
\label{sec:s6v88-action}

\begin{theorem}[Half-slab action activity compiler]
\label{thm:s6v88-action-compiler}
Suppose the positive half-slab kernel has the exact form
\begin{equation}
 F(y,x)=c_F e^{-S_+(y,x)/2},\qquad S_+\ge0.              \label{eq:s6v88-action-kernel}
\end{equation}
Let \({\cal B}^\#\subset{\mathsf X}\) be an event disseminated over
\(N\) spatial tiles.  Assume
\begin{align}
 S_+(y,x)&\ge a_sN
 \quad\text{on }{\mathsf Y}\times{\cal B}^\#,            \label{eq:s6v88-bad-action}\\
 S_+(y,x)&\le d_sN
 \quad\text{on }{\cal G}_{1/2}\times{\cal G}_0,          \label{eq:s6v88-good-action}\\
 \mu_0({\cal G}_0)&\ge e^{-b_{0,s}N},\qquad
 \mu_{1/2}({\cal G}_{1/2})\ge e^{-b_{1,s}N}.             \label{eq:s6v88-good-volume}
\end{align}
Then
\begin{equation}
 \boxed{
 \pi_0({\cal B}^\#)
 \le
 \exp\{-N(a_s-d_s-b_{0,s}-b_{1,s})\}.}                  \label{eq:s6v88-ground-disseminated}
\end{equation}
\end{theorem}

\begin{proof}
The restricted numerator in
\eqref{eq:s6v88-restricted-HS} is at most
\[
 c_F^2e^{-a_sN},
\]
because the product reference measure is normalized.  On the good
rectangle,
\[
 F\ge c_Fe^{-d_sN/2}.
\]
Apply \Cref{lem:s6v88-rectangle} with
\(d=d_s\), \(b_i=b_{i,s}\), and divide.
\end{proof}

\begin{corollary}[From one disseminated estimate to finite-set
domination]
\label{cor:s6v88-chessboard}
Assume, in addition, that the physical ground law \(\pi_0\) is
spatially reflection positive for the \(N\)-tile chessboard, and that
the star family has \(\chi\) reflection-tiling colours.  Put
\[
 \alpha_s:=a_s-d_s-b_{0,s}-b_{1,s}.
\]
If \(\alpha_s>0\), then every finite family \(E\) of star blocks
satisfies
\begin{equation}
 \boxed{
 \pi_0\left(\bigcap_{e\in E}B_e\right)
 \le\exp\left\{-{\alpha_s\over\chi}|E|\right\}.}          \label{eq:s6v88-HS-domination}
\end{equation}
\end{corollary}

\begin{proof}
Use \eqref{eq:s6v88-ground-disseminated} in
\Cref{thm:s6v86-chessboard}, then apply the largest-colour-class
argument of \Cref{cor:s6v86-colouring}.
\end{proof}

\begin{firewall}[The Hilbert--Schmidt estimate is still a probability
estimate]
\label{fire:s6v88-not-capacity}
Theorem \ref{thm:s6v88-action-compiler} estimates the physical
ground probability of a bad event even though the ground singular
function is unknown.  It does not estimate the form capacity of that
event or the unbounded conditional variance localized there.  The
V84 probability-to-capacity obstruction remains.
\end{firewall}

\subsection{Wilson boundary ownership}
\label{sec:s6v88-Wilson}

Let \({\mathscr P}_{\partial,0}\) be a collection of plaquettes whose
holonomies are measurable from the initial boundary variable \(x\)
and which are actually owned by the half-slab action.  Put
\[
 A_{\partial,0}(x)
 :=\sum_{p\in{\mathscr P}_{\partial,0}}d(U_p).
\]
For a selected boundary edge \(e\), let
\({\mathscr P}_{\partial,0}(e)\subset{\mathscr P}_{\partial,0}\) be
its chosen owned star, of cardinality \(n_e\), and define
\begin{align}
 H_e^\partial
 &:=\sum_{p\in{\mathscr P}_{\partial,0}(e)}S_{e,p},\\
 {\mathfrak F}_e^\partial
 &:=3n_e-\max_{V\in\SU(3)}\Re\Tr(VH_e^\partial),\\
 B_e^\partial(\kappa)
 &:=\{ {\mathfrak F}_e^\partial\ge\kappa\}.              \label{eq:s6v88-boundary-bad}
\end{align}
For a reflected set \(T\) of selected edges, write
\[
 B^\partial(\kappa)^\#
 :=\bigcap_{e\in T}B_e^\partial(\kappa).
\]
The proof of \Cref{lem:s6v86-F-basic} applies verbatim to this
boundary-owned star.

Suppose the half-slab action satisfies
\begin{equation}
 S_+(y,x)
 \ge\theta_\partial s\,A_{\partial,0}(x),
 \qquad \theta_\partial>0,                               \label{eq:s6v88-boundary-owner}
\end{equation}
with all other terms nonnegative.  The constant
\(\theta_\partial\) records whether a boundary plaquette is assigned
fully or fractionally to one half-slab.

\begin{lemma}[Star incidence inside the owned boundary action]
\label{lem:s6v88-boundary-incidence}
Let the disseminated initial-boundary event consist of \(N\) stars and
let each initial-boundary plaquette meet at most \(m_\partial\) selected
central edges.  Then
\begin{equation}
 S_+(y,x)
 \ge {\theta_\partial\kappa s\over m_\partial}N
 \quad\text{on }
 {\mathsf Y}\times B^\partial(\kappa)^\#.               \label{eq:s6v88-owned-cost}
\end{equation}
\end{lemma}

\begin{proof}
Apply \Cref{lem:s6v86-incidence} within the initial boundary slice:
\[
 A_{\partial,0}(x)
 \ge{\kappa N\over m_\partial}.
\]
Insert this in \eqref{eq:s6v88-boundary-owner}.
\end{proof}

\begin{lemma}[Fixed-thickness near-flat half-slab rectangle]
\label{lem:s6v88-half-tube}
Consider a local Wilson half-slab of fixed combinatorial thickness.
Assume its two boundary reference measures are product Haar after an
exact tree gauge, with at most \(\ell_0N,\ell_1N\) retained boundary
links, and that the half-slab contains at most \(qN\) plaquettes.
Then there are near-identity boundary sets
\({\cal G}_0,{\cal G}_{1/2}\) and constants \(C_i\), independent of
periodic spatial volume, such that for \(s\ge s_0\),
\begin{align}
 \mu_0({\cal G}_0)&\ge
 \exp\{-N(4\ell_0\log s+C_0)\},                          \label{eq:s6v88-G0-tube}\\
 \mu_{1/2}({\cal G}_{1/2})&\ge
 \exp\{-N(4\ell_1\log s+C_1)\},                         \label{eq:s6v88-G1-tube}\\
 S_+(y,x)&\le C_2N
 \quad\text{on }{\cal G}_{1/2}\times{\cal G}_0.          \label{eq:s6v88-S-tube}
\end{align}
\end{lemma}

\begin{proof}
Restrict every retained boundary link to the geodesic ball
\(B(I,s^{-1/2})\).  The eight-dimensional Haar volume bound used in
\Cref{lem:s6v86-tube} gives
\eqref{eq:s6v88-G0-tube}--\eqref{eq:s6v88-G1-tube}.
Every plaquette word in a fixed-thickness local slab contains a
bounded number of links.  The quadratic near-identity estimate for
\(d(U_p)\) gives \(d(U_p)\le C/s\); multiplication by \(s\) and
summation over at most \(qN\) plaquettes gives
\eqref{eq:s6v88-S-tube}.  An exact tree gauge only removes Haar
variables and cannot decrease the lower volume estimate for the
retained product chart.
\end{proof}

\begin{theorem}[Boundary-owned bare-Wilson ground activity]
\label{thm:s6v88-bare-ground}
Assume:
\begin{enumerate}[label=\textup{(\roman*)},leftmargin=3.5em]
\item the finite half-slab kernel is
      \(F=c_F e^{-S_+/2}\) with a nonnegative fixed-thickness Wilson
      action;
\item the positive boundary ownership
      \eqref{eq:s6v88-boundary-owner} holds with
      \(\inf\theta_\partial>0\);
\item the exact gauge quotient satisfies the tube hypotheses of
      \Cref{lem:s6v88-half-tube};
\item the physical singular ground law is spatially reflection
      positive, for example by
      \Cref{cor:s6v87-Euclidean-RP}.
\end{enumerate}
Then the physical ground star indicators satisfy
\begin{equation}
 \boxed{
 \pi_0\left(\bigcap_{e\in E}B_e^\partial(\kappa)\right)
 \le
 p_{\rm gr}(s,\kappa)^{|E|},\qquad
 p_{\rm gr}(s,\kappa)
 \le
 \exp\left\{
 -{\theta_\partial\kappa\over\chi m_\partial}s
 +C\log s+C'\right\}.}                                  \label{eq:s6v88-bare-ground}
\end{equation}
The constants are independent of periodic spatial volume.
\end{theorem}

\begin{proof}
Use \Cref{lem:s6v88-boundary-incidence} in
\eqref{eq:s6v88-bad-action}; thus
\[
 a_s={\theta_\partial\kappa\over m_\partial}s.
\]
The tube lemma gives \(d_s=O(1)\) and
\(b_{0,s}+b_{1,s}=O(\log s)\).  Apply
\Cref{thm:s6v88-action-compiler,cor:s6v88-chessboard}.
\end{proof}

\begin{firewall}[The full star must land on the boundary]
\label{fire:s6v88-star-landing}
The V83--V86 event uses six plaquettes in a four-dimensional one-link
star.  A three-dimensional spatial transfer boundary has only four
spatial plaquettes incident on an edge.  Therefore
\(B_e^\partial(\kappa)\) may be identified with the V83 event only
when the declared boundary variable and half-slab ownership actually
contain all six V83 staples and plaquettes.  Otherwise V88 proves
activity for the boundary-owned event
\eqref{eq:s6v88-boundary-bad}, and a separate landing theorem must
show that it contains every bad conditional geometry relevant to
GRBAT.
\end{firewall}

\begin{corollary}[Bare-ground cluster entropy is beaten]
\label{cor:s6v88-bare-cluster}
Under \Cref{thm:s6v88-bare-ground}, every fixed block degree \(D\)
and every cluster-capacity exponent \(a_s=o(s)\) obey
\[
 e^{a_s+1}D\,p_{\rm gr}(s,\kappa)<1
\]
for sufficiently large \(s\).
\end{corollary}

\begin{proof}
The negative term in
\eqref{eq:s6v88-bare-ground} is linear in \(s\), while every other
term is \(o(s)\).  Apply
\Cref{cor:s6v85-weak-cluster}.
\end{proof}

\subsection{Complete-action debits}
\label{sec:s6v88-complete}

The preceding argument is stable under a complete action only when its
signed terms are compared on the two exact regions used in the proof.

\begin{definition}[Half-slab ground-activity card, HSGAC]
\label{def:s6v88-HSGAC}
For a response-matched positive amplitude
\[
 F^{\rm ph}=c_Fe^{-S_+^{\rm ph}/2},
\]
HSGAC consists of:
\begin{enumerate}[label=\textup{(H\arabic*)},leftmargin=4em]
\item exact initial and half-boundary reference measures and gauge
      quotient;
\item spatial reflection positivity of the singular ground law;
\item a disseminated boundary-owned bad-sector lower bound
      \[
      S_+^{\rm ph}\ge
      \left({\theta_\partial\kappa\over m_\partial}s-r_s\right)N
      +c_{\Lambda,s};
      \]
\item a good rectangle with reference masses
      \(e^{-b_{0,s}N}\), \(e^{-b_{1,s}N}\) and
      \[
      S_+^{\rm ph}\le(d_s+c_{\Lambda,s}/N)N;
      \]
\item
      \(r_s+d_s+b_{0,s}+b_{1,s}=o(s)\), after the common extensive
      scalar \(c_{\Lambda,s}\) is cancelled;
\item block/reflection compatibility along the cofinal regulator
      route.
\end{enumerate}
\end{definition}

\begin{theorem}[HSGAC gives physical finite-set domination]
\label{thm:s6v88-HSGAC}
Under HSGAC,
\begin{equation}
 \pi_0^{\rm ph}
 \left(\bigcap_{e\in E}B_e^\partial(\kappa)\right)
 \le
 \exp\left[
 -{|E|\over\chi}
 \left(
 {\theta_\partial\kappa\over m_\partial}s
 -r_s-d_s-b_{0,s}-b_{1,s}
 \right)\right].                                       \label{eq:s6v88-HSGAC}
\end{equation}
Thus SFCC row (S2) holds for these physical bad stars.
\end{theorem}

\begin{proof}
Apply \Cref{thm:s6v88-action-compiler}.  The scalar
\(c_{\Lambda,s}\) multiplies the amplitude by one global factor and
cancels exactly as in \Cref{rem:s6v88-normalizer}.  HSGAC row (H5)
makes the remaining exponent positive and linear in \(s\).
Use \Cref{cor:s6v88-chessboard}.
\end{proof}

\begin{remark}[Relation to CARSC and SGRPC]
\label{rem:s6v88-cards}
CARSC postulates a local ground action; SGRPC estimates the singular
ground profile; HSGAC instead estimates one restricted squared
half-slab kernel and one variational lower bound for \(s_1\).
The three are alternative sufficient routes to the same physical
finite-set domination row.  HSGAC is closest to the actual
\textup{(YMH.1)--(YMH.2)} source data.
\end{remark}

\subsection{Audit of the attached half-slab source}
\label{sec:s6v88-source-audit}

\begin{theorem}[What the source does and does not populate]
\label{thm:s6v88-source-audit}
The attached Grand Tensor definition
\texttt{def:YM-half-slab-amplitude} explicitly supplies
\[
 F(U_+,U_0)=c\exp[-S_+(U_+,U_0)/2]
\]
for a finite Wilson slab, and
\texttt{thm:YM-half-slab-Doob} supplies the Hilbert--Schmidt identity,
the simple positive singular pair, and the physical ground marginal.
These statements populate the algebraic hypotheses of
\Cref{thm:s6v88-restricted-HS}.

The source does not, in that packet, print:
\begin{enumerate}[label=\textup{(\roman*)},leftmargin=3.5em]
\item a decomposition of \(S_+\) identifying a uniform positive
      initial-boundary plaquette owner \(\theta_\partial\);
\item an exact landing of the relevant bad conditional geometry in
      the boundary-owned event \(B_e^\partial(\kappa)\);
\item a cofinal fixed-thickness count and tree-gauge tube estimate;
\item a spatial reflection-positive midpoint-cylinder theorem for the
      resulting singular ground marginal;
\item the corresponding facts for one response-matched rather than
      bare-Wilson amplitude.
\end{enumerate}
Therefore \Cref{thm:s6v88-bare-ground} is a rigorous compiler whose
remaining geometric/source rows are explicit, not a claim that HSGAC
has already been populated.
\end{theorem}

\begin{proof}
The positive statements are the displayed formulas
\textup{(YMH.1)--(YMH.7)} in the source.  The signed-Wilsonian section
constructs effective Hessian and response identities but does not
rewrite the half-slab action \(S_+\) with the five listed
boundary/reflection/cofinal properties.  A direct search of the
half-slab packet finds no such statements.  Hence those rows cannot be
inferred from nuclearity or positivity improvement alone.
\end{proof}

\begin{assessment}[Exact advance and remaining capacity obstruction]
\label{ass:s6v88-frontier}
V88 removes pointwise control of \(u_0\) from the probability
programme.  A Wilson action cost in the squared half-slab kernel,
minus the effective stable-rank cost of the amplitude, is sufficient.
For a fixed-thickness Wilson slab these costs have the favorable scales
\[
 \text{bad action}=c sN,\qquad
 \log r_{\rm eff}=O(N\log s).
\]
The required probability activity is therefore exponential once
boundary ownership and spatial reflection positivity are verified.

This still does not control the connected-cluster form capacity.
That row must use the physical Dirichlet form or a restricted
resolvent, not another probability comparison.
\end{assessment}

\begin{programme}[V89: from restricted probability to restricted
capacity]
\label{prog:s6v88-V89}
\begin{enumerate}[leftmargin=3em]
\item Express the bad-cluster variance debit as a multiplication
      operator sandwiched by the physical resolvent or heat semigroup.
\item Prove a spectral-capacity analogue of
      \Cref{thm:s6v88-restricted-HS} which retains one Dirichlet
      energy.
\item Test whether the good-sector local gap and the restricted
      half-slab activity provide the missing smoothing exponent.
\item Track the cluster-size dependence before applying the V85
      lattice-animal sum.
\item Keep HSGAC source rows open until the explicit \(S_+\) ownership
      decomposition is supplied.
\end{enumerate}
\end{programme}

\begin{firewall}[Claim boundary after sixth-series V88]
\label{fire:s6v88-boundary}
V88 proves a restricted Hilbert--Schmidt bound for the exact physical
ground marginal; identifies its stable-rank debit; proves a positive
rectangle lower bound for the leading singular value; proves an
action-level disseminated activity compiler; derives a conditional
bare-Wilson ground activity
\[
 e^{-c s+O(\log s)};
\]
and formulates the complete HSGAC criterion.

V88 does not populate the source-specific boundary ownership, spatial
reflection, tree-gauge tube, or response-matched HSGAC rows; does not
prove a form-capacity estimate or approximate tensorization; does not
populate GMPC or HSOCC; does not construct the continuum OS theory;
and does not prove the full Yang--Mills mass gap.
\end{firewall}

\begin{theorem}[Sixth-series V88 result ledger]
\label{thm:s6v88-results}
The positive conclusions in \Cref{fire:s6v88-boundary} are exactly
those of
\Cref{thm:s6v88-restricted-HS,
cor:s6v88-stable-rank,
lem:s6v88-rectangle,
cor:s6v88-rank,
thm:s6v88-action-compiler,
cor:s6v88-chessboard,
lem:s6v88-boundary-incidence,
lem:s6v88-half-tube,
thm:s6v88-bare-ground,
cor:s6v88-bare-cluster,
thm:s6v88-HSGAC,
thm:s6v88-source-audit}.
\end{theorem}

\begin{proof}
Each conclusion is the cited result with its half-slab,
singular-value, boundary-owner, reflection-positive, fixed-thickness,
and probability-versus-capacity qualifications unchanged.
\end{proof}

\paragraph{Parent-version record.}
V88 was formed from
\texttt{Renewal\_Yang\_Mills\_Mass\_Gap\_Research\_Notes\_V87.tex}.
The V87 parent had \(39\,961\) logical source lines,
\(1\,400\,163\) bytes, and SHA--256
\[
 \texttt{8fc70da6769a4dd762f5e2c3f4f59adaf50c582a4875e571ee0671c1355a33b2}.
\]
The next compulsory companion audit is V90.

% =============================================================================
% END APPENDED SIXTH-SERIES V88 MATERIAL
% =============================================================================

% =============================================================================
% BEGIN APPENDED SIXTH-SERIES V89 MATERIAL
% =============================================================================

\section{V89: restricted transfer smoothing and a genuine
bad-set capacity inequality}
\label{sec:s6v89}

\subsection{Rarity alone cannot be upgraded to capacity}
\label{sec:s6v89-no-go}

The V84 warning that probability is not capacity has an elementary
sharp model.

\begin{proposition}[Two-state rarity/capacity separation]
\label{prop:s6v89-two-state}
Fix \(p\in(0,1)\).  On \(\{0,1\}\), let
\[
 \pi(1)=p,\qquad\pi(0)=1-p,
\]
and let a reversible continuous-time chain have edge conductance
\(q>0\), so its Dirichlet form is
\[
 {\cal E}_q(f)=q\{f(1)-f(0)\}^2.
\]
For the rare event \(B=\{1\}\), any inequality
\begin{equation}
 \int_B(f-\pi f)^2\,\dd\pi
 \le\varepsilon\operatorname{Var}_\pi(f)
   +\kappa{\cal E}_q(f)                                 \label{eq:s6v89-two-capacity}
\end{equation}
with \(\varepsilon<1-p\) requires
\begin{equation}
 \kappa\ge
 {p(1-p)\{1-p-\varepsilon\}\over q}.                    \label{eq:s6v89-kappa-lower}
\end{equation}
Thus fixed small probability \(p\), by itself, gives no uniform
capacity bound as \(q\downarrow0\).
\end{proposition}

\begin{proof}
Use the centered function
\[
 f(1)=1-p,\qquad f(0)=-p.
\]
Then
\[
 \operatorname{Var}_\pi(f)=p(1-p),\qquad
 {\cal E}_q(f)=q,\qquad
 \int_Bf^2\,\dd\pi=p(1-p)^2.
\]
Insert these identities in
\eqref{eq:s6v89-two-capacity} and solve for \(\kappa\).
\end{proof}

\begin{assessment}[What must replace probability]
\label{ass:s6v89-smoothing-needed}
The missing datum is communication out of the bad set.  In the
two-state model it is the conductance \(q\).  For the Yang--Mills
ground law it must be measured by a restricted physical transfer,
resolvent, or Dirichlet capacity.  V89 uses a restricted transfer
norm because the half-slab source makes that object exactly
computable without pointwise knowledge of the ground state.
\end{assessment}

\subsection{Restricted smoothing implies a form-capacity bound}
\label{sec:s6v89-abstract}

Let \(P\) be a self-adjoint Markov contraction on \(L^2(\pi)\), with
reversible kernel \(P(x,\dd y)\).  For a measurable event \(B\), write
\begin{align}
 \eta_P(B)&:=\|M_BP\|_{2\to2}^2,                         \label{eq:s6v89-eta}\\
 {\cal J}_{P,B}(f)
 &:=\int_B\int
       |f(x)-f(y)|^2P(x,\dd y)\pi(\dd x),                \label{eq:s6v89-JB}
\end{align}
where \(M_B\) is multiplication by \(1_B\).

\begin{theorem}[Restricted-smoothing capacity inequality]
\label{thm:s6v89-smoothing-capacity}
For every \(f\in L^2(\pi)\),
\begin{equation}
 \boxed{
 \int_B|f|^2\,\dd\pi
 \le
 2\eta_P(B)\|f\|_2^2+2{\cal J}_{P,B}(f).}                \label{eq:s6v89-capacity}
\end{equation}
Moreover,
\begin{equation}
 {\cal J}_{P,{\mathsf X}}(f)
 =2\langle f,(I-P)f\rangle.                             \label{eq:s6v89-global-J}
\end{equation}
If \(P=e^{-\delta H}\) and
\({\cal E}_P(f):=\delta^{-1}\langle f,(I-P)f\rangle\),
then for centered \(f\)
\begin{equation}
 \int_B|f|^2\,\dd\pi
 \le
 2\eta_P(B)\operatorname{Var}_\pi(f)
 +4\delta{\cal E}_P(f).                                 \label{eq:s6v89-global-capacity}
\end{equation}
\end{theorem}

\begin{proof}
Decompose
\[
 M_Bf=M_BPf+M_B(I-P)f
\]
and use \(|a+b|^2\le2|a|^2+2|b|^2\).  The first square is at most
\(\eta_P(B)\|f\|_2^2\).  Jensen's inequality for the Markov kernel
gives
\begin{align*}
 |f(x)-Pf(x)|^2
 &=
 \left|\int\{f(x)-f(y)\}P(x,\dd y)\right|^2\\
 &\le\int|f(x)-f(y)|^2P(x,\dd y).
\end{align*}
Integrate over \(x\in B\) to obtain
\eqref{eq:s6v89-capacity}.  Reversibility and invariance give
\[
 \int\!\!\int|f(x)-f(y)|^2P(x,\dd y)\pi(\dd x)
 =2\{\|f\|_2^2-\langle f,Pf\rangle\},
\]
which is \eqref{eq:s6v89-global-J}.  Enlarge \(B\) to the whole space
in the jump term and use centering for
\eqref{eq:s6v89-global-capacity}.
\end{proof}

\begin{remark}[A tunable version]
\label{rem:s6v89-tunable}
For every \(\theta>0\), the same proof with
\[
 |a+b|^2\le(1+\theta)|a|^2+(1+\theta^{-1})|b|^2
\]
gives
\[
 \int_B|f|^2\,\dd\pi
 \le(1+\theta)\eta_P(B)\|f\|_2^2
 +(1+\theta^{-1}){\cal J}_{P,B}(f).
\]
The symmetric choice \(\theta=1\) is used below.
\end{remark}

\begin{theorem}[Weighted family compiler]
\label{thm:s6v89-family}
Let \((B_i,w_i)_{i\in I}\) be a finite weighted event family and put
\begin{equation}
 {\mathscr R}:=\sum_{i\in I}w_i\,P M_{B_i}P.             \label{eq:s6v89-R}
\end{equation}
Assume, on the mean-zero space,
\begin{align}
 \|{\mathscr R}\|&\le r,                                \label{eq:s6v89-R-bound}\\
 \sum_iw_i{\cal J}_{P,B_i}(f)
 &\le2m\delta{\cal E}_P(f).                             \label{eq:s6v89-J-overlap}
\end{align}
Then the family has the weighted form-capacity bound
\begin{equation}
 \boxed{
 \sum_iw_i\int_{B_i}|f|^2\,\dd\pi
 \le
 2r\operatorname{Var}_\pi(f)
 +4m\delta{\cal E}_P(f)}                                \label{eq:s6v89-family-capacity}
\end{equation}
for every \(f\), after replacing it by \(f-\pi f\).
\end{theorem}

\begin{proof}
Sum the first decomposition in the proof of
\Cref{thm:s6v89-smoothing-capacity}.  The smoothed squares equal
\[
 \sum_iw_i\|M_{B_i}Pf\|_2^2
 =\langle f,{\mathscr R}f\rangle
 \le r\|f\|_2^2
\]
on the mean-zero space.  Use
\eqref{eq:s6v89-J-overlap} for the restricted jump terms and restore
the factor two from the elementary square inequality.
\end{proof}

\begin{corollary}[Direct insertion into the V83 compiler]
\label{cor:s6v89-V83}
If the events and weights in \Cref{thm:s6v89-family} are the bad
exterior events appearing in
\eqref{eq:s6v83-bad-debit}, then V84's capacity pair is
\begin{equation}
 \varepsilon_{\rm bad}=2r,\qquad
 \kappa_{\rm bad}=4m\delta.                             \label{eq:s6v89-pair}
\end{equation}
If \(2C_{\rm AT}r<1\), the good-exterior compiler closes with the
explicit gap \eqref{eq:s6v84-capacity-gap}.
\end{corollary}

\begin{proof}
Let \({\cal G}_i\) be the corresponding exterior sigma-field.  Since
\(B_i\in{\cal G}_i\), conditional variance with the constant
\(c=\pi f\) gives
\[
 1_{B_i}\operatorname{Var}(f\mid{\cal G}_i)
 \le
 1_{B_i}\mathbb E\!\left[
 |f-\pi f|^2\mid{\cal G}_i\right].
\]
Average, multiply by \(w_i\), and sum.  The result is bounded by the
left side of \eqref{eq:s6v89-family-capacity}.  Apply
\Cref{thm:s6v89-family,thm:s6v84-capacity-closes}.
\end{proof}

\begin{firewall}[The operator row is essential]
\label{fire:s6v89-operator-row}
Individual estimates
\(\eta_P(B_i)\ll1\) do not imply
\eqref{eq:s6v89-R-bound} without a volume loss: the norm of a sum of
positive operators can be as large as the sum of their norms.  The
cluster argument must exploit locality, an orthogonal martingale
decomposition, or a Schur row sum for the \emph{assembled} operator
\({\mathscr R}\).  Scalar probability summability does not prove this
operator statement.
\end{firewall}

\subsection{The singular ground factor cancels from restricted
transfer smoothing}
\label{sec:s6v89-Doob}

Assume now that \(F\ge0\) and retain the half-slab singular data of
V88.  Let
\begin{equation}
 Q(x,x')
 :=\int_{\mathsf Y}F(y,x)F(y,x')\,\mu_{1/2}(\dd y)      \label{eq:s6v89-Q}
\end{equation}
be the kernel of \(A^*A\).  The ground-transformed physical transfer
on \(L^2(\pi_0)\) is
\begin{equation}
 (Pf)(x)
 ={1\over s_1^2u_0(x)}
 \int Q(x,x')u_0(x')f(x')\,\mu_0(\dd x').               \label{eq:s6v89-Doob-P}
\end{equation}

\begin{lemma}[Exact restricted Doob Hilbert--Schmidt identity]
\label{lem:s6v89-Doob-HS}
For every measurable \({\cal B}\subset{\mathsf X}\),
\begin{equation}
 \boxed{
 \|M_{\cal B}P\|_{{\rm HS},L^2(\pi_0)}^2
 ={1\over s_1^4}
 \int_{\cal B}\int_{\mathsf X}|Q(x,x')|^2\,
 \mu_0(\dd x')\mu_0(\dd x).}                            \label{eq:s6v89-Doob-HS}
\end{equation}
In particular the unknown singular ground function cancels exactly.
\end{lemma}

\begin{proof}
Relative to \(\pi_0(\dd x')=u_0(x')^2\mu_0(\dd x')\), the kernel in
\eqref{eq:s6v89-Doob-P} is
\[
 p(x,x')
 ={Q(x,x')\over s_1^2u_0(x)u_0(x')}.
\]
Therefore
\begin{align*}
 \|M_{\cal B}P\|_{\rm HS}^2
 &=
 \int_{\cal B}\int
 {Q(x,x')^2\over
  s_1^4u_0(x)^2u_0(x')^2}
 \pi_0(\dd x')\pi_0(\dd x)\\
 &={1\over s_1^4}
 \int_{\cal B}\int Q(x,x')^2\,
 \mu_0(\dd x')\mu_0(\dd x).
\end{align*}
\end{proof}

\begin{lemma}[Half-slab row-square bound]
\label{lem:s6v89-row-square}
Put
\[
 R_F(x):=\int_{\mathsf Y}F(y,x)^2\,\mu_{1/2}(\dd y).
\]
Then
\begin{equation}
 |Q(x,x')|^2\le R_F(x)R_F(x')                          \label{eq:s6v89-Q-CS}
\end{equation}
and hence
\begin{equation}
 \int_{\cal B}\int|Q(x,x')|^2\,\dd\mu_0(x')\dd\mu_0(x)
 \le
 \left(\int_{\cal B}R_F\,\dd\mu_0\right)
 \|A\|_{\rm HS}^2.                                     \label{eq:s6v89-Q-restricted}
\end{equation}
\end{lemma}

\begin{proof}
Equation \eqref{eq:s6v89-Q-CS} is Cauchy--Schwarz in the intermediate
variable \(y\).  Integrate it over
\({\cal B}\times{\mathsf X}\), and use
\[
 \int R_F\,\dd\mu_0=\|A\|_{\rm HS}^2.
\]
\end{proof}

\begin{theorem}[Restricted half-slab action controls transfer
smoothing]
\label{thm:s6v89-action-smoothing}
Assume the action hypotheses of
\Cref{thm:s6v88-action-compiler}.  Put
\[
 g_s:=d_s+b_{0,s}+b_{1,s}.
\]
Then
\begin{equation}
 \boxed{
 \eta_P({\cal B}^\#)
 \le
 \|M_{{\cal B}^\#}P\|_{\rm HS}^2
 \le
 \exp\{-N(a_s-2g_s)\}.}                                 \label{eq:s6v89-action-eta}
\end{equation}
Consequently, for centered \(f\),
\begin{equation}
 \boxed{
 \int_{{\cal B}^\#}|f|^2\,\dd\pi_0
 \le
 2e^{-N(a_s-2g_s)}\operatorname{Var}_{\pi_0}(f)
 +4\delta{\cal E}_P(f).}                                \label{eq:s6v89-disseminated-capacity}
\end{equation}
\end{theorem}

\begin{proof}
The bad-action bound gives
\[
 \int_{{\cal B}^\#}R_F\,\dd\mu_0
 =
 \int_{{\cal B}^\#\times{\mathsf Y}}F^2
 \le c_F^2e^{-a_sN}.
\]
Nonnegativity of \(S_+\) gives
\(\|A\|_{\rm HS}^2\le c_F^2\).
The rectangle bound gives
\[
 s_1^2\ge c_F^2e^{-g_sN}.
\]
Insert these three estimates in
\Cref{lem:s6v89-Doob-HS,lem:s6v89-row-square}:
\[
 \|M_{{\cal B}^\#}P\|_{\rm HS}^2
 \le
 {c_F^4e^{-a_sN}\over
  c_F^4e^{-2g_sN}}
 =e^{-N(a_s-2g_s)}.
\]
The operator norm is at most the Hilbert--Schmidt norm.  Apply
\Cref{thm:s6v89-smoothing-capacity}.
\end{proof}

\begin{corollary}[Favorable Wilson scaling for the restricted
capacity]
\label{cor:s6v89-Wilson-scale}
Under the boundary-owned Wilson hypotheses of
\Cref{thm:s6v88-bare-ground},
\begin{equation}
 a_s-2g_s
 =
 {\theta_\partial\kappa\over m_\partial}s-O(\log s).
                                                               \label{eq:s6v89-Wilson-exponent}
\end{equation}
Thus the disseminated bad event has an exponentially small variance
coefficient and an energy debit \(4\delta\).
\end{corollary}

\begin{proof}
Use \Cref{lem:s6v88-boundary-incidence,lem:s6v88-half-tube} in the
definition of \(a_s\) and \(g_s\), then apply
\Cref{thm:s6v89-action-smoothing}.
\end{proof}

\begin{remark}[Capacity costs one additional stable-rank debit]
\label{rem:s6v89-double-rank}
The V88 probability exponent loses \(g_s\), whereas the transfer
smoothing exponent loses \(2g_s\).  This is not a proof artifact:
\eqref{eq:s6v89-Doob-HS} contains \(s_1^{-4}\), while the probability
bound \eqref{eq:s6v88-restricted-HS} contains \(s_1^{-2}\).  Since
\(g_s=O(\log s)\) for a fixed-thickness Wilson slab, either loss is
strictly smaller than the order-\(s\) action gain.
\end{remark}

\subsection{Connected clusters and the exact assembly row}
\label{sec:s6v89-clusters}

For a connected bad cluster \(C\), let \(\widetilde C\) be a bounded
enlargement and let \(P_C^\eta\) be a reversible local physical
transfer on the conditional cluster fibre with exterior \(\eta\).
Let \({\cal B}_C^\eta\) be the exact bad-cluster event on that fibre.

\begin{definition}[Cluster restricted-smoothing card, CRSC]
\label{def:s6v89-CRSC}
CRSC consists of:
\begin{enumerate}[label=\textup{(K\arabic*)},leftmargin=4em]
\item exact conditional local half-slab amplitudes whose Doob
      transfers are \(P_C^\eta\);
\item the uniform restricted smoothing estimate
      \begin{equation}
       \|M_{{\cal B}_C^\eta}P_C^\eta\|_{2\to2}^2
       \le A e^{-q_s|C|},                               \label{eq:s6v89-cluster-eta}
      \end{equation}
      with \(q_s=cs-O(\log s)\);
\item a canonical root for each realized cluster and the V85
      connected-set count;
\item canonical isometric embeddings of every conditional cluster
      fibre into the global \(L^2(\mu)\) fluctuation space, and an
      assembled smoothed-operator row
      \begin{equation}
       \left\|
       \sum_C w_C
       P_C^\eta M_{{\cal B}_C^\eta}P_C^\eta
       \right\|_{\mathbf1^\perp}
       \le r_s;                                         \label{eq:s6v89-cluster-R}
      \end{equation}
\item a restricted jump-energy overlap
      \begin{equation}
       \sum_Cw_C{\cal J}_{P_C^\eta,{\cal B}_C^\eta}(f)
       \le2m_{\rm cl}\delta_s{\cal E}_\mu(f);            \label{eq:s6v89-cluster-J}
      \end{equation}
\item exact identification of these local jump forms with pieces of
      the physical ground Dirichlet form.
\end{enumerate}
\end{definition}

\begin{lemma}[Rooted scalar smoothing sum]
\label{lem:s6v89-rooted-smoothing}
Under CRSC row (K2), for every root \(o\),
\begin{equation}
 \sum_{\substack{C\ni o\\C\ {\rm connected}}}
 A e^{-q_s|C|}
 \le
 {Ae^{-q_s}\over
  1-e^{1-q_s}D}                                        \label{eq:s6v89-rooted-sum}
\end{equation}
whenever \(e^{1-q_s}D<1\).
\end{lemma}

\begin{proof}
Use the V85 bound \(A_n(o)\le(eD)^{n-1}\):
\[
 \sum_{n\ge1}A(eD)^{n-1}e^{-q_sn}
 =
 {Ae^{-q_s}\over1-e^{1-q_s}D}.
\]
\end{proof}

\begin{theorem}[CRSC capacity handoff]
\label{thm:s6v89-CRSC}
If CRSC holds and the conditional-variance bad debit is dominated by
the assembled cluster \(L^2\) terms, then it has the V84 form-capacity
pair
\begin{equation}
 \boxed{
 \varepsilon_{\rm bad}=2r_s,\qquad
 \kappa_{\rm bad}=4m_{\rm cl}\delta_s.}                 \label{eq:s6v89-CRSC-pair}
\end{equation}
If \(2C_{\rm AT}r_s<1\), the V83 good-exterior compiler closes.
\end{theorem}

\begin{proof}
Apply the weighted family compiler
\Cref{thm:s6v89-family} fibrewise, integrate the exterior, and use
CRSC rows (K4)--(K6).  The assumed domination identifies its left side
with the bad conditional-variance debit.  Apply
\Cref{cor:s6v89-V83}.
\end{proof}

\begin{firewall}[The rooted scalar sum is not the operator row]
\label{fire:s6v89-scalar-not-operator}
Lemma \ref{lem:s6v89-rooted-smoothing} proves summability of the
individual restricted norms around one root.  It does not prove
\eqref{eq:s6v89-cluster-R}.  Remote local cluster operators can add on
different martingale components, and the norm of their unstructured
sum can grow with volume.  Row (K4) requires an explicit local
martingale, coloring, or Schur decomposition that assigns each
fluctuation once.
\end{firewall}

\subsection{Frontier and next audit}
\label{sec:s6v89-frontier}

\begin{theorem}[V89 capacity decision]
\label{thm:s6v89-decision}
For a positive half-slab amplitude, a bad-sector Wilson action gain
that beats twice the logarithmic stable-rank cost gives an
exponentially small restricted physical-transfer norm.  That norm
produces a true bad-set form-capacity inequality with energy debit
\(4\delta\), without estimating the singular ground function.

For connected clusters, the remaining gap is no longer an undefined
probability-to-capacity conversion.  It is the assembled operator row
\eqref{eq:s6v89-cluster-R}, together with local jump-energy incidence
\eqref{eq:s6v89-cluster-J}.
\end{theorem}

\begin{proof}
The first assertion is
\Cref{thm:s6v89-action-smoothing,cor:s6v89-Wilson-scale}.  The second
is \Cref{thm:s6v89-CRSC,fire:s6v89-scalar-not-operator}.
\end{proof}

\begin{programme}[V90: compulsory audit and cluster-operator
assembly]
\label{prog:s6v89-V90}
\begin{enumerate}[leftmargin=3em]
\item Perform the compulsory audit of the current SM and NS active
      notes, looking specifically for a local martingale, restricted
      semigroup, or operator row-sum estimate.
\item Decompose the bad conditional variance into canonically rooted
      cluster fibres and write the exact local transfer \(P_C^\eta\).
\item Test whether conditional-expectation martingale differences make
      the ranges in \eqref{eq:s6v89-cluster-R} orthogonal or only
      boundedly overlapping.
\item Prove the jump-energy incidence
      \eqref{eq:s6v89-cluster-J} from the physical block carr\'e du
      champ.
\item Keep the half-slab boundary-owner and response-matched source
      rows separate from this assembly calculation.
\end{enumerate}
\end{programme}

\begin{firewall}[Claim boundary after sixth-series V89]
\label{fire:s6v89-boundary}
V89 proves a two-state rarity/capacity no-go; proves a
restricted-smoothing capacity inequality and its weighted family
version; proves exact cancellation of the singular ground function
from a restricted Doob Hilbert--Schmidt norm; bounds that norm by
half-slab action and stable-rank estimates; obtains a conditional
disseminated Wilson capacity with exponential variance coefficient and
energy debit \(4\delta\); and formulates CRSC.

V89 does not prove the assembled cluster operator row, local
jump-energy incidence, boundary-star landing, response-matched HSGAC,
approximate tensorization, GMPC or HSOCC, the continuum OS theory, or
the full Yang--Mills mass gap.
\end{firewall}

\begin{theorem}[Sixth-series V89 result ledger]
\label{thm:s6v89-results}
The positive conclusions in \Cref{fire:s6v89-boundary} are exactly
those of
\Cref{prop:s6v89-two-state,
thm:s6v89-smoothing-capacity,
thm:s6v89-family,
cor:s6v89-V83,
lem:s6v89-Doob-HS,
lem:s6v89-row-square,
thm:s6v89-action-smoothing,
cor:s6v89-Wilson-scale,
lem:s6v89-rooted-smoothing,
thm:s6v89-CRSC,
thm:s6v89-decision}.
\end{theorem}

\begin{proof}
Each conclusion is the cited result with its transfer-time,
restricted-Hilbert--Schmidt, stable-rank, local-cluster,
operator-row, and boundary-owner qualifications unchanged.
\end{proof}

\paragraph{Parent-version record.}
V89 was formed from
\texttt{Renewal\_Yang\_Mills\_Mass\_Gap\_Research\_Notes\_V88.tex}.
The V88 parent had \(40\,584\) logical source lines,
\(1\,421\,690\) bytes, and SHA--256
\[
 \texttt{72543921f36b545914ef236d89d2b3a93522f1c8ad5525ea934245d18af995f7}.
\]
The next iteration, V90, requires the compulsory companion audit.

% =============================================================================
% END APPENDED SIXTH-SERIES V89 MATERIAL
% =============================================================================

% =============================================================================
% BEGIN APPENDED SIXTH-SERIES V90 MATERIAL
% =============================================================================

\section{Sixth-series V90: the cluster operator row from martingale
off-diagonal decay}
\label{sec:s6v90}

\subsection{Status and purpose of this iteration}

V89 separated three assertions which must not be conflated:
small classical probability of a bad cluster, small restricted
transfer norm on that cluster, and a volume-uniform bound for the sum
of all embedded restricted-transfer operators.  The first two do not
by themselves imply the third.  V90 supplies an exact functional
analytic compiler for the third assertion.  Its new input is an
off-diagonal estimate relative to an orthogonal martingale
decomposition.  This is a genuine extra estimate, not a reformulation
of the scalar cluster activity.

The compulsory five-version audit is also performed below.  It finds
one useful structural analogy in the SM V28 note---a direct
normal-fibre Hardy estimate---but no already-proved Yang--Mills
operator row.  Accordingly, no companion result is silently imported
as a populated Yang--Mills hypothesis.

\subsection{Compulsory V90 companion audit}
\label{sec:s6v90-audit}

\begin{audit}[Companion files inspected at V90]
\label{audit:s6v90-files}
The current active companion notes are
\begin{align*}
 &\texttt{Renewal\_SM\_Einstein\_Continuum\_Research\_Notes\_V29.tex},
 \\
 &\texttt{Renewal\_NS\_Stability\_Research\_Notes\_V26.tex}.
\end{align*}
At the time of the audit the SM file had \(10\,858\) logical source
lines, \(404\,421\) bytes, and SHA--256
\[
 \texttt{cd90a29df499ac632f47c8059be74db3809e002c284cad14f4d910d238304377}.
\]
The NS file had \(5\,319\) logical source lines, \(278\,283\) bytes,
and SHA--256
\[
 \texttt{82c887f8c2c69e4f28fad7c3dc8de5b9099cb5a4bf431eba1fff3e91935978ad}.
\]
The NS checksum is unchanged from the preceding companion audit.  The
SM programme has advanced from V25 to V29.
\end{audit}

\begin{assessment}[What the SM additions do and do not transfer]
\label{ass:s6v90-SM}
The SM V26--V29 additions give the following information relevant to
the present bottleneck.
\begin{enumerate}[leftmargin=3em]
\item V26 treats good/bad energy absorption and a uniform Sobolev
      route.  It explicitly exposes the failure of that route when
      the usable exponent tends to one with dimension.
\item V27 gives a rational, injective spectral repair map and keeps
      its Jacobian and reference-density charges.  This is a rigorous
      repair bookkeeping device, but not a restricted-transfer
      operator estimate.
\item V28 proves a one-dimensional normal-fibre Hardy/collar
      inequality.  Its constants are dimension-free in spectator
      coordinates.  This is a legitimate alternative capacity
      mechanism if one has a Yang--Mills fibre coordinate with a
      uniform outward drift, a controlled collar density, and bounded
      chart overlap.
\item V29 classifies conformal tails using gravitational signs.  Its
      positive cosmological and curvature-square slopes are
      model-specific and do not produce a Wilson cluster slope.
\end{enumerate}
The SU(3) obstruction in V83 is intrinsically multiwell.  None of the
audited SM notes constructs a monotone normal coordinate through all
frustrated Wilson cluster sectors, nor do they prove the martingale
off-diagonal estimate introduced below.  The V28 Hardy route is
therefore retained as an alternative, unpopulated branch; it is not
counted as closure of CRSC.
\end{assessment}

\begin{assessment}[What the NS file contributes at V90]
\label{ass:s6v90-NS}
The active NS V26 file has not changed since the V85 audit.  Its
renewal and source-localization analogies remain useful for organizing
connected objects, but it contains no physical Yang--Mills Doob
kernel, no SU(3) boundary owner, and no \(L^2\) operator row for the
cluster smoothing operators.  Thus there is no new NS import into the
V90 proof chain.
\end{assessment}

\begin{decision}[Audit consequence]
\label{dec:s6v90-audit}
The active route remains the restricted-transfer route of V89.
V90 attacks only its volume-uniform assembly step.  The possible SM
V28 normal-fibre route will be reopened only if a later Yang--Mills
iteration constructs the required monotone fibre and proves its
physical density bounds.
\end{decision}

\subsection{An operator-valued Schur lemma}
\label{sec:s6v90-Schur}

Let \({\cal H}\) be a Hilbert space with a finite orthogonal
decomposition
\begin{equation}
 {\cal H}=\bigoplus_{a\in{\cal I}}{\cal H}_a,\qquad
 \Pi_a:{\cal H}\longrightarrow{\cal H}_a.                 \label{eq:s6v90-decomp}
\end{equation}
For \(T\in{\cal B}({\cal H})\), write
\(T_{ab}:=\Pi_aT\Pi_b\).

\begin{lemma}[Operator-valued Schur test]
\label{lem:s6v90-Schur}
Suppose \(T=T^*\) and
\begin{equation}
 \sup_{a\in{\cal I}}\sum_{b\in{\cal I}}\|T_{ab}\|
 \le \rho.                                                \label{eq:s6v90-Schur-row}
\end{equation}
Then
\begin{equation}
 \|T\|_{{\cal H}\to{\cal H}}\le\rho.                      \label{eq:s6v90-Schur-conclusion}
\end{equation}
\end{lemma}

\begin{proof}
For \(f=\sum_af_a\), \(f_a=\Pi_af\), self-adjointness gives
\(\|T_{ab}\|=\|T_{ba}\|\).  Hence
\begin{align*}
 |\langle f,Tf\rangle|
 &\le\sum_{a,b}\|T_{ab}\|\,\|f_a\|\,\|f_b\|\\
 &\le {1\over2}\sum_{a,b}\|T_{ab}\|
       \bigl(\|f_a\|^2+\|f_b\|^2\bigr)\\
 &\le \rho\sum_a\|f_a\|^2
 =\rho\|f\|^2.
\end{align*}
Taking the supremum of the quadratic form over unit vectors proves
\eqref{eq:s6v90-Schur-conclusion}.
\end{proof}

\begin{remark}[Why positivity alone was insufficient]
\label{rem:s6v90-positivity}
Even when every \(R_C\) is positive and
\(\|R_C\|\le Ae^{-q|C|}\), the sum
\(\sum_CR_C\) may have a volume factor because many ranges can align.
\Cref{lem:s6v90-Schur} removes that factor only after the block rows
are summable.  This is exactly the distinction enforced by
\Cref{fire:s6v89-scalar-not-operator}.
\end{remark}

\subsection{Quasilocal connected-cluster assembly}
\label{sec:s6v90-cluster}

Let \(G=({\cal I},E)\) be the graph of martingale blocks.  For a
nonempty set \(C\subset{\cal I}\), put
\[
 d(a,C):=\min_{c\in C}d_G(a,c).
\]
Assume the graph has the uniform metric summability constant
\begin{equation}
 S_\mu:=
 \sup_{a\in{\cal I}}\sum_{b\in{\cal I}}e^{-\mu d_G(a,b)}
 <\infty                                                   \label{eq:s6v90-Smu}
\end{equation}
and the rooted connected-animal bound
\begin{equation}
 N_n(a):=
 \#\{C\ni a:C\ {\rm connected},\ |C|=n\}
 \le K^{n-1}.                                             \label{eq:s6v90-animal}
\end{equation}
Both constants are independent of the finite volume.  A bounded
degree graph admits such a \(K\), for example the deliberately rough
choice \(K=eD\) used in the preceding cluster count.

\begin{theorem}[Quasilocal cluster operator row]
\label{thm:s6v90-cluster-row}
For every nonempty connected \(C\subset{\cal I}\), let
\(R_C=R_C^*\in{\cal B}({\cal H})\).  Suppose that, uniformly in the
volume and boundary condition,
\begin{equation}
 \|\Pi_aR_C\Pi_b\|
 \le A e^{-q|C|}
       e^{-\mu\{d(a,C)+d(b,C)\}}                           \label{eq:s6v90-block-decay}
\end{equation}
for all \(a,b,C\).  If
\[
 \vartheta:=Ke^{-q}<1,
\]
then the strong finite-volume sum \(R:=\sum_CR_C\) obeys
\begin{equation}
 \boxed{\displaystyle
 \|R\|\le
 r(q,\mu):=
 {A S_\mu^2 e^{-q}\over(1-\vartheta)^2}.}                 \label{eq:s6v90-row-bound}
\end{equation}
In particular, if \(q=q_s=cs-O(\log s)\), then
\begin{equation}
 r(q_s,\mu)=e^{-cs+O(\log s)}                             \label{eq:s6v90-row-scale}
\end{equation}
provided \(A,S_\mu,K\) grow at most polynomially in \(s\).
\end{theorem}

\begin{proof}
Let \(R_{ab}=\Pi_aR\Pi_b\).  The triangle inequality and
\eqref{eq:s6v90-block-decay} give
\[
 \sum_b\|R_{ab}\|
 \le
 A\sum_Ce^{-q|C|}e^{-\mu d(a,C)}
       \sum_be^{-\mu d(b,C)}.
\]
For each \(C\),
\begin{align}
 \sum_be^{-\mu d(b,C)}
 &\le\sum_{u\in C}\sum_be^{-\mu d(b,u)}
 \le |C|S_\mu,                                           \label{eq:s6v90-bsum}\\
 e^{-\mu d(a,C)}
 &\le\sum_{v\in C}e^{-\mu d(a,v)}.                       \label{eq:s6v90-asum}
\end{align}
After inserting these estimates and summing first over the marked
vertex \(v\),
\begin{align*}
 \sum_b\|R_{ab}\|
 &\le
 AS_\mu
 \sum_v e^{-\mu d(a,v)}
 \sum_{\substack{C\ni v\\C\ {\rm connected}}}
 |C|e^{-q|C|}\\
 &\le
 AS_\mu^2
 \sum_{n\ge1}nK^{n-1}e^{-qn}\\
 &=
 {AS_\mu^2e^{-q}\over(1-Ke^{-q})^2}.
\end{align*}
The last identity is
\(\sum_{n\ge1}n\vartheta^{n-1}=(1-\vartheta)^{-2}\).
\Cref{lem:s6v90-Schur} now proves
\eqref{eq:s6v90-row-bound}.  The scale
\eqref{eq:s6v90-row-scale} follows by taking logarithms.
\end{proof}

\begin{corollary}[Strictly local special case]
\label{cor:s6v90-local}
If \(\Pi_aR_C\Pi_b=0\) unless both \(a\) and \(b\) lie in a fixed
radius enlargement \(C^{+r_0}\), and the nonzero blocks are bounded by
\(Ae^{-q|C|}\), then
\Cref{thm:s6v90-cluster-row} applies, for any fixed \(\mu>0\), after
replacing \(A\) by \(Ae^{2\mu r_0}\).
\end{corollary}

\begin{proof}
If \(a,b\in C^{+r_0}\), then
\[
 e^{-\mu\{d(a,C)+d(b,C)\}}\ge e^{-2\mu r_0},
\]
so the asserted local block bound implies
\eqref{eq:s6v90-block-decay} with \(A\) replaced by
\(Ae^{2\mu r_0}\).  If either index is outside \(C^{+r_0}\), the block
vanishes and the same estimate is automatic.
\end{proof}

\begin{firewall}[The new input in the cluster row]
\label{fire:s6v90-block-input}
The scalar restricted-smoothing estimate
\[
 \|M_{{\cal B}_C^\eta}P_C^\eta\|^2
 \le Ae^{-q_s|C|}
\]
does not prove \eqref{eq:s6v90-block-decay}.  To invoke
\Cref{thm:s6v90-cluster-row}, one must additionally show that the
embedded operator
\[
 R_C=
 w_C\,P_C^\eta M_{{\cal B}_C^\eta}P_C^\eta
\]
has exponentially decaying blocks in a specified orthogonal
decomposition.  The estimate must be uniform in volume, exterior
field, cofinal index, and admissible response source.  V90 proves the
assembly implication, not this physical off-diagonal premise.
\end{firewall}

\subsection{The canonical martingale decomposition}
\label{sec:s6v90-martingale}

Fix an ordering of the spatial link blocks and a filtration
\[
 {\cal F}_0\subset{\cal F}_1\subset\cdots\subset{\cal F}_N.
\]
Let \(E_a\) be conditional expectation onto \({\cal F}_a\) in the
physical ground measure and define
\begin{equation}
 \Pi_a:=E_a-E_{a-1},\qquad 1\le a\le N.                   \label{eq:s6v90-martingale-Pi}
\end{equation}

\begin{lemma}[Exact orthogonality]
\label{lem:s6v90-martingale}
On the mean-zero space,
\begin{equation}
 L^2_0(\mu)=\bigoplus_{a=1}^N\operatorname{Ran}\Pi_a,
 \qquad
 \Pi_a\Pi_b=\delta_{ab}\Pi_a.                             \label{eq:s6v90-martingale-orth}
\end{equation}
\end{lemma}

\begin{proof}
Conditional expectations onto nested sigma-fields are orthogonal
projections and satisfy \(E_aE_b=E_{\min(a,b)}\).  Expanding
\((E_a-E_{a-1})(E_b-E_{b-1})\) gives zero for \(a\ne b\) and
\(\Pi_a\) for \(a=b\).  The telescoping sum is
\(\sum_a\Pi_a=E_N-E_0\), which is the identity on centered functions
when \({\cal F}_N\) is the full sigma-field and
\({\cal F}_0\) is trivial.
\end{proof}

\begin{warning}[Ordering is not locality]
\label{warn:s6v90-ordering}
\Cref{lem:s6v90-martingale} is automatic, but a martingale
difference need not be spatially localized merely because its index
is attached to a spatial block.  Consequently,
\eqref{eq:s6v90-block-decay} must be proved from a conditional
connected-kernel or influence estimate; it cannot be inferred from
\eqref{eq:s6v90-martingale-orth}.
\end{warning}

\begin{definition}[Martingale off-diagonal cluster estimate (MOCE)]
\label{def:s6v90-MOCE}
MOCE at scale \(s\) is the assertion that the physical conditional
operators
\[
 R_C^\eta=
 w_C P_C^\eta M_{{\cal B}_C^\eta}P_C^\eta
\]
have canonical isometric embeddings into \(L^2_0(\mu)\) for which
\eqref{eq:s6v90-block-decay} holds with
\[
 q=q_s=cs-O(\log s),\qquad
 \inf_s\mu_s>0,
\]
and with \(A_s,S_{\mu_s},K\) subexponential in \(s\), uniformly in all
allowed \(\eta\).
\end{definition}

\begin{corollary}[MOCE closes the CRSC operator row]
\label{cor:s6v90-MOCE}
MOCE implies CRSC row (K4) with
\begin{equation}
 r_s\le
 {A_sS_{\mu_s}^2e^{-q_s}\over
  (1-Ke^{-q_s})^2}
 =e^{-cs+O(\log s)}.                                     \label{eq:s6v90-MOCE-r}
\end{equation}
\end{corollary}

\begin{proof}
Apply \Cref{thm:s6v90-cluster-row} to the decomposition
\eqref{eq:s6v90-martingale-orth}.
\end{proof}

\subsection{Jump-energy incidence as conductance ownership}
\label{sec:s6v90-jump}

For a reversible Markov kernel \(P\), let
\[
 \Gamma_P(\dd x,\dd y):=\pi(\dd x)P(x,\dd y)
\]
be its symmetric conductance measure.  A restricted jump integral is
an integral of \(|f(x)-f(y)|^2\) against a submeasure of
\(\Gamma_P\).

\begin{lemma}[Conductance-owner incidence]
\label{lem:s6v90-jump-owner}
Let \(\Gamma\) be the conductance measure of the physical comparison
kernel, and let \(\Gamma_C^B\) be the embedded conductance measure
which represents
\(w_C{\cal J}_{P_C^\eta,{\cal B}_C^\eta}\).  Suppose there are
measurable densities \(g_C\ge0\) such that
\[
 \Gamma_C^B=g_C\Gamma,\qquad
 \sum_Cg_C(x,y)\le m_{\rm cl}
 \quad\text{for \(\Gamma\)-almost every }(x,y).            \label{eq:s6v90-owner-density}
\]
Then
\begin{equation}
 \sum_Cw_C{\cal J}_{P_C^\eta,{\cal B}_C^\eta}(f)
 \le
 m_{\rm cl}{\cal J}_{P,{\mathsf X}}(f)
 =
 2m_{\rm cl}\delta\,{\cal E}_P(f).                        \label{eq:s6v90-jump-owned}
\end{equation}
\end{lemma}

\begin{proof}
By Tonelli's theorem and
\eqref{eq:s6v90-owner-density},
\begin{align*}
 \sum_Cw_C{\cal J}_{P_C^\eta,{\cal B}_C^\eta}(f)
 &=
 \int |f(x)-f(y)|^2\sum_Cg_C(x,y)\,
       \Gamma(\dd x,\dd y)\\
 &\le
 m_{\rm cl}
 \int |f(x)-f(y)|^2\Gamma(\dd x,\dd y).
\end{align*}
The last integral is
\({\cal J}_{P,{\mathsf X}}(f)
=2\langle f,(I-P)f\rangle
=2\delta{\cal E}_P(f)\) by
\eqref{eq:s6v89-global-J}.
\end{proof}

\begin{remark}[Geometric content]
\label{rem:s6v90-jump-geometric}
The hypothesis \eqref{eq:s6v90-owner-density} is verified if each
physical jump is assigned to at most \(m_{\rm cl}\) canonically rooted
cluster enlargements and every local jump conductance is genuinely a
subconductance of the physical one.  Bounded geometric overlap alone
does not prove the latter absolute-continuity assertion.
\end{remark}

\subsection{The assembled capacity compiler}
\label{sec:s6v90-capacity}

\begin{theorem}[MOCE plus conductance ownership gives cluster capacity]
\label{thm:s6v90-capacity}
Assume CRSC rows (K1)--(K3) and (K6), MOCE, and
\eqref{eq:s6v90-owner-density}.  Then for every centered
\(f\in L^2(\mu)\),
\begin{equation}
 \boxed{\displaystyle
 \sum_Cw_C\int_{{\cal B}_C^\eta}|f|^2\,\dd\mu
 \le
 2r_s\operatorname{Var}_\mu(f)
 +4m_{\rm cl}\delta_s{\cal E}_\mu(f),}                    \label{eq:s6v90-capacity}
\end{equation}
where
\[
 r_s\le
 {A_sS_{\mu_s}^2e^{-q_s}\over
  (1-Ke^{-q_s})^2}.
\]
Thus the variance coefficient is exponentially small at strong
coupling whenever \(q_s=cs-O(\log s)\).
\end{theorem}

\begin{proof}
\Cref{cor:s6v90-MOCE} supplies
\eqref{eq:s6v89-cluster-R} with the displayed \(r_s\).
\Cref{lem:s6v90-jump-owner} supplies
\eqref{eq:s6v89-cluster-J}.  Apply the weighted-family theorem
\Cref{thm:s6v89-family}.
\end{proof}

\begin{corollary}[Conditional insertion into the V84 gap compiler]
\label{cor:s6v90-gap}
If the cluster family in \Cref{thm:s6v90-capacity} is exactly the
family generated by V83's bad-exterior debit, and if the good
exterior approximate-tensorization hypotheses of V84 hold, then
\[
 \varepsilon_{\rm bad}=2r_s,\qquad
 \kappa_{\rm bad}=4m_{\rm cl}\delta_s.
\]
Whenever \(2C_{\rm AT}r_s<1\), the physical finite-volume spectral
gap is bounded below by the explicit V84 capacity-gap formula,
uniformly in volume.
\end{corollary}

\begin{proof}
This is \Cref{cor:s6v89-V83} with the two inputs populated by
\Cref{thm:s6v90-capacity}.
\end{proof}

\begin{firewall}[Conditional means conditional]
\label{fire:s6v90-capacity}
The conclusion of \Cref{cor:s6v90-gap} is not presently unconditional.
The following physical rows remain unproved:
\begin{enumerate}[leftmargin=3em]
\item boundary ownership and landing of the V83 frustrated-star
      geometry in the half-slab action;
\item response-uniform local restricted smoothing at Wilson scale;
\item MOCE, namely the block estimate
      \eqref{eq:s6v90-block-decay} for the physical embedded
      cluster operators;
\item absolute-continuous conductance ownership
      \eqref{eq:s6v90-owner-density};
\item the good-exterior approximate tensorization and the remaining
      GMPC/HSOCC continuum interfaces.
\end{enumerate}
V90 closes the combinatorial-functional implication from MOCE to a
volume-uniform row.  It does not prove MOCE.
\end{firewall}

\subsection{Where MOCE has to come from}
\label{sec:s6v90-source}

\begin{proposition}[A sufficient connected-kernel estimate]
\label{prop:s6v90-connected-kernel}
Let \(R_C\) be the physical embedded cluster operator.  Suppose there
are kernels \(R_{C,\Gamma}\) indexed by connected polymers
\(\Gamma\supset C\) such that, in operator norm,
\begin{align}
 R_C&=\sum_{\Gamma\supset C}R_{C,\Gamma},                 \label{eq:s6v90-kernel-expansion}\\
 \|\Pi_aR_{C,\Gamma}\Pi_b\|
 &\le
 A_0e^{-q_0|C|-\zeta|\Gamma\setminus C|}
 1_{\{a,b\in\Gamma^{+r_0}\}}.                            \label{eq:s6v90-kernel-bound}
\end{align}
If the animal entropy is strictly smaller than \(\zeta\), then MOCE
holds, with \(q\) equal to \(q_0\) minus a volume-independent entropy
debit and with some \(\mu>0\).
\end{proposition}

\begin{proof}
For fixed \(a,b,C\), every contributing \(\Gamma\) must connect
\(C^{+r_0}\) to the \(r_0\)-neighbourhoods of \(a\) and \(b\).
Consequently
\[
 |\Gamma\setminus C|
 \ge c_0\{d(a,C)+d(b,C)\}-c_1|C|
\]
for constants depending only on the graph geometry and \(r_0\).
Sum \eqref{eq:s6v90-kernel-bound} over connected enlargements.
The strict excess of \(\zeta\) over animal entropy leaves an
exponential factor
\[
 e^{-\mu\{d(a,C)+d(b,C)\}}
\]
and changes \(e^{-q_0|C|}\) only by
\(e^{c_2|C|}\).  This is
\eqref{eq:s6v90-block-decay} with \(q=q_0-c_2\).
\end{proof}

\begin{warning}[Expansion legitimacy]
\label{warn:s6v90-expansion}
\Cref{prop:s6v90-connected-kernel} is a sufficient theorem, not an
assertion that the singular physical Doob kernel already has the
expansion \eqref{eq:s6v90-kernel-expansion}.  Establishing absolute
operator-norm convergence, response uniformity, and compatibility
with the ground transform is the next analytic bottleneck.
\end{warning}

\begin{decision}[V90 bottleneck decision]
\label{dec:s6v90-bottleneck}
The assembled-row gap of V89 has now been reduced to two checkable
local statements:
\begin{enumerate}[leftmargin=3em]
\item a connected half-slab kernel expansion with excess decay over
      animal entropy, sufficient for MOCE;
\item a conductance submeasure/owner identity, sufficient for the
      jump-energy incidence.
\end{enumerate}
The next iteration must work upstream at the physical Doob kernel:
derive its connected covariance expansion from the positive
half-slab amplitude, or prove that such an expansion is unavailable
without an additional mixing estimate.  Merely repeating a scalar
polymer sum will not advance the proof.
\end{decision}

\begin{theorem}[Sixth-series V90 exact advance]
\label{thm:s6v90-decision}
V90 proves:
\begin{enumerate}[leftmargin=3em]
\item the operator-valued Schur test
      \Cref{lem:s6v90-Schur};
\item the volume-uniform quasilocal connected-cluster row
      \Cref{thm:s6v90-cluster-row}, with explicit constant
      \eqref{eq:s6v90-row-bound};
\item the canonical physical-measure martingale decomposition
      \Cref{lem:s6v90-martingale};
\item the conductance-owner incidence
      \Cref{lem:s6v90-jump-owner};
\item the conditional cluster capacity compiler
      \Cref{thm:s6v90-capacity};
\item the sufficient connected-kernel criterion
      \Cref{prop:s6v90-connected-kernel}.
\end{enumerate}
Thus CRSC's operator-assembly and jump-incidence rows now have exact
sufficient conditions with no volume factor.  The new primary
bottleneck is proving those conditions for the physical, singular
Doob-transformed Yang--Mills transfer.
\end{theorem}

\begin{proof}
The six conclusions are precisely the cited results.  The limitations
are those in
\Cref{fire:s6v90-block-input,fire:s6v90-capacity,
warn:s6v90-expansion}.
\end{proof}

\begin{programme}[V91: physical connected-kernel origin of MOCE]
\label{prog:s6v90-V91}
\begin{enumerate}[leftmargin=3em]
\item Write the exact normalized half-slab kernel and its two-copy
      representation for the blocks
      \(\Pi_aR_C\Pi_b\), keeping singular values explicit.
\item Separate disconnected action components which cancel under the
      martingale differences from connected components meeting
      \(a\), \(C\), and \(b\).
\item Determine whether Wilson character positivity plus the
      boundary owner yields absolute convergence with decay exceeding
      connected-animal entropy.
\item If the cancellation is insufficient, identify the least
      additional strong-coupling mixing estimate and trace it back to
      the complete action rather than postulating MOCE.
\item In parallel, formulate the exact physical conductance density
      needed in \eqref{eq:s6v90-owner-density}.
\end{enumerate}
\end{programme}

\begin{firewall}[Claim boundary after sixth-series V90]
\label{fire:s6v90-boundary}
V90 performs the compulsory SM/NS audit; proves an operator-valued
Schur test; proves a volume-uniform cluster-operator assembly theorem
from exponential martingale off-diagonal decay; identifies the exact
orthogonal martingale decomposition; proves a conductance-owner
jump-incidence lemma; and combines these into a conditional
bad-cluster form-capacity inequality.

V90 does not prove MOCE for the physical Doob kernel, conductance
ownership, boundary-star landing, response-uniform HSGAC,
good-exterior approximate tensorization, GMPC or HSOCC, the continuum
OS theory, or the full Yang--Mills mass gap.
\end{firewall}

\begin{theorem}[Sixth-series V90 result ledger]
\label{thm:s6v90-results}
The positive conclusions in \Cref{fire:s6v90-boundary} are exactly
those of
\Cref{audit:s6v90-files,
lem:s6v90-Schur,
thm:s6v90-cluster-row,
lem:s6v90-martingale,
cor:s6v90-MOCE,
lem:s6v90-jump-owner,
thm:s6v90-capacity,
prop:s6v90-connected-kernel,
thm:s6v90-decision}.
\end{theorem}

\begin{proof}
Each conclusion is the cited result with its off-diagonal,
animal-entropy, conductance-owner, response-uniformity, and
ground-transform qualifications unchanged.
\end{proof}

\paragraph{Parent-version record.}
V90 was formed from
\texttt{Renewal\_Yang\_Mills\_Mass\_Gap\_Research\_Notes\_V89.tex}.
The V89 parent had \(41\,160\) logical source lines,
\(1\,440\,324\) bytes, and SHA--256
\[
 \texttt{3421d3cc92b63006bff54915f19c66e21213674c237eb0b7bd09320c13d01f2}.
\]
The compulsory V90 companion audit used the active SM V29 and NS V26
files recorded in \Cref{audit:s6v90-files}.

% =============================================================================
% END APPENDED SIXTH-SERIES V90 MATERIAL
% =============================================================================

% =============================================================================
% BEGIN APPENDED SIXTH-SERIES V91 MATERIAL
% =============================================================================

\section{Sixth-series V91: exact ground-transform cancellation and
the projected smoothing bootstrap}
\label{sec:s6v91}

\subsection{Upstream correction of the assembly target}

The V90 Schur theorem is correct, but its proposed application to the
\emph{uncentered} operators
\[
 P M_{G_B^c}P
\]
is stronger than the V83 bad term requires and is generally
incompatible with spectator variables.  The V83 term is not
\(\int_{G_B^c}|f|^2\); it is
\[
 \mathbb E\!\left[
 1_{G_B^c}\operatorname{Var}(f\mid{\cal F}_{B^c})\right].
\]
It therefore contains the exact local fluctuation projection
\(D_B=I-\mathbb E[\cdot\mid{\cal F}_{B^c}]\).  Retaining this
projection gives a different bootstrap in which the individual
restricted norm is factored against the \emph{total conditional
variance sum}.  No norm of a sum of uncentered operators is needed.

\subsection{The complete singular ground transform cancels
unitarily}
\label{sec:s6v91-unitary}

Retain the V89 notation
\[
 Q=A^*A,\qquad Q u_0=s_1^2u_0,\qquad
 \pi_0(\dd x)=u_0(x)^2\mu_0(\dd x),
\]
and assume \(\|u_0\|_{L^2(\mu_0)}=1\) and \(u_0>0\) almost
everywhere.  Define
\begin{equation}
 {\cal U}:L^2(\pi_0)\longrightarrow L^2(\mu_0),
 \qquad ({\cal U}f)(x):=u_0(x)f(x).                       \label{eq:s6v91-U}
\end{equation}

\begin{theorem}[Exact unitary removal of the singular ground factor]
\label{thm:s6v91-unitary}
The map \({\cal U}\) is unitary and
\begin{equation}
 {\cal U}P{\cal U}^{-1}=s_1^{-2}Q.                       \label{eq:s6v91-P-conjugate}
\end{equation}
For every measurable event \({\cal B}\),
\begin{align}
 \|M_{\cal B}P\|_{L^2(\pi_0)\to L^2(\pi_0)}
 &=
 s_1^{-2}\|M_{\cal B}Q\|_{L^2(\mu_0)\to L^2(\mu_0)},
                                                               \label{eq:s6v91-norm-exact}\\
 {\cal U}(PM_{\cal B}P){\cal U}^{-1}
 &=
 s_1^{-4}QM_{\cal B}Q.                                  \label{eq:s6v91-R-exact}
\end{align}
Thus cancellation of \(u_0\) is exact at operator level, not only
after taking the Hilbert--Schmidt norm.
\end{theorem}

\begin{proof}
The definition of \(\pi_0\) gives
\[
 \|{\cal U}f\|_{L^2(\mu_0)}^2
 =\int u_0^2|f|^2\,\dd\mu_0
 =\|f\|_{L^2(\pi_0)}^2.
\]
Surjectivity follows from \(u_0>0\) almost everywhere, with inverse
\({\cal U}^{-1}g=g/u_0\).  Substitution in
\eqref{eq:s6v89-Doob-P} yields
\[
 ({\cal U}P{\cal U}^{-1}g)(x)
 ={1\over s_1^2}\int Q(x,x')g(x')\,\mu_0(\dd x').
\]
Multiplication by \(1_{\cal B}\) commutes with \({\cal U}\).
Equations \eqref{eq:s6v91-norm-exact} and
\eqref{eq:s6v91-R-exact} follow by unitary invariance of the operator
norm and by multiplication of the conjugated factors.
\end{proof}

\begin{corollary}[Two-copy half-slab kernel]
\label{cor:s6v91-two-copy}
The operator on the right of \eqref{eq:s6v91-R-exact} has
\(\mu_0\)-kernel
\begin{equation}
 \widehat r_{\cal B}(x,x')
 ={1\over s_1^4}
 \int_{\cal B}Q(x,z)Q(z,x')\,\mu_0(\dd z).                \label{eq:s6v91-two-Q}
\end{equation}
If \(F\ge0\), this is the nonnegative two-copy expression
\begin{multline}
 \widehat r_{\cal B}(x,x')
 ={1\over s_1^4}
 \int_{\cal B}\int_{{\mathsf Y}^2}
 F(y,x)F(y,z)F(y',z)F(y',x')\\
 {}\times
 \mu_{1/2}(\dd y)\mu_{1/2}(\dd y')\mu_0(\dd z).
                                                               \label{eq:s6v91-two-copy}
\end{multline}
\end{corollary}

\begin{proof}
The kernel-composition rule gives \eqref{eq:s6v91-two-Q}.
Insert \eqref{eq:s6v89-Q} twice and use Tonelli's theorem; all
integrands are nonnegative.
\end{proof}

\begin{corollary}[V89's Hilbert--Schmidt identity recovered]
\label{cor:s6v91-HS}
Taking the Hilbert--Schmidt norm in
\eqref{eq:s6v91-norm-exact} gives
\eqref{eq:s6v89-Doob-HS}.
\end{corollary}

\begin{proof}
The Hilbert--Schmidt kernel of \(M_{\cal B}Q\) is
\(1_{\cal B}(x)Q(x,x')\).  Square and integrate against
\(\mu_0\otimes\mu_0\).
\end{proof}

\subsection{Why the uncentered MOCE premise is not the natural target}
\label{sec:s6v91-spectator}

\begin{proposition}[Spectator-identity obstruction]
\label{prop:s6v91-spectator}
Let
\({\cal H}={\cal H}_L\otimes{\cal H}_S\), let
\(P=P_L\otimes I_S\), and let
\(M_{\cal B}=M_{{\cal B},L}\otimes I_S\).  Then
\begin{equation}
 PM_{\cal B}P
 =
 (P_LM_{{\cal B},L}P_L)\otimes I_S.                      \label{eq:s6v91-spectator}
\end{equation}
If \(P_LM_{{\cal B},L}P_L\ne0\) and
\((h_b)_{b\in J}\) is any orthonormal spectator family, the
projections
\[
 \Pi_b=I_L\otimes |h_b\rangle\langle h_b|
\]
satisfy
\begin{equation}
 \|\Pi_b(PM_{\cal B}P)\Pi_b\|
 =
 \|P_LM_{{\cal B},L}P_L\|                                \label{eq:s6v91-no-decay}
\end{equation}
for every \(b\).  Hence an uncentered local smoothing operator does
not acquire off-diagonal decay merely because the spectator label is
far from \({\cal B}\).
\end{proposition}

\begin{proof}
Equation \eqref{eq:s6v91-spectator} is the tensor-product
multiplication rule.  Compressing its spectator identity to
\(\operatorname{span}\{h_b\}\) leaves the local factor unchanged,
which proves \eqref{eq:s6v91-no-decay}.
\end{proof}

\begin{decision}[Status of V90's MOCE branch]
\label{dec:s6v91-MOCE}
\Cref{thm:s6v90-cluster-row} remains a valid sufficient theorem if
MOCE is independently proved for a genuinely connected or centered
operator.  However,
\Cref{prop:s6v91-spectator} shows that MOCE should not be demanded of
the raw uncentered tensor-local operator.  The primary route is
therefore changed to retain the conditional fluctuation projection
which is already present in V83.
\end{decision}

\subsection{Exact representation of the bad conditional variance}
\label{sec:s6v91-projection}

Let \(({\mathsf X},{\cal F},\pi)\) be a probability space.  For each
block \(i\), let \({\cal G}_i\subset{\cal F}\) be its exterior
sigma-field, let
\[
 E_i f:=\mathbb E_\pi[f\mid{\cal G}_i],
 \qquad D_i:=I-E_i,
\]
and let \({\cal B}_i\in{\cal G}_i\) be the bad-exterior event.

\begin{lemma}[Bad variance is a projected norm]
\label{lem:s6v91-bad-projected}
For every \(f\in L^2(\pi)\),
\begin{align}
 \mathbb E_\pi\!\left[
 1_{{\cal B}_i}\operatorname{Var}_\pi(f\mid{\cal G}_i)
 \right]
 &=
 \|M_{{\cal B}_i}D_if\|_2^2,                             \label{eq:s6v91-bad-projected}\\
 \mathbb E_\pi[
 \operatorname{Var}_\pi(f\mid{\cal G}_i)]
 &=
 \|D_if\|_2^2.                                           \label{eq:s6v91-total-projected}
\end{align}
\end{lemma}

\begin{proof}
Conditional variance is
\(\mathbb E[|f-E_if|^2\mid{\cal G}_i]\).  Multiply by the
\({\cal G}_i\)-measurable indicator and use the tower property to get
\eqref{eq:s6v91-bad-projected}.  The second identity is the same
calculation with indicator one.
\end{proof}

\begin{remark}[The lost projection in V89]
\label{rem:s6v91-lost}
The inequality in the proof of \Cref{cor:s6v89-V83} replaced
\(D_if\) by \(f-\pi f\).  It was valid, but it discarded precisely
the block fluctuation which prevents irrelevant spectator modes from
entering the bad term.  The next theorem applies restricted smoothing
before making that relaxation.
\end{remark}

\subsection{Projected restricted smoothing}
\label{sec:s6v91-smoothing}

For each \(i\), let \(P_i\) be a self-adjoint Markov contraction on
\(L^2(\pi)\) and define
\begin{align}
 \eta_i&:=\|M_{{\cal B}_i}P_i\|_{2\to2}^2,                \label{eq:s6v91-etai}\\
 {\cal J}_i(g)&:=
 \int_{{\cal B}_i}\int
 |g(x)-g(y)|^2P_i(x,\dd y)\pi(\dd x).                    \label{eq:s6v91-Ji}
\end{align}

\begin{theorem}[Projected restricted-smoothing inequality]
\label{thm:s6v91-projected-smoothing}
For every \(f\in L^2(\pi)\),
\begin{equation}
 \boxed{\displaystyle
 \mathbb E_\pi\!\left[
 1_{{\cal B}_i}\operatorname{Var}(f\mid{\cal G}_i)\right]
 \le
 2\eta_i\|D_if\|_2^2+2{\cal J}_i(D_if).}                 \label{eq:s6v91-projected-smoothing}
\end{equation}
\end{theorem}

\begin{proof}
Apply \Cref{thm:s6v89-smoothing-capacity} to the function
\(g=D_if\), the kernel \(P_i\), and the event \({\cal B}_i\).
Use \Cref{lem:s6v91-bad-projected} on the left.
\end{proof}

\begin{corollary}[No assembled smoothing-operator row]
\label{cor:s6v91-no-row}
Let \(w_i\ge0\), put
\[
 {\cal S}(f):=\sum_iw_i\|D_if\|_2^2,
 \qquad
 {\cal A}(f):=
 \sum_iw_i
 \mathbb E[
 1_{{\cal B}_i}\operatorname{Var}(f\mid{\cal G}_i)].
\]
If \(\sup_i\eta_i\le\eta\), then
\begin{equation}
 {\cal A}(f)
 \le
 2\eta{\cal S}(f)+2\sum_iw_i{\cal J}_i(D_if).             \label{eq:s6v91-summed}
\end{equation}
No estimate of
\(\|\sum_iw_iP_iM_{{\cal B}_i}P_i\|\) occurs.
\end{corollary}

\begin{proof}
Multiply \eqref{eq:s6v91-projected-smoothing} by \(w_i\), sum, and
use \(\eta_i\le\eta\).
\end{proof}

\subsection{Closing the bad sector by bootstrapping the full variance
sum}
\label{sec:s6v91-bootstrap}

\begin{definition}[Projected jump incidence]
\label{def:s6v91-PJI}
The family has projected jump-incidence constant \(m_{\rm pr}\) at
time \(\delta\) if
\begin{equation}
 \sum_iw_i{\cal J}_i(D_if)
 \le
 2m_{\rm pr}\delta\,{\cal E}_\pi(f)                       \label{eq:s6v91-PJI}
\end{equation}
for every form-domain \(f\).
\end{definition}

\begin{theorem}[Projected smoothing bootstrap]
\label{thm:s6v91-bootstrap}
Assume:
\begin{enumerate}[label=\textup{(P\arabic*)},leftmargin=4em]
\item weighted approximate tensorization
      \begin{equation}
       \operatorname{Var}_\pi(f)
       \le C_{\rm AT}{\cal S}(f);                         \label{eq:s6v91-AT}
      \end{equation}
\item a good-exterior local gap and energy incidence giving
      \begin{equation}
       {\cal G}(f):=
       \sum_iw_i\mathbb E[
       1_{{\cal B}_i^c}\operatorname{Var}(f\mid{\cal G}_i)]
       \le {m_{\rm inc}\over\lambda_{\rm good}}
       {\cal E}_\pi(f);                                  \label{eq:s6v91-good}
      \end{equation}
\item the uniform restricted smoothing
      \(\sup_i\eta_i\le\eta<1/2\);
\item projected jump incidence \eqref{eq:s6v91-PJI}.
\end{enumerate}
Then
\begin{equation}
 \boxed{\displaystyle
 \operatorname{Var}_\pi(f)
 \le
 {C_{\rm AT}\over1-2\eta}
 \left(
 {m_{\rm inc}\over\lambda_{\rm good}}
 +4m_{\rm pr}\delta
 \right){\cal E}_\pi(f).}                                \label{eq:s6v91-Poincare}
\end{equation}
Consequently the physical spectral gap is at least
\begin{equation}
 \boxed{\displaystyle
 \lambda_{\rm phys}\ge
 {1-2\eta\over
 C_{\rm AT}
 \left(m_{\rm inc}/\lambda_{\rm good}
       +4m_{\rm pr}\delta\right)}.}                       \label{eq:s6v91-gap}
\end{equation}
\end{theorem}

\begin{proof}
By \Cref{lem:s6v91-bad-projected},
\[
 {\cal S}(f)={\cal G}(f)+{\cal A}(f).
\]
Use \eqref{eq:s6v91-good}, then
\Cref{cor:s6v91-no-row,def:s6v91-PJI}:
\[
 {\cal S}(f)
 \le
 {m_{\rm inc}\over\lambda_{\rm good}}{\cal E}_\pi(f)
 +2\eta{\cal S}(f)
 +4m_{\rm pr}\delta{\cal E}_\pi(f).
\]
Move \(2\eta{\cal S}(f)\) to the left and divide by
\(1-2\eta>0\).  Finally use
\eqref{eq:s6v91-AT}.  Inverting the resulting Poincaré constant gives
\eqref{eq:s6v91-gap}.
\end{proof}

\begin{remark}[Comparison with the V84 capacity compiler]
\label{rem:s6v91-comparison}
The V84 compiler required a direct bad debit
\({\cal A}\le\varepsilon\operatorname{Var}+\kappa{\cal E}\).
The projected bootstrap instead permits
\[
 {\cal A}\le2\eta{\cal S}+4m_{\rm pr}\delta{\cal E}
\]
and absorbs the first term at the level of the full conditional
variance sum.  Its smallness condition is \(2\eta<1\), not
\(C_{\rm AT}\varepsilon<1\).  This is not a stronger conclusion from
the same hypotheses; it uses the exact projection retained in
\eqref{eq:s6v91-bad-projected}.
\end{remark}

\subsection{A local sufficient condition for projected jump incidence}
\label{sec:s6v91-local-jump}

Suppose the physical form has a nonnegative local decomposition
\begin{equation}
 {\cal E}_\pi(f)=\sum_{a\in{\cal A}}{\cal E}_a(f).
                                                               \label{eq:s6v91-form-decomp}
\end{equation}
Associate to each block \(i\) two finite neighbourhoods
\({\cal N}_i\subset{\cal N}_i^+\subset{\cal A}\).

\begin{lemma}[Local projected-energy compiler]
\label{lem:s6v91-local-energy}
Assume
\begin{align}
 {\cal J}_i(g)
 &\le
 2\delta\sum_{a\in{\cal N}_i}{\cal E}_a(g),               \label{eq:s6v91-local-J}\\
 \sum_{a\in{\cal N}_i}{\cal E}_a(D_if)
 &\le
 L\sum_{b\in{\cal N}_i^+}{\cal E}_b(f),                  \label{eq:s6v91-D-stability}\\
 \sup_{b\in{\cal A}}
 \sum_{i:\,b\in{\cal N}_i^+}w_i
 &\le\nu.                                                 \label{eq:s6v91-neighbour-overlap}
\end{align}
Then \eqref{eq:s6v91-PJI} holds with
\begin{equation}
 m_{\rm pr}=L\nu.                                        \label{eq:s6v91-mpr}
\end{equation}
\end{lemma}

\begin{proof}
Sum \eqref{eq:s6v91-local-J}, use
\eqref{eq:s6v91-D-stability}, reverse the two finite sums, and apply
\eqref{eq:s6v91-neighbour-overlap}:
\begin{align*}
 \sum_iw_i{\cal J}_i(D_if)
 &\le
 2\delta L\sum_iw_i
 \sum_{b\in{\cal N}_i^+}{\cal E}_b(f)\\
 &\le
 2\delta L\nu\sum_b{\cal E}_b(f)
 =
 2\delta L\nu{\cal E}_\pi(f).
\end{align*}
\end{proof}

\begin{firewall}[The local projection estimate is substantive]
\label{fire:s6v91-D-stability}
Conditional expectation is an \(L^2\) contraction, but this alone
does not imply \eqref{eq:s6v91-D-stability}.  That row requires
control of gradients of the conditional law as its exterior changes,
or an equivalent commutator estimate between \(D_i\) and the physical
form.  It must be proved for the gauge quotient and the
response-matched ground measure.  V91 has only isolated its exact
local form.
\end{firewall}

\subsection{Revised acquisition card and decision}

\begin{definition}[Projected restricted-smoothing card (PRSC)]
\label{def:s6v91-PRSC}
PRSC consists of:
\begin{enumerate}[label=\textup{(R\arabic*)},leftmargin=4em]
\item exact bad exterior events
      \({\cal B}_i=G_i^c\in{\cal G}_i\) containing all V83
      separated-maximizer geometries;
\item response-uniform half-slab estimates giving
      \(\eta_i\le\eta_s=e^{-cs+O(\log s)}\);
\item a good-exterior conditional gap
      \(\lambda_{\rm good}>0\);
\item weighted approximate tensorization with volume-uniform
      \(C_{\rm AT}\);
\item a local form decomposition satisfying
      \eqref{eq:s6v91-local-J}--\eqref{eq:s6v91-neighbour-overlap}
      with volume-uniform \(L,\nu\);
\item boundary ownership, cofinal compatibility, and gauge-quotient
      compatibility for every preceding row.
\end{enumerate}
\end{definition}

\begin{theorem}[PRSC-to-finite-volume-gap handoff]
\label{thm:s6v91-PRSC}
If PRSC holds and \(2\eta_s<1\), then the physical finite-volume
Hamiltonian has the volume-uniform lower gap
\[
 {1-2\eta_s\over
 C_{\rm AT}
 \left(m_{\rm inc}/\lambda_{\rm good}
       +4L\nu\delta_s\right)}.
\]
\end{theorem}

\begin{proof}
Use \Cref{lem:s6v91-local-energy} in
\Cref{thm:s6v91-bootstrap}.
\end{proof}

\begin{decision}[V91 bottleneck decision]
\label{dec:s6v91-bottleneck}
The uncentered assembled operator row is no longer on the primary
proof path.  The remaining local analytic issue in the projected
route is \eqref{eq:s6v91-D-stability}: a response-uniform energy
stability estimate for the conditional fluctuation projection.
Going upstream, this is a derivative-of-conditional-density problem.
The next iteration will derive its exact covariance commutator and
determine which Wilson Hessian or influence bound would control it.
\end{decision}

\begin{theorem}[Sixth-series V91 exact advance]
\label{thm:s6v91-decision}
V91 proves:
\begin{enumerate}[leftmargin=3em]
\item exact unitary removal of the singular ground factor from the
      full restricted transfer operator;
\item the nonnegative two-copy half-slab kernel
      \eqref{eq:s6v91-two-copy};
\item a spectator-identity obstruction to the raw uncentered MOCE
      application;
\item the exact projected representation of bad conditional
      variance;
\item the projected restricted-smoothing inequality and its
      no-operator-row sum;
\item the bootstrap gap \eqref{eq:s6v91-gap};
\item the local projected-energy compiler
      \Cref{lem:s6v91-local-energy}.
\end{enumerate}
This replaces an unnecessarily strong global operator-row target by
a local commutator/energy-stability target.
\end{theorem}

\begin{proof}
The assertions are
\Cref{thm:s6v91-unitary,
cor:s6v91-two-copy,
prop:s6v91-spectator,
lem:s6v91-bad-projected,
thm:s6v91-projected-smoothing,
cor:s6v91-no-row,
thm:s6v91-bootstrap,
lem:s6v91-local-energy}.
The remaining physical qualifications are
\Cref{fire:s6v91-D-stability}.
\end{proof}

\begin{programme}[V92: conditional-gradient commutator]
\label{prog:s6v91-V92}
\begin{enumerate}[leftmargin=3em]
\item Differentiate the exact conditional expectation \(E_i f\)
      along an exterior link and express the derivative as the
      conditional expectation of the derivative of \(f\) plus a
      covariance with the logarithmic density derivative.
\item Bound that covariance by a conditional Poincaré estimate and
      the local Wilson force, retaining the good/bad exterior split.
\item Sum the resulting gradient neighbourhoods and compute the
      explicit overlap constant in
      \eqref{eq:s6v91-D-stability}.
\item Test the identity after the singular ground transform and under
      response-matched source insertions.
\item If a uniform covariance bound would already presuppose the
      desired gap, isolate the circular row and switch to the
      half-slab two-copy representation.
\end{enumerate}
\end{programme}

\begin{firewall}[Claim boundary after sixth-series V91]
\label{fire:s6v91-boundary}
V91 proves exact operator-level ground-factor cancellation, a
two-copy restricted kernel, a spectator no-go for the uncentered
MOCE application, the projected restricted-smoothing bootstrap, and
a local sufficient condition for projected jump incidence.

V91 does not prove the response-uniform restricted norm for every
physical bad exterior, conditional-projection energy stability,
boundary-star landing, good-exterior approximate tensorization,
GMPC or HSOCC, the continuum OS theory, or the full Yang--Mills mass
gap.
\end{firewall}

\begin{theorem}[Sixth-series V91 result ledger]
\label{thm:s6v91-results}
The positive conclusions in \Cref{fire:s6v91-boundary} are exactly
those of
\Cref{thm:s6v91-unitary,
cor:s6v91-two-copy,
cor:s6v91-HS,
prop:s6v91-spectator,
lem:s6v91-bad-projected,
thm:s6v91-projected-smoothing,
cor:s6v91-no-row,
thm:s6v91-bootstrap,
lem:s6v91-local-energy,
thm:s6v91-PRSC,
thm:s6v91-decision}.
\end{theorem}

\begin{proof}
Each conclusion is the cited result with its conditional-exterior,
projected-jump, form-locality, response-uniformity, and cofinal
qualifications unchanged.
\end{proof}

\paragraph{Parent-version record.}
V91 was formed from
\texttt{Renewal\_Yang\_Mills\_Mass\_Gap\_Research\_Notes\_V90.tex}.
The V90 parent had \(41\,819\) logical source lines,
\(1\,463\,667\) bytes, and SHA--256
\[
 \texttt{3f3a3e3a3cdbce4b4e1567dca84c290818e6d4819565be801a2c4dbf81e829d9}.
\]

% =============================================================================
% END APPENDED SIXTH-SERIES V91 MATERIAL
% =============================================================================

% =============================================================================
% BEGIN APPENDED SIXTH-SERIES V92 MATERIAL
% =============================================================================

\section{Sixth-series V92: the conditional-gradient commutator and
the ground-score obstruction}
\label{sec:s6v92}

\subsection{Purpose}

V91 reduced projected jump incidence to the local Sobolev stability
of
\[
 D_i=I-\mathbb E[\cdot\mid{\cal G}_i].
\]
V92 computes the derivative of this projection exactly.  The
calculation separates a harmless differentiated observable from a
conditional covariance with the logarithmic score.  For the physical
ground measure, that score contains a derivative of the singular
ground function.  The latter cancels from restricted transfer norms
by V91, but it does \emph{not} cancel from derivatives of spatial
conditional expectation.

\subsection{Conditional differentiation}
\label{sec:s6v92-diff}

Work first in a product chart
\[
 (z,\eta)\in{\mathsf X}_i\times{\mathsf X}_{i^c}
\]
with fixed reference measure \(\dd z\) on the block fibre.  Let
\[
 \rho_i(z\mid\eta)>0,\qquad
 \int_{{\mathsf X}_i}\rho_i(z\mid\eta)\,\dd z=1
\]
be a differentiable conditional density, and write
\[
 E_i f(\eta):=\int f(z,\eta)\rho_i(z\mid\eta)\,\dd z,
 \qquad D_i:=I-E_i.
\]
For an exterior derivation \(X_a\), define the conditional score
\begin{equation}
 \sigma_{ia}(z,\eta):=
 X_a\log\rho_i(z\mid\eta).                               \label{eq:s6v92-score}
\end{equation}
All identities below first hold for smooth cylinder functions under
a dominated-differentiation hypothesis, and then extend by closure
when the displayed quantities are square integrable.

\begin{lemma}[The score is conditionally centered]
\label{lem:s6v92-score-centered}
\begin{equation}
 E_i\sigma_{ia}=0.                                       \label{eq:s6v92-score-zero}
\end{equation}
\end{lemma}

\begin{proof}
Differentiate the normalization of \(\rho_i\):
\[
 0=X_a\int\rho_i\,\dd z
 =\int (X_a\log\rho_i)\rho_i\,\dd z.
\]
\end{proof}

\begin{theorem}[Exact conditional-gradient commutator]
\label{thm:s6v92-commutator}
For every exterior derivation \(X_a\),
\begin{align}
 X_aE_if
 &=E_i(X_af)+
   \operatorname{Cov}_i(f,\sigma_{ia}),                  \label{eq:s6v92-XE}\\
 X_aD_if
 &=D_i(X_af)-
   \operatorname{Cov}_i(f,\sigma_{ia}).                  \label{eq:s6v92-XD}
\end{align}
For every fibre derivation \(Y\) acting only on \(z\),
\begin{equation}
 YD_if=Yf.                                               \label{eq:s6v92-YD}
\end{equation}
\end{theorem}

\begin{proof}
Differentiate under the conditional integral:
\begin{align*}
 X_aE_if
 &=
 \int (X_af)\rho_i\,\dd z
 +\int f(X_a\log\rho_i)\rho_i\,\dd z\\
 &=
 E_i(X_af)+E_i(f\sigma_{ia}).
\end{align*}
By \eqref{eq:s6v92-score-zero}, the second term is the conditional
covariance.  Subtract the result from \(X_af\) to get
\eqref{eq:s6v92-XD}.  Since \(E_if\) is independent of the fibre
coordinate \(z\), applying \(Y\) gives \eqref{eq:s6v92-YD}.
\end{proof}

\begin{corollary}[Exact conditional square identity]
\label{cor:s6v92-square}
For every exterior derivation \(X_a\),
\begin{equation}
 E_i|X_aD_if|^2
 =
 \operatorname{Var}_i(X_af)
 +|\operatorname{Cov}_i(f,\sigma_{ia})|^2.               \label{eq:s6v92-square}
\end{equation}
\end{corollary}

\begin{proof}
In \eqref{eq:s6v92-XD}, the function \(D_i(X_af)\) has conditional
mean zero, while
\(\operatorname{Cov}_i(f,\sigma_{ia})\) is constant on the fibre.
The cross term therefore vanishes.  Also
\(E_i|D_i(X_af)|^2=\operatorname{Var}_i(X_af)\).
\end{proof}

\subsection{A quantitative energy-stability theorem}
\label{sec:s6v92-stability}

Let the fibre carré du champ be
\[
 {\mathfrak e}_i(f)(\eta):=
 \sum_{Y\in{\cal Y}_i}E_i|Yf|^2.
\]
Assume a conditional Poincaré inequality
\begin{equation}
 \operatorname{Var}_i(h)
 \le\lambda_i(\eta)^{-1}{\mathfrak e}_i(h)(\eta).
                                                               \label{eq:s6v92-cond-gap}
\end{equation}
For an exterior direction \(a\), put
\begin{equation}
 \alpha_{ia}(\eta)^2:=
 \operatorname{Var}_i(\sigma_{ia})(\eta).                \label{eq:s6v92-alpha}
\end{equation}

\begin{theorem}[Conditional-score control of the projection gradient]
\label{thm:s6v92-score-control}
Under \eqref{eq:s6v92-cond-gap},
\begin{equation}
 E_i|X_aD_if|^2
 \le
 E_i|X_af|^2+
 {\alpha_{ia}(\eta)^2\over\lambda_i(\eta)}
 {\mathfrak e}_i(f)(\eta).                               \label{eq:s6v92-XD-bound}
\end{equation}
\end{theorem}

\begin{proof}
The first term in \eqref{eq:s6v92-square} is at most
\(E_i|X_af|^2\).  Conditional Cauchy--Schwarz and
\eqref{eq:s6v92-cond-gap} give
\[
 |\operatorname{Cov}_i(f,\sigma_{ia})|^2
 \le
 \operatorname{Var}_i(f)\operatorname{Var}_i(\sigma_{ia})
 \le
 {\alpha_{ia}^2\over\lambda_i}{\mathfrak e}_i(f).
\]
\end{proof}

Let \(\bar\pi_i(\dd\eta)\) be the exterior marginal.  For a fixed
finite set \(\partial i\) of exterior directions, define
\begin{equation}
 \Lambda_i:=
 \operatorname*{ess\,sup}_{\eta}
 {1\over\lambda_i(\eta)}
 \sum_{a\in\partial i}\alpha_{ia}(\eta)^2.                \label{eq:s6v92-Lambda}
\end{equation}

\begin{corollary}[Local Sobolev stability of \(D_i\)]
\label{cor:s6v92-D-stability}
Let the local energy consist of the fibre derivatives
\({\cal Y}_i\) and exterior derivatives \(X_a\),
\(a\in\partial i\).  If \(\Lambda_i<\infty\), then
\begin{multline}
 \int\left\{
 \sum_{Y\in{\cal Y}_i}|YD_if|^2+
 \sum_{a\in\partial i}|X_aD_if|^2
 \right\}\,\dd\pi\\
 \le
 \int\left\{
 (1+\Lambda_i)\sum_{Y\in{\cal Y}_i}|Yf|^2+
 \sum_{a\in\partial i}|X_af|^2
 \right\}\,\dd\pi.                                      \label{eq:s6v92-local-stability}
\end{multline}
In particular, the local stability row
\eqref{eq:s6v91-D-stability} holds with
\[
 L_i\le1+\Lambda_i
\]
after the fibre and boundary terms are included in
\({\cal N}_i^+\).
\end{corollary}

\begin{proof}
Use \eqref{eq:s6v92-YD} for fibre derivatives, sum
\eqref{eq:s6v92-XD-bound} over \(a\in\partial i\), integrate over
\(\bar\pi_i\), and apply \eqref{eq:s6v92-Lambda}.
\end{proof}

\begin{corollary}[Conditional-score completion of the V91 row]
\label{cor:s6v92-PJI}
If
\[
 \sup_i\Lambda_i\le\Lambda<\infty
\]
uniformly in volume, cofinal index, response source, and admissible
exterior, and if the neighbourhood overlap is \(\nu\), then the V91
local projected-energy compiler applies with
\begin{equation}
 m_{\rm pr}\le(1+\Lambda)\nu.                            \label{eq:s6v92-mpr}
\end{equation}
\end{corollary}

\begin{proof}
Use \Cref{cor:s6v92-D-stability} as
\eqref{eq:s6v91-D-stability} and then apply
\Cref{lem:s6v91-local-energy}.
\end{proof}

\subsection{The physical score and the returned ground factor}
\label{sec:s6v92-physical}

Suppose the physical conditional density relative to product Haar
measure has the form
\begin{equation}
 \rho_i(z\mid\eta)
 ={1\over Z_i(\eta)}
 e^{-V_i(z,\eta)}u_0(z,\eta)^2.                          \label{eq:s6v92-physical-density}
\end{equation}
Here \(V_i\) contains the explicitly assigned Wilson and
response-source factors not already included in \(u_0\).

\begin{proposition}[Exact physical score decomposition]
\label{prop:s6v92-physical-score}
For every exterior direction \(X_a\),
\begin{equation}
 \sigma_{ia}
 =
 D_i^{\rm fib}\!\left(
 -X_aV_i+2X_a\log u_0
 \right),                                                \label{eq:s6v92-physical-score}
\end{equation}
where
\[
 D_i^{\rm fib}h:=h-E_ih.
\]
Consequently
\begin{equation}
 \alpha_{ia}^2
 \le
 2\operatorname{Var}_i(X_aV_i)
 +8\operatorname{Var}_i(X_a\log u_0).                    \label{eq:s6v92-alpha-split}
\end{equation}
\end{proposition}

\begin{proof}
Differentiate the logarithm of
\eqref{eq:s6v92-physical-density}.  The derivative of
\(\log Z_i\) is the conditional mean of
\(-X_aV_i+2X_a\log u_0\), because the score is centered by
\Cref{lem:s6v92-score-centered}.  This proves
\eqref{eq:s6v92-physical-score}.  Apply
\(\operatorname{Var}(r+s)\le2\operatorname{Var}(r)
+2\operatorname{Var}(s)\) with \(s=2X_a\log u_0\).
\end{proof}

\begin{remark}[The explicit Wilson part]
\label{rem:s6v92-Wilson-force}
For a fixed finite-range Wilson star, differentiation along a
unit-normalized invariant vector field gives
\[
 |X_aV_i|
 \le c_W n_\star s+K_{\rm resp},
\]
where \(n_\star\) is the bounded number of incident action terms and
\(K_{\rm resp}\) is the corresponding response-source derivative.
Thus the Wilson contribution to
\eqref{eq:s6v92-alpha-split} is at most polynomial in \(s\) when
\(K_{\rm resp}\) is.  The unpopulated term is the conditional ground
score
\[
 \operatorname{Var}_i(X_a\log u_0).
\]
\end{remark}

\begin{firewall}[Ground cancellation has a precise boundary]
\label{fire:s6v92-cancellation-boundary}
\Cref{thm:s6v91-unitary} cancels \(u_0\) from
\(M_{\cal B}P\) by conjugating the entire transfer operator.  A
spatial conditional expectation is defined by the ground density
itself.  Its derivative therefore contains
\(X_a\log u_0\) through
\eqref{eq:s6v92-physical-score}.  Operator-level ground cancellation
does not imply conditional-gradient ground cancellation.
\end{firewall}

\subsection{What the singular equation does control}
\label{sec:s6v92-singular-derivative}

Assume \(X_a\) differentiates the first boundary variable of the
kernel \(Q\), and that differentiation under the integral is valid.
The singular equation is
\[
 s_1^2u_0(x)=\int Q(x,x')u_0(x')\,\mu_0(\dd x').
\]

\begin{lemma}[Exact derivative of the singular ground function]
\label{lem:s6v92-u-derivative}
\begin{align}
 X_au_0(x)
 &=
 {1\over s_1^2}
 \int (X_a^{(1)}Q)(x,x')u_0(x')\,\mu_0(\dd x'),          \label{eq:s6v92-Xu}\\
 X_a\log u_0(x)
 &=
 {\int (X_a^{(1)}Q)(x,x')u_0(x')\,\mu_0(\dd x')
  \over s_1^2u_0(x)}.                                   \label{eq:s6v92-Xlogu}
\end{align}
\end{lemma}

\begin{proof}
Differentiate the singular equation in \(x\).
The singular value is a scalar independent of the boundary point.
Divide by \(u_0(x)>0\) for the second identity.
\end{proof}

\begin{proposition}[Global Fisher bound from the differentiated
kernel]
\label{prop:s6v92-Fisher}
If \(X_a^{(1)}Q\) is bounded on \(L^2(\mu_0)\), then
\begin{align}
 \int|X_a\log u_0|^2\,\dd\pi_0
 &=\|X_au_0\|_{L^2(\mu_0)}^2                             \label{eq:s6v92-Fisher-id}\\
 &\le
 s_1^{-4}\|X_a^{(1)}Q\|_{2\to2}^2.                       \label{eq:s6v92-Fisher-bound}
\end{align}
The operator norm may be replaced by the Hilbert--Schmidt norm when
the latter is finite.
\end{proposition}

\begin{proof}
Since \(\dd\pi_0=u_0^2\dd\mu_0\),
\[
 \int|X_a\log u_0|^2\,\dd\pi_0
 =\int|X_au_0|^2\,\dd\mu_0.
\]
Apply \eqref{eq:s6v92-Xu}, the operator norm bound, and
\(\|u_0\|_2=1\).
\end{proof}

\begin{corollary}[Only an averaged conditional influence follows]
\label{cor:s6v92-average-score}
Let \(\bar\pi_i\) be the exterior marginal.  Then
\begin{equation}
 \int
 \operatorname{Var}_i(X_a\log u_0)\,
 \bar\pi_i(\dd\eta)
 \le
 s_1^{-4}\|X_a^{(1)}Q\|_{2\to2}^2.                       \label{eq:s6v92-average-score}
\end{equation}
\end{corollary}

\begin{proof}
Conditional variance is bounded by conditional second moment.
Integrate and apply \Cref{prop:s6v92-Fisher}.
\end{proof}

\begin{lemma}[Average influence does not give the required multiplier
bound]
\label{lem:s6v92-average-not-uniform}
For a nonnegative measurable coefficient \(a(\eta)\), the inequality
\begin{equation}
 \int a(\eta)h(\eta)\,\bar\pi_i(\dd\eta)
 \le C\int h(\eta)\,\bar\pi_i(\dd\eta)                   \label{eq:s6v92-multiplier}
\end{equation}
for every nonnegative \(h\in L^1(\bar\pi_i)\) holds if and only if
\(a\le C\) almost everywhere.  Therefore an \(L^1\) bound such as
\eqref{eq:s6v92-average-score} cannot by itself supply the essential
supremum in \eqref{eq:s6v92-Lambda}.
\end{lemma}

\begin{proof}
The forward implication follows by taking \(h\) to be the indicator
of \(\{a>C+\varepsilon\}\).  The reverse implication is immediate.
\end{proof}

\begin{assessment}[No hidden closure from the singular equation]
\label{ass:s6v92-no-closure}
The differentiated singular equation gives a useful global Fisher
estimate, but the projected-energy compiler needs a multiplier bound
against the local fibre-energy density
\({\mathfrak e}_i(f)(\eta)\).  By
\Cref{lem:s6v92-average-not-uniform}, the global Fisher estimate does
not provide it.  One needs either:
\begin{enumerate}[leftmargin=3em]
\item a response-uniform conditional ground-score bound;
\item an \(L^r\) ground-score bound paired with an
      \(L^{r'}\) improvement for local energy densities; or
\item a transfer-kernel commutator argument which avoids
      differentiating the spatial conditional law.
\end{enumerate}
The second choice is the same substantive moment-improvement issue
already exposed in V84, not a free consequence of rarity.
\end{assessment}

\subsection{Ground-score influence card}

\begin{definition}[Ground-score influence card (GSIC)]
\label{def:s6v92-GSIC}
GSIC for a block family consists of:
\begin{enumerate}[label=\textup{(S\arabic*)},leftmargin=4em]
\item differentiable physical conditional densities
      \eqref{eq:s6v92-physical-density};
\item conditional Poincaré constants \(\lambda_i(\eta)>0\);
\item response-uniform bounds for
      \(\operatorname{Var}_i(X_a\log u_0)\);
\item a volume-uniform row
      \[
       \sup_i\operatorname*{ess\,sup}_{\eta}
       {\sum_{a\in\partial i}\alpha_{ia}(\eta)^2
        \over\lambda_i(\eta)}
       <\infty;
      \]
\item finite neighbourhood overlap and compatibility with the gauge
      quotient and cofinal bridges.
\end{enumerate}
\end{definition}

\begin{theorem}[GSIC populates PRSC's projected-energy row]
\label{thm:s6v92-GSIC}
GSIC implies
\eqref{eq:s6v91-D-stability} with
\[
 L=1+\sup_i\Lambda_i.
\]
Together with PRSC rows (R1)--(R4) and (R6), it yields the finite-volume
gap \eqref{eq:s6v91-gap}, with
\[
 m_{\rm pr}\le
 \nu(1+\sup_i\Lambda_i).
\]
\end{theorem}

\begin{proof}
Apply
\Cref{prop:s6v92-physical-score,
cor:s6v92-D-stability,
cor:s6v92-PJI,
thm:s6v91-bootstrap}.
\end{proof}

\begin{decision}[V92 bottleneck decision]
\label{dec:s6v92-bottleneck}
The conditional-gradient route has reached a genuine ground-score
row.  Estimating that row by a conditional Poincaré inequality risks
circularity on the V83 bad exterior, while the singular equation
currently yields only an averaged Fisher bound.  The next iteration
will pursue the third branch of
\Cref{ass:s6v92-no-closure}: an abstract commutator estimate for
\([D_i,P]\) which uses the two-copy transfer kernel and identifies
exactly whether conditional ground scores can be avoided.
\end{decision}

\begin{theorem}[Sixth-series V92 exact advance]
\label{thm:s6v92-decision}
V92 proves:
\begin{enumerate}[leftmargin=3em]
\item the exact conditional-gradient commutator
      \eqref{eq:s6v92-XD};
\item the orthogonal square identity
      \eqref{eq:s6v92-square};
\item the conditional-score energy-stability estimate
      \eqref{eq:s6v92-XD-bound};
\item the exact physical score decomposition
      \eqref{eq:s6v92-physical-score};
\item the differentiated singular equation
      \eqref{eq:s6v92-Xlogu};
\item the global Fisher bound
      \eqref{eq:s6v92-Fisher-bound};
\item the proof that averaged influence is insufficient for the
      required multiplier estimate.
\end{enumerate}
Thus the V91 local stability row is reduced to GSIC, and its remaining
ground dependence is explicit rather than hidden.
\end{theorem}

\begin{proof}
The conclusions are exactly
\Cref{thm:s6v92-commutator,
cor:s6v92-square,
thm:s6v92-score-control,
prop:s6v92-physical-score,
lem:s6v92-u-derivative,
prop:s6v92-Fisher,
lem:s6v92-average-not-uniform}.
\end{proof}

\begin{programme}[V93: transfer commutator versus ground score]
\label{prog:s6v92-V93}
\begin{enumerate}[leftmargin=3em]
\item Derive an exact decomposition of
      \(M_{{\cal B}_i}D_i f\) using \(P\), \(D_iP\), and the
      commutator \([D_i,P]\).
\item Determine whether the unitary conjugation
      \({\cal U}P{\cal U}^{-1}=s_1^{-2}Q\) also gives a ground-free
      representation of the commutator.
\item Prove a projected smoothing inequality with an explicit
      commutator debit.
\item Test the tensor-product and finite-range limits to prevent a
      hidden spectator-volume factor.
\item If the commutator still contains conjugated conditional
      projections, state the minimal kernel influence estimate rather
      than claiming cancellation.
\end{enumerate}
\end{programme}

\begin{firewall}[Claim boundary after sixth-series V92]
\label{fire:s6v92-boundary}
V92 proves the conditional-gradient and ground-score identities,
derives a sufficient local energy-stability estimate, and obtains a
global Fisher bound from the differentiated half-slab singular
equation.  It also proves why that averaged bound does not supply the
uniform multiplier needed by the projected bootstrap.

V92 does not prove GSIC, conditional-projection energy stability for
the physical response-matched measure, boundary-star landing,
good-exterior approximate tensorization, GMPC or HSOCC, the continuum
OS theory, or the full Yang--Mills mass gap.
\end{firewall}

\begin{theorem}[Sixth-series V92 result ledger]
\label{thm:s6v92-results}
The positive conclusions in \Cref{fire:s6v92-boundary} are exactly
those of
\Cref{lem:s6v92-score-centered,
thm:s6v92-commutator,
cor:s6v92-square,
thm:s6v92-score-control,
cor:s6v92-D-stability,
cor:s6v92-PJI,
prop:s6v92-physical-score,
lem:s6v92-u-derivative,
prop:s6v92-Fisher,
cor:s6v92-average-score,
lem:s6v92-average-not-uniform,
thm:s6v92-GSIC,
thm:s6v92-decision}.
\end{theorem}

\begin{proof}
Each conclusion is the cited result with its differentiability,
conditional-gap, essential-supremum, ground-score, response, and
cofinal qualifications unchanged.
\end{proof}

\paragraph{Parent-version record.}
V92 was formed from
\texttt{Renewal\_Yang\_Mills\_Mass\_Gap\_Research\_Notes\_V91.tex}.
The V91 parent had \(42\,402\) logical source lines,
\(1\,482\,928\) bytes, and SHA--256
\[
 \texttt{9ffa91f680a878963594ca352980be5cae24d327ded0bc8cf3f27da96e09b346}.
\]

% =============================================================================
% END APPENDED SIXTH-SERIES V92 MATERIAL
% =============================================================================

% =============================================================================
% BEGIN APPENDED SIXTH-SERIES V93 MATERIAL
% =============================================================================

\section{Sixth-series V93: a transfer-commutator bootstrap}
\label{sec:s6v93}

\subsection{Replacing spatial score derivatives by a transfer
commutator}

V92 showed that differentiating the physical conditional projection
exposes \(X\log u_0\).  V93 develops the alternative promised there.
Instead of differentiating \(D_i\), commute it through a positive
short-time transfer \(P_i\).  This produces a precise error
\([D_i,P_i]\).  A form-small estimate for that commutator is enough
for the good/bad bootstrap and does not require an a priori
pointwise ground-score bound.

\subsection{The algebraic commutator identity}
\label{sec:s6v93-algebra}

Let \(E_i\) be conditional expectation onto the exterior sigma-field,
\(D_i=I-E_i\), and let \({\cal B}_i\) be exterior measurable.  Then
\begin{equation}
 [M_{{\cal B}_i},D_i]=0.                                 \label{eq:s6v93-MD-commute}
\end{equation}
Let \(P_i\) be a self-adjoint positive Markov contraction:
\[
 0\le P_i\le I.
\]

\begin{lemma}[Transfer decomposition of the projected defect]
\label{lem:s6v93-decomposition}
\begin{equation}
 (I-P_i)D_i
 =
 D_i(I-P_i)+[D_i,P_i].                                  \label{eq:s6v93-defect}
\end{equation}
\end{lemma}

\begin{proof}
The right side is
\[
 D_i-D_iP_i+D_iP_i-P_iD_i=D_i-P_iD_i.
\]
\end{proof}

\begin{lemma}[Positive-transfer energy inequality]
\label{lem:s6v93-positive}
For \(0\le P_i\le I\),
\begin{equation}
 \|(I-P_i)f\|_2^2
 \le\langle f,(I-P_i)f\rangle.                           \label{eq:s6v93-positive}
\end{equation}
\end{lemma}

\begin{proof}
The spectral theorem reduces the assertion to
\((1-t)^2\le1-t\) for \(0\le t\le1\).
\end{proof}

\subsection{A projected smoothing inequality with commutator debit}
\label{sec:s6v93-projected}

Put
\[
 M_i:=M_{{\cal B}_i},\qquad
 \eta_i:=\|M_iP_i\|_{2\to2}^2.
\]

\begin{theorem}[Tunable transfer-commutator inequality]
\label{thm:s6v93-tunable}
For every \(\theta,\kappa>0\) and every \(f\in L^2(\pi)\),
\begin{multline}
 \|M_iD_if\|_2^2
 \le
 (1+\theta)\eta_i\|D_if\|_2^2\\
 +(1+\theta^{-1})(1+\kappa)
 \langle f,(I-P_i)f\rangle\\
 +(1+\theta^{-1})(1+\kappa^{-1})
 \|M_i[D_i,P_i]f\|_2^2.                                 \label{eq:s6v93-tunable}
\end{multline}
\end{theorem}

\begin{proof}
Write
\[
 M_iD_if=M_iP_iD_if+M_i(I-P_i)D_if.
\]
The weighted square inequality gives
\[
 \|M_iD_if\|^2
 \le
 (1+\theta)\|M_iP_iD_if\|^2
 +(1+\theta^{-1})\|M_i(I-P_i)D_if\|^2.
\]
The first term is at most
\((1+\theta)\eta_i\|D_if\|^2\).  By
\Cref{lem:s6v93-decomposition} and
\eqref{eq:s6v93-MD-commute},
\[
 M_i(I-P_i)D_if
 =
 M_iD_i(I-P_i)f+M_i[D_i,P_i]f.
\]
Apply the weighted square inequality with parameter \(\kappa\).
Both \(M_i\) and \(D_i\) are contractions, so the first resulting
norm is at most \(\|(I-P_i)f\|\).  Finish with
\Cref{lem:s6v93-positive}.
\end{proof}

\begin{remark}[Why this is different from V91]
\label{rem:s6v93-difference}
V91 controlled the second piece by the restricted jump form of
\(D_if\), which led to spatial derivatives of \(D_i\).  V93 moves
\(D_i\) past \(P_i\), leaving the energy of \(f\) and an explicit
commutator.  The price is a new form-small commutator row.
\end{remark}

\subsection{The summed bootstrap}
\label{sec:s6v93-sum}

\begin{definition}[Local transfer and commutator incidence]
\label{def:s6v93-incidence}
For weights \(w_i\ge0\), constants \(m_P,c_D\ge0\), and transfer time
\(\delta>0\), require
\begin{align}
 \sum_iw_i\langle f,(I-P_i)f\rangle
 &\le m_P\delta{\cal E}_\pi(f),                          \label{eq:s6v93-P-incidence}\\
 \sum_iw_i\|M_i[D_i,P_i]f\|_2^2
 &\le c_D\delta{\cal E}_\pi(f).                          \label{eq:s6v93-C-incidence}
\end{align}
\end{definition}

\begin{theorem}[Transfer-commutator good/bad bootstrap]
\label{thm:s6v93-bootstrap}
Assume weighted approximate tensorization
\[
 \operatorname{Var}_\pi(f)\le C_{\rm AT}{\cal S}(f),
 \qquad
 {\cal S}(f):=\sum_iw_i\|D_if\|_2^2,
\]
the good-exterior estimate
\[
 {\cal G}(f)\le
 {m_{\rm inc}\over\lambda_{\rm good}}{\cal E}_\pi(f),
\]
\(\sup_i\eta_i\le\eta\), and
\eqref{eq:s6v93-P-incidence}--\eqref{eq:s6v93-C-incidence}.
For \(\theta,\kappa>0\), set
\begin{align}
 a_\theta&:=1+\theta,                                    \label{eq:s6v93-a}\\
 b_{\theta,\kappa}
 &:=
 (1+\theta^{-1})
 \left\{(1+\kappa)m_P+(1+\kappa^{-1})c_D\right\}.        \label{eq:s6v93-b}
\end{align}
If \(a_\theta\eta<1\), then
\begin{equation}
 \boxed{\displaystyle
 \operatorname{Var}_\pi(f)
 \le
 {C_{\rm AT}\left(
 m_{\rm inc}/\lambda_{\rm good}
 +b_{\theta,\kappa}\delta\right)
 \over1-a_\theta\eta}
 {\cal E}_\pi(f).}                                       \label{eq:s6v93-Poincare}
\end{equation}
Hence
\begin{equation}
 \boxed{\displaystyle
 \lambda_{\rm phys}\ge
 {1-a_\theta\eta\over
 C_{\rm AT}\left(
 m_{\rm inc}/\lambda_{\rm good}
 +b_{\theta,\kappa}\delta\right)}.}                      \label{eq:s6v93-gap}
\end{equation}
\end{theorem}

\begin{proof}
Sum \eqref{eq:s6v93-tunable} with weights \(w_i\) and use the two
incidence bounds:
\[
 {\cal A}(f)
 \le
 a_\theta\eta{\cal S}(f)
 +b_{\theta,\kappa}\delta{\cal E}_\pi(f).
\]
As in V91,
\({\cal S}={\cal G}+{\cal A}\).  Insert the good estimate, absorb
\(a_\theta\eta{\cal S}\), and apply approximate tensorization.
Inversion gives \eqref{eq:s6v93-gap}.
\end{proof}

\begin{corollary}[Optimal auxiliary commutator split]
\label{cor:s6v93-opt-kappa}
For fixed \(\theta\) and positive \(m_P,c_D\), the choice
\[
 \kappa=\sqrt{c_D/m_P}
\]
minimizes the braces in \eqref{eq:s6v93-b}, giving
\begin{equation}
 b_{\theta,*}
 =
 (1+\theta^{-1})(\sqrt{m_P}+\sqrt{c_D})^2.               \label{eq:s6v93-b-opt}
\end{equation}
The evident limiting choices apply if one coefficient vanishes.
\end{corollary}

\begin{proof}
The braces equal
\[
 m_P+c_D+\kappa m_P+\kappa^{-1}c_D.
\]
Arithmetic--geometric mean minimizes the last two terms at
\(2\sqrt{m_Pc_D}\).
\end{proof}

\begin{assessment}[Gain over the factor-two smoothing split]
\label{ass:s6v93-gain}
The V91 condition was \(2\eta<1\).  The tunable V93 condition is
\((1+\theta)\eta<1\).  Thus any \(\eta<1\) can be absorbed by taking
\(0<\theta<\eta^{-1}-1\), at the cost of increasing the energy
coefficient through \(1+\theta^{-1}\).  At Wilson scale
\(\eta_s=e^{-cs+O(\log s)}\), a fixed \(\theta\) is sufficient.
\end{assessment}

\subsection{Ground conjugation of the commutator}
\label{sec:s6v93-conjugation}

Use the unitary \({\cal U}\) of V91 and put
\[
 \widehat E_i:={\cal U}E_i{\cal U}^{-1},\qquad
 \widehat D_i:=I-\widehat E_i.
\]

\begin{lemma}[Exact conjugated commutator]
\label{lem:s6v93-conjugated}
\begin{align}
 {\cal U}[D_i,P]{\cal U}^{-1}
 &=
 s_1^{-2}[\widehat D_i,Q]
 =
 -s_1^{-2}[\widehat E_i,Q],                              \label{eq:s6v93-conjugated}\\
 {\cal U}M_i[D_i,P]{\cal U}^{-1}
 &=
 -s_1^{-2}M_i[\widehat E_i,Q].                           \label{eq:s6v93-restricted-conjugated}
\end{align}
\end{lemma}

\begin{proof}
Conjugation preserves products and commutators, and
\({\cal U}P{\cal U}^{-1}=s_1^{-2}Q\) by
\Cref{thm:s6v91-unitary}.  Also
\(\widehat D_i=I-\widehat E_i\).  The event multiplier commutes with
\({\cal U}\).
\end{proof}

Assume for the next formula that
\(\mu_0=\mu_i\otimes\mu_{i^c}\).  Write \(x=(z,\eta)\), set
\begin{equation}
 Z_i(\eta):=\int u_0(z,\eta)^2\,\mu_i(\dd z),\qquad
 \phi_i^\eta(z):={u_0(z,\eta)\over\sqrt{Z_i(\eta)}}.
                                                               \label{eq:s6v93-phi}
\end{equation}

\begin{proposition}[The conjugated conditional expectation is a
fibre rank-one projection]
\label{prop:s6v93-rank-one}
For \(g\in L^2(\mu_0)\),
\begin{equation}
 (\widehat E_i g)(z,\eta)
 =
 \phi_i^\eta(z)
 \int\phi_i^\eta(z')g(z',\eta)\,\mu_i(\dd z').           \label{eq:s6v93-Ehat}
\end{equation}
Thus \(\widehat E_i\) is the direct integral of the rank-one
projections
\[
 |\phi_i^\eta\rangle\langle\phi_i^\eta|.
\]
\end{proposition}

\begin{proof}
The physical conditional density of \(z\) given \(\eta\) is
\[
 {u_0(z,\eta)^2\over Z_i(\eta)}\,\mu_i(\dd z).
\]
For \(f={\cal U}^{-1}g=g/u_0\),
\[
 E_if(\eta)
 ={1\over Z_i(\eta)}
 \int u_0(z',\eta)g(z',\eta)\,\mu_i(\dd z').
\]
Multiplication by \(u_0(z,\eta)\) gives
\eqref{eq:s6v93-Ehat}.
\end{proof}

\begin{firewall}[The commutator is not ground-free]
\label{fire:s6v93-not-ground-free}
Equation \eqref{eq:s6v93-conjugated} replaces the Doob transfer by
the ground-free operator \(Q/s_1^2\), but
\eqref{eq:s6v93-Ehat} shows that the conjugated conditional
projection still contains the normalized fibre ground vector
\(\phi_i^\eta\).  Therefore V91's unitary cancellation does not, by
itself, prove
\eqref{eq:s6v93-C-incidence}.
\end{firewall}

\subsection{Leakage of the conditional ground subspace}
\label{sec:s6v93-leakage}

\begin{lemma}[Projection commutator equals off-diagonal leakage]
\label{lem:s6v93-leakage}
Let \(E\) be an orthogonal projection, \(D=I-E\), and \(Q=Q^*\).
Then
\begin{equation}
 \|[E,Q]\|=\|DQE\|=\|EQD\|.                              \label{eq:s6v93-leakage}
\end{equation}
\end{lemma}

\begin{proof}
Relative to
\({\cal H}=\operatorname{Ran}E\oplus\operatorname{Ran}D\),
\[
 [E,Q]=
 \begin{pmatrix}
 0&EQD\\
 -DQE&0
 \end{pmatrix}.
\]
The two off-diagonal blocks are adjoints.  Squaring the absolute
value gives diagonal blocks with the same largest singular value,
which proves the norm identity.
\end{proof}

\begin{definition}[Conditional ground-subspace leakage]
\label{def:s6v93-leakage}
For the physical conditional projection, define
\begin{equation}
 \ell_i:=
 s_1^{-2}\|\widehat D_iQ\widehat E_i\|.                  \label{eq:s6v93-ell}
\end{equation}
Its restricted form version is the least \(c_i\) for which
\begin{equation}
 s_1^{-4}
 \|M_i[\widehat E_i,Q]g\|_2^2
 \le
 c_i\delta\,\widehat{\cal E}(g)                          \label{eq:s6v93-form-leakage}
\end{equation}
on the conjugated physical form domain.
\end{definition}

\begin{corollary}[Exact meaning of the commutator row]
\label{cor:s6v93-meaning}
The commutator incidence
\eqref{eq:s6v93-C-incidence} is equivalent, under \({\cal U}\), to
\begin{equation}
 s_1^{-4}\sum_iw_i
 \|M_i[\widehat E_i,Q]g\|_2^2
 \le
 c_D\delta\,\widehat{\cal E}(g),                         \label{eq:s6v93-conjugated-row}
\end{equation}
where
\[
 \widehat{\cal E}(g):=
 {\cal E}_\pi({\cal U}^{-1}g).
\]
It is therefore a summed, restricted conditional-ground-subspace
leakage estimate.
\end{corollary}

\begin{proof}
Apply \eqref{eq:s6v93-restricted-conjugated}, use unitarity, and sum.
\end{proof}

\begin{proposition}[Explicit leakage kernel]
\label{prop:s6v93-leakage-kernel}
Define
\begin{equation}
 K_i(z,\eta;\eta')
 :=
 \int Q((z,\eta),(z',\eta'))\phi_i^{\eta'}(z')\,
       \mu_i(\dd z').                                    \label{eq:s6v93-K}
\end{equation}
Then the kernel of
\(\widehat D_iQ\widehat E_i\), from the exterior coefficient of a
fibre-ground vector at \(\eta'\) to the full output at \((z,\eta)\),
is
\begin{equation}
 K_i^\perp(z,\eta;\eta')
 :=
 K_i(z,\eta;\eta')
 -\phi_i^\eta(z)
  \int\phi_i^\eta(w)K_i(w,\eta;\eta')\,\mu_i(\dd w).
                                                               \label{eq:s6v93-Kperp}
\end{equation}
\end{proposition}

\begin{proof}
Apply \(Q\) to the direct-integral rank-one formula
\eqref{eq:s6v93-Ehat}.  The first term is
\eqref{eq:s6v93-K}.  Applying \(\widehat D_i=I-\widehat E_i\) at the
output subtracts its component parallel to \(\phi_i^\eta\), which is
exactly the second term in \eqref{eq:s6v93-Kperp}.
\end{proof}

\begin{remark}[Factorized consistency check]
\label{rem:s6v93-factorized}
If \(Q=Q_i\otimes Q_{i^c}\) and \(u_0=u_i\otimes u_{i^c}\), with
\(u_i\) the leading eigenvector of \(Q_i\), then
\(\phi_i^\eta\) is independent of \(\eta\) and
\(Q_i\phi_i\) is parallel to \(\phi_i\).  Hence
\[
 K_i^\perp=0,\qquad [\widehat E_i,Q]=0.
\]
Thus the leakage criterion has the correct decoupled limit and does
not create a spectator-volume defect.
\end{remark}

\subsection{Revised acquisition card}

\begin{definition}[Transfer commutator card (TCC)]
\label{def:s6v93-TCC}
TCC consists of:
\begin{enumerate}[label=\textup{(T\arabic*)},leftmargin=4em]
\item positive self-adjoint local comparison transfers \(P_i\);
\item response-uniform restricted smoothing
      \(\|M_iP_i\|^2\le\eta_s<1\);
\item local transfer-energy incidence
      \eqref{eq:s6v93-P-incidence};
\item summed restricted leakage
      \eqref{eq:s6v93-conjugated-row};
\item good-exterior conditional gaps and weighted approximate
      tensorization;
\item boundary ownership, gauge quotient, and cofinal compatibility.
\end{enumerate}
\end{definition}

\begin{theorem}[TCC-to-gap handoff]
\label{thm:s6v93-TCC}
If TCC holds, choose
\[
 0<\theta<\eta_s^{-1}-1
\]
and \(\kappa>0\).  Then the physical finite-volume spectral gap has
the explicit volume-uniform lower bound
\eqref{eq:s6v93-gap}.
\end{theorem}

\begin{proof}
TCC supplies every premise of
\Cref{thm:s6v93-bootstrap}.
\end{proof}

\begin{decision}[V93 bottleneck decision]
\label{dec:s6v93-bottleneck}
The transfer-commutator route avoids spatial derivatives of the
conditional density, but it does not erase the conditional ground
subspace.  Its exact remaining object is the leakage kernel
\eqref{eq:s6v93-Kperp}.  The next iteration must estimate this kernel
from the half-slab action.  In particular, it must decide whether a
local interaction expansion makes \(K_i^\perp\) pay for a connected
polymer crossing the boundary of \(i\), with decay exceeding its
entropy.
\end{decision}

\begin{theorem}[Sixth-series V93 exact advance]
\label{thm:s6v93-decision}
V93 proves:
\begin{enumerate}[leftmargin=3em]
\item the tunable projected transfer inequality
      \eqref{eq:s6v93-tunable};
\item the transfer-commutator bootstrap and gap
      \eqref{eq:s6v93-gap};
\item the exact conjugated commutator
      \eqref{eq:s6v93-conjugated};
\item the fibre rank-one formula
      \eqref{eq:s6v93-Ehat};
\item the equality of commutator norm and conditional-subspace
      leakage;
\item the explicit leakage kernel
      \eqref{eq:s6v93-Kperp}.
\end{enumerate}
Thus the V92 ground-score row has a rigorously equivalent
transfer-geometric alternative.  Its unproved input is a summed
form-small bound for conditional ground-subspace leakage.
\end{theorem}

\begin{proof}
The assertions are
\Cref{thm:s6v93-tunable,
thm:s6v93-bootstrap,
lem:s6v93-conjugated,
prop:s6v93-rank-one,
lem:s6v93-leakage,
prop:s6v93-leakage-kernel}.
The remaining limitation is
\Cref{fire:s6v93-not-ground-free}.
\end{proof}

\begin{programme}[V94: action control of the leakage kernel]
\label{prog:s6v93-V94}
\begin{enumerate}[leftmargin=3em]
\item Insert the two-half-slab formula for \(Q\) into
      \eqref{eq:s6v93-Kperp}.
\item Identify the exact cancellation of components disconnected
      from the block boundary.
\item Formulate a connected-polymer expansion for the remaining
      kernel and retain the singular-value denominator \(s_1^{-2}\).
\item Prove a Hilbert--Schmidt or form bound that sums over blocks
      without a volume factor.
\item If large-\(s\) Wilson character coefficients do not give a
      convergent activity expansion, record that regime obstruction
      and return to the near-flat contour variables rather than
      invoking a strong-coupling series.
\end{enumerate}
\end{programme}

\begin{firewall}[Claim boundary after sixth-series V93]
\label{fire:s6v93-boundary}
V93 proves a transfer-commutator alternative to conditional-gradient
stability, an explicit gap compiler from form-small commutator
incidence, and exact formulas for the conjugated conditional
projection and its leakage kernel.

V93 does not bound the physical leakage kernel, prove TCC,
boundary-star landing, response-uniform restricted smoothing,
good-exterior approximate tensorization, GMPC or HSOCC, the continuum
OS theory, or the full Yang--Mills mass gap.
\end{firewall}

\begin{theorem}[Sixth-series V93 result ledger]
\label{thm:s6v93-results}
The positive conclusions in \Cref{fire:s6v93-boundary} are exactly
those of
\Cref{lem:s6v93-decomposition,
lem:s6v93-positive,
thm:s6v93-tunable,
thm:s6v93-bootstrap,
cor:s6v93-opt-kappa,
lem:s6v93-conjugated,
prop:s6v93-rank-one,
lem:s6v93-leakage,
cor:s6v93-meaning,
prop:s6v93-leakage-kernel,
thm:s6v93-TCC,
thm:s6v93-decision}.
\end{theorem}

\begin{proof}
Each conclusion is the cited result with its positivity,
commutator-incidence, conditional-ground-subspace, restricted-event,
response, and cofinal qualifications unchanged.
\end{proof}

\paragraph{Parent-version record.}
V93 was formed from
\texttt{Renewal\_Yang\_Mills\_Mass\_Gap\_Research\_Notes\_V92.tex}.
The V92 parent had \(42\,972\) logical source lines,
\(1\,501\,257\) bytes, and SHA--256
\[
 \texttt{1da7c1b448c34fb13438f8c338a268f05696359a9fbaf1af2b74c3f9565084c8}.
\]

% =============================================================================
% END APPENDED SIXTH-SERIES V93 MATERIAL
% =============================================================================

% =============================================================================
% BEGIN APPENDED SIXTH-SERIES V94 MATERIAL
% =============================================================================

\section{Sixth-series V94: internalizing frustration in an enlarged
star block}
\label{sec:s6v94}

\subsection{The upstream change of block}

The V83 obstruction arose because one link was updated while its six
staples were frozen in the exterior.  That choice made frustration an
exterior event and forced V91--V93 to compare a spatial conditional
projection with a transfer that moved the frustrating variables.
V94 enlarges the updated block to contain the link and every link in
its incident plaquettes.  The frustration event is then internal to
the fibre.  A conditional Markov transfer on that same fibre commutes
exactly with the conditional fluctuation projection.  This removes
both the ground-score term of V92 and the leakage commutator of V93
from the primary route.

\subsection{The enlarged star is uniformly finite}
\label{sec:s6v94-star}

On the four-dimensional hypercubic lattice, let \(e\) be an
unoriented spatial link and let
\begin{equation}
 \Lambda_e:=
 \bigcup_{p\ni e}E(p),                                   \label{eq:s6v94-star}
\end{equation}
where \(E(p)\) is the set of four links on the plaquette \(p\).

\begin{lemma}[Star size and overlap]
\label{lem:s6v94-star-size}
In four dimensions,
\begin{align}
 |\Lambda_e|&\le 19,                                     \label{eq:s6v94-size}\\
 \sup_b\#\{e:b\in\Lambda_e\}&\le19.                      \label{eq:s6v94-overlap}
\end{align}
The same bounds hold after passing to oriented representatives.
\end{lemma}

\begin{proof}
A link belongs to \(2(4-1)=6\) elementary plaquettes.  Each contributes
at most three links other than \(e\), giving
\[
 |\Lambda_e|\le1+6\cdot3=19.
\]
The relation \(b\in\Lambda_e\) is symmetric in \(b,e\): it says that
the two links coincide or lie on a common plaquette.  Applying the
same count with \(b\) fixed proves \eqref{eq:s6v94-overlap}.
Orientation selects representatives and does not create new
unoriented links.
\end{proof}

Recall the gauge-invariant frustration functional of V86,
\[
 {\mathfrak F}_e
 =
 18-\max_{V\in SU(3)}\Re\Tr(VH_e).
\]

\begin{lemma}[Frustration becomes an internal observable]
\label{lem:s6v94-internal}
The function \({\mathfrak F}_e\) is measurable with respect to the
links in \(\Lambda_e\).  For every threshold
\[
 0<f_*<27/2,
\]
the event
\begin{equation}
 {\cal B}_e:=
 \{{\mathfrak F}_e\ge f_*\}                              \label{eq:s6v94-Be}
\end{equation}
contains the V83 center-staple obstruction and is internal to the
\(\Lambda_e\)-fibre.
\end{lemma}

\begin{proof}
Every staple in \(H_e\) is the ordered product of the other three
links of a plaquette containing \(e\).  Hence \(H_e\), and therefore
\({\mathfrak F}_e\), depends only on \(\Lambda_e\).
\Cref{lem:s6v86-center-F} gives
\({\mathfrak F}_e=27/2\) on the center-staple obstruction.
\end{proof}

\begin{warning}[Internalization is not good-sector coercivity]
\label{warn:s6v94-good}
The complement of \({\cal B}_e\) says that the staple aggregate has a
highly aligned maximizer.  V94 does not infer from this statement
alone a quantitative Poincaré inequality on the whole enlarged star.
That finite-dimensional good-sector estimate remains a separate row.
\end{warning}

\subsection{Conditional fibre transfers commute exactly}
\label{sec:s6v94-fibre}

Disintegrate the physical ground measure over the far exterior:
\begin{equation}
 \mu(\dd x_{\Lambda_e},\dd\eta)
 =
 \mu_e^\eta(\dd x_{\Lambda_e})\,
 \bar\mu_e(\dd\eta).                                    \label{eq:s6v94-disintegration}
\end{equation}
Let \(E_e\) be conditional expectation onto the far exterior and
\[
 D_e:=I-E_e.
\]
For almost every \(\eta\), let \(P_e^\eta\) be a positive
self-adjoint Markov contraction on
\(L^2(\mu_e^\eta)\).  Its direct integral is
\begin{equation}
 (P_ef)(x,\eta)
 :=
 (P_e^\eta f(\,\cdot\,,\eta))(x).                        \label{eq:s6v94-direct-P}
\end{equation}

\begin{theorem}[Exact conditional-transfer commutation]
\label{thm:s6v94-commute}
The direct-integral transfer preserves \(\mu\), is positive and
self-adjoint, and
\begin{equation}
 E_eP_e=P_eE_e=E_e,\qquad
 [D_e,P_e]=0.                                            \label{eq:s6v94-commute}
\end{equation}
\end{theorem}

\begin{proof}
All operator properties hold fibrewise and hence after integration
over \(\bar\mu_e\).  On each fibre, \(P_e^\eta1=1\), so applying
\(P_e\) to a far-exterior-measurable function leaves it unchanged:
\(P_eE_e=E_e\).  Invariance of \(\mu_e^\eta\) gives
\[
 E_eP_ef(\eta)
 =
 \int P_e^\eta f\,\dd\mu_e^\eta
 =
 \int f\,\dd\mu_e^\eta
 =
 E_ef(\eta).
\]
Subtracting these identities from the identity operator proves
\([D_e,P_e]=0\).
\end{proof}

\begin{corollary}[The V93 leakage vanishes for the matched fibre]
\label{cor:s6v94-zero-leakage}
For the enlarged-star conditional transfer,
\[
 M_{{\cal B}_e}[D_e,P_e]=0.
\]
Thus the V93 commutator-incidence constant is \(c_D=0\).
\end{corollary}

\begin{proof}
This is immediate from \eqref{eq:s6v94-commute}.
\end{proof}

\subsection{Direct-integral restricted smoothing}
\label{sec:s6v94-restricted}

Write
\[
 {\cal B}_e^\eta:=
 \{x_{\Lambda_e}:(x_{\Lambda_e},\eta)\in{\cal B}_e\}
\]
and define
\[
 \eta_e^\eta:=
 \|M_{{\cal B}_e^\eta}P_e^\eta\|_{2\to2}^2.
\]

\begin{lemma}[Exact direct-integral norm]
\label{lem:s6v94-direct-norm}
On the standard measurable field of conditional Hilbert spaces,
\begin{equation}
 \|M_{{\cal B}_e}P_e\|_{2\to2}^2
 =
 \operatorname*{ess\,sup}_{\eta}\eta_e^\eta.             \label{eq:s6v94-direct-norm}
\end{equation}
\end{lemma}

\begin{proof}
For every \(f\),
\[
 \|M_{{\cal B}_e}P_ef\|_2^2
 =
 \int
 \|M_{{\cal B}_e^\eta}P_e^\eta f_\eta\|_{L^2(\mu_e^\eta)}^2
 \,\bar\mu_e(\dd\eta),
\]
which gives the upper bound by the essential supremum.  Conversely,
for every \(\varepsilon>0\), restrict to the positive-measure set on
which the fibre norm exceeds the essential supremum minus
\(\varepsilon\), and use a measurable family of unit vectors
approximating the fibre norms.  This gives the reverse inequality
after \(\varepsilon\downarrow0\).
\end{proof}

\begin{corollary}[Fibrewise ground cancellation]
\label{cor:s6v94-fibre-cancel}
If \(P_e^\eta\) is the Doob transform of a positive conditional
half-slab amplitude with squared transfer \(Q_e^\eta\), leading
singular value \(s_{1,e}^\eta\), and conditional ground function
\(u_{0,e}^\eta\), then
\begin{equation}
 \eta_e^\eta
 =
 (s_{1,e}^\eta)^{-4}
 \|M_{{\cal B}_e^\eta}Q_e^\eta\|_{2\to2}^2.              \label{eq:s6v94-fibre-cancel}
\end{equation}
The conditional ground function cancels exactly.
\end{corollary}

\begin{proof}
Apply \Cref{thm:s6v91-unitary} on the fibre.
\end{proof}

\begin{corollary}[Conditional action input]
\label{cor:s6v94-action-input}
Suppose the half-slab action and positive-rectangle estimates of V89
hold on every enlarged-star fibre, uniformly in \(\eta\), with
\begin{equation}
 a_s-2g_s\ge c_*s-C_*\log(2+s).                          \label{eq:s6v94-action-margin}
\end{equation}
Then
\begin{equation}
 \|M_{{\cal B}_e}P_e\|_{2\to2}^2
 \le
 \eta_s:=
 \exp\{-c_*s+C_*\log(2+s)\}.                             \label{eq:s6v94-eta}
\end{equation}
\end{corollary}

\begin{proof}
Apply \Cref{thm:s6v89-action-smoothing} fibrewise with one enlarged
star and use \Cref{lem:s6v94-direct-norm}.
\end{proof}

\begin{firewall}[The essential supremum is substantive]
\label{fire:s6v94-uniform}
An action estimate averaged over \(\eta\) does not imply
\eqref{eq:s6v94-eta}.  The half-slab boundary owner, positive
rectangle, and response-source debit must be uniform over the far
exterior, or the exceptional exteriors must be handled by another
capacity layer.  This is the exact source-specific content still
missing from \Cref{cor:s6v94-action-input}.
\end{firewall}

\subsection{Commuting projected smoothing}
\label{sec:s6v94-smoothing}

\begin{theorem}[Commuting projected-smoothing inequality]
\label{thm:s6v94-smoothing}
Let \(D\) be an orthogonal projection, let \(P\) be a positive
self-adjoint contraction, and assume \([D,P]=0\).  Let \(M\) be an
orthogonal multiplication projection and put
\(\eta=\|MP\|^2\).  Then for every \(\theta>0\),
\begin{equation}
 \boxed{\displaystyle
 \|MDf\|^2
 \le
 (1+\theta)\eta\|Df\|^2
 +(1+\theta^{-1})\langle f,(I-P)f\rangle.}               \label{eq:s6v94-smoothing}
\end{equation}
\end{theorem}

\begin{proof}
Split
\[
 MDf=MPDf+M(I-P)Df
\]
and use the weighted square inequality.  The first term is at most
\(\eta\|Df\|^2\).  For the second,
\[
 \|M(I-P)Df\|^2
 \le\|(I-P)Df\|^2
 \le\langle Df,(I-P)Df\rangle
\]
by \Cref{lem:s6v93-positive}.  Since \(D\) commutes with \(P\),
\[
 D(I-P)D=D(I-P)\le I-P
\]
in operator order.  Hence the last quadratic form is at most
\(\langle f,(I-P)f\rangle\).
\end{proof}

\subsection{The enlarged-block gap compiler}
\label{sec:s6v94-gap}

Define
\[
 {\cal S}_*(f):=\sum_ew_e\|D_ef\|_2^2,
\]
and split it internally:
\begin{align}
 {\cal G}_*(f)&:=
 \sum_ew_e\|M_{{\cal B}_e^c}D_ef\|_2^2,                 \label{eq:s6v94-G}\\
 {\cal A}_*(f)&:=
 \sum_ew_e\|M_{{\cal B}_e}D_ef\|_2^2.                   \label{eq:s6v94-A}
\end{align}
Then
\begin{equation}
 {\cal S}_*={\cal G}_*+{\cal A}_*.                       \label{eq:s6v94-split}
\end{equation}

\begin{theorem}[Enlarged-star projected bootstrap]
\label{thm:s6v94-bootstrap}
Assume:
\begin{enumerate}[label=\textup{(E\arabic*)},leftmargin=4em]
\item enlarged-star approximate tensorization
      \begin{equation}
       \operatorname{Var}_\mu(f)
       \le C_*{\cal S}_*(f);                             \label{eq:s6v94-AT}
      \end{equation}
\item good-sector coercivity and form incidence
      \begin{equation}
       {\cal G}_*(f)
       \le {m_*\over\lambda_*}{\cal E}_\mu(f);           \label{eq:s6v94-good}
      \end{equation}
\item conditional positive transfers satisfying
      \eqref{eq:s6v94-commute} and
      \[
       \sup_e\|M_{{\cal B}_e}P_e\|^2\le\eta;
      \]
\item transfer-form incidence
      \begin{equation}
       \sum_ew_e\langle f,(I-P_e)f\rangle
       \le m_P\delta{\cal E}_\mu(f).                     \label{eq:s6v94-transfer-incidence}
      \end{equation}
\end{enumerate}
For every \(\theta>0\) such that
\[
 (1+\theta)\eta<1,
\]
one has
\begin{equation}
 \boxed{\displaystyle
 \operatorname{Var}_\mu(f)
 \le
 {C_*
 \left\{m_*/\lambda_*+
 (1+\theta^{-1})m_P\delta\right\}
 \over1-(1+\theta)\eta}
 {\cal E}_\mu(f).}                                      \label{eq:s6v94-Poincare}
\end{equation}
Consequently
\begin{equation}
 \boxed{\displaystyle
 \lambda_{\rm phys}\ge
 {1-(1+\theta)\eta\over
 C_*
 \left\{m_*/\lambda_*+
 (1+\theta^{-1})m_P\delta\right\}}.}                    \label{eq:s6v94-gap}
\end{equation}
\end{theorem}

\begin{proof}
Apply \Cref{thm:s6v94-smoothing} to
\((D_e,P_e,M_{{\cal B}_e})\), multiply by \(w_e\), sum, and use
\eqref{eq:s6v94-transfer-incidence}:
\[
 {\cal A}_*(f)
 \le
 (1+\theta)\eta{\cal S}_*(f)
 +(1+\theta^{-1})m_P\delta{\cal E}_\mu(f).
\]
Combine this with
\eqref{eq:s6v94-split} and \eqref{eq:s6v94-good}, absorb the
\({\cal S}_*\) term, and apply
\eqref{eq:s6v94-AT}.  Invert the Poincaré constant for the gap.
\end{proof}

\begin{lemma}[Generator incidence implies transfer incidence]
\label{lem:s6v94-generator}
Suppose
\[
 P_e=e^{-\delta H_e},\qquad H_e\ge0,
\]
and the local generators obey
\begin{equation}
 \sum_ew_e\langle f,H_ef\rangle
 \le m_H{\cal E}_\mu(f).                                \label{eq:s6v94-H-incidence}
\end{equation}
Then \eqref{eq:s6v94-transfer-incidence} holds with \(m_P=m_H\).
\end{lemma}

\begin{proof}
The spectral inequality \(1-e^{-\delta t}\le\delta t\) gives
\[
 \langle f,(I-P_e)f\rangle
 \le\delta\langle f,H_ef\rangle.
\]
Multiply by \(w_e\), sum, and use
\eqref{eq:s6v94-H-incidence}.
\end{proof}

\begin{remark}[The overlap is genuinely bounded]
\label{rem:s6v94-overlap}
If each \(H_e\) is the restriction of the physical local form to
\(\Lambda_e\), the geometric multiplicity entering
\eqref{eq:s6v94-H-incidence} is bounded by
\eqref{eq:s6v94-overlap}.  Exact form domination is still required;
the number \(19\) controls incidence but does not prove that a
conditional half-slab generator is a subform of the physical
Hamiltonian.
\end{remark}

\subsection{Why this block choice avoids the previous two
bottlenecks}

\begin{theorem}[Ground-score and leakage bypass]
\label{thm:s6v94-bypass}
For the enlarged-star fibre construction:
\begin{enumerate}[leftmargin=3em]
\item no spatial derivative of \(E_e\) occurs in
      \Cref{thm:s6v94-smoothing}, so GSIC is unnecessary;
\item \([D_e,P_e]=0\), so the TCC leakage row is identically zero;
\item the singular ground function cancels fibrewise from the
      restricted norm by \eqref{eq:s6v94-fibre-cancel};
\item no sum of uncentered smoothing operators and no MOCE row occurs
      in \Cref{thm:s6v94-bootstrap}.
\end{enumerate}
\end{theorem}

\begin{proof}
The assertions follow respectively from
\Cref{thm:s6v94-smoothing,
thm:s6v94-commute,
cor:s6v94-fibre-cancel,
thm:s6v94-bootstrap}.
\end{proof}

\begin{firewall}[What was traded, not erased]
\label{fire:s6v94-trade}
The enlarged block removes two artificial commutator problems, but it
creates four concrete acquisition rows:
\begin{enumerate}[leftmargin=3em]
\item a uniform conditional half-slab action margin
      \eqref{eq:s6v94-action-margin};
\item good-sector coercivity \eqref{eq:s6v94-good} on the
      gauge-reduced \(19\)-link fibre;
\item enlarged-star approximate tensorization
      \eqref{eq:s6v94-AT};
\item physical generator/form incidence
      \eqref{eq:s6v94-H-incidence}.
\end{enumerate}
None is supplied merely by the geometric cardinality bound.
\end{firewall}

\subsection{Enlarged-star acquisition card}

\begin{definition}[Internal frustrated-star card (IFSC)]
\label{def:s6v94-IFSC}
IFSC consists of:
\begin{enumerate}[label=\textup{(I\arabic*)},leftmargin=4em]
\item the gauge-invariant internal event
      \eqref{eq:s6v94-Be} on the tree-gauge-reduced
      \(\Lambda_e\)-fibre;
\item exact positive conditional half-slab transfers
      \(P_e^\eta\);
\item the uniform restricted norm
      \eqref{eq:s6v94-eta};
\item good-sector coercivity \eqref{eq:s6v94-good};
\item enlarged-star approximate tensorization
      \eqref{eq:s6v94-AT};
\item generator incidence \eqref{eq:s6v94-H-incidence};
\item response-source, boundary-owner, gauge-quotient, and cofinal
      compatibility of all preceding rows.
\end{enumerate}
\end{definition}

\begin{theorem}[IFSC-to-gap handoff]
\label{thm:s6v94-IFSC}
If IFSC holds and \(\eta_s<1\), choose
\[
 0<\theta<\eta_s^{-1}-1.
\]
Then the physical finite-volume Hamiltonian has the explicit
volume-uniform gap \eqref{eq:s6v94-gap}, with \(m_P=m_H\).
\end{theorem}

\begin{proof}
Use \Cref{lem:s6v94-generator} in
\Cref{thm:s6v94-bootstrap}.
\end{proof}

\begin{decision}[V94 bottleneck decision]
\label{dec:s6v94-bottleneck}
The enlarged-star construction is now the primary route.  It bypasses
MOCE, GSIC, and the TCC leakage row.  The next compulsory-audit
iteration must compare the current companion notes against the four
rows in \Cref{fire:s6v94-trade}.  After that audit, the first direct
target is the finite-dimensional good-sector coercivity of the
tree-gauge-reduced enlarged star, while the complete-action
boundary-owner row remains separately visible.
\end{decision}

\begin{theorem}[Sixth-series V94 exact advance]
\label{thm:s6v94-decision}
V94 proves:
\begin{enumerate}[leftmargin=3em]
\item the \(19\)-link size and overlap bounds for enlarged stars;
\item internalization of the V83 gauge-invariant frustration event;
\item exact commutation of conditional fibre transfer and
      conditional fluctuation;
\item the direct-integral restricted-norm identity;
\item the commuting projected-smoothing inequality;
\item the enlarged-star gap compiler
      \eqref{eq:s6v94-gap};
\item generator-to-transfer form incidence;
\item the precise bypass of the V90--V93 operator, score, and leakage
      rows.
\end{enumerate}
\end{theorem}

\begin{proof}
The claims are
\Cref{lem:s6v94-star-size,
lem:s6v94-internal,
thm:s6v94-commute,
lem:s6v94-direct-norm,
thm:s6v94-smoothing,
thm:s6v94-bootstrap,
lem:s6v94-generator,
thm:s6v94-bypass}.
The acquired hypotheses remain delimited by
\Cref{fire:s6v94-uniform,fire:s6v94-trade}.
\end{proof}

\begin{programme}[V95: compulsory audit and enlarged-star coercivity]
\label{prog:s6v94-V95}
\begin{enumerate}[leftmargin=3em]
\item Audit the active SM and NS notes against the four IFSC rows in
      \Cref{fire:s6v94-trade}.
\item Fix a tree gauge on \(\Lambda_e\) and count the remaining compact
      SU(3) coordinates.
\item Derive a quantitative good-sector lower Hessian or path
      inequality from \({\mathfrak F}_e<f_*\), without assuming a
      unique coordinate representative.
\item Track the dependence on the far exterior and response sources.
\item If frustration below \(f_*\) does not imply coercivity, refine
      the internal event by the smallest additional
      gauge-invariant degeneracy discriminant.
\end{enumerate}
\end{programme}

\begin{firewall}[Claim boundary after sixth-series V94]
\label{fire:s6v94-boundary}
V94 proves an enlarged-star internalization and a commuting
projected-smoothing gap compiler.  This rigorously removes the
uncentered operator row, spatial ground-score, and conditional
subspace leakage from the primary proof path.

V94 does not prove IFSC's uniform action margin, good-sector
coercivity, approximate tensorization, or generator incidence; does
not populate GMPC or HSOCC; does not construct the continuum OS
theory; and does not prove the full Yang--Mills mass gap.
\end{firewall}

\begin{theorem}[Sixth-series V94 result ledger]
\label{thm:s6v94-results}
The positive conclusions in \Cref{fire:s6v94-boundary} are exactly
those of
\Cref{lem:s6v94-star-size,
lem:s6v94-internal,
thm:s6v94-commute,
cor:s6v94-zero-leakage,
lem:s6v94-direct-norm,
cor:s6v94-fibre-cancel,
cor:s6v94-action-input,
thm:s6v94-smoothing,
thm:s6v94-bootstrap,
lem:s6v94-generator,
thm:s6v94-bypass,
thm:s6v94-IFSC,
thm:s6v94-decision}.
\end{theorem}

\begin{proof}
Each conclusion is the cited result with its fibrewise,
essential-supremum, good-sector, form-incidence, boundary-owner,
response, and cofinal qualifications unchanged.
\end{proof}

\paragraph{Parent-version record.}
V94 was formed from
\texttt{Renewal\_Yang\_Mills\_Mass\_Gap\_Research\_Notes\_V93.tex}.
The V93 parent had \(43\,541\) logical source lines,
\(1\,518\,469\) bytes, and SHA--256
\[
 \texttt{0c7e7a56028cf8429a503513d4371ee9279e4aaef6acba7939fe0d6a7b39bb7b}.
\]

% =============================================================================
% END APPENDED SIXTH-SERIES V94 MATERIAL
% =============================================================================

% =============================================================================
% BEGIN APPENDED SIXTH-SERIES V95 MATERIAL
% =============================================================================

\section{Sixth-series V95: compulsory audit and quantitative
aligned-staple geometry}
\label{sec:s6v95}

\subsection{Compulsory companion audit}
\label{sec:s6v95-audit}

\begin{audit}[Companion files inspected at V95]
\label{audit:s6v95-files}
The current active companion notes are
\begin{align*}
 &\texttt{Renewal\_SM\_Einstein\_Continuum\_Research\_Notes\_V34.tex},
 \\
 &\texttt{Renewal\_NS\_Stability\_Research\_Notes\_V26.tex}.
\end{align*}
At the audit time the SM file had \(13\,253\) logical source lines,
\(486\,331\) bytes, and SHA--256
\[
 \texttt{4b4ac9f76ca2600fab3dc1b6c89394124d155bf5de08306d9b7f4f4e868ba2e5}.
\]
The NS file had \(5\,319\) logical source lines, \(278\,283\) bytes,
and SHA--256
\[
 \texttt{82c887f8c2c69e4f28fad7c3dc8de5b9099cb5a4bf431eba1fff3e91935978ad}.
\]
Thus the SM programme advanced from V29 to V34 since the V90 audit,
while the NS file is unchanged.
\end{audit}

\begin{theorem}[Relevant SM V30--V34 imports]
\label{thm:s6v95-SM-audit}
The audited SM additions establish four abstract facts relevant to
IFSC.
\begin{enumerate}[leftmargin=3em]
\item A transfer which leaves the bad carrier variable fixed has
      restricted norm one; the smoother must move the bad variable.
\item For an exact conditional heat-bath projection \(Q\) and an
      internal bad event with conditional probability \(p(\eta)\),
      \[
       \|M_{\cal B}Q\|^2
       =\operatorname*{ess\,sup}_{\eta}p(\eta).
      \]
      Its finite-time interpolation has fibre norm
      \(p+(1-p)e^{-2t}\).
\item Conditional variance is one half of the independent
      replacement square, giving an exact route to update-incidence
      estimates.
\item Orthogonal local coefficient-space summands can replace a
      dangerous sum of per-cell relative constants by their weighted
      supremum.
\end{enumerate}
These facts support V94's decision to update the entire frustrated
star and to demand a worst-far-exterior fibre estimate.
\end{theorem}

\begin{proof}
These are respectively the results titled
\emph{A fibrewise quantum semigroup cannot smooth a metric event}
(SM V30), \emph{Exact restricted norm of a heat-bath projection} and
\emph{Exact finite-time restricted norm} (SM V31),
\emph{Conditional replacement identity} (SM V32), and
\emph{Orthogonal-cell weighted incidence} (SM V34).  Their stated
positivity, essential-supremum, replacement, and fixed-loader
qualifications are retained.
\end{proof}

\begin{assessment}[What the SM audit does not populate]
\label{ass:s6v95-SM-limit}
The SM heat-bath projection can be applied abstractly to an enlarged
Yang--Mills star, but its defect form is the enlarged-star conditional
variance itself.  Domination of that defect by the physical
Hamiltonian form would already be a local Poincaré comparison and is
not supplied by the heat-bath identity.  The SM notes also contain no
SU(3) good-sector coercivity, Wilson boundary-owner estimate, or
response-matched ground law.  Hence they validate the architecture but
do not populate the four open rows in
\Cref{fire:s6v94-trade}.
\end{assessment}

\begin{assessment}[NS audit consequence]
\label{ass:s6v95-NS}
The NS V26 file is unchanged.  Its rank-one Oseen refinements do not
provide a compact-group conditional gap, a positive enlarged-star
transfer, an approximate-tensorization constant, or a Yang--Mills
form-incidence theorem.  No NS result changes the IFSC acquisition
order.
\end{assessment}

\begin{decision}[V95 audit decision]
\label{dec:s6v95-audit}
The enlarged-star route remains primary.  The SM audit strengthens
two safeguards: the event must be internal to the updated fibre, and
all activity bounds must be uniform over its far exterior.  V95 now
attacks the bare-Wilson good-sector geometry; physical ground and
response corrections remain separate.
\end{decision}

\subsection{Small frustration forces quantitative staple alignment}
\label{sec:s6v95-alignment}

Let \(S_1,\ldots,S_m\in SU(3)\), put
\[
 H:=\sum_{j=1}^mS_j,
\]
and define
\begin{align}
 \Phi_H(U)&:=\Re\Tr(UH),                                 \label{eq:s6v95-Phi}\\
 {\mathfrak F}(S_1,\ldots,S_m)
 &:=3m-\max_{V\in SU(3)}\Phi_H(V).                       \label{eq:s6v95-F}
\end{align}
For \(W\in SU(3)\), let
\[
 d(W):=3-\Re\Tr W
 ={1\over2}\|W-I\|_{\rm F}^2.                            \label{eq:s6v95-d}
\]

\begin{lemma}[Exact alignment budget]
\label{lem:s6v95-budget}
Let \(V_0\) be any maximizer of \(\Phi_H\) and put
\[
 W_j:=V_0S_j,\qquad M:=V_0H.
\]
Then
\begin{align}
 \sum_{j=1}^m d(W_j)&={\mathfrak F},                     \label{eq:s6v95-budget}\\
 \sum_{j=1}^m\|W_j-I\|_{\rm F}^2&=2{\mathfrak F},        \label{eq:s6v95-square-budget}\\
 \|M-mI\|_{\rm F}
 &\le\sqrt{2m{\mathfrak F}}.                            \label{eq:s6v95-M-close}
\end{align}
\end{lemma}

\begin{proof}
Since \(M=\sum_jW_j\),
\[
 {\mathfrak F}
 =3m-\Re\Tr M
 =\sum_j(3-\Re\Tr W_j)
 =\sum_jd(W_j).
\]
Equation \eqref{eq:s6v95-square-budget} follows from
\eqref{eq:s6v95-d}.  Finally,
\[
 \|M-mI\|_{\rm F}
 \le\sum_j\|W_j-I\|_{\rm F}
 \le\sqrt{m\sum_j\|W_j-I\|_{\rm F}^2},
\]
which is \eqref{eq:s6v95-M-close}.
\end{proof}

\begin{corollary}[Pairwise staple alignment]
\label{cor:s6v95-pairwise}
For every \(j,k\),
\begin{equation}
 \|S_j-S_k\|_{\rm F}
 \le2\sqrt{2{\mathfrak F}}.                              \label{eq:s6v95-pairwise}
\end{equation}
\end{corollary}

\begin{proof}
Unitary invariance and the triangle inequality give
\[
 \|S_j-S_k\|_{\rm F}
 =\|W_j-W_k\|_{\rm F}
 \le\|W_j-I\|_{\rm F}+\|W_k-I\|_{\rm F}.
\]
Each squared norm is at most the sum
\(2{\mathfrak F}\).
\end{proof}

\subsection{Curvature at a maximizing link}
\label{sec:s6v95-curvature}

Put
\begin{equation}
 \epsilon:=\sqrt{2m{\mathfrak F}},\qquad
 A:={1\over2}(M+M^*).                                   \label{eq:s6v95-epsilon}
\end{equation}

\begin{lemma}[Positive Hermitian polar part]
\label{lem:s6v95-A}
If \(\epsilon<m\), then
\begin{equation}
 A\succeq(m-\epsilon)I.                                 \label{eq:s6v95-A-lower}
\end{equation}
\end{lemma}

\begin{proof}
\Cref{lem:s6v95-budget} gives
\(\|M-mI\|_{\rm op}\le\epsilon\).  Taking Hermitian parts does not
increase the operator norm, so
\(\|A-mI\|_{\rm op}\le\epsilon\).  The spectral bound follows.
\end{proof}

\begin{proposition}[Strict Hessian at every maximizer]
\label{prop:s6v95-Hessian}
For \(X\in\mathfrak{su}(3)\), let
\[
 \gamma_X(t):=e^{tX}V_0.
\]
Then
\begin{align}
 {d\over dt}\Phi_H(\gamma_X(t))\bigg|_{t=0}&=0,          \label{eq:s6v95-first}\\
 -{d^2\over dt^2}\Phi_H(\gamma_X(t))\bigg|_{t=0}
 &\ge(m-\epsilon)\|X\|_{\rm F}^2.                       \label{eq:s6v95-Hessian}
\end{align}
Thus every maximizer is nondegenerate whenever
\({\mathfrak F}<m/2\).
\end{proposition}

\begin{proof}
The first derivative vanishes because \(V_0\) is a maximizer.
The second derivative is
\[
 \Re\Tr(X^2M)=\Tr(X^2A),
\]
because \(X^2\) is Hermitian and the anti-Hermitian part of \(M\)
contributes only an imaginary trace.  Since \(-X^2\succeq0\),
\Cref{lem:s6v95-A} gives
\[
 -\Tr(X^2A)
 =\Tr((-X^2)A)
 \ge(m-\epsilon)\Tr(-X^2)
 =(m-\epsilon)\|X\|_{\rm F}^2.
\]
\end{proof}

\subsection{A universal unique-maximizer threshold}
\label{sec:s6v95-unique}

Fix constants \(\rho_{\exp}>0\) and \(C_{\exp}\ge1\) with the following
standard local property of \(SU(3)\): if
\(\|W-I\|_{\rm F}\le\rho_{\exp}\), then
\[
 W=e^X,\qquad X\in\mathfrak{su}(3),\qquad
 \|X\|_{\rm F}\le C_{\exp}\|W-I\|_{\rm F}.
\]
Such constants exist because the exponential map is a diffeomorphism
near the identity.  Define
\begin{align}
 r_{\rm cvx}
 &:=
 \min\left\{1,{1\over4e(\sqrt3+1/2)}\right\},             \label{eq:s6v95-rcvx}\\
 \epsilon_0(m)
 &:=
 m\min\left\{
 {1\over2},{\rho_{\exp}\over4},
 {r_{\rm cvx}\over4C_{\exp}}\right\},                    \label{eq:s6v95-eps0}\\
 f_{\rm al}(m)
 &:={\epsilon_0(m)^2\over2m}.                            \label{eq:s6v95-fal}
\end{align}

\begin{theorem}[Small frustration gives one global matrix-Fisher well]
\label{thm:s6v95-one-well}
If
\begin{equation}
 {\mathfrak F}<f_{\rm al}(m),                            \label{eq:s6v95-smallF}
\end{equation}
then \(\Phi_H\) has a unique global maximizer \(V_0\).  Moreover,
for every \(U\in SU(3)\),
\begin{equation}
 \boxed{\displaystyle
 \Phi_H(V_0)-\Phi_H(U)
 \ge {m\over8e^2}
 \|UV_0^{-1}-I\|_{\rm F}^2.}                            \label{eq:s6v95-global-quadratic}
\end{equation}
\end{theorem}

\begin{proof}
Put \(M=V_0H=mI+E\), so
\(\|E\|_{\rm F}\le\epsilon\) by
\eqref{eq:s6v95-M-close}.  For
\(W=UV_0^{-1}\), unitary invariance gives
\begin{align}
 \Phi_H(V_0)-\Phi_H(U)
 &=\Re\Tr((I-W)M)\\
 &\ge {m\over2}\|I-W\|_{\rm F}^2
      -\epsilon\|I-W\|_{\rm F}.                         \label{eq:s6v95-far}
\end{align}
Therefore, whenever
\(\|I-W\|_{\rm F}\ge4\epsilon/m\),
\begin{equation}
 \Phi_H(V_0)-\Phi_H(U)
 \ge {m\over4}\|I-W\|_{\rm F}^2.                        \label{eq:s6v95-far2}
\end{equation}

It remains to treat
\(\|I-W\|_{\rm F}<4\epsilon/m\).  The choice
\eqref{eq:s6v95-eps0} puts \(W\) in the exponential chart and writes
\(W=e^X\) with \(\|X\|_{\rm F}\le r_{\rm cvx}\).  Set
\[
 \varphi(t):=\Re\Tr(e^{tX}M).
\]
We have \(\varphi'(0)=0\) by
\eqref{eq:s6v95-first} and
\[
 -\varphi''(0)\ge(m-\epsilon)\|X\|_{\rm F}^2.
\]
For \(0\le t\le1\),
\begin{align*}
 |\varphi''(t)-\varphi''(0)|
 &=
 |\Re\Tr(X^2(e^{tX}-I)M)|\\
 &\le
 \|X^2\|_{\rm F}\|e^{tX}-I\|_{\rm op}\|M\|_{\rm F}\\
 &\le
 e\|M\|_{\rm F}\|X\|_{\rm F}^3.
\end{align*}
Here
\(\|M\|_{\rm F}\le m\sqrt3+\epsilon
\le m(\sqrt3+1/2)\), while
\(m-\epsilon\ge m/2\).  The definition of \(r_{\rm cvx}\) therefore
gives
\[
 \varphi''(t)
 \le-{m\over4}\|X\|_{\rm F}^2.
\]
Twice integrating and using \(\varphi'(0)=0\) yields
\[
 \Phi_H(V_0)-\Phi_H(U)
 =\varphi(0)-\varphi(1)
 \ge {m\over8}\|X\|_{\rm F}^2.
\]
Since \(\|e^X-I\|_{\rm F}\le e\|X\|_{\rm F}\) for
\(\|X\|_{\rm F}\le1\),
\begin{equation}
 \Phi_H(V_0)-\Phi_H(U)
 \ge {m\over8e^2}\|I-W\|_{\rm F}^2.                     \label{eq:s6v95-near}
\end{equation}
The near and far regions cover \(SU(3)\), and
\eqref{eq:s6v95-global-quadratic} follows from
\eqref{eq:s6v95-far2,eq:s6v95-near}.  Equality in the global maximum
is possible only when \(W=I\), proving uniqueness.
\end{proof}

\begin{corollary}[Six-staple good threshold]
\label{cor:s6v95-six}
For a four-dimensional Wilson link, \(m=6\).  If the V94 threshold is
chosen so that
\[
 0<f_*<f_{\rm al}(6),
\]
then every configuration in \({\cal B}_e^c\) has a unique
matrix-Fisher maximizer and obeys
\begin{equation}
 \max_U\Re\Tr(UH_e)-\Re\Tr(VH_e)
 \ge {3\over4e^2}
 \|VV_0^{-1}-I\|_{\rm F}^2.                             \label{eq:s6v95-six-quadratic}
\end{equation}
\end{corollary}

\begin{proof}
Apply \Cref{thm:s6v95-one-well} with \(m=6\).
\end{proof}

\subsection{Bare-Wilson concentration in the good sector}
\label{sec:s6v95-concentration}

\begin{proposition}[One-link concentration from aligned staples]
\label{prop:s6v95-concentration}
Let
\[
 \nu_{s,H}(\dd U)
 ={e^{s\Phi_H(U)}\over Z_{s,H}}\,\dd U
\]
with normalized Haar measure.  Under
\eqref{eq:s6v95-smallF}, for every fixed \(\rho>0\) there are
constants \(C_\rho,C_0<\infty\), independent of \(s\ge1\) and of the
staples satisfying the same frustration threshold, such that
\begin{equation}
 \nu_{s,H}\!\left(
 \|UV_0^{-1}-I\|_{\rm F}\ge\rho\right)
 \le
 C_\rho s^4
 \exp\left\{
 -{m\rho^2\over8e^2}s+C_0
 \right\}.                                               \label{eq:s6v95-concentration}
\end{equation}
\end{proposition}

\begin{proof}
\Cref{thm:s6v95-one-well} bounds the numerator of the indicated event
by
\[
 \exp\left\{
 s\Phi_H(V_0)-{m\rho^2\over8e^2}s
 \right\}.
\]
For the denominator, use exponential coordinates
\(U=e^XV_0\) with
\(\|X\|_{\rm F}\le s^{-1/2}\).  The first derivative at the maximum
vanishes, and the second derivative is bounded in absolute value by a
uniform constant depending only on \(m\) and the fixed frustration
threshold.  Hence
\[
 \Phi_H(e^XV_0)
 \ge\Phi_H(V_0)-C\|X\|_{\rm F}^2.
\]
Haar measure in an eight-dimensional exponential chart has a
Jacobian bounded below on a fixed ball, so the ball of radius
\(s^{-1/2}\) has measure at least \(c s^{-4}\).  Therefore
\[
 Z_{s,H}\ge c s^{-4}e^{s\Phi_H(V_0)-C}.
\]
Divide the numerator bound by this denominator and absorb normalized
Haar volume into \(C_\rho\).
\end{proof}

\begin{firewall}[Concentration is not the good-sector Poincaré row]
\label{fire:s6v95-concentration}
Equation \eqref{eq:s6v95-concentration} proves a single classical
well and exponential tail for the \emph{bare matrix-Fisher
conditional}.  It does not by itself prove a Poincaré inequality for
arbitrary \(L^2\) functions, and it does not include the physical
conditional factor
\[
 u_0(U,\xi)^2
\]
identified in V87.  A response-matched ground profile can alter both
the curvature and the tail unless its variation is controlled.
\end{firewall}

\subsection{What has been populated in IFSC}

\begin{theorem}[Bare good-geometry row]
\label{thm:s6v95-bare-good}
For \(f_*<f_{\rm al}(6)\), the bare Wilson part of IFSC's good sector
has:
\begin{enumerate}[leftmargin=3em]
\item pairwise staple alignment
      \eqref{eq:s6v95-pairwise};
\item a unique one-link maximizer;
\item a volume- and exterior-independent quadratic action deficit
      \eqref{eq:s6v95-six-quadratic};
\item the concentration estimate
      \eqref{eq:s6v95-concentration}.
\end{enumerate}
Thus no additional polar-degeneracy discriminant is needed for the
bare six-staple phase once \(f_*\) is chosen below the universal
threshold.
\end{theorem}

\begin{proof}
Use
\Cref{cor:s6v95-pairwise,
thm:s6v95-one-well,
cor:s6v95-six,
prop:s6v95-concentration}.
\end{proof}

\begin{warning}[The enlarged fibre contains more than the central link]
\label{warn:s6v95-fibre}
\Cref{thm:s6v95-bare-good} controls the central one-link phase for
fixed staples.  IFSC row (I4) is a coercive inequality for arbitrary
fluctuations of the whole tree-gauge-reduced enlarged star.
Sequentially resampling the remaining links may leave the good set or
change the aligned maximizer.  A block path or conditional
factorization theorem is still required.
\end{warning}

\begin{definition}[Aligned-star coercivity card (ASCC)]
\label{def:s6v95-ASCC}
ASCC consists of:
\begin{enumerate}[label=\textup{(A\arabic*)},leftmargin=4em]
\item \(f_*<f_{\rm al}(6)\);
\item a tree-gauge chart for the enlarged star with uniformly
      controlled Haar Jacobian;
\item a measurable partition of unity subordinate to aligned
      corridors, with overlap at most \(m_{\rm cor}\);
\item a one-well Poincaré inequality with constant
      \(\lambda_{\rm cor}\) for the fluctuation about each corridor
      mean, using \eqref{eq:s6v95-six-quadratic};
\item a corridor-mean gluing and exit-capacity estimate which, after
      summing over enlarged stars, gives
      \begin{equation}
       {\cal G}_*(f)
       \le
       C_{\rm cor}{\cal E}_\mu(f)
       +\varepsilon_{\rm exit}{\cal S}_*(f),             \label{eq:s6v95-gluing}
      \end{equation}
      where \(C_{\rm cor}\) is explicitly assembled from
      \(m_{\rm cor}/\lambda_{\rm cor}\), the corridor-mean graph, and
      the exit energy debit;
\item stability of \eqref{eq:s6v95-gluing} under the physical ground
      profile, response sources, far exterior, and cofinal bridges.
\end{enumerate}
\end{definition}

\begin{theorem}[ASCC gives the generalized enlarged-star gap]
\label{thm:s6v95-ASCC}
Assume ASCC, V94 enlarged-star approximate tensorization, the
restricted norm \(\eta\), and transfer-form incidence.  If
\begin{equation}
 \varepsilon_{\rm exit}+(1+\theta)\eta<1,                \label{eq:s6v95-ASCC-window}
\end{equation}
then
\begin{equation}
 \operatorname{Var}_\mu(f)
 \le
 {C_*
 \left\{C_{\rm cor}+
 (1+\theta^{-1})m_P\delta\right\}
 \over
 1-\varepsilon_{\rm exit}-(1+\theta)\eta}
 {\cal E}_\mu(f).                                       \label{eq:s6v95-ASCC-gap}
\end{equation}
\end{theorem}

\begin{proof}
Use \eqref{eq:s6v95-gluing} for the good term and
\Cref{thm:s6v94-smoothing} plus transfer-form incidence for the bad
term:
\[
 {\cal S}_*(f)
 \le
 \left\{\varepsilon_{\rm exit}+(1+\theta)\eta\right\}
 {\cal S}_*(f)
 +\left\{C_{\rm cor}+
 (1+\theta^{-1})m_P\delta\right\}{\cal E}_\mu(f).
\]
Absorb using \eqref{eq:s6v95-ASCC-window}, then apply enlarged-star
approximate tensorization.
\end{proof}

\begin{firewall}[ASCC is a card, not a populated theorem]
\label{fire:s6v95-ASCC}
\Cref{thm:s6v95-ASCC} is an implication from the stated corridor and
mean-gluing rows.  V95 has populated only ASCC row (A1) and the bare
one-link geometric input to row (A4).  It has not constructed a
physical corridor partition, controlled the corridor means, or proved
the exit-capacity term in \eqref{eq:s6v95-gluing}.
\end{firewall}

\begin{decision}[V95 bottleneck decision]
\label{dec:s6v95-bottleneck}
The bare polar-degeneracy question is resolved by
\Cref{thm:s6v95-one-well}.  The remaining good-sector bottleneck is
not uniqueness of the central matrix-Fisher maximizer; it is transport
of arbitrary \(19\)-link fibre fluctuations through aligned corridors
while retaining the physical ground and response factors.  V96 must
start with the exact tree-gauge dimension and a sequential variance
decomposition, and must stop if a step leaves the aligned corridor.
\end{decision}

\begin{theorem}[Sixth-series V95 exact advance]
\label{thm:s6v95-decision}
V95:
\begin{enumerate}[leftmargin=3em]
\item performs the compulsory SM V34/NS V26 audit;
\item proves the exact frustration-to-alignment budget;
\item proves positive Hessian at every matrix-Fisher maximizer when
      \({\mathfrak F}<m/2\);
\item proves a universal small-frustration threshold with a unique
      global maximizer and global quadratic deficit;
\item proves exponential bare one-link concentration with the correct
      \(s^4\) eight-dimensional prefactor;
\item identifies the remaining enlarged-fibre corridor and physical
      ground-profile rows.
\end{enumerate}
\end{theorem}

\begin{proof}
The conclusions are
\Cref{audit:s6v95-files,
thm:s6v95-SM-audit,
lem:s6v95-budget,
prop:s6v95-Hessian,
thm:s6v95-one-well,
prop:s6v95-concentration,
warn:s6v95-fibre,
fire:s6v95-ASCC}.
\end{proof}

\begin{programme}[V96: tree gauge and sequential aligned corridors]
\label{prog:s6v95-V96}
\begin{enumerate}[leftmargin=3em]
\item Count vertices and independent links of
      \(\Lambda_e\), choose an explicit spanning tree, and record all
      residual boundary gauge actions.
\item Write an exact martingale variance decomposition along the
      remaining SU(3) link coordinates.
\item Determine which conditional steps inherit the alignment budget
      and which can leave \({\cal B}_e^c\).
\item For the stable steps, apply the one-well geometry of V95; for
      unstable steps, define the smallest corridor-exit event.
\item Keep the physical ground-profile and response-source
      perturbations explicit in every conditional density.
\end{enumerate}
\end{programme}

\begin{firewall}[Claim boundary after sixth-series V95]
\label{fire:s6v95-boundary}
V95 performs the compulsory companion audit and proves quantitative
bare-Wilson aligned-staple geometry, including a universal
unique-maximizer threshold, a global quadratic action deficit, and
one-link concentration.

V95 does not prove enlarged-star physical good-sector coercivity,
uniform response-matched action smoothing, enlarged-star approximate
tensorization or form incidence, GMPC or HSOCC, the continuum OS
theory, or the full Yang--Mills mass gap.
\end{firewall}

\begin{theorem}[Sixth-series V95 result ledger]
\label{thm:s6v95-results}
The positive conclusions in \Cref{fire:s6v95-boundary} are exactly
those of
\Cref{audit:s6v95-files,
thm:s6v95-SM-audit,
lem:s6v95-budget,
cor:s6v95-pairwise,
lem:s6v95-A,
prop:s6v95-Hessian,
thm:s6v95-one-well,
cor:s6v95-six,
prop:s6v95-concentration,
thm:s6v95-bare-good,
thm:s6v95-ASCC,
thm:s6v95-decision}.
\end{theorem}

\begin{proof}
Each conclusion is the cited result with its bare-Wilson,
small-frustration, local-chart, corridor, physical-ground, response,
and cofinal qualifications unchanged.
\end{proof}

\paragraph{Parent-version record.}
V95 was formed from
\texttt{Renewal\_Yang\_Mills\_Mass\_Gap\_Research\_Notes\_V94.tex}.
The V94 parent had \(44\,139\) logical source lines,
\(1\,537\,433\) bytes, and SHA--256
\[
 \texttt{7c45ca13ccb5556d35bf6fc7e5fcdc5e0cbf3614398c11c5f33dd932c1021506}.
\]
The compulsory audit used the active companion files recorded in
\Cref{audit:s6v95-files}.

% =============================================================================
% END APPENDED SIXTH-SERIES V95 MATERIAL
% =============================================================================

% =============================================================================
% BEGIN APPENDED SIXTH-SERIES V96 MATERIAL
% =============================================================================

\section{Sixth-series V96: tree-gauge reduction and common-relative
star coordinates}
\label{sec:s6v96}

\subsection{The exact enlarged-star graph}
\label{sec:s6v96-graph}

Choose coordinates so the central link joins \(0\) to \(e_1\).  For
each transverse direction \(\nu\in\{2,3,4\}\) and sign
\(\sigma\in\{-1,+1\}\), the incident plaquette has vertices
\[
 0,\quad e_1,\quad \sigma e_\nu,\quad e_1+\sigma e_\nu.
\]
The six plaquettes share only the central link.

\begin{lemma}[Exact star counts]
\label{lem:s6v96-counts}
The enlarged-star graph has
\begin{equation}
 |V(\Lambda_e)|=14,\qquad
 |E(\Lambda_e)|=19,\qquad
 b_1(\Lambda_e)=6.                                      \label{eq:s6v96-counts}
\end{equation}
Here \(b_1=|E|-|V|+1\) is the cycle rank.
\end{lemma}

\begin{proof}
There are two central endpoints and two new vertices for each of six
plaquettes, hence \(2+2\cdot6=14\) vertices.  Besides the common
central edge, each plaquette contributes three distinct edges, hence
\(1+3\cdot6=19\) edges.  The graph is connected, so its cycle rank is
\(19-14+1=6\).
\end{proof}

\begin{lemma}[An explicit spanning tree]
\label{lem:s6v96-tree}
Let \(T_e\) contain the central edge and, for every
\((\nu,\sigma)\), the two edges
\[
 (0,\sigma e_\nu),\qquad
 (\sigma e_\nu,e_1+\sigma e_\nu).
\]
Then \(T_e\) is a spanning tree with \(13\) edges.  The six omitted
edges
\[
 (e_1+\sigma e_\nu,e_1)
\]
are one chord per incident plaquette.
\end{lemma}

\begin{proof}
Every outer vertex is joined to \(0\) by the displayed arms, and
\(e_1\) is joined to \(0\) by the central edge, so the graph is
connected.  It has \(1+2\cdot6=13=14-1\) edges and therefore is a
tree.  The omitted edge closes exactly its corresponding plaquette
cycle.
\end{proof}

\subsection{Tree gauge and the boundary-fibre qualification}
\label{sec:s6v96-gauge}

\begin{proposition}[Isolated-star tree gauge]
\label{prop:s6v96-tree-gauge}
On the isolated graph \(\Lambda_e\), vertex gauge transformations set
all \(13\) tree links to the identity.  The remaining variables are
six chord holonomies
\[
 S_1,\ldots,S_6\in SU(3).
\]
The residual root gauge acts by simultaneous conjugation,
\[
 (S_1,\ldots,S_6)\longmapsto
 (gS_1g^{-1},\ldots,gS_6g^{-1}).
\]
Thus the generic gauge-reduced dimension is
\begin{equation}
 6\dim SU(3)-\dim SU(3)=40.                              \label{eq:s6v96-dimension}
\end{equation}
With a fixed root frame, the six-chord coordinate space has dimension
\(48\).
\end{proposition}

\begin{proof}
Fix a root.  Recursively along the unique tree paths, choose the gauge
at each nonroot vertex so that its incoming tree link becomes the
identity.  This uses \(14-1=13\) independent vertex gauges and leaves
only the root transformation.  A chord transforms at its two
endpoints; after transporting both endpoints to the root along the
tree, this becomes simultaneous root conjugation.  The generic
stabilizer of a six-tuple is the finite center, so the generic
dimension drops from \(6\cdot8\) to \(5\cdot8=40\).
\end{proof}

\begin{warning}[A fixed conditional fibre is not gauge invariant]
\label{warn:s6v96-fixed-fibre}
Every vertex of \(\Lambda_e\) is incident to links outside the
enlarged star.  A vertex gauge transformation which tree-gauges the
internal links also changes those exterior links.  Therefore tree
gauge is not an internal symmetry of a fibre with the far exterior
\(\eta\) held pointwise fixed.
\end{warning}

\begin{theorem}[Gauge covariance transports conditional fibres]
\label{thm:s6v96-fibre-transport}
Assume the full finite-volume measure is gauge invariant and admits
the disintegration \eqref{eq:s6v94-disintegration}.  For a lattice
gauge transformation \(g\), the conditional measures may be chosen
equivariantly so that
\begin{equation}
 \mu_e^{g\eta}(gA)=\mu_e^\eta(A)                         \label{eq:s6v96-equivariant}
\end{equation}
for almost every \(\eta\) and every measurable fibre event \(A\).
Consequently, conditional Poincaré constants, restricted transfer
norms, and fibre action margins are unchanged between the fibres
\(\eta\) and \(g\eta\), provided the conditional transfers are gauge
covariant.  A uniform essential-supremum estimate over a
gauge-saturated exterior class may therefore be proved in tree gauge.
\end{theorem}

\begin{proof}
For bounded measurable \(F\) on the fibre and \(h\) on the exterior,
gauge invariance gives
\[
 \int h(g\eta)F(gx)\,
 \mu_e^\eta(\dd x)\bar\mu_e(\dd\eta)
 =
 \int h(\eta)F(x)\,
 \mu_e^\eta(\dd x)\bar\mu_e(\dd\eta).
\]
Change the exterior variable by the measure-preserving gauge action
and use uniqueness of regular conditional probabilities up to null
sets.  This yields \eqref{eq:s6v96-equivariant}.  Pullback by \(g\) is
unitary between the two conditional \(L^2\) spaces and preserves the
gauge-invariant Dirichlet form.  Gauge covariance intertwines the
transfers and event multipliers, so all stated constants agree.
\end{proof}

\begin{firewall}[Uniformity is what licenses tree gauge]
\label{fire:s6v96-gauge}
\Cref{thm:s6v96-fibre-transport} compares different exteriors.  It
does not turn a bound proved for one preferred exterior into a
worst-exterior bound.  The admissible exterior family must be
gauge-saturated, and the estimate must remain uniform after the
exterior is transformed.
\end{firewall}

\begin{lemma}[Tree gauge has constant lattice Faddeev--Popov factor]
\label{lem:s6v96-FP}
For product Haar reference measure on the \(19\) star links, the
recursive change to tree-link gauge variables and six chord variables
has unit Jacobian.  Integration of gauge-invariant functions reduces
to product Haar measure on the six chords, followed by quotienting the
residual simultaneous conjugation.
\end{lemma}

\begin{proof}
At each tree step the new link or vertex variable is obtained by left
or right multiplication by products of previously chosen group
elements.  Haar measure is invariant under both operations, so every
step has Jacobian one.  Iterating over the tree factors the \(13\)
nonroot gauge variables as a constant group volume.  Normalized Haar
measure makes that volume one.  The remaining root action is
simultaneous conjugation.
\end{proof}

\subsection{Common and relative chord variables}
\label{sec:s6v96-common}

Fix an orientation of the six chords so that they are the six staple
matrices in the tree frame, replacing a chord by its inverse when
necessary.  Define
\begin{equation}
 C:=S_1,\qquad
 R_j:=S_1^{-1}S_j,\quad 2\le j\le6,\qquad R_1:=I.        \label{eq:s6v96-CR}
\end{equation}

\begin{lemma}[Global Haar-preserving coordinate map]
\label{lem:s6v96-CR-map}
The map
\[
 (S_1,\ldots,S_6)
 \longleftrightarrow
 (C,R_2,\ldots,R_6),
 \qquad S_j=CR_j,
\]
is a global diffeomorphism of \(SU(3)^6\) and preserves product Haar
measure.
\end{lemma}

\begin{proof}
The displayed formulas are smooth inverse maps.  For fixed \(C\),
left multiplication \(R_j\mapsto CR_j\) preserves Haar measure for
every \(j\ge2\), while \(S_1=C\).  Fubini's theorem proves preservation
of the product measure.
\end{proof}

\begin{lemma}[Frustration depends only on relative variables]
\label{lem:s6v96-relative-F}
Put
\[
 H_R:=I+\sum_{j=2}^6R_j.
\]
Then
\begin{align}
 \sum_{j=1}^6S_j&=CH_R,                                  \label{eq:s6v96-H}\\
 {\mathfrak F}(S_1,\ldots,S_6)
 &=18-\max_{W\in SU(3)}\Re\Tr(WH_R).                    \label{eq:s6v96-F-relative}
\end{align}
In particular, \({\mathfrak F}\) is independent of \(C\).
\end{lemma}

\begin{proof}
The first identity follows from \(S_j=CR_j\).  In the maximum, make
the bijective substitution \(W=UC\):
\[
 \max_U\Re\Tr(UCH_R)=\max_W\Re\Tr(WH_R).
\]
\end{proof}

\begin{proposition}[Two-sided relative-tube inclusions]
\label{prop:s6v96-tube}
For
\[
 {\cal T}_-(f_*):=
 \left\{(R_2,\ldots,R_6):
 {1\over2}\sum_{j=2}^6\|R_j-I\|_{\rm F}^2<f_*\right\}
\]
and
\[
 {\cal T}_+(f_*):=
 \left\{(R_2,\ldots,R_6):
 \max_{2\le j\le6}\|R_j-I\|_{\rm F}
 <2\sqrt{2f_*}\right\},
\]
one has
\begin{equation}
 {\cal T}_-(f_*)
 \subset
 \{{\mathfrak F}<f_*\}
 \subset
 {\cal T}_+(f_*).                                       \label{eq:s6v96-tube}
\end{equation}
\end{proposition}

\begin{proof}
For the first inclusion, use \(W=I\) as a competitor in
\eqref{eq:s6v96-F-relative}:
\[
 {\mathfrak F}
 \le
 \sum_{j=2}^6(3-\Re\Tr R_j)
 ={1\over2}\sum_{j=2}^6\|R_j-I\|_{\rm F}^2.
\]
For the second, apply \eqref{eq:s6v95-pairwise} with
\(S_1=C\) and \(S_j=CR_j\):
\[
 \|R_j-I\|_{\rm F}
 =\|S_j-S_1\|_{\rm F}
 \le2\sqrt{2{\mathfrak F}}
 <2\sqrt{2f_*}.
\]
\end{proof}

\begin{corollary}[A single small relative chart]
\label{cor:s6v96-relative-chart}
Choose \(f_*\) below both \(f_{\rm al}(6)\) and
\(\rho_{\exp}^2/8\).  Then every good relative variable has a unique
logarithm
\[
 R_j=e^{X_j},\qquad X_j\in\mathfrak{su}(3),
\]
in the fixed identity chart.  Thus the five relative modes lie in one
bounded \(40\)-dimensional Euclidean chart; only the common mode
\(C\in SU(3)\) remains global.
\end{corollary}

\begin{proof}
\Cref{prop:s6v96-tube} gives
\(\|R_j-I\|_{\rm F}<2\sqrt{2f_*}<\rho_{\exp}\).  Apply the local
logarithm property used in V95 to each \(R_j\).
\end{proof}

\subsection{The common mode has the V95 one-well phase}
\label{sec:s6v96-common-well}

The six incident plaquette factors in tree gauge have bare Wilson
phase
\begin{equation}
 \sum_{j=1}^6\Re\Tr S_j
 =
 \Re\Tr(CH_R).                                           \label{eq:s6v96-star-phase}
\end{equation}

\begin{theorem}[Conditional common-mode one-well geometry]
\label{thm:s6v96-common-well}
Fix \(R=(R_2,\ldots,R_6)\) with
\({\mathfrak F}(R)<f_{\rm al}(6)\).  Then the function
\[
 C\longmapsto\Re\Tr(CH_R)
\]
has a unique global maximizer \(C_0(R)\) and
\begin{equation}
 \Re\Tr(C_0(R)H_R)-\Re\Tr(CH_R)
 \ge {3\over4e^2}
 \|CC_0(R)^{-1}-I\|_{\rm F}^2.                          \label{eq:s6v96-common-deficit}
\end{equation}
The maximizer is equivariant under residual simultaneous conjugation.
\end{theorem}

\begin{proof}
The six matrices \(I,R_2,\ldots,R_6\) have frustration
\eqref{eq:s6v96-F-relative}.  Apply
\Cref{cor:s6v95-six}.  If all chords are conjugated by \(g\), then
\(H_R\) is conjugated by \(g\); uniqueness shows that the maximizing
common coordinate is conjugated accordingly.
\end{proof}

\begin{firewall}[Only the incident Wilson phase has been reduced]
\label{fire:s6v96-complete-action}
Links in \(\Lambda_e\) also belong to plaquettes outside the six-page
star.  After conditioning on the far exterior, those terms, the
singular ground profile, and response sources add a function of
\((C,R)\) to \eqref{eq:s6v96-star-phase}.  V96 does not prove that
this added function preserves
\eqref{eq:s6v96-common-deficit}.  Complete-action stability remains
an IFSC/ASCC row.
\end{firewall}

\subsection{Exact six-mode variance decomposition}
\label{sec:s6v96-variance}

Let \(\nu(\dd C,\dd R)\) be any conditional probability on the
tree-gauge chord variables, including all physical density factors.
Let \(\nu_R\) be its relative marginal and
\(\nu(\dd C\mid R)\) its common-mode conditional.

\begin{theorem}[Common-relative variance split]
\label{thm:s6v96-variance}
For every \(f\in L^2(\nu)\), with
\[
 \bar f(R):=\int f(C,R)\,\nu(\dd C\mid R),
\]
one has
\begin{equation}
 \operatorname{Var}_\nu(f)
 =
 \int\operatorname{Var}_{\nu(\cdot\mid R)}(f)\,
      \nu_R(\dd R)
 +\operatorname{Var}_{\nu_R}(\bar f).                   \label{eq:s6v96-variance}
\end{equation}
\end{theorem}

\begin{proof}
This is the orthogonal decomposition
\[
 f-\nu f=(f-\bar f)+(\bar f-\nu f).
\]
The two terms are orthogonal because the first has conditional mean
zero given \(R\).
\end{proof}

Order the five relatives and let
\[
 {\cal F}_1=\{\varnothing,\Omega\},\qquad
 {\cal F}_k=\sigma(R_2,\ldots,R_k),\quad2\le k\le6.
\]

\begin{corollary}[Exact five-step relative martingale]
\label{cor:s6v96-relative-martingale}
\begin{equation}
 \operatorname{Var}_{\nu_R}(\bar f)
 =
 \sum_{k=2}^6
 \left\|
 \mathbb E[\bar f\mid{\cal F}_k]
 -\mathbb E[\bar f\mid{\cal F}_{k-1}]
 \right\|_{L^2(\nu_R)}^2.                               \label{eq:s6v96-relative-martingale}
\end{equation}
\end{corollary}

\begin{proof}
The five differences are orthogonal martingale increments and
their sum is \(\bar f-\nu_R\bar f\).
\end{proof}

\begin{remark}[No hidden volume factor]
\label{rem:s6v96-no-volume}
Equations \eqref{eq:s6v96-variance} and
\eqref{eq:s6v96-relative-martingale} contain one common-mode term and
five relative increments, regardless of the ambient lattice volume.
The issue is now uniform control of these six terms, not their count.
\end{remark}

\subsection{The exact remaining good-sector rows}

\begin{definition}[Common-relative aligned-star card (CRAC)]
\label{def:s6v96-CRAC}
CRAC consists of:
\begin{enumerate}[label=\textup{(C\arabic*)},leftmargin=4em]
\item the gauge-covariant tree-fibre reduction of
      \Cref{thm:s6v96-fibre-transport};
\item a complete-action conditional Poincaré bound for the common
      mode, stable relative to \eqref{eq:s6v96-common-deficit};
\item bounds for the five martingale increments in
      \eqref{eq:s6v96-relative-martingale} inside the chart
      \({\cal T}_+(f_*)\);
\item a mean-gluing and capacity estimate for the annulus
      \({\cal T}_+(f_*)\setminus{\cal T}_-(f_*)\) and for exits from
      the aligned chart;
\item response-ground, far-exterior, residual-conjugation, and cofinal
      uniformity of the preceding rows.
\end{enumerate}
\end{definition}

\begin{theorem}[CRAC supplies ASCC's corridor row]
\label{thm:s6v96-CRAC}
If CRAC holds with common-mode constant \(C_C\), relative-increment
constants \(C_2,\ldots,C_6\), and exit pair
\((\varepsilon_{\rm exit},\kappa_{\rm exit})\), then
\eqref{eq:s6v95-gluing} holds with
\begin{equation}
 C_{\rm cor}
 \le
 C_C+\sum_{k=2}^6C_k+\kappa_{\rm exit},                 \label{eq:s6v96-Ccor}
\end{equation}
and the same \(\varepsilon_{\rm exit}\).
\end{theorem}

\begin{proof}
Apply the common-mode estimate to the first term of
\eqref{eq:s6v96-variance}.  Apply the five relative estimates to the
orthogonal terms of
\eqref{eq:s6v96-relative-martingale}.  Add the exit-capacity debit.
Orthogonality makes the constants additive and produces no cross
terms.  This is precisely \eqref{eq:s6v95-gluing}.
\end{proof}

\begin{firewall}[CRAC is not populated by the coordinate count]
\label{fire:s6v96-CRAC}
The tree gauge and common-relative transformation are exact, but they
do not prove the six conditional Poincaré estimates in CRAC.  In
particular, the conditional law of \(C\) contains complete-action and
ground-profile terms, while the marginal law of the \(R_j\) is
obtained after integrating \(C\).  The latter operation can create
nonlocal effective interactions even though the relative chart is
finite dimensional.
\end{firewall}

\begin{decision}[V96 bottleneck decision]
\label{dec:s6v96-bottleneck}
The enlarged-star fibre has been reduced exactly from \(19\) link
variables to six chord variables, then to one global common mode and
five small relative modes.  The first analytic target is now the
common-mode conditional law.  V97 will derive a perturbation-stable
one-well Poincaré theorem on \(SU(3)\) with constants stated in terms
of the complete-action oscillation and gradient/Hessian debits.  If
the physical ground profile supplies none of those debits, that
failure will be recorded before treating the five relative modes.
\end{decision}

\begin{theorem}[Sixth-series V96 exact advance]
\label{thm:s6v96-decision}
V96 proves:
\begin{enumerate}[leftmargin=3em]
\item the exact \(14\)-vertex, \(19\)-edge, six-cycle star count;
\item a spanning-tree reduction to six SU(3) chord variables;
\item the required gauge-covariant transport between conditional
      fibres;
\item a Haar-preserving common-relative coordinate map;
\item the two-sided aligned-tube inclusions
      \eqref{eq:s6v96-tube};
\item the common-mode quadratic one-well deficit;
\item the exact one-plus-five variance decomposition.
\end{enumerate}
The remaining good-sector problem is CRAC, not a \(19\)-fold
unstructured block estimate.
\end{theorem}

\begin{proof}
The conclusions are
\Cref{lem:s6v96-counts,
lem:s6v96-tree,
prop:s6v96-tree-gauge,
thm:s6v96-fibre-transport,
lem:s6v96-FP,
lem:s6v96-CR-map,
prop:s6v96-tube,
thm:s6v96-common-well,
thm:s6v96-variance,
cor:s6v96-relative-martingale}.
The complete-action limitations are
\Cref{fire:s6v96-gauge,
fire:s6v96-complete-action,
fire:s6v96-CRAC}.
\end{proof}

\begin{programme}[V97: complete-action common-mode coercivity]
\label{prog:s6v96-V97}
\begin{enumerate}[leftmargin=3em]
\item Prove a compact one-well Poincaré lemma from a quadratic phase
      deficit, local Hessian control, and a bounded perturbation.
\item Apply it to \(C\mapsto s\Re\Tr(CH_R)\) uniformly for
      \(R\in{\cal T}_+(f_*)\).
\item Insert the complete-action, response, and ground-profile
      perturbations with their exact oscillation and derivative
      constants.
\item Test whether the available half-slab estimates control those
      constants or only an averaged action.
\item Keep the five relative martingale increments for the subsequent
      iteration.
\end{enumerate}
\end{programme}

\begin{firewall}[Claim boundary after sixth-series V96]
\label{fire:s6v96-boundary}
V96 proves the exact enlarged-star graph and tree-gauge reduction,
handles the fixed-exterior gauge qualification, constructs
common-relative chord coordinates, proves aligned-tube inclusions and
the common-mode bare one-well deficit, and reduces arbitrary fibre
variance to one common plus five relative terms.

V96 does not prove CRAC, complete-action common-mode coercivity,
relative-mode coercivity, physical exit capacity, uniform
response-matched action smoothing, enlarged-star approximate
tensorization or form incidence, GMPC or HSOCC, the continuum OS
theory, or the full Yang--Mills mass gap.
\end{firewall}

\begin{theorem}[Sixth-series V96 result ledger]
\label{thm:s6v96-results}
The positive conclusions in \Cref{fire:s6v96-boundary} are exactly
those of
\Cref{lem:s6v96-counts,
lem:s6v96-tree,
prop:s6v96-tree-gauge,
thm:s6v96-fibre-transport,
lem:s6v96-FP,
lem:s6v96-CR-map,
lem:s6v96-relative-F,
prop:s6v96-tube,
cor:s6v96-relative-chart,
thm:s6v96-common-well,
thm:s6v96-variance,
cor:s6v96-relative-martingale,
thm:s6v96-CRAC,
thm:s6v96-decision}.
\end{theorem}

\begin{proof}
Each conclusion is the cited result with its gauge-saturated exterior,
tree-frame, complete-action, ground-profile, response, corridor-exit,
and cofinal qualifications unchanged.
\end{proof}

\paragraph{Parent-version record.}
V96 was formed from
\texttt{Renewal\_Yang\_Mills\_Mass\_Gap\_Research\_Notes\_V95.tex}.
The V95 parent had \(44\,784\) logical source lines,
\(1\,558\,375\) bytes, and SHA--256
\[
 \texttt{ee037018e5adf5ad93b84c7e9ebbc43672c90c9b179481beb874f8115081f08d}.
\]

% =============================================================================
% END APPENDED SIXTH-SERIES V96 MATERIAL
% =============================================================================

% =============================================================================
% BEGIN APPENDED SIXTH-SERIES V97 MATERIAL
% =============================================================================

\section{Sixth-series V97: the star-energy core and exact mean gluing}
\label{sec:s6v97}

\subsection{One internal event owns both frustration and Wilson cost}
\label{sec:s6v97-defect}

In the tree frame of V96, let \(S_1,\ldots,S_6\) be the six oriented
plaquette chord holonomies and define
\begin{equation}
 {\mathfrak D}_e(S)
 :=
 \sum_{j=1}^6\bigl(3-\Re\Tr S_j\bigr)
 ={1\over2}\sum_{j=1}^6\|S_j-I\|_{\rm F}^2.             \label{eq:s6v97-D}
\end{equation}
This is the sum of the six incident Wilson plaquette deficits and is
gauge invariant before tree fixing.

\begin{lemma}[Star energy dominates frustration]
\label{lem:s6v97-F-le-D}
\begin{equation}
 0\le{\mathfrak F}_e\le{\mathfrak D}_e.                  \label{eq:s6v97-F-le-D}
\end{equation}
Consequently, for every \(a>0\),
\begin{equation}
 \{{\mathfrak F}_e\ge a\}
 \subset
 {\cal B}_{e,a}:=\{{\mathfrak D}_e\ge a\}.               \label{eq:s6v97-containment}
\end{equation}
\end{lemma}

\begin{proof}
In the tree frame,
\[
 {\mathfrak F}_e
 =
 18-\max_{V\in SU(3)}
 \sum_{j=1}^6\Re\Tr(VS_j).
\]
Using \(V=I\) as a competitor gives
\[
 {\mathfrak F}_e
 \le
 18-\sum_j\Re\Tr S_j
 ={\mathfrak D}_e.
\]
Nonnegativity follows because every summand
\(\Re\Tr(VS_j)\le3\).
\end{proof}

\begin{corollary}[Exact bare action margin]
\label{cor:s6v97-action}
For the six owned Wilson plaquettes with density factor
\[
 \exp\left\{s\sum_{j=1}^6\Re\Tr S_j\right\},
\]
the event \({\cal B}_{e,a}\) pays the exact action deficit
\begin{equation}
 sa.                                                     \label{eq:s6v97-action}
\end{equation}
\end{corollary}

\begin{proof}
The maximal phase is \(18s\), and
\[
 18s-s\sum_j\Re\Tr S_j=s{\mathfrak D}_e\ge sa
\]
on the event.
\end{proof}

\begin{assessment}[Improvement over a separate frustration event]
\label{ass:s6v97-improvement}
The internal event \({\cal B}_{e,a}\) contains every V83 frustrated
star by \eqref{eq:s6v97-containment}, while its bare Wilson action
cost is tautological.  Its complement is a simultaneous near-flat
core for all six physical cycle variables.  Thus separate aligned
corridors are unnecessary at the bare-star level.
\end{assessment}

\subsection{Convexity of the near-flat six-chord core}
\label{sec:s6v97-core}

Equip \(SU(3)\) with the bi-invariant metric
\(\|X\|_{\rm F}^2=-\Tr(X^2)\) on
\(\mathfrak{su}(3)\), and \(SU(3)^6\) with the product metric.  Put
\begin{equation}
 {\cal K}_a:=\{S\in SU(3)^6:{\mathfrak D}_e(S)<a\}.
                                                               \label{eq:s6v97-K}
\end{equation}

\begin{lemma}[Hessian of one plaquette deficit]
\label{lem:s6v97-Hessian}
Let \(d(U)=3-\Re\Tr U\).  At \(U\in SU(3)\), in the
left-invariant direction \(X\in\mathfrak{su}(3)\),
\begin{equation}
 \operatorname{Hess}d_U(X,X)
 =
 \Tr\!\left((-X^2)\operatorname{Herm}U\right).           \label{eq:s6v97-Hessian}
\end{equation}
If \(\|U-I\|_{\rm op}\le r<1\), then
\begin{equation}
 \operatorname{Hess}d_U(X,X)
 \ge(1-r)\|X\|_{\rm F}^2.                               \label{eq:s6v97-Hessian-lower}
\end{equation}
\end{lemma}

\begin{proof}
Along \(U(t)=e^{tX}U\),
\[
 {d^2\over dt^2}d(U(t))\bigg|_{t=0}
 =-\Re\Tr(X^2U)
 =\Tr((-X^2)\operatorname{Herm}U).
\]
If \(\|U-I\|_{\rm op}\le r\), then
\(\operatorname{Herm}U\succeq(1-r)I\).  Since
\(-X^2\succeq0\), the lower bound follows.
\end{proof}

\begin{lemma}[Uniform convex core]
\label{lem:s6v97-convex-core}
There exists a universal \(a_{\rm cvx}>0\) such that, for every
\(0<a<a_{\rm cvx}\), the set \({\cal K}_a\) is contained in a
strongly convex product normal neighbourhood, has convex boundary,
and
\begin{equation}
 \operatorname{Hess}{\mathfrak D}_e
 \succeq c_a g,\qquad
 c_a:=1-\sqrt{2a}>0,                                    \label{eq:s6v97-D-convex}
\end{equation}
throughout \({\cal K}_a\).
\end{lemma}

\begin{proof}
If \({\mathfrak D}_e<a\), then
\[
 \|S_j-I\|_{\rm op}
 \le\|S_j-I\|_{\rm F}<\sqrt{2a}
\]
for every \(j\).  Choose \(a_{\rm cvx}\) so small that the closed
ball of radius \(\sqrt{2a_{\rm cvx}}\) lies in a strongly convex
normal neighbourhood of \(I\), and also
\(\sqrt{2a_{\rm cvx}}<1\).
\Cref{lem:s6v97-Hessian} then gives
\eqref{eq:s6v97-D-convex} after summing the six product directions.
Along the unique product geodesic between two points of
\({\cal K}_a\), strict convexity gives
\[
 {\mathfrak D}_e(\gamma(t))
 \le
 (1-t){\mathfrak D}_e(\gamma(0))
 +t{\mathfrak D}_e(\gamma(1))<a.
\]
Thus the sublevel is geodesically convex.  Its regular boundary
\({\mathfrak D}_e=a\) has nonnegative second fundamental form because
the Hessian of the defining function is positive on tangent
directions.
\end{proof}

\subsection{A bare core Poincaré inequality}
\label{sec:s6v97-BE}

\begin{lemma}[Convex-boundary Bakry--Émery estimate]
\label{lem:s6v97-BE}
Let \(\Omega\) be a geodesically convex domain in a Riemannian
manifold with nonnegative Ricci curvature.  Let
\[
 \nu(\dd x)=Z^{-1}e^{-V(x)}\operatorname{vol}(\dd x)
\]
on \(\Omega\), with reflecting boundary, and suppose
\(\operatorname{Hess}V\succeq\kappa g\) for some \(\kappa>0\).
Then
\begin{equation}
 \operatorname{Var}_\nu(f)
 \le {1\over\kappa}
 \int_\Omega|\nabla f|^2\,\dd\nu.                       \label{eq:s6v97-BE}
\end{equation}
\end{lemma}

\begin{proof}
Let \(L=\Delta-\nabla V\cdot\nabla\) with Neumann boundary condition.
The weighted Bochner identity gives
\[
 {1\over2}L|\nabla f|^2
 -\langle\nabla f,\nabla Lf\rangle
 =
 \|\operatorname{Hess}f\|^2+
 (\operatorname{Ric}+\operatorname{Hess}V)
 (\nabla f,\nabla f).
\]
Integrate against \(\nu\).  The boundary term has the favorable sign
because the boundary is convex.  Hence
\[
 \int(Lf)^2\,\dd\nu
 \ge\kappa\int|\nabla f|^2\,\dd\nu.
\]
Apply this to Neumann eigenfunctions and use the spectral resolution
of the self-adjoint generator.  Every nonconstant eigenvalue is at
least \(\kappa\), which is equivalent to
\eqref{eq:s6v97-BE}.
\end{proof}

\begin{theorem}[Bare six-chord core gap]
\label{thm:s6v97-bare-core}
Let \(\nu_{s,a}^{0}\) be product Haar measure on
\({\cal K}_a\), normalized with density
\[
 {d\nu_{s,a}^{0}\over dS}
 \propto e^{-s{\mathfrak D}_e(S)}.
\]
For \(0<a<a_{\rm cvx}\),
\begin{equation}
 \boxed{\displaystyle
 \operatorname{Var}_{\nu_{s,a}^{0}}(f)
 \le {1\over sc_a}
 \int_{{\cal K}_a}|\nabla f|^2\,\dd\nu_{s,a}^{0}.}       \label{eq:s6v97-bare-core}
\end{equation}
\end{theorem}

\begin{proof}
The product of bi-invariant \(SU(3)\) metrics has nonnegative Ricci
curvature.  The potential is
\(V=s{\mathfrak D}_e\), whose Hessian is at least \(sc_a g\) by
\Cref{lem:s6v97-convex-core}.  Apply
\Cref{lem:s6v97-BE}.
\end{proof}

\begin{lemma}[Bounded core perturbation]
\label{lem:s6v97-Holley}
Let
\[
 {d\nu_{s,a}^{W}\over d\nu_{s,a}^{0}}
 ={e^{W}\over\nu_{s,a}^{0}(e^W)}.
\]
If
\(\operatorname{osc}_{{\cal K}_a}W\le K\), then
\begin{equation}
 \operatorname{Var}_{\nu_{s,a}^{W}}(f)
 \le
 {e^K\over sc_a}
 \int_{{\cal K}_a}|\nabla f|^2\,\dd\nu_{s,a}^{W}.        \label{eq:s6v97-perturbed-core}
\end{equation}
\end{lemma}

\begin{proof}
The ratio between the largest and smallest densities of
\(\nu_{s,a}^{W}\) relative to \(\nu_{s,a}^{0}\) is at most \(e^K\).
For any constant \(c\),
\[
 \int|f-\nu^Wf|^2\,\dd\nu^W
 \le\int|f-c|^2\,\dd\nu^W.
\]
Choose \(c=\nu^0f\), compare the variance integral to \(\nu^0\), use
\eqref{eq:s6v97-bare-core}, and compare the energy integral back to
\(\nu^W\).  The two normalized density comparisons combine to their
ratio \(e^K\).
\end{proof}

\begin{firewall}[The physical perturbation row]
\label{fire:s6v97-W}
For a physical conditional enlarged-star law, \(W\) contains all
unowned crossing plaquettes, response terms, and the singular ground
profile relative to the six owned Wilson factors.  V97 does not prove
a volume-, exterior-, and cofinal-uniform bound
\[
 \operatorname{osc}_{{\cal K}_a}W\le K.
\]
Without such a bound, \eqref{eq:s6v97-perturbed-core} is not a
physical core gap.
\end{firewall}

\subsection{Exact gluing of the core mean to the bad term}
\label{sec:s6v97-gluing}

\begin{lemma}[Two-set mean gluing]
\label{lem:s6v97-mean}
Let \(\nu\) be a probability measure, let
\({\cal B}\) have probability \(p<1\), put
\({\cal K}={\cal B}^c\) and \(q=1-p\), and let
\(\nu g=0\).  Then
\begin{equation}
 \int_{\cal K}|g|^2\,\dd\nu
 \le
 q\operatorname{Var}_{\nu(\cdot\mid{\cal K})}(g)
 +{p\over q}\int_{\cal B}|g|^2\,\dd\nu.                 \label{eq:s6v97-mean}
\end{equation}
\end{lemma}

\begin{proof}
Let \(m_{\cal K}\) and \(m_{\cal B}\) be the two conditional means.
Since \(qm_{\cal K}+pm_{\cal B}=0\),
\[
 q|m_{\cal K}|^2
 ={p^2\over q}|m_{\cal B}|^2
 \le {p\over q}\int_{\cal B}|g|^2\,\dd\nu.
\]
Add \(q\operatorname{Var}_{\nu(\cdot\mid{\cal K})}(g)\).
\end{proof}

\begin{lemma}[Restricted norm bounds bad probability]
\label{lem:s6v97-p-le-eta}
If \(P\) is Markov and
\(\eta=\|M_{\cal B}P\|^2\), then
\begin{equation}
 \nu({\cal B})\le\eta.                                   \label{eq:s6v97-p-le-eta}
\end{equation}
\end{lemma}

\begin{proof}
Apply \(M_{\cal B}P\) to the unit constant:
\[
 \|M_{\cal B}P1\|^2=\|1_{\cal B}\|^2=\nu({\cal B}).
\]
\end{proof}

\subsection{The single-core enlarged-star bootstrap}
\label{sec:s6v97-bootstrap}

\begin{theorem}[Core Poincaré plus restricted smoothing]
\label{thm:s6v97-bootstrap}
For each enlarged star \(e\), let
\({\cal B}_e=\{{\mathfrak D}_e\ge a\}\),
\({\cal K}_e={\cal B}_e^c\), and retain the conditional projection
\(D_e\) and commuting conditional transfer \(P_e\) of V94.  Assume:
\begin{enumerate}[label=\textup{(Q\arabic*)},leftmargin=4em]
\item enlarged-star approximate tensorization with constant \(C_*\);
\item a normalized-core conditional Poincaré estimate which, after
      energy incidence, gives
      \begin{equation}
       \sum_ew_e q_e
       \operatorname{Var}_{\mu_e^\eta(\cdot\mid{\cal K}_e)}
       (D_ef)
       \le C_{\rm core}{\cal E}_\mu(f);                  \label{eq:s6v97-core-incidence}
      \end{equation}
\item
      \(\sup_e\|M_{{\cal B}_e}P_e\|^2\le\eta<1\);
\item transfer-form incidence
      \eqref{eq:s6v94-transfer-incidence}.
\end{enumerate}
For every \(\theta>0\) satisfying
\begin{equation}
 (2+\theta)\eta<1,                                      \label{eq:s6v97-window}
\end{equation}
one has
\begin{equation}
 \boxed{\displaystyle
 \operatorname{Var}_\mu(f)
 \le
 {C_*\left\{
 (1-\eta)C_{\rm core}
 +(1+\theta^{-1})m_P\delta\right\}
 \over1-(2+\theta)\eta}
 {\cal E}_\mu(f).}                                      \label{eq:s6v97-Poincare}
\end{equation}
\end{theorem}

\begin{proof}
Set
\[
 S_e=\|D_ef\|^2,\qquad
 A_e=\|M_{{\cal B}_e}D_ef\|^2.
\]
Apply \Cref{lem:s6v97-mean} on each conditional fibre to
\(g=D_ef\), whose fibre mean is zero.  By
\Cref{lem:s6v97-p-le-eta}, \(p_e\le\eta\) and
\(q_e\ge1-\eta\).  Summing and using
\eqref{eq:s6v97-core-incidence} gives
\begin{equation}
 {\cal S}_*(f)
 \le
 C_{\rm core}{\cal E}_\mu(f)
 +{1\over1-\eta}{\cal A}_*(f).                          \label{eq:s6v97-S-core-A}
\end{equation}
The commuting smoothing theorem and transfer incidence give
\[
 {\cal A}_*(f)
 \le
 (1+\theta)\eta{\cal S}_*(f)
 +(1+\theta^{-1})m_P\delta{\cal E}_\mu(f).
\]
Insert this in \eqref{eq:s6v97-S-core-A}, multiply by
\(1-\eta\), and absorb.  The remaining coefficient of
\({\cal S}_*\) is
\[
 1-\eta-(1+\theta)\eta=1-(2+\theta)\eta.
\]
Apply enlarged-star approximate tensorization to obtain
\eqref{eq:s6v97-Poincare}.
\end{proof}

\begin{corollary}[Bare-core constant]
\label{cor:s6v97-core-constant}
If the physical core law differs from the bare six-chord law by a
conditional perturbation of oscillation at most \(K\), and if its
core gradient form is owned by the physical form with incidence
\(m_{\rm core}\), then
\begin{equation}
 C_{\rm core}
 \le {m_{\rm core}e^K\over sc_a}.                        \label{eq:s6v97-Ccore}
\end{equation}
\end{corollary}

\begin{proof}
Apply \Cref{lem:s6v97-Holley} on every fibre, multiply by \(q_e w_e\),
and use the assumed gradient-form incidence.
\end{proof}

\begin{corollary}[Conditional star-energy gap compiler]
\label{cor:s6v97-gap}
Assume the complete half-slab action gives
\[
 \eta_s\le e^{-c_as+O(\log s)}
\]
for \({\cal B}_{e,a}\), and the hypotheses of
\Cref{cor:s6v97-core-constant} hold with
\(K=O(1)\) and volume-uniform incidence constants.  Then, for all
large \(s\), \eqref{eq:s6v97-window} holds and the enlarged-star
bootstrap has no exponentially bad good-sector constant.
\end{corollary}

\begin{proof}
The restricted norm tends to zero exponentially, while
\eqref{eq:s6v97-Ccore} is \(O(s^{-1})\).  Insert these estimates in
\eqref{eq:s6v97-Poincare}.
\end{proof}

\begin{firewall}[Conditional compiler versus physical theorem]
\label{fire:s6v97-physical}
The bare star has now supplied both the action deficit and the
strongly convex core.  The following physical rows remain unproved:
\begin{enumerate}[leftmargin=3em]
\item complete-action restricted smoothing for
      \({\mathfrak D}_e\ge a\), uniformly over far exteriors and
      response sources;
\item the \(O(1)\) core oscillation bound for the singular ground and
      complete-action perturbation;
\item core gradient-form incidence, enlarged-star transfer incidence,
      and enlarged-star approximate tensorization.
\end{enumerate}
\Cref{cor:s6v97-gap} is conditional on those rows.
\end{firewall}

\begin{decision}[V97 bottleneck decision]
\label{dec:s6v97-bottleneck}
The aligned-corridor mean-gluing problem of V95--V96 is removed from
the primary route.  The gauge-invariant star energy
\({\mathfrak D}_e\) gives one connected convex core and one bad event,
and \Cref{lem:s6v97-mean} glues their means through the already
controlled bad norm.  The next target is the physical perturbation
\(W\): determine from the half-slab amplitude whether its oscillation
on the fixed core can be bounded independently of volume and cofinal
scale, or whether a ground-profile counterexample forces a second
restricted-smoothing layer.
\end{decision}

\begin{theorem}[Sixth-series V97 exact advance]
\label{thm:s6v97-decision}
V97 proves:
\begin{enumerate}[leftmargin=3em]
\item the gauge-invariant domination
      \({\mathfrak F}_e\le{\mathfrak D}_e\);
\item the exact owned Wilson action cost \(sa\);
\item uniform geodesic convexity of the six-chord core;
\item the bare core gap \(sc_a\) by a convex-boundary
      Bakry--Émery argument;
\item its bounded-perturbation stability;
\item the exact two-set mean-gluing inequality;
\item the single-core enlarged-star bootstrap
      \eqref{eq:s6v97-Poincare}.
\end{enumerate}
This replaces the unpopulated multi-corridor ASCC row by three
explicit physical estimates listed in
\Cref{fire:s6v97-physical}.
\end{theorem}

\begin{proof}
The conclusions are
\Cref{lem:s6v97-F-le-D,
cor:s6v97-action,
lem:s6v97-convex-core,
thm:s6v97-bare-core,
lem:s6v97-Holley,
lem:s6v97-mean,
thm:s6v97-bootstrap}.
\end{proof}

\begin{programme}[V98: physical core perturbation]
\label{prog:s6v97-V98}
\begin{enumerate}[leftmargin=3em]
\item Write the exact physical enlarged-star conditional density
      relative to the six owned Wilson factors.
\item Separate crossing-plaquette, response-source, and singular-ground
      contributions to \(W\).
\item Bound the crossing-plaquette oscillation on
      \({\cal K}_a\) using finite incidence and the fixed core radius.
\item Test whether the half-slab singular equation controls the
      conditional ground oscillation; do not replace an averaged
      estimate by an essential supremum.
\item If the ground oscillation is unavailable, construct a nested
      ground-profile bad event and apply restricted smoothing again.
\end{enumerate}
\end{programme}

\begin{firewall}[Claim boundary after sixth-series V97]
\label{fire:s6v97-boundary}
V97 replaces frustration by an owning star-energy event, proves the
bare six-chord core is geodesically convex with Poincaré gap at least
\(sc_a\), proves bounded-perturbation stability and exact core/bad
mean gluing, and derives a new enlarged-star gap compiler.

V97 does not prove the physical core perturbation bound, uniform
complete-action restricted smoothing, enlarged-star approximate
tensorization or form incidence, GMPC or HSOCC, the continuum OS
theory, or the full Yang--Mills mass gap.
\end{firewall}

\begin{theorem}[Sixth-series V97 result ledger]
\label{thm:s6v97-results}
The positive conclusions in \Cref{fire:s6v97-boundary} are exactly
those of
\Cref{lem:s6v97-F-le-D,
cor:s6v97-action,
lem:s6v97-Hessian,
lem:s6v97-convex-core,
lem:s6v97-BE,
thm:s6v97-bare-core,
lem:s6v97-Holley,
lem:s6v97-mean,
lem:s6v97-p-le-eta,
thm:s6v97-bootstrap,
cor:s6v97-core-constant,
cor:s6v97-gap,
thm:s6v97-decision}.
\end{theorem}

\begin{proof}
Each conclusion is the cited result with its tree-frame,
small-core, convex-boundary, perturbation-oscillation,
restricted-smoothing, form-incidence, response, and cofinal
qualifications unchanged.
\end{proof}

\paragraph{Parent-version record.}
V97 was formed from
\texttt{Renewal\_Yang\_Mills\_Mass\_Gap\_Research\_Notes\_V96.tex}.
The V96 parent had \(45\,355\) logical source lines,
\(1\,577\,643\) bytes, and SHA--256
\[
 \texttt{10d38b42bd51b83086a8e00981cdbf3dc47826a81d51869008c9d7436e111bc6}.
\]

% =============================================================================
% END APPENDED SIXTH-SERIES V97 MATERIAL
% =============================================================================

% =============================================================================
% BEGIN APPENDED SIXTH-SERIES V98 MATERIAL
% =============================================================================

\section{Sixth-series V98: why crude core perturbation fails and how
a Wilson collar repairs the Hessian}
\label{sec:s6v98}

\subsection{The physical core perturbation}

On a tree-gauged enlarged-star fibre, write the conditional density
schematically as
\begin{equation}
 \mu_e^\eta(\dd S)
 ={1\over Z_e^\eta}
 \exp\{-s{\mathfrak D}_e(S)+W_e^\eta(S)\}\,\dd S,        \label{eq:s6v98-density}
\end{equation}
where \(W_e^\eta\) contains crossing Wilson plaquettes, response
sources, reference-density terms, and the singular ground profile not
included in the six owned pages.  V97's bounded-perturbation branch
requires an \(O(1)\), or at worst \(O(\log s)\), oscillation of
\(W_e^\eta\) on \({\cal K}_a\).

\begin{firewall}[Boundary-gauge scope correction to V96--V97]
\label{fire:s6v98-boundary-gauge}
Tree gauge is exact for the isolated star quotient, but the required
gauge transformation depends on the internal link configuration and
changes exterior links incident to the star vertices.  Consequently,
the six-chord density \eqref{eq:s6v98-density} is not yet the density
of a pointwise-fixed-exterior \(19\)-link fibre.  The six-mode variance
decomposition of V96 is exact for the isolated quotient or for the
inner chord conditional after suitable boundary-frame variables are
retained; it is not a decomposition of all physical buffer
fluctuations.  Every physical use of V97 must therefore add the
boundary-frame/buffer modes described below.
\end{firewall}

\subsection{A quantitative obstruction to the crude oscillation route}
\label{sec:s6v98-oscillation}

Let
\[
 {\cal P}_e^\partial
 :=
 \{p:E(p)\cap\Lambda_e\ne\varnothing\}
 \setminus\{p:p\ni e\}.
\]
Since \(|\Lambda_e|\le19\) and every four-dimensional link lies in six
plaquettes,
\begin{equation}
 |{\cal P}_e^\partial|\le114.                            \label{eq:s6v98-cross-count}
\end{equation}

\begin{lemma}[Crossing Wilson oscillation]
\label{lem:s6v98-cross-osc}
There is a geometric constant \(L_\partial<\infty\), independent of
volume and \(s\), such that on the six-chord core
\({\cal K}_a\),
\begin{equation}
 \operatorname{osc}_{{\cal K}_a}
 W_{e,\rm cross}^\eta
 \le L_\partial s\sqrt a.                               \label{eq:s6v98-cross-osc}
\end{equation}
Without using the core diameter, the uniform bound is
\[
 \operatorname{osc}W_{e,\rm cross}^\eta
 \le {9\over2}|{\cal P}_e^\partial|s.
\]
\end{lemma}

\begin{proof}
For \(U,V\in SU(3)\),
\[
 |\Re\Tr U-\Re\Tr V|
 \le\sqrt3\,\|U-V\|_{\rm F}.
\]
Changing the six chord variables inside \({\cal K}_a\) changes each
by \(O(\sqrt a)\), by \eqref{eq:s6v97-D}.  A plaquette product is
Lipschitz in each of its four unitary factors, with a universal
constant, and only the bounded set
\eqref{eq:s6v98-cross-count} occurs.  Summation gives
\eqref{eq:s6v98-cross-osc}.  Globally,
\(\Re\Tr U\in[-3/2,3]\), whose range is \(9/2\).
\end{proof}

\begin{proposition}[Core-radius/action-margin incompatibility for the
crude comparison]
\label{prop:s6v98-incompatibility}
Suppose one uses only \eqref{eq:s6v98-cross-osc} to demand
\[
 \operatorname{osc}_{{\cal K}_{a_s}}
 W_{e,\rm cross}^\eta=O(\log s).
\]
Then necessarily
\begin{equation}
 a_s=O\!\left({(\log s)^2\over s^2}\right),              \label{eq:s6v98-small-a}
\end{equation}
and the bare bad-event action margin satisfies
\begin{equation}
 sa_s=O\!\left({(\log s)^2\over s}\right)\longrightarrow0.
                                                               \label{eq:s6v98-lost-margin}
\end{equation}
Thus the same threshold cannot provide both a bounded-perturbation
core comparison and an exponentially small bad restricted norm by
this crude route.
\end{proposition}

\begin{proof}
Equation \eqref{eq:s6v98-cross-osc} forces
\(\sqrt{a_s}=O((\log s)/s)\), which is
\eqref{eq:s6v98-small-a}.  Multiplication by \(s\) gives
\eqref{eq:s6v98-lost-margin}.
\end{proof}

\begin{decision}[The crossing Wilson terms must be principal]
\label{dec:s6v98-principal}
Crossing plaquettes cannot be treated as an \(O(1)\) perturbation of
the six page action at large \(s\).  They must be included in a
principal Wilson collar, so their positive Hessian near flat
holonomy is retained.  Only response and ground-profile terms may
then be treated as perturbations.
\end{decision}

\subsection{The one-layer Wilson collar}
\label{sec:s6v98-collar}

Define
\begin{align}
 {\cal P}_e^+
 &:=
 \{p:E(p)\cap\Lambda_e\ne\varnothing\},                  \label{eq:s6v98-Pplus}\\
 \Lambda_e^+
 &:=
 \bigcup_{p\in{\cal P}_e^+}E(p),                         \label{eq:s6v98-Lplus}\\
 \widehat\Lambda_e^+
 &:=
 \Lambda_e^+\cup
 \{\ell:\ell\text{ has an endpoint in }V(\Lambda_e)\},   \label{eq:s6v98-buffer}\\
 {\mathfrak D}_e^+
 &:=
 \sum_{p\in{\cal P}_e^+}
 \bigl(3-\Re\Tr U_p\bigr).                              \label{eq:s6v98-Dplus}
\end{align}

\begin{lemma}[Uniform collar size]
\label{lem:s6v98-collar-size}
The cardinalities of
\({\cal P}_e^+\), \(\Lambda_e^+\), and
\(\widehat\Lambda_e^+\), and the number of collars containing any
fixed link or plaquette, are bounded by constants
depending only on dimension.  The crude four-dimensional bounds
\begin{equation}
 |{\cal P}_e^+|\le114,\qquad
 |\Lambda_e^+|\le361,\qquad
 |\widehat\Lambda_e^+|\le473                            \label{eq:s6v98-collar-size}
\end{equation}
are sufficient.
\end{lemma}

\begin{proof}
Each of the at most \(19\) inner links belongs to six plaquettes,
giving at most \(114\) plaquettes.  Each such plaquette adds at most
three links to the original \(19\), giving
\(19+3\cdot114=361\).  Reversing the incidence count gives the same
type of dimension-only overlap bound.  The \(14\) inner vertices have
at most eight incident unoriented links each, so adjoining every such
link adds at most \(112\), giving the last crude bound.
\end{proof}

\begin{lemma}[The buffer licenses inner tree gauge at fixed far
exterior]
\label{lem:s6v98-buffer-gauge}
In the conditional fibre whose updated links are
\(\widehat\Lambda_e^+\), gauge transformations at the \(14\) vertices
of \(\Lambda_e\) change only updated links.  Hence the \(13\) nonroot
inner vertex gauges may set the V96 inner spanning tree to the
identity without changing the far exterior.
\end{lemma}

\begin{proof}
Every link changed by a gauge transformation at an inner vertex has
that vertex as an endpoint and is included by
\eqref{eq:s6v98-buffer}.  The recursive tree-gauge construction of
\Cref{prop:s6v96-tree-gauge} uses only the \(13\) nonroot inner
vertices.  It therefore acts inside the buffer fibre.
\end{proof}

\begin{warning}[Six chords do not exhaust the buffer]
\label{warn:s6v98-buffer-modes}
After the inner tree is fixed, the buffer links not belonging to the
six page chords remain genuine conditional variables.  The V97
six-chord core estimate controls the inner conditional variance given
those variables.  A physical block Poincaré inequality also requires
coercivity and mean gluing for the remaining buffer modes.
\end{warning}

\subsection{Finite-complex Wilson Hessian}
\label{sec:s6v98-Hessian}

Let \(K\) be a finite connected lattice two-complex, fix a spanning
tree in its one-skeleton, and use the remaining link coordinates.
Let
\[
 B_K:C^1_T(K;\mathfrak{su}(3))
 \longrightarrow C^2(K;\mathfrak{su}(3))
\]
be the linearized plaquette-curvature map at the identity.  Put
\begin{equation}
 \lambda_K:=
 \inf_{X\ne0}{\|B_KX\|^2\over\|X\|^2}.                  \label{eq:s6v98-lambdaK}
\end{equation}

\begin{lemma}[Exact flat Wilson Hessian]
\label{lem:s6v98-flat-Hessian}
For the gauge-fixed Wilson defect
\[
 {\mathfrak D}_K(U):=
 \sum_{p\in K^{(2)}}(3-\Re\Tr U_p),
\]
one has at the identity connection
\begin{equation}
 \operatorname{Hess}{\mathfrak D}_K(I)
 =
 B_K^*B_K.                                               \label{eq:s6v98-flat-Hessian}
\end{equation}
Consequently, if \(B_K\) is injective,
\[
 \operatorname{Hess}{\mathfrak D}_K(I)
 \succeq\lambda_K I,\qquad \lambda_K>0.
\]
\end{lemma}

\begin{proof}
For a link perturbation \(U_\ell(t)=e^{tX_\ell}\), the first
variation of the plaquette holonomy is the signed sum
\[
 {d\over dt}U_p(t)\bigg|_{t=0}=(B_KX)_p.
\]
Since
\[
 3-\Re\Tr(e^{tY})
 ={t^2\over2}\|Y\|_{\rm F}^2+O(t^3),
\]
the second variation of the total defect is
\(\sum_p\|(B_KX)_p\|^2\).  Polarization gives
\eqref{eq:s6v98-flat-Hessian}.  In finite dimension, injectivity
implies the strictly positive minimum
\eqref{eq:s6v98-lambdaK}.
\end{proof}

\begin{definition}[Curvature reconstruction row]
\label{def:s6v98-reconstruction}
A gauge-fixed finite complex has reconstruction constant \(C_K\) on
a neighbourhood of the flat sector if
\begin{equation}
 \operatorname{dist}(U,{\cal F}_K)^2
 \le C_K{\mathfrak D}_K(U),                              \label{eq:s6v98-reconstruction}
\end{equation}
where \({\cal F}_K\) is its set of flat gauge-fixed connections.
\end{definition}

\begin{theorem}[Near-flat Wilson collar coercivity]
\label{thm:s6v98-collar-coercivity}
Suppose \(B_K\) is injective, the flat gauge-fixed connection is
unique, and \eqref{eq:s6v98-reconstruction} holds.  Then there exists
\(a_K>0\), depending only on the finite complex, such that
\begin{equation}
 {\mathfrak D}_K(U)<a_K
 \quad\Longrightarrow\quad
 \operatorname{Hess}{\mathfrak D}_K(U)
 \succeq{\lambda_K\over2}I.                             \label{eq:s6v98-collar-coercivity}
\end{equation}
\end{theorem}

\begin{proof}
In tree coordinates the Hessian is a continuous matrix-valued
function.  Hence there is a radius \(r_K>0\) on which it differs from
\(B_K^*B_K\) by at most \(\lambda_K/2\) in operator norm.
Choose \(a_K\le r_K^2/C_K\).  Reconstruction then places every
configuration with defect below \(a_K\) inside that radius.
\Cref{lem:s6v98-flat-Hessian} and the operator-norm perturbation bound
give \eqref{eq:s6v98-collar-coercivity}.
\end{proof}

\begin{remark}[The six-page star has no reconstruction mystery]
\label{rem:s6v98-page}
For the original six-page complex, tree gauge leaves one chord per
page and \(B_K\) is the identity on those six chord directions.
Thus \(\lambda_K=1\), and
\eqref{eq:s6v97-D} is itself the nonlinear reconstruction estimate.
The new acquisition is reconstruction for the one-layer collar after
its boundary conditions and residual flat sectors are declared.
\end{remark}

\subsection{Response and ground perturbations after collar ownership}

\begin{corollary}[Hessian-stable complete collar]
\label{cor:s6v98-response}
Let the conditional log density on a gauge-fixed collar core be
\[
 -s{\mathfrak D}_K(U)+R_s^\eta(U).
\]
If
\begin{equation}
 \|\operatorname{Hess}R_s^\eta\|_{\rm op}
 \le\chi_s
 <{s\lambda_K\over2}                                   \label{eq:s6v98-response-Hessian}
\end{equation}
uniformly on
\(\{{\mathfrak D}_K<a_K\}\), then the full negative log density has
Hessian at least
\begin{equation}
 {s\lambda_K\over2}-\chi_s.                             \label{eq:s6v98-full-Hessian}
\end{equation}
On a convex core with convex boundary, its conditional Poincaré gap
is at least that number.
\end{corollary}

\begin{proof}
Combine \eqref{eq:s6v98-collar-coercivity} with the Hessian bound and
apply \Cref{lem:s6v97-BE}.
\end{proof}

\begin{firewall}[The singular ground Hessian is not controlled]
\label{fire:s6v98-ground}
For the physical conditional law,
\(R_s^\eta\) includes \(2\log u_0\).  The differentiated singular
equation in V92 gives an averaged first-derivative Fisher bound, not
the uniform second-derivative estimate
\eqref{eq:s6v98-response-Hessian}.  Therefore
\Cref{cor:s6v98-response} is not yet a physical core theorem.
\end{firewall}

\begin{definition}[Wilson collar curvature card (WCCC)]
\label{def:s6v98-WCCC}
WCCC consists of:
\begin{enumerate}[label=\textup{(W\arabic*)},leftmargin=4em]
\item the buffered one-layer complex on
      \(\widehat\Lambda_e^+\), with the inner tree fixed as in
      \Cref{lem:s6v98-buffer-gauge} and a declared outer-boundary
      convention;
\item injectivity of \(B_{K_e^+}\) after removal of all residual flat
      modes;
\item the nonlinear reconstruction estimate
      \eqref{eq:s6v98-reconstruction};
\item complete-action restricted smoothing for
      \(\{{\mathfrak D}_{K_e^+}\ge a_K\}\);
\item response-ground Hessian control
      \eqref{eq:s6v98-response-Hessian}, or a further capacity split
      replacing it;
\item coercivity and mean gluing for the nonchord buffer modes;
\item collar form incidence, approximate tensorization, gauge
      covariance, and cofinal compatibility.
\end{enumerate}
\end{definition}

\begin{theorem}[WCCC supplies the V97 physical core rows]
\label{thm:s6v98-WCCC}
If WCCC holds with positive margin
\[
 \kappa_s={s\lambda_K\over2}-\chi_s>0,
\]
then the core-incidence row of V97 holds with its local Poincaré
constant replaced by \(\kappa_s\), and the collar bad event supplies
the restricted norm in the V97 bootstrap.  Hence
\eqref{eq:s6v97-Poincare} applies with the WCCC incidence constants.
\end{theorem}

\begin{proof}
Inner and buffer coercivity are
\Cref{cor:s6v98-response} and WCCC row (W6); WCCC row (W4) provides
the bad restricted norm.  Row (W7) supplies the two form incidences
and approximate tensorization required by
\Cref{thm:s6v97-bootstrap}.
\end{proof}

\begin{decision}[V98 bottleneck decision]
\label{dec:s6v98-bottleneck}
The naive treatment of crossing plaquettes as a bounded perturbation
is ruled out by
\Cref{prop:s6v98-incompatibility}.  The corrected route promotes all
one-layer crossing plaquettes into the principal Wilson collar and
adds the finite boundary-frame buffer needed to justify inner tree
gauge at fixed far exterior.  The next tasks are finite and explicit:
compute the buffered collar's gauge-fixed curvature matrix and flat
sectors, retain its nonchord modes, prove its reconstruction row, and
then confront the response-ground Hessian.  V99 will perform the
linear collar rank and boundary-mode audit before any claim of
positive \(\lambda_K\).
\end{decision}

\begin{theorem}[Sixth-series V98 exact advance]
\label{thm:s6v98-decision}
V98 proves:
\begin{enumerate}[leftmargin=3em]
\item the \(O(s\sqrt a)\) crossing-plaquette oscillation bound;
\item the incompatibility between crude bounded perturbation and a
      strong bad-event action margin;
\item uniform finite size of the one-layer collar;
\item the boundary-gauge correction and a finite buffer which
      licenses inner tree gauge at fixed far exterior;
\item the exact flat Wilson Hessian \(B_K^*B_K\);
\item a near-flat collar coercivity theorem from injectivity and
      nonlinear reconstruction;
\item the precise response-ground Hessian row needed after Wilson
      ownership.
\end{enumerate}
\end{theorem}

\begin{proof}
The assertions are
\Cref{lem:s6v98-cross-osc,
prop:s6v98-incompatibility,
lem:s6v98-collar-size,
lem:s6v98-flat-Hessian,
thm:s6v98-collar-coercivity,
cor:s6v98-response}.
\end{proof}

\begin{programme}[V99: collar cohomology and curvature rank]
\label{prog:s6v98-V99}
\begin{enumerate}[leftmargin=3em]
\item Specify the cubical one-layer two-complex \(K_e^+\) exactly.
\item Compute \(B_K\), its kernel, and the effect of tree and boundary
      gauge fixing.
\item Separate pure-gauge, boundary, and genuine flat-connection
      zero modes.
\item If \(B_K\) is not injective, add the minimal plaquette or
      boundary owner needed to remove the kernel rather than assuming
      a spectral floor.
\item Prove the corresponding finite-complex reconstruction
      inequality with explicit, dimension-only constants.
\end{enumerate}
\end{programme}

\begin{firewall}[Claim boundary after sixth-series V98]
\label{fire:s6v98-boundary}
V98 proves that crossing Wilson plaquettes cannot be treated by a
crude bounded-oscillation comparison, promotes them to a finite
principal collar, corrects the fixed-exterior scope of the isolated
tree gauge by adding a finite boundary buffer, and proves an abstract
near-flat Wilson Hessian theorem from curvature-rank and
reconstruction inputs.

V98 does not prove WCCC for the actual collar, control the physical
ground Hessian, prove complete-action restricted smoothing,
approximate tensorization or form incidence, GMPC or HSOCC, the
continuum OS theory, or the full Yang--Mills mass gap.
\end{firewall}

\begin{theorem}[Sixth-series V98 result ledger]
\label{thm:s6v98-results}
The positive conclusions in \Cref{fire:s6v98-boundary} are exactly
those of
\Cref{lem:s6v98-cross-osc,
prop:s6v98-incompatibility,
lem:s6v98-collar-size,
lem:s6v98-buffer-gauge,
lem:s6v98-flat-Hessian,
thm:s6v98-collar-coercivity,
cor:s6v98-response,
thm:s6v98-WCCC,
thm:s6v98-decision}.
\end{theorem}

\begin{proof}
Each conclusion is the cited result with its finite-collar,
tree-gauge, curvature-injectivity, reconstruction,
response-ground-Hessian, restricted-smoothing, and cofinal
qualifications unchanged.
\end{proof}

\paragraph{Parent-version record.}
V98 was formed from
\texttt{Renewal\_Yang\_Mills\_Mass\_Gap\_Research\_Notes\_V97.tex}.
The V97 parent had \(45\,912\) logical source lines,
\(1\,595\,240\) bytes, and SHA--256
\[
 \texttt{d931af0961246371e276c84f97cf558a02dea150b83eefa65d3e4f9a8b364768}.
\]

% =============================================================================
% END APPENDED SIXTH-SERIES V98 MATERIAL
% =============================================================================

% =============================================================================
% BEGIN APPENDED SIXTH-SERIES V99 MATERIAL
% =============================================================================

\section{Sixth-series V99: exact collar curvature rank and the
boundary-frame split}
\label{sec:s6v99}

\subsection{Canonical finite collar}
\label{sec:s6v99-complex}

Represent a positive lattice link by \((b;\mu)\), where
\(b\in\mathbb Z^4\) is its lower endpoint and
\(\mu\in\{0,1,2,3\}\).  Represent a positive plaquette by
\((b;\mu,\nu)\), \(\mu<\nu\), with oriented boundary
\begin{equation}
 (b;\mu)+(b+e_\mu;\nu)
 -(b+e_\nu;\mu)-(b;\nu).                                \label{eq:s6v99-boundary}
\end{equation}
Take the central link to be \((0;0)\).  Let \(P_0\) be its six
incident plaquettes, \(E_0\) their union of links, and define
\begin{equation}
 P^+:=\{p:E(p)\cap E_0\ne\varnothing\},\qquad
 E^+:=\bigcup_{p\in P^+}E(p).                            \label{eq:s6v99-P-E}
\end{equation}
This is the precise one-layer collar used below.

\begin{theorem}[Exact collar census]
\label{thm:s6v99-census}
The canonical collar has
\begin{equation}
 |P_0|=6,\quad |E_0|=19,\quad |V_0|=14,\qquad
 |P^+|=72,\quad |E^+|=123,\quad |V^+|=64.               \label{eq:s6v99-census}
\end{equation}
Moreover, every link incident to an inner vertex in \(V_0\) already
belongs to \(E^+\).  Thus the buffered set
\(\widehat\Lambda_e^+\) of V98 has exactly \(123\), not merely at
most \(473\), links.
\end{theorem}

\begin{proof}
Use \eqref{eq:s6v99-boundary} as a duplicate-free enumeration rule.
The central link has, for each of three transverse directions, one
plaquette on each side, giving six.  Their edge and vertex census is
the V96 count \(19,14\).  For every one of the \(19\) links, enumerate
its six incident plaquettes and remove repetitions lexicographically;
the resulting list has \(72\) elements.  Enumerating their four
boundary links gives \(123\) distinct links and \(64\) endpoints.
Finally, enumerating the eight unoriented links incident to each of
the \(14\) inner vertices produces \(93\) distinct links, all present
in the \(123\)-link list.

This enumeration is finite and exact: it uses only tuple addition,
lexicographic set insertion, and the four signed terms in
\eqref{eq:s6v99-boundary}.  The nonzero-minor certificate below
independently checks the resulting incidence table.
\end{proof}

\subsection{The scalar curvature matrix}
\label{sec:s6v99-rank}

Let
\[
 \partial_2:\mathbb R^{P^+}\longrightarrow\mathbb R^{E^+}
\]
be the cellular boundary map given by
\eqref{eq:s6v99-boundary}, and let
\[
 B:=\partial_2^*:\mathbb R^{E^+}\longrightarrow
 \mathbb R^{P^+}.
\]
This is the scalar version of the linearized curvature map; the
\(\mathfrak{su}(3)\)-valued map at a flat connection is
\(B\otimes I_8\) after parallel transport to a common frame.

\begin{theorem}[Exact rational collar rank]
\label{thm:s6v99-rank}
\begin{equation}
 \operatorname{rank}_{\mathbb Q}B=60,\qquad
 \dim_{\mathbb Q}\ker B=63.                              \label{eq:s6v99-rank}
\end{equation}
The graph cycle rank is
\[
 |E^+|-|V^+|+1=123-64+1=60.
\]
Hence
\begin{equation}
 H^1(K_e^+;\mathbb Q)=0,\qquad
 \dim_{\mathbb Q}H^2(K_e^+;\mathbb Q)=72-60=12.         \label{eq:s6v99-cohomology}
\end{equation}
\end{theorem}

\begin{proof}
The relation \(\partial_1\partial_2=0\) gives
\[
 \operatorname{rank}B\le
 \dim\ker\partial_1=60.
\]
For the reverse inequality, order plaquettes and links
lexicographically by their integer tuples.  Delete the \(13\) inner
tree columns of \Cref{lem:s6v96-tree}.  In the resulting
\(72\times110\) integer matrix, select the following zero-based row
and column indices:
\begin{align*}
{\cal R}={}&
0,1,2,3,4,5,6,7,8,10,11,14,15,17,16,20,21,22,24,26,\\
&28,30,31,32,40,41,44,35,39,47,43,27,34,51,46,42,9,12,\\
&56,13,19,18,23,60,49,25,29,63,50,45,33,66,36,57,38,\\
&68,53,48,67,52,\\[0.3em]
{\cal C}={}&
0,2,4,6,7,8,13,15,16,18,19,21,22,23,25,27,29,30,32,33,\\
&34,36,38,39,40,41,42,44,46,47,48,51,53,55,56,57,62,64,\\
&66,67,68,69,73,75,76,77,78,81,82,83,84,85,86,87,88,89,\\
&90,91,94,95.
\end{align*}
The resulting \(60\times60\) matrix \(A\) has exact Bareiss
determinant
\begin{equation}
 \det A=-1.                                              \label{eq:s6v99-det}
\end{equation}
Equivalently, Gaussian elimination modulo \(1000003\) gives
\(\det A\equiv1000002\).  Since the integer determinant is nonzero,
\(\operatorname{rank}_{\mathbb Q}B\ge60\), proving equality.
The kernel and cohomology dimensions follow from the census and
rank-nullity.
\end{proof}

\begin{remark}[Why the certificate is rigorous]
\label{rem:s6v99-certificate}
The lists \({\cal R},{\cal C}\), the boundary rule
\eqref{eq:s6v99-boundary}, and the integer identity
\eqref{eq:s6v99-det} constitute a finite independently reproducible
certificate.  The modular residue alone already proves nonvanishing;
the exact Bareiss calculation records the stronger unimodularity.
\end{remark}

\subsection{The nullspace is entirely gauge}
\label{sec:s6v99-gauge-kernel}

Let
\[
 d_0:\mathbb R^{V^+}\longrightarrow\mathbb R^{E^+}
\]
be the scalar vertex-to-link coboundary.  Since the collar graph is
connected,
\[
 \dim\operatorname{im}d_0=|V^+|-1=63.
\]

\begin{corollary}[No genuine linear flat mode]
\label{cor:s6v99-no-flat}
\begin{equation}
 \ker B=\operatorname{im}d_0.                            \label{eq:s6v99-pure-gauge}
\end{equation}
Thus all \(63\) scalar zero modes of the full flat Wilson Hessian are
pure gauge.  For \(\mathfrak{su}(3)\)-valued links, the curvature rank
is \(60\cdot8=480\) and the gauge kernel has dimension
\(63\cdot8=504\).
\end{corollary}

\begin{proof}
Gauge invariance gives
\(\operatorname{im}d_0\subset\ker B\).  Both spaces have dimension
\(63\) by \Cref{thm:s6v99-rank}; hence they are equal.  Tensor with
\(\mathfrak{su}(3)\) for the final dimensions.
\end{proof}

\begin{corollary}[What remains after the permitted inner tree gauge]
\label{cor:s6v99-inner-tree}
After deleting the \(13\) inner-tree link columns, the curvature
matrix has
\begin{equation}
 72\text{ rows},\qquad110\text{ columns},\qquad
 \operatorname{rank}60,\qquad\operatorname{nullity}50.  \label{eq:s6v99-inner-rank}
\end{equation}
The \(50\) remaining scalar null directions are exactly the residual
boundary-frame gauge directions.
\end{corollary}

\begin{proof}
The unimodular minor in
\eqref{eq:s6v99-det} proves rank \(60\) even after the inner-tree
columns are deleted.  Rank-nullity gives \(50\).  Fixing \(13\)
independent tree gauges removes \(13\) dimensions from the
\(63\)-dimensional gauge image, leaving \(50\); inclusion and equal
dimension identify the kernel.
\end{proof}

\begin{firewall}[Linear cohomology is not the obstruction]
\label{fire:s6v99-not-H1}
The WCCC rank premise is now resolved: there is no genuine linear
curvature zero mode.  The obstruction is that the remaining \(50\)
gauge directions are based at the outer boundary.  Gauging them away
would change the pointwise-fixed far exterior.  They must instead be
held as conditional frame variables or controlled by additional
boundary action.
\end{firewall}

\subsection{A full spanning-tree conditional split}
\label{sec:s6v99-full-tree}

\begin{theorem}[Unimodular chord/frame decomposition]
\label{thm:s6v99-chord-frame}
The \(60\) link columns \({\cal C}\) in
\Cref{thm:s6v99-rank} may be taken as curvature-chord variables.
The complementary \(63\) links---the \(13\) inner-tree links together
with the other \(50\) post-inner-tree links---form a spanning tree
\(T^+\) of the \(64\)-vertex collar graph.  Thus:
\begin{enumerate}[leftmargin=3em]
\item the \(13\) inner tree links may be set to identity by inner
      vertex gauge transformations;
\item the additional \(50\) tree links are retained as conditioned
      boundary-frame variables;
\item conditional on those frame variables, the \(60\) chord
      variables have injective linearized curvature, with a
      unimodular selected plaquette minor.
\end{enumerate}
\end{theorem}

\begin{proof}
The complement has \(123-60=63=64-1\) edges and contains the inner
tree.  Direct adjacency traversal using
\eqref{eq:s6v99-boundary} reaches all \(64\) vertices.  A connected
\(64\)-vertex graph with \(63\) edges is a spanning tree.  The
unimodular minor proves injectivity on the complementary chord
coordinates.
\end{proof}

\begin{remark}[Conditioning is not forbidden gauge fixing]
\label{rem:s6v99-conditioning}
The \(50\) outer tree links are not set to the identity.  Their actual
values are included in the exterior sigma-field of the chord update.
Tree parallel transport merely trivializes Lie-algebra fibres and
defines based chord holonomies; by Haar invariance this coordinate
change has unit Jacobian.  Therefore the construction preserves the
pointwise far exterior.
\end{remark}

\subsection{An explicit linear spectral floor}
\label{sec:s6v99-floor}

\begin{proposition}[Certified curvature singular-value bound]
\label{prop:s6v99-floor}
The minor \(A\) in \eqref{eq:s6v99-det} has \(102\) nonzero entries,
all equal to \(1\) or \(-1\), and hence
\(\|A\|_{\rm F}^2=102\).  Its least singular value obeys
\begin{equation}
 \sigma_{\min}(A)^2\ge102^{-59}.                         \label{eq:s6v99-floor}
\end{equation}
Consequently the gauge-fixed flat Wilson Hessian on the selected
\(60\) chord variables has the explicit, dimension-only lower floor
\(102^{-59}\).
\end{proposition}

\begin{proof}
Let
\(\sigma_1\ge\cdots\ge\sigma_{60}>0\) be the singular values.
Unimodularity gives
\[
 \prod_{j=1}^{60}\sigma_j=|\det A|=1,
\]
while every \(\sigma_j\le\|A\|_{\rm op}\le\|A\|_{\rm F}=\sqrt{102}\).
Therefore
\[
 \sigma_{60}
 ={1\over\prod_{j<60}\sigma_j}
 \ge(\sqrt{102})^{-59}.
\]
Squaring proves \eqref{eq:s6v99-floor}.  The full curvature matrix
contains this row restriction, so its quadratic form is no smaller.
\end{proof}

\begin{assessment}[The floor is crude but sufficient in principle]
\label{ass:s6v99-floor}
The number \(102^{-59}\) is not intended as a sharp physical constant.
Its role is to prove that the linear collar floor is a genuine
positive number independent of lattice volume.  Numerical singular
values or a better combinatorial basis may improve it without
changing the proof architecture.
\end{assessment}

\subsection{Local nonlinear reconstruction}

\begin{theorem}[Uniform local chord reconstruction]
\label{thm:s6v99-local-reconstruction}
Fix the \(63\) tree-link values and let \(U^{\rm fl}\) be a flat chord
configuration compatible with them.  In tree-parallel-transport
coordinates there are constants \(r_{\rm loc},C_{\rm loc}>0\),
depending only on the finite collar geometry, such that
\begin{equation}
 \operatorname{dist}(U,U^{\rm fl})<r_{\rm loc}
 \quad\Longrightarrow\quad
 \operatorname{dist}(U,U^{\rm fl})^2
 \le C_{\rm loc}{\mathfrak D}_{K_e^+}(U).                \label{eq:s6v99-local-reconstruction}
\end{equation}
The constants may be chosen uniformly in the conditioned tree-link
values.
\end{theorem}

\begin{proof}
Use the selected \(60\) plaquette holonomies as a map from the
\(60\) chord copies of \(SU(3)\) to \(SU(3)^{60}\).  At a flat
configuration its derivative, after parallel transport to the root,
contains \(A\otimes I_8\) and is invertible by
\eqref{eq:s6v99-det}.  The quantitative inverse function theorem
therefore gives a locally Lipschitz inverse.  Near the identity,
\[
 3-\Re\Tr V\asymp\|\log V\|_{\rm F}^2,
\]
so squared plaquette-log distance is bounded by the Wilson defect.
All conditioned tree values lie in a compact product of \(SU(3)\);
parallel transports are isometries and the higher derivatives are
uniformly bounded.  A finite subcover makes
\(r_{\rm loc},C_{\rm loc}\) uniform.
\end{proof}

\begin{firewall}[Nonlinear flat sectors remain to be classified]
\label{fire:s6v99-flat-sectors}
The unimodular derivative proves local uniqueness near every flat
chord configuration.  It does not by itself prove that, for fixed
tree frames, there is only one global solution of all \(72\)
noncommutative plaquette equations.  Distinct isolated flat sectors,
if present, have zero Wilson defect and cannot be suppressed by an
action estimate.  A nonabelian presentation or cellular contraction
of \(K_e^+\) is required before the local reconstruction theorem can
be promoted to V98's global row.
\end{firewall}

\subsection{Revised collar acquisition card}

\begin{definition}[Chord-conditioned Wilson collar card (CCWCC)]
\label{def:s6v99-CCWCC}
CCWCC consists of:
\begin{enumerate}[label=\textup{(C\arabic*)},leftmargin=4em]
\item the exact collar census and full spanning tree of V99;
\item conditioning on the \(50\) outer frame links and updating the
      \(60\) curvature chords;
\item exclusion or controlled gluing of all nontrivial nonlinear flat
      sectors;
\item global small-defect reconstruction from
      \eqref{eq:s6v99-local-reconstruction};
\item complete-action restricted smoothing outside the reconstructed
      core;
\item response-ground stability inside each core;
\item a family of shifted chord updates covering every physical
      fluctuation, with approximate tensorization and form incidence.
\end{enumerate}
\end{definition}

\begin{theorem}[CCWCC populates WCCC's curvature rows]
\label{thm:s6v99-CCWCC}
CCWCC rows (C1)--(C4) imply WCCC rows (W1)--(W3), with a flat
curvature floor at least \(102^{-59}\).  If all CCWCC rows hold, the
V98-to-V97 handoff applies and yields the conditional gap
\eqref{eq:s6v97-Poincare}.
\end{theorem}

\begin{proof}
The census and boundary convention are
\Cref{thm:s6v99-census,thm:s6v99-chord-frame}.  Curvature injectivity
and its floor are
\Cref{thm:s6v99-rank,prop:s6v99-floor}.  Local reconstruction is
\Cref{thm:s6v99-local-reconstruction}; rows (C3)--(C4) globalize it.
The remaining CCWCC rows are exactly the smoothing, perturbation, and
incidence hypotheses of
\Cref{thm:s6v98-WCCC,thm:s6v97-bootstrap}.
\end{proof}

\begin{decision}[V99 bottleneck decision]
\label{dec:s6v99-bottleneck}
The collar's linear rank is now fully resolved: it has no genuine
linear flat mode, and the apparent \(50\)-dimensional kernel is
exactly the boundary-frame gauge space.  Conditioning on a
full-tree frame yields \(60\) chord variables with a certified
positive curvature floor.  The next compulsory-audit iteration must
classify the nonlinear flat sectors and examine whether the companion
programmes contain a usable finite-complex contraction or
frame-conditioned tensorization result.
\end{decision}

\begin{theorem}[Sixth-series V99 exact advance]
\label{thm:s6v99-decision}
V99 proves:
\begin{enumerate}[leftmargin=3em]
\item the exact \(64\)-vertex, \(123\)-link, \(72\)-plaquette collar
      census;
\item the exact curvature rank \(60\) and
      \(H^1(K_e^+;\mathbb Q)=0\);
\item identification of the entire linear kernel as gauge;
\item the exact \(50\)-mode boundary-frame remainder after inner tree
      fixing;
\item a full-tree conditioning split into \(63\) frames and \(60\)
      curvature chords;
\item a unimodular rank certificate and explicit floor
      \(102^{-59}\);
\item uniform local nonlinear chord reconstruction.
\end{enumerate}
\end{theorem}

\begin{proof}
The assertions are
\Cref{thm:s6v99-census,
thm:s6v99-rank,
cor:s6v99-no-flat,
cor:s6v99-inner-tree,
thm:s6v99-chord-frame,
prop:s6v99-floor,
thm:s6v99-local-reconstruction}.
\end{proof}

\begin{programme}[V100: compulsory audit and nonlinear flat sectors]
\label{prog:s6v99-V100}
\begin{enumerate}[leftmargin=3em]
\item Perform the compulsory current SM/NS audit.
\item Build the nonabelian presentation with \(60\) chord generators
      and \(72\) plaquette relators.
\item Eliminate generators using the collar's cubical attachment
      order or exhibit a residual flat representation.
\item If the presentation is trivial, promote local reconstruction
      to a global small-defect theorem by compactness.
\item Freeze the completed V100 as the next \(\texttt{\_f6}\)
      archive and restart at a compact V1 summary, as required by the
      hundred-version rule.
\end{enumerate}
\end{programme}

\begin{firewall}[Claim boundary after sixth-series V99]
\label{fire:s6v99-boundary}
V99 resolves the finite collar's linear curvature rank and
boundary-frame kernel, supplies a unimodular certificate and an
explicit spectral floor, and proves local nonlinear reconstruction
on a frame-conditioned \(60\)-chord fibre.

V99 does not classify nonlinear flat sectors, prove global
reconstruction, control response-ground perturbations, prove
complete-action restricted smoothing, approximate tensorization or
form incidence, GMPC or HSOCC, the continuum OS theory, or the full
Yang--Mills mass gap.
\end{firewall}

\begin{theorem}[Sixth-series V99 result ledger]
\label{thm:s6v99-results}
The positive conclusions in \Cref{fire:s6v99-boundary} are exactly
those of
\Cref{thm:s6v99-census,
thm:s6v99-rank,
cor:s6v99-no-flat,
cor:s6v99-inner-tree,
thm:s6v99-chord-frame,
prop:s6v99-floor,
thm:s6v99-local-reconstruction,
thm:s6v99-CCWCC,
thm:s6v99-decision}.
\end{theorem}

\begin{proof}
Each conclusion is the cited result with its finite enumeration,
frame-conditioning, nonlinear-flat-sector, response-ground,
restricted-smoothing, tensorization, and cofinal qualifications
unchanged.
\end{proof}

\paragraph{Parent-version record.}
V99 was formed from
\texttt{Renewal\_Yang\_Mills\_Mass\_Gap\_Research\_Notes\_V98.tex}.
The V98 parent had \(46\,397\) logical source lines,
\(1\,612\,762\) bytes, and SHA--256
\[
 \texttt{8510f84b31a4c5607217e5a86e5f219bce68570be2e37b70084f1b0542a5ab8f}.
\]

% =============================================================================
% END APPENDED SIXTH-SERIES V99 MATERIAL
% =============================================================================

% =============================================================================
% BEGIN APPENDED SIXTH-SERIES V100 MATERIAL
% =============================================================================

\section{Sixth-series V100: nonabelian collar contraction and global
small-defect reconstruction}
\label{sec:s6v100}

This is the compulsory five-version audit and the hundred-version
rollover note.  Its main mathematical task is the nonlinear question left
open in V99: whether the \(72\) collar plaquette equations admit an
isolated flat sector not visible to the linearized curvature matrix.  The
answer is no.  A finite nonabelian Tietze certificate contracts all
\(60\) chord generators.  This proves simple connectedness of the collar
two-complex, uniqueness of the flat chord completion for every fixed tree
frame, and the global small-defect reconstruction row required by CCWCC.

\subsection{Compulsory SM/NS audit}
\label{sec:s6v100-audit}

\begin{audit}[Current companion files read at V100]
\label{audit:s6v100-files}
The audit used the active files
\begin{enumerate}[label=\textup{(\roman*)},leftmargin=3.8em]
\item
\texttt{Renewal\_SM\_Einstein\_Continuum\_Research\_Notes\_V39.tex},
with \(15\,715\) logical source lines, \(572\,813\) bytes, and SHA--256
\[
 \texttt{9BDA6D4E1112E20FB770F08533CDBEB5B8A8FA1D8DB28F847651378CE658D7C3};
\]
\item
\texttt{Renewal\_NS\_Stability\_Research\_Notes\_V26.tex},
with \(5\,319\) logical source lines, \(278\,283\) bytes, and SHA--256
\[
 \texttt{82C887F8C2C69E4F28FAD7C3DC8DE5B9099CB5A4BF431EBA1FFF3E91935978AD}.
\]
\end{enumerate}
Both files have one live document terminator.  The NS file is unchanged
from the V95 audit.  The SM file has advanced from V34 to V39.
\end{audit}

\begin{theorem}[Transferable conclusions of the V100 audit]
\label{thm:s6v100-audit}
The companion audit changes the Yang--Mills direction in exactly the
following way.
\begin{enumerate}[leftmargin=3em]
\item SM V35 proves the Hilbert-valued commuting projected-smoothing
      inequality for an event internal to its conditional block.  This
      confirms the V94 algebra but supplies no complete-action collar
      restricted norm.
\item SM V36--V38 supply bounded-patch incidence and finite-horizon
      response compilers conditional on physical locality and infrared
      source order.  They do not establish those physical rows for the
      Yang--Mills collar.
\item SM V39 proves that a genuine unitary conjugation tangent is an
      exact range-one source and that fully covariant factors cancel.
      It also corrects the V38 claim: finite-horizon inversion controls
      the equation residual, not Hilbert-norm convergence without
      additional infrared information.
\item Consequently full-tree gauge transport may be used to classify
      flat collar configurations and to choose coordinates.  It may not
      be used to identify pointwise-fixed exterior fibres or to delete a
      physical response term unless the full conditional operator is
      proved covariant.
\item NS V26 supplies no new collar contraction, conditional
      tensorization, or response-ground theorem.
\end{enumerate}
\end{theorem}

\begin{proof}
The first item is the direct-integral identity
\(EP=PE=E\) and the resulting projected split proved in SM V35.  The
second is the bounded-overlap, graph-incidence, and finite-horizon heat
response chain of SM V36--V38.  The third and fourth are the exact
decomposition
\[
 Q\dot H\Omega=-HQG\Omega+QV\Omega
\]
and its covariance firewall in SM V39.  The hashes in
\Cref{audit:s6v100-files} fix the audited sources.  The NS conclusion is
its unchanged V26 terminal ledger.
\end{proof}

\begin{assessment}[Audit-guided restriction]
\label{ass:s6v100-audit}
The nonabelian flat-sector calculation below is purely a statement about
the finite collar complex and is invariant under full-tree gauge
transport.  No such transport is charged as a symmetry of the
pointwise-fixed conditional fibre.  This is the exact distinction required
by the SM V39 audit.
\end{assessment}

\subsection{The based nonabelian presentation}
\label{sec:s6v100-presentation}

Retain the lexicographic plaquette and link ordering and the full spanning
tree \(T^+\) of \Cref{thm:s6v99-chord-frame}.  Collapse \(T^+\) to a
basepoint.  Order the remaining chord generators
\[
 c_0,\ldots,c_{59}
\]
by the column order \({\cal C}\) in \Cref{thm:s6v99-rank}, and order the
plaquettes
\[
 p_0,\ldots,p_{71}
\]
lexicographically by \((b;\mu,\nu)\).  For each plaquette, read its
oriented boundary by \eqref{eq:s6v99-boundary} and delete the tree
letters.  Denote the resulting freely reduced word by \(w_r\).

\begin{lemma}[Presentation after maximal-tree collapse]
\label{lem:s6v100-presentation}
The fundamental group of the collar two-complex has presentation
\begin{equation}
 \pi_1(K_e^+)
 \cong
 \left\langle c_0,\ldots,c_{59}\ \middle|\
 w_0,\ldots,w_{71}\right\rangle .                       \label{eq:s6v100-presentation}
\end{equation}
Initially the \(72\) nonempty reduced relators have length histogram
\begin{equation}
 34\text{ of length }1,\qquad
 12\text{ of length }2,\qquad
 26\text{ of length }3.                                 \label{eq:s6v100-histogram}
\end{equation}
\end{lemma}

\begin{proof}
Collapsing a maximal tree in the one-skeleton is a homotopy equivalence.
The \(123-63=60\) non-tree edges give the generators, and the attaching
map of each of the \(72\) square two-cells gives its oriented boundary
word with tree letters deleted.  This is
\eqref{eq:s6v100-presentation}.  Direct enumeration from
\eqref{eq:s6v99-boundary} gives
\eqref{eq:s6v100-histogram}.
\end{proof}

\subsection{A finite Tietze contraction certificate}
\label{sec:s6v100-Tietze}

For a word \(w\), write
\[
 w\big|_{c_{g_1}=\cdots=c_{g_j}=1}
\]
for substitution followed by free and cyclic reduction.  All indices in
the next theorem are zero based.

\begin{theorem}[Nonabelian singleton-relator certificate]
\label{thm:s6v100-certificate}
There are ordered triples \((r_j,g_j,\epsilon_j)\),
\(1\le j\le60\), such that
\begin{equation}
 w_{r_j}\big|_{c_{g_1}=\cdots=c_{g_{j-1}}=1}
 =c_{g_j}^{\epsilon_j},\qquad
 \epsilon_j\in\{-1,+1\}.                                \label{eq:s6v100-singleton}
\end{equation}
The complete certificate is
\begin{align*}
&(5,3,+),(0,0,+),(1,1,+),(2,2,+),(3,4,-),(4,5,-),\\
&(6,6,+),(8,7,+),(9,8,+),(7,36,+),(11,9,+),(12,10,+),\\
&(10,37,+),(13,39,+),(16,11,-),(17,12,-),(18,13,+),\\
&(15,41,+),(19,14,+),(14,40,+),(20,15,+),(22,16,+),\\
&(23,17,+),(21,42,+),(25,45,+),(27,18,+),(28,19,-),\\
&(29,20,+),(26,46,+),(30,21,+),(33,50,+),(34,22,+),\\
&(35,23,+),(36,52,+),(37,53,+),(38,54,+),(44,26,-),\\
&(45,27,+),(32,49,+),(47,29,-),(48,30,-),(41,25,-),\\
&(43,57,+),(49,31,+),(24,44,+),(50,32,+),(31,48,+),\\
&(51,33,-),(52,34,-),(39,28,-),(46,59,+),(53,35,-),\\
&(40,24,-),(42,56,+),(66,51,+),(56,38,+),(60,43,+),\\
&(63,47,+),(67,58,-),(68,55,-).
\end{align*}
The generator indices \(g_j\) are a permutation of
\(\{0,\ldots,59\}\).  After the sixty substitutions, every remaining
relator is the empty word.
\end{theorem}

\begin{proof}
For each displayed triple, form the four signed boundary letters in
\eqref{eq:s6v99-boundary}, delete the \(63\) tree edges specified in
\Cref{thm:s6v99-chord-frame}, replace the selected chord columns by
\(c_0,\ldots,c_{59}\), and set the preceding displayed generators to the
identity.  Free reduction gives the singleton word in
\eqref{eq:s6v100-singleton}.  The calculation never forms a nontrivial
replacement word: each selected relator has reduced length one at its
step.  Its maximum intermediate word length is therefore three.

This finite check is reproduced by
\texttt{scripts/verify\_v100\_collar.py}, a \(220\)-line,
\(7\,575\)-byte integer-word verifier with SHA--256
\[
 \texttt{4B159312712BB7245213467E9BD9899A7CFF849487C33212C2EFE600BB01E1EB}.
\]
The verifier independently recovers
\[
 (|P_0|,|E_0|,|P^+|,|E^+|,|V^+|)
 =(6,19,72,123,64),
\]
checks that the \(63\)-edge complement has one component and cycle rank
zero, checks the histogram \eqref{eq:s6v100-histogram}, and reports
sixty eliminated generators, zero residual generators, and zero residual
relators.  Every operation is tuple enumeration, equality comparison, or
free-word cancellation, so the displayed data are an independently
checkable finite certificate.
\end{proof}

\begin{theorem}[The collar presentation is trivial]
\label{thm:s6v100-pi1}
\begin{equation}
 \boxed{\pi_1(K_e^+)=\{1\}.}                             \label{eq:s6v100-pi1}
\end{equation}
\end{theorem}

\begin{proof}
At the first step,
\eqref{eq:s6v100-singleton} makes \(c_{g_1}=1\).  This is a valid
Tietze elimination: substitute the identity for that generator in every
other relator and delete the singleton relation.  Inductively the
\(j\)-th step makes \(c_{g_j}=1\).  Since the \(g_j\) are a permutation
of all sixty generator indices, the presented group has no nonidentity
generator.  The final empty-relator check shows that no inconsistent
extra condition was introduced.  Apply
\Cref{lem:s6v100-presentation}.
\end{proof}

\begin{remark}[Why the argument is genuinely nonabelian]
\label{rem:s6v100-nonabelian}
The proof uses words in a free group and Tietze substitutions.  It does
not commute letters or pass to an exponent-sum matrix.  Thus it rules out
flat holonomy representations into every group, including nonabelian
\(SU(3)\); the V99 rational \(H^1\) calculation alone could not do this.
\end{remark}

\subsection{Classification of all flat collar sectors}
\label{sec:s6v100-flat}

\begin{theorem}[Unique flat chord completion at every tree frame]
\label{thm:s6v100-flat}
Fix arbitrary \(SU(3)\) values on the \(63\) links of \(T^+\), with a
root frame fixed.  There is exactly one assignment of the \(60\) chord
links for which all \(72\) plaquette holonomies are the identity.
Equivalently, every flat collar connection is gauge equivalent to the
trivial connection, and fixing the tree-link values selects exactly one
representative of that gauge orbit.
\end{theorem}

\begin{proof}
On a tree there is a unique vertex gauge transformation equal to the
identity at the root which sends all tree-link values to the identity:
construct it successively along the unique root paths.  Apply this
coordinate transformation to a flat connection.  The transformed chord
values define a homomorphism
\[
 \pi_1(K_e^+)\longrightarrow SU(3)
\]
by the presentation \eqref{eq:s6v100-presentation}.  The domain is
trivial by \Cref{thm:s6v100-pi1}, so every transformed chord is the
identity.  Undoing the uniquely determined tree gauge gives existence
and uniqueness of the chord values at the prescribed tree frame.
\end{proof}

\begin{corollary}[No nonlinear flat sector]
\label{cor:s6v100-no-sector}
The isolated-flat-sector firewall
\Cref{fire:s6v99-flat-sectors} is discharged.  For fixed tree frames the
zero set of the collar Wilson defect consists of the single point
\(U^{\rm fl}\).
\end{corollary}

\begin{proof}
For \(V\in SU(3)\),
\[
 3-\Re\Tr V\ge0,
\]
with equality exactly when \(V=I\).  Hence the Wilson defect vanishes
exactly when all plaquette holonomies are the identity.  Apply
\Cref{thm:s6v100-flat}.
\end{proof}

\subsection{Uniform global small-defect reconstruction}
\label{sec:s6v100-global}

Use a bi-invariant product distance in the tree-parallel-transport chord
coordinates.  Let \(U^{\rm fl}(\tau)\) denote the unique flat chord
completion at tree frame \(\tau\).

\begin{theorem}[Global collar reconstruction below a fixed defect]
\label{thm:s6v100-global}
There are constants \(a_{\rm glob},C_{\rm glob}>0\), depending only on
the \(64\)-vertex collar geometry, such that for every tree frame
\(\tau\) and every chord assignment \(U\),
\begin{equation}
 {\mathfrak D}_{K_e^+}(U)<a_{\rm glob}
 \quad\Longrightarrow\quad
 \operatorname{dist}\bigl(U,U^{\rm fl}(\tau)\bigr)^2
 \le C_{\rm glob}{\mathfrak D}_{K_e^+}(U).               \label{eq:s6v100-global}
\end{equation}
The constants are uniform in the finite volume, central-link location,
and conditioned tree-frame values.
\end{theorem}

\begin{proof}
Tree gauge is a product of left and right translations, hence an
isometry for a bi-invariant product metric; plaquette traces are
unchanged.  It is therefore enough to put all tree links at the identity.
The chord space \(X=SU(3)^{60}\) is compact, and by
\Cref{cor:s6v100-no-sector} the continuous nonnegative defect has the
unique zero \(U^{\rm fl}\).

Let \(r_{\rm loc}\) and \(C_{\rm loc}\) be the uniform constants from
\Cref{thm:s6v99-local-reconstruction}.  On the compact set
\[
 X_{\rm far}
 :=\{U:\operatorname{dist}(U,U^{\rm fl})\ge r_{\rm loc}/2\}
\]
the defect is strictly positive, so it has a minimum \(m_{\rm far}>0\).
Set
\[
 a_{\rm glob}:={m_{\rm far}\over2},
 \qquad C_{\rm glob}:=C_{\rm loc}.
\]
If the defect is below \(a_{\rm glob}\), the point is not in
\(X_{\rm far}\), and the V99 local reconstruction estimate applies.
Transporting back along the tree proves the uniform assertion.
\end{proof}

\begin{corollary}[CCWCC curvature rows are complete]
\label{cor:s6v100-CCWCC}
Rows (C1)--(C4) of the chord-conditioned Wilson collar card are proved.
Consequently WCCC's finite-collar census, curvature-floor, flat-sector,
and global-reconstruction rows hold, with the explicit linear floor
\(102^{-59}\).
\end{corollary}

\begin{proof}
Rows (C1)--(C2) are
\Cref{thm:s6v99-census,thm:s6v99-chord-frame,
prop:s6v99-floor}.  Row (C3) is
\Cref{thm:s6v100-flat}, and row (C4) is
\Cref{thm:s6v100-global}.  Invoke
\Cref{thm:s6v99-CCWCC}.
\end{proof}

\begin{firewall}[What simple connectedness does not prove]
\label{fire:s6v100-remaining}
The Tietze contraction classifies zero-curvature configurations; it does
not estimate the full conditional density away from zero curvature.
In particular it does not prove:
\begin{enumerate}[leftmargin=3em]
\item complete-action restricted smoothing when plaquettes outside
      \(P^+\) meet an updated chord;
\item response-ground stability for the physical conditional Hamiltonian;
\item a covering family of chord updates with approximate tensorization
      and physical Dirichlet-form incidence.
\end{enumerate}
These are CCWCC rows (C5)--(C7).  They cannot be inferred from the
unimodular minor or from the SM unitary-tangent identity.
\end{firewall}

\begin{decision}[V100 bottleneck decision]
\label{dec:s6v100-bottleneck}
The nonlinear topology is no longer the bottleneck.  The next series must
go upstream of CCWCC (C5): first determine whether any finite link block
can contain every Wilson plaquette meeting its updated variables.  If
finite action closure is impossible, replace that row by an explicit
boundary-layer influence or multiscale conditional mechanism rather than
enlarging the collar indefinitely.
\end{decision}

\begin{theorem}[Sixth-series V100 exact advance]
\label{thm:s6v100-decision}
V100 proves:
\begin{enumerate}[leftmargin=3em]
\item a nonabelian \(60\)-step Tietze certificate for the collar;
\item \(\pi_1(K_e^+)=1\);
\item uniqueness of the flat chord completion for arbitrary fixed tree
      frames;
\item uniform global small-defect reconstruction;
\item completion of CCWCC rows (C1)--(C4).
\end{enumerate}
The remaining rows are precisely the analytic complete-action,
response-ground, and tensorization/incidence rows in
\Cref{fire:s6v100-remaining}.
\end{theorem}

\begin{proof}
Use \Cref{thm:s6v100-certificate,thm:s6v100-pi1,
thm:s6v100-flat,thm:s6v100-global,cor:s6v100-CCWCC}.
\end{proof}

\begin{programme}[Post-rollover V1]
\label{prog:s6v100-next}
Freeze this note as the sixth hundred-version archive and restart with a
compact V1 containing the programme context, an honest major-results
ledger, the current theorem interface, and the open mass-gap gates.  The
first new attack will test finite complete-action closure and formulate
the smallest valid replacement for CCWCC (C5).
\end{programme}

\begin{firewall}[Claim boundary after sixth-series V100]
\label{fire:s6v100-boundary}
V100 resolves the finite collar's nonabelian flat-sector problem and
globalizes its small-defect reconstruction.  It does not prove CCWCC
(C5)--(C7), the conditional gap, GMPC or HSOCC, the Osterwalder--Schrader
continuum construction, or the full Yang--Mills mass gap.
\end{firewall}

\paragraph{Parent-version record.}
V100 was formed from
\texttt{Renewal\_Yang\_Mills\_Mass\_Gap\_Research\_Notes\_V99.tex}.
The V99 parent had \(46\,872\) logical source lines,
\(1\,630\,579\) bytes, and SHA--256
\[
 \texttt{FC4EBC96016C0AEFE6D85AF3146B5876637700CBA6540DB499BB37134E43DC10}.
\]

% =============================================================================
% END APPENDED SIXTH-SERIES V100 MATERIAL
% =============================================================================

\end{document}
